\section{Gráficos Tridimensionales}

\textit{Matplotlib} permite realizar gráficos de tres dimensiones (\(x,y,z\)) mediante \texttt{subplot\_kw{'projection' : '3d'}} (kw significa keywords). 

\subsection{Configuración del Entorno 3D}

La activación se realiza mediante el argumento \texttt{subplot\_kw} al inicializar la figura, donde se debe colocar:
\begin{itemize}
\item \texttt{projection='3d'}: Informa a \textit{Pyplot} que el objeto \texttt{ax} debe tener propiedades de profundidad, rotación y un tercer eje de coordenadas.
\item \texttt{set\_zlabel()}: Método exclusivo de esta proyección para etiquetar la profundidad.
\end{itemize}

\begin{pycode}{Creación de una hélice 3D}
import matplotlib.pyplot as plt
import numpy as np

# Datos
t = np.linspace(0, 4 * np.pi, 500)
z = t
x = np.sin(t)
y = np.cos(t)

# Inicialización del lienzo 3D
fig, ax = plt.subplots(subplot_kw={"projection": "3d"}, figsize=(5, 5)) 

# Trazado de línea en el espacio
ax.plot3D(x, y, z, color='purple', linewidth=2)

ax.set_title("Trayectoria Helicoidal en 3D")
ax.set_xlabel("Eje X")
ax.set_ylabel("Eje Y")
ax.set_zlabel("Eje Z (Tiempo/Profundidad)")

plt.show()
\end{pycode}

El gráfico de hélice se ve en la figura \ref{fig:grafico_Tridimensional}

\begin{figure}[ht]
    \centering
    %% Creator: Matplotlib, PGF backend
%%
%% To include the figure in your LaTeX document, write
%%   \input{<filename>.pgf}
%%
%% Make sure the required packages are loaded in your preamble
%%   \usepackage{pgf}
%%
%% Also ensure that all the required font packages are loaded; for instance,
%% the lmodern package is sometimes necessary when using math font.
%%   \usepackage{lmodern}
%%
%% Figures using additional raster images can only be included by \input if
%% they are in the same directory as the main LaTeX file. For loading figures
%% from other directories you can use the `import` package
%%   \usepackage{import}
%%
%% and then include the figures with
%%   \import{<path to file>}{<filename>.pgf}
%%
%% Matplotlib used the following preamble
%%   \def\mathdefault#1{#1}
%%   \everymath=\expandafter{\the\everymath\displaystyle}
%%   \IfFileExists{scrextend.sty}{
%%     \usepackage[fontsize=10.000000pt]{scrextend}
%%   }{
%%     \renewcommand{\normalsize}{\fontsize{10.000000}{12.000000}\selectfont}
%%     \normalsize
%%   }
%%   \usepackage{fontspec}
%%   \setmainfont{CMU Serif}
%%   \makeatletter\@ifpackageloaded{underscore}{}{\usepackage[strings]{underscore}}\makeatother
%%
\begingroup%
\makeatletter%
\begin{pgfpicture}%
\pgfpathrectangle{\pgfpointorigin}{\pgfqpoint{4.580936in}{4.678828in}}%
\pgfusepath{use as bounding box, clip}%
\begin{pgfscope}%
\pgfsetbuttcap%
\pgfsetmiterjoin%
\definecolor{currentfill}{rgb}{1.000000,1.000000,1.000000}%
\pgfsetfillcolor{currentfill}%
\pgfsetlinewidth{0.000000pt}%
\definecolor{currentstroke}{rgb}{1.000000,1.000000,1.000000}%
\pgfsetstrokecolor{currentstroke}%
\pgfsetdash{}{0pt}%
\pgfpathmoveto{\pgfqpoint{0.000000in}{0.000000in}}%
\pgfpathlineto{\pgfqpoint{4.580936in}{0.000000in}}%
\pgfpathlineto{\pgfqpoint{4.580936in}{4.678828in}}%
\pgfpathlineto{\pgfqpoint{0.000000in}{4.678828in}}%
\pgfpathlineto{\pgfqpoint{0.000000in}{0.000000in}}%
\pgfpathclose%
\pgfusepath{fill}%
\end{pgfscope}%
\begin{pgfscope}%
\pgfsetbuttcap%
\pgfsetmiterjoin%
\definecolor{currentfill}{rgb}{1.000000,1.000000,1.000000}%
\pgfsetfillcolor{currentfill}%
\pgfsetlinewidth{0.000000pt}%
\definecolor{currentstroke}{rgb}{0.000000,0.000000,0.000000}%
\pgfsetstrokecolor{currentstroke}%
\pgfsetstrokeopacity{0.000000}%
\pgfsetdash{}{0pt}%
\pgfpathmoveto{\pgfqpoint{0.100000in}{0.134025in}}%
\pgfpathlineto{\pgfqpoint{4.345803in}{0.134025in}}%
\pgfpathlineto{\pgfqpoint{4.345803in}{4.379828in}}%
\pgfpathlineto{\pgfqpoint{0.100000in}{4.379828in}}%
\pgfpathlineto{\pgfqpoint{0.100000in}{0.134025in}}%
\pgfpathclose%
\pgfusepath{fill}%
\end{pgfscope}%
\begin{pgfscope}%
\pgfsetbuttcap%
\pgfsetmiterjoin%
\definecolor{currentfill}{rgb}{0.950000,0.950000,0.950000}%
\pgfsetfillcolor{currentfill}%
\pgfsetfillopacity{0.500000}%
\pgfsetlinewidth{1.003750pt}%
\definecolor{currentstroke}{rgb}{0.950000,0.950000,0.950000}%
\pgfsetstrokecolor{currentstroke}%
\pgfsetstrokeopacity{0.500000}%
\pgfsetdash{}{0pt}%
\pgfpathmoveto{\pgfqpoint{0.420587in}{1.180904in}}%
\pgfpathlineto{\pgfqpoint{1.822690in}{2.356174in}}%
\pgfpathlineto{\pgfqpoint{1.803200in}{4.051124in}}%
\pgfpathlineto{\pgfqpoint{0.333999in}{2.978968in}}%
\pgfusepath{stroke,fill}%
\end{pgfscope}%
\begin{pgfscope}%
\pgfsetbuttcap%
\pgfsetmiterjoin%
\definecolor{currentfill}{rgb}{0.900000,0.900000,0.900000}%
\pgfsetfillcolor{currentfill}%
\pgfsetfillopacity{0.500000}%
\pgfsetlinewidth{1.003750pt}%
\definecolor{currentstroke}{rgb}{0.900000,0.900000,0.900000}%
\pgfsetstrokecolor{currentstroke}%
\pgfsetstrokeopacity{0.500000}%
\pgfsetdash{}{0pt}%
\pgfpathmoveto{\pgfqpoint{1.822690in}{2.356174in}}%
\pgfpathlineto{\pgfqpoint{4.072565in}{1.702224in}}%
\pgfpathlineto{\pgfqpoint{4.152855in}{3.455555in}}%
\pgfpathlineto{\pgfqpoint{1.803200in}{4.051124in}}%
\pgfusepath{stroke,fill}%
\end{pgfscope}%
\begin{pgfscope}%
\pgfsetbuttcap%
\pgfsetmiterjoin%
\definecolor{currentfill}{rgb}{0.925000,0.925000,0.925000}%
\pgfsetfillcolor{currentfill}%
\pgfsetfillopacity{0.500000}%
\pgfsetlinewidth{1.003750pt}%
\definecolor{currentstroke}{rgb}{0.925000,0.925000,0.925000}%
\pgfsetstrokecolor{currentstroke}%
\pgfsetstrokeopacity{0.500000}%
\pgfsetdash{}{0pt}%
\pgfpathmoveto{\pgfqpoint{0.420587in}{1.180904in}}%
\pgfpathlineto{\pgfqpoint{2.805565in}{0.401970in}}%
\pgfpathlineto{\pgfqpoint{4.072565in}{1.702224in}}%
\pgfpathlineto{\pgfqpoint{1.822690in}{2.356174in}}%
\pgfusepath{stroke,fill}%
\end{pgfscope}%
\begin{pgfscope}%
\pgfsetbuttcap%
\pgfsetroundjoin%
\pgfsetlinewidth{0.501875pt}%
\definecolor{currentstroke}{rgb}{0.501961,0.501961,0.501961}%
\pgfsetstrokecolor{currentstroke}%
\pgfsetdash{{1.850000pt}{0.800000pt}}{0.000000pt}%
\pgfpathmoveto{\pgfqpoint{0.565029in}{1.133729in}}%
\pgfpathlineto{\pgfqpoint{1.959516in}{2.316404in}}%
\pgfpathlineto{\pgfqpoint{1.945811in}{4.014976in}}%
\pgfusepath{stroke}%
\end{pgfscope}%
\begin{pgfscope}%
\pgfsetbuttcap%
\pgfsetroundjoin%
\pgfsetlinewidth{0.501875pt}%
\definecolor{currentstroke}{rgb}{0.501961,0.501961,0.501961}%
\pgfsetstrokecolor{currentstroke}%
\pgfsetdash{{1.850000pt}{0.800000pt}}{0.000000pt}%
\pgfpathmoveto{\pgfqpoint{0.814870in}{1.052131in}}%
\pgfpathlineto{\pgfqpoint{2.196011in}{2.247665in}}%
\pgfpathlineto{\pgfqpoint{2.192391in}{3.952476in}}%
\pgfusepath{stroke}%
\end{pgfscope}%
\begin{pgfscope}%
\pgfsetbuttcap%
\pgfsetroundjoin%
\pgfsetlinewidth{0.501875pt}%
\definecolor{currentstroke}{rgb}{0.501961,0.501961,0.501961}%
\pgfsetstrokecolor{currentstroke}%
\pgfsetdash{{1.850000pt}{0.800000pt}}{0.000000pt}%
\pgfpathmoveto{\pgfqpoint{1.067553in}{0.969605in}}%
\pgfpathlineto{\pgfqpoint{2.434972in}{2.178208in}}%
\pgfpathlineto{\pgfqpoint{2.441654in}{3.889295in}}%
\pgfusepath{stroke}%
\end{pgfscope}%
\begin{pgfscope}%
\pgfsetbuttcap%
\pgfsetroundjoin%
\pgfsetlinewidth{0.501875pt}%
\definecolor{currentstroke}{rgb}{0.501961,0.501961,0.501961}%
\pgfsetstrokecolor{currentstroke}%
\pgfsetdash{{1.850000pt}{0.800000pt}}{0.000000pt}%
\pgfpathmoveto{\pgfqpoint{1.323126in}{0.886135in}}%
\pgfpathlineto{\pgfqpoint{2.676440in}{2.108023in}}%
\pgfpathlineto{\pgfqpoint{2.693645in}{3.825422in}}%
\pgfusepath{stroke}%
\end{pgfscope}%
\begin{pgfscope}%
\pgfsetbuttcap%
\pgfsetroundjoin%
\pgfsetlinewidth{0.501875pt}%
\definecolor{currentstroke}{rgb}{0.501961,0.501961,0.501961}%
\pgfsetstrokecolor{currentstroke}%
\pgfsetdash{{1.850000pt}{0.800000pt}}{0.000000pt}%
\pgfpathmoveto{\pgfqpoint{1.581638in}{0.801705in}}%
\pgfpathlineto{\pgfqpoint{2.920453in}{2.037098in}}%
\pgfpathlineto{\pgfqpoint{2.948408in}{3.760847in}}%
\pgfusepath{stroke}%
\end{pgfscope}%
\begin{pgfscope}%
\pgfsetbuttcap%
\pgfsetroundjoin%
\pgfsetlinewidth{0.501875pt}%
\definecolor{currentstroke}{rgb}{0.501961,0.501961,0.501961}%
\pgfsetstrokecolor{currentstroke}%
\pgfsetdash{{1.850000pt}{0.800000pt}}{0.000000pt}%
\pgfpathmoveto{\pgfqpoint{1.843142in}{0.716297in}}%
\pgfpathlineto{\pgfqpoint{3.167051in}{1.965422in}}%
\pgfpathlineto{\pgfqpoint{3.205989in}{3.695558in}}%
\pgfusepath{stroke}%
\end{pgfscope}%
\begin{pgfscope}%
\pgfsetbuttcap%
\pgfsetroundjoin%
\pgfsetlinewidth{0.501875pt}%
\definecolor{currentstroke}{rgb}{0.501961,0.501961,0.501961}%
\pgfsetstrokecolor{currentstroke}%
\pgfsetdash{{1.850000pt}{0.800000pt}}{0.000000pt}%
\pgfpathmoveto{\pgfqpoint{2.107690in}{0.629896in}}%
\pgfpathlineto{\pgfqpoint{3.416277in}{1.892981in}}%
\pgfpathlineto{\pgfqpoint{3.466436in}{3.629542in}}%
\pgfusepath{stroke}%
\end{pgfscope}%
\begin{pgfscope}%
\pgfsetbuttcap%
\pgfsetroundjoin%
\pgfsetlinewidth{0.501875pt}%
\definecolor{currentstroke}{rgb}{0.501961,0.501961,0.501961}%
\pgfsetstrokecolor{currentstroke}%
\pgfsetdash{{1.850000pt}{0.800000pt}}{0.000000pt}%
\pgfpathmoveto{\pgfqpoint{2.375334in}{0.542484in}}%
\pgfpathlineto{\pgfqpoint{3.668172in}{1.819765in}}%
\pgfpathlineto{\pgfqpoint{3.729796in}{3.562788in}}%
\pgfusepath{stroke}%
\end{pgfscope}%
\begin{pgfscope}%
\pgfsetbuttcap%
\pgfsetroundjoin%
\pgfsetlinewidth{0.501875pt}%
\definecolor{currentstroke}{rgb}{0.501961,0.501961,0.501961}%
\pgfsetstrokecolor{currentstroke}%
\pgfsetdash{{1.850000pt}{0.800000pt}}{0.000000pt}%
\pgfpathmoveto{\pgfqpoint{2.646129in}{0.454042in}}%
\pgfpathlineto{\pgfqpoint{3.922780in}{1.745761in}}%
\pgfpathlineto{\pgfqpoint{3.996118in}{3.495283in}}%
\pgfusepath{stroke}%
\end{pgfscope}%
\begin{pgfscope}%
\pgfsetbuttcap%
\pgfsetroundjoin%
\pgfsetlinewidth{0.501875pt}%
\definecolor{currentstroke}{rgb}{0.501961,0.501961,0.501961}%
\pgfsetstrokecolor{currentstroke}%
\pgfsetdash{{1.850000pt}{0.800000pt}}{0.000000pt}%
\pgfpathmoveto{\pgfqpoint{0.435584in}{3.053101in}}%
\pgfpathlineto{\pgfqpoint{0.517191in}{1.261879in}}%
\pgfpathlineto{\pgfqpoint{2.893219in}{0.491926in}}%
\pgfusepath{stroke}%
\end{pgfscope}%
\begin{pgfscope}%
\pgfsetbuttcap%
\pgfsetroundjoin%
\pgfsetlinewidth{0.501875pt}%
\definecolor{currentstroke}{rgb}{0.501961,0.501961,0.501961}%
\pgfsetstrokecolor{currentstroke}%
\pgfsetdash{{1.850000pt}{0.800000pt}}{0.000000pt}%
\pgfpathmoveto{\pgfqpoint{0.606998in}{3.178191in}}%
\pgfpathlineto{\pgfqpoint{0.680315in}{1.398612in}}%
\pgfpathlineto{\pgfqpoint{3.041111in}{0.643699in}}%
\pgfusepath{stroke}%
\end{pgfscope}%
\begin{pgfscope}%
\pgfsetbuttcap%
\pgfsetroundjoin%
\pgfsetlinewidth{0.501875pt}%
\definecolor{currentstroke}{rgb}{0.501961,0.501961,0.501961}%
\pgfsetstrokecolor{currentstroke}%
\pgfsetdash{{1.850000pt}{0.800000pt}}{0.000000pt}%
\pgfpathmoveto{\pgfqpoint{0.775019in}{3.300805in}}%
\pgfpathlineto{\pgfqpoint{0.840349in}{1.532756in}}%
\pgfpathlineto{\pgfqpoint{3.186053in}{0.792445in}}%
\pgfusepath{stroke}%
\end{pgfscope}%
\begin{pgfscope}%
\pgfsetbuttcap%
\pgfsetroundjoin%
\pgfsetlinewidth{0.501875pt}%
\definecolor{currentstroke}{rgb}{0.501961,0.501961,0.501961}%
\pgfsetstrokecolor{currentstroke}%
\pgfsetdash{{1.850000pt}{0.800000pt}}{0.000000pt}%
\pgfpathmoveto{\pgfqpoint{0.939746in}{3.421015in}}%
\pgfpathlineto{\pgfqpoint{0.997381in}{1.664383in}}%
\pgfpathlineto{\pgfqpoint{3.328135in}{0.938256in}}%
\pgfusepath{stroke}%
\end{pgfscope}%
\begin{pgfscope}%
\pgfsetbuttcap%
\pgfsetroundjoin%
\pgfsetlinewidth{0.501875pt}%
\definecolor{currentstroke}{rgb}{0.501961,0.501961,0.501961}%
\pgfsetstrokecolor{currentstroke}%
\pgfsetdash{{1.850000pt}{0.800000pt}}{0.000000pt}%
\pgfpathmoveto{\pgfqpoint{1.101275in}{3.538891in}}%
\pgfpathlineto{\pgfqpoint{1.151494in}{1.793564in}}%
\pgfpathlineto{\pgfqpoint{3.467439in}{1.081216in}}%
\pgfusepath{stroke}%
\end{pgfscope}%
\begin{pgfscope}%
\pgfsetbuttcap%
\pgfsetroundjoin%
\pgfsetlinewidth{0.501875pt}%
\definecolor{currentstroke}{rgb}{0.501961,0.501961,0.501961}%
\pgfsetstrokecolor{currentstroke}%
\pgfsetdash{{1.850000pt}{0.800000pt}}{0.000000pt}%
\pgfpathmoveto{\pgfqpoint{1.259698in}{3.654501in}}%
\pgfpathlineto{\pgfqpoint{1.302768in}{1.920365in}}%
\pgfpathlineto{\pgfqpoint{3.604046in}{1.221408in}}%
\pgfusepath{stroke}%
\end{pgfscope}%
\begin{pgfscope}%
\pgfsetbuttcap%
\pgfsetroundjoin%
\pgfsetlinewidth{0.501875pt}%
\definecolor{currentstroke}{rgb}{0.501961,0.501961,0.501961}%
\pgfsetstrokecolor{currentstroke}%
\pgfsetdash{{1.850000pt}{0.800000pt}}{0.000000pt}%
\pgfpathmoveto{\pgfqpoint{1.415104in}{3.767910in}}%
\pgfpathlineto{\pgfqpoint{1.451283in}{2.044853in}}%
\pgfpathlineto{\pgfqpoint{3.738034in}{1.358913in}}%
\pgfusepath{stroke}%
\end{pgfscope}%
\begin{pgfscope}%
\pgfsetbuttcap%
\pgfsetroundjoin%
\pgfsetlinewidth{0.501875pt}%
\definecolor{currentstroke}{rgb}{0.501961,0.501961,0.501961}%
\pgfsetstrokecolor{currentstroke}%
\pgfsetdash{{1.850000pt}{0.800000pt}}{0.000000pt}%
\pgfpathmoveto{\pgfqpoint{1.567579in}{3.879179in}}%
\pgfpathlineto{\pgfqpoint{1.597112in}{2.167090in}}%
\pgfpathlineto{\pgfqpoint{3.869477in}{1.493806in}}%
\pgfusepath{stroke}%
\end{pgfscope}%
\begin{pgfscope}%
\pgfsetbuttcap%
\pgfsetroundjoin%
\pgfsetlinewidth{0.501875pt}%
\definecolor{currentstroke}{rgb}{0.501961,0.501961,0.501961}%
\pgfsetstrokecolor{currentstroke}%
\pgfsetdash{{1.850000pt}{0.800000pt}}{0.000000pt}%
\pgfpathmoveto{\pgfqpoint{1.717204in}{3.988368in}}%
\pgfpathlineto{\pgfqpoint{1.740328in}{2.287136in}}%
\pgfpathlineto{\pgfqpoint{3.998447in}{1.626161in}}%
\pgfusepath{stroke}%
\end{pgfscope}%
\begin{pgfscope}%
\pgfsetbuttcap%
\pgfsetroundjoin%
\pgfsetlinewidth{0.501875pt}%
\definecolor{currentstroke}{rgb}{0.501961,0.501961,0.501961}%
\pgfsetstrokecolor{currentstroke}%
\pgfsetdash{{1.850000pt}{0.800000pt}}{0.000000pt}%
\pgfpathmoveto{\pgfqpoint{4.074103in}{1.735816in}}%
\pgfpathlineto{\pgfqpoint{1.822316in}{2.388715in}}%
\pgfpathlineto{\pgfqpoint{0.418931in}{1.215296in}}%
\pgfusepath{stroke}%
\end{pgfscope}%
\begin{pgfscope}%
\pgfsetbuttcap%
\pgfsetroundjoin%
\pgfsetlinewidth{0.501875pt}%
\definecolor{currentstroke}{rgb}{0.501961,0.501961,0.501961}%
\pgfsetstrokecolor{currentstroke}%
\pgfsetdash{{1.850000pt}{0.800000pt}}{0.000000pt}%
\pgfpathmoveto{\pgfqpoint{4.085943in}{1.994355in}}%
\pgfpathlineto{\pgfqpoint{1.819437in}{2.639075in}}%
\pgfpathlineto{\pgfqpoint{0.406181in}{1.480067in}}%
\pgfusepath{stroke}%
\end{pgfscope}%
\begin{pgfscope}%
\pgfsetbuttcap%
\pgfsetroundjoin%
\pgfsetlinewidth{0.501875pt}%
\definecolor{currentstroke}{rgb}{0.501961,0.501961,0.501961}%
\pgfsetstrokecolor{currentstroke}%
\pgfsetdash{{1.850000pt}{0.800000pt}}{0.000000pt}%
\pgfpathmoveto{\pgfqpoint{4.097939in}{2.256329in}}%
\pgfpathlineto{\pgfqpoint{1.816522in}{2.892602in}}%
\pgfpathlineto{\pgfqpoint{0.393254in}{1.748490in}}%
\pgfusepath{stroke}%
\end{pgfscope}%
\begin{pgfscope}%
\pgfsetbuttcap%
\pgfsetroundjoin%
\pgfsetlinewidth{0.501875pt}%
\definecolor{currentstroke}{rgb}{0.501961,0.501961,0.501961}%
\pgfsetstrokecolor{currentstroke}%
\pgfsetdash{{1.850000pt}{0.800000pt}}{0.000000pt}%
\pgfpathmoveto{\pgfqpoint{4.110096in}{2.521807in}}%
\pgfpathlineto{\pgfqpoint{1.813569in}{3.149358in}}%
\pgfpathlineto{\pgfqpoint{0.380149in}{2.020642in}}%
\pgfusepath{stroke}%
\end{pgfscope}%
\begin{pgfscope}%
\pgfsetbuttcap%
\pgfsetroundjoin%
\pgfsetlinewidth{0.501875pt}%
\definecolor{currentstroke}{rgb}{0.501961,0.501961,0.501961}%
\pgfsetstrokecolor{currentstroke}%
\pgfsetdash{{1.850000pt}{0.800000pt}}{0.000000pt}%
\pgfpathmoveto{\pgfqpoint{4.122417in}{2.790860in}}%
\pgfpathlineto{\pgfqpoint{1.810579in}{3.409403in}}%
\pgfpathlineto{\pgfqpoint{0.366859in}{2.296601in}}%
\pgfusepath{stroke}%
\end{pgfscope}%
\begin{pgfscope}%
\pgfsetbuttcap%
\pgfsetroundjoin%
\pgfsetlinewidth{0.501875pt}%
\definecolor{currentstroke}{rgb}{0.501961,0.501961,0.501961}%
\pgfsetstrokecolor{currentstroke}%
\pgfsetdash{{1.850000pt}{0.800000pt}}{0.000000pt}%
\pgfpathmoveto{\pgfqpoint{4.134905in}{3.063561in}}%
\pgfpathlineto{\pgfqpoint{1.807550in}{3.672803in}}%
\pgfpathlineto{\pgfqpoint{0.353383in}{2.576447in}}%
\pgfusepath{stroke}%
\end{pgfscope}%
\begin{pgfscope}%
\pgfsetbuttcap%
\pgfsetroundjoin%
\pgfsetlinewidth{0.501875pt}%
\definecolor{currentstroke}{rgb}{0.501961,0.501961,0.501961}%
\pgfsetstrokecolor{currentstroke}%
\pgfsetdash{{1.850000pt}{0.800000pt}}{0.000000pt}%
\pgfpathmoveto{\pgfqpoint{4.147563in}{3.339984in}}%
\pgfpathlineto{\pgfqpoint{1.804482in}{3.939622in}}%
\pgfpathlineto{\pgfqpoint{0.339715in}{2.860262in}}%
\pgfusepath{stroke}%
\end{pgfscope}%
\begin{pgfscope}%
\pgfsetrectcap%
\pgfsetroundjoin%
\pgfsetlinewidth{0.803000pt}%
\definecolor{currentstroke}{rgb}{0.000000,0.000000,0.000000}%
\pgfsetstrokecolor{currentstroke}%
\pgfsetdash{}{0pt}%
\pgfpathmoveto{\pgfqpoint{0.420587in}{1.180904in}}%
\pgfpathlineto{\pgfqpoint{2.805565in}{0.401970in}}%
\pgfusepath{stroke}%
\end{pgfscope}%
\begin{pgfscope}%
\pgfpathrectangle{\pgfqpoint{0.100000in}{0.134025in}}{\pgfqpoint{4.245803in}{4.245803in}}%
\pgfusepath{clip}%
\pgfsetbuttcap%
\pgfsetroundjoin%
\pgfsetlinewidth{0.501875pt}%
\definecolor{currentstroke}{rgb}{0.501961,0.501961,0.501961}%
\pgfsetstrokecolor{currentstroke}%
\pgfsetstrokeopacity{0.500000}%
\pgfsetdash{{1.850000pt}{0.800000pt}}{0.000000pt}%
\pgfpathmoveto{\pgfqpoint{2.280277in}{2.314302in}}%
\pgfusepath{stroke}%
\end{pgfscope}%
\begin{pgfscope}%
\pgfsetrectcap%
\pgfsetroundjoin%
\pgfsetlinewidth{0.803000pt}%
\definecolor{currentstroke}{rgb}{0.000000,0.000000,0.000000}%
\pgfsetstrokecolor{currentstroke}%
\pgfsetdash{}{0pt}%
\pgfpathmoveto{\pgfqpoint{0.577172in}{1.144028in}}%
\pgfpathlineto{\pgfqpoint{0.540691in}{1.113088in}}%
\pgfusepath{stroke}%
\end{pgfscope}%
\begin{pgfscope}%
\definecolor{textcolor}{rgb}{0.000000,0.000000,0.000000}%
\pgfsetstrokecolor{textcolor}%
\pgfsetfillcolor{textcolor}%
\pgftext[x=0.460509in,y=0.915630in,,top]{\color{textcolor}{\rmfamily\fontsize{10.000000}{12.000000}\selectfont\catcode`\^=\active\def^{\ifmmode\sp\else\^{}\fi}\catcode`\%=\active\def%{\%}$\mathdefault{\ensuremath{-}1.00}$}}%
\end{pgfscope}%
\begin{pgfscope}%
\pgfpathrectangle{\pgfqpoint{0.100000in}{0.134025in}}{\pgfqpoint{4.245803in}{4.245803in}}%
\pgfusepath{clip}%
\pgfsetbuttcap%
\pgfsetroundjoin%
\pgfsetlinewidth{0.501875pt}%
\definecolor{currentstroke}{rgb}{0.501961,0.501961,0.501961}%
\pgfsetstrokecolor{currentstroke}%
\pgfsetstrokeopacity{0.500000}%
\pgfsetdash{{1.850000pt}{0.800000pt}}{0.000000pt}%
\pgfpathmoveto{\pgfqpoint{2.280277in}{2.314302in}}%
\pgfusepath{stroke}%
\end{pgfscope}%
\begin{pgfscope}%
\pgfsetrectcap%
\pgfsetroundjoin%
\pgfsetlinewidth{0.803000pt}%
\definecolor{currentstroke}{rgb}{0.000000,0.000000,0.000000}%
\pgfsetstrokecolor{currentstroke}%
\pgfsetdash{}{0pt}%
\pgfpathmoveto{\pgfqpoint{0.826903in}{1.062546in}}%
\pgfpathlineto{\pgfqpoint{0.790754in}{1.031255in}}%
\pgfusepath{stroke}%
\end{pgfscope}%
\begin{pgfscope}%
\definecolor{textcolor}{rgb}{0.000000,0.000000,0.000000}%
\pgfsetstrokecolor{textcolor}%
\pgfsetfillcolor{textcolor}%
\pgftext[x=0.710570in,y=0.832529in,,top]{\color{textcolor}{\rmfamily\fontsize{10.000000}{12.000000}\selectfont\catcode`\^=\active\def^{\ifmmode\sp\else\^{}\fi}\catcode`\%=\active\def%{\%}$\mathdefault{\ensuremath{-}0.75}$}}%
\end{pgfscope}%
\begin{pgfscope}%
\pgfpathrectangle{\pgfqpoint{0.100000in}{0.134025in}}{\pgfqpoint{4.245803in}{4.245803in}}%
\pgfusepath{clip}%
\pgfsetbuttcap%
\pgfsetroundjoin%
\pgfsetlinewidth{0.501875pt}%
\definecolor{currentstroke}{rgb}{0.501961,0.501961,0.501961}%
\pgfsetstrokecolor{currentstroke}%
\pgfsetstrokeopacity{0.500000}%
\pgfsetdash{{1.850000pt}{0.800000pt}}{0.000000pt}%
\pgfpathmoveto{\pgfqpoint{2.280277in}{2.314302in}}%
\pgfusepath{stroke}%
\end{pgfscope}%
\begin{pgfscope}%
\pgfsetrectcap%
\pgfsetroundjoin%
\pgfsetlinewidth{0.803000pt}%
\definecolor{currentstroke}{rgb}{0.000000,0.000000,0.000000}%
\pgfsetstrokecolor{currentstroke}%
\pgfsetdash{}{0pt}%
\pgfpathmoveto{\pgfqpoint{1.079471in}{0.980139in}}%
\pgfpathlineto{\pgfqpoint{1.043664in}{0.948491in}}%
\pgfusepath{stroke}%
\end{pgfscope}%
\begin{pgfscope}%
\definecolor{textcolor}{rgb}{0.000000,0.000000,0.000000}%
\pgfsetstrokecolor{textcolor}%
\pgfsetfillcolor{textcolor}%
\pgftext[x=0.963484in,y=0.748480in,,top]{\color{textcolor}{\rmfamily\fontsize{10.000000}{12.000000}\selectfont\catcode`\^=\active\def^{\ifmmode\sp\else\^{}\fi}\catcode`\%=\active\def%{\%}$\mathdefault{\ensuremath{-}0.50}$}}%
\end{pgfscope}%
\begin{pgfscope}%
\pgfpathrectangle{\pgfqpoint{0.100000in}{0.134025in}}{\pgfqpoint{4.245803in}{4.245803in}}%
\pgfusepath{clip}%
\pgfsetbuttcap%
\pgfsetroundjoin%
\pgfsetlinewidth{0.501875pt}%
\definecolor{currentstroke}{rgb}{0.501961,0.501961,0.501961}%
\pgfsetstrokecolor{currentstroke}%
\pgfsetstrokeopacity{0.500000}%
\pgfsetdash{{1.850000pt}{0.800000pt}}{0.000000pt}%
\pgfpathmoveto{\pgfqpoint{2.280277in}{2.314302in}}%
\pgfusepath{stroke}%
\end{pgfscope}%
\begin{pgfscope}%
\pgfsetrectcap%
\pgfsetroundjoin%
\pgfsetlinewidth{0.803000pt}%
\definecolor{currentstroke}{rgb}{0.000000,0.000000,0.000000}%
\pgfsetstrokecolor{currentstroke}%
\pgfsetdash{}{0pt}%
\pgfpathmoveto{\pgfqpoint{1.334927in}{0.896790in}}%
\pgfpathlineto{\pgfqpoint{1.299472in}{0.864778in}}%
\pgfusepath{stroke}%
\end{pgfscope}%
\begin{pgfscope}%
\definecolor{textcolor}{rgb}{0.000000,0.000000,0.000000}%
\pgfsetstrokecolor{textcolor}%
\pgfsetfillcolor{textcolor}%
\pgftext[x=1.219301in,y=0.663467in,,top]{\color{textcolor}{\rmfamily\fontsize{10.000000}{12.000000}\selectfont\catcode`\^=\active\def^{\ifmmode\sp\else\^{}\fi}\catcode`\%=\active\def%{\%}$\mathdefault{\ensuremath{-}0.25}$}}%
\end{pgfscope}%
\begin{pgfscope}%
\pgfpathrectangle{\pgfqpoint{0.100000in}{0.134025in}}{\pgfqpoint{4.245803in}{4.245803in}}%
\pgfusepath{clip}%
\pgfsetbuttcap%
\pgfsetroundjoin%
\pgfsetlinewidth{0.501875pt}%
\definecolor{currentstroke}{rgb}{0.501961,0.501961,0.501961}%
\pgfsetstrokecolor{currentstroke}%
\pgfsetstrokeopacity{0.500000}%
\pgfsetdash{{1.850000pt}{0.800000pt}}{0.000000pt}%
\pgfpathmoveto{\pgfqpoint{2.280277in}{2.314302in}}%
\pgfusepath{stroke}%
\end{pgfscope}%
\begin{pgfscope}%
\pgfsetrectcap%
\pgfsetroundjoin%
\pgfsetlinewidth{0.803000pt}%
\definecolor{currentstroke}{rgb}{0.000000,0.000000,0.000000}%
\pgfsetstrokecolor{currentstroke}%
\pgfsetdash{}{0pt}%
\pgfpathmoveto{\pgfqpoint{1.593318in}{0.812482in}}%
\pgfpathlineto{\pgfqpoint{1.558227in}{0.780102in}}%
\pgfusepath{stroke}%
\end{pgfscope}%
\begin{pgfscope}%
\definecolor{textcolor}{rgb}{0.000000,0.000000,0.000000}%
\pgfsetstrokecolor{textcolor}%
\pgfsetfillcolor{textcolor}%
\pgftext[x=1.478070in,y=0.577472in,,top]{\color{textcolor}{\rmfamily\fontsize{10.000000}{12.000000}\selectfont\catcode`\^=\active\def^{\ifmmode\sp\else\^{}\fi}\catcode`\%=\active\def%{\%}$\mathdefault{0.00}$}}%
\end{pgfscope}%
\begin{pgfscope}%
\pgfpathrectangle{\pgfqpoint{0.100000in}{0.134025in}}{\pgfqpoint{4.245803in}{4.245803in}}%
\pgfusepath{clip}%
\pgfsetbuttcap%
\pgfsetroundjoin%
\pgfsetlinewidth{0.501875pt}%
\definecolor{currentstroke}{rgb}{0.501961,0.501961,0.501961}%
\pgfsetstrokecolor{currentstroke}%
\pgfsetstrokeopacity{0.500000}%
\pgfsetdash{{1.850000pt}{0.800000pt}}{0.000000pt}%
\pgfpathmoveto{\pgfqpoint{2.280277in}{2.314302in}}%
\pgfusepath{stroke}%
\end{pgfscope}%
\begin{pgfscope}%
\pgfsetrectcap%
\pgfsetroundjoin%
\pgfsetlinewidth{0.803000pt}%
\definecolor{currentstroke}{rgb}{0.000000,0.000000,0.000000}%
\pgfsetstrokecolor{currentstroke}%
\pgfsetdash{}{0pt}%
\pgfpathmoveto{\pgfqpoint{1.854698in}{0.727200in}}%
\pgfpathlineto{\pgfqpoint{1.819980in}{0.694444in}}%
\pgfusepath{stroke}%
\end{pgfscope}%
\begin{pgfscope}%
\definecolor{textcolor}{rgb}{0.000000,0.000000,0.000000}%
\pgfsetstrokecolor{textcolor}%
\pgfsetfillcolor{textcolor}%
\pgftext[x=1.739843in,y=0.490479in,,top]{\color{textcolor}{\rmfamily\fontsize{10.000000}{12.000000}\selectfont\catcode`\^=\active\def^{\ifmmode\sp\else\^{}\fi}\catcode`\%=\active\def%{\%}$\mathdefault{0.25}$}}%
\end{pgfscope}%
\begin{pgfscope}%
\pgfpathrectangle{\pgfqpoint{0.100000in}{0.134025in}}{\pgfqpoint{4.245803in}{4.245803in}}%
\pgfusepath{clip}%
\pgfsetbuttcap%
\pgfsetroundjoin%
\pgfsetlinewidth{0.501875pt}%
\definecolor{currentstroke}{rgb}{0.501961,0.501961,0.501961}%
\pgfsetstrokecolor{currentstroke}%
\pgfsetstrokeopacity{0.500000}%
\pgfsetdash{{1.850000pt}{0.800000pt}}{0.000000pt}%
\pgfpathmoveto{\pgfqpoint{2.280277in}{2.314302in}}%
\pgfusepath{stroke}%
\end{pgfscope}%
\begin{pgfscope}%
\pgfsetrectcap%
\pgfsetroundjoin%
\pgfsetlinewidth{0.803000pt}%
\definecolor{currentstroke}{rgb}{0.000000,0.000000,0.000000}%
\pgfsetstrokecolor{currentstroke}%
\pgfsetdash{}{0pt}%
\pgfpathmoveto{\pgfqpoint{2.119117in}{0.640926in}}%
\pgfpathlineto{\pgfqpoint{2.084784in}{0.607787in}}%
\pgfusepath{stroke}%
\end{pgfscope}%
\begin{pgfscope}%
\definecolor{textcolor}{rgb}{0.000000,0.000000,0.000000}%
\pgfsetstrokecolor{textcolor}%
\pgfsetfillcolor{textcolor}%
\pgftext[x=2.004673in,y=0.402470in,,top]{\color{textcolor}{\rmfamily\fontsize{10.000000}{12.000000}\selectfont\catcode`\^=\active\def^{\ifmmode\sp\else\^{}\fi}\catcode`\%=\active\def%{\%}$\mathdefault{0.50}$}}%
\end{pgfscope}%
\begin{pgfscope}%
\pgfpathrectangle{\pgfqpoint{0.100000in}{0.134025in}}{\pgfqpoint{4.245803in}{4.245803in}}%
\pgfusepath{clip}%
\pgfsetbuttcap%
\pgfsetroundjoin%
\pgfsetlinewidth{0.501875pt}%
\definecolor{currentstroke}{rgb}{0.501961,0.501961,0.501961}%
\pgfsetstrokecolor{currentstroke}%
\pgfsetstrokeopacity{0.500000}%
\pgfsetdash{{1.850000pt}{0.800000pt}}{0.000000pt}%
\pgfpathmoveto{\pgfqpoint{2.280277in}{2.314302in}}%
\pgfusepath{stroke}%
\end{pgfscope}%
\begin{pgfscope}%
\pgfsetrectcap%
\pgfsetroundjoin%
\pgfsetlinewidth{0.803000pt}%
\definecolor{currentstroke}{rgb}{0.000000,0.000000,0.000000}%
\pgfsetstrokecolor{currentstroke}%
\pgfsetdash{}{0pt}%
\pgfpathmoveto{\pgfqpoint{2.386629in}{0.553643in}}%
\pgfpathlineto{\pgfqpoint{2.352692in}{0.520115in}}%
\pgfusepath{stroke}%
\end{pgfscope}%
\begin{pgfscope}%
\definecolor{textcolor}{rgb}{0.000000,0.000000,0.000000}%
\pgfsetstrokecolor{textcolor}%
\pgfsetfillcolor{textcolor}%
\pgftext[x=2.272613in,y=0.313428in,,top]{\color{textcolor}{\rmfamily\fontsize{10.000000}{12.000000}\selectfont\catcode`\^=\active\def^{\ifmmode\sp\else\^{}\fi}\catcode`\%=\active\def%{\%}$\mathdefault{0.75}$}}%
\end{pgfscope}%
\begin{pgfscope}%
\pgfpathrectangle{\pgfqpoint{0.100000in}{0.134025in}}{\pgfqpoint{4.245803in}{4.245803in}}%
\pgfusepath{clip}%
\pgfsetbuttcap%
\pgfsetroundjoin%
\pgfsetlinewidth{0.501875pt}%
\definecolor{currentstroke}{rgb}{0.501961,0.501961,0.501961}%
\pgfsetstrokecolor{currentstroke}%
\pgfsetstrokeopacity{0.500000}%
\pgfsetdash{{1.850000pt}{0.800000pt}}{0.000000pt}%
\pgfpathmoveto{\pgfqpoint{2.280277in}{2.314302in}}%
\pgfusepath{stroke}%
\end{pgfscope}%
\begin{pgfscope}%
\pgfsetrectcap%
\pgfsetroundjoin%
\pgfsetlinewidth{0.803000pt}%
\definecolor{currentstroke}{rgb}{0.000000,0.000000,0.000000}%
\pgfsetstrokecolor{currentstroke}%
\pgfsetdash{}{0pt}%
\pgfpathmoveto{\pgfqpoint{2.657289in}{0.465333in}}%
\pgfpathlineto{\pgfqpoint{2.623760in}{0.431409in}}%
\pgfusepath{stroke}%
\end{pgfscope}%
\begin{pgfscope}%
\definecolor{textcolor}{rgb}{0.000000,0.000000,0.000000}%
\pgfsetstrokecolor{textcolor}%
\pgfsetfillcolor{textcolor}%
\pgftext[x=2.543719in,y=0.223333in,,top]{\color{textcolor}{\rmfamily\fontsize{10.000000}{12.000000}\selectfont\catcode`\^=\active\def^{\ifmmode\sp\else\^{}\fi}\catcode`\%=\active\def%{\%}$\mathdefault{1.00}$}}%
\end{pgfscope}%
\begin{pgfscope}%
\definecolor{textcolor}{rgb}{0.000000,0.000000,0.000000}%
\pgfsetstrokecolor{textcolor}%
\pgfsetfillcolor{textcolor}%
\pgftext[x=1.355611in,y=0.312341in,,,rotate=341.912962]{\color{textcolor}{\rmfamily\fontsize{10.000000}{12.000000}\selectfont\catcode`\^=\active\def^{\ifmmode\sp\else\^{}\fi}\catcode`\%=\active\def%{\%}Eje X}}%
\end{pgfscope}%
\begin{pgfscope}%
\pgfsetrectcap%
\pgfsetroundjoin%
\pgfsetlinewidth{0.803000pt}%
\definecolor{currentstroke}{rgb}{0.000000,0.000000,0.000000}%
\pgfsetstrokecolor{currentstroke}%
\pgfsetdash{}{0pt}%
\pgfpathmoveto{\pgfqpoint{4.072565in}{1.702224in}}%
\pgfpathlineto{\pgfqpoint{2.805565in}{0.401970in}}%
\pgfusepath{stroke}%
\end{pgfscope}%
\begin{pgfscope}%
\pgfpathrectangle{\pgfqpoint{0.100000in}{0.134025in}}{\pgfqpoint{4.245803in}{4.245803in}}%
\pgfusepath{clip}%
\pgfsetbuttcap%
\pgfsetroundjoin%
\pgfsetlinewidth{0.501875pt}%
\definecolor{currentstroke}{rgb}{0.501961,0.501961,0.501961}%
\pgfsetstrokecolor{currentstroke}%
\pgfsetstrokeopacity{0.500000}%
\pgfsetdash{{1.850000pt}{0.800000pt}}{0.000000pt}%
\pgfpathmoveto{\pgfqpoint{2.280277in}{2.314302in}}%
\pgfusepath{stroke}%
\end{pgfscope}%
\begin{pgfscope}%
\pgfsetrectcap%
\pgfsetroundjoin%
\pgfsetlinewidth{0.803000pt}%
\definecolor{currentstroke}{rgb}{0.000000,0.000000,0.000000}%
\pgfsetstrokecolor{currentstroke}%
\pgfsetdash{}{0pt}%
\pgfpathmoveto{\pgfqpoint{2.873196in}{0.498414in}}%
\pgfpathlineto{\pgfqpoint{2.933317in}{0.478932in}}%
\pgfusepath{stroke}%
\end{pgfscope}%
\begin{pgfscope}%
\definecolor{textcolor}{rgb}{0.000000,0.000000,0.000000}%
\pgfsetstrokecolor{textcolor}%
\pgfsetfillcolor{textcolor}%
\pgftext[x=3.072103in,y=0.305505in,,top]{\color{textcolor}{\rmfamily\fontsize{10.000000}{12.000000}\selectfont\catcode`\^=\active\def^{\ifmmode\sp\else\^{}\fi}\catcode`\%=\active\def%{\%}$\mathdefault{\ensuremath{-}1.00}$}}%
\end{pgfscope}%
\begin{pgfscope}%
\pgfpathrectangle{\pgfqpoint{0.100000in}{0.134025in}}{\pgfqpoint{4.245803in}{4.245803in}}%
\pgfusepath{clip}%
\pgfsetbuttcap%
\pgfsetroundjoin%
\pgfsetlinewidth{0.501875pt}%
\definecolor{currentstroke}{rgb}{0.501961,0.501961,0.501961}%
\pgfsetstrokecolor{currentstroke}%
\pgfsetstrokeopacity{0.500000}%
\pgfsetdash{{1.850000pt}{0.800000pt}}{0.000000pt}%
\pgfpathmoveto{\pgfqpoint{2.280277in}{2.314302in}}%
\pgfusepath{stroke}%
\end{pgfscope}%
\begin{pgfscope}%
\pgfsetrectcap%
\pgfsetroundjoin%
\pgfsetlinewidth{0.803000pt}%
\definecolor{currentstroke}{rgb}{0.000000,0.000000,0.000000}%
\pgfsetstrokecolor{currentstroke}%
\pgfsetdash{}{0pt}%
\pgfpathmoveto{\pgfqpoint{3.021226in}{0.650057in}}%
\pgfpathlineto{\pgfqpoint{3.080930in}{0.630965in}}%
\pgfusepath{stroke}%
\end{pgfscope}%
\begin{pgfscope}%
\definecolor{textcolor}{rgb}{0.000000,0.000000,0.000000}%
\pgfsetstrokecolor{textcolor}%
\pgfsetfillcolor{textcolor}%
\pgftext[x=3.218269in,y=0.459235in,,top]{\color{textcolor}{\rmfamily\fontsize{10.000000}{12.000000}\selectfont\catcode`\^=\active\def^{\ifmmode\sp\else\^{}\fi}\catcode`\%=\active\def%{\%}$\mathdefault{\ensuremath{-}0.75}$}}%
\end{pgfscope}%
\begin{pgfscope}%
\pgfpathrectangle{\pgfqpoint{0.100000in}{0.134025in}}{\pgfqpoint{4.245803in}{4.245803in}}%
\pgfusepath{clip}%
\pgfsetbuttcap%
\pgfsetroundjoin%
\pgfsetlinewidth{0.501875pt}%
\definecolor{currentstroke}{rgb}{0.501961,0.501961,0.501961}%
\pgfsetstrokecolor{currentstroke}%
\pgfsetstrokeopacity{0.500000}%
\pgfsetdash{{1.850000pt}{0.800000pt}}{0.000000pt}%
\pgfpathmoveto{\pgfqpoint{2.280277in}{2.314302in}}%
\pgfusepath{stroke}%
\end{pgfscope}%
\begin{pgfscope}%
\pgfsetrectcap%
\pgfsetroundjoin%
\pgfsetlinewidth{0.803000pt}%
\definecolor{currentstroke}{rgb}{0.000000,0.000000,0.000000}%
\pgfsetstrokecolor{currentstroke}%
\pgfsetdash{}{0pt}%
\pgfpathmoveto{\pgfqpoint{3.166306in}{0.798678in}}%
\pgfpathlineto{\pgfqpoint{3.225598in}{0.779965in}}%
\pgfusepath{stroke}%
\end{pgfscope}%
\begin{pgfscope}%
\definecolor{textcolor}{rgb}{0.000000,0.000000,0.000000}%
\pgfsetstrokecolor{textcolor}%
\pgfsetfillcolor{textcolor}%
\pgftext[x=3.361518in,y=0.609897in,,top]{\color{textcolor}{\rmfamily\fontsize{10.000000}{12.000000}\selectfont\catcode`\^=\active\def^{\ifmmode\sp\else\^{}\fi}\catcode`\%=\active\def%{\%}$\mathdefault{\ensuremath{-}0.50}$}}%
\end{pgfscope}%
\begin{pgfscope}%
\pgfpathrectangle{\pgfqpoint{0.100000in}{0.134025in}}{\pgfqpoint{4.245803in}{4.245803in}}%
\pgfusepath{clip}%
\pgfsetbuttcap%
\pgfsetroundjoin%
\pgfsetlinewidth{0.501875pt}%
\definecolor{currentstroke}{rgb}{0.501961,0.501961,0.501961}%
\pgfsetstrokecolor{currentstroke}%
\pgfsetstrokeopacity{0.500000}%
\pgfsetdash{{1.850000pt}{0.800000pt}}{0.000000pt}%
\pgfpathmoveto{\pgfqpoint{2.280277in}{2.314302in}}%
\pgfusepath{stroke}%
\end{pgfscope}%
\begin{pgfscope}%
\pgfsetrectcap%
\pgfsetroundjoin%
\pgfsetlinewidth{0.803000pt}%
\definecolor{currentstroke}{rgb}{0.000000,0.000000,0.000000}%
\pgfsetstrokecolor{currentstroke}%
\pgfsetdash{}{0pt}%
\pgfpathmoveto{\pgfqpoint{3.308523in}{0.944366in}}%
\pgfpathlineto{\pgfqpoint{3.367408in}{0.926021in}}%
\pgfusepath{stroke}%
\end{pgfscope}%
\begin{pgfscope}%
\definecolor{textcolor}{rgb}{0.000000,0.000000,0.000000}%
\pgfsetstrokecolor{textcolor}%
\pgfsetfillcolor{textcolor}%
\pgftext[x=3.501939in,y=0.757584in,,top]{\color{textcolor}{\rmfamily\fontsize{10.000000}{12.000000}\selectfont\catcode`\^=\active\def^{\ifmmode\sp\else\^{}\fi}\catcode`\%=\active\def%{\%}$\mathdefault{\ensuremath{-}0.25}$}}%
\end{pgfscope}%
\begin{pgfscope}%
\pgfpathrectangle{\pgfqpoint{0.100000in}{0.134025in}}{\pgfqpoint{4.245803in}{4.245803in}}%
\pgfusepath{clip}%
\pgfsetbuttcap%
\pgfsetroundjoin%
\pgfsetlinewidth{0.501875pt}%
\definecolor{currentstroke}{rgb}{0.501961,0.501961,0.501961}%
\pgfsetstrokecolor{currentstroke}%
\pgfsetstrokeopacity{0.500000}%
\pgfsetdash{{1.850000pt}{0.800000pt}}{0.000000pt}%
\pgfpathmoveto{\pgfqpoint{2.280277in}{2.314302in}}%
\pgfusepath{stroke}%
\end{pgfscope}%
\begin{pgfscope}%
\pgfsetrectcap%
\pgfsetroundjoin%
\pgfsetlinewidth{0.803000pt}%
\definecolor{currentstroke}{rgb}{0.000000,0.000000,0.000000}%
\pgfsetstrokecolor{currentstroke}%
\pgfsetdash{}{0pt}%
\pgfpathmoveto{\pgfqpoint{3.447961in}{1.087207in}}%
\pgfpathlineto{\pgfqpoint{3.506443in}{1.069219in}}%
\pgfusepath{stroke}%
\end{pgfscope}%
\begin{pgfscope}%
\definecolor{textcolor}{rgb}{0.000000,0.000000,0.000000}%
\pgfsetstrokecolor{textcolor}%
\pgfsetfillcolor{textcolor}%
\pgftext[x=3.639613in,y=0.902383in,,top]{\color{textcolor}{\rmfamily\fontsize{10.000000}{12.000000}\selectfont\catcode`\^=\active\def^{\ifmmode\sp\else\^{}\fi}\catcode`\%=\active\def%{\%}$\mathdefault{0.00}$}}%
\end{pgfscope}%
\begin{pgfscope}%
\pgfpathrectangle{\pgfqpoint{0.100000in}{0.134025in}}{\pgfqpoint{4.245803in}{4.245803in}}%
\pgfusepath{clip}%
\pgfsetbuttcap%
\pgfsetroundjoin%
\pgfsetlinewidth{0.501875pt}%
\definecolor{currentstroke}{rgb}{0.501961,0.501961,0.501961}%
\pgfsetstrokecolor{currentstroke}%
\pgfsetstrokeopacity{0.500000}%
\pgfsetdash{{1.850000pt}{0.800000pt}}{0.000000pt}%
\pgfpathmoveto{\pgfqpoint{2.280277in}{2.314302in}}%
\pgfusepath{stroke}%
\end{pgfscope}%
\begin{pgfscope}%
\pgfsetrectcap%
\pgfsetroundjoin%
\pgfsetlinewidth{0.803000pt}%
\definecolor{currentstroke}{rgb}{0.000000,0.000000,0.000000}%
\pgfsetstrokecolor{currentstroke}%
\pgfsetdash{}{0pt}%
\pgfpathmoveto{\pgfqpoint{3.584700in}{1.227284in}}%
\pgfpathlineto{\pgfqpoint{3.642784in}{1.209642in}}%
\pgfusepath{stroke}%
\end{pgfscope}%
\begin{pgfscope}%
\definecolor{textcolor}{rgb}{0.000000,0.000000,0.000000}%
\pgfsetstrokecolor{textcolor}%
\pgfsetfillcolor{textcolor}%
\pgftext[x=3.774620in,y=1.044377in,,top]{\color{textcolor}{\rmfamily\fontsize{10.000000}{12.000000}\selectfont\catcode`\^=\active\def^{\ifmmode\sp\else\^{}\fi}\catcode`\%=\active\def%{\%}$\mathdefault{0.25}$}}%
\end{pgfscope}%
\begin{pgfscope}%
\pgfpathrectangle{\pgfqpoint{0.100000in}{0.134025in}}{\pgfqpoint{4.245803in}{4.245803in}}%
\pgfusepath{clip}%
\pgfsetbuttcap%
\pgfsetroundjoin%
\pgfsetlinewidth{0.501875pt}%
\definecolor{currentstroke}{rgb}{0.501961,0.501961,0.501961}%
\pgfsetstrokecolor{currentstroke}%
\pgfsetstrokeopacity{0.500000}%
\pgfsetdash{{1.850000pt}{0.800000pt}}{0.000000pt}%
\pgfpathmoveto{\pgfqpoint{2.280277in}{2.314302in}}%
\pgfusepath{stroke}%
\end{pgfscope}%
\begin{pgfscope}%
\pgfsetrectcap%
\pgfsetroundjoin%
\pgfsetlinewidth{0.803000pt}%
\definecolor{currentstroke}{rgb}{0.000000,0.000000,0.000000}%
\pgfsetstrokecolor{currentstroke}%
\pgfsetdash{}{0pt}%
\pgfpathmoveto{\pgfqpoint{3.718819in}{1.364676in}}%
\pgfpathlineto{\pgfqpoint{3.776509in}{1.347372in}}%
\pgfusepath{stroke}%
\end{pgfscope}%
\begin{pgfscope}%
\definecolor{textcolor}{rgb}{0.000000,0.000000,0.000000}%
\pgfsetstrokecolor{textcolor}%
\pgfsetfillcolor{textcolor}%
\pgftext[x=3.907038in,y=1.183647in,,top]{\color{textcolor}{\rmfamily\fontsize{10.000000}{12.000000}\selectfont\catcode`\^=\active\def^{\ifmmode\sp\else\^{}\fi}\catcode`\%=\active\def%{\%}$\mathdefault{0.50}$}}%
\end{pgfscope}%
\begin{pgfscope}%
\pgfpathrectangle{\pgfqpoint{0.100000in}{0.134025in}}{\pgfqpoint{4.245803in}{4.245803in}}%
\pgfusepath{clip}%
\pgfsetbuttcap%
\pgfsetroundjoin%
\pgfsetlinewidth{0.501875pt}%
\definecolor{currentstroke}{rgb}{0.501961,0.501961,0.501961}%
\pgfsetstrokecolor{currentstroke}%
\pgfsetstrokeopacity{0.500000}%
\pgfsetdash{{1.850000pt}{0.800000pt}}{0.000000pt}%
\pgfpathmoveto{\pgfqpoint{2.280277in}{2.314302in}}%
\pgfusepath{stroke}%
\end{pgfscope}%
\begin{pgfscope}%
\pgfsetrectcap%
\pgfsetroundjoin%
\pgfsetlinewidth{0.803000pt}%
\definecolor{currentstroke}{rgb}{0.000000,0.000000,0.000000}%
\pgfsetstrokecolor{currentstroke}%
\pgfsetdash{}{0pt}%
\pgfpathmoveto{\pgfqpoint{3.850392in}{1.499461in}}%
\pgfpathlineto{\pgfqpoint{3.907693in}{1.482483in}}%
\pgfusepath{stroke}%
\end{pgfscope}%
\begin{pgfscope}%
\definecolor{textcolor}{rgb}{0.000000,0.000000,0.000000}%
\pgfsetstrokecolor{textcolor}%
\pgfsetfillcolor{textcolor}%
\pgftext[x=4.036940in,y=1.320270in,,top]{\color{textcolor}{\rmfamily\fontsize{10.000000}{12.000000}\selectfont\catcode`\^=\active\def^{\ifmmode\sp\else\^{}\fi}\catcode`\%=\active\def%{\%}$\mathdefault{0.75}$}}%
\end{pgfscope}%
\begin{pgfscope}%
\pgfpathrectangle{\pgfqpoint{0.100000in}{0.134025in}}{\pgfqpoint{4.245803in}{4.245803in}}%
\pgfusepath{clip}%
\pgfsetbuttcap%
\pgfsetroundjoin%
\pgfsetlinewidth{0.501875pt}%
\definecolor{currentstroke}{rgb}{0.501961,0.501961,0.501961}%
\pgfsetstrokecolor{currentstroke}%
\pgfsetstrokeopacity{0.500000}%
\pgfsetdash{{1.850000pt}{0.800000pt}}{0.000000pt}%
\pgfpathmoveto{\pgfqpoint{2.280277in}{2.314302in}}%
\pgfusepath{stroke}%
\end{pgfscope}%
\begin{pgfscope}%
\pgfsetrectcap%
\pgfsetroundjoin%
\pgfsetlinewidth{0.803000pt}%
\definecolor{currentstroke}{rgb}{0.000000,0.000000,0.000000}%
\pgfsetstrokecolor{currentstroke}%
\pgfsetdash{}{0pt}%
\pgfpathmoveto{\pgfqpoint{3.979491in}{1.631710in}}%
\pgfpathlineto{\pgfqpoint{4.036406in}{1.615051in}}%
\pgfusepath{stroke}%
\end{pgfscope}%
\begin{pgfscope}%
\definecolor{textcolor}{rgb}{0.000000,0.000000,0.000000}%
\pgfsetstrokecolor{textcolor}%
\pgfsetfillcolor{textcolor}%
\pgftext[x=4.164396in,y=1.454323in,,top]{\color{textcolor}{\rmfamily\fontsize{10.000000}{12.000000}\selectfont\catcode`\^=\active\def^{\ifmmode\sp\else\^{}\fi}\catcode`\%=\active\def%{\%}$\mathdefault{1.00}$}}%
\end{pgfscope}%
\begin{pgfscope}%
\definecolor{textcolor}{rgb}{0.000000,0.000000,0.000000}%
\pgfsetstrokecolor{textcolor}%
\pgfsetfillcolor{textcolor}%
\pgftext[x=3.841956in,y=0.692225in,,,rotate=45.742112]{\color{textcolor}{\rmfamily\fontsize{10.000000}{12.000000}\selectfont\catcode`\^=\active\def^{\ifmmode\sp\else\^{}\fi}\catcode`\%=\active\def%{\%}Eje Y}}%
\end{pgfscope}%
\begin{pgfscope}%
\pgfsetrectcap%
\pgfsetroundjoin%
\pgfsetlinewidth{0.803000pt}%
\definecolor{currentstroke}{rgb}{0.000000,0.000000,0.000000}%
\pgfsetstrokecolor{currentstroke}%
\pgfsetdash{}{0pt}%
\pgfpathmoveto{\pgfqpoint{4.072565in}{1.702224in}}%
\pgfpathlineto{\pgfqpoint{4.152855in}{3.455555in}}%
\pgfusepath{stroke}%
\end{pgfscope}%
\begin{pgfscope}%
\pgfsetbuttcap%
\pgfsetroundjoin%
\pgfsetlinewidth{0.501875pt}%
\definecolor{currentstroke}{rgb}{0.501961,0.501961,0.501961}%
\pgfsetstrokecolor{currentstroke}%
\pgfsetstrokeopacity{0.500000}%
\pgfsetdash{{1.850000pt}{0.800000pt}}{0.000000pt}%
\pgfpathmoveto{\pgfqpoint{2.280277in}{2.314302in}}%
\pgfusepath{stroke}%
\end{pgfscope}%
\begin{pgfscope}%
\pgfsetrectcap%
\pgfsetroundjoin%
\pgfsetlinewidth{0.803000pt}%
\definecolor{currentstroke}{rgb}{0.000000,0.000000,0.000000}%
\pgfsetstrokecolor{currentstroke}%
\pgfsetdash{}{0pt}%
\pgfpathmoveto{\pgfqpoint{4.055204in}{1.741296in}}%
\pgfpathlineto{\pgfqpoint{4.111947in}{1.724844in}}%
\pgfusepath{stroke}%
\end{pgfscope}%
\begin{pgfscope}%
\definecolor{textcolor}{rgb}{0.000000,0.000000,0.000000}%
\pgfsetstrokecolor{textcolor}%
\pgfsetfillcolor{textcolor}%
\pgftext[x=4.327868in,y=1.771954in,,top]{\color{textcolor}{\rmfamily\fontsize{10.000000}{12.000000}\selectfont\catcode`\^=\active\def^{\ifmmode\sp\else\^{}\fi}\catcode`\%=\active\def%{\%}$\mathdefault{0}$}}%
\end{pgfscope}%
\begin{pgfscope}%
\pgfsetbuttcap%
\pgfsetroundjoin%
\pgfsetlinewidth{0.501875pt}%
\definecolor{currentstroke}{rgb}{0.501961,0.501961,0.501961}%
\pgfsetstrokecolor{currentstroke}%
\pgfsetstrokeopacity{0.500000}%
\pgfsetdash{{1.850000pt}{0.800000pt}}{0.000000pt}%
\pgfpathmoveto{\pgfqpoint{2.280277in}{2.314302in}}%
\pgfusepath{stroke}%
\end{pgfscope}%
\begin{pgfscope}%
\pgfsetrectcap%
\pgfsetroundjoin%
\pgfsetlinewidth{0.803000pt}%
\definecolor{currentstroke}{rgb}{0.000000,0.000000,0.000000}%
\pgfsetstrokecolor{currentstroke}%
\pgfsetdash{}{0pt}%
\pgfpathmoveto{\pgfqpoint{4.066914in}{1.999768in}}%
\pgfpathlineto{\pgfqpoint{4.124046in}{1.983516in}}%
\pgfusepath{stroke}%
\end{pgfscope}%
\begin{pgfscope}%
\definecolor{textcolor}{rgb}{0.000000,0.000000,0.000000}%
\pgfsetstrokecolor{textcolor}%
\pgfsetfillcolor{textcolor}%
\pgftext[x=4.341347in,y=2.030050in,,top]{\color{textcolor}{\rmfamily\fontsize{10.000000}{12.000000}\selectfont\catcode`\^=\active\def^{\ifmmode\sp\else\^{}\fi}\catcode`\%=\active\def%{\%}$\mathdefault{2}$}}%
\end{pgfscope}%
\begin{pgfscope}%
\pgfsetbuttcap%
\pgfsetroundjoin%
\pgfsetlinewidth{0.501875pt}%
\definecolor{currentstroke}{rgb}{0.501961,0.501961,0.501961}%
\pgfsetstrokecolor{currentstroke}%
\pgfsetstrokeopacity{0.500000}%
\pgfsetdash{{1.850000pt}{0.800000pt}}{0.000000pt}%
\pgfpathmoveto{\pgfqpoint{2.280277in}{2.314302in}}%
\pgfusepath{stroke}%
\end{pgfscope}%
\begin{pgfscope}%
\pgfsetrectcap%
\pgfsetroundjoin%
\pgfsetlinewidth{0.803000pt}%
\definecolor{currentstroke}{rgb}{0.000000,0.000000,0.000000}%
\pgfsetstrokecolor{currentstroke}%
\pgfsetdash{}{0pt}%
\pgfpathmoveto{\pgfqpoint{4.078779in}{2.261672in}}%
\pgfpathlineto{\pgfqpoint{4.136305in}{2.245628in}}%
\pgfusepath{stroke}%
\end{pgfscope}%
\begin{pgfscope}%
\definecolor{textcolor}{rgb}{0.000000,0.000000,0.000000}%
\pgfsetstrokecolor{textcolor}%
\pgfsetfillcolor{textcolor}%
\pgftext[x=4.355004in,y=2.291567in,,top]{\color{textcolor}{\rmfamily\fontsize{10.000000}{12.000000}\selectfont\catcode`\^=\active\def^{\ifmmode\sp\else\^{}\fi}\catcode`\%=\active\def%{\%}$\mathdefault{4}$}}%
\end{pgfscope}%
\begin{pgfscope}%
\pgfsetbuttcap%
\pgfsetroundjoin%
\pgfsetlinewidth{0.501875pt}%
\definecolor{currentstroke}{rgb}{0.501961,0.501961,0.501961}%
\pgfsetstrokecolor{currentstroke}%
\pgfsetstrokeopacity{0.500000}%
\pgfsetdash{{1.850000pt}{0.800000pt}}{0.000000pt}%
\pgfpathmoveto{\pgfqpoint{2.280277in}{2.314302in}}%
\pgfusepath{stroke}%
\end{pgfscope}%
\begin{pgfscope}%
\pgfsetrectcap%
\pgfsetroundjoin%
\pgfsetlinewidth{0.803000pt}%
\definecolor{currentstroke}{rgb}{0.000000,0.000000,0.000000}%
\pgfsetstrokecolor{currentstroke}%
\pgfsetdash{}{0pt}%
\pgfpathmoveto{\pgfqpoint{4.090803in}{2.527079in}}%
\pgfpathlineto{\pgfqpoint{4.148729in}{2.511250in}}%
\pgfusepath{stroke}%
\end{pgfscope}%
\begin{pgfscope}%
\definecolor{textcolor}{rgb}{0.000000,0.000000,0.000000}%
\pgfsetstrokecolor{textcolor}%
\pgfsetfillcolor{textcolor}%
\pgftext[x=4.368843in,y=2.556573in,,top]{\color{textcolor}{\rmfamily\fontsize{10.000000}{12.000000}\selectfont\catcode`\^=\active\def^{\ifmmode\sp\else\^{}\fi}\catcode`\%=\active\def%{\%}$\mathdefault{6}$}}%
\end{pgfscope}%
\begin{pgfscope}%
\pgfsetbuttcap%
\pgfsetroundjoin%
\pgfsetlinewidth{0.501875pt}%
\definecolor{currentstroke}{rgb}{0.501961,0.501961,0.501961}%
\pgfsetstrokecolor{currentstroke}%
\pgfsetstrokeopacity{0.500000}%
\pgfsetdash{{1.850000pt}{0.800000pt}}{0.000000pt}%
\pgfpathmoveto{\pgfqpoint{2.280277in}{2.314302in}}%
\pgfusepath{stroke}%
\end{pgfscope}%
\begin{pgfscope}%
\pgfsetrectcap%
\pgfsetroundjoin%
\pgfsetlinewidth{0.803000pt}%
\definecolor{currentstroke}{rgb}{0.000000,0.000000,0.000000}%
\pgfsetstrokecolor{currentstroke}%
\pgfsetdash{}{0pt}%
\pgfpathmoveto{\pgfqpoint{4.102989in}{2.796058in}}%
\pgfpathlineto{\pgfqpoint{4.161320in}{2.780451in}}%
\pgfusepath{stroke}%
\end{pgfscope}%
\begin{pgfscope}%
\definecolor{textcolor}{rgb}{0.000000,0.000000,0.000000}%
\pgfsetstrokecolor{textcolor}%
\pgfsetfillcolor{textcolor}%
\pgftext[x=4.382868in,y=2.825137in,,top]{\color{textcolor}{\rmfamily\fontsize{10.000000}{12.000000}\selectfont\catcode`\^=\active\def^{\ifmmode\sp\else\^{}\fi}\catcode`\%=\active\def%{\%}$\mathdefault{8}$}}%
\end{pgfscope}%
\begin{pgfscope}%
\pgfsetbuttcap%
\pgfsetroundjoin%
\pgfsetlinewidth{0.501875pt}%
\definecolor{currentstroke}{rgb}{0.501961,0.501961,0.501961}%
\pgfsetstrokecolor{currentstroke}%
\pgfsetstrokeopacity{0.500000}%
\pgfsetdash{{1.850000pt}{0.800000pt}}{0.000000pt}%
\pgfpathmoveto{\pgfqpoint{2.280277in}{2.314302in}}%
\pgfusepath{stroke}%
\end{pgfscope}%
\begin{pgfscope}%
\pgfsetrectcap%
\pgfsetroundjoin%
\pgfsetlinewidth{0.803000pt}%
\definecolor{currentstroke}{rgb}{0.000000,0.000000,0.000000}%
\pgfsetstrokecolor{currentstroke}%
\pgfsetdash{}{0pt}%
\pgfpathmoveto{\pgfqpoint{4.115340in}{3.068682in}}%
\pgfpathlineto{\pgfqpoint{4.174082in}{3.053305in}}%
\pgfusepath{stroke}%
\end{pgfscope}%
\begin{pgfscope}%
\definecolor{textcolor}{rgb}{0.000000,0.000000,0.000000}%
\pgfsetstrokecolor{textcolor}%
\pgfsetfillcolor{textcolor}%
\pgftext[x=4.397083in,y=3.097333in,,top]{\color{textcolor}{\rmfamily\fontsize{10.000000}{12.000000}\selectfont\catcode`\^=\active\def^{\ifmmode\sp\else\^{}\fi}\catcode`\%=\active\def%{\%}$\mathdefault{10}$}}%
\end{pgfscope}%
\begin{pgfscope}%
\pgfsetbuttcap%
\pgfsetroundjoin%
\pgfsetlinewidth{0.501875pt}%
\definecolor{currentstroke}{rgb}{0.501961,0.501961,0.501961}%
\pgfsetstrokecolor{currentstroke}%
\pgfsetstrokeopacity{0.500000}%
\pgfsetdash{{1.850000pt}{0.800000pt}}{0.000000pt}%
\pgfpathmoveto{\pgfqpoint{2.280277in}{2.314302in}}%
\pgfusepath{stroke}%
\end{pgfscope}%
\begin{pgfscope}%
\pgfsetrectcap%
\pgfsetroundjoin%
\pgfsetlinewidth{0.803000pt}%
\definecolor{currentstroke}{rgb}{0.000000,0.000000,0.000000}%
\pgfsetstrokecolor{currentstroke}%
\pgfsetdash{}{0pt}%
\pgfpathmoveto{\pgfqpoint{4.127859in}{3.345027in}}%
\pgfpathlineto{\pgfqpoint{4.187018in}{3.329887in}}%
\pgfusepath{stroke}%
\end{pgfscope}%
\begin{pgfscope}%
\definecolor{textcolor}{rgb}{0.000000,0.000000,0.000000}%
\pgfsetstrokecolor{textcolor}%
\pgfsetfillcolor{textcolor}%
\pgftext[x=4.411492in,y=3.373235in,,top]{\color{textcolor}{\rmfamily\fontsize{10.000000}{12.000000}\selectfont\catcode`\^=\active\def^{\ifmmode\sp\else\^{}\fi}\catcode`\%=\active\def%{\%}$\mathdefault{12}$}}%
\end{pgfscope}%
\begin{pgfscope}%
\definecolor{textcolor}{rgb}{0.000000,0.000000,0.000000}%
\pgfsetstrokecolor{textcolor}%
\pgfsetfillcolor{textcolor}%
\pgftext[x=4.673088in,y=2.634878in,,,rotate=87.378092]{\color{textcolor}{\rmfamily\fontsize{10.000000}{12.000000}\selectfont\catcode`\^=\active\def^{\ifmmode\sp\else\^{}\fi}\catcode`\%=\active\def%{\%}Eje Z (Tiempo/Profundidad)}}%
\end{pgfscope}%
\begin{pgfscope}%
\pgfpathrectangle{\pgfqpoint{0.100000in}{0.134025in}}{\pgfqpoint{4.245803in}{4.245803in}}%
\pgfusepath{clip}%
\pgfsetrectcap%
\pgfsetroundjoin%
\pgfsetlinewidth{2.007500pt}%
\definecolor{currentstroke}{rgb}{0.501961,0.000000,0.501961}%
\pgfsetstrokecolor{currentstroke}%
\pgfsetdash{}{0pt}%
\pgfpathmoveto{\pgfqpoint{2.842440in}{1.997851in}}%
\pgfpathlineto{\pgfqpoint{2.915567in}{1.984263in}}%
\pgfpathlineto{\pgfqpoint{2.962562in}{1.973652in}}%
\pgfpathlineto{\pgfqpoint{3.008020in}{1.961839in}}%
\pgfpathlineto{\pgfqpoint{3.051830in}{1.948862in}}%
\pgfpathlineto{\pgfqpoint{3.093885in}{1.934763in}}%
\pgfpathlineto{\pgfqpoint{3.134077in}{1.919585in}}%
\pgfpathlineto{\pgfqpoint{3.172306in}{1.903372in}}%
\pgfpathlineto{\pgfqpoint{3.208471in}{1.886175in}}%
\pgfpathlineto{\pgfqpoint{3.242475in}{1.868044in}}%
\pgfpathlineto{\pgfqpoint{3.274227in}{1.849032in}}%
\pgfpathlineto{\pgfqpoint{3.303637in}{1.829197in}}%
\pgfpathlineto{\pgfqpoint{3.330621in}{1.808596in}}%
\pgfpathlineto{\pgfqpoint{3.355100in}{1.787291in}}%
\pgfpathlineto{\pgfqpoint{3.376997in}{1.765345in}}%
\pgfpathlineto{\pgfqpoint{3.396243in}{1.742825in}}%
\pgfpathlineto{\pgfqpoint{3.412773in}{1.719797in}}%
\pgfpathlineto{\pgfqpoint{3.426528in}{1.696332in}}%
\pgfpathlineto{\pgfqpoint{3.437454in}{1.672501in}}%
\pgfpathlineto{\pgfqpoint{3.445506in}{1.648379in}}%
\pgfpathlineto{\pgfqpoint{3.450642in}{1.624040in}}%
\pgfpathlineto{\pgfqpoint{3.452829in}{1.599561in}}%
\pgfpathlineto{\pgfqpoint{3.452041in}{1.575021in}}%
\pgfpathlineto{\pgfqpoint{3.448259in}{1.550498in}}%
\pgfpathlineto{\pgfqpoint{3.441471in}{1.526074in}}%
\pgfpathlineto{\pgfqpoint{3.431675in}{1.501829in}}%
\pgfpathlineto{\pgfqpoint{3.418874in}{1.477844in}}%
\pgfpathlineto{\pgfqpoint{3.403082in}{1.454203in}}%
\pgfpathlineto{\pgfqpoint{3.384321in}{1.430986in}}%
\pgfpathlineto{\pgfqpoint{3.362620in}{1.408275in}}%
\pgfpathlineto{\pgfqpoint{3.338018in}{1.386153in}}%
\pgfpathlineto{\pgfqpoint{3.310564in}{1.364699in}}%
\pgfpathlineto{\pgfqpoint{3.280312in}{1.343993in}}%
\pgfpathlineto{\pgfqpoint{3.247329in}{1.324114in}}%
\pgfpathlineto{\pgfqpoint{3.211689in}{1.305138in}}%
\pgfpathlineto{\pgfqpoint{3.173474in}{1.287140in}}%
\pgfpathlineto{\pgfqpoint{3.132777in}{1.270193in}}%
\pgfpathlineto{\pgfqpoint{3.089697in}{1.254367in}}%
\pgfpathlineto{\pgfqpoint{3.044343in}{1.239730in}}%
\pgfpathlineto{\pgfqpoint{2.996830in}{1.226345in}}%
\pgfpathlineto{\pgfqpoint{2.947284in}{1.214273in}}%
\pgfpathlineto{\pgfqpoint{2.895835in}{1.203573in}}%
\pgfpathlineto{\pgfqpoint{2.842622in}{1.194296in}}%
\pgfpathlineto{\pgfqpoint{2.787789in}{1.186494in}}%
\pgfpathlineto{\pgfqpoint{2.731489in}{1.180210in}}%
\pgfpathlineto{\pgfqpoint{2.673877in}{1.175484in}}%
\pgfpathlineto{\pgfqpoint{2.615116in}{1.172353in}}%
\pgfpathlineto{\pgfqpoint{2.555371in}{1.170847in}}%
\pgfpathlineto{\pgfqpoint{2.494813in}{1.170991in}}%
\pgfpathlineto{\pgfqpoint{2.433614in}{1.172805in}}%
\pgfpathlineto{\pgfqpoint{2.371952in}{1.176304in}}%
\pgfpathlineto{\pgfqpoint{2.310002in}{1.181498in}}%
\pgfpathlineto{\pgfqpoint{2.247944in}{1.188390in}}%
\pgfpathlineto{\pgfqpoint{2.185958in}{1.196979in}}%
\pgfpathlineto{\pgfqpoint{2.124223in}{1.207257in}}%
\pgfpathlineto{\pgfqpoint{2.062916in}{1.219213in}}%
\pgfpathlineto{\pgfqpoint{2.002214in}{1.232827in}}%
\pgfpathlineto{\pgfqpoint{1.942292in}{1.248077in}}%
\pgfpathlineto{\pgfqpoint{1.883320in}{1.264934in}}%
\pgfpathlineto{\pgfqpoint{1.825466in}{1.283365in}}%
\pgfpathlineto{\pgfqpoint{1.768895in}{1.303331in}}%
\pgfpathlineto{\pgfqpoint{1.713763in}{1.324790in}}%
\pgfpathlineto{\pgfqpoint{1.660226in}{1.347693in}}%
\pgfpathlineto{\pgfqpoint{1.608429in}{1.371989in}}%
\pgfpathlineto{\pgfqpoint{1.558513in}{1.397622in}}%
\pgfpathlineto{\pgfqpoint{1.510613in}{1.424533in}}%
\pgfpathlineto{\pgfqpoint{1.464856in}{1.452658in}}%
\pgfpathlineto{\pgfqpoint{1.421360in}{1.481932in}}%
\pgfpathlineto{\pgfqpoint{1.380236in}{1.512284in}}%
\pgfpathlineto{\pgfqpoint{1.341588in}{1.543645in}}%
\pgfpathlineto{\pgfqpoint{1.305511in}{1.575939in}}%
\pgfpathlineto{\pgfqpoint{1.272089in}{1.609091in}}%
\pgfpathlineto{\pgfqpoint{1.241401in}{1.643023in}}%
\pgfpathlineto{\pgfqpoint{1.213514in}{1.677656in}}%
\pgfpathlineto{\pgfqpoint{1.188488in}{1.712910in}}%
\pgfpathlineto{\pgfqpoint{1.166374in}{1.748705in}}%
\pgfpathlineto{\pgfqpoint{1.147214in}{1.784958in}}%
\pgfpathlineto{\pgfqpoint{1.131040in}{1.821588in}}%
\pgfpathlineto{\pgfqpoint{1.117877in}{1.858514in}}%
\pgfpathlineto{\pgfqpoint{1.107740in}{1.895655in}}%
\pgfpathlineto{\pgfqpoint{1.100637in}{1.932928in}}%
\pgfpathlineto{\pgfqpoint{1.096567in}{1.970256in}}%
\pgfpathlineto{\pgfqpoint{1.095521in}{2.007558in}}%
\pgfpathlineto{\pgfqpoint{1.097483in}{2.044756in}}%
\pgfpathlineto{\pgfqpoint{1.102428in}{2.081775in}}%
\pgfpathlineto{\pgfqpoint{1.110324in}{2.118539in}}%
\pgfpathlineto{\pgfqpoint{1.121134in}{2.154975in}}%
\pgfpathlineto{\pgfqpoint{1.134812in}{2.191012in}}%
\pgfpathlineto{\pgfqpoint{1.151307in}{2.226580in}}%
\pgfpathlineto{\pgfqpoint{1.170561in}{2.261611in}}%
\pgfpathlineto{\pgfqpoint{1.192510in}{2.296041in}}%
\pgfpathlineto{\pgfqpoint{1.217087in}{2.329807in}}%
\pgfpathlineto{\pgfqpoint{1.244216in}{2.362849in}}%
\pgfpathlineto{\pgfqpoint{1.273818in}{2.395107in}}%
\pgfpathlineto{\pgfqpoint{1.305810in}{2.426528in}}%
\pgfpathlineto{\pgfqpoint{1.340103in}{2.457057in}}%
\pgfpathlineto{\pgfqpoint{1.376606in}{2.486645in}}%
\pgfpathlineto{\pgfqpoint{1.415222in}{2.515243in}}%
\pgfpathlineto{\pgfqpoint{1.455852in}{2.542808in}}%
\pgfpathlineto{\pgfqpoint{1.498393in}{2.569295in}}%
\pgfpathlineto{\pgfqpoint{1.542739in}{2.594667in}}%
\pgfpathlineto{\pgfqpoint{1.588783in}{2.618885in}}%
\pgfpathlineto{\pgfqpoint{1.636412in}{2.641917in}}%
\pgfpathlineto{\pgfqpoint{1.685513in}{2.663729in}}%
\pgfpathlineto{\pgfqpoint{1.735971in}{2.684293in}}%
\pgfpathlineto{\pgfqpoint{1.787669in}{2.703584in}}%
\pgfpathlineto{\pgfqpoint{1.840487in}{2.721578in}}%
\pgfpathlineto{\pgfqpoint{1.894304in}{2.738254in}}%
\pgfpathlineto{\pgfqpoint{1.948999in}{2.753594in}}%
\pgfpathlineto{\pgfqpoint{2.004448in}{2.767584in}}%
\pgfpathlineto{\pgfqpoint{2.088765in}{2.786007in}}%
\pgfpathlineto{\pgfqpoint{2.174077in}{2.801331in}}%
\pgfpathlineto{\pgfqpoint{2.259959in}{2.813538in}}%
\pgfpathlineto{\pgfqpoint{2.345985in}{2.822623in}}%
\pgfpathlineto{\pgfqpoint{2.431727in}{2.828594in}}%
\pgfpathlineto{\pgfqpoint{2.516757in}{2.831473in}}%
\pgfpathlineto{\pgfqpoint{2.600649in}{2.831295in}}%
\pgfpathlineto{\pgfqpoint{2.682983in}{2.828108in}}%
\pgfpathlineto{\pgfqpoint{2.763338in}{2.821973in}}%
\pgfpathlineto{\pgfqpoint{2.841302in}{2.812965in}}%
\pgfpathlineto{\pgfqpoint{2.916468in}{2.801172in}}%
\pgfpathlineto{\pgfqpoint{2.988439in}{2.786693in}}%
\pgfpathlineto{\pgfqpoint{3.056826in}{2.769644in}}%
\pgfpathlineto{\pgfqpoint{3.100239in}{2.756912in}}%
\pgfpathlineto{\pgfqpoint{3.141784in}{2.743135in}}%
\pgfpathlineto{\pgfqpoint{3.181353in}{2.728355in}}%
\pgfpathlineto{\pgfqpoint{3.218845in}{2.712619in}}%
\pgfpathlineto{\pgfqpoint{3.254161in}{2.695975in}}%
\pgfpathlineto{\pgfqpoint{3.287203in}{2.678474in}}%
\pgfpathlineto{\pgfqpoint{3.317882in}{2.660169in}}%
\pgfpathlineto{\pgfqpoint{3.346108in}{2.641116in}}%
\pgfpathlineto{\pgfqpoint{3.371800in}{2.621374in}}%
\pgfpathlineto{\pgfqpoint{3.394878in}{2.601001in}}%
\pgfpathlineto{\pgfqpoint{3.415269in}{2.580061in}}%
\pgfpathlineto{\pgfqpoint{3.432905in}{2.558618in}}%
\pgfpathlineto{\pgfqpoint{3.447724in}{2.536739in}}%
\pgfpathlineto{\pgfqpoint{3.459670in}{2.514492in}}%
\pgfpathlineto{\pgfqpoint{3.468692in}{2.491948in}}%
\pgfpathlineto{\pgfqpoint{3.474747in}{2.469178in}}%
\pgfpathlineto{\pgfqpoint{3.477799in}{2.446255in}}%
\pgfpathlineto{\pgfqpoint{3.477818in}{2.423254in}}%
\pgfpathlineto{\pgfqpoint{3.474782in}{2.400252in}}%
\pgfpathlineto{\pgfqpoint{3.468677in}{2.377325in}}%
\pgfpathlineto{\pgfqpoint{3.459497in}{2.354551in}}%
\pgfpathlineto{\pgfqpoint{3.447245in}{2.332007in}}%
\pgfpathlineto{\pgfqpoint{3.431930in}{2.309774in}}%
\pgfpathlineto{\pgfqpoint{3.413571in}{2.287929in}}%
\pgfpathlineto{\pgfqpoint{3.392198in}{2.266551in}}%
\pgfpathlineto{\pgfqpoint{3.367846in}{2.245719in}}%
\pgfpathlineto{\pgfqpoint{3.340562in}{2.225510in}}%
\pgfpathlineto{\pgfqpoint{3.310400in}{2.206002in}}%
\pgfpathlineto{\pgfqpoint{3.277425in}{2.187269in}}%
\pgfpathlineto{\pgfqpoint{3.241710in}{2.169386in}}%
\pgfpathlineto{\pgfqpoint{3.203337in}{2.152425in}}%
\pgfpathlineto{\pgfqpoint{3.162397in}{2.136458in}}%
\pgfpathlineto{\pgfqpoint{3.118989in}{2.121551in}}%
\pgfpathlineto{\pgfqpoint{3.073223in}{2.107771in}}%
\pgfpathlineto{\pgfqpoint{3.025216in}{2.095180in}}%
\pgfpathlineto{\pgfqpoint{2.975091in}{2.083838in}}%
\pgfpathlineto{\pgfqpoint{2.922982in}{2.073800in}}%
\pgfpathlineto{\pgfqpoint{2.869029in}{2.065119in}}%
\pgfpathlineto{\pgfqpoint{2.813378in}{2.057843in}}%
\pgfpathlineto{\pgfqpoint{2.756183in}{2.052017in}}%
\pgfpathlineto{\pgfqpoint{2.697602in}{2.047679in}}%
\pgfpathlineto{\pgfqpoint{2.637800in}{2.044867in}}%
\pgfpathlineto{\pgfqpoint{2.576947in}{2.043609in}}%
\pgfpathlineto{\pgfqpoint{2.515215in}{2.043931in}}%
\pgfpathlineto{\pgfqpoint{2.452781in}{2.045854in}}%
\pgfpathlineto{\pgfqpoint{2.389826in}{2.049394in}}%
\pgfpathlineto{\pgfqpoint{2.326531in}{2.054560in}}%
\pgfpathlineto{\pgfqpoint{2.263078in}{2.061356in}}%
\pgfpathlineto{\pgfqpoint{2.199652in}{2.069784in}}%
\pgfpathlineto{\pgfqpoint{2.136437in}{2.079836in}}%
\pgfpathlineto{\pgfqpoint{2.073616in}{2.091503in}}%
\pgfpathlineto{\pgfqpoint{2.011369in}{2.104766in}}%
\pgfpathlineto{\pgfqpoint{1.949877in}{2.119606in}}%
\pgfpathlineto{\pgfqpoint{1.889316in}{2.135995in}}%
\pgfpathlineto{\pgfqpoint{1.829858in}{2.153903in}}%
\pgfpathlineto{\pgfqpoint{1.771672in}{2.173292in}}%
\pgfpathlineto{\pgfqpoint{1.714923in}{2.194124in}}%
\pgfpathlineto{\pgfqpoint{1.659768in}{2.216351in}}%
\pgfpathlineto{\pgfqpoint{1.606359in}{2.239926in}}%
\pgfpathlineto{\pgfqpoint{1.554844in}{2.264796in}}%
\pgfpathlineto{\pgfqpoint{1.505359in}{2.290903in}}%
\pgfpathlineto{\pgfqpoint{1.458038in}{2.318187in}}%
\pgfpathlineto{\pgfqpoint{1.413004in}{2.346586in}}%
\pgfpathlineto{\pgfqpoint{1.370372in}{2.376034in}}%
\pgfpathlineto{\pgfqpoint{1.330251in}{2.406461in}}%
\pgfpathlineto{\pgfqpoint{1.292737in}{2.437798in}}%
\pgfpathlineto{\pgfqpoint{1.257922in}{2.469971in}}%
\pgfpathlineto{\pgfqpoint{1.225886in}{2.502906in}}%
\pgfpathlineto{\pgfqpoint{1.196701in}{2.536527in}}%
\pgfpathlineto{\pgfqpoint{1.170430in}{2.570758in}}%
\pgfpathlineto{\pgfqpoint{1.147126in}{2.605521in}}%
\pgfpathlineto{\pgfqpoint{1.126834in}{2.640737in}}%
\pgfpathlineto{\pgfqpoint{1.109589in}{2.676329in}}%
\pgfpathlineto{\pgfqpoint{1.095417in}{2.712218in}}%
\pgfpathlineto{\pgfqpoint{1.084337in}{2.748327in}}%
\pgfpathlineto{\pgfqpoint{1.076358in}{2.784577in}}%
\pgfpathlineto{\pgfqpoint{1.071479in}{2.820891in}}%
\pgfpathlineto{\pgfqpoint{1.069694in}{2.857194in}}%
\pgfpathlineto{\pgfqpoint{1.070987in}{2.893411in}}%
\pgfpathlineto{\pgfqpoint{1.075334in}{2.929469in}}%
\pgfpathlineto{\pgfqpoint{1.082705in}{2.965295in}}%
\pgfpathlineto{\pgfqpoint{1.093061in}{3.000818in}}%
\pgfpathlineto{\pgfqpoint{1.106358in}{3.035971in}}%
\pgfpathlineto{\pgfqpoint{1.122544in}{3.070687in}}%
\pgfpathlineto{\pgfqpoint{1.141562in}{3.104901in}}%
\pgfpathlineto{\pgfqpoint{1.163347in}{3.138550in}}%
\pgfpathlineto{\pgfqpoint{1.187830in}{3.171574in}}%
\pgfpathlineto{\pgfqpoint{1.214936in}{3.203914in}}%
\pgfpathlineto{\pgfqpoint{1.244586in}{3.235516in}}%
\pgfpathlineto{\pgfqpoint{1.276695in}{3.266326in}}%
\pgfpathlineto{\pgfqpoint{1.311173in}{3.296292in}}%
\pgfpathlineto{\pgfqpoint{1.347927in}{3.325366in}}%
\pgfpathlineto{\pgfqpoint{1.386860in}{3.353504in}}%
\pgfpathlineto{\pgfqpoint{1.427872in}{3.380660in}}%
\pgfpathlineto{\pgfqpoint{1.470858in}{3.406795in}}%
\pgfpathlineto{\pgfqpoint{1.515710in}{3.431871in}}%
\pgfpathlineto{\pgfqpoint{1.562320in}{3.455851in}}%
\pgfpathlineto{\pgfqpoint{1.610573in}{3.478704in}}%
\pgfpathlineto{\pgfqpoint{1.660355in}{3.500398in}}%
\pgfpathlineto{\pgfqpoint{1.711549in}{3.520906in}}%
\pgfpathlineto{\pgfqpoint{1.764036in}{3.540202in}}%
\pgfpathlineto{\pgfqpoint{1.817693in}{3.558265in}}%
\pgfpathlineto{\pgfqpoint{1.872400in}{3.575074in}}%
\pgfpathlineto{\pgfqpoint{1.956155in}{3.597898in}}%
\pgfpathlineto{\pgfqpoint{2.041568in}{3.617815in}}%
\pgfpathlineto{\pgfqpoint{2.128210in}{3.634787in}}%
\pgfpathlineto{\pgfqpoint{2.215651in}{3.648793in}}%
\pgfpathlineto{\pgfqpoint{2.303456in}{3.659823in}}%
\pgfpathlineto{\pgfqpoint{2.391190in}{3.667879in}}%
\pgfpathlineto{\pgfqpoint{2.478418in}{3.672976in}}%
\pgfpathlineto{\pgfqpoint{2.564706in}{3.675141in}}%
\pgfpathlineto{\pgfqpoint{2.649623in}{3.674416in}}%
\pgfpathlineto{\pgfqpoint{2.732739in}{3.670852in}}%
\pgfpathlineto{\pgfqpoint{2.813631in}{3.664514in}}%
\pgfpathlineto{\pgfqpoint{2.866117in}{3.658786in}}%
\pgfpathlineto{\pgfqpoint{2.866117in}{3.658786in}}%
\pgfusepath{stroke}%
\end{pgfscope}%
\begin{pgfscope}%
\definecolor{textcolor}{rgb}{0.000000,0.000000,0.000000}%
\pgfsetstrokecolor{textcolor}%
\pgfsetfillcolor{textcolor}%
\pgftext[x=2.222902in,y=4.463161in,,base]{\color{textcolor}{\rmfamily\fontsize{12.000000}{14.400000}\selectfont\catcode`\^=\active\def^{\ifmmode\sp\else\^{}\fi}\catcode`\%=\active\def%{\%}Trayectoria Helicoidal en 3D}}%
\end{pgfscope}%
\end{pgfpicture}%
\makeatother%
\endgroup%
 
    \caption{Gráfico 3D}
    \label{fig:grafico_Tridimensional}
\end{figure}

\subsubsection{Dispersión 3D (\textit{ax.scatter})}

El gráfico de dispersión en tres dimensiones se utiliza para identificar agrupamientos (\textit{clusters}), tendencias espaciales o correlaciones entre tres variables independientes. Al añadir un mapa de colores (\textit{colormap}), podemos visualizar una cuarta dimensión de datos de forma intuitiva.

\emph{¿Cuándo se usa?}
\begin{itemize}
\item Detección de valores atípicos (\textit{Outliers}): Es más fácil ver un punto "lejos de la nube" en un espacio 3D que en tablas planas.
\item Análisis de carteras: Por ejemplo, comparar Riesgo, Retorno y Plazo de diferentes activos financieros.
\item Ingeniería y logística: Para mapear coordenadas geográficas (X,Y) junto con una variable de altitud o costo (Z).
\end{itemize}

\begin{pycode}{Dispersión 3D Matemática}
import matplotlib.pyplot as plt
import numpy as np

# Generación de datos
rng = np.random.default_rng(42)
n = 600
x, y, z = rng.normal(0, 1, (3, n))
distancia = np.sqrt(x**2 + y**2 + z**2) # Cuarta dimensión (Color)

# Entorno 3D
fig, ax = plt.subplots(subplot_kw={"projection": "3d"}, figsize=(5, 5))

# Creación del gráfico
sc = ax.scatter(x, y, z, c=distancia, cmap="magma", s=10, alpha=0.8)

# Configuración de barra de colores
cb = fig.colorbar(sc, ax=ax, shrink=0.5, pad=0.1) 
cb.set_label("Intensidad / Distancia")

ax.set_title("Nube de Puntos Multivariable")
plt.show()
\end{pycode}

\emph{Análisis del código}
\begin{itemize}
\item \texttt{c=c}: Define el color de cada punto. En este caso, el color depende de la distancia radial (qué tan lejos está el punto del centro 0,0,0).
\item \texttt{cmap='mapa'}: Es el "mapa de colores". Algunos mapas son: \textit{Viridis}, \textit{plasma}, \textit{magma}, \textit{inferno} y \textit{cividis}.
\item \texttt{s=15}: (\textit{Size}) Define el tamaño de los puntos. En 3D, puntos muy grandes pueden solaparse y "tapar" la profundidad.
\item \texttt{alpha=0.9}: Controla la transparencia. Un poco de transparencia ayuda a ver puntos que están "detrás" de otros en la nube.
\item \texttt{fig.colorbar()}: Añade la leyenda lateral que explica qué significa cada color. Sin esto, el color sería solo decorativo y no informativo.
\item \texttt{ax.set\_box\_aspect((1, 1, 1))}: Fuerza a que los tres ejes tengan la misma escala visual (un cubo perfecto). Esto evita que el gráfico se vea "estirado" si un eje tiene valores más grandes que otros.
\end{itemize}

\begin{nota}
    Al usar \texttt{shrink} (escala) y \texttt{pad} (espaciado) en \texttt{colorbar}, ajustamos su tamaño y su distancia respecto al gráfico principal. \texttt{pad=0.1} Define la fracción de espacio entre el gráfico y la barra de colores..
\end{nota}

El gráfico de dispersión 3D se puede observar en la figura \ref{fig:grafico_3D_scatter}

\begin{figure}[ht]
    \centering
    %% Creator: Matplotlib, PGF backend
%%
%% To include the figure in your LaTeX document, write
%%   \input{<filename>.pgf}
%%
%% Make sure the required packages are loaded in your preamble
%%   \usepackage{pgf}
%%
%% Also ensure that all the required font packages are loaded; for instance,
%% the lmodern package is sometimes necessary when using math font.
%%   \usepackage{lmodern}
%%
%% Figures using additional raster images can only be included by \input if
%% they are in the same directory as the main LaTeX file. For loading figures
%% from other directories you can use the `import` package
%%   \usepackage{import}
%%
%% and then include the figures with
%%   \import{<path to file>}{<filename>.pgf}
%%
%% Matplotlib used the following preamble
%%   \def\mathdefault#1{#1}
%%   \everymath=\expandafter{\the\everymath\displaystyle}
%%   \IfFileExists{scrextend.sty}{
%%     \usepackage[fontsize=10.000000pt]{scrextend}
%%   }{
%%     \renewcommand{\normalsize}{\fontsize{10.000000}{12.000000}\selectfont}
%%     \normalsize
%%   }
%%   \usepackage{fontspec}
%%   \setmainfont{CMU Serif}
%%   \makeatletter\@ifpackageloaded{underscore}{}{\usepackage[strings]{underscore}}\makeatother
%%
\begingroup%
\makeatletter%
\begin{pgfpicture}%
\pgfpathrectangle{\pgfpointorigin}{\pgfqpoint{3.950556in}{3.414292in}}%
\pgfusepath{use as bounding box, clip}%
\begin{pgfscope}%
\pgfsetbuttcap%
\pgfsetmiterjoin%
\definecolor{currentfill}{rgb}{1.000000,1.000000,1.000000}%
\pgfsetfillcolor{currentfill}%
\pgfsetlinewidth{0.000000pt}%
\definecolor{currentstroke}{rgb}{1.000000,1.000000,1.000000}%
\pgfsetstrokecolor{currentstroke}%
\pgfsetdash{}{0pt}%
\pgfpathmoveto{\pgfqpoint{0.000000in}{0.000000in}}%
\pgfpathlineto{\pgfqpoint{3.950556in}{0.000000in}}%
\pgfpathlineto{\pgfqpoint{3.950556in}{3.414292in}}%
\pgfpathlineto{\pgfqpoint{0.000000in}{3.414292in}}%
\pgfpathlineto{\pgfqpoint{0.000000in}{0.000000in}}%
\pgfpathclose%
\pgfusepath{fill}%
\end{pgfscope}%
\begin{pgfscope}%
\pgfsetbuttcap%
\pgfsetmiterjoin%
\definecolor{currentfill}{rgb}{1.000000,1.000000,1.000000}%
\pgfsetfillcolor{currentfill}%
\pgfsetlinewidth{0.000000pt}%
\definecolor{currentstroke}{rgb}{0.000000,0.000000,0.000000}%
\pgfsetstrokecolor{currentstroke}%
\pgfsetstrokeopacity{0.000000}%
\pgfsetdash{}{0pt}%
\pgfpathmoveto{\pgfqpoint{0.100000in}{0.209042in}}%
\pgfpathlineto{\pgfqpoint{3.006250in}{0.209042in}}%
\pgfpathlineto{\pgfqpoint{3.006250in}{3.115292in}}%
\pgfpathlineto{\pgfqpoint{0.100000in}{3.115292in}}%
\pgfpathlineto{\pgfqpoint{0.100000in}{0.209042in}}%
\pgfpathclose%
\pgfusepath{fill}%
\end{pgfscope}%
\begin{pgfscope}%
\pgfsetbuttcap%
\pgfsetmiterjoin%
\definecolor{currentfill}{rgb}{0.950000,0.950000,0.950000}%
\pgfsetfillcolor{currentfill}%
\pgfsetfillopacity{0.500000}%
\pgfsetlinewidth{1.003750pt}%
\definecolor{currentstroke}{rgb}{0.950000,0.950000,0.950000}%
\pgfsetstrokecolor{currentstroke}%
\pgfsetstrokeopacity{0.500000}%
\pgfsetdash{}{0pt}%
\pgfpathmoveto{\pgfqpoint{0.319442in}{0.925630in}}%
\pgfpathlineto{\pgfqpoint{1.279181in}{1.730102in}}%
\pgfpathlineto{\pgfqpoint{1.265839in}{2.890295in}}%
\pgfpathlineto{\pgfqpoint{0.260172in}{2.156405in}}%
\pgfusepath{stroke,fill}%
\end{pgfscope}%
\begin{pgfscope}%
\pgfsetbuttcap%
\pgfsetmiterjoin%
\definecolor{currentfill}{rgb}{0.900000,0.900000,0.900000}%
\pgfsetfillcolor{currentfill}%
\pgfsetfillopacity{0.500000}%
\pgfsetlinewidth{1.003750pt}%
\definecolor{currentstroke}{rgb}{0.900000,0.900000,0.900000}%
\pgfsetstrokecolor{currentstroke}%
\pgfsetstrokeopacity{0.500000}%
\pgfsetdash{}{0pt}%
\pgfpathmoveto{\pgfqpoint{1.279181in}{1.730102in}}%
\pgfpathlineto{\pgfqpoint{2.819219in}{1.282474in}}%
\pgfpathlineto{\pgfqpoint{2.874177in}{2.482628in}}%
\pgfpathlineto{\pgfqpoint{1.265839in}{2.890295in}}%
\pgfusepath{stroke,fill}%
\end{pgfscope}%
\begin{pgfscope}%
\pgfsetbuttcap%
\pgfsetmiterjoin%
\definecolor{currentfill}{rgb}{0.925000,0.925000,0.925000}%
\pgfsetfillcolor{currentfill}%
\pgfsetfillopacity{0.500000}%
\pgfsetlinewidth{1.003750pt}%
\definecolor{currentstroke}{rgb}{0.925000,0.925000,0.925000}%
\pgfsetstrokecolor{currentstroke}%
\pgfsetstrokeopacity{0.500000}%
\pgfsetdash{}{0pt}%
\pgfpathmoveto{\pgfqpoint{0.319442in}{0.925630in}}%
\pgfpathlineto{\pgfqpoint{1.951958in}{0.392451in}}%
\pgfpathlineto{\pgfqpoint{2.819219in}{1.282474in}}%
\pgfpathlineto{\pgfqpoint{1.279181in}{1.730102in}}%
\pgfusepath{stroke,fill}%
\end{pgfscope}%
\begin{pgfscope}%
\pgfsetbuttcap%
\pgfsetroundjoin%
\pgfsetlinewidth{0.501875pt}%
\definecolor{currentstroke}{rgb}{0.501961,0.501961,0.501961}%
\pgfsetstrokecolor{currentstroke}%
\pgfsetdash{{1.850000pt}{0.800000pt}}{0.000000pt}%
\pgfpathmoveto{\pgfqpoint{0.643195in}{0.819892in}}%
\pgfpathlineto{\pgfqpoint{1.585677in}{1.641016in}}%
\pgfpathlineto{\pgfqpoint{1.585387in}{2.809299in}}%
\pgfusepath{stroke}%
\end{pgfscope}%
\begin{pgfscope}%
\pgfsetbuttcap%
\pgfsetroundjoin%
\pgfsetlinewidth{0.501875pt}%
\definecolor{currentstroke}{rgb}{0.501961,0.501961,0.501961}%
\pgfsetstrokecolor{currentstroke}%
\pgfsetdash{{1.850000pt}{0.800000pt}}{0.000000pt}%
\pgfpathmoveto{\pgfqpoint{1.120318in}{0.664064in}}%
\pgfpathlineto{\pgfqpoint{2.036388in}{1.510012in}}%
\pgfpathlineto{\pgfqpoint{2.055780in}{2.690068in}}%
\pgfusepath{stroke}%
\end{pgfscope}%
\begin{pgfscope}%
\pgfsetbuttcap%
\pgfsetroundjoin%
\pgfsetlinewidth{0.501875pt}%
\definecolor{currentstroke}{rgb}{0.501961,0.501961,0.501961}%
\pgfsetstrokecolor{currentstroke}%
\pgfsetdash{{1.850000pt}{0.800000pt}}{0.000000pt}%
\pgfpathmoveto{\pgfqpoint{1.612620in}{0.503278in}}%
\pgfpathlineto{\pgfqpoint{2.500222in}{1.375194in}}%
\pgfpathlineto{\pgfqpoint{2.540475in}{2.567212in}}%
\pgfusepath{stroke}%
\end{pgfscope}%
\begin{pgfscope}%
\pgfsetbuttcap%
\pgfsetroundjoin%
\pgfsetlinewidth{0.501875pt}%
\definecolor{currentstroke}{rgb}{0.501961,0.501961,0.501961}%
\pgfsetstrokecolor{currentstroke}%
\pgfsetdash{{1.850000pt}{0.800000pt}}{0.000000pt}%
\pgfpathmoveto{\pgfqpoint{0.280686in}{2.171375in}}%
\pgfpathlineto{\pgfqpoint{0.338946in}{0.941979in}}%
\pgfpathlineto{\pgfqpoint{1.969659in}{0.410617in}}%
\pgfusepath{stroke}%
\end{pgfscope}%
\begin{pgfscope}%
\pgfsetbuttcap%
\pgfsetroundjoin%
\pgfsetlinewidth{0.501875pt}%
\definecolor{currentstroke}{rgb}{0.501961,0.501961,0.501961}%
\pgfsetstrokecolor{currentstroke}%
\pgfsetdash{{1.850000pt}{0.800000pt}}{0.000000pt}%
\pgfpathmoveto{\pgfqpoint{0.554259in}{2.371016in}}%
\pgfpathlineto{\pgfqpoint{0.599344in}{1.160250in}}%
\pgfpathlineto{\pgfqpoint{2.205681in}{0.652834in}}%
\pgfusepath{stroke}%
\end{pgfscope}%
\begin{pgfscope}%
\pgfsetbuttcap%
\pgfsetroundjoin%
\pgfsetlinewidth{0.501875pt}%
\definecolor{currentstroke}{rgb}{0.501961,0.501961,0.501961}%
\pgfsetstrokecolor{currentstroke}%
\pgfsetdash{{1.850000pt}{0.800000pt}}{0.000000pt}%
\pgfpathmoveto{\pgfqpoint{0.815418in}{2.561598in}}%
\pgfpathlineto{\pgfqpoint{0.848429in}{1.369038in}}%
\pgfpathlineto{\pgfqpoint{2.430920in}{0.883984in}}%
\pgfusepath{stroke}%
\end{pgfscope}%
\begin{pgfscope}%
\pgfsetbuttcap%
\pgfsetroundjoin%
\pgfsetlinewidth{0.501875pt}%
\definecolor{currentstroke}{rgb}{0.501961,0.501961,0.501961}%
\pgfsetstrokecolor{currentstroke}%
\pgfsetdash{{1.850000pt}{0.800000pt}}{0.000000pt}%
\pgfpathmoveto{\pgfqpoint{1.064990in}{2.743724in}}%
\pgfpathlineto{\pgfqpoint{1.086921in}{1.568947in}}%
\pgfpathlineto{\pgfqpoint{2.646098in}{1.104809in}}%
\pgfusepath{stroke}%
\end{pgfscope}%
\begin{pgfscope}%
\pgfsetbuttcap%
\pgfsetroundjoin%
\pgfsetlinewidth{0.501875pt}%
\definecolor{currentstroke}{rgb}{0.501961,0.501961,0.501961}%
\pgfsetstrokecolor{currentstroke}%
\pgfsetdash{{1.850000pt}{0.800000pt}}{0.000000pt}%
\pgfpathmoveto{\pgfqpoint{2.830136in}{1.520883in}}%
\pgfpathlineto{\pgfqpoint{1.276526in}{1.960963in}}%
\pgfpathlineto{\pgfqpoint{0.307684in}{1.169792in}}%
\pgfusepath{stroke}%
\end{pgfscope}%
\begin{pgfscope}%
\pgfsetbuttcap%
\pgfsetroundjoin%
\pgfsetlinewidth{0.501875pt}%
\definecolor{currentstroke}{rgb}{0.501961,0.501961,0.501961}%
\pgfsetstrokecolor{currentstroke}%
\pgfsetdash{{1.850000pt}{0.800000pt}}{0.000000pt}%
\pgfpathmoveto{\pgfqpoint{2.846265in}{1.873094in}}%
\pgfpathlineto{\pgfqpoint{1.272608in}{2.301669in}}%
\pgfpathlineto{\pgfqpoint{0.290299in}{1.530801in}}%
\pgfusepath{stroke}%
\end{pgfscope}%
\begin{pgfscope}%
\pgfsetbuttcap%
\pgfsetroundjoin%
\pgfsetlinewidth{0.501875pt}%
\definecolor{currentstroke}{rgb}{0.501961,0.501961,0.501961}%
\pgfsetstrokecolor{currentstroke}%
\pgfsetdash{{1.850000pt}{0.800000pt}}{0.000000pt}%
\pgfpathmoveto{\pgfqpoint{2.862820in}{2.234606in}}%
\pgfpathlineto{\pgfqpoint{1.268592in}{2.650933in}}%
\pgfpathlineto{\pgfqpoint{0.272437in}{1.901715in}}%
\pgfusepath{stroke}%
\end{pgfscope}%
\begin{pgfscope}%
\pgfsetrectcap%
\pgfsetroundjoin%
\pgfsetlinewidth{0.803000pt}%
\definecolor{currentstroke}{rgb}{0.000000,0.000000,0.000000}%
\pgfsetstrokecolor{currentstroke}%
\pgfsetdash{}{0pt}%
\pgfpathmoveto{\pgfqpoint{0.319442in}{0.925630in}}%
\pgfpathlineto{\pgfqpoint{1.951958in}{0.392451in}}%
\pgfusepath{stroke}%
\end{pgfscope}%
\begin{pgfscope}%
\pgfpathrectangle{\pgfqpoint{0.100000in}{0.209042in}}{\pgfqpoint{2.906250in}{2.906250in}}%
\pgfusepath{clip}%
\pgfsetbuttcap%
\pgfsetroundjoin%
\pgfsetlinewidth{0.501875pt}%
\definecolor{currentstroke}{rgb}{0.501961,0.501961,0.501961}%
\pgfsetstrokecolor{currentstroke}%
\pgfsetstrokeopacity{0.500000}%
\pgfsetdash{{1.850000pt}{0.800000pt}}{0.000000pt}%
\pgfpathmoveto{\pgfqpoint{1.592399in}{1.701441in}}%
\pgfusepath{stroke}%
\end{pgfscope}%
\begin{pgfscope}%
\pgfsetrectcap%
\pgfsetroundjoin%
\pgfsetlinewidth{0.803000pt}%
\definecolor{currentstroke}{rgb}{0.000000,0.000000,0.000000}%
\pgfsetstrokecolor{currentstroke}%
\pgfsetdash{}{0pt}%
\pgfpathmoveto{\pgfqpoint{0.651407in}{0.827047in}}%
\pgfpathlineto{\pgfqpoint{0.626736in}{0.805552in}}%
\pgfusepath{stroke}%
\end{pgfscope}%
\begin{pgfscope}%
\definecolor{textcolor}{rgb}{0.000000,0.000000,0.000000}%
\pgfsetstrokecolor{textcolor}%
\pgfsetfillcolor{textcolor}%
\pgftext[x=0.538809in,y=0.599486in,,top]{\color{textcolor}{\rmfamily\fontsize{10.000000}{12.000000}\selectfont\catcode`\^=\active\def^{\ifmmode\sp\else\^{}\fi}\catcode`\%=\active\def%{\%}$\mathdefault{\ensuremath{-}2}$}}%
\end{pgfscope}%
\begin{pgfscope}%
\pgfpathrectangle{\pgfqpoint{0.100000in}{0.209042in}}{\pgfqpoint{2.906250in}{2.906250in}}%
\pgfusepath{clip}%
\pgfsetbuttcap%
\pgfsetroundjoin%
\pgfsetlinewidth{0.501875pt}%
\definecolor{currentstroke}{rgb}{0.501961,0.501961,0.501961}%
\pgfsetstrokecolor{currentstroke}%
\pgfsetstrokeopacity{0.500000}%
\pgfsetdash{{1.850000pt}{0.800000pt}}{0.000000pt}%
\pgfpathmoveto{\pgfqpoint{1.592399in}{1.701441in}}%
\pgfusepath{stroke}%
\end{pgfscope}%
\begin{pgfscope}%
\pgfsetrectcap%
\pgfsetroundjoin%
\pgfsetlinewidth{0.803000pt}%
\definecolor{currentstroke}{rgb}{0.000000,0.000000,0.000000}%
\pgfsetstrokecolor{currentstroke}%
\pgfsetdash{}{0pt}%
\pgfpathmoveto{\pgfqpoint{1.128310in}{0.671445in}}%
\pgfpathlineto{\pgfqpoint{1.104298in}{0.649271in}}%
\pgfusepath{stroke}%
\end{pgfscope}%
\begin{pgfscope}%
\definecolor{textcolor}{rgb}{0.000000,0.000000,0.000000}%
\pgfsetstrokecolor{textcolor}%
\pgfsetfillcolor{textcolor}%
\pgftext[x=1.016600in,y=0.439435in,,top]{\color{textcolor}{\rmfamily\fontsize{10.000000}{12.000000}\selectfont\catcode`\^=\active\def^{\ifmmode\sp\else\^{}\fi}\catcode`\%=\active\def%{\%}$\mathdefault{0}$}}%
\end{pgfscope}%
\begin{pgfscope}%
\pgfpathrectangle{\pgfqpoint{0.100000in}{0.209042in}}{\pgfqpoint{2.906250in}{2.906250in}}%
\pgfusepath{clip}%
\pgfsetbuttcap%
\pgfsetroundjoin%
\pgfsetlinewidth{0.501875pt}%
\definecolor{currentstroke}{rgb}{0.501961,0.501961,0.501961}%
\pgfsetstrokecolor{currentstroke}%
\pgfsetstrokeopacity{0.500000}%
\pgfsetdash{{1.850000pt}{0.800000pt}}{0.000000pt}%
\pgfpathmoveto{\pgfqpoint{1.592399in}{1.701441in}}%
\pgfusepath{stroke}%
\end{pgfscope}%
\begin{pgfscope}%
\pgfsetrectcap%
\pgfsetroundjoin%
\pgfsetlinewidth{0.803000pt}%
\definecolor{currentstroke}{rgb}{0.000000,0.000000,0.000000}%
\pgfsetstrokecolor{currentstroke}%
\pgfsetdash{}{0pt}%
\pgfpathmoveto{\pgfqpoint{1.620374in}{0.510895in}}%
\pgfpathlineto{\pgfqpoint{1.597078in}{0.488010in}}%
\pgfusepath{stroke}%
\end{pgfscope}%
\begin{pgfscope}%
\definecolor{textcolor}{rgb}{0.000000,0.000000,0.000000}%
\pgfsetstrokecolor{textcolor}%
\pgfsetfillcolor{textcolor}%
\pgftext[x=1.509667in,y=0.274267in,,top]{\color{textcolor}{\rmfamily\fontsize{10.000000}{12.000000}\selectfont\catcode`\^=\active\def^{\ifmmode\sp\else\^{}\fi}\catcode`\%=\active\def%{\%}$\mathdefault{2}$}}%
\end{pgfscope}%
\begin{pgfscope}%
\pgfsetrectcap%
\pgfsetroundjoin%
\pgfsetlinewidth{0.803000pt}%
\definecolor{currentstroke}{rgb}{0.000000,0.000000,0.000000}%
\pgfsetstrokecolor{currentstroke}%
\pgfsetdash{}{0pt}%
\pgfpathmoveto{\pgfqpoint{2.819219in}{1.282474in}}%
\pgfpathlineto{\pgfqpoint{1.951958in}{0.392451in}}%
\pgfusepath{stroke}%
\end{pgfscope}%
\begin{pgfscope}%
\pgfpathrectangle{\pgfqpoint{0.100000in}{0.209042in}}{\pgfqpoint{2.906250in}{2.906250in}}%
\pgfusepath{clip}%
\pgfsetbuttcap%
\pgfsetroundjoin%
\pgfsetlinewidth{0.501875pt}%
\definecolor{currentstroke}{rgb}{0.501961,0.501961,0.501961}%
\pgfsetstrokecolor{currentstroke}%
\pgfsetstrokeopacity{0.500000}%
\pgfsetdash{{1.850000pt}{0.800000pt}}{0.000000pt}%
\pgfpathmoveto{\pgfqpoint{1.592399in}{1.701441in}}%
\pgfusepath{stroke}%
\end{pgfscope}%
\begin{pgfscope}%
\pgfsetrectcap%
\pgfsetroundjoin%
\pgfsetlinewidth{0.803000pt}%
\definecolor{currentstroke}{rgb}{0.000000,0.000000,0.000000}%
\pgfsetstrokecolor{currentstroke}%
\pgfsetdash{}{0pt}%
\pgfpathmoveto{\pgfqpoint{1.955914in}{0.415095in}}%
\pgfpathlineto{\pgfqpoint{1.997185in}{0.401647in}}%
\pgfusepath{stroke}%
\end{pgfscope}%
\begin{pgfscope}%
\definecolor{textcolor}{rgb}{0.000000,0.000000,0.000000}%
\pgfsetstrokecolor{textcolor}%
\pgfsetfillcolor{textcolor}%
\pgftext[x=2.149307in,y=0.223333in,,top]{\color{textcolor}{\rmfamily\fontsize{10.000000}{12.000000}\selectfont\catcode`\^=\active\def^{\ifmmode\sp\else\^{}\fi}\catcode`\%=\active\def%{\%}$\mathdefault{\ensuremath{-}4}$}}%
\end{pgfscope}%
\begin{pgfscope}%
\pgfpathrectangle{\pgfqpoint{0.100000in}{0.209042in}}{\pgfqpoint{2.906250in}{2.906250in}}%
\pgfusepath{clip}%
\pgfsetbuttcap%
\pgfsetroundjoin%
\pgfsetlinewidth{0.501875pt}%
\definecolor{currentstroke}{rgb}{0.501961,0.501961,0.501961}%
\pgfsetstrokecolor{currentstroke}%
\pgfsetstrokeopacity{0.500000}%
\pgfsetdash{{1.850000pt}{0.800000pt}}{0.000000pt}%
\pgfpathmoveto{\pgfqpoint{1.592399in}{1.701441in}}%
\pgfusepath{stroke}%
\end{pgfscope}%
\begin{pgfscope}%
\pgfsetrectcap%
\pgfsetroundjoin%
\pgfsetlinewidth{0.803000pt}%
\definecolor{currentstroke}{rgb}{0.000000,0.000000,0.000000}%
\pgfsetstrokecolor{currentstroke}%
\pgfsetdash{}{0pt}%
\pgfpathmoveto{\pgfqpoint{2.192158in}{0.657105in}}%
\pgfpathlineto{\pgfqpoint{2.232762in}{0.644279in}}%
\pgfusepath{stroke}%
\end{pgfscope}%
\begin{pgfscope}%
\definecolor{textcolor}{rgb}{0.000000,0.000000,0.000000}%
\pgfsetstrokecolor{textcolor}%
\pgfsetfillcolor{textcolor}%
\pgftext[x=2.381302in,y=0.470113in,,top]{\color{textcolor}{\rmfamily\fontsize{10.000000}{12.000000}\selectfont\catcode`\^=\active\def^{\ifmmode\sp\else\^{}\fi}\catcode`\%=\active\def%{\%}$\mathdefault{\ensuremath{-}2}$}}%
\end{pgfscope}%
\begin{pgfscope}%
\pgfpathrectangle{\pgfqpoint{0.100000in}{0.209042in}}{\pgfqpoint{2.906250in}{2.906250in}}%
\pgfusepath{clip}%
\pgfsetbuttcap%
\pgfsetroundjoin%
\pgfsetlinewidth{0.501875pt}%
\definecolor{currentstroke}{rgb}{0.501961,0.501961,0.501961}%
\pgfsetstrokecolor{currentstroke}%
\pgfsetstrokeopacity{0.500000}%
\pgfsetdash{{1.850000pt}{0.800000pt}}{0.000000pt}%
\pgfpathmoveto{\pgfqpoint{1.592399in}{1.701441in}}%
\pgfusepath{stroke}%
\end{pgfscope}%
\begin{pgfscope}%
\pgfsetrectcap%
\pgfsetroundjoin%
\pgfsetlinewidth{0.803000pt}%
\definecolor{currentstroke}{rgb}{0.000000,0.000000,0.000000}%
\pgfsetstrokecolor{currentstroke}%
\pgfsetdash{}{0pt}%
\pgfpathmoveto{\pgfqpoint{2.417612in}{0.888063in}}%
\pgfpathlineto{\pgfqpoint{2.457568in}{0.875816in}}%
\pgfusepath{stroke}%
\end{pgfscope}%
\begin{pgfscope}%
\definecolor{textcolor}{rgb}{0.000000,0.000000,0.000000}%
\pgfsetstrokecolor{textcolor}%
\pgfsetfillcolor{textcolor}%
\pgftext[x=2.602691in,y=0.705609in,,top]{\color{textcolor}{\rmfamily\fontsize{10.000000}{12.000000}\selectfont\catcode`\^=\active\def^{\ifmmode\sp\else\^{}\fi}\catcode`\%=\active\def%{\%}$\mathdefault{0}$}}%
\end{pgfscope}%
\begin{pgfscope}%
\pgfpathrectangle{\pgfqpoint{0.100000in}{0.209042in}}{\pgfqpoint{2.906250in}{2.906250in}}%
\pgfusepath{clip}%
\pgfsetbuttcap%
\pgfsetroundjoin%
\pgfsetlinewidth{0.501875pt}%
\definecolor{currentstroke}{rgb}{0.501961,0.501961,0.501961}%
\pgfsetstrokecolor{currentstroke}%
\pgfsetstrokeopacity{0.500000}%
\pgfsetdash{{1.850000pt}{0.800000pt}}{0.000000pt}%
\pgfpathmoveto{\pgfqpoint{1.592399in}{1.701441in}}%
\pgfusepath{stroke}%
\end{pgfscope}%
\begin{pgfscope}%
\pgfsetrectcap%
\pgfsetroundjoin%
\pgfsetlinewidth{0.803000pt}%
\definecolor{currentstroke}{rgb}{0.000000,0.000000,0.000000}%
\pgfsetstrokecolor{currentstroke}%
\pgfsetdash{}{0pt}%
\pgfpathmoveto{\pgfqpoint{2.633000in}{1.108708in}}%
\pgfpathlineto{\pgfqpoint{2.672324in}{1.097002in}}%
\pgfusepath{stroke}%
\end{pgfscope}%
\begin{pgfscope}%
\definecolor{textcolor}{rgb}{0.000000,0.000000,0.000000}%
\pgfsetstrokecolor{textcolor}%
\pgfsetfillcolor{textcolor}%
\pgftext[x=2.814183in,y=0.930578in,,top]{\color{textcolor}{\rmfamily\fontsize{10.000000}{12.000000}\selectfont\catcode`\^=\active\def^{\ifmmode\sp\else\^{}\fi}\catcode`\%=\active\def%{\%}$\mathdefault{2}$}}%
\end{pgfscope}%
\begin{pgfscope}%
\pgfsetrectcap%
\pgfsetroundjoin%
\pgfsetlinewidth{0.803000pt}%
\definecolor{currentstroke}{rgb}{0.000000,0.000000,0.000000}%
\pgfsetstrokecolor{currentstroke}%
\pgfsetdash{}{0pt}%
\pgfpathmoveto{\pgfqpoint{2.819219in}{1.282474in}}%
\pgfpathlineto{\pgfqpoint{2.874177in}{2.482628in}}%
\pgfusepath{stroke}%
\end{pgfscope}%
\begin{pgfscope}%
\pgfsetbuttcap%
\pgfsetroundjoin%
\pgfsetlinewidth{0.501875pt}%
\definecolor{currentstroke}{rgb}{0.501961,0.501961,0.501961}%
\pgfsetstrokecolor{currentstroke}%
\pgfsetstrokeopacity{0.500000}%
\pgfsetdash{{1.850000pt}{0.800000pt}}{0.000000pt}%
\pgfpathmoveto{\pgfqpoint{1.592399in}{1.701441in}}%
\pgfusepath{stroke}%
\end{pgfscope}%
\begin{pgfscope}%
\pgfsetrectcap%
\pgfsetroundjoin%
\pgfsetlinewidth{0.803000pt}%
\definecolor{currentstroke}{rgb}{0.000000,0.000000,0.000000}%
\pgfsetstrokecolor{currentstroke}%
\pgfsetdash{}{0pt}%
\pgfpathmoveto{\pgfqpoint{2.817092in}{1.524578in}}%
\pgfpathlineto{\pgfqpoint{2.856256in}{1.513484in}}%
\pgfusepath{stroke}%
\end{pgfscope}%
\begin{pgfscope}%
\definecolor{textcolor}{rgb}{0.000000,0.000000,0.000000}%
\pgfsetstrokecolor{textcolor}%
\pgfsetfillcolor{textcolor}%
\pgftext[x=3.085586in,y=1.556438in,,top]{\color{textcolor}{\rmfamily\fontsize{10.000000}{12.000000}\selectfont\catcode`\^=\active\def^{\ifmmode\sp\else\^{}\fi}\catcode`\%=\active\def%{\%}$\mathdefault{\ensuremath{-}2}$}}%
\end{pgfscope}%
\begin{pgfscope}%
\pgfsetbuttcap%
\pgfsetroundjoin%
\pgfsetlinewidth{0.501875pt}%
\definecolor{currentstroke}{rgb}{0.501961,0.501961,0.501961}%
\pgfsetstrokecolor{currentstroke}%
\pgfsetstrokeopacity{0.500000}%
\pgfsetdash{{1.850000pt}{0.800000pt}}{0.000000pt}%
\pgfpathmoveto{\pgfqpoint{1.592399in}{1.701441in}}%
\pgfusepath{stroke}%
\end{pgfscope}%
\begin{pgfscope}%
\pgfsetrectcap%
\pgfsetroundjoin%
\pgfsetlinewidth{0.803000pt}%
\definecolor{currentstroke}{rgb}{0.000000,0.000000,0.000000}%
\pgfsetstrokecolor{currentstroke}%
\pgfsetdash{}{0pt}%
\pgfpathmoveto{\pgfqpoint{2.833044in}{1.876695in}}%
\pgfpathlineto{\pgfqpoint{2.872739in}{1.865884in}}%
\pgfusepath{stroke}%
\end{pgfscope}%
\begin{pgfscope}%
\definecolor{textcolor}{rgb}{0.000000,0.000000,0.000000}%
\pgfsetstrokecolor{textcolor}%
\pgfsetfillcolor{textcolor}%
\pgftext[x=3.104968in,y=1.907740in,,top]{\color{textcolor}{\rmfamily\fontsize{10.000000}{12.000000}\selectfont\catcode`\^=\active\def^{\ifmmode\sp\else\^{}\fi}\catcode`\%=\active\def%{\%}$\mathdefault{0}$}}%
\end{pgfscope}%
\begin{pgfscope}%
\pgfsetbuttcap%
\pgfsetroundjoin%
\pgfsetlinewidth{0.501875pt}%
\definecolor{currentstroke}{rgb}{0.501961,0.501961,0.501961}%
\pgfsetstrokecolor{currentstroke}%
\pgfsetstrokeopacity{0.500000}%
\pgfsetdash{{1.850000pt}{0.800000pt}}{0.000000pt}%
\pgfpathmoveto{\pgfqpoint{1.592399in}{1.701441in}}%
\pgfusepath{stroke}%
\end{pgfscope}%
\begin{pgfscope}%
\pgfsetrectcap%
\pgfsetroundjoin%
\pgfsetlinewidth{0.803000pt}%
\definecolor{currentstroke}{rgb}{0.000000,0.000000,0.000000}%
\pgfsetstrokecolor{currentstroke}%
\pgfsetdash{}{0pt}%
\pgfpathmoveto{\pgfqpoint{2.849417in}{2.238106in}}%
\pgfpathlineto{\pgfqpoint{2.889657in}{2.227598in}}%
\pgfusepath{stroke}%
\end{pgfscope}%
\begin{pgfscope}%
\definecolor{textcolor}{rgb}{0.000000,0.000000,0.000000}%
\pgfsetstrokecolor{textcolor}%
\pgfsetfillcolor{textcolor}%
\pgftext[x=3.124859in,y=2.268281in,,top]{\color{textcolor}{\rmfamily\fontsize{10.000000}{12.000000}\selectfont\catcode`\^=\active\def^{\ifmmode\sp\else\^{}\fi}\catcode`\%=\active\def%{\%}$\mathdefault{2}$}}%
\end{pgfscope}%
\begin{pgfscope}%
\pgfpathrectangle{\pgfqpoint{0.100000in}{0.209042in}}{\pgfqpoint{2.906250in}{2.906250in}}%
\pgfusepath{clip}%
\pgfsetbuttcap%
\pgfsetroundjoin%
\definecolor{currentfill}{rgb}{0.973381,0.461520,0.361965}%
\pgfsetfillcolor{currentfill}%
\pgfsetfillopacity{0.800000}%
\pgfsetlinewidth{1.003750pt}%
\definecolor{currentstroke}{rgb}{0.973381,0.461520,0.361965}%
\pgfsetstrokecolor{currentstroke}%
\pgfsetstrokeopacity{0.800000}%
\pgfsetdash{}{0pt}%
\pgfpathmoveto{\pgfqpoint{1.551565in}{1.661211in}}%
\pgfpathcurveto{\pgfqpoint{1.557388in}{1.661211in}}{\pgfqpoint{1.562975in}{1.663525in}}{\pgfqpoint{1.567093in}{1.667643in}}%
\pgfpathcurveto{\pgfqpoint{1.571211in}{1.671761in}}{\pgfqpoint{1.573525in}{1.677347in}}{\pgfqpoint{1.573525in}{1.683171in}}%
\pgfpathcurveto{\pgfqpoint{1.573525in}{1.688995in}}{\pgfqpoint{1.571211in}{1.694582in}}{\pgfqpoint{1.567093in}{1.698700in}}%
\pgfpathcurveto{\pgfqpoint{1.562975in}{1.702818in}}{\pgfqpoint{1.557388in}{1.705132in}}{\pgfqpoint{1.551565in}{1.705132in}}%
\pgfpathcurveto{\pgfqpoint{1.545741in}{1.705132in}}{\pgfqpoint{1.540154in}{1.702818in}}{\pgfqpoint{1.536036in}{1.698700in}}%
\pgfpathcurveto{\pgfqpoint{1.531918in}{1.694582in}}{\pgfqpoint{1.529604in}{1.688995in}}{\pgfqpoint{1.529604in}{1.683171in}}%
\pgfpathcurveto{\pgfqpoint{1.529604in}{1.677347in}}{\pgfqpoint{1.531918in}{1.671761in}}{\pgfqpoint{1.536036in}{1.667643in}}%
\pgfpathcurveto{\pgfqpoint{1.540154in}{1.663525in}}{\pgfqpoint{1.545741in}{1.661211in}}{\pgfqpoint{1.551565in}{1.661211in}}%
\pgfpathlineto{\pgfqpoint{1.551565in}{1.661211in}}%
\pgfpathclose%
\pgfusepath{stroke,fill}%
\end{pgfscope}%
\begin{pgfscope}%
\pgfpathrectangle{\pgfqpoint{0.100000in}{0.209042in}}{\pgfqpoint{2.906250in}{2.906250in}}%
\pgfusepath{clip}%
\pgfsetbuttcap%
\pgfsetroundjoin%
\definecolor{currentfill}{rgb}{0.955849,0.404400,0.360619}%
\pgfsetfillcolor{currentfill}%
\pgfsetfillopacity{0.800000}%
\pgfsetlinewidth{1.003750pt}%
\definecolor{currentstroke}{rgb}{0.955849,0.404400,0.360619}%
\pgfsetstrokecolor{currentstroke}%
\pgfsetstrokeopacity{0.800000}%
\pgfsetdash{}{0pt}%
\pgfpathmoveto{\pgfqpoint{1.421233in}{1.754757in}}%
\pgfpathcurveto{\pgfqpoint{1.427057in}{1.754757in}}{\pgfqpoint{1.432643in}{1.757071in}}{\pgfqpoint{1.436761in}{1.761189in}}%
\pgfpathcurveto{\pgfqpoint{1.440879in}{1.765307in}}{\pgfqpoint{1.443193in}{1.770893in}}{\pgfqpoint{1.443193in}{1.776717in}}%
\pgfpathcurveto{\pgfqpoint{1.443193in}{1.782541in}}{\pgfqpoint{1.440879in}{1.788127in}}{\pgfqpoint{1.436761in}{1.792245in}}%
\pgfpathcurveto{\pgfqpoint{1.432643in}{1.796364in}}{\pgfqpoint{1.427057in}{1.798677in}}{\pgfqpoint{1.421233in}{1.798677in}}%
\pgfpathcurveto{\pgfqpoint{1.415409in}{1.798677in}}{\pgfqpoint{1.409823in}{1.796364in}}{\pgfqpoint{1.405704in}{1.792245in}}%
\pgfpathcurveto{\pgfqpoint{1.401586in}{1.788127in}}{\pgfqpoint{1.399272in}{1.782541in}}{\pgfqpoint{1.399272in}{1.776717in}}%
\pgfpathcurveto{\pgfqpoint{1.399272in}{1.770893in}}{\pgfqpoint{1.401586in}{1.765307in}}{\pgfqpoint{1.405704in}{1.761189in}}%
\pgfpathcurveto{\pgfqpoint{1.409823in}{1.757071in}}{\pgfqpoint{1.415409in}{1.754757in}}{\pgfqpoint{1.421233in}{1.754757in}}%
\pgfpathlineto{\pgfqpoint{1.421233in}{1.754757in}}%
\pgfpathclose%
\pgfusepath{stroke,fill}%
\end{pgfscope}%
\begin{pgfscope}%
\pgfpathrectangle{\pgfqpoint{0.100000in}{0.209042in}}{\pgfqpoint{2.906250in}{2.906250in}}%
\pgfusepath{clip}%
\pgfsetbuttcap%
\pgfsetroundjoin%
\definecolor{currentfill}{rgb}{0.950210,0.390820,0.362468}%
\pgfsetfillcolor{currentfill}%
\pgfsetfillopacity{0.800000}%
\pgfsetlinewidth{1.003750pt}%
\definecolor{currentstroke}{rgb}{0.950210,0.390820,0.362468}%
\pgfsetstrokecolor{currentstroke}%
\pgfsetstrokeopacity{0.800000}%
\pgfsetdash{}{0pt}%
\pgfpathmoveto{\pgfqpoint{1.823833in}{1.790437in}}%
\pgfpathcurveto{\pgfqpoint{1.829657in}{1.790437in}}{\pgfqpoint{1.835243in}{1.792751in}}{\pgfqpoint{1.839362in}{1.796869in}}%
\pgfpathcurveto{\pgfqpoint{1.843480in}{1.800987in}}{\pgfqpoint{1.845794in}{1.806573in}}{\pgfqpoint{1.845794in}{1.812397in}}%
\pgfpathcurveto{\pgfqpoint{1.845794in}{1.818221in}}{\pgfqpoint{1.843480in}{1.823807in}}{\pgfqpoint{1.839362in}{1.827925in}}%
\pgfpathcurveto{\pgfqpoint{1.835243in}{1.832044in}}{\pgfqpoint{1.829657in}{1.834357in}}{\pgfqpoint{1.823833in}{1.834357in}}%
\pgfpathcurveto{\pgfqpoint{1.818009in}{1.834357in}}{\pgfqpoint{1.812423in}{1.832044in}}{\pgfqpoint{1.808305in}{1.827925in}}%
\pgfpathcurveto{\pgfqpoint{1.804187in}{1.823807in}}{\pgfqpoint{1.801873in}{1.818221in}}{\pgfqpoint{1.801873in}{1.812397in}}%
\pgfpathcurveto{\pgfqpoint{1.801873in}{1.806573in}}{\pgfqpoint{1.804187in}{1.800987in}}{\pgfqpoint{1.808305in}{1.796869in}}%
\pgfpathcurveto{\pgfqpoint{1.812423in}{1.792751in}}{\pgfqpoint{1.818009in}{1.790437in}}{\pgfqpoint{1.823833in}{1.790437in}}%
\pgfpathlineto{\pgfqpoint{1.823833in}{1.790437in}}%
\pgfpathclose%
\pgfusepath{stroke,fill}%
\end{pgfscope}%
\begin{pgfscope}%
\pgfpathrectangle{\pgfqpoint{0.100000in}{0.209042in}}{\pgfqpoint{2.906250in}{2.906250in}}%
\pgfusepath{clip}%
\pgfsetbuttcap%
\pgfsetroundjoin%
\definecolor{currentfill}{rgb}{0.981000,0.498428,0.369734}%
\pgfsetfillcolor{currentfill}%
\pgfsetfillopacity{0.800000}%
\pgfsetlinewidth{1.003750pt}%
\definecolor{currentstroke}{rgb}{0.981000,0.498428,0.369734}%
\pgfsetstrokecolor{currentstroke}%
\pgfsetstrokeopacity{0.800000}%
\pgfsetdash{}{0pt}%
\pgfpathmoveto{\pgfqpoint{1.684028in}{2.123576in}}%
\pgfpathcurveto{\pgfqpoint{1.689852in}{2.123576in}}{\pgfqpoint{1.695438in}{2.125890in}}{\pgfqpoint{1.699556in}{2.130008in}}%
\pgfpathcurveto{\pgfqpoint{1.703674in}{2.134126in}}{\pgfqpoint{1.705988in}{2.139713in}}{\pgfqpoint{1.705988in}{2.145537in}}%
\pgfpathcurveto{\pgfqpoint{1.705988in}{2.151360in}}{\pgfqpoint{1.703674in}{2.156947in}}{\pgfqpoint{1.699556in}{2.161065in}}%
\pgfpathcurveto{\pgfqpoint{1.695438in}{2.165183in}}{\pgfqpoint{1.689852in}{2.167497in}}{\pgfqpoint{1.684028in}{2.167497in}}%
\pgfpathcurveto{\pgfqpoint{1.678204in}{2.167497in}}{\pgfqpoint{1.672618in}{2.165183in}}{\pgfqpoint{1.668500in}{2.161065in}}%
\pgfpathcurveto{\pgfqpoint{1.664381in}{2.156947in}}{\pgfqpoint{1.662068in}{2.151360in}}{\pgfqpoint{1.662068in}{2.145537in}}%
\pgfpathcurveto{\pgfqpoint{1.662068in}{2.139713in}}{\pgfqpoint{1.664381in}{2.134126in}}{\pgfqpoint{1.668500in}{2.130008in}}%
\pgfpathcurveto{\pgfqpoint{1.672618in}{2.125890in}}{\pgfqpoint{1.678204in}{2.123576in}}{\pgfqpoint{1.684028in}{2.123576in}}%
\pgfpathlineto{\pgfqpoint{1.684028in}{2.123576in}}%
\pgfpathclose%
\pgfusepath{stroke,fill}%
\end{pgfscope}%
\begin{pgfscope}%
\pgfpathrectangle{\pgfqpoint{0.100000in}{0.209042in}}{\pgfqpoint{2.906250in}{2.906250in}}%
\pgfusepath{clip}%
\pgfsetbuttcap%
\pgfsetroundjoin%
\definecolor{currentfill}{rgb}{0.950210,0.390820,0.362468}%
\pgfsetfillcolor{currentfill}%
\pgfsetfillopacity{0.800000}%
\pgfsetlinewidth{1.003750pt}%
\definecolor{currentstroke}{rgb}{0.950210,0.390820,0.362468}%
\pgfsetstrokecolor{currentstroke}%
\pgfsetstrokeopacity{0.800000}%
\pgfsetdash{}{0pt}%
\pgfpathmoveto{\pgfqpoint{1.276930in}{1.632785in}}%
\pgfpathcurveto{\pgfqpoint{1.282754in}{1.632785in}}{\pgfqpoint{1.288340in}{1.635098in}}{\pgfqpoint{1.292458in}{1.639217in}}%
\pgfpathcurveto{\pgfqpoint{1.296576in}{1.643335in}}{\pgfqpoint{1.298890in}{1.648921in}}{\pgfqpoint{1.298890in}{1.654745in}}%
\pgfpathcurveto{\pgfqpoint{1.298890in}{1.660569in}}{\pgfqpoint{1.296576in}{1.666155in}}{\pgfqpoint{1.292458in}{1.670273in}}%
\pgfpathcurveto{\pgfqpoint{1.288340in}{1.674391in}}{\pgfqpoint{1.282754in}{1.676705in}}{\pgfqpoint{1.276930in}{1.676705in}}%
\pgfpathcurveto{\pgfqpoint{1.271106in}{1.676705in}}{\pgfqpoint{1.265520in}{1.674391in}}{\pgfqpoint{1.261402in}{1.670273in}}%
\pgfpathcurveto{\pgfqpoint{1.257283in}{1.666155in}}{\pgfqpoint{1.254970in}{1.660569in}}{\pgfqpoint{1.254970in}{1.654745in}}%
\pgfpathcurveto{\pgfqpoint{1.254970in}{1.648921in}}{\pgfqpoint{1.257283in}{1.643335in}}{\pgfqpoint{1.261402in}{1.639217in}}%
\pgfpathcurveto{\pgfqpoint{1.265520in}{1.635098in}}{\pgfqpoint{1.271106in}{1.632785in}}{\pgfqpoint{1.276930in}{1.632785in}}%
\pgfpathlineto{\pgfqpoint{1.276930in}{1.632785in}}%
\pgfpathclose%
\pgfusepath{stroke,fill}%
\end{pgfscope}%
\begin{pgfscope}%
\pgfpathrectangle{\pgfqpoint{0.100000in}{0.209042in}}{\pgfqpoint{2.906250in}{2.906250in}}%
\pgfusepath{clip}%
\pgfsetbuttcap%
\pgfsetroundjoin%
\definecolor{currentfill}{rgb}{0.879464,0.296125,0.403118}%
\pgfsetfillcolor{currentfill}%
\pgfsetfillopacity{0.800000}%
\pgfsetlinewidth{1.003750pt}%
\definecolor{currentstroke}{rgb}{0.879464,0.296125,0.403118}%
\pgfsetstrokecolor{currentstroke}%
\pgfsetstrokeopacity{0.800000}%
\pgfsetdash{}{0pt}%
\pgfpathmoveto{\pgfqpoint{1.518688in}{1.650829in}}%
\pgfpathcurveto{\pgfqpoint{1.524512in}{1.650829in}}{\pgfqpoint{1.530098in}{1.653143in}}{\pgfqpoint{1.534216in}{1.657261in}}%
\pgfpathcurveto{\pgfqpoint{1.538334in}{1.661379in}}{\pgfqpoint{1.540648in}{1.666965in}}{\pgfqpoint{1.540648in}{1.672789in}}%
\pgfpathcurveto{\pgfqpoint{1.540648in}{1.678613in}}{\pgfqpoint{1.538334in}{1.684199in}}{\pgfqpoint{1.534216in}{1.688317in}}%
\pgfpathcurveto{\pgfqpoint{1.530098in}{1.692435in}}{\pgfqpoint{1.524512in}{1.694749in}}{\pgfqpoint{1.518688in}{1.694749in}}%
\pgfpathcurveto{\pgfqpoint{1.512864in}{1.694749in}}{\pgfqpoint{1.507277in}{1.692435in}}{\pgfqpoint{1.503159in}{1.688317in}}%
\pgfpathcurveto{\pgfqpoint{1.499041in}{1.684199in}}{\pgfqpoint{1.496727in}{1.678613in}}{\pgfqpoint{1.496727in}{1.672789in}}%
\pgfpathcurveto{\pgfqpoint{1.496727in}{1.666965in}}{\pgfqpoint{1.499041in}{1.661379in}}{\pgfqpoint{1.503159in}{1.657261in}}%
\pgfpathcurveto{\pgfqpoint{1.507277in}{1.653143in}}{\pgfqpoint{1.512864in}{1.650829in}}{\pgfqpoint{1.518688in}{1.650829in}}%
\pgfpathlineto{\pgfqpoint{1.518688in}{1.650829in}}%
\pgfpathclose%
\pgfusepath{stroke,fill}%
\end{pgfscope}%
\begin{pgfscope}%
\pgfpathrectangle{\pgfqpoint{0.100000in}{0.209042in}}{\pgfqpoint{2.906250in}{2.906250in}}%
\pgfusepath{clip}%
\pgfsetbuttcap%
\pgfsetroundjoin%
\definecolor{currentfill}{rgb}{0.840636,0.268953,0.424666}%
\pgfsetfillcolor{currentfill}%
\pgfsetfillopacity{0.800000}%
\pgfsetlinewidth{1.003750pt}%
\definecolor{currentstroke}{rgb}{0.840636,0.268953,0.424666}%
\pgfsetstrokecolor{currentstroke}%
\pgfsetstrokeopacity{0.800000}%
\pgfsetdash{}{0pt}%
\pgfpathmoveto{\pgfqpoint{1.590782in}{1.741958in}}%
\pgfpathcurveto{\pgfqpoint{1.596606in}{1.741958in}}{\pgfqpoint{1.602192in}{1.744272in}}{\pgfqpoint{1.606310in}{1.748390in}}%
\pgfpathcurveto{\pgfqpoint{1.610428in}{1.752509in}}{\pgfqpoint{1.612742in}{1.758095in}}{\pgfqpoint{1.612742in}{1.763919in}}%
\pgfpathcurveto{\pgfqpoint{1.612742in}{1.769743in}}{\pgfqpoint{1.610428in}{1.775329in}}{\pgfqpoint{1.606310in}{1.779447in}}%
\pgfpathcurveto{\pgfqpoint{1.602192in}{1.783565in}}{\pgfqpoint{1.596606in}{1.785879in}}{\pgfqpoint{1.590782in}{1.785879in}}%
\pgfpathcurveto{\pgfqpoint{1.584958in}{1.785879in}}{\pgfqpoint{1.579372in}{1.783565in}}{\pgfqpoint{1.575254in}{1.779447in}}%
\pgfpathcurveto{\pgfqpoint{1.571136in}{1.775329in}}{\pgfqpoint{1.568822in}{1.769743in}}{\pgfqpoint{1.568822in}{1.763919in}}%
\pgfpathcurveto{\pgfqpoint{1.568822in}{1.758095in}}{\pgfqpoint{1.571136in}{1.752509in}}{\pgfqpoint{1.575254in}{1.748390in}}%
\pgfpathcurveto{\pgfqpoint{1.579372in}{1.744272in}}{\pgfqpoint{1.584958in}{1.741958in}}{\pgfqpoint{1.590782in}{1.741958in}}%
\pgfpathlineto{\pgfqpoint{1.590782in}{1.741958in}}%
\pgfpathclose%
\pgfusepath{stroke,fill}%
\end{pgfscope}%
\begin{pgfscope}%
\pgfpathrectangle{\pgfqpoint{0.100000in}{0.209042in}}{\pgfqpoint{2.906250in}{2.906250in}}%
\pgfusepath{clip}%
\pgfsetbuttcap%
\pgfsetroundjoin%
\definecolor{currentfill}{rgb}{0.933606,0.358764,0.370541}%
\pgfsetfillcolor{currentfill}%
\pgfsetfillopacity{0.800000}%
\pgfsetlinewidth{1.003750pt}%
\definecolor{currentstroke}{rgb}{0.933606,0.358764,0.370541}%
\pgfsetstrokecolor{currentstroke}%
\pgfsetstrokeopacity{0.800000}%
\pgfsetdash{}{0pt}%
\pgfpathmoveto{\pgfqpoint{1.870853in}{1.748739in}}%
\pgfpathcurveto{\pgfqpoint{1.876677in}{1.748739in}}{\pgfqpoint{1.882263in}{1.751053in}}{\pgfqpoint{1.886381in}{1.755171in}}%
\pgfpathcurveto{\pgfqpoint{1.890500in}{1.759289in}}{\pgfqpoint{1.892813in}{1.764875in}}{\pgfqpoint{1.892813in}{1.770699in}}%
\pgfpathcurveto{\pgfqpoint{1.892813in}{1.776523in}}{\pgfqpoint{1.890500in}{1.782109in}}{\pgfqpoint{1.886381in}{1.786227in}}%
\pgfpathcurveto{\pgfqpoint{1.882263in}{1.790345in}}{\pgfqpoint{1.876677in}{1.792659in}}{\pgfqpoint{1.870853in}{1.792659in}}%
\pgfpathcurveto{\pgfqpoint{1.865029in}{1.792659in}}{\pgfqpoint{1.859443in}{1.790345in}}{\pgfqpoint{1.855325in}{1.786227in}}%
\pgfpathcurveto{\pgfqpoint{1.851207in}{1.782109in}}{\pgfqpoint{1.848893in}{1.776523in}}{\pgfqpoint{1.848893in}{1.770699in}}%
\pgfpathcurveto{\pgfqpoint{1.848893in}{1.764875in}}{\pgfqpoint{1.851207in}{1.759289in}}{\pgfqpoint{1.855325in}{1.755171in}}%
\pgfpathcurveto{\pgfqpoint{1.859443in}{1.751053in}}{\pgfqpoint{1.865029in}{1.748739in}}{\pgfqpoint{1.870853in}{1.748739in}}%
\pgfpathlineto{\pgfqpoint{1.870853in}{1.748739in}}%
\pgfpathclose%
\pgfusepath{stroke,fill}%
\end{pgfscope}%
\begin{pgfscope}%
\pgfpathrectangle{\pgfqpoint{0.100000in}{0.209042in}}{\pgfqpoint{2.906250in}{2.906250in}}%
\pgfusepath{clip}%
\pgfsetbuttcap%
\pgfsetroundjoin%
\definecolor{currentfill}{rgb}{0.994309,0.616999,0.422950}%
\pgfsetfillcolor{currentfill}%
\pgfsetfillopacity{0.800000}%
\pgfsetlinewidth{1.003750pt}%
\definecolor{currentstroke}{rgb}{0.994309,0.616999,0.422950}%
\pgfsetstrokecolor{currentstroke}%
\pgfsetstrokeopacity{0.800000}%
\pgfsetdash{}{0pt}%
\pgfpathmoveto{\pgfqpoint{1.751774in}{1.344272in}}%
\pgfpathcurveto{\pgfqpoint{1.757598in}{1.344272in}}{\pgfqpoint{1.763184in}{1.346586in}}{\pgfqpoint{1.767302in}{1.350704in}}%
\pgfpathcurveto{\pgfqpoint{1.771420in}{1.354822in}}{\pgfqpoint{1.773734in}{1.360408in}}{\pgfqpoint{1.773734in}{1.366232in}}%
\pgfpathcurveto{\pgfqpoint{1.773734in}{1.372056in}}{\pgfqpoint{1.771420in}{1.377642in}}{\pgfqpoint{1.767302in}{1.381761in}}%
\pgfpathcurveto{\pgfqpoint{1.763184in}{1.385879in}}{\pgfqpoint{1.757598in}{1.388193in}}{\pgfqpoint{1.751774in}{1.388193in}}%
\pgfpathcurveto{\pgfqpoint{1.745950in}{1.388193in}}{\pgfqpoint{1.740364in}{1.385879in}}{\pgfqpoint{1.736246in}{1.381761in}}%
\pgfpathcurveto{\pgfqpoint{1.732128in}{1.377642in}}{\pgfqpoint{1.729814in}{1.372056in}}{\pgfqpoint{1.729814in}{1.366232in}}%
\pgfpathcurveto{\pgfqpoint{1.729814in}{1.360408in}}{\pgfqpoint{1.732128in}{1.354822in}}{\pgfqpoint{1.736246in}{1.350704in}}%
\pgfpathcurveto{\pgfqpoint{1.740364in}{1.346586in}}{\pgfqpoint{1.745950in}{1.344272in}}{\pgfqpoint{1.751774in}{1.344272in}}%
\pgfpathlineto{\pgfqpoint{1.751774in}{1.344272in}}%
\pgfpathclose%
\pgfusepath{stroke,fill}%
\end{pgfscope}%
\begin{pgfscope}%
\pgfpathrectangle{\pgfqpoint{0.100000in}{0.209042in}}{\pgfqpoint{2.906250in}{2.906250in}}%
\pgfusepath{clip}%
\pgfsetbuttcap%
\pgfsetroundjoin%
\definecolor{currentfill}{rgb}{0.953099,0.397563,0.361438}%
\pgfsetfillcolor{currentfill}%
\pgfsetfillopacity{0.800000}%
\pgfsetlinewidth{1.003750pt}%
\definecolor{currentstroke}{rgb}{0.953099,0.397563,0.361438}%
\pgfsetstrokecolor{currentstroke}%
\pgfsetstrokeopacity{0.800000}%
\pgfsetdash{}{0pt}%
\pgfpathmoveto{\pgfqpoint{1.190865in}{1.704843in}}%
\pgfpathcurveto{\pgfqpoint{1.196689in}{1.704843in}}{\pgfqpoint{1.202275in}{1.707157in}}{\pgfqpoint{1.206393in}{1.711275in}}%
\pgfpathcurveto{\pgfqpoint{1.210512in}{1.715393in}}{\pgfqpoint{1.212825in}{1.720979in}}{\pgfqpoint{1.212825in}{1.726803in}}%
\pgfpathcurveto{\pgfqpoint{1.212825in}{1.732627in}}{\pgfqpoint{1.210512in}{1.738213in}}{\pgfqpoint{1.206393in}{1.742331in}}%
\pgfpathcurveto{\pgfqpoint{1.202275in}{1.746450in}}{\pgfqpoint{1.196689in}{1.748763in}}{\pgfqpoint{1.190865in}{1.748763in}}%
\pgfpathcurveto{\pgfqpoint{1.185041in}{1.748763in}}{\pgfqpoint{1.179455in}{1.746450in}}{\pgfqpoint{1.175337in}{1.742331in}}%
\pgfpathcurveto{\pgfqpoint{1.171219in}{1.738213in}}{\pgfqpoint{1.168905in}{1.732627in}}{\pgfqpoint{1.168905in}{1.726803in}}%
\pgfpathcurveto{\pgfqpoint{1.168905in}{1.720979in}}{\pgfqpoint{1.171219in}{1.715393in}}{\pgfqpoint{1.175337in}{1.711275in}}%
\pgfpathcurveto{\pgfqpoint{1.179455in}{1.707157in}}{\pgfqpoint{1.185041in}{1.704843in}}{\pgfqpoint{1.190865in}{1.704843in}}%
\pgfpathlineto{\pgfqpoint{1.190865in}{1.704843in}}%
\pgfpathclose%
\pgfusepath{stroke,fill}%
\end{pgfscope}%
\begin{pgfscope}%
\pgfpathrectangle{\pgfqpoint{0.100000in}{0.209042in}}{\pgfqpoint{2.906250in}{2.906250in}}%
\pgfusepath{clip}%
\pgfsetbuttcap%
\pgfsetroundjoin%
\definecolor{currentfill}{rgb}{0.937221,0.364929,0.368567}%
\pgfsetfillcolor{currentfill}%
\pgfsetfillopacity{0.800000}%
\pgfsetlinewidth{1.003750pt}%
\definecolor{currentstroke}{rgb}{0.937221,0.364929,0.368567}%
\pgfsetstrokecolor{currentstroke}%
\pgfsetstrokeopacity{0.800000}%
\pgfsetdash{}{0pt}%
\pgfpathmoveto{\pgfqpoint{1.250921in}{1.960688in}}%
\pgfpathcurveto{\pgfqpoint{1.256745in}{1.960688in}}{\pgfqpoint{1.262331in}{1.963001in}}{\pgfqpoint{1.266449in}{1.967120in}}%
\pgfpathcurveto{\pgfqpoint{1.270567in}{1.971238in}}{\pgfqpoint{1.272881in}{1.976824in}}{\pgfqpoint{1.272881in}{1.982648in}}%
\pgfpathcurveto{\pgfqpoint{1.272881in}{1.988472in}}{\pgfqpoint{1.270567in}{1.994058in}}{\pgfqpoint{1.266449in}{1.998176in}}%
\pgfpathcurveto{\pgfqpoint{1.262331in}{2.002294in}}{\pgfqpoint{1.256745in}{2.004608in}}{\pgfqpoint{1.250921in}{2.004608in}}%
\pgfpathcurveto{\pgfqpoint{1.245097in}{2.004608in}}{\pgfqpoint{1.239511in}{2.002294in}}{\pgfqpoint{1.235393in}{1.998176in}}%
\pgfpathcurveto{\pgfqpoint{1.231274in}{1.994058in}}{\pgfqpoint{1.228961in}{1.988472in}}{\pgfqpoint{1.228961in}{1.982648in}}%
\pgfpathcurveto{\pgfqpoint{1.228961in}{1.976824in}}{\pgfqpoint{1.231274in}{1.971238in}}{\pgfqpoint{1.235393in}{1.967120in}}%
\pgfpathcurveto{\pgfqpoint{1.239511in}{1.963001in}}{\pgfqpoint{1.245097in}{1.960688in}}{\pgfqpoint{1.250921in}{1.960688in}}%
\pgfpathlineto{\pgfqpoint{1.250921in}{1.960688in}}%
\pgfpathclose%
\pgfusepath{stroke,fill}%
\end{pgfscope}%
\begin{pgfscope}%
\pgfpathrectangle{\pgfqpoint{0.100000in}{0.209042in}}{\pgfqpoint{2.906250in}{2.906250in}}%
\pgfusepath{clip}%
\pgfsetbuttcap%
\pgfsetroundjoin%
\definecolor{currentfill}{rgb}{0.798608,0.247040,0.444848}%
\pgfsetfillcolor{currentfill}%
\pgfsetfillopacity{0.800000}%
\pgfsetlinewidth{1.003750pt}%
\definecolor{currentstroke}{rgb}{0.798608,0.247040,0.444848}%
\pgfsetstrokecolor{currentstroke}%
\pgfsetstrokeopacity{0.800000}%
\pgfsetdash{}{0pt}%
\pgfpathmoveto{\pgfqpoint{1.479799in}{1.829948in}}%
\pgfpathcurveto{\pgfqpoint{1.485623in}{1.829948in}}{\pgfqpoint{1.491209in}{1.832262in}}{\pgfqpoint{1.495327in}{1.836380in}}%
\pgfpathcurveto{\pgfqpoint{1.499446in}{1.840498in}}{\pgfqpoint{1.501759in}{1.846084in}}{\pgfqpoint{1.501759in}{1.851908in}}%
\pgfpathcurveto{\pgfqpoint{1.501759in}{1.857732in}}{\pgfqpoint{1.499446in}{1.863318in}}{\pgfqpoint{1.495327in}{1.867436in}}%
\pgfpathcurveto{\pgfqpoint{1.491209in}{1.871554in}}{\pgfqpoint{1.485623in}{1.873868in}}{\pgfqpoint{1.479799in}{1.873868in}}%
\pgfpathcurveto{\pgfqpoint{1.473975in}{1.873868in}}{\pgfqpoint{1.468389in}{1.871554in}}{\pgfqpoint{1.464271in}{1.867436in}}%
\pgfpathcurveto{\pgfqpoint{1.460153in}{1.863318in}}{\pgfqpoint{1.457839in}{1.857732in}}{\pgfqpoint{1.457839in}{1.851908in}}%
\pgfpathcurveto{\pgfqpoint{1.457839in}{1.846084in}}{\pgfqpoint{1.460153in}{1.840498in}}{\pgfqpoint{1.464271in}{1.836380in}}%
\pgfpathcurveto{\pgfqpoint{1.468389in}{1.832262in}}{\pgfqpoint{1.473975in}{1.829948in}}{\pgfqpoint{1.479799in}{1.829948in}}%
\pgfpathlineto{\pgfqpoint{1.479799in}{1.829948in}}%
\pgfpathclose%
\pgfusepath{stroke,fill}%
\end{pgfscope}%
\begin{pgfscope}%
\pgfpathrectangle{\pgfqpoint{0.100000in}{0.209042in}}{\pgfqpoint{2.906250in}{2.906250in}}%
\pgfusepath{clip}%
\pgfsetbuttcap%
\pgfsetroundjoin%
\definecolor{currentfill}{rgb}{0.828886,0.262229,0.430644}%
\pgfsetfillcolor{currentfill}%
\pgfsetfillopacity{0.800000}%
\pgfsetlinewidth{1.003750pt}%
\definecolor{currentstroke}{rgb}{0.828886,0.262229,0.430644}%
\pgfsetstrokecolor{currentstroke}%
\pgfsetstrokeopacity{0.800000}%
\pgfsetdash{}{0pt}%
\pgfpathmoveto{\pgfqpoint{1.496904in}{1.642144in}}%
\pgfpathcurveto{\pgfqpoint{1.502728in}{1.642144in}}{\pgfqpoint{1.508314in}{1.644457in}}{\pgfqpoint{1.512432in}{1.648576in}}%
\pgfpathcurveto{\pgfqpoint{1.516550in}{1.652694in}}{\pgfqpoint{1.518864in}{1.658280in}}{\pgfqpoint{1.518864in}{1.664104in}}%
\pgfpathcurveto{\pgfqpoint{1.518864in}{1.669928in}}{\pgfqpoint{1.516550in}{1.675514in}}{\pgfqpoint{1.512432in}{1.679632in}}%
\pgfpathcurveto{\pgfqpoint{1.508314in}{1.683750in}}{\pgfqpoint{1.502728in}{1.686064in}}{\pgfqpoint{1.496904in}{1.686064in}}%
\pgfpathcurveto{\pgfqpoint{1.491080in}{1.686064in}}{\pgfqpoint{1.485494in}{1.683750in}}{\pgfqpoint{1.481376in}{1.679632in}}%
\pgfpathcurveto{\pgfqpoint{1.477257in}{1.675514in}}{\pgfqpoint{1.474944in}{1.669928in}}{\pgfqpoint{1.474944in}{1.664104in}}%
\pgfpathcurveto{\pgfqpoint{1.474944in}{1.658280in}}{\pgfqpoint{1.477257in}{1.652694in}}{\pgfqpoint{1.481376in}{1.648576in}}%
\pgfpathcurveto{\pgfqpoint{1.485494in}{1.644457in}}{\pgfqpoint{1.491080in}{1.642144in}}{\pgfqpoint{1.496904in}{1.642144in}}%
\pgfpathlineto{\pgfqpoint{1.496904in}{1.642144in}}%
\pgfpathclose%
\pgfusepath{stroke,fill}%
\end{pgfscope}%
\begin{pgfscope}%
\pgfpathrectangle{\pgfqpoint{0.100000in}{0.209042in}}{\pgfqpoint{2.906250in}{2.906250in}}%
\pgfusepath{clip}%
\pgfsetbuttcap%
\pgfsetroundjoin%
\definecolor{currentfill}{rgb}{0.846416,0.272473,0.421631}%
\pgfsetfillcolor{currentfill}%
\pgfsetfillopacity{0.800000}%
\pgfsetlinewidth{1.003750pt}%
\definecolor{currentstroke}{rgb}{0.846416,0.272473,0.421631}%
\pgfsetstrokecolor{currentstroke}%
\pgfsetstrokeopacity{0.800000}%
\pgfsetdash{}{0pt}%
\pgfpathmoveto{\pgfqpoint{1.382803in}{1.903832in}}%
\pgfpathcurveto{\pgfqpoint{1.388627in}{1.903832in}}{\pgfqpoint{1.394213in}{1.906146in}}{\pgfqpoint{1.398331in}{1.910264in}}%
\pgfpathcurveto{\pgfqpoint{1.402449in}{1.914382in}}{\pgfqpoint{1.404763in}{1.919968in}}{\pgfqpoint{1.404763in}{1.925792in}}%
\pgfpathcurveto{\pgfqpoint{1.404763in}{1.931616in}}{\pgfqpoint{1.402449in}{1.937202in}}{\pgfqpoint{1.398331in}{1.941320in}}%
\pgfpathcurveto{\pgfqpoint{1.394213in}{1.945438in}}{\pgfqpoint{1.388627in}{1.947752in}}{\pgfqpoint{1.382803in}{1.947752in}}%
\pgfpathcurveto{\pgfqpoint{1.376979in}{1.947752in}}{\pgfqpoint{1.371393in}{1.945438in}}{\pgfqpoint{1.367275in}{1.941320in}}%
\pgfpathcurveto{\pgfqpoint{1.363156in}{1.937202in}}{\pgfqpoint{1.360843in}{1.931616in}}{\pgfqpoint{1.360843in}{1.925792in}}%
\pgfpathcurveto{\pgfqpoint{1.360843in}{1.919968in}}{\pgfqpoint{1.363156in}{1.914382in}}{\pgfqpoint{1.367275in}{1.910264in}}%
\pgfpathcurveto{\pgfqpoint{1.371393in}{1.906146in}}{\pgfqpoint{1.376979in}{1.903832in}}{\pgfqpoint{1.382803in}{1.903832in}}%
\pgfpathlineto{\pgfqpoint{1.382803in}{1.903832in}}%
\pgfpathclose%
\pgfusepath{stroke,fill}%
\end{pgfscope}%
\begin{pgfscope}%
\pgfpathrectangle{\pgfqpoint{0.100000in}{0.209042in}}{\pgfqpoint{2.906250in}{2.906250in}}%
\pgfusepath{clip}%
\pgfsetbuttcap%
\pgfsetroundjoin%
\definecolor{currentfill}{rgb}{0.834791,0.265540,0.427671}%
\pgfsetfillcolor{currentfill}%
\pgfsetfillopacity{0.800000}%
\pgfsetlinewidth{1.003750pt}%
\definecolor{currentstroke}{rgb}{0.834791,0.265540,0.427671}%
\pgfsetstrokecolor{currentstroke}%
\pgfsetstrokeopacity{0.800000}%
\pgfsetdash{}{0pt}%
\pgfpathmoveto{\pgfqpoint{1.432475in}{1.921498in}}%
\pgfpathcurveto{\pgfqpoint{1.438299in}{1.921498in}}{\pgfqpoint{1.443885in}{1.923812in}}{\pgfqpoint{1.448004in}{1.927930in}}%
\pgfpathcurveto{\pgfqpoint{1.452122in}{1.932049in}}{\pgfqpoint{1.454436in}{1.937635in}}{\pgfqpoint{1.454436in}{1.943459in}}%
\pgfpathcurveto{\pgfqpoint{1.454436in}{1.949283in}}{\pgfqpoint{1.452122in}{1.954869in}}{\pgfqpoint{1.448004in}{1.958987in}}%
\pgfpathcurveto{\pgfqpoint{1.443885in}{1.963105in}}{\pgfqpoint{1.438299in}{1.965419in}}{\pgfqpoint{1.432475in}{1.965419in}}%
\pgfpathcurveto{\pgfqpoint{1.426651in}{1.965419in}}{\pgfqpoint{1.421065in}{1.963105in}}{\pgfqpoint{1.416947in}{1.958987in}}%
\pgfpathcurveto{\pgfqpoint{1.412829in}{1.954869in}}{\pgfqpoint{1.410515in}{1.949283in}}{\pgfqpoint{1.410515in}{1.943459in}}%
\pgfpathcurveto{\pgfqpoint{1.410515in}{1.937635in}}{\pgfqpoint{1.412829in}{1.932049in}}{\pgfqpoint{1.416947in}{1.927930in}}%
\pgfpathcurveto{\pgfqpoint{1.421065in}{1.923812in}}{\pgfqpoint{1.426651in}{1.921498in}}{\pgfqpoint{1.432475in}{1.921498in}}%
\pgfpathlineto{\pgfqpoint{1.432475in}{1.921498in}}%
\pgfpathclose%
\pgfusepath{stroke,fill}%
\end{pgfscope}%
\begin{pgfscope}%
\pgfpathrectangle{\pgfqpoint{0.100000in}{0.209042in}}{\pgfqpoint{2.906250in}{2.906250in}}%
\pgfusepath{clip}%
\pgfsetbuttcap%
\pgfsetroundjoin%
\definecolor{currentfill}{rgb}{0.816914,0.255895,0.436461}%
\pgfsetfillcolor{currentfill}%
\pgfsetfillopacity{0.800000}%
\pgfsetlinewidth{1.003750pt}%
\definecolor{currentstroke}{rgb}{0.816914,0.255895,0.436461}%
\pgfsetstrokecolor{currentstroke}%
\pgfsetstrokeopacity{0.800000}%
\pgfsetdash{}{0pt}%
\pgfpathmoveto{\pgfqpoint{1.672791in}{1.657444in}}%
\pgfpathcurveto{\pgfqpoint{1.678615in}{1.657444in}}{\pgfqpoint{1.684201in}{1.659758in}}{\pgfqpoint{1.688319in}{1.663876in}}%
\pgfpathcurveto{\pgfqpoint{1.692437in}{1.667994in}}{\pgfqpoint{1.694751in}{1.673580in}}{\pgfqpoint{1.694751in}{1.679404in}}%
\pgfpathcurveto{\pgfqpoint{1.694751in}{1.685228in}}{\pgfqpoint{1.692437in}{1.690815in}}{\pgfqpoint{1.688319in}{1.694933in}}%
\pgfpathcurveto{\pgfqpoint{1.684201in}{1.699051in}}{\pgfqpoint{1.678615in}{1.701365in}}{\pgfqpoint{1.672791in}{1.701365in}}%
\pgfpathcurveto{\pgfqpoint{1.666967in}{1.701365in}}{\pgfqpoint{1.661381in}{1.699051in}}{\pgfqpoint{1.657263in}{1.694933in}}%
\pgfpathcurveto{\pgfqpoint{1.653144in}{1.690815in}}{\pgfqpoint{1.650831in}{1.685228in}}{\pgfqpoint{1.650831in}{1.679404in}}%
\pgfpathcurveto{\pgfqpoint{1.650831in}{1.673580in}}{\pgfqpoint{1.653144in}{1.667994in}}{\pgfqpoint{1.657263in}{1.663876in}}%
\pgfpathcurveto{\pgfqpoint{1.661381in}{1.659758in}}{\pgfqpoint{1.666967in}{1.657444in}}{\pgfqpoint{1.672791in}{1.657444in}}%
\pgfpathlineto{\pgfqpoint{1.672791in}{1.657444in}}%
\pgfpathclose%
\pgfusepath{stroke,fill}%
\end{pgfscope}%
\begin{pgfscope}%
\pgfpathrectangle{\pgfqpoint{0.100000in}{0.209042in}}{\pgfqpoint{2.906250in}{2.906250in}}%
\pgfusepath{clip}%
\pgfsetbuttcap%
\pgfsetroundjoin%
\definecolor{currentfill}{rgb}{0.857763,0.279857,0.415496}%
\pgfsetfillcolor{currentfill}%
\pgfsetfillopacity{0.800000}%
\pgfsetlinewidth{1.003750pt}%
\definecolor{currentstroke}{rgb}{0.857763,0.279857,0.415496}%
\pgfsetstrokecolor{currentstroke}%
\pgfsetstrokeopacity{0.800000}%
\pgfsetdash{}{0pt}%
\pgfpathmoveto{\pgfqpoint{1.800943in}{1.866850in}}%
\pgfpathcurveto{\pgfqpoint{1.806767in}{1.866850in}}{\pgfqpoint{1.812353in}{1.869164in}}{\pgfqpoint{1.816471in}{1.873282in}}%
\pgfpathcurveto{\pgfqpoint{1.820589in}{1.877400in}}{\pgfqpoint{1.822903in}{1.882986in}}{\pgfqpoint{1.822903in}{1.888810in}}%
\pgfpathcurveto{\pgfqpoint{1.822903in}{1.894634in}}{\pgfqpoint{1.820589in}{1.900220in}}{\pgfqpoint{1.816471in}{1.904339in}}%
\pgfpathcurveto{\pgfqpoint{1.812353in}{1.908457in}}{\pgfqpoint{1.806767in}{1.910771in}}{\pgfqpoint{1.800943in}{1.910771in}}%
\pgfpathcurveto{\pgfqpoint{1.795119in}{1.910771in}}{\pgfqpoint{1.789533in}{1.908457in}}{\pgfqpoint{1.785415in}{1.904339in}}%
\pgfpathcurveto{\pgfqpoint{1.781297in}{1.900220in}}{\pgfqpoint{1.778983in}{1.894634in}}{\pgfqpoint{1.778983in}{1.888810in}}%
\pgfpathcurveto{\pgfqpoint{1.778983in}{1.882986in}}{\pgfqpoint{1.781297in}{1.877400in}}{\pgfqpoint{1.785415in}{1.873282in}}%
\pgfpathcurveto{\pgfqpoint{1.789533in}{1.869164in}}{\pgfqpoint{1.795119in}{1.866850in}}{\pgfqpoint{1.800943in}{1.866850in}}%
\pgfpathlineto{\pgfqpoint{1.800943in}{1.866850in}}%
\pgfpathclose%
\pgfusepath{stroke,fill}%
\end{pgfscope}%
\begin{pgfscope}%
\pgfpathrectangle{\pgfqpoint{0.100000in}{0.209042in}}{\pgfqpoint{2.906250in}{2.906250in}}%
\pgfusepath{clip}%
\pgfsetbuttcap%
\pgfsetroundjoin%
\definecolor{currentfill}{rgb}{0.761077,0.231214,0.460162}%
\pgfsetfillcolor{currentfill}%
\pgfsetfillopacity{0.800000}%
\pgfsetlinewidth{1.003750pt}%
\definecolor{currentstroke}{rgb}{0.761077,0.231214,0.460162}%
\pgfsetstrokecolor{currentstroke}%
\pgfsetstrokeopacity{0.800000}%
\pgfsetdash{}{0pt}%
\pgfpathmoveto{\pgfqpoint{1.580165in}{1.734283in}}%
\pgfpathcurveto{\pgfqpoint{1.585989in}{1.734283in}}{\pgfqpoint{1.591575in}{1.736596in}}{\pgfqpoint{1.595694in}{1.740715in}}%
\pgfpathcurveto{\pgfqpoint{1.599812in}{1.744833in}}{\pgfqpoint{1.602126in}{1.750419in}}{\pgfqpoint{1.602126in}{1.756243in}}%
\pgfpathcurveto{\pgfqpoint{1.602126in}{1.762067in}}{\pgfqpoint{1.599812in}{1.767653in}}{\pgfqpoint{1.595694in}{1.771771in}}%
\pgfpathcurveto{\pgfqpoint{1.591575in}{1.775889in}}{\pgfqpoint{1.585989in}{1.778203in}}{\pgfqpoint{1.580165in}{1.778203in}}%
\pgfpathcurveto{\pgfqpoint{1.574341in}{1.778203in}}{\pgfqpoint{1.568755in}{1.775889in}}{\pgfqpoint{1.564637in}{1.771771in}}%
\pgfpathcurveto{\pgfqpoint{1.560519in}{1.767653in}}{\pgfqpoint{1.558205in}{1.762067in}}{\pgfqpoint{1.558205in}{1.756243in}}%
\pgfpathcurveto{\pgfqpoint{1.558205in}{1.750419in}}{\pgfqpoint{1.560519in}{1.744833in}}{\pgfqpoint{1.564637in}{1.740715in}}%
\pgfpathcurveto{\pgfqpoint{1.568755in}{1.736596in}}{\pgfqpoint{1.574341in}{1.734283in}}{\pgfqpoint{1.580165in}{1.734283in}}%
\pgfpathlineto{\pgfqpoint{1.580165in}{1.734283in}}%
\pgfpathclose%
\pgfusepath{stroke,fill}%
\end{pgfscope}%
\begin{pgfscope}%
\pgfpathrectangle{\pgfqpoint{0.100000in}{0.209042in}}{\pgfqpoint{2.906250in}{2.906250in}}%
\pgfusepath{clip}%
\pgfsetbuttcap%
\pgfsetroundjoin%
\definecolor{currentfill}{rgb}{0.852126,0.276106,0.418573}%
\pgfsetfillcolor{currentfill}%
\pgfsetfillopacity{0.800000}%
\pgfsetlinewidth{1.003750pt}%
\definecolor{currentstroke}{rgb}{0.852126,0.276106,0.418573}%
\pgfsetstrokecolor{currentstroke}%
\pgfsetstrokeopacity{0.800000}%
\pgfsetdash{}{0pt}%
\pgfpathmoveto{\pgfqpoint{1.298221in}{1.869513in}}%
\pgfpathcurveto{\pgfqpoint{1.304045in}{1.869513in}}{\pgfqpoint{1.309631in}{1.871827in}}{\pgfqpoint{1.313749in}{1.875945in}}%
\pgfpathcurveto{\pgfqpoint{1.317867in}{1.880063in}}{\pgfqpoint{1.320181in}{1.885649in}}{\pgfqpoint{1.320181in}{1.891473in}}%
\pgfpathcurveto{\pgfqpoint{1.320181in}{1.897297in}}{\pgfqpoint{1.317867in}{1.902883in}}{\pgfqpoint{1.313749in}{1.907001in}}%
\pgfpathcurveto{\pgfqpoint{1.309631in}{1.911119in}}{\pgfqpoint{1.304045in}{1.913433in}}{\pgfqpoint{1.298221in}{1.913433in}}%
\pgfpathcurveto{\pgfqpoint{1.292397in}{1.913433in}}{\pgfqpoint{1.286811in}{1.911119in}}{\pgfqpoint{1.282692in}{1.907001in}}%
\pgfpathcurveto{\pgfqpoint{1.278574in}{1.902883in}}{\pgfqpoint{1.276260in}{1.897297in}}{\pgfqpoint{1.276260in}{1.891473in}}%
\pgfpathcurveto{\pgfqpoint{1.276260in}{1.885649in}}{\pgfqpoint{1.278574in}{1.880063in}}{\pgfqpoint{1.282692in}{1.875945in}}%
\pgfpathcurveto{\pgfqpoint{1.286811in}{1.871827in}}{\pgfqpoint{1.292397in}{1.869513in}}{\pgfqpoint{1.298221in}{1.869513in}}%
\pgfpathlineto{\pgfqpoint{1.298221in}{1.869513in}}%
\pgfpathclose%
\pgfusepath{stroke,fill}%
\end{pgfscope}%
\begin{pgfscope}%
\pgfpathrectangle{\pgfqpoint{0.100000in}{0.209042in}}{\pgfqpoint{2.906250in}{2.906250in}}%
\pgfusepath{clip}%
\pgfsetbuttcap%
\pgfsetroundjoin%
\definecolor{currentfill}{rgb}{0.735616,0.221713,0.469180}%
\pgfsetfillcolor{currentfill}%
\pgfsetfillopacity{0.800000}%
\pgfsetlinewidth{1.003750pt}%
\definecolor{currentstroke}{rgb}{0.735616,0.221713,0.469180}%
\pgfsetstrokecolor{currentstroke}%
\pgfsetstrokeopacity{0.800000}%
\pgfsetdash{}{0pt}%
\pgfpathmoveto{\pgfqpoint{1.641688in}{1.838212in}}%
\pgfpathcurveto{\pgfqpoint{1.647512in}{1.838212in}}{\pgfqpoint{1.653099in}{1.840526in}}{\pgfqpoint{1.657217in}{1.844644in}}%
\pgfpathcurveto{\pgfqpoint{1.661335in}{1.848762in}}{\pgfqpoint{1.663649in}{1.854349in}}{\pgfqpoint{1.663649in}{1.860172in}}%
\pgfpathcurveto{\pgfqpoint{1.663649in}{1.865996in}}{\pgfqpoint{1.661335in}{1.871583in}}{\pgfqpoint{1.657217in}{1.875701in}}%
\pgfpathcurveto{\pgfqpoint{1.653099in}{1.879819in}}{\pgfqpoint{1.647512in}{1.882133in}}{\pgfqpoint{1.641688in}{1.882133in}}%
\pgfpathcurveto{\pgfqpoint{1.635865in}{1.882133in}}{\pgfqpoint{1.630278in}{1.879819in}}{\pgfqpoint{1.626160in}{1.875701in}}%
\pgfpathcurveto{\pgfqpoint{1.622042in}{1.871583in}}{\pgfqpoint{1.619728in}{1.865996in}}{\pgfqpoint{1.619728in}{1.860172in}}%
\pgfpathcurveto{\pgfqpoint{1.619728in}{1.854349in}}{\pgfqpoint{1.622042in}{1.848762in}}{\pgfqpoint{1.626160in}{1.844644in}}%
\pgfpathcurveto{\pgfqpoint{1.630278in}{1.840526in}}{\pgfqpoint{1.635865in}{1.838212in}}{\pgfqpoint{1.641688in}{1.838212in}}%
\pgfpathlineto{\pgfqpoint{1.641688in}{1.838212in}}%
\pgfpathclose%
\pgfusepath{stroke,fill}%
\end{pgfscope}%
\begin{pgfscope}%
\pgfpathrectangle{\pgfqpoint{0.100000in}{0.209042in}}{\pgfqpoint{2.906250in}{2.906250in}}%
\pgfusepath{clip}%
\pgfsetbuttcap%
\pgfsetroundjoin%
\definecolor{currentfill}{rgb}{0.950210,0.390820,0.362468}%
\pgfsetfillcolor{currentfill}%
\pgfsetfillopacity{0.800000}%
\pgfsetlinewidth{1.003750pt}%
\definecolor{currentstroke}{rgb}{0.950210,0.390820,0.362468}%
\pgfsetstrokecolor{currentstroke}%
\pgfsetstrokeopacity{0.800000}%
\pgfsetdash{}{0pt}%
\pgfpathmoveto{\pgfqpoint{1.336401in}{1.451895in}}%
\pgfpathcurveto{\pgfqpoint{1.342225in}{1.451895in}}{\pgfqpoint{1.347811in}{1.454209in}}{\pgfqpoint{1.351929in}{1.458327in}}%
\pgfpathcurveto{\pgfqpoint{1.356047in}{1.462445in}}{\pgfqpoint{1.358361in}{1.468031in}}{\pgfqpoint{1.358361in}{1.473855in}}%
\pgfpathcurveto{\pgfqpoint{1.358361in}{1.479679in}}{\pgfqpoint{1.356047in}{1.485265in}}{\pgfqpoint{1.351929in}{1.489384in}}%
\pgfpathcurveto{\pgfqpoint{1.347811in}{1.493502in}}{\pgfqpoint{1.342225in}{1.495816in}}{\pgfqpoint{1.336401in}{1.495816in}}%
\pgfpathcurveto{\pgfqpoint{1.330577in}{1.495816in}}{\pgfqpoint{1.324991in}{1.493502in}}{\pgfqpoint{1.320872in}{1.489384in}}%
\pgfpathcurveto{\pgfqpoint{1.316754in}{1.485265in}}{\pgfqpoint{1.314440in}{1.479679in}}{\pgfqpoint{1.314440in}{1.473855in}}%
\pgfpathcurveto{\pgfqpoint{1.314440in}{1.468031in}}{\pgfqpoint{1.316754in}{1.462445in}}{\pgfqpoint{1.320872in}{1.458327in}}%
\pgfpathcurveto{\pgfqpoint{1.324991in}{1.454209in}}{\pgfqpoint{1.330577in}{1.451895in}}{\pgfqpoint{1.336401in}{1.451895in}}%
\pgfpathlineto{\pgfqpoint{1.336401in}{1.451895in}}%
\pgfpathclose%
\pgfusepath{stroke,fill}%
\end{pgfscope}%
\begin{pgfscope}%
\pgfpathrectangle{\pgfqpoint{0.100000in}{0.209042in}}{\pgfqpoint{2.906250in}{2.906250in}}%
\pgfusepath{clip}%
\pgfsetbuttcap%
\pgfsetroundjoin%
\definecolor{currentfill}{rgb}{0.748378,0.226377,0.464794}%
\pgfsetfillcolor{currentfill}%
\pgfsetfillopacity{0.800000}%
\pgfsetlinewidth{1.003750pt}%
\definecolor{currentstroke}{rgb}{0.748378,0.226377,0.464794}%
\pgfsetstrokecolor{currentstroke}%
\pgfsetstrokeopacity{0.800000}%
\pgfsetdash{}{0pt}%
\pgfpathmoveto{\pgfqpoint{1.434551in}{1.856929in}}%
\pgfpathcurveto{\pgfqpoint{1.440375in}{1.856929in}}{\pgfqpoint{1.445961in}{1.859243in}}{\pgfqpoint{1.450079in}{1.863361in}}%
\pgfpathcurveto{\pgfqpoint{1.454197in}{1.867479in}}{\pgfqpoint{1.456511in}{1.873065in}}{\pgfqpoint{1.456511in}{1.878889in}}%
\pgfpathcurveto{\pgfqpoint{1.456511in}{1.884713in}}{\pgfqpoint{1.454197in}{1.890299in}}{\pgfqpoint{1.450079in}{1.894417in}}%
\pgfpathcurveto{\pgfqpoint{1.445961in}{1.898535in}}{\pgfqpoint{1.440375in}{1.900849in}}{\pgfqpoint{1.434551in}{1.900849in}}%
\pgfpathcurveto{\pgfqpoint{1.428727in}{1.900849in}}{\pgfqpoint{1.423141in}{1.898535in}}{\pgfqpoint{1.419022in}{1.894417in}}%
\pgfpathcurveto{\pgfqpoint{1.414904in}{1.890299in}}{\pgfqpoint{1.412590in}{1.884713in}}{\pgfqpoint{1.412590in}{1.878889in}}%
\pgfpathcurveto{\pgfqpoint{1.412590in}{1.873065in}}{\pgfqpoint{1.414904in}{1.867479in}}{\pgfqpoint{1.419022in}{1.863361in}}%
\pgfpathcurveto{\pgfqpoint{1.423141in}{1.859243in}}{\pgfqpoint{1.428727in}{1.856929in}}{\pgfqpoint{1.434551in}{1.856929in}}%
\pgfpathlineto{\pgfqpoint{1.434551in}{1.856929in}}%
\pgfpathclose%
\pgfusepath{stroke,fill}%
\end{pgfscope}%
\begin{pgfscope}%
\pgfpathrectangle{\pgfqpoint{0.100000in}{0.209042in}}{\pgfqpoint{2.906250in}{2.906250in}}%
\pgfusepath{clip}%
\pgfsetbuttcap%
\pgfsetroundjoin%
\definecolor{currentfill}{rgb}{0.816914,0.255895,0.436461}%
\pgfsetfillcolor{currentfill}%
\pgfsetfillopacity{0.800000}%
\pgfsetlinewidth{1.003750pt}%
\definecolor{currentstroke}{rgb}{0.816914,0.255895,0.436461}%
\pgfsetstrokecolor{currentstroke}%
\pgfsetstrokeopacity{0.800000}%
\pgfsetdash{}{0pt}%
\pgfpathmoveto{\pgfqpoint{1.685815in}{1.502585in}}%
\pgfpathcurveto{\pgfqpoint{1.691638in}{1.502585in}}{\pgfqpoint{1.697225in}{1.504899in}}{\pgfqpoint{1.701343in}{1.509017in}}%
\pgfpathcurveto{\pgfqpoint{1.705461in}{1.513135in}}{\pgfqpoint{1.707775in}{1.518721in}}{\pgfqpoint{1.707775in}{1.524545in}}%
\pgfpathcurveto{\pgfqpoint{1.707775in}{1.530369in}}{\pgfqpoint{1.705461in}{1.535955in}}{\pgfqpoint{1.701343in}{1.540073in}}%
\pgfpathcurveto{\pgfqpoint{1.697225in}{1.544192in}}{\pgfqpoint{1.691638in}{1.546505in}}{\pgfqpoint{1.685815in}{1.546505in}}%
\pgfpathcurveto{\pgfqpoint{1.679991in}{1.546505in}}{\pgfqpoint{1.674404in}{1.544192in}}{\pgfqpoint{1.670286in}{1.540073in}}%
\pgfpathcurveto{\pgfqpoint{1.666168in}{1.535955in}}{\pgfqpoint{1.663854in}{1.530369in}}{\pgfqpoint{1.663854in}{1.524545in}}%
\pgfpathcurveto{\pgfqpoint{1.663854in}{1.518721in}}{\pgfqpoint{1.666168in}{1.513135in}}{\pgfqpoint{1.670286in}{1.509017in}}%
\pgfpathcurveto{\pgfqpoint{1.674404in}{1.504899in}}{\pgfqpoint{1.679991in}{1.502585in}}{\pgfqpoint{1.685815in}{1.502585in}}%
\pgfpathlineto{\pgfqpoint{1.685815in}{1.502585in}}%
\pgfpathclose%
\pgfusepath{stroke,fill}%
\end{pgfscope}%
\begin{pgfscope}%
\pgfpathrectangle{\pgfqpoint{0.100000in}{0.209042in}}{\pgfqpoint{2.906250in}{2.906250in}}%
\pgfusepath{clip}%
\pgfsetbuttcap%
\pgfsetroundjoin%
\definecolor{currentfill}{rgb}{0.652056,0.193986,0.491611}%
\pgfsetfillcolor{currentfill}%
\pgfsetfillopacity{0.800000}%
\pgfsetlinewidth{1.003750pt}%
\definecolor{currentstroke}{rgb}{0.652056,0.193986,0.491611}%
\pgfsetstrokecolor{currentstroke}%
\pgfsetstrokeopacity{0.800000}%
\pgfsetdash{}{0pt}%
\pgfpathmoveto{\pgfqpoint{1.599733in}{1.714305in}}%
\pgfpathcurveto{\pgfqpoint{1.605557in}{1.714305in}}{\pgfqpoint{1.611143in}{1.716619in}}{\pgfqpoint{1.615261in}{1.720737in}}%
\pgfpathcurveto{\pgfqpoint{1.619379in}{1.724855in}}{\pgfqpoint{1.621693in}{1.730441in}}{\pgfqpoint{1.621693in}{1.736265in}}%
\pgfpathcurveto{\pgfqpoint{1.621693in}{1.742089in}}{\pgfqpoint{1.619379in}{1.747675in}}{\pgfqpoint{1.615261in}{1.751794in}}%
\pgfpathcurveto{\pgfqpoint{1.611143in}{1.755912in}}{\pgfqpoint{1.605557in}{1.758226in}}{\pgfqpoint{1.599733in}{1.758226in}}%
\pgfpathcurveto{\pgfqpoint{1.593909in}{1.758226in}}{\pgfqpoint{1.588323in}{1.755912in}}{\pgfqpoint{1.584204in}{1.751794in}}%
\pgfpathcurveto{\pgfqpoint{1.580086in}{1.747675in}}{\pgfqpoint{1.577772in}{1.742089in}}{\pgfqpoint{1.577772in}{1.736265in}}%
\pgfpathcurveto{\pgfqpoint{1.577772in}{1.730441in}}{\pgfqpoint{1.580086in}{1.724855in}}{\pgfqpoint{1.584204in}{1.720737in}}%
\pgfpathcurveto{\pgfqpoint{1.588323in}{1.716619in}}{\pgfqpoint{1.593909in}{1.714305in}}{\pgfqpoint{1.599733in}{1.714305in}}%
\pgfpathlineto{\pgfqpoint{1.599733in}{1.714305in}}%
\pgfpathclose%
\pgfusepath{stroke,fill}%
\end{pgfscope}%
\begin{pgfscope}%
\pgfpathrectangle{\pgfqpoint{0.100000in}{0.209042in}}{\pgfqpoint{2.906250in}{2.906250in}}%
\pgfusepath{clip}%
\pgfsetbuttcap%
\pgfsetroundjoin%
\definecolor{currentfill}{rgb}{0.652056,0.193986,0.491611}%
\pgfsetfillcolor{currentfill}%
\pgfsetfillopacity{0.800000}%
\pgfsetlinewidth{1.003750pt}%
\definecolor{currentstroke}{rgb}{0.652056,0.193986,0.491611}%
\pgfsetstrokecolor{currentstroke}%
\pgfsetstrokeopacity{0.800000}%
\pgfsetdash{}{0pt}%
\pgfpathmoveto{\pgfqpoint{1.520605in}{1.680582in}}%
\pgfpathcurveto{\pgfqpoint{1.526429in}{1.680582in}}{\pgfqpoint{1.532015in}{1.682896in}}{\pgfqpoint{1.536134in}{1.687014in}}%
\pgfpathcurveto{\pgfqpoint{1.540252in}{1.691133in}}{\pgfqpoint{1.542566in}{1.696719in}}{\pgfqpoint{1.542566in}{1.702543in}}%
\pgfpathcurveto{\pgfqpoint{1.542566in}{1.708367in}}{\pgfqpoint{1.540252in}{1.713953in}}{\pgfqpoint{1.536134in}{1.718071in}}%
\pgfpathcurveto{\pgfqpoint{1.532015in}{1.722189in}}{\pgfqpoint{1.526429in}{1.724503in}}{\pgfqpoint{1.520605in}{1.724503in}}%
\pgfpathcurveto{\pgfqpoint{1.514781in}{1.724503in}}{\pgfqpoint{1.509195in}{1.722189in}}{\pgfqpoint{1.505077in}{1.718071in}}%
\pgfpathcurveto{\pgfqpoint{1.500959in}{1.713953in}}{\pgfqpoint{1.498645in}{1.708367in}}{\pgfqpoint{1.498645in}{1.702543in}}%
\pgfpathcurveto{\pgfqpoint{1.498645in}{1.696719in}}{\pgfqpoint{1.500959in}{1.691133in}}{\pgfqpoint{1.505077in}{1.687014in}}%
\pgfpathcurveto{\pgfqpoint{1.509195in}{1.682896in}}{\pgfqpoint{1.514781in}{1.680582in}}{\pgfqpoint{1.520605in}{1.680582in}}%
\pgfpathlineto{\pgfqpoint{1.520605in}{1.680582in}}%
\pgfpathclose%
\pgfusepath{stroke,fill}%
\end{pgfscope}%
\begin{pgfscope}%
\pgfpathrectangle{\pgfqpoint{0.100000in}{0.209042in}}{\pgfqpoint{2.906250in}{2.906250in}}%
\pgfusepath{clip}%
\pgfsetbuttcap%
\pgfsetroundjoin%
\definecolor{currentfill}{rgb}{0.987053,0.991438,0.749504}%
\pgfsetfillcolor{currentfill}%
\pgfsetfillopacity{0.800000}%
\pgfsetlinewidth{1.003750pt}%
\definecolor{currentstroke}{rgb}{0.987053,0.991438,0.749504}%
\pgfsetstrokecolor{currentstroke}%
\pgfsetstrokeopacity{0.800000}%
\pgfsetdash{}{0pt}%
\pgfpathmoveto{\pgfqpoint{1.738196in}{2.504930in}}%
\pgfpathcurveto{\pgfqpoint{1.744020in}{2.504930in}}{\pgfqpoint{1.749606in}{2.507244in}}{\pgfqpoint{1.753724in}{2.511362in}}%
\pgfpathcurveto{\pgfqpoint{1.757842in}{2.515480in}}{\pgfqpoint{1.760156in}{2.521066in}}{\pgfqpoint{1.760156in}{2.526890in}}%
\pgfpathcurveto{\pgfqpoint{1.760156in}{2.532714in}}{\pgfqpoint{1.757842in}{2.538300in}}{\pgfqpoint{1.753724in}{2.542419in}}%
\pgfpathcurveto{\pgfqpoint{1.749606in}{2.546537in}}{\pgfqpoint{1.744020in}{2.548851in}}{\pgfqpoint{1.738196in}{2.548851in}}%
\pgfpathcurveto{\pgfqpoint{1.732372in}{2.548851in}}{\pgfqpoint{1.726786in}{2.546537in}}{\pgfqpoint{1.722668in}{2.542419in}}%
\pgfpathcurveto{\pgfqpoint{1.718550in}{2.538300in}}{\pgfqpoint{1.716236in}{2.532714in}}{\pgfqpoint{1.716236in}{2.526890in}}%
\pgfpathcurveto{\pgfqpoint{1.716236in}{2.521066in}}{\pgfqpoint{1.718550in}{2.515480in}}{\pgfqpoint{1.722668in}{2.511362in}}%
\pgfpathcurveto{\pgfqpoint{1.726786in}{2.507244in}}{\pgfqpoint{1.732372in}{2.504930in}}{\pgfqpoint{1.738196in}{2.504930in}}%
\pgfpathlineto{\pgfqpoint{1.738196in}{2.504930in}}%
\pgfpathclose%
\pgfusepath{stroke,fill}%
\end{pgfscope}%
\begin{pgfscope}%
\pgfpathrectangle{\pgfqpoint{0.100000in}{0.209042in}}{\pgfqpoint{2.906250in}{2.906250in}}%
\pgfusepath{clip}%
\pgfsetbuttcap%
\pgfsetroundjoin%
\definecolor{currentfill}{rgb}{0.613617,0.181811,0.498536}%
\pgfsetfillcolor{currentfill}%
\pgfsetfillopacity{0.800000}%
\pgfsetlinewidth{1.003750pt}%
\definecolor{currentstroke}{rgb}{0.613617,0.181811,0.498536}%
\pgfsetstrokecolor{currentstroke}%
\pgfsetstrokeopacity{0.800000}%
\pgfsetdash{}{0pt}%
\pgfpathmoveto{\pgfqpoint{1.668800in}{1.738793in}}%
\pgfpathcurveto{\pgfqpoint{1.674624in}{1.738793in}}{\pgfqpoint{1.680210in}{1.741107in}}{\pgfqpoint{1.684328in}{1.745225in}}%
\pgfpathcurveto{\pgfqpoint{1.688446in}{1.749343in}}{\pgfqpoint{1.690760in}{1.754929in}}{\pgfqpoint{1.690760in}{1.760753in}}%
\pgfpathcurveto{\pgfqpoint{1.690760in}{1.766577in}}{\pgfqpoint{1.688446in}{1.772163in}}{\pgfqpoint{1.684328in}{1.776282in}}%
\pgfpathcurveto{\pgfqpoint{1.680210in}{1.780400in}}{\pgfqpoint{1.674624in}{1.782714in}}{\pgfqpoint{1.668800in}{1.782714in}}%
\pgfpathcurveto{\pgfqpoint{1.662976in}{1.782714in}}{\pgfqpoint{1.657390in}{1.780400in}}{\pgfqpoint{1.653272in}{1.776282in}}%
\pgfpathcurveto{\pgfqpoint{1.649153in}{1.772163in}}{\pgfqpoint{1.646840in}{1.766577in}}{\pgfqpoint{1.646840in}{1.760753in}}%
\pgfpathcurveto{\pgfqpoint{1.646840in}{1.754929in}}{\pgfqpoint{1.649153in}{1.749343in}}{\pgfqpoint{1.653272in}{1.745225in}}%
\pgfpathcurveto{\pgfqpoint{1.657390in}{1.741107in}}{\pgfqpoint{1.662976in}{1.738793in}}{\pgfqpoint{1.668800in}{1.738793in}}%
\pgfpathlineto{\pgfqpoint{1.668800in}{1.738793in}}%
\pgfpathclose%
\pgfusepath{stroke,fill}%
\end{pgfscope}%
\begin{pgfscope}%
\pgfpathrectangle{\pgfqpoint{0.100000in}{0.209042in}}{\pgfqpoint{2.906250in}{2.906250in}}%
\pgfusepath{clip}%
\pgfsetbuttcap%
\pgfsetroundjoin%
\definecolor{currentfill}{rgb}{0.703532,0.210638,0.479049}%
\pgfsetfillcolor{currentfill}%
\pgfsetfillopacity{0.800000}%
\pgfsetlinewidth{1.003750pt}%
\definecolor{currentstroke}{rgb}{0.703532,0.210638,0.479049}%
\pgfsetstrokecolor{currentstroke}%
\pgfsetstrokeopacity{0.800000}%
\pgfsetdash{}{0pt}%
\pgfpathmoveto{\pgfqpoint{1.437202in}{1.583695in}}%
\pgfpathcurveto{\pgfqpoint{1.443026in}{1.583695in}}{\pgfqpoint{1.448612in}{1.586009in}}{\pgfqpoint{1.452730in}{1.590127in}}%
\pgfpathcurveto{\pgfqpoint{1.456848in}{1.594245in}}{\pgfqpoint{1.459162in}{1.599831in}}{\pgfqpoint{1.459162in}{1.605655in}}%
\pgfpathcurveto{\pgfqpoint{1.459162in}{1.611479in}}{\pgfqpoint{1.456848in}{1.617065in}}{\pgfqpoint{1.452730in}{1.621183in}}%
\pgfpathcurveto{\pgfqpoint{1.448612in}{1.625302in}}{\pgfqpoint{1.443026in}{1.627615in}}{\pgfqpoint{1.437202in}{1.627615in}}%
\pgfpathcurveto{\pgfqpoint{1.431378in}{1.627615in}}{\pgfqpoint{1.425792in}{1.625302in}}{\pgfqpoint{1.421673in}{1.621183in}}%
\pgfpathcurveto{\pgfqpoint{1.417555in}{1.617065in}}{\pgfqpoint{1.415241in}{1.611479in}}{\pgfqpoint{1.415241in}{1.605655in}}%
\pgfpathcurveto{\pgfqpoint{1.415241in}{1.599831in}}{\pgfqpoint{1.417555in}{1.594245in}}{\pgfqpoint{1.421673in}{1.590127in}}%
\pgfpathcurveto{\pgfqpoint{1.425792in}{1.586009in}}{\pgfqpoint{1.431378in}{1.583695in}}{\pgfqpoint{1.437202in}{1.583695in}}%
\pgfpathlineto{\pgfqpoint{1.437202in}{1.583695in}}%
\pgfpathclose%
\pgfusepath{stroke,fill}%
\end{pgfscope}%
\begin{pgfscope}%
\pgfpathrectangle{\pgfqpoint{0.100000in}{0.209042in}}{\pgfqpoint{2.906250in}{2.906250in}}%
\pgfusepath{clip}%
\pgfsetbuttcap%
\pgfsetroundjoin%
\definecolor{currentfill}{rgb}{0.804752,0.249911,0.442102}%
\pgfsetfillcolor{currentfill}%
\pgfsetfillopacity{0.800000}%
\pgfsetlinewidth{1.003750pt}%
\definecolor{currentstroke}{rgb}{0.804752,0.249911,0.442102}%
\pgfsetstrokecolor{currentstroke}%
\pgfsetstrokeopacity{0.800000}%
\pgfsetdash{}{0pt}%
\pgfpathmoveto{\pgfqpoint{1.618842in}{2.070233in}}%
\pgfpathcurveto{\pgfqpoint{1.624665in}{2.070233in}}{\pgfqpoint{1.630252in}{2.072547in}}{\pgfqpoint{1.634370in}{2.076665in}}%
\pgfpathcurveto{\pgfqpoint{1.638488in}{2.080783in}}{\pgfqpoint{1.640802in}{2.086369in}}{\pgfqpoint{1.640802in}{2.092193in}}%
\pgfpathcurveto{\pgfqpoint{1.640802in}{2.098017in}}{\pgfqpoint{1.638488in}{2.103603in}}{\pgfqpoint{1.634370in}{2.107722in}}%
\pgfpathcurveto{\pgfqpoint{1.630252in}{2.111840in}}{\pgfqpoint{1.624665in}{2.114154in}}{\pgfqpoint{1.618842in}{2.114154in}}%
\pgfpathcurveto{\pgfqpoint{1.613018in}{2.114154in}}{\pgfqpoint{1.607431in}{2.111840in}}{\pgfqpoint{1.603313in}{2.107722in}}%
\pgfpathcurveto{\pgfqpoint{1.599195in}{2.103603in}}{\pgfqpoint{1.596881in}{2.098017in}}{\pgfqpoint{1.596881in}{2.092193in}}%
\pgfpathcurveto{\pgfqpoint{1.596881in}{2.086369in}}{\pgfqpoint{1.599195in}{2.080783in}}{\pgfqpoint{1.603313in}{2.076665in}}%
\pgfpathcurveto{\pgfqpoint{1.607431in}{2.072547in}}{\pgfqpoint{1.613018in}{2.070233in}}{\pgfqpoint{1.618842in}{2.070233in}}%
\pgfpathlineto{\pgfqpoint{1.618842in}{2.070233in}}%
\pgfpathclose%
\pgfusepath{stroke,fill}%
\end{pgfscope}%
\begin{pgfscope}%
\pgfpathrectangle{\pgfqpoint{0.100000in}{0.209042in}}{\pgfqpoint{2.906250in}{2.906250in}}%
\pgfusepath{clip}%
\pgfsetbuttcap%
\pgfsetroundjoin%
\definecolor{currentfill}{rgb}{0.664915,0.198075,0.488836}%
\pgfsetfillcolor{currentfill}%
\pgfsetfillopacity{0.800000}%
\pgfsetlinewidth{1.003750pt}%
\definecolor{currentstroke}{rgb}{0.664915,0.198075,0.488836}%
\pgfsetstrokecolor{currentstroke}%
\pgfsetstrokeopacity{0.800000}%
\pgfsetdash{}{0pt}%
\pgfpathmoveto{\pgfqpoint{1.754162in}{1.873445in}}%
\pgfpathcurveto{\pgfqpoint{1.759986in}{1.873445in}}{\pgfqpoint{1.765573in}{1.875759in}}{\pgfqpoint{1.769691in}{1.879877in}}%
\pgfpathcurveto{\pgfqpoint{1.773809in}{1.883995in}}{\pgfqpoint{1.776123in}{1.889581in}}{\pgfqpoint{1.776123in}{1.895405in}}%
\pgfpathcurveto{\pgfqpoint{1.776123in}{1.901229in}}{\pgfqpoint{1.773809in}{1.906815in}}{\pgfqpoint{1.769691in}{1.910933in}}%
\pgfpathcurveto{\pgfqpoint{1.765573in}{1.915052in}}{\pgfqpoint{1.759986in}{1.917365in}}{\pgfqpoint{1.754162in}{1.917365in}}%
\pgfpathcurveto{\pgfqpoint{1.748339in}{1.917365in}}{\pgfqpoint{1.742752in}{1.915052in}}{\pgfqpoint{1.738634in}{1.910933in}}%
\pgfpathcurveto{\pgfqpoint{1.734516in}{1.906815in}}{\pgfqpoint{1.732202in}{1.901229in}}{\pgfqpoint{1.732202in}{1.895405in}}%
\pgfpathcurveto{\pgfqpoint{1.732202in}{1.889581in}}{\pgfqpoint{1.734516in}{1.883995in}}{\pgfqpoint{1.738634in}{1.879877in}}%
\pgfpathcurveto{\pgfqpoint{1.742752in}{1.875759in}}{\pgfqpoint{1.748339in}{1.873445in}}{\pgfqpoint{1.754162in}{1.873445in}}%
\pgfpathlineto{\pgfqpoint{1.754162in}{1.873445in}}%
\pgfpathclose%
\pgfusepath{stroke,fill}%
\end{pgfscope}%
\begin{pgfscope}%
\pgfpathrectangle{\pgfqpoint{0.100000in}{0.209042in}}{\pgfqpoint{2.906250in}{2.906250in}}%
\pgfusepath{clip}%
\pgfsetbuttcap%
\pgfsetroundjoin%
\definecolor{currentfill}{rgb}{0.944006,0.377643,0.365136}%
\pgfsetfillcolor{currentfill}%
\pgfsetfillopacity{0.800000}%
\pgfsetlinewidth{1.003750pt}%
\definecolor{currentstroke}{rgb}{0.944006,0.377643,0.365136}%
\pgfsetstrokecolor{currentstroke}%
\pgfsetstrokeopacity{0.800000}%
\pgfsetdash{}{0pt}%
\pgfpathmoveto{\pgfqpoint{1.030600in}{1.693937in}}%
\pgfpathcurveto{\pgfqpoint{1.036424in}{1.693937in}}{\pgfqpoint{1.042010in}{1.696250in}}{\pgfqpoint{1.046128in}{1.700369in}}%
\pgfpathcurveto{\pgfqpoint{1.050246in}{1.704487in}}{\pgfqpoint{1.052560in}{1.710073in}}{\pgfqpoint{1.052560in}{1.715897in}}%
\pgfpathcurveto{\pgfqpoint{1.052560in}{1.721721in}}{\pgfqpoint{1.050246in}{1.727307in}}{\pgfqpoint{1.046128in}{1.731425in}}%
\pgfpathcurveto{\pgfqpoint{1.042010in}{1.735543in}}{\pgfqpoint{1.036424in}{1.737857in}}{\pgfqpoint{1.030600in}{1.737857in}}%
\pgfpathcurveto{\pgfqpoint{1.024776in}{1.737857in}}{\pgfqpoint{1.019190in}{1.735543in}}{\pgfqpoint{1.015071in}{1.731425in}}%
\pgfpathcurveto{\pgfqpoint{1.010953in}{1.727307in}}{\pgfqpoint{1.008639in}{1.721721in}}{\pgfqpoint{1.008639in}{1.715897in}}%
\pgfpathcurveto{\pgfqpoint{1.008639in}{1.710073in}}{\pgfqpoint{1.010953in}{1.704487in}}{\pgfqpoint{1.015071in}{1.700369in}}%
\pgfpathcurveto{\pgfqpoint{1.019190in}{1.696250in}}{\pgfqpoint{1.024776in}{1.693937in}}{\pgfqpoint{1.030600in}{1.693937in}}%
\pgfpathlineto{\pgfqpoint{1.030600in}{1.693937in}}%
\pgfpathclose%
\pgfusepath{stroke,fill}%
\end{pgfscope}%
\begin{pgfscope}%
\pgfpathrectangle{\pgfqpoint{0.100000in}{0.209042in}}{\pgfqpoint{2.906250in}{2.906250in}}%
\pgfusepath{clip}%
\pgfsetbuttcap%
\pgfsetroundjoin%
\definecolor{currentfill}{rgb}{0.690661,0.206384,0.482558}%
\pgfsetfillcolor{currentfill}%
\pgfsetfillopacity{0.800000}%
\pgfsetlinewidth{1.003750pt}%
\definecolor{currentstroke}{rgb}{0.690661,0.206384,0.482558}%
\pgfsetstrokecolor{currentstroke}%
\pgfsetstrokeopacity{0.800000}%
\pgfsetdash{}{0pt}%
\pgfpathmoveto{\pgfqpoint{1.615374in}{1.530482in}}%
\pgfpathcurveto{\pgfqpoint{1.621198in}{1.530482in}}{\pgfqpoint{1.626784in}{1.532796in}}{\pgfqpoint{1.630902in}{1.536914in}}%
\pgfpathcurveto{\pgfqpoint{1.635021in}{1.541032in}}{\pgfqpoint{1.637334in}{1.546618in}}{\pgfqpoint{1.637334in}{1.552442in}}%
\pgfpathcurveto{\pgfqpoint{1.637334in}{1.558266in}}{\pgfqpoint{1.635021in}{1.563852in}}{\pgfqpoint{1.630902in}{1.567971in}}%
\pgfpathcurveto{\pgfqpoint{1.626784in}{1.572089in}}{\pgfqpoint{1.621198in}{1.574403in}}{\pgfqpoint{1.615374in}{1.574403in}}%
\pgfpathcurveto{\pgfqpoint{1.609550in}{1.574403in}}{\pgfqpoint{1.603964in}{1.572089in}}{\pgfqpoint{1.599846in}{1.567971in}}%
\pgfpathcurveto{\pgfqpoint{1.595728in}{1.563852in}}{\pgfqpoint{1.593414in}{1.558266in}}{\pgfqpoint{1.593414in}{1.552442in}}%
\pgfpathcurveto{\pgfqpoint{1.593414in}{1.546618in}}{\pgfqpoint{1.595728in}{1.541032in}}{\pgfqpoint{1.599846in}{1.536914in}}%
\pgfpathcurveto{\pgfqpoint{1.603964in}{1.532796in}}{\pgfqpoint{1.609550in}{1.530482in}}{\pgfqpoint{1.615374in}{1.530482in}}%
\pgfpathlineto{\pgfqpoint{1.615374in}{1.530482in}}%
\pgfpathclose%
\pgfusepath{stroke,fill}%
\end{pgfscope}%
\begin{pgfscope}%
\pgfpathrectangle{\pgfqpoint{0.100000in}{0.209042in}}{\pgfqpoint{2.906250in}{2.906250in}}%
\pgfusepath{clip}%
\pgfsetbuttcap%
\pgfsetroundjoin%
\definecolor{currentfill}{rgb}{0.798608,0.247040,0.444848}%
\pgfsetfillcolor{currentfill}%
\pgfsetfillopacity{0.800000}%
\pgfsetlinewidth{1.003750pt}%
\definecolor{currentstroke}{rgb}{0.798608,0.247040,0.444848}%
\pgfsetstrokecolor{currentstroke}%
\pgfsetstrokeopacity{0.800000}%
\pgfsetdash{}{0pt}%
\pgfpathmoveto{\pgfqpoint{1.543524in}{1.431575in}}%
\pgfpathcurveto{\pgfqpoint{1.549348in}{1.431575in}}{\pgfqpoint{1.554934in}{1.433889in}}{\pgfqpoint{1.559052in}{1.438007in}}%
\pgfpathcurveto{\pgfqpoint{1.563170in}{1.442125in}}{\pgfqpoint{1.565484in}{1.447711in}}{\pgfqpoint{1.565484in}{1.453535in}}%
\pgfpathcurveto{\pgfqpoint{1.565484in}{1.459359in}}{\pgfqpoint{1.563170in}{1.464945in}}{\pgfqpoint{1.559052in}{1.469063in}}%
\pgfpathcurveto{\pgfqpoint{1.554934in}{1.473182in}}{\pgfqpoint{1.549348in}{1.475496in}}{\pgfqpoint{1.543524in}{1.475496in}}%
\pgfpathcurveto{\pgfqpoint{1.537700in}{1.475496in}}{\pgfqpoint{1.532113in}{1.473182in}}{\pgfqpoint{1.527995in}{1.469063in}}%
\pgfpathcurveto{\pgfqpoint{1.523877in}{1.464945in}}{\pgfqpoint{1.521563in}{1.459359in}}{\pgfqpoint{1.521563in}{1.453535in}}%
\pgfpathcurveto{\pgfqpoint{1.521563in}{1.447711in}}{\pgfqpoint{1.523877in}{1.442125in}}{\pgfqpoint{1.527995in}{1.438007in}}%
\pgfpathcurveto{\pgfqpoint{1.532113in}{1.433889in}}{\pgfqpoint{1.537700in}{1.431575in}}{\pgfqpoint{1.543524in}{1.431575in}}%
\pgfpathlineto{\pgfqpoint{1.543524in}{1.431575in}}%
\pgfpathclose%
\pgfusepath{stroke,fill}%
\end{pgfscope}%
\begin{pgfscope}%
\pgfpathrectangle{\pgfqpoint{0.100000in}{0.209042in}}{\pgfqpoint{2.906250in}{2.906250in}}%
\pgfusepath{clip}%
\pgfsetbuttcap%
\pgfsetroundjoin%
\definecolor{currentfill}{rgb}{0.716387,0.214982,0.475290}%
\pgfsetfillcolor{currentfill}%
\pgfsetfillopacity{0.800000}%
\pgfsetlinewidth{1.003750pt}%
\definecolor{currentstroke}{rgb}{0.716387,0.214982,0.475290}%
\pgfsetstrokecolor{currentstroke}%
\pgfsetstrokeopacity{0.800000}%
\pgfsetdash{}{0pt}%
\pgfpathmoveto{\pgfqpoint{1.834177in}{1.567167in}}%
\pgfpathcurveto{\pgfqpoint{1.840000in}{1.567167in}}{\pgfqpoint{1.845587in}{1.569481in}}{\pgfqpoint{1.849705in}{1.573599in}}%
\pgfpathcurveto{\pgfqpoint{1.853823in}{1.577717in}}{\pgfqpoint{1.856137in}{1.583303in}}{\pgfqpoint{1.856137in}{1.589127in}}%
\pgfpathcurveto{\pgfqpoint{1.856137in}{1.594951in}}{\pgfqpoint{1.853823in}{1.600537in}}{\pgfqpoint{1.849705in}{1.604655in}}%
\pgfpathcurveto{\pgfqpoint{1.845587in}{1.608774in}}{\pgfqpoint{1.840000in}{1.611087in}}{\pgfqpoint{1.834177in}{1.611087in}}%
\pgfpathcurveto{\pgfqpoint{1.828353in}{1.611087in}}{\pgfqpoint{1.822766in}{1.608774in}}{\pgfqpoint{1.818648in}{1.604655in}}%
\pgfpathcurveto{\pgfqpoint{1.814530in}{1.600537in}}{\pgfqpoint{1.812216in}{1.594951in}}{\pgfqpoint{1.812216in}{1.589127in}}%
\pgfpathcurveto{\pgfqpoint{1.812216in}{1.583303in}}{\pgfqpoint{1.814530in}{1.577717in}}{\pgfqpoint{1.818648in}{1.573599in}}%
\pgfpathcurveto{\pgfqpoint{1.822766in}{1.569481in}}{\pgfqpoint{1.828353in}{1.567167in}}{\pgfqpoint{1.834177in}{1.567167in}}%
\pgfpathlineto{\pgfqpoint{1.834177in}{1.567167in}}%
\pgfpathclose%
\pgfusepath{stroke,fill}%
\end{pgfscope}%
\begin{pgfscope}%
\pgfpathrectangle{\pgfqpoint{0.100000in}{0.209042in}}{\pgfqpoint{2.906250in}{2.906250in}}%
\pgfusepath{clip}%
\pgfsetbuttcap%
\pgfsetroundjoin%
\definecolor{currentfill}{rgb}{0.671349,0.200133,0.487358}%
\pgfsetfillcolor{currentfill}%
\pgfsetfillopacity{0.800000}%
\pgfsetlinewidth{1.003750pt}%
\definecolor{currentstroke}{rgb}{0.671349,0.200133,0.487358}%
\pgfsetstrokecolor{currentstroke}%
\pgfsetstrokeopacity{0.800000}%
\pgfsetdash{}{0pt}%
\pgfpathmoveto{\pgfqpoint{1.810304in}{1.602641in}}%
\pgfpathcurveto{\pgfqpoint{1.816128in}{1.602641in}}{\pgfqpoint{1.821714in}{1.604955in}}{\pgfqpoint{1.825833in}{1.609073in}}%
\pgfpathcurveto{\pgfqpoint{1.829951in}{1.613191in}}{\pgfqpoint{1.832265in}{1.618777in}}{\pgfqpoint{1.832265in}{1.624601in}}%
\pgfpathcurveto{\pgfqpoint{1.832265in}{1.630425in}}{\pgfqpoint{1.829951in}{1.636011in}}{\pgfqpoint{1.825833in}{1.640129in}}%
\pgfpathcurveto{\pgfqpoint{1.821714in}{1.644247in}}{\pgfqpoint{1.816128in}{1.646561in}}{\pgfqpoint{1.810304in}{1.646561in}}%
\pgfpathcurveto{\pgfqpoint{1.804480in}{1.646561in}}{\pgfqpoint{1.798894in}{1.644247in}}{\pgfqpoint{1.794776in}{1.640129in}}%
\pgfpathcurveto{\pgfqpoint{1.790658in}{1.636011in}}{\pgfqpoint{1.788344in}{1.630425in}}{\pgfqpoint{1.788344in}{1.624601in}}%
\pgfpathcurveto{\pgfqpoint{1.788344in}{1.618777in}}{\pgfqpoint{1.790658in}{1.613191in}}{\pgfqpoint{1.794776in}{1.609073in}}%
\pgfpathcurveto{\pgfqpoint{1.798894in}{1.604955in}}{\pgfqpoint{1.804480in}{1.602641in}}{\pgfqpoint{1.810304in}{1.602641in}}%
\pgfpathlineto{\pgfqpoint{1.810304in}{1.602641in}}%
\pgfpathclose%
\pgfusepath{stroke,fill}%
\end{pgfscope}%
\begin{pgfscope}%
\pgfpathrectangle{\pgfqpoint{0.100000in}{0.209042in}}{\pgfqpoint{2.906250in}{2.906250in}}%
\pgfusepath{clip}%
\pgfsetbuttcap%
\pgfsetroundjoin%
\definecolor{currentfill}{rgb}{0.600868,0.177743,0.500394}%
\pgfsetfillcolor{currentfill}%
\pgfsetfillopacity{0.800000}%
\pgfsetlinewidth{1.003750pt}%
\definecolor{currentstroke}{rgb}{0.600868,0.177743,0.500394}%
\pgfsetstrokecolor{currentstroke}%
\pgfsetstrokeopacity{0.800000}%
\pgfsetdash{}{0pt}%
\pgfpathmoveto{\pgfqpoint{1.738502in}{1.858015in}}%
\pgfpathcurveto{\pgfqpoint{1.744326in}{1.858015in}}{\pgfqpoint{1.749913in}{1.860329in}}{\pgfqpoint{1.754031in}{1.864447in}}%
\pgfpathcurveto{\pgfqpoint{1.758149in}{1.868565in}}{\pgfqpoint{1.760463in}{1.874151in}}{\pgfqpoint{1.760463in}{1.879975in}}%
\pgfpathcurveto{\pgfqpoint{1.760463in}{1.885799in}}{\pgfqpoint{1.758149in}{1.891385in}}{\pgfqpoint{1.754031in}{1.895504in}}%
\pgfpathcurveto{\pgfqpoint{1.749913in}{1.899622in}}{\pgfqpoint{1.744326in}{1.901936in}}{\pgfqpoint{1.738502in}{1.901936in}}%
\pgfpathcurveto{\pgfqpoint{1.732678in}{1.901936in}}{\pgfqpoint{1.727092in}{1.899622in}}{\pgfqpoint{1.722974in}{1.895504in}}%
\pgfpathcurveto{\pgfqpoint{1.718856in}{1.891385in}}{\pgfqpoint{1.716542in}{1.885799in}}{\pgfqpoint{1.716542in}{1.879975in}}%
\pgfpathcurveto{\pgfqpoint{1.716542in}{1.874151in}}{\pgfqpoint{1.718856in}{1.868565in}}{\pgfqpoint{1.722974in}{1.864447in}}%
\pgfpathcurveto{\pgfqpoint{1.727092in}{1.860329in}}{\pgfqpoint{1.732678in}{1.858015in}}{\pgfqpoint{1.738502in}{1.858015in}}%
\pgfpathlineto{\pgfqpoint{1.738502in}{1.858015in}}%
\pgfpathclose%
\pgfusepath{stroke,fill}%
\end{pgfscope}%
\begin{pgfscope}%
\pgfpathrectangle{\pgfqpoint{0.100000in}{0.209042in}}{\pgfqpoint{2.906250in}{2.906250in}}%
\pgfusepath{clip}%
\pgfsetbuttcap%
\pgfsetroundjoin%
\definecolor{currentfill}{rgb}{0.709962,0.212797,0.477201}%
\pgfsetfillcolor{currentfill}%
\pgfsetfillopacity{0.800000}%
\pgfsetlinewidth{1.003750pt}%
\definecolor{currentstroke}{rgb}{0.709962,0.212797,0.477201}%
\pgfsetstrokecolor{currentstroke}%
\pgfsetstrokeopacity{0.800000}%
\pgfsetdash{}{0pt}%
\pgfpathmoveto{\pgfqpoint{1.247262in}{1.877594in}}%
\pgfpathcurveto{\pgfqpoint{1.253086in}{1.877594in}}{\pgfqpoint{1.258672in}{1.879908in}}{\pgfqpoint{1.262790in}{1.884026in}}%
\pgfpathcurveto{\pgfqpoint{1.266908in}{1.888144in}}{\pgfqpoint{1.269222in}{1.893730in}}{\pgfqpoint{1.269222in}{1.899554in}}%
\pgfpathcurveto{\pgfqpoint{1.269222in}{1.905378in}}{\pgfqpoint{1.266908in}{1.910964in}}{\pgfqpoint{1.262790in}{1.915082in}}%
\pgfpathcurveto{\pgfqpoint{1.258672in}{1.919201in}}{\pgfqpoint{1.253086in}{1.921514in}}{\pgfqpoint{1.247262in}{1.921514in}}%
\pgfpathcurveto{\pgfqpoint{1.241438in}{1.921514in}}{\pgfqpoint{1.235852in}{1.919201in}}{\pgfqpoint{1.231734in}{1.915082in}}%
\pgfpathcurveto{\pgfqpoint{1.227616in}{1.910964in}}{\pgfqpoint{1.225302in}{1.905378in}}{\pgfqpoint{1.225302in}{1.899554in}}%
\pgfpathcurveto{\pgfqpoint{1.225302in}{1.893730in}}{\pgfqpoint{1.227616in}{1.888144in}}{\pgfqpoint{1.231734in}{1.884026in}}%
\pgfpathcurveto{\pgfqpoint{1.235852in}{1.879908in}}{\pgfqpoint{1.241438in}{1.877594in}}{\pgfqpoint{1.247262in}{1.877594in}}%
\pgfpathlineto{\pgfqpoint{1.247262in}{1.877594in}}%
\pgfpathclose%
\pgfusepath{stroke,fill}%
\end{pgfscope}%
\begin{pgfscope}%
\pgfpathrectangle{\pgfqpoint{0.100000in}{0.209042in}}{\pgfqpoint{2.906250in}{2.906250in}}%
\pgfusepath{clip}%
\pgfsetbuttcap%
\pgfsetroundjoin%
\definecolor{currentfill}{rgb}{0.594508,0.175701,0.501241}%
\pgfsetfillcolor{currentfill}%
\pgfsetfillopacity{0.800000}%
\pgfsetlinewidth{1.003750pt}%
\definecolor{currentstroke}{rgb}{0.594508,0.175701,0.501241}%
\pgfsetstrokecolor{currentstroke}%
\pgfsetstrokeopacity{0.800000}%
\pgfsetdash{}{0pt}%
\pgfpathmoveto{\pgfqpoint{1.723967in}{1.887133in}}%
\pgfpathcurveto{\pgfqpoint{1.729791in}{1.887133in}}{\pgfqpoint{1.735377in}{1.889447in}}{\pgfqpoint{1.739495in}{1.893565in}}%
\pgfpathcurveto{\pgfqpoint{1.743614in}{1.897683in}}{\pgfqpoint{1.745927in}{1.903270in}}{\pgfqpoint{1.745927in}{1.909093in}}%
\pgfpathcurveto{\pgfqpoint{1.745927in}{1.914917in}}{\pgfqpoint{1.743614in}{1.920504in}}{\pgfqpoint{1.739495in}{1.924622in}}%
\pgfpathcurveto{\pgfqpoint{1.735377in}{1.928740in}}{\pgfqpoint{1.729791in}{1.931054in}}{\pgfqpoint{1.723967in}{1.931054in}}%
\pgfpathcurveto{\pgfqpoint{1.718143in}{1.931054in}}{\pgfqpoint{1.712557in}{1.928740in}}{\pgfqpoint{1.708439in}{1.924622in}}%
\pgfpathcurveto{\pgfqpoint{1.704321in}{1.920504in}}{\pgfqpoint{1.702007in}{1.914917in}}{\pgfqpoint{1.702007in}{1.909093in}}%
\pgfpathcurveto{\pgfqpoint{1.702007in}{1.903270in}}{\pgfqpoint{1.704321in}{1.897683in}}{\pgfqpoint{1.708439in}{1.893565in}}%
\pgfpathcurveto{\pgfqpoint{1.712557in}{1.889447in}}{\pgfqpoint{1.718143in}{1.887133in}}{\pgfqpoint{1.723967in}{1.887133in}}%
\pgfpathlineto{\pgfqpoint{1.723967in}{1.887133in}}%
\pgfpathclose%
\pgfusepath{stroke,fill}%
\end{pgfscope}%
\begin{pgfscope}%
\pgfpathrectangle{\pgfqpoint{0.100000in}{0.209042in}}{\pgfqpoint{2.906250in}{2.906250in}}%
\pgfusepath{clip}%
\pgfsetbuttcap%
\pgfsetroundjoin%
\definecolor{currentfill}{rgb}{0.840636,0.268953,0.424666}%
\pgfsetfillcolor{currentfill}%
\pgfsetfillopacity{0.800000}%
\pgfsetlinewidth{1.003750pt}%
\definecolor{currentstroke}{rgb}{0.840636,0.268953,0.424666}%
\pgfsetstrokecolor{currentstroke}%
\pgfsetstrokeopacity{0.800000}%
\pgfsetdash{}{0pt}%
\pgfpathmoveto{\pgfqpoint{1.429335in}{1.398705in}}%
\pgfpathcurveto{\pgfqpoint{1.435159in}{1.398705in}}{\pgfqpoint{1.440745in}{1.401019in}}{\pgfqpoint{1.444863in}{1.405137in}}%
\pgfpathcurveto{\pgfqpoint{1.448981in}{1.409256in}}{\pgfqpoint{1.451295in}{1.414842in}}{\pgfqpoint{1.451295in}{1.420666in}}%
\pgfpathcurveto{\pgfqpoint{1.451295in}{1.426490in}}{\pgfqpoint{1.448981in}{1.432076in}}{\pgfqpoint{1.444863in}{1.436194in}}%
\pgfpathcurveto{\pgfqpoint{1.440745in}{1.440312in}}{\pgfqpoint{1.435159in}{1.442626in}}{\pgfqpoint{1.429335in}{1.442626in}}%
\pgfpathcurveto{\pgfqpoint{1.423511in}{1.442626in}}{\pgfqpoint{1.417925in}{1.440312in}}{\pgfqpoint{1.413807in}{1.436194in}}%
\pgfpathcurveto{\pgfqpoint{1.409688in}{1.432076in}}{\pgfqpoint{1.407375in}{1.426490in}}{\pgfqpoint{1.407375in}{1.420666in}}%
\pgfpathcurveto{\pgfqpoint{1.407375in}{1.414842in}}{\pgfqpoint{1.409688in}{1.409256in}}{\pgfqpoint{1.413807in}{1.405137in}}%
\pgfpathcurveto{\pgfqpoint{1.417925in}{1.401019in}}{\pgfqpoint{1.423511in}{1.398705in}}{\pgfqpoint{1.429335in}{1.398705in}}%
\pgfpathlineto{\pgfqpoint{1.429335in}{1.398705in}}%
\pgfpathclose%
\pgfusepath{stroke,fill}%
\end{pgfscope}%
\begin{pgfscope}%
\pgfpathrectangle{\pgfqpoint{0.100000in}{0.209042in}}{\pgfqpoint{2.906250in}{2.906250in}}%
\pgfusepath{clip}%
\pgfsetbuttcap%
\pgfsetroundjoin%
\definecolor{currentfill}{rgb}{0.709962,0.212797,0.477201}%
\pgfsetfillcolor{currentfill}%
\pgfsetfillopacity{0.800000}%
\pgfsetlinewidth{1.003750pt}%
\definecolor{currentstroke}{rgb}{0.709962,0.212797,0.477201}%
\pgfsetstrokecolor{currentstroke}%
\pgfsetstrokeopacity{0.800000}%
\pgfsetdash{}{0pt}%
\pgfpathmoveto{\pgfqpoint{1.750231in}{2.006422in}}%
\pgfpathcurveto{\pgfqpoint{1.756055in}{2.006422in}}{\pgfqpoint{1.761641in}{2.008736in}}{\pgfqpoint{1.765759in}{2.012854in}}%
\pgfpathcurveto{\pgfqpoint{1.769878in}{2.016972in}}{\pgfqpoint{1.772191in}{2.022558in}}{\pgfqpoint{1.772191in}{2.028382in}}%
\pgfpathcurveto{\pgfqpoint{1.772191in}{2.034206in}}{\pgfqpoint{1.769878in}{2.039792in}}{\pgfqpoint{1.765759in}{2.043910in}}%
\pgfpathcurveto{\pgfqpoint{1.761641in}{2.048028in}}{\pgfqpoint{1.756055in}{2.050342in}}{\pgfqpoint{1.750231in}{2.050342in}}%
\pgfpathcurveto{\pgfqpoint{1.744407in}{2.050342in}}{\pgfqpoint{1.738821in}{2.048028in}}{\pgfqpoint{1.734703in}{2.043910in}}%
\pgfpathcurveto{\pgfqpoint{1.730585in}{2.039792in}}{\pgfqpoint{1.728271in}{2.034206in}}{\pgfqpoint{1.728271in}{2.028382in}}%
\pgfpathcurveto{\pgfqpoint{1.728271in}{2.022558in}}{\pgfqpoint{1.730585in}{2.016972in}}{\pgfqpoint{1.734703in}{2.012854in}}%
\pgfpathcurveto{\pgfqpoint{1.738821in}{2.008736in}}{\pgfqpoint{1.744407in}{2.006422in}}{\pgfqpoint{1.750231in}{2.006422in}}%
\pgfpathlineto{\pgfqpoint{1.750231in}{2.006422in}}%
\pgfpathclose%
\pgfusepath{stroke,fill}%
\end{pgfscope}%
\begin{pgfscope}%
\pgfpathrectangle{\pgfqpoint{0.100000in}{0.209042in}}{\pgfqpoint{2.906250in}{2.906250in}}%
\pgfusepath{clip}%
\pgfsetbuttcap%
\pgfsetroundjoin%
\definecolor{currentfill}{rgb}{0.632805,0.187893,0.495332}%
\pgfsetfillcolor{currentfill}%
\pgfsetfillopacity{0.800000}%
\pgfsetlinewidth{1.003750pt}%
\definecolor{currentstroke}{rgb}{0.632805,0.187893,0.495332}%
\pgfsetstrokecolor{currentstroke}%
\pgfsetstrokeopacity{0.800000}%
\pgfsetdash{}{0pt}%
\pgfpathmoveto{\pgfqpoint{1.420099in}{1.967223in}}%
\pgfpathcurveto{\pgfqpoint{1.425923in}{1.967223in}}{\pgfqpoint{1.431509in}{1.969536in}}{\pgfqpoint{1.435627in}{1.973655in}}%
\pgfpathcurveto{\pgfqpoint{1.439746in}{1.977773in}}{\pgfqpoint{1.442059in}{1.983359in}}{\pgfqpoint{1.442059in}{1.989183in}}%
\pgfpathcurveto{\pgfqpoint{1.442059in}{1.995007in}}{\pgfqpoint{1.439746in}{2.000593in}}{\pgfqpoint{1.435627in}{2.004711in}}%
\pgfpathcurveto{\pgfqpoint{1.431509in}{2.008829in}}{\pgfqpoint{1.425923in}{2.011143in}}{\pgfqpoint{1.420099in}{2.011143in}}%
\pgfpathcurveto{\pgfqpoint{1.414275in}{2.011143in}}{\pgfqpoint{1.408689in}{2.008829in}}{\pgfqpoint{1.404571in}{2.004711in}}%
\pgfpathcurveto{\pgfqpoint{1.400453in}{2.000593in}}{\pgfqpoint{1.398139in}{1.995007in}}{\pgfqpoint{1.398139in}{1.989183in}}%
\pgfpathcurveto{\pgfqpoint{1.398139in}{1.983359in}}{\pgfqpoint{1.400453in}{1.977773in}}{\pgfqpoint{1.404571in}{1.973655in}}%
\pgfpathcurveto{\pgfqpoint{1.408689in}{1.969536in}}{\pgfqpoint{1.414275in}{1.967223in}}{\pgfqpoint{1.420099in}{1.967223in}}%
\pgfpathlineto{\pgfqpoint{1.420099in}{1.967223in}}%
\pgfpathclose%
\pgfusepath{stroke,fill}%
\end{pgfscope}%
\begin{pgfscope}%
\pgfpathrectangle{\pgfqpoint{0.100000in}{0.209042in}}{\pgfqpoint{2.906250in}{2.906250in}}%
\pgfusepath{clip}%
\pgfsetbuttcap%
\pgfsetroundjoin%
\definecolor{currentfill}{rgb}{0.494258,0.141462,0.507988}%
\pgfsetfillcolor{currentfill}%
\pgfsetfillopacity{0.800000}%
\pgfsetlinewidth{1.003750pt}%
\definecolor{currentstroke}{rgb}{0.494258,0.141462,0.507988}%
\pgfsetstrokecolor{currentstroke}%
\pgfsetstrokeopacity{0.800000}%
\pgfsetdash{}{0pt}%
\pgfpathmoveto{\pgfqpoint{1.545115in}{1.829897in}}%
\pgfpathcurveto{\pgfqpoint{1.550939in}{1.829897in}}{\pgfqpoint{1.556525in}{1.832210in}}{\pgfqpoint{1.560644in}{1.836329in}}%
\pgfpathcurveto{\pgfqpoint{1.564762in}{1.840447in}}{\pgfqpoint{1.567076in}{1.846033in}}{\pgfqpoint{1.567076in}{1.851857in}}%
\pgfpathcurveto{\pgfqpoint{1.567076in}{1.857681in}}{\pgfqpoint{1.564762in}{1.863267in}}{\pgfqpoint{1.560644in}{1.867385in}}%
\pgfpathcurveto{\pgfqpoint{1.556525in}{1.871503in}}{\pgfqpoint{1.550939in}{1.873817in}}{\pgfqpoint{1.545115in}{1.873817in}}%
\pgfpathcurveto{\pgfqpoint{1.539291in}{1.873817in}}{\pgfqpoint{1.533705in}{1.871503in}}{\pgfqpoint{1.529587in}{1.867385in}}%
\pgfpathcurveto{\pgfqpoint{1.525469in}{1.863267in}}{\pgfqpoint{1.523155in}{1.857681in}}{\pgfqpoint{1.523155in}{1.851857in}}%
\pgfpathcurveto{\pgfqpoint{1.523155in}{1.846033in}}{\pgfqpoint{1.525469in}{1.840447in}}{\pgfqpoint{1.529587in}{1.836329in}}%
\pgfpathcurveto{\pgfqpoint{1.533705in}{1.832210in}}{\pgfqpoint{1.539291in}{1.829897in}}{\pgfqpoint{1.545115in}{1.829897in}}%
\pgfpathlineto{\pgfqpoint{1.545115in}{1.829897in}}%
\pgfpathclose%
\pgfusepath{stroke,fill}%
\end{pgfscope}%
\begin{pgfscope}%
\pgfpathrectangle{\pgfqpoint{0.100000in}{0.209042in}}{\pgfqpoint{2.906250in}{2.906250in}}%
\pgfusepath{clip}%
\pgfsetbuttcap%
\pgfsetroundjoin%
\definecolor{currentfill}{rgb}{0.613617,0.181811,0.498536}%
\pgfsetfillcolor{currentfill}%
\pgfsetfillopacity{0.800000}%
\pgfsetlinewidth{1.003750pt}%
\definecolor{currentstroke}{rgb}{0.613617,0.181811,0.498536}%
\pgfsetstrokecolor{currentstroke}%
\pgfsetstrokeopacity{0.800000}%
\pgfsetdash{}{0pt}%
\pgfpathmoveto{\pgfqpoint{1.513232in}{1.537318in}}%
\pgfpathcurveto{\pgfqpoint{1.519056in}{1.537318in}}{\pgfqpoint{1.524642in}{1.539632in}}{\pgfqpoint{1.528760in}{1.543750in}}%
\pgfpathcurveto{\pgfqpoint{1.532878in}{1.547868in}}{\pgfqpoint{1.535192in}{1.553454in}}{\pgfqpoint{1.535192in}{1.559278in}}%
\pgfpathcurveto{\pgfqpoint{1.535192in}{1.565102in}}{\pgfqpoint{1.532878in}{1.570689in}}{\pgfqpoint{1.528760in}{1.574807in}}%
\pgfpathcurveto{\pgfqpoint{1.524642in}{1.578925in}}{\pgfqpoint{1.519056in}{1.581239in}}{\pgfqpoint{1.513232in}{1.581239in}}%
\pgfpathcurveto{\pgfqpoint{1.507408in}{1.581239in}}{\pgfqpoint{1.501822in}{1.578925in}}{\pgfqpoint{1.497704in}{1.574807in}}%
\pgfpathcurveto{\pgfqpoint{1.493586in}{1.570689in}}{\pgfqpoint{1.491272in}{1.565102in}}{\pgfqpoint{1.491272in}{1.559278in}}%
\pgfpathcurveto{\pgfqpoint{1.491272in}{1.553454in}}{\pgfqpoint{1.493586in}{1.547868in}}{\pgfqpoint{1.497704in}{1.543750in}}%
\pgfpathcurveto{\pgfqpoint{1.501822in}{1.539632in}}{\pgfqpoint{1.507408in}{1.537318in}}{\pgfqpoint{1.513232in}{1.537318in}}%
\pgfpathlineto{\pgfqpoint{1.513232in}{1.537318in}}%
\pgfpathclose%
\pgfusepath{stroke,fill}%
\end{pgfscope}%
\begin{pgfscope}%
\pgfpathrectangle{\pgfqpoint{0.100000in}{0.209042in}}{\pgfqpoint{2.906250in}{2.906250in}}%
\pgfusepath{clip}%
\pgfsetbuttcap%
\pgfsetroundjoin%
\definecolor{currentfill}{rgb}{0.500438,0.143719,0.507920}%
\pgfsetfillcolor{currentfill}%
\pgfsetfillopacity{0.800000}%
\pgfsetlinewidth{1.003750pt}%
\definecolor{currentstroke}{rgb}{0.500438,0.143719,0.507920}%
\pgfsetstrokecolor{currentstroke}%
\pgfsetstrokeopacity{0.800000}%
\pgfsetdash{}{0pt}%
\pgfpathmoveto{\pgfqpoint{1.464049in}{1.777672in}}%
\pgfpathcurveto{\pgfqpoint{1.469873in}{1.777672in}}{\pgfqpoint{1.475459in}{1.779986in}}{\pgfqpoint{1.479577in}{1.784104in}}%
\pgfpathcurveto{\pgfqpoint{1.483695in}{1.788222in}}{\pgfqpoint{1.486009in}{1.793808in}}{\pgfqpoint{1.486009in}{1.799632in}}%
\pgfpathcurveto{\pgfqpoint{1.486009in}{1.805456in}}{\pgfqpoint{1.483695in}{1.811042in}}{\pgfqpoint{1.479577in}{1.815160in}}%
\pgfpathcurveto{\pgfqpoint{1.475459in}{1.819278in}}{\pgfqpoint{1.469873in}{1.821592in}}{\pgfqpoint{1.464049in}{1.821592in}}%
\pgfpathcurveto{\pgfqpoint{1.458225in}{1.821592in}}{\pgfqpoint{1.452639in}{1.819278in}}{\pgfqpoint{1.448520in}{1.815160in}}%
\pgfpathcurveto{\pgfqpoint{1.444402in}{1.811042in}}{\pgfqpoint{1.442088in}{1.805456in}}{\pgfqpoint{1.442088in}{1.799632in}}%
\pgfpathcurveto{\pgfqpoint{1.442088in}{1.793808in}}{\pgfqpoint{1.444402in}{1.788222in}}{\pgfqpoint{1.448520in}{1.784104in}}%
\pgfpathcurveto{\pgfqpoint{1.452639in}{1.779986in}}{\pgfqpoint{1.458225in}{1.777672in}}{\pgfqpoint{1.464049in}{1.777672in}}%
\pgfpathlineto{\pgfqpoint{1.464049in}{1.777672in}}%
\pgfpathclose%
\pgfusepath{stroke,fill}%
\end{pgfscope}%
\begin{pgfscope}%
\pgfpathrectangle{\pgfqpoint{0.100000in}{0.209042in}}{\pgfqpoint{2.906250in}{2.906250in}}%
\pgfusepath{clip}%
\pgfsetbuttcap%
\pgfsetroundjoin%
\definecolor{currentfill}{rgb}{0.588158,0.173652,0.502035}%
\pgfsetfillcolor{currentfill}%
\pgfsetfillopacity{0.800000}%
\pgfsetlinewidth{1.003750pt}%
\definecolor{currentstroke}{rgb}{0.588158,0.173652,0.502035}%
\pgfsetstrokecolor{currentstroke}%
\pgfsetstrokeopacity{0.800000}%
\pgfsetdash{}{0pt}%
\pgfpathmoveto{\pgfqpoint{1.325854in}{1.839656in}}%
\pgfpathcurveto{\pgfqpoint{1.331678in}{1.839656in}}{\pgfqpoint{1.337264in}{1.841970in}}{\pgfqpoint{1.341382in}{1.846088in}}%
\pgfpathcurveto{\pgfqpoint{1.345500in}{1.850206in}}{\pgfqpoint{1.347814in}{1.855792in}}{\pgfqpoint{1.347814in}{1.861616in}}%
\pgfpathcurveto{\pgfqpoint{1.347814in}{1.867440in}}{\pgfqpoint{1.345500in}{1.873026in}}{\pgfqpoint{1.341382in}{1.877144in}}%
\pgfpathcurveto{\pgfqpoint{1.337264in}{1.881262in}}{\pgfqpoint{1.331678in}{1.883576in}}{\pgfqpoint{1.325854in}{1.883576in}}%
\pgfpathcurveto{\pgfqpoint{1.320030in}{1.883576in}}{\pgfqpoint{1.314444in}{1.881262in}}{\pgfqpoint{1.310326in}{1.877144in}}%
\pgfpathcurveto{\pgfqpoint{1.306208in}{1.873026in}}{\pgfqpoint{1.303894in}{1.867440in}}{\pgfqpoint{1.303894in}{1.861616in}}%
\pgfpathcurveto{\pgfqpoint{1.303894in}{1.855792in}}{\pgfqpoint{1.306208in}{1.850206in}}{\pgfqpoint{1.310326in}{1.846088in}}%
\pgfpathcurveto{\pgfqpoint{1.314444in}{1.841970in}}{\pgfqpoint{1.320030in}{1.839656in}}{\pgfqpoint{1.325854in}{1.839656in}}%
\pgfpathlineto{\pgfqpoint{1.325854in}{1.839656in}}%
\pgfpathclose%
\pgfusepath{stroke,fill}%
\end{pgfscope}%
\begin{pgfscope}%
\pgfpathrectangle{\pgfqpoint{0.100000in}{0.209042in}}{\pgfqpoint{2.906250in}{2.906250in}}%
\pgfusepath{clip}%
\pgfsetbuttcap%
\pgfsetroundjoin%
\definecolor{currentfill}{rgb}{0.620005,0.183840,0.497524}%
\pgfsetfillcolor{currentfill}%
\pgfsetfillopacity{0.800000}%
\pgfsetlinewidth{1.003750pt}%
\definecolor{currentstroke}{rgb}{0.620005,0.183840,0.497524}%
\pgfsetstrokecolor{currentstroke}%
\pgfsetstrokeopacity{0.800000}%
\pgfsetdash{}{0pt}%
\pgfpathmoveto{\pgfqpoint{1.691721in}{1.985592in}}%
\pgfpathcurveto{\pgfqpoint{1.697545in}{1.985592in}}{\pgfqpoint{1.703132in}{1.987906in}}{\pgfqpoint{1.707250in}{1.992024in}}%
\pgfpathcurveto{\pgfqpoint{1.711368in}{1.996142in}}{\pgfqpoint{1.713682in}{2.001728in}}{\pgfqpoint{1.713682in}{2.007552in}}%
\pgfpathcurveto{\pgfqpoint{1.713682in}{2.013376in}}{\pgfqpoint{1.711368in}{2.018962in}}{\pgfqpoint{1.707250in}{2.023081in}}%
\pgfpathcurveto{\pgfqpoint{1.703132in}{2.027199in}}{\pgfqpoint{1.697545in}{2.029513in}}{\pgfqpoint{1.691721in}{2.029513in}}%
\pgfpathcurveto{\pgfqpoint{1.685897in}{2.029513in}}{\pgfqpoint{1.680311in}{2.027199in}}{\pgfqpoint{1.676193in}{2.023081in}}%
\pgfpathcurveto{\pgfqpoint{1.672075in}{2.018962in}}{\pgfqpoint{1.669761in}{2.013376in}}{\pgfqpoint{1.669761in}{2.007552in}}%
\pgfpathcurveto{\pgfqpoint{1.669761in}{2.001728in}}{\pgfqpoint{1.672075in}{1.996142in}}{\pgfqpoint{1.676193in}{1.992024in}}%
\pgfpathcurveto{\pgfqpoint{1.680311in}{1.987906in}}{\pgfqpoint{1.685897in}{1.985592in}}{\pgfqpoint{1.691721in}{1.985592in}}%
\pgfpathlineto{\pgfqpoint{1.691721in}{1.985592in}}%
\pgfpathclose%
\pgfusepath{stroke,fill}%
\end{pgfscope}%
\begin{pgfscope}%
\pgfpathrectangle{\pgfqpoint{0.100000in}{0.209042in}}{\pgfqpoint{2.906250in}{2.906250in}}%
\pgfusepath{clip}%
\pgfsetbuttcap%
\pgfsetroundjoin%
\definecolor{currentfill}{rgb}{0.600868,0.177743,0.500394}%
\pgfsetfillcolor{currentfill}%
\pgfsetfillopacity{0.800000}%
\pgfsetlinewidth{1.003750pt}%
\definecolor{currentstroke}{rgb}{0.600868,0.177743,0.500394}%
\pgfsetstrokecolor{currentstroke}%
\pgfsetstrokeopacity{0.800000}%
\pgfsetdash{}{0pt}%
\pgfpathmoveto{\pgfqpoint{1.296093in}{1.826216in}}%
\pgfpathcurveto{\pgfqpoint{1.301917in}{1.826216in}}{\pgfqpoint{1.307503in}{1.828530in}}{\pgfqpoint{1.311621in}{1.832648in}}%
\pgfpathcurveto{\pgfqpoint{1.315740in}{1.836766in}}{\pgfqpoint{1.318053in}{1.842352in}}{\pgfqpoint{1.318053in}{1.848176in}}%
\pgfpathcurveto{\pgfqpoint{1.318053in}{1.854000in}}{\pgfqpoint{1.315740in}{1.859586in}}{\pgfqpoint{1.311621in}{1.863705in}}%
\pgfpathcurveto{\pgfqpoint{1.307503in}{1.867823in}}{\pgfqpoint{1.301917in}{1.870137in}}{\pgfqpoint{1.296093in}{1.870137in}}%
\pgfpathcurveto{\pgfqpoint{1.290269in}{1.870137in}}{\pgfqpoint{1.284683in}{1.867823in}}{\pgfqpoint{1.280565in}{1.863705in}}%
\pgfpathcurveto{\pgfqpoint{1.276447in}{1.859586in}}{\pgfqpoint{1.274133in}{1.854000in}}{\pgfqpoint{1.274133in}{1.848176in}}%
\pgfpathcurveto{\pgfqpoint{1.274133in}{1.842352in}}{\pgfqpoint{1.276447in}{1.836766in}}{\pgfqpoint{1.280565in}{1.832648in}}%
\pgfpathcurveto{\pgfqpoint{1.284683in}{1.828530in}}{\pgfqpoint{1.290269in}{1.826216in}}{\pgfqpoint{1.296093in}{1.826216in}}%
\pgfpathlineto{\pgfqpoint{1.296093in}{1.826216in}}%
\pgfpathclose%
\pgfusepath{stroke,fill}%
\end{pgfscope}%
\begin{pgfscope}%
\pgfpathrectangle{\pgfqpoint{0.100000in}{0.209042in}}{\pgfqpoint{2.906250in}{2.906250in}}%
\pgfusepath{clip}%
\pgfsetbuttcap%
\pgfsetroundjoin%
\definecolor{currentfill}{rgb}{0.652056,0.193986,0.491611}%
\pgfsetfillcolor{currentfill}%
\pgfsetfillopacity{0.800000}%
\pgfsetlinewidth{1.003750pt}%
\definecolor{currentstroke}{rgb}{0.652056,0.193986,0.491611}%
\pgfsetstrokecolor{currentstroke}%
\pgfsetstrokeopacity{0.800000}%
\pgfsetdash{}{0pt}%
\pgfpathmoveto{\pgfqpoint{1.475309in}{1.493874in}}%
\pgfpathcurveto{\pgfqpoint{1.481133in}{1.493874in}}{\pgfqpoint{1.486719in}{1.496188in}}{\pgfqpoint{1.490837in}{1.500306in}}%
\pgfpathcurveto{\pgfqpoint{1.494955in}{1.504424in}}{\pgfqpoint{1.497269in}{1.510011in}}{\pgfqpoint{1.497269in}{1.515834in}}%
\pgfpathcurveto{\pgfqpoint{1.497269in}{1.521658in}}{\pgfqpoint{1.494955in}{1.527245in}}{\pgfqpoint{1.490837in}{1.531363in}}%
\pgfpathcurveto{\pgfqpoint{1.486719in}{1.535481in}}{\pgfqpoint{1.481133in}{1.537795in}}{\pgfqpoint{1.475309in}{1.537795in}}%
\pgfpathcurveto{\pgfqpoint{1.469485in}{1.537795in}}{\pgfqpoint{1.463899in}{1.535481in}}{\pgfqpoint{1.459781in}{1.531363in}}%
\pgfpathcurveto{\pgfqpoint{1.455663in}{1.527245in}}{\pgfqpoint{1.453349in}{1.521658in}}{\pgfqpoint{1.453349in}{1.515834in}}%
\pgfpathcurveto{\pgfqpoint{1.453349in}{1.510011in}}{\pgfqpoint{1.455663in}{1.504424in}}{\pgfqpoint{1.459781in}{1.500306in}}%
\pgfpathcurveto{\pgfqpoint{1.463899in}{1.496188in}}{\pgfqpoint{1.469485in}{1.493874in}}{\pgfqpoint{1.475309in}{1.493874in}}%
\pgfpathlineto{\pgfqpoint{1.475309in}{1.493874in}}%
\pgfpathclose%
\pgfusepath{stroke,fill}%
\end{pgfscope}%
\begin{pgfscope}%
\pgfpathrectangle{\pgfqpoint{0.100000in}{0.209042in}}{\pgfqpoint{2.906250in}{2.906250in}}%
\pgfusepath{clip}%
\pgfsetbuttcap%
\pgfsetroundjoin%
\definecolor{currentfill}{rgb}{0.868793,0.287728,0.409303}%
\pgfsetfillcolor{currentfill}%
\pgfsetfillopacity{0.800000}%
\pgfsetlinewidth{1.003750pt}%
\definecolor{currentstroke}{rgb}{0.868793,0.287728,0.409303}%
\pgfsetstrokecolor{currentstroke}%
\pgfsetstrokeopacity{0.800000}%
\pgfsetdash{}{0pt}%
\pgfpathmoveto{\pgfqpoint{1.053749in}{1.713851in}}%
\pgfpathcurveto{\pgfqpoint{1.059573in}{1.713851in}}{\pgfqpoint{1.065160in}{1.716165in}}{\pgfqpoint{1.069278in}{1.720283in}}%
\pgfpathcurveto{\pgfqpoint{1.073396in}{1.724401in}}{\pgfqpoint{1.075710in}{1.729987in}}{\pgfqpoint{1.075710in}{1.735811in}}%
\pgfpathcurveto{\pgfqpoint{1.075710in}{1.741635in}}{\pgfqpoint{1.073396in}{1.747221in}}{\pgfqpoint{1.069278in}{1.751339in}}%
\pgfpathcurveto{\pgfqpoint{1.065160in}{1.755457in}}{\pgfqpoint{1.059573in}{1.757771in}}{\pgfqpoint{1.053749in}{1.757771in}}%
\pgfpathcurveto{\pgfqpoint{1.047925in}{1.757771in}}{\pgfqpoint{1.042339in}{1.755457in}}{\pgfqpoint{1.038221in}{1.751339in}}%
\pgfpathcurveto{\pgfqpoint{1.034103in}{1.747221in}}{\pgfqpoint{1.031789in}{1.741635in}}{\pgfqpoint{1.031789in}{1.735811in}}%
\pgfpathcurveto{\pgfqpoint{1.031789in}{1.729987in}}{\pgfqpoint{1.034103in}{1.724401in}}{\pgfqpoint{1.038221in}{1.720283in}}%
\pgfpathcurveto{\pgfqpoint{1.042339in}{1.716165in}}{\pgfqpoint{1.047925in}{1.713851in}}{\pgfqpoint{1.053749in}{1.713851in}}%
\pgfpathlineto{\pgfqpoint{1.053749in}{1.713851in}}%
\pgfpathclose%
\pgfusepath{stroke,fill}%
\end{pgfscope}%
\begin{pgfscope}%
\pgfpathrectangle{\pgfqpoint{0.100000in}{0.209042in}}{\pgfqpoint{2.906250in}{2.906250in}}%
\pgfusepath{clip}%
\pgfsetbuttcap%
\pgfsetroundjoin%
\definecolor{currentfill}{rgb}{0.494258,0.141462,0.507988}%
\pgfsetfillcolor{currentfill}%
\pgfsetfillopacity{0.800000}%
\pgfsetlinewidth{1.003750pt}%
\definecolor{currentstroke}{rgb}{0.494258,0.141462,0.507988}%
\pgfsetstrokecolor{currentstroke}%
\pgfsetstrokeopacity{0.800000}%
\pgfsetdash{}{0pt}%
\pgfpathmoveto{\pgfqpoint{1.417877in}{1.740293in}}%
\pgfpathcurveto{\pgfqpoint{1.423701in}{1.740293in}}{\pgfqpoint{1.429287in}{1.742607in}}{\pgfqpoint{1.433405in}{1.746725in}}%
\pgfpathcurveto{\pgfqpoint{1.437523in}{1.750843in}}{\pgfqpoint{1.439837in}{1.756429in}}{\pgfqpoint{1.439837in}{1.762253in}}%
\pgfpathcurveto{\pgfqpoint{1.439837in}{1.768077in}}{\pgfqpoint{1.437523in}{1.773663in}}{\pgfqpoint{1.433405in}{1.777781in}}%
\pgfpathcurveto{\pgfqpoint{1.429287in}{1.781899in}}{\pgfqpoint{1.423701in}{1.784213in}}{\pgfqpoint{1.417877in}{1.784213in}}%
\pgfpathcurveto{\pgfqpoint{1.412053in}{1.784213in}}{\pgfqpoint{1.406466in}{1.781899in}}{\pgfqpoint{1.402348in}{1.777781in}}%
\pgfpathcurveto{\pgfqpoint{1.398230in}{1.773663in}}{\pgfqpoint{1.395916in}{1.768077in}}{\pgfqpoint{1.395916in}{1.762253in}}%
\pgfpathcurveto{\pgfqpoint{1.395916in}{1.756429in}}{\pgfqpoint{1.398230in}{1.750843in}}{\pgfqpoint{1.402348in}{1.746725in}}%
\pgfpathcurveto{\pgfqpoint{1.406466in}{1.742607in}}{\pgfqpoint{1.412053in}{1.740293in}}{\pgfqpoint{1.417877in}{1.740293in}}%
\pgfpathlineto{\pgfqpoint{1.417877in}{1.740293in}}%
\pgfpathclose%
\pgfusepath{stroke,fill}%
\end{pgfscope}%
\begin{pgfscope}%
\pgfpathrectangle{\pgfqpoint{0.100000in}{0.209042in}}{\pgfqpoint{2.906250in}{2.906250in}}%
\pgfusepath{clip}%
\pgfsetbuttcap%
\pgfsetroundjoin%
\definecolor{currentfill}{rgb}{0.512831,0.148179,0.507648}%
\pgfsetfillcolor{currentfill}%
\pgfsetfillopacity{0.800000}%
\pgfsetlinewidth{1.003750pt}%
\definecolor{currentstroke}{rgb}{0.512831,0.148179,0.507648}%
\pgfsetstrokecolor{currentstroke}%
\pgfsetstrokeopacity{0.800000}%
\pgfsetdash{}{0pt}%
\pgfpathmoveto{\pgfqpoint{1.409036in}{1.677997in}}%
\pgfpathcurveto{\pgfqpoint{1.414859in}{1.677997in}}{\pgfqpoint{1.420446in}{1.680311in}}{\pgfqpoint{1.424564in}{1.684429in}}%
\pgfpathcurveto{\pgfqpoint{1.428682in}{1.688547in}}{\pgfqpoint{1.430996in}{1.694134in}}{\pgfqpoint{1.430996in}{1.699958in}}%
\pgfpathcurveto{\pgfqpoint{1.430996in}{1.705781in}}{\pgfqpoint{1.428682in}{1.711368in}}{\pgfqpoint{1.424564in}{1.715486in}}%
\pgfpathcurveto{\pgfqpoint{1.420446in}{1.719604in}}{\pgfqpoint{1.414859in}{1.721918in}}{\pgfqpoint{1.409036in}{1.721918in}}%
\pgfpathcurveto{\pgfqpoint{1.403212in}{1.721918in}}{\pgfqpoint{1.397625in}{1.719604in}}{\pgfqpoint{1.393507in}{1.715486in}}%
\pgfpathcurveto{\pgfqpoint{1.389389in}{1.711368in}}{\pgfqpoint{1.387075in}{1.705781in}}{\pgfqpoint{1.387075in}{1.699958in}}%
\pgfpathcurveto{\pgfqpoint{1.387075in}{1.694134in}}{\pgfqpoint{1.389389in}{1.688547in}}{\pgfqpoint{1.393507in}{1.684429in}}%
\pgfpathcurveto{\pgfqpoint{1.397625in}{1.680311in}}{\pgfqpoint{1.403212in}{1.677997in}}{\pgfqpoint{1.409036in}{1.677997in}}%
\pgfpathlineto{\pgfqpoint{1.409036in}{1.677997in}}%
\pgfpathclose%
\pgfusepath{stroke,fill}%
\end{pgfscope}%
\begin{pgfscope}%
\pgfpathrectangle{\pgfqpoint{0.100000in}{0.209042in}}{\pgfqpoint{2.906250in}{2.906250in}}%
\pgfusepath{clip}%
\pgfsetbuttcap%
\pgfsetroundjoin%
\definecolor{currentfill}{rgb}{0.512831,0.148179,0.507648}%
\pgfsetfillcolor{currentfill}%
\pgfsetfillopacity{0.800000}%
\pgfsetlinewidth{1.003750pt}%
\definecolor{currentstroke}{rgb}{0.512831,0.148179,0.507648}%
\pgfsetstrokecolor{currentstroke}%
\pgfsetstrokeopacity{0.800000}%
\pgfsetdash{}{0pt}%
\pgfpathmoveto{\pgfqpoint{1.778572in}{1.809720in}}%
\pgfpathcurveto{\pgfqpoint{1.784396in}{1.809720in}}{\pgfqpoint{1.789983in}{1.812034in}}{\pgfqpoint{1.794101in}{1.816152in}}%
\pgfpathcurveto{\pgfqpoint{1.798219in}{1.820271in}}{\pgfqpoint{1.800533in}{1.825857in}}{\pgfqpoint{1.800533in}{1.831681in}}%
\pgfpathcurveto{\pgfqpoint{1.800533in}{1.837505in}}{\pgfqpoint{1.798219in}{1.843091in}}{\pgfqpoint{1.794101in}{1.847209in}}%
\pgfpathcurveto{\pgfqpoint{1.789983in}{1.851327in}}{\pgfqpoint{1.784396in}{1.853641in}}{\pgfqpoint{1.778572in}{1.853641in}}%
\pgfpathcurveto{\pgfqpoint{1.772748in}{1.853641in}}{\pgfqpoint{1.767162in}{1.851327in}}{\pgfqpoint{1.763044in}{1.847209in}}%
\pgfpathcurveto{\pgfqpoint{1.758926in}{1.843091in}}{\pgfqpoint{1.756612in}{1.837505in}}{\pgfqpoint{1.756612in}{1.831681in}}%
\pgfpathcurveto{\pgfqpoint{1.756612in}{1.825857in}}{\pgfqpoint{1.758926in}{1.820271in}}{\pgfqpoint{1.763044in}{1.816152in}}%
\pgfpathcurveto{\pgfqpoint{1.767162in}{1.812034in}}{\pgfqpoint{1.772748in}{1.809720in}}{\pgfqpoint{1.778572in}{1.809720in}}%
\pgfpathlineto{\pgfqpoint{1.778572in}{1.809720in}}%
\pgfpathclose%
\pgfusepath{stroke,fill}%
\end{pgfscope}%
\begin{pgfscope}%
\pgfpathrectangle{\pgfqpoint{0.100000in}{0.209042in}}{\pgfqpoint{2.906250in}{2.906250in}}%
\pgfusepath{clip}%
\pgfsetbuttcap%
\pgfsetroundjoin%
\definecolor{currentfill}{rgb}{0.481929,0.136891,0.507989}%
\pgfsetfillcolor{currentfill}%
\pgfsetfillopacity{0.800000}%
\pgfsetlinewidth{1.003750pt}%
\definecolor{currentstroke}{rgb}{0.481929,0.136891,0.507989}%
\pgfsetstrokecolor{currentstroke}%
\pgfsetstrokeopacity{0.800000}%
\pgfsetdash{}{0pt}%
\pgfpathmoveto{\pgfqpoint{1.447913in}{1.668357in}}%
\pgfpathcurveto{\pgfqpoint{1.453737in}{1.668357in}}{\pgfqpoint{1.459323in}{1.670671in}}{\pgfqpoint{1.463441in}{1.674789in}}%
\pgfpathcurveto{\pgfqpoint{1.467559in}{1.678907in}}{\pgfqpoint{1.469873in}{1.684493in}}{\pgfqpoint{1.469873in}{1.690317in}}%
\pgfpathcurveto{\pgfqpoint{1.469873in}{1.696141in}}{\pgfqpoint{1.467559in}{1.701727in}}{\pgfqpoint{1.463441in}{1.705845in}}%
\pgfpathcurveto{\pgfqpoint{1.459323in}{1.709964in}}{\pgfqpoint{1.453737in}{1.712277in}}{\pgfqpoint{1.447913in}{1.712277in}}%
\pgfpathcurveto{\pgfqpoint{1.442089in}{1.712277in}}{\pgfqpoint{1.436503in}{1.709964in}}{\pgfqpoint{1.432385in}{1.705845in}}%
\pgfpathcurveto{\pgfqpoint{1.428267in}{1.701727in}}{\pgfqpoint{1.425953in}{1.696141in}}{\pgfqpoint{1.425953in}{1.690317in}}%
\pgfpathcurveto{\pgfqpoint{1.425953in}{1.684493in}}{\pgfqpoint{1.428267in}{1.678907in}}{\pgfqpoint{1.432385in}{1.674789in}}%
\pgfpathcurveto{\pgfqpoint{1.436503in}{1.670671in}}{\pgfqpoint{1.442089in}{1.668357in}}{\pgfqpoint{1.447913in}{1.668357in}}%
\pgfpathlineto{\pgfqpoint{1.447913in}{1.668357in}}%
\pgfpathclose%
\pgfusepath{stroke,fill}%
\end{pgfscope}%
\begin{pgfscope}%
\pgfpathrectangle{\pgfqpoint{0.100000in}{0.209042in}}{\pgfqpoint{2.906250in}{2.906250in}}%
\pgfusepath{clip}%
\pgfsetbuttcap%
\pgfsetroundjoin%
\definecolor{currentfill}{rgb}{0.457386,0.127522,0.507448}%
\pgfsetfillcolor{currentfill}%
\pgfsetfillopacity{0.800000}%
\pgfsetlinewidth{1.003750pt}%
\definecolor{currentstroke}{rgb}{0.457386,0.127522,0.507448}%
\pgfsetstrokecolor{currentstroke}%
\pgfsetstrokeopacity{0.800000}%
\pgfsetdash{}{0pt}%
\pgfpathmoveto{\pgfqpoint{1.556448in}{1.647903in}}%
\pgfpathcurveto{\pgfqpoint{1.562272in}{1.647903in}}{\pgfqpoint{1.567858in}{1.650217in}}{\pgfqpoint{1.571976in}{1.654335in}}%
\pgfpathcurveto{\pgfqpoint{1.576094in}{1.658454in}}{\pgfqpoint{1.578408in}{1.664040in}}{\pgfqpoint{1.578408in}{1.669864in}}%
\pgfpathcurveto{\pgfqpoint{1.578408in}{1.675688in}}{\pgfqpoint{1.576094in}{1.681274in}}{\pgfqpoint{1.571976in}{1.685392in}}%
\pgfpathcurveto{\pgfqpoint{1.567858in}{1.689510in}}{\pgfqpoint{1.562272in}{1.691824in}}{\pgfqpoint{1.556448in}{1.691824in}}%
\pgfpathcurveto{\pgfqpoint{1.550624in}{1.691824in}}{\pgfqpoint{1.545038in}{1.689510in}}{\pgfqpoint{1.540920in}{1.685392in}}%
\pgfpathcurveto{\pgfqpoint{1.536801in}{1.681274in}}{\pgfqpoint{1.534488in}{1.675688in}}{\pgfqpoint{1.534488in}{1.669864in}}%
\pgfpathcurveto{\pgfqpoint{1.534488in}{1.664040in}}{\pgfqpoint{1.536801in}{1.658454in}}{\pgfqpoint{1.540920in}{1.654335in}}%
\pgfpathcurveto{\pgfqpoint{1.545038in}{1.650217in}}{\pgfqpoint{1.550624in}{1.647903in}}{\pgfqpoint{1.556448in}{1.647903in}}%
\pgfpathlineto{\pgfqpoint{1.556448in}{1.647903in}}%
\pgfpathclose%
\pgfusepath{stroke,fill}%
\end{pgfscope}%
\begin{pgfscope}%
\pgfpathrectangle{\pgfqpoint{0.100000in}{0.209042in}}{\pgfqpoint{2.906250in}{2.906250in}}%
\pgfusepath{clip}%
\pgfsetbuttcap%
\pgfsetroundjoin%
\definecolor{currentfill}{rgb}{0.671349,0.200133,0.487358}%
\pgfsetfillcolor{currentfill}%
\pgfsetfillopacity{0.800000}%
\pgfsetlinewidth{1.003750pt}%
\definecolor{currentstroke}{rgb}{0.671349,0.200133,0.487358}%
\pgfsetstrokecolor{currentstroke}%
\pgfsetstrokeopacity{0.800000}%
\pgfsetdash{}{0pt}%
\pgfpathmoveto{\pgfqpoint{1.893190in}{1.559868in}}%
\pgfpathcurveto{\pgfqpoint{1.899014in}{1.559868in}}{\pgfqpoint{1.904600in}{1.562182in}}{\pgfqpoint{1.908718in}{1.566300in}}%
\pgfpathcurveto{\pgfqpoint{1.912836in}{1.570418in}}{\pgfqpoint{1.915150in}{1.576004in}}{\pgfqpoint{1.915150in}{1.581828in}}%
\pgfpathcurveto{\pgfqpoint{1.915150in}{1.587652in}}{\pgfqpoint{1.912836in}{1.593238in}}{\pgfqpoint{1.908718in}{1.597356in}}%
\pgfpathcurveto{\pgfqpoint{1.904600in}{1.601474in}}{\pgfqpoint{1.899014in}{1.603788in}}{\pgfqpoint{1.893190in}{1.603788in}}%
\pgfpathcurveto{\pgfqpoint{1.887366in}{1.603788in}}{\pgfqpoint{1.881780in}{1.601474in}}{\pgfqpoint{1.877662in}{1.597356in}}%
\pgfpathcurveto{\pgfqpoint{1.873544in}{1.593238in}}{\pgfqpoint{1.871230in}{1.587652in}}{\pgfqpoint{1.871230in}{1.581828in}}%
\pgfpathcurveto{\pgfqpoint{1.871230in}{1.576004in}}{\pgfqpoint{1.873544in}{1.570418in}}{\pgfqpoint{1.877662in}{1.566300in}}%
\pgfpathcurveto{\pgfqpoint{1.881780in}{1.562182in}}{\pgfqpoint{1.887366in}{1.559868in}}{\pgfqpoint{1.893190in}{1.559868in}}%
\pgfpathlineto{\pgfqpoint{1.893190in}{1.559868in}}%
\pgfpathclose%
\pgfusepath{stroke,fill}%
\end{pgfscope}%
\begin{pgfscope}%
\pgfpathrectangle{\pgfqpoint{0.100000in}{0.209042in}}{\pgfqpoint{2.906250in}{2.906250in}}%
\pgfusepath{clip}%
\pgfsetbuttcap%
\pgfsetroundjoin%
\definecolor{currentfill}{rgb}{0.786212,0.241514,0.450184}%
\pgfsetfillcolor{currentfill}%
\pgfsetfillopacity{0.800000}%
\pgfsetlinewidth{1.003750pt}%
\definecolor{currentstroke}{rgb}{0.786212,0.241514,0.450184}%
\pgfsetstrokecolor{currentstroke}%
\pgfsetstrokeopacity{0.800000}%
\pgfsetdash{}{0pt}%
\pgfpathmoveto{\pgfqpoint{1.112431in}{1.835837in}}%
\pgfpathcurveto{\pgfqpoint{1.118255in}{1.835837in}}{\pgfqpoint{1.123842in}{1.838150in}}{\pgfqpoint{1.127960in}{1.842269in}}%
\pgfpathcurveto{\pgfqpoint{1.132078in}{1.846387in}}{\pgfqpoint{1.134392in}{1.851973in}}{\pgfqpoint{1.134392in}{1.857797in}}%
\pgfpathcurveto{\pgfqpoint{1.134392in}{1.863621in}}{\pgfqpoint{1.132078in}{1.869207in}}{\pgfqpoint{1.127960in}{1.873325in}}%
\pgfpathcurveto{\pgfqpoint{1.123842in}{1.877443in}}{\pgfqpoint{1.118255in}{1.879757in}}{\pgfqpoint{1.112431in}{1.879757in}}%
\pgfpathcurveto{\pgfqpoint{1.106607in}{1.879757in}}{\pgfqpoint{1.101021in}{1.877443in}}{\pgfqpoint{1.096903in}{1.873325in}}%
\pgfpathcurveto{\pgfqpoint{1.092785in}{1.869207in}}{\pgfqpoint{1.090471in}{1.863621in}}{\pgfqpoint{1.090471in}{1.857797in}}%
\pgfpathcurveto{\pgfqpoint{1.090471in}{1.851973in}}{\pgfqpoint{1.092785in}{1.846387in}}{\pgfqpoint{1.096903in}{1.842269in}}%
\pgfpathcurveto{\pgfqpoint{1.101021in}{1.838150in}}{\pgfqpoint{1.106607in}{1.835837in}}{\pgfqpoint{1.112431in}{1.835837in}}%
\pgfpathlineto{\pgfqpoint{1.112431in}{1.835837in}}%
\pgfpathclose%
\pgfusepath{stroke,fill}%
\end{pgfscope}%
\begin{pgfscope}%
\pgfpathrectangle{\pgfqpoint{0.100000in}{0.209042in}}{\pgfqpoint{2.906250in}{2.906250in}}%
\pgfusepath{clip}%
\pgfsetbuttcap%
\pgfsetroundjoin%
\definecolor{currentfill}{rgb}{0.494258,0.141462,0.507988}%
\pgfsetfillcolor{currentfill}%
\pgfsetfillopacity{0.800000}%
\pgfsetlinewidth{1.003750pt}%
\definecolor{currentstroke}{rgb}{0.494258,0.141462,0.507988}%
\pgfsetstrokecolor{currentstroke}%
\pgfsetstrokeopacity{0.800000}%
\pgfsetdash{}{0pt}%
\pgfpathmoveto{\pgfqpoint{1.607843in}{1.582950in}}%
\pgfpathcurveto{\pgfqpoint{1.613667in}{1.582950in}}{\pgfqpoint{1.619253in}{1.585264in}}{\pgfqpoint{1.623371in}{1.589382in}}%
\pgfpathcurveto{\pgfqpoint{1.627489in}{1.593501in}}{\pgfqpoint{1.629803in}{1.599087in}}{\pgfqpoint{1.629803in}{1.604911in}}%
\pgfpathcurveto{\pgfqpoint{1.629803in}{1.610735in}}{\pgfqpoint{1.627489in}{1.616321in}}{\pgfqpoint{1.623371in}{1.620439in}}%
\pgfpathcurveto{\pgfqpoint{1.619253in}{1.624557in}}{\pgfqpoint{1.613667in}{1.626871in}}{\pgfqpoint{1.607843in}{1.626871in}}%
\pgfpathcurveto{\pgfqpoint{1.602019in}{1.626871in}}{\pgfqpoint{1.596433in}{1.624557in}}{\pgfqpoint{1.592315in}{1.620439in}}%
\pgfpathcurveto{\pgfqpoint{1.588196in}{1.616321in}}{\pgfqpoint{1.585883in}{1.610735in}}{\pgfqpoint{1.585883in}{1.604911in}}%
\pgfpathcurveto{\pgfqpoint{1.585883in}{1.599087in}}{\pgfqpoint{1.588196in}{1.593501in}}{\pgfqpoint{1.592315in}{1.589382in}}%
\pgfpathcurveto{\pgfqpoint{1.596433in}{1.585264in}}{\pgfqpoint{1.602019in}{1.582950in}}{\pgfqpoint{1.607843in}{1.582950in}}%
\pgfpathlineto{\pgfqpoint{1.607843in}{1.582950in}}%
\pgfpathclose%
\pgfusepath{stroke,fill}%
\end{pgfscope}%
\begin{pgfscope}%
\pgfpathrectangle{\pgfqpoint{0.100000in}{0.209042in}}{\pgfqpoint{2.906250in}{2.906250in}}%
\pgfusepath{clip}%
\pgfsetbuttcap%
\pgfsetroundjoin%
\definecolor{currentfill}{rgb}{0.874176,0.291859,0.406205}%
\pgfsetfillcolor{currentfill}%
\pgfsetfillopacity{0.800000}%
\pgfsetlinewidth{1.003750pt}%
\definecolor{currentstroke}{rgb}{0.874176,0.291859,0.406205}%
\pgfsetstrokecolor{currentstroke}%
\pgfsetstrokeopacity{0.800000}%
\pgfsetdash{}{0pt}%
\pgfpathmoveto{\pgfqpoint{1.792366in}{1.325309in}}%
\pgfpathcurveto{\pgfqpoint{1.798190in}{1.325309in}}{\pgfqpoint{1.803776in}{1.327623in}}{\pgfqpoint{1.807894in}{1.331741in}}%
\pgfpathcurveto{\pgfqpoint{1.812013in}{1.335859in}}{\pgfqpoint{1.814326in}{1.341445in}}{\pgfqpoint{1.814326in}{1.347269in}}%
\pgfpathcurveto{\pgfqpoint{1.814326in}{1.353093in}}{\pgfqpoint{1.812013in}{1.358679in}}{\pgfqpoint{1.807894in}{1.362797in}}%
\pgfpathcurveto{\pgfqpoint{1.803776in}{1.366916in}}{\pgfqpoint{1.798190in}{1.369229in}}{\pgfqpoint{1.792366in}{1.369229in}}%
\pgfpathcurveto{\pgfqpoint{1.786542in}{1.369229in}}{\pgfqpoint{1.780956in}{1.366916in}}{\pgfqpoint{1.776838in}{1.362797in}}%
\pgfpathcurveto{\pgfqpoint{1.772720in}{1.358679in}}{\pgfqpoint{1.770406in}{1.353093in}}{\pgfqpoint{1.770406in}{1.347269in}}%
\pgfpathcurveto{\pgfqpoint{1.770406in}{1.341445in}}{\pgfqpoint{1.772720in}{1.335859in}}{\pgfqpoint{1.776838in}{1.331741in}}%
\pgfpathcurveto{\pgfqpoint{1.780956in}{1.327623in}}{\pgfqpoint{1.786542in}{1.325309in}}{\pgfqpoint{1.792366in}{1.325309in}}%
\pgfpathlineto{\pgfqpoint{1.792366in}{1.325309in}}%
\pgfpathclose%
\pgfusepath{stroke,fill}%
\end{pgfscope}%
\begin{pgfscope}%
\pgfpathrectangle{\pgfqpoint{0.100000in}{0.209042in}}{\pgfqpoint{2.906250in}{2.906250in}}%
\pgfusepath{clip}%
\pgfsetbuttcap%
\pgfsetroundjoin%
\definecolor{currentfill}{rgb}{0.581819,0.171596,0.502777}%
\pgfsetfillcolor{currentfill}%
\pgfsetfillopacity{0.800000}%
\pgfsetlinewidth{1.003750pt}%
\definecolor{currentstroke}{rgb}{0.581819,0.171596,0.502777}%
\pgfsetstrokecolor{currentstroke}%
\pgfsetstrokeopacity{0.800000}%
\pgfsetdash{}{0pt}%
\pgfpathmoveto{\pgfqpoint{1.471691in}{1.517987in}}%
\pgfpathcurveto{\pgfqpoint{1.477515in}{1.517987in}}{\pgfqpoint{1.483101in}{1.520301in}}{\pgfqpoint{1.487219in}{1.524419in}}%
\pgfpathcurveto{\pgfqpoint{1.491337in}{1.528537in}}{\pgfqpoint{1.493651in}{1.534123in}}{\pgfqpoint{1.493651in}{1.539947in}}%
\pgfpathcurveto{\pgfqpoint{1.493651in}{1.545771in}}{\pgfqpoint{1.491337in}{1.551357in}}{\pgfqpoint{1.487219in}{1.555475in}}%
\pgfpathcurveto{\pgfqpoint{1.483101in}{1.559594in}}{\pgfqpoint{1.477515in}{1.561907in}}{\pgfqpoint{1.471691in}{1.561907in}}%
\pgfpathcurveto{\pgfqpoint{1.465867in}{1.561907in}}{\pgfqpoint{1.460281in}{1.559594in}}{\pgfqpoint{1.456162in}{1.555475in}}%
\pgfpathcurveto{\pgfqpoint{1.452044in}{1.551357in}}{\pgfqpoint{1.449730in}{1.545771in}}{\pgfqpoint{1.449730in}{1.539947in}}%
\pgfpathcurveto{\pgfqpoint{1.449730in}{1.534123in}}{\pgfqpoint{1.452044in}{1.528537in}}{\pgfqpoint{1.456162in}{1.524419in}}%
\pgfpathcurveto{\pgfqpoint{1.460281in}{1.520301in}}{\pgfqpoint{1.465867in}{1.517987in}}{\pgfqpoint{1.471691in}{1.517987in}}%
\pgfpathlineto{\pgfqpoint{1.471691in}{1.517987in}}%
\pgfpathclose%
\pgfusepath{stroke,fill}%
\end{pgfscope}%
\begin{pgfscope}%
\pgfpathrectangle{\pgfqpoint{0.100000in}{0.209042in}}{\pgfqpoint{2.906250in}{2.906250in}}%
\pgfusepath{clip}%
\pgfsetbuttcap%
\pgfsetroundjoin%
\definecolor{currentfill}{rgb}{0.414709,0.110431,0.504662}%
\pgfsetfillcolor{currentfill}%
\pgfsetfillopacity{0.800000}%
\pgfsetlinewidth{1.003750pt}%
\definecolor{currentstroke}{rgb}{0.414709,0.110431,0.504662}%
\pgfsetstrokecolor{currentstroke}%
\pgfsetstrokeopacity{0.800000}%
\pgfsetdash{}{0pt}%
\pgfpathmoveto{\pgfqpoint{1.528939in}{1.830336in}}%
\pgfpathcurveto{\pgfqpoint{1.534763in}{1.830336in}}{\pgfqpoint{1.540349in}{1.832649in}}{\pgfqpoint{1.544467in}{1.836768in}}%
\pgfpathcurveto{\pgfqpoint{1.548586in}{1.840886in}}{\pgfqpoint{1.550899in}{1.846472in}}{\pgfqpoint{1.550899in}{1.852296in}}%
\pgfpathcurveto{\pgfqpoint{1.550899in}{1.858120in}}{\pgfqpoint{1.548586in}{1.863706in}}{\pgfqpoint{1.544467in}{1.867824in}}%
\pgfpathcurveto{\pgfqpoint{1.540349in}{1.871942in}}{\pgfqpoint{1.534763in}{1.874256in}}{\pgfqpoint{1.528939in}{1.874256in}}%
\pgfpathcurveto{\pgfqpoint{1.523115in}{1.874256in}}{\pgfqpoint{1.517529in}{1.871942in}}{\pgfqpoint{1.513411in}{1.867824in}}%
\pgfpathcurveto{\pgfqpoint{1.509293in}{1.863706in}}{\pgfqpoint{1.506979in}{1.858120in}}{\pgfqpoint{1.506979in}{1.852296in}}%
\pgfpathcurveto{\pgfqpoint{1.506979in}{1.846472in}}{\pgfqpoint{1.509293in}{1.840886in}}{\pgfqpoint{1.513411in}{1.836768in}}%
\pgfpathcurveto{\pgfqpoint{1.517529in}{1.832649in}}{\pgfqpoint{1.523115in}{1.830336in}}{\pgfqpoint{1.528939in}{1.830336in}}%
\pgfpathlineto{\pgfqpoint{1.528939in}{1.830336in}}%
\pgfpathclose%
\pgfusepath{stroke,fill}%
\end{pgfscope}%
\begin{pgfscope}%
\pgfpathrectangle{\pgfqpoint{0.100000in}{0.209042in}}{\pgfqpoint{2.906250in}{2.906250in}}%
\pgfusepath{clip}%
\pgfsetbuttcap%
\pgfsetroundjoin%
\definecolor{currentfill}{rgb}{0.488088,0.139186,0.508011}%
\pgfsetfillcolor{currentfill}%
\pgfsetfillopacity{0.800000}%
\pgfsetlinewidth{1.003750pt}%
\definecolor{currentstroke}{rgb}{0.488088,0.139186,0.508011}%
\pgfsetstrokecolor{currentstroke}%
\pgfsetstrokeopacity{0.800000}%
\pgfsetdash{}{0pt}%
\pgfpathmoveto{\pgfqpoint{1.809256in}{1.719791in}}%
\pgfpathcurveto{\pgfqpoint{1.815080in}{1.719791in}}{\pgfqpoint{1.820666in}{1.722104in}}{\pgfqpoint{1.824784in}{1.726223in}}%
\pgfpathcurveto{\pgfqpoint{1.828902in}{1.730341in}}{\pgfqpoint{1.831216in}{1.735927in}}{\pgfqpoint{1.831216in}{1.741751in}}%
\pgfpathcurveto{\pgfqpoint{1.831216in}{1.747575in}}{\pgfqpoint{1.828902in}{1.753161in}}{\pgfqpoint{1.824784in}{1.757279in}}%
\pgfpathcurveto{\pgfqpoint{1.820666in}{1.761397in}}{\pgfqpoint{1.815080in}{1.763711in}}{\pgfqpoint{1.809256in}{1.763711in}}%
\pgfpathcurveto{\pgfqpoint{1.803432in}{1.763711in}}{\pgfqpoint{1.797846in}{1.761397in}}{\pgfqpoint{1.793728in}{1.757279in}}%
\pgfpathcurveto{\pgfqpoint{1.789610in}{1.753161in}}{\pgfqpoint{1.787296in}{1.747575in}}{\pgfqpoint{1.787296in}{1.741751in}}%
\pgfpathcurveto{\pgfqpoint{1.787296in}{1.735927in}}{\pgfqpoint{1.789610in}{1.730341in}}{\pgfqpoint{1.793728in}{1.726223in}}%
\pgfpathcurveto{\pgfqpoint{1.797846in}{1.722104in}}{\pgfqpoint{1.803432in}{1.719791in}}{\pgfqpoint{1.809256in}{1.719791in}}%
\pgfpathlineto{\pgfqpoint{1.809256in}{1.719791in}}%
\pgfpathclose%
\pgfusepath{stroke,fill}%
\end{pgfscope}%
\begin{pgfscope}%
\pgfpathrectangle{\pgfqpoint{0.100000in}{0.209042in}}{\pgfqpoint{2.906250in}{2.906250in}}%
\pgfusepath{clip}%
\pgfsetbuttcap%
\pgfsetroundjoin%
\definecolor{currentfill}{rgb}{0.469640,0.132245,0.507809}%
\pgfsetfillcolor{currentfill}%
\pgfsetfillopacity{0.800000}%
\pgfsetlinewidth{1.003750pt}%
\definecolor{currentstroke}{rgb}{0.469640,0.132245,0.507809}%
\pgfsetstrokecolor{currentstroke}%
\pgfsetstrokeopacity{0.800000}%
\pgfsetdash{}{0pt}%
\pgfpathmoveto{\pgfqpoint{1.382932in}{1.729259in}}%
\pgfpathcurveto{\pgfqpoint{1.388756in}{1.729259in}}{\pgfqpoint{1.394342in}{1.731573in}}{\pgfqpoint{1.398460in}{1.735691in}}%
\pgfpathcurveto{\pgfqpoint{1.402578in}{1.739809in}}{\pgfqpoint{1.404892in}{1.745396in}}{\pgfqpoint{1.404892in}{1.751220in}}%
\pgfpathcurveto{\pgfqpoint{1.404892in}{1.757044in}}{\pgfqpoint{1.402578in}{1.762630in}}{\pgfqpoint{1.398460in}{1.766748in}}%
\pgfpathcurveto{\pgfqpoint{1.394342in}{1.770866in}}{\pgfqpoint{1.388756in}{1.773180in}}{\pgfqpoint{1.382932in}{1.773180in}}%
\pgfpathcurveto{\pgfqpoint{1.377108in}{1.773180in}}{\pgfqpoint{1.371522in}{1.770866in}}{\pgfqpoint{1.367404in}{1.766748in}}%
\pgfpathcurveto{\pgfqpoint{1.363286in}{1.762630in}}{\pgfqpoint{1.360972in}{1.757044in}}{\pgfqpoint{1.360972in}{1.751220in}}%
\pgfpathcurveto{\pgfqpoint{1.360972in}{1.745396in}}{\pgfqpoint{1.363286in}{1.739809in}}{\pgfqpoint{1.367404in}{1.735691in}}%
\pgfpathcurveto{\pgfqpoint{1.371522in}{1.731573in}}{\pgfqpoint{1.377108in}{1.729259in}}{\pgfqpoint{1.382932in}{1.729259in}}%
\pgfpathlineto{\pgfqpoint{1.382932in}{1.729259in}}%
\pgfpathclose%
\pgfusepath{stroke,fill}%
\end{pgfscope}%
\begin{pgfscope}%
\pgfpathrectangle{\pgfqpoint{0.100000in}{0.209042in}}{\pgfqpoint{2.906250in}{2.906250in}}%
\pgfusepath{clip}%
\pgfsetbuttcap%
\pgfsetroundjoin%
\definecolor{currentfill}{rgb}{0.639216,0.189921,0.494150}%
\pgfsetfillcolor{currentfill}%
\pgfsetfillopacity{0.800000}%
\pgfsetlinewidth{1.003750pt}%
\definecolor{currentstroke}{rgb}{0.639216,0.189921,0.494150}%
\pgfsetstrokecolor{currentstroke}%
\pgfsetstrokeopacity{0.800000}%
\pgfsetdash{}{0pt}%
\pgfpathmoveto{\pgfqpoint{1.586486in}{1.445370in}}%
\pgfpathcurveto{\pgfqpoint{1.592310in}{1.445370in}}{\pgfqpoint{1.597896in}{1.447683in}}{\pgfqpoint{1.602014in}{1.451802in}}%
\pgfpathcurveto{\pgfqpoint{1.606132in}{1.455920in}}{\pgfqpoint{1.608446in}{1.461506in}}{\pgfqpoint{1.608446in}{1.467330in}}%
\pgfpathcurveto{\pgfqpoint{1.608446in}{1.473154in}}{\pgfqpoint{1.606132in}{1.478740in}}{\pgfqpoint{1.602014in}{1.482858in}}%
\pgfpathcurveto{\pgfqpoint{1.597896in}{1.486976in}}{\pgfqpoint{1.592310in}{1.489290in}}{\pgfqpoint{1.586486in}{1.489290in}}%
\pgfpathcurveto{\pgfqpoint{1.580662in}{1.489290in}}{\pgfqpoint{1.575076in}{1.486976in}}{\pgfqpoint{1.570958in}{1.482858in}}%
\pgfpathcurveto{\pgfqpoint{1.566839in}{1.478740in}}{\pgfqpoint{1.564525in}{1.473154in}}{\pgfqpoint{1.564525in}{1.467330in}}%
\pgfpathcurveto{\pgfqpoint{1.564525in}{1.461506in}}{\pgfqpoint{1.566839in}{1.455920in}}{\pgfqpoint{1.570958in}{1.451802in}}%
\pgfpathcurveto{\pgfqpoint{1.575076in}{1.447683in}}{\pgfqpoint{1.580662in}{1.445370in}}{\pgfqpoint{1.586486in}{1.445370in}}%
\pgfpathlineto{\pgfqpoint{1.586486in}{1.445370in}}%
\pgfpathclose%
\pgfusepath{stroke,fill}%
\end{pgfscope}%
\begin{pgfscope}%
\pgfpathrectangle{\pgfqpoint{0.100000in}{0.209042in}}{\pgfqpoint{2.906250in}{2.906250in}}%
\pgfusepath{clip}%
\pgfsetbuttcap%
\pgfsetroundjoin%
\definecolor{currentfill}{rgb}{0.834791,0.265540,0.427671}%
\pgfsetfillcolor{currentfill}%
\pgfsetfillopacity{0.800000}%
\pgfsetlinewidth{1.003750pt}%
\definecolor{currentstroke}{rgb}{0.834791,0.265540,0.427671}%
\pgfsetstrokecolor{currentstroke}%
\pgfsetstrokeopacity{0.800000}%
\pgfsetdash{}{0pt}%
\pgfpathmoveto{\pgfqpoint{1.323125in}{1.384280in}}%
\pgfpathcurveto{\pgfqpoint{1.328949in}{1.384280in}}{\pgfqpoint{1.334535in}{1.386594in}}{\pgfqpoint{1.338653in}{1.390712in}}%
\pgfpathcurveto{\pgfqpoint{1.342771in}{1.394831in}}{\pgfqpoint{1.345085in}{1.400417in}}{\pgfqpoint{1.345085in}{1.406241in}}%
\pgfpathcurveto{\pgfqpoint{1.345085in}{1.412065in}}{\pgfqpoint{1.342771in}{1.417651in}}{\pgfqpoint{1.338653in}{1.421769in}}%
\pgfpathcurveto{\pgfqpoint{1.334535in}{1.425887in}}{\pgfqpoint{1.328949in}{1.428201in}}{\pgfqpoint{1.323125in}{1.428201in}}%
\pgfpathcurveto{\pgfqpoint{1.317301in}{1.428201in}}{\pgfqpoint{1.311715in}{1.425887in}}{\pgfqpoint{1.307597in}{1.421769in}}%
\pgfpathcurveto{\pgfqpoint{1.303479in}{1.417651in}}{\pgfqpoint{1.301165in}{1.412065in}}{\pgfqpoint{1.301165in}{1.406241in}}%
\pgfpathcurveto{\pgfqpoint{1.301165in}{1.400417in}}{\pgfqpoint{1.303479in}{1.394831in}}{\pgfqpoint{1.307597in}{1.390712in}}%
\pgfpathcurveto{\pgfqpoint{1.311715in}{1.386594in}}{\pgfqpoint{1.317301in}{1.384280in}}{\pgfqpoint{1.323125in}{1.384280in}}%
\pgfpathlineto{\pgfqpoint{1.323125in}{1.384280in}}%
\pgfpathclose%
\pgfusepath{stroke,fill}%
\end{pgfscope}%
\begin{pgfscope}%
\pgfpathrectangle{\pgfqpoint{0.100000in}{0.209042in}}{\pgfqpoint{2.906250in}{2.906250in}}%
\pgfusepath{clip}%
\pgfsetbuttcap%
\pgfsetroundjoin%
\definecolor{currentfill}{rgb}{0.494258,0.141462,0.507988}%
\pgfsetfillcolor{currentfill}%
\pgfsetfillopacity{0.800000}%
\pgfsetlinewidth{1.003750pt}%
\definecolor{currentstroke}{rgb}{0.494258,0.141462,0.507988}%
\pgfsetstrokecolor{currentstroke}%
\pgfsetstrokeopacity{0.800000}%
\pgfsetdash{}{0pt}%
\pgfpathmoveto{\pgfqpoint{1.651399in}{1.561036in}}%
\pgfpathcurveto{\pgfqpoint{1.657223in}{1.561036in}}{\pgfqpoint{1.662810in}{1.563350in}}{\pgfqpoint{1.666928in}{1.567468in}}%
\pgfpathcurveto{\pgfqpoint{1.671046in}{1.571586in}}{\pgfqpoint{1.673360in}{1.577172in}}{\pgfqpoint{1.673360in}{1.582996in}}%
\pgfpathcurveto{\pgfqpoint{1.673360in}{1.588820in}}{\pgfqpoint{1.671046in}{1.594406in}}{\pgfqpoint{1.666928in}{1.598524in}}%
\pgfpathcurveto{\pgfqpoint{1.662810in}{1.602643in}}{\pgfqpoint{1.657223in}{1.604956in}}{\pgfqpoint{1.651399in}{1.604956in}}%
\pgfpathcurveto{\pgfqpoint{1.645575in}{1.604956in}}{\pgfqpoint{1.639989in}{1.602643in}}{\pgfqpoint{1.635871in}{1.598524in}}%
\pgfpathcurveto{\pgfqpoint{1.631753in}{1.594406in}}{\pgfqpoint{1.629439in}{1.588820in}}{\pgfqpoint{1.629439in}{1.582996in}}%
\pgfpathcurveto{\pgfqpoint{1.629439in}{1.577172in}}{\pgfqpoint{1.631753in}{1.571586in}}{\pgfqpoint{1.635871in}{1.567468in}}%
\pgfpathcurveto{\pgfqpoint{1.639989in}{1.563350in}}{\pgfqpoint{1.645575in}{1.561036in}}{\pgfqpoint{1.651399in}{1.561036in}}%
\pgfpathlineto{\pgfqpoint{1.651399in}{1.561036in}}%
\pgfpathclose%
\pgfusepath{stroke,fill}%
\end{pgfscope}%
\begin{pgfscope}%
\pgfpathrectangle{\pgfqpoint{0.100000in}{0.209042in}}{\pgfqpoint{2.906250in}{2.906250in}}%
\pgfusepath{clip}%
\pgfsetbuttcap%
\pgfsetroundjoin%
\definecolor{currentfill}{rgb}{0.677786,0.202203,0.485819}%
\pgfsetfillcolor{currentfill}%
\pgfsetfillopacity{0.800000}%
\pgfsetlinewidth{1.003750pt}%
\definecolor{currentstroke}{rgb}{0.677786,0.202203,0.485819}%
\pgfsetstrokecolor{currentstroke}%
\pgfsetstrokeopacity{0.800000}%
\pgfsetdash{}{0pt}%
\pgfpathmoveto{\pgfqpoint{1.326394in}{2.018961in}}%
\pgfpathcurveto{\pgfqpoint{1.332217in}{2.018961in}}{\pgfqpoint{1.337804in}{2.021275in}}{\pgfqpoint{1.341922in}{2.025393in}}%
\pgfpathcurveto{\pgfqpoint{1.346040in}{2.029511in}}{\pgfqpoint{1.348354in}{2.035098in}}{\pgfqpoint{1.348354in}{2.040922in}}%
\pgfpathcurveto{\pgfqpoint{1.348354in}{2.046746in}}{\pgfqpoint{1.346040in}{2.052332in}}{\pgfqpoint{1.341922in}{2.056450in}}%
\pgfpathcurveto{\pgfqpoint{1.337804in}{2.060568in}}{\pgfqpoint{1.332217in}{2.062882in}}{\pgfqpoint{1.326394in}{2.062882in}}%
\pgfpathcurveto{\pgfqpoint{1.320570in}{2.062882in}}{\pgfqpoint{1.314983in}{2.060568in}}{\pgfqpoint{1.310865in}{2.056450in}}%
\pgfpathcurveto{\pgfqpoint{1.306747in}{2.052332in}}{\pgfqpoint{1.304433in}{2.046746in}}{\pgfqpoint{1.304433in}{2.040922in}}%
\pgfpathcurveto{\pgfqpoint{1.304433in}{2.035098in}}{\pgfqpoint{1.306747in}{2.029511in}}{\pgfqpoint{1.310865in}{2.025393in}}%
\pgfpathcurveto{\pgfqpoint{1.314983in}{2.021275in}}{\pgfqpoint{1.320570in}{2.018961in}}{\pgfqpoint{1.326394in}{2.018961in}}%
\pgfpathlineto{\pgfqpoint{1.326394in}{2.018961in}}%
\pgfpathclose%
\pgfusepath{stroke,fill}%
\end{pgfscope}%
\begin{pgfscope}%
\pgfpathrectangle{\pgfqpoint{0.100000in}{0.209042in}}{\pgfqpoint{2.906250in}{2.906250in}}%
\pgfusepath{clip}%
\pgfsetbuttcap%
\pgfsetroundjoin%
\definecolor{currentfill}{rgb}{0.613617,0.181811,0.498536}%
\pgfsetfillcolor{currentfill}%
\pgfsetfillopacity{0.800000}%
\pgfsetlinewidth{1.003750pt}%
\definecolor{currentstroke}{rgb}{0.613617,0.181811,0.498536}%
\pgfsetstrokecolor{currentstroke}%
\pgfsetstrokeopacity{0.800000}%
\pgfsetdash{}{0pt}%
\pgfpathmoveto{\pgfqpoint{1.526058in}{1.464788in}}%
\pgfpathcurveto{\pgfqpoint{1.531882in}{1.464788in}}{\pgfqpoint{1.537469in}{1.467102in}}{\pgfqpoint{1.541587in}{1.471220in}}%
\pgfpathcurveto{\pgfqpoint{1.545705in}{1.475338in}}{\pgfqpoint{1.548019in}{1.480924in}}{\pgfqpoint{1.548019in}{1.486748in}}%
\pgfpathcurveto{\pgfqpoint{1.548019in}{1.492572in}}{\pgfqpoint{1.545705in}{1.498158in}}{\pgfqpoint{1.541587in}{1.502276in}}%
\pgfpathcurveto{\pgfqpoint{1.537469in}{1.506394in}}{\pgfqpoint{1.531882in}{1.508708in}}{\pgfqpoint{1.526058in}{1.508708in}}%
\pgfpathcurveto{\pgfqpoint{1.520234in}{1.508708in}}{\pgfqpoint{1.514648in}{1.506394in}}{\pgfqpoint{1.510530in}{1.502276in}}%
\pgfpathcurveto{\pgfqpoint{1.506412in}{1.498158in}}{\pgfqpoint{1.504098in}{1.492572in}}{\pgfqpoint{1.504098in}{1.486748in}}%
\pgfpathcurveto{\pgfqpoint{1.504098in}{1.480924in}}{\pgfqpoint{1.506412in}{1.475338in}}{\pgfqpoint{1.510530in}{1.471220in}}%
\pgfpathcurveto{\pgfqpoint{1.514648in}{1.467102in}}{\pgfqpoint{1.520234in}{1.464788in}}{\pgfqpoint{1.526058in}{1.464788in}}%
\pgfpathlineto{\pgfqpoint{1.526058in}{1.464788in}}%
\pgfpathclose%
\pgfusepath{stroke,fill}%
\end{pgfscope}%
\begin{pgfscope}%
\pgfpathrectangle{\pgfqpoint{0.100000in}{0.209042in}}{\pgfqpoint{2.906250in}{2.906250in}}%
\pgfusepath{clip}%
\pgfsetbuttcap%
\pgfsetroundjoin%
\definecolor{currentfill}{rgb}{0.690661,0.206384,0.482558}%
\pgfsetfillcolor{currentfill}%
\pgfsetfillopacity{0.800000}%
\pgfsetlinewidth{1.003750pt}%
\definecolor{currentstroke}{rgb}{0.690661,0.206384,0.482558}%
\pgfsetstrokecolor{currentstroke}%
\pgfsetstrokeopacity{0.800000}%
\pgfsetdash{}{0pt}%
\pgfpathmoveto{\pgfqpoint{1.682903in}{2.080730in}}%
\pgfpathcurveto{\pgfqpoint{1.688727in}{2.080730in}}{\pgfqpoint{1.694313in}{2.083044in}}{\pgfqpoint{1.698431in}{2.087162in}}%
\pgfpathcurveto{\pgfqpoint{1.702549in}{2.091280in}}{\pgfqpoint{1.704863in}{2.096866in}}{\pgfqpoint{1.704863in}{2.102690in}}%
\pgfpathcurveto{\pgfqpoint{1.704863in}{2.108514in}}{\pgfqpoint{1.702549in}{2.114100in}}{\pgfqpoint{1.698431in}{2.118218in}}%
\pgfpathcurveto{\pgfqpoint{1.694313in}{2.122336in}}{\pgfqpoint{1.688727in}{2.124650in}}{\pgfqpoint{1.682903in}{2.124650in}}%
\pgfpathcurveto{\pgfqpoint{1.677079in}{2.124650in}}{\pgfqpoint{1.671493in}{2.122336in}}{\pgfqpoint{1.667375in}{2.118218in}}%
\pgfpathcurveto{\pgfqpoint{1.663256in}{2.114100in}}{\pgfqpoint{1.660943in}{2.108514in}}{\pgfqpoint{1.660943in}{2.102690in}}%
\pgfpathcurveto{\pgfqpoint{1.660943in}{2.096866in}}{\pgfqpoint{1.663256in}{2.091280in}}{\pgfqpoint{1.667375in}{2.087162in}}%
\pgfpathcurveto{\pgfqpoint{1.671493in}{2.083044in}}{\pgfqpoint{1.677079in}{2.080730in}}{\pgfqpoint{1.682903in}{2.080730in}}%
\pgfpathlineto{\pgfqpoint{1.682903in}{2.080730in}}%
\pgfpathclose%
\pgfusepath{stroke,fill}%
\end{pgfscope}%
\begin{pgfscope}%
\pgfpathrectangle{\pgfqpoint{0.100000in}{0.209042in}}{\pgfqpoint{2.906250in}{2.906250in}}%
\pgfusepath{clip}%
\pgfsetbuttcap%
\pgfsetroundjoin%
\definecolor{currentfill}{rgb}{0.488088,0.139186,0.508011}%
\pgfsetfillcolor{currentfill}%
\pgfsetfillopacity{0.800000}%
\pgfsetlinewidth{1.003750pt}%
\definecolor{currentstroke}{rgb}{0.488088,0.139186,0.508011}%
\pgfsetstrokecolor{currentstroke}%
\pgfsetstrokeopacity{0.800000}%
\pgfsetdash{}{0pt}%
\pgfpathmoveto{\pgfqpoint{1.607321in}{1.544905in}}%
\pgfpathcurveto{\pgfqpoint{1.613145in}{1.544905in}}{\pgfqpoint{1.618731in}{1.547218in}}{\pgfqpoint{1.622849in}{1.551337in}}%
\pgfpathcurveto{\pgfqpoint{1.626967in}{1.555455in}}{\pgfqpoint{1.629281in}{1.561041in}}{\pgfqpoint{1.629281in}{1.566865in}}%
\pgfpathcurveto{\pgfqpoint{1.629281in}{1.572689in}}{\pgfqpoint{1.626967in}{1.578275in}}{\pgfqpoint{1.622849in}{1.582393in}}%
\pgfpathcurveto{\pgfqpoint{1.618731in}{1.586511in}}{\pgfqpoint{1.613145in}{1.588825in}}{\pgfqpoint{1.607321in}{1.588825in}}%
\pgfpathcurveto{\pgfqpoint{1.601497in}{1.588825in}}{\pgfqpoint{1.595911in}{1.586511in}}{\pgfqpoint{1.591792in}{1.582393in}}%
\pgfpathcurveto{\pgfqpoint{1.587674in}{1.578275in}}{\pgfqpoint{1.585360in}{1.572689in}}{\pgfqpoint{1.585360in}{1.566865in}}%
\pgfpathcurveto{\pgfqpoint{1.585360in}{1.561041in}}{\pgfqpoint{1.587674in}{1.555455in}}{\pgfqpoint{1.591792in}{1.551337in}}%
\pgfpathcurveto{\pgfqpoint{1.595911in}{1.547218in}}{\pgfqpoint{1.601497in}{1.544905in}}{\pgfqpoint{1.607321in}{1.544905in}}%
\pgfpathlineto{\pgfqpoint{1.607321in}{1.544905in}}%
\pgfpathclose%
\pgfusepath{stroke,fill}%
\end{pgfscope}%
\begin{pgfscope}%
\pgfpathrectangle{\pgfqpoint{0.100000in}{0.209042in}}{\pgfqpoint{2.906250in}{2.906250in}}%
\pgfusepath{clip}%
\pgfsetbuttcap%
\pgfsetroundjoin%
\definecolor{currentfill}{rgb}{0.607238,0.179779,0.499492}%
\pgfsetfillcolor{currentfill}%
\pgfsetfillopacity{0.800000}%
\pgfsetlinewidth{1.003750pt}%
\definecolor{currentstroke}{rgb}{0.607238,0.179779,0.499492}%
\pgfsetstrokecolor{currentstroke}%
\pgfsetstrokeopacity{0.800000}%
\pgfsetdash{}{0pt}%
\pgfpathmoveto{\pgfqpoint{1.251794in}{1.875678in}}%
\pgfpathcurveto{\pgfqpoint{1.257618in}{1.875678in}}{\pgfqpoint{1.263204in}{1.877992in}}{\pgfqpoint{1.267322in}{1.882110in}}%
\pgfpathcurveto{\pgfqpoint{1.271440in}{1.886228in}}{\pgfqpoint{1.273754in}{1.891815in}}{\pgfqpoint{1.273754in}{1.897638in}}%
\pgfpathcurveto{\pgfqpoint{1.273754in}{1.903462in}}{\pgfqpoint{1.271440in}{1.909049in}}{\pgfqpoint{1.267322in}{1.913167in}}%
\pgfpathcurveto{\pgfqpoint{1.263204in}{1.917285in}}{\pgfqpoint{1.257618in}{1.919599in}}{\pgfqpoint{1.251794in}{1.919599in}}%
\pgfpathcurveto{\pgfqpoint{1.245970in}{1.919599in}}{\pgfqpoint{1.240384in}{1.917285in}}{\pgfqpoint{1.236266in}{1.913167in}}%
\pgfpathcurveto{\pgfqpoint{1.232147in}{1.909049in}}{\pgfqpoint{1.229834in}{1.903462in}}{\pgfqpoint{1.229834in}{1.897638in}}%
\pgfpathcurveto{\pgfqpoint{1.229834in}{1.891815in}}{\pgfqpoint{1.232147in}{1.886228in}}{\pgfqpoint{1.236266in}{1.882110in}}%
\pgfpathcurveto{\pgfqpoint{1.240384in}{1.877992in}}{\pgfqpoint{1.245970in}{1.875678in}}{\pgfqpoint{1.251794in}{1.875678in}}%
\pgfpathlineto{\pgfqpoint{1.251794in}{1.875678in}}%
\pgfpathclose%
\pgfusepath{stroke,fill}%
\end{pgfscope}%
\begin{pgfscope}%
\pgfpathrectangle{\pgfqpoint{0.100000in}{0.209042in}}{\pgfqpoint{2.906250in}{2.906250in}}%
\pgfusepath{clip}%
\pgfsetbuttcap%
\pgfsetroundjoin%
\definecolor{currentfill}{rgb}{0.488088,0.139186,0.508011}%
\pgfsetfillcolor{currentfill}%
\pgfsetfillopacity{0.800000}%
\pgfsetlinewidth{1.003750pt}%
\definecolor{currentstroke}{rgb}{0.488088,0.139186,0.508011}%
\pgfsetstrokecolor{currentstroke}%
\pgfsetstrokeopacity{0.800000}%
\pgfsetdash{}{0pt}%
\pgfpathmoveto{\pgfqpoint{1.856560in}{1.739663in}}%
\pgfpathcurveto{\pgfqpoint{1.862384in}{1.739663in}}{\pgfqpoint{1.867970in}{1.741977in}}{\pgfqpoint{1.872088in}{1.746095in}}%
\pgfpathcurveto{\pgfqpoint{1.876207in}{1.750214in}}{\pgfqpoint{1.878520in}{1.755800in}}{\pgfqpoint{1.878520in}{1.761624in}}%
\pgfpathcurveto{\pgfqpoint{1.878520in}{1.767448in}}{\pgfqpoint{1.876207in}{1.773034in}}{\pgfqpoint{1.872088in}{1.777152in}}%
\pgfpathcurveto{\pgfqpoint{1.867970in}{1.781270in}}{\pgfqpoint{1.862384in}{1.783584in}}{\pgfqpoint{1.856560in}{1.783584in}}%
\pgfpathcurveto{\pgfqpoint{1.850736in}{1.783584in}}{\pgfqpoint{1.845150in}{1.781270in}}{\pgfqpoint{1.841032in}{1.777152in}}%
\pgfpathcurveto{\pgfqpoint{1.836914in}{1.773034in}}{\pgfqpoint{1.834600in}{1.767448in}}{\pgfqpoint{1.834600in}{1.761624in}}%
\pgfpathcurveto{\pgfqpoint{1.834600in}{1.755800in}}{\pgfqpoint{1.836914in}{1.750214in}}{\pgfqpoint{1.841032in}{1.746095in}}%
\pgfpathcurveto{\pgfqpoint{1.845150in}{1.741977in}}{\pgfqpoint{1.850736in}{1.739663in}}{\pgfqpoint{1.856560in}{1.739663in}}%
\pgfpathlineto{\pgfqpoint{1.856560in}{1.739663in}}%
\pgfpathclose%
\pgfusepath{stroke,fill}%
\end{pgfscope}%
\begin{pgfscope}%
\pgfpathrectangle{\pgfqpoint{0.100000in}{0.209042in}}{\pgfqpoint{2.906250in}{2.906250in}}%
\pgfusepath{clip}%
\pgfsetbuttcap%
\pgfsetroundjoin%
\definecolor{currentfill}{rgb}{0.506629,0.145958,0.507806}%
\pgfsetfillcolor{currentfill}%
\pgfsetfillopacity{0.800000}%
\pgfsetlinewidth{1.003750pt}%
\definecolor{currentstroke}{rgb}{0.506629,0.145958,0.507806}%
\pgfsetstrokecolor{currentstroke}%
\pgfsetstrokeopacity{0.800000}%
\pgfsetdash{}{0pt}%
\pgfpathmoveto{\pgfqpoint{1.734454in}{1.947067in}}%
\pgfpathcurveto{\pgfqpoint{1.740278in}{1.947067in}}{\pgfqpoint{1.745864in}{1.949381in}}{\pgfqpoint{1.749982in}{1.953499in}}%
\pgfpathcurveto{\pgfqpoint{1.754100in}{1.957617in}}{\pgfqpoint{1.756414in}{1.963203in}}{\pgfqpoint{1.756414in}{1.969027in}}%
\pgfpathcurveto{\pgfqpoint{1.756414in}{1.974851in}}{\pgfqpoint{1.754100in}{1.980437in}}{\pgfqpoint{1.749982in}{1.984555in}}%
\pgfpathcurveto{\pgfqpoint{1.745864in}{1.988673in}}{\pgfqpoint{1.740278in}{1.990987in}}{\pgfqpoint{1.734454in}{1.990987in}}%
\pgfpathcurveto{\pgfqpoint{1.728630in}{1.990987in}}{\pgfqpoint{1.723043in}{1.988673in}}{\pgfqpoint{1.718925in}{1.984555in}}%
\pgfpathcurveto{\pgfqpoint{1.714807in}{1.980437in}}{\pgfqpoint{1.712493in}{1.974851in}}{\pgfqpoint{1.712493in}{1.969027in}}%
\pgfpathcurveto{\pgfqpoint{1.712493in}{1.963203in}}{\pgfqpoint{1.714807in}{1.957617in}}{\pgfqpoint{1.718925in}{1.953499in}}%
\pgfpathcurveto{\pgfqpoint{1.723043in}{1.949381in}}{\pgfqpoint{1.728630in}{1.947067in}}{\pgfqpoint{1.734454in}{1.947067in}}%
\pgfpathlineto{\pgfqpoint{1.734454in}{1.947067in}}%
\pgfpathclose%
\pgfusepath{stroke,fill}%
\end{pgfscope}%
\begin{pgfscope}%
\pgfpathrectangle{\pgfqpoint{0.100000in}{0.209042in}}{\pgfqpoint{2.906250in}{2.906250in}}%
\pgfusepath{clip}%
\pgfsetbuttcap%
\pgfsetroundjoin%
\definecolor{currentfill}{rgb}{0.816914,0.255895,0.436461}%
\pgfsetfillcolor{currentfill}%
\pgfsetfillopacity{0.800000}%
\pgfsetlinewidth{1.003750pt}%
\definecolor{currentstroke}{rgb}{0.816914,0.255895,0.436461}%
\pgfsetstrokecolor{currentstroke}%
\pgfsetstrokeopacity{0.800000}%
\pgfsetdash{}{0pt}%
\pgfpathmoveto{\pgfqpoint{1.952191in}{2.050812in}}%
\pgfpathcurveto{\pgfqpoint{1.958015in}{2.050812in}}{\pgfqpoint{1.963601in}{2.053125in}}{\pgfqpoint{1.967719in}{2.057244in}}%
\pgfpathcurveto{\pgfqpoint{1.971837in}{2.061362in}}{\pgfqpoint{1.974151in}{2.066948in}}{\pgfqpoint{1.974151in}{2.072772in}}%
\pgfpathcurveto{\pgfqpoint{1.974151in}{2.078596in}}{\pgfqpoint{1.971837in}{2.084182in}}{\pgfqpoint{1.967719in}{2.088300in}}%
\pgfpathcurveto{\pgfqpoint{1.963601in}{2.092418in}}{\pgfqpoint{1.958015in}{2.094732in}}{\pgfqpoint{1.952191in}{2.094732in}}%
\pgfpathcurveto{\pgfqpoint{1.946367in}{2.094732in}}{\pgfqpoint{1.940781in}{2.092418in}}{\pgfqpoint{1.936663in}{2.088300in}}%
\pgfpathcurveto{\pgfqpoint{1.932545in}{2.084182in}}{\pgfqpoint{1.930231in}{2.078596in}}{\pgfqpoint{1.930231in}{2.072772in}}%
\pgfpathcurveto{\pgfqpoint{1.930231in}{2.066948in}}{\pgfqpoint{1.932545in}{2.061362in}}{\pgfqpoint{1.936663in}{2.057244in}}%
\pgfpathcurveto{\pgfqpoint{1.940781in}{2.053125in}}{\pgfqpoint{1.946367in}{2.050812in}}{\pgfqpoint{1.952191in}{2.050812in}}%
\pgfpathlineto{\pgfqpoint{1.952191in}{2.050812in}}%
\pgfpathclose%
\pgfusepath{stroke,fill}%
\end{pgfscope}%
\begin{pgfscope}%
\pgfpathrectangle{\pgfqpoint{0.100000in}{0.209042in}}{\pgfqpoint{2.906250in}{2.906250in}}%
\pgfusepath{clip}%
\pgfsetbuttcap%
\pgfsetroundjoin%
\definecolor{currentfill}{rgb}{0.544015,0.159033,0.506159}%
\pgfsetfillcolor{currentfill}%
\pgfsetfillopacity{0.800000}%
\pgfsetlinewidth{1.003750pt}%
\definecolor{currentstroke}{rgb}{0.544015,0.159033,0.506159}%
\pgfsetstrokecolor{currentstroke}%
\pgfsetstrokeopacity{0.800000}%
\pgfsetdash{}{0pt}%
\pgfpathmoveto{\pgfqpoint{1.308364in}{1.614591in}}%
\pgfpathcurveto{\pgfqpoint{1.314188in}{1.614591in}}{\pgfqpoint{1.319774in}{1.616905in}}{\pgfqpoint{1.323892in}{1.621023in}}%
\pgfpathcurveto{\pgfqpoint{1.328010in}{1.625141in}}{\pgfqpoint{1.330324in}{1.630727in}}{\pgfqpoint{1.330324in}{1.636551in}}%
\pgfpathcurveto{\pgfqpoint{1.330324in}{1.642375in}}{\pgfqpoint{1.328010in}{1.647961in}}{\pgfqpoint{1.323892in}{1.652080in}}%
\pgfpathcurveto{\pgfqpoint{1.319774in}{1.656198in}}{\pgfqpoint{1.314188in}{1.658512in}}{\pgfqpoint{1.308364in}{1.658512in}}%
\pgfpathcurveto{\pgfqpoint{1.302540in}{1.658512in}}{\pgfqpoint{1.296954in}{1.656198in}}{\pgfqpoint{1.292836in}{1.652080in}}%
\pgfpathcurveto{\pgfqpoint{1.288718in}{1.647961in}}{\pgfqpoint{1.286404in}{1.642375in}}{\pgfqpoint{1.286404in}{1.636551in}}%
\pgfpathcurveto{\pgfqpoint{1.286404in}{1.630727in}}{\pgfqpoint{1.288718in}{1.625141in}}{\pgfqpoint{1.292836in}{1.621023in}}%
\pgfpathcurveto{\pgfqpoint{1.296954in}{1.616905in}}{\pgfqpoint{1.302540in}{1.614591in}}{\pgfqpoint{1.308364in}{1.614591in}}%
\pgfpathlineto{\pgfqpoint{1.308364in}{1.614591in}}%
\pgfpathclose%
\pgfusepath{stroke,fill}%
\end{pgfscope}%
\begin{pgfscope}%
\pgfpathrectangle{\pgfqpoint{0.100000in}{0.209042in}}{\pgfqpoint{2.906250in}{2.906250in}}%
\pgfusepath{clip}%
\pgfsetbuttcap%
\pgfsetroundjoin%
\definecolor{currentfill}{rgb}{0.488088,0.139186,0.508011}%
\pgfsetfillcolor{currentfill}%
\pgfsetfillopacity{0.800000}%
\pgfsetlinewidth{1.003750pt}%
\definecolor{currentstroke}{rgb}{0.488088,0.139186,0.508011}%
\pgfsetstrokecolor{currentstroke}%
\pgfsetstrokeopacity{0.800000}%
\pgfsetdash{}{0pt}%
\pgfpathmoveto{\pgfqpoint{1.531836in}{1.968692in}}%
\pgfpathcurveto{\pgfqpoint{1.537660in}{1.968692in}}{\pgfqpoint{1.543246in}{1.971006in}}{\pgfqpoint{1.547364in}{1.975124in}}%
\pgfpathcurveto{\pgfqpoint{1.551482in}{1.979242in}}{\pgfqpoint{1.553796in}{1.984828in}}{\pgfqpoint{1.553796in}{1.990652in}}%
\pgfpathcurveto{\pgfqpoint{1.553796in}{1.996476in}}{\pgfqpoint{1.551482in}{2.002062in}}{\pgfqpoint{1.547364in}{2.006181in}}%
\pgfpathcurveto{\pgfqpoint{1.543246in}{2.010299in}}{\pgfqpoint{1.537660in}{2.012613in}}{\pgfqpoint{1.531836in}{2.012613in}}%
\pgfpathcurveto{\pgfqpoint{1.526012in}{2.012613in}}{\pgfqpoint{1.520426in}{2.010299in}}{\pgfqpoint{1.516308in}{2.006181in}}%
\pgfpathcurveto{\pgfqpoint{1.512189in}{2.002062in}}{\pgfqpoint{1.509876in}{1.996476in}}{\pgfqpoint{1.509876in}{1.990652in}}%
\pgfpathcurveto{\pgfqpoint{1.509876in}{1.984828in}}{\pgfqpoint{1.512189in}{1.979242in}}{\pgfqpoint{1.516308in}{1.975124in}}%
\pgfpathcurveto{\pgfqpoint{1.520426in}{1.971006in}}{\pgfqpoint{1.526012in}{1.968692in}}{\pgfqpoint{1.531836in}{1.968692in}}%
\pgfpathlineto{\pgfqpoint{1.531836in}{1.968692in}}%
\pgfpathclose%
\pgfusepath{stroke,fill}%
\end{pgfscope}%
\begin{pgfscope}%
\pgfpathrectangle{\pgfqpoint{0.100000in}{0.209042in}}{\pgfqpoint{2.906250in}{2.906250in}}%
\pgfusepath{clip}%
\pgfsetbuttcap%
\pgfsetroundjoin%
\definecolor{currentfill}{rgb}{0.457386,0.127522,0.507448}%
\pgfsetfillcolor{currentfill}%
\pgfsetfillopacity{0.800000}%
\pgfsetlinewidth{1.003750pt}%
\definecolor{currentstroke}{rgb}{0.457386,0.127522,0.507448}%
\pgfsetstrokecolor{currentstroke}%
\pgfsetstrokeopacity{0.800000}%
\pgfsetdash{}{0pt}%
\pgfpathmoveto{\pgfqpoint{1.841301in}{1.717218in}}%
\pgfpathcurveto{\pgfqpoint{1.847125in}{1.717218in}}{\pgfqpoint{1.852711in}{1.719532in}}{\pgfqpoint{1.856829in}{1.723650in}}%
\pgfpathcurveto{\pgfqpoint{1.860947in}{1.727768in}}{\pgfqpoint{1.863261in}{1.733354in}}{\pgfqpoint{1.863261in}{1.739178in}}%
\pgfpathcurveto{\pgfqpoint{1.863261in}{1.745002in}}{\pgfqpoint{1.860947in}{1.750588in}}{\pgfqpoint{1.856829in}{1.754706in}}%
\pgfpathcurveto{\pgfqpoint{1.852711in}{1.758824in}}{\pgfqpoint{1.847125in}{1.761138in}}{\pgfqpoint{1.841301in}{1.761138in}}%
\pgfpathcurveto{\pgfqpoint{1.835477in}{1.761138in}}{\pgfqpoint{1.829891in}{1.758824in}}{\pgfqpoint{1.825773in}{1.754706in}}%
\pgfpathcurveto{\pgfqpoint{1.821655in}{1.750588in}}{\pgfqpoint{1.819341in}{1.745002in}}{\pgfqpoint{1.819341in}{1.739178in}}%
\pgfpathcurveto{\pgfqpoint{1.819341in}{1.733354in}}{\pgfqpoint{1.821655in}{1.727768in}}{\pgfqpoint{1.825773in}{1.723650in}}%
\pgfpathcurveto{\pgfqpoint{1.829891in}{1.719532in}}{\pgfqpoint{1.835477in}{1.717218in}}{\pgfqpoint{1.841301in}{1.717218in}}%
\pgfpathlineto{\pgfqpoint{1.841301in}{1.717218in}}%
\pgfpathclose%
\pgfusepath{stroke,fill}%
\end{pgfscope}%
\begin{pgfscope}%
\pgfpathrectangle{\pgfqpoint{0.100000in}{0.209042in}}{\pgfqpoint{2.906250in}{2.906250in}}%
\pgfusepath{clip}%
\pgfsetbuttcap%
\pgfsetroundjoin%
\definecolor{currentfill}{rgb}{0.420791,0.112920,0.505215}%
\pgfsetfillcolor{currentfill}%
\pgfsetfillopacity{0.800000}%
\pgfsetlinewidth{1.003750pt}%
\definecolor{currentstroke}{rgb}{0.420791,0.112920,0.505215}%
\pgfsetstrokecolor{currentstroke}%
\pgfsetstrokeopacity{0.800000}%
\pgfsetdash{}{0pt}%
\pgfpathmoveto{\pgfqpoint{1.453490in}{1.887512in}}%
\pgfpathcurveto{\pgfqpoint{1.459314in}{1.887512in}}{\pgfqpoint{1.464900in}{1.889826in}}{\pgfqpoint{1.469018in}{1.893944in}}%
\pgfpathcurveto{\pgfqpoint{1.473136in}{1.898062in}}{\pgfqpoint{1.475450in}{1.903649in}}{\pgfqpoint{1.475450in}{1.909473in}}%
\pgfpathcurveto{\pgfqpoint{1.475450in}{1.915297in}}{\pgfqpoint{1.473136in}{1.920883in}}{\pgfqpoint{1.469018in}{1.925001in}}%
\pgfpathcurveto{\pgfqpoint{1.464900in}{1.929119in}}{\pgfqpoint{1.459314in}{1.931433in}}{\pgfqpoint{1.453490in}{1.931433in}}%
\pgfpathcurveto{\pgfqpoint{1.447666in}{1.931433in}}{\pgfqpoint{1.442080in}{1.929119in}}{\pgfqpoint{1.437961in}{1.925001in}}%
\pgfpathcurveto{\pgfqpoint{1.433843in}{1.920883in}}{\pgfqpoint{1.431529in}{1.915297in}}{\pgfqpoint{1.431529in}{1.909473in}}%
\pgfpathcurveto{\pgfqpoint{1.431529in}{1.903649in}}{\pgfqpoint{1.433843in}{1.898062in}}{\pgfqpoint{1.437961in}{1.893944in}}%
\pgfpathcurveto{\pgfqpoint{1.442080in}{1.889826in}}{\pgfqpoint{1.447666in}{1.887512in}}{\pgfqpoint{1.453490in}{1.887512in}}%
\pgfpathlineto{\pgfqpoint{1.453490in}{1.887512in}}%
\pgfpathclose%
\pgfusepath{stroke,fill}%
\end{pgfscope}%
\begin{pgfscope}%
\pgfpathrectangle{\pgfqpoint{0.100000in}{0.209042in}}{\pgfqpoint{2.906250in}{2.906250in}}%
\pgfusepath{clip}%
\pgfsetbuttcap%
\pgfsetroundjoin%
\definecolor{currentfill}{rgb}{0.457386,0.127522,0.507448}%
\pgfsetfillcolor{currentfill}%
\pgfsetfillopacity{0.800000}%
\pgfsetlinewidth{1.003750pt}%
\definecolor{currentstroke}{rgb}{0.457386,0.127522,0.507448}%
\pgfsetstrokecolor{currentstroke}%
\pgfsetstrokeopacity{0.800000}%
\pgfsetdash{}{0pt}%
\pgfpathmoveto{\pgfqpoint{1.790407in}{1.604278in}}%
\pgfpathcurveto{\pgfqpoint{1.796230in}{1.604278in}}{\pgfqpoint{1.801817in}{1.606592in}}{\pgfqpoint{1.805935in}{1.610710in}}%
\pgfpathcurveto{\pgfqpoint{1.810053in}{1.614828in}}{\pgfqpoint{1.812367in}{1.620415in}}{\pgfqpoint{1.812367in}{1.626239in}}%
\pgfpathcurveto{\pgfqpoint{1.812367in}{1.632062in}}{\pgfqpoint{1.810053in}{1.637649in}}{\pgfqpoint{1.805935in}{1.641767in}}%
\pgfpathcurveto{\pgfqpoint{1.801817in}{1.645885in}}{\pgfqpoint{1.796230in}{1.648199in}}{\pgfqpoint{1.790407in}{1.648199in}}%
\pgfpathcurveto{\pgfqpoint{1.784583in}{1.648199in}}{\pgfqpoint{1.778996in}{1.645885in}}{\pgfqpoint{1.774878in}{1.641767in}}%
\pgfpathcurveto{\pgfqpoint{1.770760in}{1.637649in}}{\pgfqpoint{1.768446in}{1.632062in}}{\pgfqpoint{1.768446in}{1.626239in}}%
\pgfpathcurveto{\pgfqpoint{1.768446in}{1.620415in}}{\pgfqpoint{1.770760in}{1.614828in}}{\pgfqpoint{1.774878in}{1.610710in}}%
\pgfpathcurveto{\pgfqpoint{1.778996in}{1.606592in}}{\pgfqpoint{1.784583in}{1.604278in}}{\pgfqpoint{1.790407in}{1.604278in}}%
\pgfpathlineto{\pgfqpoint{1.790407in}{1.604278in}}%
\pgfpathclose%
\pgfusepath{stroke,fill}%
\end{pgfscope}%
\begin{pgfscope}%
\pgfpathrectangle{\pgfqpoint{0.100000in}{0.209042in}}{\pgfqpoint{2.906250in}{2.906250in}}%
\pgfusepath{clip}%
\pgfsetbuttcap%
\pgfsetroundjoin%
\definecolor{currentfill}{rgb}{0.463508,0.129893,0.507652}%
\pgfsetfillcolor{currentfill}%
\pgfsetfillopacity{0.800000}%
\pgfsetlinewidth{1.003750pt}%
\definecolor{currentstroke}{rgb}{0.463508,0.129893,0.507652}%
\pgfsetstrokecolor{currentstroke}%
\pgfsetstrokeopacity{0.800000}%
\pgfsetdash{}{0pt}%
\pgfpathmoveto{\pgfqpoint{1.327257in}{1.797064in}}%
\pgfpathcurveto{\pgfqpoint{1.333081in}{1.797064in}}{\pgfqpoint{1.338667in}{1.799378in}}{\pgfqpoint{1.342785in}{1.803496in}}%
\pgfpathcurveto{\pgfqpoint{1.346903in}{1.807614in}}{\pgfqpoint{1.349217in}{1.813200in}}{\pgfqpoint{1.349217in}{1.819024in}}%
\pgfpathcurveto{\pgfqpoint{1.349217in}{1.824848in}}{\pgfqpoint{1.346903in}{1.830435in}}{\pgfqpoint{1.342785in}{1.834553in}}%
\pgfpathcurveto{\pgfqpoint{1.338667in}{1.838671in}}{\pgfqpoint{1.333081in}{1.840985in}}{\pgfqpoint{1.327257in}{1.840985in}}%
\pgfpathcurveto{\pgfqpoint{1.321433in}{1.840985in}}{\pgfqpoint{1.315847in}{1.838671in}}{\pgfqpoint{1.311728in}{1.834553in}}%
\pgfpathcurveto{\pgfqpoint{1.307610in}{1.830435in}}{\pgfqpoint{1.305296in}{1.824848in}}{\pgfqpoint{1.305296in}{1.819024in}}%
\pgfpathcurveto{\pgfqpoint{1.305296in}{1.813200in}}{\pgfqpoint{1.307610in}{1.807614in}}{\pgfqpoint{1.311728in}{1.803496in}}%
\pgfpathcurveto{\pgfqpoint{1.315847in}{1.799378in}}{\pgfqpoint{1.321433in}{1.797064in}}{\pgfqpoint{1.327257in}{1.797064in}}%
\pgfpathlineto{\pgfqpoint{1.327257in}{1.797064in}}%
\pgfpathclose%
\pgfusepath{stroke,fill}%
\end{pgfscope}%
\begin{pgfscope}%
\pgfpathrectangle{\pgfqpoint{0.100000in}{0.209042in}}{\pgfqpoint{2.906250in}{2.906250in}}%
\pgfusepath{clip}%
\pgfsetbuttcap%
\pgfsetroundjoin%
\definecolor{currentfill}{rgb}{0.933606,0.358764,0.370541}%
\pgfsetfillcolor{currentfill}%
\pgfsetfillopacity{0.800000}%
\pgfsetlinewidth{1.003750pt}%
\definecolor{currentstroke}{rgb}{0.933606,0.358764,0.370541}%
\pgfsetstrokecolor{currentstroke}%
\pgfsetstrokeopacity{0.800000}%
\pgfsetdash{}{0pt}%
\pgfpathmoveto{\pgfqpoint{1.709582in}{2.257130in}}%
\pgfpathcurveto{\pgfqpoint{1.715406in}{2.257130in}}{\pgfqpoint{1.720992in}{2.259444in}}{\pgfqpoint{1.725111in}{2.263562in}}%
\pgfpathcurveto{\pgfqpoint{1.729229in}{2.267680in}}{\pgfqpoint{1.731543in}{2.273266in}}{\pgfqpoint{1.731543in}{2.279090in}}%
\pgfpathcurveto{\pgfqpoint{1.731543in}{2.284914in}}{\pgfqpoint{1.729229in}{2.290500in}}{\pgfqpoint{1.725111in}{2.294618in}}%
\pgfpathcurveto{\pgfqpoint{1.720992in}{2.298736in}}{\pgfqpoint{1.715406in}{2.301050in}}{\pgfqpoint{1.709582in}{2.301050in}}%
\pgfpathcurveto{\pgfqpoint{1.703758in}{2.301050in}}{\pgfqpoint{1.698172in}{2.298736in}}{\pgfqpoint{1.694054in}{2.294618in}}%
\pgfpathcurveto{\pgfqpoint{1.689936in}{2.290500in}}{\pgfqpoint{1.687622in}{2.284914in}}{\pgfqpoint{1.687622in}{2.279090in}}%
\pgfpathcurveto{\pgfqpoint{1.687622in}{2.273266in}}{\pgfqpoint{1.689936in}{2.267680in}}{\pgfqpoint{1.694054in}{2.263562in}}%
\pgfpathcurveto{\pgfqpoint{1.698172in}{2.259444in}}{\pgfqpoint{1.703758in}{2.257130in}}{\pgfqpoint{1.709582in}{2.257130in}}%
\pgfpathlineto{\pgfqpoint{1.709582in}{2.257130in}}%
\pgfpathclose%
\pgfusepath{stroke,fill}%
\end{pgfscope}%
\begin{pgfscope}%
\pgfpathrectangle{\pgfqpoint{0.100000in}{0.209042in}}{\pgfqpoint{2.906250in}{2.906250in}}%
\pgfusepath{clip}%
\pgfsetbuttcap%
\pgfsetroundjoin%
\definecolor{currentfill}{rgb}{0.645633,0.191952,0.492910}%
\pgfsetfillcolor{currentfill}%
\pgfsetfillopacity{0.800000}%
\pgfsetlinewidth{1.003750pt}%
\definecolor{currentstroke}{rgb}{0.645633,0.191952,0.492910}%
\pgfsetstrokecolor{currentstroke}%
\pgfsetstrokeopacity{0.800000}%
\pgfsetdash{}{0pt}%
\pgfpathmoveto{\pgfqpoint{1.796526in}{1.431931in}}%
\pgfpathcurveto{\pgfqpoint{1.802350in}{1.431931in}}{\pgfqpoint{1.807936in}{1.434245in}}{\pgfqpoint{1.812054in}{1.438363in}}%
\pgfpathcurveto{\pgfqpoint{1.816173in}{1.442481in}}{\pgfqpoint{1.818486in}{1.448067in}}{\pgfqpoint{1.818486in}{1.453891in}}%
\pgfpathcurveto{\pgfqpoint{1.818486in}{1.459715in}}{\pgfqpoint{1.816173in}{1.465301in}}{\pgfqpoint{1.812054in}{1.469419in}}%
\pgfpathcurveto{\pgfqpoint{1.807936in}{1.473537in}}{\pgfqpoint{1.802350in}{1.475851in}}{\pgfqpoint{1.796526in}{1.475851in}}%
\pgfpathcurveto{\pgfqpoint{1.790702in}{1.475851in}}{\pgfqpoint{1.785116in}{1.473537in}}{\pgfqpoint{1.780998in}{1.469419in}}%
\pgfpathcurveto{\pgfqpoint{1.776880in}{1.465301in}}{\pgfqpoint{1.774566in}{1.459715in}}{\pgfqpoint{1.774566in}{1.453891in}}%
\pgfpathcurveto{\pgfqpoint{1.774566in}{1.448067in}}{\pgfqpoint{1.776880in}{1.442481in}}{\pgfqpoint{1.780998in}{1.438363in}}%
\pgfpathcurveto{\pgfqpoint{1.785116in}{1.434245in}}{\pgfqpoint{1.790702in}{1.431931in}}{\pgfqpoint{1.796526in}{1.431931in}}%
\pgfpathlineto{\pgfqpoint{1.796526in}{1.431931in}}%
\pgfpathclose%
\pgfusepath{stroke,fill}%
\end{pgfscope}%
\begin{pgfscope}%
\pgfpathrectangle{\pgfqpoint{0.100000in}{0.209042in}}{\pgfqpoint{2.906250in}{2.906250in}}%
\pgfusepath{clip}%
\pgfsetbuttcap%
\pgfsetroundjoin%
\definecolor{currentfill}{rgb}{0.506629,0.145958,0.507806}%
\pgfsetfillcolor{currentfill}%
\pgfsetfillopacity{0.800000}%
\pgfsetlinewidth{1.003750pt}%
\definecolor{currentstroke}{rgb}{0.506629,0.145958,0.507806}%
\pgfsetstrokecolor{currentstroke}%
\pgfsetstrokeopacity{0.800000}%
\pgfsetdash{}{0pt}%
\pgfpathmoveto{\pgfqpoint{1.292176in}{1.831127in}}%
\pgfpathcurveto{\pgfqpoint{1.298000in}{1.831127in}}{\pgfqpoint{1.303586in}{1.833441in}}{\pgfqpoint{1.307704in}{1.837559in}}%
\pgfpathcurveto{\pgfqpoint{1.311822in}{1.841677in}}{\pgfqpoint{1.314136in}{1.847263in}}{\pgfqpoint{1.314136in}{1.853087in}}%
\pgfpathcurveto{\pgfqpoint{1.314136in}{1.858911in}}{\pgfqpoint{1.311822in}{1.864497in}}{\pgfqpoint{1.307704in}{1.868615in}}%
\pgfpathcurveto{\pgfqpoint{1.303586in}{1.872734in}}{\pgfqpoint{1.298000in}{1.875047in}}{\pgfqpoint{1.292176in}{1.875047in}}%
\pgfpathcurveto{\pgfqpoint{1.286352in}{1.875047in}}{\pgfqpoint{1.280765in}{1.872734in}}{\pgfqpoint{1.276647in}{1.868615in}}%
\pgfpathcurveto{\pgfqpoint{1.272529in}{1.864497in}}{\pgfqpoint{1.270215in}{1.858911in}}{\pgfqpoint{1.270215in}{1.853087in}}%
\pgfpathcurveto{\pgfqpoint{1.270215in}{1.847263in}}{\pgfqpoint{1.272529in}{1.841677in}}{\pgfqpoint{1.276647in}{1.837559in}}%
\pgfpathcurveto{\pgfqpoint{1.280765in}{1.833441in}}{\pgfqpoint{1.286352in}{1.831127in}}{\pgfqpoint{1.292176in}{1.831127in}}%
\pgfpathlineto{\pgfqpoint{1.292176in}{1.831127in}}%
\pgfpathclose%
\pgfusepath{stroke,fill}%
\end{pgfscope}%
\begin{pgfscope}%
\pgfpathrectangle{\pgfqpoint{0.100000in}{0.209042in}}{\pgfqpoint{2.906250in}{2.906250in}}%
\pgfusepath{clip}%
\pgfsetbuttcap%
\pgfsetroundjoin%
\definecolor{currentfill}{rgb}{0.445163,0.122724,0.506901}%
\pgfsetfillcolor{currentfill}%
\pgfsetfillopacity{0.800000}%
\pgfsetlinewidth{1.003750pt}%
\definecolor{currentstroke}{rgb}{0.445163,0.122724,0.506901}%
\pgfsetstrokecolor{currentstroke}%
\pgfsetstrokeopacity{0.800000}%
\pgfsetdash{}{0pt}%
\pgfpathmoveto{\pgfqpoint{1.804523in}{1.852433in}}%
\pgfpathcurveto{\pgfqpoint{1.810347in}{1.852433in}}{\pgfqpoint{1.815933in}{1.854747in}}{\pgfqpoint{1.820052in}{1.858865in}}%
\pgfpathcurveto{\pgfqpoint{1.824170in}{1.862983in}}{\pgfqpoint{1.826484in}{1.868569in}}{\pgfqpoint{1.826484in}{1.874393in}}%
\pgfpathcurveto{\pgfqpoint{1.826484in}{1.880217in}}{\pgfqpoint{1.824170in}{1.885803in}}{\pgfqpoint{1.820052in}{1.889921in}}%
\pgfpathcurveto{\pgfqpoint{1.815933in}{1.894039in}}{\pgfqpoint{1.810347in}{1.896353in}}{\pgfqpoint{1.804523in}{1.896353in}}%
\pgfpathcurveto{\pgfqpoint{1.798699in}{1.896353in}}{\pgfqpoint{1.793113in}{1.894039in}}{\pgfqpoint{1.788995in}{1.889921in}}%
\pgfpathcurveto{\pgfqpoint{1.784877in}{1.885803in}}{\pgfqpoint{1.782563in}{1.880217in}}{\pgfqpoint{1.782563in}{1.874393in}}%
\pgfpathcurveto{\pgfqpoint{1.782563in}{1.868569in}}{\pgfqpoint{1.784877in}{1.862983in}}{\pgfqpoint{1.788995in}{1.858865in}}%
\pgfpathcurveto{\pgfqpoint{1.793113in}{1.854747in}}{\pgfqpoint{1.798699in}{1.852433in}}{\pgfqpoint{1.804523in}{1.852433in}}%
\pgfpathlineto{\pgfqpoint{1.804523in}{1.852433in}}%
\pgfpathclose%
\pgfusepath{stroke,fill}%
\end{pgfscope}%
\begin{pgfscope}%
\pgfpathrectangle{\pgfqpoint{0.100000in}{0.209042in}}{\pgfqpoint{2.906250in}{2.906250in}}%
\pgfusepath{clip}%
\pgfsetbuttcap%
\pgfsetroundjoin%
\definecolor{currentfill}{rgb}{0.512831,0.148179,0.507648}%
\pgfsetfillcolor{currentfill}%
\pgfsetfillopacity{0.800000}%
\pgfsetlinewidth{1.003750pt}%
\definecolor{currentstroke}{rgb}{0.512831,0.148179,0.507648}%
\pgfsetstrokecolor{currentstroke}%
\pgfsetstrokeopacity{0.800000}%
\pgfsetdash{}{0pt}%
\pgfpathmoveto{\pgfqpoint{1.276125in}{1.709461in}}%
\pgfpathcurveto{\pgfqpoint{1.281949in}{1.709461in}}{\pgfqpoint{1.287535in}{1.711775in}}{\pgfqpoint{1.291654in}{1.715894in}}%
\pgfpathcurveto{\pgfqpoint{1.295772in}{1.720012in}}{\pgfqpoint{1.298086in}{1.725598in}}{\pgfqpoint{1.298086in}{1.731422in}}%
\pgfpathcurveto{\pgfqpoint{1.298086in}{1.737246in}}{\pgfqpoint{1.295772in}{1.742832in}}{\pgfqpoint{1.291654in}{1.746950in}}%
\pgfpathcurveto{\pgfqpoint{1.287535in}{1.751068in}}{\pgfqpoint{1.281949in}{1.753382in}}{\pgfqpoint{1.276125in}{1.753382in}}%
\pgfpathcurveto{\pgfqpoint{1.270301in}{1.753382in}}{\pgfqpoint{1.264715in}{1.751068in}}{\pgfqpoint{1.260597in}{1.746950in}}%
\pgfpathcurveto{\pgfqpoint{1.256479in}{1.742832in}}{\pgfqpoint{1.254165in}{1.737246in}}{\pgfqpoint{1.254165in}{1.731422in}}%
\pgfpathcurveto{\pgfqpoint{1.254165in}{1.725598in}}{\pgfqpoint{1.256479in}{1.720012in}}{\pgfqpoint{1.260597in}{1.715894in}}%
\pgfpathcurveto{\pgfqpoint{1.264715in}{1.711775in}}{\pgfqpoint{1.270301in}{1.709461in}}{\pgfqpoint{1.276125in}{1.709461in}}%
\pgfpathlineto{\pgfqpoint{1.276125in}{1.709461in}}%
\pgfpathclose%
\pgfusepath{stroke,fill}%
\end{pgfscope}%
\begin{pgfscope}%
\pgfpathrectangle{\pgfqpoint{0.100000in}{0.209042in}}{\pgfqpoint{2.906250in}{2.906250in}}%
\pgfusepath{clip}%
\pgfsetbuttcap%
\pgfsetroundjoin%
\definecolor{currentfill}{rgb}{0.973381,0.461520,0.361965}%
\pgfsetfillcolor{currentfill}%
\pgfsetfillopacity{0.800000}%
\pgfsetlinewidth{1.003750pt}%
\definecolor{currentstroke}{rgb}{0.973381,0.461520,0.361965}%
\pgfsetstrokecolor{currentstroke}%
\pgfsetstrokeopacity{0.800000}%
\pgfsetdash{}{0pt}%
\pgfpathmoveto{\pgfqpoint{1.269631in}{2.280105in}}%
\pgfpathcurveto{\pgfqpoint{1.275455in}{2.280105in}}{\pgfqpoint{1.281041in}{2.282418in}}{\pgfqpoint{1.285159in}{2.286537in}}%
\pgfpathcurveto{\pgfqpoint{1.289277in}{2.290655in}}{\pgfqpoint{1.291591in}{2.296241in}}{\pgfqpoint{1.291591in}{2.302065in}}%
\pgfpathcurveto{\pgfqpoint{1.291591in}{2.307889in}}{\pgfqpoint{1.289277in}{2.313475in}}{\pgfqpoint{1.285159in}{2.317593in}}%
\pgfpathcurveto{\pgfqpoint{1.281041in}{2.321711in}}{\pgfqpoint{1.275455in}{2.324025in}}{\pgfqpoint{1.269631in}{2.324025in}}%
\pgfpathcurveto{\pgfqpoint{1.263807in}{2.324025in}}{\pgfqpoint{1.258221in}{2.321711in}}{\pgfqpoint{1.254103in}{2.317593in}}%
\pgfpathcurveto{\pgfqpoint{1.249985in}{2.313475in}}{\pgfqpoint{1.247671in}{2.307889in}}{\pgfqpoint{1.247671in}{2.302065in}}%
\pgfpathcurveto{\pgfqpoint{1.247671in}{2.296241in}}{\pgfqpoint{1.249985in}{2.290655in}}{\pgfqpoint{1.254103in}{2.286537in}}%
\pgfpathcurveto{\pgfqpoint{1.258221in}{2.282418in}}{\pgfqpoint{1.263807in}{2.280105in}}{\pgfqpoint{1.269631in}{2.280105in}}%
\pgfpathlineto{\pgfqpoint{1.269631in}{2.280105in}}%
\pgfpathclose%
\pgfusepath{stroke,fill}%
\end{pgfscope}%
\begin{pgfscope}%
\pgfpathrectangle{\pgfqpoint{0.100000in}{0.209042in}}{\pgfqpoint{2.906250in}{2.906250in}}%
\pgfusepath{clip}%
\pgfsetbuttcap%
\pgfsetroundjoin%
\definecolor{currentfill}{rgb}{0.408629,0.107930,0.504052}%
\pgfsetfillcolor{currentfill}%
\pgfsetfillopacity{0.800000}%
\pgfsetlinewidth{1.003750pt}%
\definecolor{currentstroke}{rgb}{0.408629,0.107930,0.504052}%
\pgfsetstrokecolor{currentstroke}%
\pgfsetstrokeopacity{0.800000}%
\pgfsetdash{}{0pt}%
\pgfpathmoveto{\pgfqpoint{1.364458in}{1.764037in}}%
\pgfpathcurveto{\pgfqpoint{1.370282in}{1.764037in}}{\pgfqpoint{1.375868in}{1.766351in}}{\pgfqpoint{1.379987in}{1.770469in}}%
\pgfpathcurveto{\pgfqpoint{1.384105in}{1.774587in}}{\pgfqpoint{1.386419in}{1.780173in}}{\pgfqpoint{1.386419in}{1.785997in}}%
\pgfpathcurveto{\pgfqpoint{1.386419in}{1.791821in}}{\pgfqpoint{1.384105in}{1.797408in}}{\pgfqpoint{1.379987in}{1.801526in}}%
\pgfpathcurveto{\pgfqpoint{1.375868in}{1.805644in}}{\pgfqpoint{1.370282in}{1.807958in}}{\pgfqpoint{1.364458in}{1.807958in}}%
\pgfpathcurveto{\pgfqpoint{1.358634in}{1.807958in}}{\pgfqpoint{1.353048in}{1.805644in}}{\pgfqpoint{1.348930in}{1.801526in}}%
\pgfpathcurveto{\pgfqpoint{1.344812in}{1.797408in}}{\pgfqpoint{1.342498in}{1.791821in}}{\pgfqpoint{1.342498in}{1.785997in}}%
\pgfpathcurveto{\pgfqpoint{1.342498in}{1.780173in}}{\pgfqpoint{1.344812in}{1.774587in}}{\pgfqpoint{1.348930in}{1.770469in}}%
\pgfpathcurveto{\pgfqpoint{1.353048in}{1.766351in}}{\pgfqpoint{1.358634in}{1.764037in}}{\pgfqpoint{1.364458in}{1.764037in}}%
\pgfpathlineto{\pgfqpoint{1.364458in}{1.764037in}}%
\pgfpathclose%
\pgfusepath{stroke,fill}%
\end{pgfscope}%
\begin{pgfscope}%
\pgfpathrectangle{\pgfqpoint{0.100000in}{0.209042in}}{\pgfqpoint{2.906250in}{2.906250in}}%
\pgfusepath{clip}%
\pgfsetbuttcap%
\pgfsetroundjoin%
\definecolor{currentfill}{rgb}{0.291366,0.064553,0.475462}%
\pgfsetfillcolor{currentfill}%
\pgfsetfillopacity{0.800000}%
\pgfsetlinewidth{1.003750pt}%
\definecolor{currentstroke}{rgb}{0.291366,0.064553,0.475462}%
\pgfsetstrokecolor{currentstroke}%
\pgfsetstrokeopacity{0.800000}%
\pgfsetdash{}{0pt}%
\pgfpathmoveto{\pgfqpoint{1.637706in}{1.760885in}}%
\pgfpathcurveto{\pgfqpoint{1.643530in}{1.760885in}}{\pgfqpoint{1.649117in}{1.763199in}}{\pgfqpoint{1.653235in}{1.767317in}}%
\pgfpathcurveto{\pgfqpoint{1.657353in}{1.771435in}}{\pgfqpoint{1.659667in}{1.777022in}}{\pgfqpoint{1.659667in}{1.782846in}}%
\pgfpathcurveto{\pgfqpoint{1.659667in}{1.788669in}}{\pgfqpoint{1.657353in}{1.794256in}}{\pgfqpoint{1.653235in}{1.798374in}}%
\pgfpathcurveto{\pgfqpoint{1.649117in}{1.802492in}}{\pgfqpoint{1.643530in}{1.804806in}}{\pgfqpoint{1.637706in}{1.804806in}}%
\pgfpathcurveto{\pgfqpoint{1.631882in}{1.804806in}}{\pgfqpoint{1.626296in}{1.802492in}}{\pgfqpoint{1.622178in}{1.798374in}}%
\pgfpathcurveto{\pgfqpoint{1.618060in}{1.794256in}}{\pgfqpoint{1.615746in}{1.788669in}}{\pgfqpoint{1.615746in}{1.782846in}}%
\pgfpathcurveto{\pgfqpoint{1.615746in}{1.777022in}}{\pgfqpoint{1.618060in}{1.771435in}}{\pgfqpoint{1.622178in}{1.767317in}}%
\pgfpathcurveto{\pgfqpoint{1.626296in}{1.763199in}}{\pgfqpoint{1.631882in}{1.760885in}}{\pgfqpoint{1.637706in}{1.760885in}}%
\pgfpathlineto{\pgfqpoint{1.637706in}{1.760885in}}%
\pgfpathclose%
\pgfusepath{stroke,fill}%
\end{pgfscope}%
\begin{pgfscope}%
\pgfpathrectangle{\pgfqpoint{0.100000in}{0.209042in}}{\pgfqpoint{2.906250in}{2.906250in}}%
\pgfusepath{clip}%
\pgfsetbuttcap%
\pgfsetroundjoin%
\definecolor{currentfill}{rgb}{0.457386,0.127522,0.507448}%
\pgfsetfillcolor{currentfill}%
\pgfsetfillopacity{0.800000}%
\pgfsetlinewidth{1.003750pt}%
\definecolor{currentstroke}{rgb}{0.457386,0.127522,0.507448}%
\pgfsetstrokecolor{currentstroke}%
\pgfsetstrokeopacity{0.800000}%
\pgfsetdash{}{0pt}%
\pgfpathmoveto{\pgfqpoint{1.872012in}{1.716606in}}%
\pgfpathcurveto{\pgfqpoint{1.877836in}{1.716606in}}{\pgfqpoint{1.883422in}{1.718920in}}{\pgfqpoint{1.887540in}{1.723038in}}%
\pgfpathcurveto{\pgfqpoint{1.891658in}{1.727156in}}{\pgfqpoint{1.893972in}{1.732742in}}{\pgfqpoint{1.893972in}{1.738566in}}%
\pgfpathcurveto{\pgfqpoint{1.893972in}{1.744390in}}{\pgfqpoint{1.891658in}{1.749976in}}{\pgfqpoint{1.887540in}{1.754095in}}%
\pgfpathcurveto{\pgfqpoint{1.883422in}{1.758213in}}{\pgfqpoint{1.877836in}{1.760527in}}{\pgfqpoint{1.872012in}{1.760527in}}%
\pgfpathcurveto{\pgfqpoint{1.866188in}{1.760527in}}{\pgfqpoint{1.860602in}{1.758213in}}{\pgfqpoint{1.856484in}{1.754095in}}%
\pgfpathcurveto{\pgfqpoint{1.852366in}{1.749976in}}{\pgfqpoint{1.850052in}{1.744390in}}{\pgfqpoint{1.850052in}{1.738566in}}%
\pgfpathcurveto{\pgfqpoint{1.850052in}{1.732742in}}{\pgfqpoint{1.852366in}{1.727156in}}{\pgfqpoint{1.856484in}{1.723038in}}%
\pgfpathcurveto{\pgfqpoint{1.860602in}{1.718920in}}{\pgfqpoint{1.866188in}{1.716606in}}{\pgfqpoint{1.872012in}{1.716606in}}%
\pgfpathlineto{\pgfqpoint{1.872012in}{1.716606in}}%
\pgfpathclose%
\pgfusepath{stroke,fill}%
\end{pgfscope}%
\begin{pgfscope}%
\pgfpathrectangle{\pgfqpoint{0.100000in}{0.209042in}}{\pgfqpoint{2.906250in}{2.906250in}}%
\pgfusepath{clip}%
\pgfsetbuttcap%
\pgfsetroundjoin%
\definecolor{currentfill}{rgb}{0.960949,0.418323,0.359630}%
\pgfsetfillcolor{currentfill}%
\pgfsetfillopacity{0.800000}%
\pgfsetlinewidth{1.003750pt}%
\definecolor{currentstroke}{rgb}{0.960949,0.418323,0.359630}%
\pgfsetstrokecolor{currentstroke}%
\pgfsetstrokeopacity{0.800000}%
\pgfsetdash{}{0pt}%
\pgfpathmoveto{\pgfqpoint{1.668348in}{1.173981in}}%
\pgfpathcurveto{\pgfqpoint{1.674172in}{1.173981in}}{\pgfqpoint{1.679758in}{1.176294in}}{\pgfqpoint{1.683876in}{1.180413in}}%
\pgfpathcurveto{\pgfqpoint{1.687994in}{1.184531in}}{\pgfqpoint{1.690308in}{1.190117in}}{\pgfqpoint{1.690308in}{1.195941in}}%
\pgfpathcurveto{\pgfqpoint{1.690308in}{1.201765in}}{\pgfqpoint{1.687994in}{1.207351in}}{\pgfqpoint{1.683876in}{1.211469in}}%
\pgfpathcurveto{\pgfqpoint{1.679758in}{1.215587in}}{\pgfqpoint{1.674172in}{1.217901in}}{\pgfqpoint{1.668348in}{1.217901in}}%
\pgfpathcurveto{\pgfqpoint{1.662524in}{1.217901in}}{\pgfqpoint{1.656938in}{1.215587in}}{\pgfqpoint{1.652820in}{1.211469in}}%
\pgfpathcurveto{\pgfqpoint{1.648702in}{1.207351in}}{\pgfqpoint{1.646388in}{1.201765in}}{\pgfqpoint{1.646388in}{1.195941in}}%
\pgfpathcurveto{\pgfqpoint{1.646388in}{1.190117in}}{\pgfqpoint{1.648702in}{1.184531in}}{\pgfqpoint{1.652820in}{1.180413in}}%
\pgfpathcurveto{\pgfqpoint{1.656938in}{1.176294in}}{\pgfqpoint{1.662524in}{1.173981in}}{\pgfqpoint{1.668348in}{1.173981in}}%
\pgfpathlineto{\pgfqpoint{1.668348in}{1.173981in}}%
\pgfpathclose%
\pgfusepath{stroke,fill}%
\end{pgfscope}%
\begin{pgfscope}%
\pgfpathrectangle{\pgfqpoint{0.100000in}{0.209042in}}{\pgfqpoint{2.906250in}{2.906250in}}%
\pgfusepath{clip}%
\pgfsetbuttcap%
\pgfsetroundjoin%
\definecolor{currentfill}{rgb}{0.475780,0.134577,0.507921}%
\pgfsetfillcolor{currentfill}%
\pgfsetfillopacity{0.800000}%
\pgfsetlinewidth{1.003750pt}%
\definecolor{currentstroke}{rgb}{0.475780,0.134577,0.507921}%
\pgfsetstrokecolor{currentstroke}%
\pgfsetstrokeopacity{0.800000}%
\pgfsetdash{}{0pt}%
\pgfpathmoveto{\pgfqpoint{1.888642in}{1.678240in}}%
\pgfpathcurveto{\pgfqpoint{1.894466in}{1.678240in}}{\pgfqpoint{1.900052in}{1.680554in}}{\pgfqpoint{1.904170in}{1.684672in}}%
\pgfpathcurveto{\pgfqpoint{1.908289in}{1.688790in}}{\pgfqpoint{1.910602in}{1.694376in}}{\pgfqpoint{1.910602in}{1.700200in}}%
\pgfpathcurveto{\pgfqpoint{1.910602in}{1.706024in}}{\pgfqpoint{1.908289in}{1.711610in}}{\pgfqpoint{1.904170in}{1.715728in}}%
\pgfpathcurveto{\pgfqpoint{1.900052in}{1.719847in}}{\pgfqpoint{1.894466in}{1.722160in}}{\pgfqpoint{1.888642in}{1.722160in}}%
\pgfpathcurveto{\pgfqpoint{1.882818in}{1.722160in}}{\pgfqpoint{1.877232in}{1.719847in}}{\pgfqpoint{1.873114in}{1.715728in}}%
\pgfpathcurveto{\pgfqpoint{1.868996in}{1.711610in}}{\pgfqpoint{1.866682in}{1.706024in}}{\pgfqpoint{1.866682in}{1.700200in}}%
\pgfpathcurveto{\pgfqpoint{1.866682in}{1.694376in}}{\pgfqpoint{1.868996in}{1.688790in}}{\pgfqpoint{1.873114in}{1.684672in}}%
\pgfpathcurveto{\pgfqpoint{1.877232in}{1.680554in}}{\pgfqpoint{1.882818in}{1.678240in}}{\pgfqpoint{1.888642in}{1.678240in}}%
\pgfpathlineto{\pgfqpoint{1.888642in}{1.678240in}}%
\pgfpathclose%
\pgfusepath{stroke,fill}%
\end{pgfscope}%
\begin{pgfscope}%
\pgfpathrectangle{\pgfqpoint{0.100000in}{0.209042in}}{\pgfqpoint{2.906250in}{2.906250in}}%
\pgfusepath{clip}%
\pgfsetbuttcap%
\pgfsetroundjoin%
\definecolor{currentfill}{rgb}{0.304081,0.067835,0.480812}%
\pgfsetfillcolor{currentfill}%
\pgfsetfillopacity{0.800000}%
\pgfsetlinewidth{1.003750pt}%
\definecolor{currentstroke}{rgb}{0.304081,0.067835,0.480812}%
\pgfsetstrokecolor{currentstroke}%
\pgfsetstrokeopacity{0.800000}%
\pgfsetdash{}{0pt}%
\pgfpathmoveto{\pgfqpoint{1.657073in}{1.807452in}}%
\pgfpathcurveto{\pgfqpoint{1.662897in}{1.807452in}}{\pgfqpoint{1.668483in}{1.809766in}}{\pgfqpoint{1.672601in}{1.813884in}}%
\pgfpathcurveto{\pgfqpoint{1.676719in}{1.818002in}}{\pgfqpoint{1.679033in}{1.823588in}}{\pgfqpoint{1.679033in}{1.829412in}}%
\pgfpathcurveto{\pgfqpoint{1.679033in}{1.835236in}}{\pgfqpoint{1.676719in}{1.840822in}}{\pgfqpoint{1.672601in}{1.844941in}}%
\pgfpathcurveto{\pgfqpoint{1.668483in}{1.849059in}}{\pgfqpoint{1.662897in}{1.851373in}}{\pgfqpoint{1.657073in}{1.851373in}}%
\pgfpathcurveto{\pgfqpoint{1.651249in}{1.851373in}}{\pgfqpoint{1.645663in}{1.849059in}}{\pgfqpoint{1.641545in}{1.844941in}}%
\pgfpathcurveto{\pgfqpoint{1.637426in}{1.840822in}}{\pgfqpoint{1.635113in}{1.835236in}}{\pgfqpoint{1.635113in}{1.829412in}}%
\pgfpathcurveto{\pgfqpoint{1.635113in}{1.823588in}}{\pgfqpoint{1.637426in}{1.818002in}}{\pgfqpoint{1.641545in}{1.813884in}}%
\pgfpathcurveto{\pgfqpoint{1.645663in}{1.809766in}}{\pgfqpoint{1.651249in}{1.807452in}}{\pgfqpoint{1.657073in}{1.807452in}}%
\pgfpathlineto{\pgfqpoint{1.657073in}{1.807452in}}%
\pgfpathclose%
\pgfusepath{stroke,fill}%
\end{pgfscope}%
\begin{pgfscope}%
\pgfpathrectangle{\pgfqpoint{0.100000in}{0.209042in}}{\pgfqpoint{2.906250in}{2.906250in}}%
\pgfusepath{clip}%
\pgfsetbuttcap%
\pgfsetroundjoin%
\definecolor{currentfill}{rgb}{0.506629,0.145958,0.507806}%
\pgfsetfillcolor{currentfill}%
\pgfsetfillopacity{0.800000}%
\pgfsetlinewidth{1.003750pt}%
\definecolor{currentstroke}{rgb}{0.506629,0.145958,0.507806}%
\pgfsetstrokecolor{currentstroke}%
\pgfsetstrokeopacity{0.800000}%
\pgfsetdash{}{0pt}%
\pgfpathmoveto{\pgfqpoint{1.416717in}{1.518251in}}%
\pgfpathcurveto{\pgfqpoint{1.422541in}{1.518251in}}{\pgfqpoint{1.428127in}{1.520565in}}{\pgfqpoint{1.432245in}{1.524683in}}%
\pgfpathcurveto{\pgfqpoint{1.436364in}{1.528801in}}{\pgfqpoint{1.438677in}{1.534387in}}{\pgfqpoint{1.438677in}{1.540211in}}%
\pgfpathcurveto{\pgfqpoint{1.438677in}{1.546035in}}{\pgfqpoint{1.436364in}{1.551621in}}{\pgfqpoint{1.432245in}{1.555740in}}%
\pgfpathcurveto{\pgfqpoint{1.428127in}{1.559858in}}{\pgfqpoint{1.422541in}{1.562172in}}{\pgfqpoint{1.416717in}{1.562172in}}%
\pgfpathcurveto{\pgfqpoint{1.410893in}{1.562172in}}{\pgfqpoint{1.405307in}{1.559858in}}{\pgfqpoint{1.401189in}{1.555740in}}%
\pgfpathcurveto{\pgfqpoint{1.397071in}{1.551621in}}{\pgfqpoint{1.394757in}{1.546035in}}{\pgfqpoint{1.394757in}{1.540211in}}%
\pgfpathcurveto{\pgfqpoint{1.394757in}{1.534387in}}{\pgfqpoint{1.397071in}{1.528801in}}{\pgfqpoint{1.401189in}{1.524683in}}%
\pgfpathcurveto{\pgfqpoint{1.405307in}{1.520565in}}{\pgfqpoint{1.410893in}{1.518251in}}{\pgfqpoint{1.416717in}{1.518251in}}%
\pgfpathlineto{\pgfqpoint{1.416717in}{1.518251in}}%
\pgfpathclose%
\pgfusepath{stroke,fill}%
\end{pgfscope}%
\begin{pgfscope}%
\pgfpathrectangle{\pgfqpoint{0.100000in}{0.209042in}}{\pgfqpoint{2.906250in}{2.906250in}}%
\pgfusepath{clip}%
\pgfsetbuttcap%
\pgfsetroundjoin%
\definecolor{currentfill}{rgb}{0.703532,0.210638,0.479049}%
\pgfsetfillcolor{currentfill}%
\pgfsetfillopacity{0.800000}%
\pgfsetlinewidth{1.003750pt}%
\definecolor{currentstroke}{rgb}{0.703532,0.210638,0.479049}%
\pgfsetstrokecolor{currentstroke}%
\pgfsetstrokeopacity{0.800000}%
\pgfsetdash{}{0pt}%
\pgfpathmoveto{\pgfqpoint{1.937311in}{1.450887in}}%
\pgfpathcurveto{\pgfqpoint{1.943135in}{1.450887in}}{\pgfqpoint{1.948721in}{1.453201in}}{\pgfqpoint{1.952839in}{1.457319in}}%
\pgfpathcurveto{\pgfqpoint{1.956957in}{1.461437in}}{\pgfqpoint{1.959271in}{1.467023in}}{\pgfqpoint{1.959271in}{1.472847in}}%
\pgfpathcurveto{\pgfqpoint{1.959271in}{1.478671in}}{\pgfqpoint{1.956957in}{1.484257in}}{\pgfqpoint{1.952839in}{1.488375in}}%
\pgfpathcurveto{\pgfqpoint{1.948721in}{1.492493in}}{\pgfqpoint{1.943135in}{1.494807in}}{\pgfqpoint{1.937311in}{1.494807in}}%
\pgfpathcurveto{\pgfqpoint{1.931487in}{1.494807in}}{\pgfqpoint{1.925901in}{1.492493in}}{\pgfqpoint{1.921782in}{1.488375in}}%
\pgfpathcurveto{\pgfqpoint{1.917664in}{1.484257in}}{\pgfqpoint{1.915350in}{1.478671in}}{\pgfqpoint{1.915350in}{1.472847in}}%
\pgfpathcurveto{\pgfqpoint{1.915350in}{1.467023in}}{\pgfqpoint{1.917664in}{1.461437in}}{\pgfqpoint{1.921782in}{1.457319in}}%
\pgfpathcurveto{\pgfqpoint{1.925901in}{1.453201in}}{\pgfqpoint{1.931487in}{1.450887in}}{\pgfqpoint{1.937311in}{1.450887in}}%
\pgfpathlineto{\pgfqpoint{1.937311in}{1.450887in}}%
\pgfpathclose%
\pgfusepath{stroke,fill}%
\end{pgfscope}%
\begin{pgfscope}%
\pgfpathrectangle{\pgfqpoint{0.100000in}{0.209042in}}{\pgfqpoint{2.906250in}{2.906250in}}%
\pgfusepath{clip}%
\pgfsetbuttcap%
\pgfsetroundjoin%
\definecolor{currentfill}{rgb}{0.494258,0.141462,0.507988}%
\pgfsetfillcolor{currentfill}%
\pgfsetfillopacity{0.800000}%
\pgfsetlinewidth{1.003750pt}%
\definecolor{currentstroke}{rgb}{0.494258,0.141462,0.507988}%
\pgfsetstrokecolor{currentstroke}%
\pgfsetstrokeopacity{0.800000}%
\pgfsetdash{}{0pt}%
\pgfpathmoveto{\pgfqpoint{1.566989in}{2.002963in}}%
\pgfpathcurveto{\pgfqpoint{1.572813in}{2.002963in}}{\pgfqpoint{1.578399in}{2.005276in}}{\pgfqpoint{1.582518in}{2.009395in}}%
\pgfpathcurveto{\pgfqpoint{1.586636in}{2.013513in}}{\pgfqpoint{1.588950in}{2.019099in}}{\pgfqpoint{1.588950in}{2.024923in}}%
\pgfpathcurveto{\pgfqpoint{1.588950in}{2.030747in}}{\pgfqpoint{1.586636in}{2.036333in}}{\pgfqpoint{1.582518in}{2.040451in}}%
\pgfpathcurveto{\pgfqpoint{1.578399in}{2.044569in}}{\pgfqpoint{1.572813in}{2.046883in}}{\pgfqpoint{1.566989in}{2.046883in}}%
\pgfpathcurveto{\pgfqpoint{1.561165in}{2.046883in}}{\pgfqpoint{1.555579in}{2.044569in}}{\pgfqpoint{1.551461in}{2.040451in}}%
\pgfpathcurveto{\pgfqpoint{1.547343in}{2.036333in}}{\pgfqpoint{1.545029in}{2.030747in}}{\pgfqpoint{1.545029in}{2.024923in}}%
\pgfpathcurveto{\pgfqpoint{1.545029in}{2.019099in}}{\pgfqpoint{1.547343in}{2.013513in}}{\pgfqpoint{1.551461in}{2.009395in}}%
\pgfpathcurveto{\pgfqpoint{1.555579in}{2.005276in}}{\pgfqpoint{1.561165in}{2.002963in}}{\pgfqpoint{1.566989in}{2.002963in}}%
\pgfpathlineto{\pgfqpoint{1.566989in}{2.002963in}}%
\pgfpathclose%
\pgfusepath{stroke,fill}%
\end{pgfscope}%
\begin{pgfscope}%
\pgfpathrectangle{\pgfqpoint{0.100000in}{0.209042in}}{\pgfqpoint{2.906250in}{2.906250in}}%
\pgfusepath{clip}%
\pgfsetbuttcap%
\pgfsetroundjoin%
\definecolor{currentfill}{rgb}{0.432967,0.117855,0.506160}%
\pgfsetfillcolor{currentfill}%
\pgfsetfillopacity{0.800000}%
\pgfsetlinewidth{1.003750pt}%
\definecolor{currentstroke}{rgb}{0.432967,0.117855,0.506160}%
\pgfsetstrokecolor{currentstroke}%
\pgfsetstrokeopacity{0.800000}%
\pgfsetdash{}{0pt}%
\pgfpathmoveto{\pgfqpoint{1.851975in}{1.797781in}}%
\pgfpathcurveto{\pgfqpoint{1.857799in}{1.797781in}}{\pgfqpoint{1.863385in}{1.800095in}}{\pgfqpoint{1.867503in}{1.804213in}}%
\pgfpathcurveto{\pgfqpoint{1.871621in}{1.808331in}}{\pgfqpoint{1.873935in}{1.813917in}}{\pgfqpoint{1.873935in}{1.819741in}}%
\pgfpathcurveto{\pgfqpoint{1.873935in}{1.825565in}}{\pgfqpoint{1.871621in}{1.831151in}}{\pgfqpoint{1.867503in}{1.835270in}}%
\pgfpathcurveto{\pgfqpoint{1.863385in}{1.839388in}}{\pgfqpoint{1.857799in}{1.841702in}}{\pgfqpoint{1.851975in}{1.841702in}}%
\pgfpathcurveto{\pgfqpoint{1.846151in}{1.841702in}}{\pgfqpoint{1.840565in}{1.839388in}}{\pgfqpoint{1.836447in}{1.835270in}}%
\pgfpathcurveto{\pgfqpoint{1.832329in}{1.831151in}}{\pgfqpoint{1.830015in}{1.825565in}}{\pgfqpoint{1.830015in}{1.819741in}}%
\pgfpathcurveto{\pgfqpoint{1.830015in}{1.813917in}}{\pgfqpoint{1.832329in}{1.808331in}}{\pgfqpoint{1.836447in}{1.804213in}}%
\pgfpathcurveto{\pgfqpoint{1.840565in}{1.800095in}}{\pgfqpoint{1.846151in}{1.797781in}}{\pgfqpoint{1.851975in}{1.797781in}}%
\pgfpathlineto{\pgfqpoint{1.851975in}{1.797781in}}%
\pgfpathclose%
\pgfusepath{stroke,fill}%
\end{pgfscope}%
\begin{pgfscope}%
\pgfpathrectangle{\pgfqpoint{0.100000in}{0.209042in}}{\pgfqpoint{2.906250in}{2.906250in}}%
\pgfusepath{clip}%
\pgfsetbuttcap%
\pgfsetroundjoin%
\definecolor{currentfill}{rgb}{0.512831,0.148179,0.507648}%
\pgfsetfillcolor{currentfill}%
\pgfsetfillopacity{0.800000}%
\pgfsetlinewidth{1.003750pt}%
\definecolor{currentstroke}{rgb}{0.512831,0.148179,0.507648}%
\pgfsetstrokecolor{currentstroke}%
\pgfsetstrokeopacity{0.800000}%
\pgfsetdash{}{0pt}%
\pgfpathmoveto{\pgfqpoint{1.572450in}{1.465019in}}%
\pgfpathcurveto{\pgfqpoint{1.578274in}{1.465019in}}{\pgfqpoint{1.583860in}{1.467333in}}{\pgfqpoint{1.587978in}{1.471451in}}%
\pgfpathcurveto{\pgfqpoint{1.592096in}{1.475569in}}{\pgfqpoint{1.594410in}{1.481156in}}{\pgfqpoint{1.594410in}{1.486980in}}%
\pgfpathcurveto{\pgfqpoint{1.594410in}{1.492804in}}{\pgfqpoint{1.592096in}{1.498390in}}{\pgfqpoint{1.587978in}{1.502508in}}%
\pgfpathcurveto{\pgfqpoint{1.583860in}{1.506626in}}{\pgfqpoint{1.578274in}{1.508940in}}{\pgfqpoint{1.572450in}{1.508940in}}%
\pgfpathcurveto{\pgfqpoint{1.566626in}{1.508940in}}{\pgfqpoint{1.561040in}{1.506626in}}{\pgfqpoint{1.556922in}{1.502508in}}%
\pgfpathcurveto{\pgfqpoint{1.552803in}{1.498390in}}{\pgfqpoint{1.550490in}{1.492804in}}{\pgfqpoint{1.550490in}{1.486980in}}%
\pgfpathcurveto{\pgfqpoint{1.550490in}{1.481156in}}{\pgfqpoint{1.552803in}{1.475569in}}{\pgfqpoint{1.556922in}{1.471451in}}%
\pgfpathcurveto{\pgfqpoint{1.561040in}{1.467333in}}{\pgfqpoint{1.566626in}{1.465019in}}{\pgfqpoint{1.572450in}{1.465019in}}%
\pgfpathlineto{\pgfqpoint{1.572450in}{1.465019in}}%
\pgfpathclose%
\pgfusepath{stroke,fill}%
\end{pgfscope}%
\begin{pgfscope}%
\pgfpathrectangle{\pgfqpoint{0.100000in}{0.209042in}}{\pgfqpoint{2.906250in}{2.906250in}}%
\pgfusepath{clip}%
\pgfsetbuttcap%
\pgfsetroundjoin%
\definecolor{currentfill}{rgb}{0.359898,0.087831,0.496778}%
\pgfsetfillcolor{currentfill}%
\pgfsetfillopacity{0.800000}%
\pgfsetlinewidth{1.003750pt}%
\definecolor{currentstroke}{rgb}{0.359898,0.087831,0.496778}%
\pgfsetstrokecolor{currentstroke}%
\pgfsetstrokeopacity{0.800000}%
\pgfsetdash{}{0pt}%
\pgfpathmoveto{\pgfqpoint{1.398055in}{1.786024in}}%
\pgfpathcurveto{\pgfqpoint{1.403879in}{1.786024in}}{\pgfqpoint{1.409465in}{1.788337in}}{\pgfqpoint{1.413583in}{1.792456in}}%
\pgfpathcurveto{\pgfqpoint{1.417702in}{1.796574in}}{\pgfqpoint{1.420015in}{1.802160in}}{\pgfqpoint{1.420015in}{1.807984in}}%
\pgfpathcurveto{\pgfqpoint{1.420015in}{1.813808in}}{\pgfqpoint{1.417702in}{1.819394in}}{\pgfqpoint{1.413583in}{1.823512in}}%
\pgfpathcurveto{\pgfqpoint{1.409465in}{1.827630in}}{\pgfqpoint{1.403879in}{1.829944in}}{\pgfqpoint{1.398055in}{1.829944in}}%
\pgfpathcurveto{\pgfqpoint{1.392231in}{1.829944in}}{\pgfqpoint{1.386645in}{1.827630in}}{\pgfqpoint{1.382527in}{1.823512in}}%
\pgfpathcurveto{\pgfqpoint{1.378409in}{1.819394in}}{\pgfqpoint{1.376095in}{1.813808in}}{\pgfqpoint{1.376095in}{1.807984in}}%
\pgfpathcurveto{\pgfqpoint{1.376095in}{1.802160in}}{\pgfqpoint{1.378409in}{1.796574in}}{\pgfqpoint{1.382527in}{1.792456in}}%
\pgfpathcurveto{\pgfqpoint{1.386645in}{1.788337in}}{\pgfqpoint{1.392231in}{1.786024in}}{\pgfqpoint{1.398055in}{1.786024in}}%
\pgfpathlineto{\pgfqpoint{1.398055in}{1.786024in}}%
\pgfpathclose%
\pgfusepath{stroke,fill}%
\end{pgfscope}%
\begin{pgfscope}%
\pgfpathrectangle{\pgfqpoint{0.100000in}{0.209042in}}{\pgfqpoint{2.906250in}{2.906250in}}%
\pgfusepath{clip}%
\pgfsetbuttcap%
\pgfsetroundjoin%
\definecolor{currentfill}{rgb}{0.475780,0.134577,0.507921}%
\pgfsetfillcolor{currentfill}%
\pgfsetfillopacity{0.800000}%
\pgfsetlinewidth{1.003750pt}%
\definecolor{currentstroke}{rgb}{0.475780,0.134577,0.507921}%
\pgfsetstrokecolor{currentstroke}%
\pgfsetstrokeopacity{0.800000}%
\pgfsetdash{}{0pt}%
\pgfpathmoveto{\pgfqpoint{1.379486in}{1.930193in}}%
\pgfpathcurveto{\pgfqpoint{1.385310in}{1.930193in}}{\pgfqpoint{1.390896in}{1.932507in}}{\pgfqpoint{1.395014in}{1.936625in}}%
\pgfpathcurveto{\pgfqpoint{1.399132in}{1.940743in}}{\pgfqpoint{1.401446in}{1.946330in}}{\pgfqpoint{1.401446in}{1.952153in}}%
\pgfpathcurveto{\pgfqpoint{1.401446in}{1.957977in}}{\pgfqpoint{1.399132in}{1.963564in}}{\pgfqpoint{1.395014in}{1.967682in}}%
\pgfpathcurveto{\pgfqpoint{1.390896in}{1.971800in}}{\pgfqpoint{1.385310in}{1.974114in}}{\pgfqpoint{1.379486in}{1.974114in}}%
\pgfpathcurveto{\pgfqpoint{1.373662in}{1.974114in}}{\pgfqpoint{1.368076in}{1.971800in}}{\pgfqpoint{1.363958in}{1.967682in}}%
\pgfpathcurveto{\pgfqpoint{1.359840in}{1.963564in}}{\pgfqpoint{1.357526in}{1.957977in}}{\pgfqpoint{1.357526in}{1.952153in}}%
\pgfpathcurveto{\pgfqpoint{1.357526in}{1.946330in}}{\pgfqpoint{1.359840in}{1.940743in}}{\pgfqpoint{1.363958in}{1.936625in}}%
\pgfpathcurveto{\pgfqpoint{1.368076in}{1.932507in}}{\pgfqpoint{1.373662in}{1.930193in}}{\pgfqpoint{1.379486in}{1.930193in}}%
\pgfpathlineto{\pgfqpoint{1.379486in}{1.930193in}}%
\pgfpathclose%
\pgfusepath{stroke,fill}%
\end{pgfscope}%
\begin{pgfscope}%
\pgfpathrectangle{\pgfqpoint{0.100000in}{0.209042in}}{\pgfqpoint{2.906250in}{2.906250in}}%
\pgfusepath{clip}%
\pgfsetbuttcap%
\pgfsetroundjoin%
\definecolor{currentfill}{rgb}{0.588158,0.173652,0.502035}%
\pgfsetfillcolor{currentfill}%
\pgfsetfillopacity{0.800000}%
\pgfsetlinewidth{1.003750pt}%
\definecolor{currentstroke}{rgb}{0.588158,0.173652,0.502035}%
\pgfsetstrokecolor{currentstroke}%
\pgfsetstrokeopacity{0.800000}%
\pgfsetdash{}{0pt}%
\pgfpathmoveto{\pgfqpoint{1.337253in}{1.500171in}}%
\pgfpathcurveto{\pgfqpoint{1.343076in}{1.500171in}}{\pgfqpoint{1.348663in}{1.502484in}}{\pgfqpoint{1.352781in}{1.506603in}}%
\pgfpathcurveto{\pgfqpoint{1.356899in}{1.510721in}}{\pgfqpoint{1.359213in}{1.516307in}}{\pgfqpoint{1.359213in}{1.522131in}}%
\pgfpathcurveto{\pgfqpoint{1.359213in}{1.527955in}}{\pgfqpoint{1.356899in}{1.533541in}}{\pgfqpoint{1.352781in}{1.537659in}}%
\pgfpathcurveto{\pgfqpoint{1.348663in}{1.541777in}}{\pgfqpoint{1.343076in}{1.544091in}}{\pgfqpoint{1.337253in}{1.544091in}}%
\pgfpathcurveto{\pgfqpoint{1.331429in}{1.544091in}}{\pgfqpoint{1.325842in}{1.541777in}}{\pgfqpoint{1.321724in}{1.537659in}}%
\pgfpathcurveto{\pgfqpoint{1.317606in}{1.533541in}}{\pgfqpoint{1.315292in}{1.527955in}}{\pgfqpoint{1.315292in}{1.522131in}}%
\pgfpathcurveto{\pgfqpoint{1.315292in}{1.516307in}}{\pgfqpoint{1.317606in}{1.510721in}}{\pgfqpoint{1.321724in}{1.506603in}}%
\pgfpathcurveto{\pgfqpoint{1.325842in}{1.502484in}}{\pgfqpoint{1.331429in}{1.500171in}}{\pgfqpoint{1.337253in}{1.500171in}}%
\pgfpathlineto{\pgfqpoint{1.337253in}{1.500171in}}%
\pgfpathclose%
\pgfusepath{stroke,fill}%
\end{pgfscope}%
\begin{pgfscope}%
\pgfpathrectangle{\pgfqpoint{0.100000in}{0.209042in}}{\pgfqpoint{2.906250in}{2.906250in}}%
\pgfusepath{clip}%
\pgfsetbuttcap%
\pgfsetroundjoin%
\definecolor{currentfill}{rgb}{0.742004,0.224025,0.467018}%
\pgfsetfillcolor{currentfill}%
\pgfsetfillopacity{0.800000}%
\pgfsetlinewidth{1.003750pt}%
\definecolor{currentstroke}{rgb}{0.742004,0.224025,0.467018}%
\pgfsetstrokecolor{currentstroke}%
\pgfsetstrokeopacity{0.800000}%
\pgfsetdash{}{0pt}%
\pgfpathmoveto{\pgfqpoint{1.080434in}{1.740479in}}%
\pgfpathcurveto{\pgfqpoint{1.086258in}{1.740479in}}{\pgfqpoint{1.091845in}{1.742793in}}{\pgfqpoint{1.095963in}{1.746911in}}%
\pgfpathcurveto{\pgfqpoint{1.100081in}{1.751029in}}{\pgfqpoint{1.102395in}{1.756615in}}{\pgfqpoint{1.102395in}{1.762439in}}%
\pgfpathcurveto{\pgfqpoint{1.102395in}{1.768263in}}{\pgfqpoint{1.100081in}{1.773849in}}{\pgfqpoint{1.095963in}{1.777967in}}%
\pgfpathcurveto{\pgfqpoint{1.091845in}{1.782085in}}{\pgfqpoint{1.086258in}{1.784399in}}{\pgfqpoint{1.080434in}{1.784399in}}%
\pgfpathcurveto{\pgfqpoint{1.074611in}{1.784399in}}{\pgfqpoint{1.069024in}{1.782085in}}{\pgfqpoint{1.064906in}{1.777967in}}%
\pgfpathcurveto{\pgfqpoint{1.060788in}{1.773849in}}{\pgfqpoint{1.058474in}{1.768263in}}{\pgfqpoint{1.058474in}{1.762439in}}%
\pgfpathcurveto{\pgfqpoint{1.058474in}{1.756615in}}{\pgfqpoint{1.060788in}{1.751029in}}{\pgfqpoint{1.064906in}{1.746911in}}%
\pgfpathcurveto{\pgfqpoint{1.069024in}{1.742793in}}{\pgfqpoint{1.074611in}{1.740479in}}{\pgfqpoint{1.080434in}{1.740479in}}%
\pgfpathlineto{\pgfqpoint{1.080434in}{1.740479in}}%
\pgfpathclose%
\pgfusepath{stroke,fill}%
\end{pgfscope}%
\begin{pgfscope}%
\pgfpathrectangle{\pgfqpoint{0.100000in}{0.209042in}}{\pgfqpoint{2.906250in}{2.906250in}}%
\pgfusepath{clip}%
\pgfsetbuttcap%
\pgfsetroundjoin%
\definecolor{currentfill}{rgb}{0.852126,0.276106,0.418573}%
\pgfsetfillcolor{currentfill}%
\pgfsetfillopacity{0.800000}%
\pgfsetlinewidth{1.003750pt}%
\definecolor{currentstroke}{rgb}{0.852126,0.276106,0.418573}%
\pgfsetstrokecolor{currentstroke}%
\pgfsetstrokeopacity{0.800000}%
\pgfsetdash{}{0pt}%
\pgfpathmoveto{\pgfqpoint{2.177812in}{1.619328in}}%
\pgfpathcurveto{\pgfqpoint{2.183636in}{1.619328in}}{\pgfqpoint{2.189222in}{1.621642in}}{\pgfqpoint{2.193340in}{1.625760in}}%
\pgfpathcurveto{\pgfqpoint{2.197458in}{1.629879in}}{\pgfqpoint{2.199772in}{1.635465in}}{\pgfqpoint{2.199772in}{1.641289in}}%
\pgfpathcurveto{\pgfqpoint{2.199772in}{1.647113in}}{\pgfqpoint{2.197458in}{1.652699in}}{\pgfqpoint{2.193340in}{1.656817in}}%
\pgfpathcurveto{\pgfqpoint{2.189222in}{1.660935in}}{\pgfqpoint{2.183636in}{1.663249in}}{\pgfqpoint{2.177812in}{1.663249in}}%
\pgfpathcurveto{\pgfqpoint{2.171988in}{1.663249in}}{\pgfqpoint{2.166402in}{1.660935in}}{\pgfqpoint{2.162284in}{1.656817in}}%
\pgfpathcurveto{\pgfqpoint{2.158165in}{1.652699in}}{\pgfqpoint{2.155852in}{1.647113in}}{\pgfqpoint{2.155852in}{1.641289in}}%
\pgfpathcurveto{\pgfqpoint{2.155852in}{1.635465in}}{\pgfqpoint{2.158165in}{1.629879in}}{\pgfqpoint{2.162284in}{1.625760in}}%
\pgfpathcurveto{\pgfqpoint{2.166402in}{1.621642in}}{\pgfqpoint{2.171988in}{1.619328in}}{\pgfqpoint{2.177812in}{1.619328in}}%
\pgfpathlineto{\pgfqpoint{2.177812in}{1.619328in}}%
\pgfpathclose%
\pgfusepath{stroke,fill}%
\end{pgfscope}%
\begin{pgfscope}%
\pgfpathrectangle{\pgfqpoint{0.100000in}{0.209042in}}{\pgfqpoint{2.906250in}{2.906250in}}%
\pgfusepath{clip}%
\pgfsetbuttcap%
\pgfsetroundjoin%
\definecolor{currentfill}{rgb}{0.316654,0.071690,0.485380}%
\pgfsetfillcolor{currentfill}%
\pgfsetfillopacity{0.800000}%
\pgfsetlinewidth{1.003750pt}%
\definecolor{currentstroke}{rgb}{0.316654,0.071690,0.485380}%
\pgfsetstrokecolor{currentstroke}%
\pgfsetstrokeopacity{0.800000}%
\pgfsetdash{}{0pt}%
\pgfpathmoveto{\pgfqpoint{1.667340in}{1.853815in}}%
\pgfpathcurveto{\pgfqpoint{1.673164in}{1.853815in}}{\pgfqpoint{1.678750in}{1.856129in}}{\pgfqpoint{1.682868in}{1.860247in}}%
\pgfpathcurveto{\pgfqpoint{1.686986in}{1.864366in}}{\pgfqpoint{1.689300in}{1.869952in}}{\pgfqpoint{1.689300in}{1.875776in}}%
\pgfpathcurveto{\pgfqpoint{1.689300in}{1.881600in}}{\pgfqpoint{1.686986in}{1.887186in}}{\pgfqpoint{1.682868in}{1.891304in}}%
\pgfpathcurveto{\pgfqpoint{1.678750in}{1.895422in}}{\pgfqpoint{1.673164in}{1.897736in}}{\pgfqpoint{1.667340in}{1.897736in}}%
\pgfpathcurveto{\pgfqpoint{1.661516in}{1.897736in}}{\pgfqpoint{1.655929in}{1.895422in}}{\pgfqpoint{1.651811in}{1.891304in}}%
\pgfpathcurveto{\pgfqpoint{1.647693in}{1.887186in}}{\pgfqpoint{1.645379in}{1.881600in}}{\pgfqpoint{1.645379in}{1.875776in}}%
\pgfpathcurveto{\pgfqpoint{1.645379in}{1.869952in}}{\pgfqpoint{1.647693in}{1.864366in}}{\pgfqpoint{1.651811in}{1.860247in}}%
\pgfpathcurveto{\pgfqpoint{1.655929in}{1.856129in}}{\pgfqpoint{1.661516in}{1.853815in}}{\pgfqpoint{1.667340in}{1.853815in}}%
\pgfpathlineto{\pgfqpoint{1.667340in}{1.853815in}}%
\pgfpathclose%
\pgfusepath{stroke,fill}%
\end{pgfscope}%
\begin{pgfscope}%
\pgfpathrectangle{\pgfqpoint{0.100000in}{0.209042in}}{\pgfqpoint{2.906250in}{2.906250in}}%
\pgfusepath{clip}%
\pgfsetbuttcap%
\pgfsetroundjoin%
\definecolor{currentfill}{rgb}{0.353773,0.085373,0.495501}%
\pgfsetfillcolor{currentfill}%
\pgfsetfillopacity{0.800000}%
\pgfsetlinewidth{1.003750pt}%
\definecolor{currentstroke}{rgb}{0.353773,0.085373,0.495501}%
\pgfsetstrokecolor{currentstroke}%
\pgfsetstrokeopacity{0.800000}%
\pgfsetdash{}{0pt}%
\pgfpathmoveto{\pgfqpoint{1.800833in}{1.708307in}}%
\pgfpathcurveto{\pgfqpoint{1.806657in}{1.708307in}}{\pgfqpoint{1.812243in}{1.710621in}}{\pgfqpoint{1.816361in}{1.714739in}}%
\pgfpathcurveto{\pgfqpoint{1.820479in}{1.718858in}}{\pgfqpoint{1.822793in}{1.724444in}}{\pgfqpoint{1.822793in}{1.730268in}}%
\pgfpathcurveto{\pgfqpoint{1.822793in}{1.736092in}}{\pgfqpoint{1.820479in}{1.741678in}}{\pgfqpoint{1.816361in}{1.745796in}}%
\pgfpathcurveto{\pgfqpoint{1.812243in}{1.749914in}}{\pgfqpoint{1.806657in}{1.752228in}}{\pgfqpoint{1.800833in}{1.752228in}}%
\pgfpathcurveto{\pgfqpoint{1.795009in}{1.752228in}}{\pgfqpoint{1.789423in}{1.749914in}}{\pgfqpoint{1.785305in}{1.745796in}}%
\pgfpathcurveto{\pgfqpoint{1.781186in}{1.741678in}}{\pgfqpoint{1.778873in}{1.736092in}}{\pgfqpoint{1.778873in}{1.730268in}}%
\pgfpathcurveto{\pgfqpoint{1.778873in}{1.724444in}}{\pgfqpoint{1.781186in}{1.718858in}}{\pgfqpoint{1.785305in}{1.714739in}}%
\pgfpathcurveto{\pgfqpoint{1.789423in}{1.710621in}}{\pgfqpoint{1.795009in}{1.708307in}}{\pgfqpoint{1.800833in}{1.708307in}}%
\pgfpathlineto{\pgfqpoint{1.800833in}{1.708307in}}%
\pgfpathclose%
\pgfusepath{stroke,fill}%
\end{pgfscope}%
\begin{pgfscope}%
\pgfpathrectangle{\pgfqpoint{0.100000in}{0.209042in}}{\pgfqpoint{2.906250in}{2.906250in}}%
\pgfusepath{clip}%
\pgfsetbuttcap%
\pgfsetroundjoin%
\definecolor{currentfill}{rgb}{0.258857,0.059706,0.457710}%
\pgfsetfillcolor{currentfill}%
\pgfsetfillopacity{0.800000}%
\pgfsetlinewidth{1.003750pt}%
\definecolor{currentstroke}{rgb}{0.258857,0.059706,0.457710}%
\pgfsetstrokecolor{currentstroke}%
\pgfsetstrokeopacity{0.800000}%
\pgfsetdash{}{0pt}%
\pgfpathmoveto{\pgfqpoint{1.583797in}{1.668390in}}%
\pgfpathcurveto{\pgfqpoint{1.589621in}{1.668390in}}{\pgfqpoint{1.595207in}{1.670704in}}{\pgfqpoint{1.599325in}{1.674822in}}%
\pgfpathcurveto{\pgfqpoint{1.603443in}{1.678940in}}{\pgfqpoint{1.605757in}{1.684526in}}{\pgfqpoint{1.605757in}{1.690350in}}%
\pgfpathcurveto{\pgfqpoint{1.605757in}{1.696174in}}{\pgfqpoint{1.603443in}{1.701760in}}{\pgfqpoint{1.599325in}{1.705878in}}%
\pgfpathcurveto{\pgfqpoint{1.595207in}{1.709996in}}{\pgfqpoint{1.589621in}{1.712310in}}{\pgfqpoint{1.583797in}{1.712310in}}%
\pgfpathcurveto{\pgfqpoint{1.577973in}{1.712310in}}{\pgfqpoint{1.572387in}{1.709996in}}{\pgfqpoint{1.568269in}{1.705878in}}%
\pgfpathcurveto{\pgfqpoint{1.564151in}{1.701760in}}{\pgfqpoint{1.561837in}{1.696174in}}{\pgfqpoint{1.561837in}{1.690350in}}%
\pgfpathcurveto{\pgfqpoint{1.561837in}{1.684526in}}{\pgfqpoint{1.564151in}{1.678940in}}{\pgfqpoint{1.568269in}{1.674822in}}%
\pgfpathcurveto{\pgfqpoint{1.572387in}{1.670704in}}{\pgfqpoint{1.577973in}{1.668390in}}{\pgfqpoint{1.583797in}{1.668390in}}%
\pgfpathlineto{\pgfqpoint{1.583797in}{1.668390in}}%
\pgfpathclose%
\pgfusepath{stroke,fill}%
\end{pgfscope}%
\begin{pgfscope}%
\pgfpathrectangle{\pgfqpoint{0.100000in}{0.209042in}}{\pgfqpoint{2.906250in}{2.906250in}}%
\pgfusepath{clip}%
\pgfsetbuttcap%
\pgfsetroundjoin%
\definecolor{currentfill}{rgb}{0.278493,0.061978,0.469190}%
\pgfsetfillcolor{currentfill}%
\pgfsetfillopacity{0.800000}%
\pgfsetlinewidth{1.003750pt}%
\definecolor{currentstroke}{rgb}{0.278493,0.061978,0.469190}%
\pgfsetstrokecolor{currentstroke}%
\pgfsetstrokeopacity{0.800000}%
\pgfsetdash{}{0pt}%
\pgfpathmoveto{\pgfqpoint{1.501451in}{1.804793in}}%
\pgfpathcurveto{\pgfqpoint{1.507275in}{1.804793in}}{\pgfqpoint{1.512861in}{1.807107in}}{\pgfqpoint{1.516979in}{1.811225in}}%
\pgfpathcurveto{\pgfqpoint{1.521098in}{1.815343in}}{\pgfqpoint{1.523411in}{1.820929in}}{\pgfqpoint{1.523411in}{1.826753in}}%
\pgfpathcurveto{\pgfqpoint{1.523411in}{1.832577in}}{\pgfqpoint{1.521098in}{1.838163in}}{\pgfqpoint{1.516979in}{1.842282in}}%
\pgfpathcurveto{\pgfqpoint{1.512861in}{1.846400in}}{\pgfqpoint{1.507275in}{1.848714in}}{\pgfqpoint{1.501451in}{1.848714in}}%
\pgfpathcurveto{\pgfqpoint{1.495627in}{1.848714in}}{\pgfqpoint{1.490041in}{1.846400in}}{\pgfqpoint{1.485923in}{1.842282in}}%
\pgfpathcurveto{\pgfqpoint{1.481805in}{1.838163in}}{\pgfqpoint{1.479491in}{1.832577in}}{\pgfqpoint{1.479491in}{1.826753in}}%
\pgfpathcurveto{\pgfqpoint{1.479491in}{1.820929in}}{\pgfqpoint{1.481805in}{1.815343in}}{\pgfqpoint{1.485923in}{1.811225in}}%
\pgfpathcurveto{\pgfqpoint{1.490041in}{1.807107in}}{\pgfqpoint{1.495627in}{1.804793in}}{\pgfqpoint{1.501451in}{1.804793in}}%
\pgfpathlineto{\pgfqpoint{1.501451in}{1.804793in}}%
\pgfpathclose%
\pgfusepath{stroke,fill}%
\end{pgfscope}%
\begin{pgfscope}%
\pgfpathrectangle{\pgfqpoint{0.100000in}{0.209042in}}{\pgfqpoint{2.906250in}{2.906250in}}%
\pgfusepath{clip}%
\pgfsetbuttcap%
\pgfsetroundjoin%
\definecolor{currentfill}{rgb}{0.284951,0.063168,0.472451}%
\pgfsetfillcolor{currentfill}%
\pgfsetfillopacity{0.800000}%
\pgfsetlinewidth{1.003750pt}%
\definecolor{currentstroke}{rgb}{0.284951,0.063168,0.472451}%
\pgfsetstrokecolor{currentstroke}%
\pgfsetstrokeopacity{0.800000}%
\pgfsetdash{}{0pt}%
\pgfpathmoveto{\pgfqpoint{1.580085in}{1.622527in}}%
\pgfpathcurveto{\pgfqpoint{1.585909in}{1.622527in}}{\pgfqpoint{1.591495in}{1.624841in}}{\pgfqpoint{1.595614in}{1.628959in}}%
\pgfpathcurveto{\pgfqpoint{1.599732in}{1.633077in}}{\pgfqpoint{1.602046in}{1.638664in}}{\pgfqpoint{1.602046in}{1.644488in}}%
\pgfpathcurveto{\pgfqpoint{1.602046in}{1.650312in}}{\pgfqpoint{1.599732in}{1.655898in}}{\pgfqpoint{1.595614in}{1.660016in}}%
\pgfpathcurveto{\pgfqpoint{1.591495in}{1.664134in}}{\pgfqpoint{1.585909in}{1.666448in}}{\pgfqpoint{1.580085in}{1.666448in}}%
\pgfpathcurveto{\pgfqpoint{1.574261in}{1.666448in}}{\pgfqpoint{1.568675in}{1.664134in}}{\pgfqpoint{1.564557in}{1.660016in}}%
\pgfpathcurveto{\pgfqpoint{1.560439in}{1.655898in}}{\pgfqpoint{1.558125in}{1.650312in}}{\pgfqpoint{1.558125in}{1.644488in}}%
\pgfpathcurveto{\pgfqpoint{1.558125in}{1.638664in}}{\pgfqpoint{1.560439in}{1.633077in}}{\pgfqpoint{1.564557in}{1.628959in}}%
\pgfpathcurveto{\pgfqpoint{1.568675in}{1.624841in}}{\pgfqpoint{1.574261in}{1.622527in}}{\pgfqpoint{1.580085in}{1.622527in}}%
\pgfpathlineto{\pgfqpoint{1.580085in}{1.622527in}}%
\pgfpathclose%
\pgfusepath{stroke,fill}%
\end{pgfscope}%
\begin{pgfscope}%
\pgfpathrectangle{\pgfqpoint{0.100000in}{0.209042in}}{\pgfqpoint{2.906250in}{2.906250in}}%
\pgfusepath{clip}%
\pgfsetbuttcap%
\pgfsetroundjoin%
\definecolor{currentfill}{rgb}{0.992196,0.587502,0.406299}%
\pgfsetfillcolor{currentfill}%
\pgfsetfillopacity{0.800000}%
\pgfsetlinewidth{1.003750pt}%
\definecolor{currentstroke}{rgb}{0.992196,0.587502,0.406299}%
\pgfsetstrokecolor{currentstroke}%
\pgfsetstrokeopacity{0.800000}%
\pgfsetdash{}{0pt}%
\pgfpathmoveto{\pgfqpoint{2.404503in}{1.743969in}}%
\pgfpathcurveto{\pgfqpoint{2.410327in}{1.743969in}}{\pgfqpoint{2.415913in}{1.746283in}}{\pgfqpoint{2.420031in}{1.750401in}}%
\pgfpathcurveto{\pgfqpoint{2.424150in}{1.754519in}}{\pgfqpoint{2.426463in}{1.760106in}}{\pgfqpoint{2.426463in}{1.765930in}}%
\pgfpathcurveto{\pgfqpoint{2.426463in}{1.771753in}}{\pgfqpoint{2.424150in}{1.777340in}}{\pgfqpoint{2.420031in}{1.781458in}}%
\pgfpathcurveto{\pgfqpoint{2.415913in}{1.785576in}}{\pgfqpoint{2.410327in}{1.787890in}}{\pgfqpoint{2.404503in}{1.787890in}}%
\pgfpathcurveto{\pgfqpoint{2.398679in}{1.787890in}}{\pgfqpoint{2.393093in}{1.785576in}}{\pgfqpoint{2.388975in}{1.781458in}}%
\pgfpathcurveto{\pgfqpoint{2.384857in}{1.777340in}}{\pgfqpoint{2.382543in}{1.771753in}}{\pgfqpoint{2.382543in}{1.765930in}}%
\pgfpathcurveto{\pgfqpoint{2.382543in}{1.760106in}}{\pgfqpoint{2.384857in}{1.754519in}}{\pgfqpoint{2.388975in}{1.750401in}}%
\pgfpathcurveto{\pgfqpoint{2.393093in}{1.746283in}}{\pgfqpoint{2.398679in}{1.743969in}}{\pgfqpoint{2.404503in}{1.743969in}}%
\pgfpathlineto{\pgfqpoint{2.404503in}{1.743969in}}%
\pgfpathclose%
\pgfusepath{stroke,fill}%
\end{pgfscope}%
\begin{pgfscope}%
\pgfpathrectangle{\pgfqpoint{0.100000in}{0.209042in}}{\pgfqpoint{2.906250in}{2.906250in}}%
\pgfusepath{clip}%
\pgfsetbuttcap%
\pgfsetroundjoin%
\definecolor{currentfill}{rgb}{0.469640,0.132245,0.507809}%
\pgfsetfillcolor{currentfill}%
\pgfsetfillopacity{0.800000}%
\pgfsetlinewidth{1.003750pt}%
\definecolor{currentstroke}{rgb}{0.469640,0.132245,0.507809}%
\pgfsetstrokecolor{currentstroke}%
\pgfsetstrokeopacity{0.800000}%
\pgfsetdash{}{0pt}%
\pgfpathmoveto{\pgfqpoint{1.362768in}{1.562721in}}%
\pgfpathcurveto{\pgfqpoint{1.368592in}{1.562721in}}{\pgfqpoint{1.374178in}{1.565035in}}{\pgfqpoint{1.378296in}{1.569153in}}%
\pgfpathcurveto{\pgfqpoint{1.382415in}{1.573271in}}{\pgfqpoint{1.384728in}{1.578858in}}{\pgfqpoint{1.384728in}{1.584682in}}%
\pgfpathcurveto{\pgfqpoint{1.384728in}{1.590505in}}{\pgfqpoint{1.382415in}{1.596092in}}{\pgfqpoint{1.378296in}{1.600210in}}%
\pgfpathcurveto{\pgfqpoint{1.374178in}{1.604328in}}{\pgfqpoint{1.368592in}{1.606642in}}{\pgfqpoint{1.362768in}{1.606642in}}%
\pgfpathcurveto{\pgfqpoint{1.356944in}{1.606642in}}{\pgfqpoint{1.351358in}{1.604328in}}{\pgfqpoint{1.347240in}{1.600210in}}%
\pgfpathcurveto{\pgfqpoint{1.343122in}{1.596092in}}{\pgfqpoint{1.340808in}{1.590505in}}{\pgfqpoint{1.340808in}{1.584682in}}%
\pgfpathcurveto{\pgfqpoint{1.340808in}{1.578858in}}{\pgfqpoint{1.343122in}{1.573271in}}{\pgfqpoint{1.347240in}{1.569153in}}%
\pgfpathcurveto{\pgfqpoint{1.351358in}{1.565035in}}{\pgfqpoint{1.356944in}{1.562721in}}{\pgfqpoint{1.362768in}{1.562721in}}%
\pgfpathlineto{\pgfqpoint{1.362768in}{1.562721in}}%
\pgfpathclose%
\pgfusepath{stroke,fill}%
\end{pgfscope}%
\begin{pgfscope}%
\pgfpathrectangle{\pgfqpoint{0.100000in}{0.209042in}}{\pgfqpoint{2.906250in}{2.906250in}}%
\pgfusepath{clip}%
\pgfsetbuttcap%
\pgfsetroundjoin%
\definecolor{currentfill}{rgb}{0.786212,0.241514,0.450184}%
\pgfsetfillcolor{currentfill}%
\pgfsetfillopacity{0.800000}%
\pgfsetlinewidth{1.003750pt}%
\definecolor{currentstroke}{rgb}{0.786212,0.241514,0.450184}%
\pgfsetstrokecolor{currentstroke}%
\pgfsetstrokeopacity{0.800000}%
\pgfsetdash{}{0pt}%
\pgfpathmoveto{\pgfqpoint{1.523729in}{1.293688in}}%
\pgfpathcurveto{\pgfqpoint{1.529553in}{1.293688in}}{\pgfqpoint{1.535139in}{1.296002in}}{\pgfqpoint{1.539257in}{1.300120in}}%
\pgfpathcurveto{\pgfqpoint{1.543376in}{1.304238in}}{\pgfqpoint{1.545689in}{1.309824in}}{\pgfqpoint{1.545689in}{1.315648in}}%
\pgfpathcurveto{\pgfqpoint{1.545689in}{1.321472in}}{\pgfqpoint{1.543376in}{1.327058in}}{\pgfqpoint{1.539257in}{1.331176in}}%
\pgfpathcurveto{\pgfqpoint{1.535139in}{1.335295in}}{\pgfqpoint{1.529553in}{1.337608in}}{\pgfqpoint{1.523729in}{1.337608in}}%
\pgfpathcurveto{\pgfqpoint{1.517905in}{1.337608in}}{\pgfqpoint{1.512319in}{1.335295in}}{\pgfqpoint{1.508201in}{1.331176in}}%
\pgfpathcurveto{\pgfqpoint{1.504083in}{1.327058in}}{\pgfqpoint{1.501769in}{1.321472in}}{\pgfqpoint{1.501769in}{1.315648in}}%
\pgfpathcurveto{\pgfqpoint{1.501769in}{1.309824in}}{\pgfqpoint{1.504083in}{1.304238in}}{\pgfqpoint{1.508201in}{1.300120in}}%
\pgfpathcurveto{\pgfqpoint{1.512319in}{1.296002in}}{\pgfqpoint{1.517905in}{1.293688in}}{\pgfqpoint{1.523729in}{1.293688in}}%
\pgfpathlineto{\pgfqpoint{1.523729in}{1.293688in}}%
\pgfpathclose%
\pgfusepath{stroke,fill}%
\end{pgfscope}%
\begin{pgfscope}%
\pgfpathrectangle{\pgfqpoint{0.100000in}{0.209042in}}{\pgfqpoint{2.906250in}{2.906250in}}%
\pgfusepath{clip}%
\pgfsetbuttcap%
\pgfsetroundjoin%
\definecolor{currentfill}{rgb}{0.506629,0.145958,0.507806}%
\pgfsetfillcolor{currentfill}%
\pgfsetfillopacity{0.800000}%
\pgfsetlinewidth{1.003750pt}%
\definecolor{currentstroke}{rgb}{0.506629,0.145958,0.507806}%
\pgfsetstrokecolor{currentstroke}%
\pgfsetstrokeopacity{0.800000}%
\pgfsetdash{}{0pt}%
\pgfpathmoveto{\pgfqpoint{1.312521in}{1.575246in}}%
\pgfpathcurveto{\pgfqpoint{1.318344in}{1.575246in}}{\pgfqpoint{1.323931in}{1.577560in}}{\pgfqpoint{1.328049in}{1.581678in}}%
\pgfpathcurveto{\pgfqpoint{1.332167in}{1.585796in}}{\pgfqpoint{1.334481in}{1.591382in}}{\pgfqpoint{1.334481in}{1.597206in}}%
\pgfpathcurveto{\pgfqpoint{1.334481in}{1.603030in}}{\pgfqpoint{1.332167in}{1.608616in}}{\pgfqpoint{1.328049in}{1.612734in}}%
\pgfpathcurveto{\pgfqpoint{1.323931in}{1.616853in}}{\pgfqpoint{1.318344in}{1.619166in}}{\pgfqpoint{1.312521in}{1.619166in}}%
\pgfpathcurveto{\pgfqpoint{1.306697in}{1.619166in}}{\pgfqpoint{1.301110in}{1.616853in}}{\pgfqpoint{1.296992in}{1.612734in}}%
\pgfpathcurveto{\pgfqpoint{1.292874in}{1.608616in}}{\pgfqpoint{1.290560in}{1.603030in}}{\pgfqpoint{1.290560in}{1.597206in}}%
\pgfpathcurveto{\pgfqpoint{1.290560in}{1.591382in}}{\pgfqpoint{1.292874in}{1.585796in}}{\pgfqpoint{1.296992in}{1.581678in}}%
\pgfpathcurveto{\pgfqpoint{1.301110in}{1.577560in}}{\pgfqpoint{1.306697in}{1.575246in}}{\pgfqpoint{1.312521in}{1.575246in}}%
\pgfpathlineto{\pgfqpoint{1.312521in}{1.575246in}}%
\pgfpathclose%
\pgfusepath{stroke,fill}%
\end{pgfscope}%
\begin{pgfscope}%
\pgfpathrectangle{\pgfqpoint{0.100000in}{0.209042in}}{\pgfqpoint{2.906250in}{2.906250in}}%
\pgfusepath{clip}%
\pgfsetbuttcap%
\pgfsetroundjoin%
\definecolor{currentfill}{rgb}{0.265447,0.060237,0.461840}%
\pgfsetfillcolor{currentfill}%
\pgfsetfillopacity{0.800000}%
\pgfsetlinewidth{1.003750pt}%
\definecolor{currentstroke}{rgb}{0.265447,0.060237,0.461840}%
\pgfsetstrokecolor{currentstroke}%
\pgfsetstrokeopacity{0.800000}%
\pgfsetdash{}{0pt}%
\pgfpathmoveto{\pgfqpoint{1.669488in}{1.648853in}}%
\pgfpathcurveto{\pgfqpoint{1.675312in}{1.648853in}}{\pgfqpoint{1.680898in}{1.651167in}}{\pgfqpoint{1.685016in}{1.655285in}}%
\pgfpathcurveto{\pgfqpoint{1.689134in}{1.659403in}}{\pgfqpoint{1.691448in}{1.664990in}}{\pgfqpoint{1.691448in}{1.670813in}}%
\pgfpathcurveto{\pgfqpoint{1.691448in}{1.676637in}}{\pgfqpoint{1.689134in}{1.682224in}}{\pgfqpoint{1.685016in}{1.686342in}}%
\pgfpathcurveto{\pgfqpoint{1.680898in}{1.690460in}}{\pgfqpoint{1.675312in}{1.692774in}}{\pgfqpoint{1.669488in}{1.692774in}}%
\pgfpathcurveto{\pgfqpoint{1.663664in}{1.692774in}}{\pgfqpoint{1.658078in}{1.690460in}}{\pgfqpoint{1.653960in}{1.686342in}}%
\pgfpathcurveto{\pgfqpoint{1.649842in}{1.682224in}}{\pgfqpoint{1.647528in}{1.676637in}}{\pgfqpoint{1.647528in}{1.670813in}}%
\pgfpathcurveto{\pgfqpoint{1.647528in}{1.664990in}}{\pgfqpoint{1.649842in}{1.659403in}}{\pgfqpoint{1.653960in}{1.655285in}}%
\pgfpathcurveto{\pgfqpoint{1.658078in}{1.651167in}}{\pgfqpoint{1.663664in}{1.648853in}}{\pgfqpoint{1.669488in}{1.648853in}}%
\pgfpathlineto{\pgfqpoint{1.669488in}{1.648853in}}%
\pgfpathclose%
\pgfusepath{stroke,fill}%
\end{pgfscope}%
\begin{pgfscope}%
\pgfpathrectangle{\pgfqpoint{0.100000in}{0.209042in}}{\pgfqpoint{2.906250in}{2.906250in}}%
\pgfusepath{clip}%
\pgfsetbuttcap%
\pgfsetroundjoin%
\definecolor{currentfill}{rgb}{0.372116,0.092816,0.499053}%
\pgfsetfillcolor{currentfill}%
\pgfsetfillopacity{0.800000}%
\pgfsetlinewidth{1.003750pt}%
\definecolor{currentstroke}{rgb}{0.372116,0.092816,0.499053}%
\pgfsetstrokecolor{currentstroke}%
\pgfsetstrokeopacity{0.800000}%
\pgfsetdash{}{0pt}%
\pgfpathmoveto{\pgfqpoint{1.445350in}{1.587511in}}%
\pgfpathcurveto{\pgfqpoint{1.451174in}{1.587511in}}{\pgfqpoint{1.456761in}{1.589825in}}{\pgfqpoint{1.460879in}{1.593943in}}%
\pgfpathcurveto{\pgfqpoint{1.464997in}{1.598061in}}{\pgfqpoint{1.467311in}{1.603648in}}{\pgfqpoint{1.467311in}{1.609472in}}%
\pgfpathcurveto{\pgfqpoint{1.467311in}{1.615296in}}{\pgfqpoint{1.464997in}{1.620882in}}{\pgfqpoint{1.460879in}{1.625000in}}%
\pgfpathcurveto{\pgfqpoint{1.456761in}{1.629118in}}{\pgfqpoint{1.451174in}{1.631432in}}{\pgfqpoint{1.445350in}{1.631432in}}%
\pgfpathcurveto{\pgfqpoint{1.439526in}{1.631432in}}{\pgfqpoint{1.433940in}{1.629118in}}{\pgfqpoint{1.429822in}{1.625000in}}%
\pgfpathcurveto{\pgfqpoint{1.425704in}{1.620882in}}{\pgfqpoint{1.423390in}{1.615296in}}{\pgfqpoint{1.423390in}{1.609472in}}%
\pgfpathcurveto{\pgfqpoint{1.423390in}{1.603648in}}{\pgfqpoint{1.425704in}{1.598061in}}{\pgfqpoint{1.429822in}{1.593943in}}%
\pgfpathcurveto{\pgfqpoint{1.433940in}{1.589825in}}{\pgfqpoint{1.439526in}{1.587511in}}{\pgfqpoint{1.445350in}{1.587511in}}%
\pgfpathlineto{\pgfqpoint{1.445350in}{1.587511in}}%
\pgfpathclose%
\pgfusepath{stroke,fill}%
\end{pgfscope}%
\begin{pgfscope}%
\pgfpathrectangle{\pgfqpoint{0.100000in}{0.209042in}}{\pgfqpoint{2.906250in}{2.906250in}}%
\pgfusepath{clip}%
\pgfsetbuttcap%
\pgfsetroundjoin%
\definecolor{currentfill}{rgb}{0.632805,0.187893,0.495332}%
\pgfsetfillcolor{currentfill}%
\pgfsetfillopacity{0.800000}%
\pgfsetlinewidth{1.003750pt}%
\definecolor{currentstroke}{rgb}{0.632805,0.187893,0.495332}%
\pgfsetstrokecolor{currentstroke}%
\pgfsetstrokeopacity{0.800000}%
\pgfsetdash{}{0pt}%
\pgfpathmoveto{\pgfqpoint{1.266767in}{2.000551in}}%
\pgfpathcurveto{\pgfqpoint{1.272591in}{2.000551in}}{\pgfqpoint{1.278177in}{2.002864in}}{\pgfqpoint{1.282295in}{2.006983in}}%
\pgfpathcurveto{\pgfqpoint{1.286413in}{2.011101in}}{\pgfqpoint{1.288727in}{2.016687in}}{\pgfqpoint{1.288727in}{2.022511in}}%
\pgfpathcurveto{\pgfqpoint{1.288727in}{2.028335in}}{\pgfqpoint{1.286413in}{2.033921in}}{\pgfqpoint{1.282295in}{2.038039in}}%
\pgfpathcurveto{\pgfqpoint{1.278177in}{2.042157in}}{\pgfqpoint{1.272591in}{2.044471in}}{\pgfqpoint{1.266767in}{2.044471in}}%
\pgfpathcurveto{\pgfqpoint{1.260943in}{2.044471in}}{\pgfqpoint{1.255357in}{2.042157in}}{\pgfqpoint{1.251239in}{2.038039in}}%
\pgfpathcurveto{\pgfqpoint{1.247120in}{2.033921in}}{\pgfqpoint{1.244807in}{2.028335in}}{\pgfqpoint{1.244807in}{2.022511in}}%
\pgfpathcurveto{\pgfqpoint{1.244807in}{2.016687in}}{\pgfqpoint{1.247120in}{2.011101in}}{\pgfqpoint{1.251239in}{2.006983in}}%
\pgfpathcurveto{\pgfqpoint{1.255357in}{2.002864in}}{\pgfqpoint{1.260943in}{2.000551in}}{\pgfqpoint{1.266767in}{2.000551in}}%
\pgfpathlineto{\pgfqpoint{1.266767in}{2.000551in}}%
\pgfpathclose%
\pgfusepath{stroke,fill}%
\end{pgfscope}%
\begin{pgfscope}%
\pgfpathrectangle{\pgfqpoint{0.100000in}{0.209042in}}{\pgfqpoint{2.906250in}{2.906250in}}%
\pgfusepath{clip}%
\pgfsetbuttcap%
\pgfsetroundjoin%
\definecolor{currentfill}{rgb}{0.258857,0.059706,0.457710}%
\pgfsetfillcolor{currentfill}%
\pgfsetfillopacity{0.800000}%
\pgfsetlinewidth{1.003750pt}%
\definecolor{currentstroke}{rgb}{0.258857,0.059706,0.457710}%
\pgfsetstrokecolor{currentstroke}%
\pgfsetstrokeopacity{0.800000}%
\pgfsetdash{}{0pt}%
\pgfpathmoveto{\pgfqpoint{1.565969in}{1.833534in}}%
\pgfpathcurveto{\pgfqpoint{1.571793in}{1.833534in}}{\pgfqpoint{1.577379in}{1.835848in}}{\pgfqpoint{1.581497in}{1.839966in}}%
\pgfpathcurveto{\pgfqpoint{1.585615in}{1.844084in}}{\pgfqpoint{1.587929in}{1.849671in}}{\pgfqpoint{1.587929in}{1.855495in}}%
\pgfpathcurveto{\pgfqpoint{1.587929in}{1.861318in}}{\pgfqpoint{1.585615in}{1.866905in}}{\pgfqpoint{1.581497in}{1.871023in}}%
\pgfpathcurveto{\pgfqpoint{1.577379in}{1.875141in}}{\pgfqpoint{1.571793in}{1.877455in}}{\pgfqpoint{1.565969in}{1.877455in}}%
\pgfpathcurveto{\pgfqpoint{1.560145in}{1.877455in}}{\pgfqpoint{1.554559in}{1.875141in}}{\pgfqpoint{1.550441in}{1.871023in}}%
\pgfpathcurveto{\pgfqpoint{1.546322in}{1.866905in}}{\pgfqpoint{1.544008in}{1.861318in}}{\pgfqpoint{1.544008in}{1.855495in}}%
\pgfpathcurveto{\pgfqpoint{1.544008in}{1.849671in}}{\pgfqpoint{1.546322in}{1.844084in}}{\pgfqpoint{1.550441in}{1.839966in}}%
\pgfpathcurveto{\pgfqpoint{1.554559in}{1.835848in}}{\pgfqpoint{1.560145in}{1.833534in}}{\pgfqpoint{1.565969in}{1.833534in}}%
\pgfpathlineto{\pgfqpoint{1.565969in}{1.833534in}}%
\pgfpathclose%
\pgfusepath{stroke,fill}%
\end{pgfscope}%
\begin{pgfscope}%
\pgfpathrectangle{\pgfqpoint{0.100000in}{0.209042in}}{\pgfqpoint{2.906250in}{2.906250in}}%
\pgfusepath{clip}%
\pgfsetbuttcap%
\pgfsetroundjoin%
\definecolor{currentfill}{rgb}{0.329114,0.075972,0.489287}%
\pgfsetfillcolor{currentfill}%
\pgfsetfillopacity{0.800000}%
\pgfsetlinewidth{1.003750pt}%
\definecolor{currentstroke}{rgb}{0.329114,0.075972,0.489287}%
\pgfsetstrokecolor{currentstroke}%
\pgfsetstrokeopacity{0.800000}%
\pgfsetdash{}{0pt}%
\pgfpathmoveto{\pgfqpoint{1.691965in}{1.880841in}}%
\pgfpathcurveto{\pgfqpoint{1.697789in}{1.880841in}}{\pgfqpoint{1.703375in}{1.883154in}}{\pgfqpoint{1.707493in}{1.887273in}}%
\pgfpathcurveto{\pgfqpoint{1.711611in}{1.891391in}}{\pgfqpoint{1.713925in}{1.896977in}}{\pgfqpoint{1.713925in}{1.902801in}}%
\pgfpathcurveto{\pgfqpoint{1.713925in}{1.908625in}}{\pgfqpoint{1.711611in}{1.914211in}}{\pgfqpoint{1.707493in}{1.918329in}}%
\pgfpathcurveto{\pgfqpoint{1.703375in}{1.922447in}}{\pgfqpoint{1.697789in}{1.924761in}}{\pgfqpoint{1.691965in}{1.924761in}}%
\pgfpathcurveto{\pgfqpoint{1.686141in}{1.924761in}}{\pgfqpoint{1.680555in}{1.922447in}}{\pgfqpoint{1.676437in}{1.918329in}}%
\pgfpathcurveto{\pgfqpoint{1.672318in}{1.914211in}}{\pgfqpoint{1.670005in}{1.908625in}}{\pgfqpoint{1.670005in}{1.902801in}}%
\pgfpathcurveto{\pgfqpoint{1.670005in}{1.896977in}}{\pgfqpoint{1.672318in}{1.891391in}}{\pgfqpoint{1.676437in}{1.887273in}}%
\pgfpathcurveto{\pgfqpoint{1.680555in}{1.883154in}}{\pgfqpoint{1.686141in}{1.880841in}}{\pgfqpoint{1.691965in}{1.880841in}}%
\pgfpathlineto{\pgfqpoint{1.691965in}{1.880841in}}%
\pgfpathclose%
\pgfusepath{stroke,fill}%
\end{pgfscope}%
\begin{pgfscope}%
\pgfpathrectangle{\pgfqpoint{0.100000in}{0.209042in}}{\pgfqpoint{2.906250in}{2.906250in}}%
\pgfusepath{clip}%
\pgfsetbuttcap%
\pgfsetroundjoin%
\definecolor{currentfill}{rgb}{0.291366,0.064553,0.475462}%
\pgfsetfillcolor{currentfill}%
\pgfsetfillopacity{0.800000}%
\pgfsetlinewidth{1.003750pt}%
\definecolor{currentstroke}{rgb}{0.291366,0.064553,0.475462}%
\pgfsetstrokecolor{currentstroke}%
\pgfsetstrokeopacity{0.800000}%
\pgfsetdash{}{0pt}%
\pgfpathmoveto{\pgfqpoint{1.707388in}{1.837097in}}%
\pgfpathcurveto{\pgfqpoint{1.713212in}{1.837097in}}{\pgfqpoint{1.718798in}{1.839411in}}{\pgfqpoint{1.722916in}{1.843529in}}%
\pgfpathcurveto{\pgfqpoint{1.727035in}{1.847647in}}{\pgfqpoint{1.729348in}{1.853233in}}{\pgfqpoint{1.729348in}{1.859057in}}%
\pgfpathcurveto{\pgfqpoint{1.729348in}{1.864881in}}{\pgfqpoint{1.727035in}{1.870467in}}{\pgfqpoint{1.722916in}{1.874585in}}%
\pgfpathcurveto{\pgfqpoint{1.718798in}{1.878703in}}{\pgfqpoint{1.713212in}{1.881017in}}{\pgfqpoint{1.707388in}{1.881017in}}%
\pgfpathcurveto{\pgfqpoint{1.701564in}{1.881017in}}{\pgfqpoint{1.695978in}{1.878703in}}{\pgfqpoint{1.691860in}{1.874585in}}%
\pgfpathcurveto{\pgfqpoint{1.687742in}{1.870467in}}{\pgfqpoint{1.685428in}{1.864881in}}{\pgfqpoint{1.685428in}{1.859057in}}%
\pgfpathcurveto{\pgfqpoint{1.685428in}{1.853233in}}{\pgfqpoint{1.687742in}{1.847647in}}{\pgfqpoint{1.691860in}{1.843529in}}%
\pgfpathcurveto{\pgfqpoint{1.695978in}{1.839411in}}{\pgfqpoint{1.701564in}{1.837097in}}{\pgfqpoint{1.707388in}{1.837097in}}%
\pgfpathlineto{\pgfqpoint{1.707388in}{1.837097in}}%
\pgfpathclose%
\pgfusepath{stroke,fill}%
\end{pgfscope}%
\begin{pgfscope}%
\pgfpathrectangle{\pgfqpoint{0.100000in}{0.209042in}}{\pgfqpoint{2.906250in}{2.906250in}}%
\pgfusepath{clip}%
\pgfsetbuttcap%
\pgfsetroundjoin%
\definecolor{currentfill}{rgb}{0.716387,0.214982,0.475290}%
\pgfsetfillcolor{currentfill}%
\pgfsetfillopacity{0.800000}%
\pgfsetlinewidth{1.003750pt}%
\definecolor{currentstroke}{rgb}{0.716387,0.214982,0.475290}%
\pgfsetstrokecolor{currentstroke}%
\pgfsetstrokeopacity{0.800000}%
\pgfsetdash{}{0pt}%
\pgfpathmoveto{\pgfqpoint{1.520362in}{1.330768in}}%
\pgfpathcurveto{\pgfqpoint{1.526186in}{1.330768in}}{\pgfqpoint{1.531772in}{1.333082in}}{\pgfqpoint{1.535890in}{1.337200in}}%
\pgfpathcurveto{\pgfqpoint{1.540008in}{1.341319in}}{\pgfqpoint{1.542322in}{1.346905in}}{\pgfqpoint{1.542322in}{1.352729in}}%
\pgfpathcurveto{\pgfqpoint{1.542322in}{1.358553in}}{\pgfqpoint{1.540008in}{1.364139in}}{\pgfqpoint{1.535890in}{1.368257in}}%
\pgfpathcurveto{\pgfqpoint{1.531772in}{1.372375in}}{\pgfqpoint{1.526186in}{1.374689in}}{\pgfqpoint{1.520362in}{1.374689in}}%
\pgfpathcurveto{\pgfqpoint{1.514538in}{1.374689in}}{\pgfqpoint{1.508952in}{1.372375in}}{\pgfqpoint{1.504834in}{1.368257in}}%
\pgfpathcurveto{\pgfqpoint{1.500715in}{1.364139in}}{\pgfqpoint{1.498402in}{1.358553in}}{\pgfqpoint{1.498402in}{1.352729in}}%
\pgfpathcurveto{\pgfqpoint{1.498402in}{1.346905in}}{\pgfqpoint{1.500715in}{1.341319in}}{\pgfqpoint{1.504834in}{1.337200in}}%
\pgfpathcurveto{\pgfqpoint{1.508952in}{1.333082in}}{\pgfqpoint{1.514538in}{1.330768in}}{\pgfqpoint{1.520362in}{1.330768in}}%
\pgfpathlineto{\pgfqpoint{1.520362in}{1.330768in}}%
\pgfpathclose%
\pgfusepath{stroke,fill}%
\end{pgfscope}%
\begin{pgfscope}%
\pgfpathrectangle{\pgfqpoint{0.100000in}{0.209042in}}{\pgfqpoint{2.906250in}{2.906250in}}%
\pgfusepath{clip}%
\pgfsetbuttcap%
\pgfsetroundjoin%
\definecolor{currentfill}{rgb}{0.432967,0.117855,0.506160}%
\pgfsetfillcolor{currentfill}%
\pgfsetfillopacity{0.800000}%
\pgfsetlinewidth{1.003750pt}%
\definecolor{currentstroke}{rgb}{0.432967,0.117855,0.506160}%
\pgfsetstrokecolor{currentstroke}%
\pgfsetstrokeopacity{0.800000}%
\pgfsetdash{}{0pt}%
\pgfpathmoveto{\pgfqpoint{1.505788in}{1.505255in}}%
\pgfpathcurveto{\pgfqpoint{1.511612in}{1.505255in}}{\pgfqpoint{1.517198in}{1.507569in}}{\pgfqpoint{1.521316in}{1.511687in}}%
\pgfpathcurveto{\pgfqpoint{1.525434in}{1.515805in}}{\pgfqpoint{1.527748in}{1.521391in}}{\pgfqpoint{1.527748in}{1.527215in}}%
\pgfpathcurveto{\pgfqpoint{1.527748in}{1.533039in}}{\pgfqpoint{1.525434in}{1.538625in}}{\pgfqpoint{1.521316in}{1.542744in}}%
\pgfpathcurveto{\pgfqpoint{1.517198in}{1.546862in}}{\pgfqpoint{1.511612in}{1.549176in}}{\pgfqpoint{1.505788in}{1.549176in}}%
\pgfpathcurveto{\pgfqpoint{1.499964in}{1.549176in}}{\pgfqpoint{1.494378in}{1.546862in}}{\pgfqpoint{1.490259in}{1.542744in}}%
\pgfpathcurveto{\pgfqpoint{1.486141in}{1.538625in}}{\pgfqpoint{1.483827in}{1.533039in}}{\pgfqpoint{1.483827in}{1.527215in}}%
\pgfpathcurveto{\pgfqpoint{1.483827in}{1.521391in}}{\pgfqpoint{1.486141in}{1.515805in}}{\pgfqpoint{1.490259in}{1.511687in}}%
\pgfpathcurveto{\pgfqpoint{1.494378in}{1.507569in}}{\pgfqpoint{1.499964in}{1.505255in}}{\pgfqpoint{1.505788in}{1.505255in}}%
\pgfpathlineto{\pgfqpoint{1.505788in}{1.505255in}}%
\pgfpathclose%
\pgfusepath{stroke,fill}%
\end{pgfscope}%
\begin{pgfscope}%
\pgfpathrectangle{\pgfqpoint{0.100000in}{0.209042in}}{\pgfqpoint{2.906250in}{2.906250in}}%
\pgfusepath{clip}%
\pgfsetbuttcap%
\pgfsetroundjoin%
\definecolor{currentfill}{rgb}{0.198177,0.063862,0.404009}%
\pgfsetfillcolor{currentfill}%
\pgfsetfillopacity{0.800000}%
\pgfsetlinewidth{1.003750pt}%
\definecolor{currentstroke}{rgb}{0.198177,0.063862,0.404009}%
\pgfsetstrokecolor{currentstroke}%
\pgfsetstrokeopacity{0.800000}%
\pgfsetdash{}{0pt}%
\pgfpathmoveto{\pgfqpoint{1.583102in}{1.731679in}}%
\pgfpathcurveto{\pgfqpoint{1.588926in}{1.731679in}}{\pgfqpoint{1.594512in}{1.733992in}}{\pgfqpoint{1.598630in}{1.738111in}}%
\pgfpathcurveto{\pgfqpoint{1.602749in}{1.742229in}}{\pgfqpoint{1.605062in}{1.747815in}}{\pgfqpoint{1.605062in}{1.753639in}}%
\pgfpathcurveto{\pgfqpoint{1.605062in}{1.759463in}}{\pgfqpoint{1.602749in}{1.765049in}}{\pgfqpoint{1.598630in}{1.769167in}}%
\pgfpathcurveto{\pgfqpoint{1.594512in}{1.773285in}}{\pgfqpoint{1.588926in}{1.775599in}}{\pgfqpoint{1.583102in}{1.775599in}}%
\pgfpathcurveto{\pgfqpoint{1.577278in}{1.775599in}}{\pgfqpoint{1.571692in}{1.773285in}}{\pgfqpoint{1.567574in}{1.769167in}}%
\pgfpathcurveto{\pgfqpoint{1.563456in}{1.765049in}}{\pgfqpoint{1.561142in}{1.759463in}}{\pgfqpoint{1.561142in}{1.753639in}}%
\pgfpathcurveto{\pgfqpoint{1.561142in}{1.747815in}}{\pgfqpoint{1.563456in}{1.742229in}}{\pgfqpoint{1.567574in}{1.738111in}}%
\pgfpathcurveto{\pgfqpoint{1.571692in}{1.733992in}}{\pgfqpoint{1.577278in}{1.731679in}}{\pgfqpoint{1.583102in}{1.731679in}}%
\pgfpathlineto{\pgfqpoint{1.583102in}{1.731679in}}%
\pgfpathclose%
\pgfusepath{stroke,fill}%
\end{pgfscope}%
\begin{pgfscope}%
\pgfpathrectangle{\pgfqpoint{0.100000in}{0.209042in}}{\pgfqpoint{2.906250in}{2.906250in}}%
\pgfusepath{clip}%
\pgfsetbuttcap%
\pgfsetroundjoin%
\definecolor{currentfill}{rgb}{0.218512,0.061158,0.425392}%
\pgfsetfillcolor{currentfill}%
\pgfsetfillopacity{0.800000}%
\pgfsetlinewidth{1.003750pt}%
\definecolor{currentstroke}{rgb}{0.218512,0.061158,0.425392}%
\pgfsetstrokecolor{currentstroke}%
\pgfsetstrokeopacity{0.800000}%
\pgfsetdash{}{0pt}%
\pgfpathmoveto{\pgfqpoint{1.547074in}{1.689603in}}%
\pgfpathcurveto{\pgfqpoint{1.552898in}{1.689603in}}{\pgfqpoint{1.558484in}{1.691917in}}{\pgfqpoint{1.562603in}{1.696035in}}%
\pgfpathcurveto{\pgfqpoint{1.566721in}{1.700153in}}{\pgfqpoint{1.569035in}{1.705739in}}{\pgfqpoint{1.569035in}{1.711563in}}%
\pgfpathcurveto{\pgfqpoint{1.569035in}{1.717387in}}{\pgfqpoint{1.566721in}{1.722973in}}{\pgfqpoint{1.562603in}{1.727092in}}%
\pgfpathcurveto{\pgfqpoint{1.558484in}{1.731210in}}{\pgfqpoint{1.552898in}{1.733524in}}{\pgfqpoint{1.547074in}{1.733524in}}%
\pgfpathcurveto{\pgfqpoint{1.541250in}{1.733524in}}{\pgfqpoint{1.535664in}{1.731210in}}{\pgfqpoint{1.531546in}{1.727092in}}%
\pgfpathcurveto{\pgfqpoint{1.527428in}{1.722973in}}{\pgfqpoint{1.525114in}{1.717387in}}{\pgfqpoint{1.525114in}{1.711563in}}%
\pgfpathcurveto{\pgfqpoint{1.525114in}{1.705739in}}{\pgfqpoint{1.527428in}{1.700153in}}{\pgfqpoint{1.531546in}{1.696035in}}%
\pgfpathcurveto{\pgfqpoint{1.535664in}{1.691917in}}{\pgfqpoint{1.541250in}{1.689603in}}{\pgfqpoint{1.547074in}{1.689603in}}%
\pgfpathlineto{\pgfqpoint{1.547074in}{1.689603in}}%
\pgfpathclose%
\pgfusepath{stroke,fill}%
\end{pgfscope}%
\begin{pgfscope}%
\pgfpathrectangle{\pgfqpoint{0.100000in}{0.209042in}}{\pgfqpoint{2.906250in}{2.906250in}}%
\pgfusepath{clip}%
\pgfsetbuttcap%
\pgfsetroundjoin%
\definecolor{currentfill}{rgb}{0.384299,0.097855,0.501002}%
\pgfsetfillcolor{currentfill}%
\pgfsetfillopacity{0.800000}%
\pgfsetlinewidth{1.003750pt}%
\definecolor{currentstroke}{rgb}{0.384299,0.097855,0.501002}%
\pgfsetstrokecolor{currentstroke}%
\pgfsetstrokeopacity{0.800000}%
\pgfsetdash{}{0pt}%
\pgfpathmoveto{\pgfqpoint{1.841971in}{1.637267in}}%
\pgfpathcurveto{\pgfqpoint{1.847795in}{1.637267in}}{\pgfqpoint{1.853381in}{1.639580in}}{\pgfqpoint{1.857499in}{1.643699in}}%
\pgfpathcurveto{\pgfqpoint{1.861617in}{1.647817in}}{\pgfqpoint{1.863931in}{1.653403in}}{\pgfqpoint{1.863931in}{1.659227in}}%
\pgfpathcurveto{\pgfqpoint{1.863931in}{1.665051in}}{\pgfqpoint{1.861617in}{1.670637in}}{\pgfqpoint{1.857499in}{1.674755in}}%
\pgfpathcurveto{\pgfqpoint{1.853381in}{1.678873in}}{\pgfqpoint{1.847795in}{1.681187in}}{\pgfqpoint{1.841971in}{1.681187in}}%
\pgfpathcurveto{\pgfqpoint{1.836147in}{1.681187in}}{\pgfqpoint{1.830561in}{1.678873in}}{\pgfqpoint{1.826443in}{1.674755in}}%
\pgfpathcurveto{\pgfqpoint{1.822325in}{1.670637in}}{\pgfqpoint{1.820011in}{1.665051in}}{\pgfqpoint{1.820011in}{1.659227in}}%
\pgfpathcurveto{\pgfqpoint{1.820011in}{1.653403in}}{\pgfqpoint{1.822325in}{1.647817in}}{\pgfqpoint{1.826443in}{1.643699in}}%
\pgfpathcurveto{\pgfqpoint{1.830561in}{1.639580in}}{\pgfqpoint{1.836147in}{1.637267in}}{\pgfqpoint{1.841971in}{1.637267in}}%
\pgfpathlineto{\pgfqpoint{1.841971in}{1.637267in}}%
\pgfpathclose%
\pgfusepath{stroke,fill}%
\end{pgfscope}%
\begin{pgfscope}%
\pgfpathrectangle{\pgfqpoint{0.100000in}{0.209042in}}{\pgfqpoint{2.906250in}{2.906250in}}%
\pgfusepath{clip}%
\pgfsetbuttcap%
\pgfsetroundjoin%
\definecolor{currentfill}{rgb}{0.329114,0.075972,0.489287}%
\pgfsetfillcolor{currentfill}%
\pgfsetfillopacity{0.800000}%
\pgfsetlinewidth{1.003750pt}%
\definecolor{currentstroke}{rgb}{0.329114,0.075972,0.489287}%
\pgfsetstrokecolor{currentstroke}%
\pgfsetstrokeopacity{0.800000}%
\pgfsetdash{}{0pt}%
\pgfpathmoveto{\pgfqpoint{1.609339in}{1.560557in}}%
\pgfpathcurveto{\pgfqpoint{1.615163in}{1.560557in}}{\pgfqpoint{1.620749in}{1.562871in}}{\pgfqpoint{1.624867in}{1.566989in}}%
\pgfpathcurveto{\pgfqpoint{1.628986in}{1.571107in}}{\pgfqpoint{1.631299in}{1.576693in}}{\pgfqpoint{1.631299in}{1.582517in}}%
\pgfpathcurveto{\pgfqpoint{1.631299in}{1.588341in}}{\pgfqpoint{1.628986in}{1.593927in}}{\pgfqpoint{1.624867in}{1.598046in}}%
\pgfpathcurveto{\pgfqpoint{1.620749in}{1.602164in}}{\pgfqpoint{1.615163in}{1.604478in}}{\pgfqpoint{1.609339in}{1.604478in}}%
\pgfpathcurveto{\pgfqpoint{1.603515in}{1.604478in}}{\pgfqpoint{1.597929in}{1.602164in}}{\pgfqpoint{1.593811in}{1.598046in}}%
\pgfpathcurveto{\pgfqpoint{1.589693in}{1.593927in}}{\pgfqpoint{1.587379in}{1.588341in}}{\pgfqpoint{1.587379in}{1.582517in}}%
\pgfpathcurveto{\pgfqpoint{1.587379in}{1.576693in}}{\pgfqpoint{1.589693in}{1.571107in}}{\pgfqpoint{1.593811in}{1.566989in}}%
\pgfpathcurveto{\pgfqpoint{1.597929in}{1.562871in}}{\pgfqpoint{1.603515in}{1.560557in}}{\pgfqpoint{1.609339in}{1.560557in}}%
\pgfpathlineto{\pgfqpoint{1.609339in}{1.560557in}}%
\pgfpathclose%
\pgfusepath{stroke,fill}%
\end{pgfscope}%
\begin{pgfscope}%
\pgfpathrectangle{\pgfqpoint{0.100000in}{0.209042in}}{\pgfqpoint{2.906250in}{2.906250in}}%
\pgfusepath{clip}%
\pgfsetbuttcap%
\pgfsetroundjoin%
\definecolor{currentfill}{rgb}{0.512831,0.148179,0.507648}%
\pgfsetfillcolor{currentfill}%
\pgfsetfillopacity{0.800000}%
\pgfsetlinewidth{1.003750pt}%
\definecolor{currentstroke}{rgb}{0.512831,0.148179,0.507648}%
\pgfsetstrokecolor{currentstroke}%
\pgfsetstrokeopacity{0.800000}%
\pgfsetdash{}{0pt}%
\pgfpathmoveto{\pgfqpoint{1.236003in}{1.703705in}}%
\pgfpathcurveto{\pgfqpoint{1.241827in}{1.703705in}}{\pgfqpoint{1.247414in}{1.706018in}}{\pgfqpoint{1.251532in}{1.710137in}}%
\pgfpathcurveto{\pgfqpoint{1.255650in}{1.714255in}}{\pgfqpoint{1.257964in}{1.719841in}}{\pgfqpoint{1.257964in}{1.725665in}}%
\pgfpathcurveto{\pgfqpoint{1.257964in}{1.731489in}}{\pgfqpoint{1.255650in}{1.737075in}}{\pgfqpoint{1.251532in}{1.741193in}}%
\pgfpathcurveto{\pgfqpoint{1.247414in}{1.745311in}}{\pgfqpoint{1.241827in}{1.747625in}}{\pgfqpoint{1.236003in}{1.747625in}}%
\pgfpathcurveto{\pgfqpoint{1.230180in}{1.747625in}}{\pgfqpoint{1.224593in}{1.745311in}}{\pgfqpoint{1.220475in}{1.741193in}}%
\pgfpathcurveto{\pgfqpoint{1.216357in}{1.737075in}}{\pgfqpoint{1.214043in}{1.731489in}}{\pgfqpoint{1.214043in}{1.725665in}}%
\pgfpathcurveto{\pgfqpoint{1.214043in}{1.719841in}}{\pgfqpoint{1.216357in}{1.714255in}}{\pgfqpoint{1.220475in}{1.710137in}}%
\pgfpathcurveto{\pgfqpoint{1.224593in}{1.706018in}}{\pgfqpoint{1.230180in}{1.703705in}}{\pgfqpoint{1.236003in}{1.703705in}}%
\pgfpathlineto{\pgfqpoint{1.236003in}{1.703705in}}%
\pgfpathclose%
\pgfusepath{stroke,fill}%
\end{pgfscope}%
\begin{pgfscope}%
\pgfpathrectangle{\pgfqpoint{0.100000in}{0.209042in}}{\pgfqpoint{2.906250in}{2.906250in}}%
\pgfusepath{clip}%
\pgfsetbuttcap%
\pgfsetroundjoin%
\definecolor{currentfill}{rgb}{0.258857,0.059706,0.457710}%
\pgfsetfillcolor{currentfill}%
\pgfsetfillopacity{0.800000}%
\pgfsetlinewidth{1.003750pt}%
\definecolor{currentstroke}{rgb}{0.258857,0.059706,0.457710}%
\pgfsetstrokecolor{currentstroke}%
\pgfsetstrokeopacity{0.800000}%
\pgfsetdash{}{0pt}%
\pgfpathmoveto{\pgfqpoint{1.603893in}{1.616252in}}%
\pgfpathcurveto{\pgfqpoint{1.609717in}{1.616252in}}{\pgfqpoint{1.615303in}{1.618566in}}{\pgfqpoint{1.619422in}{1.622684in}}%
\pgfpathcurveto{\pgfqpoint{1.623540in}{1.626802in}}{\pgfqpoint{1.625854in}{1.632388in}}{\pgfqpoint{1.625854in}{1.638212in}}%
\pgfpathcurveto{\pgfqpoint{1.625854in}{1.644036in}}{\pgfqpoint{1.623540in}{1.649622in}}{\pgfqpoint{1.619422in}{1.653740in}}%
\pgfpathcurveto{\pgfqpoint{1.615303in}{1.657858in}}{\pgfqpoint{1.609717in}{1.660172in}}{\pgfqpoint{1.603893in}{1.660172in}}%
\pgfpathcurveto{\pgfqpoint{1.598069in}{1.660172in}}{\pgfqpoint{1.592483in}{1.657858in}}{\pgfqpoint{1.588365in}{1.653740in}}%
\pgfpathcurveto{\pgfqpoint{1.584247in}{1.649622in}}{\pgfqpoint{1.581933in}{1.644036in}}{\pgfqpoint{1.581933in}{1.638212in}}%
\pgfpathcurveto{\pgfqpoint{1.581933in}{1.632388in}}{\pgfqpoint{1.584247in}{1.626802in}}{\pgfqpoint{1.588365in}{1.622684in}}%
\pgfpathcurveto{\pgfqpoint{1.592483in}{1.618566in}}{\pgfqpoint{1.598069in}{1.616252in}}{\pgfqpoint{1.603893in}{1.616252in}}%
\pgfpathlineto{\pgfqpoint{1.603893in}{1.616252in}}%
\pgfpathclose%
\pgfusepath{stroke,fill}%
\end{pgfscope}%
\begin{pgfscope}%
\pgfpathrectangle{\pgfqpoint{0.100000in}{0.209042in}}{\pgfqpoint{2.906250in}{2.906250in}}%
\pgfusepath{clip}%
\pgfsetbuttcap%
\pgfsetroundjoin%
\definecolor{currentfill}{rgb}{0.697098,0.208501,0.480835}%
\pgfsetfillcolor{currentfill}%
\pgfsetfillopacity{0.800000}%
\pgfsetlinewidth{1.003750pt}%
\definecolor{currentstroke}{rgb}{0.697098,0.208501,0.480835}%
\pgfsetstrokecolor{currentstroke}%
\pgfsetstrokeopacity{0.800000}%
\pgfsetdash{}{0pt}%
\pgfpathmoveto{\pgfqpoint{1.095695in}{1.768030in}}%
\pgfpathcurveto{\pgfqpoint{1.101519in}{1.768030in}}{\pgfqpoint{1.107105in}{1.770344in}}{\pgfqpoint{1.111223in}{1.774462in}}%
\pgfpathcurveto{\pgfqpoint{1.115342in}{1.778580in}}{\pgfqpoint{1.117655in}{1.784166in}}{\pgfqpoint{1.117655in}{1.789990in}}%
\pgfpathcurveto{\pgfqpoint{1.117655in}{1.795814in}}{\pgfqpoint{1.115342in}{1.801400in}}{\pgfqpoint{1.111223in}{1.805518in}}%
\pgfpathcurveto{\pgfqpoint{1.107105in}{1.809637in}}{\pgfqpoint{1.101519in}{1.811950in}}{\pgfqpoint{1.095695in}{1.811950in}}%
\pgfpathcurveto{\pgfqpoint{1.089871in}{1.811950in}}{\pgfqpoint{1.084285in}{1.809637in}}{\pgfqpoint{1.080167in}{1.805518in}}%
\pgfpathcurveto{\pgfqpoint{1.076049in}{1.801400in}}{\pgfqpoint{1.073735in}{1.795814in}}{\pgfqpoint{1.073735in}{1.789990in}}%
\pgfpathcurveto{\pgfqpoint{1.073735in}{1.784166in}}{\pgfqpoint{1.076049in}{1.778580in}}{\pgfqpoint{1.080167in}{1.774462in}}%
\pgfpathcurveto{\pgfqpoint{1.084285in}{1.770344in}}{\pgfqpoint{1.089871in}{1.768030in}}{\pgfqpoint{1.095695in}{1.768030in}}%
\pgfpathlineto{\pgfqpoint{1.095695in}{1.768030in}}%
\pgfpathclose%
\pgfusepath{stroke,fill}%
\end{pgfscope}%
\begin{pgfscope}%
\pgfpathrectangle{\pgfqpoint{0.100000in}{0.209042in}}{\pgfqpoint{2.906250in}{2.906250in}}%
\pgfusepath{clip}%
\pgfsetbuttcap%
\pgfsetroundjoin%
\definecolor{currentfill}{rgb}{0.626401,0.185867,0.496456}%
\pgfsetfillcolor{currentfill}%
\pgfsetfillopacity{0.800000}%
\pgfsetlinewidth{1.003750pt}%
\definecolor{currentstroke}{rgb}{0.626401,0.185867,0.496456}%
\pgfsetstrokecolor{currentstroke}%
\pgfsetstrokeopacity{0.800000}%
\pgfsetdash{}{0pt}%
\pgfpathmoveto{\pgfqpoint{1.814346in}{2.050252in}}%
\pgfpathcurveto{\pgfqpoint{1.820170in}{2.050252in}}{\pgfqpoint{1.825756in}{2.052566in}}{\pgfqpoint{1.829874in}{2.056684in}}%
\pgfpathcurveto{\pgfqpoint{1.833992in}{2.060802in}}{\pgfqpoint{1.836306in}{2.066388in}}{\pgfqpoint{1.836306in}{2.072212in}}%
\pgfpathcurveto{\pgfqpoint{1.836306in}{2.078036in}}{\pgfqpoint{1.833992in}{2.083622in}}{\pgfqpoint{1.829874in}{2.087740in}}%
\pgfpathcurveto{\pgfqpoint{1.825756in}{2.091859in}}{\pgfqpoint{1.820170in}{2.094173in}}{\pgfqpoint{1.814346in}{2.094173in}}%
\pgfpathcurveto{\pgfqpoint{1.808522in}{2.094173in}}{\pgfqpoint{1.802936in}{2.091859in}}{\pgfqpoint{1.798817in}{2.087740in}}%
\pgfpathcurveto{\pgfqpoint{1.794699in}{2.083622in}}{\pgfqpoint{1.792385in}{2.078036in}}{\pgfqpoint{1.792385in}{2.072212in}}%
\pgfpathcurveto{\pgfqpoint{1.792385in}{2.066388in}}{\pgfqpoint{1.794699in}{2.060802in}}{\pgfqpoint{1.798817in}{2.056684in}}%
\pgfpathcurveto{\pgfqpoint{1.802936in}{2.052566in}}{\pgfqpoint{1.808522in}{2.050252in}}{\pgfqpoint{1.814346in}{2.050252in}}%
\pgfpathlineto{\pgfqpoint{1.814346in}{2.050252in}}%
\pgfpathclose%
\pgfusepath{stroke,fill}%
\end{pgfscope}%
\begin{pgfscope}%
\pgfpathrectangle{\pgfqpoint{0.100000in}{0.209042in}}{\pgfqpoint{2.906250in}{2.906250in}}%
\pgfusepath{clip}%
\pgfsetbuttcap%
\pgfsetroundjoin%
\definecolor{currentfill}{rgb}{0.481929,0.136891,0.507989}%
\pgfsetfillcolor{currentfill}%
\pgfsetfillopacity{0.800000}%
\pgfsetlinewidth{1.003750pt}%
\definecolor{currentstroke}{rgb}{0.481929,0.136891,0.507989}%
\pgfsetstrokecolor{currentstroke}%
\pgfsetstrokeopacity{0.800000}%
\pgfsetdash{}{0pt}%
\pgfpathmoveto{\pgfqpoint{1.663058in}{1.450174in}}%
\pgfpathcurveto{\pgfqpoint{1.668882in}{1.450174in}}{\pgfqpoint{1.674468in}{1.452488in}}{\pgfqpoint{1.678587in}{1.456606in}}%
\pgfpathcurveto{\pgfqpoint{1.682705in}{1.460724in}}{\pgfqpoint{1.685019in}{1.466311in}}{\pgfqpoint{1.685019in}{1.472135in}}%
\pgfpathcurveto{\pgfqpoint{1.685019in}{1.477958in}}{\pgfqpoint{1.682705in}{1.483545in}}{\pgfqpoint{1.678587in}{1.487663in}}%
\pgfpathcurveto{\pgfqpoint{1.674468in}{1.491781in}}{\pgfqpoint{1.668882in}{1.494095in}}{\pgfqpoint{1.663058in}{1.494095in}}%
\pgfpathcurveto{\pgfqpoint{1.657234in}{1.494095in}}{\pgfqpoint{1.651648in}{1.491781in}}{\pgfqpoint{1.647530in}{1.487663in}}%
\pgfpathcurveto{\pgfqpoint{1.643412in}{1.483545in}}{\pgfqpoint{1.641098in}{1.477958in}}{\pgfqpoint{1.641098in}{1.472135in}}%
\pgfpathcurveto{\pgfqpoint{1.641098in}{1.466311in}}{\pgfqpoint{1.643412in}{1.460724in}}{\pgfqpoint{1.647530in}{1.456606in}}%
\pgfpathcurveto{\pgfqpoint{1.651648in}{1.452488in}}{\pgfqpoint{1.657234in}{1.450174in}}{\pgfqpoint{1.663058in}{1.450174in}}%
\pgfpathlineto{\pgfqpoint{1.663058in}{1.450174in}}%
\pgfpathclose%
\pgfusepath{stroke,fill}%
\end{pgfscope}%
\begin{pgfscope}%
\pgfpathrectangle{\pgfqpoint{0.100000in}{0.209042in}}{\pgfqpoint{2.906250in}{2.906250in}}%
\pgfusepath{clip}%
\pgfsetbuttcap%
\pgfsetroundjoin%
\definecolor{currentfill}{rgb}{0.537755,0.156894,0.506551}%
\pgfsetfillcolor{currentfill}%
\pgfsetfillopacity{0.800000}%
\pgfsetlinewidth{1.003750pt}%
\definecolor{currentstroke}{rgb}{0.537755,0.156894,0.506551}%
\pgfsetstrokecolor{currentstroke}%
\pgfsetstrokeopacity{0.800000}%
\pgfsetdash{}{0pt}%
\pgfpathmoveto{\pgfqpoint{1.296506in}{1.950616in}}%
\pgfpathcurveto{\pgfqpoint{1.302330in}{1.950616in}}{\pgfqpoint{1.307917in}{1.952930in}}{\pgfqpoint{1.312035in}{1.957048in}}%
\pgfpathcurveto{\pgfqpoint{1.316153in}{1.961166in}}{\pgfqpoint{1.318467in}{1.966753in}}{\pgfqpoint{1.318467in}{1.972576in}}%
\pgfpathcurveto{\pgfqpoint{1.318467in}{1.978400in}}{\pgfqpoint{1.316153in}{1.983987in}}{\pgfqpoint{1.312035in}{1.988105in}}%
\pgfpathcurveto{\pgfqpoint{1.307917in}{1.992223in}}{\pgfqpoint{1.302330in}{1.994537in}}{\pgfqpoint{1.296506in}{1.994537in}}%
\pgfpathcurveto{\pgfqpoint{1.290683in}{1.994537in}}{\pgfqpoint{1.285096in}{1.992223in}}{\pgfqpoint{1.280978in}{1.988105in}}%
\pgfpathcurveto{\pgfqpoint{1.276860in}{1.983987in}}{\pgfqpoint{1.274546in}{1.978400in}}{\pgfqpoint{1.274546in}{1.972576in}}%
\pgfpathcurveto{\pgfqpoint{1.274546in}{1.966753in}}{\pgfqpoint{1.276860in}{1.961166in}}{\pgfqpoint{1.280978in}{1.957048in}}%
\pgfpathcurveto{\pgfqpoint{1.285096in}{1.952930in}}{\pgfqpoint{1.290683in}{1.950616in}}{\pgfqpoint{1.296506in}{1.950616in}}%
\pgfpathlineto{\pgfqpoint{1.296506in}{1.950616in}}%
\pgfpathclose%
\pgfusepath{stroke,fill}%
\end{pgfscope}%
\begin{pgfscope}%
\pgfpathrectangle{\pgfqpoint{0.100000in}{0.209042in}}{\pgfqpoint{2.906250in}{2.906250in}}%
\pgfusepath{clip}%
\pgfsetbuttcap%
\pgfsetroundjoin%
\definecolor{currentfill}{rgb}{0.218512,0.061158,0.425392}%
\pgfsetfillcolor{currentfill}%
\pgfsetfillopacity{0.800000}%
\pgfsetlinewidth{1.003750pt}%
\definecolor{currentstroke}{rgb}{0.218512,0.061158,0.425392}%
\pgfsetstrokecolor{currentstroke}%
\pgfsetstrokeopacity{0.800000}%
\pgfsetdash{}{0pt}%
\pgfpathmoveto{\pgfqpoint{1.512317in}{1.669812in}}%
\pgfpathcurveto{\pgfqpoint{1.518141in}{1.669812in}}{\pgfqpoint{1.523727in}{1.672126in}}{\pgfqpoint{1.527845in}{1.676244in}}%
\pgfpathcurveto{\pgfqpoint{1.531963in}{1.680362in}}{\pgfqpoint{1.534277in}{1.685948in}}{\pgfqpoint{1.534277in}{1.691772in}}%
\pgfpathcurveto{\pgfqpoint{1.534277in}{1.697596in}}{\pgfqpoint{1.531963in}{1.703182in}}{\pgfqpoint{1.527845in}{1.707300in}}%
\pgfpathcurveto{\pgfqpoint{1.523727in}{1.711418in}}{\pgfqpoint{1.518141in}{1.713732in}}{\pgfqpoint{1.512317in}{1.713732in}}%
\pgfpathcurveto{\pgfqpoint{1.506493in}{1.713732in}}{\pgfqpoint{1.500907in}{1.711418in}}{\pgfqpoint{1.496789in}{1.707300in}}%
\pgfpathcurveto{\pgfqpoint{1.492671in}{1.703182in}}{\pgfqpoint{1.490357in}{1.697596in}}{\pgfqpoint{1.490357in}{1.691772in}}%
\pgfpathcurveto{\pgfqpoint{1.490357in}{1.685948in}}{\pgfqpoint{1.492671in}{1.680362in}}{\pgfqpoint{1.496789in}{1.676244in}}%
\pgfpathcurveto{\pgfqpoint{1.500907in}{1.672126in}}{\pgfqpoint{1.506493in}{1.669812in}}{\pgfqpoint{1.512317in}{1.669812in}}%
\pgfpathlineto{\pgfqpoint{1.512317in}{1.669812in}}%
\pgfpathclose%
\pgfusepath{stroke,fill}%
\end{pgfscope}%
\begin{pgfscope}%
\pgfpathrectangle{\pgfqpoint{0.100000in}{0.209042in}}{\pgfqpoint{2.906250in}{2.906250in}}%
\pgfusepath{clip}%
\pgfsetbuttcap%
\pgfsetroundjoin%
\definecolor{currentfill}{rgb}{0.265447,0.060237,0.461840}%
\pgfsetfillcolor{currentfill}%
\pgfsetfillopacity{0.800000}%
\pgfsetlinewidth{1.003750pt}%
\definecolor{currentstroke}{rgb}{0.265447,0.060237,0.461840}%
\pgfsetstrokecolor{currentstroke}%
\pgfsetstrokeopacity{0.800000}%
\pgfsetdash{}{0pt}%
\pgfpathmoveto{\pgfqpoint{1.758229in}{1.669916in}}%
\pgfpathcurveto{\pgfqpoint{1.764053in}{1.669916in}}{\pgfqpoint{1.769639in}{1.672229in}}{\pgfqpoint{1.773758in}{1.676348in}}%
\pgfpathcurveto{\pgfqpoint{1.777876in}{1.680466in}}{\pgfqpoint{1.780190in}{1.686052in}}{\pgfqpoint{1.780190in}{1.691876in}}%
\pgfpathcurveto{\pgfqpoint{1.780190in}{1.697700in}}{\pgfqpoint{1.777876in}{1.703286in}}{\pgfqpoint{1.773758in}{1.707404in}}%
\pgfpathcurveto{\pgfqpoint{1.769639in}{1.711522in}}{\pgfqpoint{1.764053in}{1.713836in}}{\pgfqpoint{1.758229in}{1.713836in}}%
\pgfpathcurveto{\pgfqpoint{1.752405in}{1.713836in}}{\pgfqpoint{1.746819in}{1.711522in}}{\pgfqpoint{1.742701in}{1.707404in}}%
\pgfpathcurveto{\pgfqpoint{1.738583in}{1.703286in}}{\pgfqpoint{1.736269in}{1.697700in}}{\pgfqpoint{1.736269in}{1.691876in}}%
\pgfpathcurveto{\pgfqpoint{1.736269in}{1.686052in}}{\pgfqpoint{1.738583in}{1.680466in}}{\pgfqpoint{1.742701in}{1.676348in}}%
\pgfpathcurveto{\pgfqpoint{1.746819in}{1.672229in}}{\pgfqpoint{1.752405in}{1.669916in}}{\pgfqpoint{1.758229in}{1.669916in}}%
\pgfpathlineto{\pgfqpoint{1.758229in}{1.669916in}}%
\pgfpathclose%
\pgfusepath{stroke,fill}%
\end{pgfscope}%
\begin{pgfscope}%
\pgfpathrectangle{\pgfqpoint{0.100000in}{0.209042in}}{\pgfqpoint{2.906250in}{2.906250in}}%
\pgfusepath{clip}%
\pgfsetbuttcap%
\pgfsetroundjoin%
\definecolor{currentfill}{rgb}{0.265447,0.060237,0.461840}%
\pgfsetfillcolor{currentfill}%
\pgfsetfillopacity{0.800000}%
\pgfsetlinewidth{1.003750pt}%
\definecolor{currentstroke}{rgb}{0.265447,0.060237,0.461840}%
\pgfsetstrokecolor{currentstroke}%
\pgfsetstrokeopacity{0.800000}%
\pgfsetdash{}{0pt}%
\pgfpathmoveto{\pgfqpoint{1.746210in}{1.647685in}}%
\pgfpathcurveto{\pgfqpoint{1.752034in}{1.647685in}}{\pgfqpoint{1.757620in}{1.649999in}}{\pgfqpoint{1.761739in}{1.654117in}}%
\pgfpathcurveto{\pgfqpoint{1.765857in}{1.658235in}}{\pgfqpoint{1.768171in}{1.663822in}}{\pgfqpoint{1.768171in}{1.669646in}}%
\pgfpathcurveto{\pgfqpoint{1.768171in}{1.675469in}}{\pgfqpoint{1.765857in}{1.681056in}}{\pgfqpoint{1.761739in}{1.685174in}}%
\pgfpathcurveto{\pgfqpoint{1.757620in}{1.689292in}}{\pgfqpoint{1.752034in}{1.691606in}}{\pgfqpoint{1.746210in}{1.691606in}}%
\pgfpathcurveto{\pgfqpoint{1.740386in}{1.691606in}}{\pgfqpoint{1.734800in}{1.689292in}}{\pgfqpoint{1.730682in}{1.685174in}}%
\pgfpathcurveto{\pgfqpoint{1.726564in}{1.681056in}}{\pgfqpoint{1.724250in}{1.675469in}}{\pgfqpoint{1.724250in}{1.669646in}}%
\pgfpathcurveto{\pgfqpoint{1.724250in}{1.663822in}}{\pgfqpoint{1.726564in}{1.658235in}}{\pgfqpoint{1.730682in}{1.654117in}}%
\pgfpathcurveto{\pgfqpoint{1.734800in}{1.649999in}}{\pgfqpoint{1.740386in}{1.647685in}}{\pgfqpoint{1.746210in}{1.647685in}}%
\pgfpathlineto{\pgfqpoint{1.746210in}{1.647685in}}%
\pgfpathclose%
\pgfusepath{stroke,fill}%
\end{pgfscope}%
\begin{pgfscope}%
\pgfpathrectangle{\pgfqpoint{0.100000in}{0.209042in}}{\pgfqpoint{2.906250in}{2.906250in}}%
\pgfusepath{clip}%
\pgfsetbuttcap%
\pgfsetroundjoin%
\definecolor{currentfill}{rgb}{0.697098,0.208501,0.480835}%
\pgfsetfillcolor{currentfill}%
\pgfsetfillopacity{0.800000}%
\pgfsetlinewidth{1.003750pt}%
\definecolor{currentstroke}{rgb}{0.697098,0.208501,0.480835}%
\pgfsetstrokecolor{currentstroke}%
\pgfsetstrokeopacity{0.800000}%
\pgfsetdash{}{0pt}%
\pgfpathmoveto{\pgfqpoint{1.912213in}{2.046243in}}%
\pgfpathcurveto{\pgfqpoint{1.918037in}{2.046243in}}{\pgfqpoint{1.923623in}{2.048557in}}{\pgfqpoint{1.927741in}{2.052675in}}%
\pgfpathcurveto{\pgfqpoint{1.931859in}{2.056793in}}{\pgfqpoint{1.934173in}{2.062379in}}{\pgfqpoint{1.934173in}{2.068203in}}%
\pgfpathcurveto{\pgfqpoint{1.934173in}{2.074027in}}{\pgfqpoint{1.931859in}{2.079613in}}{\pgfqpoint{1.927741in}{2.083731in}}%
\pgfpathcurveto{\pgfqpoint{1.923623in}{2.087849in}}{\pgfqpoint{1.918037in}{2.090163in}}{\pgfqpoint{1.912213in}{2.090163in}}%
\pgfpathcurveto{\pgfqpoint{1.906389in}{2.090163in}}{\pgfqpoint{1.900803in}{2.087849in}}{\pgfqpoint{1.896685in}{2.083731in}}%
\pgfpathcurveto{\pgfqpoint{1.892566in}{2.079613in}}{\pgfqpoint{1.890253in}{2.074027in}}{\pgfqpoint{1.890253in}{2.068203in}}%
\pgfpathcurveto{\pgfqpoint{1.890253in}{2.062379in}}{\pgfqpoint{1.892566in}{2.056793in}}{\pgfqpoint{1.896685in}{2.052675in}}%
\pgfpathcurveto{\pgfqpoint{1.900803in}{2.048557in}}{\pgfqpoint{1.906389in}{2.046243in}}{\pgfqpoint{1.912213in}{2.046243in}}%
\pgfpathlineto{\pgfqpoint{1.912213in}{2.046243in}}%
\pgfpathclose%
\pgfusepath{stroke,fill}%
\end{pgfscope}%
\begin{pgfscope}%
\pgfpathrectangle{\pgfqpoint{0.100000in}{0.209042in}}{\pgfqpoint{2.906250in}{2.906250in}}%
\pgfusepath{clip}%
\pgfsetbuttcap%
\pgfsetroundjoin%
\definecolor{currentfill}{rgb}{0.171713,0.067305,0.370771}%
\pgfsetfillcolor{currentfill}%
\pgfsetfillopacity{0.800000}%
\pgfsetlinewidth{1.003750pt}%
\definecolor{currentstroke}{rgb}{0.171713,0.067305,0.370771}%
\pgfsetstrokecolor{currentstroke}%
\pgfsetstrokeopacity{0.800000}%
\pgfsetdash{}{0pt}%
\pgfpathmoveto{\pgfqpoint{1.595477in}{1.687180in}}%
\pgfpathcurveto{\pgfqpoint{1.601301in}{1.687180in}}{\pgfqpoint{1.606887in}{1.689494in}}{\pgfqpoint{1.611005in}{1.693612in}}%
\pgfpathcurveto{\pgfqpoint{1.615123in}{1.697730in}}{\pgfqpoint{1.617437in}{1.703316in}}{\pgfqpoint{1.617437in}{1.709140in}}%
\pgfpathcurveto{\pgfqpoint{1.617437in}{1.714964in}}{\pgfqpoint{1.615123in}{1.720550in}}{\pgfqpoint{1.611005in}{1.724669in}}%
\pgfpathcurveto{\pgfqpoint{1.606887in}{1.728787in}}{\pgfqpoint{1.601301in}{1.731101in}}{\pgfqpoint{1.595477in}{1.731101in}}%
\pgfpathcurveto{\pgfqpoint{1.589653in}{1.731101in}}{\pgfqpoint{1.584067in}{1.728787in}}{\pgfqpoint{1.579948in}{1.724669in}}%
\pgfpathcurveto{\pgfqpoint{1.575830in}{1.720550in}}{\pgfqpoint{1.573516in}{1.714964in}}{\pgfqpoint{1.573516in}{1.709140in}}%
\pgfpathcurveto{\pgfqpoint{1.573516in}{1.703316in}}{\pgfqpoint{1.575830in}{1.697730in}}{\pgfqpoint{1.579948in}{1.693612in}}%
\pgfpathcurveto{\pgfqpoint{1.584067in}{1.689494in}}{\pgfqpoint{1.589653in}{1.687180in}}{\pgfqpoint{1.595477in}{1.687180in}}%
\pgfpathlineto{\pgfqpoint{1.595477in}{1.687180in}}%
\pgfpathclose%
\pgfusepath{stroke,fill}%
\end{pgfscope}%
\begin{pgfscope}%
\pgfpathrectangle{\pgfqpoint{0.100000in}{0.209042in}}{\pgfqpoint{2.906250in}{2.906250in}}%
\pgfusepath{clip}%
\pgfsetbuttcap%
\pgfsetroundjoin%
\definecolor{currentfill}{rgb}{0.550287,0.161158,0.505719}%
\pgfsetfillcolor{currentfill}%
\pgfsetfillopacity{0.800000}%
\pgfsetlinewidth{1.003750pt}%
\definecolor{currentstroke}{rgb}{0.550287,0.161158,0.505719}%
\pgfsetstrokecolor{currentstroke}%
\pgfsetstrokeopacity{0.800000}%
\pgfsetdash{}{0pt}%
\pgfpathmoveto{\pgfqpoint{1.647017in}{2.059067in}}%
\pgfpathcurveto{\pgfqpoint{1.652840in}{2.059067in}}{\pgfqpoint{1.658427in}{2.061381in}}{\pgfqpoint{1.662545in}{2.065499in}}%
\pgfpathcurveto{\pgfqpoint{1.666663in}{2.069617in}}{\pgfqpoint{1.668977in}{2.075203in}}{\pgfqpoint{1.668977in}{2.081027in}}%
\pgfpathcurveto{\pgfqpoint{1.668977in}{2.086851in}}{\pgfqpoint{1.666663in}{2.092438in}}{\pgfqpoint{1.662545in}{2.096556in}}%
\pgfpathcurveto{\pgfqpoint{1.658427in}{2.100674in}}{\pgfqpoint{1.652840in}{2.102988in}}{\pgfqpoint{1.647017in}{2.102988in}}%
\pgfpathcurveto{\pgfqpoint{1.641193in}{2.102988in}}{\pgfqpoint{1.635606in}{2.100674in}}{\pgfqpoint{1.631488in}{2.096556in}}%
\pgfpathcurveto{\pgfqpoint{1.627370in}{2.092438in}}{\pgfqpoint{1.625056in}{2.086851in}}{\pgfqpoint{1.625056in}{2.081027in}}%
\pgfpathcurveto{\pgfqpoint{1.625056in}{2.075203in}}{\pgfqpoint{1.627370in}{2.069617in}}{\pgfqpoint{1.631488in}{2.065499in}}%
\pgfpathcurveto{\pgfqpoint{1.635606in}{2.061381in}}{\pgfqpoint{1.641193in}{2.059067in}}{\pgfqpoint{1.647017in}{2.059067in}}%
\pgfpathlineto{\pgfqpoint{1.647017in}{2.059067in}}%
\pgfpathclose%
\pgfusepath{stroke,fill}%
\end{pgfscope}%
\begin{pgfscope}%
\pgfpathrectangle{\pgfqpoint{0.100000in}{0.209042in}}{\pgfqpoint{2.906250in}{2.906250in}}%
\pgfusepath{clip}%
\pgfsetbuttcap%
\pgfsetroundjoin%
\definecolor{currentfill}{rgb}{0.426877,0.115395,0.505714}%
\pgfsetfillcolor{currentfill}%
\pgfsetfillopacity{0.800000}%
\pgfsetlinewidth{1.003750pt}%
\definecolor{currentstroke}{rgb}{0.426877,0.115395,0.505714}%
\pgfsetstrokecolor{currentstroke}%
\pgfsetstrokeopacity{0.800000}%
\pgfsetdash{}{0pt}%
\pgfpathmoveto{\pgfqpoint{1.758987in}{1.501719in}}%
\pgfpathcurveto{\pgfqpoint{1.764811in}{1.501719in}}{\pgfqpoint{1.770397in}{1.504033in}}{\pgfqpoint{1.774515in}{1.508151in}}%
\pgfpathcurveto{\pgfqpoint{1.778634in}{1.512269in}}{\pgfqpoint{1.780947in}{1.517856in}}{\pgfqpoint{1.780947in}{1.523679in}}%
\pgfpathcurveto{\pgfqpoint{1.780947in}{1.529503in}}{\pgfqpoint{1.778634in}{1.535090in}}{\pgfqpoint{1.774515in}{1.539208in}}%
\pgfpathcurveto{\pgfqpoint{1.770397in}{1.543326in}}{\pgfqpoint{1.764811in}{1.545640in}}{\pgfqpoint{1.758987in}{1.545640in}}%
\pgfpathcurveto{\pgfqpoint{1.753163in}{1.545640in}}{\pgfqpoint{1.747577in}{1.543326in}}{\pgfqpoint{1.743459in}{1.539208in}}%
\pgfpathcurveto{\pgfqpoint{1.739341in}{1.535090in}}{\pgfqpoint{1.737027in}{1.529503in}}{\pgfqpoint{1.737027in}{1.523679in}}%
\pgfpathcurveto{\pgfqpoint{1.737027in}{1.517856in}}{\pgfqpoint{1.739341in}{1.512269in}}{\pgfqpoint{1.743459in}{1.508151in}}%
\pgfpathcurveto{\pgfqpoint{1.747577in}{1.504033in}}{\pgfqpoint{1.753163in}{1.501719in}}{\pgfqpoint{1.758987in}{1.501719in}}%
\pgfpathlineto{\pgfqpoint{1.758987in}{1.501719in}}%
\pgfpathclose%
\pgfusepath{stroke,fill}%
\end{pgfscope}%
\begin{pgfscope}%
\pgfpathrectangle{\pgfqpoint{0.100000in}{0.209042in}}{\pgfqpoint{2.906250in}{2.906250in}}%
\pgfusepath{clip}%
\pgfsetbuttcap%
\pgfsetroundjoin%
\definecolor{currentfill}{rgb}{0.232077,0.059889,0.437695}%
\pgfsetfillcolor{currentfill}%
\pgfsetfillopacity{0.800000}%
\pgfsetlinewidth{1.003750pt}%
\definecolor{currentstroke}{rgb}{0.232077,0.059889,0.437695}%
\pgfsetstrokecolor{currentstroke}%
\pgfsetstrokeopacity{0.800000}%
\pgfsetdash{}{0pt}%
\pgfpathmoveto{\pgfqpoint{1.548353in}{1.623030in}}%
\pgfpathcurveto{\pgfqpoint{1.554177in}{1.623030in}}{\pgfqpoint{1.559763in}{1.625344in}}{\pgfqpoint{1.563881in}{1.629462in}}%
\pgfpathcurveto{\pgfqpoint{1.567999in}{1.633580in}}{\pgfqpoint{1.570313in}{1.639167in}}{\pgfqpoint{1.570313in}{1.644991in}}%
\pgfpathcurveto{\pgfqpoint{1.570313in}{1.650815in}}{\pgfqpoint{1.567999in}{1.656401in}}{\pgfqpoint{1.563881in}{1.660519in}}%
\pgfpathcurveto{\pgfqpoint{1.559763in}{1.664637in}}{\pgfqpoint{1.554177in}{1.666951in}}{\pgfqpoint{1.548353in}{1.666951in}}%
\pgfpathcurveto{\pgfqpoint{1.542529in}{1.666951in}}{\pgfqpoint{1.536942in}{1.664637in}}{\pgfqpoint{1.532824in}{1.660519in}}%
\pgfpathcurveto{\pgfqpoint{1.528706in}{1.656401in}}{\pgfqpoint{1.526392in}{1.650815in}}{\pgfqpoint{1.526392in}{1.644991in}}%
\pgfpathcurveto{\pgfqpoint{1.526392in}{1.639167in}}{\pgfqpoint{1.528706in}{1.633580in}}{\pgfqpoint{1.532824in}{1.629462in}}%
\pgfpathcurveto{\pgfqpoint{1.536942in}{1.625344in}}{\pgfqpoint{1.542529in}{1.623030in}}{\pgfqpoint{1.548353in}{1.623030in}}%
\pgfpathlineto{\pgfqpoint{1.548353in}{1.623030in}}%
\pgfpathclose%
\pgfusepath{stroke,fill}%
\end{pgfscope}%
\begin{pgfscope}%
\pgfpathrectangle{\pgfqpoint{0.100000in}{0.209042in}}{\pgfqpoint{2.906250in}{2.906250in}}%
\pgfusepath{clip}%
\pgfsetbuttcap%
\pgfsetroundjoin%
\definecolor{currentfill}{rgb}{0.519045,0.150383,0.507443}%
\pgfsetfillcolor{currentfill}%
\pgfsetfillopacity{0.800000}%
\pgfsetlinewidth{1.003750pt}%
\definecolor{currentstroke}{rgb}{0.519045,0.150383,0.507443}%
\pgfsetstrokecolor{currentstroke}%
\pgfsetstrokeopacity{0.800000}%
\pgfsetdash{}{0pt}%
\pgfpathmoveto{\pgfqpoint{1.597707in}{2.046844in}}%
\pgfpathcurveto{\pgfqpoint{1.603531in}{2.046844in}}{\pgfqpoint{1.609117in}{2.049158in}}{\pgfqpoint{1.613236in}{2.053276in}}%
\pgfpathcurveto{\pgfqpoint{1.617354in}{2.057394in}}{\pgfqpoint{1.619668in}{2.062981in}}{\pgfqpoint{1.619668in}{2.068805in}}%
\pgfpathcurveto{\pgfqpoint{1.619668in}{2.074629in}}{\pgfqpoint{1.617354in}{2.080215in}}{\pgfqpoint{1.613236in}{2.084333in}}%
\pgfpathcurveto{\pgfqpoint{1.609117in}{2.088451in}}{\pgfqpoint{1.603531in}{2.090765in}}{\pgfqpoint{1.597707in}{2.090765in}}%
\pgfpathcurveto{\pgfqpoint{1.591883in}{2.090765in}}{\pgfqpoint{1.586297in}{2.088451in}}{\pgfqpoint{1.582179in}{2.084333in}}%
\pgfpathcurveto{\pgfqpoint{1.578061in}{2.080215in}}{\pgfqpoint{1.575747in}{2.074629in}}{\pgfqpoint{1.575747in}{2.068805in}}%
\pgfpathcurveto{\pgfqpoint{1.575747in}{2.062981in}}{\pgfqpoint{1.578061in}{2.057394in}}{\pgfqpoint{1.582179in}{2.053276in}}%
\pgfpathcurveto{\pgfqpoint{1.586297in}{2.049158in}}{\pgfqpoint{1.591883in}{2.046844in}}{\pgfqpoint{1.597707in}{2.046844in}}%
\pgfpathlineto{\pgfqpoint{1.597707in}{2.046844in}}%
\pgfpathclose%
\pgfusepath{stroke,fill}%
\end{pgfscope}%
\begin{pgfscope}%
\pgfpathrectangle{\pgfqpoint{0.100000in}{0.209042in}}{\pgfqpoint{2.906250in}{2.906250in}}%
\pgfusepath{clip}%
\pgfsetbuttcap%
\pgfsetroundjoin%
\definecolor{currentfill}{rgb}{0.396467,0.102902,0.502658}%
\pgfsetfillcolor{currentfill}%
\pgfsetfillopacity{0.800000}%
\pgfsetlinewidth{1.003750pt}%
\definecolor{currentstroke}{rgb}{0.396467,0.102902,0.502658}%
\pgfsetstrokecolor{currentstroke}%
\pgfsetstrokeopacity{0.800000}%
\pgfsetdash{}{0pt}%
\pgfpathmoveto{\pgfqpoint{1.363973in}{1.882410in}}%
\pgfpathcurveto{\pgfqpoint{1.369797in}{1.882410in}}{\pgfqpoint{1.375384in}{1.884723in}}{\pgfqpoint{1.379502in}{1.888842in}}%
\pgfpathcurveto{\pgfqpoint{1.383620in}{1.892960in}}{\pgfqpoint{1.385934in}{1.898546in}}{\pgfqpoint{1.385934in}{1.904370in}}%
\pgfpathcurveto{\pgfqpoint{1.385934in}{1.910194in}}{\pgfqpoint{1.383620in}{1.915780in}}{\pgfqpoint{1.379502in}{1.919898in}}%
\pgfpathcurveto{\pgfqpoint{1.375384in}{1.924016in}}{\pgfqpoint{1.369797in}{1.926330in}}{\pgfqpoint{1.363973in}{1.926330in}}%
\pgfpathcurveto{\pgfqpoint{1.358149in}{1.926330in}}{\pgfqpoint{1.352563in}{1.924016in}}{\pgfqpoint{1.348445in}{1.919898in}}%
\pgfpathcurveto{\pgfqpoint{1.344327in}{1.915780in}}{\pgfqpoint{1.342013in}{1.910194in}}{\pgfqpoint{1.342013in}{1.904370in}}%
\pgfpathcurveto{\pgfqpoint{1.342013in}{1.898546in}}{\pgfqpoint{1.344327in}{1.892960in}}{\pgfqpoint{1.348445in}{1.888842in}}%
\pgfpathcurveto{\pgfqpoint{1.352563in}{1.884723in}}{\pgfqpoint{1.358149in}{1.882410in}}{\pgfqpoint{1.363973in}{1.882410in}}%
\pgfpathlineto{\pgfqpoint{1.363973in}{1.882410in}}%
\pgfpathclose%
\pgfusepath{stroke,fill}%
\end{pgfscope}%
\begin{pgfscope}%
\pgfpathrectangle{\pgfqpoint{0.100000in}{0.209042in}}{\pgfqpoint{2.906250in}{2.906250in}}%
\pgfusepath{clip}%
\pgfsetbuttcap%
\pgfsetroundjoin%
\definecolor{currentfill}{rgb}{0.716387,0.214982,0.475290}%
\pgfsetfillcolor{currentfill}%
\pgfsetfillopacity{0.800000}%
\pgfsetlinewidth{1.003750pt}%
\definecolor{currentstroke}{rgb}{0.716387,0.214982,0.475290}%
\pgfsetstrokecolor{currentstroke}%
\pgfsetstrokeopacity{0.800000}%
\pgfsetdash{}{0pt}%
\pgfpathmoveto{\pgfqpoint{1.088544in}{1.658102in}}%
\pgfpathcurveto{\pgfqpoint{1.094368in}{1.658102in}}{\pgfqpoint{1.099955in}{1.660415in}}{\pgfqpoint{1.104073in}{1.664534in}}%
\pgfpathcurveto{\pgfqpoint{1.108191in}{1.668652in}}{\pgfqpoint{1.110505in}{1.674238in}}{\pgfqpoint{1.110505in}{1.680062in}}%
\pgfpathcurveto{\pgfqpoint{1.110505in}{1.685886in}}{\pgfqpoint{1.108191in}{1.691472in}}{\pgfqpoint{1.104073in}{1.695590in}}%
\pgfpathcurveto{\pgfqpoint{1.099955in}{1.699708in}}{\pgfqpoint{1.094368in}{1.702022in}}{\pgfqpoint{1.088544in}{1.702022in}}%
\pgfpathcurveto{\pgfqpoint{1.082720in}{1.702022in}}{\pgfqpoint{1.077134in}{1.699708in}}{\pgfqpoint{1.073016in}{1.695590in}}%
\pgfpathcurveto{\pgfqpoint{1.068898in}{1.691472in}}{\pgfqpoint{1.066584in}{1.685886in}}{\pgfqpoint{1.066584in}{1.680062in}}%
\pgfpathcurveto{\pgfqpoint{1.066584in}{1.674238in}}{\pgfqpoint{1.068898in}{1.668652in}}{\pgfqpoint{1.073016in}{1.664534in}}%
\pgfpathcurveto{\pgfqpoint{1.077134in}{1.660415in}}{\pgfqpoint{1.082720in}{1.658102in}}{\pgfqpoint{1.088544in}{1.658102in}}%
\pgfpathlineto{\pgfqpoint{1.088544in}{1.658102in}}%
\pgfpathclose%
\pgfusepath{stroke,fill}%
\end{pgfscope}%
\begin{pgfscope}%
\pgfpathrectangle{\pgfqpoint{0.100000in}{0.209042in}}{\pgfqpoint{2.906250in}{2.906250in}}%
\pgfusepath{clip}%
\pgfsetbuttcap%
\pgfsetroundjoin%
\definecolor{currentfill}{rgb}{0.198177,0.063862,0.404009}%
\pgfsetfillcolor{currentfill}%
\pgfsetfillopacity{0.800000}%
\pgfsetlinewidth{1.003750pt}%
\definecolor{currentstroke}{rgb}{0.198177,0.063862,0.404009}%
\pgfsetstrokecolor{currentstroke}%
\pgfsetstrokeopacity{0.800000}%
\pgfsetdash{}{0pt}%
\pgfpathmoveto{\pgfqpoint{1.653547in}{1.637326in}}%
\pgfpathcurveto{\pgfqpoint{1.659371in}{1.637326in}}{\pgfqpoint{1.664957in}{1.639640in}}{\pgfqpoint{1.669075in}{1.643758in}}%
\pgfpathcurveto{\pgfqpoint{1.673193in}{1.647876in}}{\pgfqpoint{1.675507in}{1.653462in}}{\pgfqpoint{1.675507in}{1.659286in}}%
\pgfpathcurveto{\pgfqpoint{1.675507in}{1.665110in}}{\pgfqpoint{1.673193in}{1.670696in}}{\pgfqpoint{1.669075in}{1.674814in}}%
\pgfpathcurveto{\pgfqpoint{1.664957in}{1.678932in}}{\pgfqpoint{1.659371in}{1.681246in}}{\pgfqpoint{1.653547in}{1.681246in}}%
\pgfpathcurveto{\pgfqpoint{1.647723in}{1.681246in}}{\pgfqpoint{1.642137in}{1.678932in}}{\pgfqpoint{1.638019in}{1.674814in}}%
\pgfpathcurveto{\pgfqpoint{1.633901in}{1.670696in}}{\pgfqpoint{1.631587in}{1.665110in}}{\pgfqpoint{1.631587in}{1.659286in}}%
\pgfpathcurveto{\pgfqpoint{1.631587in}{1.653462in}}{\pgfqpoint{1.633901in}{1.647876in}}{\pgfqpoint{1.638019in}{1.643758in}}%
\pgfpathcurveto{\pgfqpoint{1.642137in}{1.639640in}}{\pgfqpoint{1.647723in}{1.637326in}}{\pgfqpoint{1.653547in}{1.637326in}}%
\pgfpathlineto{\pgfqpoint{1.653547in}{1.637326in}}%
\pgfpathclose%
\pgfusepath{stroke,fill}%
\end{pgfscope}%
\begin{pgfscope}%
\pgfpathrectangle{\pgfqpoint{0.100000in}{0.209042in}}{\pgfqpoint{2.906250in}{2.906250in}}%
\pgfusepath{clip}%
\pgfsetbuttcap%
\pgfsetroundjoin%
\definecolor{currentfill}{rgb}{0.359898,0.087831,0.496778}%
\pgfsetfillcolor{currentfill}%
\pgfsetfillopacity{0.800000}%
\pgfsetlinewidth{1.003750pt}%
\definecolor{currentstroke}{rgb}{0.359898,0.087831,0.496778}%
\pgfsetstrokecolor{currentstroke}%
\pgfsetstrokeopacity{0.800000}%
\pgfsetdash{}{0pt}%
\pgfpathmoveto{\pgfqpoint{1.372067in}{1.630197in}}%
\pgfpathcurveto{\pgfqpoint{1.377890in}{1.630197in}}{\pgfqpoint{1.383477in}{1.632510in}}{\pgfqpoint{1.387595in}{1.636629in}}%
\pgfpathcurveto{\pgfqpoint{1.391713in}{1.640747in}}{\pgfqpoint{1.394027in}{1.646333in}}{\pgfqpoint{1.394027in}{1.652157in}}%
\pgfpathcurveto{\pgfqpoint{1.394027in}{1.657981in}}{\pgfqpoint{1.391713in}{1.663567in}}{\pgfqpoint{1.387595in}{1.667685in}}%
\pgfpathcurveto{\pgfqpoint{1.383477in}{1.671803in}}{\pgfqpoint{1.377890in}{1.674117in}}{\pgfqpoint{1.372067in}{1.674117in}}%
\pgfpathcurveto{\pgfqpoint{1.366243in}{1.674117in}}{\pgfqpoint{1.360656in}{1.671803in}}{\pgfqpoint{1.356538in}{1.667685in}}%
\pgfpathcurveto{\pgfqpoint{1.352420in}{1.663567in}}{\pgfqpoint{1.350106in}{1.657981in}}{\pgfqpoint{1.350106in}{1.652157in}}%
\pgfpathcurveto{\pgfqpoint{1.350106in}{1.646333in}}{\pgfqpoint{1.352420in}{1.640747in}}{\pgfqpoint{1.356538in}{1.636629in}}%
\pgfpathcurveto{\pgfqpoint{1.360656in}{1.632510in}}{\pgfqpoint{1.366243in}{1.630197in}}{\pgfqpoint{1.372067in}{1.630197in}}%
\pgfpathlineto{\pgfqpoint{1.372067in}{1.630197in}}%
\pgfpathclose%
\pgfusepath{stroke,fill}%
\end{pgfscope}%
\begin{pgfscope}%
\pgfpathrectangle{\pgfqpoint{0.100000in}{0.209042in}}{\pgfqpoint{2.906250in}{2.906250in}}%
\pgfusepath{clip}%
\pgfsetbuttcap%
\pgfsetroundjoin%
\definecolor{currentfill}{rgb}{0.420791,0.112920,0.505215}%
\pgfsetfillcolor{currentfill}%
\pgfsetfillopacity{0.800000}%
\pgfsetlinewidth{1.003750pt}%
\definecolor{currentstroke}{rgb}{0.420791,0.112920,0.505215}%
\pgfsetstrokecolor{currentstroke}%
\pgfsetstrokeopacity{0.800000}%
\pgfsetdash{}{0pt}%
\pgfpathmoveto{\pgfqpoint{1.918190in}{1.671426in}}%
\pgfpathcurveto{\pgfqpoint{1.924014in}{1.671426in}}{\pgfqpoint{1.929600in}{1.673740in}}{\pgfqpoint{1.933718in}{1.677858in}}%
\pgfpathcurveto{\pgfqpoint{1.937837in}{1.681976in}}{\pgfqpoint{1.940150in}{1.687563in}}{\pgfqpoint{1.940150in}{1.693386in}}%
\pgfpathcurveto{\pgfqpoint{1.940150in}{1.699210in}}{\pgfqpoint{1.937837in}{1.704797in}}{\pgfqpoint{1.933718in}{1.708915in}}%
\pgfpathcurveto{\pgfqpoint{1.929600in}{1.713033in}}{\pgfqpoint{1.924014in}{1.715347in}}{\pgfqpoint{1.918190in}{1.715347in}}%
\pgfpathcurveto{\pgfqpoint{1.912366in}{1.715347in}}{\pgfqpoint{1.906780in}{1.713033in}}{\pgfqpoint{1.902662in}{1.708915in}}%
\pgfpathcurveto{\pgfqpoint{1.898544in}{1.704797in}}{\pgfqpoint{1.896230in}{1.699210in}}{\pgfqpoint{1.896230in}{1.693386in}}%
\pgfpathcurveto{\pgfqpoint{1.896230in}{1.687563in}}{\pgfqpoint{1.898544in}{1.681976in}}{\pgfqpoint{1.902662in}{1.677858in}}%
\pgfpathcurveto{\pgfqpoint{1.906780in}{1.673740in}}{\pgfqpoint{1.912366in}{1.671426in}}{\pgfqpoint{1.918190in}{1.671426in}}%
\pgfpathlineto{\pgfqpoint{1.918190in}{1.671426in}}%
\pgfpathclose%
\pgfusepath{stroke,fill}%
\end{pgfscope}%
\begin{pgfscope}%
\pgfpathrectangle{\pgfqpoint{0.100000in}{0.209042in}}{\pgfqpoint{2.906250in}{2.906250in}}%
\pgfusepath{clip}%
\pgfsetbuttcap%
\pgfsetroundjoin%
\definecolor{currentfill}{rgb}{0.469640,0.132245,0.507809}%
\pgfsetfillcolor{currentfill}%
\pgfsetfillopacity{0.800000}%
\pgfsetlinewidth{1.003750pt}%
\definecolor{currentstroke}{rgb}{0.469640,0.132245,0.507809}%
\pgfsetstrokecolor{currentstroke}%
\pgfsetstrokeopacity{0.800000}%
\pgfsetdash{}{0pt}%
\pgfpathmoveto{\pgfqpoint{1.310104in}{1.907844in}}%
\pgfpathcurveto{\pgfqpoint{1.315928in}{1.907844in}}{\pgfqpoint{1.321514in}{1.910158in}}{\pgfqpoint{1.325632in}{1.914276in}}%
\pgfpathcurveto{\pgfqpoint{1.329750in}{1.918394in}}{\pgfqpoint{1.332064in}{1.923980in}}{\pgfqpoint{1.332064in}{1.929804in}}%
\pgfpathcurveto{\pgfqpoint{1.332064in}{1.935628in}}{\pgfqpoint{1.329750in}{1.941214in}}{\pgfqpoint{1.325632in}{1.945333in}}%
\pgfpathcurveto{\pgfqpoint{1.321514in}{1.949451in}}{\pgfqpoint{1.315928in}{1.951765in}}{\pgfqpoint{1.310104in}{1.951765in}}%
\pgfpathcurveto{\pgfqpoint{1.304280in}{1.951765in}}{\pgfqpoint{1.298694in}{1.949451in}}{\pgfqpoint{1.294576in}{1.945333in}}%
\pgfpathcurveto{\pgfqpoint{1.290457in}{1.941214in}}{\pgfqpoint{1.288144in}{1.935628in}}{\pgfqpoint{1.288144in}{1.929804in}}%
\pgfpathcurveto{\pgfqpoint{1.288144in}{1.923980in}}{\pgfqpoint{1.290457in}{1.918394in}}{\pgfqpoint{1.294576in}{1.914276in}}%
\pgfpathcurveto{\pgfqpoint{1.298694in}{1.910158in}}{\pgfqpoint{1.304280in}{1.907844in}}{\pgfqpoint{1.310104in}{1.907844in}}%
\pgfpathlineto{\pgfqpoint{1.310104in}{1.907844in}}%
\pgfpathclose%
\pgfusepath{stroke,fill}%
\end{pgfscope}%
\begin{pgfscope}%
\pgfpathrectangle{\pgfqpoint{0.100000in}{0.209042in}}{\pgfqpoint{2.906250in}{2.906250in}}%
\pgfusepath{clip}%
\pgfsetbuttcap%
\pgfsetroundjoin%
\definecolor{currentfill}{rgb}{0.218512,0.061158,0.425392}%
\pgfsetfillcolor{currentfill}%
\pgfsetfillopacity{0.800000}%
\pgfsetlinewidth{1.003750pt}%
\definecolor{currentstroke}{rgb}{0.218512,0.061158,0.425392}%
\pgfsetstrokecolor{currentstroke}%
\pgfsetstrokeopacity{0.800000}%
\pgfsetdash{}{0pt}%
\pgfpathmoveto{\pgfqpoint{1.525986in}{1.840964in}}%
\pgfpathcurveto{\pgfqpoint{1.531810in}{1.840964in}}{\pgfqpoint{1.537396in}{1.843278in}}{\pgfqpoint{1.541514in}{1.847396in}}%
\pgfpathcurveto{\pgfqpoint{1.545632in}{1.851514in}}{\pgfqpoint{1.547946in}{1.857100in}}{\pgfqpoint{1.547946in}{1.862924in}}%
\pgfpathcurveto{\pgfqpoint{1.547946in}{1.868748in}}{\pgfqpoint{1.545632in}{1.874334in}}{\pgfqpoint{1.541514in}{1.878453in}}%
\pgfpathcurveto{\pgfqpoint{1.537396in}{1.882571in}}{\pgfqpoint{1.531810in}{1.884885in}}{\pgfqpoint{1.525986in}{1.884885in}}%
\pgfpathcurveto{\pgfqpoint{1.520162in}{1.884885in}}{\pgfqpoint{1.514576in}{1.882571in}}{\pgfqpoint{1.510458in}{1.878453in}}%
\pgfpathcurveto{\pgfqpoint{1.506340in}{1.874334in}}{\pgfqpoint{1.504026in}{1.868748in}}{\pgfqpoint{1.504026in}{1.862924in}}%
\pgfpathcurveto{\pgfqpoint{1.504026in}{1.857100in}}{\pgfqpoint{1.506340in}{1.851514in}}{\pgfqpoint{1.510458in}{1.847396in}}%
\pgfpathcurveto{\pgfqpoint{1.514576in}{1.843278in}}{\pgfqpoint{1.520162in}{1.840964in}}{\pgfqpoint{1.525986in}{1.840964in}}%
\pgfpathlineto{\pgfqpoint{1.525986in}{1.840964in}}%
\pgfpathclose%
\pgfusepath{stroke,fill}%
\end{pgfscope}%
\begin{pgfscope}%
\pgfpathrectangle{\pgfqpoint{0.100000in}{0.209042in}}{\pgfqpoint{2.906250in}{2.906250in}}%
\pgfusepath{clip}%
\pgfsetbuttcap%
\pgfsetroundjoin%
\definecolor{currentfill}{rgb}{0.531507,0.154739,0.506895}%
\pgfsetfillcolor{currentfill}%
\pgfsetfillopacity{0.800000}%
\pgfsetlinewidth{1.003750pt}%
\definecolor{currentstroke}{rgb}{0.531507,0.154739,0.506895}%
\pgfsetstrokecolor{currentstroke}%
\pgfsetstrokeopacity{0.800000}%
\pgfsetdash{}{0pt}%
\pgfpathmoveto{\pgfqpoint{1.273050in}{1.936612in}}%
\pgfpathcurveto{\pgfqpoint{1.278874in}{1.936612in}}{\pgfqpoint{1.284460in}{1.938926in}}{\pgfqpoint{1.288578in}{1.943044in}}%
\pgfpathcurveto{\pgfqpoint{1.292696in}{1.947162in}}{\pgfqpoint{1.295010in}{1.952749in}}{\pgfqpoint{1.295010in}{1.958572in}}%
\pgfpathcurveto{\pgfqpoint{1.295010in}{1.964396in}}{\pgfqpoint{1.292696in}{1.969983in}}{\pgfqpoint{1.288578in}{1.974101in}}%
\pgfpathcurveto{\pgfqpoint{1.284460in}{1.978219in}}{\pgfqpoint{1.278874in}{1.980533in}}{\pgfqpoint{1.273050in}{1.980533in}}%
\pgfpathcurveto{\pgfqpoint{1.267226in}{1.980533in}}{\pgfqpoint{1.261640in}{1.978219in}}{\pgfqpoint{1.257522in}{1.974101in}}%
\pgfpathcurveto{\pgfqpoint{1.253404in}{1.969983in}}{\pgfqpoint{1.251090in}{1.964396in}}{\pgfqpoint{1.251090in}{1.958572in}}%
\pgfpathcurveto{\pgfqpoint{1.251090in}{1.952749in}}{\pgfqpoint{1.253404in}{1.947162in}}{\pgfqpoint{1.257522in}{1.943044in}}%
\pgfpathcurveto{\pgfqpoint{1.261640in}{1.938926in}}{\pgfqpoint{1.267226in}{1.936612in}}{\pgfqpoint{1.273050in}{1.936612in}}%
\pgfpathlineto{\pgfqpoint{1.273050in}{1.936612in}}%
\pgfpathclose%
\pgfusepath{stroke,fill}%
\end{pgfscope}%
\begin{pgfscope}%
\pgfpathrectangle{\pgfqpoint{0.100000in}{0.209042in}}{\pgfqpoint{2.906250in}{2.906250in}}%
\pgfusepath{clip}%
\pgfsetbuttcap%
\pgfsetroundjoin%
\definecolor{currentfill}{rgb}{0.353773,0.085373,0.495501}%
\pgfsetfillcolor{currentfill}%
\pgfsetfillopacity{0.800000}%
\pgfsetlinewidth{1.003750pt}%
\definecolor{currentstroke}{rgb}{0.353773,0.085373,0.495501}%
\pgfsetstrokecolor{currentstroke}%
\pgfsetstrokeopacity{0.800000}%
\pgfsetdash{}{0pt}%
\pgfpathmoveto{\pgfqpoint{1.865838in}{1.780952in}}%
\pgfpathcurveto{\pgfqpoint{1.871662in}{1.780952in}}{\pgfqpoint{1.877248in}{1.783266in}}{\pgfqpoint{1.881366in}{1.787384in}}%
\pgfpathcurveto{\pgfqpoint{1.885484in}{1.791502in}}{\pgfqpoint{1.887798in}{1.797088in}}{\pgfqpoint{1.887798in}{1.802912in}}%
\pgfpathcurveto{\pgfqpoint{1.887798in}{1.808736in}}{\pgfqpoint{1.885484in}{1.814322in}}{\pgfqpoint{1.881366in}{1.818441in}}%
\pgfpathcurveto{\pgfqpoint{1.877248in}{1.822559in}}{\pgfqpoint{1.871662in}{1.824873in}}{\pgfqpoint{1.865838in}{1.824873in}}%
\pgfpathcurveto{\pgfqpoint{1.860014in}{1.824873in}}{\pgfqpoint{1.854428in}{1.822559in}}{\pgfqpoint{1.850310in}{1.818441in}}%
\pgfpathcurveto{\pgfqpoint{1.846191in}{1.814322in}}{\pgfqpoint{1.843878in}{1.808736in}}{\pgfqpoint{1.843878in}{1.802912in}}%
\pgfpathcurveto{\pgfqpoint{1.843878in}{1.797088in}}{\pgfqpoint{1.846191in}{1.791502in}}{\pgfqpoint{1.850310in}{1.787384in}}%
\pgfpathcurveto{\pgfqpoint{1.854428in}{1.783266in}}{\pgfqpoint{1.860014in}{1.780952in}}{\pgfqpoint{1.865838in}{1.780952in}}%
\pgfpathlineto{\pgfqpoint{1.865838in}{1.780952in}}%
\pgfpathclose%
\pgfusepath{stroke,fill}%
\end{pgfscope}%
\begin{pgfscope}%
\pgfpathrectangle{\pgfqpoint{0.100000in}{0.209042in}}{\pgfqpoint{2.906250in}{2.906250in}}%
\pgfusepath{clip}%
\pgfsetbuttcap%
\pgfsetroundjoin%
\definecolor{currentfill}{rgb}{0.556571,0.163269,0.505230}%
\pgfsetfillcolor{currentfill}%
\pgfsetfillopacity{0.800000}%
\pgfsetlinewidth{1.003750pt}%
\definecolor{currentstroke}{rgb}{0.556571,0.163269,0.505230}%
\pgfsetstrokecolor{currentstroke}%
\pgfsetstrokeopacity{0.800000}%
\pgfsetdash{}{0pt}%
\pgfpathmoveto{\pgfqpoint{1.620623in}{1.385677in}}%
\pgfpathcurveto{\pgfqpoint{1.626447in}{1.385677in}}{\pgfqpoint{1.632033in}{1.387991in}}{\pgfqpoint{1.636151in}{1.392109in}}%
\pgfpathcurveto{\pgfqpoint{1.640269in}{1.396227in}}{\pgfqpoint{1.642583in}{1.401813in}}{\pgfqpoint{1.642583in}{1.407637in}}%
\pgfpathcurveto{\pgfqpoint{1.642583in}{1.413461in}}{\pgfqpoint{1.640269in}{1.419047in}}{\pgfqpoint{1.636151in}{1.423165in}}%
\pgfpathcurveto{\pgfqpoint{1.632033in}{1.427283in}}{\pgfqpoint{1.626447in}{1.429597in}}{\pgfqpoint{1.620623in}{1.429597in}}%
\pgfpathcurveto{\pgfqpoint{1.614799in}{1.429597in}}{\pgfqpoint{1.609213in}{1.427283in}}{\pgfqpoint{1.605095in}{1.423165in}}%
\pgfpathcurveto{\pgfqpoint{1.600977in}{1.419047in}}{\pgfqpoint{1.598663in}{1.413461in}}{\pgfqpoint{1.598663in}{1.407637in}}%
\pgfpathcurveto{\pgfqpoint{1.598663in}{1.401813in}}{\pgfqpoint{1.600977in}{1.396227in}}{\pgfqpoint{1.605095in}{1.392109in}}%
\pgfpathcurveto{\pgfqpoint{1.609213in}{1.387991in}}{\pgfqpoint{1.614799in}{1.385677in}}{\pgfqpoint{1.620623in}{1.385677in}}%
\pgfpathlineto{\pgfqpoint{1.620623in}{1.385677in}}%
\pgfpathclose%
\pgfusepath{stroke,fill}%
\end{pgfscope}%
\begin{pgfscope}%
\pgfpathrectangle{\pgfqpoint{0.100000in}{0.209042in}}{\pgfqpoint{2.906250in}{2.906250in}}%
\pgfusepath{clip}%
\pgfsetbuttcap%
\pgfsetroundjoin%
\definecolor{currentfill}{rgb}{0.322899,0.073782,0.487408}%
\pgfsetfillcolor{currentfill}%
\pgfsetfillopacity{0.800000}%
\pgfsetlinewidth{1.003750pt}%
\definecolor{currentstroke}{rgb}{0.322899,0.073782,0.487408}%
\pgfsetstrokecolor{currentstroke}%
\pgfsetstrokeopacity{0.800000}%
\pgfsetdash{}{0pt}%
\pgfpathmoveto{\pgfqpoint{1.388711in}{1.638899in}}%
\pgfpathcurveto{\pgfqpoint{1.394535in}{1.638899in}}{\pgfqpoint{1.400121in}{1.641213in}}{\pgfqpoint{1.404239in}{1.645331in}}%
\pgfpathcurveto{\pgfqpoint{1.408357in}{1.649449in}}{\pgfqpoint{1.410671in}{1.655035in}}{\pgfqpoint{1.410671in}{1.660859in}}%
\pgfpathcurveto{\pgfqpoint{1.410671in}{1.666683in}}{\pgfqpoint{1.408357in}{1.672269in}}{\pgfqpoint{1.404239in}{1.676388in}}%
\pgfpathcurveto{\pgfqpoint{1.400121in}{1.680506in}}{\pgfqpoint{1.394535in}{1.682820in}}{\pgfqpoint{1.388711in}{1.682820in}}%
\pgfpathcurveto{\pgfqpoint{1.382887in}{1.682820in}}{\pgfqpoint{1.377301in}{1.680506in}}{\pgfqpoint{1.373183in}{1.676388in}}%
\pgfpathcurveto{\pgfqpoint{1.369065in}{1.672269in}}{\pgfqpoint{1.366751in}{1.666683in}}{\pgfqpoint{1.366751in}{1.660859in}}%
\pgfpathcurveto{\pgfqpoint{1.366751in}{1.655035in}}{\pgfqpoint{1.369065in}{1.649449in}}{\pgfqpoint{1.373183in}{1.645331in}}%
\pgfpathcurveto{\pgfqpoint{1.377301in}{1.641213in}}{\pgfqpoint{1.382887in}{1.638899in}}{\pgfqpoint{1.388711in}{1.638899in}}%
\pgfpathlineto{\pgfqpoint{1.388711in}{1.638899in}}%
\pgfpathclose%
\pgfusepath{stroke,fill}%
\end{pgfscope}%
\begin{pgfscope}%
\pgfpathrectangle{\pgfqpoint{0.100000in}{0.209042in}}{\pgfqpoint{2.906250in}{2.906250in}}%
\pgfusepath{clip}%
\pgfsetbuttcap%
\pgfsetroundjoin%
\definecolor{currentfill}{rgb}{0.414709,0.110431,0.504662}%
\pgfsetfillcolor{currentfill}%
\pgfsetfillopacity{0.800000}%
\pgfsetlinewidth{1.003750pt}%
\definecolor{currentstroke}{rgb}{0.414709,0.110431,0.504662}%
\pgfsetstrokecolor{currentstroke}%
\pgfsetstrokeopacity{0.800000}%
\pgfsetdash{}{0pt}%
\pgfpathmoveto{\pgfqpoint{1.591911in}{1.469334in}}%
\pgfpathcurveto{\pgfqpoint{1.597735in}{1.469334in}}{\pgfqpoint{1.603321in}{1.471648in}}{\pgfqpoint{1.607439in}{1.475766in}}%
\pgfpathcurveto{\pgfqpoint{1.611557in}{1.479884in}}{\pgfqpoint{1.613871in}{1.485471in}}{\pgfqpoint{1.613871in}{1.491295in}}%
\pgfpathcurveto{\pgfqpoint{1.613871in}{1.497118in}}{\pgfqpoint{1.611557in}{1.502705in}}{\pgfqpoint{1.607439in}{1.506823in}}%
\pgfpathcurveto{\pgfqpoint{1.603321in}{1.510941in}}{\pgfqpoint{1.597735in}{1.513255in}}{\pgfqpoint{1.591911in}{1.513255in}}%
\pgfpathcurveto{\pgfqpoint{1.586087in}{1.513255in}}{\pgfqpoint{1.580501in}{1.510941in}}{\pgfqpoint{1.576383in}{1.506823in}}%
\pgfpathcurveto{\pgfqpoint{1.572265in}{1.502705in}}{\pgfqpoint{1.569951in}{1.497118in}}{\pgfqpoint{1.569951in}{1.491295in}}%
\pgfpathcurveto{\pgfqpoint{1.569951in}{1.485471in}}{\pgfqpoint{1.572265in}{1.479884in}}{\pgfqpoint{1.576383in}{1.475766in}}%
\pgfpathcurveto{\pgfqpoint{1.580501in}{1.471648in}}{\pgfqpoint{1.586087in}{1.469334in}}{\pgfqpoint{1.591911in}{1.469334in}}%
\pgfpathlineto{\pgfqpoint{1.591911in}{1.469334in}}%
\pgfpathclose%
\pgfusepath{stroke,fill}%
\end{pgfscope}%
\begin{pgfscope}%
\pgfpathrectangle{\pgfqpoint{0.100000in}{0.209042in}}{\pgfqpoint{2.906250in}{2.906250in}}%
\pgfusepath{clip}%
\pgfsetbuttcap%
\pgfsetroundjoin%
\definecolor{currentfill}{rgb}{0.384299,0.097855,0.501002}%
\pgfsetfillcolor{currentfill}%
\pgfsetfillopacity{0.800000}%
\pgfsetlinewidth{1.003750pt}%
\definecolor{currentstroke}{rgb}{0.384299,0.097855,0.501002}%
\pgfsetstrokecolor{currentstroke}%
\pgfsetstrokeopacity{0.800000}%
\pgfsetdash{}{0pt}%
\pgfpathmoveto{\pgfqpoint{1.387532in}{1.574852in}}%
\pgfpathcurveto{\pgfqpoint{1.393356in}{1.574852in}}{\pgfqpoint{1.398942in}{1.577166in}}{\pgfqpoint{1.403060in}{1.581284in}}%
\pgfpathcurveto{\pgfqpoint{1.407179in}{1.585402in}}{\pgfqpoint{1.409492in}{1.590989in}}{\pgfqpoint{1.409492in}{1.596813in}}%
\pgfpathcurveto{\pgfqpoint{1.409492in}{1.602636in}}{\pgfqpoint{1.407179in}{1.608223in}}{\pgfqpoint{1.403060in}{1.612341in}}%
\pgfpathcurveto{\pgfqpoint{1.398942in}{1.616459in}}{\pgfqpoint{1.393356in}{1.618773in}}{\pgfqpoint{1.387532in}{1.618773in}}%
\pgfpathcurveto{\pgfqpoint{1.381708in}{1.618773in}}{\pgfqpoint{1.376122in}{1.616459in}}{\pgfqpoint{1.372004in}{1.612341in}}%
\pgfpathcurveto{\pgfqpoint{1.367886in}{1.608223in}}{\pgfqpoint{1.365572in}{1.602636in}}{\pgfqpoint{1.365572in}{1.596813in}}%
\pgfpathcurveto{\pgfqpoint{1.365572in}{1.590989in}}{\pgfqpoint{1.367886in}{1.585402in}}{\pgfqpoint{1.372004in}{1.581284in}}%
\pgfpathcurveto{\pgfqpoint{1.376122in}{1.577166in}}{\pgfqpoint{1.381708in}{1.574852in}}{\pgfqpoint{1.387532in}{1.574852in}}%
\pgfpathlineto{\pgfqpoint{1.387532in}{1.574852in}}%
\pgfpathclose%
\pgfusepath{stroke,fill}%
\end{pgfscope}%
\begin{pgfscope}%
\pgfpathrectangle{\pgfqpoint{0.100000in}{0.209042in}}{\pgfqpoint{2.906250in}{2.906250in}}%
\pgfusepath{clip}%
\pgfsetbuttcap%
\pgfsetroundjoin%
\definecolor{currentfill}{rgb}{0.463508,0.129893,0.507652}%
\pgfsetfillcolor{currentfill}%
\pgfsetfillopacity{0.800000}%
\pgfsetlinewidth{1.003750pt}%
\definecolor{currentstroke}{rgb}{0.463508,0.129893,0.507652}%
\pgfsetstrokecolor{currentstroke}%
\pgfsetstrokeopacity{0.800000}%
\pgfsetdash{}{0pt}%
\pgfpathmoveto{\pgfqpoint{1.707786in}{1.447589in}}%
\pgfpathcurveto{\pgfqpoint{1.713610in}{1.447589in}}{\pgfqpoint{1.719196in}{1.449902in}}{\pgfqpoint{1.723314in}{1.454021in}}%
\pgfpathcurveto{\pgfqpoint{1.727433in}{1.458139in}}{\pgfqpoint{1.729746in}{1.463725in}}{\pgfqpoint{1.729746in}{1.469549in}}%
\pgfpathcurveto{\pgfqpoint{1.729746in}{1.475373in}}{\pgfqpoint{1.727433in}{1.480959in}}{\pgfqpoint{1.723314in}{1.485077in}}%
\pgfpathcurveto{\pgfqpoint{1.719196in}{1.489195in}}{\pgfqpoint{1.713610in}{1.491509in}}{\pgfqpoint{1.707786in}{1.491509in}}%
\pgfpathcurveto{\pgfqpoint{1.701962in}{1.491509in}}{\pgfqpoint{1.696376in}{1.489195in}}{\pgfqpoint{1.692258in}{1.485077in}}%
\pgfpathcurveto{\pgfqpoint{1.688140in}{1.480959in}}{\pgfqpoint{1.685826in}{1.475373in}}{\pgfqpoint{1.685826in}{1.469549in}}%
\pgfpathcurveto{\pgfqpoint{1.685826in}{1.463725in}}{\pgfqpoint{1.688140in}{1.458139in}}{\pgfqpoint{1.692258in}{1.454021in}}%
\pgfpathcurveto{\pgfqpoint{1.696376in}{1.449902in}}{\pgfqpoint{1.701962in}{1.447589in}}{\pgfqpoint{1.707786in}{1.447589in}}%
\pgfpathlineto{\pgfqpoint{1.707786in}{1.447589in}}%
\pgfpathclose%
\pgfusepath{stroke,fill}%
\end{pgfscope}%
\begin{pgfscope}%
\pgfpathrectangle{\pgfqpoint{0.100000in}{0.209042in}}{\pgfqpoint{2.906250in}{2.906250in}}%
\pgfusepath{clip}%
\pgfsetbuttcap%
\pgfsetroundjoin%
\definecolor{currentfill}{rgb}{0.353773,0.085373,0.495501}%
\pgfsetfillcolor{currentfill}%
\pgfsetfillopacity{0.800000}%
\pgfsetlinewidth{1.003750pt}%
\definecolor{currentstroke}{rgb}{0.353773,0.085373,0.495501}%
\pgfsetstrokecolor{currentstroke}%
\pgfsetstrokeopacity{0.800000}%
\pgfsetdash{}{0pt}%
\pgfpathmoveto{\pgfqpoint{1.871760in}{1.655318in}}%
\pgfpathcurveto{\pgfqpoint{1.877584in}{1.655318in}}{\pgfqpoint{1.883170in}{1.657632in}}{\pgfqpoint{1.887288in}{1.661750in}}%
\pgfpathcurveto{\pgfqpoint{1.891406in}{1.665869in}}{\pgfqpoint{1.893720in}{1.671455in}}{\pgfqpoint{1.893720in}{1.677279in}}%
\pgfpathcurveto{\pgfqpoint{1.893720in}{1.683103in}}{\pgfqpoint{1.891406in}{1.688689in}}{\pgfqpoint{1.887288in}{1.692807in}}%
\pgfpathcurveto{\pgfqpoint{1.883170in}{1.696925in}}{\pgfqpoint{1.877584in}{1.699239in}}{\pgfqpoint{1.871760in}{1.699239in}}%
\pgfpathcurveto{\pgfqpoint{1.865936in}{1.699239in}}{\pgfqpoint{1.860350in}{1.696925in}}{\pgfqpoint{1.856231in}{1.692807in}}%
\pgfpathcurveto{\pgfqpoint{1.852113in}{1.688689in}}{\pgfqpoint{1.849799in}{1.683103in}}{\pgfqpoint{1.849799in}{1.677279in}}%
\pgfpathcurveto{\pgfqpoint{1.849799in}{1.671455in}}{\pgfqpoint{1.852113in}{1.665869in}}{\pgfqpoint{1.856231in}{1.661750in}}%
\pgfpathcurveto{\pgfqpoint{1.860350in}{1.657632in}}{\pgfqpoint{1.865936in}{1.655318in}}{\pgfqpoint{1.871760in}{1.655318in}}%
\pgfpathlineto{\pgfqpoint{1.871760in}{1.655318in}}%
\pgfpathclose%
\pgfusepath{stroke,fill}%
\end{pgfscope}%
\begin{pgfscope}%
\pgfpathrectangle{\pgfqpoint{0.100000in}{0.209042in}}{\pgfqpoint{2.906250in}{2.906250in}}%
\pgfusepath{clip}%
\pgfsetbuttcap%
\pgfsetroundjoin%
\definecolor{currentfill}{rgb}{0.265447,0.060237,0.461840}%
\pgfsetfillcolor{currentfill}%
\pgfsetfillopacity{0.800000}%
\pgfsetlinewidth{1.003750pt}%
\definecolor{currentstroke}{rgb}{0.265447,0.060237,0.461840}%
\pgfsetstrokecolor{currentstroke}%
\pgfsetstrokeopacity{0.800000}%
\pgfsetdash{}{0pt}%
\pgfpathmoveto{\pgfqpoint{1.487121in}{1.871892in}}%
\pgfpathcurveto{\pgfqpoint{1.492945in}{1.871892in}}{\pgfqpoint{1.498531in}{1.874206in}}{\pgfqpoint{1.502649in}{1.878324in}}%
\pgfpathcurveto{\pgfqpoint{1.506767in}{1.882442in}}{\pgfqpoint{1.509081in}{1.888029in}}{\pgfqpoint{1.509081in}{1.893853in}}%
\pgfpathcurveto{\pgfqpoint{1.509081in}{1.899676in}}{\pgfqpoint{1.506767in}{1.905263in}}{\pgfqpoint{1.502649in}{1.909381in}}%
\pgfpathcurveto{\pgfqpoint{1.498531in}{1.913499in}}{\pgfqpoint{1.492945in}{1.915813in}}{\pgfqpoint{1.487121in}{1.915813in}}%
\pgfpathcurveto{\pgfqpoint{1.481297in}{1.915813in}}{\pgfqpoint{1.475711in}{1.913499in}}{\pgfqpoint{1.471593in}{1.909381in}}%
\pgfpathcurveto{\pgfqpoint{1.467475in}{1.905263in}}{\pgfqpoint{1.465161in}{1.899676in}}{\pgfqpoint{1.465161in}{1.893853in}}%
\pgfpathcurveto{\pgfqpoint{1.465161in}{1.888029in}}{\pgfqpoint{1.467475in}{1.882442in}}{\pgfqpoint{1.471593in}{1.878324in}}%
\pgfpathcurveto{\pgfqpoint{1.475711in}{1.874206in}}{\pgfqpoint{1.481297in}{1.871892in}}{\pgfqpoint{1.487121in}{1.871892in}}%
\pgfpathlineto{\pgfqpoint{1.487121in}{1.871892in}}%
\pgfpathclose%
\pgfusepath{stroke,fill}%
\end{pgfscope}%
\begin{pgfscope}%
\pgfpathrectangle{\pgfqpoint{0.100000in}{0.209042in}}{\pgfqpoint{2.906250in}{2.906250in}}%
\pgfusepath{clip}%
\pgfsetbuttcap%
\pgfsetroundjoin%
\definecolor{currentfill}{rgb}{0.645633,0.191952,0.492910}%
\pgfsetfillcolor{currentfill}%
\pgfsetfillopacity{0.800000}%
\pgfsetlinewidth{1.003750pt}%
\definecolor{currentstroke}{rgb}{0.645633,0.191952,0.492910}%
\pgfsetstrokecolor{currentstroke}%
\pgfsetstrokeopacity{0.800000}%
\pgfsetdash{}{0pt}%
\pgfpathmoveto{\pgfqpoint{1.347363in}{2.079448in}}%
\pgfpathcurveto{\pgfqpoint{1.353187in}{2.079448in}}{\pgfqpoint{1.358773in}{2.081762in}}{\pgfqpoint{1.362891in}{2.085880in}}%
\pgfpathcurveto{\pgfqpoint{1.367010in}{2.089998in}}{\pgfqpoint{1.369323in}{2.095585in}}{\pgfqpoint{1.369323in}{2.101408in}}%
\pgfpathcurveto{\pgfqpoint{1.369323in}{2.107232in}}{\pgfqpoint{1.367010in}{2.112819in}}{\pgfqpoint{1.362891in}{2.116937in}}%
\pgfpathcurveto{\pgfqpoint{1.358773in}{2.121055in}}{\pgfqpoint{1.353187in}{2.123369in}}{\pgfqpoint{1.347363in}{2.123369in}}%
\pgfpathcurveto{\pgfqpoint{1.341539in}{2.123369in}}{\pgfqpoint{1.335953in}{2.121055in}}{\pgfqpoint{1.331835in}{2.116937in}}%
\pgfpathcurveto{\pgfqpoint{1.327717in}{2.112819in}}{\pgfqpoint{1.325403in}{2.107232in}}{\pgfqpoint{1.325403in}{2.101408in}}%
\pgfpathcurveto{\pgfqpoint{1.325403in}{2.095585in}}{\pgfqpoint{1.327717in}{2.089998in}}{\pgfqpoint{1.331835in}{2.085880in}}%
\pgfpathcurveto{\pgfqpoint{1.335953in}{2.081762in}}{\pgfqpoint{1.341539in}{2.079448in}}{\pgfqpoint{1.347363in}{2.079448in}}%
\pgfpathlineto{\pgfqpoint{1.347363in}{2.079448in}}%
\pgfpathclose%
\pgfusepath{stroke,fill}%
\end{pgfscope}%
\begin{pgfscope}%
\pgfpathrectangle{\pgfqpoint{0.100000in}{0.209042in}}{\pgfqpoint{2.906250in}{2.906250in}}%
\pgfusepath{clip}%
\pgfsetbuttcap%
\pgfsetroundjoin%
\definecolor{currentfill}{rgb}{0.271994,0.060994,0.465660}%
\pgfsetfillcolor{currentfill}%
\pgfsetfillopacity{0.800000}%
\pgfsetlinewidth{1.003750pt}%
\definecolor{currentstroke}{rgb}{0.271994,0.060994,0.465660}%
\pgfsetstrokecolor{currentstroke}%
\pgfsetstrokeopacity{0.800000}%
\pgfsetdash{}{0pt}%
\pgfpathmoveto{\pgfqpoint{1.395511in}{1.705760in}}%
\pgfpathcurveto{\pgfqpoint{1.401335in}{1.705760in}}{\pgfqpoint{1.406921in}{1.708074in}}{\pgfqpoint{1.411039in}{1.712192in}}%
\pgfpathcurveto{\pgfqpoint{1.415158in}{1.716310in}}{\pgfqpoint{1.417471in}{1.721897in}}{\pgfqpoint{1.417471in}{1.727720in}}%
\pgfpathcurveto{\pgfqpoint{1.417471in}{1.733544in}}{\pgfqpoint{1.415158in}{1.739131in}}{\pgfqpoint{1.411039in}{1.743249in}}%
\pgfpathcurveto{\pgfqpoint{1.406921in}{1.747367in}}{\pgfqpoint{1.401335in}{1.749681in}}{\pgfqpoint{1.395511in}{1.749681in}}%
\pgfpathcurveto{\pgfqpoint{1.389687in}{1.749681in}}{\pgfqpoint{1.384101in}{1.747367in}}{\pgfqpoint{1.379983in}{1.743249in}}%
\pgfpathcurveto{\pgfqpoint{1.375865in}{1.739131in}}{\pgfqpoint{1.373551in}{1.733544in}}{\pgfqpoint{1.373551in}{1.727720in}}%
\pgfpathcurveto{\pgfqpoint{1.373551in}{1.721897in}}{\pgfqpoint{1.375865in}{1.716310in}}{\pgfqpoint{1.379983in}{1.712192in}}%
\pgfpathcurveto{\pgfqpoint{1.384101in}{1.708074in}}{\pgfqpoint{1.389687in}{1.705760in}}{\pgfqpoint{1.395511in}{1.705760in}}%
\pgfpathlineto{\pgfqpoint{1.395511in}{1.705760in}}%
\pgfpathclose%
\pgfusepath{stroke,fill}%
\end{pgfscope}%
\begin{pgfscope}%
\pgfpathrectangle{\pgfqpoint{0.100000in}{0.209042in}}{\pgfqpoint{2.906250in}{2.906250in}}%
\pgfusepath{clip}%
\pgfsetbuttcap%
\pgfsetroundjoin%
\definecolor{currentfill}{rgb}{0.408629,0.107930,0.504052}%
\pgfsetfillcolor{currentfill}%
\pgfsetfillopacity{0.800000}%
\pgfsetlinewidth{1.003750pt}%
\definecolor{currentstroke}{rgb}{0.408629,0.107930,0.504052}%
\pgfsetstrokecolor{currentstroke}%
\pgfsetstrokeopacity{0.800000}%
\pgfsetdash{}{0pt}%
\pgfpathmoveto{\pgfqpoint{1.419396in}{1.522780in}}%
\pgfpathcurveto{\pgfqpoint{1.425220in}{1.522780in}}{\pgfqpoint{1.430806in}{1.525093in}}{\pgfqpoint{1.434925in}{1.529212in}}%
\pgfpathcurveto{\pgfqpoint{1.439043in}{1.533330in}}{\pgfqpoint{1.441357in}{1.538916in}}{\pgfqpoint{1.441357in}{1.544740in}}%
\pgfpathcurveto{\pgfqpoint{1.441357in}{1.550564in}}{\pgfqpoint{1.439043in}{1.556150in}}{\pgfqpoint{1.434925in}{1.560268in}}%
\pgfpathcurveto{\pgfqpoint{1.430806in}{1.564386in}}{\pgfqpoint{1.425220in}{1.566700in}}{\pgfqpoint{1.419396in}{1.566700in}}%
\pgfpathcurveto{\pgfqpoint{1.413572in}{1.566700in}}{\pgfqpoint{1.407986in}{1.564386in}}{\pgfqpoint{1.403868in}{1.560268in}}%
\pgfpathcurveto{\pgfqpoint{1.399750in}{1.556150in}}{\pgfqpoint{1.397436in}{1.550564in}}{\pgfqpoint{1.397436in}{1.544740in}}%
\pgfpathcurveto{\pgfqpoint{1.397436in}{1.538916in}}{\pgfqpoint{1.399750in}{1.533330in}}{\pgfqpoint{1.403868in}{1.529212in}}%
\pgfpathcurveto{\pgfqpoint{1.407986in}{1.525093in}}{\pgfqpoint{1.413572in}{1.522780in}}{\pgfqpoint{1.419396in}{1.522780in}}%
\pgfpathlineto{\pgfqpoint{1.419396in}{1.522780in}}%
\pgfpathclose%
\pgfusepath{stroke,fill}%
\end{pgfscope}%
\begin{pgfscope}%
\pgfpathrectangle{\pgfqpoint{0.100000in}{0.209042in}}{\pgfqpoint{2.906250in}{2.906250in}}%
\pgfusepath{clip}%
\pgfsetbuttcap%
\pgfsetroundjoin%
\definecolor{currentfill}{rgb}{0.378211,0.095332,0.500067}%
\pgfsetfillcolor{currentfill}%
\pgfsetfillopacity{0.800000}%
\pgfsetlinewidth{1.003750pt}%
\definecolor{currentstroke}{rgb}{0.378211,0.095332,0.500067}%
\pgfsetstrokecolor{currentstroke}%
\pgfsetstrokeopacity{0.800000}%
\pgfsetdash{}{0pt}%
\pgfpathmoveto{\pgfqpoint{1.837591in}{1.561120in}}%
\pgfpathcurveto{\pgfqpoint{1.843415in}{1.561120in}}{\pgfqpoint{1.849001in}{1.563434in}}{\pgfqpoint{1.853119in}{1.567552in}}%
\pgfpathcurveto{\pgfqpoint{1.857237in}{1.571670in}}{\pgfqpoint{1.859551in}{1.577257in}}{\pgfqpoint{1.859551in}{1.583081in}}%
\pgfpathcurveto{\pgfqpoint{1.859551in}{1.588904in}}{\pgfqpoint{1.857237in}{1.594491in}}{\pgfqpoint{1.853119in}{1.598609in}}%
\pgfpathcurveto{\pgfqpoint{1.849001in}{1.602727in}}{\pgfqpoint{1.843415in}{1.605041in}}{\pgfqpoint{1.837591in}{1.605041in}}%
\pgfpathcurveto{\pgfqpoint{1.831767in}{1.605041in}}{\pgfqpoint{1.826181in}{1.602727in}}{\pgfqpoint{1.822063in}{1.598609in}}%
\pgfpathcurveto{\pgfqpoint{1.817945in}{1.594491in}}{\pgfqpoint{1.815631in}{1.588904in}}{\pgfqpoint{1.815631in}{1.583081in}}%
\pgfpathcurveto{\pgfqpoint{1.815631in}{1.577257in}}{\pgfqpoint{1.817945in}{1.571670in}}{\pgfqpoint{1.822063in}{1.567552in}}%
\pgfpathcurveto{\pgfqpoint{1.826181in}{1.563434in}}{\pgfqpoint{1.831767in}{1.561120in}}{\pgfqpoint{1.837591in}{1.561120in}}%
\pgfpathlineto{\pgfqpoint{1.837591in}{1.561120in}}%
\pgfpathclose%
\pgfusepath{stroke,fill}%
\end{pgfscope}%
\begin{pgfscope}%
\pgfpathrectangle{\pgfqpoint{0.100000in}{0.209042in}}{\pgfqpoint{2.906250in}{2.906250in}}%
\pgfusepath{clip}%
\pgfsetbuttcap%
\pgfsetroundjoin%
\definecolor{currentfill}{rgb}{0.531507,0.154739,0.506895}%
\pgfsetfillcolor{currentfill}%
\pgfsetfillopacity{0.800000}%
\pgfsetlinewidth{1.003750pt}%
\definecolor{currentstroke}{rgb}{0.531507,0.154739,0.506895}%
\pgfsetstrokecolor{currentstroke}%
\pgfsetstrokeopacity{0.800000}%
\pgfsetdash{}{0pt}%
\pgfpathmoveto{\pgfqpoint{2.022136in}{1.720776in}}%
\pgfpathcurveto{\pgfqpoint{2.027960in}{1.720776in}}{\pgfqpoint{2.033546in}{1.723090in}}{\pgfqpoint{2.037664in}{1.727208in}}%
\pgfpathcurveto{\pgfqpoint{2.041782in}{1.731326in}}{\pgfqpoint{2.044096in}{1.736913in}}{\pgfqpoint{2.044096in}{1.742736in}}%
\pgfpathcurveto{\pgfqpoint{2.044096in}{1.748560in}}{\pgfqpoint{2.041782in}{1.754147in}}{\pgfqpoint{2.037664in}{1.758265in}}%
\pgfpathcurveto{\pgfqpoint{2.033546in}{1.762383in}}{\pgfqpoint{2.027960in}{1.764697in}}{\pgfqpoint{2.022136in}{1.764697in}}%
\pgfpathcurveto{\pgfqpoint{2.016312in}{1.764697in}}{\pgfqpoint{2.010726in}{1.762383in}}{\pgfqpoint{2.006608in}{1.758265in}}%
\pgfpathcurveto{\pgfqpoint{2.002489in}{1.754147in}}{\pgfqpoint{2.000176in}{1.748560in}}{\pgfqpoint{2.000176in}{1.742736in}}%
\pgfpathcurveto{\pgfqpoint{2.000176in}{1.736913in}}{\pgfqpoint{2.002489in}{1.731326in}}{\pgfqpoint{2.006608in}{1.727208in}}%
\pgfpathcurveto{\pgfqpoint{2.010726in}{1.723090in}}{\pgfqpoint{2.016312in}{1.720776in}}{\pgfqpoint{2.022136in}{1.720776in}}%
\pgfpathlineto{\pgfqpoint{2.022136in}{1.720776in}}%
\pgfpathclose%
\pgfusepath{stroke,fill}%
\end{pgfscope}%
\begin{pgfscope}%
\pgfpathrectangle{\pgfqpoint{0.100000in}{0.209042in}}{\pgfqpoint{2.906250in}{2.906250in}}%
\pgfusepath{clip}%
\pgfsetbuttcap%
\pgfsetroundjoin%
\definecolor{currentfill}{rgb}{0.359898,0.087831,0.496778}%
\pgfsetfillcolor{currentfill}%
\pgfsetfillopacity{0.800000}%
\pgfsetlinewidth{1.003750pt}%
\definecolor{currentstroke}{rgb}{0.359898,0.087831,0.496778}%
\pgfsetstrokecolor{currentstroke}%
\pgfsetstrokeopacity{0.800000}%
\pgfsetdash{}{0pt}%
\pgfpathmoveto{\pgfqpoint{1.324044in}{1.789954in}}%
\pgfpathcurveto{\pgfqpoint{1.329868in}{1.789954in}}{\pgfqpoint{1.335454in}{1.792268in}}{\pgfqpoint{1.339572in}{1.796386in}}%
\pgfpathcurveto{\pgfqpoint{1.343690in}{1.800505in}}{\pgfqpoint{1.346004in}{1.806091in}}{\pgfqpoint{1.346004in}{1.811915in}}%
\pgfpathcurveto{\pgfqpoint{1.346004in}{1.817739in}}{\pgfqpoint{1.343690in}{1.823325in}}{\pgfqpoint{1.339572in}{1.827443in}}%
\pgfpathcurveto{\pgfqpoint{1.335454in}{1.831561in}}{\pgfqpoint{1.329868in}{1.833875in}}{\pgfqpoint{1.324044in}{1.833875in}}%
\pgfpathcurveto{\pgfqpoint{1.318220in}{1.833875in}}{\pgfqpoint{1.312634in}{1.831561in}}{\pgfqpoint{1.308516in}{1.827443in}}%
\pgfpathcurveto{\pgfqpoint{1.304397in}{1.823325in}}{\pgfqpoint{1.302084in}{1.817739in}}{\pgfqpoint{1.302084in}{1.811915in}}%
\pgfpathcurveto{\pgfqpoint{1.302084in}{1.806091in}}{\pgfqpoint{1.304397in}{1.800505in}}{\pgfqpoint{1.308516in}{1.796386in}}%
\pgfpathcurveto{\pgfqpoint{1.312634in}{1.792268in}}{\pgfqpoint{1.318220in}{1.789954in}}{\pgfqpoint{1.324044in}{1.789954in}}%
\pgfpathlineto{\pgfqpoint{1.324044in}{1.789954in}}%
\pgfpathclose%
\pgfusepath{stroke,fill}%
\end{pgfscope}%
\begin{pgfscope}%
\pgfpathrectangle{\pgfqpoint{0.100000in}{0.209042in}}{\pgfqpoint{2.906250in}{2.906250in}}%
\pgfusepath{clip}%
\pgfsetbuttcap%
\pgfsetroundjoin%
\definecolor{currentfill}{rgb}{0.697098,0.208501,0.480835}%
\pgfsetfillcolor{currentfill}%
\pgfsetfillopacity{0.800000}%
\pgfsetlinewidth{1.003750pt}%
\definecolor{currentstroke}{rgb}{0.697098,0.208501,0.480835}%
\pgfsetstrokecolor{currentstroke}%
\pgfsetstrokeopacity{0.800000}%
\pgfsetdash{}{0pt}%
\pgfpathmoveto{\pgfqpoint{2.083544in}{1.526250in}}%
\pgfpathcurveto{\pgfqpoint{2.089368in}{1.526250in}}{\pgfqpoint{2.094954in}{1.528564in}}{\pgfqpoint{2.099073in}{1.532682in}}%
\pgfpathcurveto{\pgfqpoint{2.103191in}{1.536800in}}{\pgfqpoint{2.105505in}{1.542386in}}{\pgfqpoint{2.105505in}{1.548210in}}%
\pgfpathcurveto{\pgfqpoint{2.105505in}{1.554034in}}{\pgfqpoint{2.103191in}{1.559620in}}{\pgfqpoint{2.099073in}{1.563738in}}%
\pgfpathcurveto{\pgfqpoint{2.094954in}{1.567856in}}{\pgfqpoint{2.089368in}{1.570170in}}{\pgfqpoint{2.083544in}{1.570170in}}%
\pgfpathcurveto{\pgfqpoint{2.077720in}{1.570170in}}{\pgfqpoint{2.072134in}{1.567856in}}{\pgfqpoint{2.068016in}{1.563738in}}%
\pgfpathcurveto{\pgfqpoint{2.063898in}{1.559620in}}{\pgfqpoint{2.061584in}{1.554034in}}{\pgfqpoint{2.061584in}{1.548210in}}%
\pgfpathcurveto{\pgfqpoint{2.061584in}{1.542386in}}{\pgfqpoint{2.063898in}{1.536800in}}{\pgfqpoint{2.068016in}{1.532682in}}%
\pgfpathcurveto{\pgfqpoint{2.072134in}{1.528564in}}{\pgfqpoint{2.077720in}{1.526250in}}{\pgfqpoint{2.083544in}{1.526250in}}%
\pgfpathlineto{\pgfqpoint{2.083544in}{1.526250in}}%
\pgfpathclose%
\pgfusepath{stroke,fill}%
\end{pgfscope}%
\begin{pgfscope}%
\pgfpathrectangle{\pgfqpoint{0.100000in}{0.209042in}}{\pgfqpoint{2.906250in}{2.906250in}}%
\pgfusepath{clip}%
\pgfsetbuttcap%
\pgfsetroundjoin%
\definecolor{currentfill}{rgb}{0.963310,0.425390,0.359469}%
\pgfsetfillcolor{currentfill}%
\pgfsetfillopacity{0.800000}%
\pgfsetlinewidth{1.003750pt}%
\definecolor{currentstroke}{rgb}{0.963310,0.425390,0.359469}%
\pgfsetstrokecolor{currentstroke}%
\pgfsetstrokeopacity{0.800000}%
\pgfsetdash{}{0pt}%
\pgfpathmoveto{\pgfqpoint{1.478462in}{1.139218in}}%
\pgfpathcurveto{\pgfqpoint{1.484286in}{1.139218in}}{\pgfqpoint{1.489872in}{1.141531in}}{\pgfqpoint{1.493990in}{1.145650in}}%
\pgfpathcurveto{\pgfqpoint{1.498108in}{1.149768in}}{\pgfqpoint{1.500422in}{1.155354in}}{\pgfqpoint{1.500422in}{1.161178in}}%
\pgfpathcurveto{\pgfqpoint{1.500422in}{1.167002in}}{\pgfqpoint{1.498108in}{1.172588in}}{\pgfqpoint{1.493990in}{1.176706in}}%
\pgfpathcurveto{\pgfqpoint{1.489872in}{1.180824in}}{\pgfqpoint{1.484286in}{1.183138in}}{\pgfqpoint{1.478462in}{1.183138in}}%
\pgfpathcurveto{\pgfqpoint{1.472638in}{1.183138in}}{\pgfqpoint{1.467052in}{1.180824in}}{\pgfqpoint{1.462934in}{1.176706in}}%
\pgfpathcurveto{\pgfqpoint{1.458816in}{1.172588in}}{\pgfqpoint{1.456502in}{1.167002in}}{\pgfqpoint{1.456502in}{1.161178in}}%
\pgfpathcurveto{\pgfqpoint{1.456502in}{1.155354in}}{\pgfqpoint{1.458816in}{1.149768in}}{\pgfqpoint{1.462934in}{1.145650in}}%
\pgfpathcurveto{\pgfqpoint{1.467052in}{1.141531in}}{\pgfqpoint{1.472638in}{1.139218in}}{\pgfqpoint{1.478462in}{1.139218in}}%
\pgfpathlineto{\pgfqpoint{1.478462in}{1.139218in}}%
\pgfpathclose%
\pgfusepath{stroke,fill}%
\end{pgfscope}%
\begin{pgfscope}%
\pgfpathrectangle{\pgfqpoint{0.100000in}{0.209042in}}{\pgfqpoint{2.906250in}{2.906250in}}%
\pgfusepath{clip}%
\pgfsetbuttcap%
\pgfsetroundjoin%
\definecolor{currentfill}{rgb}{0.798608,0.247040,0.444848}%
\pgfsetfillcolor{currentfill}%
\pgfsetfillopacity{0.800000}%
\pgfsetlinewidth{1.003750pt}%
\definecolor{currentstroke}{rgb}{0.798608,0.247040,0.444848}%
\pgfsetstrokecolor{currentstroke}%
\pgfsetstrokeopacity{0.800000}%
\pgfsetdash{}{0pt}%
\pgfpathmoveto{\pgfqpoint{1.108419in}{1.485964in}}%
\pgfpathcurveto{\pgfqpoint{1.114243in}{1.485964in}}{\pgfqpoint{1.119829in}{1.488278in}}{\pgfqpoint{1.123947in}{1.492396in}}%
\pgfpathcurveto{\pgfqpoint{1.128065in}{1.496514in}}{\pgfqpoint{1.130379in}{1.502101in}}{\pgfqpoint{1.130379in}{1.507925in}}%
\pgfpathcurveto{\pgfqpoint{1.130379in}{1.513749in}}{\pgfqpoint{1.128065in}{1.519335in}}{\pgfqpoint{1.123947in}{1.523453in}}%
\pgfpathcurveto{\pgfqpoint{1.119829in}{1.527571in}}{\pgfqpoint{1.114243in}{1.529885in}}{\pgfqpoint{1.108419in}{1.529885in}}%
\pgfpathcurveto{\pgfqpoint{1.102595in}{1.529885in}}{\pgfqpoint{1.097008in}{1.527571in}}{\pgfqpoint{1.092890in}{1.523453in}}%
\pgfpathcurveto{\pgfqpoint{1.088772in}{1.519335in}}{\pgfqpoint{1.086458in}{1.513749in}}{\pgfqpoint{1.086458in}{1.507925in}}%
\pgfpathcurveto{\pgfqpoint{1.086458in}{1.502101in}}{\pgfqpoint{1.088772in}{1.496514in}}{\pgfqpoint{1.092890in}{1.492396in}}%
\pgfpathcurveto{\pgfqpoint{1.097008in}{1.488278in}}{\pgfqpoint{1.102595in}{1.485964in}}{\pgfqpoint{1.108419in}{1.485964in}}%
\pgfpathlineto{\pgfqpoint{1.108419in}{1.485964in}}%
\pgfpathclose%
\pgfusepath{stroke,fill}%
\end{pgfscope}%
\begin{pgfscope}%
\pgfpathrectangle{\pgfqpoint{0.100000in}{0.209042in}}{\pgfqpoint{2.906250in}{2.906250in}}%
\pgfusepath{clip}%
\pgfsetbuttcap%
\pgfsetroundjoin%
\definecolor{currentfill}{rgb}{0.198177,0.063862,0.404009}%
\pgfsetfillcolor{currentfill}%
\pgfsetfillopacity{0.800000}%
\pgfsetlinewidth{1.003750pt}%
\definecolor{currentstroke}{rgb}{0.198177,0.063862,0.404009}%
\pgfsetstrokecolor{currentstroke}%
\pgfsetstrokeopacity{0.800000}%
\pgfsetdash{}{0pt}%
\pgfpathmoveto{\pgfqpoint{1.464517in}{1.803518in}}%
\pgfpathcurveto{\pgfqpoint{1.470341in}{1.803518in}}{\pgfqpoint{1.475927in}{1.805832in}}{\pgfqpoint{1.480045in}{1.809950in}}%
\pgfpathcurveto{\pgfqpoint{1.484164in}{1.814068in}}{\pgfqpoint{1.486477in}{1.819654in}}{\pgfqpoint{1.486477in}{1.825478in}}%
\pgfpathcurveto{\pgfqpoint{1.486477in}{1.831302in}}{\pgfqpoint{1.484164in}{1.836889in}}{\pgfqpoint{1.480045in}{1.841007in}}%
\pgfpathcurveto{\pgfqpoint{1.475927in}{1.845125in}}{\pgfqpoint{1.470341in}{1.847439in}}{\pgfqpoint{1.464517in}{1.847439in}}%
\pgfpathcurveto{\pgfqpoint{1.458693in}{1.847439in}}{\pgfqpoint{1.453107in}{1.845125in}}{\pgfqpoint{1.448989in}{1.841007in}}%
\pgfpathcurveto{\pgfqpoint{1.444871in}{1.836889in}}{\pgfqpoint{1.442557in}{1.831302in}}{\pgfqpoint{1.442557in}{1.825478in}}%
\pgfpathcurveto{\pgfqpoint{1.442557in}{1.819654in}}{\pgfqpoint{1.444871in}{1.814068in}}{\pgfqpoint{1.448989in}{1.809950in}}%
\pgfpathcurveto{\pgfqpoint{1.453107in}{1.805832in}}{\pgfqpoint{1.458693in}{1.803518in}}{\pgfqpoint{1.464517in}{1.803518in}}%
\pgfpathlineto{\pgfqpoint{1.464517in}{1.803518in}}%
\pgfpathclose%
\pgfusepath{stroke,fill}%
\end{pgfscope}%
\begin{pgfscope}%
\pgfpathrectangle{\pgfqpoint{0.100000in}{0.209042in}}{\pgfqpoint{2.906250in}{2.906250in}}%
\pgfusepath{clip}%
\pgfsetbuttcap%
\pgfsetroundjoin%
\definecolor{currentfill}{rgb}{0.963310,0.425390,0.359469}%
\pgfsetfillcolor{currentfill}%
\pgfsetfillopacity{0.800000}%
\pgfsetlinewidth{1.003750pt}%
\definecolor{currentstroke}{rgb}{0.963310,0.425390,0.359469}%
\pgfsetstrokecolor{currentstroke}%
\pgfsetstrokeopacity{0.800000}%
\pgfsetdash{}{0pt}%
\pgfpathmoveto{\pgfqpoint{1.554857in}{2.331603in}}%
\pgfpathcurveto{\pgfqpoint{1.560681in}{2.331603in}}{\pgfqpoint{1.566268in}{2.333917in}}{\pgfqpoint{1.570386in}{2.338035in}}%
\pgfpathcurveto{\pgfqpoint{1.574504in}{2.342153in}}{\pgfqpoint{1.576818in}{2.347740in}}{\pgfqpoint{1.576818in}{2.353563in}}%
\pgfpathcurveto{\pgfqpoint{1.576818in}{2.359387in}}{\pgfqpoint{1.574504in}{2.364974in}}{\pgfqpoint{1.570386in}{2.369092in}}%
\pgfpathcurveto{\pgfqpoint{1.566268in}{2.373210in}}{\pgfqpoint{1.560681in}{2.375524in}}{\pgfqpoint{1.554857in}{2.375524in}}%
\pgfpathcurveto{\pgfqpoint{1.549033in}{2.375524in}}{\pgfqpoint{1.543447in}{2.373210in}}{\pgfqpoint{1.539329in}{2.369092in}}%
\pgfpathcurveto{\pgfqpoint{1.535211in}{2.364974in}}{\pgfqpoint{1.532897in}{2.359387in}}{\pgfqpoint{1.532897in}{2.353563in}}%
\pgfpathcurveto{\pgfqpoint{1.532897in}{2.347740in}}{\pgfqpoint{1.535211in}{2.342153in}}{\pgfqpoint{1.539329in}{2.338035in}}%
\pgfpathcurveto{\pgfqpoint{1.543447in}{2.333917in}}{\pgfqpoint{1.549033in}{2.331603in}}{\pgfqpoint{1.554857in}{2.331603in}}%
\pgfpathlineto{\pgfqpoint{1.554857in}{2.331603in}}%
\pgfpathclose%
\pgfusepath{stroke,fill}%
\end{pgfscope}%
\begin{pgfscope}%
\pgfpathrectangle{\pgfqpoint{0.100000in}{0.209042in}}{\pgfqpoint{2.906250in}{2.906250in}}%
\pgfusepath{clip}%
\pgfsetbuttcap%
\pgfsetroundjoin%
\definecolor{currentfill}{rgb}{0.258857,0.059706,0.457710}%
\pgfsetfillcolor{currentfill}%
\pgfsetfillopacity{0.800000}%
\pgfsetlinewidth{1.003750pt}%
\definecolor{currentstroke}{rgb}{0.258857,0.059706,0.457710}%
\pgfsetstrokecolor{currentstroke}%
\pgfsetstrokeopacity{0.800000}%
\pgfsetdash{}{0pt}%
\pgfpathmoveto{\pgfqpoint{1.804323in}{1.790017in}}%
\pgfpathcurveto{\pgfqpoint{1.810147in}{1.790017in}}{\pgfqpoint{1.815733in}{1.792331in}}{\pgfqpoint{1.819852in}{1.796449in}}%
\pgfpathcurveto{\pgfqpoint{1.823970in}{1.800568in}}{\pgfqpoint{1.826284in}{1.806154in}}{\pgfqpoint{1.826284in}{1.811978in}}%
\pgfpathcurveto{\pgfqpoint{1.826284in}{1.817802in}}{\pgfqpoint{1.823970in}{1.823388in}}{\pgfqpoint{1.819852in}{1.827506in}}%
\pgfpathcurveto{\pgfqpoint{1.815733in}{1.831624in}}{\pgfqpoint{1.810147in}{1.833938in}}{\pgfqpoint{1.804323in}{1.833938in}}%
\pgfpathcurveto{\pgfqpoint{1.798499in}{1.833938in}}{\pgfqpoint{1.792913in}{1.831624in}}{\pgfqpoint{1.788795in}{1.827506in}}%
\pgfpathcurveto{\pgfqpoint{1.784677in}{1.823388in}}{\pgfqpoint{1.782363in}{1.817802in}}{\pgfqpoint{1.782363in}{1.811978in}}%
\pgfpathcurveto{\pgfqpoint{1.782363in}{1.806154in}}{\pgfqpoint{1.784677in}{1.800568in}}{\pgfqpoint{1.788795in}{1.796449in}}%
\pgfpathcurveto{\pgfqpoint{1.792913in}{1.792331in}}{\pgfqpoint{1.798499in}{1.790017in}}{\pgfqpoint{1.804323in}{1.790017in}}%
\pgfpathlineto{\pgfqpoint{1.804323in}{1.790017in}}%
\pgfpathclose%
\pgfusepath{stroke,fill}%
\end{pgfscope}%
\begin{pgfscope}%
\pgfpathrectangle{\pgfqpoint{0.100000in}{0.209042in}}{\pgfqpoint{2.906250in}{2.906250in}}%
\pgfusepath{clip}%
\pgfsetbuttcap%
\pgfsetroundjoin%
\definecolor{currentfill}{rgb}{0.378211,0.095332,0.500067}%
\pgfsetfillcolor{currentfill}%
\pgfsetfillopacity{0.800000}%
\pgfsetlinewidth{1.003750pt}%
\definecolor{currentstroke}{rgb}{0.378211,0.095332,0.500067}%
\pgfsetstrokecolor{currentstroke}%
\pgfsetstrokeopacity{0.800000}%
\pgfsetdash{}{0pt}%
\pgfpathmoveto{\pgfqpoint{1.914515in}{1.673447in}}%
\pgfpathcurveto{\pgfqpoint{1.920339in}{1.673447in}}{\pgfqpoint{1.925925in}{1.675761in}}{\pgfqpoint{1.930044in}{1.679879in}}%
\pgfpathcurveto{\pgfqpoint{1.934162in}{1.683997in}}{\pgfqpoint{1.936476in}{1.689583in}}{\pgfqpoint{1.936476in}{1.695407in}}%
\pgfpathcurveto{\pgfqpoint{1.936476in}{1.701231in}}{\pgfqpoint{1.934162in}{1.706817in}}{\pgfqpoint{1.930044in}{1.710935in}}%
\pgfpathcurveto{\pgfqpoint{1.925925in}{1.715053in}}{\pgfqpoint{1.920339in}{1.717367in}}{\pgfqpoint{1.914515in}{1.717367in}}%
\pgfpathcurveto{\pgfqpoint{1.908691in}{1.717367in}}{\pgfqpoint{1.903105in}{1.715053in}}{\pgfqpoint{1.898987in}{1.710935in}}%
\pgfpathcurveto{\pgfqpoint{1.894869in}{1.706817in}}{\pgfqpoint{1.892555in}{1.701231in}}{\pgfqpoint{1.892555in}{1.695407in}}%
\pgfpathcurveto{\pgfqpoint{1.892555in}{1.689583in}}{\pgfqpoint{1.894869in}{1.683997in}}{\pgfqpoint{1.898987in}{1.679879in}}%
\pgfpathcurveto{\pgfqpoint{1.903105in}{1.675761in}}{\pgfqpoint{1.908691in}{1.673447in}}{\pgfqpoint{1.914515in}{1.673447in}}%
\pgfpathlineto{\pgfqpoint{1.914515in}{1.673447in}}%
\pgfpathclose%
\pgfusepath{stroke,fill}%
\end{pgfscope}%
\begin{pgfscope}%
\pgfpathrectangle{\pgfqpoint{0.100000in}{0.209042in}}{\pgfqpoint{2.906250in}{2.906250in}}%
\pgfusepath{clip}%
\pgfsetbuttcap%
\pgfsetroundjoin%
\definecolor{currentfill}{rgb}{0.965549,0.432519,0.359529}%
\pgfsetfillcolor{currentfill}%
\pgfsetfillopacity{0.800000}%
\pgfsetlinewidth{1.003750pt}%
\definecolor{currentstroke}{rgb}{0.965549,0.432519,0.359529}%
\pgfsetstrokecolor{currentstroke}%
\pgfsetstrokeopacity{0.800000}%
\pgfsetdash{}{0pt}%
\pgfpathmoveto{\pgfqpoint{1.799677in}{1.129295in}}%
\pgfpathcurveto{\pgfqpoint{1.805501in}{1.129295in}}{\pgfqpoint{1.811087in}{1.131609in}}{\pgfqpoint{1.815205in}{1.135727in}}%
\pgfpathcurveto{\pgfqpoint{1.819324in}{1.139845in}}{\pgfqpoint{1.821638in}{1.145431in}}{\pgfqpoint{1.821638in}{1.151255in}}%
\pgfpathcurveto{\pgfqpoint{1.821638in}{1.157079in}}{\pgfqpoint{1.819324in}{1.162665in}}{\pgfqpoint{1.815205in}{1.166783in}}%
\pgfpathcurveto{\pgfqpoint{1.811087in}{1.170901in}}{\pgfqpoint{1.805501in}{1.173215in}}{\pgfqpoint{1.799677in}{1.173215in}}%
\pgfpathcurveto{\pgfqpoint{1.793853in}{1.173215in}}{\pgfqpoint{1.788267in}{1.170901in}}{\pgfqpoint{1.784149in}{1.166783in}}%
\pgfpathcurveto{\pgfqpoint{1.780031in}{1.162665in}}{\pgfqpoint{1.777717in}{1.157079in}}{\pgfqpoint{1.777717in}{1.151255in}}%
\pgfpathcurveto{\pgfqpoint{1.777717in}{1.145431in}}{\pgfqpoint{1.780031in}{1.139845in}}{\pgfqpoint{1.784149in}{1.135727in}}%
\pgfpathcurveto{\pgfqpoint{1.788267in}{1.131609in}}{\pgfqpoint{1.793853in}{1.129295in}}{\pgfqpoint{1.799677in}{1.129295in}}%
\pgfpathlineto{\pgfqpoint{1.799677in}{1.129295in}}%
\pgfpathclose%
\pgfusepath{stroke,fill}%
\end{pgfscope}%
\begin{pgfscope}%
\pgfpathrectangle{\pgfqpoint{0.100000in}{0.209042in}}{\pgfqpoint{2.906250in}{2.906250in}}%
\pgfusepath{clip}%
\pgfsetbuttcap%
\pgfsetroundjoin%
\definecolor{currentfill}{rgb}{0.884651,0.300530,0.400047}%
\pgfsetfillcolor{currentfill}%
\pgfsetfillopacity{0.800000}%
\pgfsetlinewidth{1.003750pt}%
\definecolor{currentstroke}{rgb}{0.884651,0.300530,0.400047}%
\pgfsetstrokecolor{currentstroke}%
\pgfsetstrokeopacity{0.800000}%
\pgfsetdash{}{0pt}%
\pgfpathmoveto{\pgfqpoint{1.461782in}{1.208779in}}%
\pgfpathcurveto{\pgfqpoint{1.467606in}{1.208779in}}{\pgfqpoint{1.473192in}{1.211092in}}{\pgfqpoint{1.477311in}{1.215211in}}%
\pgfpathcurveto{\pgfqpoint{1.481429in}{1.219329in}}{\pgfqpoint{1.483743in}{1.224915in}}{\pgfqpoint{1.483743in}{1.230739in}}%
\pgfpathcurveto{\pgfqpoint{1.483743in}{1.236563in}}{\pgfqpoint{1.481429in}{1.242149in}}{\pgfqpoint{1.477311in}{1.246267in}}%
\pgfpathcurveto{\pgfqpoint{1.473192in}{1.250385in}}{\pgfqpoint{1.467606in}{1.252699in}}{\pgfqpoint{1.461782in}{1.252699in}}%
\pgfpathcurveto{\pgfqpoint{1.455958in}{1.252699in}}{\pgfqpoint{1.450372in}{1.250385in}}{\pgfqpoint{1.446254in}{1.246267in}}%
\pgfpathcurveto{\pgfqpoint{1.442136in}{1.242149in}}{\pgfqpoint{1.439822in}{1.236563in}}{\pgfqpoint{1.439822in}{1.230739in}}%
\pgfpathcurveto{\pgfqpoint{1.439822in}{1.224915in}}{\pgfqpoint{1.442136in}{1.219329in}}{\pgfqpoint{1.446254in}{1.215211in}}%
\pgfpathcurveto{\pgfqpoint{1.450372in}{1.211092in}}{\pgfqpoint{1.455958in}{1.208779in}}{\pgfqpoint{1.461782in}{1.208779in}}%
\pgfpathlineto{\pgfqpoint{1.461782in}{1.208779in}}%
\pgfpathclose%
\pgfusepath{stroke,fill}%
\end{pgfscope}%
\begin{pgfscope}%
\pgfpathrectangle{\pgfqpoint{0.100000in}{0.209042in}}{\pgfqpoint{2.906250in}{2.906250in}}%
\pgfusepath{clip}%
\pgfsetbuttcap%
\pgfsetroundjoin%
\definecolor{currentfill}{rgb}{0.129380,0.067935,0.305443}%
\pgfsetfillcolor{currentfill}%
\pgfsetfillopacity{0.800000}%
\pgfsetlinewidth{1.003750pt}%
\definecolor{currentstroke}{rgb}{0.129380,0.067935,0.305443}%
\pgfsetstrokecolor{currentstroke}%
\pgfsetstrokeopacity{0.800000}%
\pgfsetdash{}{0pt}%
\pgfpathmoveto{\pgfqpoint{1.688634in}{1.783873in}}%
\pgfpathcurveto{\pgfqpoint{1.694458in}{1.783873in}}{\pgfqpoint{1.700044in}{1.786187in}}{\pgfqpoint{1.704162in}{1.790305in}}%
\pgfpathcurveto{\pgfqpoint{1.708280in}{1.794423in}}{\pgfqpoint{1.710594in}{1.800009in}}{\pgfqpoint{1.710594in}{1.805833in}}%
\pgfpathcurveto{\pgfqpoint{1.710594in}{1.811657in}}{\pgfqpoint{1.708280in}{1.817243in}}{\pgfqpoint{1.704162in}{1.821361in}}%
\pgfpathcurveto{\pgfqpoint{1.700044in}{1.825480in}}{\pgfqpoint{1.694458in}{1.827794in}}{\pgfqpoint{1.688634in}{1.827794in}}%
\pgfpathcurveto{\pgfqpoint{1.682810in}{1.827794in}}{\pgfqpoint{1.677224in}{1.825480in}}{\pgfqpoint{1.673106in}{1.821361in}}%
\pgfpathcurveto{\pgfqpoint{1.668987in}{1.817243in}}{\pgfqpoint{1.666674in}{1.811657in}}{\pgfqpoint{1.666674in}{1.805833in}}%
\pgfpathcurveto{\pgfqpoint{1.666674in}{1.800009in}}{\pgfqpoint{1.668987in}{1.794423in}}{\pgfqpoint{1.673106in}{1.790305in}}%
\pgfpathcurveto{\pgfqpoint{1.677224in}{1.786187in}}{\pgfqpoint{1.682810in}{1.783873in}}{\pgfqpoint{1.688634in}{1.783873in}}%
\pgfpathlineto{\pgfqpoint{1.688634in}{1.783873in}}%
\pgfpathclose%
\pgfusepath{stroke,fill}%
\end{pgfscope}%
\begin{pgfscope}%
\pgfpathrectangle{\pgfqpoint{0.100000in}{0.209042in}}{\pgfqpoint{2.906250in}{2.906250in}}%
\pgfusepath{clip}%
\pgfsetbuttcap%
\pgfsetroundjoin%
\definecolor{currentfill}{rgb}{0.773695,0.236249,0.455289}%
\pgfsetfillcolor{currentfill}%
\pgfsetfillopacity{0.800000}%
\pgfsetlinewidth{1.003750pt}%
\definecolor{currentstroke}{rgb}{0.773695,0.236249,0.455289}%
\pgfsetstrokecolor{currentstroke}%
\pgfsetstrokeopacity{0.800000}%
\pgfsetdash{}{0pt}%
\pgfpathmoveto{\pgfqpoint{1.563620in}{2.199703in}}%
\pgfpathcurveto{\pgfqpoint{1.569444in}{2.199703in}}{\pgfqpoint{1.575030in}{2.202017in}}{\pgfqpoint{1.579148in}{2.206135in}}%
\pgfpathcurveto{\pgfqpoint{1.583266in}{2.210253in}}{\pgfqpoint{1.585580in}{2.215840in}}{\pgfqpoint{1.585580in}{2.221663in}}%
\pgfpathcurveto{\pgfqpoint{1.585580in}{2.227487in}}{\pgfqpoint{1.583266in}{2.233074in}}{\pgfqpoint{1.579148in}{2.237192in}}%
\pgfpathcurveto{\pgfqpoint{1.575030in}{2.241310in}}{\pgfqpoint{1.569444in}{2.243624in}}{\pgfqpoint{1.563620in}{2.243624in}}%
\pgfpathcurveto{\pgfqpoint{1.557796in}{2.243624in}}{\pgfqpoint{1.552210in}{2.241310in}}{\pgfqpoint{1.548092in}{2.237192in}}%
\pgfpathcurveto{\pgfqpoint{1.543973in}{2.233074in}}{\pgfqpoint{1.541660in}{2.227487in}}{\pgfqpoint{1.541660in}{2.221663in}}%
\pgfpathcurveto{\pgfqpoint{1.541660in}{2.215840in}}{\pgfqpoint{1.543973in}{2.210253in}}{\pgfqpoint{1.548092in}{2.206135in}}%
\pgfpathcurveto{\pgfqpoint{1.552210in}{2.202017in}}{\pgfqpoint{1.557796in}{2.199703in}}{\pgfqpoint{1.563620in}{2.199703in}}%
\pgfpathlineto{\pgfqpoint{1.563620in}{2.199703in}}%
\pgfpathclose%
\pgfusepath{stroke,fill}%
\end{pgfscope}%
\begin{pgfscope}%
\pgfpathrectangle{\pgfqpoint{0.100000in}{0.209042in}}{\pgfqpoint{2.906250in}{2.906250in}}%
\pgfusepath{clip}%
\pgfsetbuttcap%
\pgfsetroundjoin%
\definecolor{currentfill}{rgb}{0.316654,0.071690,0.485380}%
\pgfsetfillcolor{currentfill}%
\pgfsetfillopacity{0.800000}%
\pgfsetlinewidth{1.003750pt}%
\definecolor{currentstroke}{rgb}{0.316654,0.071690,0.485380}%
\pgfsetstrokecolor{currentstroke}%
\pgfsetstrokeopacity{0.800000}%
\pgfsetdash{}{0pt}%
\pgfpathmoveto{\pgfqpoint{1.498627in}{1.931815in}}%
\pgfpathcurveto{\pgfqpoint{1.504451in}{1.931815in}}{\pgfqpoint{1.510037in}{1.934129in}}{\pgfqpoint{1.514155in}{1.938247in}}%
\pgfpathcurveto{\pgfqpoint{1.518273in}{1.942365in}}{\pgfqpoint{1.520587in}{1.947951in}}{\pgfqpoint{1.520587in}{1.953775in}}%
\pgfpathcurveto{\pgfqpoint{1.520587in}{1.959599in}}{\pgfqpoint{1.518273in}{1.965186in}}{\pgfqpoint{1.514155in}{1.969304in}}%
\pgfpathcurveto{\pgfqpoint{1.510037in}{1.973422in}}{\pgfqpoint{1.504451in}{1.975736in}}{\pgfqpoint{1.498627in}{1.975736in}}%
\pgfpathcurveto{\pgfqpoint{1.492803in}{1.975736in}}{\pgfqpoint{1.487217in}{1.973422in}}{\pgfqpoint{1.483099in}{1.969304in}}%
\pgfpathcurveto{\pgfqpoint{1.478981in}{1.965186in}}{\pgfqpoint{1.476667in}{1.959599in}}{\pgfqpoint{1.476667in}{1.953775in}}%
\pgfpathcurveto{\pgfqpoint{1.476667in}{1.947951in}}{\pgfqpoint{1.478981in}{1.942365in}}{\pgfqpoint{1.483099in}{1.938247in}}%
\pgfpathcurveto{\pgfqpoint{1.487217in}{1.934129in}}{\pgfqpoint{1.492803in}{1.931815in}}{\pgfqpoint{1.498627in}{1.931815in}}%
\pgfpathlineto{\pgfqpoint{1.498627in}{1.931815in}}%
\pgfpathclose%
\pgfusepath{stroke,fill}%
\end{pgfscope}%
\begin{pgfscope}%
\pgfpathrectangle{\pgfqpoint{0.100000in}{0.209042in}}{\pgfqpoint{2.906250in}{2.906250in}}%
\pgfusepath{clip}%
\pgfsetbuttcap%
\pgfsetroundjoin%
\definecolor{currentfill}{rgb}{0.773695,0.236249,0.455289}%
\pgfsetfillcolor{currentfill}%
\pgfsetfillopacity{0.800000}%
\pgfsetlinewidth{1.003750pt}%
\definecolor{currentstroke}{rgb}{0.773695,0.236249,0.455289}%
\pgfsetstrokecolor{currentstroke}%
\pgfsetstrokeopacity{0.800000}%
\pgfsetdash{}{0pt}%
\pgfpathmoveto{\pgfqpoint{2.022488in}{2.040722in}}%
\pgfpathcurveto{\pgfqpoint{2.028312in}{2.040722in}}{\pgfqpoint{2.033898in}{2.043036in}}{\pgfqpoint{2.038016in}{2.047154in}}%
\pgfpathcurveto{\pgfqpoint{2.042134in}{2.051272in}}{\pgfqpoint{2.044448in}{2.056858in}}{\pgfqpoint{2.044448in}{2.062682in}}%
\pgfpathcurveto{\pgfqpoint{2.044448in}{2.068506in}}{\pgfqpoint{2.042134in}{2.074092in}}{\pgfqpoint{2.038016in}{2.078211in}}%
\pgfpathcurveto{\pgfqpoint{2.033898in}{2.082329in}}{\pgfqpoint{2.028312in}{2.084643in}}{\pgfqpoint{2.022488in}{2.084643in}}%
\pgfpathcurveto{\pgfqpoint{2.016664in}{2.084643in}}{\pgfqpoint{2.011078in}{2.082329in}}{\pgfqpoint{2.006959in}{2.078211in}}%
\pgfpathcurveto{\pgfqpoint{2.002841in}{2.074092in}}{\pgfqpoint{2.000527in}{2.068506in}}{\pgfqpoint{2.000527in}{2.062682in}}%
\pgfpathcurveto{\pgfqpoint{2.000527in}{2.056858in}}{\pgfqpoint{2.002841in}{2.051272in}}{\pgfqpoint{2.006959in}{2.047154in}}%
\pgfpathcurveto{\pgfqpoint{2.011078in}{2.043036in}}{\pgfqpoint{2.016664in}{2.040722in}}{\pgfqpoint{2.022488in}{2.040722in}}%
\pgfpathlineto{\pgfqpoint{2.022488in}{2.040722in}}%
\pgfpathclose%
\pgfusepath{stroke,fill}%
\end{pgfscope}%
\begin{pgfscope}%
\pgfpathrectangle{\pgfqpoint{0.100000in}{0.209042in}}{\pgfqpoint{2.906250in}{2.906250in}}%
\pgfusepath{clip}%
\pgfsetbuttcap%
\pgfsetroundjoin%
\definecolor{currentfill}{rgb}{0.989363,0.557873,0.391671}%
\pgfsetfillcolor{currentfill}%
\pgfsetfillopacity{0.800000}%
\pgfsetlinewidth{1.003750pt}%
\definecolor{currentstroke}{rgb}{0.989363,0.557873,0.391671}%
\pgfsetstrokecolor{currentstroke}%
\pgfsetstrokeopacity{0.800000}%
\pgfsetdash{}{0pt}%
\pgfpathmoveto{\pgfqpoint{0.791317in}{1.838972in}}%
\pgfpathcurveto{\pgfqpoint{0.797141in}{1.838972in}}{\pgfqpoint{0.802727in}{1.841286in}}{\pgfqpoint{0.806845in}{1.845404in}}%
\pgfpathcurveto{\pgfqpoint{0.810964in}{1.849522in}}{\pgfqpoint{0.813277in}{1.855109in}}{\pgfqpoint{0.813277in}{1.860933in}}%
\pgfpathcurveto{\pgfqpoint{0.813277in}{1.866756in}}{\pgfqpoint{0.810964in}{1.872343in}}{\pgfqpoint{0.806845in}{1.876461in}}%
\pgfpathcurveto{\pgfqpoint{0.802727in}{1.880579in}}{\pgfqpoint{0.797141in}{1.882893in}}{\pgfqpoint{0.791317in}{1.882893in}}%
\pgfpathcurveto{\pgfqpoint{0.785493in}{1.882893in}}{\pgfqpoint{0.779907in}{1.880579in}}{\pgfqpoint{0.775789in}{1.876461in}}%
\pgfpathcurveto{\pgfqpoint{0.771671in}{1.872343in}}{\pgfqpoint{0.769357in}{1.866756in}}{\pgfqpoint{0.769357in}{1.860933in}}%
\pgfpathcurveto{\pgfqpoint{0.769357in}{1.855109in}}{\pgfqpoint{0.771671in}{1.849522in}}{\pgfqpoint{0.775789in}{1.845404in}}%
\pgfpathcurveto{\pgfqpoint{0.779907in}{1.841286in}}{\pgfqpoint{0.785493in}{1.838972in}}{\pgfqpoint{0.791317in}{1.838972in}}%
\pgfpathlineto{\pgfqpoint{0.791317in}{1.838972in}}%
\pgfpathclose%
\pgfusepath{stroke,fill}%
\end{pgfscope}%
\begin{pgfscope}%
\pgfpathrectangle{\pgfqpoint{0.100000in}{0.209042in}}{\pgfqpoint{2.906250in}{2.906250in}}%
\pgfusepath{clip}%
\pgfsetbuttcap%
\pgfsetroundjoin%
\definecolor{currentfill}{rgb}{0.716387,0.214982,0.475290}%
\pgfsetfillcolor{currentfill}%
\pgfsetfillopacity{0.800000}%
\pgfsetlinewidth{1.003750pt}%
\definecolor{currentstroke}{rgb}{0.716387,0.214982,0.475290}%
\pgfsetstrokecolor{currentstroke}%
\pgfsetstrokeopacity{0.800000}%
\pgfsetdash{}{0pt}%
\pgfpathmoveto{\pgfqpoint{1.415545in}{2.148045in}}%
\pgfpathcurveto{\pgfqpoint{1.421369in}{2.148045in}}{\pgfqpoint{1.426955in}{2.150359in}}{\pgfqpoint{1.431074in}{2.154477in}}%
\pgfpathcurveto{\pgfqpoint{1.435192in}{2.158595in}}{\pgfqpoint{1.437506in}{2.164181in}}{\pgfqpoint{1.437506in}{2.170005in}}%
\pgfpathcurveto{\pgfqpoint{1.437506in}{2.175829in}}{\pgfqpoint{1.435192in}{2.181415in}}{\pgfqpoint{1.431074in}{2.185533in}}%
\pgfpathcurveto{\pgfqpoint{1.426955in}{2.189651in}}{\pgfqpoint{1.421369in}{2.191965in}}{\pgfqpoint{1.415545in}{2.191965in}}%
\pgfpathcurveto{\pgfqpoint{1.409721in}{2.191965in}}{\pgfqpoint{1.404135in}{2.189651in}}{\pgfqpoint{1.400017in}{2.185533in}}%
\pgfpathcurveto{\pgfqpoint{1.395899in}{2.181415in}}{\pgfqpoint{1.393585in}{2.175829in}}{\pgfqpoint{1.393585in}{2.170005in}}%
\pgfpathcurveto{\pgfqpoint{1.393585in}{2.164181in}}{\pgfqpoint{1.395899in}{2.158595in}}{\pgfqpoint{1.400017in}{2.154477in}}%
\pgfpathcurveto{\pgfqpoint{1.404135in}{2.150359in}}{\pgfqpoint{1.409721in}{2.148045in}}{\pgfqpoint{1.415545in}{2.148045in}}%
\pgfpathlineto{\pgfqpoint{1.415545in}{2.148045in}}%
\pgfpathclose%
\pgfusepath{stroke,fill}%
\end{pgfscope}%
\begin{pgfscope}%
\pgfpathrectangle{\pgfqpoint{0.100000in}{0.209042in}}{\pgfqpoint{2.906250in}{2.906250in}}%
\pgfusepath{clip}%
\pgfsetbuttcap%
\pgfsetroundjoin%
\definecolor{currentfill}{rgb}{0.569172,0.167454,0.504105}%
\pgfsetfillcolor{currentfill}%
\pgfsetfillopacity{0.800000}%
\pgfsetlinewidth{1.003750pt}%
\definecolor{currentstroke}{rgb}{0.569172,0.167454,0.504105}%
\pgfsetstrokecolor{currentstroke}%
\pgfsetstrokeopacity{0.800000}%
\pgfsetdash{}{0pt}%
\pgfpathmoveto{\pgfqpoint{1.729533in}{1.364166in}}%
\pgfpathcurveto{\pgfqpoint{1.735357in}{1.364166in}}{\pgfqpoint{1.740944in}{1.366480in}}{\pgfqpoint{1.745062in}{1.370598in}}%
\pgfpathcurveto{\pgfqpoint{1.749180in}{1.374716in}}{\pgfqpoint{1.751494in}{1.380302in}}{\pgfqpoint{1.751494in}{1.386126in}}%
\pgfpathcurveto{\pgfqpoint{1.751494in}{1.391950in}}{\pgfqpoint{1.749180in}{1.397536in}}{\pgfqpoint{1.745062in}{1.401654in}}%
\pgfpathcurveto{\pgfqpoint{1.740944in}{1.405772in}}{\pgfqpoint{1.735357in}{1.408086in}}{\pgfqpoint{1.729533in}{1.408086in}}%
\pgfpathcurveto{\pgfqpoint{1.723709in}{1.408086in}}{\pgfqpoint{1.718123in}{1.405772in}}{\pgfqpoint{1.714005in}{1.401654in}}%
\pgfpathcurveto{\pgfqpoint{1.709887in}{1.397536in}}{\pgfqpoint{1.707573in}{1.391950in}}{\pgfqpoint{1.707573in}{1.386126in}}%
\pgfpathcurveto{\pgfqpoint{1.707573in}{1.380302in}}{\pgfqpoint{1.709887in}{1.374716in}}{\pgfqpoint{1.714005in}{1.370598in}}%
\pgfpathcurveto{\pgfqpoint{1.718123in}{1.366480in}}{\pgfqpoint{1.723709in}{1.364166in}}{\pgfqpoint{1.729533in}{1.364166in}}%
\pgfpathlineto{\pgfqpoint{1.729533in}{1.364166in}}%
\pgfpathclose%
\pgfusepath{stroke,fill}%
\end{pgfscope}%
\begin{pgfscope}%
\pgfpathrectangle{\pgfqpoint{0.100000in}{0.209042in}}{\pgfqpoint{2.906250in}{2.906250in}}%
\pgfusepath{clip}%
\pgfsetbuttcap%
\pgfsetroundjoin%
\definecolor{currentfill}{rgb}{0.469640,0.132245,0.507809}%
\pgfsetfillcolor{currentfill}%
\pgfsetfillopacity{0.800000}%
\pgfsetlinewidth{1.003750pt}%
\definecolor{currentstroke}{rgb}{0.469640,0.132245,0.507809}%
\pgfsetstrokecolor{currentstroke}%
\pgfsetstrokeopacity{0.800000}%
\pgfsetdash{}{0pt}%
\pgfpathmoveto{\pgfqpoint{1.232976in}{1.671663in}}%
\pgfpathcurveto{\pgfqpoint{1.238800in}{1.671663in}}{\pgfqpoint{1.244386in}{1.673977in}}{\pgfqpoint{1.248504in}{1.678095in}}%
\pgfpathcurveto{\pgfqpoint{1.252622in}{1.682213in}}{\pgfqpoint{1.254936in}{1.687800in}}{\pgfqpoint{1.254936in}{1.693623in}}%
\pgfpathcurveto{\pgfqpoint{1.254936in}{1.699447in}}{\pgfqpoint{1.252622in}{1.705034in}}{\pgfqpoint{1.248504in}{1.709152in}}%
\pgfpathcurveto{\pgfqpoint{1.244386in}{1.713270in}}{\pgfqpoint{1.238800in}{1.715584in}}{\pgfqpoint{1.232976in}{1.715584in}}%
\pgfpathcurveto{\pgfqpoint{1.227152in}{1.715584in}}{\pgfqpoint{1.221566in}{1.713270in}}{\pgfqpoint{1.217447in}{1.709152in}}%
\pgfpathcurveto{\pgfqpoint{1.213329in}{1.705034in}}{\pgfqpoint{1.211015in}{1.699447in}}{\pgfqpoint{1.211015in}{1.693623in}}%
\pgfpathcurveto{\pgfqpoint{1.211015in}{1.687800in}}{\pgfqpoint{1.213329in}{1.682213in}}{\pgfqpoint{1.217447in}{1.678095in}}%
\pgfpathcurveto{\pgfqpoint{1.221566in}{1.673977in}}{\pgfqpoint{1.227152in}{1.671663in}}{\pgfqpoint{1.232976in}{1.671663in}}%
\pgfpathlineto{\pgfqpoint{1.232976in}{1.671663in}}%
\pgfpathclose%
\pgfusepath{stroke,fill}%
\end{pgfscope}%
\begin{pgfscope}%
\pgfpathrectangle{\pgfqpoint{0.100000in}{0.209042in}}{\pgfqpoint{2.906250in}{2.906250in}}%
\pgfusepath{clip}%
\pgfsetbuttcap%
\pgfsetroundjoin%
\definecolor{currentfill}{rgb}{0.366012,0.090314,0.497960}%
\pgfsetfillcolor{currentfill}%
\pgfsetfillopacity{0.800000}%
\pgfsetlinewidth{1.003750pt}%
\definecolor{currentstroke}{rgb}{0.366012,0.090314,0.497960}%
\pgfsetstrokecolor{currentstroke}%
\pgfsetstrokeopacity{0.800000}%
\pgfsetdash{}{0pt}%
\pgfpathmoveto{\pgfqpoint{1.744698in}{1.945632in}}%
\pgfpathcurveto{\pgfqpoint{1.750521in}{1.945632in}}{\pgfqpoint{1.756108in}{1.947946in}}{\pgfqpoint{1.760226in}{1.952064in}}%
\pgfpathcurveto{\pgfqpoint{1.764344in}{1.956182in}}{\pgfqpoint{1.766658in}{1.961768in}}{\pgfqpoint{1.766658in}{1.967592in}}%
\pgfpathcurveto{\pgfqpoint{1.766658in}{1.973416in}}{\pgfqpoint{1.764344in}{1.979002in}}{\pgfqpoint{1.760226in}{1.983121in}}%
\pgfpathcurveto{\pgfqpoint{1.756108in}{1.987239in}}{\pgfqpoint{1.750521in}{1.989553in}}{\pgfqpoint{1.744698in}{1.989553in}}%
\pgfpathcurveto{\pgfqpoint{1.738874in}{1.989553in}}{\pgfqpoint{1.733287in}{1.987239in}}{\pgfqpoint{1.729169in}{1.983121in}}%
\pgfpathcurveto{\pgfqpoint{1.725051in}{1.979002in}}{\pgfqpoint{1.722737in}{1.973416in}}{\pgfqpoint{1.722737in}{1.967592in}}%
\pgfpathcurveto{\pgfqpoint{1.722737in}{1.961768in}}{\pgfqpoint{1.725051in}{1.956182in}}{\pgfqpoint{1.729169in}{1.952064in}}%
\pgfpathcurveto{\pgfqpoint{1.733287in}{1.947946in}}{\pgfqpoint{1.738874in}{1.945632in}}{\pgfqpoint{1.744698in}{1.945632in}}%
\pgfpathlineto{\pgfqpoint{1.744698in}{1.945632in}}%
\pgfpathclose%
\pgfusepath{stroke,fill}%
\end{pgfscope}%
\begin{pgfscope}%
\pgfpathrectangle{\pgfqpoint{0.100000in}{0.209042in}}{\pgfqpoint{2.906250in}{2.906250in}}%
\pgfusepath{clip}%
\pgfsetbuttcap%
\pgfsetroundjoin%
\definecolor{currentfill}{rgb}{0.445163,0.122724,0.506901}%
\pgfsetfillcolor{currentfill}%
\pgfsetfillopacity{0.800000}%
\pgfsetlinewidth{1.003750pt}%
\definecolor{currentstroke}{rgb}{0.445163,0.122724,0.506901}%
\pgfsetstrokecolor{currentstroke}%
\pgfsetstrokeopacity{0.800000}%
\pgfsetdash{}{0pt}%
\pgfpathmoveto{\pgfqpoint{1.280701in}{1.875762in}}%
\pgfpathcurveto{\pgfqpoint{1.286525in}{1.875762in}}{\pgfqpoint{1.292111in}{1.878076in}}{\pgfqpoint{1.296229in}{1.882194in}}%
\pgfpathcurveto{\pgfqpoint{1.300348in}{1.886312in}}{\pgfqpoint{1.302661in}{1.891898in}}{\pgfqpoint{1.302661in}{1.897722in}}%
\pgfpathcurveto{\pgfqpoint{1.302661in}{1.903546in}}{\pgfqpoint{1.300348in}{1.909132in}}{\pgfqpoint{1.296229in}{1.913250in}}%
\pgfpathcurveto{\pgfqpoint{1.292111in}{1.917369in}}{\pgfqpoint{1.286525in}{1.919682in}}{\pgfqpoint{1.280701in}{1.919682in}}%
\pgfpathcurveto{\pgfqpoint{1.274877in}{1.919682in}}{\pgfqpoint{1.269291in}{1.917369in}}{\pgfqpoint{1.265173in}{1.913250in}}%
\pgfpathcurveto{\pgfqpoint{1.261055in}{1.909132in}}{\pgfqpoint{1.258741in}{1.903546in}}{\pgfqpoint{1.258741in}{1.897722in}}%
\pgfpathcurveto{\pgfqpoint{1.258741in}{1.891898in}}{\pgfqpoint{1.261055in}{1.886312in}}{\pgfqpoint{1.265173in}{1.882194in}}%
\pgfpathcurveto{\pgfqpoint{1.269291in}{1.878076in}}{\pgfqpoint{1.274877in}{1.875762in}}{\pgfqpoint{1.280701in}{1.875762in}}%
\pgfpathlineto{\pgfqpoint{1.280701in}{1.875762in}}%
\pgfpathclose%
\pgfusepath{stroke,fill}%
\end{pgfscope}%
\begin{pgfscope}%
\pgfpathrectangle{\pgfqpoint{0.100000in}{0.209042in}}{\pgfqpoint{2.906250in}{2.906250in}}%
\pgfusepath{clip}%
\pgfsetbuttcap%
\pgfsetroundjoin%
\definecolor{currentfill}{rgb}{0.639216,0.189921,0.494150}%
\pgfsetfillcolor{currentfill}%
\pgfsetfillopacity{0.800000}%
\pgfsetlinewidth{1.003750pt}%
\definecolor{currentstroke}{rgb}{0.639216,0.189921,0.494150}%
\pgfsetstrokecolor{currentstroke}%
\pgfsetstrokeopacity{0.800000}%
\pgfsetdash{}{0pt}%
\pgfpathmoveto{\pgfqpoint{1.785979in}{2.094385in}}%
\pgfpathcurveto{\pgfqpoint{1.791802in}{2.094385in}}{\pgfqpoint{1.797389in}{2.096699in}}{\pgfqpoint{1.801507in}{2.100817in}}%
\pgfpathcurveto{\pgfqpoint{1.805625in}{2.104935in}}{\pgfqpoint{1.807939in}{2.110521in}}{\pgfqpoint{1.807939in}{2.116345in}}%
\pgfpathcurveto{\pgfqpoint{1.807939in}{2.122169in}}{\pgfqpoint{1.805625in}{2.127755in}}{\pgfqpoint{1.801507in}{2.131873in}}%
\pgfpathcurveto{\pgfqpoint{1.797389in}{2.135991in}}{\pgfqpoint{1.791802in}{2.138305in}}{\pgfqpoint{1.785979in}{2.138305in}}%
\pgfpathcurveto{\pgfqpoint{1.780155in}{2.138305in}}{\pgfqpoint{1.774568in}{2.135991in}}{\pgfqpoint{1.770450in}{2.131873in}}%
\pgfpathcurveto{\pgfqpoint{1.766332in}{2.127755in}}{\pgfqpoint{1.764018in}{2.122169in}}{\pgfqpoint{1.764018in}{2.116345in}}%
\pgfpathcurveto{\pgfqpoint{1.764018in}{2.110521in}}{\pgfqpoint{1.766332in}{2.104935in}}{\pgfqpoint{1.770450in}{2.100817in}}%
\pgfpathcurveto{\pgfqpoint{1.774568in}{2.096699in}}{\pgfqpoint{1.780155in}{2.094385in}}{\pgfqpoint{1.785979in}{2.094385in}}%
\pgfpathlineto{\pgfqpoint{1.785979in}{2.094385in}}%
\pgfpathclose%
\pgfusepath{stroke,fill}%
\end{pgfscope}%
\begin{pgfscope}%
\pgfpathrectangle{\pgfqpoint{0.100000in}{0.209042in}}{\pgfqpoint{2.906250in}{2.906250in}}%
\pgfusepath{clip}%
\pgfsetbuttcap%
\pgfsetroundjoin%
\definecolor{currentfill}{rgb}{0.899552,0.314616,0.391037}%
\pgfsetfillcolor{currentfill}%
\pgfsetfillopacity{0.800000}%
\pgfsetlinewidth{1.003750pt}%
\definecolor{currentstroke}{rgb}{0.899552,0.314616,0.391037}%
\pgfsetstrokecolor{currentstroke}%
\pgfsetstrokeopacity{0.800000}%
\pgfsetdash{}{0pt}%
\pgfpathmoveto{\pgfqpoint{1.824583in}{2.233989in}}%
\pgfpathcurveto{\pgfqpoint{1.830406in}{2.233989in}}{\pgfqpoint{1.835993in}{2.236303in}}{\pgfqpoint{1.840111in}{2.240421in}}%
\pgfpathcurveto{\pgfqpoint{1.844229in}{2.244539in}}{\pgfqpoint{1.846543in}{2.250125in}}{\pgfqpoint{1.846543in}{2.255949in}}%
\pgfpathcurveto{\pgfqpoint{1.846543in}{2.261773in}}{\pgfqpoint{1.844229in}{2.267359in}}{\pgfqpoint{1.840111in}{2.271477in}}%
\pgfpathcurveto{\pgfqpoint{1.835993in}{2.275595in}}{\pgfqpoint{1.830406in}{2.277909in}}{\pgfqpoint{1.824583in}{2.277909in}}%
\pgfpathcurveto{\pgfqpoint{1.818759in}{2.277909in}}{\pgfqpoint{1.813172in}{2.275595in}}{\pgfqpoint{1.809054in}{2.271477in}}%
\pgfpathcurveto{\pgfqpoint{1.804936in}{2.267359in}}{\pgfqpoint{1.802622in}{2.261773in}}{\pgfqpoint{1.802622in}{2.255949in}}%
\pgfpathcurveto{\pgfqpoint{1.802622in}{2.250125in}}{\pgfqpoint{1.804936in}{2.244539in}}{\pgfqpoint{1.809054in}{2.240421in}}%
\pgfpathcurveto{\pgfqpoint{1.813172in}{2.236303in}}{\pgfqpoint{1.818759in}{2.233989in}}{\pgfqpoint{1.824583in}{2.233989in}}%
\pgfpathlineto{\pgfqpoint{1.824583in}{2.233989in}}%
\pgfpathclose%
\pgfusepath{stroke,fill}%
\end{pgfscope}%
\begin{pgfscope}%
\pgfpathrectangle{\pgfqpoint{0.100000in}{0.209042in}}{\pgfqpoint{2.906250in}{2.906250in}}%
\pgfusepath{clip}%
\pgfsetbuttcap%
\pgfsetroundjoin%
\definecolor{currentfill}{rgb}{0.550287,0.161158,0.505719}%
\pgfsetfillcolor{currentfill}%
\pgfsetfillopacity{0.800000}%
\pgfsetlinewidth{1.003750pt}%
\definecolor{currentstroke}{rgb}{0.550287,0.161158,0.505719}%
\pgfsetstrokecolor{currentstroke}%
\pgfsetstrokeopacity{0.800000}%
\pgfsetdash{}{0pt}%
\pgfpathmoveto{\pgfqpoint{2.012375in}{1.556180in}}%
\pgfpathcurveto{\pgfqpoint{2.018199in}{1.556180in}}{\pgfqpoint{2.023785in}{1.558494in}}{\pgfqpoint{2.027903in}{1.562612in}}%
\pgfpathcurveto{\pgfqpoint{2.032021in}{1.566730in}}{\pgfqpoint{2.034335in}{1.572317in}}{\pgfqpoint{2.034335in}{1.578141in}}%
\pgfpathcurveto{\pgfqpoint{2.034335in}{1.583964in}}{\pgfqpoint{2.032021in}{1.589551in}}{\pgfqpoint{2.027903in}{1.593669in}}%
\pgfpathcurveto{\pgfqpoint{2.023785in}{1.597787in}}{\pgfqpoint{2.018199in}{1.600101in}}{\pgfqpoint{2.012375in}{1.600101in}}%
\pgfpathcurveto{\pgfqpoint{2.006551in}{1.600101in}}{\pgfqpoint{2.000965in}{1.597787in}}{\pgfqpoint{1.996847in}{1.593669in}}%
\pgfpathcurveto{\pgfqpoint{1.992728in}{1.589551in}}{\pgfqpoint{1.990415in}{1.583964in}}{\pgfqpoint{1.990415in}{1.578141in}}%
\pgfpathcurveto{\pgfqpoint{1.990415in}{1.572317in}}{\pgfqpoint{1.992728in}{1.566730in}}{\pgfqpoint{1.996847in}{1.562612in}}%
\pgfpathcurveto{\pgfqpoint{2.000965in}{1.558494in}}{\pgfqpoint{2.006551in}{1.556180in}}{\pgfqpoint{2.012375in}{1.556180in}}%
\pgfpathlineto{\pgfqpoint{2.012375in}{1.556180in}}%
\pgfpathclose%
\pgfusepath{stroke,fill}%
\end{pgfscope}%
\begin{pgfscope}%
\pgfpathrectangle{\pgfqpoint{0.100000in}{0.209042in}}{\pgfqpoint{2.906250in}{2.906250in}}%
\pgfusepath{clip}%
\pgfsetbuttcap%
\pgfsetroundjoin%
\definecolor{currentfill}{rgb}{0.420791,0.112920,0.505215}%
\pgfsetfillcolor{currentfill}%
\pgfsetfillopacity{0.800000}%
\pgfsetlinewidth{1.003750pt}%
\definecolor{currentstroke}{rgb}{0.420791,0.112920,0.505215}%
\pgfsetstrokecolor{currentstroke}%
\pgfsetstrokeopacity{0.800000}%
\pgfsetdash{}{0pt}%
\pgfpathmoveto{\pgfqpoint{1.296719in}{1.867014in}}%
\pgfpathcurveto{\pgfqpoint{1.302543in}{1.867014in}}{\pgfqpoint{1.308129in}{1.869328in}}{\pgfqpoint{1.312247in}{1.873446in}}%
\pgfpathcurveto{\pgfqpoint{1.316365in}{1.877564in}}{\pgfqpoint{1.318679in}{1.883150in}}{\pgfqpoint{1.318679in}{1.888974in}}%
\pgfpathcurveto{\pgfqpoint{1.318679in}{1.894798in}}{\pgfqpoint{1.316365in}{1.900384in}}{\pgfqpoint{1.312247in}{1.904502in}}%
\pgfpathcurveto{\pgfqpoint{1.308129in}{1.908621in}}{\pgfqpoint{1.302543in}{1.910935in}}{\pgfqpoint{1.296719in}{1.910935in}}%
\pgfpathcurveto{\pgfqpoint{1.290895in}{1.910935in}}{\pgfqpoint{1.285309in}{1.908621in}}{\pgfqpoint{1.281191in}{1.904502in}}%
\pgfpathcurveto{\pgfqpoint{1.277072in}{1.900384in}}{\pgfqpoint{1.274759in}{1.894798in}}{\pgfqpoint{1.274759in}{1.888974in}}%
\pgfpathcurveto{\pgfqpoint{1.274759in}{1.883150in}}{\pgfqpoint{1.277072in}{1.877564in}}{\pgfqpoint{1.281191in}{1.873446in}}%
\pgfpathcurveto{\pgfqpoint{1.285309in}{1.869328in}}{\pgfqpoint{1.290895in}{1.867014in}}{\pgfqpoint{1.296719in}{1.867014in}}%
\pgfpathlineto{\pgfqpoint{1.296719in}{1.867014in}}%
\pgfpathclose%
\pgfusepath{stroke,fill}%
\end{pgfscope}%
\begin{pgfscope}%
\pgfpathrectangle{\pgfqpoint{0.100000in}{0.209042in}}{\pgfqpoint{2.906250in}{2.906250in}}%
\pgfusepath{clip}%
\pgfsetbuttcap%
\pgfsetroundjoin%
\definecolor{currentfill}{rgb}{0.146785,0.068738,0.334011}%
\pgfsetfillcolor{currentfill}%
\pgfsetfillopacity{0.800000}%
\pgfsetlinewidth{1.003750pt}%
\definecolor{currentstroke}{rgb}{0.146785,0.068738,0.334011}%
\pgfsetstrokecolor{currentstroke}%
\pgfsetstrokeopacity{0.800000}%
\pgfsetdash{}{0pt}%
\pgfpathmoveto{\pgfqpoint{1.487469in}{1.803211in}}%
\pgfpathcurveto{\pgfqpoint{1.493293in}{1.803211in}}{\pgfqpoint{1.498879in}{1.805525in}}{\pgfqpoint{1.502997in}{1.809643in}}%
\pgfpathcurveto{\pgfqpoint{1.507115in}{1.813761in}}{\pgfqpoint{1.509429in}{1.819347in}}{\pgfqpoint{1.509429in}{1.825171in}}%
\pgfpathcurveto{\pgfqpoint{1.509429in}{1.830995in}}{\pgfqpoint{1.507115in}{1.836581in}}{\pgfqpoint{1.502997in}{1.840699in}}%
\pgfpathcurveto{\pgfqpoint{1.498879in}{1.844817in}}{\pgfqpoint{1.493293in}{1.847131in}}{\pgfqpoint{1.487469in}{1.847131in}}%
\pgfpathcurveto{\pgfqpoint{1.481645in}{1.847131in}}{\pgfqpoint{1.476059in}{1.844817in}}{\pgfqpoint{1.471940in}{1.840699in}}%
\pgfpathcurveto{\pgfqpoint{1.467822in}{1.836581in}}{\pgfqpoint{1.465508in}{1.830995in}}{\pgfqpoint{1.465508in}{1.825171in}}%
\pgfpathcurveto{\pgfqpoint{1.465508in}{1.819347in}}{\pgfqpoint{1.467822in}{1.813761in}}{\pgfqpoint{1.471940in}{1.809643in}}%
\pgfpathcurveto{\pgfqpoint{1.476059in}{1.805525in}}{\pgfqpoint{1.481645in}{1.803211in}}{\pgfqpoint{1.487469in}{1.803211in}}%
\pgfpathlineto{\pgfqpoint{1.487469in}{1.803211in}}%
\pgfpathclose%
\pgfusepath{stroke,fill}%
\end{pgfscope}%
\begin{pgfscope}%
\pgfpathrectangle{\pgfqpoint{0.100000in}{0.209042in}}{\pgfqpoint{2.906250in}{2.906250in}}%
\pgfusepath{clip}%
\pgfsetbuttcap%
\pgfsetroundjoin%
\definecolor{currentfill}{rgb}{0.316654,0.071690,0.485380}%
\pgfsetfillcolor{currentfill}%
\pgfsetfillopacity{0.800000}%
\pgfsetlinewidth{1.003750pt}%
\definecolor{currentstroke}{rgb}{0.316654,0.071690,0.485380}%
\pgfsetstrokecolor{currentstroke}%
\pgfsetstrokeopacity{0.800000}%
\pgfsetdash{}{0pt}%
\pgfpathmoveto{\pgfqpoint{1.746360in}{1.913847in}}%
\pgfpathcurveto{\pgfqpoint{1.752184in}{1.913847in}}{\pgfqpoint{1.757770in}{1.916161in}}{\pgfqpoint{1.761889in}{1.920279in}}%
\pgfpathcurveto{\pgfqpoint{1.766007in}{1.924397in}}{\pgfqpoint{1.768321in}{1.929983in}}{\pgfqpoint{1.768321in}{1.935807in}}%
\pgfpathcurveto{\pgfqpoint{1.768321in}{1.941631in}}{\pgfqpoint{1.766007in}{1.947217in}}{\pgfqpoint{1.761889in}{1.951335in}}%
\pgfpathcurveto{\pgfqpoint{1.757770in}{1.955453in}}{\pgfqpoint{1.752184in}{1.957767in}}{\pgfqpoint{1.746360in}{1.957767in}}%
\pgfpathcurveto{\pgfqpoint{1.740536in}{1.957767in}}{\pgfqpoint{1.734950in}{1.955453in}}{\pgfqpoint{1.730832in}{1.951335in}}%
\pgfpathcurveto{\pgfqpoint{1.726714in}{1.947217in}}{\pgfqpoint{1.724400in}{1.941631in}}{\pgfqpoint{1.724400in}{1.935807in}}%
\pgfpathcurveto{\pgfqpoint{1.724400in}{1.929983in}}{\pgfqpoint{1.726714in}{1.924397in}}{\pgfqpoint{1.730832in}{1.920279in}}%
\pgfpathcurveto{\pgfqpoint{1.734950in}{1.916161in}}{\pgfqpoint{1.740536in}{1.913847in}}{\pgfqpoint{1.746360in}{1.913847in}}%
\pgfpathlineto{\pgfqpoint{1.746360in}{1.913847in}}%
\pgfpathclose%
\pgfusepath{stroke,fill}%
\end{pgfscope}%
\begin{pgfscope}%
\pgfpathrectangle{\pgfqpoint{0.100000in}{0.209042in}}{\pgfqpoint{2.906250in}{2.906250in}}%
\pgfusepath{clip}%
\pgfsetbuttcap%
\pgfsetroundjoin%
\definecolor{currentfill}{rgb}{0.384299,0.097855,0.501002}%
\pgfsetfillcolor{currentfill}%
\pgfsetfillopacity{0.800000}%
\pgfsetlinewidth{1.003750pt}%
\definecolor{currentstroke}{rgb}{0.384299,0.097855,0.501002}%
\pgfsetstrokecolor{currentstroke}%
\pgfsetstrokeopacity{0.800000}%
\pgfsetdash{}{0pt}%
\pgfpathmoveto{\pgfqpoint{1.506058in}{1.981314in}}%
\pgfpathcurveto{\pgfqpoint{1.511881in}{1.981314in}}{\pgfqpoint{1.517468in}{1.983628in}}{\pgfqpoint{1.521586in}{1.987746in}}%
\pgfpathcurveto{\pgfqpoint{1.525704in}{1.991865in}}{\pgfqpoint{1.528018in}{1.997451in}}{\pgfqpoint{1.528018in}{2.003275in}}%
\pgfpathcurveto{\pgfqpoint{1.528018in}{2.009099in}}{\pgfqpoint{1.525704in}{2.014685in}}{\pgfqpoint{1.521586in}{2.018803in}}%
\pgfpathcurveto{\pgfqpoint{1.517468in}{2.022921in}}{\pgfqpoint{1.511881in}{2.025235in}}{\pgfqpoint{1.506058in}{2.025235in}}%
\pgfpathcurveto{\pgfqpoint{1.500234in}{2.025235in}}{\pgfqpoint{1.494647in}{2.022921in}}{\pgfqpoint{1.490529in}{2.018803in}}%
\pgfpathcurveto{\pgfqpoint{1.486411in}{2.014685in}}{\pgfqpoint{1.484097in}{2.009099in}}{\pgfqpoint{1.484097in}{2.003275in}}%
\pgfpathcurveto{\pgfqpoint{1.484097in}{1.997451in}}{\pgfqpoint{1.486411in}{1.991865in}}{\pgfqpoint{1.490529in}{1.987746in}}%
\pgfpathcurveto{\pgfqpoint{1.494647in}{1.983628in}}{\pgfqpoint{1.500234in}{1.981314in}}{\pgfqpoint{1.506058in}{1.981314in}}%
\pgfpathlineto{\pgfqpoint{1.506058in}{1.981314in}}%
\pgfpathclose%
\pgfusepath{stroke,fill}%
\end{pgfscope}%
\begin{pgfscope}%
\pgfpathrectangle{\pgfqpoint{0.100000in}{0.209042in}}{\pgfqpoint{2.906250in}{2.906250in}}%
\pgfusepath{clip}%
\pgfsetbuttcap%
\pgfsetroundjoin%
\definecolor{currentfill}{rgb}{0.537755,0.156894,0.506551}%
\pgfsetfillcolor{currentfill}%
\pgfsetfillopacity{0.800000}%
\pgfsetlinewidth{1.003750pt}%
\definecolor{currentstroke}{rgb}{0.537755,0.156894,0.506551}%
\pgfsetstrokecolor{currentstroke}%
\pgfsetstrokeopacity{0.800000}%
\pgfsetdash{}{0pt}%
\pgfpathmoveto{\pgfqpoint{1.690082in}{1.373377in}}%
\pgfpathcurveto{\pgfqpoint{1.695905in}{1.373377in}}{\pgfqpoint{1.701492in}{1.375691in}}{\pgfqpoint{1.705610in}{1.379809in}}%
\pgfpathcurveto{\pgfqpoint{1.709728in}{1.383927in}}{\pgfqpoint{1.712042in}{1.389514in}}{\pgfqpoint{1.712042in}{1.395337in}}%
\pgfpathcurveto{\pgfqpoint{1.712042in}{1.401161in}}{\pgfqpoint{1.709728in}{1.406748in}}{\pgfqpoint{1.705610in}{1.410866in}}%
\pgfpathcurveto{\pgfqpoint{1.701492in}{1.414984in}}{\pgfqpoint{1.695905in}{1.417298in}}{\pgfqpoint{1.690082in}{1.417298in}}%
\pgfpathcurveto{\pgfqpoint{1.684258in}{1.417298in}}{\pgfqpoint{1.678671in}{1.414984in}}{\pgfqpoint{1.674553in}{1.410866in}}%
\pgfpathcurveto{\pgfqpoint{1.670435in}{1.406748in}}{\pgfqpoint{1.668121in}{1.401161in}}{\pgfqpoint{1.668121in}{1.395337in}}%
\pgfpathcurveto{\pgfqpoint{1.668121in}{1.389514in}}{\pgfqpoint{1.670435in}{1.383927in}}{\pgfqpoint{1.674553in}{1.379809in}}%
\pgfpathcurveto{\pgfqpoint{1.678671in}{1.375691in}}{\pgfqpoint{1.684258in}{1.373377in}}{\pgfqpoint{1.690082in}{1.373377in}}%
\pgfpathlineto{\pgfqpoint{1.690082in}{1.373377in}}%
\pgfpathclose%
\pgfusepath{stroke,fill}%
\end{pgfscope}%
\begin{pgfscope}%
\pgfpathrectangle{\pgfqpoint{0.100000in}{0.209042in}}{\pgfqpoint{2.906250in}{2.906250in}}%
\pgfusepath{clip}%
\pgfsetbuttcap%
\pgfsetroundjoin%
\definecolor{currentfill}{rgb}{0.113094,0.065492,0.276784}%
\pgfsetfillcolor{currentfill}%
\pgfsetfillopacity{0.800000}%
\pgfsetlinewidth{1.003750pt}%
\definecolor{currentstroke}{rgb}{0.113094,0.065492,0.276784}%
\pgfsetstrokecolor{currentstroke}%
\pgfsetstrokeopacity{0.800000}%
\pgfsetdash{}{0pt}%
\pgfpathmoveto{\pgfqpoint{1.699185in}{1.647386in}}%
\pgfpathcurveto{\pgfqpoint{1.705009in}{1.647386in}}{\pgfqpoint{1.710595in}{1.649700in}}{\pgfqpoint{1.714713in}{1.653818in}}%
\pgfpathcurveto{\pgfqpoint{1.718831in}{1.657936in}}{\pgfqpoint{1.721145in}{1.663523in}}{\pgfqpoint{1.721145in}{1.669346in}}%
\pgfpathcurveto{\pgfqpoint{1.721145in}{1.675170in}}{\pgfqpoint{1.718831in}{1.680757in}}{\pgfqpoint{1.714713in}{1.684875in}}%
\pgfpathcurveto{\pgfqpoint{1.710595in}{1.688993in}}{\pgfqpoint{1.705009in}{1.691307in}}{\pgfqpoint{1.699185in}{1.691307in}}%
\pgfpathcurveto{\pgfqpoint{1.693361in}{1.691307in}}{\pgfqpoint{1.687774in}{1.688993in}}{\pgfqpoint{1.683656in}{1.684875in}}%
\pgfpathcurveto{\pgfqpoint{1.679538in}{1.680757in}}{\pgfqpoint{1.677224in}{1.675170in}}{\pgfqpoint{1.677224in}{1.669346in}}%
\pgfpathcurveto{\pgfqpoint{1.677224in}{1.663523in}}{\pgfqpoint{1.679538in}{1.657936in}}{\pgfqpoint{1.683656in}{1.653818in}}%
\pgfpathcurveto{\pgfqpoint{1.687774in}{1.649700in}}{\pgfqpoint{1.693361in}{1.647386in}}{\pgfqpoint{1.699185in}{1.647386in}}%
\pgfpathlineto{\pgfqpoint{1.699185in}{1.647386in}}%
\pgfpathclose%
\pgfusepath{stroke,fill}%
\end{pgfscope}%
\begin{pgfscope}%
\pgfpathrectangle{\pgfqpoint{0.100000in}{0.209042in}}{\pgfqpoint{2.906250in}{2.906250in}}%
\pgfusepath{clip}%
\pgfsetbuttcap%
\pgfsetroundjoin%
\definecolor{currentfill}{rgb}{0.494258,0.141462,0.507988}%
\pgfsetfillcolor{currentfill}%
\pgfsetfillopacity{0.800000}%
\pgfsetlinewidth{1.003750pt}%
\definecolor{currentstroke}{rgb}{0.494258,0.141462,0.507988}%
\pgfsetstrokecolor{currentstroke}%
\pgfsetstrokeopacity{0.800000}%
\pgfsetdash{}{0pt}%
\pgfpathmoveto{\pgfqpoint{1.714474in}{1.400014in}}%
\pgfpathcurveto{\pgfqpoint{1.720298in}{1.400014in}}{\pgfqpoint{1.725884in}{1.402328in}}{\pgfqpoint{1.730002in}{1.406446in}}%
\pgfpathcurveto{\pgfqpoint{1.734121in}{1.410564in}}{\pgfqpoint{1.736435in}{1.416150in}}{\pgfqpoint{1.736435in}{1.421974in}}%
\pgfpathcurveto{\pgfqpoint{1.736435in}{1.427798in}}{\pgfqpoint{1.734121in}{1.433384in}}{\pgfqpoint{1.730002in}{1.437502in}}%
\pgfpathcurveto{\pgfqpoint{1.725884in}{1.441621in}}{\pgfqpoint{1.720298in}{1.443934in}}{\pgfqpoint{1.714474in}{1.443934in}}%
\pgfpathcurveto{\pgfqpoint{1.708650in}{1.443934in}}{\pgfqpoint{1.703064in}{1.441621in}}{\pgfqpoint{1.698946in}{1.437502in}}%
\pgfpathcurveto{\pgfqpoint{1.694828in}{1.433384in}}{\pgfqpoint{1.692514in}{1.427798in}}{\pgfqpoint{1.692514in}{1.421974in}}%
\pgfpathcurveto{\pgfqpoint{1.692514in}{1.416150in}}{\pgfqpoint{1.694828in}{1.410564in}}{\pgfqpoint{1.698946in}{1.406446in}}%
\pgfpathcurveto{\pgfqpoint{1.703064in}{1.402328in}}{\pgfqpoint{1.708650in}{1.400014in}}{\pgfqpoint{1.714474in}{1.400014in}}%
\pgfpathlineto{\pgfqpoint{1.714474in}{1.400014in}}%
\pgfpathclose%
\pgfusepath{stroke,fill}%
\end{pgfscope}%
\begin{pgfscope}%
\pgfpathrectangle{\pgfqpoint{0.100000in}{0.209042in}}{\pgfqpoint{2.906250in}{2.906250in}}%
\pgfusepath{clip}%
\pgfsetbuttcap%
\pgfsetroundjoin%
\definecolor{currentfill}{rgb}{0.140858,0.068654,0.324538}%
\pgfsetfillcolor{currentfill}%
\pgfsetfillopacity{0.800000}%
\pgfsetlinewidth{1.003750pt}%
\definecolor{currentstroke}{rgb}{0.140858,0.068654,0.324538}%
\pgfsetstrokecolor{currentstroke}%
\pgfsetstrokeopacity{0.800000}%
\pgfsetdash{}{0pt}%
\pgfpathmoveto{\pgfqpoint{1.730623in}{1.636837in}}%
\pgfpathcurveto{\pgfqpoint{1.736447in}{1.636837in}}{\pgfqpoint{1.742033in}{1.639151in}}{\pgfqpoint{1.746151in}{1.643269in}}%
\pgfpathcurveto{\pgfqpoint{1.750270in}{1.647387in}}{\pgfqpoint{1.752583in}{1.652974in}}{\pgfqpoint{1.752583in}{1.658798in}}%
\pgfpathcurveto{\pgfqpoint{1.752583in}{1.664621in}}{\pgfqpoint{1.750270in}{1.670208in}}{\pgfqpoint{1.746151in}{1.674326in}}%
\pgfpathcurveto{\pgfqpoint{1.742033in}{1.678444in}}{\pgfqpoint{1.736447in}{1.680758in}}{\pgfqpoint{1.730623in}{1.680758in}}%
\pgfpathcurveto{\pgfqpoint{1.724799in}{1.680758in}}{\pgfqpoint{1.719213in}{1.678444in}}{\pgfqpoint{1.715095in}{1.674326in}}%
\pgfpathcurveto{\pgfqpoint{1.710977in}{1.670208in}}{\pgfqpoint{1.708663in}{1.664621in}}{\pgfqpoint{1.708663in}{1.658798in}}%
\pgfpathcurveto{\pgfqpoint{1.708663in}{1.652974in}}{\pgfqpoint{1.710977in}{1.647387in}}{\pgfqpoint{1.715095in}{1.643269in}}%
\pgfpathcurveto{\pgfqpoint{1.719213in}{1.639151in}}{\pgfqpoint{1.724799in}{1.636837in}}{\pgfqpoint{1.730623in}{1.636837in}}%
\pgfpathlineto{\pgfqpoint{1.730623in}{1.636837in}}%
\pgfpathclose%
\pgfusepath{stroke,fill}%
\end{pgfscope}%
\begin{pgfscope}%
\pgfpathrectangle{\pgfqpoint{0.100000in}{0.209042in}}{\pgfqpoint{2.906250in}{2.906250in}}%
\pgfusepath{clip}%
\pgfsetbuttcap%
\pgfsetroundjoin%
\definecolor{currentfill}{rgb}{0.512831,0.148179,0.507648}%
\pgfsetfillcolor{currentfill}%
\pgfsetfillopacity{0.800000}%
\pgfsetlinewidth{1.003750pt}%
\definecolor{currentstroke}{rgb}{0.512831,0.148179,0.507648}%
\pgfsetstrokecolor{currentstroke}%
\pgfsetstrokeopacity{0.800000}%
\pgfsetdash{}{0pt}%
\pgfpathmoveto{\pgfqpoint{1.193753in}{1.749605in}}%
\pgfpathcurveto{\pgfqpoint{1.199577in}{1.749605in}}{\pgfqpoint{1.205163in}{1.751919in}}{\pgfqpoint{1.209281in}{1.756037in}}%
\pgfpathcurveto{\pgfqpoint{1.213399in}{1.760155in}}{\pgfqpoint{1.215713in}{1.765741in}}{\pgfqpoint{1.215713in}{1.771565in}}%
\pgfpathcurveto{\pgfqpoint{1.215713in}{1.777389in}}{\pgfqpoint{1.213399in}{1.782975in}}{\pgfqpoint{1.209281in}{1.787093in}}%
\pgfpathcurveto{\pgfqpoint{1.205163in}{1.791212in}}{\pgfqpoint{1.199577in}{1.793525in}}{\pgfqpoint{1.193753in}{1.793525in}}%
\pgfpathcurveto{\pgfqpoint{1.187929in}{1.793525in}}{\pgfqpoint{1.182343in}{1.791212in}}{\pgfqpoint{1.178225in}{1.787093in}}%
\pgfpathcurveto{\pgfqpoint{1.174107in}{1.782975in}}{\pgfqpoint{1.171793in}{1.777389in}}{\pgfqpoint{1.171793in}{1.771565in}}%
\pgfpathcurveto{\pgfqpoint{1.171793in}{1.765741in}}{\pgfqpoint{1.174107in}{1.760155in}}{\pgfqpoint{1.178225in}{1.756037in}}%
\pgfpathcurveto{\pgfqpoint{1.182343in}{1.751919in}}{\pgfqpoint{1.187929in}{1.749605in}}{\pgfqpoint{1.193753in}{1.749605in}}%
\pgfpathlineto{\pgfqpoint{1.193753in}{1.749605in}}%
\pgfpathclose%
\pgfusepath{stroke,fill}%
\end{pgfscope}%
\begin{pgfscope}%
\pgfpathrectangle{\pgfqpoint{0.100000in}{0.209042in}}{\pgfqpoint{2.906250in}{2.906250in}}%
\pgfusepath{clip}%
\pgfsetbuttcap%
\pgfsetroundjoin%
\definecolor{currentfill}{rgb}{0.562866,0.165368,0.504692}%
\pgfsetfillcolor{currentfill}%
\pgfsetfillopacity{0.800000}%
\pgfsetlinewidth{1.003750pt}%
\definecolor{currentstroke}{rgb}{0.562866,0.165368,0.504692}%
\pgfsetstrokecolor{currentstroke}%
\pgfsetstrokeopacity{0.800000}%
\pgfsetdash{}{0pt}%
\pgfpathmoveto{\pgfqpoint{1.884142in}{2.003774in}}%
\pgfpathcurveto{\pgfqpoint{1.889966in}{2.003774in}}{\pgfqpoint{1.895552in}{2.006088in}}{\pgfqpoint{1.899670in}{2.010206in}}%
\pgfpathcurveto{\pgfqpoint{1.903788in}{2.014324in}}{\pgfqpoint{1.906102in}{2.019910in}}{\pgfqpoint{1.906102in}{2.025734in}}%
\pgfpathcurveto{\pgfqpoint{1.906102in}{2.031558in}}{\pgfqpoint{1.903788in}{2.037144in}}{\pgfqpoint{1.899670in}{2.041262in}}%
\pgfpathcurveto{\pgfqpoint{1.895552in}{2.045380in}}{\pgfqpoint{1.889966in}{2.047694in}}{\pgfqpoint{1.884142in}{2.047694in}}%
\pgfpathcurveto{\pgfqpoint{1.878318in}{2.047694in}}{\pgfqpoint{1.872732in}{2.045380in}}{\pgfqpoint{1.868614in}{2.041262in}}%
\pgfpathcurveto{\pgfqpoint{1.864496in}{2.037144in}}{\pgfqpoint{1.862182in}{2.031558in}}{\pgfqpoint{1.862182in}{2.025734in}}%
\pgfpathcurveto{\pgfqpoint{1.862182in}{2.019910in}}{\pgfqpoint{1.864496in}{2.014324in}}{\pgfqpoint{1.868614in}{2.010206in}}%
\pgfpathcurveto{\pgfqpoint{1.872732in}{2.006088in}}{\pgfqpoint{1.878318in}{2.003774in}}{\pgfqpoint{1.884142in}{2.003774in}}%
\pgfpathlineto{\pgfqpoint{1.884142in}{2.003774in}}%
\pgfpathclose%
\pgfusepath{stroke,fill}%
\end{pgfscope}%
\begin{pgfscope}%
\pgfpathrectangle{\pgfqpoint{0.100000in}{0.209042in}}{\pgfqpoint{2.906250in}{2.906250in}}%
\pgfusepath{clip}%
\pgfsetbuttcap%
\pgfsetroundjoin%
\definecolor{currentfill}{rgb}{0.171713,0.067305,0.370771}%
\pgfsetfillcolor{currentfill}%
\pgfsetfillopacity{0.800000}%
\pgfsetlinewidth{1.003750pt}%
\definecolor{currentstroke}{rgb}{0.171713,0.067305,0.370771}%
\pgfsetstrokecolor{currentstroke}%
\pgfsetstrokeopacity{0.800000}%
\pgfsetdash{}{0pt}%
\pgfpathmoveto{\pgfqpoint{1.509869in}{1.606341in}}%
\pgfpathcurveto{\pgfqpoint{1.515693in}{1.606341in}}{\pgfqpoint{1.521280in}{1.608655in}}{\pgfqpoint{1.525398in}{1.612773in}}%
\pgfpathcurveto{\pgfqpoint{1.529516in}{1.616892in}}{\pgfqpoint{1.531830in}{1.622478in}}{\pgfqpoint{1.531830in}{1.628302in}}%
\pgfpathcurveto{\pgfqpoint{1.531830in}{1.634126in}}{\pgfqpoint{1.529516in}{1.639712in}}{\pgfqpoint{1.525398in}{1.643830in}}%
\pgfpathcurveto{\pgfqpoint{1.521280in}{1.647948in}}{\pgfqpoint{1.515693in}{1.650262in}}{\pgfqpoint{1.509869in}{1.650262in}}%
\pgfpathcurveto{\pgfqpoint{1.504045in}{1.650262in}}{\pgfqpoint{1.498459in}{1.647948in}}{\pgfqpoint{1.494341in}{1.643830in}}%
\pgfpathcurveto{\pgfqpoint{1.490223in}{1.639712in}}{\pgfqpoint{1.487909in}{1.634126in}}{\pgfqpoint{1.487909in}{1.628302in}}%
\pgfpathcurveto{\pgfqpoint{1.487909in}{1.622478in}}{\pgfqpoint{1.490223in}{1.616892in}}{\pgfqpoint{1.494341in}{1.612773in}}%
\pgfpathcurveto{\pgfqpoint{1.498459in}{1.608655in}}{\pgfqpoint{1.504045in}{1.606341in}}{\pgfqpoint{1.509869in}{1.606341in}}%
\pgfpathlineto{\pgfqpoint{1.509869in}{1.606341in}}%
\pgfpathclose%
\pgfusepath{stroke,fill}%
\end{pgfscope}%
\begin{pgfscope}%
\pgfpathrectangle{\pgfqpoint{0.100000in}{0.209042in}}{\pgfqpoint{2.906250in}{2.906250in}}%
\pgfusepath{clip}%
\pgfsetbuttcap%
\pgfsetroundjoin%
\definecolor{currentfill}{rgb}{0.390384,0.100379,0.501864}%
\pgfsetfillcolor{currentfill}%
\pgfsetfillopacity{0.800000}%
\pgfsetlinewidth{1.003750pt}%
\definecolor{currentstroke}{rgb}{0.390384,0.100379,0.501864}%
\pgfsetstrokecolor{currentstroke}%
\pgfsetstrokeopacity{0.800000}%
\pgfsetdash{}{0pt}%
\pgfpathmoveto{\pgfqpoint{1.456234in}{1.487335in}}%
\pgfpathcurveto{\pgfqpoint{1.462058in}{1.487335in}}{\pgfqpoint{1.467644in}{1.489649in}}{\pgfqpoint{1.471762in}{1.493767in}}%
\pgfpathcurveto{\pgfqpoint{1.475881in}{1.497885in}}{\pgfqpoint{1.478194in}{1.503471in}}{\pgfqpoint{1.478194in}{1.509295in}}%
\pgfpathcurveto{\pgfqpoint{1.478194in}{1.515119in}}{\pgfqpoint{1.475881in}{1.520705in}}{\pgfqpoint{1.471762in}{1.524824in}}%
\pgfpathcurveto{\pgfqpoint{1.467644in}{1.528942in}}{\pgfqpoint{1.462058in}{1.531256in}}{\pgfqpoint{1.456234in}{1.531256in}}%
\pgfpathcurveto{\pgfqpoint{1.450410in}{1.531256in}}{\pgfqpoint{1.444824in}{1.528942in}}{\pgfqpoint{1.440706in}{1.524824in}}%
\pgfpathcurveto{\pgfqpoint{1.436588in}{1.520705in}}{\pgfqpoint{1.434274in}{1.515119in}}{\pgfqpoint{1.434274in}{1.509295in}}%
\pgfpathcurveto{\pgfqpoint{1.434274in}{1.503471in}}{\pgfqpoint{1.436588in}{1.497885in}}{\pgfqpoint{1.440706in}{1.493767in}}%
\pgfpathcurveto{\pgfqpoint{1.444824in}{1.489649in}}{\pgfqpoint{1.450410in}{1.487335in}}{\pgfqpoint{1.456234in}{1.487335in}}%
\pgfpathlineto{\pgfqpoint{1.456234in}{1.487335in}}%
\pgfpathclose%
\pgfusepath{stroke,fill}%
\end{pgfscope}%
\begin{pgfscope}%
\pgfpathrectangle{\pgfqpoint{0.100000in}{0.209042in}}{\pgfqpoint{2.906250in}{2.906250in}}%
\pgfusepath{clip}%
\pgfsetbuttcap%
\pgfsetroundjoin%
\definecolor{currentfill}{rgb}{0.074257,0.052017,0.202660}%
\pgfsetfillcolor{currentfill}%
\pgfsetfillopacity{0.800000}%
\pgfsetlinewidth{1.003750pt}%
\definecolor{currentstroke}{rgb}{0.074257,0.052017,0.202660}%
\pgfsetstrokecolor{currentstroke}%
\pgfsetstrokeopacity{0.800000}%
\pgfsetdash{}{0pt}%
\pgfpathmoveto{\pgfqpoint{1.586954in}{1.791194in}}%
\pgfpathcurveto{\pgfqpoint{1.592778in}{1.791194in}}{\pgfqpoint{1.598364in}{1.793507in}}{\pgfqpoint{1.602482in}{1.797626in}}%
\pgfpathcurveto{\pgfqpoint{1.606600in}{1.801744in}}{\pgfqpoint{1.608914in}{1.807330in}}{\pgfqpoint{1.608914in}{1.813154in}}%
\pgfpathcurveto{\pgfqpoint{1.608914in}{1.818978in}}{\pgfqpoint{1.606600in}{1.824564in}}{\pgfqpoint{1.602482in}{1.828682in}}%
\pgfpathcurveto{\pgfqpoint{1.598364in}{1.832800in}}{\pgfqpoint{1.592778in}{1.835114in}}{\pgfqpoint{1.586954in}{1.835114in}}%
\pgfpathcurveto{\pgfqpoint{1.581130in}{1.835114in}}{\pgfqpoint{1.575544in}{1.832800in}}{\pgfqpoint{1.571426in}{1.828682in}}%
\pgfpathcurveto{\pgfqpoint{1.567308in}{1.824564in}}{\pgfqpoint{1.564994in}{1.818978in}}{\pgfqpoint{1.564994in}{1.813154in}}%
\pgfpathcurveto{\pgfqpoint{1.564994in}{1.807330in}}{\pgfqpoint{1.567308in}{1.801744in}}{\pgfqpoint{1.571426in}{1.797626in}}%
\pgfpathcurveto{\pgfqpoint{1.575544in}{1.793507in}}{\pgfqpoint{1.581130in}{1.791194in}}{\pgfqpoint{1.586954in}{1.791194in}}%
\pgfpathlineto{\pgfqpoint{1.586954in}{1.791194in}}%
\pgfpathclose%
\pgfusepath{stroke,fill}%
\end{pgfscope}%
\begin{pgfscope}%
\pgfpathrectangle{\pgfqpoint{0.100000in}{0.209042in}}{\pgfqpoint{2.906250in}{2.906250in}}%
\pgfusepath{clip}%
\pgfsetbuttcap%
\pgfsetroundjoin%
\definecolor{currentfill}{rgb}{0.232077,0.059889,0.437695}%
\pgfsetfillcolor{currentfill}%
\pgfsetfillopacity{0.800000}%
\pgfsetlinewidth{1.003750pt}%
\definecolor{currentstroke}{rgb}{0.232077,0.059889,0.437695}%
\pgfsetstrokecolor{currentstroke}%
\pgfsetstrokeopacity{0.800000}%
\pgfsetdash{}{0pt}%
\pgfpathmoveto{\pgfqpoint{1.836506in}{1.730365in}}%
\pgfpathcurveto{\pgfqpoint{1.842330in}{1.730365in}}{\pgfqpoint{1.847917in}{1.732679in}}{\pgfqpoint{1.852035in}{1.736797in}}%
\pgfpathcurveto{\pgfqpoint{1.856153in}{1.740915in}}{\pgfqpoint{1.858467in}{1.746501in}}{\pgfqpoint{1.858467in}{1.752325in}}%
\pgfpathcurveto{\pgfqpoint{1.858467in}{1.758149in}}{\pgfqpoint{1.856153in}{1.763735in}}{\pgfqpoint{1.852035in}{1.767853in}}%
\pgfpathcurveto{\pgfqpoint{1.847917in}{1.771971in}}{\pgfqpoint{1.842330in}{1.774285in}}{\pgfqpoint{1.836506in}{1.774285in}}%
\pgfpathcurveto{\pgfqpoint{1.830682in}{1.774285in}}{\pgfqpoint{1.825096in}{1.771971in}}{\pgfqpoint{1.820978in}{1.767853in}}%
\pgfpathcurveto{\pgfqpoint{1.816860in}{1.763735in}}{\pgfqpoint{1.814546in}{1.758149in}}{\pgfqpoint{1.814546in}{1.752325in}}%
\pgfpathcurveto{\pgfqpoint{1.814546in}{1.746501in}}{\pgfqpoint{1.816860in}{1.740915in}}{\pgfqpoint{1.820978in}{1.736797in}}%
\pgfpathcurveto{\pgfqpoint{1.825096in}{1.732679in}}{\pgfqpoint{1.830682in}{1.730365in}}{\pgfqpoint{1.836506in}{1.730365in}}%
\pgfpathlineto{\pgfqpoint{1.836506in}{1.730365in}}%
\pgfpathclose%
\pgfusepath{stroke,fill}%
\end{pgfscope}%
\begin{pgfscope}%
\pgfpathrectangle{\pgfqpoint{0.100000in}{0.209042in}}{\pgfqpoint{2.906250in}{2.906250in}}%
\pgfusepath{clip}%
\pgfsetbuttcap%
\pgfsetroundjoin%
\definecolor{currentfill}{rgb}{0.913354,0.330052,0.382563}%
\pgfsetfillcolor{currentfill}%
\pgfsetfillopacity{0.800000}%
\pgfsetlinewidth{1.003750pt}%
\definecolor{currentstroke}{rgb}{0.913354,0.330052,0.382563}%
\pgfsetstrokecolor{currentstroke}%
\pgfsetstrokeopacity{0.800000}%
\pgfsetdash{}{0pt}%
\pgfpathmoveto{\pgfqpoint{1.461706in}{1.175630in}}%
\pgfpathcurveto{\pgfqpoint{1.467530in}{1.175630in}}{\pgfqpoint{1.473116in}{1.177944in}}{\pgfqpoint{1.477234in}{1.182062in}}%
\pgfpathcurveto{\pgfqpoint{1.481352in}{1.186180in}}{\pgfqpoint{1.483666in}{1.191766in}}{\pgfqpoint{1.483666in}{1.197590in}}%
\pgfpathcurveto{\pgfqpoint{1.483666in}{1.203414in}}{\pgfqpoint{1.481352in}{1.209000in}}{\pgfqpoint{1.477234in}{1.213118in}}%
\pgfpathcurveto{\pgfqpoint{1.473116in}{1.217236in}}{\pgfqpoint{1.467530in}{1.219550in}}{\pgfqpoint{1.461706in}{1.219550in}}%
\pgfpathcurveto{\pgfqpoint{1.455882in}{1.219550in}}{\pgfqpoint{1.450296in}{1.217236in}}{\pgfqpoint{1.446178in}{1.213118in}}%
\pgfpathcurveto{\pgfqpoint{1.442059in}{1.209000in}}{\pgfqpoint{1.439746in}{1.203414in}}{\pgfqpoint{1.439746in}{1.197590in}}%
\pgfpathcurveto{\pgfqpoint{1.439746in}{1.191766in}}{\pgfqpoint{1.442059in}{1.186180in}}{\pgfqpoint{1.446178in}{1.182062in}}%
\pgfpathcurveto{\pgfqpoint{1.450296in}{1.177944in}}{\pgfqpoint{1.455882in}{1.175630in}}{\pgfqpoint{1.461706in}{1.175630in}}%
\pgfpathlineto{\pgfqpoint{1.461706in}{1.175630in}}%
\pgfpathclose%
\pgfusepath{stroke,fill}%
\end{pgfscope}%
\begin{pgfscope}%
\pgfpathrectangle{\pgfqpoint{0.100000in}{0.209042in}}{\pgfqpoint{2.906250in}{2.906250in}}%
\pgfusepath{clip}%
\pgfsetbuttcap%
\pgfsetroundjoin%
\definecolor{currentfill}{rgb}{0.218512,0.061158,0.425392}%
\pgfsetfillcolor{currentfill}%
\pgfsetfillopacity{0.800000}%
\pgfsetlinewidth{1.003750pt}%
\definecolor{currentstroke}{rgb}{0.218512,0.061158,0.425392}%
\pgfsetstrokecolor{currentstroke}%
\pgfsetstrokeopacity{0.800000}%
\pgfsetdash{}{0pt}%
\pgfpathmoveto{\pgfqpoint{1.592240in}{1.542201in}}%
\pgfpathcurveto{\pgfqpoint{1.598064in}{1.542201in}}{\pgfqpoint{1.603650in}{1.544515in}}{\pgfqpoint{1.607768in}{1.548633in}}%
\pgfpathcurveto{\pgfqpoint{1.611886in}{1.552751in}}{\pgfqpoint{1.614200in}{1.558337in}}{\pgfqpoint{1.614200in}{1.564161in}}%
\pgfpathcurveto{\pgfqpoint{1.614200in}{1.569985in}}{\pgfqpoint{1.611886in}{1.575571in}}{\pgfqpoint{1.607768in}{1.579690in}}%
\pgfpathcurveto{\pgfqpoint{1.603650in}{1.583808in}}{\pgfqpoint{1.598064in}{1.586122in}}{\pgfqpoint{1.592240in}{1.586122in}}%
\pgfpathcurveto{\pgfqpoint{1.586416in}{1.586122in}}{\pgfqpoint{1.580830in}{1.583808in}}{\pgfqpoint{1.576711in}{1.579690in}}%
\pgfpathcurveto{\pgfqpoint{1.572593in}{1.575571in}}{\pgfqpoint{1.570279in}{1.569985in}}{\pgfqpoint{1.570279in}{1.564161in}}%
\pgfpathcurveto{\pgfqpoint{1.570279in}{1.558337in}}{\pgfqpoint{1.572593in}{1.552751in}}{\pgfqpoint{1.576711in}{1.548633in}}%
\pgfpathcurveto{\pgfqpoint{1.580830in}{1.544515in}}{\pgfqpoint{1.586416in}{1.542201in}}{\pgfqpoint{1.592240in}{1.542201in}}%
\pgfpathlineto{\pgfqpoint{1.592240in}{1.542201in}}%
\pgfpathclose%
\pgfusepath{stroke,fill}%
\end{pgfscope}%
\begin{pgfscope}%
\pgfpathrectangle{\pgfqpoint{0.100000in}{0.209042in}}{\pgfqpoint{2.906250in}{2.906250in}}%
\pgfusepath{clip}%
\pgfsetbuttcap%
\pgfsetroundjoin%
\definecolor{currentfill}{rgb}{0.043830,0.033830,0.141886}%
\pgfsetfillcolor{currentfill}%
\pgfsetfillopacity{0.800000}%
\pgfsetlinewidth{1.003750pt}%
\definecolor{currentstroke}{rgb}{0.043830,0.033830,0.141886}%
\pgfsetstrokecolor{currentstroke}%
\pgfsetstrokeopacity{0.800000}%
\pgfsetdash{}{0pt}%
\pgfpathmoveto{\pgfqpoint{1.639425in}{1.731435in}}%
\pgfpathcurveto{\pgfqpoint{1.645249in}{1.731435in}}{\pgfqpoint{1.650835in}{1.733749in}}{\pgfqpoint{1.654953in}{1.737867in}}%
\pgfpathcurveto{\pgfqpoint{1.659072in}{1.741985in}}{\pgfqpoint{1.661385in}{1.747571in}}{\pgfqpoint{1.661385in}{1.753395in}}%
\pgfpathcurveto{\pgfqpoint{1.661385in}{1.759219in}}{\pgfqpoint{1.659072in}{1.764805in}}{\pgfqpoint{1.654953in}{1.768923in}}%
\pgfpathcurveto{\pgfqpoint{1.650835in}{1.773041in}}{\pgfqpoint{1.645249in}{1.775355in}}{\pgfqpoint{1.639425in}{1.775355in}}%
\pgfpathcurveto{\pgfqpoint{1.633601in}{1.775355in}}{\pgfqpoint{1.628015in}{1.773041in}}{\pgfqpoint{1.623897in}{1.768923in}}%
\pgfpathcurveto{\pgfqpoint{1.619779in}{1.764805in}}{\pgfqpoint{1.617465in}{1.759219in}}{\pgfqpoint{1.617465in}{1.753395in}}%
\pgfpathcurveto{\pgfqpoint{1.617465in}{1.747571in}}{\pgfqpoint{1.619779in}{1.741985in}}{\pgfqpoint{1.623897in}{1.737867in}}%
\pgfpathcurveto{\pgfqpoint{1.628015in}{1.733749in}}{\pgfqpoint{1.633601in}{1.731435in}}{\pgfqpoint{1.639425in}{1.731435in}}%
\pgfpathlineto{\pgfqpoint{1.639425in}{1.731435in}}%
\pgfpathclose%
\pgfusepath{stroke,fill}%
\end{pgfscope}%
\begin{pgfscope}%
\pgfpathrectangle{\pgfqpoint{0.100000in}{0.209042in}}{\pgfqpoint{2.906250in}{2.906250in}}%
\pgfusepath{clip}%
\pgfsetbuttcap%
\pgfsetroundjoin%
\definecolor{currentfill}{rgb}{0.378211,0.095332,0.500067}%
\pgfsetfillcolor{currentfill}%
\pgfsetfillopacity{0.800000}%
\pgfsetlinewidth{1.003750pt}%
\definecolor{currentstroke}{rgb}{0.378211,0.095332,0.500067}%
\pgfsetstrokecolor{currentstroke}%
\pgfsetstrokeopacity{0.800000}%
\pgfsetdash{}{0pt}%
\pgfpathmoveto{\pgfqpoint{1.528695in}{1.984596in}}%
\pgfpathcurveto{\pgfqpoint{1.534519in}{1.984596in}}{\pgfqpoint{1.540105in}{1.986910in}}{\pgfqpoint{1.544223in}{1.991028in}}%
\pgfpathcurveto{\pgfqpoint{1.548341in}{1.995146in}}{\pgfqpoint{1.550655in}{2.000732in}}{\pgfqpoint{1.550655in}{2.006556in}}%
\pgfpathcurveto{\pgfqpoint{1.550655in}{2.012380in}}{\pgfqpoint{1.548341in}{2.017966in}}{\pgfqpoint{1.544223in}{2.022084in}}%
\pgfpathcurveto{\pgfqpoint{1.540105in}{2.026203in}}{\pgfqpoint{1.534519in}{2.028516in}}{\pgfqpoint{1.528695in}{2.028516in}}%
\pgfpathcurveto{\pgfqpoint{1.522871in}{2.028516in}}{\pgfqpoint{1.517285in}{2.026203in}}{\pgfqpoint{1.513167in}{2.022084in}}%
\pgfpathcurveto{\pgfqpoint{1.509048in}{2.017966in}}{\pgfqpoint{1.506735in}{2.012380in}}{\pgfqpoint{1.506735in}{2.006556in}}%
\pgfpathcurveto{\pgfqpoint{1.506735in}{2.000732in}}{\pgfqpoint{1.509048in}{1.995146in}}{\pgfqpoint{1.513167in}{1.991028in}}%
\pgfpathcurveto{\pgfqpoint{1.517285in}{1.986910in}}{\pgfqpoint{1.522871in}{1.984596in}}{\pgfqpoint{1.528695in}{1.984596in}}%
\pgfpathlineto{\pgfqpoint{1.528695in}{1.984596in}}%
\pgfpathclose%
\pgfusepath{stroke,fill}%
\end{pgfscope}%
\begin{pgfscope}%
\pgfpathrectangle{\pgfqpoint{0.100000in}{0.209042in}}{\pgfqpoint{2.906250in}{2.906250in}}%
\pgfusepath{clip}%
\pgfsetbuttcap%
\pgfsetroundjoin%
\definecolor{currentfill}{rgb}{0.198177,0.063862,0.404009}%
\pgfsetfillcolor{currentfill}%
\pgfsetfillopacity{0.800000}%
\pgfsetlinewidth{1.003750pt}%
\definecolor{currentstroke}{rgb}{0.198177,0.063862,0.404009}%
\pgfsetstrokecolor{currentstroke}%
\pgfsetstrokeopacity{0.800000}%
\pgfsetdash{}{0pt}%
\pgfpathmoveto{\pgfqpoint{1.769439in}{1.599796in}}%
\pgfpathcurveto{\pgfqpoint{1.775263in}{1.599796in}}{\pgfqpoint{1.780849in}{1.602110in}}{\pgfqpoint{1.784967in}{1.606228in}}%
\pgfpathcurveto{\pgfqpoint{1.789085in}{1.610346in}}{\pgfqpoint{1.791399in}{1.615932in}}{\pgfqpoint{1.791399in}{1.621756in}}%
\pgfpathcurveto{\pgfqpoint{1.791399in}{1.627580in}}{\pgfqpoint{1.789085in}{1.633166in}}{\pgfqpoint{1.784967in}{1.637284in}}%
\pgfpathcurveto{\pgfqpoint{1.780849in}{1.641403in}}{\pgfqpoint{1.775263in}{1.643716in}}{\pgfqpoint{1.769439in}{1.643716in}}%
\pgfpathcurveto{\pgfqpoint{1.763615in}{1.643716in}}{\pgfqpoint{1.758029in}{1.641403in}}{\pgfqpoint{1.753911in}{1.637284in}}%
\pgfpathcurveto{\pgfqpoint{1.749793in}{1.633166in}}{\pgfqpoint{1.747479in}{1.627580in}}{\pgfqpoint{1.747479in}{1.621756in}}%
\pgfpathcurveto{\pgfqpoint{1.747479in}{1.615932in}}{\pgfqpoint{1.749793in}{1.610346in}}{\pgfqpoint{1.753911in}{1.606228in}}%
\pgfpathcurveto{\pgfqpoint{1.758029in}{1.602110in}}{\pgfqpoint{1.763615in}{1.599796in}}{\pgfqpoint{1.769439in}{1.599796in}}%
\pgfpathlineto{\pgfqpoint{1.769439in}{1.599796in}}%
\pgfpathclose%
\pgfusepath{stroke,fill}%
\end{pgfscope}%
\begin{pgfscope}%
\pgfpathrectangle{\pgfqpoint{0.100000in}{0.209042in}}{\pgfqpoint{2.906250in}{2.906250in}}%
\pgfusepath{clip}%
\pgfsetbuttcap%
\pgfsetroundjoin%
\definecolor{currentfill}{rgb}{0.506629,0.145958,0.507806}%
\pgfsetfillcolor{currentfill}%
\pgfsetfillopacity{0.800000}%
\pgfsetlinewidth{1.003750pt}%
\definecolor{currentstroke}{rgb}{0.506629,0.145958,0.507806}%
\pgfsetstrokecolor{currentstroke}%
\pgfsetstrokeopacity{0.800000}%
\pgfsetdash{}{0pt}%
\pgfpathmoveto{\pgfqpoint{1.893590in}{1.962710in}}%
\pgfpathcurveto{\pgfqpoint{1.899414in}{1.962710in}}{\pgfqpoint{1.905000in}{1.965024in}}{\pgfqpoint{1.909118in}{1.969142in}}%
\pgfpathcurveto{\pgfqpoint{1.913236in}{1.973260in}}{\pgfqpoint{1.915550in}{1.978847in}}{\pgfqpoint{1.915550in}{1.984670in}}%
\pgfpathcurveto{\pgfqpoint{1.915550in}{1.990494in}}{\pgfqpoint{1.913236in}{1.996081in}}{\pgfqpoint{1.909118in}{2.000199in}}%
\pgfpathcurveto{\pgfqpoint{1.905000in}{2.004317in}}{\pgfqpoint{1.899414in}{2.006631in}}{\pgfqpoint{1.893590in}{2.006631in}}%
\pgfpathcurveto{\pgfqpoint{1.887766in}{2.006631in}}{\pgfqpoint{1.882180in}{2.004317in}}{\pgfqpoint{1.878062in}{2.000199in}}%
\pgfpathcurveto{\pgfqpoint{1.873944in}{1.996081in}}{\pgfqpoint{1.871630in}{1.990494in}}{\pgfqpoint{1.871630in}{1.984670in}}%
\pgfpathcurveto{\pgfqpoint{1.871630in}{1.978847in}}{\pgfqpoint{1.873944in}{1.973260in}}{\pgfqpoint{1.878062in}{1.969142in}}%
\pgfpathcurveto{\pgfqpoint{1.882180in}{1.965024in}}{\pgfqpoint{1.887766in}{1.962710in}}{\pgfqpoint{1.893590in}{1.962710in}}%
\pgfpathlineto{\pgfqpoint{1.893590in}{1.962710in}}%
\pgfpathclose%
\pgfusepath{stroke,fill}%
\end{pgfscope}%
\begin{pgfscope}%
\pgfpathrectangle{\pgfqpoint{0.100000in}{0.209042in}}{\pgfqpoint{2.906250in}{2.906250in}}%
\pgfusepath{clip}%
\pgfsetbuttcap%
\pgfsetroundjoin%
\definecolor{currentfill}{rgb}{0.039608,0.031090,0.133515}%
\pgfsetfillcolor{currentfill}%
\pgfsetfillopacity{0.800000}%
\pgfsetlinewidth{1.003750pt}%
\definecolor{currentstroke}{rgb}{0.039608,0.031090,0.133515}%
\pgfsetstrokecolor{currentstroke}%
\pgfsetstrokeopacity{0.800000}%
\pgfsetdash{}{0pt}%
\pgfpathmoveto{\pgfqpoint{1.659333in}{1.722100in}}%
\pgfpathcurveto{\pgfqpoint{1.665157in}{1.722100in}}{\pgfqpoint{1.670743in}{1.724414in}}{\pgfqpoint{1.674861in}{1.728532in}}%
\pgfpathcurveto{\pgfqpoint{1.678979in}{1.732650in}}{\pgfqpoint{1.681293in}{1.738237in}}{\pgfqpoint{1.681293in}{1.744061in}}%
\pgfpathcurveto{\pgfqpoint{1.681293in}{1.749884in}}{\pgfqpoint{1.678979in}{1.755471in}}{\pgfqpoint{1.674861in}{1.759589in}}%
\pgfpathcurveto{\pgfqpoint{1.670743in}{1.763707in}}{\pgfqpoint{1.665157in}{1.766021in}}{\pgfqpoint{1.659333in}{1.766021in}}%
\pgfpathcurveto{\pgfqpoint{1.653509in}{1.766021in}}{\pgfqpoint{1.647923in}{1.763707in}}{\pgfqpoint{1.643804in}{1.759589in}}%
\pgfpathcurveto{\pgfqpoint{1.639686in}{1.755471in}}{\pgfqpoint{1.637372in}{1.749884in}}{\pgfqpoint{1.637372in}{1.744061in}}%
\pgfpathcurveto{\pgfqpoint{1.637372in}{1.738237in}}{\pgfqpoint{1.639686in}{1.732650in}}{\pgfqpoint{1.643804in}{1.728532in}}%
\pgfpathcurveto{\pgfqpoint{1.647923in}{1.724414in}}{\pgfqpoint{1.653509in}{1.722100in}}{\pgfqpoint{1.659333in}{1.722100in}}%
\pgfpathlineto{\pgfqpoint{1.659333in}{1.722100in}}%
\pgfpathclose%
\pgfusepath{stroke,fill}%
\end{pgfscope}%
\begin{pgfscope}%
\pgfpathrectangle{\pgfqpoint{0.100000in}{0.209042in}}{\pgfqpoint{2.906250in}{2.906250in}}%
\pgfusepath{clip}%
\pgfsetbuttcap%
\pgfsetroundjoin%
\definecolor{currentfill}{rgb}{0.463508,0.129893,0.507652}%
\pgfsetfillcolor{currentfill}%
\pgfsetfillopacity{0.800000}%
\pgfsetlinewidth{1.003750pt}%
\definecolor{currentstroke}{rgb}{0.463508,0.129893,0.507652}%
\pgfsetstrokecolor{currentstroke}%
\pgfsetstrokeopacity{0.800000}%
\pgfsetdash{}{0pt}%
\pgfpathmoveto{\pgfqpoint{1.836649in}{1.972359in}}%
\pgfpathcurveto{\pgfqpoint{1.842473in}{1.972359in}}{\pgfqpoint{1.848059in}{1.974673in}}{\pgfqpoint{1.852177in}{1.978791in}}%
\pgfpathcurveto{\pgfqpoint{1.856295in}{1.982909in}}{\pgfqpoint{1.858609in}{1.988495in}}{\pgfqpoint{1.858609in}{1.994319in}}%
\pgfpathcurveto{\pgfqpoint{1.858609in}{2.000143in}}{\pgfqpoint{1.856295in}{2.005729in}}{\pgfqpoint{1.852177in}{2.009848in}}%
\pgfpathcurveto{\pgfqpoint{1.848059in}{2.013966in}}{\pgfqpoint{1.842473in}{2.016280in}}{\pgfqpoint{1.836649in}{2.016280in}}%
\pgfpathcurveto{\pgfqpoint{1.830825in}{2.016280in}}{\pgfqpoint{1.825239in}{2.013966in}}{\pgfqpoint{1.821121in}{2.009848in}}%
\pgfpathcurveto{\pgfqpoint{1.817002in}{2.005729in}}{\pgfqpoint{1.814689in}{2.000143in}}{\pgfqpoint{1.814689in}{1.994319in}}%
\pgfpathcurveto{\pgfqpoint{1.814689in}{1.988495in}}{\pgfqpoint{1.817002in}{1.982909in}}{\pgfqpoint{1.821121in}{1.978791in}}%
\pgfpathcurveto{\pgfqpoint{1.825239in}{1.974673in}}{\pgfqpoint{1.830825in}{1.972359in}}{\pgfqpoint{1.836649in}{1.972359in}}%
\pgfpathlineto{\pgfqpoint{1.836649in}{1.972359in}}%
\pgfpathclose%
\pgfusepath{stroke,fill}%
\end{pgfscope}%
\begin{pgfscope}%
\pgfpathrectangle{\pgfqpoint{0.100000in}{0.209042in}}{\pgfqpoint{2.906250in}{2.906250in}}%
\pgfusepath{clip}%
\pgfsetbuttcap%
\pgfsetroundjoin%
\definecolor{currentfill}{rgb}{0.123833,0.067295,0.295879}%
\pgfsetfillcolor{currentfill}%
\pgfsetfillopacity{0.800000}%
\pgfsetlinewidth{1.003750pt}%
\definecolor{currentstroke}{rgb}{0.123833,0.067295,0.295879}%
\pgfsetstrokecolor{currentstroke}%
\pgfsetstrokeopacity{0.800000}%
\pgfsetdash{}{0pt}%
\pgfpathmoveto{\pgfqpoint{1.569743in}{1.844491in}}%
\pgfpathcurveto{\pgfqpoint{1.575567in}{1.844491in}}{\pgfqpoint{1.581153in}{1.846805in}}{\pgfqpoint{1.585271in}{1.850923in}}%
\pgfpathcurveto{\pgfqpoint{1.589389in}{1.855041in}}{\pgfqpoint{1.591703in}{1.860628in}}{\pgfqpoint{1.591703in}{1.866452in}}%
\pgfpathcurveto{\pgfqpoint{1.591703in}{1.872276in}}{\pgfqpoint{1.589389in}{1.877862in}}{\pgfqpoint{1.585271in}{1.881980in}}%
\pgfpathcurveto{\pgfqpoint{1.581153in}{1.886098in}}{\pgfqpoint{1.575567in}{1.888412in}}{\pgfqpoint{1.569743in}{1.888412in}}%
\pgfpathcurveto{\pgfqpoint{1.563919in}{1.888412in}}{\pgfqpoint{1.558333in}{1.886098in}}{\pgfqpoint{1.554215in}{1.881980in}}%
\pgfpathcurveto{\pgfqpoint{1.550097in}{1.877862in}}{\pgfqpoint{1.547783in}{1.872276in}}{\pgfqpoint{1.547783in}{1.866452in}}%
\pgfpathcurveto{\pgfqpoint{1.547783in}{1.860628in}}{\pgfqpoint{1.550097in}{1.855041in}}{\pgfqpoint{1.554215in}{1.850923in}}%
\pgfpathcurveto{\pgfqpoint{1.558333in}{1.846805in}}{\pgfqpoint{1.563919in}{1.844491in}}{\pgfqpoint{1.569743in}{1.844491in}}%
\pgfpathlineto{\pgfqpoint{1.569743in}{1.844491in}}%
\pgfpathclose%
\pgfusepath{stroke,fill}%
\end{pgfscope}%
\begin{pgfscope}%
\pgfpathrectangle{\pgfqpoint{0.100000in}{0.209042in}}{\pgfqpoint{2.906250in}{2.906250in}}%
\pgfusepath{clip}%
\pgfsetbuttcap%
\pgfsetroundjoin%
\definecolor{currentfill}{rgb}{0.588158,0.173652,0.502035}%
\pgfsetfillcolor{currentfill}%
\pgfsetfillopacity{0.800000}%
\pgfsetlinewidth{1.003750pt}%
\definecolor{currentstroke}{rgb}{0.588158,0.173652,0.502035}%
\pgfsetstrokecolor{currentstroke}%
\pgfsetstrokeopacity{0.800000}%
\pgfsetdash{}{0pt}%
\pgfpathmoveto{\pgfqpoint{2.061245in}{1.828244in}}%
\pgfpathcurveto{\pgfqpoint{2.067069in}{1.828244in}}{\pgfqpoint{2.072655in}{1.830558in}}{\pgfqpoint{2.076773in}{1.834676in}}%
\pgfpathcurveto{\pgfqpoint{2.080891in}{1.838794in}}{\pgfqpoint{2.083205in}{1.844380in}}{\pgfqpoint{2.083205in}{1.850204in}}%
\pgfpathcurveto{\pgfqpoint{2.083205in}{1.856028in}}{\pgfqpoint{2.080891in}{1.861614in}}{\pgfqpoint{2.076773in}{1.865733in}}%
\pgfpathcurveto{\pgfqpoint{2.072655in}{1.869851in}}{\pgfqpoint{2.067069in}{1.872165in}}{\pgfqpoint{2.061245in}{1.872165in}}%
\pgfpathcurveto{\pgfqpoint{2.055421in}{1.872165in}}{\pgfqpoint{2.049835in}{1.869851in}}{\pgfqpoint{2.045717in}{1.865733in}}%
\pgfpathcurveto{\pgfqpoint{2.041599in}{1.861614in}}{\pgfqpoint{2.039285in}{1.856028in}}{\pgfqpoint{2.039285in}{1.850204in}}%
\pgfpathcurveto{\pgfqpoint{2.039285in}{1.844380in}}{\pgfqpoint{2.041599in}{1.838794in}}{\pgfqpoint{2.045717in}{1.834676in}}%
\pgfpathcurveto{\pgfqpoint{2.049835in}{1.830558in}}{\pgfqpoint{2.055421in}{1.828244in}}{\pgfqpoint{2.061245in}{1.828244in}}%
\pgfpathlineto{\pgfqpoint{2.061245in}{1.828244in}}%
\pgfpathclose%
\pgfusepath{stroke,fill}%
\end{pgfscope}%
\begin{pgfscope}%
\pgfpathrectangle{\pgfqpoint{0.100000in}{0.209042in}}{\pgfqpoint{2.906250in}{2.906250in}}%
\pgfusepath{clip}%
\pgfsetbuttcap%
\pgfsetroundjoin%
\definecolor{currentfill}{rgb}{0.463508,0.129893,0.507652}%
\pgfsetfillcolor{currentfill}%
\pgfsetfillopacity{0.800000}%
\pgfsetlinewidth{1.003750pt}%
\definecolor{currentstroke}{rgb}{0.463508,0.129893,0.507652}%
\pgfsetstrokecolor{currentstroke}%
\pgfsetstrokeopacity{0.800000}%
\pgfsetdash{}{0pt}%
\pgfpathmoveto{\pgfqpoint{1.956297in}{1.859983in}}%
\pgfpathcurveto{\pgfqpoint{1.962121in}{1.859983in}}{\pgfqpoint{1.967707in}{1.862297in}}{\pgfqpoint{1.971826in}{1.866415in}}%
\pgfpathcurveto{\pgfqpoint{1.975944in}{1.870534in}}{\pgfqpoint{1.978258in}{1.876120in}}{\pgfqpoint{1.978258in}{1.881944in}}%
\pgfpathcurveto{\pgfqpoint{1.978258in}{1.887768in}}{\pgfqpoint{1.975944in}{1.893354in}}{\pgfqpoint{1.971826in}{1.897472in}}%
\pgfpathcurveto{\pgfqpoint{1.967707in}{1.901590in}}{\pgfqpoint{1.962121in}{1.903904in}}{\pgfqpoint{1.956297in}{1.903904in}}%
\pgfpathcurveto{\pgfqpoint{1.950473in}{1.903904in}}{\pgfqpoint{1.944887in}{1.901590in}}{\pgfqpoint{1.940769in}{1.897472in}}%
\pgfpathcurveto{\pgfqpoint{1.936651in}{1.893354in}}{\pgfqpoint{1.934337in}{1.887768in}}{\pgfqpoint{1.934337in}{1.881944in}}%
\pgfpathcurveto{\pgfqpoint{1.934337in}{1.876120in}}{\pgfqpoint{1.936651in}{1.870534in}}{\pgfqpoint{1.940769in}{1.866415in}}%
\pgfpathcurveto{\pgfqpoint{1.944887in}{1.862297in}}{\pgfqpoint{1.950473in}{1.859983in}}{\pgfqpoint{1.956297in}{1.859983in}}%
\pgfpathlineto{\pgfqpoint{1.956297in}{1.859983in}}%
\pgfpathclose%
\pgfusepath{stroke,fill}%
\end{pgfscope}%
\begin{pgfscope}%
\pgfpathrectangle{\pgfqpoint{0.100000in}{0.209042in}}{\pgfqpoint{2.906250in}{2.906250in}}%
\pgfusepath{clip}%
\pgfsetbuttcap%
\pgfsetroundjoin%
\definecolor{currentfill}{rgb}{0.097833,0.061531,0.248477}%
\pgfsetfillcolor{currentfill}%
\pgfsetfillopacity{0.800000}%
\pgfsetlinewidth{1.003750pt}%
\definecolor{currentstroke}{rgb}{0.097833,0.061531,0.248477}%
\pgfsetstrokecolor{currentstroke}%
\pgfsetstrokeopacity{0.800000}%
\pgfsetdash{}{0pt}%
\pgfpathmoveto{\pgfqpoint{1.625819in}{1.606242in}}%
\pgfpathcurveto{\pgfqpoint{1.631643in}{1.606242in}}{\pgfqpoint{1.637229in}{1.608556in}}{\pgfqpoint{1.641347in}{1.612674in}}%
\pgfpathcurveto{\pgfqpoint{1.645465in}{1.616793in}}{\pgfqpoint{1.647779in}{1.622379in}}{\pgfqpoint{1.647779in}{1.628203in}}%
\pgfpathcurveto{\pgfqpoint{1.647779in}{1.634027in}}{\pgfqpoint{1.645465in}{1.639613in}}{\pgfqpoint{1.641347in}{1.643731in}}%
\pgfpathcurveto{\pgfqpoint{1.637229in}{1.647849in}}{\pgfqpoint{1.631643in}{1.650163in}}{\pgfqpoint{1.625819in}{1.650163in}}%
\pgfpathcurveto{\pgfqpoint{1.619995in}{1.650163in}}{\pgfqpoint{1.614409in}{1.647849in}}{\pgfqpoint{1.610291in}{1.643731in}}%
\pgfpathcurveto{\pgfqpoint{1.606172in}{1.639613in}}{\pgfqpoint{1.603859in}{1.634027in}}{\pgfqpoint{1.603859in}{1.628203in}}%
\pgfpathcurveto{\pgfqpoint{1.603859in}{1.622379in}}{\pgfqpoint{1.606172in}{1.616793in}}{\pgfqpoint{1.610291in}{1.612674in}}%
\pgfpathcurveto{\pgfqpoint{1.614409in}{1.608556in}}{\pgfqpoint{1.619995in}{1.606242in}}{\pgfqpoint{1.625819in}{1.606242in}}%
\pgfpathlineto{\pgfqpoint{1.625819in}{1.606242in}}%
\pgfpathclose%
\pgfusepath{stroke,fill}%
\end{pgfscope}%
\begin{pgfscope}%
\pgfpathrectangle{\pgfqpoint{0.100000in}{0.209042in}}{\pgfqpoint{2.906250in}{2.906250in}}%
\pgfusepath{clip}%
\pgfsetbuttcap%
\pgfsetroundjoin%
\definecolor{currentfill}{rgb}{0.165308,0.067911,0.361816}%
\pgfsetfillcolor{currentfill}%
\pgfsetfillopacity{0.800000}%
\pgfsetlinewidth{1.003750pt}%
\definecolor{currentstroke}{rgb}{0.165308,0.067911,0.361816}%
\pgfsetstrokecolor{currentstroke}%
\pgfsetstrokeopacity{0.800000}%
\pgfsetdash{}{0pt}%
\pgfpathmoveto{\pgfqpoint{1.808649in}{1.709138in}}%
\pgfpathcurveto{\pgfqpoint{1.814473in}{1.709138in}}{\pgfqpoint{1.820059in}{1.711452in}}{\pgfqpoint{1.824177in}{1.715570in}}%
\pgfpathcurveto{\pgfqpoint{1.828295in}{1.719688in}}{\pgfqpoint{1.830609in}{1.725275in}}{\pgfqpoint{1.830609in}{1.731099in}}%
\pgfpathcurveto{\pgfqpoint{1.830609in}{1.736922in}}{\pgfqpoint{1.828295in}{1.742509in}}{\pgfqpoint{1.824177in}{1.746627in}}%
\pgfpathcurveto{\pgfqpoint{1.820059in}{1.750745in}}{\pgfqpoint{1.814473in}{1.753059in}}{\pgfqpoint{1.808649in}{1.753059in}}%
\pgfpathcurveto{\pgfqpoint{1.802825in}{1.753059in}}{\pgfqpoint{1.797239in}{1.750745in}}{\pgfqpoint{1.793121in}{1.746627in}}%
\pgfpathcurveto{\pgfqpoint{1.789002in}{1.742509in}}{\pgfqpoint{1.786689in}{1.736922in}}{\pgfqpoint{1.786689in}{1.731099in}}%
\pgfpathcurveto{\pgfqpoint{1.786689in}{1.725275in}}{\pgfqpoint{1.789002in}{1.719688in}}{\pgfqpoint{1.793121in}{1.715570in}}%
\pgfpathcurveto{\pgfqpoint{1.797239in}{1.711452in}}{\pgfqpoint{1.802825in}{1.709138in}}{\pgfqpoint{1.808649in}{1.709138in}}%
\pgfpathlineto{\pgfqpoint{1.808649in}{1.709138in}}%
\pgfpathclose%
\pgfusepath{stroke,fill}%
\end{pgfscope}%
\begin{pgfscope}%
\pgfpathrectangle{\pgfqpoint{0.100000in}{0.209042in}}{\pgfqpoint{2.906250in}{2.906250in}}%
\pgfusepath{clip}%
\pgfsetbuttcap%
\pgfsetroundjoin%
\definecolor{currentfill}{rgb}{0.929845,0.352734,0.372677}%
\pgfsetfillcolor{currentfill}%
\pgfsetfillopacity{0.800000}%
\pgfsetlinewidth{1.003750pt}%
\definecolor{currentstroke}{rgb}{0.929845,0.352734,0.372677}%
\pgfsetstrokecolor{currentstroke}%
\pgfsetstrokeopacity{0.800000}%
\pgfsetdash{}{0pt}%
\pgfpathmoveto{\pgfqpoint{2.328200in}{1.625048in}}%
\pgfpathcurveto{\pgfqpoint{2.334024in}{1.625048in}}{\pgfqpoint{2.339610in}{1.627362in}}{\pgfqpoint{2.343728in}{1.631480in}}%
\pgfpathcurveto{\pgfqpoint{2.347846in}{1.635598in}}{\pgfqpoint{2.350160in}{1.641184in}}{\pgfqpoint{2.350160in}{1.647008in}}%
\pgfpathcurveto{\pgfqpoint{2.350160in}{1.652832in}}{\pgfqpoint{2.347846in}{1.658418in}}{\pgfqpoint{2.343728in}{1.662537in}}%
\pgfpathcurveto{\pgfqpoint{2.339610in}{1.666655in}}{\pgfqpoint{2.334024in}{1.668969in}}{\pgfqpoint{2.328200in}{1.668969in}}%
\pgfpathcurveto{\pgfqpoint{2.322376in}{1.668969in}}{\pgfqpoint{2.316790in}{1.666655in}}{\pgfqpoint{2.312672in}{1.662537in}}%
\pgfpathcurveto{\pgfqpoint{2.308554in}{1.658418in}}{\pgfqpoint{2.306240in}{1.652832in}}{\pgfqpoint{2.306240in}{1.647008in}}%
\pgfpathcurveto{\pgfqpoint{2.306240in}{1.641184in}}{\pgfqpoint{2.308554in}{1.635598in}}{\pgfqpoint{2.312672in}{1.631480in}}%
\pgfpathcurveto{\pgfqpoint{2.316790in}{1.627362in}}{\pgfqpoint{2.322376in}{1.625048in}}{\pgfqpoint{2.328200in}{1.625048in}}%
\pgfpathlineto{\pgfqpoint{2.328200in}{1.625048in}}%
\pgfpathclose%
\pgfusepath{stroke,fill}%
\end{pgfscope}%
\begin{pgfscope}%
\pgfpathrectangle{\pgfqpoint{0.100000in}{0.209042in}}{\pgfqpoint{2.906250in}{2.906250in}}%
\pgfusepath{clip}%
\pgfsetbuttcap%
\pgfsetroundjoin%
\definecolor{currentfill}{rgb}{0.445163,0.122724,0.506901}%
\pgfsetfillcolor{currentfill}%
\pgfsetfillopacity{0.800000}%
\pgfsetlinewidth{1.003750pt}%
\definecolor{currentstroke}{rgb}{0.445163,0.122724,0.506901}%
\pgfsetstrokecolor{currentstroke}%
\pgfsetstrokeopacity{0.800000}%
\pgfsetdash{}{0pt}%
\pgfpathmoveto{\pgfqpoint{1.448252in}{2.006265in}}%
\pgfpathcurveto{\pgfqpoint{1.454076in}{2.006265in}}{\pgfqpoint{1.459662in}{2.008579in}}{\pgfqpoint{1.463780in}{2.012697in}}%
\pgfpathcurveto{\pgfqpoint{1.467898in}{2.016816in}}{\pgfqpoint{1.470212in}{2.022402in}}{\pgfqpoint{1.470212in}{2.028226in}}%
\pgfpathcurveto{\pgfqpoint{1.470212in}{2.034050in}}{\pgfqpoint{1.467898in}{2.039636in}}{\pgfqpoint{1.463780in}{2.043754in}}%
\pgfpathcurveto{\pgfqpoint{1.459662in}{2.047872in}}{\pgfqpoint{1.454076in}{2.050186in}}{\pgfqpoint{1.448252in}{2.050186in}}%
\pgfpathcurveto{\pgfqpoint{1.442428in}{2.050186in}}{\pgfqpoint{1.436842in}{2.047872in}}{\pgfqpoint{1.432724in}{2.043754in}}%
\pgfpathcurveto{\pgfqpoint{1.428606in}{2.039636in}}{\pgfqpoint{1.426292in}{2.034050in}}{\pgfqpoint{1.426292in}{2.028226in}}%
\pgfpathcurveto{\pgfqpoint{1.426292in}{2.022402in}}{\pgfqpoint{1.428606in}{2.016816in}}{\pgfqpoint{1.432724in}{2.012697in}}%
\pgfpathcurveto{\pgfqpoint{1.436842in}{2.008579in}}{\pgfqpoint{1.442428in}{2.006265in}}{\pgfqpoint{1.448252in}{2.006265in}}%
\pgfpathlineto{\pgfqpoint{1.448252in}{2.006265in}}%
\pgfpathclose%
\pgfusepath{stroke,fill}%
\end{pgfscope}%
\begin{pgfscope}%
\pgfpathrectangle{\pgfqpoint{0.100000in}{0.209042in}}{\pgfqpoint{2.906250in}{2.906250in}}%
\pgfusepath{clip}%
\pgfsetbuttcap%
\pgfsetroundjoin%
\definecolor{currentfill}{rgb}{0.402548,0.105420,0.503386}%
\pgfsetfillcolor{currentfill}%
\pgfsetfillopacity{0.800000}%
\pgfsetlinewidth{1.003750pt}%
\definecolor{currentstroke}{rgb}{0.402548,0.105420,0.503386}%
\pgfsetstrokecolor{currentstroke}%
\pgfsetstrokeopacity{0.800000}%
\pgfsetdash{}{0pt}%
\pgfpathmoveto{\pgfqpoint{1.962120in}{1.634827in}}%
\pgfpathcurveto{\pgfqpoint{1.967944in}{1.634827in}}{\pgfqpoint{1.973530in}{1.637141in}}{\pgfqpoint{1.977648in}{1.641259in}}%
\pgfpathcurveto{\pgfqpoint{1.981766in}{1.645377in}}{\pgfqpoint{1.984080in}{1.650963in}}{\pgfqpoint{1.984080in}{1.656787in}}%
\pgfpathcurveto{\pgfqpoint{1.984080in}{1.662611in}}{\pgfqpoint{1.981766in}{1.668197in}}{\pgfqpoint{1.977648in}{1.672316in}}%
\pgfpathcurveto{\pgfqpoint{1.973530in}{1.676434in}}{\pgfqpoint{1.967944in}{1.678748in}}{\pgfqpoint{1.962120in}{1.678748in}}%
\pgfpathcurveto{\pgfqpoint{1.956296in}{1.678748in}}{\pgfqpoint{1.950710in}{1.676434in}}{\pgfqpoint{1.946592in}{1.672316in}}%
\pgfpathcurveto{\pgfqpoint{1.942473in}{1.668197in}}{\pgfqpoint{1.940160in}{1.662611in}}{\pgfqpoint{1.940160in}{1.656787in}}%
\pgfpathcurveto{\pgfqpoint{1.940160in}{1.650963in}}{\pgfqpoint{1.942473in}{1.645377in}}{\pgfqpoint{1.946592in}{1.641259in}}%
\pgfpathcurveto{\pgfqpoint{1.950710in}{1.637141in}}{\pgfqpoint{1.956296in}{1.634827in}}{\pgfqpoint{1.962120in}{1.634827in}}%
\pgfpathlineto{\pgfqpoint{1.962120in}{1.634827in}}%
\pgfpathclose%
\pgfusepath{stroke,fill}%
\end{pgfscope}%
\begin{pgfscope}%
\pgfpathrectangle{\pgfqpoint{0.100000in}{0.209042in}}{\pgfqpoint{2.906250in}{2.906250in}}%
\pgfusepath{clip}%
\pgfsetbuttcap%
\pgfsetroundjoin%
\definecolor{currentfill}{rgb}{0.488088,0.139186,0.508011}%
\pgfsetfillcolor{currentfill}%
\pgfsetfillopacity{0.800000}%
\pgfsetlinewidth{1.003750pt}%
\definecolor{currentstroke}{rgb}{0.488088,0.139186,0.508011}%
\pgfsetstrokecolor{currentstroke}%
\pgfsetstrokeopacity{0.800000}%
\pgfsetdash{}{0pt}%
\pgfpathmoveto{\pgfqpoint{1.596102in}{1.381326in}}%
\pgfpathcurveto{\pgfqpoint{1.601926in}{1.381326in}}{\pgfqpoint{1.607512in}{1.383639in}}{\pgfqpoint{1.611630in}{1.387758in}}%
\pgfpathcurveto{\pgfqpoint{1.615748in}{1.391876in}}{\pgfqpoint{1.618062in}{1.397462in}}{\pgfqpoint{1.618062in}{1.403286in}}%
\pgfpathcurveto{\pgfqpoint{1.618062in}{1.409110in}}{\pgfqpoint{1.615748in}{1.414696in}}{\pgfqpoint{1.611630in}{1.418814in}}%
\pgfpathcurveto{\pgfqpoint{1.607512in}{1.422932in}}{\pgfqpoint{1.601926in}{1.425246in}}{\pgfqpoint{1.596102in}{1.425246in}}%
\pgfpathcurveto{\pgfqpoint{1.590278in}{1.425246in}}{\pgfqpoint{1.584692in}{1.422932in}}{\pgfqpoint{1.580574in}{1.418814in}}%
\pgfpathcurveto{\pgfqpoint{1.576456in}{1.414696in}}{\pgfqpoint{1.574142in}{1.409110in}}{\pgfqpoint{1.574142in}{1.403286in}}%
\pgfpathcurveto{\pgfqpoint{1.574142in}{1.397462in}}{\pgfqpoint{1.576456in}{1.391876in}}{\pgfqpoint{1.580574in}{1.387758in}}%
\pgfpathcurveto{\pgfqpoint{1.584692in}{1.383639in}}{\pgfqpoint{1.590278in}{1.381326in}}{\pgfqpoint{1.596102in}{1.381326in}}%
\pgfpathlineto{\pgfqpoint{1.596102in}{1.381326in}}%
\pgfpathclose%
\pgfusepath{stroke,fill}%
\end{pgfscope}%
\begin{pgfscope}%
\pgfpathrectangle{\pgfqpoint{0.100000in}{0.209042in}}{\pgfqpoint{2.906250in}{2.906250in}}%
\pgfusepath{clip}%
\pgfsetbuttcap%
\pgfsetroundjoin%
\definecolor{currentfill}{rgb}{0.152839,0.068637,0.343404}%
\pgfsetfillcolor{currentfill}%
\pgfsetfillopacity{0.800000}%
\pgfsetlinewidth{1.003750pt}%
\definecolor{currentstroke}{rgb}{0.152839,0.068637,0.343404}%
\pgfsetstrokecolor{currentstroke}%
\pgfsetstrokeopacity{0.800000}%
\pgfsetdash{}{0pt}%
\pgfpathmoveto{\pgfqpoint{1.439051in}{1.780091in}}%
\pgfpathcurveto{\pgfqpoint{1.444875in}{1.780091in}}{\pgfqpoint{1.450461in}{1.782405in}}{\pgfqpoint{1.454579in}{1.786523in}}%
\pgfpathcurveto{\pgfqpoint{1.458697in}{1.790642in}}{\pgfqpoint{1.461011in}{1.796228in}}{\pgfqpoint{1.461011in}{1.802052in}}%
\pgfpathcurveto{\pgfqpoint{1.461011in}{1.807876in}}{\pgfqpoint{1.458697in}{1.813462in}}{\pgfqpoint{1.454579in}{1.817580in}}%
\pgfpathcurveto{\pgfqpoint{1.450461in}{1.821698in}}{\pgfqpoint{1.444875in}{1.824012in}}{\pgfqpoint{1.439051in}{1.824012in}}%
\pgfpathcurveto{\pgfqpoint{1.433227in}{1.824012in}}{\pgfqpoint{1.427641in}{1.821698in}}{\pgfqpoint{1.423523in}{1.817580in}}%
\pgfpathcurveto{\pgfqpoint{1.419405in}{1.813462in}}{\pgfqpoint{1.417091in}{1.807876in}}{\pgfqpoint{1.417091in}{1.802052in}}%
\pgfpathcurveto{\pgfqpoint{1.417091in}{1.796228in}}{\pgfqpoint{1.419405in}{1.790642in}}{\pgfqpoint{1.423523in}{1.786523in}}%
\pgfpathcurveto{\pgfqpoint{1.427641in}{1.782405in}}{\pgfqpoint{1.433227in}{1.780091in}}{\pgfqpoint{1.439051in}{1.780091in}}%
\pgfpathlineto{\pgfqpoint{1.439051in}{1.780091in}}%
\pgfpathclose%
\pgfusepath{stroke,fill}%
\end{pgfscope}%
\begin{pgfscope}%
\pgfpathrectangle{\pgfqpoint{0.100000in}{0.209042in}}{\pgfqpoint{2.906250in}{2.906250in}}%
\pgfusepath{clip}%
\pgfsetbuttcap%
\pgfsetroundjoin%
\definecolor{currentfill}{rgb}{0.304081,0.067835,0.480812}%
\pgfsetfillcolor{currentfill}%
\pgfsetfillopacity{0.800000}%
\pgfsetlinewidth{1.003750pt}%
\definecolor{currentstroke}{rgb}{0.304081,0.067835,0.480812}%
\pgfsetstrokecolor{currentstroke}%
\pgfsetstrokeopacity{0.800000}%
\pgfsetdash{}{0pt}%
\pgfpathmoveto{\pgfqpoint{1.820259in}{1.873142in}}%
\pgfpathcurveto{\pgfqpoint{1.826083in}{1.873142in}}{\pgfqpoint{1.831670in}{1.875456in}}{\pgfqpoint{1.835788in}{1.879574in}}%
\pgfpathcurveto{\pgfqpoint{1.839906in}{1.883692in}}{\pgfqpoint{1.842220in}{1.889278in}}{\pgfqpoint{1.842220in}{1.895102in}}%
\pgfpathcurveto{\pgfqpoint{1.842220in}{1.900926in}}{\pgfqpoint{1.839906in}{1.906512in}}{\pgfqpoint{1.835788in}{1.910631in}}%
\pgfpathcurveto{\pgfqpoint{1.831670in}{1.914749in}}{\pgfqpoint{1.826083in}{1.917063in}}{\pgfqpoint{1.820259in}{1.917063in}}%
\pgfpathcurveto{\pgfqpoint{1.814435in}{1.917063in}}{\pgfqpoint{1.808849in}{1.914749in}}{\pgfqpoint{1.804731in}{1.910631in}}%
\pgfpathcurveto{\pgfqpoint{1.800613in}{1.906512in}}{\pgfqpoint{1.798299in}{1.900926in}}{\pgfqpoint{1.798299in}{1.895102in}}%
\pgfpathcurveto{\pgfqpoint{1.798299in}{1.889278in}}{\pgfqpoint{1.800613in}{1.883692in}}{\pgfqpoint{1.804731in}{1.879574in}}%
\pgfpathcurveto{\pgfqpoint{1.808849in}{1.875456in}}{\pgfqpoint{1.814435in}{1.873142in}}{\pgfqpoint{1.820259in}{1.873142in}}%
\pgfpathlineto{\pgfqpoint{1.820259in}{1.873142in}}%
\pgfpathclose%
\pgfusepath{stroke,fill}%
\end{pgfscope}%
\begin{pgfscope}%
\pgfpathrectangle{\pgfqpoint{0.100000in}{0.209042in}}{\pgfqpoint{2.906250in}{2.906250in}}%
\pgfusepath{clip}%
\pgfsetbuttcap%
\pgfsetroundjoin%
\definecolor{currentfill}{rgb}{0.378211,0.095332,0.500067}%
\pgfsetfillcolor{currentfill}%
\pgfsetfillopacity{0.800000}%
\pgfsetlinewidth{1.003750pt}%
\definecolor{currentstroke}{rgb}{0.378211,0.095332,0.500067}%
\pgfsetstrokecolor{currentstroke}%
\pgfsetstrokeopacity{0.800000}%
\pgfsetdash{}{0pt}%
\pgfpathmoveto{\pgfqpoint{1.674962in}{1.443328in}}%
\pgfpathcurveto{\pgfqpoint{1.680786in}{1.443328in}}{\pgfqpoint{1.686373in}{1.445642in}}{\pgfqpoint{1.690491in}{1.449760in}}%
\pgfpathcurveto{\pgfqpoint{1.694609in}{1.453878in}}{\pgfqpoint{1.696923in}{1.459465in}}{\pgfqpoint{1.696923in}{1.465289in}}%
\pgfpathcurveto{\pgfqpoint{1.696923in}{1.471113in}}{\pgfqpoint{1.694609in}{1.476699in}}{\pgfqpoint{1.690491in}{1.480817in}}%
\pgfpathcurveto{\pgfqpoint{1.686373in}{1.484935in}}{\pgfqpoint{1.680786in}{1.487249in}}{\pgfqpoint{1.674962in}{1.487249in}}%
\pgfpathcurveto{\pgfqpoint{1.669139in}{1.487249in}}{\pgfqpoint{1.663552in}{1.484935in}}{\pgfqpoint{1.659434in}{1.480817in}}%
\pgfpathcurveto{\pgfqpoint{1.655316in}{1.476699in}}{\pgfqpoint{1.653002in}{1.471113in}}{\pgfqpoint{1.653002in}{1.465289in}}%
\pgfpathcurveto{\pgfqpoint{1.653002in}{1.459465in}}{\pgfqpoint{1.655316in}{1.453878in}}{\pgfqpoint{1.659434in}{1.449760in}}%
\pgfpathcurveto{\pgfqpoint{1.663552in}{1.445642in}}{\pgfqpoint{1.669139in}{1.443328in}}{\pgfqpoint{1.674962in}{1.443328in}}%
\pgfpathlineto{\pgfqpoint{1.674962in}{1.443328in}}%
\pgfpathclose%
\pgfusepath{stroke,fill}%
\end{pgfscope}%
\begin{pgfscope}%
\pgfpathrectangle{\pgfqpoint{0.100000in}{0.209042in}}{\pgfqpoint{2.906250in}{2.906250in}}%
\pgfusepath{clip}%
\pgfsetbuttcap%
\pgfsetroundjoin%
\definecolor{currentfill}{rgb}{0.140858,0.068654,0.324538}%
\pgfsetfillcolor{currentfill}%
\pgfsetfillopacity{0.800000}%
\pgfsetlinewidth{1.003750pt}%
\definecolor{currentstroke}{rgb}{0.140858,0.068654,0.324538}%
\pgfsetstrokecolor{currentstroke}%
\pgfsetstrokeopacity{0.800000}%
\pgfsetdash{}{0pt}%
\pgfpathmoveto{\pgfqpoint{1.729385in}{1.595502in}}%
\pgfpathcurveto{\pgfqpoint{1.735208in}{1.595502in}}{\pgfqpoint{1.740795in}{1.597816in}}{\pgfqpoint{1.744913in}{1.601934in}}%
\pgfpathcurveto{\pgfqpoint{1.749031in}{1.606052in}}{\pgfqpoint{1.751345in}{1.611638in}}{\pgfqpoint{1.751345in}{1.617462in}}%
\pgfpathcurveto{\pgfqpoint{1.751345in}{1.623286in}}{\pgfqpoint{1.749031in}{1.628872in}}{\pgfqpoint{1.744913in}{1.632990in}}%
\pgfpathcurveto{\pgfqpoint{1.740795in}{1.637108in}}{\pgfqpoint{1.735208in}{1.639422in}}{\pgfqpoint{1.729385in}{1.639422in}}%
\pgfpathcurveto{\pgfqpoint{1.723561in}{1.639422in}}{\pgfqpoint{1.717974in}{1.637108in}}{\pgfqpoint{1.713856in}{1.632990in}}%
\pgfpathcurveto{\pgfqpoint{1.709738in}{1.628872in}}{\pgfqpoint{1.707424in}{1.623286in}}{\pgfqpoint{1.707424in}{1.617462in}}%
\pgfpathcurveto{\pgfqpoint{1.707424in}{1.611638in}}{\pgfqpoint{1.709738in}{1.606052in}}{\pgfqpoint{1.713856in}{1.601934in}}%
\pgfpathcurveto{\pgfqpoint{1.717974in}{1.597816in}}{\pgfqpoint{1.723561in}{1.595502in}}{\pgfqpoint{1.729385in}{1.595502in}}%
\pgfpathlineto{\pgfqpoint{1.729385in}{1.595502in}}%
\pgfpathclose%
\pgfusepath{stroke,fill}%
\end{pgfscope}%
\begin{pgfscope}%
\pgfpathrectangle{\pgfqpoint{0.100000in}{0.209042in}}{\pgfqpoint{2.906250in}{2.906250in}}%
\pgfusepath{clip}%
\pgfsetbuttcap%
\pgfsetroundjoin%
\definecolor{currentfill}{rgb}{0.171713,0.067305,0.370771}%
\pgfsetfillcolor{currentfill}%
\pgfsetfillopacity{0.800000}%
\pgfsetlinewidth{1.003750pt}%
\definecolor{currentstroke}{rgb}{0.171713,0.067305,0.370771}%
\pgfsetstrokecolor{currentstroke}%
\pgfsetstrokeopacity{0.800000}%
\pgfsetdash{}{0pt}%
\pgfpathmoveto{\pgfqpoint{1.677744in}{1.873586in}}%
\pgfpathcurveto{\pgfqpoint{1.683568in}{1.873586in}}{\pgfqpoint{1.689154in}{1.875900in}}{\pgfqpoint{1.693272in}{1.880018in}}%
\pgfpathcurveto{\pgfqpoint{1.697390in}{1.884136in}}{\pgfqpoint{1.699704in}{1.889723in}}{\pgfqpoint{1.699704in}{1.895546in}}%
\pgfpathcurveto{\pgfqpoint{1.699704in}{1.901370in}}{\pgfqpoint{1.697390in}{1.906957in}}{\pgfqpoint{1.693272in}{1.911075in}}%
\pgfpathcurveto{\pgfqpoint{1.689154in}{1.915193in}}{\pgfqpoint{1.683568in}{1.917507in}}{\pgfqpoint{1.677744in}{1.917507in}}%
\pgfpathcurveto{\pgfqpoint{1.671920in}{1.917507in}}{\pgfqpoint{1.666334in}{1.915193in}}{\pgfqpoint{1.662216in}{1.911075in}}%
\pgfpathcurveto{\pgfqpoint{1.658098in}{1.906957in}}{\pgfqpoint{1.655784in}{1.901370in}}{\pgfqpoint{1.655784in}{1.895546in}}%
\pgfpathcurveto{\pgfqpoint{1.655784in}{1.889723in}}{\pgfqpoint{1.658098in}{1.884136in}}{\pgfqpoint{1.662216in}{1.880018in}}%
\pgfpathcurveto{\pgfqpoint{1.666334in}{1.875900in}}{\pgfqpoint{1.671920in}{1.873586in}}{\pgfqpoint{1.677744in}{1.873586in}}%
\pgfpathlineto{\pgfqpoint{1.677744in}{1.873586in}}%
\pgfpathclose%
\pgfusepath{stroke,fill}%
\end{pgfscope}%
\begin{pgfscope}%
\pgfpathrectangle{\pgfqpoint{0.100000in}{0.209042in}}{\pgfqpoint{2.906250in}{2.906250in}}%
\pgfusepath{clip}%
\pgfsetbuttcap%
\pgfsetroundjoin%
\definecolor{currentfill}{rgb}{0.146785,0.068738,0.334011}%
\pgfsetfillcolor{currentfill}%
\pgfsetfillopacity{0.800000}%
\pgfsetlinewidth{1.003750pt}%
\definecolor{currentstroke}{rgb}{0.146785,0.068738,0.334011}%
\pgfsetstrokecolor{currentstroke}%
\pgfsetstrokeopacity{0.800000}%
\pgfsetdash{}{0pt}%
\pgfpathmoveto{\pgfqpoint{1.712872in}{1.841927in}}%
\pgfpathcurveto{\pgfqpoint{1.718696in}{1.841927in}}{\pgfqpoint{1.724282in}{1.844241in}}{\pgfqpoint{1.728400in}{1.848359in}}%
\pgfpathcurveto{\pgfqpoint{1.732518in}{1.852477in}}{\pgfqpoint{1.734832in}{1.858063in}}{\pgfqpoint{1.734832in}{1.863887in}}%
\pgfpathcurveto{\pgfqpoint{1.734832in}{1.869711in}}{\pgfqpoint{1.732518in}{1.875297in}}{\pgfqpoint{1.728400in}{1.879415in}}%
\pgfpathcurveto{\pgfqpoint{1.724282in}{1.883533in}}{\pgfqpoint{1.718696in}{1.885847in}}{\pgfqpoint{1.712872in}{1.885847in}}%
\pgfpathcurveto{\pgfqpoint{1.707048in}{1.885847in}}{\pgfqpoint{1.701462in}{1.883533in}}{\pgfqpoint{1.697344in}{1.879415in}}%
\pgfpathcurveto{\pgfqpoint{1.693226in}{1.875297in}}{\pgfqpoint{1.690912in}{1.869711in}}{\pgfqpoint{1.690912in}{1.863887in}}%
\pgfpathcurveto{\pgfqpoint{1.690912in}{1.858063in}}{\pgfqpoint{1.693226in}{1.852477in}}{\pgfqpoint{1.697344in}{1.848359in}}%
\pgfpathcurveto{\pgfqpoint{1.701462in}{1.844241in}}{\pgfqpoint{1.707048in}{1.841927in}}{\pgfqpoint{1.712872in}{1.841927in}}%
\pgfpathlineto{\pgfqpoint{1.712872in}{1.841927in}}%
\pgfpathclose%
\pgfusepath{stroke,fill}%
\end{pgfscope}%
\begin{pgfscope}%
\pgfpathrectangle{\pgfqpoint{0.100000in}{0.209042in}}{\pgfqpoint{2.906250in}{2.906250in}}%
\pgfusepath{clip}%
\pgfsetbuttcap%
\pgfsetroundjoin%
\definecolor{currentfill}{rgb}{0.947180,0.384178,0.363701}%
\pgfsetfillcolor{currentfill}%
\pgfsetfillopacity{0.800000}%
\pgfsetlinewidth{1.003750pt}%
\definecolor{currentstroke}{rgb}{0.947180,0.384178,0.363701}%
\pgfsetstrokecolor{currentstroke}%
\pgfsetstrokeopacity{0.800000}%
\pgfsetdash{}{0pt}%
\pgfpathmoveto{\pgfqpoint{1.642197in}{1.114953in}}%
\pgfpathcurveto{\pgfqpoint{1.648021in}{1.114953in}}{\pgfqpoint{1.653607in}{1.117266in}}{\pgfqpoint{1.657726in}{1.121385in}}%
\pgfpathcurveto{\pgfqpoint{1.661844in}{1.125503in}}{\pgfqpoint{1.664158in}{1.131089in}}{\pgfqpoint{1.664158in}{1.136913in}}%
\pgfpathcurveto{\pgfqpoint{1.664158in}{1.142737in}}{\pgfqpoint{1.661844in}{1.148323in}}{\pgfqpoint{1.657726in}{1.152441in}}%
\pgfpathcurveto{\pgfqpoint{1.653607in}{1.156559in}}{\pgfqpoint{1.648021in}{1.158873in}}{\pgfqpoint{1.642197in}{1.158873in}}%
\pgfpathcurveto{\pgfqpoint{1.636373in}{1.158873in}}{\pgfqpoint{1.630787in}{1.156559in}}{\pgfqpoint{1.626669in}{1.152441in}}%
\pgfpathcurveto{\pgfqpoint{1.622551in}{1.148323in}}{\pgfqpoint{1.620237in}{1.142737in}}{\pgfqpoint{1.620237in}{1.136913in}}%
\pgfpathcurveto{\pgfqpoint{1.620237in}{1.131089in}}{\pgfqpoint{1.622551in}{1.125503in}}{\pgfqpoint{1.626669in}{1.121385in}}%
\pgfpathcurveto{\pgfqpoint{1.630787in}{1.117266in}}{\pgfqpoint{1.636373in}{1.114953in}}{\pgfqpoint{1.642197in}{1.114953in}}%
\pgfpathlineto{\pgfqpoint{1.642197in}{1.114953in}}%
\pgfpathclose%
\pgfusepath{stroke,fill}%
\end{pgfscope}%
\begin{pgfscope}%
\pgfpathrectangle{\pgfqpoint{0.100000in}{0.209042in}}{\pgfqpoint{2.906250in}{2.906250in}}%
\pgfusepath{clip}%
\pgfsetbuttcap%
\pgfsetroundjoin%
\definecolor{currentfill}{rgb}{0.211718,0.061992,0.418647}%
\pgfsetfillcolor{currentfill}%
\pgfsetfillopacity{0.800000}%
\pgfsetlinewidth{1.003750pt}%
\definecolor{currentstroke}{rgb}{0.211718,0.061992,0.418647}%
\pgfsetstrokecolor{currentstroke}%
\pgfsetstrokeopacity{0.800000}%
\pgfsetdash{}{0pt}%
\pgfpathmoveto{\pgfqpoint{1.419062in}{1.621120in}}%
\pgfpathcurveto{\pgfqpoint{1.424886in}{1.621120in}}{\pgfqpoint{1.430472in}{1.623434in}}{\pgfqpoint{1.434590in}{1.627552in}}%
\pgfpathcurveto{\pgfqpoint{1.438708in}{1.631670in}}{\pgfqpoint{1.441022in}{1.637256in}}{\pgfqpoint{1.441022in}{1.643080in}}%
\pgfpathcurveto{\pgfqpoint{1.441022in}{1.648904in}}{\pgfqpoint{1.438708in}{1.654490in}}{\pgfqpoint{1.434590in}{1.658608in}}%
\pgfpathcurveto{\pgfqpoint{1.430472in}{1.662726in}}{\pgfqpoint{1.424886in}{1.665040in}}{\pgfqpoint{1.419062in}{1.665040in}}%
\pgfpathcurveto{\pgfqpoint{1.413238in}{1.665040in}}{\pgfqpoint{1.407652in}{1.662726in}}{\pgfqpoint{1.403534in}{1.658608in}}%
\pgfpathcurveto{\pgfqpoint{1.399415in}{1.654490in}}{\pgfqpoint{1.397102in}{1.648904in}}{\pgfqpoint{1.397102in}{1.643080in}}%
\pgfpathcurveto{\pgfqpoint{1.397102in}{1.637256in}}{\pgfqpoint{1.399415in}{1.631670in}}{\pgfqpoint{1.403534in}{1.627552in}}%
\pgfpathcurveto{\pgfqpoint{1.407652in}{1.623434in}}{\pgfqpoint{1.413238in}{1.621120in}}{\pgfqpoint{1.419062in}{1.621120in}}%
\pgfpathlineto{\pgfqpoint{1.419062in}{1.621120in}}%
\pgfpathclose%
\pgfusepath{stroke,fill}%
\end{pgfscope}%
\begin{pgfscope}%
\pgfpathrectangle{\pgfqpoint{0.100000in}{0.209042in}}{\pgfqpoint{2.906250in}{2.906250in}}%
\pgfusepath{clip}%
\pgfsetbuttcap%
\pgfsetroundjoin%
\definecolor{currentfill}{rgb}{0.211718,0.061992,0.418647}%
\pgfsetfillcolor{currentfill}%
\pgfsetfillopacity{0.800000}%
\pgfsetlinewidth{1.003750pt}%
\definecolor{currentstroke}{rgb}{0.211718,0.061992,0.418647}%
\pgfsetstrokecolor{currentstroke}%
\pgfsetstrokeopacity{0.800000}%
\pgfsetdash{}{0pt}%
\pgfpathmoveto{\pgfqpoint{1.573136in}{1.905749in}}%
\pgfpathcurveto{\pgfqpoint{1.578960in}{1.905749in}}{\pgfqpoint{1.584546in}{1.908063in}}{\pgfqpoint{1.588665in}{1.912181in}}%
\pgfpathcurveto{\pgfqpoint{1.592783in}{1.916299in}}{\pgfqpoint{1.595097in}{1.921885in}}{\pgfqpoint{1.595097in}{1.927709in}}%
\pgfpathcurveto{\pgfqpoint{1.595097in}{1.933533in}}{\pgfqpoint{1.592783in}{1.939119in}}{\pgfqpoint{1.588665in}{1.943237in}}%
\pgfpathcurveto{\pgfqpoint{1.584546in}{1.947356in}}{\pgfqpoint{1.578960in}{1.949669in}}{\pgfqpoint{1.573136in}{1.949669in}}%
\pgfpathcurveto{\pgfqpoint{1.567312in}{1.949669in}}{\pgfqpoint{1.561726in}{1.947356in}}{\pgfqpoint{1.557608in}{1.943237in}}%
\pgfpathcurveto{\pgfqpoint{1.553490in}{1.939119in}}{\pgfqpoint{1.551176in}{1.933533in}}{\pgfqpoint{1.551176in}{1.927709in}}%
\pgfpathcurveto{\pgfqpoint{1.551176in}{1.921885in}}{\pgfqpoint{1.553490in}{1.916299in}}{\pgfqpoint{1.557608in}{1.912181in}}%
\pgfpathcurveto{\pgfqpoint{1.561726in}{1.908063in}}{\pgfqpoint{1.567312in}{1.905749in}}{\pgfqpoint{1.573136in}{1.905749in}}%
\pgfpathlineto{\pgfqpoint{1.573136in}{1.905749in}}%
\pgfpathclose%
\pgfusepath{stroke,fill}%
\end{pgfscope}%
\begin{pgfscope}%
\pgfpathrectangle{\pgfqpoint{0.100000in}{0.209042in}}{\pgfqpoint{2.906250in}{2.906250in}}%
\pgfusepath{clip}%
\pgfsetbuttcap%
\pgfsetroundjoin%
\definecolor{currentfill}{rgb}{0.211718,0.061992,0.418647}%
\pgfsetfillcolor{currentfill}%
\pgfsetfillopacity{0.800000}%
\pgfsetlinewidth{1.003750pt}%
\definecolor{currentstroke}{rgb}{0.211718,0.061992,0.418647}%
\pgfsetstrokecolor{currentstroke}%
\pgfsetstrokeopacity{0.800000}%
\pgfsetdash{}{0pt}%
\pgfpathmoveto{\pgfqpoint{1.684270in}{1.529372in}}%
\pgfpathcurveto{\pgfqpoint{1.690094in}{1.529372in}}{\pgfqpoint{1.695680in}{1.531686in}}{\pgfqpoint{1.699798in}{1.535804in}}%
\pgfpathcurveto{\pgfqpoint{1.703917in}{1.539923in}}{\pgfqpoint{1.706230in}{1.545509in}}{\pgfqpoint{1.706230in}{1.551333in}}%
\pgfpathcurveto{\pgfqpoint{1.706230in}{1.557157in}}{\pgfqpoint{1.703917in}{1.562743in}}{\pgfqpoint{1.699798in}{1.566861in}}%
\pgfpathcurveto{\pgfqpoint{1.695680in}{1.570979in}}{\pgfqpoint{1.690094in}{1.573293in}}{\pgfqpoint{1.684270in}{1.573293in}}%
\pgfpathcurveto{\pgfqpoint{1.678446in}{1.573293in}}{\pgfqpoint{1.672860in}{1.570979in}}{\pgfqpoint{1.668742in}{1.566861in}}%
\pgfpathcurveto{\pgfqpoint{1.664624in}{1.562743in}}{\pgfqpoint{1.662310in}{1.557157in}}{\pgfqpoint{1.662310in}{1.551333in}}%
\pgfpathcurveto{\pgfqpoint{1.662310in}{1.545509in}}{\pgfqpoint{1.664624in}{1.539923in}}{\pgfqpoint{1.668742in}{1.535804in}}%
\pgfpathcurveto{\pgfqpoint{1.672860in}{1.531686in}}{\pgfqpoint{1.678446in}{1.529372in}}{\pgfqpoint{1.684270in}{1.529372in}}%
\pgfpathlineto{\pgfqpoint{1.684270in}{1.529372in}}%
\pgfpathclose%
\pgfusepath{stroke,fill}%
\end{pgfscope}%
\begin{pgfscope}%
\pgfpathrectangle{\pgfqpoint{0.100000in}{0.209042in}}{\pgfqpoint{2.906250in}{2.906250in}}%
\pgfusepath{clip}%
\pgfsetbuttcap%
\pgfsetroundjoin%
\definecolor{currentfill}{rgb}{0.359898,0.087831,0.496778}%
\pgfsetfillcolor{currentfill}%
\pgfsetfillopacity{0.800000}%
\pgfsetlinewidth{1.003750pt}%
\definecolor{currentstroke}{rgb}{0.359898,0.087831,0.496778}%
\pgfsetstrokecolor{currentstroke}%
\pgfsetstrokeopacity{0.800000}%
\pgfsetdash{}{0pt}%
\pgfpathmoveto{\pgfqpoint{1.309416in}{1.630696in}}%
\pgfpathcurveto{\pgfqpoint{1.315239in}{1.630696in}}{\pgfqpoint{1.320826in}{1.633010in}}{\pgfqpoint{1.324944in}{1.637128in}}%
\pgfpathcurveto{\pgfqpoint{1.329062in}{1.641246in}}{\pgfqpoint{1.331376in}{1.646832in}}{\pgfqpoint{1.331376in}{1.652656in}}%
\pgfpathcurveto{\pgfqpoint{1.331376in}{1.658480in}}{\pgfqpoint{1.329062in}{1.664066in}}{\pgfqpoint{1.324944in}{1.668184in}}%
\pgfpathcurveto{\pgfqpoint{1.320826in}{1.672303in}}{\pgfqpoint{1.315239in}{1.674616in}}{\pgfqpoint{1.309416in}{1.674616in}}%
\pgfpathcurveto{\pgfqpoint{1.303592in}{1.674616in}}{\pgfqpoint{1.298005in}{1.672303in}}{\pgfqpoint{1.293887in}{1.668184in}}%
\pgfpathcurveto{\pgfqpoint{1.289769in}{1.664066in}}{\pgfqpoint{1.287455in}{1.658480in}}{\pgfqpoint{1.287455in}{1.652656in}}%
\pgfpathcurveto{\pgfqpoint{1.287455in}{1.646832in}}{\pgfqpoint{1.289769in}{1.641246in}}{\pgfqpoint{1.293887in}{1.637128in}}%
\pgfpathcurveto{\pgfqpoint{1.298005in}{1.633010in}}{\pgfqpoint{1.303592in}{1.630696in}}{\pgfqpoint{1.309416in}{1.630696in}}%
\pgfpathlineto{\pgfqpoint{1.309416in}{1.630696in}}%
\pgfpathclose%
\pgfusepath{stroke,fill}%
\end{pgfscope}%
\begin{pgfscope}%
\pgfpathrectangle{\pgfqpoint{0.100000in}{0.209042in}}{\pgfqpoint{2.906250in}{2.906250in}}%
\pgfusepath{clip}%
\pgfsetbuttcap%
\pgfsetroundjoin%
\definecolor{currentfill}{rgb}{0.031696,0.025765,0.116965}%
\pgfsetfillcolor{currentfill}%
\pgfsetfillopacity{0.800000}%
\pgfsetlinewidth{1.003750pt}%
\definecolor{currentstroke}{rgb}{0.031696,0.025765,0.116965}%
\pgfsetstrokecolor{currentstroke}%
\pgfsetstrokeopacity{0.800000}%
\pgfsetdash{}{0pt}%
\pgfpathmoveto{\pgfqpoint{1.688857in}{1.698162in}}%
\pgfpathcurveto{\pgfqpoint{1.694681in}{1.698162in}}{\pgfqpoint{1.700267in}{1.700476in}}{\pgfqpoint{1.704385in}{1.704594in}}%
\pgfpathcurveto{\pgfqpoint{1.708503in}{1.708712in}}{\pgfqpoint{1.710817in}{1.714298in}}{\pgfqpoint{1.710817in}{1.720122in}}%
\pgfpathcurveto{\pgfqpoint{1.710817in}{1.725946in}}{\pgfqpoint{1.708503in}{1.731532in}}{\pgfqpoint{1.704385in}{1.735650in}}%
\pgfpathcurveto{\pgfqpoint{1.700267in}{1.739768in}}{\pgfqpoint{1.694681in}{1.742082in}}{\pgfqpoint{1.688857in}{1.742082in}}%
\pgfpathcurveto{\pgfqpoint{1.683033in}{1.742082in}}{\pgfqpoint{1.677447in}{1.739768in}}{\pgfqpoint{1.673329in}{1.735650in}}%
\pgfpathcurveto{\pgfqpoint{1.669210in}{1.731532in}}{\pgfqpoint{1.666897in}{1.725946in}}{\pgfqpoint{1.666897in}{1.720122in}}%
\pgfpathcurveto{\pgfqpoint{1.666897in}{1.714298in}}{\pgfqpoint{1.669210in}{1.708712in}}{\pgfqpoint{1.673329in}{1.704594in}}%
\pgfpathcurveto{\pgfqpoint{1.677447in}{1.700476in}}{\pgfqpoint{1.683033in}{1.698162in}}{\pgfqpoint{1.688857in}{1.698162in}}%
\pgfpathlineto{\pgfqpoint{1.688857in}{1.698162in}}%
\pgfpathclose%
\pgfusepath{stroke,fill}%
\end{pgfscope}%
\begin{pgfscope}%
\pgfpathrectangle{\pgfqpoint{0.100000in}{0.209042in}}{\pgfqpoint{2.906250in}{2.906250in}}%
\pgfusepath{clip}%
\pgfsetbuttcap%
\pgfsetroundjoin%
\definecolor{currentfill}{rgb}{0.569172,0.167454,0.504105}%
\pgfsetfillcolor{currentfill}%
\pgfsetfillopacity{0.800000}%
\pgfsetlinewidth{1.003750pt}%
\definecolor{currentstroke}{rgb}{0.569172,0.167454,0.504105}%
\pgfsetstrokecolor{currentstroke}%
\pgfsetstrokeopacity{0.800000}%
\pgfsetdash{}{0pt}%
\pgfpathmoveto{\pgfqpoint{1.832093in}{1.366234in}}%
\pgfpathcurveto{\pgfqpoint{1.837917in}{1.366234in}}{\pgfqpoint{1.843504in}{1.368548in}}{\pgfqpoint{1.847622in}{1.372666in}}%
\pgfpathcurveto{\pgfqpoint{1.851740in}{1.376784in}}{\pgfqpoint{1.854054in}{1.382370in}}{\pgfqpoint{1.854054in}{1.388194in}}%
\pgfpathcurveto{\pgfqpoint{1.854054in}{1.394018in}}{\pgfqpoint{1.851740in}{1.399605in}}{\pgfqpoint{1.847622in}{1.403723in}}%
\pgfpathcurveto{\pgfqpoint{1.843504in}{1.407841in}}{\pgfqpoint{1.837917in}{1.410155in}}{\pgfqpoint{1.832093in}{1.410155in}}%
\pgfpathcurveto{\pgfqpoint{1.826269in}{1.410155in}}{\pgfqpoint{1.820683in}{1.407841in}}{\pgfqpoint{1.816565in}{1.403723in}}%
\pgfpathcurveto{\pgfqpoint{1.812447in}{1.399605in}}{\pgfqpoint{1.810133in}{1.394018in}}{\pgfqpoint{1.810133in}{1.388194in}}%
\pgfpathcurveto{\pgfqpoint{1.810133in}{1.382370in}}{\pgfqpoint{1.812447in}{1.376784in}}{\pgfqpoint{1.816565in}{1.372666in}}%
\pgfpathcurveto{\pgfqpoint{1.820683in}{1.368548in}}{\pgfqpoint{1.826269in}{1.366234in}}{\pgfqpoint{1.832093in}{1.366234in}}%
\pgfpathlineto{\pgfqpoint{1.832093in}{1.366234in}}%
\pgfpathclose%
\pgfusepath{stroke,fill}%
\end{pgfscope}%
\begin{pgfscope}%
\pgfpathrectangle{\pgfqpoint{0.100000in}{0.209042in}}{\pgfqpoint{2.906250in}{2.906250in}}%
\pgfusepath{clip}%
\pgfsetbuttcap%
\pgfsetroundjoin%
\definecolor{currentfill}{rgb}{0.703532,0.210638,0.479049}%
\pgfsetfillcolor{currentfill}%
\pgfsetfillopacity{0.800000}%
\pgfsetlinewidth{1.003750pt}%
\definecolor{currentstroke}{rgb}{0.703532,0.210638,0.479049}%
\pgfsetstrokecolor{currentstroke}%
\pgfsetstrokeopacity{0.800000}%
\pgfsetdash{}{0pt}%
\pgfpathmoveto{\pgfqpoint{1.076762in}{1.639678in}}%
\pgfpathcurveto{\pgfqpoint{1.082586in}{1.639678in}}{\pgfqpoint{1.088172in}{1.641992in}}{\pgfqpoint{1.092290in}{1.646110in}}%
\pgfpathcurveto{\pgfqpoint{1.096408in}{1.650228in}}{\pgfqpoint{1.098722in}{1.655815in}}{\pgfqpoint{1.098722in}{1.661639in}}%
\pgfpathcurveto{\pgfqpoint{1.098722in}{1.667463in}}{\pgfqpoint{1.096408in}{1.673049in}}{\pgfqpoint{1.092290in}{1.677167in}}%
\pgfpathcurveto{\pgfqpoint{1.088172in}{1.681285in}}{\pgfqpoint{1.082586in}{1.683599in}}{\pgfqpoint{1.076762in}{1.683599in}}%
\pgfpathcurveto{\pgfqpoint{1.070938in}{1.683599in}}{\pgfqpoint{1.065351in}{1.681285in}}{\pgfqpoint{1.061233in}{1.677167in}}%
\pgfpathcurveto{\pgfqpoint{1.057115in}{1.673049in}}{\pgfqpoint{1.054801in}{1.667463in}}{\pgfqpoint{1.054801in}{1.661639in}}%
\pgfpathcurveto{\pgfqpoint{1.054801in}{1.655815in}}{\pgfqpoint{1.057115in}{1.650228in}}{\pgfqpoint{1.061233in}{1.646110in}}%
\pgfpathcurveto{\pgfqpoint{1.065351in}{1.641992in}}{\pgfqpoint{1.070938in}{1.639678in}}{\pgfqpoint{1.076762in}{1.639678in}}%
\pgfpathlineto{\pgfqpoint{1.076762in}{1.639678in}}%
\pgfpathclose%
\pgfusepath{stroke,fill}%
\end{pgfscope}%
\begin{pgfscope}%
\pgfpathrectangle{\pgfqpoint{0.100000in}{0.209042in}}{\pgfqpoint{2.906250in}{2.906250in}}%
\pgfusepath{clip}%
\pgfsetbuttcap%
\pgfsetroundjoin%
\definecolor{currentfill}{rgb}{0.043830,0.033830,0.141886}%
\pgfsetfillcolor{currentfill}%
\pgfsetfillopacity{0.800000}%
\pgfsetlinewidth{1.003750pt}%
\definecolor{currentstroke}{rgb}{0.043830,0.033830,0.141886}%
\pgfsetstrokecolor{currentstroke}%
\pgfsetstrokeopacity{0.800000}%
\pgfsetdash{}{0pt}%
\pgfpathmoveto{\pgfqpoint{1.672949in}{1.776169in}}%
\pgfpathcurveto{\pgfqpoint{1.678773in}{1.776169in}}{\pgfqpoint{1.684359in}{1.778483in}}{\pgfqpoint{1.688477in}{1.782601in}}%
\pgfpathcurveto{\pgfqpoint{1.692595in}{1.786719in}}{\pgfqpoint{1.694909in}{1.792306in}}{\pgfqpoint{1.694909in}{1.798130in}}%
\pgfpathcurveto{\pgfqpoint{1.694909in}{1.803954in}}{\pgfqpoint{1.692595in}{1.809540in}}{\pgfqpoint{1.688477in}{1.813658in}}%
\pgfpathcurveto{\pgfqpoint{1.684359in}{1.817776in}}{\pgfqpoint{1.678773in}{1.820090in}}{\pgfqpoint{1.672949in}{1.820090in}}%
\pgfpathcurveto{\pgfqpoint{1.667125in}{1.820090in}}{\pgfqpoint{1.661539in}{1.817776in}}{\pgfqpoint{1.657420in}{1.813658in}}%
\pgfpathcurveto{\pgfqpoint{1.653302in}{1.809540in}}{\pgfqpoint{1.650988in}{1.803954in}}{\pgfqpoint{1.650988in}{1.798130in}}%
\pgfpathcurveto{\pgfqpoint{1.650988in}{1.792306in}}{\pgfqpoint{1.653302in}{1.786719in}}{\pgfqpoint{1.657420in}{1.782601in}}%
\pgfpathcurveto{\pgfqpoint{1.661539in}{1.778483in}}{\pgfqpoint{1.667125in}{1.776169in}}{\pgfqpoint{1.672949in}{1.776169in}}%
\pgfpathlineto{\pgfqpoint{1.672949in}{1.776169in}}%
\pgfpathclose%
\pgfusepath{stroke,fill}%
\end{pgfscope}%
\begin{pgfscope}%
\pgfpathrectangle{\pgfqpoint{0.100000in}{0.209042in}}{\pgfqpoint{2.906250in}{2.906250in}}%
\pgfusepath{clip}%
\pgfsetbuttcap%
\pgfsetroundjoin%
\definecolor{currentfill}{rgb}{0.232077,0.059889,0.437695}%
\pgfsetfillcolor{currentfill}%
\pgfsetfillopacity{0.800000}%
\pgfsetlinewidth{1.003750pt}%
\definecolor{currentstroke}{rgb}{0.232077,0.059889,0.437695}%
\pgfsetstrokecolor{currentstroke}%
\pgfsetstrokeopacity{0.800000}%
\pgfsetdash{}{0pt}%
\pgfpathmoveto{\pgfqpoint{1.370311in}{1.737893in}}%
\pgfpathcurveto{\pgfqpoint{1.376135in}{1.737893in}}{\pgfqpoint{1.381721in}{1.740207in}}{\pgfqpoint{1.385839in}{1.744325in}}%
\pgfpathcurveto{\pgfqpoint{1.389957in}{1.748443in}}{\pgfqpoint{1.392271in}{1.754029in}}{\pgfqpoint{1.392271in}{1.759853in}}%
\pgfpathcurveto{\pgfqpoint{1.392271in}{1.765677in}}{\pgfqpoint{1.389957in}{1.771263in}}{\pgfqpoint{1.385839in}{1.775381in}}%
\pgfpathcurveto{\pgfqpoint{1.381721in}{1.779499in}}{\pgfqpoint{1.376135in}{1.781813in}}{\pgfqpoint{1.370311in}{1.781813in}}%
\pgfpathcurveto{\pgfqpoint{1.364487in}{1.781813in}}{\pgfqpoint{1.358901in}{1.779499in}}{\pgfqpoint{1.354783in}{1.775381in}}%
\pgfpathcurveto{\pgfqpoint{1.350665in}{1.771263in}}{\pgfqpoint{1.348351in}{1.765677in}}{\pgfqpoint{1.348351in}{1.759853in}}%
\pgfpathcurveto{\pgfqpoint{1.348351in}{1.754029in}}{\pgfqpoint{1.350665in}{1.748443in}}{\pgfqpoint{1.354783in}{1.744325in}}%
\pgfpathcurveto{\pgfqpoint{1.358901in}{1.740207in}}{\pgfqpoint{1.364487in}{1.737893in}}{\pgfqpoint{1.370311in}{1.737893in}}%
\pgfpathlineto{\pgfqpoint{1.370311in}{1.737893in}}%
\pgfpathclose%
\pgfusepath{stroke,fill}%
\end{pgfscope}%
\begin{pgfscope}%
\pgfpathrectangle{\pgfqpoint{0.100000in}{0.209042in}}{\pgfqpoint{2.906250in}{2.906250in}}%
\pgfusepath{clip}%
\pgfsetbuttcap%
\pgfsetroundjoin%
\definecolor{currentfill}{rgb}{0.347636,0.082946,0.494121}%
\pgfsetfillcolor{currentfill}%
\pgfsetfillopacity{0.800000}%
\pgfsetlinewidth{1.003750pt}%
\definecolor{currentstroke}{rgb}{0.347636,0.082946,0.494121}%
\pgfsetstrokecolor{currentstroke}%
\pgfsetstrokeopacity{0.800000}%
\pgfsetdash{}{0pt}%
\pgfpathmoveto{\pgfqpoint{1.728470in}{1.462537in}}%
\pgfpathcurveto{\pgfqpoint{1.734294in}{1.462537in}}{\pgfqpoint{1.739880in}{1.464851in}}{\pgfqpoint{1.743998in}{1.468969in}}%
\pgfpathcurveto{\pgfqpoint{1.748116in}{1.473087in}}{\pgfqpoint{1.750430in}{1.478673in}}{\pgfqpoint{1.750430in}{1.484497in}}%
\pgfpathcurveto{\pgfqpoint{1.750430in}{1.490321in}}{\pgfqpoint{1.748116in}{1.495907in}}{\pgfqpoint{1.743998in}{1.500025in}}%
\pgfpathcurveto{\pgfqpoint{1.739880in}{1.504143in}}{\pgfqpoint{1.734294in}{1.506457in}}{\pgfqpoint{1.728470in}{1.506457in}}%
\pgfpathcurveto{\pgfqpoint{1.722646in}{1.506457in}}{\pgfqpoint{1.717060in}{1.504143in}}{\pgfqpoint{1.712942in}{1.500025in}}%
\pgfpathcurveto{\pgfqpoint{1.708823in}{1.495907in}}{\pgfqpoint{1.706510in}{1.490321in}}{\pgfqpoint{1.706510in}{1.484497in}}%
\pgfpathcurveto{\pgfqpoint{1.706510in}{1.478673in}}{\pgfqpoint{1.708823in}{1.473087in}}{\pgfqpoint{1.712942in}{1.468969in}}%
\pgfpathcurveto{\pgfqpoint{1.717060in}{1.464851in}}{\pgfqpoint{1.722646in}{1.462537in}}{\pgfqpoint{1.728470in}{1.462537in}}%
\pgfpathlineto{\pgfqpoint{1.728470in}{1.462537in}}%
\pgfpathclose%
\pgfusepath{stroke,fill}%
\end{pgfscope}%
\begin{pgfscope}%
\pgfpathrectangle{\pgfqpoint{0.100000in}{0.209042in}}{\pgfqpoint{2.906250in}{2.906250in}}%
\pgfusepath{clip}%
\pgfsetbuttcap%
\pgfsetroundjoin%
\definecolor{currentfill}{rgb}{0.191460,0.064818,0.396152}%
\pgfsetfillcolor{currentfill}%
\pgfsetfillopacity{0.800000}%
\pgfsetlinewidth{1.003750pt}%
\definecolor{currentstroke}{rgb}{0.191460,0.064818,0.396152}%
\pgfsetstrokecolor{currentstroke}%
\pgfsetstrokeopacity{0.800000}%
\pgfsetdash{}{0pt}%
\pgfpathmoveto{\pgfqpoint{1.740157in}{1.559010in}}%
\pgfpathcurveto{\pgfqpoint{1.745981in}{1.559010in}}{\pgfqpoint{1.751567in}{1.561324in}}{\pgfqpoint{1.755685in}{1.565442in}}%
\pgfpathcurveto{\pgfqpoint{1.759803in}{1.569560in}}{\pgfqpoint{1.762117in}{1.575147in}}{\pgfqpoint{1.762117in}{1.580970in}}%
\pgfpathcurveto{\pgfqpoint{1.762117in}{1.586794in}}{\pgfqpoint{1.759803in}{1.592381in}}{\pgfqpoint{1.755685in}{1.596499in}}%
\pgfpathcurveto{\pgfqpoint{1.751567in}{1.600617in}}{\pgfqpoint{1.745981in}{1.602931in}}{\pgfqpoint{1.740157in}{1.602931in}}%
\pgfpathcurveto{\pgfqpoint{1.734333in}{1.602931in}}{\pgfqpoint{1.728747in}{1.600617in}}{\pgfqpoint{1.724629in}{1.596499in}}%
\pgfpathcurveto{\pgfqpoint{1.720511in}{1.592381in}}{\pgfqpoint{1.718197in}{1.586794in}}{\pgfqpoint{1.718197in}{1.580970in}}%
\pgfpathcurveto{\pgfqpoint{1.718197in}{1.575147in}}{\pgfqpoint{1.720511in}{1.569560in}}{\pgfqpoint{1.724629in}{1.565442in}}%
\pgfpathcurveto{\pgfqpoint{1.728747in}{1.561324in}}{\pgfqpoint{1.734333in}{1.559010in}}{\pgfqpoint{1.740157in}{1.559010in}}%
\pgfpathlineto{\pgfqpoint{1.740157in}{1.559010in}}%
\pgfpathclose%
\pgfusepath{stroke,fill}%
\end{pgfscope}%
\begin{pgfscope}%
\pgfpathrectangle{\pgfqpoint{0.100000in}{0.209042in}}{\pgfqpoint{2.906250in}{2.906250in}}%
\pgfusepath{clip}%
\pgfsetbuttcap%
\pgfsetroundjoin%
\definecolor{currentfill}{rgb}{0.232077,0.059889,0.437695}%
\pgfsetfillcolor{currentfill}%
\pgfsetfillopacity{0.800000}%
\pgfsetlinewidth{1.003750pt}%
\definecolor{currentstroke}{rgb}{0.232077,0.059889,0.437695}%
\pgfsetstrokecolor{currentstroke}%
\pgfsetstrokeopacity{0.800000}%
\pgfsetdash{}{0pt}%
\pgfpathmoveto{\pgfqpoint{1.763813in}{1.544646in}}%
\pgfpathcurveto{\pgfqpoint{1.769637in}{1.544646in}}{\pgfqpoint{1.775223in}{1.546960in}}{\pgfqpoint{1.779341in}{1.551078in}}%
\pgfpathcurveto{\pgfqpoint{1.783459in}{1.555196in}}{\pgfqpoint{1.785773in}{1.560782in}}{\pgfqpoint{1.785773in}{1.566606in}}%
\pgfpathcurveto{\pgfqpoint{1.785773in}{1.572430in}}{\pgfqpoint{1.783459in}{1.578016in}}{\pgfqpoint{1.779341in}{1.582134in}}%
\pgfpathcurveto{\pgfqpoint{1.775223in}{1.586252in}}{\pgfqpoint{1.769637in}{1.588566in}}{\pgfqpoint{1.763813in}{1.588566in}}%
\pgfpathcurveto{\pgfqpoint{1.757989in}{1.588566in}}{\pgfqpoint{1.752403in}{1.586252in}}{\pgfqpoint{1.748285in}{1.582134in}}%
\pgfpathcurveto{\pgfqpoint{1.744167in}{1.578016in}}{\pgfqpoint{1.741853in}{1.572430in}}{\pgfqpoint{1.741853in}{1.566606in}}%
\pgfpathcurveto{\pgfqpoint{1.741853in}{1.560782in}}{\pgfqpoint{1.744167in}{1.555196in}}{\pgfqpoint{1.748285in}{1.551078in}}%
\pgfpathcurveto{\pgfqpoint{1.752403in}{1.546960in}}{\pgfqpoint{1.757989in}{1.544646in}}{\pgfqpoint{1.763813in}{1.544646in}}%
\pgfpathlineto{\pgfqpoint{1.763813in}{1.544646in}}%
\pgfpathclose%
\pgfusepath{stroke,fill}%
\end{pgfscope}%
\begin{pgfscope}%
\pgfpathrectangle{\pgfqpoint{0.100000in}{0.209042in}}{\pgfqpoint{2.906250in}{2.906250in}}%
\pgfusepath{clip}%
\pgfsetbuttcap%
\pgfsetroundjoin%
\definecolor{currentfill}{rgb}{0.265447,0.060237,0.461840}%
\pgfsetfillcolor{currentfill}%
\pgfsetfillopacity{0.800000}%
\pgfsetlinewidth{1.003750pt}%
\definecolor{currentstroke}{rgb}{0.265447,0.060237,0.461840}%
\pgfsetstrokecolor{currentstroke}%
\pgfsetstrokeopacity{0.800000}%
\pgfsetdash{}{0pt}%
\pgfpathmoveto{\pgfqpoint{1.533002in}{1.925436in}}%
\pgfpathcurveto{\pgfqpoint{1.538826in}{1.925436in}}{\pgfqpoint{1.544412in}{1.927750in}}{\pgfqpoint{1.548531in}{1.931868in}}%
\pgfpathcurveto{\pgfqpoint{1.552649in}{1.935986in}}{\pgfqpoint{1.554963in}{1.941572in}}{\pgfqpoint{1.554963in}{1.947396in}}%
\pgfpathcurveto{\pgfqpoint{1.554963in}{1.953220in}}{\pgfqpoint{1.552649in}{1.958806in}}{\pgfqpoint{1.548531in}{1.962924in}}%
\pgfpathcurveto{\pgfqpoint{1.544412in}{1.967042in}}{\pgfqpoint{1.538826in}{1.969356in}}{\pgfqpoint{1.533002in}{1.969356in}}%
\pgfpathcurveto{\pgfqpoint{1.527178in}{1.969356in}}{\pgfqpoint{1.521592in}{1.967042in}}{\pgfqpoint{1.517474in}{1.962924in}}%
\pgfpathcurveto{\pgfqpoint{1.513356in}{1.958806in}}{\pgfqpoint{1.511042in}{1.953220in}}{\pgfqpoint{1.511042in}{1.947396in}}%
\pgfpathcurveto{\pgfqpoint{1.511042in}{1.941572in}}{\pgfqpoint{1.513356in}{1.935986in}}{\pgfqpoint{1.517474in}{1.931868in}}%
\pgfpathcurveto{\pgfqpoint{1.521592in}{1.927750in}}{\pgfqpoint{1.527178in}{1.925436in}}{\pgfqpoint{1.533002in}{1.925436in}}%
\pgfpathlineto{\pgfqpoint{1.533002in}{1.925436in}}%
\pgfpathclose%
\pgfusepath{stroke,fill}%
\end{pgfscope}%
\begin{pgfscope}%
\pgfpathrectangle{\pgfqpoint{0.100000in}{0.209042in}}{\pgfqpoint{2.906250in}{2.906250in}}%
\pgfusepath{clip}%
\pgfsetbuttcap%
\pgfsetroundjoin%
\definecolor{currentfill}{rgb}{0.322899,0.073782,0.487408}%
\pgfsetfillcolor{currentfill}%
\pgfsetfillopacity{0.800000}%
\pgfsetlinewidth{1.003750pt}%
\definecolor{currentstroke}{rgb}{0.322899,0.073782,0.487408}%
\pgfsetstrokecolor{currentstroke}%
\pgfsetstrokeopacity{0.800000}%
\pgfsetdash{}{0pt}%
\pgfpathmoveto{\pgfqpoint{1.784440in}{1.917646in}}%
\pgfpathcurveto{\pgfqpoint{1.790264in}{1.917646in}}{\pgfqpoint{1.795850in}{1.919960in}}{\pgfqpoint{1.799968in}{1.924078in}}%
\pgfpathcurveto{\pgfqpoint{1.804086in}{1.928196in}}{\pgfqpoint{1.806400in}{1.933782in}}{\pgfqpoint{1.806400in}{1.939606in}}%
\pgfpathcurveto{\pgfqpoint{1.806400in}{1.945430in}}{\pgfqpoint{1.804086in}{1.951016in}}{\pgfqpoint{1.799968in}{1.955134in}}%
\pgfpathcurveto{\pgfqpoint{1.795850in}{1.959252in}}{\pgfqpoint{1.790264in}{1.961566in}}{\pgfqpoint{1.784440in}{1.961566in}}%
\pgfpathcurveto{\pgfqpoint{1.778616in}{1.961566in}}{\pgfqpoint{1.773030in}{1.959252in}}{\pgfqpoint{1.768912in}{1.955134in}}%
\pgfpathcurveto{\pgfqpoint{1.764793in}{1.951016in}}{\pgfqpoint{1.762480in}{1.945430in}}{\pgfqpoint{1.762480in}{1.939606in}}%
\pgfpathcurveto{\pgfqpoint{1.762480in}{1.933782in}}{\pgfqpoint{1.764793in}{1.928196in}}{\pgfqpoint{1.768912in}{1.924078in}}%
\pgfpathcurveto{\pgfqpoint{1.773030in}{1.919960in}}{\pgfqpoint{1.778616in}{1.917646in}}{\pgfqpoint{1.784440in}{1.917646in}}%
\pgfpathlineto{\pgfqpoint{1.784440in}{1.917646in}}%
\pgfpathclose%
\pgfusepath{stroke,fill}%
\end{pgfscope}%
\begin{pgfscope}%
\pgfpathrectangle{\pgfqpoint{0.100000in}{0.209042in}}{\pgfqpoint{2.906250in}{2.906250in}}%
\pgfusepath{clip}%
\pgfsetbuttcap%
\pgfsetroundjoin%
\definecolor{currentfill}{rgb}{0.488088,0.139186,0.508011}%
\pgfsetfillcolor{currentfill}%
\pgfsetfillopacity{0.800000}%
\pgfsetlinewidth{1.003750pt}%
\definecolor{currentstroke}{rgb}{0.488088,0.139186,0.508011}%
\pgfsetstrokecolor{currentstroke}%
\pgfsetstrokeopacity{0.800000}%
\pgfsetdash{}{0pt}%
\pgfpathmoveto{\pgfqpoint{1.995486in}{1.840809in}}%
\pgfpathcurveto{\pgfqpoint{2.001310in}{1.840809in}}{\pgfqpoint{2.006897in}{1.843122in}}{\pgfqpoint{2.011015in}{1.847241in}}%
\pgfpathcurveto{\pgfqpoint{2.015133in}{1.851359in}}{\pgfqpoint{2.017447in}{1.856945in}}{\pgfqpoint{2.017447in}{1.862769in}}%
\pgfpathcurveto{\pgfqpoint{2.017447in}{1.868593in}}{\pgfqpoint{2.015133in}{1.874179in}}{\pgfqpoint{2.011015in}{1.878297in}}%
\pgfpathcurveto{\pgfqpoint{2.006897in}{1.882415in}}{\pgfqpoint{2.001310in}{1.884729in}}{\pgfqpoint{1.995486in}{1.884729in}}%
\pgfpathcurveto{\pgfqpoint{1.989663in}{1.884729in}}{\pgfqpoint{1.984076in}{1.882415in}}{\pgfqpoint{1.979958in}{1.878297in}}%
\pgfpathcurveto{\pgfqpoint{1.975840in}{1.874179in}}{\pgfqpoint{1.973526in}{1.868593in}}{\pgfqpoint{1.973526in}{1.862769in}}%
\pgfpathcurveto{\pgfqpoint{1.973526in}{1.856945in}}{\pgfqpoint{1.975840in}{1.851359in}}{\pgfqpoint{1.979958in}{1.847241in}}%
\pgfpathcurveto{\pgfqpoint{1.984076in}{1.843122in}}{\pgfqpoint{1.989663in}{1.840809in}}{\pgfqpoint{1.995486in}{1.840809in}}%
\pgfpathlineto{\pgfqpoint{1.995486in}{1.840809in}}%
\pgfpathclose%
\pgfusepath{stroke,fill}%
\end{pgfscope}%
\begin{pgfscope}%
\pgfpathrectangle{\pgfqpoint{0.100000in}{0.209042in}}{\pgfqpoint{2.906250in}{2.906250in}}%
\pgfusepath{clip}%
\pgfsetbuttcap%
\pgfsetroundjoin%
\definecolor{currentfill}{rgb}{0.420791,0.112920,0.505215}%
\pgfsetfillcolor{currentfill}%
\pgfsetfillopacity{0.800000}%
\pgfsetlinewidth{1.003750pt}%
\definecolor{currentstroke}{rgb}{0.420791,0.112920,0.505215}%
\pgfsetstrokecolor{currentstroke}%
\pgfsetstrokeopacity{0.800000}%
\pgfsetdash{}{0pt}%
\pgfpathmoveto{\pgfqpoint{1.518989in}{1.426634in}}%
\pgfpathcurveto{\pgfqpoint{1.524813in}{1.426634in}}{\pgfqpoint{1.530399in}{1.428948in}}{\pgfqpoint{1.534517in}{1.433066in}}%
\pgfpathcurveto{\pgfqpoint{1.538635in}{1.437185in}}{\pgfqpoint{1.540949in}{1.442771in}}{\pgfqpoint{1.540949in}{1.448595in}}%
\pgfpathcurveto{\pgfqpoint{1.540949in}{1.454419in}}{\pgfqpoint{1.538635in}{1.460005in}}{\pgfqpoint{1.534517in}{1.464123in}}%
\pgfpathcurveto{\pgfqpoint{1.530399in}{1.468241in}}{\pgfqpoint{1.524813in}{1.470555in}}{\pgfqpoint{1.518989in}{1.470555in}}%
\pgfpathcurveto{\pgfqpoint{1.513165in}{1.470555in}}{\pgfqpoint{1.507579in}{1.468241in}}{\pgfqpoint{1.503461in}{1.464123in}}%
\pgfpathcurveto{\pgfqpoint{1.499343in}{1.460005in}}{\pgfqpoint{1.497029in}{1.454419in}}{\pgfqpoint{1.497029in}{1.448595in}}%
\pgfpathcurveto{\pgfqpoint{1.497029in}{1.442771in}}{\pgfqpoint{1.499343in}{1.437185in}}{\pgfqpoint{1.503461in}{1.433066in}}%
\pgfpathcurveto{\pgfqpoint{1.507579in}{1.428948in}}{\pgfqpoint{1.513165in}{1.426634in}}{\pgfqpoint{1.518989in}{1.426634in}}%
\pgfpathlineto{\pgfqpoint{1.518989in}{1.426634in}}%
\pgfpathclose%
\pgfusepath{stroke,fill}%
\end{pgfscope}%
\begin{pgfscope}%
\pgfpathrectangle{\pgfqpoint{0.100000in}{0.209042in}}{\pgfqpoint{2.906250in}{2.906250in}}%
\pgfusepath{clip}%
\pgfsetbuttcap%
\pgfsetroundjoin%
\definecolor{currentfill}{rgb}{0.531507,0.154739,0.506895}%
\pgfsetfillcolor{currentfill}%
\pgfsetfillopacity{0.800000}%
\pgfsetlinewidth{1.003750pt}%
\definecolor{currentstroke}{rgb}{0.531507,0.154739,0.506895}%
\pgfsetstrokecolor{currentstroke}%
\pgfsetstrokeopacity{0.800000}%
\pgfsetdash{}{0pt}%
\pgfpathmoveto{\pgfqpoint{1.453863in}{1.382994in}}%
\pgfpathcurveto{\pgfqpoint{1.459687in}{1.382994in}}{\pgfqpoint{1.465273in}{1.385308in}}{\pgfqpoint{1.469391in}{1.389426in}}%
\pgfpathcurveto{\pgfqpoint{1.473509in}{1.393545in}}{\pgfqpoint{1.475823in}{1.399131in}}{\pgfqpoint{1.475823in}{1.404955in}}%
\pgfpathcurveto{\pgfqpoint{1.475823in}{1.410779in}}{\pgfqpoint{1.473509in}{1.416365in}}{\pgfqpoint{1.469391in}{1.420483in}}%
\pgfpathcurveto{\pgfqpoint{1.465273in}{1.424601in}}{\pgfqpoint{1.459687in}{1.426915in}}{\pgfqpoint{1.453863in}{1.426915in}}%
\pgfpathcurveto{\pgfqpoint{1.448039in}{1.426915in}}{\pgfqpoint{1.442453in}{1.424601in}}{\pgfqpoint{1.438335in}{1.420483in}}%
\pgfpathcurveto{\pgfqpoint{1.434216in}{1.416365in}}{\pgfqpoint{1.431903in}{1.410779in}}{\pgfqpoint{1.431903in}{1.404955in}}%
\pgfpathcurveto{\pgfqpoint{1.431903in}{1.399131in}}{\pgfqpoint{1.434216in}{1.393545in}}{\pgfqpoint{1.438335in}{1.389426in}}%
\pgfpathcurveto{\pgfqpoint{1.442453in}{1.385308in}}{\pgfqpoint{1.448039in}{1.382994in}}{\pgfqpoint{1.453863in}{1.382994in}}%
\pgfpathlineto{\pgfqpoint{1.453863in}{1.382994in}}%
\pgfpathclose%
\pgfusepath{stroke,fill}%
\end{pgfscope}%
\begin{pgfscope}%
\pgfpathrectangle{\pgfqpoint{0.100000in}{0.209042in}}{\pgfqpoint{2.906250in}{2.906250in}}%
\pgfusepath{clip}%
\pgfsetbuttcap%
\pgfsetroundjoin%
\definecolor{currentfill}{rgb}{0.773695,0.236249,0.455289}%
\pgfsetfillcolor{currentfill}%
\pgfsetfillopacity{0.800000}%
\pgfsetlinewidth{1.003750pt}%
\definecolor{currentstroke}{rgb}{0.773695,0.236249,0.455289}%
\pgfsetstrokecolor{currentstroke}%
\pgfsetstrokeopacity{0.800000}%
\pgfsetdash{}{0pt}%
\pgfpathmoveto{\pgfqpoint{1.098546in}{1.986065in}}%
\pgfpathcurveto{\pgfqpoint{1.104370in}{1.986065in}}{\pgfqpoint{1.109957in}{1.988379in}}{\pgfqpoint{1.114075in}{1.992497in}}%
\pgfpathcurveto{\pgfqpoint{1.118193in}{1.996615in}}{\pgfqpoint{1.120507in}{2.002201in}}{\pgfqpoint{1.120507in}{2.008025in}}%
\pgfpathcurveto{\pgfqpoint{1.120507in}{2.013849in}}{\pgfqpoint{1.118193in}{2.019435in}}{\pgfqpoint{1.114075in}{2.023553in}}%
\pgfpathcurveto{\pgfqpoint{1.109957in}{2.027672in}}{\pgfqpoint{1.104370in}{2.029985in}}{\pgfqpoint{1.098546in}{2.029985in}}%
\pgfpathcurveto{\pgfqpoint{1.092723in}{2.029985in}}{\pgfqpoint{1.087136in}{2.027672in}}{\pgfqpoint{1.083018in}{2.023553in}}%
\pgfpathcurveto{\pgfqpoint{1.078900in}{2.019435in}}{\pgfqpoint{1.076586in}{2.013849in}}{\pgfqpoint{1.076586in}{2.008025in}}%
\pgfpathcurveto{\pgfqpoint{1.076586in}{2.002201in}}{\pgfqpoint{1.078900in}{1.996615in}}{\pgfqpoint{1.083018in}{1.992497in}}%
\pgfpathcurveto{\pgfqpoint{1.087136in}{1.988379in}}{\pgfqpoint{1.092723in}{1.986065in}}{\pgfqpoint{1.098546in}{1.986065in}}%
\pgfpathlineto{\pgfqpoint{1.098546in}{1.986065in}}%
\pgfpathclose%
\pgfusepath{stroke,fill}%
\end{pgfscope}%
\begin{pgfscope}%
\pgfpathrectangle{\pgfqpoint{0.100000in}{0.209042in}}{\pgfqpoint{2.906250in}{2.906250in}}%
\pgfusepath{clip}%
\pgfsetbuttcap%
\pgfsetroundjoin%
\definecolor{currentfill}{rgb}{0.420791,0.112920,0.505215}%
\pgfsetfillcolor{currentfill}%
\pgfsetfillopacity{0.800000}%
\pgfsetlinewidth{1.003750pt}%
\definecolor{currentstroke}{rgb}{0.420791,0.112920,0.505215}%
\pgfsetstrokecolor{currentstroke}%
\pgfsetstrokeopacity{0.800000}%
\pgfsetdash{}{0pt}%
\pgfpathmoveto{\pgfqpoint{1.587640in}{1.411148in}}%
\pgfpathcurveto{\pgfqpoint{1.593464in}{1.411148in}}{\pgfqpoint{1.599050in}{1.413461in}}{\pgfqpoint{1.603168in}{1.417580in}}%
\pgfpathcurveto{\pgfqpoint{1.607286in}{1.421698in}}{\pgfqpoint{1.609600in}{1.427284in}}{\pgfqpoint{1.609600in}{1.433108in}}%
\pgfpathcurveto{\pgfqpoint{1.609600in}{1.438932in}}{\pgfqpoint{1.607286in}{1.444518in}}{\pgfqpoint{1.603168in}{1.448636in}}%
\pgfpathcurveto{\pgfqpoint{1.599050in}{1.452754in}}{\pgfqpoint{1.593464in}{1.455068in}}{\pgfqpoint{1.587640in}{1.455068in}}%
\pgfpathcurveto{\pgfqpoint{1.581816in}{1.455068in}}{\pgfqpoint{1.576230in}{1.452754in}}{\pgfqpoint{1.572112in}{1.448636in}}%
\pgfpathcurveto{\pgfqpoint{1.567994in}{1.444518in}}{\pgfqpoint{1.565680in}{1.438932in}}{\pgfqpoint{1.565680in}{1.433108in}}%
\pgfpathcurveto{\pgfqpoint{1.565680in}{1.427284in}}{\pgfqpoint{1.567994in}{1.421698in}}{\pgfqpoint{1.572112in}{1.417580in}}%
\pgfpathcurveto{\pgfqpoint{1.576230in}{1.413461in}}{\pgfqpoint{1.581816in}{1.411148in}}{\pgfqpoint{1.587640in}{1.411148in}}%
\pgfpathlineto{\pgfqpoint{1.587640in}{1.411148in}}%
\pgfpathclose%
\pgfusepath{stroke,fill}%
\end{pgfscope}%
\begin{pgfscope}%
\pgfpathrectangle{\pgfqpoint{0.100000in}{0.209042in}}{\pgfqpoint{2.906250in}{2.906250in}}%
\pgfusepath{clip}%
\pgfsetbuttcap%
\pgfsetroundjoin%
\definecolor{currentfill}{rgb}{0.353773,0.085373,0.495501}%
\pgfsetfillcolor{currentfill}%
\pgfsetfillopacity{0.800000}%
\pgfsetlinewidth{1.003750pt}%
\definecolor{currentstroke}{rgb}{0.353773,0.085373,0.495501}%
\pgfsetstrokecolor{currentstroke}%
\pgfsetstrokeopacity{0.800000}%
\pgfsetdash{}{0pt}%
\pgfpathmoveto{\pgfqpoint{1.288825in}{1.754160in}}%
\pgfpathcurveto{\pgfqpoint{1.294649in}{1.754160in}}{\pgfqpoint{1.300235in}{1.756474in}}{\pgfqpoint{1.304353in}{1.760592in}}%
\pgfpathcurveto{\pgfqpoint{1.308471in}{1.764710in}}{\pgfqpoint{1.310785in}{1.770296in}}{\pgfqpoint{1.310785in}{1.776120in}}%
\pgfpathcurveto{\pgfqpoint{1.310785in}{1.781944in}}{\pgfqpoint{1.308471in}{1.787530in}}{\pgfqpoint{1.304353in}{1.791648in}}%
\pgfpathcurveto{\pgfqpoint{1.300235in}{1.795766in}}{\pgfqpoint{1.294649in}{1.798080in}}{\pgfqpoint{1.288825in}{1.798080in}}%
\pgfpathcurveto{\pgfqpoint{1.283001in}{1.798080in}}{\pgfqpoint{1.277415in}{1.795766in}}{\pgfqpoint{1.273296in}{1.791648in}}%
\pgfpathcurveto{\pgfqpoint{1.269178in}{1.787530in}}{\pgfqpoint{1.266864in}{1.781944in}}{\pgfqpoint{1.266864in}{1.776120in}}%
\pgfpathcurveto{\pgfqpoint{1.266864in}{1.770296in}}{\pgfqpoint{1.269178in}{1.764710in}}{\pgfqpoint{1.273296in}{1.760592in}}%
\pgfpathcurveto{\pgfqpoint{1.277415in}{1.756474in}}{\pgfqpoint{1.283001in}{1.754160in}}{\pgfqpoint{1.288825in}{1.754160in}}%
\pgfpathlineto{\pgfqpoint{1.288825in}{1.754160in}}%
\pgfpathclose%
\pgfusepath{stroke,fill}%
\end{pgfscope}%
\begin{pgfscope}%
\pgfpathrectangle{\pgfqpoint{0.100000in}{0.209042in}}{\pgfqpoint{2.906250in}{2.906250in}}%
\pgfusepath{clip}%
\pgfsetbuttcap%
\pgfsetroundjoin%
\definecolor{currentfill}{rgb}{0.531507,0.154739,0.506895}%
\pgfsetfillcolor{currentfill}%
\pgfsetfillopacity{0.800000}%
\pgfsetlinewidth{1.003750pt}%
\definecolor{currentstroke}{rgb}{0.531507,0.154739,0.506895}%
\pgfsetstrokecolor{currentstroke}%
\pgfsetstrokeopacity{0.800000}%
\pgfsetdash{}{0pt}%
\pgfpathmoveto{\pgfqpoint{1.931775in}{1.958371in}}%
\pgfpathcurveto{\pgfqpoint{1.937599in}{1.958371in}}{\pgfqpoint{1.943185in}{1.960685in}}{\pgfqpoint{1.947303in}{1.964803in}}%
\pgfpathcurveto{\pgfqpoint{1.951421in}{1.968921in}}{\pgfqpoint{1.953735in}{1.974507in}}{\pgfqpoint{1.953735in}{1.980331in}}%
\pgfpathcurveto{\pgfqpoint{1.953735in}{1.986155in}}{\pgfqpoint{1.951421in}{1.991741in}}{\pgfqpoint{1.947303in}{1.995860in}}%
\pgfpathcurveto{\pgfqpoint{1.943185in}{1.999978in}}{\pgfqpoint{1.937599in}{2.002292in}}{\pgfqpoint{1.931775in}{2.002292in}}%
\pgfpathcurveto{\pgfqpoint{1.925951in}{2.002292in}}{\pgfqpoint{1.920365in}{1.999978in}}{\pgfqpoint{1.916247in}{1.995860in}}%
\pgfpathcurveto{\pgfqpoint{1.912128in}{1.991741in}}{\pgfqpoint{1.909815in}{1.986155in}}{\pgfqpoint{1.909815in}{1.980331in}}%
\pgfpathcurveto{\pgfqpoint{1.909815in}{1.974507in}}{\pgfqpoint{1.912128in}{1.968921in}}{\pgfqpoint{1.916247in}{1.964803in}}%
\pgfpathcurveto{\pgfqpoint{1.920365in}{1.960685in}}{\pgfqpoint{1.925951in}{1.958371in}}{\pgfqpoint{1.931775in}{1.958371in}}%
\pgfpathlineto{\pgfqpoint{1.931775in}{1.958371in}}%
\pgfpathclose%
\pgfusepath{stroke,fill}%
\end{pgfscope}%
\begin{pgfscope}%
\pgfpathrectangle{\pgfqpoint{0.100000in}{0.209042in}}{\pgfqpoint{2.906250in}{2.906250in}}%
\pgfusepath{clip}%
\pgfsetbuttcap%
\pgfsetroundjoin%
\definecolor{currentfill}{rgb}{0.083446,0.056225,0.220755}%
\pgfsetfillcolor{currentfill}%
\pgfsetfillopacity{0.800000}%
\pgfsetlinewidth{1.003750pt}%
\definecolor{currentstroke}{rgb}{0.083446,0.056225,0.220755}%
\pgfsetstrokecolor{currentstroke}%
\pgfsetstrokeopacity{0.800000}%
\pgfsetdash{}{0pt}%
\pgfpathmoveto{\pgfqpoint{1.702267in}{1.611961in}}%
\pgfpathcurveto{\pgfqpoint{1.708091in}{1.611961in}}{\pgfqpoint{1.713677in}{1.614275in}}{\pgfqpoint{1.717795in}{1.618393in}}%
\pgfpathcurveto{\pgfqpoint{1.721913in}{1.622511in}}{\pgfqpoint{1.724227in}{1.628098in}}{\pgfqpoint{1.724227in}{1.633921in}}%
\pgfpathcurveto{\pgfqpoint{1.724227in}{1.639745in}}{\pgfqpoint{1.721913in}{1.645332in}}{\pgfqpoint{1.717795in}{1.649450in}}%
\pgfpathcurveto{\pgfqpoint{1.713677in}{1.653568in}}{\pgfqpoint{1.708091in}{1.655882in}}{\pgfqpoint{1.702267in}{1.655882in}}%
\pgfpathcurveto{\pgfqpoint{1.696443in}{1.655882in}}{\pgfqpoint{1.690857in}{1.653568in}}{\pgfqpoint{1.686739in}{1.649450in}}%
\pgfpathcurveto{\pgfqpoint{1.682621in}{1.645332in}}{\pgfqpoint{1.680307in}{1.639745in}}{\pgfqpoint{1.680307in}{1.633921in}}%
\pgfpathcurveto{\pgfqpoint{1.680307in}{1.628098in}}{\pgfqpoint{1.682621in}{1.622511in}}{\pgfqpoint{1.686739in}{1.618393in}}%
\pgfpathcurveto{\pgfqpoint{1.690857in}{1.614275in}}{\pgfqpoint{1.696443in}{1.611961in}}{\pgfqpoint{1.702267in}{1.611961in}}%
\pgfpathlineto{\pgfqpoint{1.702267in}{1.611961in}}%
\pgfpathclose%
\pgfusepath{stroke,fill}%
\end{pgfscope}%
\begin{pgfscope}%
\pgfpathrectangle{\pgfqpoint{0.100000in}{0.209042in}}{\pgfqpoint{2.906250in}{2.906250in}}%
\pgfusepath{clip}%
\pgfsetbuttcap%
\pgfsetroundjoin%
\definecolor{currentfill}{rgb}{0.102815,0.063010,0.257854}%
\pgfsetfillcolor{currentfill}%
\pgfsetfillopacity{0.800000}%
\pgfsetlinewidth{1.003750pt}%
\definecolor{currentstroke}{rgb}{0.102815,0.063010,0.257854}%
\pgfsetstrokecolor{currentstroke}%
\pgfsetstrokeopacity{0.800000}%
\pgfsetdash{}{0pt}%
\pgfpathmoveto{\pgfqpoint{1.547926in}{1.835933in}}%
\pgfpathcurveto{\pgfqpoint{1.553750in}{1.835933in}}{\pgfqpoint{1.559337in}{1.838247in}}{\pgfqpoint{1.563455in}{1.842365in}}%
\pgfpathcurveto{\pgfqpoint{1.567573in}{1.846484in}}{\pgfqpoint{1.569887in}{1.852070in}}{\pgfqpoint{1.569887in}{1.857894in}}%
\pgfpathcurveto{\pgfqpoint{1.569887in}{1.863718in}}{\pgfqpoint{1.567573in}{1.869304in}}{\pgfqpoint{1.563455in}{1.873422in}}%
\pgfpathcurveto{\pgfqpoint{1.559337in}{1.877540in}}{\pgfqpoint{1.553750in}{1.879854in}}{\pgfqpoint{1.547926in}{1.879854in}}%
\pgfpathcurveto{\pgfqpoint{1.542103in}{1.879854in}}{\pgfqpoint{1.536516in}{1.877540in}}{\pgfqpoint{1.532398in}{1.873422in}}%
\pgfpathcurveto{\pgfqpoint{1.528280in}{1.869304in}}{\pgfqpoint{1.525966in}{1.863718in}}{\pgfqpoint{1.525966in}{1.857894in}}%
\pgfpathcurveto{\pgfqpoint{1.525966in}{1.852070in}}{\pgfqpoint{1.528280in}{1.846484in}}{\pgfqpoint{1.532398in}{1.842365in}}%
\pgfpathcurveto{\pgfqpoint{1.536516in}{1.838247in}}{\pgfqpoint{1.542103in}{1.835933in}}{\pgfqpoint{1.547926in}{1.835933in}}%
\pgfpathlineto{\pgfqpoint{1.547926in}{1.835933in}}%
\pgfpathclose%
\pgfusepath{stroke,fill}%
\end{pgfscope}%
\begin{pgfscope}%
\pgfpathrectangle{\pgfqpoint{0.100000in}{0.209042in}}{\pgfqpoint{2.906250in}{2.906250in}}%
\pgfusepath{clip}%
\pgfsetbuttcap%
\pgfsetroundjoin%
\definecolor{currentfill}{rgb}{0.955849,0.404400,0.360619}%
\pgfsetfillcolor{currentfill}%
\pgfsetfillopacity{0.800000}%
\pgfsetlinewidth{1.003750pt}%
\definecolor{currentstroke}{rgb}{0.955849,0.404400,0.360619}%
\pgfsetstrokecolor{currentstroke}%
\pgfsetstrokeopacity{0.800000}%
\pgfsetdash{}{0pt}%
\pgfpathmoveto{\pgfqpoint{2.223019in}{1.323738in}}%
\pgfpathcurveto{\pgfqpoint{2.228843in}{1.323738in}}{\pgfqpoint{2.234430in}{1.326052in}}{\pgfqpoint{2.238548in}{1.330170in}}%
\pgfpathcurveto{\pgfqpoint{2.242666in}{1.334288in}}{\pgfqpoint{2.244980in}{1.339874in}}{\pgfqpoint{2.244980in}{1.345698in}}%
\pgfpathcurveto{\pgfqpoint{2.244980in}{1.351522in}}{\pgfqpoint{2.242666in}{1.357108in}}{\pgfqpoint{2.238548in}{1.361227in}}%
\pgfpathcurveto{\pgfqpoint{2.234430in}{1.365345in}}{\pgfqpoint{2.228843in}{1.367659in}}{\pgfqpoint{2.223019in}{1.367659in}}%
\pgfpathcurveto{\pgfqpoint{2.217195in}{1.367659in}}{\pgfqpoint{2.211609in}{1.365345in}}{\pgfqpoint{2.207491in}{1.361227in}}%
\pgfpathcurveto{\pgfqpoint{2.203373in}{1.357108in}}{\pgfqpoint{2.201059in}{1.351522in}}{\pgfqpoint{2.201059in}{1.345698in}}%
\pgfpathcurveto{\pgfqpoint{2.201059in}{1.339874in}}{\pgfqpoint{2.203373in}{1.334288in}}{\pgfqpoint{2.207491in}{1.330170in}}%
\pgfpathcurveto{\pgfqpoint{2.211609in}{1.326052in}}{\pgfqpoint{2.217195in}{1.323738in}}{\pgfqpoint{2.223019in}{1.323738in}}%
\pgfpathlineto{\pgfqpoint{2.223019in}{1.323738in}}%
\pgfpathclose%
\pgfusepath{stroke,fill}%
\end{pgfscope}%
\begin{pgfscope}%
\pgfpathrectangle{\pgfqpoint{0.100000in}{0.209042in}}{\pgfqpoint{2.906250in}{2.906250in}}%
\pgfusepath{clip}%
\pgfsetbuttcap%
\pgfsetroundjoin%
\definecolor{currentfill}{rgb}{0.506629,0.145958,0.507806}%
\pgfsetfillcolor{currentfill}%
\pgfsetfillopacity{0.800000}%
\pgfsetlinewidth{1.003750pt}%
\definecolor{currentstroke}{rgb}{0.506629,0.145958,0.507806}%
\pgfsetstrokecolor{currentstroke}%
\pgfsetstrokeopacity{0.800000}%
\pgfsetdash{}{0pt}%
\pgfpathmoveto{\pgfqpoint{1.291034in}{1.498546in}}%
\pgfpathcurveto{\pgfqpoint{1.296858in}{1.498546in}}{\pgfqpoint{1.302444in}{1.500860in}}{\pgfqpoint{1.306562in}{1.504978in}}%
\pgfpathcurveto{\pgfqpoint{1.310681in}{1.509096in}}{\pgfqpoint{1.312994in}{1.514683in}}{\pgfqpoint{1.312994in}{1.520507in}}%
\pgfpathcurveto{\pgfqpoint{1.312994in}{1.526330in}}{\pgfqpoint{1.310681in}{1.531917in}}{\pgfqpoint{1.306562in}{1.536035in}}%
\pgfpathcurveto{\pgfqpoint{1.302444in}{1.540153in}}{\pgfqpoint{1.296858in}{1.542467in}}{\pgfqpoint{1.291034in}{1.542467in}}%
\pgfpathcurveto{\pgfqpoint{1.285210in}{1.542467in}}{\pgfqpoint{1.279624in}{1.540153in}}{\pgfqpoint{1.275506in}{1.536035in}}%
\pgfpathcurveto{\pgfqpoint{1.271388in}{1.531917in}}{\pgfqpoint{1.269074in}{1.526330in}}{\pgfqpoint{1.269074in}{1.520507in}}%
\pgfpathcurveto{\pgfqpoint{1.269074in}{1.514683in}}{\pgfqpoint{1.271388in}{1.509096in}}{\pgfqpoint{1.275506in}{1.504978in}}%
\pgfpathcurveto{\pgfqpoint{1.279624in}{1.500860in}}{\pgfqpoint{1.285210in}{1.498546in}}{\pgfqpoint{1.291034in}{1.498546in}}%
\pgfpathlineto{\pgfqpoint{1.291034in}{1.498546in}}%
\pgfpathclose%
\pgfusepath{stroke,fill}%
\end{pgfscope}%
\begin{pgfscope}%
\pgfpathrectangle{\pgfqpoint{0.100000in}{0.209042in}}{\pgfqpoint{2.906250in}{2.906250in}}%
\pgfusepath{clip}%
\pgfsetbuttcap%
\pgfsetroundjoin%
\definecolor{currentfill}{rgb}{0.329114,0.075972,0.489287}%
\pgfsetfillcolor{currentfill}%
\pgfsetfillopacity{0.800000}%
\pgfsetlinewidth{1.003750pt}%
\definecolor{currentstroke}{rgb}{0.329114,0.075972,0.489287}%
\pgfsetstrokecolor{currentstroke}%
\pgfsetstrokeopacity{0.800000}%
\pgfsetdash{}{0pt}%
\pgfpathmoveto{\pgfqpoint{1.882427in}{1.563622in}}%
\pgfpathcurveto{\pgfqpoint{1.888251in}{1.563622in}}{\pgfqpoint{1.893838in}{1.565935in}}{\pgfqpoint{1.897956in}{1.570054in}}%
\pgfpathcurveto{\pgfqpoint{1.902074in}{1.574172in}}{\pgfqpoint{1.904388in}{1.579758in}}{\pgfqpoint{1.904388in}{1.585582in}}%
\pgfpathcurveto{\pgfqpoint{1.904388in}{1.591406in}}{\pgfqpoint{1.902074in}{1.596992in}}{\pgfqpoint{1.897956in}{1.601110in}}%
\pgfpathcurveto{\pgfqpoint{1.893838in}{1.605228in}}{\pgfqpoint{1.888251in}{1.607542in}}{\pgfqpoint{1.882427in}{1.607542in}}%
\pgfpathcurveto{\pgfqpoint{1.876604in}{1.607542in}}{\pgfqpoint{1.871017in}{1.605228in}}{\pgfqpoint{1.866899in}{1.601110in}}%
\pgfpathcurveto{\pgfqpoint{1.862781in}{1.596992in}}{\pgfqpoint{1.860467in}{1.591406in}}{\pgfqpoint{1.860467in}{1.585582in}}%
\pgfpathcurveto{\pgfqpoint{1.860467in}{1.579758in}}{\pgfqpoint{1.862781in}{1.574172in}}{\pgfqpoint{1.866899in}{1.570054in}}%
\pgfpathcurveto{\pgfqpoint{1.871017in}{1.565935in}}{\pgfqpoint{1.876604in}{1.563622in}}{\pgfqpoint{1.882427in}{1.563622in}}%
\pgfpathlineto{\pgfqpoint{1.882427in}{1.563622in}}%
\pgfpathclose%
\pgfusepath{stroke,fill}%
\end{pgfscope}%
\begin{pgfscope}%
\pgfpathrectangle{\pgfqpoint{0.100000in}{0.209042in}}{\pgfqpoint{2.906250in}{2.906250in}}%
\pgfusepath{clip}%
\pgfsetbuttcap%
\pgfsetroundjoin%
\definecolor{currentfill}{rgb}{0.069764,0.049726,0.193735}%
\pgfsetfillcolor{currentfill}%
\pgfsetfillopacity{0.800000}%
\pgfsetlinewidth{1.003750pt}%
\definecolor{currentstroke}{rgb}{0.069764,0.049726,0.193735}%
\pgfsetstrokecolor{currentstroke}%
\pgfsetstrokeopacity{0.800000}%
\pgfsetdash{}{0pt}%
\pgfpathmoveto{\pgfqpoint{1.701899in}{1.793694in}}%
\pgfpathcurveto{\pgfqpoint{1.707723in}{1.793694in}}{\pgfqpoint{1.713309in}{1.796008in}}{\pgfqpoint{1.717427in}{1.800126in}}%
\pgfpathcurveto{\pgfqpoint{1.721545in}{1.804245in}}{\pgfqpoint{1.723859in}{1.809831in}}{\pgfqpoint{1.723859in}{1.815655in}}%
\pgfpathcurveto{\pgfqpoint{1.723859in}{1.821479in}}{\pgfqpoint{1.721545in}{1.827065in}}{\pgfqpoint{1.717427in}{1.831183in}}%
\pgfpathcurveto{\pgfqpoint{1.713309in}{1.835301in}}{\pgfqpoint{1.707723in}{1.837615in}}{\pgfqpoint{1.701899in}{1.837615in}}%
\pgfpathcurveto{\pgfqpoint{1.696075in}{1.837615in}}{\pgfqpoint{1.690489in}{1.835301in}}{\pgfqpoint{1.686371in}{1.831183in}}%
\pgfpathcurveto{\pgfqpoint{1.682252in}{1.827065in}}{\pgfqpoint{1.679939in}{1.821479in}}{\pgfqpoint{1.679939in}{1.815655in}}%
\pgfpathcurveto{\pgfqpoint{1.679939in}{1.809831in}}{\pgfqpoint{1.682252in}{1.804245in}}{\pgfqpoint{1.686371in}{1.800126in}}%
\pgfpathcurveto{\pgfqpoint{1.690489in}{1.796008in}}{\pgfqpoint{1.696075in}{1.793694in}}{\pgfqpoint{1.701899in}{1.793694in}}%
\pgfpathlineto{\pgfqpoint{1.701899in}{1.793694in}}%
\pgfpathclose%
\pgfusepath{stroke,fill}%
\end{pgfscope}%
\begin{pgfscope}%
\pgfpathrectangle{\pgfqpoint{0.100000in}{0.209042in}}{\pgfqpoint{2.906250in}{2.906250in}}%
\pgfusepath{clip}%
\pgfsetbuttcap%
\pgfsetroundjoin%
\definecolor{currentfill}{rgb}{0.322899,0.073782,0.487408}%
\pgfsetfillcolor{currentfill}%
\pgfsetfillopacity{0.800000}%
\pgfsetlinewidth{1.003750pt}%
\definecolor{currentstroke}{rgb}{0.322899,0.073782,0.487408}%
\pgfsetstrokecolor{currentstroke}%
\pgfsetstrokeopacity{0.800000}%
\pgfsetdash{}{0pt}%
\pgfpathmoveto{\pgfqpoint{1.913354in}{1.618438in}}%
\pgfpathcurveto{\pgfqpoint{1.919178in}{1.618438in}}{\pgfqpoint{1.924764in}{1.620752in}}{\pgfqpoint{1.928882in}{1.624870in}}%
\pgfpathcurveto{\pgfqpoint{1.933000in}{1.628988in}}{\pgfqpoint{1.935314in}{1.634574in}}{\pgfqpoint{1.935314in}{1.640398in}}%
\pgfpathcurveto{\pgfqpoint{1.935314in}{1.646222in}}{\pgfqpoint{1.933000in}{1.651808in}}{\pgfqpoint{1.928882in}{1.655926in}}%
\pgfpathcurveto{\pgfqpoint{1.924764in}{1.660044in}}{\pgfqpoint{1.919178in}{1.662358in}}{\pgfqpoint{1.913354in}{1.662358in}}%
\pgfpathcurveto{\pgfqpoint{1.907530in}{1.662358in}}{\pgfqpoint{1.901944in}{1.660044in}}{\pgfqpoint{1.897826in}{1.655926in}}%
\pgfpathcurveto{\pgfqpoint{1.893707in}{1.651808in}}{\pgfqpoint{1.891393in}{1.646222in}}{\pgfqpoint{1.891393in}{1.640398in}}%
\pgfpathcurveto{\pgfqpoint{1.891393in}{1.634574in}}{\pgfqpoint{1.893707in}{1.628988in}}{\pgfqpoint{1.897826in}{1.624870in}}%
\pgfpathcurveto{\pgfqpoint{1.901944in}{1.620752in}}{\pgfqpoint{1.907530in}{1.618438in}}{\pgfqpoint{1.913354in}{1.618438in}}%
\pgfpathlineto{\pgfqpoint{1.913354in}{1.618438in}}%
\pgfpathclose%
\pgfusepath{stroke,fill}%
\end{pgfscope}%
\begin{pgfscope}%
\pgfpathrectangle{\pgfqpoint{0.100000in}{0.209042in}}{\pgfqpoint{2.906250in}{2.906250in}}%
\pgfusepath{clip}%
\pgfsetbuttcap%
\pgfsetroundjoin%
\definecolor{currentfill}{rgb}{0.225302,0.060445,0.431742}%
\pgfsetfillcolor{currentfill}%
\pgfsetfillopacity{0.800000}%
\pgfsetlinewidth{1.003750pt}%
\definecolor{currentstroke}{rgb}{0.225302,0.060445,0.431742}%
\pgfsetstrokecolor{currentstroke}%
\pgfsetstrokeopacity{0.800000}%
\pgfsetdash{}{0pt}%
\pgfpathmoveto{\pgfqpoint{1.793757in}{1.848170in}}%
\pgfpathcurveto{\pgfqpoint{1.799580in}{1.848170in}}{\pgfqpoint{1.805167in}{1.850484in}}{\pgfqpoint{1.809285in}{1.854602in}}%
\pgfpathcurveto{\pgfqpoint{1.813403in}{1.858721in}}{\pgfqpoint{1.815717in}{1.864307in}}{\pgfqpoint{1.815717in}{1.870131in}}%
\pgfpathcurveto{\pgfqpoint{1.815717in}{1.875955in}}{\pgfqpoint{1.813403in}{1.881541in}}{\pgfqpoint{1.809285in}{1.885659in}}%
\pgfpathcurveto{\pgfqpoint{1.805167in}{1.889777in}}{\pgfqpoint{1.799580in}{1.892091in}}{\pgfqpoint{1.793757in}{1.892091in}}%
\pgfpathcurveto{\pgfqpoint{1.787933in}{1.892091in}}{\pgfqpoint{1.782346in}{1.889777in}}{\pgfqpoint{1.778228in}{1.885659in}}%
\pgfpathcurveto{\pgfqpoint{1.774110in}{1.881541in}}{\pgfqpoint{1.771796in}{1.875955in}}{\pgfqpoint{1.771796in}{1.870131in}}%
\pgfpathcurveto{\pgfqpoint{1.771796in}{1.864307in}}{\pgfqpoint{1.774110in}{1.858721in}}{\pgfqpoint{1.778228in}{1.854602in}}%
\pgfpathcurveto{\pgfqpoint{1.782346in}{1.850484in}}{\pgfqpoint{1.787933in}{1.848170in}}{\pgfqpoint{1.793757in}{1.848170in}}%
\pgfpathlineto{\pgfqpoint{1.793757in}{1.848170in}}%
\pgfpathclose%
\pgfusepath{stroke,fill}%
\end{pgfscope}%
\begin{pgfscope}%
\pgfpathrectangle{\pgfqpoint{0.100000in}{0.209042in}}{\pgfqpoint{2.906250in}{2.906250in}}%
\pgfusepath{clip}%
\pgfsetbuttcap%
\pgfsetroundjoin%
\definecolor{currentfill}{rgb}{0.191460,0.064818,0.396152}%
\pgfsetfillcolor{currentfill}%
\pgfsetfillopacity{0.800000}%
\pgfsetlinewidth{1.003750pt}%
\definecolor{currentstroke}{rgb}{0.191460,0.064818,0.396152}%
\pgfsetstrokecolor{currentstroke}%
\pgfsetstrokeopacity{0.800000}%
\pgfsetdash{}{0pt}%
\pgfpathmoveto{\pgfqpoint{1.409016in}{1.796170in}}%
\pgfpathcurveto{\pgfqpoint{1.414840in}{1.796170in}}{\pgfqpoint{1.420426in}{1.798484in}}{\pgfqpoint{1.424544in}{1.802602in}}%
\pgfpathcurveto{\pgfqpoint{1.428662in}{1.806720in}}{\pgfqpoint{1.430976in}{1.812306in}}{\pgfqpoint{1.430976in}{1.818130in}}%
\pgfpathcurveto{\pgfqpoint{1.430976in}{1.823954in}}{\pgfqpoint{1.428662in}{1.829540in}}{\pgfqpoint{1.424544in}{1.833659in}}%
\pgfpathcurveto{\pgfqpoint{1.420426in}{1.837777in}}{\pgfqpoint{1.414840in}{1.840091in}}{\pgfqpoint{1.409016in}{1.840091in}}%
\pgfpathcurveto{\pgfqpoint{1.403192in}{1.840091in}}{\pgfqpoint{1.397606in}{1.837777in}}{\pgfqpoint{1.393487in}{1.833659in}}%
\pgfpathcurveto{\pgfqpoint{1.389369in}{1.829540in}}{\pgfqpoint{1.387055in}{1.823954in}}{\pgfqpoint{1.387055in}{1.818130in}}%
\pgfpathcurveto{\pgfqpoint{1.387055in}{1.812306in}}{\pgfqpoint{1.389369in}{1.806720in}}{\pgfqpoint{1.393487in}{1.802602in}}%
\pgfpathcurveto{\pgfqpoint{1.397606in}{1.798484in}}{\pgfqpoint{1.403192in}{1.796170in}}{\pgfqpoint{1.409016in}{1.796170in}}%
\pgfpathlineto{\pgfqpoint{1.409016in}{1.796170in}}%
\pgfpathclose%
\pgfusepath{stroke,fill}%
\end{pgfscope}%
\begin{pgfscope}%
\pgfpathrectangle{\pgfqpoint{0.100000in}{0.209042in}}{\pgfqpoint{2.906250in}{2.906250in}}%
\pgfusepath{clip}%
\pgfsetbuttcap%
\pgfsetroundjoin%
\definecolor{currentfill}{rgb}{0.191460,0.064818,0.396152}%
\pgfsetfillcolor{currentfill}%
\pgfsetfillopacity{0.800000}%
\pgfsetlinewidth{1.003750pt}%
\definecolor{currentstroke}{rgb}{0.191460,0.064818,0.396152}%
\pgfsetstrokecolor{currentstroke}%
\pgfsetstrokeopacity{0.800000}%
\pgfsetdash{}{0pt}%
\pgfpathmoveto{\pgfqpoint{1.674938in}{1.530240in}}%
\pgfpathcurveto{\pgfqpoint{1.680762in}{1.530240in}}{\pgfqpoint{1.686348in}{1.532554in}}{\pgfqpoint{1.690466in}{1.536672in}}%
\pgfpathcurveto{\pgfqpoint{1.694585in}{1.540790in}}{\pgfqpoint{1.696898in}{1.546376in}}{\pgfqpoint{1.696898in}{1.552200in}}%
\pgfpathcurveto{\pgfqpoint{1.696898in}{1.558024in}}{\pgfqpoint{1.694585in}{1.563610in}}{\pgfqpoint{1.690466in}{1.567729in}}%
\pgfpathcurveto{\pgfqpoint{1.686348in}{1.571847in}}{\pgfqpoint{1.680762in}{1.574161in}}{\pgfqpoint{1.674938in}{1.574161in}}%
\pgfpathcurveto{\pgfqpoint{1.669114in}{1.574161in}}{\pgfqpoint{1.663528in}{1.571847in}}{\pgfqpoint{1.659410in}{1.567729in}}%
\pgfpathcurveto{\pgfqpoint{1.655292in}{1.563610in}}{\pgfqpoint{1.652978in}{1.558024in}}{\pgfqpoint{1.652978in}{1.552200in}}%
\pgfpathcurveto{\pgfqpoint{1.652978in}{1.546376in}}{\pgfqpoint{1.655292in}{1.540790in}}{\pgfqpoint{1.659410in}{1.536672in}}%
\pgfpathcurveto{\pgfqpoint{1.663528in}{1.532554in}}{\pgfqpoint{1.669114in}{1.530240in}}{\pgfqpoint{1.674938in}{1.530240in}}%
\pgfpathlineto{\pgfqpoint{1.674938in}{1.530240in}}%
\pgfpathclose%
\pgfusepath{stroke,fill}%
\end{pgfscope}%
\begin{pgfscope}%
\pgfpathrectangle{\pgfqpoint{0.100000in}{0.209042in}}{\pgfqpoint{2.906250in}{2.906250in}}%
\pgfusepath{clip}%
\pgfsetbuttcap%
\pgfsetroundjoin%
\definecolor{currentfill}{rgb}{0.123833,0.067295,0.295879}%
\pgfsetfillcolor{currentfill}%
\pgfsetfillopacity{0.800000}%
\pgfsetlinewidth{1.003750pt}%
\definecolor{currentstroke}{rgb}{0.123833,0.067295,0.295879}%
\pgfsetstrokecolor{currentstroke}%
\pgfsetstrokeopacity{0.800000}%
\pgfsetdash{}{0pt}%
\pgfpathmoveto{\pgfqpoint{1.755533in}{1.806628in}}%
\pgfpathcurveto{\pgfqpoint{1.761357in}{1.806628in}}{\pgfqpoint{1.766943in}{1.808942in}}{\pgfqpoint{1.771062in}{1.813060in}}%
\pgfpathcurveto{\pgfqpoint{1.775180in}{1.817178in}}{\pgfqpoint{1.777494in}{1.822765in}}{\pgfqpoint{1.777494in}{1.828588in}}%
\pgfpathcurveto{\pgfqpoint{1.777494in}{1.834412in}}{\pgfqpoint{1.775180in}{1.839999in}}{\pgfqpoint{1.771062in}{1.844117in}}%
\pgfpathcurveto{\pgfqpoint{1.766943in}{1.848235in}}{\pgfqpoint{1.761357in}{1.850549in}}{\pgfqpoint{1.755533in}{1.850549in}}%
\pgfpathcurveto{\pgfqpoint{1.749709in}{1.850549in}}{\pgfqpoint{1.744123in}{1.848235in}}{\pgfqpoint{1.740005in}{1.844117in}}%
\pgfpathcurveto{\pgfqpoint{1.735887in}{1.839999in}}{\pgfqpoint{1.733573in}{1.834412in}}{\pgfqpoint{1.733573in}{1.828588in}}%
\pgfpathcurveto{\pgfqpoint{1.733573in}{1.822765in}}{\pgfqpoint{1.735887in}{1.817178in}}{\pgfqpoint{1.740005in}{1.813060in}}%
\pgfpathcurveto{\pgfqpoint{1.744123in}{1.808942in}}{\pgfqpoint{1.749709in}{1.806628in}}{\pgfqpoint{1.755533in}{1.806628in}}%
\pgfpathlineto{\pgfqpoint{1.755533in}{1.806628in}}%
\pgfpathclose%
\pgfusepath{stroke,fill}%
\end{pgfscope}%
\begin{pgfscope}%
\pgfpathrectangle{\pgfqpoint{0.100000in}{0.209042in}}{\pgfqpoint{2.906250in}{2.906250in}}%
\pgfusepath{clip}%
\pgfsetbuttcap%
\pgfsetroundjoin%
\definecolor{currentfill}{rgb}{0.426877,0.115395,0.505714}%
\pgfsetfillcolor{currentfill}%
\pgfsetfillopacity{0.800000}%
\pgfsetlinewidth{1.003750pt}%
\definecolor{currentstroke}{rgb}{0.426877,0.115395,0.505714}%
\pgfsetstrokecolor{currentstroke}%
\pgfsetstrokeopacity{0.800000}%
\pgfsetdash{}{0pt}%
\pgfpathmoveto{\pgfqpoint{1.474562in}{1.435039in}}%
\pgfpathcurveto{\pgfqpoint{1.480386in}{1.435039in}}{\pgfqpoint{1.485972in}{1.437353in}}{\pgfqpoint{1.490090in}{1.441471in}}%
\pgfpathcurveto{\pgfqpoint{1.494209in}{1.445589in}}{\pgfqpoint{1.496522in}{1.451175in}}{\pgfqpoint{1.496522in}{1.456999in}}%
\pgfpathcurveto{\pgfqpoint{1.496522in}{1.462823in}}{\pgfqpoint{1.494209in}{1.468409in}}{\pgfqpoint{1.490090in}{1.472527in}}%
\pgfpathcurveto{\pgfqpoint{1.485972in}{1.476646in}}{\pgfqpoint{1.480386in}{1.478959in}}{\pgfqpoint{1.474562in}{1.478959in}}%
\pgfpathcurveto{\pgfqpoint{1.468738in}{1.478959in}}{\pgfqpoint{1.463152in}{1.476646in}}{\pgfqpoint{1.459034in}{1.472527in}}%
\pgfpathcurveto{\pgfqpoint{1.454916in}{1.468409in}}{\pgfqpoint{1.452602in}{1.462823in}}{\pgfqpoint{1.452602in}{1.456999in}}%
\pgfpathcurveto{\pgfqpoint{1.452602in}{1.451175in}}{\pgfqpoint{1.454916in}{1.445589in}}{\pgfqpoint{1.459034in}{1.441471in}}%
\pgfpathcurveto{\pgfqpoint{1.463152in}{1.437353in}}{\pgfqpoint{1.468738in}{1.435039in}}{\pgfqpoint{1.474562in}{1.435039in}}%
\pgfpathlineto{\pgfqpoint{1.474562in}{1.435039in}}%
\pgfpathclose%
\pgfusepath{stroke,fill}%
\end{pgfscope}%
\begin{pgfscope}%
\pgfpathrectangle{\pgfqpoint{0.100000in}{0.209042in}}{\pgfqpoint{2.906250in}{2.906250in}}%
\pgfusepath{clip}%
\pgfsetbuttcap%
\pgfsetroundjoin%
\definecolor{currentfill}{rgb}{0.211718,0.061992,0.418647}%
\pgfsetfillcolor{currentfill}%
\pgfsetfillopacity{0.800000}%
\pgfsetlinewidth{1.003750pt}%
\definecolor{currentstroke}{rgb}{0.211718,0.061992,0.418647}%
\pgfsetstrokecolor{currentstroke}%
\pgfsetstrokeopacity{0.800000}%
\pgfsetdash{}{0pt}%
\pgfpathmoveto{\pgfqpoint{1.378429in}{1.731809in}}%
\pgfpathcurveto{\pgfqpoint{1.384252in}{1.731809in}}{\pgfqpoint{1.389839in}{1.734123in}}{\pgfqpoint{1.393957in}{1.738241in}}%
\pgfpathcurveto{\pgfqpoint{1.398075in}{1.742359in}}{\pgfqpoint{1.400389in}{1.747945in}}{\pgfqpoint{1.400389in}{1.753769in}}%
\pgfpathcurveto{\pgfqpoint{1.400389in}{1.759593in}}{\pgfqpoint{1.398075in}{1.765179in}}{\pgfqpoint{1.393957in}{1.769297in}}%
\pgfpathcurveto{\pgfqpoint{1.389839in}{1.773416in}}{\pgfqpoint{1.384252in}{1.775729in}}{\pgfqpoint{1.378429in}{1.775729in}}%
\pgfpathcurveto{\pgfqpoint{1.372605in}{1.775729in}}{\pgfqpoint{1.367018in}{1.773416in}}{\pgfqpoint{1.362900in}{1.769297in}}%
\pgfpathcurveto{\pgfqpoint{1.358782in}{1.765179in}}{\pgfqpoint{1.356468in}{1.759593in}}{\pgfqpoint{1.356468in}{1.753769in}}%
\pgfpathcurveto{\pgfqpoint{1.356468in}{1.747945in}}{\pgfqpoint{1.358782in}{1.742359in}}{\pgfqpoint{1.362900in}{1.738241in}}%
\pgfpathcurveto{\pgfqpoint{1.367018in}{1.734123in}}{\pgfqpoint{1.372605in}{1.731809in}}{\pgfqpoint{1.378429in}{1.731809in}}%
\pgfpathlineto{\pgfqpoint{1.378429in}{1.731809in}}%
\pgfpathclose%
\pgfusepath{stroke,fill}%
\end{pgfscope}%
\begin{pgfscope}%
\pgfpathrectangle{\pgfqpoint{0.100000in}{0.209042in}}{\pgfqpoint{2.906250in}{2.906250in}}%
\pgfusepath{clip}%
\pgfsetbuttcap%
\pgfsetroundjoin%
\definecolor{currentfill}{rgb}{0.304081,0.067835,0.480812}%
\pgfsetfillcolor{currentfill}%
\pgfsetfillopacity{0.800000}%
\pgfsetlinewidth{1.003750pt}%
\definecolor{currentstroke}{rgb}{0.304081,0.067835,0.480812}%
\pgfsetstrokecolor{currentstroke}%
\pgfsetstrokeopacity{0.800000}%
\pgfsetdash{}{0pt}%
\pgfpathmoveto{\pgfqpoint{1.907920in}{1.634771in}}%
\pgfpathcurveto{\pgfqpoint{1.913744in}{1.634771in}}{\pgfqpoint{1.919330in}{1.637085in}}{\pgfqpoint{1.923448in}{1.641203in}}%
\pgfpathcurveto{\pgfqpoint{1.927566in}{1.645321in}}{\pgfqpoint{1.929880in}{1.650907in}}{\pgfqpoint{1.929880in}{1.656731in}}%
\pgfpathcurveto{\pgfqpoint{1.929880in}{1.662555in}}{\pgfqpoint{1.927566in}{1.668141in}}{\pgfqpoint{1.923448in}{1.672260in}}%
\pgfpathcurveto{\pgfqpoint{1.919330in}{1.676378in}}{\pgfqpoint{1.913744in}{1.678692in}}{\pgfqpoint{1.907920in}{1.678692in}}%
\pgfpathcurveto{\pgfqpoint{1.902096in}{1.678692in}}{\pgfqpoint{1.896510in}{1.676378in}}{\pgfqpoint{1.892392in}{1.672260in}}%
\pgfpathcurveto{\pgfqpoint{1.888274in}{1.668141in}}{\pgfqpoint{1.885960in}{1.662555in}}{\pgfqpoint{1.885960in}{1.656731in}}%
\pgfpathcurveto{\pgfqpoint{1.885960in}{1.650907in}}{\pgfqpoint{1.888274in}{1.645321in}}{\pgfqpoint{1.892392in}{1.641203in}}%
\pgfpathcurveto{\pgfqpoint{1.896510in}{1.637085in}}{\pgfqpoint{1.902096in}{1.634771in}}{\pgfqpoint{1.907920in}{1.634771in}}%
\pgfpathlineto{\pgfqpoint{1.907920in}{1.634771in}}%
\pgfpathclose%
\pgfusepath{stroke,fill}%
\end{pgfscope}%
\begin{pgfscope}%
\pgfpathrectangle{\pgfqpoint{0.100000in}{0.209042in}}{\pgfqpoint{2.906250in}{2.906250in}}%
\pgfusepath{clip}%
\pgfsetbuttcap%
\pgfsetroundjoin%
\definecolor{currentfill}{rgb}{0.009426,0.008022,0.052844}%
\pgfsetfillcolor{currentfill}%
\pgfsetfillopacity{0.800000}%
\pgfsetlinewidth{1.003750pt}%
\definecolor{currentstroke}{rgb}{0.009426,0.008022,0.052844}%
\pgfsetstrokecolor{currentstroke}%
\pgfsetstrokeopacity{0.800000}%
\pgfsetdash{}{0pt}%
\pgfpathmoveto{\pgfqpoint{1.589792in}{1.672326in}}%
\pgfpathcurveto{\pgfqpoint{1.595616in}{1.672326in}}{\pgfqpoint{1.601202in}{1.674640in}}{\pgfqpoint{1.605320in}{1.678758in}}%
\pgfpathcurveto{\pgfqpoint{1.609439in}{1.682876in}}{\pgfqpoint{1.611752in}{1.688462in}}{\pgfqpoint{1.611752in}{1.694286in}}%
\pgfpathcurveto{\pgfqpoint{1.611752in}{1.700110in}}{\pgfqpoint{1.609439in}{1.705697in}}{\pgfqpoint{1.605320in}{1.709815in}}%
\pgfpathcurveto{\pgfqpoint{1.601202in}{1.713933in}}{\pgfqpoint{1.595616in}{1.716247in}}{\pgfqpoint{1.589792in}{1.716247in}}%
\pgfpathcurveto{\pgfqpoint{1.583968in}{1.716247in}}{\pgfqpoint{1.578382in}{1.713933in}}{\pgfqpoint{1.574264in}{1.709815in}}%
\pgfpathcurveto{\pgfqpoint{1.570146in}{1.705697in}}{\pgfqpoint{1.567832in}{1.700110in}}{\pgfqpoint{1.567832in}{1.694286in}}%
\pgfpathcurveto{\pgfqpoint{1.567832in}{1.688462in}}{\pgfqpoint{1.570146in}{1.682876in}}{\pgfqpoint{1.574264in}{1.678758in}}%
\pgfpathcurveto{\pgfqpoint{1.578382in}{1.674640in}}{\pgfqpoint{1.583968in}{1.672326in}}{\pgfqpoint{1.589792in}{1.672326in}}%
\pgfpathlineto{\pgfqpoint{1.589792in}{1.672326in}}%
\pgfpathclose%
\pgfusepath{stroke,fill}%
\end{pgfscope}%
\begin{pgfscope}%
\pgfpathrectangle{\pgfqpoint{0.100000in}{0.209042in}}{\pgfqpoint{2.906250in}{2.906250in}}%
\pgfusepath{clip}%
\pgfsetbuttcap%
\pgfsetroundjoin%
\definecolor{currentfill}{rgb}{0.420791,0.112920,0.505215}%
\pgfsetfillcolor{currentfill}%
\pgfsetfillopacity{0.800000}%
\pgfsetlinewidth{1.003750pt}%
\definecolor{currentstroke}{rgb}{0.420791,0.112920,0.505215}%
\pgfsetstrokecolor{currentstroke}%
\pgfsetstrokeopacity{0.800000}%
\pgfsetdash{}{0pt}%
\pgfpathmoveto{\pgfqpoint{1.997355in}{1.691738in}}%
\pgfpathcurveto{\pgfqpoint{2.003179in}{1.691738in}}{\pgfqpoint{2.008765in}{1.694052in}}{\pgfqpoint{2.012883in}{1.698170in}}%
\pgfpathcurveto{\pgfqpoint{2.017001in}{1.702288in}}{\pgfqpoint{2.019315in}{1.707874in}}{\pgfqpoint{2.019315in}{1.713698in}}%
\pgfpathcurveto{\pgfqpoint{2.019315in}{1.719522in}}{\pgfqpoint{2.017001in}{1.725108in}}{\pgfqpoint{2.012883in}{1.729226in}}%
\pgfpathcurveto{\pgfqpoint{2.008765in}{1.733345in}}{\pgfqpoint{2.003179in}{1.735658in}}{\pgfqpoint{1.997355in}{1.735658in}}%
\pgfpathcurveto{\pgfqpoint{1.991531in}{1.735658in}}{\pgfqpoint{1.985945in}{1.733345in}}{\pgfqpoint{1.981826in}{1.729226in}}%
\pgfpathcurveto{\pgfqpoint{1.977708in}{1.725108in}}{\pgfqpoint{1.975394in}{1.719522in}}{\pgfqpoint{1.975394in}{1.713698in}}%
\pgfpathcurveto{\pgfqpoint{1.975394in}{1.707874in}}{\pgfqpoint{1.977708in}{1.702288in}}{\pgfqpoint{1.981826in}{1.698170in}}%
\pgfpathcurveto{\pgfqpoint{1.985945in}{1.694052in}}{\pgfqpoint{1.991531in}{1.691738in}}{\pgfqpoint{1.997355in}{1.691738in}}%
\pgfpathlineto{\pgfqpoint{1.997355in}{1.691738in}}%
\pgfpathclose%
\pgfusepath{stroke,fill}%
\end{pgfscope}%
\begin{pgfscope}%
\pgfpathrectangle{\pgfqpoint{0.100000in}{0.209042in}}{\pgfqpoint{2.906250in}{2.906250in}}%
\pgfusepath{clip}%
\pgfsetbuttcap%
\pgfsetroundjoin%
\definecolor{currentfill}{rgb}{0.709962,0.212797,0.477201}%
\pgfsetfillcolor{currentfill}%
\pgfsetfillopacity{0.800000}%
\pgfsetlinewidth{1.003750pt}%
\definecolor{currentstroke}{rgb}{0.709962,0.212797,0.477201}%
\pgfsetstrokecolor{currentstroke}%
\pgfsetstrokeopacity{0.800000}%
\pgfsetdash{}{0pt}%
\pgfpathmoveto{\pgfqpoint{1.657408in}{1.252335in}}%
\pgfpathcurveto{\pgfqpoint{1.663232in}{1.252335in}}{\pgfqpoint{1.668818in}{1.254649in}}{\pgfqpoint{1.672936in}{1.258767in}}%
\pgfpathcurveto{\pgfqpoint{1.677055in}{1.262885in}}{\pgfqpoint{1.679368in}{1.268471in}}{\pgfqpoint{1.679368in}{1.274295in}}%
\pgfpathcurveto{\pgfqpoint{1.679368in}{1.280119in}}{\pgfqpoint{1.677055in}{1.285705in}}{\pgfqpoint{1.672936in}{1.289823in}}%
\pgfpathcurveto{\pgfqpoint{1.668818in}{1.293941in}}{\pgfqpoint{1.663232in}{1.296255in}}{\pgfqpoint{1.657408in}{1.296255in}}%
\pgfpathcurveto{\pgfqpoint{1.651584in}{1.296255in}}{\pgfqpoint{1.645998in}{1.293941in}}{\pgfqpoint{1.641880in}{1.289823in}}%
\pgfpathcurveto{\pgfqpoint{1.637762in}{1.285705in}}{\pgfqpoint{1.635448in}{1.280119in}}{\pgfqpoint{1.635448in}{1.274295in}}%
\pgfpathcurveto{\pgfqpoint{1.635448in}{1.268471in}}{\pgfqpoint{1.637762in}{1.262885in}}{\pgfqpoint{1.641880in}{1.258767in}}%
\pgfpathcurveto{\pgfqpoint{1.645998in}{1.254649in}}{\pgfqpoint{1.651584in}{1.252335in}}{\pgfqpoint{1.657408in}{1.252335in}}%
\pgfpathlineto{\pgfqpoint{1.657408in}{1.252335in}}%
\pgfpathclose%
\pgfusepath{stroke,fill}%
\end{pgfscope}%
\begin{pgfscope}%
\pgfpathrectangle{\pgfqpoint{0.100000in}{0.209042in}}{\pgfqpoint{2.906250in}{2.906250in}}%
\pgfusepath{clip}%
\pgfsetbuttcap%
\pgfsetroundjoin%
\definecolor{currentfill}{rgb}{0.475780,0.134577,0.507921}%
\pgfsetfillcolor{currentfill}%
\pgfsetfillopacity{0.800000}%
\pgfsetlinewidth{1.003750pt}%
\definecolor{currentstroke}{rgb}{0.475780,0.134577,0.507921}%
\pgfsetstrokecolor{currentstroke}%
\pgfsetstrokeopacity{0.800000}%
\pgfsetdash{}{0pt}%
\pgfpathmoveto{\pgfqpoint{1.522488in}{2.044101in}}%
\pgfpathcurveto{\pgfqpoint{1.528312in}{2.044101in}}{\pgfqpoint{1.533898in}{2.046415in}}{\pgfqpoint{1.538016in}{2.050533in}}%
\pgfpathcurveto{\pgfqpoint{1.542134in}{2.054651in}}{\pgfqpoint{1.544448in}{2.060237in}}{\pgfqpoint{1.544448in}{2.066061in}}%
\pgfpathcurveto{\pgfqpoint{1.544448in}{2.071885in}}{\pgfqpoint{1.542134in}{2.077472in}}{\pgfqpoint{1.538016in}{2.081590in}}%
\pgfpathcurveto{\pgfqpoint{1.533898in}{2.085708in}}{\pgfqpoint{1.528312in}{2.088022in}}{\pgfqpoint{1.522488in}{2.088022in}}%
\pgfpathcurveto{\pgfqpoint{1.516664in}{2.088022in}}{\pgfqpoint{1.511078in}{2.085708in}}{\pgfqpoint{1.506960in}{2.081590in}}%
\pgfpathcurveto{\pgfqpoint{1.502841in}{2.077472in}}{\pgfqpoint{1.500528in}{2.071885in}}{\pgfqpoint{1.500528in}{2.066061in}}%
\pgfpathcurveto{\pgfqpoint{1.500528in}{2.060237in}}{\pgfqpoint{1.502841in}{2.054651in}}{\pgfqpoint{1.506960in}{2.050533in}}%
\pgfpathcurveto{\pgfqpoint{1.511078in}{2.046415in}}{\pgfqpoint{1.516664in}{2.044101in}}{\pgfqpoint{1.522488in}{2.044101in}}%
\pgfpathlineto{\pgfqpoint{1.522488in}{2.044101in}}%
\pgfpathclose%
\pgfusepath{stroke,fill}%
\end{pgfscope}%
\begin{pgfscope}%
\pgfpathrectangle{\pgfqpoint{0.100000in}{0.209042in}}{\pgfqpoint{2.906250in}{2.906250in}}%
\pgfusepath{clip}%
\pgfsetbuttcap%
\pgfsetroundjoin%
\definecolor{currentfill}{rgb}{0.191460,0.064818,0.396152}%
\pgfsetfillcolor{currentfill}%
\pgfsetfillopacity{0.800000}%
\pgfsetlinewidth{1.003750pt}%
\definecolor{currentstroke}{rgb}{0.191460,0.064818,0.396152}%
\pgfsetstrokecolor{currentstroke}%
\pgfsetstrokeopacity{0.800000}%
\pgfsetdash{}{0pt}%
\pgfpathmoveto{\pgfqpoint{1.825559in}{1.611871in}}%
\pgfpathcurveto{\pgfqpoint{1.831383in}{1.611871in}}{\pgfqpoint{1.836969in}{1.614185in}}{\pgfqpoint{1.841087in}{1.618303in}}%
\pgfpathcurveto{\pgfqpoint{1.845205in}{1.622421in}}{\pgfqpoint{1.847519in}{1.628007in}}{\pgfqpoint{1.847519in}{1.633831in}}%
\pgfpathcurveto{\pgfqpoint{1.847519in}{1.639655in}}{\pgfqpoint{1.845205in}{1.645241in}}{\pgfqpoint{1.841087in}{1.649359in}}%
\pgfpathcurveto{\pgfqpoint{1.836969in}{1.653478in}}{\pgfqpoint{1.831383in}{1.655791in}}{\pgfqpoint{1.825559in}{1.655791in}}%
\pgfpathcurveto{\pgfqpoint{1.819735in}{1.655791in}}{\pgfqpoint{1.814149in}{1.653478in}}{\pgfqpoint{1.810031in}{1.649359in}}%
\pgfpathcurveto{\pgfqpoint{1.805913in}{1.645241in}}{\pgfqpoint{1.803599in}{1.639655in}}{\pgfqpoint{1.803599in}{1.633831in}}%
\pgfpathcurveto{\pgfqpoint{1.803599in}{1.628007in}}{\pgfqpoint{1.805913in}{1.622421in}}{\pgfqpoint{1.810031in}{1.618303in}}%
\pgfpathcurveto{\pgfqpoint{1.814149in}{1.614185in}}{\pgfqpoint{1.819735in}{1.611871in}}{\pgfqpoint{1.825559in}{1.611871in}}%
\pgfpathlineto{\pgfqpoint{1.825559in}{1.611871in}}%
\pgfpathclose%
\pgfusepath{stroke,fill}%
\end{pgfscope}%
\begin{pgfscope}%
\pgfpathrectangle{\pgfqpoint{0.100000in}{0.209042in}}{\pgfqpoint{2.906250in}{2.906250in}}%
\pgfusepath{clip}%
\pgfsetbuttcap%
\pgfsetroundjoin%
\definecolor{currentfill}{rgb}{0.378211,0.095332,0.500067}%
\pgfsetfillcolor{currentfill}%
\pgfsetfillopacity{0.800000}%
\pgfsetlinewidth{1.003750pt}%
\definecolor{currentstroke}{rgb}{0.378211,0.095332,0.500067}%
\pgfsetstrokecolor{currentstroke}%
\pgfsetstrokeopacity{0.800000}%
\pgfsetdash{}{0pt}%
\pgfpathmoveto{\pgfqpoint{1.898795in}{1.872098in}}%
\pgfpathcurveto{\pgfqpoint{1.904619in}{1.872098in}}{\pgfqpoint{1.910205in}{1.874411in}}{\pgfqpoint{1.914323in}{1.878530in}}%
\pgfpathcurveto{\pgfqpoint{1.918442in}{1.882648in}}{\pgfqpoint{1.920755in}{1.888234in}}{\pgfqpoint{1.920755in}{1.894058in}}%
\pgfpathcurveto{\pgfqpoint{1.920755in}{1.899882in}}{\pgfqpoint{1.918442in}{1.905468in}}{\pgfqpoint{1.914323in}{1.909586in}}%
\pgfpathcurveto{\pgfqpoint{1.910205in}{1.913704in}}{\pgfqpoint{1.904619in}{1.916018in}}{\pgfqpoint{1.898795in}{1.916018in}}%
\pgfpathcurveto{\pgfqpoint{1.892971in}{1.916018in}}{\pgfqpoint{1.887385in}{1.913704in}}{\pgfqpoint{1.883267in}{1.909586in}}%
\pgfpathcurveto{\pgfqpoint{1.879149in}{1.905468in}}{\pgfqpoint{1.876835in}{1.899882in}}{\pgfqpoint{1.876835in}{1.894058in}}%
\pgfpathcurveto{\pgfqpoint{1.876835in}{1.888234in}}{\pgfqpoint{1.879149in}{1.882648in}}{\pgfqpoint{1.883267in}{1.878530in}}%
\pgfpathcurveto{\pgfqpoint{1.887385in}{1.874411in}}{\pgfqpoint{1.892971in}{1.872098in}}{\pgfqpoint{1.898795in}{1.872098in}}%
\pgfpathlineto{\pgfqpoint{1.898795in}{1.872098in}}%
\pgfpathclose%
\pgfusepath{stroke,fill}%
\end{pgfscope}%
\begin{pgfscope}%
\pgfpathrectangle{\pgfqpoint{0.100000in}{0.209042in}}{\pgfqpoint{2.906250in}{2.906250in}}%
\pgfusepath{clip}%
\pgfsetbuttcap%
\pgfsetroundjoin%
\definecolor{currentfill}{rgb}{0.463508,0.129893,0.507652}%
\pgfsetfillcolor{currentfill}%
\pgfsetfillopacity{0.800000}%
\pgfsetlinewidth{1.003750pt}%
\definecolor{currentstroke}{rgb}{0.463508,0.129893,0.507652}%
\pgfsetstrokecolor{currentstroke}%
\pgfsetstrokeopacity{0.800000}%
\pgfsetdash{}{0pt}%
\pgfpathmoveto{\pgfqpoint{1.569532in}{1.385177in}}%
\pgfpathcurveto{\pgfqpoint{1.575356in}{1.385177in}}{\pgfqpoint{1.580943in}{1.387491in}}{\pgfqpoint{1.585061in}{1.391609in}}%
\pgfpathcurveto{\pgfqpoint{1.589179in}{1.395728in}}{\pgfqpoint{1.591493in}{1.401314in}}{\pgfqpoint{1.591493in}{1.407138in}}%
\pgfpathcurveto{\pgfqpoint{1.591493in}{1.412962in}}{\pgfqpoint{1.589179in}{1.418548in}}{\pgfqpoint{1.585061in}{1.422666in}}%
\pgfpathcurveto{\pgfqpoint{1.580943in}{1.426784in}}{\pgfqpoint{1.575356in}{1.429098in}}{\pgfqpoint{1.569532in}{1.429098in}}%
\pgfpathcurveto{\pgfqpoint{1.563709in}{1.429098in}}{\pgfqpoint{1.558122in}{1.426784in}}{\pgfqpoint{1.554004in}{1.422666in}}%
\pgfpathcurveto{\pgfqpoint{1.549886in}{1.418548in}}{\pgfqpoint{1.547572in}{1.412962in}}{\pgfqpoint{1.547572in}{1.407138in}}%
\pgfpathcurveto{\pgfqpoint{1.547572in}{1.401314in}}{\pgfqpoint{1.549886in}{1.395728in}}{\pgfqpoint{1.554004in}{1.391609in}}%
\pgfpathcurveto{\pgfqpoint{1.558122in}{1.387491in}}{\pgfqpoint{1.563709in}{1.385177in}}{\pgfqpoint{1.569532in}{1.385177in}}%
\pgfpathlineto{\pgfqpoint{1.569532in}{1.385177in}}%
\pgfpathclose%
\pgfusepath{stroke,fill}%
\end{pgfscope}%
\begin{pgfscope}%
\pgfpathrectangle{\pgfqpoint{0.100000in}{0.209042in}}{\pgfqpoint{2.906250in}{2.906250in}}%
\pgfusepath{clip}%
\pgfsetbuttcap%
\pgfsetroundjoin%
\definecolor{currentfill}{rgb}{0.347636,0.082946,0.494121}%
\pgfsetfillcolor{currentfill}%
\pgfsetfillopacity{0.800000}%
\pgfsetlinewidth{1.003750pt}%
\definecolor{currentstroke}{rgb}{0.347636,0.082946,0.494121}%
\pgfsetstrokecolor{currentstroke}%
\pgfsetstrokeopacity{0.800000}%
\pgfsetdash{}{0pt}%
\pgfpathmoveto{\pgfqpoint{1.804724in}{1.926956in}}%
\pgfpathcurveto{\pgfqpoint{1.810548in}{1.926956in}}{\pgfqpoint{1.816134in}{1.929270in}}{\pgfqpoint{1.820252in}{1.933388in}}%
\pgfpathcurveto{\pgfqpoint{1.824370in}{1.937506in}}{\pgfqpoint{1.826684in}{1.943092in}}{\pgfqpoint{1.826684in}{1.948916in}}%
\pgfpathcurveto{\pgfqpoint{1.826684in}{1.954740in}}{\pgfqpoint{1.824370in}{1.960326in}}{\pgfqpoint{1.820252in}{1.964445in}}%
\pgfpathcurveto{\pgfqpoint{1.816134in}{1.968563in}}{\pgfqpoint{1.810548in}{1.970877in}}{\pgfqpoint{1.804724in}{1.970877in}}%
\pgfpathcurveto{\pgfqpoint{1.798900in}{1.970877in}}{\pgfqpoint{1.793314in}{1.968563in}}{\pgfqpoint{1.789196in}{1.964445in}}%
\pgfpathcurveto{\pgfqpoint{1.785078in}{1.960326in}}{\pgfqpoint{1.782764in}{1.954740in}}{\pgfqpoint{1.782764in}{1.948916in}}%
\pgfpathcurveto{\pgfqpoint{1.782764in}{1.943092in}}{\pgfqpoint{1.785078in}{1.937506in}}{\pgfqpoint{1.789196in}{1.933388in}}%
\pgfpathcurveto{\pgfqpoint{1.793314in}{1.929270in}}{\pgfqpoint{1.798900in}{1.926956in}}{\pgfqpoint{1.804724in}{1.926956in}}%
\pgfpathlineto{\pgfqpoint{1.804724in}{1.926956in}}%
\pgfpathclose%
\pgfusepath{stroke,fill}%
\end{pgfscope}%
\begin{pgfscope}%
\pgfpathrectangle{\pgfqpoint{0.100000in}{0.209042in}}{\pgfqpoint{2.906250in}{2.906250in}}%
\pgfusepath{clip}%
\pgfsetbuttcap%
\pgfsetroundjoin%
\definecolor{currentfill}{rgb}{0.594508,0.175701,0.501241}%
\pgfsetfillcolor{currentfill}%
\pgfsetfillopacity{0.800000}%
\pgfsetlinewidth{1.003750pt}%
\definecolor{currentstroke}{rgb}{0.594508,0.175701,0.501241}%
\pgfsetstrokecolor{currentstroke}%
\pgfsetstrokeopacity{0.800000}%
\pgfsetdash{}{0pt}%
\pgfpathmoveto{\pgfqpoint{1.387531in}{2.072605in}}%
\pgfpathcurveto{\pgfqpoint{1.393355in}{2.072605in}}{\pgfqpoint{1.398941in}{2.074919in}}{\pgfqpoint{1.403059in}{2.079037in}}%
\pgfpathcurveto{\pgfqpoint{1.407177in}{2.083155in}}{\pgfqpoint{1.409491in}{2.088741in}}{\pgfqpoint{1.409491in}{2.094565in}}%
\pgfpathcurveto{\pgfqpoint{1.409491in}{2.100389in}}{\pgfqpoint{1.407177in}{2.105975in}}{\pgfqpoint{1.403059in}{2.110093in}}%
\pgfpathcurveto{\pgfqpoint{1.398941in}{2.114211in}}{\pgfqpoint{1.393355in}{2.116525in}}{\pgfqpoint{1.387531in}{2.116525in}}%
\pgfpathcurveto{\pgfqpoint{1.381707in}{2.116525in}}{\pgfqpoint{1.376121in}{2.114211in}}{\pgfqpoint{1.372003in}{2.110093in}}%
\pgfpathcurveto{\pgfqpoint{1.367884in}{2.105975in}}{\pgfqpoint{1.365571in}{2.100389in}}{\pgfqpoint{1.365571in}{2.094565in}}%
\pgfpathcurveto{\pgfqpoint{1.365571in}{2.088741in}}{\pgfqpoint{1.367884in}{2.083155in}}{\pgfqpoint{1.372003in}{2.079037in}}%
\pgfpathcurveto{\pgfqpoint{1.376121in}{2.074919in}}{\pgfqpoint{1.381707in}{2.072605in}}{\pgfqpoint{1.387531in}{2.072605in}}%
\pgfpathlineto{\pgfqpoint{1.387531in}{2.072605in}}%
\pgfpathclose%
\pgfusepath{stroke,fill}%
\end{pgfscope}%
\begin{pgfscope}%
\pgfpathrectangle{\pgfqpoint{0.100000in}{0.209042in}}{\pgfqpoint{2.906250in}{2.906250in}}%
\pgfusepath{clip}%
\pgfsetbuttcap%
\pgfsetroundjoin%
\definecolor{currentfill}{rgb}{0.123833,0.067295,0.295879}%
\pgfsetfillcolor{currentfill}%
\pgfsetfillopacity{0.800000}%
\pgfsetlinewidth{1.003750pt}%
\definecolor{currentstroke}{rgb}{0.123833,0.067295,0.295879}%
\pgfsetstrokecolor{currentstroke}%
\pgfsetstrokeopacity{0.800000}%
\pgfsetdash{}{0pt}%
\pgfpathmoveto{\pgfqpoint{1.503462in}{1.835325in}}%
\pgfpathcurveto{\pgfqpoint{1.509286in}{1.835325in}}{\pgfqpoint{1.514872in}{1.837638in}}{\pgfqpoint{1.518990in}{1.841757in}}%
\pgfpathcurveto{\pgfqpoint{1.523108in}{1.845875in}}{\pgfqpoint{1.525422in}{1.851461in}}{\pgfqpoint{1.525422in}{1.857285in}}%
\pgfpathcurveto{\pgfqpoint{1.525422in}{1.863109in}}{\pgfqpoint{1.523108in}{1.868695in}}{\pgfqpoint{1.518990in}{1.872813in}}%
\pgfpathcurveto{\pgfqpoint{1.514872in}{1.876931in}}{\pgfqpoint{1.509286in}{1.879245in}}{\pgfqpoint{1.503462in}{1.879245in}}%
\pgfpathcurveto{\pgfqpoint{1.497638in}{1.879245in}}{\pgfqpoint{1.492052in}{1.876931in}}{\pgfqpoint{1.487934in}{1.872813in}}%
\pgfpathcurveto{\pgfqpoint{1.483816in}{1.868695in}}{\pgfqpoint{1.481502in}{1.863109in}}{\pgfqpoint{1.481502in}{1.857285in}}%
\pgfpathcurveto{\pgfqpoint{1.481502in}{1.851461in}}{\pgfqpoint{1.483816in}{1.845875in}}{\pgfqpoint{1.487934in}{1.841757in}}%
\pgfpathcurveto{\pgfqpoint{1.492052in}{1.837638in}}{\pgfqpoint{1.497638in}{1.835325in}}{\pgfqpoint{1.503462in}{1.835325in}}%
\pgfpathlineto{\pgfqpoint{1.503462in}{1.835325in}}%
\pgfpathclose%
\pgfusepath{stroke,fill}%
\end{pgfscope}%
\begin{pgfscope}%
\pgfpathrectangle{\pgfqpoint{0.100000in}{0.209042in}}{\pgfqpoint{2.906250in}{2.906250in}}%
\pgfusepath{clip}%
\pgfsetbuttcap%
\pgfsetroundjoin%
\definecolor{currentfill}{rgb}{0.645633,0.191952,0.492910}%
\pgfsetfillcolor{currentfill}%
\pgfsetfillopacity{0.800000}%
\pgfsetlinewidth{1.003750pt}%
\definecolor{currentstroke}{rgb}{0.645633,0.191952,0.492910}%
\pgfsetstrokecolor{currentstroke}%
\pgfsetstrokeopacity{0.800000}%
\pgfsetdash{}{0pt}%
\pgfpathmoveto{\pgfqpoint{1.229890in}{2.012276in}}%
\pgfpathcurveto{\pgfqpoint{1.235714in}{2.012276in}}{\pgfqpoint{1.241300in}{2.014590in}}{\pgfqpoint{1.245418in}{2.018708in}}%
\pgfpathcurveto{\pgfqpoint{1.249536in}{2.022826in}}{\pgfqpoint{1.251850in}{2.028412in}}{\pgfqpoint{1.251850in}{2.034236in}}%
\pgfpathcurveto{\pgfqpoint{1.251850in}{2.040060in}}{\pgfqpoint{1.249536in}{2.045646in}}{\pgfqpoint{1.245418in}{2.049765in}}%
\pgfpathcurveto{\pgfqpoint{1.241300in}{2.053883in}}{\pgfqpoint{1.235714in}{2.056197in}}{\pgfqpoint{1.229890in}{2.056197in}}%
\pgfpathcurveto{\pgfqpoint{1.224066in}{2.056197in}}{\pgfqpoint{1.218479in}{2.053883in}}{\pgfqpoint{1.214361in}{2.049765in}}%
\pgfpathcurveto{\pgfqpoint{1.210243in}{2.045646in}}{\pgfqpoint{1.207929in}{2.040060in}}{\pgfqpoint{1.207929in}{2.034236in}}%
\pgfpathcurveto{\pgfqpoint{1.207929in}{2.028412in}}{\pgfqpoint{1.210243in}{2.022826in}}{\pgfqpoint{1.214361in}{2.018708in}}%
\pgfpathcurveto{\pgfqpoint{1.218479in}{2.014590in}}{\pgfqpoint{1.224066in}{2.012276in}}{\pgfqpoint{1.229890in}{2.012276in}}%
\pgfpathlineto{\pgfqpoint{1.229890in}{2.012276in}}%
\pgfpathclose%
\pgfusepath{stroke,fill}%
\end{pgfscope}%
\begin{pgfscope}%
\pgfpathrectangle{\pgfqpoint{0.100000in}{0.209042in}}{\pgfqpoint{2.906250in}{2.906250in}}%
\pgfusepath{clip}%
\pgfsetbuttcap%
\pgfsetroundjoin%
\definecolor{currentfill}{rgb}{0.043830,0.033830,0.141886}%
\pgfsetfillcolor{currentfill}%
\pgfsetfillopacity{0.800000}%
\pgfsetlinewidth{1.003750pt}%
\definecolor{currentstroke}{rgb}{0.043830,0.033830,0.141886}%
\pgfsetstrokecolor{currentstroke}%
\pgfsetstrokeopacity{0.800000}%
\pgfsetdash{}{0pt}%
\pgfpathmoveto{\pgfqpoint{1.711942in}{1.767372in}}%
\pgfpathcurveto{\pgfqpoint{1.717766in}{1.767372in}}{\pgfqpoint{1.723352in}{1.769686in}}{\pgfqpoint{1.727470in}{1.773804in}}%
\pgfpathcurveto{\pgfqpoint{1.731589in}{1.777923in}}{\pgfqpoint{1.733902in}{1.783509in}}{\pgfqpoint{1.733902in}{1.789333in}}%
\pgfpathcurveto{\pgfqpoint{1.733902in}{1.795157in}}{\pgfqpoint{1.731589in}{1.800743in}}{\pgfqpoint{1.727470in}{1.804861in}}%
\pgfpathcurveto{\pgfqpoint{1.723352in}{1.808979in}}{\pgfqpoint{1.717766in}{1.811293in}}{\pgfqpoint{1.711942in}{1.811293in}}%
\pgfpathcurveto{\pgfqpoint{1.706118in}{1.811293in}}{\pgfqpoint{1.700532in}{1.808979in}}{\pgfqpoint{1.696414in}{1.804861in}}%
\pgfpathcurveto{\pgfqpoint{1.692296in}{1.800743in}}{\pgfqpoint{1.689982in}{1.795157in}}{\pgfqpoint{1.689982in}{1.789333in}}%
\pgfpathcurveto{\pgfqpoint{1.689982in}{1.783509in}}{\pgfqpoint{1.692296in}{1.777923in}}{\pgfqpoint{1.696414in}{1.773804in}}%
\pgfpathcurveto{\pgfqpoint{1.700532in}{1.769686in}}{\pgfqpoint{1.706118in}{1.767372in}}{\pgfqpoint{1.711942in}{1.767372in}}%
\pgfpathlineto{\pgfqpoint{1.711942in}{1.767372in}}%
\pgfpathclose%
\pgfusepath{stroke,fill}%
\end{pgfscope}%
\begin{pgfscope}%
\pgfpathrectangle{\pgfqpoint{0.100000in}{0.209042in}}{\pgfqpoint{2.906250in}{2.906250in}}%
\pgfusepath{clip}%
\pgfsetbuttcap%
\pgfsetroundjoin%
\definecolor{currentfill}{rgb}{0.123833,0.067295,0.295879}%
\pgfsetfillcolor{currentfill}%
\pgfsetfillopacity{0.800000}%
\pgfsetlinewidth{1.003750pt}%
\definecolor{currentstroke}{rgb}{0.123833,0.067295,0.295879}%
\pgfsetstrokecolor{currentstroke}%
\pgfsetstrokeopacity{0.800000}%
\pgfsetdash{}{0pt}%
\pgfpathmoveto{\pgfqpoint{1.756538in}{1.597230in}}%
\pgfpathcurveto{\pgfqpoint{1.762362in}{1.597230in}}{\pgfqpoint{1.767949in}{1.599544in}}{\pgfqpoint{1.772067in}{1.603662in}}%
\pgfpathcurveto{\pgfqpoint{1.776185in}{1.607780in}}{\pgfqpoint{1.778499in}{1.613367in}}{\pgfqpoint{1.778499in}{1.619191in}}%
\pgfpathcurveto{\pgfqpoint{1.778499in}{1.625014in}}{\pgfqpoint{1.776185in}{1.630601in}}{\pgfqpoint{1.772067in}{1.634719in}}%
\pgfpathcurveto{\pgfqpoint{1.767949in}{1.638837in}}{\pgfqpoint{1.762362in}{1.641151in}}{\pgfqpoint{1.756538in}{1.641151in}}%
\pgfpathcurveto{\pgfqpoint{1.750715in}{1.641151in}}{\pgfqpoint{1.745128in}{1.638837in}}{\pgfqpoint{1.741010in}{1.634719in}}%
\pgfpathcurveto{\pgfqpoint{1.736892in}{1.630601in}}{\pgfqpoint{1.734578in}{1.625014in}}{\pgfqpoint{1.734578in}{1.619191in}}%
\pgfpathcurveto{\pgfqpoint{1.734578in}{1.613367in}}{\pgfqpoint{1.736892in}{1.607780in}}{\pgfqpoint{1.741010in}{1.603662in}}%
\pgfpathcurveto{\pgfqpoint{1.745128in}{1.599544in}}{\pgfqpoint{1.750715in}{1.597230in}}{\pgfqpoint{1.756538in}{1.597230in}}%
\pgfpathlineto{\pgfqpoint{1.756538in}{1.597230in}}%
\pgfpathclose%
\pgfusepath{stroke,fill}%
\end{pgfscope}%
\begin{pgfscope}%
\pgfpathrectangle{\pgfqpoint{0.100000in}{0.209042in}}{\pgfqpoint{2.906250in}{2.906250in}}%
\pgfusepath{clip}%
\pgfsetbuttcap%
\pgfsetroundjoin%
\definecolor{currentfill}{rgb}{0.031696,0.025765,0.116965}%
\pgfsetfillcolor{currentfill}%
\pgfsetfillopacity{0.800000}%
\pgfsetlinewidth{1.003750pt}%
\definecolor{currentstroke}{rgb}{0.031696,0.025765,0.116965}%
\pgfsetstrokecolor{currentstroke}%
\pgfsetstrokeopacity{0.800000}%
\pgfsetdash{}{0pt}%
\pgfpathmoveto{\pgfqpoint{1.532472in}{1.670616in}}%
\pgfpathcurveto{\pgfqpoint{1.538296in}{1.670616in}}{\pgfqpoint{1.543882in}{1.672930in}}{\pgfqpoint{1.548000in}{1.677048in}}%
\pgfpathcurveto{\pgfqpoint{1.552118in}{1.681166in}}{\pgfqpoint{1.554432in}{1.686752in}}{\pgfqpoint{1.554432in}{1.692576in}}%
\pgfpathcurveto{\pgfqpoint{1.554432in}{1.698400in}}{\pgfqpoint{1.552118in}{1.703986in}}{\pgfqpoint{1.548000in}{1.708104in}}%
\pgfpathcurveto{\pgfqpoint{1.543882in}{1.712223in}}{\pgfqpoint{1.538296in}{1.714536in}}{\pgfqpoint{1.532472in}{1.714536in}}%
\pgfpathcurveto{\pgfqpoint{1.526648in}{1.714536in}}{\pgfqpoint{1.521062in}{1.712223in}}{\pgfqpoint{1.516944in}{1.708104in}}%
\pgfpathcurveto{\pgfqpoint{1.512826in}{1.703986in}}{\pgfqpoint{1.510512in}{1.698400in}}{\pgfqpoint{1.510512in}{1.692576in}}%
\pgfpathcurveto{\pgfqpoint{1.510512in}{1.686752in}}{\pgfqpoint{1.512826in}{1.681166in}}{\pgfqpoint{1.516944in}{1.677048in}}%
\pgfpathcurveto{\pgfqpoint{1.521062in}{1.672930in}}{\pgfqpoint{1.526648in}{1.670616in}}{\pgfqpoint{1.532472in}{1.670616in}}%
\pgfpathlineto{\pgfqpoint{1.532472in}{1.670616in}}%
\pgfpathclose%
\pgfusepath{stroke,fill}%
\end{pgfscope}%
\begin{pgfscope}%
\pgfpathrectangle{\pgfqpoint{0.100000in}{0.209042in}}{\pgfqpoint{2.906250in}{2.906250in}}%
\pgfusepath{clip}%
\pgfsetbuttcap%
\pgfsetroundjoin%
\definecolor{currentfill}{rgb}{0.359898,0.087831,0.496778}%
\pgfsetfillcolor{currentfill}%
\pgfsetfillopacity{0.800000}%
\pgfsetlinewidth{1.003750pt}%
\definecolor{currentstroke}{rgb}{0.359898,0.087831,0.496778}%
\pgfsetstrokecolor{currentstroke}%
\pgfsetstrokeopacity{0.800000}%
\pgfsetdash{}{0pt}%
\pgfpathmoveto{\pgfqpoint{1.827917in}{1.479375in}}%
\pgfpathcurveto{\pgfqpoint{1.833741in}{1.479375in}}{\pgfqpoint{1.839328in}{1.481689in}}{\pgfqpoint{1.843446in}{1.485807in}}%
\pgfpathcurveto{\pgfqpoint{1.847564in}{1.489925in}}{\pgfqpoint{1.849878in}{1.495511in}}{\pgfqpoint{1.849878in}{1.501335in}}%
\pgfpathcurveto{\pgfqpoint{1.849878in}{1.507159in}}{\pgfqpoint{1.847564in}{1.512746in}}{\pgfqpoint{1.843446in}{1.516864in}}%
\pgfpathcurveto{\pgfqpoint{1.839328in}{1.520982in}}{\pgfqpoint{1.833741in}{1.523296in}}{\pgfqpoint{1.827917in}{1.523296in}}%
\pgfpathcurveto{\pgfqpoint{1.822094in}{1.523296in}}{\pgfqpoint{1.816507in}{1.520982in}}{\pgfqpoint{1.812389in}{1.516864in}}%
\pgfpathcurveto{\pgfqpoint{1.808271in}{1.512746in}}{\pgfqpoint{1.805957in}{1.507159in}}{\pgfqpoint{1.805957in}{1.501335in}}%
\pgfpathcurveto{\pgfqpoint{1.805957in}{1.495511in}}{\pgfqpoint{1.808271in}{1.489925in}}{\pgfqpoint{1.812389in}{1.485807in}}%
\pgfpathcurveto{\pgfqpoint{1.816507in}{1.481689in}}{\pgfqpoint{1.822094in}{1.479375in}}{\pgfqpoint{1.827917in}{1.479375in}}%
\pgfpathlineto{\pgfqpoint{1.827917in}{1.479375in}}%
\pgfpathclose%
\pgfusepath{stroke,fill}%
\end{pgfscope}%
\begin{pgfscope}%
\pgfpathrectangle{\pgfqpoint{0.100000in}{0.209042in}}{\pgfqpoint{2.906250in}{2.906250in}}%
\pgfusepath{clip}%
\pgfsetbuttcap%
\pgfsetroundjoin%
\definecolor{currentfill}{rgb}{0.469640,0.132245,0.507809}%
\pgfsetfillcolor{currentfill}%
\pgfsetfillopacity{0.800000}%
\pgfsetlinewidth{1.003750pt}%
\definecolor{currentstroke}{rgb}{0.469640,0.132245,0.507809}%
\pgfsetstrokecolor{currentstroke}%
\pgfsetstrokeopacity{0.800000}%
\pgfsetdash{}{0pt}%
\pgfpathmoveto{\pgfqpoint{1.844945in}{1.980148in}}%
\pgfpathcurveto{\pgfqpoint{1.850768in}{1.980148in}}{\pgfqpoint{1.856355in}{1.982462in}}{\pgfqpoint{1.860473in}{1.986580in}}%
\pgfpathcurveto{\pgfqpoint{1.864591in}{1.990698in}}{\pgfqpoint{1.866905in}{1.996284in}}{\pgfqpoint{1.866905in}{2.002108in}}%
\pgfpathcurveto{\pgfqpoint{1.866905in}{2.007932in}}{\pgfqpoint{1.864591in}{2.013518in}}{\pgfqpoint{1.860473in}{2.017637in}}%
\pgfpathcurveto{\pgfqpoint{1.856355in}{2.021755in}}{\pgfqpoint{1.850768in}{2.024069in}}{\pgfqpoint{1.844945in}{2.024069in}}%
\pgfpathcurveto{\pgfqpoint{1.839121in}{2.024069in}}{\pgfqpoint{1.833534in}{2.021755in}}{\pgfqpoint{1.829416in}{2.017637in}}%
\pgfpathcurveto{\pgfqpoint{1.825298in}{2.013518in}}{\pgfqpoint{1.822984in}{2.007932in}}{\pgfqpoint{1.822984in}{2.002108in}}%
\pgfpathcurveto{\pgfqpoint{1.822984in}{1.996284in}}{\pgfqpoint{1.825298in}{1.990698in}}{\pgfqpoint{1.829416in}{1.986580in}}%
\pgfpathcurveto{\pgfqpoint{1.833534in}{1.982462in}}{\pgfqpoint{1.839121in}{1.980148in}}{\pgfqpoint{1.844945in}{1.980148in}}%
\pgfpathlineto{\pgfqpoint{1.844945in}{1.980148in}}%
\pgfpathclose%
\pgfusepath{stroke,fill}%
\end{pgfscope}%
\begin{pgfscope}%
\pgfpathrectangle{\pgfqpoint{0.100000in}{0.209042in}}{\pgfqpoint{2.906250in}{2.906250in}}%
\pgfusepath{clip}%
\pgfsetbuttcap%
\pgfsetroundjoin%
\definecolor{currentfill}{rgb}{0.607238,0.179779,0.499492}%
\pgfsetfillcolor{currentfill}%
\pgfsetfillopacity{0.800000}%
\pgfsetlinewidth{1.003750pt}%
\definecolor{currentstroke}{rgb}{0.607238,0.179779,0.499492}%
\pgfsetstrokecolor{currentstroke}%
\pgfsetstrokeopacity{0.800000}%
\pgfsetdash{}{0pt}%
\pgfpathmoveto{\pgfqpoint{1.123938in}{1.783410in}}%
\pgfpathcurveto{\pgfqpoint{1.129762in}{1.783410in}}{\pgfqpoint{1.135348in}{1.785723in}}{\pgfqpoint{1.139466in}{1.789842in}}%
\pgfpathcurveto{\pgfqpoint{1.143584in}{1.793960in}}{\pgfqpoint{1.145898in}{1.799546in}}{\pgfqpoint{1.145898in}{1.805370in}}%
\pgfpathcurveto{\pgfqpoint{1.145898in}{1.811194in}}{\pgfqpoint{1.143584in}{1.816780in}}{\pgfqpoint{1.139466in}{1.820898in}}%
\pgfpathcurveto{\pgfqpoint{1.135348in}{1.825016in}}{\pgfqpoint{1.129762in}{1.827330in}}{\pgfqpoint{1.123938in}{1.827330in}}%
\pgfpathcurveto{\pgfqpoint{1.118114in}{1.827330in}}{\pgfqpoint{1.112527in}{1.825016in}}{\pgfqpoint{1.108409in}{1.820898in}}%
\pgfpathcurveto{\pgfqpoint{1.104291in}{1.816780in}}{\pgfqpoint{1.101977in}{1.811194in}}{\pgfqpoint{1.101977in}{1.805370in}}%
\pgfpathcurveto{\pgfqpoint{1.101977in}{1.799546in}}{\pgfqpoint{1.104291in}{1.793960in}}{\pgfqpoint{1.108409in}{1.789842in}}%
\pgfpathcurveto{\pgfqpoint{1.112527in}{1.785723in}}{\pgfqpoint{1.118114in}{1.783410in}}{\pgfqpoint{1.123938in}{1.783410in}}%
\pgfpathlineto{\pgfqpoint{1.123938in}{1.783410in}}%
\pgfpathclose%
\pgfusepath{stroke,fill}%
\end{pgfscope}%
\begin{pgfscope}%
\pgfpathrectangle{\pgfqpoint{0.100000in}{0.209042in}}{\pgfqpoint{2.906250in}{2.906250in}}%
\pgfusepath{clip}%
\pgfsetbuttcap%
\pgfsetroundjoin%
\definecolor{currentfill}{rgb}{0.519045,0.150383,0.507443}%
\pgfsetfillcolor{currentfill}%
\pgfsetfillopacity{0.800000}%
\pgfsetlinewidth{1.003750pt}%
\definecolor{currentstroke}{rgb}{0.519045,0.150383,0.507443}%
\pgfsetstrokecolor{currentstroke}%
\pgfsetstrokeopacity{0.800000}%
\pgfsetdash{}{0pt}%
\pgfpathmoveto{\pgfqpoint{2.024329in}{1.539239in}}%
\pgfpathcurveto{\pgfqpoint{2.030153in}{1.539239in}}{\pgfqpoint{2.035739in}{1.541553in}}{\pgfqpoint{2.039857in}{1.545671in}}%
\pgfpathcurveto{\pgfqpoint{2.043976in}{1.549789in}}{\pgfqpoint{2.046289in}{1.555375in}}{\pgfqpoint{2.046289in}{1.561199in}}%
\pgfpathcurveto{\pgfqpoint{2.046289in}{1.567023in}}{\pgfqpoint{2.043976in}{1.572609in}}{\pgfqpoint{2.039857in}{1.576728in}}%
\pgfpathcurveto{\pgfqpoint{2.035739in}{1.580846in}}{\pgfqpoint{2.030153in}{1.583160in}}{\pgfqpoint{2.024329in}{1.583160in}}%
\pgfpathcurveto{\pgfqpoint{2.018505in}{1.583160in}}{\pgfqpoint{2.012919in}{1.580846in}}{\pgfqpoint{2.008801in}{1.576728in}}%
\pgfpathcurveto{\pgfqpoint{2.004683in}{1.572609in}}{\pgfqpoint{2.002369in}{1.567023in}}{\pgfqpoint{2.002369in}{1.561199in}}%
\pgfpathcurveto{\pgfqpoint{2.002369in}{1.555375in}}{\pgfqpoint{2.004683in}{1.549789in}}{\pgfqpoint{2.008801in}{1.545671in}}%
\pgfpathcurveto{\pgfqpoint{2.012919in}{1.541553in}}{\pgfqpoint{2.018505in}{1.539239in}}{\pgfqpoint{2.024329in}{1.539239in}}%
\pgfpathlineto{\pgfqpoint{2.024329in}{1.539239in}}%
\pgfpathclose%
\pgfusepath{stroke,fill}%
\end{pgfscope}%
\begin{pgfscope}%
\pgfpathrectangle{\pgfqpoint{0.100000in}{0.209042in}}{\pgfqpoint{2.906250in}{2.906250in}}%
\pgfusepath{clip}%
\pgfsetbuttcap%
\pgfsetroundjoin%
\definecolor{currentfill}{rgb}{0.822926,0.259016,0.433573}%
\pgfsetfillcolor{currentfill}%
\pgfsetfillopacity{0.800000}%
\pgfsetlinewidth{1.003750pt}%
\definecolor{currentstroke}{rgb}{0.822926,0.259016,0.433573}%
\pgfsetstrokecolor{currentstroke}%
\pgfsetstrokeopacity{0.800000}%
\pgfsetdash{}{0pt}%
\pgfpathmoveto{\pgfqpoint{0.994351in}{1.826780in}}%
\pgfpathcurveto{\pgfqpoint{1.000175in}{1.826780in}}{\pgfqpoint{1.005761in}{1.829094in}}{\pgfqpoint{1.009879in}{1.833212in}}%
\pgfpathcurveto{\pgfqpoint{1.013997in}{1.837331in}}{\pgfqpoint{1.016311in}{1.842917in}}{\pgfqpoint{1.016311in}{1.848741in}}%
\pgfpathcurveto{\pgfqpoint{1.016311in}{1.854565in}}{\pgfqpoint{1.013997in}{1.860151in}}{\pgfqpoint{1.009879in}{1.864269in}}%
\pgfpathcurveto{\pgfqpoint{1.005761in}{1.868387in}}{\pgfqpoint{1.000175in}{1.870701in}}{\pgfqpoint{0.994351in}{1.870701in}}%
\pgfpathcurveto{\pgfqpoint{0.988527in}{1.870701in}}{\pgfqpoint{0.982941in}{1.868387in}}{\pgfqpoint{0.978822in}{1.864269in}}%
\pgfpathcurveto{\pgfqpoint{0.974704in}{1.860151in}}{\pgfqpoint{0.972390in}{1.854565in}}{\pgfqpoint{0.972390in}{1.848741in}}%
\pgfpathcurveto{\pgfqpoint{0.972390in}{1.842917in}}{\pgfqpoint{0.974704in}{1.837331in}}{\pgfqpoint{0.978822in}{1.833212in}}%
\pgfpathcurveto{\pgfqpoint{0.982941in}{1.829094in}}{\pgfqpoint{0.988527in}{1.826780in}}{\pgfqpoint{0.994351in}{1.826780in}}%
\pgfpathlineto{\pgfqpoint{0.994351in}{1.826780in}}%
\pgfpathclose%
\pgfusepath{stroke,fill}%
\end{pgfscope}%
\begin{pgfscope}%
\pgfpathrectangle{\pgfqpoint{0.100000in}{0.209042in}}{\pgfqpoint{2.906250in}{2.906250in}}%
\pgfusepath{clip}%
\pgfsetbuttcap%
\pgfsetroundjoin%
\definecolor{currentfill}{rgb}{0.018815,0.016026,0.084584}%
\pgfsetfillcolor{currentfill}%
\pgfsetfillopacity{0.800000}%
\pgfsetlinewidth{1.003750pt}%
\definecolor{currentstroke}{rgb}{0.018815,0.016026,0.084584}%
\pgfsetstrokecolor{currentstroke}%
\pgfsetstrokeopacity{0.800000}%
\pgfsetdash{}{0pt}%
\pgfpathmoveto{\pgfqpoint{1.544126in}{1.744518in}}%
\pgfpathcurveto{\pgfqpoint{1.549950in}{1.744518in}}{\pgfqpoint{1.555536in}{1.746832in}}{\pgfqpoint{1.559654in}{1.750950in}}%
\pgfpathcurveto{\pgfqpoint{1.563772in}{1.755068in}}{\pgfqpoint{1.566086in}{1.760655in}}{\pgfqpoint{1.566086in}{1.766478in}}%
\pgfpathcurveto{\pgfqpoint{1.566086in}{1.772302in}}{\pgfqpoint{1.563772in}{1.777889in}}{\pgfqpoint{1.559654in}{1.782007in}}%
\pgfpathcurveto{\pgfqpoint{1.555536in}{1.786125in}}{\pgfqpoint{1.549950in}{1.788439in}}{\pgfqpoint{1.544126in}{1.788439in}}%
\pgfpathcurveto{\pgfqpoint{1.538302in}{1.788439in}}{\pgfqpoint{1.532716in}{1.786125in}}{\pgfqpoint{1.528598in}{1.782007in}}%
\pgfpathcurveto{\pgfqpoint{1.524480in}{1.777889in}}{\pgfqpoint{1.522166in}{1.772302in}}{\pgfqpoint{1.522166in}{1.766478in}}%
\pgfpathcurveto{\pgfqpoint{1.522166in}{1.760655in}}{\pgfqpoint{1.524480in}{1.755068in}}{\pgfqpoint{1.528598in}{1.750950in}}%
\pgfpathcurveto{\pgfqpoint{1.532716in}{1.746832in}}{\pgfqpoint{1.538302in}{1.744518in}}{\pgfqpoint{1.544126in}{1.744518in}}%
\pgfpathlineto{\pgfqpoint{1.544126in}{1.744518in}}%
\pgfpathclose%
\pgfusepath{stroke,fill}%
\end{pgfscope}%
\begin{pgfscope}%
\pgfpathrectangle{\pgfqpoint{0.100000in}{0.209042in}}{\pgfqpoint{2.906250in}{2.906250in}}%
\pgfusepath{clip}%
\pgfsetbuttcap%
\pgfsetroundjoin%
\definecolor{currentfill}{rgb}{0.432967,0.117855,0.506160}%
\pgfsetfillcolor{currentfill}%
\pgfsetfillopacity{0.800000}%
\pgfsetlinewidth{1.003750pt}%
\definecolor{currentstroke}{rgb}{0.432967,0.117855,0.506160}%
\pgfsetstrokecolor{currentstroke}%
\pgfsetstrokeopacity{0.800000}%
\pgfsetdash{}{0pt}%
\pgfpathmoveto{\pgfqpoint{1.662210in}{1.394704in}}%
\pgfpathcurveto{\pgfqpoint{1.668034in}{1.394704in}}{\pgfqpoint{1.673620in}{1.397018in}}{\pgfqpoint{1.677739in}{1.401136in}}%
\pgfpathcurveto{\pgfqpoint{1.681857in}{1.405254in}}{\pgfqpoint{1.684171in}{1.410840in}}{\pgfqpoint{1.684171in}{1.416664in}}%
\pgfpathcurveto{\pgfqpoint{1.684171in}{1.422488in}}{\pgfqpoint{1.681857in}{1.428074in}}{\pgfqpoint{1.677739in}{1.432192in}}%
\pgfpathcurveto{\pgfqpoint{1.673620in}{1.436311in}}{\pgfqpoint{1.668034in}{1.438624in}}{\pgfqpoint{1.662210in}{1.438624in}}%
\pgfpathcurveto{\pgfqpoint{1.656386in}{1.438624in}}{\pgfqpoint{1.650800in}{1.436311in}}{\pgfqpoint{1.646682in}{1.432192in}}%
\pgfpathcurveto{\pgfqpoint{1.642564in}{1.428074in}}{\pgfqpoint{1.640250in}{1.422488in}}{\pgfqpoint{1.640250in}{1.416664in}}%
\pgfpathcurveto{\pgfqpoint{1.640250in}{1.410840in}}{\pgfqpoint{1.642564in}{1.405254in}}{\pgfqpoint{1.646682in}{1.401136in}}%
\pgfpathcurveto{\pgfqpoint{1.650800in}{1.397018in}}{\pgfqpoint{1.656386in}{1.394704in}}{\pgfqpoint{1.662210in}{1.394704in}}%
\pgfpathlineto{\pgfqpoint{1.662210in}{1.394704in}}%
\pgfpathclose%
\pgfusepath{stroke,fill}%
\end{pgfscope}%
\begin{pgfscope}%
\pgfpathrectangle{\pgfqpoint{0.100000in}{0.209042in}}{\pgfqpoint{2.906250in}{2.906250in}}%
\pgfusepath{clip}%
\pgfsetbuttcap%
\pgfsetroundjoin%
\definecolor{currentfill}{rgb}{0.097833,0.061531,0.248477}%
\pgfsetfillcolor{currentfill}%
\pgfsetfillopacity{0.800000}%
\pgfsetlinewidth{1.003750pt}%
\definecolor{currentstroke}{rgb}{0.097833,0.061531,0.248477}%
\pgfsetstrokecolor{currentstroke}%
\pgfsetstrokeopacity{0.800000}%
\pgfsetdash{}{0pt}%
\pgfpathmoveto{\pgfqpoint{1.572259in}{1.843475in}}%
\pgfpathcurveto{\pgfqpoint{1.578082in}{1.843475in}}{\pgfqpoint{1.583669in}{1.845789in}}{\pgfqpoint{1.587787in}{1.849907in}}%
\pgfpathcurveto{\pgfqpoint{1.591905in}{1.854025in}}{\pgfqpoint{1.594219in}{1.859611in}}{\pgfqpoint{1.594219in}{1.865435in}}%
\pgfpathcurveto{\pgfqpoint{1.594219in}{1.871259in}}{\pgfqpoint{1.591905in}{1.876845in}}{\pgfqpoint{1.587787in}{1.880963in}}%
\pgfpathcurveto{\pgfqpoint{1.583669in}{1.885081in}}{\pgfqpoint{1.578082in}{1.887395in}}{\pgfqpoint{1.572259in}{1.887395in}}%
\pgfpathcurveto{\pgfqpoint{1.566435in}{1.887395in}}{\pgfqpoint{1.560848in}{1.885081in}}{\pgfqpoint{1.556730in}{1.880963in}}%
\pgfpathcurveto{\pgfqpoint{1.552612in}{1.876845in}}{\pgfqpoint{1.550298in}{1.871259in}}{\pgfqpoint{1.550298in}{1.865435in}}%
\pgfpathcurveto{\pgfqpoint{1.550298in}{1.859611in}}{\pgfqpoint{1.552612in}{1.854025in}}{\pgfqpoint{1.556730in}{1.849907in}}%
\pgfpathcurveto{\pgfqpoint{1.560848in}{1.845789in}}{\pgfqpoint{1.566435in}{1.843475in}}{\pgfqpoint{1.572259in}{1.843475in}}%
\pgfpathlineto{\pgfqpoint{1.572259in}{1.843475in}}%
\pgfpathclose%
\pgfusepath{stroke,fill}%
\end{pgfscope}%
\begin{pgfscope}%
\pgfpathrectangle{\pgfqpoint{0.100000in}{0.209042in}}{\pgfqpoint{2.906250in}{2.906250in}}%
\pgfusepath{clip}%
\pgfsetbuttcap%
\pgfsetroundjoin%
\definecolor{currentfill}{rgb}{0.384299,0.097855,0.501002}%
\pgfsetfillcolor{currentfill}%
\pgfsetfillopacity{0.800000}%
\pgfsetlinewidth{1.003750pt}%
\definecolor{currentstroke}{rgb}{0.384299,0.097855,0.501002}%
\pgfsetstrokecolor{currentstroke}%
\pgfsetstrokeopacity{0.800000}%
\pgfsetdash{}{0pt}%
\pgfpathmoveto{\pgfqpoint{1.903690in}{1.515135in}}%
\pgfpathcurveto{\pgfqpoint{1.909514in}{1.515135in}}{\pgfqpoint{1.915101in}{1.517448in}}{\pgfqpoint{1.919219in}{1.521567in}}%
\pgfpathcurveto{\pgfqpoint{1.923337in}{1.525685in}}{\pgfqpoint{1.925651in}{1.531271in}}{\pgfqpoint{1.925651in}{1.537095in}}%
\pgfpathcurveto{\pgfqpoint{1.925651in}{1.542919in}}{\pgfqpoint{1.923337in}{1.548505in}}{\pgfqpoint{1.919219in}{1.552623in}}%
\pgfpathcurveto{\pgfqpoint{1.915101in}{1.556741in}}{\pgfqpoint{1.909514in}{1.559055in}}{\pgfqpoint{1.903690in}{1.559055in}}%
\pgfpathcurveto{\pgfqpoint{1.897866in}{1.559055in}}{\pgfqpoint{1.892280in}{1.556741in}}{\pgfqpoint{1.888162in}{1.552623in}}%
\pgfpathcurveto{\pgfqpoint{1.884044in}{1.548505in}}{\pgfqpoint{1.881730in}{1.542919in}}{\pgfqpoint{1.881730in}{1.537095in}}%
\pgfpathcurveto{\pgfqpoint{1.881730in}{1.531271in}}{\pgfqpoint{1.884044in}{1.525685in}}{\pgfqpoint{1.888162in}{1.521567in}}%
\pgfpathcurveto{\pgfqpoint{1.892280in}{1.517448in}}{\pgfqpoint{1.897866in}{1.515135in}}{\pgfqpoint{1.903690in}{1.515135in}}%
\pgfpathlineto{\pgfqpoint{1.903690in}{1.515135in}}%
\pgfpathclose%
\pgfusepath{stroke,fill}%
\end{pgfscope}%
\begin{pgfscope}%
\pgfpathrectangle{\pgfqpoint{0.100000in}{0.209042in}}{\pgfqpoint{2.906250in}{2.906250in}}%
\pgfusepath{clip}%
\pgfsetbuttcap%
\pgfsetroundjoin%
\definecolor{currentfill}{rgb}{0.372116,0.092816,0.499053}%
\pgfsetfillcolor{currentfill}%
\pgfsetfillopacity{0.800000}%
\pgfsetlinewidth{1.003750pt}%
\definecolor{currentstroke}{rgb}{0.372116,0.092816,0.499053}%
\pgfsetstrokecolor{currentstroke}%
\pgfsetstrokeopacity{0.800000}%
\pgfsetdash{}{0pt}%
\pgfpathmoveto{\pgfqpoint{1.967967in}{1.728576in}}%
\pgfpathcurveto{\pgfqpoint{1.973791in}{1.728576in}}{\pgfqpoint{1.979378in}{1.730890in}}{\pgfqpoint{1.983496in}{1.735008in}}%
\pgfpathcurveto{\pgfqpoint{1.987614in}{1.739126in}}{\pgfqpoint{1.989928in}{1.744712in}}{\pgfqpoint{1.989928in}{1.750536in}}%
\pgfpathcurveto{\pgfqpoint{1.989928in}{1.756360in}}{\pgfqpoint{1.987614in}{1.761946in}}{\pgfqpoint{1.983496in}{1.766064in}}%
\pgfpathcurveto{\pgfqpoint{1.979378in}{1.770182in}}{\pgfqpoint{1.973791in}{1.772496in}}{\pgfqpoint{1.967967in}{1.772496in}}%
\pgfpathcurveto{\pgfqpoint{1.962144in}{1.772496in}}{\pgfqpoint{1.956557in}{1.770182in}}{\pgfqpoint{1.952439in}{1.766064in}}%
\pgfpathcurveto{\pgfqpoint{1.948321in}{1.761946in}}{\pgfqpoint{1.946007in}{1.756360in}}{\pgfqpoint{1.946007in}{1.750536in}}%
\pgfpathcurveto{\pgfqpoint{1.946007in}{1.744712in}}{\pgfqpoint{1.948321in}{1.739126in}}{\pgfqpoint{1.952439in}{1.735008in}}%
\pgfpathcurveto{\pgfqpoint{1.956557in}{1.730890in}}{\pgfqpoint{1.962144in}{1.728576in}}{\pgfqpoint{1.967967in}{1.728576in}}%
\pgfpathlineto{\pgfqpoint{1.967967in}{1.728576in}}%
\pgfpathclose%
\pgfusepath{stroke,fill}%
\end{pgfscope}%
\begin{pgfscope}%
\pgfpathrectangle{\pgfqpoint{0.100000in}{0.209042in}}{\pgfqpoint{2.906250in}{2.906250in}}%
\pgfusepath{clip}%
\pgfsetbuttcap%
\pgfsetroundjoin%
\definecolor{currentfill}{rgb}{0.069764,0.049726,0.193735}%
\pgfsetfillcolor{currentfill}%
\pgfsetfillopacity{0.800000}%
\pgfsetlinewidth{1.003750pt}%
\definecolor{currentstroke}{rgb}{0.069764,0.049726,0.193735}%
\pgfsetstrokecolor{currentstroke}%
\pgfsetstrokeopacity{0.800000}%
\pgfsetdash{}{0pt}%
\pgfpathmoveto{\pgfqpoint{1.646202in}{1.591186in}}%
\pgfpathcurveto{\pgfqpoint{1.652026in}{1.591186in}}{\pgfqpoint{1.657612in}{1.593500in}}{\pgfqpoint{1.661730in}{1.597618in}}%
\pgfpathcurveto{\pgfqpoint{1.665848in}{1.601736in}}{\pgfqpoint{1.668162in}{1.607322in}}{\pgfqpoint{1.668162in}{1.613146in}}%
\pgfpathcurveto{\pgfqpoint{1.668162in}{1.618970in}}{\pgfqpoint{1.665848in}{1.624556in}}{\pgfqpoint{1.661730in}{1.628674in}}%
\pgfpathcurveto{\pgfqpoint{1.657612in}{1.632792in}}{\pgfqpoint{1.652026in}{1.635106in}}{\pgfqpoint{1.646202in}{1.635106in}}%
\pgfpathcurveto{\pgfqpoint{1.640378in}{1.635106in}}{\pgfqpoint{1.634792in}{1.632792in}}{\pgfqpoint{1.630674in}{1.628674in}}%
\pgfpathcurveto{\pgfqpoint{1.626556in}{1.624556in}}{\pgfqpoint{1.624242in}{1.618970in}}{\pgfqpoint{1.624242in}{1.613146in}}%
\pgfpathcurveto{\pgfqpoint{1.624242in}{1.607322in}}{\pgfqpoint{1.626556in}{1.601736in}}{\pgfqpoint{1.630674in}{1.597618in}}%
\pgfpathcurveto{\pgfqpoint{1.634792in}{1.593500in}}{\pgfqpoint{1.640378in}{1.591186in}}{\pgfqpoint{1.646202in}{1.591186in}}%
\pgfpathlineto{\pgfqpoint{1.646202in}{1.591186in}}%
\pgfpathclose%
\pgfusepath{stroke,fill}%
\end{pgfscope}%
\begin{pgfscope}%
\pgfpathrectangle{\pgfqpoint{0.100000in}{0.209042in}}{\pgfqpoint{2.906250in}{2.906250in}}%
\pgfusepath{clip}%
\pgfsetbuttcap%
\pgfsetroundjoin%
\definecolor{currentfill}{rgb}{0.335308,0.078236,0.491024}%
\pgfsetfillcolor{currentfill}%
\pgfsetfillopacity{0.800000}%
\pgfsetlinewidth{1.003750pt}%
\definecolor{currentstroke}{rgb}{0.335308,0.078236,0.491024}%
\pgfsetstrokecolor{currentstroke}%
\pgfsetstrokeopacity{0.800000}%
\pgfsetdash{}{0pt}%
\pgfpathmoveto{\pgfqpoint{1.321147in}{1.628161in}}%
\pgfpathcurveto{\pgfqpoint{1.326971in}{1.628161in}}{\pgfqpoint{1.332558in}{1.630475in}}{\pgfqpoint{1.336676in}{1.634593in}}%
\pgfpathcurveto{\pgfqpoint{1.340794in}{1.638711in}}{\pgfqpoint{1.343108in}{1.644297in}}{\pgfqpoint{1.343108in}{1.650121in}}%
\pgfpathcurveto{\pgfqpoint{1.343108in}{1.655945in}}{\pgfqpoint{1.340794in}{1.661531in}}{\pgfqpoint{1.336676in}{1.665649in}}%
\pgfpathcurveto{\pgfqpoint{1.332558in}{1.669767in}}{\pgfqpoint{1.326971in}{1.672081in}}{\pgfqpoint{1.321147in}{1.672081in}}%
\pgfpathcurveto{\pgfqpoint{1.315323in}{1.672081in}}{\pgfqpoint{1.309737in}{1.669767in}}{\pgfqpoint{1.305619in}{1.665649in}}%
\pgfpathcurveto{\pgfqpoint{1.301501in}{1.661531in}}{\pgfqpoint{1.299187in}{1.655945in}}{\pgfqpoint{1.299187in}{1.650121in}}%
\pgfpathcurveto{\pgfqpoint{1.299187in}{1.644297in}}{\pgfqpoint{1.301501in}{1.638711in}}{\pgfqpoint{1.305619in}{1.634593in}}%
\pgfpathcurveto{\pgfqpoint{1.309737in}{1.630475in}}{\pgfqpoint{1.315323in}{1.628161in}}{\pgfqpoint{1.321147in}{1.628161in}}%
\pgfpathlineto{\pgfqpoint{1.321147in}{1.628161in}}%
\pgfpathclose%
\pgfusepath{stroke,fill}%
\end{pgfscope}%
\begin{pgfscope}%
\pgfpathrectangle{\pgfqpoint{0.100000in}{0.209042in}}{\pgfqpoint{2.906250in}{2.906250in}}%
\pgfusepath{clip}%
\pgfsetbuttcap%
\pgfsetroundjoin%
\definecolor{currentfill}{rgb}{0.366012,0.090314,0.497960}%
\pgfsetfillcolor{currentfill}%
\pgfsetfillopacity{0.800000}%
\pgfsetlinewidth{1.003750pt}%
\definecolor{currentstroke}{rgb}{0.366012,0.090314,0.497960}%
\pgfsetstrokecolor{currentstroke}%
\pgfsetstrokeopacity{0.800000}%
\pgfsetdash{}{0pt}%
\pgfpathmoveto{\pgfqpoint{1.688668in}{1.981449in}}%
\pgfpathcurveto{\pgfqpoint{1.694492in}{1.981449in}}{\pgfqpoint{1.700078in}{1.983762in}}{\pgfqpoint{1.704197in}{1.987881in}}%
\pgfpathcurveto{\pgfqpoint{1.708315in}{1.991999in}}{\pgfqpoint{1.710629in}{1.997585in}}{\pgfqpoint{1.710629in}{2.003409in}}%
\pgfpathcurveto{\pgfqpoint{1.710629in}{2.009233in}}{\pgfqpoint{1.708315in}{2.014819in}}{\pgfqpoint{1.704197in}{2.018937in}}%
\pgfpathcurveto{\pgfqpoint{1.700078in}{2.023055in}}{\pgfqpoint{1.694492in}{2.025369in}}{\pgfqpoint{1.688668in}{2.025369in}}%
\pgfpathcurveto{\pgfqpoint{1.682844in}{2.025369in}}{\pgfqpoint{1.677258in}{2.023055in}}{\pgfqpoint{1.673140in}{2.018937in}}%
\pgfpathcurveto{\pgfqpoint{1.669022in}{2.014819in}}{\pgfqpoint{1.666708in}{2.009233in}}{\pgfqpoint{1.666708in}{2.003409in}}%
\pgfpathcurveto{\pgfqpoint{1.666708in}{1.997585in}}{\pgfqpoint{1.669022in}{1.991999in}}{\pgfqpoint{1.673140in}{1.987881in}}%
\pgfpathcurveto{\pgfqpoint{1.677258in}{1.983762in}}{\pgfqpoint{1.682844in}{1.981449in}}{\pgfqpoint{1.688668in}{1.981449in}}%
\pgfpathlineto{\pgfqpoint{1.688668in}{1.981449in}}%
\pgfpathclose%
\pgfusepath{stroke,fill}%
\end{pgfscope}%
\begin{pgfscope}%
\pgfpathrectangle{\pgfqpoint{0.100000in}{0.209042in}}{\pgfqpoint{2.906250in}{2.906250in}}%
\pgfusepath{clip}%
\pgfsetbuttcap%
\pgfsetroundjoin%
\definecolor{currentfill}{rgb}{0.432967,0.117855,0.506160}%
\pgfsetfillcolor{currentfill}%
\pgfsetfillopacity{0.800000}%
\pgfsetlinewidth{1.003750pt}%
\definecolor{currentstroke}{rgb}{0.432967,0.117855,0.506160}%
\pgfsetstrokecolor{currentstroke}%
\pgfsetstrokeopacity{0.800000}%
\pgfsetdash{}{0pt}%
\pgfpathmoveto{\pgfqpoint{1.859574in}{1.447067in}}%
\pgfpathcurveto{\pgfqpoint{1.865398in}{1.447067in}}{\pgfqpoint{1.870984in}{1.449381in}}{\pgfqpoint{1.875102in}{1.453499in}}%
\pgfpathcurveto{\pgfqpoint{1.879220in}{1.457617in}}{\pgfqpoint{1.881534in}{1.463203in}}{\pgfqpoint{1.881534in}{1.469027in}}%
\pgfpathcurveto{\pgfqpoint{1.881534in}{1.474851in}}{\pgfqpoint{1.879220in}{1.480437in}}{\pgfqpoint{1.875102in}{1.484555in}}%
\pgfpathcurveto{\pgfqpoint{1.870984in}{1.488674in}}{\pgfqpoint{1.865398in}{1.490987in}}{\pgfqpoint{1.859574in}{1.490987in}}%
\pgfpathcurveto{\pgfqpoint{1.853750in}{1.490987in}}{\pgfqpoint{1.848164in}{1.488674in}}{\pgfqpoint{1.844046in}{1.484555in}}%
\pgfpathcurveto{\pgfqpoint{1.839927in}{1.480437in}}{\pgfqpoint{1.837613in}{1.474851in}}{\pgfqpoint{1.837613in}{1.469027in}}%
\pgfpathcurveto{\pgfqpoint{1.837613in}{1.463203in}}{\pgfqpoint{1.839927in}{1.457617in}}{\pgfqpoint{1.844046in}{1.453499in}}%
\pgfpathcurveto{\pgfqpoint{1.848164in}{1.449381in}}{\pgfqpoint{1.853750in}{1.447067in}}{\pgfqpoint{1.859574in}{1.447067in}}%
\pgfpathlineto{\pgfqpoint{1.859574in}{1.447067in}}%
\pgfpathclose%
\pgfusepath{stroke,fill}%
\end{pgfscope}%
\begin{pgfscope}%
\pgfpathrectangle{\pgfqpoint{0.100000in}{0.209042in}}{\pgfqpoint{2.906250in}{2.906250in}}%
\pgfusepath{clip}%
\pgfsetbuttcap%
\pgfsetroundjoin%
\definecolor{currentfill}{rgb}{0.291366,0.064553,0.475462}%
\pgfsetfillcolor{currentfill}%
\pgfsetfillopacity{0.800000}%
\pgfsetlinewidth{1.003750pt}%
\definecolor{currentstroke}{rgb}{0.291366,0.064553,0.475462}%
\pgfsetstrokecolor{currentstroke}%
\pgfsetstrokeopacity{0.800000}%
\pgfsetdash{}{0pt}%
\pgfpathmoveto{\pgfqpoint{1.339029in}{1.792262in}}%
\pgfpathcurveto{\pgfqpoint{1.344853in}{1.792262in}}{\pgfqpoint{1.350439in}{1.794576in}}{\pgfqpoint{1.354557in}{1.798694in}}%
\pgfpathcurveto{\pgfqpoint{1.358675in}{1.802812in}}{\pgfqpoint{1.360989in}{1.808399in}}{\pgfqpoint{1.360989in}{1.814223in}}%
\pgfpathcurveto{\pgfqpoint{1.360989in}{1.820046in}}{\pgfqpoint{1.358675in}{1.825633in}}{\pgfqpoint{1.354557in}{1.829751in}}%
\pgfpathcurveto{\pgfqpoint{1.350439in}{1.833869in}}{\pgfqpoint{1.344853in}{1.836183in}}{\pgfqpoint{1.339029in}{1.836183in}}%
\pgfpathcurveto{\pgfqpoint{1.333205in}{1.836183in}}{\pgfqpoint{1.327619in}{1.833869in}}{\pgfqpoint{1.323501in}{1.829751in}}%
\pgfpathcurveto{\pgfqpoint{1.319383in}{1.825633in}}{\pgfqpoint{1.317069in}{1.820046in}}{\pgfqpoint{1.317069in}{1.814223in}}%
\pgfpathcurveto{\pgfqpoint{1.317069in}{1.808399in}}{\pgfqpoint{1.319383in}{1.802812in}}{\pgfqpoint{1.323501in}{1.798694in}}%
\pgfpathcurveto{\pgfqpoint{1.327619in}{1.794576in}}{\pgfqpoint{1.333205in}{1.792262in}}{\pgfqpoint{1.339029in}{1.792262in}}%
\pgfpathlineto{\pgfqpoint{1.339029in}{1.792262in}}%
\pgfpathclose%
\pgfusepath{stroke,fill}%
\end{pgfscope}%
\begin{pgfscope}%
\pgfpathrectangle{\pgfqpoint{0.100000in}{0.209042in}}{\pgfqpoint{2.906250in}{2.906250in}}%
\pgfusepath{clip}%
\pgfsetbuttcap%
\pgfsetroundjoin%
\definecolor{currentfill}{rgb}{0.316654,0.071690,0.485380}%
\pgfsetfillcolor{currentfill}%
\pgfsetfillopacity{0.800000}%
\pgfsetlinewidth{1.003750pt}%
\definecolor{currentstroke}{rgb}{0.316654,0.071690,0.485380}%
\pgfsetstrokecolor{currentstroke}%
\pgfsetstrokeopacity{0.800000}%
\pgfsetdash{}{0pt}%
\pgfpathmoveto{\pgfqpoint{1.419346in}{1.910796in}}%
\pgfpathcurveto{\pgfqpoint{1.425170in}{1.910796in}}{\pgfqpoint{1.430756in}{1.913110in}}{\pgfqpoint{1.434874in}{1.917228in}}%
\pgfpathcurveto{\pgfqpoint{1.438992in}{1.921346in}}{\pgfqpoint{1.441306in}{1.926932in}}{\pgfqpoint{1.441306in}{1.932756in}}%
\pgfpathcurveto{\pgfqpoint{1.441306in}{1.938580in}}{\pgfqpoint{1.438992in}{1.944167in}}{\pgfqpoint{1.434874in}{1.948285in}}%
\pgfpathcurveto{\pgfqpoint{1.430756in}{1.952403in}}{\pgfqpoint{1.425170in}{1.954717in}}{\pgfqpoint{1.419346in}{1.954717in}}%
\pgfpathcurveto{\pgfqpoint{1.413522in}{1.954717in}}{\pgfqpoint{1.407936in}{1.952403in}}{\pgfqpoint{1.403818in}{1.948285in}}%
\pgfpathcurveto{\pgfqpoint{1.399699in}{1.944167in}}{\pgfqpoint{1.397386in}{1.938580in}}{\pgfqpoint{1.397386in}{1.932756in}}%
\pgfpathcurveto{\pgfqpoint{1.397386in}{1.926932in}}{\pgfqpoint{1.399699in}{1.921346in}}{\pgfqpoint{1.403818in}{1.917228in}}%
\pgfpathcurveto{\pgfqpoint{1.407936in}{1.913110in}}{\pgfqpoint{1.413522in}{1.910796in}}{\pgfqpoint{1.419346in}{1.910796in}}%
\pgfpathlineto{\pgfqpoint{1.419346in}{1.910796in}}%
\pgfpathclose%
\pgfusepath{stroke,fill}%
\end{pgfscope}%
\begin{pgfscope}%
\pgfpathrectangle{\pgfqpoint{0.100000in}{0.209042in}}{\pgfqpoint{2.906250in}{2.906250in}}%
\pgfusepath{clip}%
\pgfsetbuttcap%
\pgfsetroundjoin%
\definecolor{currentfill}{rgb}{0.408629,0.107930,0.504052}%
\pgfsetfillcolor{currentfill}%
\pgfsetfillopacity{0.800000}%
\pgfsetlinewidth{1.003750pt}%
\definecolor{currentstroke}{rgb}{0.408629,0.107930,0.504052}%
\pgfsetstrokecolor{currentstroke}%
\pgfsetstrokeopacity{0.800000}%
\pgfsetdash{}{0pt}%
\pgfpathmoveto{\pgfqpoint{1.949685in}{1.545326in}}%
\pgfpathcurveto{\pgfqpoint{1.955508in}{1.545326in}}{\pgfqpoint{1.961095in}{1.547639in}}{\pgfqpoint{1.965213in}{1.551758in}}%
\pgfpathcurveto{\pgfqpoint{1.969331in}{1.555876in}}{\pgfqpoint{1.971645in}{1.561462in}}{\pgfqpoint{1.971645in}{1.567286in}}%
\pgfpathcurveto{\pgfqpoint{1.971645in}{1.573110in}}{\pgfqpoint{1.969331in}{1.578696in}}{\pgfqpoint{1.965213in}{1.582814in}}%
\pgfpathcurveto{\pgfqpoint{1.961095in}{1.586932in}}{\pgfqpoint{1.955508in}{1.589246in}}{\pgfqpoint{1.949685in}{1.589246in}}%
\pgfpathcurveto{\pgfqpoint{1.943861in}{1.589246in}}{\pgfqpoint{1.938274in}{1.586932in}}{\pgfqpoint{1.934156in}{1.582814in}}%
\pgfpathcurveto{\pgfqpoint{1.930038in}{1.578696in}}{\pgfqpoint{1.927724in}{1.573110in}}{\pgfqpoint{1.927724in}{1.567286in}}%
\pgfpathcurveto{\pgfqpoint{1.927724in}{1.561462in}}{\pgfqpoint{1.930038in}{1.555876in}}{\pgfqpoint{1.934156in}{1.551758in}}%
\pgfpathcurveto{\pgfqpoint{1.938274in}{1.547639in}}{\pgfqpoint{1.943861in}{1.545326in}}{\pgfqpoint{1.949685in}{1.545326in}}%
\pgfpathlineto{\pgfqpoint{1.949685in}{1.545326in}}%
\pgfpathclose%
\pgfusepath{stroke,fill}%
\end{pgfscope}%
\begin{pgfscope}%
\pgfpathrectangle{\pgfqpoint{0.100000in}{0.209042in}}{\pgfqpoint{2.906250in}{2.906250in}}%
\pgfusepath{clip}%
\pgfsetbuttcap%
\pgfsetroundjoin%
\definecolor{currentfill}{rgb}{0.531507,0.154739,0.506895}%
\pgfsetfillcolor{currentfill}%
\pgfsetfillopacity{0.800000}%
\pgfsetlinewidth{1.003750pt}%
\definecolor{currentstroke}{rgb}{0.531507,0.154739,0.506895}%
\pgfsetstrokecolor{currentstroke}%
\pgfsetstrokeopacity{0.800000}%
\pgfsetdash{}{0pt}%
\pgfpathmoveto{\pgfqpoint{1.170208in}{1.728815in}}%
\pgfpathcurveto{\pgfqpoint{1.176032in}{1.728815in}}{\pgfqpoint{1.181618in}{1.731129in}}{\pgfqpoint{1.185736in}{1.735247in}}%
\pgfpathcurveto{\pgfqpoint{1.189854in}{1.739365in}}{\pgfqpoint{1.192168in}{1.744951in}}{\pgfqpoint{1.192168in}{1.750775in}}%
\pgfpathcurveto{\pgfqpoint{1.192168in}{1.756599in}}{\pgfqpoint{1.189854in}{1.762185in}}{\pgfqpoint{1.185736in}{1.766304in}}%
\pgfpathcurveto{\pgfqpoint{1.181618in}{1.770422in}}{\pgfqpoint{1.176032in}{1.772736in}}{\pgfqpoint{1.170208in}{1.772736in}}%
\pgfpathcurveto{\pgfqpoint{1.164384in}{1.772736in}}{\pgfqpoint{1.158798in}{1.770422in}}{\pgfqpoint{1.154680in}{1.766304in}}%
\pgfpathcurveto{\pgfqpoint{1.150561in}{1.762185in}}{\pgfqpoint{1.148247in}{1.756599in}}{\pgfqpoint{1.148247in}{1.750775in}}%
\pgfpathcurveto{\pgfqpoint{1.148247in}{1.744951in}}{\pgfqpoint{1.150561in}{1.739365in}}{\pgfqpoint{1.154680in}{1.735247in}}%
\pgfpathcurveto{\pgfqpoint{1.158798in}{1.731129in}}{\pgfqpoint{1.164384in}{1.728815in}}{\pgfqpoint{1.170208in}{1.728815in}}%
\pgfpathlineto{\pgfqpoint{1.170208in}{1.728815in}}%
\pgfpathclose%
\pgfusepath{stroke,fill}%
\end{pgfscope}%
\begin{pgfscope}%
\pgfpathrectangle{\pgfqpoint{0.100000in}{0.209042in}}{\pgfqpoint{2.906250in}{2.906250in}}%
\pgfusepath{clip}%
\pgfsetbuttcap%
\pgfsetroundjoin%
\definecolor{currentfill}{rgb}{0.846416,0.272473,0.421631}%
\pgfsetfillcolor{currentfill}%
\pgfsetfillopacity{0.800000}%
\pgfsetlinewidth{1.003750pt}%
\definecolor{currentstroke}{rgb}{0.846416,0.272473,0.421631}%
\pgfsetstrokecolor{currentstroke}%
\pgfsetstrokeopacity{0.800000}%
\pgfsetdash{}{0pt}%
\pgfpathmoveto{\pgfqpoint{1.908578in}{1.215360in}}%
\pgfpathcurveto{\pgfqpoint{1.914402in}{1.215360in}}{\pgfqpoint{1.919988in}{1.217674in}}{\pgfqpoint{1.924106in}{1.221792in}}%
\pgfpathcurveto{\pgfqpoint{1.928224in}{1.225910in}}{\pgfqpoint{1.930538in}{1.231496in}}{\pgfqpoint{1.930538in}{1.237320in}}%
\pgfpathcurveto{\pgfqpoint{1.930538in}{1.243144in}}{\pgfqpoint{1.928224in}{1.248730in}}{\pgfqpoint{1.924106in}{1.252848in}}%
\pgfpathcurveto{\pgfqpoint{1.919988in}{1.256966in}}{\pgfqpoint{1.914402in}{1.259280in}}{\pgfqpoint{1.908578in}{1.259280in}}%
\pgfpathcurveto{\pgfqpoint{1.902754in}{1.259280in}}{\pgfqpoint{1.897167in}{1.256966in}}{\pgfqpoint{1.893049in}{1.252848in}}%
\pgfpathcurveto{\pgfqpoint{1.888931in}{1.248730in}}{\pgfqpoint{1.886617in}{1.243144in}}{\pgfqpoint{1.886617in}{1.237320in}}%
\pgfpathcurveto{\pgfqpoint{1.886617in}{1.231496in}}{\pgfqpoint{1.888931in}{1.225910in}}{\pgfqpoint{1.893049in}{1.221792in}}%
\pgfpathcurveto{\pgfqpoint{1.897167in}{1.217674in}}{\pgfqpoint{1.902754in}{1.215360in}}{\pgfqpoint{1.908578in}{1.215360in}}%
\pgfpathlineto{\pgfqpoint{1.908578in}{1.215360in}}%
\pgfpathclose%
\pgfusepath{stroke,fill}%
\end{pgfscope}%
\begin{pgfscope}%
\pgfpathrectangle{\pgfqpoint{0.100000in}{0.209042in}}{\pgfqpoint{2.906250in}{2.906250in}}%
\pgfusepath{clip}%
\pgfsetbuttcap%
\pgfsetroundjoin%
\definecolor{currentfill}{rgb}{0.684224,0.204286,0.484219}%
\pgfsetfillcolor{currentfill}%
\pgfsetfillopacity{0.800000}%
\pgfsetlinewidth{1.003750pt}%
\definecolor{currentstroke}{rgb}{0.684224,0.204286,0.484219}%
\pgfsetstrokecolor{currentstroke}%
\pgfsetstrokeopacity{0.800000}%
\pgfsetdash{}{0pt}%
\pgfpathmoveto{\pgfqpoint{1.267732in}{1.384772in}}%
\pgfpathcurveto{\pgfqpoint{1.273556in}{1.384772in}}{\pgfqpoint{1.279142in}{1.387086in}}{\pgfqpoint{1.283260in}{1.391204in}}%
\pgfpathcurveto{\pgfqpoint{1.287378in}{1.395322in}}{\pgfqpoint{1.289692in}{1.400909in}}{\pgfqpoint{1.289692in}{1.406732in}}%
\pgfpathcurveto{\pgfqpoint{1.289692in}{1.412556in}}{\pgfqpoint{1.287378in}{1.418143in}}{\pgfqpoint{1.283260in}{1.422261in}}%
\pgfpathcurveto{\pgfqpoint{1.279142in}{1.426379in}}{\pgfqpoint{1.273556in}{1.428693in}}{\pgfqpoint{1.267732in}{1.428693in}}%
\pgfpathcurveto{\pgfqpoint{1.261908in}{1.428693in}}{\pgfqpoint{1.256322in}{1.426379in}}{\pgfqpoint{1.252203in}{1.422261in}}%
\pgfpathcurveto{\pgfqpoint{1.248085in}{1.418143in}}{\pgfqpoint{1.245771in}{1.412556in}}{\pgfqpoint{1.245771in}{1.406732in}}%
\pgfpathcurveto{\pgfqpoint{1.245771in}{1.400909in}}{\pgfqpoint{1.248085in}{1.395322in}}{\pgfqpoint{1.252203in}{1.391204in}}%
\pgfpathcurveto{\pgfqpoint{1.256322in}{1.387086in}}{\pgfqpoint{1.261908in}{1.384772in}}{\pgfqpoint{1.267732in}{1.384772in}}%
\pgfpathlineto{\pgfqpoint{1.267732in}{1.384772in}}%
\pgfpathclose%
\pgfusepath{stroke,fill}%
\end{pgfscope}%
\begin{pgfscope}%
\pgfpathrectangle{\pgfqpoint{0.100000in}{0.209042in}}{\pgfqpoint{2.906250in}{2.906250in}}%
\pgfusepath{clip}%
\pgfsetbuttcap%
\pgfsetroundjoin%
\definecolor{currentfill}{rgb}{0.001462,0.000466,0.013866}%
\pgfsetfillcolor{currentfill}%
\pgfsetfillopacity{0.800000}%
\pgfsetlinewidth{1.003750pt}%
\definecolor{currentstroke}{rgb}{0.001462,0.000466,0.013866}%
\pgfsetstrokecolor{currentstroke}%
\pgfsetstrokeopacity{0.800000}%
\pgfsetdash{}{0pt}%
\pgfpathmoveto{\pgfqpoint{1.576334in}{1.702094in}}%
\pgfpathcurveto{\pgfqpoint{1.582158in}{1.702094in}}{\pgfqpoint{1.587744in}{1.704408in}}{\pgfqpoint{1.591862in}{1.708526in}}%
\pgfpathcurveto{\pgfqpoint{1.595980in}{1.712644in}}{\pgfqpoint{1.598294in}{1.718230in}}{\pgfqpoint{1.598294in}{1.724054in}}%
\pgfpathcurveto{\pgfqpoint{1.598294in}{1.729878in}}{\pgfqpoint{1.595980in}{1.735465in}}{\pgfqpoint{1.591862in}{1.739583in}}%
\pgfpathcurveto{\pgfqpoint{1.587744in}{1.743701in}}{\pgfqpoint{1.582158in}{1.746015in}}{\pgfqpoint{1.576334in}{1.746015in}}%
\pgfpathcurveto{\pgfqpoint{1.570510in}{1.746015in}}{\pgfqpoint{1.564924in}{1.743701in}}{\pgfqpoint{1.560806in}{1.739583in}}%
\pgfpathcurveto{\pgfqpoint{1.556688in}{1.735465in}}{\pgfqpoint{1.554374in}{1.729878in}}{\pgfqpoint{1.554374in}{1.724054in}}%
\pgfpathcurveto{\pgfqpoint{1.554374in}{1.718230in}}{\pgfqpoint{1.556688in}{1.712644in}}{\pgfqpoint{1.560806in}{1.708526in}}%
\pgfpathcurveto{\pgfqpoint{1.564924in}{1.704408in}}{\pgfqpoint{1.570510in}{1.702094in}}{\pgfqpoint{1.576334in}{1.702094in}}%
\pgfpathlineto{\pgfqpoint{1.576334in}{1.702094in}}%
\pgfpathclose%
\pgfusepath{stroke,fill}%
\end{pgfscope}%
\begin{pgfscope}%
\pgfpathrectangle{\pgfqpoint{0.100000in}{0.209042in}}{\pgfqpoint{2.906250in}{2.906250in}}%
\pgfusepath{clip}%
\pgfsetbuttcap%
\pgfsetroundjoin%
\definecolor{currentfill}{rgb}{0.463508,0.129893,0.507652}%
\pgfsetfillcolor{currentfill}%
\pgfsetfillopacity{0.800000}%
\pgfsetlinewidth{1.003750pt}%
\definecolor{currentstroke}{rgb}{0.463508,0.129893,0.507652}%
\pgfsetstrokecolor{currentstroke}%
\pgfsetstrokeopacity{0.800000}%
\pgfsetdash{}{0pt}%
\pgfpathmoveto{\pgfqpoint{1.272654in}{1.556193in}}%
\pgfpathcurveto{\pgfqpoint{1.278478in}{1.556193in}}{\pgfqpoint{1.284064in}{1.558507in}}{\pgfqpoint{1.288182in}{1.562625in}}%
\pgfpathcurveto{\pgfqpoint{1.292301in}{1.566743in}}{\pgfqpoint{1.294614in}{1.572329in}}{\pgfqpoint{1.294614in}{1.578153in}}%
\pgfpathcurveto{\pgfqpoint{1.294614in}{1.583977in}}{\pgfqpoint{1.292301in}{1.589563in}}{\pgfqpoint{1.288182in}{1.593681in}}%
\pgfpathcurveto{\pgfqpoint{1.284064in}{1.597799in}}{\pgfqpoint{1.278478in}{1.600113in}}{\pgfqpoint{1.272654in}{1.600113in}}%
\pgfpathcurveto{\pgfqpoint{1.266830in}{1.600113in}}{\pgfqpoint{1.261244in}{1.597799in}}{\pgfqpoint{1.257126in}{1.593681in}}%
\pgfpathcurveto{\pgfqpoint{1.253008in}{1.589563in}}{\pgfqpoint{1.250694in}{1.583977in}}{\pgfqpoint{1.250694in}{1.578153in}}%
\pgfpathcurveto{\pgfqpoint{1.250694in}{1.572329in}}{\pgfqpoint{1.253008in}{1.566743in}}{\pgfqpoint{1.257126in}{1.562625in}}%
\pgfpathcurveto{\pgfqpoint{1.261244in}{1.558507in}}{\pgfqpoint{1.266830in}{1.556193in}}{\pgfqpoint{1.272654in}{1.556193in}}%
\pgfpathlineto{\pgfqpoint{1.272654in}{1.556193in}}%
\pgfpathclose%
\pgfusepath{stroke,fill}%
\end{pgfscope}%
\begin{pgfscope}%
\pgfpathrectangle{\pgfqpoint{0.100000in}{0.209042in}}{\pgfqpoint{2.906250in}{2.906250in}}%
\pgfusepath{clip}%
\pgfsetbuttcap%
\pgfsetroundjoin%
\definecolor{currentfill}{rgb}{0.645633,0.191952,0.492910}%
\pgfsetfillcolor{currentfill}%
\pgfsetfillopacity{0.800000}%
\pgfsetlinewidth{1.003750pt}%
\definecolor{currentstroke}{rgb}{0.645633,0.191952,0.492910}%
\pgfsetstrokecolor{currentstroke}%
\pgfsetstrokeopacity{0.800000}%
\pgfsetdash{}{0pt}%
\pgfpathmoveto{\pgfqpoint{1.396095in}{1.331190in}}%
\pgfpathcurveto{\pgfqpoint{1.401919in}{1.331190in}}{\pgfqpoint{1.407505in}{1.333504in}}{\pgfqpoint{1.411624in}{1.337622in}}%
\pgfpathcurveto{\pgfqpoint{1.415742in}{1.341740in}}{\pgfqpoint{1.418056in}{1.347326in}}{\pgfqpoint{1.418056in}{1.353150in}}%
\pgfpathcurveto{\pgfqpoint{1.418056in}{1.358974in}}{\pgfqpoint{1.415742in}{1.364560in}}{\pgfqpoint{1.411624in}{1.368679in}}%
\pgfpathcurveto{\pgfqpoint{1.407505in}{1.372797in}}{\pgfqpoint{1.401919in}{1.375111in}}{\pgfqpoint{1.396095in}{1.375111in}}%
\pgfpathcurveto{\pgfqpoint{1.390271in}{1.375111in}}{\pgfqpoint{1.384685in}{1.372797in}}{\pgfqpoint{1.380567in}{1.368679in}}%
\pgfpathcurveto{\pgfqpoint{1.376449in}{1.364560in}}{\pgfqpoint{1.374135in}{1.358974in}}{\pgfqpoint{1.374135in}{1.353150in}}%
\pgfpathcurveto{\pgfqpoint{1.374135in}{1.347326in}}{\pgfqpoint{1.376449in}{1.341740in}}{\pgfqpoint{1.380567in}{1.337622in}}%
\pgfpathcurveto{\pgfqpoint{1.384685in}{1.333504in}}{\pgfqpoint{1.390271in}{1.331190in}}{\pgfqpoint{1.396095in}{1.331190in}}%
\pgfpathlineto{\pgfqpoint{1.396095in}{1.331190in}}%
\pgfpathclose%
\pgfusepath{stroke,fill}%
\end{pgfscope}%
\begin{pgfscope}%
\pgfpathrectangle{\pgfqpoint{0.100000in}{0.209042in}}{\pgfqpoint{2.906250in}{2.906250in}}%
\pgfusepath{clip}%
\pgfsetbuttcap%
\pgfsetroundjoin%
\definecolor{currentfill}{rgb}{0.271994,0.060994,0.465660}%
\pgfsetfillcolor{currentfill}%
\pgfsetfillopacity{0.800000}%
\pgfsetlinewidth{1.003750pt}%
\definecolor{currentstroke}{rgb}{0.271994,0.060994,0.465660}%
\pgfsetstrokecolor{currentstroke}%
\pgfsetstrokeopacity{0.800000}%
\pgfsetdash{}{0pt}%
\pgfpathmoveto{\pgfqpoint{1.656599in}{1.478034in}}%
\pgfpathcurveto{\pgfqpoint{1.662423in}{1.478034in}}{\pgfqpoint{1.668009in}{1.480348in}}{\pgfqpoint{1.672128in}{1.484466in}}%
\pgfpathcurveto{\pgfqpoint{1.676246in}{1.488584in}}{\pgfqpoint{1.678560in}{1.494170in}}{\pgfqpoint{1.678560in}{1.499994in}}%
\pgfpathcurveto{\pgfqpoint{1.678560in}{1.505818in}}{\pgfqpoint{1.676246in}{1.511404in}}{\pgfqpoint{1.672128in}{1.515523in}}%
\pgfpathcurveto{\pgfqpoint{1.668009in}{1.519641in}}{\pgfqpoint{1.662423in}{1.521955in}}{\pgfqpoint{1.656599in}{1.521955in}}%
\pgfpathcurveto{\pgfqpoint{1.650775in}{1.521955in}}{\pgfqpoint{1.645189in}{1.519641in}}{\pgfqpoint{1.641071in}{1.515523in}}%
\pgfpathcurveto{\pgfqpoint{1.636953in}{1.511404in}}{\pgfqpoint{1.634639in}{1.505818in}}{\pgfqpoint{1.634639in}{1.499994in}}%
\pgfpathcurveto{\pgfqpoint{1.634639in}{1.494170in}}{\pgfqpoint{1.636953in}{1.488584in}}{\pgfqpoint{1.641071in}{1.484466in}}%
\pgfpathcurveto{\pgfqpoint{1.645189in}{1.480348in}}{\pgfqpoint{1.650775in}{1.478034in}}{\pgfqpoint{1.656599in}{1.478034in}}%
\pgfpathlineto{\pgfqpoint{1.656599in}{1.478034in}}%
\pgfpathclose%
\pgfusepath{stroke,fill}%
\end{pgfscope}%
\begin{pgfscope}%
\pgfpathrectangle{\pgfqpoint{0.100000in}{0.209042in}}{\pgfqpoint{2.906250in}{2.906250in}}%
\pgfusepath{clip}%
\pgfsetbuttcap%
\pgfsetroundjoin%
\definecolor{currentfill}{rgb}{0.378211,0.095332,0.500067}%
\pgfsetfillcolor{currentfill}%
\pgfsetfillopacity{0.800000}%
\pgfsetlinewidth{1.003750pt}%
\definecolor{currentstroke}{rgb}{0.378211,0.095332,0.500067}%
\pgfsetstrokecolor{currentstroke}%
\pgfsetstrokeopacity{0.800000}%
\pgfsetdash{}{0pt}%
\pgfpathmoveto{\pgfqpoint{1.611640in}{1.995115in}}%
\pgfpathcurveto{\pgfqpoint{1.617463in}{1.995115in}}{\pgfqpoint{1.623050in}{1.997429in}}{\pgfqpoint{1.627168in}{2.001547in}}%
\pgfpathcurveto{\pgfqpoint{1.631286in}{2.005665in}}{\pgfqpoint{1.633600in}{2.011251in}}{\pgfqpoint{1.633600in}{2.017075in}}%
\pgfpathcurveto{\pgfqpoint{1.633600in}{2.022899in}}{\pgfqpoint{1.631286in}{2.028486in}}{\pgfqpoint{1.627168in}{2.032604in}}%
\pgfpathcurveto{\pgfqpoint{1.623050in}{2.036722in}}{\pgfqpoint{1.617463in}{2.039036in}}{\pgfqpoint{1.611640in}{2.039036in}}%
\pgfpathcurveto{\pgfqpoint{1.605816in}{2.039036in}}{\pgfqpoint{1.600229in}{2.036722in}}{\pgfqpoint{1.596111in}{2.032604in}}%
\pgfpathcurveto{\pgfqpoint{1.591993in}{2.028486in}}{\pgfqpoint{1.589679in}{2.022899in}}{\pgfqpoint{1.589679in}{2.017075in}}%
\pgfpathcurveto{\pgfqpoint{1.589679in}{2.011251in}}{\pgfqpoint{1.591993in}{2.005665in}}{\pgfqpoint{1.596111in}{2.001547in}}%
\pgfpathcurveto{\pgfqpoint{1.600229in}{1.997429in}}{\pgfqpoint{1.605816in}{1.995115in}}{\pgfqpoint{1.611640in}{1.995115in}}%
\pgfpathlineto{\pgfqpoint{1.611640in}{1.995115in}}%
\pgfpathclose%
\pgfusepath{stroke,fill}%
\end{pgfscope}%
\begin{pgfscope}%
\pgfpathrectangle{\pgfqpoint{0.100000in}{0.209042in}}{\pgfqpoint{2.906250in}{2.906250in}}%
\pgfusepath{clip}%
\pgfsetbuttcap%
\pgfsetroundjoin%
\definecolor{currentfill}{rgb}{0.613617,0.181811,0.498536}%
\pgfsetfillcolor{currentfill}%
\pgfsetfillopacity{0.800000}%
\pgfsetlinewidth{1.003750pt}%
\definecolor{currentstroke}{rgb}{0.613617,0.181811,0.498536}%
\pgfsetstrokecolor{currentstroke}%
\pgfsetstrokeopacity{0.800000}%
\pgfsetdash{}{0pt}%
\pgfpathmoveto{\pgfqpoint{1.142468in}{1.606040in}}%
\pgfpathcurveto{\pgfqpoint{1.148292in}{1.606040in}}{\pgfqpoint{1.153879in}{1.608354in}}{\pgfqpoint{1.157997in}{1.612472in}}%
\pgfpathcurveto{\pgfqpoint{1.162115in}{1.616590in}}{\pgfqpoint{1.164429in}{1.622177in}}{\pgfqpoint{1.164429in}{1.628000in}}%
\pgfpathcurveto{\pgfqpoint{1.164429in}{1.633824in}}{\pgfqpoint{1.162115in}{1.639411in}}{\pgfqpoint{1.157997in}{1.643529in}}%
\pgfpathcurveto{\pgfqpoint{1.153879in}{1.647647in}}{\pgfqpoint{1.148292in}{1.649961in}}{\pgfqpoint{1.142468in}{1.649961in}}%
\pgfpathcurveto{\pgfqpoint{1.136645in}{1.649961in}}{\pgfqpoint{1.131058in}{1.647647in}}{\pgfqpoint{1.126940in}{1.643529in}}%
\pgfpathcurveto{\pgfqpoint{1.122822in}{1.639411in}}{\pgfqpoint{1.120508in}{1.633824in}}{\pgfqpoint{1.120508in}{1.628000in}}%
\pgfpathcurveto{\pgfqpoint{1.120508in}{1.622177in}}{\pgfqpoint{1.122822in}{1.616590in}}{\pgfqpoint{1.126940in}{1.612472in}}%
\pgfpathcurveto{\pgfqpoint{1.131058in}{1.608354in}}{\pgfqpoint{1.136645in}{1.606040in}}{\pgfqpoint{1.142468in}{1.606040in}}%
\pgfpathlineto{\pgfqpoint{1.142468in}{1.606040in}}%
\pgfpathclose%
\pgfusepath{stroke,fill}%
\end{pgfscope}%
\begin{pgfscope}%
\pgfpathrectangle{\pgfqpoint{0.100000in}{0.209042in}}{\pgfqpoint{2.906250in}{2.906250in}}%
\pgfusepath{clip}%
\pgfsetbuttcap%
\pgfsetroundjoin%
\definecolor{currentfill}{rgb}{0.107899,0.064335,0.267289}%
\pgfsetfillcolor{currentfill}%
\pgfsetfillopacity{0.800000}%
\pgfsetlinewidth{1.003750pt}%
\definecolor{currentstroke}{rgb}{0.107899,0.064335,0.267289}%
\pgfsetstrokecolor{currentstroke}%
\pgfsetstrokeopacity{0.800000}%
\pgfsetdash{}{0pt}%
\pgfpathmoveto{\pgfqpoint{1.608366in}{1.562439in}}%
\pgfpathcurveto{\pgfqpoint{1.614190in}{1.562439in}}{\pgfqpoint{1.619776in}{1.564753in}}{\pgfqpoint{1.623894in}{1.568871in}}%
\pgfpathcurveto{\pgfqpoint{1.628012in}{1.572990in}}{\pgfqpoint{1.630326in}{1.578576in}}{\pgfqpoint{1.630326in}{1.584400in}}%
\pgfpathcurveto{\pgfqpoint{1.630326in}{1.590224in}}{\pgfqpoint{1.628012in}{1.595810in}}{\pgfqpoint{1.623894in}{1.599928in}}%
\pgfpathcurveto{\pgfqpoint{1.619776in}{1.604046in}}{\pgfqpoint{1.614190in}{1.606360in}}{\pgfqpoint{1.608366in}{1.606360in}}%
\pgfpathcurveto{\pgfqpoint{1.602542in}{1.606360in}}{\pgfqpoint{1.596956in}{1.604046in}}{\pgfqpoint{1.592837in}{1.599928in}}%
\pgfpathcurveto{\pgfqpoint{1.588719in}{1.595810in}}{\pgfqpoint{1.586405in}{1.590224in}}{\pgfqpoint{1.586405in}{1.584400in}}%
\pgfpathcurveto{\pgfqpoint{1.586405in}{1.578576in}}{\pgfqpoint{1.588719in}{1.572990in}}{\pgfqpoint{1.592837in}{1.568871in}}%
\pgfpathcurveto{\pgfqpoint{1.596956in}{1.564753in}}{\pgfqpoint{1.602542in}{1.562439in}}{\pgfqpoint{1.608366in}{1.562439in}}%
\pgfpathlineto{\pgfqpoint{1.608366in}{1.562439in}}%
\pgfpathclose%
\pgfusepath{stroke,fill}%
\end{pgfscope}%
\begin{pgfscope}%
\pgfpathrectangle{\pgfqpoint{0.100000in}{0.209042in}}{\pgfqpoint{2.906250in}{2.906250in}}%
\pgfusepath{clip}%
\pgfsetbuttcap%
\pgfsetroundjoin%
\definecolor{currentfill}{rgb}{0.056615,0.042160,0.167446}%
\pgfsetfillcolor{currentfill}%
\pgfsetfillopacity{0.800000}%
\pgfsetlinewidth{1.003750pt}%
\definecolor{currentstroke}{rgb}{0.056615,0.042160,0.167446}%
\pgfsetstrokecolor{currentstroke}%
\pgfsetstrokeopacity{0.800000}%
\pgfsetdash{}{0pt}%
\pgfpathmoveto{\pgfqpoint{1.573557in}{1.608093in}}%
\pgfpathcurveto{\pgfqpoint{1.579381in}{1.608093in}}{\pgfqpoint{1.584967in}{1.610406in}}{\pgfqpoint{1.589085in}{1.614525in}}%
\pgfpathcurveto{\pgfqpoint{1.593203in}{1.618643in}}{\pgfqpoint{1.595517in}{1.624229in}}{\pgfqpoint{1.595517in}{1.630053in}}%
\pgfpathcurveto{\pgfqpoint{1.595517in}{1.635877in}}{\pgfqpoint{1.593203in}{1.641463in}}{\pgfqpoint{1.589085in}{1.645581in}}%
\pgfpathcurveto{\pgfqpoint{1.584967in}{1.649699in}}{\pgfqpoint{1.579381in}{1.652013in}}{\pgfqpoint{1.573557in}{1.652013in}}%
\pgfpathcurveto{\pgfqpoint{1.567733in}{1.652013in}}{\pgfqpoint{1.562147in}{1.649699in}}{\pgfqpoint{1.558028in}{1.645581in}}%
\pgfpathcurveto{\pgfqpoint{1.553910in}{1.641463in}}{\pgfqpoint{1.551596in}{1.635877in}}{\pgfqpoint{1.551596in}{1.630053in}}%
\pgfpathcurveto{\pgfqpoint{1.551596in}{1.624229in}}{\pgfqpoint{1.553910in}{1.618643in}}{\pgfqpoint{1.558028in}{1.614525in}}%
\pgfpathcurveto{\pgfqpoint{1.562147in}{1.610406in}}{\pgfqpoint{1.567733in}{1.608093in}}{\pgfqpoint{1.573557in}{1.608093in}}%
\pgfpathlineto{\pgfqpoint{1.573557in}{1.608093in}}%
\pgfpathclose%
\pgfusepath{stroke,fill}%
\end{pgfscope}%
\begin{pgfscope}%
\pgfpathrectangle{\pgfqpoint{0.100000in}{0.209042in}}{\pgfqpoint{2.906250in}{2.906250in}}%
\pgfusepath{clip}%
\pgfsetbuttcap%
\pgfsetroundjoin%
\definecolor{currentfill}{rgb}{0.088155,0.058133,0.229922}%
\pgfsetfillcolor{currentfill}%
\pgfsetfillopacity{0.800000}%
\pgfsetlinewidth{1.003750pt}%
\definecolor{currentstroke}{rgb}{0.088155,0.058133,0.229922}%
\pgfsetstrokecolor{currentstroke}%
\pgfsetstrokeopacity{0.800000}%
\pgfsetdash{}{0pt}%
\pgfpathmoveto{\pgfqpoint{1.495067in}{1.635215in}}%
\pgfpathcurveto{\pgfqpoint{1.500891in}{1.635215in}}{\pgfqpoint{1.506478in}{1.637528in}}{\pgfqpoint{1.510596in}{1.641647in}}%
\pgfpathcurveto{\pgfqpoint{1.514714in}{1.645765in}}{\pgfqpoint{1.517028in}{1.651351in}}{\pgfqpoint{1.517028in}{1.657175in}}%
\pgfpathcurveto{\pgfqpoint{1.517028in}{1.662999in}}{\pgfqpoint{1.514714in}{1.668585in}}{\pgfqpoint{1.510596in}{1.672703in}}%
\pgfpathcurveto{\pgfqpoint{1.506478in}{1.676821in}}{\pgfqpoint{1.500891in}{1.679135in}}{\pgfqpoint{1.495067in}{1.679135in}}%
\pgfpathcurveto{\pgfqpoint{1.489243in}{1.679135in}}{\pgfqpoint{1.483657in}{1.676821in}}{\pgfqpoint{1.479539in}{1.672703in}}%
\pgfpathcurveto{\pgfqpoint{1.475421in}{1.668585in}}{\pgfqpoint{1.473107in}{1.662999in}}{\pgfqpoint{1.473107in}{1.657175in}}%
\pgfpathcurveto{\pgfqpoint{1.473107in}{1.651351in}}{\pgfqpoint{1.475421in}{1.645765in}}{\pgfqpoint{1.479539in}{1.641647in}}%
\pgfpathcurveto{\pgfqpoint{1.483657in}{1.637528in}}{\pgfqpoint{1.489243in}{1.635215in}}{\pgfqpoint{1.495067in}{1.635215in}}%
\pgfpathlineto{\pgfqpoint{1.495067in}{1.635215in}}%
\pgfpathclose%
\pgfusepath{stroke,fill}%
\end{pgfscope}%
\begin{pgfscope}%
\pgfpathrectangle{\pgfqpoint{0.100000in}{0.209042in}}{\pgfqpoint{2.906250in}{2.906250in}}%
\pgfusepath{clip}%
\pgfsetbuttcap%
\pgfsetroundjoin%
\definecolor{currentfill}{rgb}{0.278493,0.061978,0.469190}%
\pgfsetfillcolor{currentfill}%
\pgfsetfillopacity{0.800000}%
\pgfsetlinewidth{1.003750pt}%
\definecolor{currentstroke}{rgb}{0.278493,0.061978,0.469190}%
\pgfsetstrokecolor{currentstroke}%
\pgfsetstrokeopacity{0.800000}%
\pgfsetdash{}{0pt}%
\pgfpathmoveto{\pgfqpoint{1.767031in}{1.906477in}}%
\pgfpathcurveto{\pgfqpoint{1.772855in}{1.906477in}}{\pgfqpoint{1.778441in}{1.908790in}}{\pgfqpoint{1.782559in}{1.912909in}}%
\pgfpathcurveto{\pgfqpoint{1.786678in}{1.917027in}}{\pgfqpoint{1.788991in}{1.922613in}}{\pgfqpoint{1.788991in}{1.928437in}}%
\pgfpathcurveto{\pgfqpoint{1.788991in}{1.934261in}}{\pgfqpoint{1.786678in}{1.939847in}}{\pgfqpoint{1.782559in}{1.943965in}}%
\pgfpathcurveto{\pgfqpoint{1.778441in}{1.948083in}}{\pgfqpoint{1.772855in}{1.950397in}}{\pgfqpoint{1.767031in}{1.950397in}}%
\pgfpathcurveto{\pgfqpoint{1.761207in}{1.950397in}}{\pgfqpoint{1.755621in}{1.948083in}}{\pgfqpoint{1.751503in}{1.943965in}}%
\pgfpathcurveto{\pgfqpoint{1.747385in}{1.939847in}}{\pgfqpoint{1.745071in}{1.934261in}}{\pgfqpoint{1.745071in}{1.928437in}}%
\pgfpathcurveto{\pgfqpoint{1.745071in}{1.922613in}}{\pgfqpoint{1.747385in}{1.917027in}}{\pgfqpoint{1.751503in}{1.912909in}}%
\pgfpathcurveto{\pgfqpoint{1.755621in}{1.908790in}}{\pgfqpoint{1.761207in}{1.906477in}}{\pgfqpoint{1.767031in}{1.906477in}}%
\pgfpathlineto{\pgfqpoint{1.767031in}{1.906477in}}%
\pgfpathclose%
\pgfusepath{stroke,fill}%
\end{pgfscope}%
\begin{pgfscope}%
\pgfpathrectangle{\pgfqpoint{0.100000in}{0.209042in}}{\pgfqpoint{2.906250in}{2.906250in}}%
\pgfusepath{clip}%
\pgfsetbuttcap%
\pgfsetroundjoin%
\definecolor{currentfill}{rgb}{0.297740,0.066117,0.478243}%
\pgfsetfillcolor{currentfill}%
\pgfsetfillopacity{0.800000}%
\pgfsetlinewidth{1.003750pt}%
\definecolor{currentstroke}{rgb}{0.297740,0.066117,0.478243}%
\pgfsetstrokecolor{currentstroke}%
\pgfsetstrokeopacity{0.800000}%
\pgfsetdash{}{0pt}%
\pgfpathmoveto{\pgfqpoint{1.900021in}{1.789976in}}%
\pgfpathcurveto{\pgfqpoint{1.905845in}{1.789976in}}{\pgfqpoint{1.911431in}{1.792290in}}{\pgfqpoint{1.915549in}{1.796408in}}%
\pgfpathcurveto{\pgfqpoint{1.919667in}{1.800527in}}{\pgfqpoint{1.921981in}{1.806113in}}{\pgfqpoint{1.921981in}{1.811937in}}%
\pgfpathcurveto{\pgfqpoint{1.921981in}{1.817761in}}{\pgfqpoint{1.919667in}{1.823347in}}{\pgfqpoint{1.915549in}{1.827465in}}%
\pgfpathcurveto{\pgfqpoint{1.911431in}{1.831583in}}{\pgfqpoint{1.905845in}{1.833897in}}{\pgfqpoint{1.900021in}{1.833897in}}%
\pgfpathcurveto{\pgfqpoint{1.894197in}{1.833897in}}{\pgfqpoint{1.888611in}{1.831583in}}{\pgfqpoint{1.884493in}{1.827465in}}%
\pgfpathcurveto{\pgfqpoint{1.880374in}{1.823347in}}{\pgfqpoint{1.878061in}{1.817761in}}{\pgfqpoint{1.878061in}{1.811937in}}%
\pgfpathcurveto{\pgfqpoint{1.878061in}{1.806113in}}{\pgfqpoint{1.880374in}{1.800527in}}{\pgfqpoint{1.884493in}{1.796408in}}%
\pgfpathcurveto{\pgfqpoint{1.888611in}{1.792290in}}{\pgfqpoint{1.894197in}{1.789976in}}{\pgfqpoint{1.900021in}{1.789976in}}%
\pgfpathlineto{\pgfqpoint{1.900021in}{1.789976in}}%
\pgfpathclose%
\pgfusepath{stroke,fill}%
\end{pgfscope}%
\begin{pgfscope}%
\pgfpathrectangle{\pgfqpoint{0.100000in}{0.209042in}}{\pgfqpoint{2.906250in}{2.906250in}}%
\pgfusepath{clip}%
\pgfsetbuttcap%
\pgfsetroundjoin%
\definecolor{currentfill}{rgb}{0.402548,0.105420,0.503386}%
\pgfsetfillcolor{currentfill}%
\pgfsetfillopacity{0.800000}%
\pgfsetlinewidth{1.003750pt}%
\definecolor{currentstroke}{rgb}{0.402548,0.105420,0.503386}%
\pgfsetstrokecolor{currentstroke}%
\pgfsetstrokeopacity{0.800000}%
\pgfsetdash{}{0pt}%
\pgfpathmoveto{\pgfqpoint{1.611736in}{2.008748in}}%
\pgfpathcurveto{\pgfqpoint{1.617559in}{2.008748in}}{\pgfqpoint{1.623146in}{2.011062in}}{\pgfqpoint{1.627264in}{2.015180in}}%
\pgfpathcurveto{\pgfqpoint{1.631382in}{2.019299in}}{\pgfqpoint{1.633696in}{2.024885in}}{\pgfqpoint{1.633696in}{2.030709in}}%
\pgfpathcurveto{\pgfqpoint{1.633696in}{2.036533in}}{\pgfqpoint{1.631382in}{2.042119in}}{\pgfqpoint{1.627264in}{2.046237in}}%
\pgfpathcurveto{\pgfqpoint{1.623146in}{2.050355in}}{\pgfqpoint{1.617559in}{2.052669in}}{\pgfqpoint{1.611736in}{2.052669in}}%
\pgfpathcurveto{\pgfqpoint{1.605912in}{2.052669in}}{\pgfqpoint{1.600325in}{2.050355in}}{\pgfqpoint{1.596207in}{2.046237in}}%
\pgfpathcurveto{\pgfqpoint{1.592089in}{2.042119in}}{\pgfqpoint{1.589775in}{2.036533in}}{\pgfqpoint{1.589775in}{2.030709in}}%
\pgfpathcurveto{\pgfqpoint{1.589775in}{2.024885in}}{\pgfqpoint{1.592089in}{2.019299in}}{\pgfqpoint{1.596207in}{2.015180in}}%
\pgfpathcurveto{\pgfqpoint{1.600325in}{2.011062in}}{\pgfqpoint{1.605912in}{2.008748in}}{\pgfqpoint{1.611736in}{2.008748in}}%
\pgfpathlineto{\pgfqpoint{1.611736in}{2.008748in}}%
\pgfpathclose%
\pgfusepath{stroke,fill}%
\end{pgfscope}%
\begin{pgfscope}%
\pgfpathrectangle{\pgfqpoint{0.100000in}{0.209042in}}{\pgfqpoint{2.906250in}{2.906250in}}%
\pgfusepath{clip}%
\pgfsetbuttcap%
\pgfsetroundjoin%
\definecolor{currentfill}{rgb}{0.335308,0.078236,0.491024}%
\pgfsetfillcolor{currentfill}%
\pgfsetfillopacity{0.800000}%
\pgfsetlinewidth{1.003750pt}%
\definecolor{currentstroke}{rgb}{0.335308,0.078236,0.491024}%
\pgfsetstrokecolor{currentstroke}%
\pgfsetstrokeopacity{0.800000}%
\pgfsetdash{}{0pt}%
\pgfpathmoveto{\pgfqpoint{1.342536in}{1.847693in}}%
\pgfpathcurveto{\pgfqpoint{1.348360in}{1.847693in}}{\pgfqpoint{1.353946in}{1.850007in}}{\pgfqpoint{1.358064in}{1.854125in}}%
\pgfpathcurveto{\pgfqpoint{1.362182in}{1.858243in}}{\pgfqpoint{1.364496in}{1.863829in}}{\pgfqpoint{1.364496in}{1.869653in}}%
\pgfpathcurveto{\pgfqpoint{1.364496in}{1.875477in}}{\pgfqpoint{1.362182in}{1.881063in}}{\pgfqpoint{1.358064in}{1.885182in}}%
\pgfpathcurveto{\pgfqpoint{1.353946in}{1.889300in}}{\pgfqpoint{1.348360in}{1.891614in}}{\pgfqpoint{1.342536in}{1.891614in}}%
\pgfpathcurveto{\pgfqpoint{1.336712in}{1.891614in}}{\pgfqpoint{1.331126in}{1.889300in}}{\pgfqpoint{1.327008in}{1.885182in}}%
\pgfpathcurveto{\pgfqpoint{1.322890in}{1.881063in}}{\pgfqpoint{1.320576in}{1.875477in}}{\pgfqpoint{1.320576in}{1.869653in}}%
\pgfpathcurveto{\pgfqpoint{1.320576in}{1.863829in}}{\pgfqpoint{1.322890in}{1.858243in}}{\pgfqpoint{1.327008in}{1.854125in}}%
\pgfpathcurveto{\pgfqpoint{1.331126in}{1.850007in}}{\pgfqpoint{1.336712in}{1.847693in}}{\pgfqpoint{1.342536in}{1.847693in}}%
\pgfpathlineto{\pgfqpoint{1.342536in}{1.847693in}}%
\pgfpathclose%
\pgfusepath{stroke,fill}%
\end{pgfscope}%
\begin{pgfscope}%
\pgfpathrectangle{\pgfqpoint{0.100000in}{0.209042in}}{\pgfqpoint{2.906250in}{2.906250in}}%
\pgfusepath{clip}%
\pgfsetbuttcap%
\pgfsetroundjoin%
\definecolor{currentfill}{rgb}{0.451271,0.125132,0.507198}%
\pgfsetfillcolor{currentfill}%
\pgfsetfillopacity{0.800000}%
\pgfsetlinewidth{1.003750pt}%
\definecolor{currentstroke}{rgb}{0.451271,0.125132,0.507198}%
\pgfsetstrokecolor{currentstroke}%
\pgfsetstrokeopacity{0.800000}%
\pgfsetdash{}{0pt}%
\pgfpathmoveto{\pgfqpoint{1.908055in}{1.926891in}}%
\pgfpathcurveto{\pgfqpoint{1.913879in}{1.926891in}}{\pgfqpoint{1.919465in}{1.929205in}}{\pgfqpoint{1.923583in}{1.933323in}}%
\pgfpathcurveto{\pgfqpoint{1.927701in}{1.937441in}}{\pgfqpoint{1.930015in}{1.943027in}}{\pgfqpoint{1.930015in}{1.948851in}}%
\pgfpathcurveto{\pgfqpoint{1.930015in}{1.954675in}}{\pgfqpoint{1.927701in}{1.960261in}}{\pgfqpoint{1.923583in}{1.964379in}}%
\pgfpathcurveto{\pgfqpoint{1.919465in}{1.968497in}}{\pgfqpoint{1.913879in}{1.970811in}}{\pgfqpoint{1.908055in}{1.970811in}}%
\pgfpathcurveto{\pgfqpoint{1.902231in}{1.970811in}}{\pgfqpoint{1.896644in}{1.968497in}}{\pgfqpoint{1.892526in}{1.964379in}}%
\pgfpathcurveto{\pgfqpoint{1.888408in}{1.960261in}}{\pgfqpoint{1.886094in}{1.954675in}}{\pgfqpoint{1.886094in}{1.948851in}}%
\pgfpathcurveto{\pgfqpoint{1.886094in}{1.943027in}}{\pgfqpoint{1.888408in}{1.937441in}}{\pgfqpoint{1.892526in}{1.933323in}}%
\pgfpathcurveto{\pgfqpoint{1.896644in}{1.929205in}}{\pgfqpoint{1.902231in}{1.926891in}}{\pgfqpoint{1.908055in}{1.926891in}}%
\pgfpathlineto{\pgfqpoint{1.908055in}{1.926891in}}%
\pgfpathclose%
\pgfusepath{stroke,fill}%
\end{pgfscope}%
\begin{pgfscope}%
\pgfpathrectangle{\pgfqpoint{0.100000in}{0.209042in}}{\pgfqpoint{2.906250in}{2.906250in}}%
\pgfusepath{clip}%
\pgfsetbuttcap%
\pgfsetroundjoin%
\definecolor{currentfill}{rgb}{0.159018,0.068354,0.352688}%
\pgfsetfillcolor{currentfill}%
\pgfsetfillopacity{0.800000}%
\pgfsetlinewidth{1.003750pt}%
\definecolor{currentstroke}{rgb}{0.159018,0.068354,0.352688}%
\pgfsetstrokecolor{currentstroke}%
\pgfsetstrokeopacity{0.800000}%
\pgfsetdash{}{0pt}%
\pgfpathmoveto{\pgfqpoint{1.511231in}{1.860500in}}%
\pgfpathcurveto{\pgfqpoint{1.517055in}{1.860500in}}{\pgfqpoint{1.522642in}{1.862814in}}{\pgfqpoint{1.526760in}{1.866932in}}%
\pgfpathcurveto{\pgfqpoint{1.530878in}{1.871050in}}{\pgfqpoint{1.533192in}{1.876636in}}{\pgfqpoint{1.533192in}{1.882460in}}%
\pgfpathcurveto{\pgfqpoint{1.533192in}{1.888284in}}{\pgfqpoint{1.530878in}{1.893870in}}{\pgfqpoint{1.526760in}{1.897989in}}%
\pgfpathcurveto{\pgfqpoint{1.522642in}{1.902107in}}{\pgfqpoint{1.517055in}{1.904421in}}{\pgfqpoint{1.511231in}{1.904421in}}%
\pgfpathcurveto{\pgfqpoint{1.505407in}{1.904421in}}{\pgfqpoint{1.499821in}{1.902107in}}{\pgfqpoint{1.495703in}{1.897989in}}%
\pgfpathcurveto{\pgfqpoint{1.491585in}{1.893870in}}{\pgfqpoint{1.489271in}{1.888284in}}{\pgfqpoint{1.489271in}{1.882460in}}%
\pgfpathcurveto{\pgfqpoint{1.489271in}{1.876636in}}{\pgfqpoint{1.491585in}{1.871050in}}{\pgfqpoint{1.495703in}{1.866932in}}%
\pgfpathcurveto{\pgfqpoint{1.499821in}{1.862814in}}{\pgfqpoint{1.505407in}{1.860500in}}{\pgfqpoint{1.511231in}{1.860500in}}%
\pgfpathlineto{\pgfqpoint{1.511231in}{1.860500in}}%
\pgfpathclose%
\pgfusepath{stroke,fill}%
\end{pgfscope}%
\begin{pgfscope}%
\pgfpathrectangle{\pgfqpoint{0.100000in}{0.209042in}}{\pgfqpoint{2.906250in}{2.906250in}}%
\pgfusepath{clip}%
\pgfsetbuttcap%
\pgfsetroundjoin%
\definecolor{currentfill}{rgb}{0.258857,0.059706,0.457710}%
\pgfsetfillcolor{currentfill}%
\pgfsetfillopacity{0.800000}%
\pgfsetlinewidth{1.003750pt}%
\definecolor{currentstroke}{rgb}{0.258857,0.059706,0.457710}%
\pgfsetstrokecolor{currentstroke}%
\pgfsetstrokeopacity{0.800000}%
\pgfsetdash{}{0pt}%
\pgfpathmoveto{\pgfqpoint{1.901764in}{1.708754in}}%
\pgfpathcurveto{\pgfqpoint{1.907588in}{1.708754in}}{\pgfqpoint{1.913174in}{1.711068in}}{\pgfqpoint{1.917292in}{1.715186in}}%
\pgfpathcurveto{\pgfqpoint{1.921411in}{1.719304in}}{\pgfqpoint{1.923724in}{1.724891in}}{\pgfqpoint{1.923724in}{1.730715in}}%
\pgfpathcurveto{\pgfqpoint{1.923724in}{1.736538in}}{\pgfqpoint{1.921411in}{1.742125in}}{\pgfqpoint{1.917292in}{1.746243in}}%
\pgfpathcurveto{\pgfqpoint{1.913174in}{1.750361in}}{\pgfqpoint{1.907588in}{1.752675in}}{\pgfqpoint{1.901764in}{1.752675in}}%
\pgfpathcurveto{\pgfqpoint{1.895940in}{1.752675in}}{\pgfqpoint{1.890354in}{1.750361in}}{\pgfqpoint{1.886236in}{1.746243in}}%
\pgfpathcurveto{\pgfqpoint{1.882118in}{1.742125in}}{\pgfqpoint{1.879804in}{1.736538in}}{\pgfqpoint{1.879804in}{1.730715in}}%
\pgfpathcurveto{\pgfqpoint{1.879804in}{1.724891in}}{\pgfqpoint{1.882118in}{1.719304in}}{\pgfqpoint{1.886236in}{1.715186in}}%
\pgfpathcurveto{\pgfqpoint{1.890354in}{1.711068in}}{\pgfqpoint{1.895940in}{1.708754in}}{\pgfqpoint{1.901764in}{1.708754in}}%
\pgfpathlineto{\pgfqpoint{1.901764in}{1.708754in}}%
\pgfpathclose%
\pgfusepath{stroke,fill}%
\end{pgfscope}%
\begin{pgfscope}%
\pgfpathrectangle{\pgfqpoint{0.100000in}{0.209042in}}{\pgfqpoint{2.906250in}{2.906250in}}%
\pgfusepath{clip}%
\pgfsetbuttcap%
\pgfsetroundjoin%
\definecolor{currentfill}{rgb}{0.439062,0.120298,0.506555}%
\pgfsetfillcolor{currentfill}%
\pgfsetfillopacity{0.800000}%
\pgfsetlinewidth{1.003750pt}%
\definecolor{currentstroke}{rgb}{0.439062,0.120298,0.506555}%
\pgfsetstrokecolor{currentstroke}%
\pgfsetstrokeopacity{0.800000}%
\pgfsetdash{}{0pt}%
\pgfpathmoveto{\pgfqpoint{1.434522in}{1.435333in}}%
\pgfpathcurveto{\pgfqpoint{1.440346in}{1.435333in}}{\pgfqpoint{1.445932in}{1.437647in}}{\pgfqpoint{1.450051in}{1.441765in}}%
\pgfpathcurveto{\pgfqpoint{1.454169in}{1.445883in}}{\pgfqpoint{1.456483in}{1.451470in}}{\pgfqpoint{1.456483in}{1.457294in}}%
\pgfpathcurveto{\pgfqpoint{1.456483in}{1.463117in}}{\pgfqpoint{1.454169in}{1.468704in}}{\pgfqpoint{1.450051in}{1.472822in}}%
\pgfpathcurveto{\pgfqpoint{1.445932in}{1.476940in}}{\pgfqpoint{1.440346in}{1.479254in}}{\pgfqpoint{1.434522in}{1.479254in}}%
\pgfpathcurveto{\pgfqpoint{1.428698in}{1.479254in}}{\pgfqpoint{1.423112in}{1.476940in}}{\pgfqpoint{1.418994in}{1.472822in}}%
\pgfpathcurveto{\pgfqpoint{1.414876in}{1.468704in}}{\pgfqpoint{1.412562in}{1.463117in}}{\pgfqpoint{1.412562in}{1.457294in}}%
\pgfpathcurveto{\pgfqpoint{1.412562in}{1.451470in}}{\pgfqpoint{1.414876in}{1.445883in}}{\pgfqpoint{1.418994in}{1.441765in}}%
\pgfpathcurveto{\pgfqpoint{1.423112in}{1.437647in}}{\pgfqpoint{1.428698in}{1.435333in}}{\pgfqpoint{1.434522in}{1.435333in}}%
\pgfpathlineto{\pgfqpoint{1.434522in}{1.435333in}}%
\pgfpathclose%
\pgfusepath{stroke,fill}%
\end{pgfscope}%
\begin{pgfscope}%
\pgfpathrectangle{\pgfqpoint{0.100000in}{0.209042in}}{\pgfqpoint{2.906250in}{2.906250in}}%
\pgfusepath{clip}%
\pgfsetbuttcap%
\pgfsetroundjoin%
\definecolor{currentfill}{rgb}{0.677786,0.202203,0.485819}%
\pgfsetfillcolor{currentfill}%
\pgfsetfillopacity{0.800000}%
\pgfsetlinewidth{1.003750pt}%
\definecolor{currentstroke}{rgb}{0.677786,0.202203,0.485819}%
\pgfsetstrokecolor{currentstroke}%
\pgfsetstrokeopacity{0.800000}%
\pgfsetdash{}{0pt}%
\pgfpathmoveto{\pgfqpoint{1.263618in}{2.056653in}}%
\pgfpathcurveto{\pgfqpoint{1.269442in}{2.056653in}}{\pgfqpoint{1.275028in}{2.058967in}}{\pgfqpoint{1.279146in}{2.063085in}}%
\pgfpathcurveto{\pgfqpoint{1.283264in}{2.067204in}}{\pgfqpoint{1.285578in}{2.072790in}}{\pgfqpoint{1.285578in}{2.078614in}}%
\pgfpathcurveto{\pgfqpoint{1.285578in}{2.084438in}}{\pgfqpoint{1.283264in}{2.090024in}}{\pgfqpoint{1.279146in}{2.094142in}}%
\pgfpathcurveto{\pgfqpoint{1.275028in}{2.098260in}}{\pgfqpoint{1.269442in}{2.100574in}}{\pgfqpoint{1.263618in}{2.100574in}}%
\pgfpathcurveto{\pgfqpoint{1.257794in}{2.100574in}}{\pgfqpoint{1.252208in}{2.098260in}}{\pgfqpoint{1.248090in}{2.094142in}}%
\pgfpathcurveto{\pgfqpoint{1.243972in}{2.090024in}}{\pgfqpoint{1.241658in}{2.084438in}}{\pgfqpoint{1.241658in}{2.078614in}}%
\pgfpathcurveto{\pgfqpoint{1.241658in}{2.072790in}}{\pgfqpoint{1.243972in}{2.067204in}}{\pgfqpoint{1.248090in}{2.063085in}}%
\pgfpathcurveto{\pgfqpoint{1.252208in}{2.058967in}}{\pgfqpoint{1.257794in}{2.056653in}}{\pgfqpoint{1.263618in}{2.056653in}}%
\pgfpathlineto{\pgfqpoint{1.263618in}{2.056653in}}%
\pgfpathclose%
\pgfusepath{stroke,fill}%
\end{pgfscope}%
\begin{pgfscope}%
\pgfpathrectangle{\pgfqpoint{0.100000in}{0.209042in}}{\pgfqpoint{2.906250in}{2.906250in}}%
\pgfusepath{clip}%
\pgfsetbuttcap%
\pgfsetroundjoin%
\definecolor{currentfill}{rgb}{0.060949,0.044794,0.176129}%
\pgfsetfillcolor{currentfill}%
\pgfsetfillopacity{0.800000}%
\pgfsetlinewidth{1.003750pt}%
\definecolor{currentstroke}{rgb}{0.060949,0.044794,0.176129}%
\pgfsetstrokecolor{currentstroke}%
\pgfsetstrokeopacity{0.800000}%
\pgfsetdash{}{0pt}%
\pgfpathmoveto{\pgfqpoint{1.532515in}{1.626225in}}%
\pgfpathcurveto{\pgfqpoint{1.538339in}{1.626225in}}{\pgfqpoint{1.543925in}{1.628539in}}{\pgfqpoint{1.548043in}{1.632657in}}%
\pgfpathcurveto{\pgfqpoint{1.552162in}{1.636775in}}{\pgfqpoint{1.554475in}{1.642361in}}{\pgfqpoint{1.554475in}{1.648185in}}%
\pgfpathcurveto{\pgfqpoint{1.554475in}{1.654009in}}{\pgfqpoint{1.552162in}{1.659595in}}{\pgfqpoint{1.548043in}{1.663713in}}%
\pgfpathcurveto{\pgfqpoint{1.543925in}{1.667831in}}{\pgfqpoint{1.538339in}{1.670145in}}{\pgfqpoint{1.532515in}{1.670145in}}%
\pgfpathcurveto{\pgfqpoint{1.526691in}{1.670145in}}{\pgfqpoint{1.521105in}{1.667831in}}{\pgfqpoint{1.516987in}{1.663713in}}%
\pgfpathcurveto{\pgfqpoint{1.512869in}{1.659595in}}{\pgfqpoint{1.510555in}{1.654009in}}{\pgfqpoint{1.510555in}{1.648185in}}%
\pgfpathcurveto{\pgfqpoint{1.510555in}{1.642361in}}{\pgfqpoint{1.512869in}{1.636775in}}{\pgfqpoint{1.516987in}{1.632657in}}%
\pgfpathcurveto{\pgfqpoint{1.521105in}{1.628539in}}{\pgfqpoint{1.526691in}{1.626225in}}{\pgfqpoint{1.532515in}{1.626225in}}%
\pgfpathlineto{\pgfqpoint{1.532515in}{1.626225in}}%
\pgfpathclose%
\pgfusepath{stroke,fill}%
\end{pgfscope}%
\begin{pgfscope}%
\pgfpathrectangle{\pgfqpoint{0.100000in}{0.209042in}}{\pgfqpoint{2.906250in}{2.906250in}}%
\pgfusepath{clip}%
\pgfsetbuttcap%
\pgfsetroundjoin%
\definecolor{currentfill}{rgb}{0.950210,0.390820,0.362468}%
\pgfsetfillcolor{currentfill}%
\pgfsetfillopacity{0.800000}%
\pgfsetlinewidth{1.003750pt}%
\definecolor{currentstroke}{rgb}{0.950210,0.390820,0.362468}%
\pgfsetstrokecolor{currentstroke}%
\pgfsetstrokeopacity{0.800000}%
\pgfsetdash{}{0pt}%
\pgfpathmoveto{\pgfqpoint{1.148057in}{1.260895in}}%
\pgfpathcurveto{\pgfqpoint{1.153881in}{1.260895in}}{\pgfqpoint{1.159467in}{1.263209in}}{\pgfqpoint{1.163585in}{1.267327in}}%
\pgfpathcurveto{\pgfqpoint{1.167704in}{1.271445in}}{\pgfqpoint{1.170017in}{1.277031in}}{\pgfqpoint{1.170017in}{1.282855in}}%
\pgfpathcurveto{\pgfqpoint{1.170017in}{1.288679in}}{\pgfqpoint{1.167704in}{1.294266in}}{\pgfqpoint{1.163585in}{1.298384in}}%
\pgfpathcurveto{\pgfqpoint{1.159467in}{1.302502in}}{\pgfqpoint{1.153881in}{1.304816in}}{\pgfqpoint{1.148057in}{1.304816in}}%
\pgfpathcurveto{\pgfqpoint{1.142233in}{1.304816in}}{\pgfqpoint{1.136647in}{1.302502in}}{\pgfqpoint{1.132529in}{1.298384in}}%
\pgfpathcurveto{\pgfqpoint{1.128411in}{1.294266in}}{\pgfqpoint{1.126097in}{1.288679in}}{\pgfqpoint{1.126097in}{1.282855in}}%
\pgfpathcurveto{\pgfqpoint{1.126097in}{1.277031in}}{\pgfqpoint{1.128411in}{1.271445in}}{\pgfqpoint{1.132529in}{1.267327in}}%
\pgfpathcurveto{\pgfqpoint{1.136647in}{1.263209in}}{\pgfqpoint{1.142233in}{1.260895in}}{\pgfqpoint{1.148057in}{1.260895in}}%
\pgfpathlineto{\pgfqpoint{1.148057in}{1.260895in}}%
\pgfpathclose%
\pgfusepath{stroke,fill}%
\end{pgfscope}%
\begin{pgfscope}%
\pgfpathrectangle{\pgfqpoint{0.100000in}{0.209042in}}{\pgfqpoint{2.906250in}{2.906250in}}%
\pgfusepath{clip}%
\pgfsetbuttcap%
\pgfsetroundjoin%
\definecolor{currentfill}{rgb}{0.969680,0.446936,0.360311}%
\pgfsetfillcolor{currentfill}%
\pgfsetfillopacity{0.800000}%
\pgfsetlinewidth{1.003750pt}%
\definecolor{currentstroke}{rgb}{0.969680,0.446936,0.360311}%
\pgfsetstrokecolor{currentstroke}%
\pgfsetstrokeopacity{0.800000}%
\pgfsetdash{}{0pt}%
\pgfpathmoveto{\pgfqpoint{1.119706in}{1.244666in}}%
\pgfpathcurveto{\pgfqpoint{1.125529in}{1.244666in}}{\pgfqpoint{1.131116in}{1.246980in}}{\pgfqpoint{1.135234in}{1.251098in}}%
\pgfpathcurveto{\pgfqpoint{1.139352in}{1.255216in}}{\pgfqpoint{1.141666in}{1.260802in}}{\pgfqpoint{1.141666in}{1.266626in}}%
\pgfpathcurveto{\pgfqpoint{1.141666in}{1.272450in}}{\pgfqpoint{1.139352in}{1.278036in}}{\pgfqpoint{1.135234in}{1.282155in}}%
\pgfpathcurveto{\pgfqpoint{1.131116in}{1.286273in}}{\pgfqpoint{1.125529in}{1.288587in}}{\pgfqpoint{1.119706in}{1.288587in}}%
\pgfpathcurveto{\pgfqpoint{1.113882in}{1.288587in}}{\pgfqpoint{1.108295in}{1.286273in}}{\pgfqpoint{1.104177in}{1.282155in}}%
\pgfpathcurveto{\pgfqpoint{1.100059in}{1.278036in}}{\pgfqpoint{1.097745in}{1.272450in}}{\pgfqpoint{1.097745in}{1.266626in}}%
\pgfpathcurveto{\pgfqpoint{1.097745in}{1.260802in}}{\pgfqpoint{1.100059in}{1.255216in}}{\pgfqpoint{1.104177in}{1.251098in}}%
\pgfpathcurveto{\pgfqpoint{1.108295in}{1.246980in}}{\pgfqpoint{1.113882in}{1.244666in}}{\pgfqpoint{1.119706in}{1.244666in}}%
\pgfpathlineto{\pgfqpoint{1.119706in}{1.244666in}}%
\pgfpathclose%
\pgfusepath{stroke,fill}%
\end{pgfscope}%
\begin{pgfscope}%
\pgfpathrectangle{\pgfqpoint{0.100000in}{0.209042in}}{\pgfqpoint{2.906250in}{2.906250in}}%
\pgfusepath{clip}%
\pgfsetbuttcap%
\pgfsetroundjoin%
\definecolor{currentfill}{rgb}{0.639216,0.189921,0.494150}%
\pgfsetfillcolor{currentfill}%
\pgfsetfillopacity{0.800000}%
\pgfsetlinewidth{1.003750pt}%
\definecolor{currentstroke}{rgb}{0.639216,0.189921,0.494150}%
\pgfsetstrokecolor{currentstroke}%
\pgfsetstrokeopacity{0.800000}%
\pgfsetdash{}{0pt}%
\pgfpathmoveto{\pgfqpoint{2.038820in}{1.946269in}}%
\pgfpathcurveto{\pgfqpoint{2.044644in}{1.946269in}}{\pgfqpoint{2.050231in}{1.948583in}}{\pgfqpoint{2.054349in}{1.952701in}}%
\pgfpathcurveto{\pgfqpoint{2.058467in}{1.956819in}}{\pgfqpoint{2.060781in}{1.962406in}}{\pgfqpoint{2.060781in}{1.968229in}}%
\pgfpathcurveto{\pgfqpoint{2.060781in}{1.974053in}}{\pgfqpoint{2.058467in}{1.979640in}}{\pgfqpoint{2.054349in}{1.983758in}}%
\pgfpathcurveto{\pgfqpoint{2.050231in}{1.987876in}}{\pgfqpoint{2.044644in}{1.990190in}}{\pgfqpoint{2.038820in}{1.990190in}}%
\pgfpathcurveto{\pgfqpoint{2.032996in}{1.990190in}}{\pgfqpoint{2.027410in}{1.987876in}}{\pgfqpoint{2.023292in}{1.983758in}}%
\pgfpathcurveto{\pgfqpoint{2.019174in}{1.979640in}}{\pgfqpoint{2.016860in}{1.974053in}}{\pgfqpoint{2.016860in}{1.968229in}}%
\pgfpathcurveto{\pgfqpoint{2.016860in}{1.962406in}}{\pgfqpoint{2.019174in}{1.956819in}}{\pgfqpoint{2.023292in}{1.952701in}}%
\pgfpathcurveto{\pgfqpoint{2.027410in}{1.948583in}}{\pgfqpoint{2.032996in}{1.946269in}}{\pgfqpoint{2.038820in}{1.946269in}}%
\pgfpathlineto{\pgfqpoint{2.038820in}{1.946269in}}%
\pgfpathclose%
\pgfusepath{stroke,fill}%
\end{pgfscope}%
\begin{pgfscope}%
\pgfpathrectangle{\pgfqpoint{0.100000in}{0.209042in}}{\pgfqpoint{2.906250in}{2.906250in}}%
\pgfusepath{clip}%
\pgfsetbuttcap%
\pgfsetroundjoin%
\definecolor{currentfill}{rgb}{0.556571,0.163269,0.505230}%
\pgfsetfillcolor{currentfill}%
\pgfsetfillopacity{0.800000}%
\pgfsetlinewidth{1.003750pt}%
\definecolor{currentstroke}{rgb}{0.556571,0.163269,0.505230}%
\pgfsetstrokecolor{currentstroke}%
\pgfsetstrokeopacity{0.800000}%
\pgfsetdash{}{0pt}%
\pgfpathmoveto{\pgfqpoint{1.467254in}{1.355237in}}%
\pgfpathcurveto{\pgfqpoint{1.473078in}{1.355237in}}{\pgfqpoint{1.478665in}{1.357550in}}{\pgfqpoint{1.482783in}{1.361669in}}%
\pgfpathcurveto{\pgfqpoint{1.486901in}{1.365787in}}{\pgfqpoint{1.489215in}{1.371373in}}{\pgfqpoint{1.489215in}{1.377197in}}%
\pgfpathcurveto{\pgfqpoint{1.489215in}{1.383021in}}{\pgfqpoint{1.486901in}{1.388607in}}{\pgfqpoint{1.482783in}{1.392725in}}%
\pgfpathcurveto{\pgfqpoint{1.478665in}{1.396843in}}{\pgfqpoint{1.473078in}{1.399157in}}{\pgfqpoint{1.467254in}{1.399157in}}%
\pgfpathcurveto{\pgfqpoint{1.461431in}{1.399157in}}{\pgfqpoint{1.455844in}{1.396843in}}{\pgfqpoint{1.451726in}{1.392725in}}%
\pgfpathcurveto{\pgfqpoint{1.447608in}{1.388607in}}{\pgfqpoint{1.445294in}{1.383021in}}{\pgfqpoint{1.445294in}{1.377197in}}%
\pgfpathcurveto{\pgfqpoint{1.445294in}{1.371373in}}{\pgfqpoint{1.447608in}{1.365787in}}{\pgfqpoint{1.451726in}{1.361669in}}%
\pgfpathcurveto{\pgfqpoint{1.455844in}{1.357550in}}{\pgfqpoint{1.461431in}{1.355237in}}{\pgfqpoint{1.467254in}{1.355237in}}%
\pgfpathlineto{\pgfqpoint{1.467254in}{1.355237in}}%
\pgfpathclose%
\pgfusepath{stroke,fill}%
\end{pgfscope}%
\begin{pgfscope}%
\pgfpathrectangle{\pgfqpoint{0.100000in}{0.209042in}}{\pgfqpoint{2.906250in}{2.906250in}}%
\pgfusepath{clip}%
\pgfsetbuttcap%
\pgfsetroundjoin%
\definecolor{currentfill}{rgb}{0.232077,0.059889,0.437695}%
\pgfsetfillcolor{currentfill}%
\pgfsetfillopacity{0.800000}%
\pgfsetlinewidth{1.003750pt}%
\definecolor{currentstroke}{rgb}{0.232077,0.059889,0.437695}%
\pgfsetstrokecolor{currentstroke}%
\pgfsetstrokeopacity{0.800000}%
\pgfsetdash{}{0pt}%
\pgfpathmoveto{\pgfqpoint{1.868262in}{1.614418in}}%
\pgfpathcurveto{\pgfqpoint{1.874086in}{1.614418in}}{\pgfqpoint{1.879672in}{1.616732in}}{\pgfqpoint{1.883790in}{1.620850in}}%
\pgfpathcurveto{\pgfqpoint{1.887908in}{1.624968in}}{\pgfqpoint{1.890222in}{1.630554in}}{\pgfqpoint{1.890222in}{1.636378in}}%
\pgfpathcurveto{\pgfqpoint{1.890222in}{1.642202in}}{\pgfqpoint{1.887908in}{1.647788in}}{\pgfqpoint{1.883790in}{1.651906in}}%
\pgfpathcurveto{\pgfqpoint{1.879672in}{1.656024in}}{\pgfqpoint{1.874086in}{1.658338in}}{\pgfqpoint{1.868262in}{1.658338in}}%
\pgfpathcurveto{\pgfqpoint{1.862438in}{1.658338in}}{\pgfqpoint{1.856852in}{1.656024in}}{\pgfqpoint{1.852734in}{1.651906in}}%
\pgfpathcurveto{\pgfqpoint{1.848615in}{1.647788in}}{\pgfqpoint{1.846301in}{1.642202in}}{\pgfqpoint{1.846301in}{1.636378in}}%
\pgfpathcurveto{\pgfqpoint{1.846301in}{1.630554in}}{\pgfqpoint{1.848615in}{1.624968in}}{\pgfqpoint{1.852734in}{1.620850in}}%
\pgfpathcurveto{\pgfqpoint{1.856852in}{1.616732in}}{\pgfqpoint{1.862438in}{1.614418in}}{\pgfqpoint{1.868262in}{1.614418in}}%
\pgfpathlineto{\pgfqpoint{1.868262in}{1.614418in}}%
\pgfpathclose%
\pgfusepath{stroke,fill}%
\end{pgfscope}%
\begin{pgfscope}%
\pgfpathrectangle{\pgfqpoint{0.100000in}{0.209042in}}{\pgfqpoint{2.906250in}{2.906250in}}%
\pgfusepath{clip}%
\pgfsetbuttcap%
\pgfsetroundjoin%
\definecolor{currentfill}{rgb}{0.218512,0.061158,0.425392}%
\pgfsetfillcolor{currentfill}%
\pgfsetfillopacity{0.800000}%
\pgfsetlinewidth{1.003750pt}%
\definecolor{currentstroke}{rgb}{0.218512,0.061158,0.425392}%
\pgfsetstrokecolor{currentstroke}%
\pgfsetstrokeopacity{0.800000}%
\pgfsetdash{}{0pt}%
\pgfpathmoveto{\pgfqpoint{1.841121in}{1.588012in}}%
\pgfpathcurveto{\pgfqpoint{1.846945in}{1.588012in}}{\pgfqpoint{1.852531in}{1.590326in}}{\pgfqpoint{1.856650in}{1.594444in}}%
\pgfpathcurveto{\pgfqpoint{1.860768in}{1.598563in}}{\pgfqpoint{1.863082in}{1.604149in}}{\pgfqpoint{1.863082in}{1.609973in}}%
\pgfpathcurveto{\pgfqpoint{1.863082in}{1.615797in}}{\pgfqpoint{1.860768in}{1.621383in}}{\pgfqpoint{1.856650in}{1.625501in}}%
\pgfpathcurveto{\pgfqpoint{1.852531in}{1.629619in}}{\pgfqpoint{1.846945in}{1.631933in}}{\pgfqpoint{1.841121in}{1.631933in}}%
\pgfpathcurveto{\pgfqpoint{1.835297in}{1.631933in}}{\pgfqpoint{1.829711in}{1.629619in}}{\pgfqpoint{1.825593in}{1.625501in}}%
\pgfpathcurveto{\pgfqpoint{1.821475in}{1.621383in}}{\pgfqpoint{1.819161in}{1.615797in}}{\pgfqpoint{1.819161in}{1.609973in}}%
\pgfpathcurveto{\pgfqpoint{1.819161in}{1.604149in}}{\pgfqpoint{1.821475in}{1.598563in}}{\pgfqpoint{1.825593in}{1.594444in}}%
\pgfpathcurveto{\pgfqpoint{1.829711in}{1.590326in}}{\pgfqpoint{1.835297in}{1.588012in}}{\pgfqpoint{1.841121in}{1.588012in}}%
\pgfpathlineto{\pgfqpoint{1.841121in}{1.588012in}}%
\pgfpathclose%
\pgfusepath{stroke,fill}%
\end{pgfscope}%
\begin{pgfscope}%
\pgfpathrectangle{\pgfqpoint{0.100000in}{0.209042in}}{\pgfqpoint{2.906250in}{2.906250in}}%
\pgfusepath{clip}%
\pgfsetbuttcap%
\pgfsetroundjoin%
\definecolor{currentfill}{rgb}{0.078815,0.054184,0.211667}%
\pgfsetfillcolor{currentfill}%
\pgfsetfillopacity{0.800000}%
\pgfsetlinewidth{1.003750pt}%
\definecolor{currentstroke}{rgb}{0.078815,0.054184,0.211667}%
\pgfsetstrokecolor{currentstroke}%
\pgfsetstrokeopacity{0.800000}%
\pgfsetdash{}{0pt}%
\pgfpathmoveto{\pgfqpoint{1.701703in}{1.814097in}}%
\pgfpathcurveto{\pgfqpoint{1.707527in}{1.814097in}}{\pgfqpoint{1.713113in}{1.816411in}}{\pgfqpoint{1.717231in}{1.820529in}}%
\pgfpathcurveto{\pgfqpoint{1.721349in}{1.824648in}}{\pgfqpoint{1.723663in}{1.830234in}}{\pgfqpoint{1.723663in}{1.836058in}}%
\pgfpathcurveto{\pgfqpoint{1.723663in}{1.841882in}}{\pgfqpoint{1.721349in}{1.847468in}}{\pgfqpoint{1.717231in}{1.851586in}}%
\pgfpathcurveto{\pgfqpoint{1.713113in}{1.855704in}}{\pgfqpoint{1.707527in}{1.858018in}}{\pgfqpoint{1.701703in}{1.858018in}}%
\pgfpathcurveto{\pgfqpoint{1.695879in}{1.858018in}}{\pgfqpoint{1.690293in}{1.855704in}}{\pgfqpoint{1.686175in}{1.851586in}}%
\pgfpathcurveto{\pgfqpoint{1.682057in}{1.847468in}}{\pgfqpoint{1.679743in}{1.841882in}}{\pgfqpoint{1.679743in}{1.836058in}}%
\pgfpathcurveto{\pgfqpoint{1.679743in}{1.830234in}}{\pgfqpoint{1.682057in}{1.824648in}}{\pgfqpoint{1.686175in}{1.820529in}}%
\pgfpathcurveto{\pgfqpoint{1.690293in}{1.816411in}}{\pgfqpoint{1.695879in}{1.814097in}}{\pgfqpoint{1.701703in}{1.814097in}}%
\pgfpathlineto{\pgfqpoint{1.701703in}{1.814097in}}%
\pgfpathclose%
\pgfusepath{stroke,fill}%
\end{pgfscope}%
\begin{pgfscope}%
\pgfpathrectangle{\pgfqpoint{0.100000in}{0.209042in}}{\pgfqpoint{2.906250in}{2.906250in}}%
\pgfusepath{clip}%
\pgfsetbuttcap%
\pgfsetroundjoin%
\definecolor{currentfill}{rgb}{0.171713,0.067305,0.370771}%
\pgfsetfillcolor{currentfill}%
\pgfsetfillopacity{0.800000}%
\pgfsetlinewidth{1.003750pt}%
\definecolor{currentstroke}{rgb}{0.171713,0.067305,0.370771}%
\pgfsetstrokecolor{currentstroke}%
\pgfsetstrokeopacity{0.800000}%
\pgfsetdash{}{0pt}%
\pgfpathmoveto{\pgfqpoint{1.404004in}{1.693288in}}%
\pgfpathcurveto{\pgfqpoint{1.409828in}{1.693288in}}{\pgfqpoint{1.415414in}{1.695602in}}{\pgfqpoint{1.419532in}{1.699720in}}%
\pgfpathcurveto{\pgfqpoint{1.423650in}{1.703838in}}{\pgfqpoint{1.425964in}{1.709424in}}{\pgfqpoint{1.425964in}{1.715248in}}%
\pgfpathcurveto{\pgfqpoint{1.425964in}{1.721072in}}{\pgfqpoint{1.423650in}{1.726658in}}{\pgfqpoint{1.419532in}{1.730776in}}%
\pgfpathcurveto{\pgfqpoint{1.415414in}{1.734894in}}{\pgfqpoint{1.409828in}{1.737208in}}{\pgfqpoint{1.404004in}{1.737208in}}%
\pgfpathcurveto{\pgfqpoint{1.398180in}{1.737208in}}{\pgfqpoint{1.392594in}{1.734894in}}{\pgfqpoint{1.388476in}{1.730776in}}%
\pgfpathcurveto{\pgfqpoint{1.384358in}{1.726658in}}{\pgfqpoint{1.382044in}{1.721072in}}{\pgfqpoint{1.382044in}{1.715248in}}%
\pgfpathcurveto{\pgfqpoint{1.382044in}{1.709424in}}{\pgfqpoint{1.384358in}{1.703838in}}{\pgfqpoint{1.388476in}{1.699720in}}%
\pgfpathcurveto{\pgfqpoint{1.392594in}{1.695602in}}{\pgfqpoint{1.398180in}{1.693288in}}{\pgfqpoint{1.404004in}{1.693288in}}%
\pgfpathlineto{\pgfqpoint{1.404004in}{1.693288in}}%
\pgfpathclose%
\pgfusepath{stroke,fill}%
\end{pgfscope}%
\begin{pgfscope}%
\pgfpathrectangle{\pgfqpoint{0.100000in}{0.209042in}}{\pgfqpoint{2.906250in}{2.906250in}}%
\pgfusepath{clip}%
\pgfsetbuttcap%
\pgfsetroundjoin%
\definecolor{currentfill}{rgb}{0.983485,0.513280,0.374198}%
\pgfsetfillcolor{currentfill}%
\pgfsetfillopacity{0.800000}%
\pgfsetlinewidth{1.003750pt}%
\definecolor{currentstroke}{rgb}{0.983485,0.513280,0.374198}%
\pgfsetstrokecolor{currentstroke}%
\pgfsetstrokeopacity{0.800000}%
\pgfsetdash{}{0pt}%
\pgfpathmoveto{\pgfqpoint{2.136280in}{1.159819in}}%
\pgfpathcurveto{\pgfqpoint{2.142104in}{1.159819in}}{\pgfqpoint{2.147690in}{1.162133in}}{\pgfqpoint{2.151808in}{1.166251in}}%
\pgfpathcurveto{\pgfqpoint{2.155927in}{1.170370in}}{\pgfqpoint{2.158240in}{1.175956in}}{\pgfqpoint{2.158240in}{1.181780in}}%
\pgfpathcurveto{\pgfqpoint{2.158240in}{1.187604in}}{\pgfqpoint{2.155927in}{1.193190in}}{\pgfqpoint{2.151808in}{1.197308in}}%
\pgfpathcurveto{\pgfqpoint{2.147690in}{1.201426in}}{\pgfqpoint{2.142104in}{1.203740in}}{\pgfqpoint{2.136280in}{1.203740in}}%
\pgfpathcurveto{\pgfqpoint{2.130456in}{1.203740in}}{\pgfqpoint{2.124870in}{1.201426in}}{\pgfqpoint{2.120752in}{1.197308in}}%
\pgfpathcurveto{\pgfqpoint{2.116634in}{1.193190in}}{\pgfqpoint{2.114320in}{1.187604in}}{\pgfqpoint{2.114320in}{1.181780in}}%
\pgfpathcurveto{\pgfqpoint{2.114320in}{1.175956in}}{\pgfqpoint{2.116634in}{1.170370in}}{\pgfqpoint{2.120752in}{1.166251in}}%
\pgfpathcurveto{\pgfqpoint{2.124870in}{1.162133in}}{\pgfqpoint{2.130456in}{1.159819in}}{\pgfqpoint{2.136280in}{1.159819in}}%
\pgfpathlineto{\pgfqpoint{2.136280in}{1.159819in}}%
\pgfpathclose%
\pgfusepath{stroke,fill}%
\end{pgfscope}%
\begin{pgfscope}%
\pgfpathrectangle{\pgfqpoint{0.100000in}{0.209042in}}{\pgfqpoint{2.906250in}{2.906250in}}%
\pgfusepath{clip}%
\pgfsetbuttcap%
\pgfsetroundjoin%
\definecolor{currentfill}{rgb}{0.013708,0.011771,0.068667}%
\pgfsetfillcolor{currentfill}%
\pgfsetfillopacity{0.800000}%
\pgfsetlinewidth{1.003750pt}%
\definecolor{currentstroke}{rgb}{0.013708,0.011771,0.068667}%
\pgfsetstrokecolor{currentstroke}%
\pgfsetstrokeopacity{0.800000}%
\pgfsetdash{}{0pt}%
\pgfpathmoveto{\pgfqpoint{1.605219in}{1.772658in}}%
\pgfpathcurveto{\pgfqpoint{1.611043in}{1.772658in}}{\pgfqpoint{1.616629in}{1.774972in}}{\pgfqpoint{1.620747in}{1.779090in}}%
\pgfpathcurveto{\pgfqpoint{1.624865in}{1.783208in}}{\pgfqpoint{1.627179in}{1.788794in}}{\pgfqpoint{1.627179in}{1.794618in}}%
\pgfpathcurveto{\pgfqpoint{1.627179in}{1.800442in}}{\pgfqpoint{1.624865in}{1.806028in}}{\pgfqpoint{1.620747in}{1.810146in}}%
\pgfpathcurveto{\pgfqpoint{1.616629in}{1.814264in}}{\pgfqpoint{1.611043in}{1.816578in}}{\pgfqpoint{1.605219in}{1.816578in}}%
\pgfpathcurveto{\pgfqpoint{1.599395in}{1.816578in}}{\pgfqpoint{1.593809in}{1.814264in}}{\pgfqpoint{1.589691in}{1.810146in}}%
\pgfpathcurveto{\pgfqpoint{1.585573in}{1.806028in}}{\pgfqpoint{1.583259in}{1.800442in}}{\pgfqpoint{1.583259in}{1.794618in}}%
\pgfpathcurveto{\pgfqpoint{1.583259in}{1.788794in}}{\pgfqpoint{1.585573in}{1.783208in}}{\pgfqpoint{1.589691in}{1.779090in}}%
\pgfpathcurveto{\pgfqpoint{1.593809in}{1.774972in}}{\pgfqpoint{1.599395in}{1.772658in}}{\pgfqpoint{1.605219in}{1.772658in}}%
\pgfpathlineto{\pgfqpoint{1.605219in}{1.772658in}}%
\pgfpathclose%
\pgfusepath{stroke,fill}%
\end{pgfscope}%
\begin{pgfscope}%
\pgfpathrectangle{\pgfqpoint{0.100000in}{0.209042in}}{\pgfqpoint{2.906250in}{2.906250in}}%
\pgfusepath{clip}%
\pgfsetbuttcap%
\pgfsetroundjoin%
\definecolor{currentfill}{rgb}{0.060949,0.044794,0.176129}%
\pgfsetfillcolor{currentfill}%
\pgfsetfillopacity{0.800000}%
\pgfsetlinewidth{1.003750pt}%
\definecolor{currentstroke}{rgb}{0.060949,0.044794,0.176129}%
\pgfsetstrokecolor{currentstroke}%
\pgfsetstrokeopacity{0.800000}%
\pgfsetdash{}{0pt}%
\pgfpathmoveto{\pgfqpoint{1.488061in}{1.726048in}}%
\pgfpathcurveto{\pgfqpoint{1.493885in}{1.726048in}}{\pgfqpoint{1.499471in}{1.728362in}}{\pgfqpoint{1.503589in}{1.732480in}}%
\pgfpathcurveto{\pgfqpoint{1.507707in}{1.736598in}}{\pgfqpoint{1.510021in}{1.742185in}}{\pgfqpoint{1.510021in}{1.748009in}}%
\pgfpathcurveto{\pgfqpoint{1.510021in}{1.753832in}}{\pgfqpoint{1.507707in}{1.759419in}}{\pgfqpoint{1.503589in}{1.763537in}}%
\pgfpathcurveto{\pgfqpoint{1.499471in}{1.767655in}}{\pgfqpoint{1.493885in}{1.769969in}}{\pgfqpoint{1.488061in}{1.769969in}}%
\pgfpathcurveto{\pgfqpoint{1.482237in}{1.769969in}}{\pgfqpoint{1.476651in}{1.767655in}}{\pgfqpoint{1.472533in}{1.763537in}}%
\pgfpathcurveto{\pgfqpoint{1.468415in}{1.759419in}}{\pgfqpoint{1.466101in}{1.753832in}}{\pgfqpoint{1.466101in}{1.748009in}}%
\pgfpathcurveto{\pgfqpoint{1.466101in}{1.742185in}}{\pgfqpoint{1.468415in}{1.736598in}}{\pgfqpoint{1.472533in}{1.732480in}}%
\pgfpathcurveto{\pgfqpoint{1.476651in}{1.728362in}}{\pgfqpoint{1.482237in}{1.726048in}}{\pgfqpoint{1.488061in}{1.726048in}}%
\pgfpathlineto{\pgfqpoint{1.488061in}{1.726048in}}%
\pgfpathclose%
\pgfusepath{stroke,fill}%
\end{pgfscope}%
\begin{pgfscope}%
\pgfpathrectangle{\pgfqpoint{0.100000in}{0.209042in}}{\pgfqpoint{2.906250in}{2.906250in}}%
\pgfusepath{clip}%
\pgfsetbuttcap%
\pgfsetroundjoin%
\definecolor{currentfill}{rgb}{0.697098,0.208501,0.480835}%
\pgfsetfillcolor{currentfill}%
\pgfsetfillopacity{0.800000}%
\pgfsetlinewidth{1.003750pt}%
\definecolor{currentstroke}{rgb}{0.697098,0.208501,0.480835}%
\pgfsetstrokecolor{currentstroke}%
\pgfsetstrokeopacity{0.800000}%
\pgfsetdash{}{0pt}%
\pgfpathmoveto{\pgfqpoint{1.579651in}{2.158797in}}%
\pgfpathcurveto{\pgfqpoint{1.585474in}{2.158797in}}{\pgfqpoint{1.591061in}{2.161111in}}{\pgfqpoint{1.595179in}{2.165229in}}%
\pgfpathcurveto{\pgfqpoint{1.599297in}{2.169347in}}{\pgfqpoint{1.601611in}{2.174933in}}{\pgfqpoint{1.601611in}{2.180757in}}%
\pgfpathcurveto{\pgfqpoint{1.601611in}{2.186581in}}{\pgfqpoint{1.599297in}{2.192167in}}{\pgfqpoint{1.595179in}{2.196285in}}%
\pgfpathcurveto{\pgfqpoint{1.591061in}{2.200403in}}{\pgfqpoint{1.585474in}{2.202717in}}{\pgfqpoint{1.579651in}{2.202717in}}%
\pgfpathcurveto{\pgfqpoint{1.573827in}{2.202717in}}{\pgfqpoint{1.568240in}{2.200403in}}{\pgfqpoint{1.564122in}{2.196285in}}%
\pgfpathcurveto{\pgfqpoint{1.560004in}{2.192167in}}{\pgfqpoint{1.557690in}{2.186581in}}{\pgfqpoint{1.557690in}{2.180757in}}%
\pgfpathcurveto{\pgfqpoint{1.557690in}{2.174933in}}{\pgfqpoint{1.560004in}{2.169347in}}{\pgfqpoint{1.564122in}{2.165229in}}%
\pgfpathcurveto{\pgfqpoint{1.568240in}{2.161111in}}{\pgfqpoint{1.573827in}{2.158797in}}{\pgfqpoint{1.579651in}{2.158797in}}%
\pgfpathlineto{\pgfqpoint{1.579651in}{2.158797in}}%
\pgfpathclose%
\pgfusepath{stroke,fill}%
\end{pgfscope}%
\begin{pgfscope}%
\pgfpathrectangle{\pgfqpoint{0.100000in}{0.209042in}}{\pgfqpoint{2.906250in}{2.906250in}}%
\pgfusepath{clip}%
\pgfsetbuttcap%
\pgfsetroundjoin%
\definecolor{currentfill}{rgb}{0.118405,0.066479,0.286321}%
\pgfsetfillcolor{currentfill}%
\pgfsetfillopacity{0.800000}%
\pgfsetlinewidth{1.003750pt}%
\definecolor{currentstroke}{rgb}{0.118405,0.066479,0.286321}%
\pgfsetstrokecolor{currentstroke}%
\pgfsetstrokeopacity{0.800000}%
\pgfsetdash{}{0pt}%
\pgfpathmoveto{\pgfqpoint{1.565983in}{1.851782in}}%
\pgfpathcurveto{\pgfqpoint{1.571807in}{1.851782in}}{\pgfqpoint{1.577393in}{1.854096in}}{\pgfqpoint{1.581512in}{1.858214in}}%
\pgfpathcurveto{\pgfqpoint{1.585630in}{1.862333in}}{\pgfqpoint{1.587944in}{1.867919in}}{\pgfqpoint{1.587944in}{1.873743in}}%
\pgfpathcurveto{\pgfqpoint{1.587944in}{1.879567in}}{\pgfqpoint{1.585630in}{1.885153in}}{\pgfqpoint{1.581512in}{1.889271in}}%
\pgfpathcurveto{\pgfqpoint{1.577393in}{1.893389in}}{\pgfqpoint{1.571807in}{1.895703in}}{\pgfqpoint{1.565983in}{1.895703in}}%
\pgfpathcurveto{\pgfqpoint{1.560159in}{1.895703in}}{\pgfqpoint{1.554573in}{1.893389in}}{\pgfqpoint{1.550455in}{1.889271in}}%
\pgfpathcurveto{\pgfqpoint{1.546337in}{1.885153in}}{\pgfqpoint{1.544023in}{1.879567in}}{\pgfqpoint{1.544023in}{1.873743in}}%
\pgfpathcurveto{\pgfqpoint{1.544023in}{1.867919in}}{\pgfqpoint{1.546337in}{1.862333in}}{\pgfqpoint{1.550455in}{1.858214in}}%
\pgfpathcurveto{\pgfqpoint{1.554573in}{1.854096in}}{\pgfqpoint{1.560159in}{1.851782in}}{\pgfqpoint{1.565983in}{1.851782in}}%
\pgfpathlineto{\pgfqpoint{1.565983in}{1.851782in}}%
\pgfpathclose%
\pgfusepath{stroke,fill}%
\end{pgfscope}%
\begin{pgfscope}%
\pgfpathrectangle{\pgfqpoint{0.100000in}{0.209042in}}{\pgfqpoint{2.906250in}{2.906250in}}%
\pgfusepath{clip}%
\pgfsetbuttcap%
\pgfsetroundjoin%
\definecolor{currentfill}{rgb}{0.107899,0.064335,0.267289}%
\pgfsetfillcolor{currentfill}%
\pgfsetfillopacity{0.800000}%
\pgfsetlinewidth{1.003750pt}%
\definecolor{currentstroke}{rgb}{0.107899,0.064335,0.267289}%
\pgfsetstrokecolor{currentstroke}%
\pgfsetstrokeopacity{0.800000}%
\pgfsetdash{}{0pt}%
\pgfpathmoveto{\pgfqpoint{1.689168in}{1.838562in}}%
\pgfpathcurveto{\pgfqpoint{1.694992in}{1.838562in}}{\pgfqpoint{1.700579in}{1.840876in}}{\pgfqpoint{1.704697in}{1.844994in}}%
\pgfpathcurveto{\pgfqpoint{1.708815in}{1.849112in}}{\pgfqpoint{1.711129in}{1.854698in}}{\pgfqpoint{1.711129in}{1.860522in}}%
\pgfpathcurveto{\pgfqpoint{1.711129in}{1.866346in}}{\pgfqpoint{1.708815in}{1.871932in}}{\pgfqpoint{1.704697in}{1.876051in}}%
\pgfpathcurveto{\pgfqpoint{1.700579in}{1.880169in}}{\pgfqpoint{1.694992in}{1.882483in}}{\pgfqpoint{1.689168in}{1.882483in}}%
\pgfpathcurveto{\pgfqpoint{1.683344in}{1.882483in}}{\pgfqpoint{1.677758in}{1.880169in}}{\pgfqpoint{1.673640in}{1.876051in}}%
\pgfpathcurveto{\pgfqpoint{1.669522in}{1.871932in}}{\pgfqpoint{1.667208in}{1.866346in}}{\pgfqpoint{1.667208in}{1.860522in}}%
\pgfpathcurveto{\pgfqpoint{1.667208in}{1.854698in}}{\pgfqpoint{1.669522in}{1.849112in}}{\pgfqpoint{1.673640in}{1.844994in}}%
\pgfpathcurveto{\pgfqpoint{1.677758in}{1.840876in}}{\pgfqpoint{1.683344in}{1.838562in}}{\pgfqpoint{1.689168in}{1.838562in}}%
\pgfpathlineto{\pgfqpoint{1.689168in}{1.838562in}}%
\pgfpathclose%
\pgfusepath{stroke,fill}%
\end{pgfscope}%
\begin{pgfscope}%
\pgfpathrectangle{\pgfqpoint{0.100000in}{0.209042in}}{\pgfqpoint{2.906250in}{2.906250in}}%
\pgfusepath{clip}%
\pgfsetbuttcap%
\pgfsetroundjoin%
\definecolor{currentfill}{rgb}{0.457386,0.127522,0.507448}%
\pgfsetfillcolor{currentfill}%
\pgfsetfillopacity{0.800000}%
\pgfsetlinewidth{1.003750pt}%
\definecolor{currentstroke}{rgb}{0.457386,0.127522,0.507448}%
\pgfsetstrokecolor{currentstroke}%
\pgfsetstrokeopacity{0.800000}%
\pgfsetdash{}{0pt}%
\pgfpathmoveto{\pgfqpoint{1.776462in}{1.389874in}}%
\pgfpathcurveto{\pgfqpoint{1.782286in}{1.389874in}}{\pgfqpoint{1.787872in}{1.392188in}}{\pgfqpoint{1.791990in}{1.396306in}}%
\pgfpathcurveto{\pgfqpoint{1.796108in}{1.400424in}}{\pgfqpoint{1.798422in}{1.406010in}}{\pgfqpoint{1.798422in}{1.411834in}}%
\pgfpathcurveto{\pgfqpoint{1.798422in}{1.417658in}}{\pgfqpoint{1.796108in}{1.423244in}}{\pgfqpoint{1.791990in}{1.427362in}}%
\pgfpathcurveto{\pgfqpoint{1.787872in}{1.431480in}}{\pgfqpoint{1.782286in}{1.433794in}}{\pgfqpoint{1.776462in}{1.433794in}}%
\pgfpathcurveto{\pgfqpoint{1.770638in}{1.433794in}}{\pgfqpoint{1.765052in}{1.431480in}}{\pgfqpoint{1.760934in}{1.427362in}}%
\pgfpathcurveto{\pgfqpoint{1.756815in}{1.423244in}}{\pgfqpoint{1.754502in}{1.417658in}}{\pgfqpoint{1.754502in}{1.411834in}}%
\pgfpathcurveto{\pgfqpoint{1.754502in}{1.406010in}}{\pgfqpoint{1.756815in}{1.400424in}}{\pgfqpoint{1.760934in}{1.396306in}}%
\pgfpathcurveto{\pgfqpoint{1.765052in}{1.392188in}}{\pgfqpoint{1.770638in}{1.389874in}}{\pgfqpoint{1.776462in}{1.389874in}}%
\pgfpathlineto{\pgfqpoint{1.776462in}{1.389874in}}%
\pgfpathclose%
\pgfusepath{stroke,fill}%
\end{pgfscope}%
\begin{pgfscope}%
\pgfpathrectangle{\pgfqpoint{0.100000in}{0.209042in}}{\pgfqpoint{2.906250in}{2.906250in}}%
\pgfusepath{clip}%
\pgfsetbuttcap%
\pgfsetroundjoin%
\definecolor{currentfill}{rgb}{0.021692,0.018320,0.092610}%
\pgfsetfillcolor{currentfill}%
\pgfsetfillopacity{0.800000}%
\pgfsetlinewidth{1.003750pt}%
\definecolor{currentstroke}{rgb}{0.021692,0.018320,0.092610}%
\pgfsetstrokecolor{currentstroke}%
\pgfsetstrokeopacity{0.800000}%
\pgfsetdash{}{0pt}%
\pgfpathmoveto{\pgfqpoint{1.608123in}{1.781002in}}%
\pgfpathcurveto{\pgfqpoint{1.613947in}{1.781002in}}{\pgfqpoint{1.619534in}{1.783316in}}{\pgfqpoint{1.623652in}{1.787434in}}%
\pgfpathcurveto{\pgfqpoint{1.627770in}{1.791552in}}{\pgfqpoint{1.630084in}{1.797138in}}{\pgfqpoint{1.630084in}{1.802962in}}%
\pgfpathcurveto{\pgfqpoint{1.630084in}{1.808786in}}{\pgfqpoint{1.627770in}{1.814372in}}{\pgfqpoint{1.623652in}{1.818490in}}%
\pgfpathcurveto{\pgfqpoint{1.619534in}{1.822609in}}{\pgfqpoint{1.613947in}{1.824922in}}{\pgfqpoint{1.608123in}{1.824922in}}%
\pgfpathcurveto{\pgfqpoint{1.602299in}{1.824922in}}{\pgfqpoint{1.596713in}{1.822609in}}{\pgfqpoint{1.592595in}{1.818490in}}%
\pgfpathcurveto{\pgfqpoint{1.588477in}{1.814372in}}{\pgfqpoint{1.586163in}{1.808786in}}{\pgfqpoint{1.586163in}{1.802962in}}%
\pgfpathcurveto{\pgfqpoint{1.586163in}{1.797138in}}{\pgfqpoint{1.588477in}{1.791552in}}{\pgfqpoint{1.592595in}{1.787434in}}%
\pgfpathcurveto{\pgfqpoint{1.596713in}{1.783316in}}{\pgfqpoint{1.602299in}{1.781002in}}{\pgfqpoint{1.608123in}{1.781002in}}%
\pgfpathlineto{\pgfqpoint{1.608123in}{1.781002in}}%
\pgfpathclose%
\pgfusepath{stroke,fill}%
\end{pgfscope}%
\begin{pgfscope}%
\pgfpathrectangle{\pgfqpoint{0.100000in}{0.209042in}}{\pgfqpoint{2.906250in}{2.906250in}}%
\pgfusepath{clip}%
\pgfsetbuttcap%
\pgfsetroundjoin%
\definecolor{currentfill}{rgb}{0.007588,0.006356,0.044973}%
\pgfsetfillcolor{currentfill}%
\pgfsetfillopacity{0.800000}%
\pgfsetlinewidth{1.003750pt}%
\definecolor{currentstroke}{rgb}{0.007588,0.006356,0.044973}%
\pgfsetstrokecolor{currentstroke}%
\pgfsetstrokeopacity{0.800000}%
\pgfsetdash{}{0pt}%
\pgfpathmoveto{\pgfqpoint{1.577126in}{1.666424in}}%
\pgfpathcurveto{\pgfqpoint{1.582950in}{1.666424in}}{\pgfqpoint{1.588536in}{1.668738in}}{\pgfqpoint{1.592654in}{1.672856in}}%
\pgfpathcurveto{\pgfqpoint{1.596772in}{1.676974in}}{\pgfqpoint{1.599086in}{1.682560in}}{\pgfqpoint{1.599086in}{1.688384in}}%
\pgfpathcurveto{\pgfqpoint{1.599086in}{1.694208in}}{\pgfqpoint{1.596772in}{1.699794in}}{\pgfqpoint{1.592654in}{1.703912in}}%
\pgfpathcurveto{\pgfqpoint{1.588536in}{1.708030in}}{\pgfqpoint{1.582950in}{1.710344in}}{\pgfqpoint{1.577126in}{1.710344in}}%
\pgfpathcurveto{\pgfqpoint{1.571302in}{1.710344in}}{\pgfqpoint{1.565716in}{1.708030in}}{\pgfqpoint{1.561598in}{1.703912in}}%
\pgfpathcurveto{\pgfqpoint{1.557480in}{1.699794in}}{\pgfqpoint{1.555166in}{1.694208in}}{\pgfqpoint{1.555166in}{1.688384in}}%
\pgfpathcurveto{\pgfqpoint{1.555166in}{1.682560in}}{\pgfqpoint{1.557480in}{1.676974in}}{\pgfqpoint{1.561598in}{1.672856in}}%
\pgfpathcurveto{\pgfqpoint{1.565716in}{1.668738in}}{\pgfqpoint{1.571302in}{1.666424in}}{\pgfqpoint{1.577126in}{1.666424in}}%
\pgfpathlineto{\pgfqpoint{1.577126in}{1.666424in}}%
\pgfpathclose%
\pgfusepath{stroke,fill}%
\end{pgfscope}%
\begin{pgfscope}%
\pgfpathrectangle{\pgfqpoint{0.100000in}{0.209042in}}{\pgfqpoint{2.906250in}{2.906250in}}%
\pgfusepath{clip}%
\pgfsetbuttcap%
\pgfsetroundjoin%
\definecolor{currentfill}{rgb}{0.024792,0.020715,0.100676}%
\pgfsetfillcolor{currentfill}%
\pgfsetfillopacity{0.800000}%
\pgfsetlinewidth{1.003750pt}%
\definecolor{currentstroke}{rgb}{0.024792,0.020715,0.100676}%
\pgfsetstrokecolor{currentstroke}%
\pgfsetstrokeopacity{0.800000}%
\pgfsetdash{}{0pt}%
\pgfpathmoveto{\pgfqpoint{1.530832in}{1.726121in}}%
\pgfpathcurveto{\pgfqpoint{1.536656in}{1.726121in}}{\pgfqpoint{1.542242in}{1.728435in}}{\pgfqpoint{1.546360in}{1.732553in}}%
\pgfpathcurveto{\pgfqpoint{1.550478in}{1.736672in}}{\pgfqpoint{1.552792in}{1.742258in}}{\pgfqpoint{1.552792in}{1.748082in}}%
\pgfpathcurveto{\pgfqpoint{1.552792in}{1.753906in}}{\pgfqpoint{1.550478in}{1.759492in}}{\pgfqpoint{1.546360in}{1.763610in}}%
\pgfpathcurveto{\pgfqpoint{1.542242in}{1.767728in}}{\pgfqpoint{1.536656in}{1.770042in}}{\pgfqpoint{1.530832in}{1.770042in}}%
\pgfpathcurveto{\pgfqpoint{1.525008in}{1.770042in}}{\pgfqpoint{1.519422in}{1.767728in}}{\pgfqpoint{1.515303in}{1.763610in}}%
\pgfpathcurveto{\pgfqpoint{1.511185in}{1.759492in}}{\pgfqpoint{1.508871in}{1.753906in}}{\pgfqpoint{1.508871in}{1.748082in}}%
\pgfpathcurveto{\pgfqpoint{1.508871in}{1.742258in}}{\pgfqpoint{1.511185in}{1.736672in}}{\pgfqpoint{1.515303in}{1.732553in}}%
\pgfpathcurveto{\pgfqpoint{1.519422in}{1.728435in}}{\pgfqpoint{1.525008in}{1.726121in}}{\pgfqpoint{1.530832in}{1.726121in}}%
\pgfpathlineto{\pgfqpoint{1.530832in}{1.726121in}}%
\pgfpathclose%
\pgfusepath{stroke,fill}%
\end{pgfscope}%
\begin{pgfscope}%
\pgfpathrectangle{\pgfqpoint{0.100000in}{0.209042in}}{\pgfqpoint{2.906250in}{2.906250in}}%
\pgfusepath{clip}%
\pgfsetbuttcap%
\pgfsetroundjoin%
\definecolor{currentfill}{rgb}{0.039608,0.031090,0.133515}%
\pgfsetfillcolor{currentfill}%
\pgfsetfillopacity{0.800000}%
\pgfsetlinewidth{1.003750pt}%
\definecolor{currentstroke}{rgb}{0.039608,0.031090,0.133515}%
\pgfsetstrokecolor{currentstroke}%
\pgfsetstrokeopacity{0.800000}%
\pgfsetdash{}{0pt}%
\pgfpathmoveto{\pgfqpoint{1.692202in}{1.780518in}}%
\pgfpathcurveto{\pgfqpoint{1.698026in}{1.780518in}}{\pgfqpoint{1.703612in}{1.782832in}}{\pgfqpoint{1.707730in}{1.786950in}}%
\pgfpathcurveto{\pgfqpoint{1.711848in}{1.791068in}}{\pgfqpoint{1.714162in}{1.796654in}}{\pgfqpoint{1.714162in}{1.802478in}}%
\pgfpathcurveto{\pgfqpoint{1.714162in}{1.808302in}}{\pgfqpoint{1.711848in}{1.813888in}}{\pgfqpoint{1.707730in}{1.818007in}}%
\pgfpathcurveto{\pgfqpoint{1.703612in}{1.822125in}}{\pgfqpoint{1.698026in}{1.824439in}}{\pgfqpoint{1.692202in}{1.824439in}}%
\pgfpathcurveto{\pgfqpoint{1.686378in}{1.824439in}}{\pgfqpoint{1.680792in}{1.822125in}}{\pgfqpoint{1.676674in}{1.818007in}}%
\pgfpathcurveto{\pgfqpoint{1.672556in}{1.813888in}}{\pgfqpoint{1.670242in}{1.808302in}}{\pgfqpoint{1.670242in}{1.802478in}}%
\pgfpathcurveto{\pgfqpoint{1.670242in}{1.796654in}}{\pgfqpoint{1.672556in}{1.791068in}}{\pgfqpoint{1.676674in}{1.786950in}}%
\pgfpathcurveto{\pgfqpoint{1.680792in}{1.782832in}}{\pgfqpoint{1.686378in}{1.780518in}}{\pgfqpoint{1.692202in}{1.780518in}}%
\pgfpathlineto{\pgfqpoint{1.692202in}{1.780518in}}%
\pgfpathclose%
\pgfusepath{stroke,fill}%
\end{pgfscope}%
\begin{pgfscope}%
\pgfpathrectangle{\pgfqpoint{0.100000in}{0.209042in}}{\pgfqpoint{2.906250in}{2.906250in}}%
\pgfusepath{clip}%
\pgfsetbuttcap%
\pgfsetroundjoin%
\definecolor{currentfill}{rgb}{0.140858,0.068654,0.324538}%
\pgfsetfillcolor{currentfill}%
\pgfsetfillopacity{0.800000}%
\pgfsetlinewidth{1.003750pt}%
\definecolor{currentstroke}{rgb}{0.140858,0.068654,0.324538}%
\pgfsetstrokecolor{currentstroke}%
\pgfsetstrokeopacity{0.800000}%
\pgfsetdash{}{0pt}%
\pgfpathmoveto{\pgfqpoint{1.834126in}{1.698198in}}%
\pgfpathcurveto{\pgfqpoint{1.839950in}{1.698198in}}{\pgfqpoint{1.845536in}{1.700512in}}{\pgfqpoint{1.849654in}{1.704630in}}%
\pgfpathcurveto{\pgfqpoint{1.853772in}{1.708749in}}{\pgfqpoint{1.856086in}{1.714335in}}{\pgfqpoint{1.856086in}{1.720159in}}%
\pgfpathcurveto{\pgfqpoint{1.856086in}{1.725983in}}{\pgfqpoint{1.853772in}{1.731569in}}{\pgfqpoint{1.849654in}{1.735687in}}%
\pgfpathcurveto{\pgfqpoint{1.845536in}{1.739805in}}{\pgfqpoint{1.839950in}{1.742119in}}{\pgfqpoint{1.834126in}{1.742119in}}%
\pgfpathcurveto{\pgfqpoint{1.828302in}{1.742119in}}{\pgfqpoint{1.822716in}{1.739805in}}{\pgfqpoint{1.818598in}{1.735687in}}%
\pgfpathcurveto{\pgfqpoint{1.814479in}{1.731569in}}{\pgfqpoint{1.812166in}{1.725983in}}{\pgfqpoint{1.812166in}{1.720159in}}%
\pgfpathcurveto{\pgfqpoint{1.812166in}{1.714335in}}{\pgfqpoint{1.814479in}{1.708749in}}{\pgfqpoint{1.818598in}{1.704630in}}%
\pgfpathcurveto{\pgfqpoint{1.822716in}{1.700512in}}{\pgfqpoint{1.828302in}{1.698198in}}{\pgfqpoint{1.834126in}{1.698198in}}%
\pgfpathlineto{\pgfqpoint{1.834126in}{1.698198in}}%
\pgfpathclose%
\pgfusepath{stroke,fill}%
\end{pgfscope}%
\begin{pgfscope}%
\pgfpathrectangle{\pgfqpoint{0.100000in}{0.209042in}}{\pgfqpoint{2.906250in}{2.906250in}}%
\pgfusepath{clip}%
\pgfsetbuttcap%
\pgfsetroundjoin%
\definecolor{currentfill}{rgb}{0.002258,0.001295,0.018331}%
\pgfsetfillcolor{currentfill}%
\pgfsetfillopacity{0.800000}%
\pgfsetlinewidth{1.003750pt}%
\definecolor{currentstroke}{rgb}{0.002258,0.001295,0.018331}%
\pgfsetstrokecolor{currentstroke}%
\pgfsetstrokeopacity{0.800000}%
\pgfsetdash{}{0pt}%
\pgfpathmoveto{\pgfqpoint{1.576467in}{1.698796in}}%
\pgfpathcurveto{\pgfqpoint{1.582291in}{1.698796in}}{\pgfqpoint{1.587877in}{1.701109in}}{\pgfqpoint{1.591995in}{1.705228in}}%
\pgfpathcurveto{\pgfqpoint{1.596113in}{1.709346in}}{\pgfqpoint{1.598427in}{1.714932in}}{\pgfqpoint{1.598427in}{1.720756in}}%
\pgfpathcurveto{\pgfqpoint{1.598427in}{1.726580in}}{\pgfqpoint{1.596113in}{1.732166in}}{\pgfqpoint{1.591995in}{1.736284in}}%
\pgfpathcurveto{\pgfqpoint{1.587877in}{1.740402in}}{\pgfqpoint{1.582291in}{1.742716in}}{\pgfqpoint{1.576467in}{1.742716in}}%
\pgfpathcurveto{\pgfqpoint{1.570643in}{1.742716in}}{\pgfqpoint{1.565057in}{1.740402in}}{\pgfqpoint{1.560938in}{1.736284in}}%
\pgfpathcurveto{\pgfqpoint{1.556820in}{1.732166in}}{\pgfqpoint{1.554506in}{1.726580in}}{\pgfqpoint{1.554506in}{1.720756in}}%
\pgfpathcurveto{\pgfqpoint{1.554506in}{1.714932in}}{\pgfqpoint{1.556820in}{1.709346in}}{\pgfqpoint{1.560938in}{1.705228in}}%
\pgfpathcurveto{\pgfqpoint{1.565057in}{1.701109in}}{\pgfqpoint{1.570643in}{1.698796in}}{\pgfqpoint{1.576467in}{1.698796in}}%
\pgfpathlineto{\pgfqpoint{1.576467in}{1.698796in}}%
\pgfpathclose%
\pgfusepath{stroke,fill}%
\end{pgfscope}%
\begin{pgfscope}%
\pgfpathrectangle{\pgfqpoint{0.100000in}{0.209042in}}{\pgfqpoint{2.906250in}{2.906250in}}%
\pgfusepath{clip}%
\pgfsetbuttcap%
\pgfsetroundjoin%
\definecolor{currentfill}{rgb}{0.754737,0.228772,0.462509}%
\pgfsetfillcolor{currentfill}%
\pgfsetfillopacity{0.800000}%
\pgfsetlinewidth{1.003750pt}%
\definecolor{currentstroke}{rgb}{0.754737,0.228772,0.462509}%
\pgfsetstrokecolor{currentstroke}%
\pgfsetstrokeopacity{0.800000}%
\pgfsetdash{}{0pt}%
\pgfpathmoveto{\pgfqpoint{1.912793in}{1.264681in}}%
\pgfpathcurveto{\pgfqpoint{1.918617in}{1.264681in}}{\pgfqpoint{1.924203in}{1.266995in}}{\pgfqpoint{1.928322in}{1.271113in}}%
\pgfpathcurveto{\pgfqpoint{1.932440in}{1.275232in}}{\pgfqpoint{1.934754in}{1.280818in}}{\pgfqpoint{1.934754in}{1.286642in}}%
\pgfpathcurveto{\pgfqpoint{1.934754in}{1.292466in}}{\pgfqpoint{1.932440in}{1.298052in}}{\pgfqpoint{1.928322in}{1.302170in}}%
\pgfpathcurveto{\pgfqpoint{1.924203in}{1.306288in}}{\pgfqpoint{1.918617in}{1.308602in}}{\pgfqpoint{1.912793in}{1.308602in}}%
\pgfpathcurveto{\pgfqpoint{1.906969in}{1.308602in}}{\pgfqpoint{1.901383in}{1.306288in}}{\pgfqpoint{1.897265in}{1.302170in}}%
\pgfpathcurveto{\pgfqpoint{1.893147in}{1.298052in}}{\pgfqpoint{1.890833in}{1.292466in}}{\pgfqpoint{1.890833in}{1.286642in}}%
\pgfpathcurveto{\pgfqpoint{1.890833in}{1.280818in}}{\pgfqpoint{1.893147in}{1.275232in}}{\pgfqpoint{1.897265in}{1.271113in}}%
\pgfpathcurveto{\pgfqpoint{1.901383in}{1.266995in}}{\pgfqpoint{1.906969in}{1.264681in}}{\pgfqpoint{1.912793in}{1.264681in}}%
\pgfpathlineto{\pgfqpoint{1.912793in}{1.264681in}}%
\pgfpathclose%
\pgfusepath{stroke,fill}%
\end{pgfscope}%
\begin{pgfscope}%
\pgfpathrectangle{\pgfqpoint{0.100000in}{0.209042in}}{\pgfqpoint{2.906250in}{2.906250in}}%
\pgfusepath{clip}%
\pgfsetbuttcap%
\pgfsetroundjoin%
\definecolor{currentfill}{rgb}{0.868793,0.287728,0.409303}%
\pgfsetfillcolor{currentfill}%
\pgfsetfillopacity{0.800000}%
\pgfsetlinewidth{1.003750pt}%
\definecolor{currentstroke}{rgb}{0.868793,0.287728,0.409303}%
\pgfsetstrokecolor{currentstroke}%
\pgfsetstrokeopacity{0.800000}%
\pgfsetdash{}{0pt}%
\pgfpathmoveto{\pgfqpoint{1.637907in}{2.247650in}}%
\pgfpathcurveto{\pgfqpoint{1.643731in}{2.247650in}}{\pgfqpoint{1.649317in}{2.249964in}}{\pgfqpoint{1.653435in}{2.254082in}}%
\pgfpathcurveto{\pgfqpoint{1.657553in}{2.258200in}}{\pgfqpoint{1.659867in}{2.263786in}}{\pgfqpoint{1.659867in}{2.269610in}}%
\pgfpathcurveto{\pgfqpoint{1.659867in}{2.275434in}}{\pgfqpoint{1.657553in}{2.281020in}}{\pgfqpoint{1.653435in}{2.285139in}}%
\pgfpathcurveto{\pgfqpoint{1.649317in}{2.289257in}}{\pgfqpoint{1.643731in}{2.291571in}}{\pgfqpoint{1.637907in}{2.291571in}}%
\pgfpathcurveto{\pgfqpoint{1.632083in}{2.291571in}}{\pgfqpoint{1.626497in}{2.289257in}}{\pgfqpoint{1.622379in}{2.285139in}}%
\pgfpathcurveto{\pgfqpoint{1.618260in}{2.281020in}}{\pgfqpoint{1.615947in}{2.275434in}}{\pgfqpoint{1.615947in}{2.269610in}}%
\pgfpathcurveto{\pgfqpoint{1.615947in}{2.263786in}}{\pgfqpoint{1.618260in}{2.258200in}}{\pgfqpoint{1.622379in}{2.254082in}}%
\pgfpathcurveto{\pgfqpoint{1.626497in}{2.249964in}}{\pgfqpoint{1.632083in}{2.247650in}}{\pgfqpoint{1.637907in}{2.247650in}}%
\pgfpathlineto{\pgfqpoint{1.637907in}{2.247650in}}%
\pgfpathclose%
\pgfusepath{stroke,fill}%
\end{pgfscope}%
\begin{pgfscope}%
\pgfpathrectangle{\pgfqpoint{0.100000in}{0.209042in}}{\pgfqpoint{2.906250in}{2.906250in}}%
\pgfusepath{clip}%
\pgfsetbuttcap%
\pgfsetroundjoin%
\definecolor{currentfill}{rgb}{0.204935,0.062907,0.411514}%
\pgfsetfillcolor{currentfill}%
\pgfsetfillopacity{0.800000}%
\pgfsetlinewidth{1.003750pt}%
\definecolor{currentstroke}{rgb}{0.204935,0.062907,0.411514}%
\pgfsetstrokecolor{currentstroke}%
\pgfsetstrokeopacity{0.800000}%
\pgfsetdash{}{0pt}%
\pgfpathmoveto{\pgfqpoint{1.663458in}{1.504657in}}%
\pgfpathcurveto{\pgfqpoint{1.669282in}{1.504657in}}{\pgfqpoint{1.674868in}{1.506971in}}{\pgfqpoint{1.678986in}{1.511089in}}%
\pgfpathcurveto{\pgfqpoint{1.683104in}{1.515207in}}{\pgfqpoint{1.685418in}{1.520793in}}{\pgfqpoint{1.685418in}{1.526617in}}%
\pgfpathcurveto{\pgfqpoint{1.685418in}{1.532441in}}{\pgfqpoint{1.683104in}{1.538027in}}{\pgfqpoint{1.678986in}{1.542145in}}%
\pgfpathcurveto{\pgfqpoint{1.674868in}{1.546263in}}{\pgfqpoint{1.669282in}{1.548577in}}{\pgfqpoint{1.663458in}{1.548577in}}%
\pgfpathcurveto{\pgfqpoint{1.657634in}{1.548577in}}{\pgfqpoint{1.652048in}{1.546263in}}{\pgfqpoint{1.647930in}{1.542145in}}%
\pgfpathcurveto{\pgfqpoint{1.643812in}{1.538027in}}{\pgfqpoint{1.641498in}{1.532441in}}{\pgfqpoint{1.641498in}{1.526617in}}%
\pgfpathcurveto{\pgfqpoint{1.641498in}{1.520793in}}{\pgfqpoint{1.643812in}{1.515207in}}{\pgfqpoint{1.647930in}{1.511089in}}%
\pgfpathcurveto{\pgfqpoint{1.652048in}{1.506971in}}{\pgfqpoint{1.657634in}{1.504657in}}{\pgfqpoint{1.663458in}{1.504657in}}%
\pgfpathlineto{\pgfqpoint{1.663458in}{1.504657in}}%
\pgfpathclose%
\pgfusepath{stroke,fill}%
\end{pgfscope}%
\begin{pgfscope}%
\pgfpathrectangle{\pgfqpoint{0.100000in}{0.209042in}}{\pgfqpoint{2.906250in}{2.906250in}}%
\pgfusepath{clip}%
\pgfsetbuttcap%
\pgfsetroundjoin%
\definecolor{currentfill}{rgb}{0.378211,0.095332,0.500067}%
\pgfsetfillcolor{currentfill}%
\pgfsetfillopacity{0.800000}%
\pgfsetlinewidth{1.003750pt}%
\definecolor{currentstroke}{rgb}{0.378211,0.095332,0.500067}%
\pgfsetstrokecolor{currentstroke}%
\pgfsetstrokeopacity{0.800000}%
\pgfsetdash{}{0pt}%
\pgfpathmoveto{\pgfqpoint{1.924290in}{1.849808in}}%
\pgfpathcurveto{\pgfqpoint{1.930114in}{1.849808in}}{\pgfqpoint{1.935700in}{1.852122in}}{\pgfqpoint{1.939818in}{1.856240in}}%
\pgfpathcurveto{\pgfqpoint{1.943936in}{1.860359in}}{\pgfqpoint{1.946250in}{1.865945in}}{\pgfqpoint{1.946250in}{1.871769in}}%
\pgfpathcurveto{\pgfqpoint{1.946250in}{1.877593in}}{\pgfqpoint{1.943936in}{1.883179in}}{\pgfqpoint{1.939818in}{1.887297in}}%
\pgfpathcurveto{\pgfqpoint{1.935700in}{1.891415in}}{\pgfqpoint{1.930114in}{1.893729in}}{\pgfqpoint{1.924290in}{1.893729in}}%
\pgfpathcurveto{\pgfqpoint{1.918466in}{1.893729in}}{\pgfqpoint{1.912880in}{1.891415in}}{\pgfqpoint{1.908762in}{1.887297in}}%
\pgfpathcurveto{\pgfqpoint{1.904644in}{1.883179in}}{\pgfqpoint{1.902330in}{1.877593in}}{\pgfqpoint{1.902330in}{1.871769in}}%
\pgfpathcurveto{\pgfqpoint{1.902330in}{1.865945in}}{\pgfqpoint{1.904644in}{1.860359in}}{\pgfqpoint{1.908762in}{1.856240in}}%
\pgfpathcurveto{\pgfqpoint{1.912880in}{1.852122in}}{\pgfqpoint{1.918466in}{1.849808in}}{\pgfqpoint{1.924290in}{1.849808in}}%
\pgfpathlineto{\pgfqpoint{1.924290in}{1.849808in}}%
\pgfpathclose%
\pgfusepath{stroke,fill}%
\end{pgfscope}%
\begin{pgfscope}%
\pgfpathrectangle{\pgfqpoint{0.100000in}{0.209042in}}{\pgfqpoint{2.906250in}{2.906250in}}%
\pgfusepath{clip}%
\pgfsetbuttcap%
\pgfsetroundjoin%
\definecolor{currentfill}{rgb}{0.975082,0.468861,0.363111}%
\pgfsetfillcolor{currentfill}%
\pgfsetfillopacity{0.800000}%
\pgfsetlinewidth{1.003750pt}%
\definecolor{currentstroke}{rgb}{0.975082,0.468861,0.363111}%
\pgfsetstrokecolor{currentstroke}%
\pgfsetstrokeopacity{0.800000}%
\pgfsetdash{}{0pt}%
\pgfpathmoveto{\pgfqpoint{2.422149in}{1.623391in}}%
\pgfpathcurveto{\pgfqpoint{2.427973in}{1.623391in}}{\pgfqpoint{2.433559in}{1.625705in}}{\pgfqpoint{2.437677in}{1.629823in}}%
\pgfpathcurveto{\pgfqpoint{2.441795in}{1.633941in}}{\pgfqpoint{2.444109in}{1.639527in}}{\pgfqpoint{2.444109in}{1.645351in}}%
\pgfpathcurveto{\pgfqpoint{2.444109in}{1.651175in}}{\pgfqpoint{2.441795in}{1.656761in}}{\pgfqpoint{2.437677in}{1.660879in}}%
\pgfpathcurveto{\pgfqpoint{2.433559in}{1.664997in}}{\pgfqpoint{2.427973in}{1.667311in}}{\pgfqpoint{2.422149in}{1.667311in}}%
\pgfpathcurveto{\pgfqpoint{2.416325in}{1.667311in}}{\pgfqpoint{2.410739in}{1.664997in}}{\pgfqpoint{2.406621in}{1.660879in}}%
\pgfpathcurveto{\pgfqpoint{2.402503in}{1.656761in}}{\pgfqpoint{2.400189in}{1.651175in}}{\pgfqpoint{2.400189in}{1.645351in}}%
\pgfpathcurveto{\pgfqpoint{2.400189in}{1.639527in}}{\pgfqpoint{2.402503in}{1.633941in}}{\pgfqpoint{2.406621in}{1.629823in}}%
\pgfpathcurveto{\pgfqpoint{2.410739in}{1.625705in}}{\pgfqpoint{2.416325in}{1.623391in}}{\pgfqpoint{2.422149in}{1.623391in}}%
\pgfpathlineto{\pgfqpoint{2.422149in}{1.623391in}}%
\pgfpathclose%
\pgfusepath{stroke,fill}%
\end{pgfscope}%
\begin{pgfscope}%
\pgfpathrectangle{\pgfqpoint{0.100000in}{0.209042in}}{\pgfqpoint{2.906250in}{2.906250in}}%
\pgfusepath{clip}%
\pgfsetbuttcap%
\pgfsetroundjoin%
\definecolor{currentfill}{rgb}{0.171713,0.067305,0.370771}%
\pgfsetfillcolor{currentfill}%
\pgfsetfillopacity{0.800000}%
\pgfsetlinewidth{1.003750pt}%
\definecolor{currentstroke}{rgb}{0.171713,0.067305,0.370771}%
\pgfsetstrokecolor{currentstroke}%
\pgfsetstrokeopacity{0.800000}%
\pgfsetdash{}{0pt}%
\pgfpathmoveto{\pgfqpoint{1.833312in}{1.763750in}}%
\pgfpathcurveto{\pgfqpoint{1.839136in}{1.763750in}}{\pgfqpoint{1.844722in}{1.766064in}}{\pgfqpoint{1.848840in}{1.770182in}}%
\pgfpathcurveto{\pgfqpoint{1.852959in}{1.774300in}}{\pgfqpoint{1.855273in}{1.779887in}}{\pgfqpoint{1.855273in}{1.785711in}}%
\pgfpathcurveto{\pgfqpoint{1.855273in}{1.791535in}}{\pgfqpoint{1.852959in}{1.797121in}}{\pgfqpoint{1.848840in}{1.801239in}}%
\pgfpathcurveto{\pgfqpoint{1.844722in}{1.805357in}}{\pgfqpoint{1.839136in}{1.807671in}}{\pgfqpoint{1.833312in}{1.807671in}}%
\pgfpathcurveto{\pgfqpoint{1.827488in}{1.807671in}}{\pgfqpoint{1.821902in}{1.805357in}}{\pgfqpoint{1.817784in}{1.801239in}}%
\pgfpathcurveto{\pgfqpoint{1.813666in}{1.797121in}}{\pgfqpoint{1.811352in}{1.791535in}}{\pgfqpoint{1.811352in}{1.785711in}}%
\pgfpathcurveto{\pgfqpoint{1.811352in}{1.779887in}}{\pgfqpoint{1.813666in}{1.774300in}}{\pgfqpoint{1.817784in}{1.770182in}}%
\pgfpathcurveto{\pgfqpoint{1.821902in}{1.766064in}}{\pgfqpoint{1.827488in}{1.763750in}}{\pgfqpoint{1.833312in}{1.763750in}}%
\pgfpathlineto{\pgfqpoint{1.833312in}{1.763750in}}%
\pgfpathclose%
\pgfusepath{stroke,fill}%
\end{pgfscope}%
\begin{pgfscope}%
\pgfpathrectangle{\pgfqpoint{0.100000in}{0.209042in}}{\pgfqpoint{2.906250in}{2.906250in}}%
\pgfusepath{clip}%
\pgfsetbuttcap%
\pgfsetroundjoin%
\definecolor{currentfill}{rgb}{0.088155,0.058133,0.229922}%
\pgfsetfillcolor{currentfill}%
\pgfsetfillopacity{0.800000}%
\pgfsetlinewidth{1.003750pt}%
\definecolor{currentstroke}{rgb}{0.088155,0.058133,0.229922}%
\pgfsetstrokecolor{currentstroke}%
\pgfsetstrokeopacity{0.800000}%
\pgfsetdash{}{0pt}%
\pgfpathmoveto{\pgfqpoint{1.752154in}{1.613515in}}%
\pgfpathcurveto{\pgfqpoint{1.757978in}{1.613515in}}{\pgfqpoint{1.763564in}{1.615828in}}{\pgfqpoint{1.767682in}{1.619947in}}%
\pgfpathcurveto{\pgfqpoint{1.771800in}{1.624065in}}{\pgfqpoint{1.774114in}{1.629651in}}{\pgfqpoint{1.774114in}{1.635475in}}%
\pgfpathcurveto{\pgfqpoint{1.774114in}{1.641299in}}{\pgfqpoint{1.771800in}{1.646885in}}{\pgfqpoint{1.767682in}{1.651003in}}%
\pgfpathcurveto{\pgfqpoint{1.763564in}{1.655121in}}{\pgfqpoint{1.757978in}{1.657435in}}{\pgfqpoint{1.752154in}{1.657435in}}%
\pgfpathcurveto{\pgfqpoint{1.746330in}{1.657435in}}{\pgfqpoint{1.740744in}{1.655121in}}{\pgfqpoint{1.736626in}{1.651003in}}%
\pgfpathcurveto{\pgfqpoint{1.732507in}{1.646885in}}{\pgfqpoint{1.730194in}{1.641299in}}{\pgfqpoint{1.730194in}{1.635475in}}%
\pgfpathcurveto{\pgfqpoint{1.730194in}{1.629651in}}{\pgfqpoint{1.732507in}{1.624065in}}{\pgfqpoint{1.736626in}{1.619947in}}%
\pgfpathcurveto{\pgfqpoint{1.740744in}{1.615828in}}{\pgfqpoint{1.746330in}{1.613515in}}{\pgfqpoint{1.752154in}{1.613515in}}%
\pgfpathlineto{\pgfqpoint{1.752154in}{1.613515in}}%
\pgfpathclose%
\pgfusepath{stroke,fill}%
\end{pgfscope}%
\begin{pgfscope}%
\pgfpathrectangle{\pgfqpoint{0.100000in}{0.209042in}}{\pgfqpoint{2.906250in}{2.906250in}}%
\pgfusepath{clip}%
\pgfsetbuttcap%
\pgfsetroundjoin%
\definecolor{currentfill}{rgb}{0.211718,0.061992,0.418647}%
\pgfsetfillcolor{currentfill}%
\pgfsetfillopacity{0.800000}%
\pgfsetlinewidth{1.003750pt}%
\definecolor{currentstroke}{rgb}{0.211718,0.061992,0.418647}%
\pgfsetstrokecolor{currentstroke}%
\pgfsetstrokeopacity{0.800000}%
\pgfsetdash{}{0pt}%
\pgfpathmoveto{\pgfqpoint{1.557977in}{1.513550in}}%
\pgfpathcurveto{\pgfqpoint{1.563801in}{1.513550in}}{\pgfqpoint{1.569387in}{1.515864in}}{\pgfqpoint{1.573505in}{1.519982in}}%
\pgfpathcurveto{\pgfqpoint{1.577623in}{1.524100in}}{\pgfqpoint{1.579937in}{1.529686in}}{\pgfqpoint{1.579937in}{1.535510in}}%
\pgfpathcurveto{\pgfqpoint{1.579937in}{1.541334in}}{\pgfqpoint{1.577623in}{1.546921in}}{\pgfqpoint{1.573505in}{1.551039in}}%
\pgfpathcurveto{\pgfqpoint{1.569387in}{1.555157in}}{\pgfqpoint{1.563801in}{1.557471in}}{\pgfqpoint{1.557977in}{1.557471in}}%
\pgfpathcurveto{\pgfqpoint{1.552153in}{1.557471in}}{\pgfqpoint{1.546567in}{1.555157in}}{\pgfqpoint{1.542449in}{1.551039in}}%
\pgfpathcurveto{\pgfqpoint{1.538331in}{1.546921in}}{\pgfqpoint{1.536017in}{1.541334in}}{\pgfqpoint{1.536017in}{1.535510in}}%
\pgfpathcurveto{\pgfqpoint{1.536017in}{1.529686in}}{\pgfqpoint{1.538331in}{1.524100in}}{\pgfqpoint{1.542449in}{1.519982in}}%
\pgfpathcurveto{\pgfqpoint{1.546567in}{1.515864in}}{\pgfqpoint{1.552153in}{1.513550in}}{\pgfqpoint{1.557977in}{1.513550in}}%
\pgfpathlineto{\pgfqpoint{1.557977in}{1.513550in}}%
\pgfpathclose%
\pgfusepath{stroke,fill}%
\end{pgfscope}%
\begin{pgfscope}%
\pgfpathrectangle{\pgfqpoint{0.100000in}{0.209042in}}{\pgfqpoint{2.906250in}{2.906250in}}%
\pgfusepath{clip}%
\pgfsetbuttcap%
\pgfsetroundjoin%
\definecolor{currentfill}{rgb}{0.018815,0.016026,0.084584}%
\pgfsetfillcolor{currentfill}%
\pgfsetfillopacity{0.800000}%
\pgfsetlinewidth{1.003750pt}%
\definecolor{currentstroke}{rgb}{0.018815,0.016026,0.084584}%
\pgfsetstrokecolor{currentstroke}%
\pgfsetstrokeopacity{0.800000}%
\pgfsetdash{}{0pt}%
\pgfpathmoveto{\pgfqpoint{1.696112in}{1.745462in}}%
\pgfpathcurveto{\pgfqpoint{1.701936in}{1.745462in}}{\pgfqpoint{1.707523in}{1.747776in}}{\pgfqpoint{1.711641in}{1.751894in}}%
\pgfpathcurveto{\pgfqpoint{1.715759in}{1.756012in}}{\pgfqpoint{1.718073in}{1.761598in}}{\pgfqpoint{1.718073in}{1.767422in}}%
\pgfpathcurveto{\pgfqpoint{1.718073in}{1.773246in}}{\pgfqpoint{1.715759in}{1.778832in}}{\pgfqpoint{1.711641in}{1.782950in}}%
\pgfpathcurveto{\pgfqpoint{1.707523in}{1.787068in}}{\pgfqpoint{1.701936in}{1.789382in}}{\pgfqpoint{1.696112in}{1.789382in}}%
\pgfpathcurveto{\pgfqpoint{1.690289in}{1.789382in}}{\pgfqpoint{1.684702in}{1.787068in}}{\pgfqpoint{1.680584in}{1.782950in}}%
\pgfpathcurveto{\pgfqpoint{1.676466in}{1.778832in}}{\pgfqpoint{1.674152in}{1.773246in}}{\pgfqpoint{1.674152in}{1.767422in}}%
\pgfpathcurveto{\pgfqpoint{1.674152in}{1.761598in}}{\pgfqpoint{1.676466in}{1.756012in}}{\pgfqpoint{1.680584in}{1.751894in}}%
\pgfpathcurveto{\pgfqpoint{1.684702in}{1.747776in}}{\pgfqpoint{1.690289in}{1.745462in}}{\pgfqpoint{1.696112in}{1.745462in}}%
\pgfpathlineto{\pgfqpoint{1.696112in}{1.745462in}}%
\pgfpathclose%
\pgfusepath{stroke,fill}%
\end{pgfscope}%
\begin{pgfscope}%
\pgfpathrectangle{\pgfqpoint{0.100000in}{0.209042in}}{\pgfqpoint{2.906250in}{2.906250in}}%
\pgfusepath{clip}%
\pgfsetbuttcap%
\pgfsetroundjoin%
\definecolor{currentfill}{rgb}{0.506629,0.145958,0.507806}%
\pgfsetfillcolor{currentfill}%
\pgfsetfillopacity{0.800000}%
\pgfsetlinewidth{1.003750pt}%
\definecolor{currentstroke}{rgb}{0.506629,0.145958,0.507806}%
\pgfsetstrokecolor{currentstroke}%
\pgfsetstrokeopacity{0.800000}%
\pgfsetdash{}{0pt}%
\pgfpathmoveto{\pgfqpoint{1.305945in}{1.954916in}}%
\pgfpathcurveto{\pgfqpoint{1.311769in}{1.954916in}}{\pgfqpoint{1.317356in}{1.957230in}}{\pgfqpoint{1.321474in}{1.961348in}}%
\pgfpathcurveto{\pgfqpoint{1.325592in}{1.965467in}}{\pgfqpoint{1.327906in}{1.971053in}}{\pgfqpoint{1.327906in}{1.976877in}}%
\pgfpathcurveto{\pgfqpoint{1.327906in}{1.982701in}}{\pgfqpoint{1.325592in}{1.988287in}}{\pgfqpoint{1.321474in}{1.992405in}}%
\pgfpathcurveto{\pgfqpoint{1.317356in}{1.996523in}}{\pgfqpoint{1.311769in}{1.998837in}}{\pgfqpoint{1.305945in}{1.998837in}}%
\pgfpathcurveto{\pgfqpoint{1.300122in}{1.998837in}}{\pgfqpoint{1.294535in}{1.996523in}}{\pgfqpoint{1.290417in}{1.992405in}}%
\pgfpathcurveto{\pgfqpoint{1.286299in}{1.988287in}}{\pgfqpoint{1.283985in}{1.982701in}}{\pgfqpoint{1.283985in}{1.976877in}}%
\pgfpathcurveto{\pgfqpoint{1.283985in}{1.971053in}}{\pgfqpoint{1.286299in}{1.965467in}}{\pgfqpoint{1.290417in}{1.961348in}}%
\pgfpathcurveto{\pgfqpoint{1.294535in}{1.957230in}}{\pgfqpoint{1.300122in}{1.954916in}}{\pgfqpoint{1.305945in}{1.954916in}}%
\pgfpathlineto{\pgfqpoint{1.305945in}{1.954916in}}%
\pgfpathclose%
\pgfusepath{stroke,fill}%
\end{pgfscope}%
\begin{pgfscope}%
\pgfpathrectangle{\pgfqpoint{0.100000in}{0.209042in}}{\pgfqpoint{2.906250in}{2.906250in}}%
\pgfusepath{clip}%
\pgfsetbuttcap%
\pgfsetroundjoin%
\definecolor{currentfill}{rgb}{0.958464,0.411324,0.360014}%
\pgfsetfillcolor{currentfill}%
\pgfsetfillopacity{0.800000}%
\pgfsetlinewidth{1.003750pt}%
\definecolor{currentstroke}{rgb}{0.958464,0.411324,0.360014}%
\pgfsetstrokecolor{currentstroke}%
\pgfsetstrokeopacity{0.800000}%
\pgfsetdash{}{0pt}%
\pgfpathmoveto{\pgfqpoint{0.889483in}{1.882420in}}%
\pgfpathcurveto{\pgfqpoint{0.895307in}{1.882420in}}{\pgfqpoint{0.900893in}{1.884734in}}{\pgfqpoint{0.905011in}{1.888852in}}%
\pgfpathcurveto{\pgfqpoint{0.909129in}{1.892971in}}{\pgfqpoint{0.911443in}{1.898557in}}{\pgfqpoint{0.911443in}{1.904381in}}%
\pgfpathcurveto{\pgfqpoint{0.911443in}{1.910205in}}{\pgfqpoint{0.909129in}{1.915791in}}{\pgfqpoint{0.905011in}{1.919909in}}%
\pgfpathcurveto{\pgfqpoint{0.900893in}{1.924027in}}{\pgfqpoint{0.895307in}{1.926341in}}{\pgfqpoint{0.889483in}{1.926341in}}%
\pgfpathcurveto{\pgfqpoint{0.883659in}{1.926341in}}{\pgfqpoint{0.878073in}{1.924027in}}{\pgfqpoint{0.873955in}{1.919909in}}%
\pgfpathcurveto{\pgfqpoint{0.869836in}{1.915791in}}{\pgfqpoint{0.867523in}{1.910205in}}{\pgfqpoint{0.867523in}{1.904381in}}%
\pgfpathcurveto{\pgfqpoint{0.867523in}{1.898557in}}{\pgfqpoint{0.869836in}{1.892971in}}{\pgfqpoint{0.873955in}{1.888852in}}%
\pgfpathcurveto{\pgfqpoint{0.878073in}{1.884734in}}{\pgfqpoint{0.883659in}{1.882420in}}{\pgfqpoint{0.889483in}{1.882420in}}%
\pgfpathlineto{\pgfqpoint{0.889483in}{1.882420in}}%
\pgfpathclose%
\pgfusepath{stroke,fill}%
\end{pgfscope}%
\begin{pgfscope}%
\pgfpathrectangle{\pgfqpoint{0.100000in}{0.209042in}}{\pgfqpoint{2.906250in}{2.906250in}}%
\pgfusepath{clip}%
\pgfsetbuttcap%
\pgfsetroundjoin%
\definecolor{currentfill}{rgb}{0.451271,0.125132,0.507198}%
\pgfsetfillcolor{currentfill}%
\pgfsetfillopacity{0.800000}%
\pgfsetlinewidth{1.003750pt}%
\definecolor{currentstroke}{rgb}{0.451271,0.125132,0.507198}%
\pgfsetstrokecolor{currentstroke}%
\pgfsetstrokeopacity{0.800000}%
\pgfsetdash{}{0pt}%
\pgfpathmoveto{\pgfqpoint{1.833957in}{1.974103in}}%
\pgfpathcurveto{\pgfqpoint{1.839781in}{1.974103in}}{\pgfqpoint{1.845367in}{1.976417in}}{\pgfqpoint{1.849485in}{1.980535in}}%
\pgfpathcurveto{\pgfqpoint{1.853604in}{1.984653in}}{\pgfqpoint{1.855917in}{1.990239in}}{\pgfqpoint{1.855917in}{1.996063in}}%
\pgfpathcurveto{\pgfqpoint{1.855917in}{2.001887in}}{\pgfqpoint{1.853604in}{2.007473in}}{\pgfqpoint{1.849485in}{2.011592in}}%
\pgfpathcurveto{\pgfqpoint{1.845367in}{2.015710in}}{\pgfqpoint{1.839781in}{2.018024in}}{\pgfqpoint{1.833957in}{2.018024in}}%
\pgfpathcurveto{\pgfqpoint{1.828133in}{2.018024in}}{\pgfqpoint{1.822547in}{2.015710in}}{\pgfqpoint{1.818429in}{2.011592in}}%
\pgfpathcurveto{\pgfqpoint{1.814311in}{2.007473in}}{\pgfqpoint{1.811997in}{2.001887in}}{\pgfqpoint{1.811997in}{1.996063in}}%
\pgfpathcurveto{\pgfqpoint{1.811997in}{1.990239in}}{\pgfqpoint{1.814311in}{1.984653in}}{\pgfqpoint{1.818429in}{1.980535in}}%
\pgfpathcurveto{\pgfqpoint{1.822547in}{1.976417in}}{\pgfqpoint{1.828133in}{1.974103in}}{\pgfqpoint{1.833957in}{1.974103in}}%
\pgfpathlineto{\pgfqpoint{1.833957in}{1.974103in}}%
\pgfpathclose%
\pgfusepath{stroke,fill}%
\end{pgfscope}%
\begin{pgfscope}%
\pgfpathrectangle{\pgfqpoint{0.100000in}{0.209042in}}{\pgfqpoint{2.906250in}{2.906250in}}%
\pgfusepath{clip}%
\pgfsetbuttcap%
\pgfsetroundjoin%
\definecolor{currentfill}{rgb}{0.252220,0.059415,0.453248}%
\pgfsetfillcolor{currentfill}%
\pgfsetfillopacity{0.800000}%
\pgfsetlinewidth{1.003750pt}%
\definecolor{currentstroke}{rgb}{0.252220,0.059415,0.453248}%
\pgfsetstrokecolor{currentstroke}%
\pgfsetstrokeopacity{0.800000}%
\pgfsetdash{}{0pt}%
\pgfpathmoveto{\pgfqpoint{1.622053in}{1.920110in}}%
\pgfpathcurveto{\pgfqpoint{1.627877in}{1.920110in}}{\pgfqpoint{1.633463in}{1.922424in}}{\pgfqpoint{1.637581in}{1.926542in}}%
\pgfpathcurveto{\pgfqpoint{1.641699in}{1.930660in}}{\pgfqpoint{1.644013in}{1.936246in}}{\pgfqpoint{1.644013in}{1.942070in}}%
\pgfpathcurveto{\pgfqpoint{1.644013in}{1.947894in}}{\pgfqpoint{1.641699in}{1.953480in}}{\pgfqpoint{1.637581in}{1.957598in}}%
\pgfpathcurveto{\pgfqpoint{1.633463in}{1.961716in}}{\pgfqpoint{1.627877in}{1.964030in}}{\pgfqpoint{1.622053in}{1.964030in}}%
\pgfpathcurveto{\pgfqpoint{1.616229in}{1.964030in}}{\pgfqpoint{1.610643in}{1.961716in}}{\pgfqpoint{1.606525in}{1.957598in}}%
\pgfpathcurveto{\pgfqpoint{1.602406in}{1.953480in}}{\pgfqpoint{1.600092in}{1.947894in}}{\pgfqpoint{1.600092in}{1.942070in}}%
\pgfpathcurveto{\pgfqpoint{1.600092in}{1.936246in}}{\pgfqpoint{1.602406in}{1.930660in}}{\pgfqpoint{1.606525in}{1.926542in}}%
\pgfpathcurveto{\pgfqpoint{1.610643in}{1.922424in}}{\pgfqpoint{1.616229in}{1.920110in}}{\pgfqpoint{1.622053in}{1.920110in}}%
\pgfpathlineto{\pgfqpoint{1.622053in}{1.920110in}}%
\pgfpathclose%
\pgfusepath{stroke,fill}%
\end{pgfscope}%
\begin{pgfscope}%
\pgfpathrectangle{\pgfqpoint{0.100000in}{0.209042in}}{\pgfqpoint{2.906250in}{2.906250in}}%
\pgfusepath{clip}%
\pgfsetbuttcap%
\pgfsetroundjoin%
\definecolor{currentfill}{rgb}{0.632805,0.187893,0.495332}%
\pgfsetfillcolor{currentfill}%
\pgfsetfillopacity{0.800000}%
\pgfsetlinewidth{1.003750pt}%
\definecolor{currentstroke}{rgb}{0.632805,0.187893,0.495332}%
\pgfsetstrokecolor{currentstroke}%
\pgfsetstrokeopacity{0.800000}%
\pgfsetdash{}{0pt}%
\pgfpathmoveto{\pgfqpoint{2.150729in}{1.660824in}}%
\pgfpathcurveto{\pgfqpoint{2.156553in}{1.660824in}}{\pgfqpoint{2.162139in}{1.663137in}}{\pgfqpoint{2.166257in}{1.667256in}}%
\pgfpathcurveto{\pgfqpoint{2.170375in}{1.671374in}}{\pgfqpoint{2.172689in}{1.676960in}}{\pgfqpoint{2.172689in}{1.682784in}}%
\pgfpathcurveto{\pgfqpoint{2.172689in}{1.688608in}}{\pgfqpoint{2.170375in}{1.694194in}}{\pgfqpoint{2.166257in}{1.698312in}}%
\pgfpathcurveto{\pgfqpoint{2.162139in}{1.702430in}}{\pgfqpoint{2.156553in}{1.704744in}}{\pgfqpoint{2.150729in}{1.704744in}}%
\pgfpathcurveto{\pgfqpoint{2.144905in}{1.704744in}}{\pgfqpoint{2.139319in}{1.702430in}}{\pgfqpoint{2.135201in}{1.698312in}}%
\pgfpathcurveto{\pgfqpoint{2.131083in}{1.694194in}}{\pgfqpoint{2.128769in}{1.688608in}}{\pgfqpoint{2.128769in}{1.682784in}}%
\pgfpathcurveto{\pgfqpoint{2.128769in}{1.676960in}}{\pgfqpoint{2.131083in}{1.671374in}}{\pgfqpoint{2.135201in}{1.667256in}}%
\pgfpathcurveto{\pgfqpoint{2.139319in}{1.663137in}}{\pgfqpoint{2.144905in}{1.660824in}}{\pgfqpoint{2.150729in}{1.660824in}}%
\pgfpathlineto{\pgfqpoint{2.150729in}{1.660824in}}%
\pgfpathclose%
\pgfusepath{stroke,fill}%
\end{pgfscope}%
\begin{pgfscope}%
\pgfpathrectangle{\pgfqpoint{0.100000in}{0.209042in}}{\pgfqpoint{2.906250in}{2.906250in}}%
\pgfusepath{clip}%
\pgfsetbuttcap%
\pgfsetroundjoin%
\definecolor{currentfill}{rgb}{0.123833,0.067295,0.295879}%
\pgfsetfillcolor{currentfill}%
\pgfsetfillopacity{0.800000}%
\pgfsetlinewidth{1.003750pt}%
\definecolor{currentstroke}{rgb}{0.123833,0.067295,0.295879}%
\pgfsetstrokecolor{currentstroke}%
\pgfsetstrokeopacity{0.800000}%
\pgfsetdash{}{0pt}%
\pgfpathmoveto{\pgfqpoint{1.809415in}{1.644675in}}%
\pgfpathcurveto{\pgfqpoint{1.815239in}{1.644675in}}{\pgfqpoint{1.820825in}{1.646989in}}{\pgfqpoint{1.824943in}{1.651107in}}%
\pgfpathcurveto{\pgfqpoint{1.829061in}{1.655225in}}{\pgfqpoint{1.831375in}{1.660811in}}{\pgfqpoint{1.831375in}{1.666635in}}%
\pgfpathcurveto{\pgfqpoint{1.831375in}{1.672459in}}{\pgfqpoint{1.829061in}{1.678045in}}{\pgfqpoint{1.824943in}{1.682163in}}%
\pgfpathcurveto{\pgfqpoint{1.820825in}{1.686281in}}{\pgfqpoint{1.815239in}{1.688595in}}{\pgfqpoint{1.809415in}{1.688595in}}%
\pgfpathcurveto{\pgfqpoint{1.803591in}{1.688595in}}{\pgfqpoint{1.798005in}{1.686281in}}{\pgfqpoint{1.793887in}{1.682163in}}%
\pgfpathcurveto{\pgfqpoint{1.789769in}{1.678045in}}{\pgfqpoint{1.787455in}{1.672459in}}{\pgfqpoint{1.787455in}{1.666635in}}%
\pgfpathcurveto{\pgfqpoint{1.787455in}{1.660811in}}{\pgfqpoint{1.789769in}{1.655225in}}{\pgfqpoint{1.793887in}{1.651107in}}%
\pgfpathcurveto{\pgfqpoint{1.798005in}{1.646989in}}{\pgfqpoint{1.803591in}{1.644675in}}{\pgfqpoint{1.809415in}{1.644675in}}%
\pgfpathlineto{\pgfqpoint{1.809415in}{1.644675in}}%
\pgfpathclose%
\pgfusepath{stroke,fill}%
\end{pgfscope}%
\begin{pgfscope}%
\pgfpathrectangle{\pgfqpoint{0.100000in}{0.209042in}}{\pgfqpoint{2.906250in}{2.906250in}}%
\pgfusepath{clip}%
\pgfsetbuttcap%
\pgfsetroundjoin%
\definecolor{currentfill}{rgb}{0.028123,0.023201,0.108787}%
\pgfsetfillcolor{currentfill}%
\pgfsetfillopacity{0.800000}%
\pgfsetlinewidth{1.003750pt}%
\definecolor{currentstroke}{rgb}{0.028123,0.023201,0.108787}%
\pgfsetstrokecolor{currentstroke}%
\pgfsetstrokeopacity{0.800000}%
\pgfsetdash{}{0pt}%
\pgfpathmoveto{\pgfqpoint{1.550691in}{1.750532in}}%
\pgfpathcurveto{\pgfqpoint{1.556515in}{1.750532in}}{\pgfqpoint{1.562102in}{1.752846in}}{\pgfqpoint{1.566220in}{1.756964in}}%
\pgfpathcurveto{\pgfqpoint{1.570338in}{1.761082in}}{\pgfqpoint{1.572652in}{1.766668in}}{\pgfqpoint{1.572652in}{1.772492in}}%
\pgfpathcurveto{\pgfqpoint{1.572652in}{1.778316in}}{\pgfqpoint{1.570338in}{1.783902in}}{\pgfqpoint{1.566220in}{1.788020in}}%
\pgfpathcurveto{\pgfqpoint{1.562102in}{1.792138in}}{\pgfqpoint{1.556515in}{1.794452in}}{\pgfqpoint{1.550691in}{1.794452in}}%
\pgfpathcurveto{\pgfqpoint{1.544868in}{1.794452in}}{\pgfqpoint{1.539281in}{1.792138in}}{\pgfqpoint{1.535163in}{1.788020in}}%
\pgfpathcurveto{\pgfqpoint{1.531045in}{1.783902in}}{\pgfqpoint{1.528731in}{1.778316in}}{\pgfqpoint{1.528731in}{1.772492in}}%
\pgfpathcurveto{\pgfqpoint{1.528731in}{1.766668in}}{\pgfqpoint{1.531045in}{1.761082in}}{\pgfqpoint{1.535163in}{1.756964in}}%
\pgfpathcurveto{\pgfqpoint{1.539281in}{1.752846in}}{\pgfqpoint{1.544868in}{1.750532in}}{\pgfqpoint{1.550691in}{1.750532in}}%
\pgfpathlineto{\pgfqpoint{1.550691in}{1.750532in}}%
\pgfpathclose%
\pgfusepath{stroke,fill}%
\end{pgfscope}%
\begin{pgfscope}%
\pgfpathrectangle{\pgfqpoint{0.100000in}{0.209042in}}{\pgfqpoint{2.906250in}{2.906250in}}%
\pgfusepath{clip}%
\pgfsetbuttcap%
\pgfsetroundjoin%
\definecolor{currentfill}{rgb}{0.074257,0.052017,0.202660}%
\pgfsetfillcolor{currentfill}%
\pgfsetfillopacity{0.800000}%
\pgfsetlinewidth{1.003750pt}%
\definecolor{currentstroke}{rgb}{0.074257,0.052017,0.202660}%
\pgfsetstrokecolor{currentstroke}%
\pgfsetstrokeopacity{0.800000}%
\pgfsetdash{}{0pt}%
\pgfpathmoveto{\pgfqpoint{1.515203in}{1.632253in}}%
\pgfpathcurveto{\pgfqpoint{1.521027in}{1.632253in}}{\pgfqpoint{1.526613in}{1.634567in}}{\pgfqpoint{1.530731in}{1.638685in}}%
\pgfpathcurveto{\pgfqpoint{1.534850in}{1.642803in}}{\pgfqpoint{1.537163in}{1.648389in}}{\pgfqpoint{1.537163in}{1.654213in}}%
\pgfpathcurveto{\pgfqpoint{1.537163in}{1.660037in}}{\pgfqpoint{1.534850in}{1.665623in}}{\pgfqpoint{1.530731in}{1.669741in}}%
\pgfpathcurveto{\pgfqpoint{1.526613in}{1.673859in}}{\pgfqpoint{1.521027in}{1.676173in}}{\pgfqpoint{1.515203in}{1.676173in}}%
\pgfpathcurveto{\pgfqpoint{1.509379in}{1.676173in}}{\pgfqpoint{1.503793in}{1.673859in}}{\pgfqpoint{1.499675in}{1.669741in}}%
\pgfpathcurveto{\pgfqpoint{1.495557in}{1.665623in}}{\pgfqpoint{1.493243in}{1.660037in}}{\pgfqpoint{1.493243in}{1.654213in}}%
\pgfpathcurveto{\pgfqpoint{1.493243in}{1.648389in}}{\pgfqpoint{1.495557in}{1.642803in}}{\pgfqpoint{1.499675in}{1.638685in}}%
\pgfpathcurveto{\pgfqpoint{1.503793in}{1.634567in}}{\pgfqpoint{1.509379in}{1.632253in}}{\pgfqpoint{1.515203in}{1.632253in}}%
\pgfpathlineto{\pgfqpoint{1.515203in}{1.632253in}}%
\pgfpathclose%
\pgfusepath{stroke,fill}%
\end{pgfscope}%
\begin{pgfscope}%
\pgfpathrectangle{\pgfqpoint{0.100000in}{0.209042in}}{\pgfqpoint{2.906250in}{2.906250in}}%
\pgfusepath{clip}%
\pgfsetbuttcap%
\pgfsetroundjoin%
\definecolor{currentfill}{rgb}{0.123833,0.067295,0.295879}%
\pgfsetfillcolor{currentfill}%
\pgfsetfillopacity{0.800000}%
\pgfsetlinewidth{1.003750pt}%
\definecolor{currentstroke}{rgb}{0.123833,0.067295,0.295879}%
\pgfsetstrokecolor{currentstroke}%
\pgfsetstrokeopacity{0.800000}%
\pgfsetdash{}{0pt}%
\pgfpathmoveto{\pgfqpoint{1.567365in}{1.846080in}}%
\pgfpathcurveto{\pgfqpoint{1.573188in}{1.846080in}}{\pgfqpoint{1.578775in}{1.848394in}}{\pgfqpoint{1.582893in}{1.852512in}}%
\pgfpathcurveto{\pgfqpoint{1.587011in}{1.856630in}}{\pgfqpoint{1.589325in}{1.862216in}}{\pgfqpoint{1.589325in}{1.868040in}}%
\pgfpathcurveto{\pgfqpoint{1.589325in}{1.873864in}}{\pgfqpoint{1.587011in}{1.879450in}}{\pgfqpoint{1.582893in}{1.883568in}}%
\pgfpathcurveto{\pgfqpoint{1.578775in}{1.887686in}}{\pgfqpoint{1.573188in}{1.890000in}}{\pgfqpoint{1.567365in}{1.890000in}}%
\pgfpathcurveto{\pgfqpoint{1.561541in}{1.890000in}}{\pgfqpoint{1.555954in}{1.887686in}}{\pgfqpoint{1.551836in}{1.883568in}}%
\pgfpathcurveto{\pgfqpoint{1.547718in}{1.879450in}}{\pgfqpoint{1.545404in}{1.873864in}}{\pgfqpoint{1.545404in}{1.868040in}}%
\pgfpathcurveto{\pgfqpoint{1.545404in}{1.862216in}}{\pgfqpoint{1.547718in}{1.856630in}}{\pgfqpoint{1.551836in}{1.852512in}}%
\pgfpathcurveto{\pgfqpoint{1.555954in}{1.848394in}}{\pgfqpoint{1.561541in}{1.846080in}}{\pgfqpoint{1.567365in}{1.846080in}}%
\pgfpathlineto{\pgfqpoint{1.567365in}{1.846080in}}%
\pgfpathclose%
\pgfusepath{stroke,fill}%
\end{pgfscope}%
\begin{pgfscope}%
\pgfpathrectangle{\pgfqpoint{0.100000in}{0.209042in}}{\pgfqpoint{2.906250in}{2.906250in}}%
\pgfusepath{clip}%
\pgfsetbuttcap%
\pgfsetroundjoin%
\definecolor{currentfill}{rgb}{0.035520,0.028397,0.125209}%
\pgfsetfillcolor{currentfill}%
\pgfsetfillopacity{0.800000}%
\pgfsetlinewidth{1.003750pt}%
\definecolor{currentstroke}{rgb}{0.035520,0.028397,0.125209}%
\pgfsetstrokecolor{currentstroke}%
\pgfsetstrokeopacity{0.800000}%
\pgfsetdash{}{0pt}%
\pgfpathmoveto{\pgfqpoint{1.734407in}{1.717985in}}%
\pgfpathcurveto{\pgfqpoint{1.740231in}{1.717985in}}{\pgfqpoint{1.745817in}{1.720298in}}{\pgfqpoint{1.749935in}{1.724417in}}%
\pgfpathcurveto{\pgfqpoint{1.754053in}{1.728535in}}{\pgfqpoint{1.756367in}{1.734121in}}{\pgfqpoint{1.756367in}{1.739945in}}%
\pgfpathcurveto{\pgfqpoint{1.756367in}{1.745769in}}{\pgfqpoint{1.754053in}{1.751355in}}{\pgfqpoint{1.749935in}{1.755473in}}%
\pgfpathcurveto{\pgfqpoint{1.745817in}{1.759591in}}{\pgfqpoint{1.740231in}{1.761905in}}{\pgfqpoint{1.734407in}{1.761905in}}%
\pgfpathcurveto{\pgfqpoint{1.728583in}{1.761905in}}{\pgfqpoint{1.722997in}{1.759591in}}{\pgfqpoint{1.718879in}{1.755473in}}%
\pgfpathcurveto{\pgfqpoint{1.714761in}{1.751355in}}{\pgfqpoint{1.712447in}{1.745769in}}{\pgfqpoint{1.712447in}{1.739945in}}%
\pgfpathcurveto{\pgfqpoint{1.712447in}{1.734121in}}{\pgfqpoint{1.714761in}{1.728535in}}{\pgfqpoint{1.718879in}{1.724417in}}%
\pgfpathcurveto{\pgfqpoint{1.722997in}{1.720298in}}{\pgfqpoint{1.728583in}{1.717985in}}{\pgfqpoint{1.734407in}{1.717985in}}%
\pgfpathlineto{\pgfqpoint{1.734407in}{1.717985in}}%
\pgfpathclose%
\pgfusepath{stroke,fill}%
\end{pgfscope}%
\begin{pgfscope}%
\pgfpathrectangle{\pgfqpoint{0.100000in}{0.209042in}}{\pgfqpoint{2.906250in}{2.906250in}}%
\pgfusepath{clip}%
\pgfsetbuttcap%
\pgfsetroundjoin%
\definecolor{currentfill}{rgb}{0.159018,0.068354,0.352688}%
\pgfsetfillcolor{currentfill}%
\pgfsetfillopacity{0.800000}%
\pgfsetlinewidth{1.003750pt}%
\definecolor{currentstroke}{rgb}{0.159018,0.068354,0.352688}%
\pgfsetstrokecolor{currentstroke}%
\pgfsetstrokeopacity{0.800000}%
\pgfsetdash{}{0pt}%
\pgfpathmoveto{\pgfqpoint{1.779614in}{1.815393in}}%
\pgfpathcurveto{\pgfqpoint{1.785438in}{1.815393in}}{\pgfqpoint{1.791025in}{1.817707in}}{\pgfqpoint{1.795143in}{1.821825in}}%
\pgfpathcurveto{\pgfqpoint{1.799261in}{1.825943in}}{\pgfqpoint{1.801575in}{1.831530in}}{\pgfqpoint{1.801575in}{1.837353in}}%
\pgfpathcurveto{\pgfqpoint{1.801575in}{1.843177in}}{\pgfqpoint{1.799261in}{1.848764in}}{\pgfqpoint{1.795143in}{1.852882in}}%
\pgfpathcurveto{\pgfqpoint{1.791025in}{1.857000in}}{\pgfqpoint{1.785438in}{1.859314in}}{\pgfqpoint{1.779614in}{1.859314in}}%
\pgfpathcurveto{\pgfqpoint{1.773790in}{1.859314in}}{\pgfqpoint{1.768204in}{1.857000in}}{\pgfqpoint{1.764086in}{1.852882in}}%
\pgfpathcurveto{\pgfqpoint{1.759968in}{1.848764in}}{\pgfqpoint{1.757654in}{1.843177in}}{\pgfqpoint{1.757654in}{1.837353in}}%
\pgfpathcurveto{\pgfqpoint{1.757654in}{1.831530in}}{\pgfqpoint{1.759968in}{1.825943in}}{\pgfqpoint{1.764086in}{1.821825in}}%
\pgfpathcurveto{\pgfqpoint{1.768204in}{1.817707in}}{\pgfqpoint{1.773790in}{1.815393in}}{\pgfqpoint{1.779614in}{1.815393in}}%
\pgfpathlineto{\pgfqpoint{1.779614in}{1.815393in}}%
\pgfpathclose%
\pgfusepath{stroke,fill}%
\end{pgfscope}%
\begin{pgfscope}%
\pgfpathrectangle{\pgfqpoint{0.100000in}{0.209042in}}{\pgfqpoint{2.906250in}{2.906250in}}%
\pgfusepath{clip}%
\pgfsetbuttcap%
\pgfsetroundjoin%
\definecolor{currentfill}{rgb}{0.198177,0.063862,0.404009}%
\pgfsetfillcolor{currentfill}%
\pgfsetfillopacity{0.800000}%
\pgfsetlinewidth{1.003750pt}%
\definecolor{currentstroke}{rgb}{0.198177,0.063862,0.404009}%
\pgfsetstrokecolor{currentstroke}%
\pgfsetstrokeopacity{0.800000}%
\pgfsetdash{}{0pt}%
\pgfpathmoveto{\pgfqpoint{1.625064in}{1.507774in}}%
\pgfpathcurveto{\pgfqpoint{1.630888in}{1.507774in}}{\pgfqpoint{1.636474in}{1.510088in}}{\pgfqpoint{1.640592in}{1.514206in}}%
\pgfpathcurveto{\pgfqpoint{1.644710in}{1.518324in}}{\pgfqpoint{1.647024in}{1.523910in}}{\pgfqpoint{1.647024in}{1.529734in}}%
\pgfpathcurveto{\pgfqpoint{1.647024in}{1.535558in}}{\pgfqpoint{1.644710in}{1.541144in}}{\pgfqpoint{1.640592in}{1.545262in}}%
\pgfpathcurveto{\pgfqpoint{1.636474in}{1.549381in}}{\pgfqpoint{1.630888in}{1.551694in}}{\pgfqpoint{1.625064in}{1.551694in}}%
\pgfpathcurveto{\pgfqpoint{1.619240in}{1.551694in}}{\pgfqpoint{1.613654in}{1.549381in}}{\pgfqpoint{1.609536in}{1.545262in}}%
\pgfpathcurveto{\pgfqpoint{1.605417in}{1.541144in}}{\pgfqpoint{1.603104in}{1.535558in}}{\pgfqpoint{1.603104in}{1.529734in}}%
\pgfpathcurveto{\pgfqpoint{1.603104in}{1.523910in}}{\pgfqpoint{1.605417in}{1.518324in}}{\pgfqpoint{1.609536in}{1.514206in}}%
\pgfpathcurveto{\pgfqpoint{1.613654in}{1.510088in}}{\pgfqpoint{1.619240in}{1.507774in}}{\pgfqpoint{1.625064in}{1.507774in}}%
\pgfpathlineto{\pgfqpoint{1.625064in}{1.507774in}}%
\pgfpathclose%
\pgfusepath{stroke,fill}%
\end{pgfscope}%
\begin{pgfscope}%
\pgfpathrectangle{\pgfqpoint{0.100000in}{0.209042in}}{\pgfqpoint{2.906250in}{2.906250in}}%
\pgfusepath{clip}%
\pgfsetbuttcap%
\pgfsetroundjoin%
\definecolor{currentfill}{rgb}{0.021692,0.018320,0.092610}%
\pgfsetfillcolor{currentfill}%
\pgfsetfillopacity{0.800000}%
\pgfsetlinewidth{1.003750pt}%
\definecolor{currentstroke}{rgb}{0.021692,0.018320,0.092610}%
\pgfsetstrokecolor{currentstroke}%
\pgfsetstrokeopacity{0.800000}%
\pgfsetdash{}{0pt}%
\pgfpathmoveto{\pgfqpoint{1.718896in}{1.712887in}}%
\pgfpathcurveto{\pgfqpoint{1.724720in}{1.712887in}}{\pgfqpoint{1.730306in}{1.715201in}}{\pgfqpoint{1.734424in}{1.719319in}}%
\pgfpathcurveto{\pgfqpoint{1.738542in}{1.723437in}}{\pgfqpoint{1.740856in}{1.729024in}}{\pgfqpoint{1.740856in}{1.734847in}}%
\pgfpathcurveto{\pgfqpoint{1.740856in}{1.740671in}}{\pgfqpoint{1.738542in}{1.746258in}}{\pgfqpoint{1.734424in}{1.750376in}}%
\pgfpathcurveto{\pgfqpoint{1.730306in}{1.754494in}}{\pgfqpoint{1.724720in}{1.756808in}}{\pgfqpoint{1.718896in}{1.756808in}}%
\pgfpathcurveto{\pgfqpoint{1.713072in}{1.756808in}}{\pgfqpoint{1.707486in}{1.754494in}}{\pgfqpoint{1.703367in}{1.750376in}}%
\pgfpathcurveto{\pgfqpoint{1.699249in}{1.746258in}}{\pgfqpoint{1.696935in}{1.740671in}}{\pgfqpoint{1.696935in}{1.734847in}}%
\pgfpathcurveto{\pgfqpoint{1.696935in}{1.729024in}}{\pgfqpoint{1.699249in}{1.723437in}}{\pgfqpoint{1.703367in}{1.719319in}}%
\pgfpathcurveto{\pgfqpoint{1.707486in}{1.715201in}}{\pgfqpoint{1.713072in}{1.712887in}}{\pgfqpoint{1.718896in}{1.712887in}}%
\pgfpathlineto{\pgfqpoint{1.718896in}{1.712887in}}%
\pgfpathclose%
\pgfusepath{stroke,fill}%
\end{pgfscope}%
\begin{pgfscope}%
\pgfpathrectangle{\pgfqpoint{0.100000in}{0.209042in}}{\pgfqpoint{2.906250in}{2.906250in}}%
\pgfusepath{clip}%
\pgfsetbuttcap%
\pgfsetroundjoin%
\definecolor{currentfill}{rgb}{0.341482,0.080564,0.492631}%
\pgfsetfillcolor{currentfill}%
\pgfsetfillopacity{0.800000}%
\pgfsetlinewidth{1.003750pt}%
\definecolor{currentstroke}{rgb}{0.341482,0.080564,0.492631}%
\pgfsetstrokecolor{currentstroke}%
\pgfsetstrokeopacity{0.800000}%
\pgfsetdash{}{0pt}%
\pgfpathmoveto{\pgfqpoint{1.480886in}{1.943542in}}%
\pgfpathcurveto{\pgfqpoint{1.486710in}{1.943542in}}{\pgfqpoint{1.492296in}{1.945855in}}{\pgfqpoint{1.496414in}{1.949974in}}%
\pgfpathcurveto{\pgfqpoint{1.500532in}{1.954092in}}{\pgfqpoint{1.502846in}{1.959678in}}{\pgfqpoint{1.502846in}{1.965502in}}%
\pgfpathcurveto{\pgfqpoint{1.502846in}{1.971326in}}{\pgfqpoint{1.500532in}{1.976912in}}{\pgfqpoint{1.496414in}{1.981030in}}%
\pgfpathcurveto{\pgfqpoint{1.492296in}{1.985148in}}{\pgfqpoint{1.486710in}{1.987462in}}{\pgfqpoint{1.480886in}{1.987462in}}%
\pgfpathcurveto{\pgfqpoint{1.475062in}{1.987462in}}{\pgfqpoint{1.469476in}{1.985148in}}{\pgfqpoint{1.465358in}{1.981030in}}%
\pgfpathcurveto{\pgfqpoint{1.461240in}{1.976912in}}{\pgfqpoint{1.458926in}{1.971326in}}{\pgfqpoint{1.458926in}{1.965502in}}%
\pgfpathcurveto{\pgfqpoint{1.458926in}{1.959678in}}{\pgfqpoint{1.461240in}{1.954092in}}{\pgfqpoint{1.465358in}{1.949974in}}%
\pgfpathcurveto{\pgfqpoint{1.469476in}{1.945855in}}{\pgfqpoint{1.475062in}{1.943542in}}{\pgfqpoint{1.480886in}{1.943542in}}%
\pgfpathlineto{\pgfqpoint{1.480886in}{1.943542in}}%
\pgfpathclose%
\pgfusepath{stroke,fill}%
\end{pgfscope}%
\begin{pgfscope}%
\pgfpathrectangle{\pgfqpoint{0.100000in}{0.209042in}}{\pgfqpoint{2.906250in}{2.906250in}}%
\pgfusepath{clip}%
\pgfsetbuttcap%
\pgfsetroundjoin%
\definecolor{currentfill}{rgb}{0.135053,0.068391,0.315000}%
\pgfsetfillcolor{currentfill}%
\pgfsetfillopacity{0.800000}%
\pgfsetlinewidth{1.003750pt}%
\definecolor{currentstroke}{rgb}{0.135053,0.068391,0.315000}%
\pgfsetstrokecolor{currentstroke}%
\pgfsetstrokeopacity{0.800000}%
\pgfsetdash{}{0pt}%
\pgfpathmoveto{\pgfqpoint{1.824157in}{1.722750in}}%
\pgfpathcurveto{\pgfqpoint{1.829981in}{1.722750in}}{\pgfqpoint{1.835567in}{1.725064in}}{\pgfqpoint{1.839685in}{1.729182in}}%
\pgfpathcurveto{\pgfqpoint{1.843803in}{1.733300in}}{\pgfqpoint{1.846117in}{1.738886in}}{\pgfqpoint{1.846117in}{1.744710in}}%
\pgfpathcurveto{\pgfqpoint{1.846117in}{1.750534in}}{\pgfqpoint{1.843803in}{1.756120in}}{\pgfqpoint{1.839685in}{1.760238in}}%
\pgfpathcurveto{\pgfqpoint{1.835567in}{1.764357in}}{\pgfqpoint{1.829981in}{1.766670in}}{\pgfqpoint{1.824157in}{1.766670in}}%
\pgfpathcurveto{\pgfqpoint{1.818333in}{1.766670in}}{\pgfqpoint{1.812747in}{1.764357in}}{\pgfqpoint{1.808629in}{1.760238in}}%
\pgfpathcurveto{\pgfqpoint{1.804511in}{1.756120in}}{\pgfqpoint{1.802197in}{1.750534in}}{\pgfqpoint{1.802197in}{1.744710in}}%
\pgfpathcurveto{\pgfqpoint{1.802197in}{1.738886in}}{\pgfqpoint{1.804511in}{1.733300in}}{\pgfqpoint{1.808629in}{1.729182in}}%
\pgfpathcurveto{\pgfqpoint{1.812747in}{1.725064in}}{\pgfqpoint{1.818333in}{1.722750in}}{\pgfqpoint{1.824157in}{1.722750in}}%
\pgfpathlineto{\pgfqpoint{1.824157in}{1.722750in}}%
\pgfpathclose%
\pgfusepath{stroke,fill}%
\end{pgfscope}%
\begin{pgfscope}%
\pgfpathrectangle{\pgfqpoint{0.100000in}{0.209042in}}{\pgfqpoint{2.906250in}{2.906250in}}%
\pgfusepath{clip}%
\pgfsetbuttcap%
\pgfsetroundjoin%
\definecolor{currentfill}{rgb}{0.684224,0.204286,0.484219}%
\pgfsetfillcolor{currentfill}%
\pgfsetfillopacity{0.800000}%
\pgfsetlinewidth{1.003750pt}%
\definecolor{currentstroke}{rgb}{0.684224,0.204286,0.484219}%
\pgfsetstrokecolor{currentstroke}%
\pgfsetstrokeopacity{0.800000}%
\pgfsetdash{}{0pt}%
\pgfpathmoveto{\pgfqpoint{1.871208in}{1.287233in}}%
\pgfpathcurveto{\pgfqpoint{1.877032in}{1.287233in}}{\pgfqpoint{1.882618in}{1.289547in}}{\pgfqpoint{1.886736in}{1.293665in}}%
\pgfpathcurveto{\pgfqpoint{1.890854in}{1.297783in}}{\pgfqpoint{1.893168in}{1.303369in}}{\pgfqpoint{1.893168in}{1.309193in}}%
\pgfpathcurveto{\pgfqpoint{1.893168in}{1.315017in}}{\pgfqpoint{1.890854in}{1.320603in}}{\pgfqpoint{1.886736in}{1.324721in}}%
\pgfpathcurveto{\pgfqpoint{1.882618in}{1.328840in}}{\pgfqpoint{1.877032in}{1.331153in}}{\pgfqpoint{1.871208in}{1.331153in}}%
\pgfpathcurveto{\pgfqpoint{1.865384in}{1.331153in}}{\pgfqpoint{1.859798in}{1.328840in}}{\pgfqpoint{1.855680in}{1.324721in}}%
\pgfpathcurveto{\pgfqpoint{1.851561in}{1.320603in}}{\pgfqpoint{1.849248in}{1.315017in}}{\pgfqpoint{1.849248in}{1.309193in}}%
\pgfpathcurveto{\pgfqpoint{1.849248in}{1.303369in}}{\pgfqpoint{1.851561in}{1.297783in}}{\pgfqpoint{1.855680in}{1.293665in}}%
\pgfpathcurveto{\pgfqpoint{1.859798in}{1.289547in}}{\pgfqpoint{1.865384in}{1.287233in}}{\pgfqpoint{1.871208in}{1.287233in}}%
\pgfpathlineto{\pgfqpoint{1.871208in}{1.287233in}}%
\pgfpathclose%
\pgfusepath{stroke,fill}%
\end{pgfscope}%
\begin{pgfscope}%
\pgfpathrectangle{\pgfqpoint{0.100000in}{0.209042in}}{\pgfqpoint{2.906250in}{2.906250in}}%
\pgfusepath{clip}%
\pgfsetbuttcap%
\pgfsetroundjoin%
\definecolor{currentfill}{rgb}{0.703532,0.210638,0.479049}%
\pgfsetfillcolor{currentfill}%
\pgfsetfillopacity{0.800000}%
\pgfsetlinewidth{1.003750pt}%
\definecolor{currentstroke}{rgb}{0.703532,0.210638,0.479049}%
\pgfsetstrokecolor{currentstroke}%
\pgfsetstrokeopacity{0.800000}%
\pgfsetdash{}{0pt}%
\pgfpathmoveto{\pgfqpoint{1.401174in}{2.128642in}}%
\pgfpathcurveto{\pgfqpoint{1.406998in}{2.128642in}}{\pgfqpoint{1.412584in}{2.130956in}}{\pgfqpoint{1.416702in}{2.135074in}}%
\pgfpathcurveto{\pgfqpoint{1.420820in}{2.139192in}}{\pgfqpoint{1.423134in}{2.144778in}}{\pgfqpoint{1.423134in}{2.150602in}}%
\pgfpathcurveto{\pgfqpoint{1.423134in}{2.156426in}}{\pgfqpoint{1.420820in}{2.162012in}}{\pgfqpoint{1.416702in}{2.166130in}}%
\pgfpathcurveto{\pgfqpoint{1.412584in}{2.170249in}}{\pgfqpoint{1.406998in}{2.172562in}}{\pgfqpoint{1.401174in}{2.172562in}}%
\pgfpathcurveto{\pgfqpoint{1.395350in}{2.172562in}}{\pgfqpoint{1.389764in}{2.170249in}}{\pgfqpoint{1.385646in}{2.166130in}}%
\pgfpathcurveto{\pgfqpoint{1.381528in}{2.162012in}}{\pgfqpoint{1.379214in}{2.156426in}}{\pgfqpoint{1.379214in}{2.150602in}}%
\pgfpathcurveto{\pgfqpoint{1.379214in}{2.144778in}}{\pgfqpoint{1.381528in}{2.139192in}}{\pgfqpoint{1.385646in}{2.135074in}}%
\pgfpathcurveto{\pgfqpoint{1.389764in}{2.130956in}}{\pgfqpoint{1.395350in}{2.128642in}}{\pgfqpoint{1.401174in}{2.128642in}}%
\pgfpathlineto{\pgfqpoint{1.401174in}{2.128642in}}%
\pgfpathclose%
\pgfusepath{stroke,fill}%
\end{pgfscope}%
\begin{pgfscope}%
\pgfpathrectangle{\pgfqpoint{0.100000in}{0.209042in}}{\pgfqpoint{2.906250in}{2.906250in}}%
\pgfusepath{clip}%
\pgfsetbuttcap%
\pgfsetroundjoin%
\definecolor{currentfill}{rgb}{0.165308,0.067911,0.361816}%
\pgfsetfillcolor{currentfill}%
\pgfsetfillopacity{0.800000}%
\pgfsetlinewidth{1.003750pt}%
\definecolor{currentstroke}{rgb}{0.165308,0.067911,0.361816}%
\pgfsetstrokecolor{currentstroke}%
\pgfsetstrokeopacity{0.800000}%
\pgfsetdash{}{0pt}%
\pgfpathmoveto{\pgfqpoint{1.823422in}{1.768195in}}%
\pgfpathcurveto{\pgfqpoint{1.829245in}{1.768195in}}{\pgfqpoint{1.834832in}{1.770509in}}{\pgfqpoint{1.838950in}{1.774627in}}%
\pgfpathcurveto{\pgfqpoint{1.843068in}{1.778745in}}{\pgfqpoint{1.845382in}{1.784331in}}{\pgfqpoint{1.845382in}{1.790155in}}%
\pgfpathcurveto{\pgfqpoint{1.845382in}{1.795979in}}{\pgfqpoint{1.843068in}{1.801565in}}{\pgfqpoint{1.838950in}{1.805683in}}%
\pgfpathcurveto{\pgfqpoint{1.834832in}{1.809801in}}{\pgfqpoint{1.829245in}{1.812115in}}{\pgfqpoint{1.823422in}{1.812115in}}%
\pgfpathcurveto{\pgfqpoint{1.817598in}{1.812115in}}{\pgfqpoint{1.812011in}{1.809801in}}{\pgfqpoint{1.807893in}{1.805683in}}%
\pgfpathcurveto{\pgfqpoint{1.803775in}{1.801565in}}{\pgfqpoint{1.801461in}{1.795979in}}{\pgfqpoint{1.801461in}{1.790155in}}%
\pgfpathcurveto{\pgfqpoint{1.801461in}{1.784331in}}{\pgfqpoint{1.803775in}{1.778745in}}{\pgfqpoint{1.807893in}{1.774627in}}%
\pgfpathcurveto{\pgfqpoint{1.812011in}{1.770509in}}{\pgfqpoint{1.817598in}{1.768195in}}{\pgfqpoint{1.823422in}{1.768195in}}%
\pgfpathlineto{\pgfqpoint{1.823422in}{1.768195in}}%
\pgfpathclose%
\pgfusepath{stroke,fill}%
\end{pgfscope}%
\begin{pgfscope}%
\pgfpathrectangle{\pgfqpoint{0.100000in}{0.209042in}}{\pgfqpoint{2.906250in}{2.906250in}}%
\pgfusepath{clip}%
\pgfsetbuttcap%
\pgfsetroundjoin%
\definecolor{currentfill}{rgb}{0.396467,0.102902,0.502658}%
\pgfsetfillcolor{currentfill}%
\pgfsetfillopacity{0.800000}%
\pgfsetlinewidth{1.003750pt}%
\definecolor{currentstroke}{rgb}{0.396467,0.102902,0.502658}%
\pgfsetstrokecolor{currentstroke}%
\pgfsetstrokeopacity{0.800000}%
\pgfsetdash{}{0pt}%
\pgfpathmoveto{\pgfqpoint{1.359861in}{1.517857in}}%
\pgfpathcurveto{\pgfqpoint{1.365685in}{1.517857in}}{\pgfqpoint{1.371271in}{1.520171in}}{\pgfqpoint{1.375390in}{1.524289in}}%
\pgfpathcurveto{\pgfqpoint{1.379508in}{1.528407in}}{\pgfqpoint{1.381822in}{1.533993in}}{\pgfqpoint{1.381822in}{1.539817in}}%
\pgfpathcurveto{\pgfqpoint{1.381822in}{1.545641in}}{\pgfqpoint{1.379508in}{1.551227in}}{\pgfqpoint{1.375390in}{1.555345in}}%
\pgfpathcurveto{\pgfqpoint{1.371271in}{1.559464in}}{\pgfqpoint{1.365685in}{1.561778in}}{\pgfqpoint{1.359861in}{1.561778in}}%
\pgfpathcurveto{\pgfqpoint{1.354037in}{1.561778in}}{\pgfqpoint{1.348451in}{1.559464in}}{\pgfqpoint{1.344333in}{1.555345in}}%
\pgfpathcurveto{\pgfqpoint{1.340215in}{1.551227in}}{\pgfqpoint{1.337901in}{1.545641in}}{\pgfqpoint{1.337901in}{1.539817in}}%
\pgfpathcurveto{\pgfqpoint{1.337901in}{1.533993in}}{\pgfqpoint{1.340215in}{1.528407in}}{\pgfqpoint{1.344333in}{1.524289in}}%
\pgfpathcurveto{\pgfqpoint{1.348451in}{1.520171in}}{\pgfqpoint{1.354037in}{1.517857in}}{\pgfqpoint{1.359861in}{1.517857in}}%
\pgfpathlineto{\pgfqpoint{1.359861in}{1.517857in}}%
\pgfpathclose%
\pgfusepath{stroke,fill}%
\end{pgfscope}%
\begin{pgfscope}%
\pgfpathrectangle{\pgfqpoint{0.100000in}{0.209042in}}{\pgfqpoint{2.906250in}{2.906250in}}%
\pgfusepath{clip}%
\pgfsetbuttcap%
\pgfsetroundjoin%
\definecolor{currentfill}{rgb}{0.304081,0.067835,0.480812}%
\pgfsetfillcolor{currentfill}%
\pgfsetfillopacity{0.800000}%
\pgfsetlinewidth{1.003750pt}%
\definecolor{currentstroke}{rgb}{0.304081,0.067835,0.480812}%
\pgfsetstrokecolor{currentstroke}%
\pgfsetstrokeopacity{0.800000}%
\pgfsetdash{}{0pt}%
\pgfpathmoveto{\pgfqpoint{1.709262in}{1.459221in}}%
\pgfpathcurveto{\pgfqpoint{1.715086in}{1.459221in}}{\pgfqpoint{1.720672in}{1.461535in}}{\pgfqpoint{1.724790in}{1.465653in}}%
\pgfpathcurveto{\pgfqpoint{1.728908in}{1.469771in}}{\pgfqpoint{1.731222in}{1.475358in}}{\pgfqpoint{1.731222in}{1.481181in}}%
\pgfpathcurveto{\pgfqpoint{1.731222in}{1.487005in}}{\pgfqpoint{1.728908in}{1.492592in}}{\pgfqpoint{1.724790in}{1.496710in}}%
\pgfpathcurveto{\pgfqpoint{1.720672in}{1.500828in}}{\pgfqpoint{1.715086in}{1.503142in}}{\pgfqpoint{1.709262in}{1.503142in}}%
\pgfpathcurveto{\pgfqpoint{1.703438in}{1.503142in}}{\pgfqpoint{1.697852in}{1.500828in}}{\pgfqpoint{1.693734in}{1.496710in}}%
\pgfpathcurveto{\pgfqpoint{1.689615in}{1.492592in}}{\pgfqpoint{1.687302in}{1.487005in}}{\pgfqpoint{1.687302in}{1.481181in}}%
\pgfpathcurveto{\pgfqpoint{1.687302in}{1.475358in}}{\pgfqpoint{1.689615in}{1.469771in}}{\pgfqpoint{1.693734in}{1.465653in}}%
\pgfpathcurveto{\pgfqpoint{1.697852in}{1.461535in}}{\pgfqpoint{1.703438in}{1.459221in}}{\pgfqpoint{1.709262in}{1.459221in}}%
\pgfpathlineto{\pgfqpoint{1.709262in}{1.459221in}}%
\pgfpathclose%
\pgfusepath{stroke,fill}%
\end{pgfscope}%
\begin{pgfscope}%
\pgfpathrectangle{\pgfqpoint{0.100000in}{0.209042in}}{\pgfqpoint{2.906250in}{2.906250in}}%
\pgfusepath{clip}%
\pgfsetbuttcap%
\pgfsetroundjoin%
\definecolor{currentfill}{rgb}{0.171713,0.067305,0.370771}%
\pgfsetfillcolor{currentfill}%
\pgfsetfillopacity{0.800000}%
\pgfsetlinewidth{1.003750pt}%
\definecolor{currentstroke}{rgb}{0.171713,0.067305,0.370771}%
\pgfsetstrokecolor{currentstroke}%
\pgfsetstrokeopacity{0.800000}%
\pgfsetdash{}{0pt}%
\pgfpathmoveto{\pgfqpoint{1.436616in}{1.627151in}}%
\pgfpathcurveto{\pgfqpoint{1.442440in}{1.627151in}}{\pgfqpoint{1.448026in}{1.629465in}}{\pgfqpoint{1.452144in}{1.633583in}}%
\pgfpathcurveto{\pgfqpoint{1.456263in}{1.637701in}}{\pgfqpoint{1.458576in}{1.643287in}}{\pgfqpoint{1.458576in}{1.649111in}}%
\pgfpathcurveto{\pgfqpoint{1.458576in}{1.654935in}}{\pgfqpoint{1.456263in}{1.660521in}}{\pgfqpoint{1.452144in}{1.664640in}}%
\pgfpathcurveto{\pgfqpoint{1.448026in}{1.668758in}}{\pgfqpoint{1.442440in}{1.671072in}}{\pgfqpoint{1.436616in}{1.671072in}}%
\pgfpathcurveto{\pgfqpoint{1.430792in}{1.671072in}}{\pgfqpoint{1.425206in}{1.668758in}}{\pgfqpoint{1.421088in}{1.664640in}}%
\pgfpathcurveto{\pgfqpoint{1.416970in}{1.660521in}}{\pgfqpoint{1.414656in}{1.654935in}}{\pgfqpoint{1.414656in}{1.649111in}}%
\pgfpathcurveto{\pgfqpoint{1.414656in}{1.643287in}}{\pgfqpoint{1.416970in}{1.637701in}}{\pgfqpoint{1.421088in}{1.633583in}}%
\pgfpathcurveto{\pgfqpoint{1.425206in}{1.629465in}}{\pgfqpoint{1.430792in}{1.627151in}}{\pgfqpoint{1.436616in}{1.627151in}}%
\pgfpathlineto{\pgfqpoint{1.436616in}{1.627151in}}%
\pgfpathclose%
\pgfusepath{stroke,fill}%
\end{pgfscope}%
\begin{pgfscope}%
\pgfpathrectangle{\pgfqpoint{0.100000in}{0.209042in}}{\pgfqpoint{2.906250in}{2.906250in}}%
\pgfusepath{clip}%
\pgfsetbuttcap%
\pgfsetroundjoin%
\definecolor{currentfill}{rgb}{0.432967,0.117855,0.506160}%
\pgfsetfillcolor{currentfill}%
\pgfsetfillopacity{0.800000}%
\pgfsetlinewidth{1.003750pt}%
\definecolor{currentstroke}{rgb}{0.432967,0.117855,0.506160}%
\pgfsetstrokecolor{currentstroke}%
\pgfsetstrokeopacity{0.800000}%
\pgfsetdash{}{0pt}%
\pgfpathmoveto{\pgfqpoint{1.466579in}{1.989629in}}%
\pgfpathcurveto{\pgfqpoint{1.472403in}{1.989629in}}{\pgfqpoint{1.477989in}{1.991943in}}{\pgfqpoint{1.482107in}{1.996061in}}%
\pgfpathcurveto{\pgfqpoint{1.486225in}{2.000179in}}{\pgfqpoint{1.488539in}{2.005765in}}{\pgfqpoint{1.488539in}{2.011589in}}%
\pgfpathcurveto{\pgfqpoint{1.488539in}{2.017413in}}{\pgfqpoint{1.486225in}{2.022999in}}{\pgfqpoint{1.482107in}{2.027118in}}%
\pgfpathcurveto{\pgfqpoint{1.477989in}{2.031236in}}{\pgfqpoint{1.472403in}{2.033550in}}{\pgfqpoint{1.466579in}{2.033550in}}%
\pgfpathcurveto{\pgfqpoint{1.460755in}{2.033550in}}{\pgfqpoint{1.455169in}{2.031236in}}{\pgfqpoint{1.451050in}{2.027118in}}%
\pgfpathcurveto{\pgfqpoint{1.446932in}{2.022999in}}{\pgfqpoint{1.444618in}{2.017413in}}{\pgfqpoint{1.444618in}{2.011589in}}%
\pgfpathcurveto{\pgfqpoint{1.444618in}{2.005765in}}{\pgfqpoint{1.446932in}{2.000179in}}{\pgfqpoint{1.451050in}{1.996061in}}%
\pgfpathcurveto{\pgfqpoint{1.455169in}{1.991943in}}{\pgfqpoint{1.460755in}{1.989629in}}{\pgfqpoint{1.466579in}{1.989629in}}%
\pgfpathlineto{\pgfqpoint{1.466579in}{1.989629in}}%
\pgfpathclose%
\pgfusepath{stroke,fill}%
\end{pgfscope}%
\begin{pgfscope}%
\pgfpathrectangle{\pgfqpoint{0.100000in}{0.209042in}}{\pgfqpoint{2.906250in}{2.906250in}}%
\pgfusepath{clip}%
\pgfsetbuttcap%
\pgfsetroundjoin%
\definecolor{currentfill}{rgb}{0.414709,0.110431,0.504662}%
\pgfsetfillcolor{currentfill}%
\pgfsetfillopacity{0.800000}%
\pgfsetlinewidth{1.003750pt}%
\definecolor{currentstroke}{rgb}{0.414709,0.110431,0.504662}%
\pgfsetstrokecolor{currentstroke}%
\pgfsetstrokeopacity{0.800000}%
\pgfsetdash{}{0pt}%
\pgfpathmoveto{\pgfqpoint{1.603306in}{2.004957in}}%
\pgfpathcurveto{\pgfqpoint{1.609130in}{2.004957in}}{\pgfqpoint{1.614716in}{2.007270in}}{\pgfqpoint{1.618834in}{2.011389in}}%
\pgfpathcurveto{\pgfqpoint{1.622952in}{2.015507in}}{\pgfqpoint{1.625266in}{2.021093in}}{\pgfqpoint{1.625266in}{2.026917in}}%
\pgfpathcurveto{\pgfqpoint{1.625266in}{2.032741in}}{\pgfqpoint{1.622952in}{2.038327in}}{\pgfqpoint{1.618834in}{2.042445in}}%
\pgfpathcurveto{\pgfqpoint{1.614716in}{2.046563in}}{\pgfqpoint{1.609130in}{2.048877in}}{\pgfqpoint{1.603306in}{2.048877in}}%
\pgfpathcurveto{\pgfqpoint{1.597482in}{2.048877in}}{\pgfqpoint{1.591896in}{2.046563in}}{\pgfqpoint{1.587778in}{2.042445in}}%
\pgfpathcurveto{\pgfqpoint{1.583659in}{2.038327in}}{\pgfqpoint{1.581345in}{2.032741in}}{\pgfqpoint{1.581345in}{2.026917in}}%
\pgfpathcurveto{\pgfqpoint{1.581345in}{2.021093in}}{\pgfqpoint{1.583659in}{2.015507in}}{\pgfqpoint{1.587778in}{2.011389in}}%
\pgfpathcurveto{\pgfqpoint{1.591896in}{2.007270in}}{\pgfqpoint{1.597482in}{2.004957in}}{\pgfqpoint{1.603306in}{2.004957in}}%
\pgfpathlineto{\pgfqpoint{1.603306in}{2.004957in}}%
\pgfpathclose%
\pgfusepath{stroke,fill}%
\end{pgfscope}%
\begin{pgfscope}%
\pgfpathrectangle{\pgfqpoint{0.100000in}{0.209042in}}{\pgfqpoint{2.906250in}{2.906250in}}%
\pgfusepath{clip}%
\pgfsetbuttcap%
\pgfsetroundjoin%
\definecolor{currentfill}{rgb}{0.291366,0.064553,0.475462}%
\pgfsetfillcolor{currentfill}%
\pgfsetfillopacity{0.800000}%
\pgfsetlinewidth{1.003750pt}%
\definecolor{currentstroke}{rgb}{0.291366,0.064553,0.475462}%
\pgfsetstrokecolor{currentstroke}%
\pgfsetstrokeopacity{0.800000}%
\pgfsetdash{}{0pt}%
\pgfpathmoveto{\pgfqpoint{1.339896in}{1.715308in}}%
\pgfpathcurveto{\pgfqpoint{1.345720in}{1.715308in}}{\pgfqpoint{1.351306in}{1.717622in}}{\pgfqpoint{1.355424in}{1.721740in}}%
\pgfpathcurveto{\pgfqpoint{1.359542in}{1.725858in}}{\pgfqpoint{1.361856in}{1.731444in}}{\pgfqpoint{1.361856in}{1.737268in}}%
\pgfpathcurveto{\pgfqpoint{1.361856in}{1.743092in}}{\pgfqpoint{1.359542in}{1.748678in}}{\pgfqpoint{1.355424in}{1.752796in}}%
\pgfpathcurveto{\pgfqpoint{1.351306in}{1.756915in}}{\pgfqpoint{1.345720in}{1.759228in}}{\pgfqpoint{1.339896in}{1.759228in}}%
\pgfpathcurveto{\pgfqpoint{1.334072in}{1.759228in}}{\pgfqpoint{1.328486in}{1.756915in}}{\pgfqpoint{1.324368in}{1.752796in}}%
\pgfpathcurveto{\pgfqpoint{1.320250in}{1.748678in}}{\pgfqpoint{1.317936in}{1.743092in}}{\pgfqpoint{1.317936in}{1.737268in}}%
\pgfpathcurveto{\pgfqpoint{1.317936in}{1.731444in}}{\pgfqpoint{1.320250in}{1.725858in}}{\pgfqpoint{1.324368in}{1.721740in}}%
\pgfpathcurveto{\pgfqpoint{1.328486in}{1.717622in}}{\pgfqpoint{1.334072in}{1.715308in}}{\pgfqpoint{1.339896in}{1.715308in}}%
\pgfpathlineto{\pgfqpoint{1.339896in}{1.715308in}}%
\pgfpathclose%
\pgfusepath{stroke,fill}%
\end{pgfscope}%
\begin{pgfscope}%
\pgfpathrectangle{\pgfqpoint{0.100000in}{0.209042in}}{\pgfqpoint{2.906250in}{2.906250in}}%
\pgfusepath{clip}%
\pgfsetbuttcap%
\pgfsetroundjoin%
\definecolor{currentfill}{rgb}{0.048062,0.036607,0.150327}%
\pgfsetfillcolor{currentfill}%
\pgfsetfillopacity{0.800000}%
\pgfsetlinewidth{1.003750pt}%
\definecolor{currentstroke}{rgb}{0.048062,0.036607,0.150327}%
\pgfsetstrokecolor{currentstroke}%
\pgfsetstrokeopacity{0.800000}%
\pgfsetdash{}{0pt}%
\pgfpathmoveto{\pgfqpoint{1.529370in}{1.664939in}}%
\pgfpathcurveto{\pgfqpoint{1.535194in}{1.664939in}}{\pgfqpoint{1.540780in}{1.667253in}}{\pgfqpoint{1.544898in}{1.671371in}}%
\pgfpathcurveto{\pgfqpoint{1.549016in}{1.675489in}}{\pgfqpoint{1.551330in}{1.681076in}}{\pgfqpoint{1.551330in}{1.686899in}}%
\pgfpathcurveto{\pgfqpoint{1.551330in}{1.692723in}}{\pgfqpoint{1.549016in}{1.698310in}}{\pgfqpoint{1.544898in}{1.702428in}}%
\pgfpathcurveto{\pgfqpoint{1.540780in}{1.706546in}}{\pgfqpoint{1.535194in}{1.708860in}}{\pgfqpoint{1.529370in}{1.708860in}}%
\pgfpathcurveto{\pgfqpoint{1.523546in}{1.708860in}}{\pgfqpoint{1.517959in}{1.706546in}}{\pgfqpoint{1.513841in}{1.702428in}}%
\pgfpathcurveto{\pgfqpoint{1.509723in}{1.698310in}}{\pgfqpoint{1.507409in}{1.692723in}}{\pgfqpoint{1.507409in}{1.686899in}}%
\pgfpathcurveto{\pgfqpoint{1.507409in}{1.681076in}}{\pgfqpoint{1.509723in}{1.675489in}}{\pgfqpoint{1.513841in}{1.671371in}}%
\pgfpathcurveto{\pgfqpoint{1.517959in}{1.667253in}}{\pgfqpoint{1.523546in}{1.664939in}}{\pgfqpoint{1.529370in}{1.664939in}}%
\pgfpathlineto{\pgfqpoint{1.529370in}{1.664939in}}%
\pgfpathclose%
\pgfusepath{stroke,fill}%
\end{pgfscope}%
\begin{pgfscope}%
\pgfpathrectangle{\pgfqpoint{0.100000in}{0.209042in}}{\pgfqpoint{2.906250in}{2.906250in}}%
\pgfusepath{clip}%
\pgfsetbuttcap%
\pgfsetroundjoin%
\definecolor{currentfill}{rgb}{0.359898,0.087831,0.496778}%
\pgfsetfillcolor{currentfill}%
\pgfsetfillopacity{0.800000}%
\pgfsetlinewidth{1.003750pt}%
\definecolor{currentstroke}{rgb}{0.359898,0.087831,0.496778}%
\pgfsetstrokecolor{currentstroke}%
\pgfsetstrokeopacity{0.800000}%
\pgfsetdash{}{0pt}%
\pgfpathmoveto{\pgfqpoint{1.708007in}{1.961395in}}%
\pgfpathcurveto{\pgfqpoint{1.713831in}{1.961395in}}{\pgfqpoint{1.719417in}{1.963709in}}{\pgfqpoint{1.723535in}{1.967827in}}%
\pgfpathcurveto{\pgfqpoint{1.727653in}{1.971946in}}{\pgfqpoint{1.729967in}{1.977532in}}{\pgfqpoint{1.729967in}{1.983356in}}%
\pgfpathcurveto{\pgfqpoint{1.729967in}{1.989180in}}{\pgfqpoint{1.727653in}{1.994766in}}{\pgfqpoint{1.723535in}{1.998884in}}%
\pgfpathcurveto{\pgfqpoint{1.719417in}{2.003002in}}{\pgfqpoint{1.713831in}{2.005316in}}{\pgfqpoint{1.708007in}{2.005316in}}%
\pgfpathcurveto{\pgfqpoint{1.702183in}{2.005316in}}{\pgfqpoint{1.696597in}{2.003002in}}{\pgfqpoint{1.692479in}{1.998884in}}%
\pgfpathcurveto{\pgfqpoint{1.688361in}{1.994766in}}{\pgfqpoint{1.686047in}{1.989180in}}{\pgfqpoint{1.686047in}{1.983356in}}%
\pgfpathcurveto{\pgfqpoint{1.686047in}{1.977532in}}{\pgfqpoint{1.688361in}{1.971946in}}{\pgfqpoint{1.692479in}{1.967827in}}%
\pgfpathcurveto{\pgfqpoint{1.696597in}{1.963709in}}{\pgfqpoint{1.702183in}{1.961395in}}{\pgfqpoint{1.708007in}{1.961395in}}%
\pgfpathlineto{\pgfqpoint{1.708007in}{1.961395in}}%
\pgfpathclose%
\pgfusepath{stroke,fill}%
\end{pgfscope}%
\begin{pgfscope}%
\pgfpathrectangle{\pgfqpoint{0.100000in}{0.209042in}}{\pgfqpoint{2.906250in}{2.906250in}}%
\pgfusepath{clip}%
\pgfsetbuttcap%
\pgfsetroundjoin%
\definecolor{currentfill}{rgb}{0.588158,0.173652,0.502035}%
\pgfsetfillcolor{currentfill}%
\pgfsetfillopacity{0.800000}%
\pgfsetlinewidth{1.003750pt}%
\definecolor{currentstroke}{rgb}{0.588158,0.173652,0.502035}%
\pgfsetstrokecolor{currentstroke}%
\pgfsetstrokeopacity{0.800000}%
\pgfsetdash{}{0pt}%
\pgfpathmoveto{\pgfqpoint{1.184610in}{1.878136in}}%
\pgfpathcurveto{\pgfqpoint{1.190433in}{1.878136in}}{\pgfqpoint{1.196020in}{1.880450in}}{\pgfqpoint{1.200138in}{1.884568in}}%
\pgfpathcurveto{\pgfqpoint{1.204256in}{1.888686in}}{\pgfqpoint{1.206570in}{1.894272in}}{\pgfqpoint{1.206570in}{1.900096in}}%
\pgfpathcurveto{\pgfqpoint{1.206570in}{1.905920in}}{\pgfqpoint{1.204256in}{1.911506in}}{\pgfqpoint{1.200138in}{1.915624in}}%
\pgfpathcurveto{\pgfqpoint{1.196020in}{1.919742in}}{\pgfqpoint{1.190433in}{1.922056in}}{\pgfqpoint{1.184610in}{1.922056in}}%
\pgfpathcurveto{\pgfqpoint{1.178786in}{1.922056in}}{\pgfqpoint{1.173199in}{1.919742in}}{\pgfqpoint{1.169081in}{1.915624in}}%
\pgfpathcurveto{\pgfqpoint{1.164963in}{1.911506in}}{\pgfqpoint{1.162649in}{1.905920in}}{\pgfqpoint{1.162649in}{1.900096in}}%
\pgfpathcurveto{\pgfqpoint{1.162649in}{1.894272in}}{\pgfqpoint{1.164963in}{1.888686in}}{\pgfqpoint{1.169081in}{1.884568in}}%
\pgfpathcurveto{\pgfqpoint{1.173199in}{1.880450in}}{\pgfqpoint{1.178786in}{1.878136in}}{\pgfqpoint{1.184610in}{1.878136in}}%
\pgfpathlineto{\pgfqpoint{1.184610in}{1.878136in}}%
\pgfpathclose%
\pgfusepath{stroke,fill}%
\end{pgfscope}%
\begin{pgfscope}%
\pgfpathrectangle{\pgfqpoint{0.100000in}{0.209042in}}{\pgfqpoint{2.906250in}{2.906250in}}%
\pgfusepath{clip}%
\pgfsetbuttcap%
\pgfsetroundjoin%
\definecolor{currentfill}{rgb}{0.007588,0.006356,0.044973}%
\pgfsetfillcolor{currentfill}%
\pgfsetfillopacity{0.800000}%
\pgfsetlinewidth{1.003750pt}%
\definecolor{currentstroke}{rgb}{0.007588,0.006356,0.044973}%
\pgfsetstrokecolor{currentstroke}%
\pgfsetstrokeopacity{0.800000}%
\pgfsetdash{}{0pt}%
\pgfpathmoveto{\pgfqpoint{1.595502in}{1.680234in}}%
\pgfpathcurveto{\pgfqpoint{1.601326in}{1.680234in}}{\pgfqpoint{1.606912in}{1.682548in}}{\pgfqpoint{1.611030in}{1.686666in}}%
\pgfpathcurveto{\pgfqpoint{1.615148in}{1.690784in}}{\pgfqpoint{1.617462in}{1.696370in}}{\pgfqpoint{1.617462in}{1.702194in}}%
\pgfpathcurveto{\pgfqpoint{1.617462in}{1.708018in}}{\pgfqpoint{1.615148in}{1.713604in}}{\pgfqpoint{1.611030in}{1.717722in}}%
\pgfpathcurveto{\pgfqpoint{1.606912in}{1.721840in}}{\pgfqpoint{1.601326in}{1.724154in}}{\pgfqpoint{1.595502in}{1.724154in}}%
\pgfpathcurveto{\pgfqpoint{1.589678in}{1.724154in}}{\pgfqpoint{1.584092in}{1.721840in}}{\pgfqpoint{1.579973in}{1.717722in}}%
\pgfpathcurveto{\pgfqpoint{1.575855in}{1.713604in}}{\pgfqpoint{1.573541in}{1.708018in}}{\pgfqpoint{1.573541in}{1.702194in}}%
\pgfpathcurveto{\pgfqpoint{1.573541in}{1.696370in}}{\pgfqpoint{1.575855in}{1.690784in}}{\pgfqpoint{1.579973in}{1.686666in}}%
\pgfpathcurveto{\pgfqpoint{1.584092in}{1.682548in}}{\pgfqpoint{1.589678in}{1.680234in}}{\pgfqpoint{1.595502in}{1.680234in}}%
\pgfpathlineto{\pgfqpoint{1.595502in}{1.680234in}}%
\pgfpathclose%
\pgfusepath{stroke,fill}%
\end{pgfscope}%
\begin{pgfscope}%
\pgfpathrectangle{\pgfqpoint{0.100000in}{0.209042in}}{\pgfqpoint{2.906250in}{2.906250in}}%
\pgfusepath{clip}%
\pgfsetbuttcap%
\pgfsetroundjoin%
\definecolor{currentfill}{rgb}{0.271994,0.060994,0.465660}%
\pgfsetfillcolor{currentfill}%
\pgfsetfillopacity{0.800000}%
\pgfsetlinewidth{1.003750pt}%
\definecolor{currentstroke}{rgb}{0.271994,0.060994,0.465660}%
\pgfsetstrokecolor{currentstroke}%
\pgfsetstrokeopacity{0.800000}%
\pgfsetdash{}{0pt}%
\pgfpathmoveto{\pgfqpoint{1.466246in}{1.525539in}}%
\pgfpathcurveto{\pgfqpoint{1.472070in}{1.525539in}}{\pgfqpoint{1.477657in}{1.527853in}}{\pgfqpoint{1.481775in}{1.531971in}}%
\pgfpathcurveto{\pgfqpoint{1.485893in}{1.536089in}}{\pgfqpoint{1.488207in}{1.541675in}}{\pgfqpoint{1.488207in}{1.547499in}}%
\pgfpathcurveto{\pgfqpoint{1.488207in}{1.553323in}}{\pgfqpoint{1.485893in}{1.558909in}}{\pgfqpoint{1.481775in}{1.563027in}}%
\pgfpathcurveto{\pgfqpoint{1.477657in}{1.567146in}}{\pgfqpoint{1.472070in}{1.569459in}}{\pgfqpoint{1.466246in}{1.569459in}}%
\pgfpathcurveto{\pgfqpoint{1.460423in}{1.569459in}}{\pgfqpoint{1.454836in}{1.567146in}}{\pgfqpoint{1.450718in}{1.563027in}}%
\pgfpathcurveto{\pgfqpoint{1.446600in}{1.558909in}}{\pgfqpoint{1.444286in}{1.553323in}}{\pgfqpoint{1.444286in}{1.547499in}}%
\pgfpathcurveto{\pgfqpoint{1.444286in}{1.541675in}}{\pgfqpoint{1.446600in}{1.536089in}}{\pgfqpoint{1.450718in}{1.531971in}}%
\pgfpathcurveto{\pgfqpoint{1.454836in}{1.527853in}}{\pgfqpoint{1.460423in}{1.525539in}}{\pgfqpoint{1.466246in}{1.525539in}}%
\pgfpathlineto{\pgfqpoint{1.466246in}{1.525539in}}%
\pgfpathclose%
\pgfusepath{stroke,fill}%
\end{pgfscope}%
\begin{pgfscope}%
\pgfpathrectangle{\pgfqpoint{0.100000in}{0.209042in}}{\pgfqpoint{2.906250in}{2.906250in}}%
\pgfusepath{clip}%
\pgfsetbuttcap%
\pgfsetroundjoin%
\definecolor{currentfill}{rgb}{0.716387,0.214982,0.475290}%
\pgfsetfillcolor{currentfill}%
\pgfsetfillopacity{0.800000}%
\pgfsetlinewidth{1.003750pt}%
\definecolor{currentstroke}{rgb}{0.716387,0.214982,0.475290}%
\pgfsetstrokecolor{currentstroke}%
\pgfsetstrokeopacity{0.800000}%
\pgfsetdash{}{0pt}%
\pgfpathmoveto{\pgfqpoint{2.188235in}{1.784977in}}%
\pgfpathcurveto{\pgfqpoint{2.194059in}{1.784977in}}{\pgfqpoint{2.199645in}{1.787290in}}{\pgfqpoint{2.203763in}{1.791409in}}%
\pgfpathcurveto{\pgfqpoint{2.207881in}{1.795527in}}{\pgfqpoint{2.210195in}{1.801113in}}{\pgfqpoint{2.210195in}{1.806937in}}%
\pgfpathcurveto{\pgfqpoint{2.210195in}{1.812761in}}{\pgfqpoint{2.207881in}{1.818347in}}{\pgfqpoint{2.203763in}{1.822465in}}%
\pgfpathcurveto{\pgfqpoint{2.199645in}{1.826583in}}{\pgfqpoint{2.194059in}{1.828897in}}{\pgfqpoint{2.188235in}{1.828897in}}%
\pgfpathcurveto{\pgfqpoint{2.182411in}{1.828897in}}{\pgfqpoint{2.176825in}{1.826583in}}{\pgfqpoint{2.172706in}{1.822465in}}%
\pgfpathcurveto{\pgfqpoint{2.168588in}{1.818347in}}{\pgfqpoint{2.166274in}{1.812761in}}{\pgfqpoint{2.166274in}{1.806937in}}%
\pgfpathcurveto{\pgfqpoint{2.166274in}{1.801113in}}{\pgfqpoint{2.168588in}{1.795527in}}{\pgfqpoint{2.172706in}{1.791409in}}%
\pgfpathcurveto{\pgfqpoint{2.176825in}{1.787290in}}{\pgfqpoint{2.182411in}{1.784977in}}{\pgfqpoint{2.188235in}{1.784977in}}%
\pgfpathlineto{\pgfqpoint{2.188235in}{1.784977in}}%
\pgfpathclose%
\pgfusepath{stroke,fill}%
\end{pgfscope}%
\begin{pgfscope}%
\pgfpathrectangle{\pgfqpoint{0.100000in}{0.209042in}}{\pgfqpoint{2.906250in}{2.906250in}}%
\pgfusepath{clip}%
\pgfsetbuttcap%
\pgfsetroundjoin%
\definecolor{currentfill}{rgb}{0.005950,0.004843,0.037130}%
\pgfsetfillcolor{currentfill}%
\pgfsetfillopacity{0.800000}%
\pgfsetlinewidth{1.003750pt}%
\definecolor{currentstroke}{rgb}{0.005950,0.004843,0.037130}%
\pgfsetstrokecolor{currentstroke}%
\pgfsetstrokeopacity{0.800000}%
\pgfsetdash{}{0pt}%
\pgfpathmoveto{\pgfqpoint{1.662262in}{1.688173in}}%
\pgfpathcurveto{\pgfqpoint{1.668086in}{1.688173in}}{\pgfqpoint{1.673672in}{1.690487in}}{\pgfqpoint{1.677790in}{1.694605in}}%
\pgfpathcurveto{\pgfqpoint{1.681908in}{1.698723in}}{\pgfqpoint{1.684222in}{1.704309in}}{\pgfqpoint{1.684222in}{1.710133in}}%
\pgfpathcurveto{\pgfqpoint{1.684222in}{1.715957in}}{\pgfqpoint{1.681908in}{1.721543in}}{\pgfqpoint{1.677790in}{1.725662in}}%
\pgfpathcurveto{\pgfqpoint{1.673672in}{1.729780in}}{\pgfqpoint{1.668086in}{1.732094in}}{\pgfqpoint{1.662262in}{1.732094in}}%
\pgfpathcurveto{\pgfqpoint{1.656438in}{1.732094in}}{\pgfqpoint{1.650852in}{1.729780in}}{\pgfqpoint{1.646734in}{1.725662in}}%
\pgfpathcurveto{\pgfqpoint{1.642616in}{1.721543in}}{\pgfqpoint{1.640302in}{1.715957in}}{\pgfqpoint{1.640302in}{1.710133in}}%
\pgfpathcurveto{\pgfqpoint{1.640302in}{1.704309in}}{\pgfqpoint{1.642616in}{1.698723in}}{\pgfqpoint{1.646734in}{1.694605in}}%
\pgfpathcurveto{\pgfqpoint{1.650852in}{1.690487in}}{\pgfqpoint{1.656438in}{1.688173in}}{\pgfqpoint{1.662262in}{1.688173in}}%
\pgfpathlineto{\pgfqpoint{1.662262in}{1.688173in}}%
\pgfpathclose%
\pgfusepath{stroke,fill}%
\end{pgfscope}%
\begin{pgfscope}%
\pgfpathrectangle{\pgfqpoint{0.100000in}{0.209042in}}{\pgfqpoint{2.906250in}{2.906250in}}%
\pgfusepath{clip}%
\pgfsetbuttcap%
\pgfsetroundjoin%
\definecolor{currentfill}{rgb}{0.940687,0.371224,0.366762}%
\pgfsetfillcolor{currentfill}%
\pgfsetfillopacity{0.800000}%
\pgfsetlinewidth{1.003750pt}%
\definecolor{currentstroke}{rgb}{0.940687,0.371224,0.366762}%
\pgfsetstrokecolor{currentstroke}%
\pgfsetstrokeopacity{0.800000}%
\pgfsetdash{}{0pt}%
\pgfpathmoveto{\pgfqpoint{1.564370in}{1.105633in}}%
\pgfpathcurveto{\pgfqpoint{1.570194in}{1.105633in}}{\pgfqpoint{1.575780in}{1.107947in}}{\pgfqpoint{1.579898in}{1.112065in}}%
\pgfpathcurveto{\pgfqpoint{1.584016in}{1.116183in}}{\pgfqpoint{1.586330in}{1.121769in}}{\pgfqpoint{1.586330in}{1.127593in}}%
\pgfpathcurveto{\pgfqpoint{1.586330in}{1.133417in}}{\pgfqpoint{1.584016in}{1.139003in}}{\pgfqpoint{1.579898in}{1.143121in}}%
\pgfpathcurveto{\pgfqpoint{1.575780in}{1.147240in}}{\pgfqpoint{1.570194in}{1.149553in}}{\pgfqpoint{1.564370in}{1.149553in}}%
\pgfpathcurveto{\pgfqpoint{1.558546in}{1.149553in}}{\pgfqpoint{1.552960in}{1.147240in}}{\pgfqpoint{1.548842in}{1.143121in}}%
\pgfpathcurveto{\pgfqpoint{1.544723in}{1.139003in}}{\pgfqpoint{1.542410in}{1.133417in}}{\pgfqpoint{1.542410in}{1.127593in}}%
\pgfpathcurveto{\pgfqpoint{1.542410in}{1.121769in}}{\pgfqpoint{1.544723in}{1.116183in}}{\pgfqpoint{1.548842in}{1.112065in}}%
\pgfpathcurveto{\pgfqpoint{1.552960in}{1.107947in}}{\pgfqpoint{1.558546in}{1.105633in}}{\pgfqpoint{1.564370in}{1.105633in}}%
\pgfpathlineto{\pgfqpoint{1.564370in}{1.105633in}}%
\pgfpathclose%
\pgfusepath{stroke,fill}%
\end{pgfscope}%
\begin{pgfscope}%
\pgfpathrectangle{\pgfqpoint{0.100000in}{0.209042in}}{\pgfqpoint{2.906250in}{2.906250in}}%
\pgfusepath{clip}%
\pgfsetbuttcap%
\pgfsetroundjoin%
\definecolor{currentfill}{rgb}{0.113094,0.065492,0.276784}%
\pgfsetfillcolor{currentfill}%
\pgfsetfillopacity{0.800000}%
\pgfsetlinewidth{1.003750pt}%
\definecolor{currentstroke}{rgb}{0.113094,0.065492,0.276784}%
\pgfsetstrokecolor{currentstroke}%
\pgfsetstrokeopacity{0.800000}%
\pgfsetdash{}{0pt}%
\pgfpathmoveto{\pgfqpoint{1.769275in}{1.777220in}}%
\pgfpathcurveto{\pgfqpoint{1.775099in}{1.777220in}}{\pgfqpoint{1.780686in}{1.779534in}}{\pgfqpoint{1.784804in}{1.783652in}}%
\pgfpathcurveto{\pgfqpoint{1.788922in}{1.787771in}}{\pgfqpoint{1.791236in}{1.793357in}}{\pgfqpoint{1.791236in}{1.799181in}}%
\pgfpathcurveto{\pgfqpoint{1.791236in}{1.805005in}}{\pgfqpoint{1.788922in}{1.810591in}}{\pgfqpoint{1.784804in}{1.814709in}}%
\pgfpathcurveto{\pgfqpoint{1.780686in}{1.818827in}}{\pgfqpoint{1.775099in}{1.821141in}}{\pgfqpoint{1.769275in}{1.821141in}}%
\pgfpathcurveto{\pgfqpoint{1.763451in}{1.821141in}}{\pgfqpoint{1.757865in}{1.818827in}}{\pgfqpoint{1.753747in}{1.814709in}}%
\pgfpathcurveto{\pgfqpoint{1.749629in}{1.810591in}}{\pgfqpoint{1.747315in}{1.805005in}}{\pgfqpoint{1.747315in}{1.799181in}}%
\pgfpathcurveto{\pgfqpoint{1.747315in}{1.793357in}}{\pgfqpoint{1.749629in}{1.787771in}}{\pgfqpoint{1.753747in}{1.783652in}}%
\pgfpathcurveto{\pgfqpoint{1.757865in}{1.779534in}}{\pgfqpoint{1.763451in}{1.777220in}}{\pgfqpoint{1.769275in}{1.777220in}}%
\pgfpathlineto{\pgfqpoint{1.769275in}{1.777220in}}%
\pgfpathclose%
\pgfusepath{stroke,fill}%
\end{pgfscope}%
\begin{pgfscope}%
\pgfpathrectangle{\pgfqpoint{0.100000in}{0.209042in}}{\pgfqpoint{2.906250in}{2.906250in}}%
\pgfusepath{clip}%
\pgfsetbuttcap%
\pgfsetroundjoin%
\definecolor{currentfill}{rgb}{0.556571,0.163269,0.505230}%
\pgfsetfillcolor{currentfill}%
\pgfsetfillopacity{0.800000}%
\pgfsetlinewidth{1.003750pt}%
\definecolor{currentstroke}{rgb}{0.556571,0.163269,0.505230}%
\pgfsetstrokecolor{currentstroke}%
\pgfsetstrokeopacity{0.800000}%
\pgfsetdash{}{0pt}%
\pgfpathmoveto{\pgfqpoint{1.167438in}{1.749559in}}%
\pgfpathcurveto{\pgfqpoint{1.173262in}{1.749559in}}{\pgfqpoint{1.178848in}{1.751873in}}{\pgfqpoint{1.182966in}{1.755991in}}%
\pgfpathcurveto{\pgfqpoint{1.187085in}{1.760109in}}{\pgfqpoint{1.189398in}{1.765695in}}{\pgfqpoint{1.189398in}{1.771519in}}%
\pgfpathcurveto{\pgfqpoint{1.189398in}{1.777343in}}{\pgfqpoint{1.187085in}{1.782929in}}{\pgfqpoint{1.182966in}{1.787047in}}%
\pgfpathcurveto{\pgfqpoint{1.178848in}{1.791166in}}{\pgfqpoint{1.173262in}{1.793479in}}{\pgfqpoint{1.167438in}{1.793479in}}%
\pgfpathcurveto{\pgfqpoint{1.161614in}{1.793479in}}{\pgfqpoint{1.156028in}{1.791166in}}{\pgfqpoint{1.151910in}{1.787047in}}%
\pgfpathcurveto{\pgfqpoint{1.147792in}{1.782929in}}{\pgfqpoint{1.145478in}{1.777343in}}{\pgfqpoint{1.145478in}{1.771519in}}%
\pgfpathcurveto{\pgfqpoint{1.145478in}{1.765695in}}{\pgfqpoint{1.147792in}{1.760109in}}{\pgfqpoint{1.151910in}{1.755991in}}%
\pgfpathcurveto{\pgfqpoint{1.156028in}{1.751873in}}{\pgfqpoint{1.161614in}{1.749559in}}{\pgfqpoint{1.167438in}{1.749559in}}%
\pgfpathlineto{\pgfqpoint{1.167438in}{1.749559in}}%
\pgfpathclose%
\pgfusepath{stroke,fill}%
\end{pgfscope}%
\begin{pgfscope}%
\pgfpathrectangle{\pgfqpoint{0.100000in}{0.209042in}}{\pgfqpoint{2.906250in}{2.906250in}}%
\pgfusepath{clip}%
\pgfsetbuttcap%
\pgfsetroundjoin%
\definecolor{currentfill}{rgb}{0.028123,0.023201,0.108787}%
\pgfsetfillcolor{currentfill}%
\pgfsetfillopacity{0.800000}%
\pgfsetlinewidth{1.003750pt}%
\definecolor{currentstroke}{rgb}{0.028123,0.023201,0.108787}%
\pgfsetstrokecolor{currentstroke}%
\pgfsetstrokeopacity{0.800000}%
\pgfsetdash{}{0pt}%
\pgfpathmoveto{\pgfqpoint{1.686813in}{1.645046in}}%
\pgfpathcurveto{\pgfqpoint{1.692637in}{1.645046in}}{\pgfqpoint{1.698223in}{1.647360in}}{\pgfqpoint{1.702342in}{1.651478in}}%
\pgfpathcurveto{\pgfqpoint{1.706460in}{1.655596in}}{\pgfqpoint{1.708774in}{1.661182in}}{\pgfqpoint{1.708774in}{1.667006in}}%
\pgfpathcurveto{\pgfqpoint{1.708774in}{1.672830in}}{\pgfqpoint{1.706460in}{1.678416in}}{\pgfqpoint{1.702342in}{1.682535in}}%
\pgfpathcurveto{\pgfqpoint{1.698223in}{1.686653in}}{\pgfqpoint{1.692637in}{1.688967in}}{\pgfqpoint{1.686813in}{1.688967in}}%
\pgfpathcurveto{\pgfqpoint{1.680989in}{1.688967in}}{\pgfqpoint{1.675403in}{1.686653in}}{\pgfqpoint{1.671285in}{1.682535in}}%
\pgfpathcurveto{\pgfqpoint{1.667167in}{1.678416in}}{\pgfqpoint{1.664853in}{1.672830in}}{\pgfqpoint{1.664853in}{1.667006in}}%
\pgfpathcurveto{\pgfqpoint{1.664853in}{1.661182in}}{\pgfqpoint{1.667167in}{1.655596in}}{\pgfqpoint{1.671285in}{1.651478in}}%
\pgfpathcurveto{\pgfqpoint{1.675403in}{1.647360in}}{\pgfqpoint{1.680989in}{1.645046in}}{\pgfqpoint{1.686813in}{1.645046in}}%
\pgfpathlineto{\pgfqpoint{1.686813in}{1.645046in}}%
\pgfpathclose%
\pgfusepath{stroke,fill}%
\end{pgfscope}%
\begin{pgfscope}%
\pgfpathrectangle{\pgfqpoint{0.100000in}{0.209042in}}{\pgfqpoint{2.906250in}{2.906250in}}%
\pgfusepath{clip}%
\pgfsetbuttcap%
\pgfsetroundjoin%
\definecolor{currentfill}{rgb}{0.652056,0.193986,0.491611}%
\pgfsetfillcolor{currentfill}%
\pgfsetfillopacity{0.800000}%
\pgfsetlinewidth{1.003750pt}%
\definecolor{currentstroke}{rgb}{0.652056,0.193986,0.491611}%
\pgfsetstrokecolor{currentstroke}%
\pgfsetstrokeopacity{0.800000}%
\pgfsetdash{}{0pt}%
\pgfpathmoveto{\pgfqpoint{1.261649in}{1.411410in}}%
\pgfpathcurveto{\pgfqpoint{1.267473in}{1.411410in}}{\pgfqpoint{1.273059in}{1.413724in}}{\pgfqpoint{1.277178in}{1.417842in}}%
\pgfpathcurveto{\pgfqpoint{1.281296in}{1.421960in}}{\pgfqpoint{1.283610in}{1.427546in}}{\pgfqpoint{1.283610in}{1.433370in}}%
\pgfpathcurveto{\pgfqpoint{1.283610in}{1.439194in}}{\pgfqpoint{1.281296in}{1.444780in}}{\pgfqpoint{1.277178in}{1.448898in}}%
\pgfpathcurveto{\pgfqpoint{1.273059in}{1.453016in}}{\pgfqpoint{1.267473in}{1.455330in}}{\pgfqpoint{1.261649in}{1.455330in}}%
\pgfpathcurveto{\pgfqpoint{1.255825in}{1.455330in}}{\pgfqpoint{1.250239in}{1.453016in}}{\pgfqpoint{1.246121in}{1.448898in}}%
\pgfpathcurveto{\pgfqpoint{1.242003in}{1.444780in}}{\pgfqpoint{1.239689in}{1.439194in}}{\pgfqpoint{1.239689in}{1.433370in}}%
\pgfpathcurveto{\pgfqpoint{1.239689in}{1.427546in}}{\pgfqpoint{1.242003in}{1.421960in}}{\pgfqpoint{1.246121in}{1.417842in}}%
\pgfpathcurveto{\pgfqpoint{1.250239in}{1.413724in}}{\pgfqpoint{1.255825in}{1.411410in}}{\pgfqpoint{1.261649in}{1.411410in}}%
\pgfpathlineto{\pgfqpoint{1.261649in}{1.411410in}}%
\pgfpathclose%
\pgfusepath{stroke,fill}%
\end{pgfscope}%
\begin{pgfscope}%
\pgfpathrectangle{\pgfqpoint{0.100000in}{0.209042in}}{\pgfqpoint{2.906250in}{2.906250in}}%
\pgfusepath{clip}%
\pgfsetbuttcap%
\pgfsetroundjoin%
\definecolor{currentfill}{rgb}{0.329114,0.075972,0.489287}%
\pgfsetfillcolor{currentfill}%
\pgfsetfillopacity{0.800000}%
\pgfsetlinewidth{1.003750pt}%
\definecolor{currentstroke}{rgb}{0.329114,0.075972,0.489287}%
\pgfsetstrokecolor{currentstroke}%
\pgfsetstrokeopacity{0.800000}%
\pgfsetdash{}{0pt}%
\pgfpathmoveto{\pgfqpoint{1.845393in}{1.877193in}}%
\pgfpathcurveto{\pgfqpoint{1.851217in}{1.877193in}}{\pgfqpoint{1.856804in}{1.879507in}}{\pgfqpoint{1.860922in}{1.883625in}}%
\pgfpathcurveto{\pgfqpoint{1.865040in}{1.887743in}}{\pgfqpoint{1.867354in}{1.893330in}}{\pgfqpoint{1.867354in}{1.899154in}}%
\pgfpathcurveto{\pgfqpoint{1.867354in}{1.904978in}}{\pgfqpoint{1.865040in}{1.910564in}}{\pgfqpoint{1.860922in}{1.914682in}}%
\pgfpathcurveto{\pgfqpoint{1.856804in}{1.918800in}}{\pgfqpoint{1.851217in}{1.921114in}}{\pgfqpoint{1.845393in}{1.921114in}}%
\pgfpathcurveto{\pgfqpoint{1.839569in}{1.921114in}}{\pgfqpoint{1.833983in}{1.918800in}}{\pgfqpoint{1.829865in}{1.914682in}}%
\pgfpathcurveto{\pgfqpoint{1.825747in}{1.910564in}}{\pgfqpoint{1.823433in}{1.904978in}}{\pgfqpoint{1.823433in}{1.899154in}}%
\pgfpathcurveto{\pgfqpoint{1.823433in}{1.893330in}}{\pgfqpoint{1.825747in}{1.887743in}}{\pgfqpoint{1.829865in}{1.883625in}}%
\pgfpathcurveto{\pgfqpoint{1.833983in}{1.879507in}}{\pgfqpoint{1.839569in}{1.877193in}}{\pgfqpoint{1.845393in}{1.877193in}}%
\pgfpathlineto{\pgfqpoint{1.845393in}{1.877193in}}%
\pgfpathclose%
\pgfusepath{stroke,fill}%
\end{pgfscope}%
\begin{pgfscope}%
\pgfpathrectangle{\pgfqpoint{0.100000in}{0.209042in}}{\pgfqpoint{2.906250in}{2.906250in}}%
\pgfusepath{clip}%
\pgfsetbuttcap%
\pgfsetroundjoin%
\definecolor{currentfill}{rgb}{0.271994,0.060994,0.465660}%
\pgfsetfillcolor{currentfill}%
\pgfsetfillopacity{0.800000}%
\pgfsetlinewidth{1.003750pt}%
\definecolor{currentstroke}{rgb}{0.271994,0.060994,0.465660}%
\pgfsetstrokecolor{currentstroke}%
\pgfsetstrokeopacity{0.800000}%
\pgfsetdash{}{0pt}%
\pgfpathmoveto{\pgfqpoint{1.360729in}{1.721320in}}%
\pgfpathcurveto{\pgfqpoint{1.366553in}{1.721320in}}{\pgfqpoint{1.372140in}{1.723634in}}{\pgfqpoint{1.376258in}{1.727752in}}%
\pgfpathcurveto{\pgfqpoint{1.380376in}{1.731870in}}{\pgfqpoint{1.382690in}{1.737456in}}{\pgfqpoint{1.382690in}{1.743280in}}%
\pgfpathcurveto{\pgfqpoint{1.382690in}{1.749104in}}{\pgfqpoint{1.380376in}{1.754690in}}{\pgfqpoint{1.376258in}{1.758808in}}%
\pgfpathcurveto{\pgfqpoint{1.372140in}{1.762926in}}{\pgfqpoint{1.366553in}{1.765240in}}{\pgfqpoint{1.360729in}{1.765240in}}%
\pgfpathcurveto{\pgfqpoint{1.354906in}{1.765240in}}{\pgfqpoint{1.349319in}{1.762926in}}{\pgfqpoint{1.345201in}{1.758808in}}%
\pgfpathcurveto{\pgfqpoint{1.341083in}{1.754690in}}{\pgfqpoint{1.338769in}{1.749104in}}{\pgfqpoint{1.338769in}{1.743280in}}%
\pgfpathcurveto{\pgfqpoint{1.338769in}{1.737456in}}{\pgfqpoint{1.341083in}{1.731870in}}{\pgfqpoint{1.345201in}{1.727752in}}%
\pgfpathcurveto{\pgfqpoint{1.349319in}{1.723634in}}{\pgfqpoint{1.354906in}{1.721320in}}{\pgfqpoint{1.360729in}{1.721320in}}%
\pgfpathlineto{\pgfqpoint{1.360729in}{1.721320in}}%
\pgfpathclose%
\pgfusepath{stroke,fill}%
\end{pgfscope}%
\begin{pgfscope}%
\pgfpathrectangle{\pgfqpoint{0.100000in}{0.209042in}}{\pgfqpoint{2.906250in}{2.906250in}}%
\pgfusepath{clip}%
\pgfsetbuttcap%
\pgfsetroundjoin%
\definecolor{currentfill}{rgb}{0.469640,0.132245,0.507809}%
\pgfsetfillcolor{currentfill}%
\pgfsetfillopacity{0.800000}%
\pgfsetlinewidth{1.003750pt}%
\definecolor{currentstroke}{rgb}{0.469640,0.132245,0.507809}%
\pgfsetstrokecolor{currentstroke}%
\pgfsetstrokeopacity{0.800000}%
\pgfsetdash{}{0pt}%
\pgfpathmoveto{\pgfqpoint{1.301989in}{1.898655in}}%
\pgfpathcurveto{\pgfqpoint{1.307813in}{1.898655in}}{\pgfqpoint{1.313400in}{1.900969in}}{\pgfqpoint{1.317518in}{1.905087in}}%
\pgfpathcurveto{\pgfqpoint{1.321636in}{1.909205in}}{\pgfqpoint{1.323950in}{1.914791in}}{\pgfqpoint{1.323950in}{1.920615in}}%
\pgfpathcurveto{\pgfqpoint{1.323950in}{1.926439in}}{\pgfqpoint{1.321636in}{1.932025in}}{\pgfqpoint{1.317518in}{1.936144in}}%
\pgfpathcurveto{\pgfqpoint{1.313400in}{1.940262in}}{\pgfqpoint{1.307813in}{1.942576in}}{\pgfqpoint{1.301989in}{1.942576in}}%
\pgfpathcurveto{\pgfqpoint{1.296165in}{1.942576in}}{\pgfqpoint{1.290579in}{1.940262in}}{\pgfqpoint{1.286461in}{1.936144in}}%
\pgfpathcurveto{\pgfqpoint{1.282343in}{1.932025in}}{\pgfqpoint{1.280029in}{1.926439in}}{\pgfqpoint{1.280029in}{1.920615in}}%
\pgfpathcurveto{\pgfqpoint{1.280029in}{1.914791in}}{\pgfqpoint{1.282343in}{1.909205in}}{\pgfqpoint{1.286461in}{1.905087in}}%
\pgfpathcurveto{\pgfqpoint{1.290579in}{1.900969in}}{\pgfqpoint{1.296165in}{1.898655in}}{\pgfqpoint{1.301989in}{1.898655in}}%
\pgfpathlineto{\pgfqpoint{1.301989in}{1.898655in}}%
\pgfpathclose%
\pgfusepath{stroke,fill}%
\end{pgfscope}%
\begin{pgfscope}%
\pgfpathrectangle{\pgfqpoint{0.100000in}{0.209042in}}{\pgfqpoint{2.906250in}{2.906250in}}%
\pgfusepath{clip}%
\pgfsetbuttcap%
\pgfsetroundjoin%
\definecolor{currentfill}{rgb}{0.291366,0.064553,0.475462}%
\pgfsetfillcolor{currentfill}%
\pgfsetfillopacity{0.800000}%
\pgfsetlinewidth{1.003750pt}%
\definecolor{currentstroke}{rgb}{0.291366,0.064553,0.475462}%
\pgfsetstrokecolor{currentstroke}%
\pgfsetstrokeopacity{0.800000}%
\pgfsetdash{}{0pt}%
\pgfpathmoveto{\pgfqpoint{1.361829in}{1.643434in}}%
\pgfpathcurveto{\pgfqpoint{1.367652in}{1.643434in}}{\pgfqpoint{1.373239in}{1.645748in}}{\pgfqpoint{1.377357in}{1.649866in}}%
\pgfpathcurveto{\pgfqpoint{1.381475in}{1.653984in}}{\pgfqpoint{1.383789in}{1.659570in}}{\pgfqpoint{1.383789in}{1.665394in}}%
\pgfpathcurveto{\pgfqpoint{1.383789in}{1.671218in}}{\pgfqpoint{1.381475in}{1.676804in}}{\pgfqpoint{1.377357in}{1.680922in}}%
\pgfpathcurveto{\pgfqpoint{1.373239in}{1.685041in}}{\pgfqpoint{1.367652in}{1.687354in}}{\pgfqpoint{1.361829in}{1.687354in}}%
\pgfpathcurveto{\pgfqpoint{1.356005in}{1.687354in}}{\pgfqpoint{1.350418in}{1.685041in}}{\pgfqpoint{1.346300in}{1.680922in}}%
\pgfpathcurveto{\pgfqpoint{1.342182in}{1.676804in}}{\pgfqpoint{1.339868in}{1.671218in}}{\pgfqpoint{1.339868in}{1.665394in}}%
\pgfpathcurveto{\pgfqpoint{1.339868in}{1.659570in}}{\pgfqpoint{1.342182in}{1.653984in}}{\pgfqpoint{1.346300in}{1.649866in}}%
\pgfpathcurveto{\pgfqpoint{1.350418in}{1.645748in}}{\pgfqpoint{1.356005in}{1.643434in}}{\pgfqpoint{1.361829in}{1.643434in}}%
\pgfpathlineto{\pgfqpoint{1.361829in}{1.643434in}}%
\pgfpathclose%
\pgfusepath{stroke,fill}%
\end{pgfscope}%
\begin{pgfscope}%
\pgfpathrectangle{\pgfqpoint{0.100000in}{0.209042in}}{\pgfqpoint{2.906250in}{2.906250in}}%
\pgfusepath{clip}%
\pgfsetbuttcap%
\pgfsetroundjoin%
\definecolor{currentfill}{rgb}{0.632805,0.187893,0.495332}%
\pgfsetfillcolor{currentfill}%
\pgfsetfillopacity{0.800000}%
\pgfsetlinewidth{1.003750pt}%
\definecolor{currentstroke}{rgb}{0.632805,0.187893,0.495332}%
\pgfsetstrokecolor{currentstroke}%
\pgfsetstrokeopacity{0.800000}%
\pgfsetdash{}{0pt}%
\pgfpathmoveto{\pgfqpoint{1.121466in}{1.732997in}}%
\pgfpathcurveto{\pgfqpoint{1.127290in}{1.732997in}}{\pgfqpoint{1.132876in}{1.735311in}}{\pgfqpoint{1.136995in}{1.739429in}}%
\pgfpathcurveto{\pgfqpoint{1.141113in}{1.743548in}}{\pgfqpoint{1.143427in}{1.749134in}}{\pgfqpoint{1.143427in}{1.754958in}}%
\pgfpathcurveto{\pgfqpoint{1.143427in}{1.760782in}}{\pgfqpoint{1.141113in}{1.766368in}}{\pgfqpoint{1.136995in}{1.770486in}}%
\pgfpathcurveto{\pgfqpoint{1.132876in}{1.774604in}}{\pgfqpoint{1.127290in}{1.776918in}}{\pgfqpoint{1.121466in}{1.776918in}}%
\pgfpathcurveto{\pgfqpoint{1.115642in}{1.776918in}}{\pgfqpoint{1.110056in}{1.774604in}}{\pgfqpoint{1.105938in}{1.770486in}}%
\pgfpathcurveto{\pgfqpoint{1.101820in}{1.766368in}}{\pgfqpoint{1.099506in}{1.760782in}}{\pgfqpoint{1.099506in}{1.754958in}}%
\pgfpathcurveto{\pgfqpoint{1.099506in}{1.749134in}}{\pgfqpoint{1.101820in}{1.743548in}}{\pgfqpoint{1.105938in}{1.739429in}}%
\pgfpathcurveto{\pgfqpoint{1.110056in}{1.735311in}}{\pgfqpoint{1.115642in}{1.732997in}}{\pgfqpoint{1.121466in}{1.732997in}}%
\pgfpathlineto{\pgfqpoint{1.121466in}{1.732997in}}%
\pgfpathclose%
\pgfusepath{stroke,fill}%
\end{pgfscope}%
\begin{pgfscope}%
\pgfpathrectangle{\pgfqpoint{0.100000in}{0.209042in}}{\pgfqpoint{2.906250in}{2.906250in}}%
\pgfusepath{clip}%
\pgfsetbuttcap%
\pgfsetroundjoin%
\definecolor{currentfill}{rgb}{0.451271,0.125132,0.507198}%
\pgfsetfillcolor{currentfill}%
\pgfsetfillopacity{0.800000}%
\pgfsetlinewidth{1.003750pt}%
\definecolor{currentstroke}{rgb}{0.451271,0.125132,0.507198}%
\pgfsetstrokecolor{currentstroke}%
\pgfsetstrokeopacity{0.800000}%
\pgfsetdash{}{0pt}%
\pgfpathmoveto{\pgfqpoint{1.600894in}{2.017391in}}%
\pgfpathcurveto{\pgfqpoint{1.606718in}{2.017391in}}{\pgfqpoint{1.612304in}{2.019704in}}{\pgfqpoint{1.616422in}{2.023823in}}%
\pgfpathcurveto{\pgfqpoint{1.620541in}{2.027941in}}{\pgfqpoint{1.622854in}{2.033527in}}{\pgfqpoint{1.622854in}{2.039351in}}%
\pgfpathcurveto{\pgfqpoint{1.622854in}{2.045175in}}{\pgfqpoint{1.620541in}{2.050761in}}{\pgfqpoint{1.616422in}{2.054879in}}%
\pgfpathcurveto{\pgfqpoint{1.612304in}{2.058997in}}{\pgfqpoint{1.606718in}{2.061311in}}{\pgfqpoint{1.600894in}{2.061311in}}%
\pgfpathcurveto{\pgfqpoint{1.595070in}{2.061311in}}{\pgfqpoint{1.589484in}{2.058997in}}{\pgfqpoint{1.585366in}{2.054879in}}%
\pgfpathcurveto{\pgfqpoint{1.581248in}{2.050761in}}{\pgfqpoint{1.578934in}{2.045175in}}{\pgfqpoint{1.578934in}{2.039351in}}%
\pgfpathcurveto{\pgfqpoint{1.578934in}{2.033527in}}{\pgfqpoint{1.581248in}{2.027941in}}{\pgfqpoint{1.585366in}{2.023823in}}%
\pgfpathcurveto{\pgfqpoint{1.589484in}{2.019704in}}{\pgfqpoint{1.595070in}{2.017391in}}{\pgfqpoint{1.600894in}{2.017391in}}%
\pgfpathlineto{\pgfqpoint{1.600894in}{2.017391in}}%
\pgfpathclose%
\pgfusepath{stroke,fill}%
\end{pgfscope}%
\begin{pgfscope}%
\pgfpathrectangle{\pgfqpoint{0.100000in}{0.209042in}}{\pgfqpoint{2.906250in}{2.906250in}}%
\pgfusepath{clip}%
\pgfsetbuttcap%
\pgfsetroundjoin%
\definecolor{currentfill}{rgb}{0.588158,0.173652,0.502035}%
\pgfsetfillcolor{currentfill}%
\pgfsetfillopacity{0.800000}%
\pgfsetlinewidth{1.003750pt}%
\definecolor{currentstroke}{rgb}{0.588158,0.173652,0.502035}%
\pgfsetstrokecolor{currentstroke}%
\pgfsetstrokeopacity{0.800000}%
\pgfsetdash{}{0pt}%
\pgfpathmoveto{\pgfqpoint{1.877153in}{2.027877in}}%
\pgfpathcurveto{\pgfqpoint{1.882977in}{2.027877in}}{\pgfqpoint{1.888564in}{2.030191in}}{\pgfqpoint{1.892682in}{2.034309in}}%
\pgfpathcurveto{\pgfqpoint{1.896800in}{2.038428in}}{\pgfqpoint{1.899114in}{2.044014in}}{\pgfqpoint{1.899114in}{2.049838in}}%
\pgfpathcurveto{\pgfqpoint{1.899114in}{2.055662in}}{\pgfqpoint{1.896800in}{2.061248in}}{\pgfqpoint{1.892682in}{2.065366in}}%
\pgfpathcurveto{\pgfqpoint{1.888564in}{2.069484in}}{\pgfqpoint{1.882977in}{2.071798in}}{\pgfqpoint{1.877153in}{2.071798in}}%
\pgfpathcurveto{\pgfqpoint{1.871330in}{2.071798in}}{\pgfqpoint{1.865743in}{2.069484in}}{\pgfqpoint{1.861625in}{2.065366in}}%
\pgfpathcurveto{\pgfqpoint{1.857507in}{2.061248in}}{\pgfqpoint{1.855193in}{2.055662in}}{\pgfqpoint{1.855193in}{2.049838in}}%
\pgfpathcurveto{\pgfqpoint{1.855193in}{2.044014in}}{\pgfqpoint{1.857507in}{2.038428in}}{\pgfqpoint{1.861625in}{2.034309in}}%
\pgfpathcurveto{\pgfqpoint{1.865743in}{2.030191in}}{\pgfqpoint{1.871330in}{2.027877in}}{\pgfqpoint{1.877153in}{2.027877in}}%
\pgfpathlineto{\pgfqpoint{1.877153in}{2.027877in}}%
\pgfpathclose%
\pgfusepath{stroke,fill}%
\end{pgfscope}%
\begin{pgfscope}%
\pgfpathrectangle{\pgfqpoint{0.100000in}{0.209042in}}{\pgfqpoint{2.906250in}{2.906250in}}%
\pgfusepath{clip}%
\pgfsetbuttcap%
\pgfsetroundjoin%
\definecolor{currentfill}{rgb}{0.594508,0.175701,0.501241}%
\pgfsetfillcolor{currentfill}%
\pgfsetfillopacity{0.800000}%
\pgfsetlinewidth{1.003750pt}%
\definecolor{currentstroke}{rgb}{0.594508,0.175701,0.501241}%
\pgfsetstrokecolor{currentstroke}%
\pgfsetstrokeopacity{0.800000}%
\pgfsetdash{}{0pt}%
\pgfpathmoveto{\pgfqpoint{2.089809in}{1.523511in}}%
\pgfpathcurveto{\pgfqpoint{2.095633in}{1.523511in}}{\pgfqpoint{2.101219in}{1.525825in}}{\pgfqpoint{2.105337in}{1.529943in}}%
\pgfpathcurveto{\pgfqpoint{2.109455in}{1.534061in}}{\pgfqpoint{2.111769in}{1.539647in}}{\pgfqpoint{2.111769in}{1.545471in}}%
\pgfpathcurveto{\pgfqpoint{2.111769in}{1.551295in}}{\pgfqpoint{2.109455in}{1.556881in}}{\pgfqpoint{2.105337in}{1.561000in}}%
\pgfpathcurveto{\pgfqpoint{2.101219in}{1.565118in}}{\pgfqpoint{2.095633in}{1.567432in}}{\pgfqpoint{2.089809in}{1.567432in}}%
\pgfpathcurveto{\pgfqpoint{2.083985in}{1.567432in}}{\pgfqpoint{2.078399in}{1.565118in}}{\pgfqpoint{2.074280in}{1.561000in}}%
\pgfpathcurveto{\pgfqpoint{2.070162in}{1.556881in}}{\pgfqpoint{2.067848in}{1.551295in}}{\pgfqpoint{2.067848in}{1.545471in}}%
\pgfpathcurveto{\pgfqpoint{2.067848in}{1.539647in}}{\pgfqpoint{2.070162in}{1.534061in}}{\pgfqpoint{2.074280in}{1.529943in}}%
\pgfpathcurveto{\pgfqpoint{2.078399in}{1.525825in}}{\pgfqpoint{2.083985in}{1.523511in}}{\pgfqpoint{2.089809in}{1.523511in}}%
\pgfpathlineto{\pgfqpoint{2.089809in}{1.523511in}}%
\pgfpathclose%
\pgfusepath{stroke,fill}%
\end{pgfscope}%
\begin{pgfscope}%
\pgfpathrectangle{\pgfqpoint{0.100000in}{0.209042in}}{\pgfqpoint{2.906250in}{2.906250in}}%
\pgfusepath{clip}%
\pgfsetbuttcap%
\pgfsetroundjoin%
\definecolor{currentfill}{rgb}{0.146785,0.068738,0.334011}%
\pgfsetfillcolor{currentfill}%
\pgfsetfillopacity{0.800000}%
\pgfsetlinewidth{1.003750pt}%
\definecolor{currentstroke}{rgb}{0.146785,0.068738,0.334011}%
\pgfsetstrokecolor{currentstroke}%
\pgfsetstrokeopacity{0.800000}%
\pgfsetdash{}{0pt}%
\pgfpathmoveto{\pgfqpoint{1.815739in}{1.632500in}}%
\pgfpathcurveto{\pgfqpoint{1.821563in}{1.632500in}}{\pgfqpoint{1.827149in}{1.634813in}}{\pgfqpoint{1.831267in}{1.638932in}}%
\pgfpathcurveto{\pgfqpoint{1.835385in}{1.643050in}}{\pgfqpoint{1.837699in}{1.648636in}}{\pgfqpoint{1.837699in}{1.654460in}}%
\pgfpathcurveto{\pgfqpoint{1.837699in}{1.660284in}}{\pgfqpoint{1.835385in}{1.665870in}}{\pgfqpoint{1.831267in}{1.669988in}}%
\pgfpathcurveto{\pgfqpoint{1.827149in}{1.674106in}}{\pgfqpoint{1.821563in}{1.676420in}}{\pgfqpoint{1.815739in}{1.676420in}}%
\pgfpathcurveto{\pgfqpoint{1.809915in}{1.676420in}}{\pgfqpoint{1.804329in}{1.674106in}}{\pgfqpoint{1.800211in}{1.669988in}}%
\pgfpathcurveto{\pgfqpoint{1.796092in}{1.665870in}}{\pgfqpoint{1.793779in}{1.660284in}}{\pgfqpoint{1.793779in}{1.654460in}}%
\pgfpathcurveto{\pgfqpoint{1.793779in}{1.648636in}}{\pgfqpoint{1.796092in}{1.643050in}}{\pgfqpoint{1.800211in}{1.638932in}}%
\pgfpathcurveto{\pgfqpoint{1.804329in}{1.634813in}}{\pgfqpoint{1.809915in}{1.632500in}}{\pgfqpoint{1.815739in}{1.632500in}}%
\pgfpathlineto{\pgfqpoint{1.815739in}{1.632500in}}%
\pgfpathclose%
\pgfusepath{stroke,fill}%
\end{pgfscope}%
\begin{pgfscope}%
\pgfpathrectangle{\pgfqpoint{0.100000in}{0.209042in}}{\pgfqpoint{2.906250in}{2.906250in}}%
\pgfusepath{clip}%
\pgfsetbuttcap%
\pgfsetroundjoin%
\definecolor{currentfill}{rgb}{0.519045,0.150383,0.507443}%
\pgfsetfillcolor{currentfill}%
\pgfsetfillopacity{0.800000}%
\pgfsetlinewidth{1.003750pt}%
\definecolor{currentstroke}{rgb}{0.519045,0.150383,0.507443}%
\pgfsetstrokecolor{currentstroke}%
\pgfsetstrokeopacity{0.800000}%
\pgfsetdash{}{0pt}%
\pgfpathmoveto{\pgfqpoint{1.467309in}{1.374491in}}%
\pgfpathcurveto{\pgfqpoint{1.473133in}{1.374491in}}{\pgfqpoint{1.478719in}{1.376804in}}{\pgfqpoint{1.482837in}{1.380923in}}%
\pgfpathcurveto{\pgfqpoint{1.486955in}{1.385041in}}{\pgfqpoint{1.489269in}{1.390627in}}{\pgfqpoint{1.489269in}{1.396451in}}%
\pgfpathcurveto{\pgfqpoint{1.489269in}{1.402275in}}{\pgfqpoint{1.486955in}{1.407861in}}{\pgfqpoint{1.482837in}{1.411979in}}%
\pgfpathcurveto{\pgfqpoint{1.478719in}{1.416097in}}{\pgfqpoint{1.473133in}{1.418411in}}{\pgfqpoint{1.467309in}{1.418411in}}%
\pgfpathcurveto{\pgfqpoint{1.461485in}{1.418411in}}{\pgfqpoint{1.455899in}{1.416097in}}{\pgfqpoint{1.451781in}{1.411979in}}%
\pgfpathcurveto{\pgfqpoint{1.447662in}{1.407861in}}{\pgfqpoint{1.445349in}{1.402275in}}{\pgfqpoint{1.445349in}{1.396451in}}%
\pgfpathcurveto{\pgfqpoint{1.445349in}{1.390627in}}{\pgfqpoint{1.447662in}{1.385041in}}{\pgfqpoint{1.451781in}{1.380923in}}%
\pgfpathcurveto{\pgfqpoint{1.455899in}{1.376804in}}{\pgfqpoint{1.461485in}{1.374491in}}{\pgfqpoint{1.467309in}{1.374491in}}%
\pgfpathlineto{\pgfqpoint{1.467309in}{1.374491in}}%
\pgfpathclose%
\pgfusepath{stroke,fill}%
\end{pgfscope}%
\begin{pgfscope}%
\pgfpathrectangle{\pgfqpoint{0.100000in}{0.209042in}}{\pgfqpoint{2.906250in}{2.906250in}}%
\pgfusepath{clip}%
\pgfsetbuttcap%
\pgfsetroundjoin%
\definecolor{currentfill}{rgb}{0.102815,0.063010,0.257854}%
\pgfsetfillcolor{currentfill}%
\pgfsetfillopacity{0.800000}%
\pgfsetlinewidth{1.003750pt}%
\definecolor{currentstroke}{rgb}{0.102815,0.063010,0.257854}%
\pgfsetstrokecolor{currentstroke}%
\pgfsetstrokeopacity{0.800000}%
\pgfsetdash{}{0pt}%
\pgfpathmoveto{\pgfqpoint{1.537276in}{1.607161in}}%
\pgfpathcurveto{\pgfqpoint{1.543100in}{1.607161in}}{\pgfqpoint{1.548686in}{1.609475in}}{\pgfqpoint{1.552804in}{1.613593in}}%
\pgfpathcurveto{\pgfqpoint{1.556922in}{1.617712in}}{\pgfqpoint{1.559236in}{1.623298in}}{\pgfqpoint{1.559236in}{1.629122in}}%
\pgfpathcurveto{\pgfqpoint{1.559236in}{1.634946in}}{\pgfqpoint{1.556922in}{1.640532in}}{\pgfqpoint{1.552804in}{1.644650in}}%
\pgfpathcurveto{\pgfqpoint{1.548686in}{1.648768in}}{\pgfqpoint{1.543100in}{1.651082in}}{\pgfqpoint{1.537276in}{1.651082in}}%
\pgfpathcurveto{\pgfqpoint{1.531452in}{1.651082in}}{\pgfqpoint{1.525866in}{1.648768in}}{\pgfqpoint{1.521748in}{1.644650in}}%
\pgfpathcurveto{\pgfqpoint{1.517629in}{1.640532in}}{\pgfqpoint{1.515316in}{1.634946in}}{\pgfqpoint{1.515316in}{1.629122in}}%
\pgfpathcurveto{\pgfqpoint{1.515316in}{1.623298in}}{\pgfqpoint{1.517629in}{1.617712in}}{\pgfqpoint{1.521748in}{1.613593in}}%
\pgfpathcurveto{\pgfqpoint{1.525866in}{1.609475in}}{\pgfqpoint{1.531452in}{1.607161in}}{\pgfqpoint{1.537276in}{1.607161in}}%
\pgfpathlineto{\pgfqpoint{1.537276in}{1.607161in}}%
\pgfpathclose%
\pgfusepath{stroke,fill}%
\end{pgfscope}%
\begin{pgfscope}%
\pgfpathrectangle{\pgfqpoint{0.100000in}{0.209042in}}{\pgfqpoint{2.906250in}{2.906250in}}%
\pgfusepath{clip}%
\pgfsetbuttcap%
\pgfsetroundjoin%
\definecolor{currentfill}{rgb}{0.018815,0.016026,0.084584}%
\pgfsetfillcolor{currentfill}%
\pgfsetfillopacity{0.800000}%
\pgfsetlinewidth{1.003750pt}%
\definecolor{currentstroke}{rgb}{0.018815,0.016026,0.084584}%
\pgfsetstrokecolor{currentstroke}%
\pgfsetstrokeopacity{0.800000}%
\pgfsetdash{}{0pt}%
\pgfpathmoveto{\pgfqpoint{1.659688in}{1.675813in}}%
\pgfpathcurveto{\pgfqpoint{1.665512in}{1.675813in}}{\pgfqpoint{1.671098in}{1.678127in}}{\pgfqpoint{1.675216in}{1.682245in}}%
\pgfpathcurveto{\pgfqpoint{1.679334in}{1.686363in}}{\pgfqpoint{1.681648in}{1.691949in}}{\pgfqpoint{1.681648in}{1.697773in}}%
\pgfpathcurveto{\pgfqpoint{1.681648in}{1.703597in}}{\pgfqpoint{1.679334in}{1.709183in}}{\pgfqpoint{1.675216in}{1.713302in}}%
\pgfpathcurveto{\pgfqpoint{1.671098in}{1.717420in}}{\pgfqpoint{1.665512in}{1.719734in}}{\pgfqpoint{1.659688in}{1.719734in}}%
\pgfpathcurveto{\pgfqpoint{1.653864in}{1.719734in}}{\pgfqpoint{1.648278in}{1.717420in}}{\pgfqpoint{1.644159in}{1.713302in}}%
\pgfpathcurveto{\pgfqpoint{1.640041in}{1.709183in}}{\pgfqpoint{1.637727in}{1.703597in}}{\pgfqpoint{1.637727in}{1.697773in}}%
\pgfpathcurveto{\pgfqpoint{1.637727in}{1.691949in}}{\pgfqpoint{1.640041in}{1.686363in}}{\pgfqpoint{1.644159in}{1.682245in}}%
\pgfpathcurveto{\pgfqpoint{1.648278in}{1.678127in}}{\pgfqpoint{1.653864in}{1.675813in}}{\pgfqpoint{1.659688in}{1.675813in}}%
\pgfpathlineto{\pgfqpoint{1.659688in}{1.675813in}}%
\pgfpathclose%
\pgfusepath{stroke,fill}%
\end{pgfscope}%
\begin{pgfscope}%
\pgfpathrectangle{\pgfqpoint{0.100000in}{0.209042in}}{\pgfqpoint{2.906250in}{2.906250in}}%
\pgfusepath{clip}%
\pgfsetbuttcap%
\pgfsetroundjoin%
\definecolor{currentfill}{rgb}{0.359898,0.087831,0.496778}%
\pgfsetfillcolor{currentfill}%
\pgfsetfillopacity{0.800000}%
\pgfsetlinewidth{1.003750pt}%
\definecolor{currentstroke}{rgb}{0.359898,0.087831,0.496778}%
\pgfsetstrokecolor{currentstroke}%
\pgfsetstrokeopacity{0.800000}%
\pgfsetdash{}{0pt}%
\pgfpathmoveto{\pgfqpoint{1.432673in}{1.494358in}}%
\pgfpathcurveto{\pgfqpoint{1.438497in}{1.494358in}}{\pgfqpoint{1.444083in}{1.496672in}}{\pgfqpoint{1.448201in}{1.500790in}}%
\pgfpathcurveto{\pgfqpoint{1.452320in}{1.504908in}}{\pgfqpoint{1.454633in}{1.510495in}}{\pgfqpoint{1.454633in}{1.516318in}}%
\pgfpathcurveto{\pgfqpoint{1.454633in}{1.522142in}}{\pgfqpoint{1.452320in}{1.527729in}}{\pgfqpoint{1.448201in}{1.531847in}}%
\pgfpathcurveto{\pgfqpoint{1.444083in}{1.535965in}}{\pgfqpoint{1.438497in}{1.538279in}}{\pgfqpoint{1.432673in}{1.538279in}}%
\pgfpathcurveto{\pgfqpoint{1.426849in}{1.538279in}}{\pgfqpoint{1.421263in}{1.535965in}}{\pgfqpoint{1.417145in}{1.531847in}}%
\pgfpathcurveto{\pgfqpoint{1.413027in}{1.527729in}}{\pgfqpoint{1.410713in}{1.522142in}}{\pgfqpoint{1.410713in}{1.516318in}}%
\pgfpathcurveto{\pgfqpoint{1.410713in}{1.510495in}}{\pgfqpoint{1.413027in}{1.504908in}}{\pgfqpoint{1.417145in}{1.500790in}}%
\pgfpathcurveto{\pgfqpoint{1.421263in}{1.496672in}}{\pgfqpoint{1.426849in}{1.494358in}}{\pgfqpoint{1.432673in}{1.494358in}}%
\pgfpathlineto{\pgfqpoint{1.432673in}{1.494358in}}%
\pgfpathclose%
\pgfusepath{stroke,fill}%
\end{pgfscope}%
\begin{pgfscope}%
\pgfpathrectangle{\pgfqpoint{0.100000in}{0.209042in}}{\pgfqpoint{2.906250in}{2.906250in}}%
\pgfusepath{clip}%
\pgfsetbuttcap%
\pgfsetroundjoin%
\definecolor{currentfill}{rgb}{0.159018,0.068354,0.352688}%
\pgfsetfillcolor{currentfill}%
\pgfsetfillopacity{0.800000}%
\pgfsetlinewidth{1.003750pt}%
\definecolor{currentstroke}{rgb}{0.159018,0.068354,0.352688}%
\pgfsetstrokecolor{currentstroke}%
\pgfsetstrokeopacity{0.800000}%
\pgfsetdash{}{0pt}%
\pgfpathmoveto{\pgfqpoint{1.727007in}{1.548471in}}%
\pgfpathcurveto{\pgfqpoint{1.732831in}{1.548471in}}{\pgfqpoint{1.738417in}{1.550785in}}{\pgfqpoint{1.742535in}{1.554903in}}%
\pgfpathcurveto{\pgfqpoint{1.746653in}{1.559021in}}{\pgfqpoint{1.748967in}{1.564607in}}{\pgfqpoint{1.748967in}{1.570431in}}%
\pgfpathcurveto{\pgfqpoint{1.748967in}{1.576255in}}{\pgfqpoint{1.746653in}{1.581841in}}{\pgfqpoint{1.742535in}{1.585960in}}%
\pgfpathcurveto{\pgfqpoint{1.738417in}{1.590078in}}{\pgfqpoint{1.732831in}{1.592392in}}{\pgfqpoint{1.727007in}{1.592392in}}%
\pgfpathcurveto{\pgfqpoint{1.721183in}{1.592392in}}{\pgfqpoint{1.715597in}{1.590078in}}{\pgfqpoint{1.711479in}{1.585960in}}%
\pgfpathcurveto{\pgfqpoint{1.707361in}{1.581841in}}{\pgfqpoint{1.705047in}{1.576255in}}{\pgfqpoint{1.705047in}{1.570431in}}%
\pgfpathcurveto{\pgfqpoint{1.705047in}{1.564607in}}{\pgfqpoint{1.707361in}{1.559021in}}{\pgfqpoint{1.711479in}{1.554903in}}%
\pgfpathcurveto{\pgfqpoint{1.715597in}{1.550785in}}{\pgfqpoint{1.721183in}{1.548471in}}{\pgfqpoint{1.727007in}{1.548471in}}%
\pgfpathlineto{\pgfqpoint{1.727007in}{1.548471in}}%
\pgfpathclose%
\pgfusepath{stroke,fill}%
\end{pgfscope}%
\begin{pgfscope}%
\pgfpathrectangle{\pgfqpoint{0.100000in}{0.209042in}}{\pgfqpoint{2.906250in}{2.906250in}}%
\pgfusepath{clip}%
\pgfsetbuttcap%
\pgfsetroundjoin%
\definecolor{currentfill}{rgb}{0.048062,0.036607,0.150327}%
\pgfsetfillcolor{currentfill}%
\pgfsetfillopacity{0.800000}%
\pgfsetlinewidth{1.003750pt}%
\definecolor{currentstroke}{rgb}{0.048062,0.036607,0.150327}%
\pgfsetstrokecolor{currentstroke}%
\pgfsetstrokeopacity{0.800000}%
\pgfsetdash{}{0pt}%
\pgfpathmoveto{\pgfqpoint{1.569596in}{1.734382in}}%
\pgfpathcurveto{\pgfqpoint{1.575420in}{1.734382in}}{\pgfqpoint{1.581006in}{1.736696in}}{\pgfqpoint{1.585124in}{1.740814in}}%
\pgfpathcurveto{\pgfqpoint{1.589242in}{1.744932in}}{\pgfqpoint{1.591556in}{1.750518in}}{\pgfqpoint{1.591556in}{1.756342in}}%
\pgfpathcurveto{\pgfqpoint{1.591556in}{1.762166in}}{\pgfqpoint{1.589242in}{1.767752in}}{\pgfqpoint{1.585124in}{1.771870in}}%
\pgfpathcurveto{\pgfqpoint{1.581006in}{1.775988in}}{\pgfqpoint{1.575420in}{1.778302in}}{\pgfqpoint{1.569596in}{1.778302in}}%
\pgfpathcurveto{\pgfqpoint{1.563772in}{1.778302in}}{\pgfqpoint{1.558186in}{1.775988in}}{\pgfqpoint{1.554068in}{1.771870in}}%
\pgfpathcurveto{\pgfqpoint{1.549950in}{1.767752in}}{\pgfqpoint{1.547636in}{1.762166in}}{\pgfqpoint{1.547636in}{1.756342in}}%
\pgfpathcurveto{\pgfqpoint{1.547636in}{1.750518in}}{\pgfqpoint{1.549950in}{1.744932in}}{\pgfqpoint{1.554068in}{1.740814in}}%
\pgfpathcurveto{\pgfqpoint{1.558186in}{1.736696in}}{\pgfqpoint{1.563772in}{1.734382in}}{\pgfqpoint{1.569596in}{1.734382in}}%
\pgfpathlineto{\pgfqpoint{1.569596in}{1.734382in}}%
\pgfpathclose%
\pgfusepath{stroke,fill}%
\end{pgfscope}%
\begin{pgfscope}%
\pgfpathrectangle{\pgfqpoint{0.100000in}{0.209042in}}{\pgfqpoint{2.906250in}{2.906250in}}%
\pgfusepath{clip}%
\pgfsetbuttcap%
\pgfsetroundjoin%
\definecolor{currentfill}{rgb}{0.594508,0.175701,0.501241}%
\pgfsetfillcolor{currentfill}%
\pgfsetfillopacity{0.800000}%
\pgfsetlinewidth{1.003750pt}%
\definecolor{currentstroke}{rgb}{0.594508,0.175701,0.501241}%
\pgfsetstrokecolor{currentstroke}%
\pgfsetstrokeopacity{0.800000}%
\pgfsetdash{}{0pt}%
\pgfpathmoveto{\pgfqpoint{1.445051in}{2.068278in}}%
\pgfpathcurveto{\pgfqpoint{1.450875in}{2.068278in}}{\pgfqpoint{1.456461in}{2.070592in}}{\pgfqpoint{1.460579in}{2.074710in}}%
\pgfpathcurveto{\pgfqpoint{1.464697in}{2.078828in}}{\pgfqpoint{1.467011in}{2.084414in}}{\pgfqpoint{1.467011in}{2.090238in}}%
\pgfpathcurveto{\pgfqpoint{1.467011in}{2.096062in}}{\pgfqpoint{1.464697in}{2.101648in}}{\pgfqpoint{1.460579in}{2.105766in}}%
\pgfpathcurveto{\pgfqpoint{1.456461in}{2.109884in}}{\pgfqpoint{1.450875in}{2.112198in}}{\pgfqpoint{1.445051in}{2.112198in}}%
\pgfpathcurveto{\pgfqpoint{1.439227in}{2.112198in}}{\pgfqpoint{1.433640in}{2.109884in}}{\pgfqpoint{1.429522in}{2.105766in}}%
\pgfpathcurveto{\pgfqpoint{1.425404in}{2.101648in}}{\pgfqpoint{1.423090in}{2.096062in}}{\pgfqpoint{1.423090in}{2.090238in}}%
\pgfpathcurveto{\pgfqpoint{1.423090in}{2.084414in}}{\pgfqpoint{1.425404in}{2.078828in}}{\pgfqpoint{1.429522in}{2.074710in}}%
\pgfpathcurveto{\pgfqpoint{1.433640in}{2.070592in}}{\pgfqpoint{1.439227in}{2.068278in}}{\pgfqpoint{1.445051in}{2.068278in}}%
\pgfpathlineto{\pgfqpoint{1.445051in}{2.068278in}}%
\pgfpathclose%
\pgfusepath{stroke,fill}%
\end{pgfscope}%
\begin{pgfscope}%
\pgfpathrectangle{\pgfqpoint{0.100000in}{0.209042in}}{\pgfqpoint{2.906250in}{2.906250in}}%
\pgfusepath{clip}%
\pgfsetbuttcap%
\pgfsetroundjoin%
\definecolor{currentfill}{rgb}{0.065330,0.047318,0.184892}%
\pgfsetfillcolor{currentfill}%
\pgfsetfillopacity{0.800000}%
\pgfsetlinewidth{1.003750pt}%
\definecolor{currentstroke}{rgb}{0.065330,0.047318,0.184892}%
\pgfsetstrokecolor{currentstroke}%
\pgfsetstrokeopacity{0.800000}%
\pgfsetdash{}{0pt}%
\pgfpathmoveto{\pgfqpoint{1.535033in}{1.689455in}}%
\pgfpathcurveto{\pgfqpoint{1.540857in}{1.689455in}}{\pgfqpoint{1.546444in}{1.691769in}}{\pgfqpoint{1.550562in}{1.695887in}}%
\pgfpathcurveto{\pgfqpoint{1.554680in}{1.700005in}}{\pgfqpoint{1.556994in}{1.705591in}}{\pgfqpoint{1.556994in}{1.711415in}}%
\pgfpathcurveto{\pgfqpoint{1.556994in}{1.717239in}}{\pgfqpoint{1.554680in}{1.722825in}}{\pgfqpoint{1.550562in}{1.726944in}}%
\pgfpathcurveto{\pgfqpoint{1.546444in}{1.731062in}}{\pgfqpoint{1.540857in}{1.733376in}}{\pgfqpoint{1.535033in}{1.733376in}}%
\pgfpathcurveto{\pgfqpoint{1.529209in}{1.733376in}}{\pgfqpoint{1.523623in}{1.731062in}}{\pgfqpoint{1.519505in}{1.726944in}}%
\pgfpathcurveto{\pgfqpoint{1.515387in}{1.722825in}}{\pgfqpoint{1.513073in}{1.717239in}}{\pgfqpoint{1.513073in}{1.711415in}}%
\pgfpathcurveto{\pgfqpoint{1.513073in}{1.705591in}}{\pgfqpoint{1.515387in}{1.700005in}}{\pgfqpoint{1.519505in}{1.695887in}}%
\pgfpathcurveto{\pgfqpoint{1.523623in}{1.691769in}}{\pgfqpoint{1.529209in}{1.689455in}}{\pgfqpoint{1.535033in}{1.689455in}}%
\pgfpathlineto{\pgfqpoint{1.535033in}{1.689455in}}%
\pgfpathclose%
\pgfusepath{stroke,fill}%
\end{pgfscope}%
\begin{pgfscope}%
\pgfpathrectangle{\pgfqpoint{0.100000in}{0.209042in}}{\pgfqpoint{2.906250in}{2.906250in}}%
\pgfusepath{clip}%
\pgfsetbuttcap%
\pgfsetroundjoin%
\definecolor{currentfill}{rgb}{0.488088,0.139186,0.508011}%
\pgfsetfillcolor{currentfill}%
\pgfsetfillopacity{0.800000}%
\pgfsetlinewidth{1.003750pt}%
\definecolor{currentstroke}{rgb}{0.488088,0.139186,0.508011}%
\pgfsetstrokecolor{currentstroke}%
\pgfsetstrokeopacity{0.800000}%
\pgfsetdash{}{0pt}%
\pgfpathmoveto{\pgfqpoint{1.764986in}{1.364257in}}%
\pgfpathcurveto{\pgfqpoint{1.770810in}{1.364257in}}{\pgfqpoint{1.776396in}{1.366570in}}{\pgfqpoint{1.780514in}{1.370689in}}%
\pgfpathcurveto{\pgfqpoint{1.784633in}{1.374807in}}{\pgfqpoint{1.786946in}{1.380393in}}{\pgfqpoint{1.786946in}{1.386217in}}%
\pgfpathcurveto{\pgfqpoint{1.786946in}{1.392041in}}{\pgfqpoint{1.784633in}{1.397627in}}{\pgfqpoint{1.780514in}{1.401745in}}%
\pgfpathcurveto{\pgfqpoint{1.776396in}{1.405863in}}{\pgfqpoint{1.770810in}{1.408177in}}{\pgfqpoint{1.764986in}{1.408177in}}%
\pgfpathcurveto{\pgfqpoint{1.759162in}{1.408177in}}{\pgfqpoint{1.753576in}{1.405863in}}{\pgfqpoint{1.749458in}{1.401745in}}%
\pgfpathcurveto{\pgfqpoint{1.745340in}{1.397627in}}{\pgfqpoint{1.743026in}{1.392041in}}{\pgfqpoint{1.743026in}{1.386217in}}%
\pgfpathcurveto{\pgfqpoint{1.743026in}{1.380393in}}{\pgfqpoint{1.745340in}{1.374807in}}{\pgfqpoint{1.749458in}{1.370689in}}%
\pgfpathcurveto{\pgfqpoint{1.753576in}{1.366570in}}{\pgfqpoint{1.759162in}{1.364257in}}{\pgfqpoint{1.764986in}{1.364257in}}%
\pgfpathlineto{\pgfqpoint{1.764986in}{1.364257in}}%
\pgfpathclose%
\pgfusepath{stroke,fill}%
\end{pgfscope}%
\begin{pgfscope}%
\pgfpathrectangle{\pgfqpoint{0.100000in}{0.209042in}}{\pgfqpoint{2.906250in}{2.906250in}}%
\pgfusepath{clip}%
\pgfsetbuttcap%
\pgfsetroundjoin%
\definecolor{currentfill}{rgb}{0.457386,0.127522,0.507448}%
\pgfsetfillcolor{currentfill}%
\pgfsetfillopacity{0.800000}%
\pgfsetlinewidth{1.003750pt}%
\definecolor{currentstroke}{rgb}{0.457386,0.127522,0.507448}%
\pgfsetstrokecolor{currentstroke}%
\pgfsetstrokeopacity{0.800000}%
\pgfsetdash{}{0pt}%
\pgfpathmoveto{\pgfqpoint{1.613364in}{1.373172in}}%
\pgfpathcurveto{\pgfqpoint{1.619188in}{1.373172in}}{\pgfqpoint{1.624774in}{1.375486in}}{\pgfqpoint{1.628892in}{1.379604in}}%
\pgfpathcurveto{\pgfqpoint{1.633010in}{1.383723in}}{\pgfqpoint{1.635324in}{1.389309in}}{\pgfqpoint{1.635324in}{1.395133in}}%
\pgfpathcurveto{\pgfqpoint{1.635324in}{1.400957in}}{\pgfqpoint{1.633010in}{1.406543in}}{\pgfqpoint{1.628892in}{1.410661in}}%
\pgfpathcurveto{\pgfqpoint{1.624774in}{1.414779in}}{\pgfqpoint{1.619188in}{1.417093in}}{\pgfqpoint{1.613364in}{1.417093in}}%
\pgfpathcurveto{\pgfqpoint{1.607540in}{1.417093in}}{\pgfqpoint{1.601954in}{1.414779in}}{\pgfqpoint{1.597835in}{1.410661in}}%
\pgfpathcurveto{\pgfqpoint{1.593717in}{1.406543in}}{\pgfqpoint{1.591403in}{1.400957in}}{\pgfqpoint{1.591403in}{1.395133in}}%
\pgfpathcurveto{\pgfqpoint{1.591403in}{1.389309in}}{\pgfqpoint{1.593717in}{1.383723in}}{\pgfqpoint{1.597835in}{1.379604in}}%
\pgfpathcurveto{\pgfqpoint{1.601954in}{1.375486in}}{\pgfqpoint{1.607540in}{1.373172in}}{\pgfqpoint{1.613364in}{1.373172in}}%
\pgfpathlineto{\pgfqpoint{1.613364in}{1.373172in}}%
\pgfpathclose%
\pgfusepath{stroke,fill}%
\end{pgfscope}%
\begin{pgfscope}%
\pgfpathrectangle{\pgfqpoint{0.100000in}{0.209042in}}{\pgfqpoint{2.906250in}{2.906250in}}%
\pgfusepath{clip}%
\pgfsetbuttcap%
\pgfsetroundjoin%
\definecolor{currentfill}{rgb}{0.556571,0.163269,0.505230}%
\pgfsetfillcolor{currentfill}%
\pgfsetfillopacity{0.800000}%
\pgfsetlinewidth{1.003750pt}%
\definecolor{currentstroke}{rgb}{0.556571,0.163269,0.505230}%
\pgfsetstrokecolor{currentstroke}%
\pgfsetstrokeopacity{0.800000}%
\pgfsetdash{}{0pt}%
\pgfpathmoveto{\pgfqpoint{1.954379in}{1.404456in}}%
\pgfpathcurveto{\pgfqpoint{1.960203in}{1.404456in}}{\pgfqpoint{1.965789in}{1.406770in}}{\pgfqpoint{1.969907in}{1.410888in}}%
\pgfpathcurveto{\pgfqpoint{1.974025in}{1.415006in}}{\pgfqpoint{1.976339in}{1.420592in}}{\pgfqpoint{1.976339in}{1.426416in}}%
\pgfpathcurveto{\pgfqpoint{1.976339in}{1.432240in}}{\pgfqpoint{1.974025in}{1.437826in}}{\pgfqpoint{1.969907in}{1.441944in}}%
\pgfpathcurveto{\pgfqpoint{1.965789in}{1.446062in}}{\pgfqpoint{1.960203in}{1.448376in}}{\pgfqpoint{1.954379in}{1.448376in}}%
\pgfpathcurveto{\pgfqpoint{1.948555in}{1.448376in}}{\pgfqpoint{1.942969in}{1.446062in}}{\pgfqpoint{1.938851in}{1.441944in}}%
\pgfpathcurveto{\pgfqpoint{1.934733in}{1.437826in}}{\pgfqpoint{1.932419in}{1.432240in}}{\pgfqpoint{1.932419in}{1.426416in}}%
\pgfpathcurveto{\pgfqpoint{1.932419in}{1.420592in}}{\pgfqpoint{1.934733in}{1.415006in}}{\pgfqpoint{1.938851in}{1.410888in}}%
\pgfpathcurveto{\pgfqpoint{1.942969in}{1.406770in}}{\pgfqpoint{1.948555in}{1.404456in}}{\pgfqpoint{1.954379in}{1.404456in}}%
\pgfpathlineto{\pgfqpoint{1.954379in}{1.404456in}}%
\pgfpathclose%
\pgfusepath{stroke,fill}%
\end{pgfscope}%
\begin{pgfscope}%
\pgfpathrectangle{\pgfqpoint{0.100000in}{0.209042in}}{\pgfqpoint{2.906250in}{2.906250in}}%
\pgfusepath{clip}%
\pgfsetbuttcap%
\pgfsetroundjoin%
\definecolor{currentfill}{rgb}{0.985693,0.528148,0.379371}%
\pgfsetfillcolor{currentfill}%
\pgfsetfillopacity{0.800000}%
\pgfsetlinewidth{1.003750pt}%
\definecolor{currentstroke}{rgb}{0.985693,0.528148,0.379371}%
\pgfsetstrokecolor{currentstroke}%
\pgfsetstrokeopacity{0.800000}%
\pgfsetdash{}{0pt}%
\pgfpathmoveto{\pgfqpoint{1.029147in}{1.271366in}}%
\pgfpathcurveto{\pgfqpoint{1.034971in}{1.271366in}}{\pgfqpoint{1.040557in}{1.273679in}}{\pgfqpoint{1.044675in}{1.277798in}}%
\pgfpathcurveto{\pgfqpoint{1.048793in}{1.281916in}}{\pgfqpoint{1.051107in}{1.287502in}}{\pgfqpoint{1.051107in}{1.293326in}}%
\pgfpathcurveto{\pgfqpoint{1.051107in}{1.299150in}}{\pgfqpoint{1.048793in}{1.304736in}}{\pgfqpoint{1.044675in}{1.308854in}}%
\pgfpathcurveto{\pgfqpoint{1.040557in}{1.312972in}}{\pgfqpoint{1.034971in}{1.315286in}}{\pgfqpoint{1.029147in}{1.315286in}}%
\pgfpathcurveto{\pgfqpoint{1.023323in}{1.315286in}}{\pgfqpoint{1.017737in}{1.312972in}}{\pgfqpoint{1.013619in}{1.308854in}}%
\pgfpathcurveto{\pgfqpoint{1.009500in}{1.304736in}}{\pgfqpoint{1.007187in}{1.299150in}}{\pgfqpoint{1.007187in}{1.293326in}}%
\pgfpathcurveto{\pgfqpoint{1.007187in}{1.287502in}}{\pgfqpoint{1.009500in}{1.281916in}}{\pgfqpoint{1.013619in}{1.277798in}}%
\pgfpathcurveto{\pgfqpoint{1.017737in}{1.273679in}}{\pgfqpoint{1.023323in}{1.271366in}}{\pgfqpoint{1.029147in}{1.271366in}}%
\pgfpathlineto{\pgfqpoint{1.029147in}{1.271366in}}%
\pgfpathclose%
\pgfusepath{stroke,fill}%
\end{pgfscope}%
\begin{pgfscope}%
\pgfpathrectangle{\pgfqpoint{0.100000in}{0.209042in}}{\pgfqpoint{2.906250in}{2.906250in}}%
\pgfusepath{clip}%
\pgfsetbuttcap%
\pgfsetroundjoin%
\definecolor{currentfill}{rgb}{0.588158,0.173652,0.502035}%
\pgfsetfillcolor{currentfill}%
\pgfsetfillopacity{0.800000}%
\pgfsetlinewidth{1.003750pt}%
\definecolor{currentstroke}{rgb}{0.588158,0.173652,0.502035}%
\pgfsetstrokecolor{currentstroke}%
\pgfsetstrokeopacity{0.800000}%
\pgfsetdash{}{0pt}%
\pgfpathmoveto{\pgfqpoint{1.640472in}{1.299900in}}%
\pgfpathcurveto{\pgfqpoint{1.646296in}{1.299900in}}{\pgfqpoint{1.651882in}{1.302214in}}{\pgfqpoint{1.656000in}{1.306332in}}%
\pgfpathcurveto{\pgfqpoint{1.660118in}{1.310450in}}{\pgfqpoint{1.662432in}{1.316036in}}{\pgfqpoint{1.662432in}{1.321860in}}%
\pgfpathcurveto{\pgfqpoint{1.662432in}{1.327684in}}{\pgfqpoint{1.660118in}{1.333270in}}{\pgfqpoint{1.656000in}{1.337389in}}%
\pgfpathcurveto{\pgfqpoint{1.651882in}{1.341507in}}{\pgfqpoint{1.646296in}{1.343821in}}{\pgfqpoint{1.640472in}{1.343821in}}%
\pgfpathcurveto{\pgfqpoint{1.634648in}{1.343821in}}{\pgfqpoint{1.629062in}{1.341507in}}{\pgfqpoint{1.624943in}{1.337389in}}%
\pgfpathcurveto{\pgfqpoint{1.620825in}{1.333270in}}{\pgfqpoint{1.618511in}{1.327684in}}{\pgfqpoint{1.618511in}{1.321860in}}%
\pgfpathcurveto{\pgfqpoint{1.618511in}{1.316036in}}{\pgfqpoint{1.620825in}{1.310450in}}{\pgfqpoint{1.624943in}{1.306332in}}%
\pgfpathcurveto{\pgfqpoint{1.629062in}{1.302214in}}{\pgfqpoint{1.634648in}{1.299900in}}{\pgfqpoint{1.640472in}{1.299900in}}%
\pgfpathlineto{\pgfqpoint{1.640472in}{1.299900in}}%
\pgfpathclose%
\pgfusepath{stroke,fill}%
\end{pgfscope}%
\begin{pgfscope}%
\pgfpathrectangle{\pgfqpoint{0.100000in}{0.209042in}}{\pgfqpoint{2.906250in}{2.906250in}}%
\pgfusepath{clip}%
\pgfsetbuttcap%
\pgfsetroundjoin%
\definecolor{currentfill}{rgb}{0.123833,0.067295,0.295879}%
\pgfsetfillcolor{currentfill}%
\pgfsetfillopacity{0.800000}%
\pgfsetlinewidth{1.003750pt}%
\definecolor{currentstroke}{rgb}{0.123833,0.067295,0.295879}%
\pgfsetstrokecolor{currentstroke}%
\pgfsetstrokeopacity{0.800000}%
\pgfsetdash{}{0pt}%
\pgfpathmoveto{\pgfqpoint{1.788544in}{1.638847in}}%
\pgfpathcurveto{\pgfqpoint{1.794368in}{1.638847in}}{\pgfqpoint{1.799954in}{1.641161in}}{\pgfqpoint{1.804072in}{1.645279in}}%
\pgfpathcurveto{\pgfqpoint{1.808190in}{1.649397in}}{\pgfqpoint{1.810504in}{1.654983in}}{\pgfqpoint{1.810504in}{1.660807in}}%
\pgfpathcurveto{\pgfqpoint{1.810504in}{1.666631in}}{\pgfqpoint{1.808190in}{1.672217in}}{\pgfqpoint{1.804072in}{1.676336in}}%
\pgfpathcurveto{\pgfqpoint{1.799954in}{1.680454in}}{\pgfqpoint{1.794368in}{1.682768in}}{\pgfqpoint{1.788544in}{1.682768in}}%
\pgfpathcurveto{\pgfqpoint{1.782720in}{1.682768in}}{\pgfqpoint{1.777134in}{1.680454in}}{\pgfqpoint{1.773016in}{1.676336in}}%
\pgfpathcurveto{\pgfqpoint{1.768898in}{1.672217in}}{\pgfqpoint{1.766584in}{1.666631in}}{\pgfqpoint{1.766584in}{1.660807in}}%
\pgfpathcurveto{\pgfqpoint{1.766584in}{1.654983in}}{\pgfqpoint{1.768898in}{1.649397in}}{\pgfqpoint{1.773016in}{1.645279in}}%
\pgfpathcurveto{\pgfqpoint{1.777134in}{1.641161in}}{\pgfqpoint{1.782720in}{1.638847in}}{\pgfqpoint{1.788544in}{1.638847in}}%
\pgfpathlineto{\pgfqpoint{1.788544in}{1.638847in}}%
\pgfpathclose%
\pgfusepath{stroke,fill}%
\end{pgfscope}%
\begin{pgfscope}%
\pgfpathrectangle{\pgfqpoint{0.100000in}{0.209042in}}{\pgfqpoint{2.906250in}{2.906250in}}%
\pgfusepath{clip}%
\pgfsetbuttcap%
\pgfsetroundjoin%
\definecolor{currentfill}{rgb}{0.329114,0.075972,0.489287}%
\pgfsetfillcolor{currentfill}%
\pgfsetfillopacity{0.800000}%
\pgfsetlinewidth{1.003750pt}%
\definecolor{currentstroke}{rgb}{0.329114,0.075972,0.489287}%
\pgfsetstrokecolor{currentstroke}%
\pgfsetstrokeopacity{0.800000}%
\pgfsetdash{}{0pt}%
\pgfpathmoveto{\pgfqpoint{1.923662in}{1.594205in}}%
\pgfpathcurveto{\pgfqpoint{1.929486in}{1.594205in}}{\pgfqpoint{1.935072in}{1.596519in}}{\pgfqpoint{1.939191in}{1.600637in}}%
\pgfpathcurveto{\pgfqpoint{1.943309in}{1.604755in}}{\pgfqpoint{1.945623in}{1.610341in}}{\pgfqpoint{1.945623in}{1.616165in}}%
\pgfpathcurveto{\pgfqpoint{1.945623in}{1.621989in}}{\pgfqpoint{1.943309in}{1.627575in}}{\pgfqpoint{1.939191in}{1.631693in}}%
\pgfpathcurveto{\pgfqpoint{1.935072in}{1.635811in}}{\pgfqpoint{1.929486in}{1.638125in}}{\pgfqpoint{1.923662in}{1.638125in}}%
\pgfpathcurveto{\pgfqpoint{1.917838in}{1.638125in}}{\pgfqpoint{1.912252in}{1.635811in}}{\pgfqpoint{1.908134in}{1.631693in}}%
\pgfpathcurveto{\pgfqpoint{1.904016in}{1.627575in}}{\pgfqpoint{1.901702in}{1.621989in}}{\pgfqpoint{1.901702in}{1.616165in}}%
\pgfpathcurveto{\pgfqpoint{1.901702in}{1.610341in}}{\pgfqpoint{1.904016in}{1.604755in}}{\pgfqpoint{1.908134in}{1.600637in}}%
\pgfpathcurveto{\pgfqpoint{1.912252in}{1.596519in}}{\pgfqpoint{1.917838in}{1.594205in}}{\pgfqpoint{1.923662in}{1.594205in}}%
\pgfpathlineto{\pgfqpoint{1.923662in}{1.594205in}}%
\pgfpathclose%
\pgfusepath{stroke,fill}%
\end{pgfscope}%
\begin{pgfscope}%
\pgfpathrectangle{\pgfqpoint{0.100000in}{0.209042in}}{\pgfqpoint{2.906250in}{2.906250in}}%
\pgfusepath{clip}%
\pgfsetbuttcap%
\pgfsetroundjoin%
\definecolor{currentfill}{rgb}{0.165308,0.067911,0.361816}%
\pgfsetfillcolor{currentfill}%
\pgfsetfillopacity{0.800000}%
\pgfsetlinewidth{1.003750pt}%
\definecolor{currentstroke}{rgb}{0.165308,0.067911,0.361816}%
\pgfsetstrokecolor{currentstroke}%
\pgfsetstrokeopacity{0.800000}%
\pgfsetdash{}{0pt}%
\pgfpathmoveto{\pgfqpoint{1.721283in}{1.831253in}}%
\pgfpathcurveto{\pgfqpoint{1.727107in}{1.831253in}}{\pgfqpoint{1.732693in}{1.833567in}}{\pgfqpoint{1.736811in}{1.837685in}}%
\pgfpathcurveto{\pgfqpoint{1.740930in}{1.841803in}}{\pgfqpoint{1.743243in}{1.847389in}}{\pgfqpoint{1.743243in}{1.853213in}}%
\pgfpathcurveto{\pgfqpoint{1.743243in}{1.859037in}}{\pgfqpoint{1.740930in}{1.864623in}}{\pgfqpoint{1.736811in}{1.868741in}}%
\pgfpathcurveto{\pgfqpoint{1.732693in}{1.872860in}}{\pgfqpoint{1.727107in}{1.875174in}}{\pgfqpoint{1.721283in}{1.875174in}}%
\pgfpathcurveto{\pgfqpoint{1.715459in}{1.875174in}}{\pgfqpoint{1.709873in}{1.872860in}}{\pgfqpoint{1.705755in}{1.868741in}}%
\pgfpathcurveto{\pgfqpoint{1.701637in}{1.864623in}}{\pgfqpoint{1.699323in}{1.859037in}}{\pgfqpoint{1.699323in}{1.853213in}}%
\pgfpathcurveto{\pgfqpoint{1.699323in}{1.847389in}}{\pgfqpoint{1.701637in}{1.841803in}}{\pgfqpoint{1.705755in}{1.837685in}}%
\pgfpathcurveto{\pgfqpoint{1.709873in}{1.833567in}}{\pgfqpoint{1.715459in}{1.831253in}}{\pgfqpoint{1.721283in}{1.831253in}}%
\pgfpathlineto{\pgfqpoint{1.721283in}{1.831253in}}%
\pgfpathclose%
\pgfusepath{stroke,fill}%
\end{pgfscope}%
\begin{pgfscope}%
\pgfpathrectangle{\pgfqpoint{0.100000in}{0.209042in}}{\pgfqpoint{2.906250in}{2.906250in}}%
\pgfusepath{clip}%
\pgfsetbuttcap%
\pgfsetroundjoin%
\definecolor{currentfill}{rgb}{0.218512,0.061158,0.425392}%
\pgfsetfillcolor{currentfill}%
\pgfsetfillopacity{0.800000}%
\pgfsetlinewidth{1.003750pt}%
\definecolor{currentstroke}{rgb}{0.218512,0.061158,0.425392}%
\pgfsetstrokecolor{currentstroke}%
\pgfsetstrokeopacity{0.800000}%
\pgfsetdash{}{0pt}%
\pgfpathmoveto{\pgfqpoint{1.860296in}{1.731267in}}%
\pgfpathcurveto{\pgfqpoint{1.866120in}{1.731267in}}{\pgfqpoint{1.871706in}{1.733581in}}{\pgfqpoint{1.875825in}{1.737699in}}%
\pgfpathcurveto{\pgfqpoint{1.879943in}{1.741817in}}{\pgfqpoint{1.882257in}{1.747404in}}{\pgfqpoint{1.882257in}{1.753227in}}%
\pgfpathcurveto{\pgfqpoint{1.882257in}{1.759051in}}{\pgfqpoint{1.879943in}{1.764638in}}{\pgfqpoint{1.875825in}{1.768756in}}%
\pgfpathcurveto{\pgfqpoint{1.871706in}{1.772874in}}{\pgfqpoint{1.866120in}{1.775188in}}{\pgfqpoint{1.860296in}{1.775188in}}%
\pgfpathcurveto{\pgfqpoint{1.854472in}{1.775188in}}{\pgfqpoint{1.848886in}{1.772874in}}{\pgfqpoint{1.844768in}{1.768756in}}%
\pgfpathcurveto{\pgfqpoint{1.840650in}{1.764638in}}{\pgfqpoint{1.838336in}{1.759051in}}{\pgfqpoint{1.838336in}{1.753227in}}%
\pgfpathcurveto{\pgfqpoint{1.838336in}{1.747404in}}{\pgfqpoint{1.840650in}{1.741817in}}{\pgfqpoint{1.844768in}{1.737699in}}%
\pgfpathcurveto{\pgfqpoint{1.848886in}{1.733581in}}{\pgfqpoint{1.854472in}{1.731267in}}{\pgfqpoint{1.860296in}{1.731267in}}%
\pgfpathlineto{\pgfqpoint{1.860296in}{1.731267in}}%
\pgfpathclose%
\pgfusepath{stroke,fill}%
\end{pgfscope}%
\begin{pgfscope}%
\pgfpathrectangle{\pgfqpoint{0.100000in}{0.209042in}}{\pgfqpoint{2.906250in}{2.906250in}}%
\pgfusepath{clip}%
\pgfsetbuttcap%
\pgfsetroundjoin%
\definecolor{currentfill}{rgb}{0.742004,0.224025,0.467018}%
\pgfsetfillcolor{currentfill}%
\pgfsetfillopacity{0.800000}%
\pgfsetlinewidth{1.003750pt}%
\definecolor{currentstroke}{rgb}{0.742004,0.224025,0.467018}%
\pgfsetstrokecolor{currentstroke}%
\pgfsetstrokeopacity{0.800000}%
\pgfsetdash{}{0pt}%
\pgfpathmoveto{\pgfqpoint{1.971296in}{1.288455in}}%
\pgfpathcurveto{\pgfqpoint{1.977120in}{1.288455in}}{\pgfqpoint{1.982706in}{1.290769in}}{\pgfqpoint{1.986825in}{1.294888in}}%
\pgfpathcurveto{\pgfqpoint{1.990943in}{1.299006in}}{\pgfqpoint{1.993257in}{1.304592in}}{\pgfqpoint{1.993257in}{1.310416in}}%
\pgfpathcurveto{\pgfqpoint{1.993257in}{1.316240in}}{\pgfqpoint{1.990943in}{1.321826in}}{\pgfqpoint{1.986825in}{1.325944in}}%
\pgfpathcurveto{\pgfqpoint{1.982706in}{1.330062in}}{\pgfqpoint{1.977120in}{1.332376in}}{\pgfqpoint{1.971296in}{1.332376in}}%
\pgfpathcurveto{\pgfqpoint{1.965472in}{1.332376in}}{\pgfqpoint{1.959886in}{1.330062in}}{\pgfqpoint{1.955768in}{1.325944in}}%
\pgfpathcurveto{\pgfqpoint{1.951650in}{1.321826in}}{\pgfqpoint{1.949336in}{1.316240in}}{\pgfqpoint{1.949336in}{1.310416in}}%
\pgfpathcurveto{\pgfqpoint{1.949336in}{1.304592in}}{\pgfqpoint{1.951650in}{1.299006in}}{\pgfqpoint{1.955768in}{1.294888in}}%
\pgfpathcurveto{\pgfqpoint{1.959886in}{1.290769in}}{\pgfqpoint{1.965472in}{1.288455in}}{\pgfqpoint{1.971296in}{1.288455in}}%
\pgfpathlineto{\pgfqpoint{1.971296in}{1.288455in}}%
\pgfpathclose%
\pgfusepath{stroke,fill}%
\end{pgfscope}%
\begin{pgfscope}%
\pgfpathrectangle{\pgfqpoint{0.100000in}{0.209042in}}{\pgfqpoint{2.906250in}{2.906250in}}%
\pgfusepath{clip}%
\pgfsetbuttcap%
\pgfsetroundjoin%
\definecolor{currentfill}{rgb}{0.525270,0.152569,0.507192}%
\pgfsetfillcolor{currentfill}%
\pgfsetfillopacity{0.800000}%
\pgfsetlinewidth{1.003750pt}%
\definecolor{currentstroke}{rgb}{0.525270,0.152569,0.507192}%
\pgfsetstrokecolor{currentstroke}%
\pgfsetstrokeopacity{0.800000}%
\pgfsetdash{}{0pt}%
\pgfpathmoveto{\pgfqpoint{1.270696in}{1.517258in}}%
\pgfpathcurveto{\pgfqpoint{1.276520in}{1.517258in}}{\pgfqpoint{1.282106in}{1.519572in}}{\pgfqpoint{1.286224in}{1.523690in}}%
\pgfpathcurveto{\pgfqpoint{1.290342in}{1.527808in}}{\pgfqpoint{1.292656in}{1.533394in}}{\pgfqpoint{1.292656in}{1.539218in}}%
\pgfpathcurveto{\pgfqpoint{1.292656in}{1.545042in}}{\pgfqpoint{1.290342in}{1.550628in}}{\pgfqpoint{1.286224in}{1.554746in}}%
\pgfpathcurveto{\pgfqpoint{1.282106in}{1.558864in}}{\pgfqpoint{1.276520in}{1.561178in}}{\pgfqpoint{1.270696in}{1.561178in}}%
\pgfpathcurveto{\pgfqpoint{1.264872in}{1.561178in}}{\pgfqpoint{1.259286in}{1.558864in}}{\pgfqpoint{1.255168in}{1.554746in}}%
\pgfpathcurveto{\pgfqpoint{1.251049in}{1.550628in}}{\pgfqpoint{1.248736in}{1.545042in}}{\pgfqpoint{1.248736in}{1.539218in}}%
\pgfpathcurveto{\pgfqpoint{1.248736in}{1.533394in}}{\pgfqpoint{1.251049in}{1.527808in}}{\pgfqpoint{1.255168in}{1.523690in}}%
\pgfpathcurveto{\pgfqpoint{1.259286in}{1.519572in}}{\pgfqpoint{1.264872in}{1.517258in}}{\pgfqpoint{1.270696in}{1.517258in}}%
\pgfpathlineto{\pgfqpoint{1.270696in}{1.517258in}}%
\pgfpathclose%
\pgfusepath{stroke,fill}%
\end{pgfscope}%
\begin{pgfscope}%
\pgfpathrectangle{\pgfqpoint{0.100000in}{0.209042in}}{\pgfqpoint{2.906250in}{2.906250in}}%
\pgfusepath{clip}%
\pgfsetbuttcap%
\pgfsetroundjoin%
\definecolor{currentfill}{rgb}{0.297740,0.066117,0.478243}%
\pgfsetfillcolor{currentfill}%
\pgfsetfillopacity{0.800000}%
\pgfsetlinewidth{1.003750pt}%
\definecolor{currentstroke}{rgb}{0.297740,0.066117,0.478243}%
\pgfsetstrokecolor{currentstroke}%
\pgfsetstrokeopacity{0.800000}%
\pgfsetdash{}{0pt}%
\pgfpathmoveto{\pgfqpoint{1.792998in}{1.874720in}}%
\pgfpathcurveto{\pgfqpoint{1.798822in}{1.874720in}}{\pgfqpoint{1.804408in}{1.877034in}}{\pgfqpoint{1.808526in}{1.881152in}}%
\pgfpathcurveto{\pgfqpoint{1.812644in}{1.885270in}}{\pgfqpoint{1.814958in}{1.890856in}}{\pgfqpoint{1.814958in}{1.896680in}}%
\pgfpathcurveto{\pgfqpoint{1.814958in}{1.902504in}}{\pgfqpoint{1.812644in}{1.908090in}}{\pgfqpoint{1.808526in}{1.912208in}}%
\pgfpathcurveto{\pgfqpoint{1.804408in}{1.916327in}}{\pgfqpoint{1.798822in}{1.918640in}}{\pgfqpoint{1.792998in}{1.918640in}}%
\pgfpathcurveto{\pgfqpoint{1.787174in}{1.918640in}}{\pgfqpoint{1.781588in}{1.916327in}}{\pgfqpoint{1.777470in}{1.912208in}}%
\pgfpathcurveto{\pgfqpoint{1.773351in}{1.908090in}}{\pgfqpoint{1.771038in}{1.902504in}}{\pgfqpoint{1.771038in}{1.896680in}}%
\pgfpathcurveto{\pgfqpoint{1.771038in}{1.890856in}}{\pgfqpoint{1.773351in}{1.885270in}}{\pgfqpoint{1.777470in}{1.881152in}}%
\pgfpathcurveto{\pgfqpoint{1.781588in}{1.877034in}}{\pgfqpoint{1.787174in}{1.874720in}}{\pgfqpoint{1.792998in}{1.874720in}}%
\pgfpathlineto{\pgfqpoint{1.792998in}{1.874720in}}%
\pgfpathclose%
\pgfusepath{stroke,fill}%
\end{pgfscope}%
\begin{pgfscope}%
\pgfpathrectangle{\pgfqpoint{0.100000in}{0.209042in}}{\pgfqpoint{2.906250in}{2.906250in}}%
\pgfusepath{clip}%
\pgfsetbuttcap%
\pgfsetroundjoin%
\definecolor{currentfill}{rgb}{0.252220,0.059415,0.453248}%
\pgfsetfillcolor{currentfill}%
\pgfsetfillopacity{0.800000}%
\pgfsetlinewidth{1.003750pt}%
\definecolor{currentstroke}{rgb}{0.252220,0.059415,0.453248}%
\pgfsetstrokecolor{currentstroke}%
\pgfsetstrokeopacity{0.800000}%
\pgfsetdash{}{0pt}%
\pgfpathmoveto{\pgfqpoint{1.823599in}{1.545059in}}%
\pgfpathcurveto{\pgfqpoint{1.829423in}{1.545059in}}{\pgfqpoint{1.835009in}{1.547373in}}{\pgfqpoint{1.839128in}{1.551491in}}%
\pgfpathcurveto{\pgfqpoint{1.843246in}{1.555609in}}{\pgfqpoint{1.845560in}{1.561195in}}{\pgfqpoint{1.845560in}{1.567019in}}%
\pgfpathcurveto{\pgfqpoint{1.845560in}{1.572843in}}{\pgfqpoint{1.843246in}{1.578429in}}{\pgfqpoint{1.839128in}{1.582548in}}%
\pgfpathcurveto{\pgfqpoint{1.835009in}{1.586666in}}{\pgfqpoint{1.829423in}{1.588980in}}{\pgfqpoint{1.823599in}{1.588980in}}%
\pgfpathcurveto{\pgfqpoint{1.817775in}{1.588980in}}{\pgfqpoint{1.812189in}{1.586666in}}{\pgfqpoint{1.808071in}{1.582548in}}%
\pgfpathcurveto{\pgfqpoint{1.803953in}{1.578429in}}{\pgfqpoint{1.801639in}{1.572843in}}{\pgfqpoint{1.801639in}{1.567019in}}%
\pgfpathcurveto{\pgfqpoint{1.801639in}{1.561195in}}{\pgfqpoint{1.803953in}{1.555609in}}{\pgfqpoint{1.808071in}{1.551491in}}%
\pgfpathcurveto{\pgfqpoint{1.812189in}{1.547373in}}{\pgfqpoint{1.817775in}{1.545059in}}{\pgfqpoint{1.823599in}{1.545059in}}%
\pgfpathlineto{\pgfqpoint{1.823599in}{1.545059in}}%
\pgfpathclose%
\pgfusepath{stroke,fill}%
\end{pgfscope}%
\begin{pgfscope}%
\pgfpathrectangle{\pgfqpoint{0.100000in}{0.209042in}}{\pgfqpoint{2.906250in}{2.906250in}}%
\pgfusepath{clip}%
\pgfsetbuttcap%
\pgfsetroundjoin%
\definecolor{currentfill}{rgb}{0.408629,0.107930,0.504052}%
\pgfsetfillcolor{currentfill}%
\pgfsetfillopacity{0.800000}%
\pgfsetlinewidth{1.003750pt}%
\definecolor{currentstroke}{rgb}{0.408629,0.107930,0.504052}%
\pgfsetstrokecolor{currentstroke}%
\pgfsetstrokeopacity{0.800000}%
\pgfsetdash{}{0pt}%
\pgfpathmoveto{\pgfqpoint{1.864979in}{1.453664in}}%
\pgfpathcurveto{\pgfqpoint{1.870803in}{1.453664in}}{\pgfqpoint{1.876389in}{1.455978in}}{\pgfqpoint{1.880507in}{1.460096in}}%
\pgfpathcurveto{\pgfqpoint{1.884625in}{1.464215in}}{\pgfqpoint{1.886939in}{1.469801in}}{\pgfqpoint{1.886939in}{1.475625in}}%
\pgfpathcurveto{\pgfqpoint{1.886939in}{1.481449in}}{\pgfqpoint{1.884625in}{1.487035in}}{\pgfqpoint{1.880507in}{1.491153in}}%
\pgfpathcurveto{\pgfqpoint{1.876389in}{1.495271in}}{\pgfqpoint{1.870803in}{1.497585in}}{\pgfqpoint{1.864979in}{1.497585in}}%
\pgfpathcurveto{\pgfqpoint{1.859155in}{1.497585in}}{\pgfqpoint{1.853569in}{1.495271in}}{\pgfqpoint{1.849451in}{1.491153in}}%
\pgfpathcurveto{\pgfqpoint{1.845333in}{1.487035in}}{\pgfqpoint{1.843019in}{1.481449in}}{\pgfqpoint{1.843019in}{1.475625in}}%
\pgfpathcurveto{\pgfqpoint{1.843019in}{1.469801in}}{\pgfqpoint{1.845333in}{1.464215in}}{\pgfqpoint{1.849451in}{1.460096in}}%
\pgfpathcurveto{\pgfqpoint{1.853569in}{1.455978in}}{\pgfqpoint{1.859155in}{1.453664in}}{\pgfqpoint{1.864979in}{1.453664in}}%
\pgfpathlineto{\pgfqpoint{1.864979in}{1.453664in}}%
\pgfpathclose%
\pgfusepath{stroke,fill}%
\end{pgfscope}%
\begin{pgfscope}%
\pgfpathrectangle{\pgfqpoint{0.100000in}{0.209042in}}{\pgfqpoint{2.906250in}{2.906250in}}%
\pgfusepath{clip}%
\pgfsetbuttcap%
\pgfsetroundjoin%
\definecolor{currentfill}{rgb}{0.088155,0.058133,0.229922}%
\pgfsetfillcolor{currentfill}%
\pgfsetfillopacity{0.800000}%
\pgfsetlinewidth{1.003750pt}%
\definecolor{currentstroke}{rgb}{0.088155,0.058133,0.229922}%
\pgfsetstrokecolor{currentstroke}%
\pgfsetstrokeopacity{0.800000}%
\pgfsetdash{}{0pt}%
\pgfpathmoveto{\pgfqpoint{1.618829in}{1.781514in}}%
\pgfpathcurveto{\pgfqpoint{1.624653in}{1.781514in}}{\pgfqpoint{1.630239in}{1.783828in}}{\pgfqpoint{1.634357in}{1.787946in}}%
\pgfpathcurveto{\pgfqpoint{1.638475in}{1.792064in}}{\pgfqpoint{1.640789in}{1.797650in}}{\pgfqpoint{1.640789in}{1.803474in}}%
\pgfpathcurveto{\pgfqpoint{1.640789in}{1.809298in}}{\pgfqpoint{1.638475in}{1.814884in}}{\pgfqpoint{1.634357in}{1.819002in}}%
\pgfpathcurveto{\pgfqpoint{1.630239in}{1.823121in}}{\pgfqpoint{1.624653in}{1.825434in}}{\pgfqpoint{1.618829in}{1.825434in}}%
\pgfpathcurveto{\pgfqpoint{1.613005in}{1.825434in}}{\pgfqpoint{1.607419in}{1.823121in}}{\pgfqpoint{1.603300in}{1.819002in}}%
\pgfpathcurveto{\pgfqpoint{1.599182in}{1.814884in}}{\pgfqpoint{1.596868in}{1.809298in}}{\pgfqpoint{1.596868in}{1.803474in}}%
\pgfpathcurveto{\pgfqpoint{1.596868in}{1.797650in}}{\pgfqpoint{1.599182in}{1.792064in}}{\pgfqpoint{1.603300in}{1.787946in}}%
\pgfpathcurveto{\pgfqpoint{1.607419in}{1.783828in}}{\pgfqpoint{1.613005in}{1.781514in}}{\pgfqpoint{1.618829in}{1.781514in}}%
\pgfpathlineto{\pgfqpoint{1.618829in}{1.781514in}}%
\pgfpathclose%
\pgfusepath{stroke,fill}%
\end{pgfscope}%
\begin{pgfscope}%
\pgfpathrectangle{\pgfqpoint{0.100000in}{0.209042in}}{\pgfqpoint{2.906250in}{2.906250in}}%
\pgfusepath{clip}%
\pgfsetbuttcap%
\pgfsetroundjoin%
\definecolor{currentfill}{rgb}{0.550287,0.161158,0.505719}%
\pgfsetfillcolor{currentfill}%
\pgfsetfillopacity{0.800000}%
\pgfsetlinewidth{1.003750pt}%
\definecolor{currentstroke}{rgb}{0.550287,0.161158,0.505719}%
\pgfsetstrokecolor{currentstroke}%
\pgfsetstrokeopacity{0.800000}%
\pgfsetdash{}{0pt}%
\pgfpathmoveto{\pgfqpoint{1.279486in}{1.932894in}}%
\pgfpathcurveto{\pgfqpoint{1.285310in}{1.932894in}}{\pgfqpoint{1.290897in}{1.935207in}}{\pgfqpoint{1.295015in}{1.939326in}}%
\pgfpathcurveto{\pgfqpoint{1.299133in}{1.943444in}}{\pgfqpoint{1.301447in}{1.949030in}}{\pgfqpoint{1.301447in}{1.954854in}}%
\pgfpathcurveto{\pgfqpoint{1.301447in}{1.960678in}}{\pgfqpoint{1.299133in}{1.966264in}}{\pgfqpoint{1.295015in}{1.970382in}}%
\pgfpathcurveto{\pgfqpoint{1.290897in}{1.974500in}}{\pgfqpoint{1.285310in}{1.976814in}}{\pgfqpoint{1.279486in}{1.976814in}}%
\pgfpathcurveto{\pgfqpoint{1.273663in}{1.976814in}}{\pgfqpoint{1.268076in}{1.974500in}}{\pgfqpoint{1.263958in}{1.970382in}}%
\pgfpathcurveto{\pgfqpoint{1.259840in}{1.966264in}}{\pgfqpoint{1.257526in}{1.960678in}}{\pgfqpoint{1.257526in}{1.954854in}}%
\pgfpathcurveto{\pgfqpoint{1.257526in}{1.949030in}}{\pgfqpoint{1.259840in}{1.943444in}}{\pgfqpoint{1.263958in}{1.939326in}}%
\pgfpathcurveto{\pgfqpoint{1.268076in}{1.935207in}}{\pgfqpoint{1.273663in}{1.932894in}}{\pgfqpoint{1.279486in}{1.932894in}}%
\pgfpathlineto{\pgfqpoint{1.279486in}{1.932894in}}%
\pgfpathclose%
\pgfusepath{stroke,fill}%
\end{pgfscope}%
\begin{pgfscope}%
\pgfpathrectangle{\pgfqpoint{0.100000in}{0.209042in}}{\pgfqpoint{2.906250in}{2.906250in}}%
\pgfusepath{clip}%
\pgfsetbuttcap%
\pgfsetroundjoin%
\definecolor{currentfill}{rgb}{0.225302,0.060445,0.431742}%
\pgfsetfillcolor{currentfill}%
\pgfsetfillopacity{0.800000}%
\pgfsetlinewidth{1.003750pt}%
\definecolor{currentstroke}{rgb}{0.225302,0.060445,0.431742}%
\pgfsetstrokecolor{currentstroke}%
\pgfsetstrokeopacity{0.800000}%
\pgfsetdash{}{0pt}%
\pgfpathmoveto{\pgfqpoint{1.431959in}{1.621636in}}%
\pgfpathcurveto{\pgfqpoint{1.437782in}{1.621636in}}{\pgfqpoint{1.443369in}{1.623950in}}{\pgfqpoint{1.447487in}{1.628068in}}%
\pgfpathcurveto{\pgfqpoint{1.451605in}{1.632187in}}{\pgfqpoint{1.453919in}{1.637773in}}{\pgfqpoint{1.453919in}{1.643597in}}%
\pgfpathcurveto{\pgfqpoint{1.453919in}{1.649421in}}{\pgfqpoint{1.451605in}{1.655007in}}{\pgfqpoint{1.447487in}{1.659125in}}%
\pgfpathcurveto{\pgfqpoint{1.443369in}{1.663243in}}{\pgfqpoint{1.437782in}{1.665557in}}{\pgfqpoint{1.431959in}{1.665557in}}%
\pgfpathcurveto{\pgfqpoint{1.426135in}{1.665557in}}{\pgfqpoint{1.420548in}{1.663243in}}{\pgfqpoint{1.416430in}{1.659125in}}%
\pgfpathcurveto{\pgfqpoint{1.412312in}{1.655007in}}{\pgfqpoint{1.409998in}{1.649421in}}{\pgfqpoint{1.409998in}{1.643597in}}%
\pgfpathcurveto{\pgfqpoint{1.409998in}{1.637773in}}{\pgfqpoint{1.412312in}{1.632187in}}{\pgfqpoint{1.416430in}{1.628068in}}%
\pgfpathcurveto{\pgfqpoint{1.420548in}{1.623950in}}{\pgfqpoint{1.426135in}{1.621636in}}{\pgfqpoint{1.431959in}{1.621636in}}%
\pgfpathlineto{\pgfqpoint{1.431959in}{1.621636in}}%
\pgfpathclose%
\pgfusepath{stroke,fill}%
\end{pgfscope}%
\begin{pgfscope}%
\pgfpathrectangle{\pgfqpoint{0.100000in}{0.209042in}}{\pgfqpoint{2.906250in}{2.906250in}}%
\pgfusepath{clip}%
\pgfsetbuttcap%
\pgfsetroundjoin%
\definecolor{currentfill}{rgb}{0.165308,0.067911,0.361816}%
\pgfsetfillcolor{currentfill}%
\pgfsetfillopacity{0.800000}%
\pgfsetlinewidth{1.003750pt}%
\definecolor{currentstroke}{rgb}{0.165308,0.067911,0.361816}%
\pgfsetstrokecolor{currentstroke}%
\pgfsetstrokeopacity{0.800000}%
\pgfsetdash{}{0pt}%
\pgfpathmoveto{\pgfqpoint{1.477598in}{1.770020in}}%
\pgfpathcurveto{\pgfqpoint{1.483422in}{1.770020in}}{\pgfqpoint{1.489008in}{1.772333in}}{\pgfqpoint{1.493126in}{1.776452in}}%
\pgfpathcurveto{\pgfqpoint{1.497244in}{1.780570in}}{\pgfqpoint{1.499558in}{1.786156in}}{\pgfqpoint{1.499558in}{1.791980in}}%
\pgfpathcurveto{\pgfqpoint{1.499558in}{1.797804in}}{\pgfqpoint{1.497244in}{1.803390in}}{\pgfqpoint{1.493126in}{1.807508in}}%
\pgfpathcurveto{\pgfqpoint{1.489008in}{1.811626in}}{\pgfqpoint{1.483422in}{1.813940in}}{\pgfqpoint{1.477598in}{1.813940in}}%
\pgfpathcurveto{\pgfqpoint{1.471774in}{1.813940in}}{\pgfqpoint{1.466188in}{1.811626in}}{\pgfqpoint{1.462070in}{1.807508in}}%
\pgfpathcurveto{\pgfqpoint{1.457952in}{1.803390in}}{\pgfqpoint{1.455638in}{1.797804in}}{\pgfqpoint{1.455638in}{1.791980in}}%
\pgfpathcurveto{\pgfqpoint{1.455638in}{1.786156in}}{\pgfqpoint{1.457952in}{1.780570in}}{\pgfqpoint{1.462070in}{1.776452in}}%
\pgfpathcurveto{\pgfqpoint{1.466188in}{1.772333in}}{\pgfqpoint{1.471774in}{1.770020in}}{\pgfqpoint{1.477598in}{1.770020in}}%
\pgfpathlineto{\pgfqpoint{1.477598in}{1.770020in}}%
\pgfpathclose%
\pgfusepath{stroke,fill}%
\end{pgfscope}%
\begin{pgfscope}%
\pgfpathrectangle{\pgfqpoint{0.100000in}{0.209042in}}{\pgfqpoint{2.906250in}{2.906250in}}%
\pgfusepath{clip}%
\pgfsetbuttcap%
\pgfsetroundjoin%
\definecolor{currentfill}{rgb}{0.297740,0.066117,0.478243}%
\pgfsetfillcolor{currentfill}%
\pgfsetfillopacity{0.800000}%
\pgfsetlinewidth{1.003750pt}%
\definecolor{currentstroke}{rgb}{0.297740,0.066117,0.478243}%
\pgfsetstrokecolor{currentstroke}%
\pgfsetstrokeopacity{0.800000}%
\pgfsetdash{}{0pt}%
\pgfpathmoveto{\pgfqpoint{1.885846in}{1.789881in}}%
\pgfpathcurveto{\pgfqpoint{1.891670in}{1.789881in}}{\pgfqpoint{1.897256in}{1.792195in}}{\pgfqpoint{1.901374in}{1.796313in}}%
\pgfpathcurveto{\pgfqpoint{1.905492in}{1.800432in}}{\pgfqpoint{1.907806in}{1.806018in}}{\pgfqpoint{1.907806in}{1.811842in}}%
\pgfpathcurveto{\pgfqpoint{1.907806in}{1.817666in}}{\pgfqpoint{1.905492in}{1.823252in}}{\pgfqpoint{1.901374in}{1.827370in}}%
\pgfpathcurveto{\pgfqpoint{1.897256in}{1.831488in}}{\pgfqpoint{1.891670in}{1.833802in}}{\pgfqpoint{1.885846in}{1.833802in}}%
\pgfpathcurveto{\pgfqpoint{1.880022in}{1.833802in}}{\pgfqpoint{1.874436in}{1.831488in}}{\pgfqpoint{1.870318in}{1.827370in}}%
\pgfpathcurveto{\pgfqpoint{1.866199in}{1.823252in}}{\pgfqpoint{1.863885in}{1.817666in}}{\pgfqpoint{1.863885in}{1.811842in}}%
\pgfpathcurveto{\pgfqpoint{1.863885in}{1.806018in}}{\pgfqpoint{1.866199in}{1.800432in}}{\pgfqpoint{1.870318in}{1.796313in}}%
\pgfpathcurveto{\pgfqpoint{1.874436in}{1.792195in}}{\pgfqpoint{1.880022in}{1.789881in}}{\pgfqpoint{1.885846in}{1.789881in}}%
\pgfpathlineto{\pgfqpoint{1.885846in}{1.789881in}}%
\pgfpathclose%
\pgfusepath{stroke,fill}%
\end{pgfscope}%
\begin{pgfscope}%
\pgfpathrectangle{\pgfqpoint{0.100000in}{0.209042in}}{\pgfqpoint{2.906250in}{2.906250in}}%
\pgfusepath{clip}%
\pgfsetbuttcap%
\pgfsetroundjoin%
\definecolor{currentfill}{rgb}{0.684224,0.204286,0.484219}%
\pgfsetfillcolor{currentfill}%
\pgfsetfillopacity{0.800000}%
\pgfsetlinewidth{1.003750pt}%
\definecolor{currentstroke}{rgb}{0.684224,0.204286,0.484219}%
\pgfsetstrokecolor{currentstroke}%
\pgfsetstrokeopacity{0.800000}%
\pgfsetdash{}{0pt}%
\pgfpathmoveto{\pgfqpoint{1.395851in}{2.093397in}}%
\pgfpathcurveto{\pgfqpoint{1.401675in}{2.093397in}}{\pgfqpoint{1.407261in}{2.095711in}}{\pgfqpoint{1.411379in}{2.099829in}}%
\pgfpathcurveto{\pgfqpoint{1.415497in}{2.103947in}}{\pgfqpoint{1.417811in}{2.109534in}}{\pgfqpoint{1.417811in}{2.115358in}}%
\pgfpathcurveto{\pgfqpoint{1.417811in}{2.121181in}}{\pgfqpoint{1.415497in}{2.126768in}}{\pgfqpoint{1.411379in}{2.130886in}}%
\pgfpathcurveto{\pgfqpoint{1.407261in}{2.135004in}}{\pgfqpoint{1.401675in}{2.137318in}}{\pgfqpoint{1.395851in}{2.137318in}}%
\pgfpathcurveto{\pgfqpoint{1.390027in}{2.137318in}}{\pgfqpoint{1.384441in}{2.135004in}}{\pgfqpoint{1.380322in}{2.130886in}}%
\pgfpathcurveto{\pgfqpoint{1.376204in}{2.126768in}}{\pgfqpoint{1.373890in}{2.121181in}}{\pgfqpoint{1.373890in}{2.115358in}}%
\pgfpathcurveto{\pgfqpoint{1.373890in}{2.109534in}}{\pgfqpoint{1.376204in}{2.103947in}}{\pgfqpoint{1.380322in}{2.099829in}}%
\pgfpathcurveto{\pgfqpoint{1.384441in}{2.095711in}}{\pgfqpoint{1.390027in}{2.093397in}}{\pgfqpoint{1.395851in}{2.093397in}}%
\pgfpathlineto{\pgfqpoint{1.395851in}{2.093397in}}%
\pgfpathclose%
\pgfusepath{stroke,fill}%
\end{pgfscope}%
\begin{pgfscope}%
\pgfpathrectangle{\pgfqpoint{0.100000in}{0.209042in}}{\pgfqpoint{2.906250in}{2.906250in}}%
\pgfusepath{clip}%
\pgfsetbuttcap%
\pgfsetroundjoin%
\definecolor{currentfill}{rgb}{0.310382,0.069702,0.483186}%
\pgfsetfillcolor{currentfill}%
\pgfsetfillopacity{0.800000}%
\pgfsetlinewidth{1.003750pt}%
\definecolor{currentstroke}{rgb}{0.310382,0.069702,0.483186}%
\pgfsetstrokecolor{currentstroke}%
\pgfsetstrokeopacity{0.800000}%
\pgfsetdash{}{0pt}%
\pgfpathmoveto{\pgfqpoint{1.908503in}{1.595308in}}%
\pgfpathcurveto{\pgfqpoint{1.914327in}{1.595308in}}{\pgfqpoint{1.919913in}{1.597622in}}{\pgfqpoint{1.924031in}{1.601740in}}%
\pgfpathcurveto{\pgfqpoint{1.928149in}{1.605858in}}{\pgfqpoint{1.930463in}{1.611444in}}{\pgfqpoint{1.930463in}{1.617268in}}%
\pgfpathcurveto{\pgfqpoint{1.930463in}{1.623092in}}{\pgfqpoint{1.928149in}{1.628679in}}{\pgfqpoint{1.924031in}{1.632797in}}%
\pgfpathcurveto{\pgfqpoint{1.919913in}{1.636915in}}{\pgfqpoint{1.914327in}{1.639229in}}{\pgfqpoint{1.908503in}{1.639229in}}%
\pgfpathcurveto{\pgfqpoint{1.902679in}{1.639229in}}{\pgfqpoint{1.897093in}{1.636915in}}{\pgfqpoint{1.892974in}{1.632797in}}%
\pgfpathcurveto{\pgfqpoint{1.888856in}{1.628679in}}{\pgfqpoint{1.886542in}{1.623092in}}{\pgfqpoint{1.886542in}{1.617268in}}%
\pgfpathcurveto{\pgfqpoint{1.886542in}{1.611444in}}{\pgfqpoint{1.888856in}{1.605858in}}{\pgfqpoint{1.892974in}{1.601740in}}%
\pgfpathcurveto{\pgfqpoint{1.897093in}{1.597622in}}{\pgfqpoint{1.902679in}{1.595308in}}{\pgfqpoint{1.908503in}{1.595308in}}%
\pgfpathlineto{\pgfqpoint{1.908503in}{1.595308in}}%
\pgfpathclose%
\pgfusepath{stroke,fill}%
\end{pgfscope}%
\begin{pgfscope}%
\pgfpathrectangle{\pgfqpoint{0.100000in}{0.209042in}}{\pgfqpoint{2.906250in}{2.906250in}}%
\pgfusepath{clip}%
\pgfsetbuttcap%
\pgfsetroundjoin%
\definecolor{currentfill}{rgb}{0.519045,0.150383,0.507443}%
\pgfsetfillcolor{currentfill}%
\pgfsetfillopacity{0.800000}%
\pgfsetlinewidth{1.003750pt}%
\definecolor{currentstroke}{rgb}{0.519045,0.150383,0.507443}%
\pgfsetstrokecolor{currentstroke}%
\pgfsetstrokeopacity{0.800000}%
\pgfsetdash{}{0pt}%
\pgfpathmoveto{\pgfqpoint{1.715495in}{2.032546in}}%
\pgfpathcurveto{\pgfqpoint{1.721319in}{2.032546in}}{\pgfqpoint{1.726905in}{2.034860in}}{\pgfqpoint{1.731023in}{2.038978in}}%
\pgfpathcurveto{\pgfqpoint{1.735141in}{2.043096in}}{\pgfqpoint{1.737455in}{2.048682in}}{\pgfqpoint{1.737455in}{2.054506in}}%
\pgfpathcurveto{\pgfqpoint{1.737455in}{2.060330in}}{\pgfqpoint{1.735141in}{2.065916in}}{\pgfqpoint{1.731023in}{2.070035in}}%
\pgfpathcurveto{\pgfqpoint{1.726905in}{2.074153in}}{\pgfqpoint{1.721319in}{2.076467in}}{\pgfqpoint{1.715495in}{2.076467in}}%
\pgfpathcurveto{\pgfqpoint{1.709671in}{2.076467in}}{\pgfqpoint{1.704085in}{2.074153in}}{\pgfqpoint{1.699967in}{2.070035in}}%
\pgfpathcurveto{\pgfqpoint{1.695848in}{2.065916in}}{\pgfqpoint{1.693535in}{2.060330in}}{\pgfqpoint{1.693535in}{2.054506in}}%
\pgfpathcurveto{\pgfqpoint{1.693535in}{2.048682in}}{\pgfqpoint{1.695848in}{2.043096in}}{\pgfqpoint{1.699967in}{2.038978in}}%
\pgfpathcurveto{\pgfqpoint{1.704085in}{2.034860in}}{\pgfqpoint{1.709671in}{2.032546in}}{\pgfqpoint{1.715495in}{2.032546in}}%
\pgfpathlineto{\pgfqpoint{1.715495in}{2.032546in}}%
\pgfpathclose%
\pgfusepath{stroke,fill}%
\end{pgfscope}%
\begin{pgfscope}%
\pgfpathrectangle{\pgfqpoint{0.100000in}{0.209042in}}{\pgfqpoint{2.906250in}{2.906250in}}%
\pgfusepath{clip}%
\pgfsetbuttcap%
\pgfsetroundjoin%
\definecolor{currentfill}{rgb}{0.600868,0.177743,0.500394}%
\pgfsetfillcolor{currentfill}%
\pgfsetfillopacity{0.800000}%
\pgfsetlinewidth{1.003750pt}%
\definecolor{currentstroke}{rgb}{0.600868,0.177743,0.500394}%
\pgfsetstrokecolor{currentstroke}%
\pgfsetstrokeopacity{0.800000}%
\pgfsetdash{}{0pt}%
\pgfpathmoveto{\pgfqpoint{1.799352in}{1.307593in}}%
\pgfpathcurveto{\pgfqpoint{1.805176in}{1.307593in}}{\pgfqpoint{1.810762in}{1.309907in}}{\pgfqpoint{1.814880in}{1.314025in}}%
\pgfpathcurveto{\pgfqpoint{1.818998in}{1.318143in}}{\pgfqpoint{1.821312in}{1.323730in}}{\pgfqpoint{1.821312in}{1.329554in}}%
\pgfpathcurveto{\pgfqpoint{1.821312in}{1.335377in}}{\pgfqpoint{1.818998in}{1.340964in}}{\pgfqpoint{1.814880in}{1.345082in}}%
\pgfpathcurveto{\pgfqpoint{1.810762in}{1.349200in}}{\pgfqpoint{1.805176in}{1.351514in}}{\pgfqpoint{1.799352in}{1.351514in}}%
\pgfpathcurveto{\pgfqpoint{1.793528in}{1.351514in}}{\pgfqpoint{1.787942in}{1.349200in}}{\pgfqpoint{1.783824in}{1.345082in}}%
\pgfpathcurveto{\pgfqpoint{1.779705in}{1.340964in}}{\pgfqpoint{1.777392in}{1.335377in}}{\pgfqpoint{1.777392in}{1.329554in}}%
\pgfpathcurveto{\pgfqpoint{1.777392in}{1.323730in}}{\pgfqpoint{1.779705in}{1.318143in}}{\pgfqpoint{1.783824in}{1.314025in}}%
\pgfpathcurveto{\pgfqpoint{1.787942in}{1.309907in}}{\pgfqpoint{1.793528in}{1.307593in}}{\pgfqpoint{1.799352in}{1.307593in}}%
\pgfpathlineto{\pgfqpoint{1.799352in}{1.307593in}}%
\pgfpathclose%
\pgfusepath{stroke,fill}%
\end{pgfscope}%
\begin{pgfscope}%
\pgfpathrectangle{\pgfqpoint{0.100000in}{0.209042in}}{\pgfqpoint{2.906250in}{2.906250in}}%
\pgfusepath{clip}%
\pgfsetbuttcap%
\pgfsetroundjoin%
\definecolor{currentfill}{rgb}{0.359898,0.087831,0.496778}%
\pgfsetfillcolor{currentfill}%
\pgfsetfillopacity{0.800000}%
\pgfsetlinewidth{1.003750pt}%
\definecolor{currentstroke}{rgb}{0.359898,0.087831,0.496778}%
\pgfsetstrokecolor{currentstroke}%
\pgfsetstrokeopacity{0.800000}%
\pgfsetdash{}{0pt}%
\pgfpathmoveto{\pgfqpoint{1.458439in}{1.911293in}}%
\pgfpathcurveto{\pgfqpoint{1.464263in}{1.911293in}}{\pgfqpoint{1.469850in}{1.913607in}}{\pgfqpoint{1.473968in}{1.917725in}}%
\pgfpathcurveto{\pgfqpoint{1.478086in}{1.921843in}}{\pgfqpoint{1.480400in}{1.927429in}}{\pgfqpoint{1.480400in}{1.933253in}}%
\pgfpathcurveto{\pgfqpoint{1.480400in}{1.939077in}}{\pgfqpoint{1.478086in}{1.944664in}}{\pgfqpoint{1.473968in}{1.948782in}}%
\pgfpathcurveto{\pgfqpoint{1.469850in}{1.952900in}}{\pgfqpoint{1.464263in}{1.955214in}}{\pgfqpoint{1.458439in}{1.955214in}}%
\pgfpathcurveto{\pgfqpoint{1.452616in}{1.955214in}}{\pgfqpoint{1.447029in}{1.952900in}}{\pgfqpoint{1.442911in}{1.948782in}}%
\pgfpathcurveto{\pgfqpoint{1.438793in}{1.944664in}}{\pgfqpoint{1.436479in}{1.939077in}}{\pgfqpoint{1.436479in}{1.933253in}}%
\pgfpathcurveto{\pgfqpoint{1.436479in}{1.927429in}}{\pgfqpoint{1.438793in}{1.921843in}}{\pgfqpoint{1.442911in}{1.917725in}}%
\pgfpathcurveto{\pgfqpoint{1.447029in}{1.913607in}}{\pgfqpoint{1.452616in}{1.911293in}}{\pgfqpoint{1.458439in}{1.911293in}}%
\pgfpathlineto{\pgfqpoint{1.458439in}{1.911293in}}%
\pgfpathclose%
\pgfusepath{stroke,fill}%
\end{pgfscope}%
\begin{pgfscope}%
\pgfpathrectangle{\pgfqpoint{0.100000in}{0.209042in}}{\pgfqpoint{2.906250in}{2.906250in}}%
\pgfusepath{clip}%
\pgfsetbuttcap%
\pgfsetroundjoin%
\definecolor{currentfill}{rgb}{0.372116,0.092816,0.499053}%
\pgfsetfillcolor{currentfill}%
\pgfsetfillopacity{0.800000}%
\pgfsetlinewidth{1.003750pt}%
\definecolor{currentstroke}{rgb}{0.372116,0.092816,0.499053}%
\pgfsetstrokecolor{currentstroke}%
\pgfsetstrokeopacity{0.800000}%
\pgfsetdash{}{0pt}%
\pgfpathmoveto{\pgfqpoint{1.335165in}{1.609436in}}%
\pgfpathcurveto{\pgfqpoint{1.340989in}{1.609436in}}{\pgfqpoint{1.346575in}{1.611750in}}{\pgfqpoint{1.350693in}{1.615868in}}%
\pgfpathcurveto{\pgfqpoint{1.354811in}{1.619986in}}{\pgfqpoint{1.357125in}{1.625572in}}{\pgfqpoint{1.357125in}{1.631396in}}%
\pgfpathcurveto{\pgfqpoint{1.357125in}{1.637220in}}{\pgfqpoint{1.354811in}{1.642806in}}{\pgfqpoint{1.350693in}{1.646924in}}%
\pgfpathcurveto{\pgfqpoint{1.346575in}{1.651042in}}{\pgfqpoint{1.340989in}{1.653356in}}{\pgfqpoint{1.335165in}{1.653356in}}%
\pgfpathcurveto{\pgfqpoint{1.329341in}{1.653356in}}{\pgfqpoint{1.323755in}{1.651042in}}{\pgfqpoint{1.319636in}{1.646924in}}%
\pgfpathcurveto{\pgfqpoint{1.315518in}{1.642806in}}{\pgfqpoint{1.313204in}{1.637220in}}{\pgfqpoint{1.313204in}{1.631396in}}%
\pgfpathcurveto{\pgfqpoint{1.313204in}{1.625572in}}{\pgfqpoint{1.315518in}{1.619986in}}{\pgfqpoint{1.319636in}{1.615868in}}%
\pgfpathcurveto{\pgfqpoint{1.323755in}{1.611750in}}{\pgfqpoint{1.329341in}{1.609436in}}{\pgfqpoint{1.335165in}{1.609436in}}%
\pgfpathlineto{\pgfqpoint{1.335165in}{1.609436in}}%
\pgfpathclose%
\pgfusepath{stroke,fill}%
\end{pgfscope}%
\begin{pgfscope}%
\pgfpathrectangle{\pgfqpoint{0.100000in}{0.209042in}}{\pgfqpoint{2.906250in}{2.906250in}}%
\pgfusepath{clip}%
\pgfsetbuttcap%
\pgfsetroundjoin%
\definecolor{currentfill}{rgb}{0.078815,0.054184,0.211667}%
\pgfsetfillcolor{currentfill}%
\pgfsetfillopacity{0.800000}%
\pgfsetlinewidth{1.003750pt}%
\definecolor{currentstroke}{rgb}{0.078815,0.054184,0.211667}%
\pgfsetstrokecolor{currentstroke}%
\pgfsetstrokeopacity{0.800000}%
\pgfsetdash{}{0pt}%
\pgfpathmoveto{\pgfqpoint{1.596593in}{1.633107in}}%
\pgfpathcurveto{\pgfqpoint{1.602417in}{1.633107in}}{\pgfqpoint{1.608003in}{1.635421in}}{\pgfqpoint{1.612121in}{1.639539in}}%
\pgfpathcurveto{\pgfqpoint{1.616239in}{1.643657in}}{\pgfqpoint{1.618553in}{1.649243in}}{\pgfqpoint{1.618553in}{1.655067in}}%
\pgfpathcurveto{\pgfqpoint{1.618553in}{1.660891in}}{\pgfqpoint{1.616239in}{1.666477in}}{\pgfqpoint{1.612121in}{1.670595in}}%
\pgfpathcurveto{\pgfqpoint{1.608003in}{1.674714in}}{\pgfqpoint{1.602417in}{1.677027in}}{\pgfqpoint{1.596593in}{1.677027in}}%
\pgfpathcurveto{\pgfqpoint{1.590769in}{1.677027in}}{\pgfqpoint{1.585183in}{1.674714in}}{\pgfqpoint{1.581065in}{1.670595in}}%
\pgfpathcurveto{\pgfqpoint{1.576946in}{1.666477in}}{\pgfqpoint{1.574633in}{1.660891in}}{\pgfqpoint{1.574633in}{1.655067in}}%
\pgfpathcurveto{\pgfqpoint{1.574633in}{1.649243in}}{\pgfqpoint{1.576946in}{1.643657in}}{\pgfqpoint{1.581065in}{1.639539in}}%
\pgfpathcurveto{\pgfqpoint{1.585183in}{1.635421in}}{\pgfqpoint{1.590769in}{1.633107in}}{\pgfqpoint{1.596593in}{1.633107in}}%
\pgfpathlineto{\pgfqpoint{1.596593in}{1.633107in}}%
\pgfpathclose%
\pgfusepath{stroke,fill}%
\end{pgfscope}%
\begin{pgfscope}%
\pgfpathrectangle{\pgfqpoint{0.100000in}{0.209042in}}{\pgfqpoint{2.906250in}{2.906250in}}%
\pgfusepath{clip}%
\pgfsetbuttcap%
\pgfsetroundjoin%
\definecolor{currentfill}{rgb}{0.840636,0.268953,0.424666}%
\pgfsetfillcolor{currentfill}%
\pgfsetfillopacity{0.800000}%
\pgfsetlinewidth{1.003750pt}%
\definecolor{currentstroke}{rgb}{0.840636,0.268953,0.424666}%
\pgfsetstrokecolor{currentstroke}%
\pgfsetstrokeopacity{0.800000}%
\pgfsetdash{}{0pt}%
\pgfpathmoveto{\pgfqpoint{1.605036in}{2.211882in}}%
\pgfpathcurveto{\pgfqpoint{1.610860in}{2.211882in}}{\pgfqpoint{1.616447in}{2.214196in}}{\pgfqpoint{1.620565in}{2.218314in}}%
\pgfpathcurveto{\pgfqpoint{1.624683in}{2.222432in}}{\pgfqpoint{1.626997in}{2.228018in}}{\pgfqpoint{1.626997in}{2.233842in}}%
\pgfpathcurveto{\pgfqpoint{1.626997in}{2.239666in}}{\pgfqpoint{1.624683in}{2.245252in}}{\pgfqpoint{1.620565in}{2.249370in}}%
\pgfpathcurveto{\pgfqpoint{1.616447in}{2.253489in}}{\pgfqpoint{1.610860in}{2.255802in}}{\pgfqpoint{1.605036in}{2.255802in}}%
\pgfpathcurveto{\pgfqpoint{1.599213in}{2.255802in}}{\pgfqpoint{1.593626in}{2.253489in}}{\pgfqpoint{1.589508in}{2.249370in}}%
\pgfpathcurveto{\pgfqpoint{1.585390in}{2.245252in}}{\pgfqpoint{1.583076in}{2.239666in}}{\pgfqpoint{1.583076in}{2.233842in}}%
\pgfpathcurveto{\pgfqpoint{1.583076in}{2.228018in}}{\pgfqpoint{1.585390in}{2.222432in}}{\pgfqpoint{1.589508in}{2.218314in}}%
\pgfpathcurveto{\pgfqpoint{1.593626in}{2.214196in}}{\pgfqpoint{1.599213in}{2.211882in}}{\pgfqpoint{1.605036in}{2.211882in}}%
\pgfpathlineto{\pgfqpoint{1.605036in}{2.211882in}}%
\pgfpathclose%
\pgfusepath{stroke,fill}%
\end{pgfscope}%
\begin{pgfscope}%
\pgfpathrectangle{\pgfqpoint{0.100000in}{0.209042in}}{\pgfqpoint{2.906250in}{2.906250in}}%
\pgfusepath{clip}%
\pgfsetbuttcap%
\pgfsetroundjoin%
\definecolor{currentfill}{rgb}{0.146785,0.068738,0.334011}%
\pgfsetfillcolor{currentfill}%
\pgfsetfillopacity{0.800000}%
\pgfsetlinewidth{1.003750pt}%
\definecolor{currentstroke}{rgb}{0.146785,0.068738,0.334011}%
\pgfsetstrokecolor{currentstroke}%
\pgfsetstrokeopacity{0.800000}%
\pgfsetdash{}{0pt}%
\pgfpathmoveto{\pgfqpoint{1.536476in}{1.791534in}}%
\pgfpathcurveto{\pgfqpoint{1.542300in}{1.791534in}}{\pgfqpoint{1.547886in}{1.793848in}}{\pgfqpoint{1.552004in}{1.797966in}}%
\pgfpathcurveto{\pgfqpoint{1.556122in}{1.802085in}}{\pgfqpoint{1.558436in}{1.807671in}}{\pgfqpoint{1.558436in}{1.813495in}}%
\pgfpathcurveto{\pgfqpoint{1.558436in}{1.819319in}}{\pgfqpoint{1.556122in}{1.824905in}}{\pgfqpoint{1.552004in}{1.829023in}}%
\pgfpathcurveto{\pgfqpoint{1.547886in}{1.833141in}}{\pgfqpoint{1.542300in}{1.835455in}}{\pgfqpoint{1.536476in}{1.835455in}}%
\pgfpathcurveto{\pgfqpoint{1.530652in}{1.835455in}}{\pgfqpoint{1.525066in}{1.833141in}}{\pgfqpoint{1.520947in}{1.829023in}}%
\pgfpathcurveto{\pgfqpoint{1.516829in}{1.824905in}}{\pgfqpoint{1.514515in}{1.819319in}}{\pgfqpoint{1.514515in}{1.813495in}}%
\pgfpathcurveto{\pgfqpoint{1.514515in}{1.807671in}}{\pgfqpoint{1.516829in}{1.802085in}}{\pgfqpoint{1.520947in}{1.797966in}}%
\pgfpathcurveto{\pgfqpoint{1.525066in}{1.793848in}}{\pgfqpoint{1.530652in}{1.791534in}}{\pgfqpoint{1.536476in}{1.791534in}}%
\pgfpathlineto{\pgfqpoint{1.536476in}{1.791534in}}%
\pgfpathclose%
\pgfusepath{stroke,fill}%
\end{pgfscope}%
\begin{pgfscope}%
\pgfpathrectangle{\pgfqpoint{0.100000in}{0.209042in}}{\pgfqpoint{2.906250in}{2.906250in}}%
\pgfusepath{clip}%
\pgfsetbuttcap%
\pgfsetroundjoin%
\definecolor{currentfill}{rgb}{0.347636,0.082946,0.494121}%
\pgfsetfillcolor{currentfill}%
\pgfsetfillopacity{0.800000}%
\pgfsetlinewidth{1.003750pt}%
\definecolor{currentstroke}{rgb}{0.347636,0.082946,0.494121}%
\pgfsetstrokecolor{currentstroke}%
\pgfsetstrokeopacity{0.800000}%
\pgfsetdash{}{0pt}%
\pgfpathmoveto{\pgfqpoint{1.902267in}{1.541077in}}%
\pgfpathcurveto{\pgfqpoint{1.908090in}{1.541077in}}{\pgfqpoint{1.913677in}{1.543391in}}{\pgfqpoint{1.917795in}{1.547509in}}%
\pgfpathcurveto{\pgfqpoint{1.921913in}{1.551627in}}{\pgfqpoint{1.924227in}{1.557213in}}{\pgfqpoint{1.924227in}{1.563037in}}%
\pgfpathcurveto{\pgfqpoint{1.924227in}{1.568861in}}{\pgfqpoint{1.921913in}{1.574447in}}{\pgfqpoint{1.917795in}{1.578565in}}%
\pgfpathcurveto{\pgfqpoint{1.913677in}{1.582683in}}{\pgfqpoint{1.908090in}{1.584997in}}{\pgfqpoint{1.902267in}{1.584997in}}%
\pgfpathcurveto{\pgfqpoint{1.896443in}{1.584997in}}{\pgfqpoint{1.890856in}{1.582683in}}{\pgfqpoint{1.886738in}{1.578565in}}%
\pgfpathcurveto{\pgfqpoint{1.882620in}{1.574447in}}{\pgfqpoint{1.880306in}{1.568861in}}{\pgfqpoint{1.880306in}{1.563037in}}%
\pgfpathcurveto{\pgfqpoint{1.880306in}{1.557213in}}{\pgfqpoint{1.882620in}{1.551627in}}{\pgfqpoint{1.886738in}{1.547509in}}%
\pgfpathcurveto{\pgfqpoint{1.890856in}{1.543391in}}{\pgfqpoint{1.896443in}{1.541077in}}{\pgfqpoint{1.902267in}{1.541077in}}%
\pgfpathlineto{\pgfqpoint{1.902267in}{1.541077in}}%
\pgfpathclose%
\pgfusepath{stroke,fill}%
\end{pgfscope}%
\begin{pgfscope}%
\pgfpathrectangle{\pgfqpoint{0.100000in}{0.209042in}}{\pgfqpoint{2.906250in}{2.906250in}}%
\pgfusepath{clip}%
\pgfsetbuttcap%
\pgfsetroundjoin%
\definecolor{currentfill}{rgb}{0.469640,0.132245,0.507809}%
\pgfsetfillcolor{currentfill}%
\pgfsetfillopacity{0.800000}%
\pgfsetlinewidth{1.003750pt}%
\definecolor{currentstroke}{rgb}{0.469640,0.132245,0.507809}%
\pgfsetstrokecolor{currentstroke}%
\pgfsetstrokeopacity{0.800000}%
\pgfsetdash{}{0pt}%
\pgfpathmoveto{\pgfqpoint{1.764776in}{1.989521in}}%
\pgfpathcurveto{\pgfqpoint{1.770600in}{1.989521in}}{\pgfqpoint{1.776186in}{1.991835in}}{\pgfqpoint{1.780304in}{1.995953in}}%
\pgfpathcurveto{\pgfqpoint{1.784422in}{2.000071in}}{\pgfqpoint{1.786736in}{2.005657in}}{\pgfqpoint{1.786736in}{2.011481in}}%
\pgfpathcurveto{\pgfqpoint{1.786736in}{2.017305in}}{\pgfqpoint{1.784422in}{2.022891in}}{\pgfqpoint{1.780304in}{2.027010in}}%
\pgfpathcurveto{\pgfqpoint{1.776186in}{2.031128in}}{\pgfqpoint{1.770600in}{2.033442in}}{\pgfqpoint{1.764776in}{2.033442in}}%
\pgfpathcurveto{\pgfqpoint{1.758952in}{2.033442in}}{\pgfqpoint{1.753366in}{2.031128in}}{\pgfqpoint{1.749247in}{2.027010in}}%
\pgfpathcurveto{\pgfqpoint{1.745129in}{2.022891in}}{\pgfqpoint{1.742815in}{2.017305in}}{\pgfqpoint{1.742815in}{2.011481in}}%
\pgfpathcurveto{\pgfqpoint{1.742815in}{2.005657in}}{\pgfqpoint{1.745129in}{2.000071in}}{\pgfqpoint{1.749247in}{1.995953in}}%
\pgfpathcurveto{\pgfqpoint{1.753366in}{1.991835in}}{\pgfqpoint{1.758952in}{1.989521in}}{\pgfqpoint{1.764776in}{1.989521in}}%
\pgfpathlineto{\pgfqpoint{1.764776in}{1.989521in}}%
\pgfpathclose%
\pgfusepath{stroke,fill}%
\end{pgfscope}%
\begin{pgfscope}%
\pgfpathrectangle{\pgfqpoint{0.100000in}{0.209042in}}{\pgfqpoint{2.906250in}{2.906250in}}%
\pgfusepath{clip}%
\pgfsetbuttcap%
\pgfsetroundjoin%
\definecolor{currentfill}{rgb}{0.396467,0.102902,0.502658}%
\pgfsetfillcolor{currentfill}%
\pgfsetfillopacity{0.800000}%
\pgfsetlinewidth{1.003750pt}%
\definecolor{currentstroke}{rgb}{0.396467,0.102902,0.502658}%
\pgfsetstrokecolor{currentstroke}%
\pgfsetstrokeopacity{0.800000}%
\pgfsetdash{}{0pt}%
\pgfpathmoveto{\pgfqpoint{1.353273in}{1.554980in}}%
\pgfpathcurveto{\pgfqpoint{1.359097in}{1.554980in}}{\pgfqpoint{1.364683in}{1.557294in}}{\pgfqpoint{1.368801in}{1.561412in}}%
\pgfpathcurveto{\pgfqpoint{1.372920in}{1.565530in}}{\pgfqpoint{1.375233in}{1.571116in}}{\pgfqpoint{1.375233in}{1.576940in}}%
\pgfpathcurveto{\pgfqpoint{1.375233in}{1.582764in}}{\pgfqpoint{1.372920in}{1.588350in}}{\pgfqpoint{1.368801in}{1.592468in}}%
\pgfpathcurveto{\pgfqpoint{1.364683in}{1.596586in}}{\pgfqpoint{1.359097in}{1.598900in}}{\pgfqpoint{1.353273in}{1.598900in}}%
\pgfpathcurveto{\pgfqpoint{1.347449in}{1.598900in}}{\pgfqpoint{1.341863in}{1.596586in}}{\pgfqpoint{1.337745in}{1.592468in}}%
\pgfpathcurveto{\pgfqpoint{1.333627in}{1.588350in}}{\pgfqpoint{1.331313in}{1.582764in}}{\pgfqpoint{1.331313in}{1.576940in}}%
\pgfpathcurveto{\pgfqpoint{1.331313in}{1.571116in}}{\pgfqpoint{1.333627in}{1.565530in}}{\pgfqpoint{1.337745in}{1.561412in}}%
\pgfpathcurveto{\pgfqpoint{1.341863in}{1.557294in}}{\pgfqpoint{1.347449in}{1.554980in}}{\pgfqpoint{1.353273in}{1.554980in}}%
\pgfpathlineto{\pgfqpoint{1.353273in}{1.554980in}}%
\pgfpathclose%
\pgfusepath{stroke,fill}%
\end{pgfscope}%
\begin{pgfscope}%
\pgfpathrectangle{\pgfqpoint{0.100000in}{0.209042in}}{\pgfqpoint{2.906250in}{2.906250in}}%
\pgfusepath{clip}%
\pgfsetbuttcap%
\pgfsetroundjoin%
\definecolor{currentfill}{rgb}{0.378211,0.095332,0.500067}%
\pgfsetfillcolor{currentfill}%
\pgfsetfillopacity{0.800000}%
\pgfsetlinewidth{1.003750pt}%
\definecolor{currentstroke}{rgb}{0.378211,0.095332,0.500067}%
\pgfsetstrokecolor{currentstroke}%
\pgfsetstrokeopacity{0.800000}%
\pgfsetdash{}{0pt}%
\pgfpathmoveto{\pgfqpoint{1.321519in}{1.640967in}}%
\pgfpathcurveto{\pgfqpoint{1.327343in}{1.640967in}}{\pgfqpoint{1.332929in}{1.643281in}}{\pgfqpoint{1.337047in}{1.647399in}}%
\pgfpathcurveto{\pgfqpoint{1.341165in}{1.651517in}}{\pgfqpoint{1.343479in}{1.657103in}}{\pgfqpoint{1.343479in}{1.662927in}}%
\pgfpathcurveto{\pgfqpoint{1.343479in}{1.668751in}}{\pgfqpoint{1.341165in}{1.674337in}}{\pgfqpoint{1.337047in}{1.678456in}}%
\pgfpathcurveto{\pgfqpoint{1.332929in}{1.682574in}}{\pgfqpoint{1.327343in}{1.684888in}}{\pgfqpoint{1.321519in}{1.684888in}}%
\pgfpathcurveto{\pgfqpoint{1.315695in}{1.684888in}}{\pgfqpoint{1.310109in}{1.682574in}}{\pgfqpoint{1.305991in}{1.678456in}}%
\pgfpathcurveto{\pgfqpoint{1.301873in}{1.674337in}}{\pgfqpoint{1.299559in}{1.668751in}}{\pgfqpoint{1.299559in}{1.662927in}}%
\pgfpathcurveto{\pgfqpoint{1.299559in}{1.657103in}}{\pgfqpoint{1.301873in}{1.651517in}}{\pgfqpoint{1.305991in}{1.647399in}}%
\pgfpathcurveto{\pgfqpoint{1.310109in}{1.643281in}}{\pgfqpoint{1.315695in}{1.640967in}}{\pgfqpoint{1.321519in}{1.640967in}}%
\pgfpathlineto{\pgfqpoint{1.321519in}{1.640967in}}%
\pgfpathclose%
\pgfusepath{stroke,fill}%
\end{pgfscope}%
\begin{pgfscope}%
\pgfpathrectangle{\pgfqpoint{0.100000in}{0.209042in}}{\pgfqpoint{2.906250in}{2.906250in}}%
\pgfusepath{clip}%
\pgfsetbuttcap%
\pgfsetroundjoin%
\definecolor{currentfill}{rgb}{0.297740,0.066117,0.478243}%
\pgfsetfillcolor{currentfill}%
\pgfsetfillopacity{0.800000}%
\pgfsetlinewidth{1.003750pt}%
\definecolor{currentstroke}{rgb}{0.297740,0.066117,0.478243}%
\pgfsetstrokecolor{currentstroke}%
\pgfsetstrokeopacity{0.800000}%
\pgfsetdash{}{0pt}%
\pgfpathmoveto{\pgfqpoint{1.821953in}{1.517037in}}%
\pgfpathcurveto{\pgfqpoint{1.827777in}{1.517037in}}{\pgfqpoint{1.833363in}{1.519351in}}{\pgfqpoint{1.837481in}{1.523469in}}%
\pgfpathcurveto{\pgfqpoint{1.841599in}{1.527588in}}{\pgfqpoint{1.843913in}{1.533174in}}{\pgfqpoint{1.843913in}{1.538998in}}%
\pgfpathcurveto{\pgfqpoint{1.843913in}{1.544822in}}{\pgfqpoint{1.841599in}{1.550408in}}{\pgfqpoint{1.837481in}{1.554526in}}%
\pgfpathcurveto{\pgfqpoint{1.833363in}{1.558644in}}{\pgfqpoint{1.827777in}{1.560958in}}{\pgfqpoint{1.821953in}{1.560958in}}%
\pgfpathcurveto{\pgfqpoint{1.816129in}{1.560958in}}{\pgfqpoint{1.810543in}{1.558644in}}{\pgfqpoint{1.806424in}{1.554526in}}%
\pgfpathcurveto{\pgfqpoint{1.802306in}{1.550408in}}{\pgfqpoint{1.799992in}{1.544822in}}{\pgfqpoint{1.799992in}{1.538998in}}%
\pgfpathcurveto{\pgfqpoint{1.799992in}{1.533174in}}{\pgfqpoint{1.802306in}{1.527588in}}{\pgfqpoint{1.806424in}{1.523469in}}%
\pgfpathcurveto{\pgfqpoint{1.810543in}{1.519351in}}{\pgfqpoint{1.816129in}{1.517037in}}{\pgfqpoint{1.821953in}{1.517037in}}%
\pgfpathlineto{\pgfqpoint{1.821953in}{1.517037in}}%
\pgfpathclose%
\pgfusepath{stroke,fill}%
\end{pgfscope}%
\begin{pgfscope}%
\pgfpathrectangle{\pgfqpoint{0.100000in}{0.209042in}}{\pgfqpoint{2.906250in}{2.906250in}}%
\pgfusepath{clip}%
\pgfsetbuttcap%
\pgfsetroundjoin%
\definecolor{currentfill}{rgb}{0.594508,0.175701,0.501241}%
\pgfsetfillcolor{currentfill}%
\pgfsetfillopacity{0.800000}%
\pgfsetlinewidth{1.003750pt}%
\definecolor{currentstroke}{rgb}{0.594508,0.175701,0.501241}%
\pgfsetstrokecolor{currentstroke}%
\pgfsetstrokeopacity{0.800000}%
\pgfsetdash{}{0pt}%
\pgfpathmoveto{\pgfqpoint{1.502992in}{1.321697in}}%
\pgfpathcurveto{\pgfqpoint{1.508816in}{1.321697in}}{\pgfqpoint{1.514402in}{1.324011in}}{\pgfqpoint{1.518520in}{1.328129in}}%
\pgfpathcurveto{\pgfqpoint{1.522639in}{1.332247in}}{\pgfqpoint{1.524952in}{1.337833in}}{\pgfqpoint{1.524952in}{1.343657in}}%
\pgfpathcurveto{\pgfqpoint{1.524952in}{1.349481in}}{\pgfqpoint{1.522639in}{1.355067in}}{\pgfqpoint{1.518520in}{1.359185in}}%
\pgfpathcurveto{\pgfqpoint{1.514402in}{1.363304in}}{\pgfqpoint{1.508816in}{1.365618in}}{\pgfqpoint{1.502992in}{1.365618in}}%
\pgfpathcurveto{\pgfqpoint{1.497168in}{1.365618in}}{\pgfqpoint{1.491582in}{1.363304in}}{\pgfqpoint{1.487464in}{1.359185in}}%
\pgfpathcurveto{\pgfqpoint{1.483346in}{1.355067in}}{\pgfqpoint{1.481032in}{1.349481in}}{\pgfqpoint{1.481032in}{1.343657in}}%
\pgfpathcurveto{\pgfqpoint{1.481032in}{1.337833in}}{\pgfqpoint{1.483346in}{1.332247in}}{\pgfqpoint{1.487464in}{1.328129in}}%
\pgfpathcurveto{\pgfqpoint{1.491582in}{1.324011in}}{\pgfqpoint{1.497168in}{1.321697in}}{\pgfqpoint{1.502992in}{1.321697in}}%
\pgfpathlineto{\pgfqpoint{1.502992in}{1.321697in}}%
\pgfpathclose%
\pgfusepath{stroke,fill}%
\end{pgfscope}%
\begin{pgfscope}%
\pgfpathrectangle{\pgfqpoint{0.100000in}{0.209042in}}{\pgfqpoint{2.906250in}{2.906250in}}%
\pgfusepath{clip}%
\pgfsetbuttcap%
\pgfsetroundjoin%
\definecolor{currentfill}{rgb}{0.748378,0.226377,0.464794}%
\pgfsetfillcolor{currentfill}%
\pgfsetfillopacity{0.800000}%
\pgfsetlinewidth{1.003750pt}%
\definecolor{currentstroke}{rgb}{0.748378,0.226377,0.464794}%
\pgfsetstrokecolor{currentstroke}%
\pgfsetstrokeopacity{0.800000}%
\pgfsetdash{}{0pt}%
\pgfpathmoveto{\pgfqpoint{1.440711in}{2.138689in}}%
\pgfpathcurveto{\pgfqpoint{1.446535in}{2.138689in}}{\pgfqpoint{1.452121in}{2.141003in}}{\pgfqpoint{1.456239in}{2.145121in}}%
\pgfpathcurveto{\pgfqpoint{1.460358in}{2.149239in}}{\pgfqpoint{1.462671in}{2.154826in}}{\pgfqpoint{1.462671in}{2.160649in}}%
\pgfpathcurveto{\pgfqpoint{1.462671in}{2.166473in}}{\pgfqpoint{1.460358in}{2.172060in}}{\pgfqpoint{1.456239in}{2.176178in}}%
\pgfpathcurveto{\pgfqpoint{1.452121in}{2.180296in}}{\pgfqpoint{1.446535in}{2.182610in}}{\pgfqpoint{1.440711in}{2.182610in}}%
\pgfpathcurveto{\pgfqpoint{1.434887in}{2.182610in}}{\pgfqpoint{1.429301in}{2.180296in}}{\pgfqpoint{1.425183in}{2.176178in}}%
\pgfpathcurveto{\pgfqpoint{1.421065in}{2.172060in}}{\pgfqpoint{1.418751in}{2.166473in}}{\pgfqpoint{1.418751in}{2.160649in}}%
\pgfpathcurveto{\pgfqpoint{1.418751in}{2.154826in}}{\pgfqpoint{1.421065in}{2.149239in}}{\pgfqpoint{1.425183in}{2.145121in}}%
\pgfpathcurveto{\pgfqpoint{1.429301in}{2.141003in}}{\pgfqpoint{1.434887in}{2.138689in}}{\pgfqpoint{1.440711in}{2.138689in}}%
\pgfpathlineto{\pgfqpoint{1.440711in}{2.138689in}}%
\pgfpathclose%
\pgfusepath{stroke,fill}%
\end{pgfscope}%
\begin{pgfscope}%
\pgfpathrectangle{\pgfqpoint{0.100000in}{0.209042in}}{\pgfqpoint{2.906250in}{2.906250in}}%
\pgfusepath{clip}%
\pgfsetbuttcap%
\pgfsetroundjoin%
\definecolor{currentfill}{rgb}{0.607238,0.179779,0.499492}%
\pgfsetfillcolor{currentfill}%
\pgfsetfillopacity{0.800000}%
\pgfsetlinewidth{1.003750pt}%
\definecolor{currentstroke}{rgb}{0.607238,0.179779,0.499492}%
\pgfsetstrokecolor{currentstroke}%
\pgfsetstrokeopacity{0.800000}%
\pgfsetdash{}{0pt}%
\pgfpathmoveto{\pgfqpoint{2.007936in}{1.408522in}}%
\pgfpathcurveto{\pgfqpoint{2.013760in}{1.408522in}}{\pgfqpoint{2.019346in}{1.410836in}}{\pgfqpoint{2.023464in}{1.414954in}}%
\pgfpathcurveto{\pgfqpoint{2.027582in}{1.419072in}}{\pgfqpoint{2.029896in}{1.424658in}}{\pgfqpoint{2.029896in}{1.430482in}}%
\pgfpathcurveto{\pgfqpoint{2.029896in}{1.436306in}}{\pgfqpoint{2.027582in}{1.441892in}}{\pgfqpoint{2.023464in}{1.446010in}}%
\pgfpathcurveto{\pgfqpoint{2.019346in}{1.450129in}}{\pgfqpoint{2.013760in}{1.452442in}}{\pgfqpoint{2.007936in}{1.452442in}}%
\pgfpathcurveto{\pgfqpoint{2.002112in}{1.452442in}}{\pgfqpoint{1.996526in}{1.450129in}}{\pgfqpoint{1.992408in}{1.446010in}}%
\pgfpathcurveto{\pgfqpoint{1.988289in}{1.441892in}}{\pgfqpoint{1.985976in}{1.436306in}}{\pgfqpoint{1.985976in}{1.430482in}}%
\pgfpathcurveto{\pgfqpoint{1.985976in}{1.424658in}}{\pgfqpoint{1.988289in}{1.419072in}}{\pgfqpoint{1.992408in}{1.414954in}}%
\pgfpathcurveto{\pgfqpoint{1.996526in}{1.410836in}}{\pgfqpoint{2.002112in}{1.408522in}}{\pgfqpoint{2.007936in}{1.408522in}}%
\pgfpathlineto{\pgfqpoint{2.007936in}{1.408522in}}%
\pgfpathclose%
\pgfusepath{stroke,fill}%
\end{pgfscope}%
\begin{pgfscope}%
\pgfpathrectangle{\pgfqpoint{0.100000in}{0.209042in}}{\pgfqpoint{2.906250in}{2.906250in}}%
\pgfusepath{clip}%
\pgfsetbuttcap%
\pgfsetroundjoin%
\definecolor{currentfill}{rgb}{0.252220,0.059415,0.453248}%
\pgfsetfillcolor{currentfill}%
\pgfsetfillopacity{0.800000}%
\pgfsetlinewidth{1.003750pt}%
\definecolor{currentstroke}{rgb}{0.252220,0.059415,0.453248}%
\pgfsetstrokecolor{currentstroke}%
\pgfsetstrokeopacity{0.800000}%
\pgfsetdash{}{0pt}%
\pgfpathmoveto{\pgfqpoint{1.471729in}{1.822879in}}%
\pgfpathcurveto{\pgfqpoint{1.477553in}{1.822879in}}{\pgfqpoint{1.483139in}{1.825193in}}{\pgfqpoint{1.487257in}{1.829311in}}%
\pgfpathcurveto{\pgfqpoint{1.491375in}{1.833429in}}{\pgfqpoint{1.493689in}{1.839016in}}{\pgfqpoint{1.493689in}{1.844840in}}%
\pgfpathcurveto{\pgfqpoint{1.493689in}{1.850663in}}{\pgfqpoint{1.491375in}{1.856250in}}{\pgfqpoint{1.487257in}{1.860368in}}%
\pgfpathcurveto{\pgfqpoint{1.483139in}{1.864486in}}{\pgfqpoint{1.477553in}{1.866800in}}{\pgfqpoint{1.471729in}{1.866800in}}%
\pgfpathcurveto{\pgfqpoint{1.465905in}{1.866800in}}{\pgfqpoint{1.460319in}{1.864486in}}{\pgfqpoint{1.456201in}{1.860368in}}%
\pgfpathcurveto{\pgfqpoint{1.452083in}{1.856250in}}{\pgfqpoint{1.449769in}{1.850663in}}{\pgfqpoint{1.449769in}{1.844840in}}%
\pgfpathcurveto{\pgfqpoint{1.449769in}{1.839016in}}{\pgfqpoint{1.452083in}{1.833429in}}{\pgfqpoint{1.456201in}{1.829311in}}%
\pgfpathcurveto{\pgfqpoint{1.460319in}{1.825193in}}{\pgfqpoint{1.465905in}{1.822879in}}{\pgfqpoint{1.471729in}{1.822879in}}%
\pgfpathlineto{\pgfqpoint{1.471729in}{1.822879in}}%
\pgfpathclose%
\pgfusepath{stroke,fill}%
\end{pgfscope}%
\begin{pgfscope}%
\pgfpathrectangle{\pgfqpoint{0.100000in}{0.209042in}}{\pgfqpoint{2.906250in}{2.906250in}}%
\pgfusepath{clip}%
\pgfsetbuttcap%
\pgfsetroundjoin%
\definecolor{currentfill}{rgb}{0.341482,0.080564,0.492631}%
\pgfsetfillcolor{currentfill}%
\pgfsetfillopacity{0.800000}%
\pgfsetlinewidth{1.003750pt}%
\definecolor{currentstroke}{rgb}{0.341482,0.080564,0.492631}%
\pgfsetstrokecolor{currentstroke}%
\pgfsetstrokeopacity{0.800000}%
\pgfsetdash{}{0pt}%
\pgfpathmoveto{\pgfqpoint{1.496272in}{1.898783in}}%
\pgfpathcurveto{\pgfqpoint{1.502096in}{1.898783in}}{\pgfqpoint{1.507682in}{1.901097in}}{\pgfqpoint{1.511800in}{1.905215in}}%
\pgfpathcurveto{\pgfqpoint{1.515919in}{1.909333in}}{\pgfqpoint{1.518232in}{1.914920in}}{\pgfqpoint{1.518232in}{1.920744in}}%
\pgfpathcurveto{\pgfqpoint{1.518232in}{1.926567in}}{\pgfqpoint{1.515919in}{1.932154in}}{\pgfqpoint{1.511800in}{1.936272in}}%
\pgfpathcurveto{\pgfqpoint{1.507682in}{1.940390in}}{\pgfqpoint{1.502096in}{1.942704in}}{\pgfqpoint{1.496272in}{1.942704in}}%
\pgfpathcurveto{\pgfqpoint{1.490448in}{1.942704in}}{\pgfqpoint{1.484862in}{1.940390in}}{\pgfqpoint{1.480744in}{1.936272in}}%
\pgfpathcurveto{\pgfqpoint{1.476626in}{1.932154in}}{\pgfqpoint{1.474312in}{1.926567in}}{\pgfqpoint{1.474312in}{1.920744in}}%
\pgfpathcurveto{\pgfqpoint{1.474312in}{1.914920in}}{\pgfqpoint{1.476626in}{1.909333in}}{\pgfqpoint{1.480744in}{1.905215in}}%
\pgfpathcurveto{\pgfqpoint{1.484862in}{1.901097in}}{\pgfqpoint{1.490448in}{1.898783in}}{\pgfqpoint{1.496272in}{1.898783in}}%
\pgfpathlineto{\pgfqpoint{1.496272in}{1.898783in}}%
\pgfpathclose%
\pgfusepath{stroke,fill}%
\end{pgfscope}%
\begin{pgfscope}%
\pgfpathrectangle{\pgfqpoint{0.100000in}{0.209042in}}{\pgfqpoint{2.906250in}{2.906250in}}%
\pgfusepath{clip}%
\pgfsetbuttcap%
\pgfsetroundjoin%
\definecolor{currentfill}{rgb}{0.245543,0.059352,0.448436}%
\pgfsetfillcolor{currentfill}%
\pgfsetfillopacity{0.800000}%
\pgfsetlinewidth{1.003750pt}%
\definecolor{currentstroke}{rgb}{0.245543,0.059352,0.448436}%
\pgfsetstrokecolor{currentstroke}%
\pgfsetstrokeopacity{0.800000}%
\pgfsetdash{}{0pt}%
\pgfpathmoveto{\pgfqpoint{1.547509in}{1.532471in}}%
\pgfpathcurveto{\pgfqpoint{1.553333in}{1.532471in}}{\pgfqpoint{1.558919in}{1.534785in}}{\pgfqpoint{1.563037in}{1.538903in}}%
\pgfpathcurveto{\pgfqpoint{1.567155in}{1.543021in}}{\pgfqpoint{1.569469in}{1.548607in}}{\pgfqpoint{1.569469in}{1.554431in}}%
\pgfpathcurveto{\pgfqpoint{1.569469in}{1.560255in}}{\pgfqpoint{1.567155in}{1.565841in}}{\pgfqpoint{1.563037in}{1.569959in}}%
\pgfpathcurveto{\pgfqpoint{1.558919in}{1.574077in}}{\pgfqpoint{1.553333in}{1.576391in}}{\pgfqpoint{1.547509in}{1.576391in}}%
\pgfpathcurveto{\pgfqpoint{1.541685in}{1.576391in}}{\pgfqpoint{1.536099in}{1.574077in}}{\pgfqpoint{1.531981in}{1.569959in}}%
\pgfpathcurveto{\pgfqpoint{1.527862in}{1.565841in}}{\pgfqpoint{1.525549in}{1.560255in}}{\pgfqpoint{1.525549in}{1.554431in}}%
\pgfpathcurveto{\pgfqpoint{1.525549in}{1.548607in}}{\pgfqpoint{1.527862in}{1.543021in}}{\pgfqpoint{1.531981in}{1.538903in}}%
\pgfpathcurveto{\pgfqpoint{1.536099in}{1.534785in}}{\pgfqpoint{1.541685in}{1.532471in}}{\pgfqpoint{1.547509in}{1.532471in}}%
\pgfpathlineto{\pgfqpoint{1.547509in}{1.532471in}}%
\pgfpathclose%
\pgfusepath{stroke,fill}%
\end{pgfscope}%
\begin{pgfscope}%
\pgfpathrectangle{\pgfqpoint{0.100000in}{0.209042in}}{\pgfqpoint{2.906250in}{2.906250in}}%
\pgfusepath{clip}%
\pgfsetbuttcap%
\pgfsetroundjoin%
\definecolor{currentfill}{rgb}{0.372116,0.092816,0.499053}%
\pgfsetfillcolor{currentfill}%
\pgfsetfillopacity{0.800000}%
\pgfsetlinewidth{1.003750pt}%
\definecolor{currentstroke}{rgb}{0.372116,0.092816,0.499053}%
\pgfsetstrokecolor{currentstroke}%
\pgfsetstrokeopacity{0.800000}%
\pgfsetdash{}{0pt}%
\pgfpathmoveto{\pgfqpoint{1.955081in}{1.723255in}}%
\pgfpathcurveto{\pgfqpoint{1.960905in}{1.723255in}}{\pgfqpoint{1.966491in}{1.725569in}}{\pgfqpoint{1.970609in}{1.729687in}}%
\pgfpathcurveto{\pgfqpoint{1.974727in}{1.733805in}}{\pgfqpoint{1.977041in}{1.739392in}}{\pgfqpoint{1.977041in}{1.745216in}}%
\pgfpathcurveto{\pgfqpoint{1.977041in}{1.751039in}}{\pgfqpoint{1.974727in}{1.756626in}}{\pgfqpoint{1.970609in}{1.760744in}}%
\pgfpathcurveto{\pgfqpoint{1.966491in}{1.764862in}}{\pgfqpoint{1.960905in}{1.767176in}}{\pgfqpoint{1.955081in}{1.767176in}}%
\pgfpathcurveto{\pgfqpoint{1.949257in}{1.767176in}}{\pgfqpoint{1.943671in}{1.764862in}}{\pgfqpoint{1.939552in}{1.760744in}}%
\pgfpathcurveto{\pgfqpoint{1.935434in}{1.756626in}}{\pgfqpoint{1.933120in}{1.751039in}}{\pgfqpoint{1.933120in}{1.745216in}}%
\pgfpathcurveto{\pgfqpoint{1.933120in}{1.739392in}}{\pgfqpoint{1.935434in}{1.733805in}}{\pgfqpoint{1.939552in}{1.729687in}}%
\pgfpathcurveto{\pgfqpoint{1.943671in}{1.725569in}}{\pgfqpoint{1.949257in}{1.723255in}}{\pgfqpoint{1.955081in}{1.723255in}}%
\pgfpathlineto{\pgfqpoint{1.955081in}{1.723255in}}%
\pgfpathclose%
\pgfusepath{stroke,fill}%
\end{pgfscope}%
\begin{pgfscope}%
\pgfpathrectangle{\pgfqpoint{0.100000in}{0.209042in}}{\pgfqpoint{2.906250in}{2.906250in}}%
\pgfusepath{clip}%
\pgfsetbuttcap%
\pgfsetroundjoin%
\definecolor{currentfill}{rgb}{0.729216,0.219437,0.471279}%
\pgfsetfillcolor{currentfill}%
\pgfsetfillopacity{0.800000}%
\pgfsetlinewidth{1.003750pt}%
\definecolor{currentstroke}{rgb}{0.729216,0.219437,0.471279}%
\pgfsetstrokecolor{currentstroke}%
\pgfsetstrokeopacity{0.800000}%
\pgfsetdash{}{0pt}%
\pgfpathmoveto{\pgfqpoint{1.082014in}{1.694824in}}%
\pgfpathcurveto{\pgfqpoint{1.087838in}{1.694824in}}{\pgfqpoint{1.093424in}{1.697138in}}{\pgfqpoint{1.097542in}{1.701256in}}%
\pgfpathcurveto{\pgfqpoint{1.101660in}{1.705374in}}{\pgfqpoint{1.103974in}{1.710960in}}{\pgfqpoint{1.103974in}{1.716784in}}%
\pgfpathcurveto{\pgfqpoint{1.103974in}{1.722608in}}{\pgfqpoint{1.101660in}{1.728194in}}{\pgfqpoint{1.097542in}{1.732312in}}%
\pgfpathcurveto{\pgfqpoint{1.093424in}{1.736431in}}{\pgfqpoint{1.087838in}{1.738744in}}{\pgfqpoint{1.082014in}{1.738744in}}%
\pgfpathcurveto{\pgfqpoint{1.076190in}{1.738744in}}{\pgfqpoint{1.070604in}{1.736431in}}{\pgfqpoint{1.066485in}{1.732312in}}%
\pgfpathcurveto{\pgfqpoint{1.062367in}{1.728194in}}{\pgfqpoint{1.060053in}{1.722608in}}{\pgfqpoint{1.060053in}{1.716784in}}%
\pgfpathcurveto{\pgfqpoint{1.060053in}{1.710960in}}{\pgfqpoint{1.062367in}{1.705374in}}{\pgfqpoint{1.066485in}{1.701256in}}%
\pgfpathcurveto{\pgfqpoint{1.070604in}{1.697138in}}{\pgfqpoint{1.076190in}{1.694824in}}{\pgfqpoint{1.082014in}{1.694824in}}%
\pgfpathlineto{\pgfqpoint{1.082014in}{1.694824in}}%
\pgfpathclose%
\pgfusepath{stroke,fill}%
\end{pgfscope}%
\begin{pgfscope}%
\pgfpathrectangle{\pgfqpoint{0.100000in}{0.209042in}}{\pgfqpoint{2.906250in}{2.906250in}}%
\pgfusepath{clip}%
\pgfsetbuttcap%
\pgfsetroundjoin%
\definecolor{currentfill}{rgb}{0.304081,0.067835,0.480812}%
\pgfsetfillcolor{currentfill}%
\pgfsetfillopacity{0.800000}%
\pgfsetlinewidth{1.003750pt}%
\definecolor{currentstroke}{rgb}{0.304081,0.067835,0.480812}%
\pgfsetstrokecolor{currentstroke}%
\pgfsetstrokeopacity{0.800000}%
\pgfsetdash{}{0pt}%
\pgfpathmoveto{\pgfqpoint{1.905366in}{1.716981in}}%
\pgfpathcurveto{\pgfqpoint{1.911190in}{1.716981in}}{\pgfqpoint{1.916776in}{1.719295in}}{\pgfqpoint{1.920894in}{1.723413in}}%
\pgfpathcurveto{\pgfqpoint{1.925012in}{1.727531in}}{\pgfqpoint{1.927326in}{1.733117in}}{\pgfqpoint{1.927326in}{1.738941in}}%
\pgfpathcurveto{\pgfqpoint{1.927326in}{1.744765in}}{\pgfqpoint{1.925012in}{1.750351in}}{\pgfqpoint{1.920894in}{1.754470in}}%
\pgfpathcurveto{\pgfqpoint{1.916776in}{1.758588in}}{\pgfqpoint{1.911190in}{1.760902in}}{\pgfqpoint{1.905366in}{1.760902in}}%
\pgfpathcurveto{\pgfqpoint{1.899542in}{1.760902in}}{\pgfqpoint{1.893956in}{1.758588in}}{\pgfqpoint{1.889838in}{1.754470in}}%
\pgfpathcurveto{\pgfqpoint{1.885720in}{1.750351in}}{\pgfqpoint{1.883406in}{1.744765in}}{\pgfqpoint{1.883406in}{1.738941in}}%
\pgfpathcurveto{\pgfqpoint{1.883406in}{1.733117in}}{\pgfqpoint{1.885720in}{1.727531in}}{\pgfqpoint{1.889838in}{1.723413in}}%
\pgfpathcurveto{\pgfqpoint{1.893956in}{1.719295in}}{\pgfqpoint{1.899542in}{1.716981in}}{\pgfqpoint{1.905366in}{1.716981in}}%
\pgfpathlineto{\pgfqpoint{1.905366in}{1.716981in}}%
\pgfpathclose%
\pgfusepath{stroke,fill}%
\end{pgfscope}%
\begin{pgfscope}%
\pgfpathrectangle{\pgfqpoint{0.100000in}{0.209042in}}{\pgfqpoint{2.906250in}{2.906250in}}%
\pgfusepath{clip}%
\pgfsetbuttcap%
\pgfsetroundjoin%
\definecolor{currentfill}{rgb}{0.152839,0.068637,0.343404}%
\pgfsetfillcolor{currentfill}%
\pgfsetfillopacity{0.800000}%
\pgfsetlinewidth{1.003750pt}%
\definecolor{currentstroke}{rgb}{0.152839,0.068637,0.343404}%
\pgfsetstrokecolor{currentstroke}%
\pgfsetstrokeopacity{0.800000}%
\pgfsetdash{}{0pt}%
\pgfpathmoveto{\pgfqpoint{1.783522in}{1.714176in}}%
\pgfpathcurveto{\pgfqpoint{1.789346in}{1.714176in}}{\pgfqpoint{1.794932in}{1.716490in}}{\pgfqpoint{1.799050in}{1.720608in}}%
\pgfpathcurveto{\pgfqpoint{1.803168in}{1.724726in}}{\pgfqpoint{1.805482in}{1.730312in}}{\pgfqpoint{1.805482in}{1.736136in}}%
\pgfpathcurveto{\pgfqpoint{1.805482in}{1.741960in}}{\pgfqpoint{1.803168in}{1.747546in}}{\pgfqpoint{1.799050in}{1.751664in}}%
\pgfpathcurveto{\pgfqpoint{1.794932in}{1.755782in}}{\pgfqpoint{1.789346in}{1.758096in}}{\pgfqpoint{1.783522in}{1.758096in}}%
\pgfpathcurveto{\pgfqpoint{1.777698in}{1.758096in}}{\pgfqpoint{1.772111in}{1.755782in}}{\pgfqpoint{1.767993in}{1.751664in}}%
\pgfpathcurveto{\pgfqpoint{1.763875in}{1.747546in}}{\pgfqpoint{1.761561in}{1.741960in}}{\pgfqpoint{1.761561in}{1.736136in}}%
\pgfpathcurveto{\pgfqpoint{1.761561in}{1.730312in}}{\pgfqpoint{1.763875in}{1.724726in}}{\pgfqpoint{1.767993in}{1.720608in}}%
\pgfpathcurveto{\pgfqpoint{1.772111in}{1.716490in}}{\pgfqpoint{1.777698in}{1.714176in}}{\pgfqpoint{1.783522in}{1.714176in}}%
\pgfpathlineto{\pgfqpoint{1.783522in}{1.714176in}}%
\pgfpathclose%
\pgfusepath{stroke,fill}%
\end{pgfscope}%
\begin{pgfscope}%
\pgfpathrectangle{\pgfqpoint{0.100000in}{0.209042in}}{\pgfqpoint{2.906250in}{2.906250in}}%
\pgfusepath{clip}%
\pgfsetbuttcap%
\pgfsetroundjoin%
\definecolor{currentfill}{rgb}{0.544015,0.159033,0.506159}%
\pgfsetfillcolor{currentfill}%
\pgfsetfillopacity{0.800000}%
\pgfsetlinewidth{1.003750pt}%
\definecolor{currentstroke}{rgb}{0.544015,0.159033,0.506159}%
\pgfsetstrokecolor{currentstroke}%
\pgfsetstrokeopacity{0.800000}%
\pgfsetdash{}{0pt}%
\pgfpathmoveto{\pgfqpoint{1.226372in}{1.810994in}}%
\pgfpathcurveto{\pgfqpoint{1.232196in}{1.810994in}}{\pgfqpoint{1.237782in}{1.813308in}}{\pgfqpoint{1.241901in}{1.817426in}}%
\pgfpathcurveto{\pgfqpoint{1.246019in}{1.821544in}}{\pgfqpoint{1.248333in}{1.827130in}}{\pgfqpoint{1.248333in}{1.832954in}}%
\pgfpathcurveto{\pgfqpoint{1.248333in}{1.838778in}}{\pgfqpoint{1.246019in}{1.844365in}}{\pgfqpoint{1.241901in}{1.848483in}}%
\pgfpathcurveto{\pgfqpoint{1.237782in}{1.852601in}}{\pgfqpoint{1.232196in}{1.854915in}}{\pgfqpoint{1.226372in}{1.854915in}}%
\pgfpathcurveto{\pgfqpoint{1.220548in}{1.854915in}}{\pgfqpoint{1.214962in}{1.852601in}}{\pgfqpoint{1.210844in}{1.848483in}}%
\pgfpathcurveto{\pgfqpoint{1.206726in}{1.844365in}}{\pgfqpoint{1.204412in}{1.838778in}}{\pgfqpoint{1.204412in}{1.832954in}}%
\pgfpathcurveto{\pgfqpoint{1.204412in}{1.827130in}}{\pgfqpoint{1.206726in}{1.821544in}}{\pgfqpoint{1.210844in}{1.817426in}}%
\pgfpathcurveto{\pgfqpoint{1.214962in}{1.813308in}}{\pgfqpoint{1.220548in}{1.810994in}}{\pgfqpoint{1.226372in}{1.810994in}}%
\pgfpathlineto{\pgfqpoint{1.226372in}{1.810994in}}%
\pgfpathclose%
\pgfusepath{stroke,fill}%
\end{pgfscope}%
\begin{pgfscope}%
\pgfpathrectangle{\pgfqpoint{0.100000in}{0.209042in}}{\pgfqpoint{2.906250in}{2.906250in}}%
\pgfusepath{clip}%
\pgfsetbuttcap%
\pgfsetroundjoin%
\definecolor{currentfill}{rgb}{0.451271,0.125132,0.507198}%
\pgfsetfillcolor{currentfill}%
\pgfsetfillopacity{0.800000}%
\pgfsetlinewidth{1.003750pt}%
\definecolor{currentstroke}{rgb}{0.451271,0.125132,0.507198}%
\pgfsetstrokecolor{currentstroke}%
\pgfsetstrokeopacity{0.800000}%
\pgfsetdash{}{0pt}%
\pgfpathmoveto{\pgfqpoint{2.010069in}{1.612944in}}%
\pgfpathcurveto{\pgfqpoint{2.015893in}{1.612944in}}{\pgfqpoint{2.021479in}{1.615258in}}{\pgfqpoint{2.025597in}{1.619376in}}%
\pgfpathcurveto{\pgfqpoint{2.029715in}{1.623494in}}{\pgfqpoint{2.032029in}{1.629080in}}{\pgfqpoint{2.032029in}{1.634904in}}%
\pgfpathcurveto{\pgfqpoint{2.032029in}{1.640728in}}{\pgfqpoint{2.029715in}{1.646314in}}{\pgfqpoint{2.025597in}{1.650432in}}%
\pgfpathcurveto{\pgfqpoint{2.021479in}{1.654550in}}{\pgfqpoint{2.015893in}{1.656864in}}{\pgfqpoint{2.010069in}{1.656864in}}%
\pgfpathcurveto{\pgfqpoint{2.004245in}{1.656864in}}{\pgfqpoint{1.998659in}{1.654550in}}{\pgfqpoint{1.994540in}{1.650432in}}%
\pgfpathcurveto{\pgfqpoint{1.990422in}{1.646314in}}{\pgfqpoint{1.988108in}{1.640728in}}{\pgfqpoint{1.988108in}{1.634904in}}%
\pgfpathcurveto{\pgfqpoint{1.988108in}{1.629080in}}{\pgfqpoint{1.990422in}{1.623494in}}{\pgfqpoint{1.994540in}{1.619376in}}%
\pgfpathcurveto{\pgfqpoint{1.998659in}{1.615258in}}{\pgfqpoint{2.004245in}{1.612944in}}{\pgfqpoint{2.010069in}{1.612944in}}%
\pgfpathlineto{\pgfqpoint{2.010069in}{1.612944in}}%
\pgfpathclose%
\pgfusepath{stroke,fill}%
\end{pgfscope}%
\begin{pgfscope}%
\pgfpathrectangle{\pgfqpoint{0.100000in}{0.209042in}}{\pgfqpoint{2.906250in}{2.906250in}}%
\pgfusepath{clip}%
\pgfsetbuttcap%
\pgfsetroundjoin%
\definecolor{currentfill}{rgb}{0.353773,0.085373,0.495501}%
\pgfsetfillcolor{currentfill}%
\pgfsetfillopacity{0.800000}%
\pgfsetlinewidth{1.003750pt}%
\definecolor{currentstroke}{rgb}{0.353773,0.085373,0.495501}%
\pgfsetstrokecolor{currentstroke}%
\pgfsetstrokeopacity{0.800000}%
\pgfsetdash{}{0pt}%
\pgfpathmoveto{\pgfqpoint{1.676473in}{1.924827in}}%
\pgfpathcurveto{\pgfqpoint{1.682297in}{1.924827in}}{\pgfqpoint{1.687883in}{1.927141in}}{\pgfqpoint{1.692001in}{1.931259in}}%
\pgfpathcurveto{\pgfqpoint{1.696119in}{1.935377in}}{\pgfqpoint{1.698433in}{1.940963in}}{\pgfqpoint{1.698433in}{1.946787in}}%
\pgfpathcurveto{\pgfqpoint{1.698433in}{1.952611in}}{\pgfqpoint{1.696119in}{1.958197in}}{\pgfqpoint{1.692001in}{1.962315in}}%
\pgfpathcurveto{\pgfqpoint{1.687883in}{1.966433in}}{\pgfqpoint{1.682297in}{1.968747in}}{\pgfqpoint{1.676473in}{1.968747in}}%
\pgfpathcurveto{\pgfqpoint{1.670649in}{1.968747in}}{\pgfqpoint{1.665063in}{1.966433in}}{\pgfqpoint{1.660945in}{1.962315in}}%
\pgfpathcurveto{\pgfqpoint{1.656827in}{1.958197in}}{\pgfqpoint{1.654513in}{1.952611in}}{\pgfqpoint{1.654513in}{1.946787in}}%
\pgfpathcurveto{\pgfqpoint{1.654513in}{1.940963in}}{\pgfqpoint{1.656827in}{1.935377in}}{\pgfqpoint{1.660945in}{1.931259in}}%
\pgfpathcurveto{\pgfqpoint{1.665063in}{1.927141in}}{\pgfqpoint{1.670649in}{1.924827in}}{\pgfqpoint{1.676473in}{1.924827in}}%
\pgfpathlineto{\pgfqpoint{1.676473in}{1.924827in}}%
\pgfpathclose%
\pgfusepath{stroke,fill}%
\end{pgfscope}%
\begin{pgfscope}%
\pgfpathrectangle{\pgfqpoint{0.100000in}{0.209042in}}{\pgfqpoint{2.906250in}{2.906250in}}%
\pgfusepath{clip}%
\pgfsetbuttcap%
\pgfsetroundjoin%
\definecolor{currentfill}{rgb}{0.594508,0.175701,0.501241}%
\pgfsetfillcolor{currentfill}%
\pgfsetfillopacity{0.800000}%
\pgfsetlinewidth{1.003750pt}%
\definecolor{currentstroke}{rgb}{0.594508,0.175701,0.501241}%
\pgfsetstrokecolor{currentstroke}%
\pgfsetstrokeopacity{0.800000}%
\pgfsetdash{}{0pt}%
\pgfpathmoveto{\pgfqpoint{1.181951in}{1.779962in}}%
\pgfpathcurveto{\pgfqpoint{1.187775in}{1.779962in}}{\pgfqpoint{1.193361in}{1.782276in}}{\pgfqpoint{1.197479in}{1.786394in}}%
\pgfpathcurveto{\pgfqpoint{1.201597in}{1.790513in}}{\pgfqpoint{1.203911in}{1.796099in}}{\pgfqpoint{1.203911in}{1.801923in}}%
\pgfpathcurveto{\pgfqpoint{1.203911in}{1.807747in}}{\pgfqpoint{1.201597in}{1.813333in}}{\pgfqpoint{1.197479in}{1.817451in}}%
\pgfpathcurveto{\pgfqpoint{1.193361in}{1.821569in}}{\pgfqpoint{1.187775in}{1.823883in}}{\pgfqpoint{1.181951in}{1.823883in}}%
\pgfpathcurveto{\pgfqpoint{1.176127in}{1.823883in}}{\pgfqpoint{1.170541in}{1.821569in}}{\pgfqpoint{1.166423in}{1.817451in}}%
\pgfpathcurveto{\pgfqpoint{1.162305in}{1.813333in}}{\pgfqpoint{1.159991in}{1.807747in}}{\pgfqpoint{1.159991in}{1.801923in}}%
\pgfpathcurveto{\pgfqpoint{1.159991in}{1.796099in}}{\pgfqpoint{1.162305in}{1.790513in}}{\pgfqpoint{1.166423in}{1.786394in}}%
\pgfpathcurveto{\pgfqpoint{1.170541in}{1.782276in}}{\pgfqpoint{1.176127in}{1.779962in}}{\pgfqpoint{1.181951in}{1.779962in}}%
\pgfpathlineto{\pgfqpoint{1.181951in}{1.779962in}}%
\pgfpathclose%
\pgfusepath{stroke,fill}%
\end{pgfscope}%
\begin{pgfscope}%
\pgfpathrectangle{\pgfqpoint{0.100000in}{0.209042in}}{\pgfqpoint{2.906250in}{2.906250in}}%
\pgfusepath{clip}%
\pgfsetbuttcap%
\pgfsetroundjoin%
\definecolor{currentfill}{rgb}{0.353773,0.085373,0.495501}%
\pgfsetfillcolor{currentfill}%
\pgfsetfillopacity{0.800000}%
\pgfsetlinewidth{1.003750pt}%
\definecolor{currentstroke}{rgb}{0.353773,0.085373,0.495501}%
\pgfsetstrokecolor{currentstroke}%
\pgfsetstrokeopacity{0.800000}%
\pgfsetdash{}{0pt}%
\pgfpathmoveto{\pgfqpoint{1.946842in}{1.672522in}}%
\pgfpathcurveto{\pgfqpoint{1.952666in}{1.672522in}}{\pgfqpoint{1.958252in}{1.674836in}}{\pgfqpoint{1.962370in}{1.678954in}}%
\pgfpathcurveto{\pgfqpoint{1.966488in}{1.683073in}}{\pgfqpoint{1.968802in}{1.688659in}}{\pgfqpoint{1.968802in}{1.694483in}}%
\pgfpathcurveto{\pgfqpoint{1.968802in}{1.700307in}}{\pgfqpoint{1.966488in}{1.705893in}}{\pgfqpoint{1.962370in}{1.710011in}}%
\pgfpathcurveto{\pgfqpoint{1.958252in}{1.714129in}}{\pgfqpoint{1.952666in}{1.716443in}}{\pgfqpoint{1.946842in}{1.716443in}}%
\pgfpathcurveto{\pgfqpoint{1.941018in}{1.716443in}}{\pgfqpoint{1.935432in}{1.714129in}}{\pgfqpoint{1.931313in}{1.710011in}}%
\pgfpathcurveto{\pgfqpoint{1.927195in}{1.705893in}}{\pgfqpoint{1.924881in}{1.700307in}}{\pgfqpoint{1.924881in}{1.694483in}}%
\pgfpathcurveto{\pgfqpoint{1.924881in}{1.688659in}}{\pgfqpoint{1.927195in}{1.683073in}}{\pgfqpoint{1.931313in}{1.678954in}}%
\pgfpathcurveto{\pgfqpoint{1.935432in}{1.674836in}}{\pgfqpoint{1.941018in}{1.672522in}}{\pgfqpoint{1.946842in}{1.672522in}}%
\pgfpathlineto{\pgfqpoint{1.946842in}{1.672522in}}%
\pgfpathclose%
\pgfusepath{stroke,fill}%
\end{pgfscope}%
\begin{pgfscope}%
\pgfpathrectangle{\pgfqpoint{0.100000in}{0.209042in}}{\pgfqpoint{2.906250in}{2.906250in}}%
\pgfusepath{clip}%
\pgfsetbuttcap%
\pgfsetroundjoin%
\definecolor{currentfill}{rgb}{0.550287,0.161158,0.505719}%
\pgfsetfillcolor{currentfill}%
\pgfsetfillopacity{0.800000}%
\pgfsetlinewidth{1.003750pt}%
\definecolor{currentstroke}{rgb}{0.550287,0.161158,0.505719}%
\pgfsetstrokecolor{currentstroke}%
\pgfsetstrokeopacity{0.800000}%
\pgfsetdash{}{0pt}%
\pgfpathmoveto{\pgfqpoint{1.847458in}{1.997380in}}%
\pgfpathcurveto{\pgfqpoint{1.853282in}{1.997380in}}{\pgfqpoint{1.858868in}{1.999694in}}{\pgfqpoint{1.862987in}{2.003812in}}%
\pgfpathcurveto{\pgfqpoint{1.867105in}{2.007930in}}{\pgfqpoint{1.869419in}{2.013516in}}{\pgfqpoint{1.869419in}{2.019340in}}%
\pgfpathcurveto{\pgfqpoint{1.869419in}{2.025164in}}{\pgfqpoint{1.867105in}{2.030750in}}{\pgfqpoint{1.862987in}{2.034868in}}%
\pgfpathcurveto{\pgfqpoint{1.858868in}{2.038986in}}{\pgfqpoint{1.853282in}{2.041300in}}{\pgfqpoint{1.847458in}{2.041300in}}%
\pgfpathcurveto{\pgfqpoint{1.841634in}{2.041300in}}{\pgfqpoint{1.836048in}{2.038986in}}{\pgfqpoint{1.831930in}{2.034868in}}%
\pgfpathcurveto{\pgfqpoint{1.827812in}{2.030750in}}{\pgfqpoint{1.825498in}{2.025164in}}{\pgfqpoint{1.825498in}{2.019340in}}%
\pgfpathcurveto{\pgfqpoint{1.825498in}{2.013516in}}{\pgfqpoint{1.827812in}{2.007930in}}{\pgfqpoint{1.831930in}{2.003812in}}%
\pgfpathcurveto{\pgfqpoint{1.836048in}{1.999694in}}{\pgfqpoint{1.841634in}{1.997380in}}{\pgfqpoint{1.847458in}{1.997380in}}%
\pgfpathlineto{\pgfqpoint{1.847458in}{1.997380in}}%
\pgfpathclose%
\pgfusepath{stroke,fill}%
\end{pgfscope}%
\begin{pgfscope}%
\pgfpathrectangle{\pgfqpoint{0.100000in}{0.209042in}}{\pgfqpoint{2.906250in}{2.906250in}}%
\pgfusepath{clip}%
\pgfsetbuttcap%
\pgfsetroundjoin%
\definecolor{currentfill}{rgb}{0.863320,0.283729,0.412403}%
\pgfsetfillcolor{currentfill}%
\pgfsetfillopacity{0.800000}%
\pgfsetlinewidth{1.003750pt}%
\definecolor{currentstroke}{rgb}{0.863320,0.283729,0.412403}%
\pgfsetstrokecolor{currentstroke}%
\pgfsetstrokeopacity{0.800000}%
\pgfsetdash{}{0pt}%
\pgfpathmoveto{\pgfqpoint{1.472595in}{1.181064in}}%
\pgfpathcurveto{\pgfqpoint{1.478419in}{1.181064in}}{\pgfqpoint{1.484005in}{1.183378in}}{\pgfqpoint{1.488124in}{1.187496in}}%
\pgfpathcurveto{\pgfqpoint{1.492242in}{1.191614in}}{\pgfqpoint{1.494556in}{1.197201in}}{\pgfqpoint{1.494556in}{1.203025in}}%
\pgfpathcurveto{\pgfqpoint{1.494556in}{1.208848in}}{\pgfqpoint{1.492242in}{1.214435in}}{\pgfqpoint{1.488124in}{1.218553in}}%
\pgfpathcurveto{\pgfqpoint{1.484005in}{1.222671in}}{\pgfqpoint{1.478419in}{1.224985in}}{\pgfqpoint{1.472595in}{1.224985in}}%
\pgfpathcurveto{\pgfqpoint{1.466771in}{1.224985in}}{\pgfqpoint{1.461185in}{1.222671in}}{\pgfqpoint{1.457067in}{1.218553in}}%
\pgfpathcurveto{\pgfqpoint{1.452949in}{1.214435in}}{\pgfqpoint{1.450635in}{1.208848in}}{\pgfqpoint{1.450635in}{1.203025in}}%
\pgfpathcurveto{\pgfqpoint{1.450635in}{1.197201in}}{\pgfqpoint{1.452949in}{1.191614in}}{\pgfqpoint{1.457067in}{1.187496in}}%
\pgfpathcurveto{\pgfqpoint{1.461185in}{1.183378in}}{\pgfqpoint{1.466771in}{1.181064in}}{\pgfqpoint{1.472595in}{1.181064in}}%
\pgfpathlineto{\pgfqpoint{1.472595in}{1.181064in}}%
\pgfpathclose%
\pgfusepath{stroke,fill}%
\end{pgfscope}%
\begin{pgfscope}%
\pgfpathrectangle{\pgfqpoint{0.100000in}{0.209042in}}{\pgfqpoint{2.906250in}{2.906250in}}%
\pgfusepath{clip}%
\pgfsetbuttcap%
\pgfsetroundjoin%
\definecolor{currentfill}{rgb}{0.426877,0.115395,0.505714}%
\pgfsetfillcolor{currentfill}%
\pgfsetfillopacity{0.800000}%
\pgfsetlinewidth{1.003750pt}%
\definecolor{currentstroke}{rgb}{0.426877,0.115395,0.505714}%
\pgfsetstrokecolor{currentstroke}%
\pgfsetstrokeopacity{0.800000}%
\pgfsetdash{}{0pt}%
\pgfpathmoveto{\pgfqpoint{1.567458in}{1.408366in}}%
\pgfpathcurveto{\pgfqpoint{1.573282in}{1.408366in}}{\pgfqpoint{1.578868in}{1.410680in}}{\pgfqpoint{1.582986in}{1.414798in}}%
\pgfpathcurveto{\pgfqpoint{1.587104in}{1.418916in}}{\pgfqpoint{1.589418in}{1.424502in}}{\pgfqpoint{1.589418in}{1.430326in}}%
\pgfpathcurveto{\pgfqpoint{1.589418in}{1.436150in}}{\pgfqpoint{1.587104in}{1.441736in}}{\pgfqpoint{1.582986in}{1.445855in}}%
\pgfpathcurveto{\pgfqpoint{1.578868in}{1.449973in}}{\pgfqpoint{1.573282in}{1.452287in}}{\pgfqpoint{1.567458in}{1.452287in}}%
\pgfpathcurveto{\pgfqpoint{1.561634in}{1.452287in}}{\pgfqpoint{1.556048in}{1.449973in}}{\pgfqpoint{1.551929in}{1.445855in}}%
\pgfpathcurveto{\pgfqpoint{1.547811in}{1.441736in}}{\pgfqpoint{1.545497in}{1.436150in}}{\pgfqpoint{1.545497in}{1.430326in}}%
\pgfpathcurveto{\pgfqpoint{1.545497in}{1.424502in}}{\pgfqpoint{1.547811in}{1.418916in}}{\pgfqpoint{1.551929in}{1.414798in}}%
\pgfpathcurveto{\pgfqpoint{1.556048in}{1.410680in}}{\pgfqpoint{1.561634in}{1.408366in}}{\pgfqpoint{1.567458in}{1.408366in}}%
\pgfpathlineto{\pgfqpoint{1.567458in}{1.408366in}}%
\pgfpathclose%
\pgfusepath{stroke,fill}%
\end{pgfscope}%
\begin{pgfscope}%
\pgfpathrectangle{\pgfqpoint{0.100000in}{0.209042in}}{\pgfqpoint{2.906250in}{2.906250in}}%
\pgfusepath{clip}%
\pgfsetbuttcap%
\pgfsetroundjoin%
\definecolor{currentfill}{rgb}{0.469640,0.132245,0.507809}%
\pgfsetfillcolor{currentfill}%
\pgfsetfillopacity{0.800000}%
\pgfsetlinewidth{1.003750pt}%
\definecolor{currentstroke}{rgb}{0.469640,0.132245,0.507809}%
\pgfsetstrokecolor{currentstroke}%
\pgfsetstrokeopacity{0.800000}%
\pgfsetdash{}{0pt}%
\pgfpathmoveto{\pgfqpoint{2.028708in}{1.641652in}}%
\pgfpathcurveto{\pgfqpoint{2.034532in}{1.641652in}}{\pgfqpoint{2.040118in}{1.643966in}}{\pgfqpoint{2.044237in}{1.648084in}}%
\pgfpathcurveto{\pgfqpoint{2.048355in}{1.652202in}}{\pgfqpoint{2.050669in}{1.657788in}}{\pgfqpoint{2.050669in}{1.663612in}}%
\pgfpathcurveto{\pgfqpoint{2.050669in}{1.669436in}}{\pgfqpoint{2.048355in}{1.675022in}}{\pgfqpoint{2.044237in}{1.679140in}}%
\pgfpathcurveto{\pgfqpoint{2.040118in}{1.683259in}}{\pgfqpoint{2.034532in}{1.685572in}}{\pgfqpoint{2.028708in}{1.685572in}}%
\pgfpathcurveto{\pgfqpoint{2.022884in}{1.685572in}}{\pgfqpoint{2.017298in}{1.683259in}}{\pgfqpoint{2.013180in}{1.679140in}}%
\pgfpathcurveto{\pgfqpoint{2.009062in}{1.675022in}}{\pgfqpoint{2.006748in}{1.669436in}}{\pgfqpoint{2.006748in}{1.663612in}}%
\pgfpathcurveto{\pgfqpoint{2.006748in}{1.657788in}}{\pgfqpoint{2.009062in}{1.652202in}}{\pgfqpoint{2.013180in}{1.648084in}}%
\pgfpathcurveto{\pgfqpoint{2.017298in}{1.643966in}}{\pgfqpoint{2.022884in}{1.641652in}}{\pgfqpoint{2.028708in}{1.641652in}}%
\pgfpathlineto{\pgfqpoint{2.028708in}{1.641652in}}%
\pgfpathclose%
\pgfusepath{stroke,fill}%
\end{pgfscope}%
\begin{pgfscope}%
\pgfpathrectangle{\pgfqpoint{0.100000in}{0.209042in}}{\pgfqpoint{2.906250in}{2.906250in}}%
\pgfusepath{clip}%
\pgfsetbuttcap%
\pgfsetroundjoin%
\definecolor{currentfill}{rgb}{0.500438,0.143719,0.507920}%
\pgfsetfillcolor{currentfill}%
\pgfsetfillopacity{0.800000}%
\pgfsetlinewidth{1.003750pt}%
\definecolor{currentstroke}{rgb}{0.500438,0.143719,0.507920}%
\pgfsetstrokecolor{currentstroke}%
\pgfsetstrokeopacity{0.800000}%
\pgfsetdash{}{0pt}%
\pgfpathmoveto{\pgfqpoint{1.388619in}{1.442360in}}%
\pgfpathcurveto{\pgfqpoint{1.394443in}{1.442360in}}{\pgfqpoint{1.400029in}{1.444674in}}{\pgfqpoint{1.404147in}{1.448792in}}%
\pgfpathcurveto{\pgfqpoint{1.408265in}{1.452910in}}{\pgfqpoint{1.410579in}{1.458496in}}{\pgfqpoint{1.410579in}{1.464320in}}%
\pgfpathcurveto{\pgfqpoint{1.410579in}{1.470144in}}{\pgfqpoint{1.408265in}{1.475730in}}{\pgfqpoint{1.404147in}{1.479849in}}%
\pgfpathcurveto{\pgfqpoint{1.400029in}{1.483967in}}{\pgfqpoint{1.394443in}{1.486281in}}{\pgfqpoint{1.388619in}{1.486281in}}%
\pgfpathcurveto{\pgfqpoint{1.382795in}{1.486281in}}{\pgfqpoint{1.377209in}{1.483967in}}{\pgfqpoint{1.373090in}{1.479849in}}%
\pgfpathcurveto{\pgfqpoint{1.368972in}{1.475730in}}{\pgfqpoint{1.366658in}{1.470144in}}{\pgfqpoint{1.366658in}{1.464320in}}%
\pgfpathcurveto{\pgfqpoint{1.366658in}{1.458496in}}{\pgfqpoint{1.368972in}{1.452910in}}{\pgfqpoint{1.373090in}{1.448792in}}%
\pgfpathcurveto{\pgfqpoint{1.377209in}{1.444674in}}{\pgfqpoint{1.382795in}{1.442360in}}{\pgfqpoint{1.388619in}{1.442360in}}%
\pgfpathlineto{\pgfqpoint{1.388619in}{1.442360in}}%
\pgfpathclose%
\pgfusepath{stroke,fill}%
\end{pgfscope}%
\begin{pgfscope}%
\pgfpathrectangle{\pgfqpoint{0.100000in}{0.209042in}}{\pgfqpoint{2.906250in}{2.906250in}}%
\pgfusepath{clip}%
\pgfsetbuttcap%
\pgfsetroundjoin%
\definecolor{currentfill}{rgb}{0.645633,0.191952,0.492910}%
\pgfsetfillcolor{currentfill}%
\pgfsetfillopacity{0.800000}%
\pgfsetlinewidth{1.003750pt}%
\definecolor{currentstroke}{rgb}{0.645633,0.191952,0.492910}%
\pgfsetstrokecolor{currentstroke}%
\pgfsetstrokeopacity{0.800000}%
\pgfsetdash{}{0pt}%
\pgfpathmoveto{\pgfqpoint{2.150834in}{1.618310in}}%
\pgfpathcurveto{\pgfqpoint{2.156658in}{1.618310in}}{\pgfqpoint{2.162245in}{1.620624in}}{\pgfqpoint{2.166363in}{1.624742in}}%
\pgfpathcurveto{\pgfqpoint{2.170481in}{1.628860in}}{\pgfqpoint{2.172795in}{1.634447in}}{\pgfqpoint{2.172795in}{1.640271in}}%
\pgfpathcurveto{\pgfqpoint{2.172795in}{1.646095in}}{\pgfqpoint{2.170481in}{1.651681in}}{\pgfqpoint{2.166363in}{1.655799in}}%
\pgfpathcurveto{\pgfqpoint{2.162245in}{1.659917in}}{\pgfqpoint{2.156658in}{1.662231in}}{\pgfqpoint{2.150834in}{1.662231in}}%
\pgfpathcurveto{\pgfqpoint{2.145010in}{1.662231in}}{\pgfqpoint{2.139424in}{1.659917in}}{\pgfqpoint{2.135306in}{1.655799in}}%
\pgfpathcurveto{\pgfqpoint{2.131188in}{1.651681in}}{\pgfqpoint{2.128874in}{1.646095in}}{\pgfqpoint{2.128874in}{1.640271in}}%
\pgfpathcurveto{\pgfqpoint{2.128874in}{1.634447in}}{\pgfqpoint{2.131188in}{1.628860in}}{\pgfqpoint{2.135306in}{1.624742in}}%
\pgfpathcurveto{\pgfqpoint{2.139424in}{1.620624in}}{\pgfqpoint{2.145010in}{1.618310in}}{\pgfqpoint{2.150834in}{1.618310in}}%
\pgfpathlineto{\pgfqpoint{2.150834in}{1.618310in}}%
\pgfpathclose%
\pgfusepath{stroke,fill}%
\end{pgfscope}%
\begin{pgfscope}%
\pgfpathrectangle{\pgfqpoint{0.100000in}{0.209042in}}{\pgfqpoint{2.906250in}{2.906250in}}%
\pgfusepath{clip}%
\pgfsetbuttcap%
\pgfsetroundjoin%
\definecolor{currentfill}{rgb}{0.481929,0.136891,0.507989}%
\pgfsetfillcolor{currentfill}%
\pgfsetfillopacity{0.800000}%
\pgfsetlinewidth{1.003750pt}%
\definecolor{currentstroke}{rgb}{0.481929,0.136891,0.507989}%
\pgfsetstrokecolor{currentstroke}%
\pgfsetstrokeopacity{0.800000}%
\pgfsetdash{}{0pt}%
\pgfpathmoveto{\pgfqpoint{2.027973in}{1.728664in}}%
\pgfpathcurveto{\pgfqpoint{2.033797in}{1.728664in}}{\pgfqpoint{2.039383in}{1.730978in}}{\pgfqpoint{2.043501in}{1.735096in}}%
\pgfpathcurveto{\pgfqpoint{2.047619in}{1.739214in}}{\pgfqpoint{2.049933in}{1.744800in}}{\pgfqpoint{2.049933in}{1.750624in}}%
\pgfpathcurveto{\pgfqpoint{2.049933in}{1.756448in}}{\pgfqpoint{2.047619in}{1.762034in}}{\pgfqpoint{2.043501in}{1.766153in}}%
\pgfpathcurveto{\pgfqpoint{2.039383in}{1.770271in}}{\pgfqpoint{2.033797in}{1.772585in}}{\pgfqpoint{2.027973in}{1.772585in}}%
\pgfpathcurveto{\pgfqpoint{2.022149in}{1.772585in}}{\pgfqpoint{2.016563in}{1.770271in}}{\pgfqpoint{2.012445in}{1.766153in}}%
\pgfpathcurveto{\pgfqpoint{2.008327in}{1.762034in}}{\pgfqpoint{2.006013in}{1.756448in}}{\pgfqpoint{2.006013in}{1.750624in}}%
\pgfpathcurveto{\pgfqpoint{2.006013in}{1.744800in}}{\pgfqpoint{2.008327in}{1.739214in}}{\pgfqpoint{2.012445in}{1.735096in}}%
\pgfpathcurveto{\pgfqpoint{2.016563in}{1.730978in}}{\pgfqpoint{2.022149in}{1.728664in}}{\pgfqpoint{2.027973in}{1.728664in}}%
\pgfpathlineto{\pgfqpoint{2.027973in}{1.728664in}}%
\pgfpathclose%
\pgfusepath{stroke,fill}%
\end{pgfscope}%
\begin{pgfscope}%
\pgfpathrectangle{\pgfqpoint{0.100000in}{0.209042in}}{\pgfqpoint{2.906250in}{2.906250in}}%
\pgfusepath{clip}%
\pgfsetbuttcap%
\pgfsetroundjoin%
\definecolor{currentfill}{rgb}{0.384299,0.097855,0.501002}%
\pgfsetfillcolor{currentfill}%
\pgfsetfillopacity{0.800000}%
\pgfsetlinewidth{1.003750pt}%
\definecolor{currentstroke}{rgb}{0.384299,0.097855,0.501002}%
\pgfsetstrokecolor{currentstroke}%
\pgfsetstrokeopacity{0.800000}%
\pgfsetdash{}{0pt}%
\pgfpathmoveto{\pgfqpoint{1.445496in}{1.493440in}}%
\pgfpathcurveto{\pgfqpoint{1.451320in}{1.493440in}}{\pgfqpoint{1.456906in}{1.495754in}}{\pgfqpoint{1.461024in}{1.499872in}}%
\pgfpathcurveto{\pgfqpoint{1.465143in}{1.503990in}}{\pgfqpoint{1.467456in}{1.509576in}}{\pgfqpoint{1.467456in}{1.515400in}}%
\pgfpathcurveto{\pgfqpoint{1.467456in}{1.521224in}}{\pgfqpoint{1.465143in}{1.526810in}}{\pgfqpoint{1.461024in}{1.530928in}}%
\pgfpathcurveto{\pgfqpoint{1.456906in}{1.535046in}}{\pgfqpoint{1.451320in}{1.537360in}}{\pgfqpoint{1.445496in}{1.537360in}}%
\pgfpathcurveto{\pgfqpoint{1.439672in}{1.537360in}}{\pgfqpoint{1.434086in}{1.535046in}}{\pgfqpoint{1.429968in}{1.530928in}}%
\pgfpathcurveto{\pgfqpoint{1.425850in}{1.526810in}}{\pgfqpoint{1.423536in}{1.521224in}}{\pgfqpoint{1.423536in}{1.515400in}}%
\pgfpathcurveto{\pgfqpoint{1.423536in}{1.509576in}}{\pgfqpoint{1.425850in}{1.503990in}}{\pgfqpoint{1.429968in}{1.499872in}}%
\pgfpathcurveto{\pgfqpoint{1.434086in}{1.495754in}}{\pgfqpoint{1.439672in}{1.493440in}}{\pgfqpoint{1.445496in}{1.493440in}}%
\pgfpathlineto{\pgfqpoint{1.445496in}{1.493440in}}%
\pgfpathclose%
\pgfusepath{stroke,fill}%
\end{pgfscope}%
\begin{pgfscope}%
\pgfpathrectangle{\pgfqpoint{0.100000in}{0.209042in}}{\pgfqpoint{2.906250in}{2.906250in}}%
\pgfusepath{clip}%
\pgfsetbuttcap%
\pgfsetroundjoin%
\definecolor{currentfill}{rgb}{0.519045,0.150383,0.507443}%
\pgfsetfillcolor{currentfill}%
\pgfsetfillopacity{0.800000}%
\pgfsetlinewidth{1.003750pt}%
\definecolor{currentstroke}{rgb}{0.519045,0.150383,0.507443}%
\pgfsetstrokecolor{currentstroke}%
\pgfsetstrokeopacity{0.800000}%
\pgfsetdash{}{0pt}%
\pgfpathmoveto{\pgfqpoint{1.240744in}{1.633071in}}%
\pgfpathcurveto{\pgfqpoint{1.246568in}{1.633071in}}{\pgfqpoint{1.252154in}{1.635385in}}{\pgfqpoint{1.256272in}{1.639503in}}%
\pgfpathcurveto{\pgfqpoint{1.260390in}{1.643621in}}{\pgfqpoint{1.262704in}{1.649207in}}{\pgfqpoint{1.262704in}{1.655031in}}%
\pgfpathcurveto{\pgfqpoint{1.262704in}{1.660855in}}{\pgfqpoint{1.260390in}{1.666441in}}{\pgfqpoint{1.256272in}{1.670559in}}%
\pgfpathcurveto{\pgfqpoint{1.252154in}{1.674678in}}{\pgfqpoint{1.246568in}{1.676991in}}{\pgfqpoint{1.240744in}{1.676991in}}%
\pgfpathcurveto{\pgfqpoint{1.234920in}{1.676991in}}{\pgfqpoint{1.229334in}{1.674678in}}{\pgfqpoint{1.225215in}{1.670559in}}%
\pgfpathcurveto{\pgfqpoint{1.221097in}{1.666441in}}{\pgfqpoint{1.218783in}{1.660855in}}{\pgfqpoint{1.218783in}{1.655031in}}%
\pgfpathcurveto{\pgfqpoint{1.218783in}{1.649207in}}{\pgfqpoint{1.221097in}{1.643621in}}{\pgfqpoint{1.225215in}{1.639503in}}%
\pgfpathcurveto{\pgfqpoint{1.229334in}{1.635385in}}{\pgfqpoint{1.234920in}{1.633071in}}{\pgfqpoint{1.240744in}{1.633071in}}%
\pgfpathlineto{\pgfqpoint{1.240744in}{1.633071in}}%
\pgfpathclose%
\pgfusepath{stroke,fill}%
\end{pgfscope}%
\begin{pgfscope}%
\pgfpathrectangle{\pgfqpoint{0.100000in}{0.209042in}}{\pgfqpoint{2.906250in}{2.906250in}}%
\pgfusepath{clip}%
\pgfsetbuttcap%
\pgfsetroundjoin%
\definecolor{currentfill}{rgb}{0.297740,0.066117,0.478243}%
\pgfsetfillcolor{currentfill}%
\pgfsetfillopacity{0.800000}%
\pgfsetlinewidth{1.003750pt}%
\definecolor{currentstroke}{rgb}{0.297740,0.066117,0.478243}%
\pgfsetstrokecolor{currentstroke}%
\pgfsetstrokeopacity{0.800000}%
\pgfsetdash{}{0pt}%
\pgfpathmoveto{\pgfqpoint{1.662587in}{1.887547in}}%
\pgfpathcurveto{\pgfqpoint{1.668411in}{1.887547in}}{\pgfqpoint{1.673997in}{1.889861in}}{\pgfqpoint{1.678115in}{1.893979in}}%
\pgfpathcurveto{\pgfqpoint{1.682233in}{1.898097in}}{\pgfqpoint{1.684547in}{1.903683in}}{\pgfqpoint{1.684547in}{1.909507in}}%
\pgfpathcurveto{\pgfqpoint{1.684547in}{1.915331in}}{\pgfqpoint{1.682233in}{1.920917in}}{\pgfqpoint{1.678115in}{1.925036in}}%
\pgfpathcurveto{\pgfqpoint{1.673997in}{1.929154in}}{\pgfqpoint{1.668411in}{1.931468in}}{\pgfqpoint{1.662587in}{1.931468in}}%
\pgfpathcurveto{\pgfqpoint{1.656763in}{1.931468in}}{\pgfqpoint{1.651177in}{1.929154in}}{\pgfqpoint{1.647058in}{1.925036in}}%
\pgfpathcurveto{\pgfqpoint{1.642940in}{1.920917in}}{\pgfqpoint{1.640626in}{1.915331in}}{\pgfqpoint{1.640626in}{1.909507in}}%
\pgfpathcurveto{\pgfqpoint{1.640626in}{1.903683in}}{\pgfqpoint{1.642940in}{1.898097in}}{\pgfqpoint{1.647058in}{1.893979in}}%
\pgfpathcurveto{\pgfqpoint{1.651177in}{1.889861in}}{\pgfqpoint{1.656763in}{1.887547in}}{\pgfqpoint{1.662587in}{1.887547in}}%
\pgfpathlineto{\pgfqpoint{1.662587in}{1.887547in}}%
\pgfpathclose%
\pgfusepath{stroke,fill}%
\end{pgfscope}%
\begin{pgfscope}%
\pgfpathrectangle{\pgfqpoint{0.100000in}{0.209042in}}{\pgfqpoint{2.906250in}{2.906250in}}%
\pgfusepath{clip}%
\pgfsetbuttcap%
\pgfsetroundjoin%
\definecolor{currentfill}{rgb}{0.451271,0.125132,0.507198}%
\pgfsetfillcolor{currentfill}%
\pgfsetfillopacity{0.800000}%
\pgfsetlinewidth{1.003750pt}%
\definecolor{currentstroke}{rgb}{0.451271,0.125132,0.507198}%
\pgfsetstrokecolor{currentstroke}%
\pgfsetstrokeopacity{0.800000}%
\pgfsetdash{}{0pt}%
\pgfpathmoveto{\pgfqpoint{2.015391in}{1.686399in}}%
\pgfpathcurveto{\pgfqpoint{2.021215in}{1.686399in}}{\pgfqpoint{2.026801in}{1.688713in}}{\pgfqpoint{2.030920in}{1.692831in}}%
\pgfpathcurveto{\pgfqpoint{2.035038in}{1.696949in}}{\pgfqpoint{2.037352in}{1.702536in}}{\pgfqpoint{2.037352in}{1.708360in}}%
\pgfpathcurveto{\pgfqpoint{2.037352in}{1.714183in}}{\pgfqpoint{2.035038in}{1.719770in}}{\pgfqpoint{2.030920in}{1.723888in}}%
\pgfpathcurveto{\pgfqpoint{2.026801in}{1.728006in}}{\pgfqpoint{2.021215in}{1.730320in}}{\pgfqpoint{2.015391in}{1.730320in}}%
\pgfpathcurveto{\pgfqpoint{2.009567in}{1.730320in}}{\pgfqpoint{2.003981in}{1.728006in}}{\pgfqpoint{1.999863in}{1.723888in}}%
\pgfpathcurveto{\pgfqpoint{1.995745in}{1.719770in}}{\pgfqpoint{1.993431in}{1.714183in}}{\pgfqpoint{1.993431in}{1.708360in}}%
\pgfpathcurveto{\pgfqpoint{1.993431in}{1.702536in}}{\pgfqpoint{1.995745in}{1.696949in}}{\pgfqpoint{1.999863in}{1.692831in}}%
\pgfpathcurveto{\pgfqpoint{2.003981in}{1.688713in}}{\pgfqpoint{2.009567in}{1.686399in}}{\pgfqpoint{2.015391in}{1.686399in}}%
\pgfpathlineto{\pgfqpoint{2.015391in}{1.686399in}}%
\pgfpathclose%
\pgfusepath{stroke,fill}%
\end{pgfscope}%
\begin{pgfscope}%
\pgfpathrectangle{\pgfqpoint{0.100000in}{0.209042in}}{\pgfqpoint{2.906250in}{2.906250in}}%
\pgfusepath{clip}%
\pgfsetbuttcap%
\pgfsetroundjoin%
\definecolor{currentfill}{rgb}{0.384299,0.097855,0.501002}%
\pgfsetfillcolor{currentfill}%
\pgfsetfillopacity{0.800000}%
\pgfsetlinewidth{1.003750pt}%
\definecolor{currentstroke}{rgb}{0.384299,0.097855,0.501002}%
\pgfsetstrokecolor{currentstroke}%
\pgfsetstrokeopacity{0.800000}%
\pgfsetdash{}{0pt}%
\pgfpathmoveto{\pgfqpoint{1.936019in}{1.561911in}}%
\pgfpathcurveto{\pgfqpoint{1.941843in}{1.561911in}}{\pgfqpoint{1.947429in}{1.564225in}}{\pgfqpoint{1.951547in}{1.568343in}}%
\pgfpathcurveto{\pgfqpoint{1.955665in}{1.572462in}}{\pgfqpoint{1.957979in}{1.578048in}}{\pgfqpoint{1.957979in}{1.583872in}}%
\pgfpathcurveto{\pgfqpoint{1.957979in}{1.589696in}}{\pgfqpoint{1.955665in}{1.595282in}}{\pgfqpoint{1.951547in}{1.599400in}}%
\pgfpathcurveto{\pgfqpoint{1.947429in}{1.603518in}}{\pgfqpoint{1.941843in}{1.605832in}}{\pgfqpoint{1.936019in}{1.605832in}}%
\pgfpathcurveto{\pgfqpoint{1.930195in}{1.605832in}}{\pgfqpoint{1.924609in}{1.603518in}}{\pgfqpoint{1.920490in}{1.599400in}}%
\pgfpathcurveto{\pgfqpoint{1.916372in}{1.595282in}}{\pgfqpoint{1.914058in}{1.589696in}}{\pgfqpoint{1.914058in}{1.583872in}}%
\pgfpathcurveto{\pgfqpoint{1.914058in}{1.578048in}}{\pgfqpoint{1.916372in}{1.572462in}}{\pgfqpoint{1.920490in}{1.568343in}}%
\pgfpathcurveto{\pgfqpoint{1.924609in}{1.564225in}}{\pgfqpoint{1.930195in}{1.561911in}}{\pgfqpoint{1.936019in}{1.561911in}}%
\pgfpathlineto{\pgfqpoint{1.936019in}{1.561911in}}%
\pgfpathclose%
\pgfusepath{stroke,fill}%
\end{pgfscope}%
\begin{pgfscope}%
\pgfpathrectangle{\pgfqpoint{0.100000in}{0.209042in}}{\pgfqpoint{2.906250in}{2.906250in}}%
\pgfusepath{clip}%
\pgfsetbuttcap%
\pgfsetroundjoin%
\definecolor{currentfill}{rgb}{0.671349,0.200133,0.487358}%
\pgfsetfillcolor{currentfill}%
\pgfsetfillopacity{0.800000}%
\pgfsetlinewidth{1.003750pt}%
\definecolor{currentstroke}{rgb}{0.671349,0.200133,0.487358}%
\pgfsetstrokecolor{currentstroke}%
\pgfsetstrokeopacity{0.800000}%
\pgfsetdash{}{0pt}%
\pgfpathmoveto{\pgfqpoint{1.858799in}{1.287909in}}%
\pgfpathcurveto{\pgfqpoint{1.864623in}{1.287909in}}{\pgfqpoint{1.870209in}{1.290223in}}{\pgfqpoint{1.874328in}{1.294341in}}%
\pgfpathcurveto{\pgfqpoint{1.878446in}{1.298460in}}{\pgfqpoint{1.880760in}{1.304046in}}{\pgfqpoint{1.880760in}{1.309870in}}%
\pgfpathcurveto{\pgfqpoint{1.880760in}{1.315694in}}{\pgfqpoint{1.878446in}{1.321280in}}{\pgfqpoint{1.874328in}{1.325398in}}%
\pgfpathcurveto{\pgfqpoint{1.870209in}{1.329516in}}{\pgfqpoint{1.864623in}{1.331830in}}{\pgfqpoint{1.858799in}{1.331830in}}%
\pgfpathcurveto{\pgfqpoint{1.852975in}{1.331830in}}{\pgfqpoint{1.847389in}{1.329516in}}{\pgfqpoint{1.843271in}{1.325398in}}%
\pgfpathcurveto{\pgfqpoint{1.839153in}{1.321280in}}{\pgfqpoint{1.836839in}{1.315694in}}{\pgfqpoint{1.836839in}{1.309870in}}%
\pgfpathcurveto{\pgfqpoint{1.836839in}{1.304046in}}{\pgfqpoint{1.839153in}{1.298460in}}{\pgfqpoint{1.843271in}{1.294341in}}%
\pgfpathcurveto{\pgfqpoint{1.847389in}{1.290223in}}{\pgfqpoint{1.852975in}{1.287909in}}{\pgfqpoint{1.858799in}{1.287909in}}%
\pgfpathlineto{\pgfqpoint{1.858799in}{1.287909in}}%
\pgfpathclose%
\pgfusepath{stroke,fill}%
\end{pgfscope}%
\begin{pgfscope}%
\pgfpathrectangle{\pgfqpoint{0.100000in}{0.209042in}}{\pgfqpoint{2.906250in}{2.906250in}}%
\pgfusepath{clip}%
\pgfsetbuttcap%
\pgfsetroundjoin%
\definecolor{currentfill}{rgb}{0.664915,0.198075,0.488836}%
\pgfsetfillcolor{currentfill}%
\pgfsetfillopacity{0.800000}%
\pgfsetlinewidth{1.003750pt}%
\definecolor{currentstroke}{rgb}{0.664915,0.198075,0.488836}%
\pgfsetstrokecolor{currentstroke}%
\pgfsetstrokeopacity{0.800000}%
\pgfsetdash{}{0pt}%
\pgfpathmoveto{\pgfqpoint{1.672059in}{2.101575in}}%
\pgfpathcurveto{\pgfqpoint{1.677883in}{2.101575in}}{\pgfqpoint{1.683469in}{2.103889in}}{\pgfqpoint{1.687588in}{2.108007in}}%
\pgfpathcurveto{\pgfqpoint{1.691706in}{2.112125in}}{\pgfqpoint{1.694020in}{2.117712in}}{\pgfqpoint{1.694020in}{2.123536in}}%
\pgfpathcurveto{\pgfqpoint{1.694020in}{2.129360in}}{\pgfqpoint{1.691706in}{2.134946in}}{\pgfqpoint{1.687588in}{2.139064in}}%
\pgfpathcurveto{\pgfqpoint{1.683469in}{2.143182in}}{\pgfqpoint{1.677883in}{2.145496in}}{\pgfqpoint{1.672059in}{2.145496in}}%
\pgfpathcurveto{\pgfqpoint{1.666235in}{2.145496in}}{\pgfqpoint{1.660649in}{2.143182in}}{\pgfqpoint{1.656531in}{2.139064in}}%
\pgfpathcurveto{\pgfqpoint{1.652413in}{2.134946in}}{\pgfqpoint{1.650099in}{2.129360in}}{\pgfqpoint{1.650099in}{2.123536in}}%
\pgfpathcurveto{\pgfqpoint{1.650099in}{2.117712in}}{\pgfqpoint{1.652413in}{2.112125in}}{\pgfqpoint{1.656531in}{2.108007in}}%
\pgfpathcurveto{\pgfqpoint{1.660649in}{2.103889in}}{\pgfqpoint{1.666235in}{2.101575in}}{\pgfqpoint{1.672059in}{2.101575in}}%
\pgfpathlineto{\pgfqpoint{1.672059in}{2.101575in}}%
\pgfpathclose%
\pgfusepath{stroke,fill}%
\end{pgfscope}%
\begin{pgfscope}%
\pgfpathrectangle{\pgfqpoint{0.100000in}{0.209042in}}{\pgfqpoint{2.906250in}{2.906250in}}%
\pgfusepath{clip}%
\pgfsetbuttcap%
\pgfsetroundjoin%
\definecolor{currentfill}{rgb}{0.451271,0.125132,0.507198}%
\pgfsetfillcolor{currentfill}%
\pgfsetfillopacity{0.800000}%
\pgfsetlinewidth{1.003750pt}%
\definecolor{currentstroke}{rgb}{0.451271,0.125132,0.507198}%
\pgfsetstrokecolor{currentstroke}%
\pgfsetstrokeopacity{0.800000}%
\pgfsetdash{}{0pt}%
\pgfpathmoveto{\pgfqpoint{1.289109in}{1.752428in}}%
\pgfpathcurveto{\pgfqpoint{1.294933in}{1.752428in}}{\pgfqpoint{1.300519in}{1.754742in}}{\pgfqpoint{1.304637in}{1.758860in}}%
\pgfpathcurveto{\pgfqpoint{1.308755in}{1.762978in}}{\pgfqpoint{1.311069in}{1.768564in}}{\pgfqpoint{1.311069in}{1.774388in}}%
\pgfpathcurveto{\pgfqpoint{1.311069in}{1.780212in}}{\pgfqpoint{1.308755in}{1.785799in}}{\pgfqpoint{1.304637in}{1.789917in}}%
\pgfpathcurveto{\pgfqpoint{1.300519in}{1.794035in}}{\pgfqpoint{1.294933in}{1.796349in}}{\pgfqpoint{1.289109in}{1.796349in}}%
\pgfpathcurveto{\pgfqpoint{1.283285in}{1.796349in}}{\pgfqpoint{1.277699in}{1.794035in}}{\pgfqpoint{1.273581in}{1.789917in}}%
\pgfpathcurveto{\pgfqpoint{1.269463in}{1.785799in}}{\pgfqpoint{1.267149in}{1.780212in}}{\pgfqpoint{1.267149in}{1.774388in}}%
\pgfpathcurveto{\pgfqpoint{1.267149in}{1.768564in}}{\pgfqpoint{1.269463in}{1.762978in}}{\pgfqpoint{1.273581in}{1.758860in}}%
\pgfpathcurveto{\pgfqpoint{1.277699in}{1.754742in}}{\pgfqpoint{1.283285in}{1.752428in}}{\pgfqpoint{1.289109in}{1.752428in}}%
\pgfpathlineto{\pgfqpoint{1.289109in}{1.752428in}}%
\pgfpathclose%
\pgfusepath{stroke,fill}%
\end{pgfscope}%
\begin{pgfscope}%
\pgfpathrectangle{\pgfqpoint{0.100000in}{0.209042in}}{\pgfqpoint{2.906250in}{2.906250in}}%
\pgfusepath{clip}%
\pgfsetbuttcap%
\pgfsetroundjoin%
\definecolor{currentfill}{rgb}{0.238826,0.059517,0.443256}%
\pgfsetfillcolor{currentfill}%
\pgfsetfillopacity{0.800000}%
\pgfsetlinewidth{1.003750pt}%
\definecolor{currentstroke}{rgb}{0.238826,0.059517,0.443256}%
\pgfsetstrokecolor{currentstroke}%
\pgfsetstrokeopacity{0.800000}%
\pgfsetdash{}{0pt}%
\pgfpathmoveto{\pgfqpoint{1.798672in}{1.780069in}}%
\pgfpathcurveto{\pgfqpoint{1.804496in}{1.780069in}}{\pgfqpoint{1.810082in}{1.782383in}}{\pgfqpoint{1.814200in}{1.786501in}}%
\pgfpathcurveto{\pgfqpoint{1.818318in}{1.790619in}}{\pgfqpoint{1.820632in}{1.796205in}}{\pgfqpoint{1.820632in}{1.802029in}}%
\pgfpathcurveto{\pgfqpoint{1.820632in}{1.807853in}}{\pgfqpoint{1.818318in}{1.813439in}}{\pgfqpoint{1.814200in}{1.817557in}}%
\pgfpathcurveto{\pgfqpoint{1.810082in}{1.821676in}}{\pgfqpoint{1.804496in}{1.823989in}}{\pgfqpoint{1.798672in}{1.823989in}}%
\pgfpathcurveto{\pgfqpoint{1.792848in}{1.823989in}}{\pgfqpoint{1.787262in}{1.821676in}}{\pgfqpoint{1.783144in}{1.817557in}}%
\pgfpathcurveto{\pgfqpoint{1.779025in}{1.813439in}}{\pgfqpoint{1.776712in}{1.807853in}}{\pgfqpoint{1.776712in}{1.802029in}}%
\pgfpathcurveto{\pgfqpoint{1.776712in}{1.796205in}}{\pgfqpoint{1.779025in}{1.790619in}}{\pgfqpoint{1.783144in}{1.786501in}}%
\pgfpathcurveto{\pgfqpoint{1.787262in}{1.782383in}}{\pgfqpoint{1.792848in}{1.780069in}}{\pgfqpoint{1.798672in}{1.780069in}}%
\pgfpathlineto{\pgfqpoint{1.798672in}{1.780069in}}%
\pgfpathclose%
\pgfusepath{stroke,fill}%
\end{pgfscope}%
\begin{pgfscope}%
\pgfpathrectangle{\pgfqpoint{0.100000in}{0.209042in}}{\pgfqpoint{2.906250in}{2.906250in}}%
\pgfusepath{clip}%
\pgfsetbuttcap%
\pgfsetroundjoin%
\definecolor{currentfill}{rgb}{0.588158,0.173652,0.502035}%
\pgfsetfillcolor{currentfill}%
\pgfsetfillopacity{0.800000}%
\pgfsetlinewidth{1.003750pt}%
\definecolor{currentstroke}{rgb}{0.588158,0.173652,0.502035}%
\pgfsetstrokecolor{currentstroke}%
\pgfsetstrokeopacity{0.800000}%
\pgfsetdash{}{0pt}%
\pgfpathmoveto{\pgfqpoint{1.347652in}{1.409600in}}%
\pgfpathcurveto{\pgfqpoint{1.353476in}{1.409600in}}{\pgfqpoint{1.359062in}{1.411914in}}{\pgfqpoint{1.363180in}{1.416032in}}%
\pgfpathcurveto{\pgfqpoint{1.367298in}{1.420150in}}{\pgfqpoint{1.369612in}{1.425736in}}{\pgfqpoint{1.369612in}{1.431560in}}%
\pgfpathcurveto{\pgfqpoint{1.369612in}{1.437384in}}{\pgfqpoint{1.367298in}{1.442970in}}{\pgfqpoint{1.363180in}{1.447088in}}%
\pgfpathcurveto{\pgfqpoint{1.359062in}{1.451207in}}{\pgfqpoint{1.353476in}{1.453520in}}{\pgfqpoint{1.347652in}{1.453520in}}%
\pgfpathcurveto{\pgfqpoint{1.341828in}{1.453520in}}{\pgfqpoint{1.336242in}{1.451207in}}{\pgfqpoint{1.332123in}{1.447088in}}%
\pgfpathcurveto{\pgfqpoint{1.328005in}{1.442970in}}{\pgfqpoint{1.325691in}{1.437384in}}{\pgfqpoint{1.325691in}{1.431560in}}%
\pgfpathcurveto{\pgfqpoint{1.325691in}{1.425736in}}{\pgfqpoint{1.328005in}{1.420150in}}{\pgfqpoint{1.332123in}{1.416032in}}%
\pgfpathcurveto{\pgfqpoint{1.336242in}{1.411914in}}{\pgfqpoint{1.341828in}{1.409600in}}{\pgfqpoint{1.347652in}{1.409600in}}%
\pgfpathlineto{\pgfqpoint{1.347652in}{1.409600in}}%
\pgfpathclose%
\pgfusepath{stroke,fill}%
\end{pgfscope}%
\begin{pgfscope}%
\pgfpathrectangle{\pgfqpoint{0.100000in}{0.209042in}}{\pgfqpoint{2.906250in}{2.906250in}}%
\pgfusepath{clip}%
\pgfsetbuttcap%
\pgfsetroundjoin%
\definecolor{currentfill}{rgb}{0.420791,0.112920,0.505215}%
\pgfsetfillcolor{currentfill}%
\pgfsetfillopacity{0.800000}%
\pgfsetlinewidth{1.003750pt}%
\definecolor{currentstroke}{rgb}{0.420791,0.112920,0.505215}%
\pgfsetstrokecolor{currentstroke}%
\pgfsetstrokeopacity{0.800000}%
\pgfsetdash{}{0pt}%
\pgfpathmoveto{\pgfqpoint{1.448735in}{1.916845in}}%
\pgfpathcurveto{\pgfqpoint{1.454559in}{1.916845in}}{\pgfqpoint{1.460145in}{1.919159in}}{\pgfqpoint{1.464263in}{1.923277in}}%
\pgfpathcurveto{\pgfqpoint{1.468382in}{1.927395in}}{\pgfqpoint{1.470695in}{1.932981in}}{\pgfqpoint{1.470695in}{1.938805in}}%
\pgfpathcurveto{\pgfqpoint{1.470695in}{1.944629in}}{\pgfqpoint{1.468382in}{1.950215in}}{\pgfqpoint{1.464263in}{1.954333in}}%
\pgfpathcurveto{\pgfqpoint{1.460145in}{1.958452in}}{\pgfqpoint{1.454559in}{1.960766in}}{\pgfqpoint{1.448735in}{1.960766in}}%
\pgfpathcurveto{\pgfqpoint{1.442911in}{1.960766in}}{\pgfqpoint{1.437325in}{1.958452in}}{\pgfqpoint{1.433207in}{1.954333in}}%
\pgfpathcurveto{\pgfqpoint{1.429089in}{1.950215in}}{\pgfqpoint{1.426775in}{1.944629in}}{\pgfqpoint{1.426775in}{1.938805in}}%
\pgfpathcurveto{\pgfqpoint{1.426775in}{1.932981in}}{\pgfqpoint{1.429089in}{1.927395in}}{\pgfqpoint{1.433207in}{1.923277in}}%
\pgfpathcurveto{\pgfqpoint{1.437325in}{1.919159in}}{\pgfqpoint{1.442911in}{1.916845in}}{\pgfqpoint{1.448735in}{1.916845in}}%
\pgfpathlineto{\pgfqpoint{1.448735in}{1.916845in}}%
\pgfpathclose%
\pgfusepath{stroke,fill}%
\end{pgfscope}%
\begin{pgfscope}%
\pgfpathrectangle{\pgfqpoint{0.100000in}{0.209042in}}{\pgfqpoint{2.906250in}{2.906250in}}%
\pgfusepath{clip}%
\pgfsetbuttcap%
\pgfsetroundjoin%
\definecolor{currentfill}{rgb}{0.140858,0.068654,0.324538}%
\pgfsetfillcolor{currentfill}%
\pgfsetfillopacity{0.800000}%
\pgfsetlinewidth{1.003750pt}%
\definecolor{currentstroke}{rgb}{0.140858,0.068654,0.324538}%
\pgfsetstrokecolor{currentstroke}%
\pgfsetstrokeopacity{0.800000}%
\pgfsetdash{}{0pt}%
\pgfpathmoveto{\pgfqpoint{1.738738in}{1.690008in}}%
\pgfpathcurveto{\pgfqpoint{1.744562in}{1.690008in}}{\pgfqpoint{1.750148in}{1.692322in}}{\pgfqpoint{1.754267in}{1.696440in}}%
\pgfpathcurveto{\pgfqpoint{1.758385in}{1.700559in}}{\pgfqpoint{1.760699in}{1.706145in}}{\pgfqpoint{1.760699in}{1.711969in}}%
\pgfpathcurveto{\pgfqpoint{1.760699in}{1.717793in}}{\pgfqpoint{1.758385in}{1.723379in}}{\pgfqpoint{1.754267in}{1.727497in}}%
\pgfpathcurveto{\pgfqpoint{1.750148in}{1.731615in}}{\pgfqpoint{1.744562in}{1.733929in}}{\pgfqpoint{1.738738in}{1.733929in}}%
\pgfpathcurveto{\pgfqpoint{1.732914in}{1.733929in}}{\pgfqpoint{1.727328in}{1.731615in}}{\pgfqpoint{1.723210in}{1.727497in}}%
\pgfpathcurveto{\pgfqpoint{1.719092in}{1.723379in}}{\pgfqpoint{1.716778in}{1.717793in}}{\pgfqpoint{1.716778in}{1.711969in}}%
\pgfpathcurveto{\pgfqpoint{1.716778in}{1.706145in}}{\pgfqpoint{1.719092in}{1.700559in}}{\pgfqpoint{1.723210in}{1.696440in}}%
\pgfpathcurveto{\pgfqpoint{1.727328in}{1.692322in}}{\pgfqpoint{1.732914in}{1.690008in}}{\pgfqpoint{1.738738in}{1.690008in}}%
\pgfpathlineto{\pgfqpoint{1.738738in}{1.690008in}}%
\pgfpathclose%
\pgfusepath{stroke,fill}%
\end{pgfscope}%
\begin{pgfscope}%
\pgfpathrectangle{\pgfqpoint{0.100000in}{0.209042in}}{\pgfqpoint{2.906250in}{2.906250in}}%
\pgfusepath{clip}%
\pgfsetbuttcap%
\pgfsetroundjoin%
\definecolor{currentfill}{rgb}{0.500438,0.143719,0.507920}%
\pgfsetfillcolor{currentfill}%
\pgfsetfillopacity{0.800000}%
\pgfsetlinewidth{1.003750pt}%
\definecolor{currentstroke}{rgb}{0.500438,0.143719,0.507920}%
\pgfsetstrokecolor{currentstroke}%
\pgfsetstrokeopacity{0.800000}%
\pgfsetdash{}{0pt}%
\pgfpathmoveto{\pgfqpoint{1.267248in}{1.789105in}}%
\pgfpathcurveto{\pgfqpoint{1.273071in}{1.789105in}}{\pgfqpoint{1.278658in}{1.791419in}}{\pgfqpoint{1.282776in}{1.795537in}}%
\pgfpathcurveto{\pgfqpoint{1.286894in}{1.799655in}}{\pgfqpoint{1.289208in}{1.805242in}}{\pgfqpoint{1.289208in}{1.811066in}}%
\pgfpathcurveto{\pgfqpoint{1.289208in}{1.816889in}}{\pgfqpoint{1.286894in}{1.822476in}}{\pgfqpoint{1.282776in}{1.826594in}}%
\pgfpathcurveto{\pgfqpoint{1.278658in}{1.830712in}}{\pgfqpoint{1.273071in}{1.833026in}}{\pgfqpoint{1.267248in}{1.833026in}}%
\pgfpathcurveto{\pgfqpoint{1.261424in}{1.833026in}}{\pgfqpoint{1.255837in}{1.830712in}}{\pgfqpoint{1.251719in}{1.826594in}}%
\pgfpathcurveto{\pgfqpoint{1.247601in}{1.822476in}}{\pgfqpoint{1.245287in}{1.816889in}}{\pgfqpoint{1.245287in}{1.811066in}}%
\pgfpathcurveto{\pgfqpoint{1.245287in}{1.805242in}}{\pgfqpoint{1.247601in}{1.799655in}}{\pgfqpoint{1.251719in}{1.795537in}}%
\pgfpathcurveto{\pgfqpoint{1.255837in}{1.791419in}}{\pgfqpoint{1.261424in}{1.789105in}}{\pgfqpoint{1.267248in}{1.789105in}}%
\pgfpathlineto{\pgfqpoint{1.267248in}{1.789105in}}%
\pgfpathclose%
\pgfusepath{stroke,fill}%
\end{pgfscope}%
\begin{pgfscope}%
\pgfpathrectangle{\pgfqpoint{0.100000in}{0.209042in}}{\pgfqpoint{2.906250in}{2.906250in}}%
\pgfusepath{clip}%
\pgfsetbuttcap%
\pgfsetroundjoin%
\definecolor{currentfill}{rgb}{0.204935,0.062907,0.411514}%
\pgfsetfillcolor{currentfill}%
\pgfsetfillopacity{0.800000}%
\pgfsetlinewidth{1.003750pt}%
\definecolor{currentstroke}{rgb}{0.204935,0.062907,0.411514}%
\pgfsetstrokecolor{currentstroke}%
\pgfsetstrokeopacity{0.800000}%
\pgfsetdash{}{0pt}%
\pgfpathmoveto{\pgfqpoint{1.557267in}{1.581882in}}%
\pgfpathcurveto{\pgfqpoint{1.563091in}{1.581882in}}{\pgfqpoint{1.568677in}{1.584196in}}{\pgfqpoint{1.572795in}{1.588314in}}%
\pgfpathcurveto{\pgfqpoint{1.576913in}{1.592432in}}{\pgfqpoint{1.579227in}{1.598018in}}{\pgfqpoint{1.579227in}{1.603842in}}%
\pgfpathcurveto{\pgfqpoint{1.579227in}{1.609666in}}{\pgfqpoint{1.576913in}{1.615252in}}{\pgfqpoint{1.572795in}{1.619370in}}%
\pgfpathcurveto{\pgfqpoint{1.568677in}{1.623489in}}{\pgfqpoint{1.563091in}{1.625802in}}{\pgfqpoint{1.557267in}{1.625802in}}%
\pgfpathcurveto{\pgfqpoint{1.551443in}{1.625802in}}{\pgfqpoint{1.545857in}{1.623489in}}{\pgfqpoint{1.541739in}{1.619370in}}%
\pgfpathcurveto{\pgfqpoint{1.537620in}{1.615252in}}{\pgfqpoint{1.535307in}{1.609666in}}{\pgfqpoint{1.535307in}{1.603842in}}%
\pgfpathcurveto{\pgfqpoint{1.535307in}{1.598018in}}{\pgfqpoint{1.537620in}{1.592432in}}{\pgfqpoint{1.541739in}{1.588314in}}%
\pgfpathcurveto{\pgfqpoint{1.545857in}{1.584196in}}{\pgfqpoint{1.551443in}{1.581882in}}{\pgfqpoint{1.557267in}{1.581882in}}%
\pgfpathlineto{\pgfqpoint{1.557267in}{1.581882in}}%
\pgfpathclose%
\pgfusepath{stroke,fill}%
\end{pgfscope}%
\begin{pgfscope}%
\pgfpathrectangle{\pgfqpoint{0.100000in}{0.209042in}}{\pgfqpoint{2.906250in}{2.906250in}}%
\pgfusepath{clip}%
\pgfsetbuttcap%
\pgfsetroundjoin%
\definecolor{currentfill}{rgb}{0.123833,0.067295,0.295879}%
\pgfsetfillcolor{currentfill}%
\pgfsetfillopacity{0.800000}%
\pgfsetlinewidth{1.003750pt}%
\definecolor{currentstroke}{rgb}{0.123833,0.067295,0.295879}%
\pgfsetstrokecolor{currentstroke}%
\pgfsetstrokeopacity{0.800000}%
\pgfsetdash{}{0pt}%
\pgfpathmoveto{\pgfqpoint{1.687187in}{1.667493in}}%
\pgfpathcurveto{\pgfqpoint{1.693011in}{1.667493in}}{\pgfqpoint{1.698597in}{1.669807in}}{\pgfqpoint{1.702715in}{1.673925in}}%
\pgfpathcurveto{\pgfqpoint{1.706833in}{1.678044in}}{\pgfqpoint{1.709147in}{1.683630in}}{\pgfqpoint{1.709147in}{1.689454in}}%
\pgfpathcurveto{\pgfqpoint{1.709147in}{1.695278in}}{\pgfqpoint{1.706833in}{1.700864in}}{\pgfqpoint{1.702715in}{1.704982in}}%
\pgfpathcurveto{\pgfqpoint{1.698597in}{1.709100in}}{\pgfqpoint{1.693011in}{1.711414in}}{\pgfqpoint{1.687187in}{1.711414in}}%
\pgfpathcurveto{\pgfqpoint{1.681363in}{1.711414in}}{\pgfqpoint{1.675777in}{1.709100in}}{\pgfqpoint{1.671659in}{1.704982in}}%
\pgfpathcurveto{\pgfqpoint{1.667540in}{1.700864in}}{\pgfqpoint{1.665227in}{1.695278in}}{\pgfqpoint{1.665227in}{1.689454in}}%
\pgfpathcurveto{\pgfqpoint{1.665227in}{1.683630in}}{\pgfqpoint{1.667540in}{1.678044in}}{\pgfqpoint{1.671659in}{1.673925in}}%
\pgfpathcurveto{\pgfqpoint{1.675777in}{1.669807in}}{\pgfqpoint{1.681363in}{1.667493in}}{\pgfqpoint{1.687187in}{1.667493in}}%
\pgfpathlineto{\pgfqpoint{1.687187in}{1.667493in}}%
\pgfpathclose%
\pgfusepath{stroke,fill}%
\end{pgfscope}%
\begin{pgfscope}%
\pgfpathrectangle{\pgfqpoint{0.100000in}{0.209042in}}{\pgfqpoint{2.906250in}{2.906250in}}%
\pgfusepath{clip}%
\pgfsetbuttcap%
\pgfsetroundjoin%
\definecolor{currentfill}{rgb}{0.984622,0.520713,0.376698}%
\pgfsetfillcolor{currentfill}%
\pgfsetfillopacity{0.800000}%
\pgfsetlinewidth{1.003750pt}%
\definecolor{currentstroke}{rgb}{0.984622,0.520713,0.376698}%
\pgfsetstrokecolor{currentstroke}%
\pgfsetstrokeopacity{0.800000}%
\pgfsetdash{}{0pt}%
\pgfpathmoveto{\pgfqpoint{1.336019in}{1.088571in}}%
\pgfpathcurveto{\pgfqpoint{1.341843in}{1.088571in}}{\pgfqpoint{1.347429in}{1.090884in}}{\pgfqpoint{1.351547in}{1.095003in}}%
\pgfpathcurveto{\pgfqpoint{1.355665in}{1.099121in}}{\pgfqpoint{1.357979in}{1.104707in}}{\pgfqpoint{1.357979in}{1.110531in}}%
\pgfpathcurveto{\pgfqpoint{1.357979in}{1.116355in}}{\pgfqpoint{1.355665in}{1.121941in}}{\pgfqpoint{1.351547in}{1.126059in}}%
\pgfpathcurveto{\pgfqpoint{1.347429in}{1.130177in}}{\pgfqpoint{1.341843in}{1.132491in}}{\pgfqpoint{1.336019in}{1.132491in}}%
\pgfpathcurveto{\pgfqpoint{1.330195in}{1.132491in}}{\pgfqpoint{1.324609in}{1.130177in}}{\pgfqpoint{1.320491in}{1.126059in}}%
\pgfpathcurveto{\pgfqpoint{1.316372in}{1.121941in}}{\pgfqpoint{1.314059in}{1.116355in}}{\pgfqpoint{1.314059in}{1.110531in}}%
\pgfpathcurveto{\pgfqpoint{1.314059in}{1.104707in}}{\pgfqpoint{1.316372in}{1.099121in}}{\pgfqpoint{1.320491in}{1.095003in}}%
\pgfpathcurveto{\pgfqpoint{1.324609in}{1.090884in}}{\pgfqpoint{1.330195in}{1.088571in}}{\pgfqpoint{1.336019in}{1.088571in}}%
\pgfpathlineto{\pgfqpoint{1.336019in}{1.088571in}}%
\pgfpathclose%
\pgfusepath{stroke,fill}%
\end{pgfscope}%
\begin{pgfscope}%
\pgfpathrectangle{\pgfqpoint{0.100000in}{0.209042in}}{\pgfqpoint{2.906250in}{2.906250in}}%
\pgfusepath{clip}%
\pgfsetbuttcap%
\pgfsetroundjoin%
\definecolor{currentfill}{rgb}{0.225302,0.060445,0.431742}%
\pgfsetfillcolor{currentfill}%
\pgfsetfillopacity{0.800000}%
\pgfsetlinewidth{1.003750pt}%
\definecolor{currentstroke}{rgb}{0.225302,0.060445,0.431742}%
\pgfsetstrokecolor{currentstroke}%
\pgfsetstrokeopacity{0.800000}%
\pgfsetdash{}{0pt}%
\pgfpathmoveto{\pgfqpoint{1.740077in}{1.559423in}}%
\pgfpathcurveto{\pgfqpoint{1.745901in}{1.559423in}}{\pgfqpoint{1.751488in}{1.561737in}}{\pgfqpoint{1.755606in}{1.565855in}}%
\pgfpathcurveto{\pgfqpoint{1.759724in}{1.569973in}}{\pgfqpoint{1.762038in}{1.575560in}}{\pgfqpoint{1.762038in}{1.581383in}}%
\pgfpathcurveto{\pgfqpoint{1.762038in}{1.587207in}}{\pgfqpoint{1.759724in}{1.592794in}}{\pgfqpoint{1.755606in}{1.596912in}}%
\pgfpathcurveto{\pgfqpoint{1.751488in}{1.601030in}}{\pgfqpoint{1.745901in}{1.603344in}}{\pgfqpoint{1.740077in}{1.603344in}}%
\pgfpathcurveto{\pgfqpoint{1.734254in}{1.603344in}}{\pgfqpoint{1.728667in}{1.601030in}}{\pgfqpoint{1.724549in}{1.596912in}}%
\pgfpathcurveto{\pgfqpoint{1.720431in}{1.592794in}}{\pgfqpoint{1.718117in}{1.587207in}}{\pgfqpoint{1.718117in}{1.581383in}}%
\pgfpathcurveto{\pgfqpoint{1.718117in}{1.575560in}}{\pgfqpoint{1.720431in}{1.569973in}}{\pgfqpoint{1.724549in}{1.565855in}}%
\pgfpathcurveto{\pgfqpoint{1.728667in}{1.561737in}}{\pgfqpoint{1.734254in}{1.559423in}}{\pgfqpoint{1.740077in}{1.559423in}}%
\pgfpathlineto{\pgfqpoint{1.740077in}{1.559423in}}%
\pgfpathclose%
\pgfusepath{stroke,fill}%
\end{pgfscope}%
\begin{pgfscope}%
\pgfpathrectangle{\pgfqpoint{0.100000in}{0.209042in}}{\pgfqpoint{2.906250in}{2.906250in}}%
\pgfusepath{clip}%
\pgfsetbuttcap%
\pgfsetroundjoin%
\definecolor{currentfill}{rgb}{0.245543,0.059352,0.448436}%
\pgfsetfillcolor{currentfill}%
\pgfsetfillopacity{0.800000}%
\pgfsetlinewidth{1.003750pt}%
\definecolor{currentstroke}{rgb}{0.245543,0.059352,0.448436}%
\pgfsetstrokecolor{currentstroke}%
\pgfsetstrokeopacity{0.800000}%
\pgfsetdash{}{0pt}%
\pgfpathmoveto{\pgfqpoint{1.514647in}{1.794681in}}%
\pgfpathcurveto{\pgfqpoint{1.520471in}{1.794681in}}{\pgfqpoint{1.526058in}{1.796995in}}{\pgfqpoint{1.530176in}{1.801113in}}%
\pgfpathcurveto{\pgfqpoint{1.534294in}{1.805231in}}{\pgfqpoint{1.536608in}{1.810817in}}{\pgfqpoint{1.536608in}{1.816641in}}%
\pgfpathcurveto{\pgfqpoint{1.536608in}{1.822465in}}{\pgfqpoint{1.534294in}{1.828051in}}{\pgfqpoint{1.530176in}{1.832169in}}%
\pgfpathcurveto{\pgfqpoint{1.526058in}{1.836287in}}{\pgfqpoint{1.520471in}{1.838601in}}{\pgfqpoint{1.514647in}{1.838601in}}%
\pgfpathcurveto{\pgfqpoint{1.508824in}{1.838601in}}{\pgfqpoint{1.503237in}{1.836287in}}{\pgfqpoint{1.499119in}{1.832169in}}%
\pgfpathcurveto{\pgfqpoint{1.495001in}{1.828051in}}{\pgfqpoint{1.492687in}{1.822465in}}{\pgfqpoint{1.492687in}{1.816641in}}%
\pgfpathcurveto{\pgfqpoint{1.492687in}{1.810817in}}{\pgfqpoint{1.495001in}{1.805231in}}{\pgfqpoint{1.499119in}{1.801113in}}%
\pgfpathcurveto{\pgfqpoint{1.503237in}{1.796995in}}{\pgfqpoint{1.508824in}{1.794681in}}{\pgfqpoint{1.514647in}{1.794681in}}%
\pgfpathlineto{\pgfqpoint{1.514647in}{1.794681in}}%
\pgfpathclose%
\pgfusepath{stroke,fill}%
\end{pgfscope}%
\begin{pgfscope}%
\pgfpathrectangle{\pgfqpoint{0.100000in}{0.209042in}}{\pgfqpoint{2.906250in}{2.906250in}}%
\pgfusepath{clip}%
\pgfsetbuttcap%
\pgfsetroundjoin%
\definecolor{currentfill}{rgb}{0.525270,0.152569,0.507192}%
\pgfsetfillcolor{currentfill}%
\pgfsetfillopacity{0.800000}%
\pgfsetlinewidth{1.003750pt}%
\definecolor{currentstroke}{rgb}{0.525270,0.152569,0.507192}%
\pgfsetstrokecolor{currentstroke}%
\pgfsetstrokeopacity{0.800000}%
\pgfsetdash{}{0pt}%
\pgfpathmoveto{\pgfqpoint{1.726429in}{1.348408in}}%
\pgfpathcurveto{\pgfqpoint{1.732253in}{1.348408in}}{\pgfqpoint{1.737839in}{1.350722in}}{\pgfqpoint{1.741957in}{1.354840in}}%
\pgfpathcurveto{\pgfqpoint{1.746076in}{1.358958in}}{\pgfqpoint{1.748389in}{1.364544in}}{\pgfqpoint{1.748389in}{1.370368in}}%
\pgfpathcurveto{\pgfqpoint{1.748389in}{1.376192in}}{\pgfqpoint{1.746076in}{1.381778in}}{\pgfqpoint{1.741957in}{1.385897in}}%
\pgfpathcurveto{\pgfqpoint{1.737839in}{1.390015in}}{\pgfqpoint{1.732253in}{1.392329in}}{\pgfqpoint{1.726429in}{1.392329in}}%
\pgfpathcurveto{\pgfqpoint{1.720605in}{1.392329in}}{\pgfqpoint{1.715019in}{1.390015in}}{\pgfqpoint{1.710901in}{1.385897in}}%
\pgfpathcurveto{\pgfqpoint{1.706783in}{1.381778in}}{\pgfqpoint{1.704469in}{1.376192in}}{\pgfqpoint{1.704469in}{1.370368in}}%
\pgfpathcurveto{\pgfqpoint{1.704469in}{1.364544in}}{\pgfqpoint{1.706783in}{1.358958in}}{\pgfqpoint{1.710901in}{1.354840in}}%
\pgfpathcurveto{\pgfqpoint{1.715019in}{1.350722in}}{\pgfqpoint{1.720605in}{1.348408in}}{\pgfqpoint{1.726429in}{1.348408in}}%
\pgfpathlineto{\pgfqpoint{1.726429in}{1.348408in}}%
\pgfpathclose%
\pgfusepath{stroke,fill}%
\end{pgfscope}%
\begin{pgfscope}%
\pgfpathrectangle{\pgfqpoint{0.100000in}{0.209042in}}{\pgfqpoint{2.906250in}{2.906250in}}%
\pgfusepath{clip}%
\pgfsetbuttcap%
\pgfsetroundjoin%
\definecolor{currentfill}{rgb}{0.908884,0.324755,0.385308}%
\pgfsetfillcolor{currentfill}%
\pgfsetfillopacity{0.800000}%
\pgfsetlinewidth{1.003750pt}%
\definecolor{currentstroke}{rgb}{0.908884,0.324755,0.385308}%
\pgfsetstrokecolor{currentstroke}%
\pgfsetstrokeopacity{0.800000}%
\pgfsetdash{}{0pt}%
\pgfpathmoveto{\pgfqpoint{1.235364in}{2.144845in}}%
\pgfpathcurveto{\pgfqpoint{1.241188in}{2.144845in}}{\pgfqpoint{1.246774in}{2.147158in}}{\pgfqpoint{1.250892in}{2.151277in}}%
\pgfpathcurveto{\pgfqpoint{1.255010in}{2.155395in}}{\pgfqpoint{1.257324in}{2.160981in}}{\pgfqpoint{1.257324in}{2.166805in}}%
\pgfpathcurveto{\pgfqpoint{1.257324in}{2.172629in}}{\pgfqpoint{1.255010in}{2.178215in}}{\pgfqpoint{1.250892in}{2.182333in}}%
\pgfpathcurveto{\pgfqpoint{1.246774in}{2.186451in}}{\pgfqpoint{1.241188in}{2.188765in}}{\pgfqpoint{1.235364in}{2.188765in}}%
\pgfpathcurveto{\pgfqpoint{1.229540in}{2.188765in}}{\pgfqpoint{1.223954in}{2.186451in}}{\pgfqpoint{1.219836in}{2.182333in}}%
\pgfpathcurveto{\pgfqpoint{1.215718in}{2.178215in}}{\pgfqpoint{1.213404in}{2.172629in}}{\pgfqpoint{1.213404in}{2.166805in}}%
\pgfpathcurveto{\pgfqpoint{1.213404in}{2.160981in}}{\pgfqpoint{1.215718in}{2.155395in}}{\pgfqpoint{1.219836in}{2.151277in}}%
\pgfpathcurveto{\pgfqpoint{1.223954in}{2.147158in}}{\pgfqpoint{1.229540in}{2.144845in}}{\pgfqpoint{1.235364in}{2.144845in}}%
\pgfpathlineto{\pgfqpoint{1.235364in}{2.144845in}}%
\pgfpathclose%
\pgfusepath{stroke,fill}%
\end{pgfscope}%
\begin{pgfscope}%
\pgfpathrectangle{\pgfqpoint{0.100000in}{0.209042in}}{\pgfqpoint{2.906250in}{2.906250in}}%
\pgfusepath{clip}%
\pgfsetbuttcap%
\pgfsetroundjoin%
\definecolor{currentfill}{rgb}{0.265447,0.060237,0.461840}%
\pgfsetfillcolor{currentfill}%
\pgfsetfillopacity{0.800000}%
\pgfsetlinewidth{1.003750pt}%
\definecolor{currentstroke}{rgb}{0.265447,0.060237,0.461840}%
\pgfsetstrokecolor{currentstroke}%
\pgfsetstrokeopacity{0.800000}%
\pgfsetdash{}{0pt}%
\pgfpathmoveto{\pgfqpoint{1.798933in}{1.785773in}}%
\pgfpathcurveto{\pgfqpoint{1.804757in}{1.785773in}}{\pgfqpoint{1.810343in}{1.788086in}}{\pgfqpoint{1.814461in}{1.792205in}}%
\pgfpathcurveto{\pgfqpoint{1.818580in}{1.796323in}}{\pgfqpoint{1.820893in}{1.801909in}}{\pgfqpoint{1.820893in}{1.807733in}}%
\pgfpathcurveto{\pgfqpoint{1.820893in}{1.813557in}}{\pgfqpoint{1.818580in}{1.819143in}}{\pgfqpoint{1.814461in}{1.823261in}}%
\pgfpathcurveto{\pgfqpoint{1.810343in}{1.827379in}}{\pgfqpoint{1.804757in}{1.829693in}}{\pgfqpoint{1.798933in}{1.829693in}}%
\pgfpathcurveto{\pgfqpoint{1.793109in}{1.829693in}}{\pgfqpoint{1.787523in}{1.827379in}}{\pgfqpoint{1.783405in}{1.823261in}}%
\pgfpathcurveto{\pgfqpoint{1.779287in}{1.819143in}}{\pgfqpoint{1.776973in}{1.813557in}}{\pgfqpoint{1.776973in}{1.807733in}}%
\pgfpathcurveto{\pgfqpoint{1.776973in}{1.801909in}}{\pgfqpoint{1.779287in}{1.796323in}}{\pgfqpoint{1.783405in}{1.792205in}}%
\pgfpathcurveto{\pgfqpoint{1.787523in}{1.788086in}}{\pgfqpoint{1.793109in}{1.785773in}}{\pgfqpoint{1.798933in}{1.785773in}}%
\pgfpathlineto{\pgfqpoint{1.798933in}{1.785773in}}%
\pgfpathclose%
\pgfusepath{stroke,fill}%
\end{pgfscope}%
\begin{pgfscope}%
\pgfpathrectangle{\pgfqpoint{0.100000in}{0.209042in}}{\pgfqpoint{2.906250in}{2.906250in}}%
\pgfusepath{clip}%
\pgfsetbuttcap%
\pgfsetroundjoin%
\definecolor{currentfill}{rgb}{0.297740,0.066117,0.478243}%
\pgfsetfillcolor{currentfill}%
\pgfsetfillopacity{0.800000}%
\pgfsetlinewidth{1.003750pt}%
\definecolor{currentstroke}{rgb}{0.297740,0.066117,0.478243}%
\pgfsetstrokecolor{currentstroke}%
\pgfsetstrokeopacity{0.800000}%
\pgfsetdash{}{0pt}%
\pgfpathmoveto{\pgfqpoint{1.426406in}{1.655375in}}%
\pgfpathcurveto{\pgfqpoint{1.432229in}{1.655375in}}{\pgfqpoint{1.437816in}{1.657689in}}{\pgfqpoint{1.441934in}{1.661808in}}%
\pgfpathcurveto{\pgfqpoint{1.446052in}{1.665926in}}{\pgfqpoint{1.448366in}{1.671512in}}{\pgfqpoint{1.448366in}{1.677336in}}%
\pgfpathcurveto{\pgfqpoint{1.448366in}{1.683160in}}{\pgfqpoint{1.446052in}{1.688746in}}{\pgfqpoint{1.441934in}{1.692864in}}%
\pgfpathcurveto{\pgfqpoint{1.437816in}{1.696982in}}{\pgfqpoint{1.432229in}{1.699296in}}{\pgfqpoint{1.426406in}{1.699296in}}%
\pgfpathcurveto{\pgfqpoint{1.420582in}{1.699296in}}{\pgfqpoint{1.414995in}{1.696982in}}{\pgfqpoint{1.410877in}{1.692864in}}%
\pgfpathcurveto{\pgfqpoint{1.406759in}{1.688746in}}{\pgfqpoint{1.404445in}{1.683160in}}{\pgfqpoint{1.404445in}{1.677336in}}%
\pgfpathcurveto{\pgfqpoint{1.404445in}{1.671512in}}{\pgfqpoint{1.406759in}{1.665926in}}{\pgfqpoint{1.410877in}{1.661808in}}%
\pgfpathcurveto{\pgfqpoint{1.414995in}{1.657689in}}{\pgfqpoint{1.420582in}{1.655375in}}{\pgfqpoint{1.426406in}{1.655375in}}%
\pgfpathlineto{\pgfqpoint{1.426406in}{1.655375in}}%
\pgfpathclose%
\pgfusepath{stroke,fill}%
\end{pgfscope}%
\begin{pgfscope}%
\pgfpathrectangle{\pgfqpoint{0.100000in}{0.209042in}}{\pgfqpoint{2.906250in}{2.906250in}}%
\pgfusepath{clip}%
\pgfsetbuttcap%
\pgfsetroundjoin%
\definecolor{currentfill}{rgb}{0.384299,0.097855,0.501002}%
\pgfsetfillcolor{currentfill}%
\pgfsetfillopacity{0.800000}%
\pgfsetlinewidth{1.003750pt}%
\definecolor{currentstroke}{rgb}{0.384299,0.097855,0.501002}%
\pgfsetstrokecolor{currentstroke}%
\pgfsetstrokeopacity{0.800000}%
\pgfsetdash{}{0pt}%
\pgfpathmoveto{\pgfqpoint{1.818345in}{1.470207in}}%
\pgfpathcurveto{\pgfqpoint{1.824169in}{1.470207in}}{\pgfqpoint{1.829755in}{1.472521in}}{\pgfqpoint{1.833874in}{1.476639in}}%
\pgfpathcurveto{\pgfqpoint{1.837992in}{1.480757in}}{\pgfqpoint{1.840306in}{1.486343in}}{\pgfqpoint{1.840306in}{1.492167in}}%
\pgfpathcurveto{\pgfqpoint{1.840306in}{1.497991in}}{\pgfqpoint{1.837992in}{1.503577in}}{\pgfqpoint{1.833874in}{1.507696in}}%
\pgfpathcurveto{\pgfqpoint{1.829755in}{1.511814in}}{\pgfqpoint{1.824169in}{1.514128in}}{\pgfqpoint{1.818345in}{1.514128in}}%
\pgfpathcurveto{\pgfqpoint{1.812521in}{1.514128in}}{\pgfqpoint{1.806935in}{1.511814in}}{\pgfqpoint{1.802817in}{1.507696in}}%
\pgfpathcurveto{\pgfqpoint{1.798699in}{1.503577in}}{\pgfqpoint{1.796385in}{1.497991in}}{\pgfqpoint{1.796385in}{1.492167in}}%
\pgfpathcurveto{\pgfqpoint{1.796385in}{1.486343in}}{\pgfqpoint{1.798699in}{1.480757in}}{\pgfqpoint{1.802817in}{1.476639in}}%
\pgfpathcurveto{\pgfqpoint{1.806935in}{1.472521in}}{\pgfqpoint{1.812521in}{1.470207in}}{\pgfqpoint{1.818345in}{1.470207in}}%
\pgfpathlineto{\pgfqpoint{1.818345in}{1.470207in}}%
\pgfpathclose%
\pgfusepath{stroke,fill}%
\end{pgfscope}%
\begin{pgfscope}%
\pgfpathrectangle{\pgfqpoint{0.100000in}{0.209042in}}{\pgfqpoint{2.906250in}{2.906250in}}%
\pgfusepath{clip}%
\pgfsetbuttcap%
\pgfsetroundjoin%
\definecolor{currentfill}{rgb}{0.146785,0.068738,0.334011}%
\pgfsetfillcolor{currentfill}%
\pgfsetfillopacity{0.800000}%
\pgfsetlinewidth{1.003750pt}%
\definecolor{currentstroke}{rgb}{0.146785,0.068738,0.334011}%
\pgfsetstrokecolor{currentstroke}%
\pgfsetstrokeopacity{0.800000}%
\pgfsetdash{}{0pt}%
\pgfpathmoveto{\pgfqpoint{1.650680in}{1.723855in}}%
\pgfpathcurveto{\pgfqpoint{1.656504in}{1.723855in}}{\pgfqpoint{1.662090in}{1.726169in}}{\pgfqpoint{1.666208in}{1.730287in}}%
\pgfpathcurveto{\pgfqpoint{1.670326in}{1.734405in}}{\pgfqpoint{1.672640in}{1.739991in}}{\pgfqpoint{1.672640in}{1.745815in}}%
\pgfpathcurveto{\pgfqpoint{1.672640in}{1.751639in}}{\pgfqpoint{1.670326in}{1.757225in}}{\pgfqpoint{1.666208in}{1.761343in}}%
\pgfpathcurveto{\pgfqpoint{1.662090in}{1.765462in}}{\pgfqpoint{1.656504in}{1.767775in}}{\pgfqpoint{1.650680in}{1.767775in}}%
\pgfpathcurveto{\pgfqpoint{1.644856in}{1.767775in}}{\pgfqpoint{1.639270in}{1.765462in}}{\pgfqpoint{1.635152in}{1.761343in}}%
\pgfpathcurveto{\pgfqpoint{1.631034in}{1.757225in}}{\pgfqpoint{1.628720in}{1.751639in}}{\pgfqpoint{1.628720in}{1.745815in}}%
\pgfpathcurveto{\pgfqpoint{1.628720in}{1.739991in}}{\pgfqpoint{1.631034in}{1.734405in}}{\pgfqpoint{1.635152in}{1.730287in}}%
\pgfpathcurveto{\pgfqpoint{1.639270in}{1.726169in}}{\pgfqpoint{1.644856in}{1.723855in}}{\pgfqpoint{1.650680in}{1.723855in}}%
\pgfpathlineto{\pgfqpoint{1.650680in}{1.723855in}}%
\pgfpathclose%
\pgfusepath{stroke,fill}%
\end{pgfscope}%
\begin{pgfscope}%
\pgfpathrectangle{\pgfqpoint{0.100000in}{0.209042in}}{\pgfqpoint{2.906250in}{2.906250in}}%
\pgfusepath{clip}%
\pgfsetbuttcap%
\pgfsetroundjoin%
\definecolor{currentfill}{rgb}{0.969680,0.446936,0.360311}%
\pgfsetfillcolor{currentfill}%
\pgfsetfillopacity{0.800000}%
\pgfsetlinewidth{1.003750pt}%
\definecolor{currentstroke}{rgb}{0.969680,0.446936,0.360311}%
\pgfsetstrokecolor{currentstroke}%
\pgfsetstrokeopacity{0.800000}%
\pgfsetdash{}{0pt}%
\pgfpathmoveto{\pgfqpoint{1.029143in}{1.349228in}}%
\pgfpathcurveto{\pgfqpoint{1.034967in}{1.349228in}}{\pgfqpoint{1.040553in}{1.351542in}}{\pgfqpoint{1.044671in}{1.355660in}}%
\pgfpathcurveto{\pgfqpoint{1.048789in}{1.359778in}}{\pgfqpoint{1.051103in}{1.365365in}}{\pgfqpoint{1.051103in}{1.371189in}}%
\pgfpathcurveto{\pgfqpoint{1.051103in}{1.377012in}}{\pgfqpoint{1.048789in}{1.382599in}}{\pgfqpoint{1.044671in}{1.386717in}}%
\pgfpathcurveto{\pgfqpoint{1.040553in}{1.390835in}}{\pgfqpoint{1.034967in}{1.393149in}}{\pgfqpoint{1.029143in}{1.393149in}}%
\pgfpathcurveto{\pgfqpoint{1.023319in}{1.393149in}}{\pgfqpoint{1.017733in}{1.390835in}}{\pgfqpoint{1.013615in}{1.386717in}}%
\pgfpathcurveto{\pgfqpoint{1.009496in}{1.382599in}}{\pgfqpoint{1.007183in}{1.377012in}}{\pgfqpoint{1.007183in}{1.371189in}}%
\pgfpathcurveto{\pgfqpoint{1.007183in}{1.365365in}}{\pgfqpoint{1.009496in}{1.359778in}}{\pgfqpoint{1.013615in}{1.355660in}}%
\pgfpathcurveto{\pgfqpoint{1.017733in}{1.351542in}}{\pgfqpoint{1.023319in}{1.349228in}}{\pgfqpoint{1.029143in}{1.349228in}}%
\pgfpathlineto{\pgfqpoint{1.029143in}{1.349228in}}%
\pgfpathclose%
\pgfusepath{stroke,fill}%
\end{pgfscope}%
\begin{pgfscope}%
\pgfpathrectangle{\pgfqpoint{0.100000in}{0.209042in}}{\pgfqpoint{2.906250in}{2.906250in}}%
\pgfusepath{clip}%
\pgfsetbuttcap%
\pgfsetroundjoin%
\definecolor{currentfill}{rgb}{0.284951,0.063168,0.472451}%
\pgfsetfillcolor{currentfill}%
\pgfsetfillopacity{0.800000}%
\pgfsetlinewidth{1.003750pt}%
\definecolor{currentstroke}{rgb}{0.284951,0.063168,0.472451}%
\pgfsetstrokecolor{currentstroke}%
\pgfsetstrokeopacity{0.800000}%
\pgfsetdash{}{0pt}%
\pgfpathmoveto{\pgfqpoint{1.614496in}{1.853070in}}%
\pgfpathcurveto{\pgfqpoint{1.620320in}{1.853070in}}{\pgfqpoint{1.625907in}{1.855384in}}{\pgfqpoint{1.630025in}{1.859502in}}%
\pgfpathcurveto{\pgfqpoint{1.634143in}{1.863620in}}{\pgfqpoint{1.636457in}{1.869206in}}{\pgfqpoint{1.636457in}{1.875030in}}%
\pgfpathcurveto{\pgfqpoint{1.636457in}{1.880854in}}{\pgfqpoint{1.634143in}{1.886440in}}{\pgfqpoint{1.630025in}{1.890558in}}%
\pgfpathcurveto{\pgfqpoint{1.625907in}{1.894677in}}{\pgfqpoint{1.620320in}{1.896990in}}{\pgfqpoint{1.614496in}{1.896990in}}%
\pgfpathcurveto{\pgfqpoint{1.608673in}{1.896990in}}{\pgfqpoint{1.603086in}{1.894677in}}{\pgfqpoint{1.598968in}{1.890558in}}%
\pgfpathcurveto{\pgfqpoint{1.594850in}{1.886440in}}{\pgfqpoint{1.592536in}{1.880854in}}{\pgfqpoint{1.592536in}{1.875030in}}%
\pgfpathcurveto{\pgfqpoint{1.592536in}{1.869206in}}{\pgfqpoint{1.594850in}{1.863620in}}{\pgfqpoint{1.598968in}{1.859502in}}%
\pgfpathcurveto{\pgfqpoint{1.603086in}{1.855384in}}{\pgfqpoint{1.608673in}{1.853070in}}{\pgfqpoint{1.614496in}{1.853070in}}%
\pgfpathlineto{\pgfqpoint{1.614496in}{1.853070in}}%
\pgfpathclose%
\pgfusepath{stroke,fill}%
\end{pgfscope}%
\begin{pgfscope}%
\pgfpathrectangle{\pgfqpoint{0.100000in}{0.209042in}}{\pgfqpoint{2.906250in}{2.906250in}}%
\pgfusepath{clip}%
\pgfsetbuttcap%
\pgfsetroundjoin%
\definecolor{currentfill}{rgb}{0.506629,0.145958,0.507806}%
\pgfsetfillcolor{currentfill}%
\pgfsetfillopacity{0.800000}%
\pgfsetlinewidth{1.003750pt}%
\definecolor{currentstroke}{rgb}{0.506629,0.145958,0.507806}%
\pgfsetstrokecolor{currentstroke}%
\pgfsetstrokeopacity{0.800000}%
\pgfsetdash{}{0pt}%
\pgfpathmoveto{\pgfqpoint{1.934798in}{1.446715in}}%
\pgfpathcurveto{\pgfqpoint{1.940622in}{1.446715in}}{\pgfqpoint{1.946208in}{1.449029in}}{\pgfqpoint{1.950326in}{1.453147in}}%
\pgfpathcurveto{\pgfqpoint{1.954444in}{1.457266in}}{\pgfqpoint{1.956758in}{1.462852in}}{\pgfqpoint{1.956758in}{1.468676in}}%
\pgfpathcurveto{\pgfqpoint{1.956758in}{1.474500in}}{\pgfqpoint{1.954444in}{1.480086in}}{\pgfqpoint{1.950326in}{1.484204in}}%
\pgfpathcurveto{\pgfqpoint{1.946208in}{1.488322in}}{\pgfqpoint{1.940622in}{1.490636in}}{\pgfqpoint{1.934798in}{1.490636in}}%
\pgfpathcurveto{\pgfqpoint{1.928974in}{1.490636in}}{\pgfqpoint{1.923388in}{1.488322in}}{\pgfqpoint{1.919270in}{1.484204in}}%
\pgfpathcurveto{\pgfqpoint{1.915152in}{1.480086in}}{\pgfqpoint{1.912838in}{1.474500in}}{\pgfqpoint{1.912838in}{1.468676in}}%
\pgfpathcurveto{\pgfqpoint{1.912838in}{1.462852in}}{\pgfqpoint{1.915152in}{1.457266in}}{\pgfqpoint{1.919270in}{1.453147in}}%
\pgfpathcurveto{\pgfqpoint{1.923388in}{1.449029in}}{\pgfqpoint{1.928974in}{1.446715in}}{\pgfqpoint{1.934798in}{1.446715in}}%
\pgfpathlineto{\pgfqpoint{1.934798in}{1.446715in}}%
\pgfpathclose%
\pgfusepath{stroke,fill}%
\end{pgfscope}%
\begin{pgfscope}%
\pgfpathrectangle{\pgfqpoint{0.100000in}{0.209042in}}{\pgfqpoint{2.906250in}{2.906250in}}%
\pgfusepath{clip}%
\pgfsetbuttcap%
\pgfsetroundjoin%
\definecolor{currentfill}{rgb}{0.390384,0.100379,0.501864}%
\pgfsetfillcolor{currentfill}%
\pgfsetfillopacity{0.800000}%
\pgfsetlinewidth{1.003750pt}%
\definecolor{currentstroke}{rgb}{0.390384,0.100379,0.501864}%
\pgfsetstrokecolor{currentstroke}%
\pgfsetstrokeopacity{0.800000}%
\pgfsetdash{}{0pt}%
\pgfpathmoveto{\pgfqpoint{1.775351in}{1.895778in}}%
\pgfpathcurveto{\pgfqpoint{1.781175in}{1.895778in}}{\pgfqpoint{1.786761in}{1.898092in}}{\pgfqpoint{1.790879in}{1.902210in}}%
\pgfpathcurveto{\pgfqpoint{1.794997in}{1.906328in}}{\pgfqpoint{1.797311in}{1.911915in}}{\pgfqpoint{1.797311in}{1.917739in}}%
\pgfpathcurveto{\pgfqpoint{1.797311in}{1.923562in}}{\pgfqpoint{1.794997in}{1.929149in}}{\pgfqpoint{1.790879in}{1.933267in}}%
\pgfpathcurveto{\pgfqpoint{1.786761in}{1.937385in}}{\pgfqpoint{1.781175in}{1.939699in}}{\pgfqpoint{1.775351in}{1.939699in}}%
\pgfpathcurveto{\pgfqpoint{1.769527in}{1.939699in}}{\pgfqpoint{1.763941in}{1.937385in}}{\pgfqpoint{1.759823in}{1.933267in}}%
\pgfpathcurveto{\pgfqpoint{1.755705in}{1.929149in}}{\pgfqpoint{1.753391in}{1.923562in}}{\pgfqpoint{1.753391in}{1.917739in}}%
\pgfpathcurveto{\pgfqpoint{1.753391in}{1.911915in}}{\pgfqpoint{1.755705in}{1.906328in}}{\pgfqpoint{1.759823in}{1.902210in}}%
\pgfpathcurveto{\pgfqpoint{1.763941in}{1.898092in}}{\pgfqpoint{1.769527in}{1.895778in}}{\pgfqpoint{1.775351in}{1.895778in}}%
\pgfpathlineto{\pgfqpoint{1.775351in}{1.895778in}}%
\pgfpathclose%
\pgfusepath{stroke,fill}%
\end{pgfscope}%
\begin{pgfscope}%
\pgfpathrectangle{\pgfqpoint{0.100000in}{0.209042in}}{\pgfqpoint{2.906250in}{2.906250in}}%
\pgfusepath{clip}%
\pgfsetbuttcap%
\pgfsetroundjoin%
\definecolor{currentfill}{rgb}{0.211718,0.061992,0.418647}%
\pgfsetfillcolor{currentfill}%
\pgfsetfillopacity{0.800000}%
\pgfsetlinewidth{1.003750pt}%
\definecolor{currentstroke}{rgb}{0.211718,0.061992,0.418647}%
\pgfsetstrokecolor{currentstroke}%
\pgfsetstrokeopacity{0.800000}%
\pgfsetdash{}{0pt}%
\pgfpathmoveto{\pgfqpoint{1.783035in}{1.644308in}}%
\pgfpathcurveto{\pgfqpoint{1.788859in}{1.644308in}}{\pgfqpoint{1.794445in}{1.646622in}}{\pgfqpoint{1.798563in}{1.650740in}}%
\pgfpathcurveto{\pgfqpoint{1.802681in}{1.654858in}}{\pgfqpoint{1.804995in}{1.660444in}}{\pgfqpoint{1.804995in}{1.666268in}}%
\pgfpathcurveto{\pgfqpoint{1.804995in}{1.672092in}}{\pgfqpoint{1.802681in}{1.677678in}}{\pgfqpoint{1.798563in}{1.681797in}}%
\pgfpathcurveto{\pgfqpoint{1.794445in}{1.685915in}}{\pgfqpoint{1.788859in}{1.688229in}}{\pgfqpoint{1.783035in}{1.688229in}}%
\pgfpathcurveto{\pgfqpoint{1.777211in}{1.688229in}}{\pgfqpoint{1.771625in}{1.685915in}}{\pgfqpoint{1.767507in}{1.681797in}}%
\pgfpathcurveto{\pgfqpoint{1.763389in}{1.677678in}}{\pgfqpoint{1.761075in}{1.672092in}}{\pgfqpoint{1.761075in}{1.666268in}}%
\pgfpathcurveto{\pgfqpoint{1.761075in}{1.660444in}}{\pgfqpoint{1.763389in}{1.654858in}}{\pgfqpoint{1.767507in}{1.650740in}}%
\pgfpathcurveto{\pgfqpoint{1.771625in}{1.646622in}}{\pgfqpoint{1.777211in}{1.644308in}}{\pgfqpoint{1.783035in}{1.644308in}}%
\pgfpathlineto{\pgfqpoint{1.783035in}{1.644308in}}%
\pgfpathclose%
\pgfusepath{stroke,fill}%
\end{pgfscope}%
\begin{pgfscope}%
\pgfpathrectangle{\pgfqpoint{0.100000in}{0.209042in}}{\pgfqpoint{2.906250in}{2.906250in}}%
\pgfusepath{clip}%
\pgfsetbuttcap%
\pgfsetroundjoin%
\definecolor{currentfill}{rgb}{0.804752,0.249911,0.442102}%
\pgfsetfillcolor{currentfill}%
\pgfsetfillopacity{0.800000}%
\pgfsetlinewidth{1.003750pt}%
\definecolor{currentstroke}{rgb}{0.804752,0.249911,0.442102}%
\pgfsetstrokecolor{currentstroke}%
\pgfsetstrokeopacity{0.800000}%
\pgfsetdash{}{0pt}%
\pgfpathmoveto{\pgfqpoint{1.508299in}{1.214238in}}%
\pgfpathcurveto{\pgfqpoint{1.514123in}{1.214238in}}{\pgfqpoint{1.519709in}{1.216552in}}{\pgfqpoint{1.523827in}{1.220670in}}%
\pgfpathcurveto{\pgfqpoint{1.527945in}{1.224788in}}{\pgfqpoint{1.530259in}{1.230374in}}{\pgfqpoint{1.530259in}{1.236198in}}%
\pgfpathcurveto{\pgfqpoint{1.530259in}{1.242022in}}{\pgfqpoint{1.527945in}{1.247608in}}{\pgfqpoint{1.523827in}{1.251726in}}%
\pgfpathcurveto{\pgfqpoint{1.519709in}{1.255845in}}{\pgfqpoint{1.514123in}{1.258158in}}{\pgfqpoint{1.508299in}{1.258158in}}%
\pgfpathcurveto{\pgfqpoint{1.502475in}{1.258158in}}{\pgfqpoint{1.496889in}{1.255845in}}{\pgfqpoint{1.492771in}{1.251726in}}%
\pgfpathcurveto{\pgfqpoint{1.488652in}{1.247608in}}{\pgfqpoint{1.486339in}{1.242022in}}{\pgfqpoint{1.486339in}{1.236198in}}%
\pgfpathcurveto{\pgfqpoint{1.486339in}{1.230374in}}{\pgfqpoint{1.488652in}{1.224788in}}{\pgfqpoint{1.492771in}{1.220670in}}%
\pgfpathcurveto{\pgfqpoint{1.496889in}{1.216552in}}{\pgfqpoint{1.502475in}{1.214238in}}{\pgfqpoint{1.508299in}{1.214238in}}%
\pgfpathlineto{\pgfqpoint{1.508299in}{1.214238in}}%
\pgfpathclose%
\pgfusepath{stroke,fill}%
\end{pgfscope}%
\begin{pgfscope}%
\pgfpathrectangle{\pgfqpoint{0.100000in}{0.209042in}}{\pgfqpoint{2.906250in}{2.906250in}}%
\pgfusepath{clip}%
\pgfsetbuttcap%
\pgfsetroundjoin%
\definecolor{currentfill}{rgb}{0.159018,0.068354,0.352688}%
\pgfsetfillcolor{currentfill}%
\pgfsetfillopacity{0.800000}%
\pgfsetlinewidth{1.003750pt}%
\definecolor{currentstroke}{rgb}{0.159018,0.068354,0.352688}%
\pgfsetstrokecolor{currentstroke}%
\pgfsetstrokeopacity{0.800000}%
\pgfsetdash{}{0pt}%
\pgfpathmoveto{\pgfqpoint{1.617163in}{1.673476in}}%
\pgfpathcurveto{\pgfqpoint{1.622987in}{1.673476in}}{\pgfqpoint{1.628573in}{1.675790in}}{\pgfqpoint{1.632692in}{1.679908in}}%
\pgfpathcurveto{\pgfqpoint{1.636810in}{1.684026in}}{\pgfqpoint{1.639124in}{1.689612in}}{\pgfqpoint{1.639124in}{1.695436in}}%
\pgfpathcurveto{\pgfqpoint{1.639124in}{1.701260in}}{\pgfqpoint{1.636810in}{1.706846in}}{\pgfqpoint{1.632692in}{1.710964in}}%
\pgfpathcurveto{\pgfqpoint{1.628573in}{1.715083in}}{\pgfqpoint{1.622987in}{1.717396in}}{\pgfqpoint{1.617163in}{1.717396in}}%
\pgfpathcurveto{\pgfqpoint{1.611339in}{1.717396in}}{\pgfqpoint{1.605753in}{1.715083in}}{\pgfqpoint{1.601635in}{1.710964in}}%
\pgfpathcurveto{\pgfqpoint{1.597517in}{1.706846in}}{\pgfqpoint{1.595203in}{1.701260in}}{\pgfqpoint{1.595203in}{1.695436in}}%
\pgfpathcurveto{\pgfqpoint{1.595203in}{1.689612in}}{\pgfqpoint{1.597517in}{1.684026in}}{\pgfqpoint{1.601635in}{1.679908in}}%
\pgfpathcurveto{\pgfqpoint{1.605753in}{1.675790in}}{\pgfqpoint{1.611339in}{1.673476in}}{\pgfqpoint{1.617163in}{1.673476in}}%
\pgfpathlineto{\pgfqpoint{1.617163in}{1.673476in}}%
\pgfpathclose%
\pgfusepath{stroke,fill}%
\end{pgfscope}%
\begin{pgfscope}%
\pgfpathrectangle{\pgfqpoint{0.100000in}{0.209042in}}{\pgfqpoint{2.906250in}{2.906250in}}%
\pgfusepath{clip}%
\pgfsetbuttcap%
\pgfsetroundjoin%
\definecolor{currentfill}{rgb}{0.322899,0.073782,0.487408}%
\pgfsetfillcolor{currentfill}%
\pgfsetfillopacity{0.800000}%
\pgfsetlinewidth{1.003750pt}%
\definecolor{currentstroke}{rgb}{0.322899,0.073782,0.487408}%
\pgfsetstrokecolor{currentstroke}%
\pgfsetstrokeopacity{0.800000}%
\pgfsetdash{}{0pt}%
\pgfpathmoveto{\pgfqpoint{1.416026in}{1.668120in}}%
\pgfpathcurveto{\pgfqpoint{1.421850in}{1.668120in}}{\pgfqpoint{1.427436in}{1.670434in}}{\pgfqpoint{1.431554in}{1.674552in}}%
\pgfpathcurveto{\pgfqpoint{1.435672in}{1.678670in}}{\pgfqpoint{1.437986in}{1.684257in}}{\pgfqpoint{1.437986in}{1.690080in}}%
\pgfpathcurveto{\pgfqpoint{1.437986in}{1.695904in}}{\pgfqpoint{1.435672in}{1.701491in}}{\pgfqpoint{1.431554in}{1.705609in}}%
\pgfpathcurveto{\pgfqpoint{1.427436in}{1.709727in}}{\pgfqpoint{1.421850in}{1.712041in}}{\pgfqpoint{1.416026in}{1.712041in}}%
\pgfpathcurveto{\pgfqpoint{1.410202in}{1.712041in}}{\pgfqpoint{1.404615in}{1.709727in}}{\pgfqpoint{1.400497in}{1.705609in}}%
\pgfpathcurveto{\pgfqpoint{1.396379in}{1.701491in}}{\pgfqpoint{1.394065in}{1.695904in}}{\pgfqpoint{1.394065in}{1.690080in}}%
\pgfpathcurveto{\pgfqpoint{1.394065in}{1.684257in}}{\pgfqpoint{1.396379in}{1.678670in}}{\pgfqpoint{1.400497in}{1.674552in}}%
\pgfpathcurveto{\pgfqpoint{1.404615in}{1.670434in}}{\pgfqpoint{1.410202in}{1.668120in}}{\pgfqpoint{1.416026in}{1.668120in}}%
\pgfpathlineto{\pgfqpoint{1.416026in}{1.668120in}}%
\pgfpathclose%
\pgfusepath{stroke,fill}%
\end{pgfscope}%
\begin{pgfscope}%
\pgfpathrectangle{\pgfqpoint{0.100000in}{0.209042in}}{\pgfqpoint{2.906250in}{2.906250in}}%
\pgfusepath{clip}%
\pgfsetbuttcap%
\pgfsetroundjoin%
\definecolor{currentfill}{rgb}{0.639216,0.189921,0.494150}%
\pgfsetfillcolor{currentfill}%
\pgfsetfillopacity{0.800000}%
\pgfsetlinewidth{1.003750pt}%
\definecolor{currentstroke}{rgb}{0.639216,0.189921,0.494150}%
\pgfsetstrokecolor{currentstroke}%
\pgfsetstrokeopacity{0.800000}%
\pgfsetdash{}{0pt}%
\pgfpathmoveto{\pgfqpoint{1.521162in}{1.306983in}}%
\pgfpathcurveto{\pgfqpoint{1.526986in}{1.306983in}}{\pgfqpoint{1.532572in}{1.309296in}}{\pgfqpoint{1.536690in}{1.313415in}}%
\pgfpathcurveto{\pgfqpoint{1.540809in}{1.317533in}}{\pgfqpoint{1.543122in}{1.323119in}}{\pgfqpoint{1.543122in}{1.328943in}}%
\pgfpathcurveto{\pgfqpoint{1.543122in}{1.334767in}}{\pgfqpoint{1.540809in}{1.340353in}}{\pgfqpoint{1.536690in}{1.344471in}}%
\pgfpathcurveto{\pgfqpoint{1.532572in}{1.348589in}}{\pgfqpoint{1.526986in}{1.350903in}}{\pgfqpoint{1.521162in}{1.350903in}}%
\pgfpathcurveto{\pgfqpoint{1.515338in}{1.350903in}}{\pgfqpoint{1.509752in}{1.348589in}}{\pgfqpoint{1.505634in}{1.344471in}}%
\pgfpathcurveto{\pgfqpoint{1.501516in}{1.340353in}}{\pgfqpoint{1.499202in}{1.334767in}}{\pgfqpoint{1.499202in}{1.328943in}}%
\pgfpathcurveto{\pgfqpoint{1.499202in}{1.323119in}}{\pgfqpoint{1.501516in}{1.317533in}}{\pgfqpoint{1.505634in}{1.313415in}}%
\pgfpathcurveto{\pgfqpoint{1.509752in}{1.309296in}}{\pgfqpoint{1.515338in}{1.306983in}}{\pgfqpoint{1.521162in}{1.306983in}}%
\pgfpathlineto{\pgfqpoint{1.521162in}{1.306983in}}%
\pgfpathclose%
\pgfusepath{stroke,fill}%
\end{pgfscope}%
\begin{pgfscope}%
\pgfpathrectangle{\pgfqpoint{0.100000in}{0.209042in}}{\pgfqpoint{2.906250in}{2.906250in}}%
\pgfusepath{clip}%
\pgfsetbuttcap%
\pgfsetroundjoin%
\definecolor{currentfill}{rgb}{0.671349,0.200133,0.487358}%
\pgfsetfillcolor{currentfill}%
\pgfsetfillopacity{0.800000}%
\pgfsetlinewidth{1.003750pt}%
\definecolor{currentstroke}{rgb}{0.671349,0.200133,0.487358}%
\pgfsetstrokecolor{currentstroke}%
\pgfsetstrokeopacity{0.800000}%
\pgfsetdash{}{0pt}%
\pgfpathmoveto{\pgfqpoint{2.117871in}{1.805305in}}%
\pgfpathcurveto{\pgfqpoint{2.123695in}{1.805305in}}{\pgfqpoint{2.129281in}{1.807619in}}{\pgfqpoint{2.133399in}{1.811737in}}%
\pgfpathcurveto{\pgfqpoint{2.137517in}{1.815855in}}{\pgfqpoint{2.139831in}{1.821442in}}{\pgfqpoint{2.139831in}{1.827266in}}%
\pgfpathcurveto{\pgfqpoint{2.139831in}{1.833089in}}{\pgfqpoint{2.137517in}{1.838676in}}{\pgfqpoint{2.133399in}{1.842794in}}%
\pgfpathcurveto{\pgfqpoint{2.129281in}{1.846912in}}{\pgfqpoint{2.123695in}{1.849226in}}{\pgfqpoint{2.117871in}{1.849226in}}%
\pgfpathcurveto{\pgfqpoint{2.112047in}{1.849226in}}{\pgfqpoint{2.106461in}{1.846912in}}{\pgfqpoint{2.102343in}{1.842794in}}%
\pgfpathcurveto{\pgfqpoint{2.098225in}{1.838676in}}{\pgfqpoint{2.095911in}{1.833089in}}{\pgfqpoint{2.095911in}{1.827266in}}%
\pgfpathcurveto{\pgfqpoint{2.095911in}{1.821442in}}{\pgfqpoint{2.098225in}{1.815855in}}{\pgfqpoint{2.102343in}{1.811737in}}%
\pgfpathcurveto{\pgfqpoint{2.106461in}{1.807619in}}{\pgfqpoint{2.112047in}{1.805305in}}{\pgfqpoint{2.117871in}{1.805305in}}%
\pgfpathlineto{\pgfqpoint{2.117871in}{1.805305in}}%
\pgfpathclose%
\pgfusepath{stroke,fill}%
\end{pgfscope}%
\begin{pgfscope}%
\pgfpathrectangle{\pgfqpoint{0.100000in}{0.209042in}}{\pgfqpoint{2.906250in}{2.906250in}}%
\pgfusepath{clip}%
\pgfsetbuttcap%
\pgfsetroundjoin%
\definecolor{currentfill}{rgb}{0.322899,0.073782,0.487408}%
\pgfsetfillcolor{currentfill}%
\pgfsetfillopacity{0.800000}%
\pgfsetlinewidth{1.003750pt}%
\definecolor{currentstroke}{rgb}{0.322899,0.073782,0.487408}%
\pgfsetstrokecolor{currentstroke}%
\pgfsetstrokeopacity{0.800000}%
\pgfsetdash{}{0pt}%
\pgfpathmoveto{\pgfqpoint{1.455059in}{1.775796in}}%
\pgfpathcurveto{\pgfqpoint{1.460883in}{1.775796in}}{\pgfqpoint{1.466469in}{1.778110in}}{\pgfqpoint{1.470588in}{1.782228in}}%
\pgfpathcurveto{\pgfqpoint{1.474706in}{1.786346in}}{\pgfqpoint{1.477020in}{1.791933in}}{\pgfqpoint{1.477020in}{1.797756in}}%
\pgfpathcurveto{\pgfqpoint{1.477020in}{1.803580in}}{\pgfqpoint{1.474706in}{1.809167in}}{\pgfqpoint{1.470588in}{1.813285in}}%
\pgfpathcurveto{\pgfqpoint{1.466469in}{1.817403in}}{\pgfqpoint{1.460883in}{1.819717in}}{\pgfqpoint{1.455059in}{1.819717in}}%
\pgfpathcurveto{\pgfqpoint{1.449235in}{1.819717in}}{\pgfqpoint{1.443649in}{1.817403in}}{\pgfqpoint{1.439531in}{1.813285in}}%
\pgfpathcurveto{\pgfqpoint{1.435413in}{1.809167in}}{\pgfqpoint{1.433099in}{1.803580in}}{\pgfqpoint{1.433099in}{1.797756in}}%
\pgfpathcurveto{\pgfqpoint{1.433099in}{1.791933in}}{\pgfqpoint{1.435413in}{1.786346in}}{\pgfqpoint{1.439531in}{1.782228in}}%
\pgfpathcurveto{\pgfqpoint{1.443649in}{1.778110in}}{\pgfqpoint{1.449235in}{1.775796in}}{\pgfqpoint{1.455059in}{1.775796in}}%
\pgfpathlineto{\pgfqpoint{1.455059in}{1.775796in}}%
\pgfpathclose%
\pgfusepath{stroke,fill}%
\end{pgfscope}%
\begin{pgfscope}%
\pgfpathrectangle{\pgfqpoint{0.100000in}{0.209042in}}{\pgfqpoint{2.906250in}{2.906250in}}%
\pgfusepath{clip}%
\pgfsetbuttcap%
\pgfsetroundjoin%
\definecolor{currentfill}{rgb}{0.457386,0.127522,0.507448}%
\pgfsetfillcolor{currentfill}%
\pgfsetfillopacity{0.800000}%
\pgfsetlinewidth{1.003750pt}%
\definecolor{currentstroke}{rgb}{0.457386,0.127522,0.507448}%
\pgfsetstrokecolor{currentstroke}%
\pgfsetstrokeopacity{0.800000}%
\pgfsetdash{}{0pt}%
\pgfpathmoveto{\pgfqpoint{1.324767in}{1.769481in}}%
\pgfpathcurveto{\pgfqpoint{1.330591in}{1.769481in}}{\pgfqpoint{1.336177in}{1.771795in}}{\pgfqpoint{1.340295in}{1.775913in}}%
\pgfpathcurveto{\pgfqpoint{1.344413in}{1.780031in}}{\pgfqpoint{1.346727in}{1.785617in}}{\pgfqpoint{1.346727in}{1.791441in}}%
\pgfpathcurveto{\pgfqpoint{1.346727in}{1.797265in}}{\pgfqpoint{1.344413in}{1.802852in}}{\pgfqpoint{1.340295in}{1.806970in}}%
\pgfpathcurveto{\pgfqpoint{1.336177in}{1.811088in}}{\pgfqpoint{1.330591in}{1.813402in}}{\pgfqpoint{1.324767in}{1.813402in}}%
\pgfpathcurveto{\pgfqpoint{1.318943in}{1.813402in}}{\pgfqpoint{1.313357in}{1.811088in}}{\pgfqpoint{1.309239in}{1.806970in}}%
\pgfpathcurveto{\pgfqpoint{1.305121in}{1.802852in}}{\pgfqpoint{1.302807in}{1.797265in}}{\pgfqpoint{1.302807in}{1.791441in}}%
\pgfpathcurveto{\pgfqpoint{1.302807in}{1.785617in}}{\pgfqpoint{1.305121in}{1.780031in}}{\pgfqpoint{1.309239in}{1.775913in}}%
\pgfpathcurveto{\pgfqpoint{1.313357in}{1.771795in}}{\pgfqpoint{1.318943in}{1.769481in}}{\pgfqpoint{1.324767in}{1.769481in}}%
\pgfpathlineto{\pgfqpoint{1.324767in}{1.769481in}}%
\pgfpathclose%
\pgfusepath{stroke,fill}%
\end{pgfscope}%
\begin{pgfscope}%
\pgfpathrectangle{\pgfqpoint{0.100000in}{0.209042in}}{\pgfqpoint{2.906250in}{2.906250in}}%
\pgfusepath{clip}%
\pgfsetbuttcap%
\pgfsetroundjoin%
\definecolor{currentfill}{rgb}{0.600868,0.177743,0.500394}%
\pgfsetfillcolor{currentfill}%
\pgfsetfillopacity{0.800000}%
\pgfsetlinewidth{1.003750pt}%
\definecolor{currentstroke}{rgb}{0.600868,0.177743,0.500394}%
\pgfsetstrokecolor{currentstroke}%
\pgfsetstrokeopacity{0.800000}%
\pgfsetdash{}{0pt}%
\pgfpathmoveto{\pgfqpoint{2.049348in}{1.483435in}}%
\pgfpathcurveto{\pgfqpoint{2.055172in}{1.483435in}}{\pgfqpoint{2.060758in}{1.485749in}}{\pgfqpoint{2.064876in}{1.489867in}}%
\pgfpathcurveto{\pgfqpoint{2.068994in}{1.493985in}}{\pgfqpoint{2.071308in}{1.499571in}}{\pgfqpoint{2.071308in}{1.505395in}}%
\pgfpathcurveto{\pgfqpoint{2.071308in}{1.511219in}}{\pgfqpoint{2.068994in}{1.516805in}}{\pgfqpoint{2.064876in}{1.520924in}}%
\pgfpathcurveto{\pgfqpoint{2.060758in}{1.525042in}}{\pgfqpoint{2.055172in}{1.527356in}}{\pgfqpoint{2.049348in}{1.527356in}}%
\pgfpathcurveto{\pgfqpoint{2.043524in}{1.527356in}}{\pgfqpoint{2.037938in}{1.525042in}}{\pgfqpoint{2.033820in}{1.520924in}}%
\pgfpathcurveto{\pgfqpoint{2.029702in}{1.516805in}}{\pgfqpoint{2.027388in}{1.511219in}}{\pgfqpoint{2.027388in}{1.505395in}}%
\pgfpathcurveto{\pgfqpoint{2.027388in}{1.499571in}}{\pgfqpoint{2.029702in}{1.493985in}}{\pgfqpoint{2.033820in}{1.489867in}}%
\pgfpathcurveto{\pgfqpoint{2.037938in}{1.485749in}}{\pgfqpoint{2.043524in}{1.483435in}}{\pgfqpoint{2.049348in}{1.483435in}}%
\pgfpathlineto{\pgfqpoint{2.049348in}{1.483435in}}%
\pgfpathclose%
\pgfusepath{stroke,fill}%
\end{pgfscope}%
\begin{pgfscope}%
\pgfpathrectangle{\pgfqpoint{0.100000in}{0.209042in}}{\pgfqpoint{2.906250in}{2.906250in}}%
\pgfusepath{clip}%
\pgfsetbuttcap%
\pgfsetroundjoin%
\definecolor{currentfill}{rgb}{0.512831,0.148179,0.507648}%
\pgfsetfillcolor{currentfill}%
\pgfsetfillopacity{0.800000}%
\pgfsetlinewidth{1.003750pt}%
\definecolor{currentstroke}{rgb}{0.512831,0.148179,0.507648}%
\pgfsetstrokecolor{currentstroke}%
\pgfsetstrokeopacity{0.800000}%
\pgfsetdash{}{0pt}%
\pgfpathmoveto{\pgfqpoint{1.771555in}{1.970565in}}%
\pgfpathcurveto{\pgfqpoint{1.777379in}{1.970565in}}{\pgfqpoint{1.782966in}{1.972879in}}{\pgfqpoint{1.787084in}{1.976997in}}%
\pgfpathcurveto{\pgfqpoint{1.791202in}{1.981116in}}{\pgfqpoint{1.793516in}{1.986702in}}{\pgfqpoint{1.793516in}{1.992526in}}%
\pgfpathcurveto{\pgfqpoint{1.793516in}{1.998350in}}{\pgfqpoint{1.791202in}{2.003936in}}{\pgfqpoint{1.787084in}{2.008054in}}%
\pgfpathcurveto{\pgfqpoint{1.782966in}{2.012172in}}{\pgfqpoint{1.777379in}{2.014486in}}{\pgfqpoint{1.771555in}{2.014486in}}%
\pgfpathcurveto{\pgfqpoint{1.765732in}{2.014486in}}{\pgfqpoint{1.760145in}{2.012172in}}{\pgfqpoint{1.756027in}{2.008054in}}%
\pgfpathcurveto{\pgfqpoint{1.751909in}{2.003936in}}{\pgfqpoint{1.749595in}{1.998350in}}{\pgfqpoint{1.749595in}{1.992526in}}%
\pgfpathcurveto{\pgfqpoint{1.749595in}{1.986702in}}{\pgfqpoint{1.751909in}{1.981116in}}{\pgfqpoint{1.756027in}{1.976997in}}%
\pgfpathcurveto{\pgfqpoint{1.760145in}{1.972879in}}{\pgfqpoint{1.765732in}{1.970565in}}{\pgfqpoint{1.771555in}{1.970565in}}%
\pgfpathlineto{\pgfqpoint{1.771555in}{1.970565in}}%
\pgfpathclose%
\pgfusepath{stroke,fill}%
\end{pgfscope}%
\begin{pgfscope}%
\pgfpathrectangle{\pgfqpoint{0.100000in}{0.209042in}}{\pgfqpoint{2.906250in}{2.906250in}}%
\pgfusepath{clip}%
\pgfsetbuttcap%
\pgfsetroundjoin%
\definecolor{currentfill}{rgb}{0.889731,0.305079,0.397002}%
\pgfsetfillcolor{currentfill}%
\pgfsetfillopacity{0.800000}%
\pgfsetlinewidth{1.003750pt}%
\definecolor{currentstroke}{rgb}{0.889731,0.305079,0.397002}%
\pgfsetstrokecolor{currentstroke}%
\pgfsetstrokeopacity{0.800000}%
\pgfsetdash{}{0pt}%
\pgfpathmoveto{\pgfqpoint{1.013453in}{1.829518in}}%
\pgfpathcurveto{\pgfqpoint{1.019277in}{1.829518in}}{\pgfqpoint{1.024863in}{1.831832in}}{\pgfqpoint{1.028981in}{1.835950in}}%
\pgfpathcurveto{\pgfqpoint{1.033099in}{1.840068in}}{\pgfqpoint{1.035413in}{1.845655in}}{\pgfqpoint{1.035413in}{1.851478in}}%
\pgfpathcurveto{\pgfqpoint{1.035413in}{1.857302in}}{\pgfqpoint{1.033099in}{1.862889in}}{\pgfqpoint{1.028981in}{1.867007in}}%
\pgfpathcurveto{\pgfqpoint{1.024863in}{1.871125in}}{\pgfqpoint{1.019277in}{1.873439in}}{\pgfqpoint{1.013453in}{1.873439in}}%
\pgfpathcurveto{\pgfqpoint{1.007629in}{1.873439in}}{\pgfqpoint{1.002043in}{1.871125in}}{\pgfqpoint{0.997925in}{1.867007in}}%
\pgfpathcurveto{\pgfqpoint{0.993806in}{1.862889in}}{\pgfqpoint{0.991493in}{1.857302in}}{\pgfqpoint{0.991493in}{1.851478in}}%
\pgfpathcurveto{\pgfqpoint{0.991493in}{1.845655in}}{\pgfqpoint{0.993806in}{1.840068in}}{\pgfqpoint{0.997925in}{1.835950in}}%
\pgfpathcurveto{\pgfqpoint{1.002043in}{1.831832in}}{\pgfqpoint{1.007629in}{1.829518in}}{\pgfqpoint{1.013453in}{1.829518in}}%
\pgfpathlineto{\pgfqpoint{1.013453in}{1.829518in}}%
\pgfpathclose%
\pgfusepath{stroke,fill}%
\end{pgfscope}%
\begin{pgfscope}%
\pgfpathrectangle{\pgfqpoint{0.100000in}{0.209042in}}{\pgfqpoint{2.906250in}{2.906250in}}%
\pgfusepath{clip}%
\pgfsetbuttcap%
\pgfsetroundjoin%
\definecolor{currentfill}{rgb}{0.347636,0.082946,0.494121}%
\pgfsetfillcolor{currentfill}%
\pgfsetfillopacity{0.800000}%
\pgfsetlinewidth{1.003750pt}%
\definecolor{currentstroke}{rgb}{0.347636,0.082946,0.494121}%
\pgfsetstrokecolor{currentstroke}%
\pgfsetstrokeopacity{0.800000}%
\pgfsetdash{}{0pt}%
\pgfpathmoveto{\pgfqpoint{1.415772in}{1.735324in}}%
\pgfpathcurveto{\pgfqpoint{1.421596in}{1.735324in}}{\pgfqpoint{1.427182in}{1.737638in}}{\pgfqpoint{1.431300in}{1.741756in}}%
\pgfpathcurveto{\pgfqpoint{1.435418in}{1.745874in}}{\pgfqpoint{1.437732in}{1.751461in}}{\pgfqpoint{1.437732in}{1.757285in}}%
\pgfpathcurveto{\pgfqpoint{1.437732in}{1.763109in}}{\pgfqpoint{1.435418in}{1.768695in}}{\pgfqpoint{1.431300in}{1.772813in}}%
\pgfpathcurveto{\pgfqpoint{1.427182in}{1.776931in}}{\pgfqpoint{1.421596in}{1.779245in}}{\pgfqpoint{1.415772in}{1.779245in}}%
\pgfpathcurveto{\pgfqpoint{1.409948in}{1.779245in}}{\pgfqpoint{1.404362in}{1.776931in}}{\pgfqpoint{1.400243in}{1.772813in}}%
\pgfpathcurveto{\pgfqpoint{1.396125in}{1.768695in}}{\pgfqpoint{1.393811in}{1.763109in}}{\pgfqpoint{1.393811in}{1.757285in}}%
\pgfpathcurveto{\pgfqpoint{1.393811in}{1.751461in}}{\pgfqpoint{1.396125in}{1.745874in}}{\pgfqpoint{1.400243in}{1.741756in}}%
\pgfpathcurveto{\pgfqpoint{1.404362in}{1.737638in}}{\pgfqpoint{1.409948in}{1.735324in}}{\pgfqpoint{1.415772in}{1.735324in}}%
\pgfpathlineto{\pgfqpoint{1.415772in}{1.735324in}}%
\pgfpathclose%
\pgfusepath{stroke,fill}%
\end{pgfscope}%
\begin{pgfscope}%
\pgfpathrectangle{\pgfqpoint{0.100000in}{0.209042in}}{\pgfqpoint{2.906250in}{2.906250in}}%
\pgfusepath{clip}%
\pgfsetbuttcap%
\pgfsetroundjoin%
\definecolor{currentfill}{rgb}{0.439062,0.120298,0.506555}%
\pgfsetfillcolor{currentfill}%
\pgfsetfillopacity{0.800000}%
\pgfsetlinewidth{1.003750pt}%
\definecolor{currentstroke}{rgb}{0.439062,0.120298,0.506555}%
\pgfsetstrokecolor{currentstroke}%
\pgfsetstrokeopacity{0.800000}%
\pgfsetdash{}{0pt}%
\pgfpathmoveto{\pgfqpoint{1.960306in}{1.742430in}}%
\pgfpathcurveto{\pgfqpoint{1.966130in}{1.742430in}}{\pgfqpoint{1.971716in}{1.744744in}}{\pgfqpoint{1.975834in}{1.748862in}}%
\pgfpathcurveto{\pgfqpoint{1.979952in}{1.752980in}}{\pgfqpoint{1.982266in}{1.758566in}}{\pgfqpoint{1.982266in}{1.764390in}}%
\pgfpathcurveto{\pgfqpoint{1.982266in}{1.770214in}}{\pgfqpoint{1.979952in}{1.775800in}}{\pgfqpoint{1.975834in}{1.779918in}}%
\pgfpathcurveto{\pgfqpoint{1.971716in}{1.784036in}}{\pgfqpoint{1.966130in}{1.786350in}}{\pgfqpoint{1.960306in}{1.786350in}}%
\pgfpathcurveto{\pgfqpoint{1.954482in}{1.786350in}}{\pgfqpoint{1.948896in}{1.784036in}}{\pgfqpoint{1.944778in}{1.779918in}}%
\pgfpathcurveto{\pgfqpoint{1.940659in}{1.775800in}}{\pgfqpoint{1.938346in}{1.770214in}}{\pgfqpoint{1.938346in}{1.764390in}}%
\pgfpathcurveto{\pgfqpoint{1.938346in}{1.758566in}}{\pgfqpoint{1.940659in}{1.752980in}}{\pgfqpoint{1.944778in}{1.748862in}}%
\pgfpathcurveto{\pgfqpoint{1.948896in}{1.744744in}}{\pgfqpoint{1.954482in}{1.742430in}}{\pgfqpoint{1.960306in}{1.742430in}}%
\pgfpathlineto{\pgfqpoint{1.960306in}{1.742430in}}%
\pgfpathclose%
\pgfusepath{stroke,fill}%
\end{pgfscope}%
\begin{pgfscope}%
\pgfpathrectangle{\pgfqpoint{0.100000in}{0.209042in}}{\pgfqpoint{2.906250in}{2.906250in}}%
\pgfusepath{clip}%
\pgfsetbuttcap%
\pgfsetroundjoin%
\definecolor{currentfill}{rgb}{0.378211,0.095332,0.500067}%
\pgfsetfillcolor{currentfill}%
\pgfsetfillopacity{0.800000}%
\pgfsetlinewidth{1.003750pt}%
\definecolor{currentstroke}{rgb}{0.378211,0.095332,0.500067}%
\pgfsetstrokecolor{currentstroke}%
\pgfsetstrokeopacity{0.800000}%
\pgfsetdash{}{0pt}%
\pgfpathmoveto{\pgfqpoint{1.687077in}{1.892680in}}%
\pgfpathcurveto{\pgfqpoint{1.692900in}{1.892680in}}{\pgfqpoint{1.698487in}{1.894994in}}{\pgfqpoint{1.702605in}{1.899112in}}%
\pgfpathcurveto{\pgfqpoint{1.706723in}{1.903230in}}{\pgfqpoint{1.709037in}{1.908816in}}{\pgfqpoint{1.709037in}{1.914640in}}%
\pgfpathcurveto{\pgfqpoint{1.709037in}{1.920464in}}{\pgfqpoint{1.706723in}{1.926050in}}{\pgfqpoint{1.702605in}{1.930169in}}%
\pgfpathcurveto{\pgfqpoint{1.698487in}{1.934287in}}{\pgfqpoint{1.692900in}{1.936601in}}{\pgfqpoint{1.687077in}{1.936601in}}%
\pgfpathcurveto{\pgfqpoint{1.681253in}{1.936601in}}{\pgfqpoint{1.675666in}{1.934287in}}{\pgfqpoint{1.671548in}{1.930169in}}%
\pgfpathcurveto{\pgfqpoint{1.667430in}{1.926050in}}{\pgfqpoint{1.665116in}{1.920464in}}{\pgfqpoint{1.665116in}{1.914640in}}%
\pgfpathcurveto{\pgfqpoint{1.665116in}{1.908816in}}{\pgfqpoint{1.667430in}{1.903230in}}{\pgfqpoint{1.671548in}{1.899112in}}%
\pgfpathcurveto{\pgfqpoint{1.675666in}{1.894994in}}{\pgfqpoint{1.681253in}{1.892680in}}{\pgfqpoint{1.687077in}{1.892680in}}%
\pgfpathlineto{\pgfqpoint{1.687077in}{1.892680in}}%
\pgfpathclose%
\pgfusepath{stroke,fill}%
\end{pgfscope}%
\begin{pgfscope}%
\pgfpathrectangle{\pgfqpoint{0.100000in}{0.209042in}}{\pgfqpoint{2.906250in}{2.906250in}}%
\pgfusepath{clip}%
\pgfsetbuttcap%
\pgfsetroundjoin%
\definecolor{currentfill}{rgb}{0.329114,0.075972,0.489287}%
\pgfsetfillcolor{currentfill}%
\pgfsetfillopacity{0.800000}%
\pgfsetlinewidth{1.003750pt}%
\definecolor{currentstroke}{rgb}{0.329114,0.075972,0.489287}%
\pgfsetstrokecolor{currentstroke}%
\pgfsetstrokeopacity{0.800000}%
\pgfsetdash{}{0pt}%
\pgfpathmoveto{\pgfqpoint{1.821462in}{1.553092in}}%
\pgfpathcurveto{\pgfqpoint{1.827286in}{1.553092in}}{\pgfqpoint{1.832872in}{1.555406in}}{\pgfqpoint{1.836990in}{1.559524in}}%
\pgfpathcurveto{\pgfqpoint{1.841108in}{1.563642in}}{\pgfqpoint{1.843422in}{1.569228in}}{\pgfqpoint{1.843422in}{1.575052in}}%
\pgfpathcurveto{\pgfqpoint{1.843422in}{1.580876in}}{\pgfqpoint{1.841108in}{1.586462in}}{\pgfqpoint{1.836990in}{1.590580in}}%
\pgfpathcurveto{\pgfqpoint{1.832872in}{1.594699in}}{\pgfqpoint{1.827286in}{1.597012in}}{\pgfqpoint{1.821462in}{1.597012in}}%
\pgfpathcurveto{\pgfqpoint{1.815638in}{1.597012in}}{\pgfqpoint{1.810052in}{1.594699in}}{\pgfqpoint{1.805933in}{1.590580in}}%
\pgfpathcurveto{\pgfqpoint{1.801815in}{1.586462in}}{\pgfqpoint{1.799501in}{1.580876in}}{\pgfqpoint{1.799501in}{1.575052in}}%
\pgfpathcurveto{\pgfqpoint{1.799501in}{1.569228in}}{\pgfqpoint{1.801815in}{1.563642in}}{\pgfqpoint{1.805933in}{1.559524in}}%
\pgfpathcurveto{\pgfqpoint{1.810052in}{1.555406in}}{\pgfqpoint{1.815638in}{1.553092in}}{\pgfqpoint{1.821462in}{1.553092in}}%
\pgfpathlineto{\pgfqpoint{1.821462in}{1.553092in}}%
\pgfpathclose%
\pgfusepath{stroke,fill}%
\end{pgfscope}%
\begin{pgfscope}%
\pgfpathrectangle{\pgfqpoint{0.100000in}{0.209042in}}{\pgfqpoint{2.906250in}{2.906250in}}%
\pgfusepath{clip}%
\pgfsetbuttcap%
\pgfsetroundjoin%
\definecolor{currentfill}{rgb}{0.238826,0.059517,0.443256}%
\pgfsetfillcolor{currentfill}%
\pgfsetfillopacity{0.800000}%
\pgfsetlinewidth{1.003750pt}%
\definecolor{currentstroke}{rgb}{0.238826,0.059517,0.443256}%
\pgfsetstrokecolor{currentstroke}%
\pgfsetstrokeopacity{0.800000}%
\pgfsetdash{}{0pt}%
\pgfpathmoveto{\pgfqpoint{1.706632in}{1.758701in}}%
\pgfpathcurveto{\pgfqpoint{1.712456in}{1.758701in}}{\pgfqpoint{1.718042in}{1.761015in}}{\pgfqpoint{1.722160in}{1.765133in}}%
\pgfpathcurveto{\pgfqpoint{1.726279in}{1.769251in}}{\pgfqpoint{1.728593in}{1.774837in}}{\pgfqpoint{1.728593in}{1.780661in}}%
\pgfpathcurveto{\pgfqpoint{1.728593in}{1.786485in}}{\pgfqpoint{1.726279in}{1.792071in}}{\pgfqpoint{1.722160in}{1.796189in}}%
\pgfpathcurveto{\pgfqpoint{1.718042in}{1.800307in}}{\pgfqpoint{1.712456in}{1.802621in}}{\pgfqpoint{1.706632in}{1.802621in}}%
\pgfpathcurveto{\pgfqpoint{1.700808in}{1.802621in}}{\pgfqpoint{1.695222in}{1.800307in}}{\pgfqpoint{1.691104in}{1.796189in}}%
\pgfpathcurveto{\pgfqpoint{1.686986in}{1.792071in}}{\pgfqpoint{1.684672in}{1.786485in}}{\pgfqpoint{1.684672in}{1.780661in}}%
\pgfpathcurveto{\pgfqpoint{1.684672in}{1.774837in}}{\pgfqpoint{1.686986in}{1.769251in}}{\pgfqpoint{1.691104in}{1.765133in}}%
\pgfpathcurveto{\pgfqpoint{1.695222in}{1.761015in}}{\pgfqpoint{1.700808in}{1.758701in}}{\pgfqpoint{1.706632in}{1.758701in}}%
\pgfpathlineto{\pgfqpoint{1.706632in}{1.758701in}}%
\pgfpathclose%
\pgfusepath{stroke,fill}%
\end{pgfscope}%
\begin{pgfscope}%
\pgfpathrectangle{\pgfqpoint{0.100000in}{0.209042in}}{\pgfqpoint{2.906250in}{2.906250in}}%
\pgfusepath{clip}%
\pgfsetbuttcap%
\pgfsetroundjoin%
\definecolor{currentfill}{rgb}{0.297740,0.066117,0.478243}%
\pgfsetfillcolor{currentfill}%
\pgfsetfillopacity{0.800000}%
\pgfsetlinewidth{1.003750pt}%
\definecolor{currentstroke}{rgb}{0.297740,0.066117,0.478243}%
\pgfsetstrokecolor{currentstroke}%
\pgfsetstrokeopacity{0.800000}%
\pgfsetdash{}{0pt}%
\pgfpathmoveto{\pgfqpoint{1.667163in}{1.525472in}}%
\pgfpathcurveto{\pgfqpoint{1.672987in}{1.525472in}}{\pgfqpoint{1.678573in}{1.527785in}}{\pgfqpoint{1.682691in}{1.531904in}}%
\pgfpathcurveto{\pgfqpoint{1.686809in}{1.536022in}}{\pgfqpoint{1.689123in}{1.541608in}}{\pgfqpoint{1.689123in}{1.547432in}}%
\pgfpathcurveto{\pgfqpoint{1.689123in}{1.553256in}}{\pgfqpoint{1.686809in}{1.558842in}}{\pgfqpoint{1.682691in}{1.562960in}}%
\pgfpathcurveto{\pgfqpoint{1.678573in}{1.567078in}}{\pgfqpoint{1.672987in}{1.569392in}}{\pgfqpoint{1.667163in}{1.569392in}}%
\pgfpathcurveto{\pgfqpoint{1.661339in}{1.569392in}}{\pgfqpoint{1.655753in}{1.567078in}}{\pgfqpoint{1.651635in}{1.562960in}}%
\pgfpathcurveto{\pgfqpoint{1.647517in}{1.558842in}}{\pgfqpoint{1.645203in}{1.553256in}}{\pgfqpoint{1.645203in}{1.547432in}}%
\pgfpathcurveto{\pgfqpoint{1.645203in}{1.541608in}}{\pgfqpoint{1.647517in}{1.536022in}}{\pgfqpoint{1.651635in}{1.531904in}}%
\pgfpathcurveto{\pgfqpoint{1.655753in}{1.527785in}}{\pgfqpoint{1.661339in}{1.525472in}}{\pgfqpoint{1.667163in}{1.525472in}}%
\pgfpathlineto{\pgfqpoint{1.667163in}{1.525472in}}%
\pgfpathclose%
\pgfusepath{stroke,fill}%
\end{pgfscope}%
\begin{pgfscope}%
\pgfpathrectangle{\pgfqpoint{0.100000in}{0.209042in}}{\pgfqpoint{2.906250in}{2.906250in}}%
\pgfusepath{clip}%
\pgfsetbuttcap%
\pgfsetroundjoin%
\definecolor{currentfill}{rgb}{0.537755,0.156894,0.506551}%
\pgfsetfillcolor{currentfill}%
\pgfsetfillopacity{0.800000}%
\pgfsetlinewidth{1.003750pt}%
\definecolor{currentstroke}{rgb}{0.537755,0.156894,0.506551}%
\pgfsetstrokecolor{currentstroke}%
\pgfsetstrokeopacity{0.800000}%
\pgfsetdash{}{0pt}%
\pgfpathmoveto{\pgfqpoint{1.687463in}{1.996290in}}%
\pgfpathcurveto{\pgfqpoint{1.693287in}{1.996290in}}{\pgfqpoint{1.698873in}{1.998604in}}{\pgfqpoint{1.702991in}{2.002722in}}%
\pgfpathcurveto{\pgfqpoint{1.707109in}{2.006840in}}{\pgfqpoint{1.709423in}{2.012426in}}{\pgfqpoint{1.709423in}{2.018250in}}%
\pgfpathcurveto{\pgfqpoint{1.709423in}{2.024074in}}{\pgfqpoint{1.707109in}{2.029661in}}{\pgfqpoint{1.702991in}{2.033779in}}%
\pgfpathcurveto{\pgfqpoint{1.698873in}{2.037897in}}{\pgfqpoint{1.693287in}{2.040211in}}{\pgfqpoint{1.687463in}{2.040211in}}%
\pgfpathcurveto{\pgfqpoint{1.681639in}{2.040211in}}{\pgfqpoint{1.676053in}{2.037897in}}{\pgfqpoint{1.671935in}{2.033779in}}%
\pgfpathcurveto{\pgfqpoint{1.667816in}{2.029661in}}{\pgfqpoint{1.665503in}{2.024074in}}{\pgfqpoint{1.665503in}{2.018250in}}%
\pgfpathcurveto{\pgfqpoint{1.665503in}{2.012426in}}{\pgfqpoint{1.667816in}{2.006840in}}{\pgfqpoint{1.671935in}{2.002722in}}%
\pgfpathcurveto{\pgfqpoint{1.676053in}{1.998604in}}{\pgfqpoint{1.681639in}{1.996290in}}{\pgfqpoint{1.687463in}{1.996290in}}%
\pgfpathlineto{\pgfqpoint{1.687463in}{1.996290in}}%
\pgfpathclose%
\pgfusepath{stroke,fill}%
\end{pgfscope}%
\begin{pgfscope}%
\pgfpathrectangle{\pgfqpoint{0.100000in}{0.209042in}}{\pgfqpoint{2.906250in}{2.906250in}}%
\pgfusepath{clip}%
\pgfsetbuttcap%
\pgfsetroundjoin%
\definecolor{currentfill}{rgb}{0.500438,0.143719,0.507920}%
\pgfsetfillcolor{currentfill}%
\pgfsetfillopacity{0.800000}%
\pgfsetlinewidth{1.003750pt}%
\definecolor{currentstroke}{rgb}{0.500438,0.143719,0.507920}%
\pgfsetstrokecolor{currentstroke}%
\pgfsetstrokeopacity{0.800000}%
\pgfsetdash{}{0pt}%
\pgfpathmoveto{\pgfqpoint{1.748926in}{1.381754in}}%
\pgfpathcurveto{\pgfqpoint{1.754750in}{1.381754in}}{\pgfqpoint{1.760336in}{1.384068in}}{\pgfqpoint{1.764455in}{1.388186in}}%
\pgfpathcurveto{\pgfqpoint{1.768573in}{1.392304in}}{\pgfqpoint{1.770887in}{1.397890in}}{\pgfqpoint{1.770887in}{1.403714in}}%
\pgfpathcurveto{\pgfqpoint{1.770887in}{1.409538in}}{\pgfqpoint{1.768573in}{1.415124in}}{\pgfqpoint{1.764455in}{1.419242in}}%
\pgfpathcurveto{\pgfqpoint{1.760336in}{1.423360in}}{\pgfqpoint{1.754750in}{1.425674in}}{\pgfqpoint{1.748926in}{1.425674in}}%
\pgfpathcurveto{\pgfqpoint{1.743102in}{1.425674in}}{\pgfqpoint{1.737516in}{1.423360in}}{\pgfqpoint{1.733398in}{1.419242in}}%
\pgfpathcurveto{\pgfqpoint{1.729280in}{1.415124in}}{\pgfqpoint{1.726966in}{1.409538in}}{\pgfqpoint{1.726966in}{1.403714in}}%
\pgfpathcurveto{\pgfqpoint{1.726966in}{1.397890in}}{\pgfqpoint{1.729280in}{1.392304in}}{\pgfqpoint{1.733398in}{1.388186in}}%
\pgfpathcurveto{\pgfqpoint{1.737516in}{1.384068in}}{\pgfqpoint{1.743102in}{1.381754in}}{\pgfqpoint{1.748926in}{1.381754in}}%
\pgfpathlineto{\pgfqpoint{1.748926in}{1.381754in}}%
\pgfpathclose%
\pgfusepath{stroke,fill}%
\end{pgfscope}%
\begin{pgfscope}%
\pgfpathrectangle{\pgfqpoint{0.100000in}{0.209042in}}{\pgfqpoint{2.906250in}{2.906250in}}%
\pgfusepath{clip}%
\pgfsetbuttcap%
\pgfsetroundjoin%
\definecolor{currentfill}{rgb}{0.297740,0.066117,0.478243}%
\pgfsetfillcolor{currentfill}%
\pgfsetfillopacity{0.800000}%
\pgfsetlinewidth{1.003750pt}%
\definecolor{currentstroke}{rgb}{0.297740,0.066117,0.478243}%
\pgfsetstrokecolor{currentstroke}%
\pgfsetstrokeopacity{0.800000}%
\pgfsetdash{}{0pt}%
\pgfpathmoveto{\pgfqpoint{1.779824in}{1.555611in}}%
\pgfpathcurveto{\pgfqpoint{1.785648in}{1.555611in}}{\pgfqpoint{1.791234in}{1.557925in}}{\pgfqpoint{1.795352in}{1.562043in}}%
\pgfpathcurveto{\pgfqpoint{1.799470in}{1.566161in}}{\pgfqpoint{1.801784in}{1.571747in}}{\pgfqpoint{1.801784in}{1.577571in}}%
\pgfpathcurveto{\pgfqpoint{1.801784in}{1.583395in}}{\pgfqpoint{1.799470in}{1.588982in}}{\pgfqpoint{1.795352in}{1.593100in}}%
\pgfpathcurveto{\pgfqpoint{1.791234in}{1.597218in}}{\pgfqpoint{1.785648in}{1.599532in}}{\pgfqpoint{1.779824in}{1.599532in}}%
\pgfpathcurveto{\pgfqpoint{1.774000in}{1.599532in}}{\pgfqpoint{1.768414in}{1.597218in}}{\pgfqpoint{1.764296in}{1.593100in}}%
\pgfpathcurveto{\pgfqpoint{1.760178in}{1.588982in}}{\pgfqpoint{1.757864in}{1.583395in}}{\pgfqpoint{1.757864in}{1.577571in}}%
\pgfpathcurveto{\pgfqpoint{1.757864in}{1.571747in}}{\pgfqpoint{1.760178in}{1.566161in}}{\pgfqpoint{1.764296in}{1.562043in}}%
\pgfpathcurveto{\pgfqpoint{1.768414in}{1.557925in}}{\pgfqpoint{1.774000in}{1.555611in}}{\pgfqpoint{1.779824in}{1.555611in}}%
\pgfpathlineto{\pgfqpoint{1.779824in}{1.555611in}}%
\pgfpathclose%
\pgfusepath{stroke,fill}%
\end{pgfscope}%
\begin{pgfscope}%
\pgfpathrectangle{\pgfqpoint{0.100000in}{0.209042in}}{\pgfqpoint{2.906250in}{2.906250in}}%
\pgfusepath{clip}%
\pgfsetbuttcap%
\pgfsetroundjoin%
\definecolor{currentfill}{rgb}{0.252220,0.059415,0.453248}%
\pgfsetfillcolor{currentfill}%
\pgfsetfillopacity{0.800000}%
\pgfsetlinewidth{1.003750pt}%
\definecolor{currentstroke}{rgb}{0.252220,0.059415,0.453248}%
\pgfsetstrokecolor{currentstroke}%
\pgfsetstrokeopacity{0.800000}%
\pgfsetdash{}{0pt}%
\pgfpathmoveto{\pgfqpoint{1.684195in}{1.780802in}}%
\pgfpathcurveto{\pgfqpoint{1.690019in}{1.780802in}}{\pgfqpoint{1.695606in}{1.783116in}}{\pgfqpoint{1.699724in}{1.787234in}}%
\pgfpathcurveto{\pgfqpoint{1.703842in}{1.791352in}}{\pgfqpoint{1.706156in}{1.796938in}}{\pgfqpoint{1.706156in}{1.802762in}}%
\pgfpathcurveto{\pgfqpoint{1.706156in}{1.808586in}}{\pgfqpoint{1.703842in}{1.814172in}}{\pgfqpoint{1.699724in}{1.818290in}}%
\pgfpathcurveto{\pgfqpoint{1.695606in}{1.822409in}}{\pgfqpoint{1.690019in}{1.824722in}}{\pgfqpoint{1.684195in}{1.824722in}}%
\pgfpathcurveto{\pgfqpoint{1.678372in}{1.824722in}}{\pgfqpoint{1.672785in}{1.822409in}}{\pgfqpoint{1.668667in}{1.818290in}}%
\pgfpathcurveto{\pgfqpoint{1.664549in}{1.814172in}}{\pgfqpoint{1.662235in}{1.808586in}}{\pgfqpoint{1.662235in}{1.802762in}}%
\pgfpathcurveto{\pgfqpoint{1.662235in}{1.796938in}}{\pgfqpoint{1.664549in}{1.791352in}}{\pgfqpoint{1.668667in}{1.787234in}}%
\pgfpathcurveto{\pgfqpoint{1.672785in}{1.783116in}}{\pgfqpoint{1.678372in}{1.780802in}}{\pgfqpoint{1.684195in}{1.780802in}}%
\pgfpathlineto{\pgfqpoint{1.684195in}{1.780802in}}%
\pgfpathclose%
\pgfusepath{stroke,fill}%
\end{pgfscope}%
\begin{pgfscope}%
\pgfpathrectangle{\pgfqpoint{0.100000in}{0.209042in}}{\pgfqpoint{2.906250in}{2.906250in}}%
\pgfusepath{clip}%
\pgfsetbuttcap%
\pgfsetroundjoin%
\definecolor{currentfill}{rgb}{0.463508,0.129893,0.507652}%
\pgfsetfillcolor{currentfill}%
\pgfsetfillopacity{0.800000}%
\pgfsetlinewidth{1.003750pt}%
\definecolor{currentstroke}{rgb}{0.463508,0.129893,0.507652}%
\pgfsetstrokecolor{currentstroke}%
\pgfsetstrokeopacity{0.800000}%
\pgfsetdash{}{0pt}%
\pgfpathmoveto{\pgfqpoint{1.471310in}{1.913201in}}%
\pgfpathcurveto{\pgfqpoint{1.477134in}{1.913201in}}{\pgfqpoint{1.482720in}{1.915515in}}{\pgfqpoint{1.486838in}{1.919633in}}%
\pgfpathcurveto{\pgfqpoint{1.490956in}{1.923751in}}{\pgfqpoint{1.493270in}{1.929338in}}{\pgfqpoint{1.493270in}{1.935162in}}%
\pgfpathcurveto{\pgfqpoint{1.493270in}{1.940986in}}{\pgfqpoint{1.490956in}{1.946572in}}{\pgfqpoint{1.486838in}{1.950690in}}%
\pgfpathcurveto{\pgfqpoint{1.482720in}{1.954808in}}{\pgfqpoint{1.477134in}{1.957122in}}{\pgfqpoint{1.471310in}{1.957122in}}%
\pgfpathcurveto{\pgfqpoint{1.465486in}{1.957122in}}{\pgfqpoint{1.459900in}{1.954808in}}{\pgfqpoint{1.455782in}{1.950690in}}%
\pgfpathcurveto{\pgfqpoint{1.451664in}{1.946572in}}{\pgfqpoint{1.449350in}{1.940986in}}{\pgfqpoint{1.449350in}{1.935162in}}%
\pgfpathcurveto{\pgfqpoint{1.449350in}{1.929338in}}{\pgfqpoint{1.451664in}{1.923751in}}{\pgfqpoint{1.455782in}{1.919633in}}%
\pgfpathcurveto{\pgfqpoint{1.459900in}{1.915515in}}{\pgfqpoint{1.465486in}{1.913201in}}{\pgfqpoint{1.471310in}{1.913201in}}%
\pgfpathlineto{\pgfqpoint{1.471310in}{1.913201in}}%
\pgfpathclose%
\pgfusepath{stroke,fill}%
\end{pgfscope}%
\begin{pgfscope}%
\pgfpathrectangle{\pgfqpoint{0.100000in}{0.209042in}}{\pgfqpoint{2.906250in}{2.906250in}}%
\pgfusepath{clip}%
\pgfsetbuttcap%
\pgfsetroundjoin%
\definecolor{currentfill}{rgb}{0.983485,0.513280,0.374198}%
\pgfsetfillcolor{currentfill}%
\pgfsetfillopacity{0.800000}%
\pgfsetlinewidth{1.003750pt}%
\definecolor{currentstroke}{rgb}{0.983485,0.513280,0.374198}%
\pgfsetstrokecolor{currentstroke}%
\pgfsetstrokeopacity{0.800000}%
\pgfsetdash{}{0pt}%
\pgfpathmoveto{\pgfqpoint{1.845836in}{1.040563in}}%
\pgfpathcurveto{\pgfqpoint{1.851660in}{1.040563in}}{\pgfqpoint{1.857246in}{1.042877in}}{\pgfqpoint{1.861364in}{1.046995in}}%
\pgfpathcurveto{\pgfqpoint{1.865482in}{1.051113in}}{\pgfqpoint{1.867796in}{1.056700in}}{\pgfqpoint{1.867796in}{1.062524in}}%
\pgfpathcurveto{\pgfqpoint{1.867796in}{1.068347in}}{\pgfqpoint{1.865482in}{1.073934in}}{\pgfqpoint{1.861364in}{1.078052in}}%
\pgfpathcurveto{\pgfqpoint{1.857246in}{1.082170in}}{\pgfqpoint{1.851660in}{1.084484in}}{\pgfqpoint{1.845836in}{1.084484in}}%
\pgfpathcurveto{\pgfqpoint{1.840012in}{1.084484in}}{\pgfqpoint{1.834426in}{1.082170in}}{\pgfqpoint{1.830308in}{1.078052in}}%
\pgfpathcurveto{\pgfqpoint{1.826190in}{1.073934in}}{\pgfqpoint{1.823876in}{1.068347in}}{\pgfqpoint{1.823876in}{1.062524in}}%
\pgfpathcurveto{\pgfqpoint{1.823876in}{1.056700in}}{\pgfqpoint{1.826190in}{1.051113in}}{\pgfqpoint{1.830308in}{1.046995in}}%
\pgfpathcurveto{\pgfqpoint{1.834426in}{1.042877in}}{\pgfqpoint{1.840012in}{1.040563in}}{\pgfqpoint{1.845836in}{1.040563in}}%
\pgfpathlineto{\pgfqpoint{1.845836in}{1.040563in}}%
\pgfpathclose%
\pgfusepath{stroke,fill}%
\end{pgfscope}%
\begin{pgfscope}%
\pgfpathrectangle{\pgfqpoint{0.100000in}{0.209042in}}{\pgfqpoint{2.906250in}{2.906250in}}%
\pgfusepath{clip}%
\pgfsetbuttcap%
\pgfsetroundjoin%
\definecolor{currentfill}{rgb}{0.709962,0.212797,0.477201}%
\pgfsetfillcolor{currentfill}%
\pgfsetfillopacity{0.800000}%
\pgfsetlinewidth{1.003750pt}%
\definecolor{currentstroke}{rgb}{0.709962,0.212797,0.477201}%
\pgfsetstrokecolor{currentstroke}%
\pgfsetstrokeopacity{0.800000}%
\pgfsetdash{}{0pt}%
\pgfpathmoveto{\pgfqpoint{2.160739in}{1.543961in}}%
\pgfpathcurveto{\pgfqpoint{2.166563in}{1.543961in}}{\pgfqpoint{2.172149in}{1.546275in}}{\pgfqpoint{2.176268in}{1.550393in}}%
\pgfpathcurveto{\pgfqpoint{2.180386in}{1.554512in}}{\pgfqpoint{2.182700in}{1.560098in}}{\pgfqpoint{2.182700in}{1.565922in}}%
\pgfpathcurveto{\pgfqpoint{2.182700in}{1.571746in}}{\pgfqpoint{2.180386in}{1.577332in}}{\pgfqpoint{2.176268in}{1.581450in}}%
\pgfpathcurveto{\pgfqpoint{2.172149in}{1.585568in}}{\pgfqpoint{2.166563in}{1.587882in}}{\pgfqpoint{2.160739in}{1.587882in}}%
\pgfpathcurveto{\pgfqpoint{2.154915in}{1.587882in}}{\pgfqpoint{2.149329in}{1.585568in}}{\pgfqpoint{2.145211in}{1.581450in}}%
\pgfpathcurveto{\pgfqpoint{2.141093in}{1.577332in}}{\pgfqpoint{2.138779in}{1.571746in}}{\pgfqpoint{2.138779in}{1.565922in}}%
\pgfpathcurveto{\pgfqpoint{2.138779in}{1.560098in}}{\pgfqpoint{2.141093in}{1.554512in}}{\pgfqpoint{2.145211in}{1.550393in}}%
\pgfpathcurveto{\pgfqpoint{2.149329in}{1.546275in}}{\pgfqpoint{2.154915in}{1.543961in}}{\pgfqpoint{2.160739in}{1.543961in}}%
\pgfpathlineto{\pgfqpoint{2.160739in}{1.543961in}}%
\pgfpathclose%
\pgfusepath{stroke,fill}%
\end{pgfscope}%
\begin{pgfscope}%
\pgfpathrectangle{\pgfqpoint{0.100000in}{0.209042in}}{\pgfqpoint{2.906250in}{2.906250in}}%
\pgfusepath{clip}%
\pgfsetbuttcap%
\pgfsetroundjoin%
\definecolor{currentfill}{rgb}{0.632805,0.187893,0.495332}%
\pgfsetfillcolor{currentfill}%
\pgfsetfillopacity{0.800000}%
\pgfsetlinewidth{1.003750pt}%
\definecolor{currentstroke}{rgb}{0.632805,0.187893,0.495332}%
\pgfsetstrokecolor{currentstroke}%
\pgfsetstrokeopacity{0.800000}%
\pgfsetdash{}{0pt}%
\pgfpathmoveto{\pgfqpoint{1.551883in}{2.046190in}}%
\pgfpathcurveto{\pgfqpoint{1.557707in}{2.046190in}}{\pgfqpoint{1.563294in}{2.048504in}}{\pgfqpoint{1.567412in}{2.052622in}}%
\pgfpathcurveto{\pgfqpoint{1.571530in}{2.056740in}}{\pgfqpoint{1.573844in}{2.062326in}}{\pgfqpoint{1.573844in}{2.068150in}}%
\pgfpathcurveto{\pgfqpoint{1.573844in}{2.073974in}}{\pgfqpoint{1.571530in}{2.079560in}}{\pgfqpoint{1.567412in}{2.083679in}}%
\pgfpathcurveto{\pgfqpoint{1.563294in}{2.087797in}}{\pgfqpoint{1.557707in}{2.090111in}}{\pgfqpoint{1.551883in}{2.090111in}}%
\pgfpathcurveto{\pgfqpoint{1.546060in}{2.090111in}}{\pgfqpoint{1.540473in}{2.087797in}}{\pgfqpoint{1.536355in}{2.083679in}}%
\pgfpathcurveto{\pgfqpoint{1.532237in}{2.079560in}}{\pgfqpoint{1.529923in}{2.073974in}}{\pgfqpoint{1.529923in}{2.068150in}}%
\pgfpathcurveto{\pgfqpoint{1.529923in}{2.062326in}}{\pgfqpoint{1.532237in}{2.056740in}}{\pgfqpoint{1.536355in}{2.052622in}}%
\pgfpathcurveto{\pgfqpoint{1.540473in}{2.048504in}}{\pgfqpoint{1.546060in}{2.046190in}}{\pgfqpoint{1.551883in}{2.046190in}}%
\pgfpathlineto{\pgfqpoint{1.551883in}{2.046190in}}%
\pgfpathclose%
\pgfusepath{stroke,fill}%
\end{pgfscope}%
\begin{pgfscope}%
\pgfpathrectangle{\pgfqpoint{0.100000in}{0.209042in}}{\pgfqpoint{2.906250in}{2.906250in}}%
\pgfusepath{clip}%
\pgfsetbuttcap%
\pgfsetroundjoin%
\definecolor{currentfill}{rgb}{0.297740,0.066117,0.478243}%
\pgfsetfillcolor{currentfill}%
\pgfsetfillopacity{0.800000}%
\pgfsetlinewidth{1.003750pt}%
\definecolor{currentstroke}{rgb}{0.297740,0.066117,0.478243}%
\pgfsetstrokecolor{currentstroke}%
\pgfsetstrokeopacity{0.800000}%
\pgfsetdash{}{0pt}%
\pgfpathmoveto{\pgfqpoint{1.755429in}{1.789886in}}%
\pgfpathcurveto{\pgfqpoint{1.761253in}{1.789886in}}{\pgfqpoint{1.766839in}{1.792200in}}{\pgfqpoint{1.770958in}{1.796318in}}%
\pgfpathcurveto{\pgfqpoint{1.775076in}{1.800436in}}{\pgfqpoint{1.777390in}{1.806023in}}{\pgfqpoint{1.777390in}{1.811847in}}%
\pgfpathcurveto{\pgfqpoint{1.777390in}{1.817671in}}{\pgfqpoint{1.775076in}{1.823257in}}{\pgfqpoint{1.770958in}{1.827375in}}%
\pgfpathcurveto{\pgfqpoint{1.766839in}{1.831493in}}{\pgfqpoint{1.761253in}{1.833807in}}{\pgfqpoint{1.755429in}{1.833807in}}%
\pgfpathcurveto{\pgfqpoint{1.749605in}{1.833807in}}{\pgfqpoint{1.744019in}{1.831493in}}{\pgfqpoint{1.739901in}{1.827375in}}%
\pgfpathcurveto{\pgfqpoint{1.735783in}{1.823257in}}{\pgfqpoint{1.733469in}{1.817671in}}{\pgfqpoint{1.733469in}{1.811847in}}%
\pgfpathcurveto{\pgfqpoint{1.733469in}{1.806023in}}{\pgfqpoint{1.735783in}{1.800436in}}{\pgfqpoint{1.739901in}{1.796318in}}%
\pgfpathcurveto{\pgfqpoint{1.744019in}{1.792200in}}{\pgfqpoint{1.749605in}{1.789886in}}{\pgfqpoint{1.755429in}{1.789886in}}%
\pgfpathlineto{\pgfqpoint{1.755429in}{1.789886in}}%
\pgfpathclose%
\pgfusepath{stroke,fill}%
\end{pgfscope}%
\begin{pgfscope}%
\pgfpathrectangle{\pgfqpoint{0.100000in}{0.209042in}}{\pgfqpoint{2.906250in}{2.906250in}}%
\pgfusepath{clip}%
\pgfsetbuttcap%
\pgfsetroundjoin%
\definecolor{currentfill}{rgb}{0.519045,0.150383,0.507443}%
\pgfsetfillcolor{currentfill}%
\pgfsetfillopacity{0.800000}%
\pgfsetlinewidth{1.003750pt}%
\definecolor{currentstroke}{rgb}{0.519045,0.150383,0.507443}%
\pgfsetstrokecolor{currentstroke}%
\pgfsetstrokeopacity{0.800000}%
\pgfsetdash{}{0pt}%
\pgfpathmoveto{\pgfqpoint{1.981875in}{1.819627in}}%
\pgfpathcurveto{\pgfqpoint{1.987699in}{1.819627in}}{\pgfqpoint{1.993285in}{1.821941in}}{\pgfqpoint{1.997404in}{1.826059in}}%
\pgfpathcurveto{\pgfqpoint{2.001522in}{1.830177in}}{\pgfqpoint{2.003836in}{1.835764in}}{\pgfqpoint{2.003836in}{1.841588in}}%
\pgfpathcurveto{\pgfqpoint{2.003836in}{1.847411in}}{\pgfqpoint{2.001522in}{1.852998in}}{\pgfqpoint{1.997404in}{1.857116in}}%
\pgfpathcurveto{\pgfqpoint{1.993285in}{1.861234in}}{\pgfqpoint{1.987699in}{1.863548in}}{\pgfqpoint{1.981875in}{1.863548in}}%
\pgfpathcurveto{\pgfqpoint{1.976051in}{1.863548in}}{\pgfqpoint{1.970465in}{1.861234in}}{\pgfqpoint{1.966347in}{1.857116in}}%
\pgfpathcurveto{\pgfqpoint{1.962229in}{1.852998in}}{\pgfqpoint{1.959915in}{1.847411in}}{\pgfqpoint{1.959915in}{1.841588in}}%
\pgfpathcurveto{\pgfqpoint{1.959915in}{1.835764in}}{\pgfqpoint{1.962229in}{1.830177in}}{\pgfqpoint{1.966347in}{1.826059in}}%
\pgfpathcurveto{\pgfqpoint{1.970465in}{1.821941in}}{\pgfqpoint{1.976051in}{1.819627in}}{\pgfqpoint{1.981875in}{1.819627in}}%
\pgfpathlineto{\pgfqpoint{1.981875in}{1.819627in}}%
\pgfpathclose%
\pgfusepath{stroke,fill}%
\end{pgfscope}%
\begin{pgfscope}%
\pgfpathrectangle{\pgfqpoint{0.100000in}{0.209042in}}{\pgfqpoint{2.906250in}{2.906250in}}%
\pgfusepath{clip}%
\pgfsetbuttcap%
\pgfsetroundjoin%
\definecolor{currentfill}{rgb}{0.218512,0.061158,0.425392}%
\pgfsetfillcolor{currentfill}%
\pgfsetfillopacity{0.800000}%
\pgfsetlinewidth{1.003750pt}%
\definecolor{currentstroke}{rgb}{0.218512,0.061158,0.425392}%
\pgfsetstrokecolor{currentstroke}%
\pgfsetstrokeopacity{0.800000}%
\pgfsetdash{}{0pt}%
\pgfpathmoveto{\pgfqpoint{1.636665in}{1.639775in}}%
\pgfpathcurveto{\pgfqpoint{1.642489in}{1.639775in}}{\pgfqpoint{1.648075in}{1.642089in}}{\pgfqpoint{1.652193in}{1.646207in}}%
\pgfpathcurveto{\pgfqpoint{1.656311in}{1.650325in}}{\pgfqpoint{1.658625in}{1.655911in}}{\pgfqpoint{1.658625in}{1.661735in}}%
\pgfpathcurveto{\pgfqpoint{1.658625in}{1.667559in}}{\pgfqpoint{1.656311in}{1.673145in}}{\pgfqpoint{1.652193in}{1.677263in}}%
\pgfpathcurveto{\pgfqpoint{1.648075in}{1.681382in}}{\pgfqpoint{1.642489in}{1.683695in}}{\pgfqpoint{1.636665in}{1.683695in}}%
\pgfpathcurveto{\pgfqpoint{1.630841in}{1.683695in}}{\pgfqpoint{1.625255in}{1.681382in}}{\pgfqpoint{1.621137in}{1.677263in}}%
\pgfpathcurveto{\pgfqpoint{1.617018in}{1.673145in}}{\pgfqpoint{1.614705in}{1.667559in}}{\pgfqpoint{1.614705in}{1.661735in}}%
\pgfpathcurveto{\pgfqpoint{1.614705in}{1.655911in}}{\pgfqpoint{1.617018in}{1.650325in}}{\pgfqpoint{1.621137in}{1.646207in}}%
\pgfpathcurveto{\pgfqpoint{1.625255in}{1.642089in}}{\pgfqpoint{1.630841in}{1.639775in}}{\pgfqpoint{1.636665in}{1.639775in}}%
\pgfpathlineto{\pgfqpoint{1.636665in}{1.639775in}}%
\pgfpathclose%
\pgfusepath{stroke,fill}%
\end{pgfscope}%
\begin{pgfscope}%
\pgfpathrectangle{\pgfqpoint{0.100000in}{0.209042in}}{\pgfqpoint{2.906250in}{2.906250in}}%
\pgfusepath{clip}%
\pgfsetbuttcap%
\pgfsetroundjoin%
\definecolor{currentfill}{rgb}{0.512831,0.148179,0.507648}%
\pgfsetfillcolor{currentfill}%
\pgfsetfillopacity{0.800000}%
\pgfsetlinewidth{1.003750pt}%
\definecolor{currentstroke}{rgb}{0.512831,0.148179,0.507648}%
\pgfsetstrokecolor{currentstroke}%
\pgfsetstrokeopacity{0.800000}%
\pgfsetdash{}{0pt}%
\pgfpathmoveto{\pgfqpoint{1.861359in}{1.412524in}}%
\pgfpathcurveto{\pgfqpoint{1.867183in}{1.412524in}}{\pgfqpoint{1.872769in}{1.414838in}}{\pgfqpoint{1.876887in}{1.418956in}}%
\pgfpathcurveto{\pgfqpoint{1.881005in}{1.423074in}}{\pgfqpoint{1.883319in}{1.428660in}}{\pgfqpoint{1.883319in}{1.434484in}}%
\pgfpathcurveto{\pgfqpoint{1.883319in}{1.440308in}}{\pgfqpoint{1.881005in}{1.445894in}}{\pgfqpoint{1.876887in}{1.450012in}}%
\pgfpathcurveto{\pgfqpoint{1.872769in}{1.454131in}}{\pgfqpoint{1.867183in}{1.456444in}}{\pgfqpoint{1.861359in}{1.456444in}}%
\pgfpathcurveto{\pgfqpoint{1.855535in}{1.456444in}}{\pgfqpoint{1.849949in}{1.454131in}}{\pgfqpoint{1.845831in}{1.450012in}}%
\pgfpathcurveto{\pgfqpoint{1.841713in}{1.445894in}}{\pgfqpoint{1.839399in}{1.440308in}}{\pgfqpoint{1.839399in}{1.434484in}}%
\pgfpathcurveto{\pgfqpoint{1.839399in}{1.428660in}}{\pgfqpoint{1.841713in}{1.423074in}}{\pgfqpoint{1.845831in}{1.418956in}}%
\pgfpathcurveto{\pgfqpoint{1.849949in}{1.414838in}}{\pgfqpoint{1.855535in}{1.412524in}}{\pgfqpoint{1.861359in}{1.412524in}}%
\pgfpathlineto{\pgfqpoint{1.861359in}{1.412524in}}%
\pgfpathclose%
\pgfusepath{stroke,fill}%
\end{pgfscope}%
\begin{pgfscope}%
\pgfpathrectangle{\pgfqpoint{0.100000in}{0.209042in}}{\pgfqpoint{2.906250in}{2.906250in}}%
\pgfusepath{clip}%
\pgfsetbuttcap%
\pgfsetroundjoin%
\definecolor{currentfill}{rgb}{0.868793,0.287728,0.409303}%
\pgfsetfillcolor{currentfill}%
\pgfsetfillopacity{0.800000}%
\pgfsetlinewidth{1.003750pt}%
\definecolor{currentstroke}{rgb}{0.868793,0.287728,0.409303}%
\pgfsetstrokecolor{currentstroke}%
\pgfsetstrokeopacity{0.800000}%
\pgfsetdash{}{0pt}%
\pgfpathmoveto{\pgfqpoint{2.177054in}{1.943435in}}%
\pgfpathcurveto{\pgfqpoint{2.182878in}{1.943435in}}{\pgfqpoint{2.188464in}{1.945749in}}{\pgfqpoint{2.192582in}{1.949867in}}%
\pgfpathcurveto{\pgfqpoint{2.196700in}{1.953986in}}{\pgfqpoint{2.199014in}{1.959572in}}{\pgfqpoint{2.199014in}{1.965396in}}%
\pgfpathcurveto{\pgfqpoint{2.199014in}{1.971220in}}{\pgfqpoint{2.196700in}{1.976806in}}{\pgfqpoint{2.192582in}{1.980924in}}%
\pgfpathcurveto{\pgfqpoint{2.188464in}{1.985042in}}{\pgfqpoint{2.182878in}{1.987356in}}{\pgfqpoint{2.177054in}{1.987356in}}%
\pgfpathcurveto{\pgfqpoint{2.171230in}{1.987356in}}{\pgfqpoint{2.165644in}{1.985042in}}{\pgfqpoint{2.161526in}{1.980924in}}%
\pgfpathcurveto{\pgfqpoint{2.157407in}{1.976806in}}{\pgfqpoint{2.155093in}{1.971220in}}{\pgfqpoint{2.155093in}{1.965396in}}%
\pgfpathcurveto{\pgfqpoint{2.155093in}{1.959572in}}{\pgfqpoint{2.157407in}{1.953986in}}{\pgfqpoint{2.161526in}{1.949867in}}%
\pgfpathcurveto{\pgfqpoint{2.165644in}{1.945749in}}{\pgfqpoint{2.171230in}{1.943435in}}{\pgfqpoint{2.177054in}{1.943435in}}%
\pgfpathlineto{\pgfqpoint{2.177054in}{1.943435in}}%
\pgfpathclose%
\pgfusepath{stroke,fill}%
\end{pgfscope}%
\begin{pgfscope}%
\pgfpathrectangle{\pgfqpoint{0.100000in}{0.209042in}}{\pgfqpoint{2.906250in}{2.906250in}}%
\pgfusepath{clip}%
\pgfsetbuttcap%
\pgfsetroundjoin%
\definecolor{currentfill}{rgb}{0.304081,0.067835,0.480812}%
\pgfsetfillcolor{currentfill}%
\pgfsetfillopacity{0.800000}%
\pgfsetlinewidth{1.003750pt}%
\definecolor{currentstroke}{rgb}{0.304081,0.067835,0.480812}%
\pgfsetstrokecolor{currentstroke}%
\pgfsetstrokeopacity{0.800000}%
\pgfsetdash{}{0pt}%
\pgfpathmoveto{\pgfqpoint{1.702795in}{1.537538in}}%
\pgfpathcurveto{\pgfqpoint{1.708619in}{1.537538in}}{\pgfqpoint{1.714205in}{1.539852in}}{\pgfqpoint{1.718323in}{1.543970in}}%
\pgfpathcurveto{\pgfqpoint{1.722442in}{1.548088in}}{\pgfqpoint{1.724755in}{1.553674in}}{\pgfqpoint{1.724755in}{1.559498in}}%
\pgfpathcurveto{\pgfqpoint{1.724755in}{1.565322in}}{\pgfqpoint{1.722442in}{1.570908in}}{\pgfqpoint{1.718323in}{1.575027in}}%
\pgfpathcurveto{\pgfqpoint{1.714205in}{1.579145in}}{\pgfqpoint{1.708619in}{1.581459in}}{\pgfqpoint{1.702795in}{1.581459in}}%
\pgfpathcurveto{\pgfqpoint{1.696971in}{1.581459in}}{\pgfqpoint{1.691385in}{1.579145in}}{\pgfqpoint{1.687267in}{1.575027in}}%
\pgfpathcurveto{\pgfqpoint{1.683149in}{1.570908in}}{\pgfqpoint{1.680835in}{1.565322in}}{\pgfqpoint{1.680835in}{1.559498in}}%
\pgfpathcurveto{\pgfqpoint{1.680835in}{1.553674in}}{\pgfqpoint{1.683149in}{1.548088in}}{\pgfqpoint{1.687267in}{1.543970in}}%
\pgfpathcurveto{\pgfqpoint{1.691385in}{1.539852in}}{\pgfqpoint{1.696971in}{1.537538in}}{\pgfqpoint{1.702795in}{1.537538in}}%
\pgfpathlineto{\pgfqpoint{1.702795in}{1.537538in}}%
\pgfpathclose%
\pgfusepath{stroke,fill}%
\end{pgfscope}%
\begin{pgfscope}%
\pgfpathrectangle{\pgfqpoint{0.100000in}{0.209042in}}{\pgfqpoint{2.906250in}{2.906250in}}%
\pgfusepath{clip}%
\pgfsetbuttcap%
\pgfsetroundjoin%
\definecolor{currentfill}{rgb}{0.408629,0.107930,0.504052}%
\pgfsetfillcolor{currentfill}%
\pgfsetfillopacity{0.800000}%
\pgfsetlinewidth{1.003750pt}%
\definecolor{currentstroke}{rgb}{0.408629,0.107930,0.504052}%
\pgfsetstrokecolor{currentstroke}%
\pgfsetstrokeopacity{0.800000}%
\pgfsetdash{}{0pt}%
\pgfpathmoveto{\pgfqpoint{1.477183in}{1.513621in}}%
\pgfpathcurveto{\pgfqpoint{1.483007in}{1.513621in}}{\pgfqpoint{1.488593in}{1.515935in}}{\pgfqpoint{1.492711in}{1.520053in}}%
\pgfpathcurveto{\pgfqpoint{1.496829in}{1.524171in}}{\pgfqpoint{1.499143in}{1.529758in}}{\pgfqpoint{1.499143in}{1.535582in}}%
\pgfpathcurveto{\pgfqpoint{1.499143in}{1.541405in}}{\pgfqpoint{1.496829in}{1.546992in}}{\pgfqpoint{1.492711in}{1.551110in}}%
\pgfpathcurveto{\pgfqpoint{1.488593in}{1.555228in}}{\pgfqpoint{1.483007in}{1.557542in}}{\pgfqpoint{1.477183in}{1.557542in}}%
\pgfpathcurveto{\pgfqpoint{1.471359in}{1.557542in}}{\pgfqpoint{1.465773in}{1.555228in}}{\pgfqpoint{1.461654in}{1.551110in}}%
\pgfpathcurveto{\pgfqpoint{1.457536in}{1.546992in}}{\pgfqpoint{1.455222in}{1.541405in}}{\pgfqpoint{1.455222in}{1.535582in}}%
\pgfpathcurveto{\pgfqpoint{1.455222in}{1.529758in}}{\pgfqpoint{1.457536in}{1.524171in}}{\pgfqpoint{1.461654in}{1.520053in}}%
\pgfpathcurveto{\pgfqpoint{1.465773in}{1.515935in}}{\pgfqpoint{1.471359in}{1.513621in}}{\pgfqpoint{1.477183in}{1.513621in}}%
\pgfpathlineto{\pgfqpoint{1.477183in}{1.513621in}}%
\pgfpathclose%
\pgfusepath{stroke,fill}%
\end{pgfscope}%
\begin{pgfscope}%
\pgfpathrectangle{\pgfqpoint{0.100000in}{0.209042in}}{\pgfqpoint{2.906250in}{2.906250in}}%
\pgfusepath{clip}%
\pgfsetbuttcap%
\pgfsetroundjoin%
\definecolor{currentfill}{rgb}{0.804752,0.249911,0.442102}%
\pgfsetfillcolor{currentfill}%
\pgfsetfillopacity{0.800000}%
\pgfsetlinewidth{1.003750pt}%
\definecolor{currentstroke}{rgb}{0.804752,0.249911,0.442102}%
\pgfsetstrokecolor{currentstroke}%
\pgfsetstrokeopacity{0.800000}%
\pgfsetdash{}{0pt}%
\pgfpathmoveto{\pgfqpoint{1.773574in}{2.131341in}}%
\pgfpathcurveto{\pgfqpoint{1.779398in}{2.131341in}}{\pgfqpoint{1.784984in}{2.133655in}}{\pgfqpoint{1.789103in}{2.137773in}}%
\pgfpathcurveto{\pgfqpoint{1.793221in}{2.141892in}}{\pgfqpoint{1.795535in}{2.147478in}}{\pgfqpoint{1.795535in}{2.153302in}}%
\pgfpathcurveto{\pgfqpoint{1.795535in}{2.159126in}}{\pgfqpoint{1.793221in}{2.164712in}}{\pgfqpoint{1.789103in}{2.168830in}}%
\pgfpathcurveto{\pgfqpoint{1.784984in}{2.172948in}}{\pgfqpoint{1.779398in}{2.175262in}}{\pgfqpoint{1.773574in}{2.175262in}}%
\pgfpathcurveto{\pgfqpoint{1.767750in}{2.175262in}}{\pgfqpoint{1.762164in}{2.172948in}}{\pgfqpoint{1.758046in}{2.168830in}}%
\pgfpathcurveto{\pgfqpoint{1.753928in}{2.164712in}}{\pgfqpoint{1.751614in}{2.159126in}}{\pgfqpoint{1.751614in}{2.153302in}}%
\pgfpathcurveto{\pgfqpoint{1.751614in}{2.147478in}}{\pgfqpoint{1.753928in}{2.141892in}}{\pgfqpoint{1.758046in}{2.137773in}}%
\pgfpathcurveto{\pgfqpoint{1.762164in}{2.133655in}}{\pgfqpoint{1.767750in}{2.131341in}}{\pgfqpoint{1.773574in}{2.131341in}}%
\pgfpathlineto{\pgfqpoint{1.773574in}{2.131341in}}%
\pgfpathclose%
\pgfusepath{stroke,fill}%
\end{pgfscope}%
\begin{pgfscope}%
\pgfpathrectangle{\pgfqpoint{0.100000in}{0.209042in}}{\pgfqpoint{2.906250in}{2.906250in}}%
\pgfusepath{clip}%
\pgfsetbuttcap%
\pgfsetroundjoin%
\definecolor{currentfill}{rgb}{0.792427,0.244242,0.447543}%
\pgfsetfillcolor{currentfill}%
\pgfsetfillopacity{0.800000}%
\pgfsetlinewidth{1.003750pt}%
\definecolor{currentstroke}{rgb}{0.792427,0.244242,0.447543}%
\pgfsetstrokecolor{currentstroke}%
\pgfsetstrokeopacity{0.800000}%
\pgfsetdash{}{0pt}%
\pgfpathmoveto{\pgfqpoint{1.903967in}{1.243175in}}%
\pgfpathcurveto{\pgfqpoint{1.909791in}{1.243175in}}{\pgfqpoint{1.915377in}{1.245489in}}{\pgfqpoint{1.919495in}{1.249607in}}%
\pgfpathcurveto{\pgfqpoint{1.923613in}{1.253725in}}{\pgfqpoint{1.925927in}{1.259311in}}{\pgfqpoint{1.925927in}{1.265135in}}%
\pgfpathcurveto{\pgfqpoint{1.925927in}{1.270959in}}{\pgfqpoint{1.923613in}{1.276545in}}{\pgfqpoint{1.919495in}{1.280663in}}%
\pgfpathcurveto{\pgfqpoint{1.915377in}{1.284781in}}{\pgfqpoint{1.909791in}{1.287095in}}{\pgfqpoint{1.903967in}{1.287095in}}%
\pgfpathcurveto{\pgfqpoint{1.898143in}{1.287095in}}{\pgfqpoint{1.892556in}{1.284781in}}{\pgfqpoint{1.888438in}{1.280663in}}%
\pgfpathcurveto{\pgfqpoint{1.884320in}{1.276545in}}{\pgfqpoint{1.882006in}{1.270959in}}{\pgfqpoint{1.882006in}{1.265135in}}%
\pgfpathcurveto{\pgfqpoint{1.882006in}{1.259311in}}{\pgfqpoint{1.884320in}{1.253725in}}{\pgfqpoint{1.888438in}{1.249607in}}%
\pgfpathcurveto{\pgfqpoint{1.892556in}{1.245489in}}{\pgfqpoint{1.898143in}{1.243175in}}{\pgfqpoint{1.903967in}{1.243175in}}%
\pgfpathlineto{\pgfqpoint{1.903967in}{1.243175in}}%
\pgfpathclose%
\pgfusepath{stroke,fill}%
\end{pgfscope}%
\begin{pgfscope}%
\pgfpathrectangle{\pgfqpoint{0.100000in}{0.209042in}}{\pgfqpoint{2.906250in}{2.906250in}}%
\pgfusepath{clip}%
\pgfsetbuttcap%
\pgfsetroundjoin%
\definecolor{currentfill}{rgb}{0.284951,0.063168,0.472451}%
\pgfsetfillcolor{currentfill}%
\pgfsetfillopacity{0.800000}%
\pgfsetlinewidth{1.003750pt}%
\definecolor{currentstroke}{rgb}{0.284951,0.063168,0.472451}%
\pgfsetstrokecolor{currentstroke}%
\pgfsetstrokeopacity{0.800000}%
\pgfsetdash{}{0pt}%
\pgfpathmoveto{\pgfqpoint{1.789155in}{1.637266in}}%
\pgfpathcurveto{\pgfqpoint{1.794979in}{1.637266in}}{\pgfqpoint{1.800565in}{1.639580in}}{\pgfqpoint{1.804683in}{1.643698in}}%
\pgfpathcurveto{\pgfqpoint{1.808801in}{1.647816in}}{\pgfqpoint{1.811115in}{1.653402in}}{\pgfqpoint{1.811115in}{1.659226in}}%
\pgfpathcurveto{\pgfqpoint{1.811115in}{1.665050in}}{\pgfqpoint{1.808801in}{1.670636in}}{\pgfqpoint{1.804683in}{1.674754in}}%
\pgfpathcurveto{\pgfqpoint{1.800565in}{1.678872in}}{\pgfqpoint{1.794979in}{1.681186in}}{\pgfqpoint{1.789155in}{1.681186in}}%
\pgfpathcurveto{\pgfqpoint{1.783331in}{1.681186in}}{\pgfqpoint{1.777745in}{1.678872in}}{\pgfqpoint{1.773627in}{1.674754in}}%
\pgfpathcurveto{\pgfqpoint{1.769508in}{1.670636in}}{\pgfqpoint{1.767195in}{1.665050in}}{\pgfqpoint{1.767195in}{1.659226in}}%
\pgfpathcurveto{\pgfqpoint{1.767195in}{1.653402in}}{\pgfqpoint{1.769508in}{1.647816in}}{\pgfqpoint{1.773627in}{1.643698in}}%
\pgfpathcurveto{\pgfqpoint{1.777745in}{1.639580in}}{\pgfqpoint{1.783331in}{1.637266in}}{\pgfqpoint{1.789155in}{1.637266in}}%
\pgfpathlineto{\pgfqpoint{1.789155in}{1.637266in}}%
\pgfpathclose%
\pgfusepath{stroke,fill}%
\end{pgfscope}%
\begin{pgfscope}%
\pgfpathrectangle{\pgfqpoint{0.100000in}{0.209042in}}{\pgfqpoint{2.906250in}{2.906250in}}%
\pgfusepath{clip}%
\pgfsetbuttcap%
\pgfsetroundjoin%
\definecolor{currentfill}{rgb}{0.488088,0.139186,0.508011}%
\pgfsetfillcolor{currentfill}%
\pgfsetfillopacity{0.800000}%
\pgfsetlinewidth{1.003750pt}%
\definecolor{currentstroke}{rgb}{0.488088,0.139186,0.508011}%
\pgfsetstrokecolor{currentstroke}%
\pgfsetstrokeopacity{0.800000}%
\pgfsetdash{}{0pt}%
\pgfpathmoveto{\pgfqpoint{1.995832in}{1.702075in}}%
\pgfpathcurveto{\pgfqpoint{2.001656in}{1.702075in}}{\pgfqpoint{2.007242in}{1.704389in}}{\pgfqpoint{2.011361in}{1.708507in}}%
\pgfpathcurveto{\pgfqpoint{2.015479in}{1.712625in}}{\pgfqpoint{2.017793in}{1.718212in}}{\pgfqpoint{2.017793in}{1.724036in}}%
\pgfpathcurveto{\pgfqpoint{2.017793in}{1.729859in}}{\pgfqpoint{2.015479in}{1.735446in}}{\pgfqpoint{2.011361in}{1.739564in}}%
\pgfpathcurveto{\pgfqpoint{2.007242in}{1.743682in}}{\pgfqpoint{2.001656in}{1.745996in}}{\pgfqpoint{1.995832in}{1.745996in}}%
\pgfpathcurveto{\pgfqpoint{1.990008in}{1.745996in}}{\pgfqpoint{1.984422in}{1.743682in}}{\pgfqpoint{1.980304in}{1.739564in}}%
\pgfpathcurveto{\pgfqpoint{1.976186in}{1.735446in}}{\pgfqpoint{1.973872in}{1.729859in}}{\pgfqpoint{1.973872in}{1.724036in}}%
\pgfpathcurveto{\pgfqpoint{1.973872in}{1.718212in}}{\pgfqpoint{1.976186in}{1.712625in}}{\pgfqpoint{1.980304in}{1.708507in}}%
\pgfpathcurveto{\pgfqpoint{1.984422in}{1.704389in}}{\pgfqpoint{1.990008in}{1.702075in}}{\pgfqpoint{1.995832in}{1.702075in}}%
\pgfpathlineto{\pgfqpoint{1.995832in}{1.702075in}}%
\pgfpathclose%
\pgfusepath{stroke,fill}%
\end{pgfscope}%
\begin{pgfscope}%
\pgfpathrectangle{\pgfqpoint{0.100000in}{0.209042in}}{\pgfqpoint{2.906250in}{2.906250in}}%
\pgfusepath{clip}%
\pgfsetbuttcap%
\pgfsetroundjoin%
\definecolor{currentfill}{rgb}{0.994738,0.624350,0.427397}%
\pgfsetfillcolor{currentfill}%
\pgfsetfillopacity{0.800000}%
\pgfsetlinewidth{1.003750pt}%
\definecolor{currentstroke}{rgb}{0.994738,0.624350,0.427397}%
\pgfsetstrokecolor{currentstroke}%
\pgfsetstrokeopacity{0.800000}%
\pgfsetdash{}{0pt}%
\pgfpathmoveto{\pgfqpoint{1.558733in}{0.991881in}}%
\pgfpathcurveto{\pgfqpoint{1.564557in}{0.991881in}}{\pgfqpoint{1.570143in}{0.994195in}}{\pgfqpoint{1.574261in}{0.998313in}}%
\pgfpathcurveto{\pgfqpoint{1.578379in}{1.002431in}}{\pgfqpoint{1.580693in}{1.008017in}}{\pgfqpoint{1.580693in}{1.013841in}}%
\pgfpathcurveto{\pgfqpoint{1.580693in}{1.019665in}}{\pgfqpoint{1.578379in}{1.025252in}}{\pgfqpoint{1.574261in}{1.029370in}}%
\pgfpathcurveto{\pgfqpoint{1.570143in}{1.033488in}}{\pgfqpoint{1.564557in}{1.035802in}}{\pgfqpoint{1.558733in}{1.035802in}}%
\pgfpathcurveto{\pgfqpoint{1.552909in}{1.035802in}}{\pgfqpoint{1.547323in}{1.033488in}}{\pgfqpoint{1.543205in}{1.029370in}}%
\pgfpathcurveto{\pgfqpoint{1.539086in}{1.025252in}}{\pgfqpoint{1.536773in}{1.019665in}}{\pgfqpoint{1.536773in}{1.013841in}}%
\pgfpathcurveto{\pgfqpoint{1.536773in}{1.008017in}}{\pgfqpoint{1.539086in}{1.002431in}}{\pgfqpoint{1.543205in}{0.998313in}}%
\pgfpathcurveto{\pgfqpoint{1.547323in}{0.994195in}}{\pgfqpoint{1.552909in}{0.991881in}}{\pgfqpoint{1.558733in}{0.991881in}}%
\pgfpathlineto{\pgfqpoint{1.558733in}{0.991881in}}%
\pgfpathclose%
\pgfusepath{stroke,fill}%
\end{pgfscope}%
\begin{pgfscope}%
\pgfpathrectangle{\pgfqpoint{0.100000in}{0.209042in}}{\pgfqpoint{2.906250in}{2.906250in}}%
\pgfusepath{clip}%
\pgfsetbuttcap%
\pgfsetroundjoin%
\definecolor{currentfill}{rgb}{0.420791,0.112920,0.505215}%
\pgfsetfillcolor{currentfill}%
\pgfsetfillopacity{0.800000}%
\pgfsetlinewidth{1.003750pt}%
\definecolor{currentstroke}{rgb}{0.420791,0.112920,0.505215}%
\pgfsetstrokecolor{currentstroke}%
\pgfsetstrokeopacity{0.800000}%
\pgfsetdash{}{0pt}%
\pgfpathmoveto{\pgfqpoint{1.470646in}{1.518087in}}%
\pgfpathcurveto{\pgfqpoint{1.476470in}{1.518087in}}{\pgfqpoint{1.482056in}{1.520401in}}{\pgfqpoint{1.486174in}{1.524519in}}%
\pgfpathcurveto{\pgfqpoint{1.490292in}{1.528637in}}{\pgfqpoint{1.492606in}{1.534223in}}{\pgfqpoint{1.492606in}{1.540047in}}%
\pgfpathcurveto{\pgfqpoint{1.492606in}{1.545871in}}{\pgfqpoint{1.490292in}{1.551457in}}{\pgfqpoint{1.486174in}{1.555575in}}%
\pgfpathcurveto{\pgfqpoint{1.482056in}{1.559694in}}{\pgfqpoint{1.476470in}{1.562007in}}{\pgfqpoint{1.470646in}{1.562007in}}%
\pgfpathcurveto{\pgfqpoint{1.464822in}{1.562007in}}{\pgfqpoint{1.459236in}{1.559694in}}{\pgfqpoint{1.455118in}{1.555575in}}%
\pgfpathcurveto{\pgfqpoint{1.450999in}{1.551457in}}{\pgfqpoint{1.448685in}{1.545871in}}{\pgfqpoint{1.448685in}{1.540047in}}%
\pgfpathcurveto{\pgfqpoint{1.448685in}{1.534223in}}{\pgfqpoint{1.450999in}{1.528637in}}{\pgfqpoint{1.455118in}{1.524519in}}%
\pgfpathcurveto{\pgfqpoint{1.459236in}{1.520401in}}{\pgfqpoint{1.464822in}{1.518087in}}{\pgfqpoint{1.470646in}{1.518087in}}%
\pgfpathlineto{\pgfqpoint{1.470646in}{1.518087in}}%
\pgfpathclose%
\pgfusepath{stroke,fill}%
\end{pgfscope}%
\begin{pgfscope}%
\pgfpathrectangle{\pgfqpoint{0.100000in}{0.209042in}}{\pgfqpoint{2.906250in}{2.906250in}}%
\pgfusepath{clip}%
\pgfsetbuttcap%
\pgfsetroundjoin%
\definecolor{currentfill}{rgb}{0.402548,0.105420,0.503386}%
\pgfsetfillcolor{currentfill}%
\pgfsetfillopacity{0.800000}%
\pgfsetlinewidth{1.003750pt}%
\definecolor{currentstroke}{rgb}{0.402548,0.105420,0.503386}%
\pgfsetstrokecolor{currentstroke}%
\pgfsetstrokeopacity{0.800000}%
\pgfsetdash{}{0pt}%
\pgfpathmoveto{\pgfqpoint{1.449989in}{1.796118in}}%
\pgfpathcurveto{\pgfqpoint{1.455813in}{1.796118in}}{\pgfqpoint{1.461399in}{1.798432in}}{\pgfqpoint{1.465517in}{1.802550in}}%
\pgfpathcurveto{\pgfqpoint{1.469635in}{1.806668in}}{\pgfqpoint{1.471949in}{1.812254in}}{\pgfqpoint{1.471949in}{1.818078in}}%
\pgfpathcurveto{\pgfqpoint{1.471949in}{1.823902in}}{\pgfqpoint{1.469635in}{1.829488in}}{\pgfqpoint{1.465517in}{1.833606in}}%
\pgfpathcurveto{\pgfqpoint{1.461399in}{1.837724in}}{\pgfqpoint{1.455813in}{1.840038in}}{\pgfqpoint{1.449989in}{1.840038in}}%
\pgfpathcurveto{\pgfqpoint{1.444165in}{1.840038in}}{\pgfqpoint{1.438579in}{1.837724in}}{\pgfqpoint{1.434461in}{1.833606in}}%
\pgfpathcurveto{\pgfqpoint{1.430343in}{1.829488in}}{\pgfqpoint{1.428029in}{1.823902in}}{\pgfqpoint{1.428029in}{1.818078in}}%
\pgfpathcurveto{\pgfqpoint{1.428029in}{1.812254in}}{\pgfqpoint{1.430343in}{1.806668in}}{\pgfqpoint{1.434461in}{1.802550in}}%
\pgfpathcurveto{\pgfqpoint{1.438579in}{1.798432in}}{\pgfqpoint{1.444165in}{1.796118in}}{\pgfqpoint{1.449989in}{1.796118in}}%
\pgfpathlineto{\pgfqpoint{1.449989in}{1.796118in}}%
\pgfpathclose%
\pgfusepath{stroke,fill}%
\end{pgfscope}%
\begin{pgfscope}%
\pgfpathrectangle{\pgfqpoint{0.100000in}{0.209042in}}{\pgfqpoint{2.906250in}{2.906250in}}%
\pgfusepath{clip}%
\pgfsetbuttcap%
\pgfsetroundjoin%
\definecolor{currentfill}{rgb}{0.341482,0.080564,0.492631}%
\pgfsetfillcolor{currentfill}%
\pgfsetfillopacity{0.800000}%
\pgfsetlinewidth{1.003750pt}%
\definecolor{currentstroke}{rgb}{0.341482,0.080564,0.492631}%
\pgfsetstrokecolor{currentstroke}%
\pgfsetstrokeopacity{0.800000}%
\pgfsetdash{}{0pt}%
\pgfpathmoveto{\pgfqpoint{1.712818in}{1.524617in}}%
\pgfpathcurveto{\pgfqpoint{1.718642in}{1.524617in}}{\pgfqpoint{1.724228in}{1.526931in}}{\pgfqpoint{1.728346in}{1.531049in}}%
\pgfpathcurveto{\pgfqpoint{1.732465in}{1.535167in}}{\pgfqpoint{1.734778in}{1.540753in}}{\pgfqpoint{1.734778in}{1.546577in}}%
\pgfpathcurveto{\pgfqpoint{1.734778in}{1.552401in}}{\pgfqpoint{1.732465in}{1.557988in}}{\pgfqpoint{1.728346in}{1.562106in}}%
\pgfpathcurveto{\pgfqpoint{1.724228in}{1.566224in}}{\pgfqpoint{1.718642in}{1.568538in}}{\pgfqpoint{1.712818in}{1.568538in}}%
\pgfpathcurveto{\pgfqpoint{1.706994in}{1.568538in}}{\pgfqpoint{1.701408in}{1.566224in}}{\pgfqpoint{1.697290in}{1.562106in}}%
\pgfpathcurveto{\pgfqpoint{1.693172in}{1.557988in}}{\pgfqpoint{1.690858in}{1.552401in}}{\pgfqpoint{1.690858in}{1.546577in}}%
\pgfpathcurveto{\pgfqpoint{1.690858in}{1.540753in}}{\pgfqpoint{1.693172in}{1.535167in}}{\pgfqpoint{1.697290in}{1.531049in}}%
\pgfpathcurveto{\pgfqpoint{1.701408in}{1.526931in}}{\pgfqpoint{1.706994in}{1.524617in}}{\pgfqpoint{1.712818in}{1.524617in}}%
\pgfpathlineto{\pgfqpoint{1.712818in}{1.524617in}}%
\pgfpathclose%
\pgfusepath{stroke,fill}%
\end{pgfscope}%
\begin{pgfscope}%
\pgfpathrectangle{\pgfqpoint{0.100000in}{0.209042in}}{\pgfqpoint{2.906250in}{2.906250in}}%
\pgfusepath{clip}%
\pgfsetbuttcap%
\pgfsetroundjoin%
\definecolor{currentfill}{rgb}{0.722805,0.217194,0.473316}%
\pgfsetfillcolor{currentfill}%
\pgfsetfillopacity{0.800000}%
\pgfsetlinewidth{1.003750pt}%
\definecolor{currentstroke}{rgb}{0.722805,0.217194,0.473316}%
\pgfsetstrokecolor{currentstroke}%
\pgfsetstrokeopacity{0.800000}%
\pgfsetdash{}{0pt}%
\pgfpathmoveto{\pgfqpoint{1.286824in}{1.977053in}}%
\pgfpathcurveto{\pgfqpoint{1.292648in}{1.977053in}}{\pgfqpoint{1.298234in}{1.979367in}}{\pgfqpoint{1.302353in}{1.983485in}}%
\pgfpathcurveto{\pgfqpoint{1.306471in}{1.987604in}}{\pgfqpoint{1.308785in}{1.993190in}}{\pgfqpoint{1.308785in}{1.999014in}}%
\pgfpathcurveto{\pgfqpoint{1.308785in}{2.004838in}}{\pgfqpoint{1.306471in}{2.010424in}}{\pgfqpoint{1.302353in}{2.014542in}}%
\pgfpathcurveto{\pgfqpoint{1.298234in}{2.018660in}}{\pgfqpoint{1.292648in}{2.020974in}}{\pgfqpoint{1.286824in}{2.020974in}}%
\pgfpathcurveto{\pgfqpoint{1.281000in}{2.020974in}}{\pgfqpoint{1.275414in}{2.018660in}}{\pgfqpoint{1.271296in}{2.014542in}}%
\pgfpathcurveto{\pgfqpoint{1.267178in}{2.010424in}}{\pgfqpoint{1.264864in}{2.004838in}}{\pgfqpoint{1.264864in}{1.999014in}}%
\pgfpathcurveto{\pgfqpoint{1.264864in}{1.993190in}}{\pgfqpoint{1.267178in}{1.987604in}}{\pgfqpoint{1.271296in}{1.983485in}}%
\pgfpathcurveto{\pgfqpoint{1.275414in}{1.979367in}}{\pgfqpoint{1.281000in}{1.977053in}}{\pgfqpoint{1.286824in}{1.977053in}}%
\pgfpathlineto{\pgfqpoint{1.286824in}{1.977053in}}%
\pgfpathclose%
\pgfusepath{stroke,fill}%
\end{pgfscope}%
\begin{pgfscope}%
\pgfpathrectangle{\pgfqpoint{0.100000in}{0.209042in}}{\pgfqpoint{2.906250in}{2.906250in}}%
\pgfusepath{clip}%
\pgfsetbuttcap%
\pgfsetroundjoin%
\definecolor{currentfill}{rgb}{0.664915,0.198075,0.488836}%
\pgfsetfillcolor{currentfill}%
\pgfsetfillopacity{0.800000}%
\pgfsetlinewidth{1.003750pt}%
\definecolor{currentstroke}{rgb}{0.664915,0.198075,0.488836}%
\pgfsetstrokecolor{currentstroke}%
\pgfsetstrokeopacity{0.800000}%
\pgfsetdash{}{0pt}%
\pgfpathmoveto{\pgfqpoint{1.190781in}{1.651499in}}%
\pgfpathcurveto{\pgfqpoint{1.196604in}{1.651499in}}{\pgfqpoint{1.202191in}{1.653813in}}{\pgfqpoint{1.206309in}{1.657931in}}%
\pgfpathcurveto{\pgfqpoint{1.210427in}{1.662049in}}{\pgfqpoint{1.212741in}{1.667636in}}{\pgfqpoint{1.212741in}{1.673459in}}%
\pgfpathcurveto{\pgfqpoint{1.212741in}{1.679283in}}{\pgfqpoint{1.210427in}{1.684870in}}{\pgfqpoint{1.206309in}{1.688988in}}%
\pgfpathcurveto{\pgfqpoint{1.202191in}{1.693106in}}{\pgfqpoint{1.196604in}{1.695420in}}{\pgfqpoint{1.190781in}{1.695420in}}%
\pgfpathcurveto{\pgfqpoint{1.184957in}{1.695420in}}{\pgfqpoint{1.179370in}{1.693106in}}{\pgfqpoint{1.175252in}{1.688988in}}%
\pgfpathcurveto{\pgfqpoint{1.171134in}{1.684870in}}{\pgfqpoint{1.168820in}{1.679283in}}{\pgfqpoint{1.168820in}{1.673459in}}%
\pgfpathcurveto{\pgfqpoint{1.168820in}{1.667636in}}{\pgfqpoint{1.171134in}{1.662049in}}{\pgfqpoint{1.175252in}{1.657931in}}%
\pgfpathcurveto{\pgfqpoint{1.179370in}{1.653813in}}{\pgfqpoint{1.184957in}{1.651499in}}{\pgfqpoint{1.190781in}{1.651499in}}%
\pgfpathlineto{\pgfqpoint{1.190781in}{1.651499in}}%
\pgfpathclose%
\pgfusepath{stroke,fill}%
\end{pgfscope}%
\begin{pgfscope}%
\pgfpathrectangle{\pgfqpoint{0.100000in}{0.209042in}}{\pgfqpoint{2.906250in}{2.906250in}}%
\pgfusepath{clip}%
\pgfsetbuttcap%
\pgfsetroundjoin%
\definecolor{currentfill}{rgb}{0.697098,0.208501,0.480835}%
\pgfsetfillcolor{currentfill}%
\pgfsetfillopacity{0.800000}%
\pgfsetlinewidth{1.003750pt}%
\definecolor{currentstroke}{rgb}{0.697098,0.208501,0.480835}%
\pgfsetstrokecolor{currentstroke}%
\pgfsetstrokeopacity{0.800000}%
\pgfsetdash{}{0pt}%
\pgfpathmoveto{\pgfqpoint{1.193742in}{1.575914in}}%
\pgfpathcurveto{\pgfqpoint{1.199566in}{1.575914in}}{\pgfqpoint{1.205152in}{1.578228in}}{\pgfqpoint{1.209270in}{1.582346in}}%
\pgfpathcurveto{\pgfqpoint{1.213389in}{1.586464in}}{\pgfqpoint{1.215702in}{1.592051in}}{\pgfqpoint{1.215702in}{1.597874in}}%
\pgfpathcurveto{\pgfqpoint{1.215702in}{1.603698in}}{\pgfqpoint{1.213389in}{1.609285in}}{\pgfqpoint{1.209270in}{1.613403in}}%
\pgfpathcurveto{\pgfqpoint{1.205152in}{1.617521in}}{\pgfqpoint{1.199566in}{1.619835in}}{\pgfqpoint{1.193742in}{1.619835in}}%
\pgfpathcurveto{\pgfqpoint{1.187918in}{1.619835in}}{\pgfqpoint{1.182332in}{1.617521in}}{\pgfqpoint{1.178214in}{1.613403in}}%
\pgfpathcurveto{\pgfqpoint{1.174096in}{1.609285in}}{\pgfqpoint{1.171782in}{1.603698in}}{\pgfqpoint{1.171782in}{1.597874in}}%
\pgfpathcurveto{\pgfqpoint{1.171782in}{1.592051in}}{\pgfqpoint{1.174096in}{1.586464in}}{\pgfqpoint{1.178214in}{1.582346in}}%
\pgfpathcurveto{\pgfqpoint{1.182332in}{1.578228in}}{\pgfqpoint{1.187918in}{1.575914in}}{\pgfqpoint{1.193742in}{1.575914in}}%
\pgfpathlineto{\pgfqpoint{1.193742in}{1.575914in}}%
\pgfpathclose%
\pgfusepath{stroke,fill}%
\end{pgfscope}%
\begin{pgfscope}%
\pgfpathrectangle{\pgfqpoint{0.100000in}{0.209042in}}{\pgfqpoint{2.906250in}{2.906250in}}%
\pgfusepath{clip}%
\pgfsetbuttcap%
\pgfsetroundjoin%
\definecolor{currentfill}{rgb}{0.310382,0.069702,0.483186}%
\pgfsetfillcolor{currentfill}%
\pgfsetfillopacity{0.800000}%
\pgfsetlinewidth{1.003750pt}%
\definecolor{currentstroke}{rgb}{0.310382,0.069702,0.483186}%
\pgfsetstrokecolor{currentstroke}%
\pgfsetstrokeopacity{0.800000}%
\pgfsetdash{}{0pt}%
\pgfpathmoveto{\pgfqpoint{1.551479in}{1.682484in}}%
\pgfpathcurveto{\pgfqpoint{1.557303in}{1.682484in}}{\pgfqpoint{1.562889in}{1.684798in}}{\pgfqpoint{1.567007in}{1.688916in}}%
\pgfpathcurveto{\pgfqpoint{1.571125in}{1.693034in}}{\pgfqpoint{1.573439in}{1.698620in}}{\pgfqpoint{1.573439in}{1.704444in}}%
\pgfpathcurveto{\pgfqpoint{1.573439in}{1.710268in}}{\pgfqpoint{1.571125in}{1.715854in}}{\pgfqpoint{1.567007in}{1.719972in}}%
\pgfpathcurveto{\pgfqpoint{1.562889in}{1.724090in}}{\pgfqpoint{1.557303in}{1.726404in}}{\pgfqpoint{1.551479in}{1.726404in}}%
\pgfpathcurveto{\pgfqpoint{1.545655in}{1.726404in}}{\pgfqpoint{1.540068in}{1.724090in}}{\pgfqpoint{1.535950in}{1.719972in}}%
\pgfpathcurveto{\pgfqpoint{1.531832in}{1.715854in}}{\pgfqpoint{1.529518in}{1.710268in}}{\pgfqpoint{1.529518in}{1.704444in}}%
\pgfpathcurveto{\pgfqpoint{1.529518in}{1.698620in}}{\pgfqpoint{1.531832in}{1.693034in}}{\pgfqpoint{1.535950in}{1.688916in}}%
\pgfpathcurveto{\pgfqpoint{1.540068in}{1.684798in}}{\pgfqpoint{1.545655in}{1.682484in}}{\pgfqpoint{1.551479in}{1.682484in}}%
\pgfpathlineto{\pgfqpoint{1.551479in}{1.682484in}}%
\pgfpathclose%
\pgfusepath{stroke,fill}%
\end{pgfscope}%
\begin{pgfscope}%
\pgfpathrectangle{\pgfqpoint{0.100000in}{0.209042in}}{\pgfqpoint{2.906250in}{2.906250in}}%
\pgfusepath{clip}%
\pgfsetbuttcap%
\pgfsetroundjoin%
\definecolor{currentfill}{rgb}{0.329114,0.075972,0.489287}%
\pgfsetfillcolor{currentfill}%
\pgfsetfillopacity{0.800000}%
\pgfsetlinewidth{1.003750pt}%
\definecolor{currentstroke}{rgb}{0.329114,0.075972,0.489287}%
\pgfsetstrokecolor{currentstroke}%
\pgfsetstrokeopacity{0.800000}%
\pgfsetdash{}{0pt}%
\pgfpathmoveto{\pgfqpoint{1.557312in}{1.745727in}}%
\pgfpathcurveto{\pgfqpoint{1.563136in}{1.745727in}}{\pgfqpoint{1.568722in}{1.748041in}}{\pgfqpoint{1.572840in}{1.752159in}}%
\pgfpathcurveto{\pgfqpoint{1.576959in}{1.756277in}}{\pgfqpoint{1.579272in}{1.761863in}}{\pgfqpoint{1.579272in}{1.767687in}}%
\pgfpathcurveto{\pgfqpoint{1.579272in}{1.773511in}}{\pgfqpoint{1.576959in}{1.779097in}}{\pgfqpoint{1.572840in}{1.783216in}}%
\pgfpathcurveto{\pgfqpoint{1.568722in}{1.787334in}}{\pgfqpoint{1.563136in}{1.789648in}}{\pgfqpoint{1.557312in}{1.789648in}}%
\pgfpathcurveto{\pgfqpoint{1.551488in}{1.789648in}}{\pgfqpoint{1.545902in}{1.787334in}}{\pgfqpoint{1.541784in}{1.783216in}}%
\pgfpathcurveto{\pgfqpoint{1.537666in}{1.779097in}}{\pgfqpoint{1.535352in}{1.773511in}}{\pgfqpoint{1.535352in}{1.767687in}}%
\pgfpathcurveto{\pgfqpoint{1.535352in}{1.761863in}}{\pgfqpoint{1.537666in}{1.756277in}}{\pgfqpoint{1.541784in}{1.752159in}}%
\pgfpathcurveto{\pgfqpoint{1.545902in}{1.748041in}}{\pgfqpoint{1.551488in}{1.745727in}}{\pgfqpoint{1.557312in}{1.745727in}}%
\pgfpathlineto{\pgfqpoint{1.557312in}{1.745727in}}%
\pgfpathclose%
\pgfusepath{stroke,fill}%
\end{pgfscope}%
\begin{pgfscope}%
\pgfpathrectangle{\pgfqpoint{0.100000in}{0.209042in}}{\pgfqpoint{2.906250in}{2.906250in}}%
\pgfusepath{clip}%
\pgfsetbuttcap%
\pgfsetroundjoin%
\definecolor{currentfill}{rgb}{0.322899,0.073782,0.487408}%
\pgfsetfillcolor{currentfill}%
\pgfsetfillopacity{0.800000}%
\pgfsetlinewidth{1.003750pt}%
\definecolor{currentstroke}{rgb}{0.322899,0.073782,0.487408}%
\pgfsetstrokecolor{currentstroke}%
\pgfsetstrokeopacity{0.800000}%
\pgfsetdash{}{0pt}%
\pgfpathmoveto{\pgfqpoint{1.772995in}{1.709976in}}%
\pgfpathcurveto{\pgfqpoint{1.778819in}{1.709976in}}{\pgfqpoint{1.784405in}{1.712289in}}{\pgfqpoint{1.788523in}{1.716408in}}%
\pgfpathcurveto{\pgfqpoint{1.792641in}{1.720526in}}{\pgfqpoint{1.794955in}{1.726112in}}{\pgfqpoint{1.794955in}{1.731936in}}%
\pgfpathcurveto{\pgfqpoint{1.794955in}{1.737760in}}{\pgfqpoint{1.792641in}{1.743346in}}{\pgfqpoint{1.788523in}{1.747464in}}%
\pgfpathcurveto{\pgfqpoint{1.784405in}{1.751582in}}{\pgfqpoint{1.778819in}{1.753896in}}{\pgfqpoint{1.772995in}{1.753896in}}%
\pgfpathcurveto{\pgfqpoint{1.767171in}{1.753896in}}{\pgfqpoint{1.761585in}{1.751582in}}{\pgfqpoint{1.757467in}{1.747464in}}%
\pgfpathcurveto{\pgfqpoint{1.753348in}{1.743346in}}{\pgfqpoint{1.751035in}{1.737760in}}{\pgfqpoint{1.751035in}{1.731936in}}%
\pgfpathcurveto{\pgfqpoint{1.751035in}{1.726112in}}{\pgfqpoint{1.753348in}{1.720526in}}{\pgfqpoint{1.757467in}{1.716408in}}%
\pgfpathcurveto{\pgfqpoint{1.761585in}{1.712289in}}{\pgfqpoint{1.767171in}{1.709976in}}{\pgfqpoint{1.772995in}{1.709976in}}%
\pgfpathlineto{\pgfqpoint{1.772995in}{1.709976in}}%
\pgfpathclose%
\pgfusepath{stroke,fill}%
\end{pgfscope}%
\begin{pgfscope}%
\pgfpathrectangle{\pgfqpoint{0.100000in}{0.209042in}}{\pgfqpoint{2.906250in}{2.906250in}}%
\pgfusepath{clip}%
\pgfsetbuttcap%
\pgfsetroundjoin%
\definecolor{currentfill}{rgb}{0.506629,0.145958,0.507806}%
\pgfsetfillcolor{currentfill}%
\pgfsetfillopacity{0.800000}%
\pgfsetlinewidth{1.003750pt}%
\definecolor{currentstroke}{rgb}{0.506629,0.145958,0.507806}%
\pgfsetstrokecolor{currentstroke}%
\pgfsetstrokeopacity{0.800000}%
\pgfsetdash{}{0pt}%
\pgfpathmoveto{\pgfqpoint{1.759967in}{1.921366in}}%
\pgfpathcurveto{\pgfqpoint{1.765791in}{1.921366in}}{\pgfqpoint{1.771377in}{1.923680in}}{\pgfqpoint{1.775495in}{1.927798in}}%
\pgfpathcurveto{\pgfqpoint{1.779613in}{1.931916in}}{\pgfqpoint{1.781927in}{1.937502in}}{\pgfqpoint{1.781927in}{1.943326in}}%
\pgfpathcurveto{\pgfqpoint{1.781927in}{1.949150in}}{\pgfqpoint{1.779613in}{1.954736in}}{\pgfqpoint{1.775495in}{1.958854in}}%
\pgfpathcurveto{\pgfqpoint{1.771377in}{1.962973in}}{\pgfqpoint{1.765791in}{1.965287in}}{\pgfqpoint{1.759967in}{1.965287in}}%
\pgfpathcurveto{\pgfqpoint{1.754143in}{1.965287in}}{\pgfqpoint{1.748557in}{1.962973in}}{\pgfqpoint{1.744439in}{1.958854in}}%
\pgfpathcurveto{\pgfqpoint{1.740320in}{1.954736in}}{\pgfqpoint{1.738006in}{1.949150in}}{\pgfqpoint{1.738006in}{1.943326in}}%
\pgfpathcurveto{\pgfqpoint{1.738006in}{1.937502in}}{\pgfqpoint{1.740320in}{1.931916in}}{\pgfqpoint{1.744439in}{1.927798in}}%
\pgfpathcurveto{\pgfqpoint{1.748557in}{1.923680in}}{\pgfqpoint{1.754143in}{1.921366in}}{\pgfqpoint{1.759967in}{1.921366in}}%
\pgfpathlineto{\pgfqpoint{1.759967in}{1.921366in}}%
\pgfpathclose%
\pgfusepath{stroke,fill}%
\end{pgfscope}%
\begin{pgfscope}%
\pgfpathrectangle{\pgfqpoint{0.100000in}{0.209042in}}{\pgfqpoint{2.906250in}{2.906250in}}%
\pgfusepath{clip}%
\pgfsetbuttcap%
\pgfsetroundjoin%
\definecolor{currentfill}{rgb}{0.684224,0.204286,0.484219}%
\pgfsetfillcolor{currentfill}%
\pgfsetfillopacity{0.800000}%
\pgfsetlinewidth{1.003750pt}%
\definecolor{currentstroke}{rgb}{0.684224,0.204286,0.484219}%
\pgfsetstrokecolor{currentstroke}%
\pgfsetstrokeopacity{0.800000}%
\pgfsetdash{}{0pt}%
\pgfpathmoveto{\pgfqpoint{1.777714in}{2.038545in}}%
\pgfpathcurveto{\pgfqpoint{1.783538in}{2.038545in}}{\pgfqpoint{1.789124in}{2.040858in}}{\pgfqpoint{1.793242in}{2.044977in}}%
\pgfpathcurveto{\pgfqpoint{1.797360in}{2.049095in}}{\pgfqpoint{1.799674in}{2.054681in}}{\pgfqpoint{1.799674in}{2.060505in}}%
\pgfpathcurveto{\pgfqpoint{1.799674in}{2.066329in}}{\pgfqpoint{1.797360in}{2.071915in}}{\pgfqpoint{1.793242in}{2.076033in}}%
\pgfpathcurveto{\pgfqpoint{1.789124in}{2.080151in}}{\pgfqpoint{1.783538in}{2.082465in}}{\pgfqpoint{1.777714in}{2.082465in}}%
\pgfpathcurveto{\pgfqpoint{1.771890in}{2.082465in}}{\pgfqpoint{1.766304in}{2.080151in}}{\pgfqpoint{1.762185in}{2.076033in}}%
\pgfpathcurveto{\pgfqpoint{1.758067in}{2.071915in}}{\pgfqpoint{1.755753in}{2.066329in}}{\pgfqpoint{1.755753in}{2.060505in}}%
\pgfpathcurveto{\pgfqpoint{1.755753in}{2.054681in}}{\pgfqpoint{1.758067in}{2.049095in}}{\pgfqpoint{1.762185in}{2.044977in}}%
\pgfpathcurveto{\pgfqpoint{1.766304in}{2.040858in}}{\pgfqpoint{1.771890in}{2.038545in}}{\pgfqpoint{1.777714in}{2.038545in}}%
\pgfpathlineto{\pgfqpoint{1.777714in}{2.038545in}}%
\pgfpathclose%
\pgfusepath{stroke,fill}%
\end{pgfscope}%
\begin{pgfscope}%
\pgfpathrectangle{\pgfqpoint{0.100000in}{0.209042in}}{\pgfqpoint{2.906250in}{2.906250in}}%
\pgfusepath{clip}%
\pgfsetbuttcap%
\pgfsetroundjoin%
\definecolor{currentfill}{rgb}{0.432967,0.117855,0.506160}%
\pgfsetfillcolor{currentfill}%
\pgfsetfillopacity{0.800000}%
\pgfsetlinewidth{1.003750pt}%
\definecolor{currentstroke}{rgb}{0.432967,0.117855,0.506160}%
\pgfsetstrokecolor{currentstroke}%
\pgfsetstrokeopacity{0.800000}%
\pgfsetdash{}{0pt}%
\pgfpathmoveto{\pgfqpoint{1.684637in}{1.459704in}}%
\pgfpathcurveto{\pgfqpoint{1.690461in}{1.459704in}}{\pgfqpoint{1.696047in}{1.462018in}}{\pgfqpoint{1.700166in}{1.466136in}}%
\pgfpathcurveto{\pgfqpoint{1.704284in}{1.470254in}}{\pgfqpoint{1.706598in}{1.475840in}}{\pgfqpoint{1.706598in}{1.481664in}}%
\pgfpathcurveto{\pgfqpoint{1.706598in}{1.487488in}}{\pgfqpoint{1.704284in}{1.493075in}}{\pgfqpoint{1.700166in}{1.497193in}}%
\pgfpathcurveto{\pgfqpoint{1.696047in}{1.501311in}}{\pgfqpoint{1.690461in}{1.503625in}}{\pgfqpoint{1.684637in}{1.503625in}}%
\pgfpathcurveto{\pgfqpoint{1.678813in}{1.503625in}}{\pgfqpoint{1.673227in}{1.501311in}}{\pgfqpoint{1.669109in}{1.497193in}}%
\pgfpathcurveto{\pgfqpoint{1.664991in}{1.493075in}}{\pgfqpoint{1.662677in}{1.487488in}}{\pgfqpoint{1.662677in}{1.481664in}}%
\pgfpathcurveto{\pgfqpoint{1.662677in}{1.475840in}}{\pgfqpoint{1.664991in}{1.470254in}}{\pgfqpoint{1.669109in}{1.466136in}}%
\pgfpathcurveto{\pgfqpoint{1.673227in}{1.462018in}}{\pgfqpoint{1.678813in}{1.459704in}}{\pgfqpoint{1.684637in}{1.459704in}}%
\pgfpathlineto{\pgfqpoint{1.684637in}{1.459704in}}%
\pgfpathclose%
\pgfusepath{stroke,fill}%
\end{pgfscope}%
\begin{pgfscope}%
\pgfpathrectangle{\pgfqpoint{0.100000in}{0.209042in}}{\pgfqpoint{2.906250in}{2.906250in}}%
\pgfusepath{clip}%
\pgfsetbuttcap%
\pgfsetroundjoin%
\definecolor{currentfill}{rgb}{0.481929,0.136891,0.507989}%
\pgfsetfillcolor{currentfill}%
\pgfsetfillopacity{0.800000}%
\pgfsetlinewidth{1.003750pt}%
\definecolor{currentstroke}{rgb}{0.481929,0.136891,0.507989}%
\pgfsetstrokecolor{currentstroke}%
\pgfsetstrokeopacity{0.800000}%
\pgfsetdash{}{0pt}%
\pgfpathmoveto{\pgfqpoint{1.819996in}{1.872712in}}%
\pgfpathcurveto{\pgfqpoint{1.825820in}{1.872712in}}{\pgfqpoint{1.831406in}{1.875026in}}{\pgfqpoint{1.835524in}{1.879144in}}%
\pgfpathcurveto{\pgfqpoint{1.839642in}{1.883262in}}{\pgfqpoint{1.841956in}{1.888848in}}{\pgfqpoint{1.841956in}{1.894672in}}%
\pgfpathcurveto{\pgfqpoint{1.841956in}{1.900496in}}{\pgfqpoint{1.839642in}{1.906082in}}{\pgfqpoint{1.835524in}{1.910200in}}%
\pgfpathcurveto{\pgfqpoint{1.831406in}{1.914318in}}{\pgfqpoint{1.825820in}{1.916632in}}{\pgfqpoint{1.819996in}{1.916632in}}%
\pgfpathcurveto{\pgfqpoint{1.814172in}{1.916632in}}{\pgfqpoint{1.808586in}{1.914318in}}{\pgfqpoint{1.804468in}{1.910200in}}%
\pgfpathcurveto{\pgfqpoint{1.800350in}{1.906082in}}{\pgfqpoint{1.798036in}{1.900496in}}{\pgfqpoint{1.798036in}{1.894672in}}%
\pgfpathcurveto{\pgfqpoint{1.798036in}{1.888848in}}{\pgfqpoint{1.800350in}{1.883262in}}{\pgfqpoint{1.804468in}{1.879144in}}%
\pgfpathcurveto{\pgfqpoint{1.808586in}{1.875026in}}{\pgfqpoint{1.814172in}{1.872712in}}{\pgfqpoint{1.819996in}{1.872712in}}%
\pgfpathlineto{\pgfqpoint{1.819996in}{1.872712in}}%
\pgfpathclose%
\pgfusepath{stroke,fill}%
\end{pgfscope}%
\begin{pgfscope}%
\pgfpathrectangle{\pgfqpoint{0.100000in}{0.209042in}}{\pgfqpoint{2.906250in}{2.906250in}}%
\pgfusepath{clip}%
\pgfsetbuttcap%
\pgfsetroundjoin%
\definecolor{currentfill}{rgb}{0.879464,0.296125,0.403118}%
\pgfsetfillcolor{currentfill}%
\pgfsetfillopacity{0.800000}%
\pgfsetlinewidth{1.003750pt}%
\definecolor{currentstroke}{rgb}{0.879464,0.296125,0.403118}%
\pgfsetstrokecolor{currentstroke}%
\pgfsetstrokeopacity{0.800000}%
\pgfsetdash{}{0pt}%
\pgfpathmoveto{\pgfqpoint{1.777208in}{1.166657in}}%
\pgfpathcurveto{\pgfqpoint{1.783032in}{1.166657in}}{\pgfqpoint{1.788619in}{1.168970in}}{\pgfqpoint{1.792737in}{1.173089in}}%
\pgfpathcurveto{\pgfqpoint{1.796855in}{1.177207in}}{\pgfqpoint{1.799169in}{1.182793in}}{\pgfqpoint{1.799169in}{1.188617in}}%
\pgfpathcurveto{\pgfqpoint{1.799169in}{1.194441in}}{\pgfqpoint{1.796855in}{1.200027in}}{\pgfqpoint{1.792737in}{1.204145in}}%
\pgfpathcurveto{\pgfqpoint{1.788619in}{1.208263in}}{\pgfqpoint{1.783032in}{1.210577in}}{\pgfqpoint{1.777208in}{1.210577in}}%
\pgfpathcurveto{\pgfqpoint{1.771385in}{1.210577in}}{\pgfqpoint{1.765798in}{1.208263in}}{\pgfqpoint{1.761680in}{1.204145in}}%
\pgfpathcurveto{\pgfqpoint{1.757562in}{1.200027in}}{\pgfqpoint{1.755248in}{1.194441in}}{\pgfqpoint{1.755248in}{1.188617in}}%
\pgfpathcurveto{\pgfqpoint{1.755248in}{1.182793in}}{\pgfqpoint{1.757562in}{1.177207in}}{\pgfqpoint{1.761680in}{1.173089in}}%
\pgfpathcurveto{\pgfqpoint{1.765798in}{1.168970in}}{\pgfqpoint{1.771385in}{1.166657in}}{\pgfqpoint{1.777208in}{1.166657in}}%
\pgfpathlineto{\pgfqpoint{1.777208in}{1.166657in}}%
\pgfpathclose%
\pgfusepath{stroke,fill}%
\end{pgfscope}%
\begin{pgfscope}%
\pgfpathrectangle{\pgfqpoint{0.100000in}{0.209042in}}{\pgfqpoint{2.906250in}{2.906250in}}%
\pgfusepath{clip}%
\pgfsetbuttcap%
\pgfsetroundjoin%
\definecolor{currentfill}{rgb}{0.341482,0.080564,0.492631}%
\pgfsetfillcolor{currentfill}%
\pgfsetfillopacity{0.800000}%
\pgfsetlinewidth{1.003750pt}%
\definecolor{currentstroke}{rgb}{0.341482,0.080564,0.492631}%
\pgfsetstrokecolor{currentstroke}%
\pgfsetstrokeopacity{0.800000}%
\pgfsetdash{}{0pt}%
\pgfpathmoveto{\pgfqpoint{1.529167in}{1.643202in}}%
\pgfpathcurveto{\pgfqpoint{1.534990in}{1.643202in}}{\pgfqpoint{1.540577in}{1.645516in}}{\pgfqpoint{1.544695in}{1.649634in}}%
\pgfpathcurveto{\pgfqpoint{1.548813in}{1.653752in}}{\pgfqpoint{1.551127in}{1.659338in}}{\pgfqpoint{1.551127in}{1.665162in}}%
\pgfpathcurveto{\pgfqpoint{1.551127in}{1.670986in}}{\pgfqpoint{1.548813in}{1.676572in}}{\pgfqpoint{1.544695in}{1.680690in}}%
\pgfpathcurveto{\pgfqpoint{1.540577in}{1.684808in}}{\pgfqpoint{1.534990in}{1.687122in}}{\pgfqpoint{1.529167in}{1.687122in}}%
\pgfpathcurveto{\pgfqpoint{1.523343in}{1.687122in}}{\pgfqpoint{1.517756in}{1.684808in}}{\pgfqpoint{1.513638in}{1.680690in}}%
\pgfpathcurveto{\pgfqpoint{1.509520in}{1.676572in}}{\pgfqpoint{1.507206in}{1.670986in}}{\pgfqpoint{1.507206in}{1.665162in}}%
\pgfpathcurveto{\pgfqpoint{1.507206in}{1.659338in}}{\pgfqpoint{1.509520in}{1.653752in}}{\pgfqpoint{1.513638in}{1.649634in}}%
\pgfpathcurveto{\pgfqpoint{1.517756in}{1.645516in}}{\pgfqpoint{1.523343in}{1.643202in}}{\pgfqpoint{1.529167in}{1.643202in}}%
\pgfpathlineto{\pgfqpoint{1.529167in}{1.643202in}}%
\pgfpathclose%
\pgfusepath{stroke,fill}%
\end{pgfscope}%
\begin{pgfscope}%
\pgfpathrectangle{\pgfqpoint{0.100000in}{0.209042in}}{\pgfqpoint{2.906250in}{2.906250in}}%
\pgfusepath{clip}%
\pgfsetbuttcap%
\pgfsetroundjoin%
\definecolor{currentfill}{rgb}{0.664915,0.198075,0.488836}%
\pgfsetfillcolor{currentfill}%
\pgfsetfillopacity{0.800000}%
\pgfsetlinewidth{1.003750pt}%
\definecolor{currentstroke}{rgb}{0.664915,0.198075,0.488836}%
\pgfsetstrokecolor{currentstroke}%
\pgfsetstrokeopacity{0.800000}%
\pgfsetdash{}{0pt}%
\pgfpathmoveto{\pgfqpoint{2.116596in}{1.678529in}}%
\pgfpathcurveto{\pgfqpoint{2.122419in}{1.678529in}}{\pgfqpoint{2.128006in}{1.680843in}}{\pgfqpoint{2.132124in}{1.684961in}}%
\pgfpathcurveto{\pgfqpoint{2.136242in}{1.689079in}}{\pgfqpoint{2.138556in}{1.694665in}}{\pgfqpoint{2.138556in}{1.700489in}}%
\pgfpathcurveto{\pgfqpoint{2.138556in}{1.706313in}}{\pgfqpoint{2.136242in}{1.711899in}}{\pgfqpoint{2.132124in}{1.716018in}}%
\pgfpathcurveto{\pgfqpoint{2.128006in}{1.720136in}}{\pgfqpoint{2.122419in}{1.722450in}}{\pgfqpoint{2.116596in}{1.722450in}}%
\pgfpathcurveto{\pgfqpoint{2.110772in}{1.722450in}}{\pgfqpoint{2.105185in}{1.720136in}}{\pgfqpoint{2.101067in}{1.716018in}}%
\pgfpathcurveto{\pgfqpoint{2.096949in}{1.711899in}}{\pgfqpoint{2.094635in}{1.706313in}}{\pgfqpoint{2.094635in}{1.700489in}}%
\pgfpathcurveto{\pgfqpoint{2.094635in}{1.694665in}}{\pgfqpoint{2.096949in}{1.689079in}}{\pgfqpoint{2.101067in}{1.684961in}}%
\pgfpathcurveto{\pgfqpoint{2.105185in}{1.680843in}}{\pgfqpoint{2.110772in}{1.678529in}}{\pgfqpoint{2.116596in}{1.678529in}}%
\pgfpathlineto{\pgfqpoint{2.116596in}{1.678529in}}%
\pgfpathclose%
\pgfusepath{stroke,fill}%
\end{pgfscope}%
\begin{pgfscope}%
\pgfpathrectangle{\pgfqpoint{0.100000in}{0.209042in}}{\pgfqpoint{2.906250in}{2.906250in}}%
\pgfusepath{clip}%
\pgfsetbuttcap%
\pgfsetroundjoin%
\definecolor{currentfill}{rgb}{0.652056,0.193986,0.491611}%
\pgfsetfillcolor{currentfill}%
\pgfsetfillopacity{0.800000}%
\pgfsetlinewidth{1.003750pt}%
\definecolor{currentstroke}{rgb}{0.652056,0.193986,0.491611}%
\pgfsetstrokecolor{currentstroke}%
\pgfsetstrokeopacity{0.800000}%
\pgfsetdash{}{0pt}%
\pgfpathmoveto{\pgfqpoint{1.957270in}{1.383118in}}%
\pgfpathcurveto{\pgfqpoint{1.963094in}{1.383118in}}{\pgfqpoint{1.968680in}{1.385432in}}{\pgfqpoint{1.972799in}{1.389550in}}%
\pgfpathcurveto{\pgfqpoint{1.976917in}{1.393668in}}{\pgfqpoint{1.979231in}{1.399254in}}{\pgfqpoint{1.979231in}{1.405078in}}%
\pgfpathcurveto{\pgfqpoint{1.979231in}{1.410902in}}{\pgfqpoint{1.976917in}{1.416488in}}{\pgfqpoint{1.972799in}{1.420607in}}%
\pgfpathcurveto{\pgfqpoint{1.968680in}{1.424725in}}{\pgfqpoint{1.963094in}{1.427039in}}{\pgfqpoint{1.957270in}{1.427039in}}%
\pgfpathcurveto{\pgfqpoint{1.951446in}{1.427039in}}{\pgfqpoint{1.945860in}{1.424725in}}{\pgfqpoint{1.941742in}{1.420607in}}%
\pgfpathcurveto{\pgfqpoint{1.937624in}{1.416488in}}{\pgfqpoint{1.935310in}{1.410902in}}{\pgfqpoint{1.935310in}{1.405078in}}%
\pgfpathcurveto{\pgfqpoint{1.935310in}{1.399254in}}{\pgfqpoint{1.937624in}{1.393668in}}{\pgfqpoint{1.941742in}{1.389550in}}%
\pgfpathcurveto{\pgfqpoint{1.945860in}{1.385432in}}{\pgfqpoint{1.951446in}{1.383118in}}{\pgfqpoint{1.957270in}{1.383118in}}%
\pgfpathlineto{\pgfqpoint{1.957270in}{1.383118in}}%
\pgfpathclose%
\pgfusepath{stroke,fill}%
\end{pgfscope}%
\begin{pgfscope}%
\pgfpathrectangle{\pgfqpoint{0.100000in}{0.209042in}}{\pgfqpoint{2.906250in}{2.906250in}}%
\pgfusepath{clip}%
\pgfsetbuttcap%
\pgfsetroundjoin%
\definecolor{currentfill}{rgb}{0.390384,0.100379,0.501864}%
\pgfsetfillcolor{currentfill}%
\pgfsetfillopacity{0.800000}%
\pgfsetlinewidth{1.003750pt}%
\definecolor{currentstroke}{rgb}{0.390384,0.100379,0.501864}%
\pgfsetstrokecolor{currentstroke}%
\pgfsetstrokeopacity{0.800000}%
\pgfsetdash{}{0pt}%
\pgfpathmoveto{\pgfqpoint{1.586610in}{1.823334in}}%
\pgfpathcurveto{\pgfqpoint{1.592434in}{1.823334in}}{\pgfqpoint{1.598020in}{1.825648in}}{\pgfqpoint{1.602138in}{1.829766in}}%
\pgfpathcurveto{\pgfqpoint{1.606256in}{1.833884in}}{\pgfqpoint{1.608570in}{1.839470in}}{\pgfqpoint{1.608570in}{1.845294in}}%
\pgfpathcurveto{\pgfqpoint{1.608570in}{1.851118in}}{\pgfqpoint{1.606256in}{1.856704in}}{\pgfqpoint{1.602138in}{1.860822in}}%
\pgfpathcurveto{\pgfqpoint{1.598020in}{1.864940in}}{\pgfqpoint{1.592434in}{1.867254in}}{\pgfqpoint{1.586610in}{1.867254in}}%
\pgfpathcurveto{\pgfqpoint{1.580786in}{1.867254in}}{\pgfqpoint{1.575199in}{1.864940in}}{\pgfqpoint{1.571081in}{1.860822in}}%
\pgfpathcurveto{\pgfqpoint{1.566963in}{1.856704in}}{\pgfqpoint{1.564649in}{1.851118in}}{\pgfqpoint{1.564649in}{1.845294in}}%
\pgfpathcurveto{\pgfqpoint{1.564649in}{1.839470in}}{\pgfqpoint{1.566963in}{1.833884in}}{\pgfqpoint{1.571081in}{1.829766in}}%
\pgfpathcurveto{\pgfqpoint{1.575199in}{1.825648in}}{\pgfqpoint{1.580786in}{1.823334in}}{\pgfqpoint{1.586610in}{1.823334in}}%
\pgfpathlineto{\pgfqpoint{1.586610in}{1.823334in}}%
\pgfpathclose%
\pgfusepath{stroke,fill}%
\end{pgfscope}%
\begin{pgfscope}%
\pgfpathrectangle{\pgfqpoint{0.100000in}{0.209042in}}{\pgfqpoint{2.906250in}{2.906250in}}%
\pgfusepath{clip}%
\pgfsetbuttcap%
\pgfsetroundjoin%
\definecolor{currentfill}{rgb}{0.481929,0.136891,0.507989}%
\pgfsetfillcolor{currentfill}%
\pgfsetfillopacity{0.800000}%
\pgfsetlinewidth{1.003750pt}%
\definecolor{currentstroke}{rgb}{0.481929,0.136891,0.507989}%
\pgfsetstrokecolor{currentstroke}%
\pgfsetstrokeopacity{0.800000}%
\pgfsetdash{}{0pt}%
\pgfpathmoveto{\pgfqpoint{1.444529in}{1.833921in}}%
\pgfpathcurveto{\pgfqpoint{1.450353in}{1.833921in}}{\pgfqpoint{1.455939in}{1.836235in}}{\pgfqpoint{1.460057in}{1.840353in}}%
\pgfpathcurveto{\pgfqpoint{1.464175in}{1.844471in}}{\pgfqpoint{1.466489in}{1.850057in}}{\pgfqpoint{1.466489in}{1.855881in}}%
\pgfpathcurveto{\pgfqpoint{1.466489in}{1.861705in}}{\pgfqpoint{1.464175in}{1.867291in}}{\pgfqpoint{1.460057in}{1.871410in}}%
\pgfpathcurveto{\pgfqpoint{1.455939in}{1.875528in}}{\pgfqpoint{1.450353in}{1.877842in}}{\pgfqpoint{1.444529in}{1.877842in}}%
\pgfpathcurveto{\pgfqpoint{1.438705in}{1.877842in}}{\pgfqpoint{1.433119in}{1.875528in}}{\pgfqpoint{1.429001in}{1.871410in}}%
\pgfpathcurveto{\pgfqpoint{1.424883in}{1.867291in}}{\pgfqpoint{1.422569in}{1.861705in}}{\pgfqpoint{1.422569in}{1.855881in}}%
\pgfpathcurveto{\pgfqpoint{1.422569in}{1.850057in}}{\pgfqpoint{1.424883in}{1.844471in}}{\pgfqpoint{1.429001in}{1.840353in}}%
\pgfpathcurveto{\pgfqpoint{1.433119in}{1.836235in}}{\pgfqpoint{1.438705in}{1.833921in}}{\pgfqpoint{1.444529in}{1.833921in}}%
\pgfpathlineto{\pgfqpoint{1.444529in}{1.833921in}}%
\pgfpathclose%
\pgfusepath{stroke,fill}%
\end{pgfscope}%
\begin{pgfscope}%
\pgfpathrectangle{\pgfqpoint{0.100000in}{0.209042in}}{\pgfqpoint{2.906250in}{2.906250in}}%
\pgfusepath{clip}%
\pgfsetbuttcap%
\pgfsetroundjoin%
\definecolor{currentfill}{rgb}{0.408629,0.107930,0.504052}%
\pgfsetfillcolor{currentfill}%
\pgfsetfillopacity{0.800000}%
\pgfsetlinewidth{1.003750pt}%
\definecolor{currentstroke}{rgb}{0.408629,0.107930,0.504052}%
\pgfsetstrokecolor{currentstroke}%
\pgfsetstrokeopacity{0.800000}%
\pgfsetdash{}{0pt}%
\pgfpathmoveto{\pgfqpoint{1.823435in}{1.556550in}}%
\pgfpathcurveto{\pgfqpoint{1.829258in}{1.556550in}}{\pgfqpoint{1.834845in}{1.558864in}}{\pgfqpoint{1.838963in}{1.562982in}}%
\pgfpathcurveto{\pgfqpoint{1.843081in}{1.567100in}}{\pgfqpoint{1.845395in}{1.572687in}}{\pgfqpoint{1.845395in}{1.578511in}}%
\pgfpathcurveto{\pgfqpoint{1.845395in}{1.584334in}}{\pgfqpoint{1.843081in}{1.589921in}}{\pgfqpoint{1.838963in}{1.594039in}}%
\pgfpathcurveto{\pgfqpoint{1.834845in}{1.598157in}}{\pgfqpoint{1.829258in}{1.600471in}}{\pgfqpoint{1.823435in}{1.600471in}}%
\pgfpathcurveto{\pgfqpoint{1.817611in}{1.600471in}}{\pgfqpoint{1.812024in}{1.598157in}}{\pgfqpoint{1.807906in}{1.594039in}}%
\pgfpathcurveto{\pgfqpoint{1.803788in}{1.589921in}}{\pgfqpoint{1.801474in}{1.584334in}}{\pgfqpoint{1.801474in}{1.578511in}}%
\pgfpathcurveto{\pgfqpoint{1.801474in}{1.572687in}}{\pgfqpoint{1.803788in}{1.567100in}}{\pgfqpoint{1.807906in}{1.562982in}}%
\pgfpathcurveto{\pgfqpoint{1.812024in}{1.558864in}}{\pgfqpoint{1.817611in}{1.556550in}}{\pgfqpoint{1.823435in}{1.556550in}}%
\pgfpathlineto{\pgfqpoint{1.823435in}{1.556550in}}%
\pgfpathclose%
\pgfusepath{stroke,fill}%
\end{pgfscope}%
\begin{pgfscope}%
\pgfpathrectangle{\pgfqpoint{0.100000in}{0.209042in}}{\pgfqpoint{2.906250in}{2.906250in}}%
\pgfusepath{clip}%
\pgfsetbuttcap%
\pgfsetroundjoin%
\definecolor{currentfill}{rgb}{0.366012,0.090314,0.497960}%
\pgfsetfillcolor{currentfill}%
\pgfsetfillopacity{0.800000}%
\pgfsetlinewidth{1.003750pt}%
\definecolor{currentstroke}{rgb}{0.366012,0.090314,0.497960}%
\pgfsetstrokecolor{currentstroke}%
\pgfsetstrokeopacity{0.800000}%
\pgfsetdash{}{0pt}%
\pgfpathmoveto{\pgfqpoint{1.658075in}{1.781380in}}%
\pgfpathcurveto{\pgfqpoint{1.663899in}{1.781380in}}{\pgfqpoint{1.669485in}{1.783694in}}{\pgfqpoint{1.673603in}{1.787812in}}%
\pgfpathcurveto{\pgfqpoint{1.677721in}{1.791930in}}{\pgfqpoint{1.680035in}{1.797516in}}{\pgfqpoint{1.680035in}{1.803340in}}%
\pgfpathcurveto{\pgfqpoint{1.680035in}{1.809164in}}{\pgfqpoint{1.677721in}{1.814750in}}{\pgfqpoint{1.673603in}{1.818869in}}%
\pgfpathcurveto{\pgfqpoint{1.669485in}{1.822987in}}{\pgfqpoint{1.663899in}{1.825301in}}{\pgfqpoint{1.658075in}{1.825301in}}%
\pgfpathcurveto{\pgfqpoint{1.652251in}{1.825301in}}{\pgfqpoint{1.646665in}{1.822987in}}{\pgfqpoint{1.642546in}{1.818869in}}%
\pgfpathcurveto{\pgfqpoint{1.638428in}{1.814750in}}{\pgfqpoint{1.636114in}{1.809164in}}{\pgfqpoint{1.636114in}{1.803340in}}%
\pgfpathcurveto{\pgfqpoint{1.636114in}{1.797516in}}{\pgfqpoint{1.638428in}{1.791930in}}{\pgfqpoint{1.642546in}{1.787812in}}%
\pgfpathcurveto{\pgfqpoint{1.646665in}{1.783694in}}{\pgfqpoint{1.652251in}{1.781380in}}{\pgfqpoint{1.658075in}{1.781380in}}%
\pgfpathlineto{\pgfqpoint{1.658075in}{1.781380in}}%
\pgfpathclose%
\pgfusepath{stroke,fill}%
\end{pgfscope}%
\begin{pgfscope}%
\pgfpathrectangle{\pgfqpoint{0.100000in}{0.209042in}}{\pgfqpoint{2.906250in}{2.906250in}}%
\pgfusepath{clip}%
\pgfsetbuttcap%
\pgfsetroundjoin%
\definecolor{currentfill}{rgb}{0.445163,0.122724,0.506901}%
\pgfsetfillcolor{currentfill}%
\pgfsetfillopacity{0.800000}%
\pgfsetlinewidth{1.003750pt}%
\definecolor{currentstroke}{rgb}{0.445163,0.122724,0.506901}%
\pgfsetstrokecolor{currentstroke}%
\pgfsetstrokeopacity{0.800000}%
\pgfsetdash{}{0pt}%
\pgfpathmoveto{\pgfqpoint{1.653599in}{1.864334in}}%
\pgfpathcurveto{\pgfqpoint{1.659423in}{1.864334in}}{\pgfqpoint{1.665009in}{1.866648in}}{\pgfqpoint{1.669127in}{1.870766in}}%
\pgfpathcurveto{\pgfqpoint{1.673245in}{1.874884in}}{\pgfqpoint{1.675559in}{1.880471in}}{\pgfqpoint{1.675559in}{1.886295in}}%
\pgfpathcurveto{\pgfqpoint{1.675559in}{1.892118in}}{\pgfqpoint{1.673245in}{1.897705in}}{\pgfqpoint{1.669127in}{1.901823in}}%
\pgfpathcurveto{\pgfqpoint{1.665009in}{1.905941in}}{\pgfqpoint{1.659423in}{1.908255in}}{\pgfqpoint{1.653599in}{1.908255in}}%
\pgfpathcurveto{\pgfqpoint{1.647775in}{1.908255in}}{\pgfqpoint{1.642189in}{1.905941in}}{\pgfqpoint{1.638071in}{1.901823in}}%
\pgfpathcurveto{\pgfqpoint{1.633953in}{1.897705in}}{\pgfqpoint{1.631639in}{1.892118in}}{\pgfqpoint{1.631639in}{1.886295in}}%
\pgfpathcurveto{\pgfqpoint{1.631639in}{1.880471in}}{\pgfqpoint{1.633953in}{1.874884in}}{\pgfqpoint{1.638071in}{1.870766in}}%
\pgfpathcurveto{\pgfqpoint{1.642189in}{1.866648in}}{\pgfqpoint{1.647775in}{1.864334in}}{\pgfqpoint{1.653599in}{1.864334in}}%
\pgfpathlineto{\pgfqpoint{1.653599in}{1.864334in}}%
\pgfpathclose%
\pgfusepath{stroke,fill}%
\end{pgfscope}%
\begin{pgfscope}%
\pgfpathrectangle{\pgfqpoint{0.100000in}{0.209042in}}{\pgfqpoint{2.906250in}{2.906250in}}%
\pgfusepath{clip}%
\pgfsetbuttcap%
\pgfsetroundjoin%
\definecolor{currentfill}{rgb}{0.709962,0.212797,0.477201}%
\pgfsetfillcolor{currentfill}%
\pgfsetfillopacity{0.800000}%
\pgfsetlinewidth{1.003750pt}%
\definecolor{currentstroke}{rgb}{0.709962,0.212797,0.477201}%
\pgfsetstrokecolor{currentstroke}%
\pgfsetstrokeopacity{0.800000}%
\pgfsetdash{}{0pt}%
\pgfpathmoveto{\pgfqpoint{2.057095in}{1.872895in}}%
\pgfpathcurveto{\pgfqpoint{2.062919in}{1.872895in}}{\pgfqpoint{2.068505in}{1.875209in}}{\pgfqpoint{2.072623in}{1.879327in}}%
\pgfpathcurveto{\pgfqpoint{2.076741in}{1.883445in}}{\pgfqpoint{2.079055in}{1.889031in}}{\pgfqpoint{2.079055in}{1.894855in}}%
\pgfpathcurveto{\pgfqpoint{2.079055in}{1.900679in}}{\pgfqpoint{2.076741in}{1.906265in}}{\pgfqpoint{2.072623in}{1.910383in}}%
\pgfpathcurveto{\pgfqpoint{2.068505in}{1.914502in}}{\pgfqpoint{2.062919in}{1.916815in}}{\pgfqpoint{2.057095in}{1.916815in}}%
\pgfpathcurveto{\pgfqpoint{2.051271in}{1.916815in}}{\pgfqpoint{2.045684in}{1.914502in}}{\pgfqpoint{2.041566in}{1.910383in}}%
\pgfpathcurveto{\pgfqpoint{2.037448in}{1.906265in}}{\pgfqpoint{2.035134in}{1.900679in}}{\pgfqpoint{2.035134in}{1.894855in}}%
\pgfpathcurveto{\pgfqpoint{2.035134in}{1.889031in}}{\pgfqpoint{2.037448in}{1.883445in}}{\pgfqpoint{2.041566in}{1.879327in}}%
\pgfpathcurveto{\pgfqpoint{2.045684in}{1.875209in}}{\pgfqpoint{2.051271in}{1.872895in}}{\pgfqpoint{2.057095in}{1.872895in}}%
\pgfpathlineto{\pgfqpoint{2.057095in}{1.872895in}}%
\pgfpathclose%
\pgfusepath{stroke,fill}%
\end{pgfscope}%
\begin{pgfscope}%
\pgfpathrectangle{\pgfqpoint{0.100000in}{0.209042in}}{\pgfqpoint{2.906250in}{2.906250in}}%
\pgfusepath{clip}%
\pgfsetbuttcap%
\pgfsetroundjoin%
\definecolor{currentfill}{rgb}{0.378211,0.095332,0.500067}%
\pgfsetfillcolor{currentfill}%
\pgfsetfillopacity{0.800000}%
\pgfsetlinewidth{1.003750pt}%
\definecolor{currentstroke}{rgb}{0.378211,0.095332,0.500067}%
\pgfsetstrokecolor{currentstroke}%
\pgfsetstrokeopacity{0.800000}%
\pgfsetdash{}{0pt}%
\pgfpathmoveto{\pgfqpoint{1.511463in}{1.695599in}}%
\pgfpathcurveto{\pgfqpoint{1.517287in}{1.695599in}}{\pgfqpoint{1.522873in}{1.697912in}}{\pgfqpoint{1.526991in}{1.702031in}}%
\pgfpathcurveto{\pgfqpoint{1.531109in}{1.706149in}}{\pgfqpoint{1.533423in}{1.711735in}}{\pgfqpoint{1.533423in}{1.717559in}}%
\pgfpathcurveto{\pgfqpoint{1.533423in}{1.723383in}}{\pgfqpoint{1.531109in}{1.728969in}}{\pgfqpoint{1.526991in}{1.733087in}}%
\pgfpathcurveto{\pgfqpoint{1.522873in}{1.737205in}}{\pgfqpoint{1.517287in}{1.739519in}}{\pgfqpoint{1.511463in}{1.739519in}}%
\pgfpathcurveto{\pgfqpoint{1.505639in}{1.739519in}}{\pgfqpoint{1.500053in}{1.737205in}}{\pgfqpoint{1.495935in}{1.733087in}}%
\pgfpathcurveto{\pgfqpoint{1.491816in}{1.728969in}}{\pgfqpoint{1.489503in}{1.723383in}}{\pgfqpoint{1.489503in}{1.717559in}}%
\pgfpathcurveto{\pgfqpoint{1.489503in}{1.711735in}}{\pgfqpoint{1.491816in}{1.706149in}}{\pgfqpoint{1.495935in}{1.702031in}}%
\pgfpathcurveto{\pgfqpoint{1.500053in}{1.697912in}}{\pgfqpoint{1.505639in}{1.695599in}}{\pgfqpoint{1.511463in}{1.695599in}}%
\pgfpathlineto{\pgfqpoint{1.511463in}{1.695599in}}%
\pgfpathclose%
\pgfusepath{stroke,fill}%
\end{pgfscope}%
\begin{pgfscope}%
\pgfpathrectangle{\pgfqpoint{0.100000in}{0.209042in}}{\pgfqpoint{2.906250in}{2.906250in}}%
\pgfusepath{clip}%
\pgfsetbuttcap%
\pgfsetroundjoin%
\definecolor{currentfill}{rgb}{0.729216,0.219437,0.471279}%
\pgfsetfillcolor{currentfill}%
\pgfsetfillopacity{0.800000}%
\pgfsetlinewidth{1.003750pt}%
\definecolor{currentstroke}{rgb}{0.729216,0.219437,0.471279}%
\pgfsetstrokecolor{currentstroke}%
\pgfsetstrokeopacity{0.800000}%
\pgfsetdash{}{0pt}%
\pgfpathmoveto{\pgfqpoint{2.056145in}{1.399780in}}%
\pgfpathcurveto{\pgfqpoint{2.061969in}{1.399780in}}{\pgfqpoint{2.067555in}{1.402094in}}{\pgfqpoint{2.071674in}{1.406212in}}%
\pgfpathcurveto{\pgfqpoint{2.075792in}{1.410330in}}{\pgfqpoint{2.078106in}{1.415916in}}{\pgfqpoint{2.078106in}{1.421740in}}%
\pgfpathcurveto{\pgfqpoint{2.078106in}{1.427564in}}{\pgfqpoint{2.075792in}{1.433150in}}{\pgfqpoint{2.071674in}{1.437269in}}%
\pgfpathcurveto{\pgfqpoint{2.067555in}{1.441387in}}{\pgfqpoint{2.061969in}{1.443701in}}{\pgfqpoint{2.056145in}{1.443701in}}%
\pgfpathcurveto{\pgfqpoint{2.050321in}{1.443701in}}{\pgfqpoint{2.044735in}{1.441387in}}{\pgfqpoint{2.040617in}{1.437269in}}%
\pgfpathcurveto{\pgfqpoint{2.036499in}{1.433150in}}{\pgfqpoint{2.034185in}{1.427564in}}{\pgfqpoint{2.034185in}{1.421740in}}%
\pgfpathcurveto{\pgfqpoint{2.034185in}{1.415916in}}{\pgfqpoint{2.036499in}{1.410330in}}{\pgfqpoint{2.040617in}{1.406212in}}%
\pgfpathcurveto{\pgfqpoint{2.044735in}{1.402094in}}{\pgfqpoint{2.050321in}{1.399780in}}{\pgfqpoint{2.056145in}{1.399780in}}%
\pgfpathlineto{\pgfqpoint{2.056145in}{1.399780in}}%
\pgfpathclose%
\pgfusepath{stroke,fill}%
\end{pgfscope}%
\begin{pgfscope}%
\pgfpathrectangle{\pgfqpoint{0.100000in}{0.209042in}}{\pgfqpoint{2.906250in}{2.906250in}}%
\pgfusepath{clip}%
\pgfsetbuttcap%
\pgfsetroundjoin%
\definecolor{currentfill}{rgb}{0.439062,0.120298,0.506555}%
\pgfsetfillcolor{currentfill}%
\pgfsetfillopacity{0.800000}%
\pgfsetlinewidth{1.003750pt}%
\definecolor{currentstroke}{rgb}{0.439062,0.120298,0.506555}%
\pgfsetstrokecolor{currentstroke}%
\pgfsetstrokeopacity{0.800000}%
\pgfsetdash{}{0pt}%
\pgfpathmoveto{\pgfqpoint{1.783372in}{1.494259in}}%
\pgfpathcurveto{\pgfqpoint{1.789196in}{1.494259in}}{\pgfqpoint{1.794782in}{1.496573in}}{\pgfqpoint{1.798900in}{1.500691in}}%
\pgfpathcurveto{\pgfqpoint{1.803018in}{1.504810in}}{\pgfqpoint{1.805332in}{1.510396in}}{\pgfqpoint{1.805332in}{1.516220in}}%
\pgfpathcurveto{\pgfqpoint{1.805332in}{1.522044in}}{\pgfqpoint{1.803018in}{1.527630in}}{\pgfqpoint{1.798900in}{1.531748in}}%
\pgfpathcurveto{\pgfqpoint{1.794782in}{1.535866in}}{\pgfqpoint{1.789196in}{1.538180in}}{\pgfqpoint{1.783372in}{1.538180in}}%
\pgfpathcurveto{\pgfqpoint{1.777548in}{1.538180in}}{\pgfqpoint{1.771962in}{1.535866in}}{\pgfqpoint{1.767843in}{1.531748in}}%
\pgfpathcurveto{\pgfqpoint{1.763725in}{1.527630in}}{\pgfqpoint{1.761411in}{1.522044in}}{\pgfqpoint{1.761411in}{1.516220in}}%
\pgfpathcurveto{\pgfqpoint{1.761411in}{1.510396in}}{\pgfqpoint{1.763725in}{1.504810in}}{\pgfqpoint{1.767843in}{1.500691in}}%
\pgfpathcurveto{\pgfqpoint{1.771962in}{1.496573in}}{\pgfqpoint{1.777548in}{1.494259in}}{\pgfqpoint{1.783372in}{1.494259in}}%
\pgfpathlineto{\pgfqpoint{1.783372in}{1.494259in}}%
\pgfpathclose%
\pgfusepath{stroke,fill}%
\end{pgfscope}%
\begin{pgfscope}%
\pgfpathrectangle{\pgfqpoint{0.100000in}{0.209042in}}{\pgfqpoint{2.906250in}{2.906250in}}%
\pgfusepath{clip}%
\pgfsetbuttcap%
\pgfsetroundjoin%
\definecolor{currentfill}{rgb}{0.481929,0.136891,0.507989}%
\pgfsetfillcolor{currentfill}%
\pgfsetfillopacity{0.800000}%
\pgfsetlinewidth{1.003750pt}%
\definecolor{currentstroke}{rgb}{0.481929,0.136891,0.507989}%
\pgfsetstrokecolor{currentstroke}%
\pgfsetstrokeopacity{0.800000}%
\pgfsetdash{}{0pt}%
\pgfpathmoveto{\pgfqpoint{1.494050in}{1.488063in}}%
\pgfpathcurveto{\pgfqpoint{1.499874in}{1.488063in}}{\pgfqpoint{1.505460in}{1.490376in}}{\pgfqpoint{1.509578in}{1.494495in}}%
\pgfpathcurveto{\pgfqpoint{1.513696in}{1.498613in}}{\pgfqpoint{1.516010in}{1.504199in}}{\pgfqpoint{1.516010in}{1.510023in}}%
\pgfpathcurveto{\pgfqpoint{1.516010in}{1.515847in}}{\pgfqpoint{1.513696in}{1.521433in}}{\pgfqpoint{1.509578in}{1.525551in}}%
\pgfpathcurveto{\pgfqpoint{1.505460in}{1.529669in}}{\pgfqpoint{1.499874in}{1.531983in}}{\pgfqpoint{1.494050in}{1.531983in}}%
\pgfpathcurveto{\pgfqpoint{1.488226in}{1.531983in}}{\pgfqpoint{1.482640in}{1.529669in}}{\pgfqpoint{1.478521in}{1.525551in}}%
\pgfpathcurveto{\pgfqpoint{1.474403in}{1.521433in}}{\pgfqpoint{1.472089in}{1.515847in}}{\pgfqpoint{1.472089in}{1.510023in}}%
\pgfpathcurveto{\pgfqpoint{1.472089in}{1.504199in}}{\pgfqpoint{1.474403in}{1.498613in}}{\pgfqpoint{1.478521in}{1.494495in}}%
\pgfpathcurveto{\pgfqpoint{1.482640in}{1.490376in}}{\pgfqpoint{1.488226in}{1.488063in}}{\pgfqpoint{1.494050in}{1.488063in}}%
\pgfpathlineto{\pgfqpoint{1.494050in}{1.488063in}}%
\pgfpathclose%
\pgfusepath{stroke,fill}%
\end{pgfscope}%
\begin{pgfscope}%
\pgfpathrectangle{\pgfqpoint{0.100000in}{0.209042in}}{\pgfqpoint{2.906250in}{2.906250in}}%
\pgfusepath{clip}%
\pgfsetbuttcap%
\pgfsetroundjoin%
\definecolor{currentfill}{rgb}{0.353773,0.085373,0.495501}%
\pgfsetfillcolor{currentfill}%
\pgfsetfillopacity{0.800000}%
\pgfsetlinewidth{1.003750pt}%
\definecolor{currentstroke}{rgb}{0.353773,0.085373,0.495501}%
\pgfsetstrokecolor{currentstroke}%
\pgfsetstrokeopacity{0.800000}%
\pgfsetdash{}{0pt}%
\pgfpathmoveto{\pgfqpoint{1.719085in}{1.592171in}}%
\pgfpathcurveto{\pgfqpoint{1.724908in}{1.592171in}}{\pgfqpoint{1.730495in}{1.594484in}}{\pgfqpoint{1.734613in}{1.598603in}}%
\pgfpathcurveto{\pgfqpoint{1.738731in}{1.602721in}}{\pgfqpoint{1.741045in}{1.608307in}}{\pgfqpoint{1.741045in}{1.614131in}}%
\pgfpathcurveto{\pgfqpoint{1.741045in}{1.619955in}}{\pgfqpoint{1.738731in}{1.625541in}}{\pgfqpoint{1.734613in}{1.629659in}}%
\pgfpathcurveto{\pgfqpoint{1.730495in}{1.633777in}}{\pgfqpoint{1.724908in}{1.636091in}}{\pgfqpoint{1.719085in}{1.636091in}}%
\pgfpathcurveto{\pgfqpoint{1.713261in}{1.636091in}}{\pgfqpoint{1.707674in}{1.633777in}}{\pgfqpoint{1.703556in}{1.629659in}}%
\pgfpathcurveto{\pgfqpoint{1.699438in}{1.625541in}}{\pgfqpoint{1.697124in}{1.619955in}}{\pgfqpoint{1.697124in}{1.614131in}}%
\pgfpathcurveto{\pgfqpoint{1.697124in}{1.608307in}}{\pgfqpoint{1.699438in}{1.602721in}}{\pgfqpoint{1.703556in}{1.598603in}}%
\pgfpathcurveto{\pgfqpoint{1.707674in}{1.594484in}}{\pgfqpoint{1.713261in}{1.592171in}}{\pgfqpoint{1.719085in}{1.592171in}}%
\pgfpathlineto{\pgfqpoint{1.719085in}{1.592171in}}%
\pgfpathclose%
\pgfusepath{stroke,fill}%
\end{pgfscope}%
\begin{pgfscope}%
\pgfpathrectangle{\pgfqpoint{0.100000in}{0.209042in}}{\pgfqpoint{2.906250in}{2.906250in}}%
\pgfusepath{clip}%
\pgfsetbuttcap%
\pgfsetroundjoin%
\definecolor{currentfill}{rgb}{0.426877,0.115395,0.505714}%
\pgfsetfillcolor{currentfill}%
\pgfsetfillopacity{0.800000}%
\pgfsetlinewidth{1.003750pt}%
\definecolor{currentstroke}{rgb}{0.426877,0.115395,0.505714}%
\pgfsetstrokecolor{currentstroke}%
\pgfsetstrokeopacity{0.800000}%
\pgfsetdash{}{0pt}%
\pgfpathmoveto{\pgfqpoint{1.842476in}{1.760008in}}%
\pgfpathcurveto{\pgfqpoint{1.848299in}{1.760008in}}{\pgfqpoint{1.853886in}{1.762322in}}{\pgfqpoint{1.858004in}{1.766440in}}%
\pgfpathcurveto{\pgfqpoint{1.862122in}{1.770558in}}{\pgfqpoint{1.864436in}{1.776145in}}{\pgfqpoint{1.864436in}{1.781969in}}%
\pgfpathcurveto{\pgfqpoint{1.864436in}{1.787793in}}{\pgfqpoint{1.862122in}{1.793379in}}{\pgfqpoint{1.858004in}{1.797497in}}%
\pgfpathcurveto{\pgfqpoint{1.853886in}{1.801615in}}{\pgfqpoint{1.848299in}{1.803929in}}{\pgfqpoint{1.842476in}{1.803929in}}%
\pgfpathcurveto{\pgfqpoint{1.836652in}{1.803929in}}{\pgfqpoint{1.831065in}{1.801615in}}{\pgfqpoint{1.826947in}{1.797497in}}%
\pgfpathcurveto{\pgfqpoint{1.822829in}{1.793379in}}{\pgfqpoint{1.820515in}{1.787793in}}{\pgfqpoint{1.820515in}{1.781969in}}%
\pgfpathcurveto{\pgfqpoint{1.820515in}{1.776145in}}{\pgfqpoint{1.822829in}{1.770558in}}{\pgfqpoint{1.826947in}{1.766440in}}%
\pgfpathcurveto{\pgfqpoint{1.831065in}{1.762322in}}{\pgfqpoint{1.836652in}{1.760008in}}{\pgfqpoint{1.842476in}{1.760008in}}%
\pgfpathlineto{\pgfqpoint{1.842476in}{1.760008in}}%
\pgfpathclose%
\pgfusepath{stroke,fill}%
\end{pgfscope}%
\begin{pgfscope}%
\pgfpathrectangle{\pgfqpoint{0.100000in}{0.209042in}}{\pgfqpoint{2.906250in}{2.906250in}}%
\pgfusepath{clip}%
\pgfsetbuttcap%
\pgfsetroundjoin%
\definecolor{currentfill}{rgb}{0.779968,0.238851,0.452765}%
\pgfsetfillcolor{currentfill}%
\pgfsetfillopacity{0.800000}%
\pgfsetlinewidth{1.003750pt}%
\definecolor{currentstroke}{rgb}{0.779968,0.238851,0.452765}%
\pgfsetstrokecolor{currentstroke}%
\pgfsetstrokeopacity{0.800000}%
\pgfsetdash{}{0pt}%
\pgfpathmoveto{\pgfqpoint{1.153932in}{1.582380in}}%
\pgfpathcurveto{\pgfqpoint{1.159755in}{1.582380in}}{\pgfqpoint{1.165342in}{1.584693in}}{\pgfqpoint{1.169460in}{1.588812in}}%
\pgfpathcurveto{\pgfqpoint{1.173578in}{1.592930in}}{\pgfqpoint{1.175892in}{1.598516in}}{\pgfqpoint{1.175892in}{1.604340in}}%
\pgfpathcurveto{\pgfqpoint{1.175892in}{1.610164in}}{\pgfqpoint{1.173578in}{1.615750in}}{\pgfqpoint{1.169460in}{1.619868in}}%
\pgfpathcurveto{\pgfqpoint{1.165342in}{1.623986in}}{\pgfqpoint{1.159755in}{1.626300in}}{\pgfqpoint{1.153932in}{1.626300in}}%
\pgfpathcurveto{\pgfqpoint{1.148108in}{1.626300in}}{\pgfqpoint{1.142521in}{1.623986in}}{\pgfqpoint{1.138403in}{1.619868in}}%
\pgfpathcurveto{\pgfqpoint{1.134285in}{1.615750in}}{\pgfqpoint{1.131971in}{1.610164in}}{\pgfqpoint{1.131971in}{1.604340in}}%
\pgfpathcurveto{\pgfqpoint{1.131971in}{1.598516in}}{\pgfqpoint{1.134285in}{1.592930in}}{\pgfqpoint{1.138403in}{1.588812in}}%
\pgfpathcurveto{\pgfqpoint{1.142521in}{1.584693in}}{\pgfqpoint{1.148108in}{1.582380in}}{\pgfqpoint{1.153932in}{1.582380in}}%
\pgfpathlineto{\pgfqpoint{1.153932in}{1.582380in}}%
\pgfpathclose%
\pgfusepath{stroke,fill}%
\end{pgfscope}%
\begin{pgfscope}%
\pgfpathrectangle{\pgfqpoint{0.100000in}{0.209042in}}{\pgfqpoint{2.906250in}{2.906250in}}%
\pgfusepath{clip}%
\pgfsetbuttcap%
\pgfsetroundjoin%
\definecolor{currentfill}{rgb}{0.594508,0.175701,0.501241}%
\pgfsetfillcolor{currentfill}%
\pgfsetfillopacity{0.800000}%
\pgfsetlinewidth{1.003750pt}%
\definecolor{currentstroke}{rgb}{0.594508,0.175701,0.501241}%
\pgfsetstrokecolor{currentstroke}%
\pgfsetstrokeopacity{0.800000}%
\pgfsetdash{}{0pt}%
\pgfpathmoveto{\pgfqpoint{2.041146in}{1.600945in}}%
\pgfpathcurveto{\pgfqpoint{2.046970in}{1.600945in}}{\pgfqpoint{2.052556in}{1.603259in}}{\pgfqpoint{2.056674in}{1.607377in}}%
\pgfpathcurveto{\pgfqpoint{2.060792in}{1.611495in}}{\pgfqpoint{2.063106in}{1.617082in}}{\pgfqpoint{2.063106in}{1.622906in}}%
\pgfpathcurveto{\pgfqpoint{2.063106in}{1.628730in}}{\pgfqpoint{2.060792in}{1.634316in}}{\pgfqpoint{2.056674in}{1.638434in}}%
\pgfpathcurveto{\pgfqpoint{2.052556in}{1.642552in}}{\pgfqpoint{2.046970in}{1.644866in}}{\pgfqpoint{2.041146in}{1.644866in}}%
\pgfpathcurveto{\pgfqpoint{2.035322in}{1.644866in}}{\pgfqpoint{2.029736in}{1.642552in}}{\pgfqpoint{2.025618in}{1.638434in}}%
\pgfpathcurveto{\pgfqpoint{2.021500in}{1.634316in}}{\pgfqpoint{2.019186in}{1.628730in}}{\pgfqpoint{2.019186in}{1.622906in}}%
\pgfpathcurveto{\pgfqpoint{2.019186in}{1.617082in}}{\pgfqpoint{2.021500in}{1.611495in}}{\pgfqpoint{2.025618in}{1.607377in}}%
\pgfpathcurveto{\pgfqpoint{2.029736in}{1.603259in}}{\pgfqpoint{2.035322in}{1.600945in}}{\pgfqpoint{2.041146in}{1.600945in}}%
\pgfpathlineto{\pgfqpoint{2.041146in}{1.600945in}}%
\pgfpathclose%
\pgfusepath{stroke,fill}%
\end{pgfscope}%
\begin{pgfscope}%
\pgfpathrectangle{\pgfqpoint{0.100000in}{0.209042in}}{\pgfqpoint{2.906250in}{2.906250in}}%
\pgfusepath{clip}%
\pgfsetbuttcap%
\pgfsetroundjoin%
\definecolor{currentfill}{rgb}{0.562866,0.165368,0.504692}%
\pgfsetfillcolor{currentfill}%
\pgfsetfillopacity{0.800000}%
\pgfsetlinewidth{1.003750pt}%
\definecolor{currentstroke}{rgb}{0.562866,0.165368,0.504692}%
\pgfsetstrokecolor{currentstroke}%
\pgfsetstrokeopacity{0.800000}%
\pgfsetdash{}{0pt}%
\pgfpathmoveto{\pgfqpoint{1.940160in}{1.833759in}}%
\pgfpathcurveto{\pgfqpoint{1.945984in}{1.833759in}}{\pgfqpoint{1.951570in}{1.836073in}}{\pgfqpoint{1.955689in}{1.840191in}}%
\pgfpathcurveto{\pgfqpoint{1.959807in}{1.844309in}}{\pgfqpoint{1.962121in}{1.849895in}}{\pgfqpoint{1.962121in}{1.855719in}}%
\pgfpathcurveto{\pgfqpoint{1.962121in}{1.861543in}}{\pgfqpoint{1.959807in}{1.867129in}}{\pgfqpoint{1.955689in}{1.871247in}}%
\pgfpathcurveto{\pgfqpoint{1.951570in}{1.875366in}}{\pgfqpoint{1.945984in}{1.877679in}}{\pgfqpoint{1.940160in}{1.877679in}}%
\pgfpathcurveto{\pgfqpoint{1.934336in}{1.877679in}}{\pgfqpoint{1.928750in}{1.875366in}}{\pgfqpoint{1.924632in}{1.871247in}}%
\pgfpathcurveto{\pgfqpoint{1.920514in}{1.867129in}}{\pgfqpoint{1.918200in}{1.861543in}}{\pgfqpoint{1.918200in}{1.855719in}}%
\pgfpathcurveto{\pgfqpoint{1.918200in}{1.849895in}}{\pgfqpoint{1.920514in}{1.844309in}}{\pgfqpoint{1.924632in}{1.840191in}}%
\pgfpathcurveto{\pgfqpoint{1.928750in}{1.836073in}}{\pgfqpoint{1.934336in}{1.833759in}}{\pgfqpoint{1.940160in}{1.833759in}}%
\pgfpathlineto{\pgfqpoint{1.940160in}{1.833759in}}%
\pgfpathclose%
\pgfusepath{stroke,fill}%
\end{pgfscope}%
\begin{pgfscope}%
\pgfpathrectangle{\pgfqpoint{0.100000in}{0.209042in}}{\pgfqpoint{2.906250in}{2.906250in}}%
\pgfusepath{clip}%
\pgfsetbuttcap%
\pgfsetroundjoin%
\definecolor{currentfill}{rgb}{0.347636,0.082946,0.494121}%
\pgfsetfillcolor{currentfill}%
\pgfsetfillopacity{0.800000}%
\pgfsetlinewidth{1.003750pt}%
\definecolor{currentstroke}{rgb}{0.347636,0.082946,0.494121}%
\pgfsetstrokecolor{currentstroke}%
\pgfsetstrokeopacity{0.800000}%
\pgfsetdash{}{0pt}%
\pgfpathmoveto{\pgfqpoint{1.601564in}{1.670298in}}%
\pgfpathcurveto{\pgfqpoint{1.607388in}{1.670298in}}{\pgfqpoint{1.612974in}{1.672612in}}{\pgfqpoint{1.617093in}{1.676730in}}%
\pgfpathcurveto{\pgfqpoint{1.621211in}{1.680848in}}{\pgfqpoint{1.623525in}{1.686434in}}{\pgfqpoint{1.623525in}{1.692258in}}%
\pgfpathcurveto{\pgfqpoint{1.623525in}{1.698082in}}{\pgfqpoint{1.621211in}{1.703668in}}{\pgfqpoint{1.617093in}{1.707786in}}%
\pgfpathcurveto{\pgfqpoint{1.612974in}{1.711905in}}{\pgfqpoint{1.607388in}{1.714218in}}{\pgfqpoint{1.601564in}{1.714218in}}%
\pgfpathcurveto{\pgfqpoint{1.595740in}{1.714218in}}{\pgfqpoint{1.590154in}{1.711905in}}{\pgfqpoint{1.586036in}{1.707786in}}%
\pgfpathcurveto{\pgfqpoint{1.581918in}{1.703668in}}{\pgfqpoint{1.579604in}{1.698082in}}{\pgfqpoint{1.579604in}{1.692258in}}%
\pgfpathcurveto{\pgfqpoint{1.579604in}{1.686434in}}{\pgfqpoint{1.581918in}{1.680848in}}{\pgfqpoint{1.586036in}{1.676730in}}%
\pgfpathcurveto{\pgfqpoint{1.590154in}{1.672612in}}{\pgfqpoint{1.595740in}{1.670298in}}{\pgfqpoint{1.601564in}{1.670298in}}%
\pgfpathlineto{\pgfqpoint{1.601564in}{1.670298in}}%
\pgfpathclose%
\pgfusepath{stroke,fill}%
\end{pgfscope}%
\begin{pgfscope}%
\pgfpathrectangle{\pgfqpoint{0.100000in}{0.209042in}}{\pgfqpoint{2.906250in}{2.906250in}}%
\pgfusepath{clip}%
\pgfsetbuttcap%
\pgfsetroundjoin%
\definecolor{currentfill}{rgb}{0.366012,0.090314,0.497960}%
\pgfsetfillcolor{currentfill}%
\pgfsetfillopacity{0.800000}%
\pgfsetlinewidth{1.003750pt}%
\definecolor{currentstroke}{rgb}{0.366012,0.090314,0.497960}%
\pgfsetstrokecolor{currentstroke}%
\pgfsetstrokeopacity{0.800000}%
\pgfsetdash{}{0pt}%
\pgfpathmoveto{\pgfqpoint{1.595790in}{1.611351in}}%
\pgfpathcurveto{\pgfqpoint{1.601614in}{1.611351in}}{\pgfqpoint{1.607200in}{1.613665in}}{\pgfqpoint{1.611318in}{1.617783in}}%
\pgfpathcurveto{\pgfqpoint{1.615436in}{1.621901in}}{\pgfqpoint{1.617750in}{1.627487in}}{\pgfqpoint{1.617750in}{1.633311in}}%
\pgfpathcurveto{\pgfqpoint{1.617750in}{1.639135in}}{\pgfqpoint{1.615436in}{1.644721in}}{\pgfqpoint{1.611318in}{1.648839in}}%
\pgfpathcurveto{\pgfqpoint{1.607200in}{1.652957in}}{\pgfqpoint{1.601614in}{1.655271in}}{\pgfqpoint{1.595790in}{1.655271in}}%
\pgfpathcurveto{\pgfqpoint{1.589966in}{1.655271in}}{\pgfqpoint{1.584380in}{1.652957in}}{\pgfqpoint{1.580262in}{1.648839in}}%
\pgfpathcurveto{\pgfqpoint{1.576143in}{1.644721in}}{\pgfqpoint{1.573830in}{1.639135in}}{\pgfqpoint{1.573830in}{1.633311in}}%
\pgfpathcurveto{\pgfqpoint{1.573830in}{1.627487in}}{\pgfqpoint{1.576143in}{1.621901in}}{\pgfqpoint{1.580262in}{1.617783in}}%
\pgfpathcurveto{\pgfqpoint{1.584380in}{1.613665in}}{\pgfqpoint{1.589966in}{1.611351in}}{\pgfqpoint{1.595790in}{1.611351in}}%
\pgfpathlineto{\pgfqpoint{1.595790in}{1.611351in}}%
\pgfpathclose%
\pgfusepath{stroke,fill}%
\end{pgfscope}%
\begin{pgfscope}%
\pgfpathrectangle{\pgfqpoint{0.100000in}{0.209042in}}{\pgfqpoint{2.906250in}{2.906250in}}%
\pgfusepath{clip}%
\pgfsetbuttcap%
\pgfsetroundjoin%
\definecolor{currentfill}{rgb}{0.469640,0.132245,0.507809}%
\pgfsetfillcolor{currentfill}%
\pgfsetfillopacity{0.800000}%
\pgfsetlinewidth{1.003750pt}%
\definecolor{currentstroke}{rgb}{0.469640,0.132245,0.507809}%
\pgfsetstrokecolor{currentstroke}%
\pgfsetstrokeopacity{0.800000}%
\pgfsetdash{}{0pt}%
\pgfpathmoveto{\pgfqpoint{1.777467in}{1.836967in}}%
\pgfpathcurveto{\pgfqpoint{1.783291in}{1.836967in}}{\pgfqpoint{1.788877in}{1.839280in}}{\pgfqpoint{1.792995in}{1.843399in}}%
\pgfpathcurveto{\pgfqpoint{1.797113in}{1.847517in}}{\pgfqpoint{1.799427in}{1.853103in}}{\pgfqpoint{1.799427in}{1.858927in}}%
\pgfpathcurveto{\pgfqpoint{1.799427in}{1.864751in}}{\pgfqpoint{1.797113in}{1.870337in}}{\pgfqpoint{1.792995in}{1.874455in}}%
\pgfpathcurveto{\pgfqpoint{1.788877in}{1.878573in}}{\pgfqpoint{1.783291in}{1.880887in}}{\pgfqpoint{1.777467in}{1.880887in}}%
\pgfpathcurveto{\pgfqpoint{1.771643in}{1.880887in}}{\pgfqpoint{1.766057in}{1.878573in}}{\pgfqpoint{1.761939in}{1.874455in}}%
\pgfpathcurveto{\pgfqpoint{1.757821in}{1.870337in}}{\pgfqpoint{1.755507in}{1.864751in}}{\pgfqpoint{1.755507in}{1.858927in}}%
\pgfpathcurveto{\pgfqpoint{1.755507in}{1.853103in}}{\pgfqpoint{1.757821in}{1.847517in}}{\pgfqpoint{1.761939in}{1.843399in}}%
\pgfpathcurveto{\pgfqpoint{1.766057in}{1.839280in}}{\pgfqpoint{1.771643in}{1.836967in}}{\pgfqpoint{1.777467in}{1.836967in}}%
\pgfpathlineto{\pgfqpoint{1.777467in}{1.836967in}}%
\pgfpathclose%
\pgfusepath{stroke,fill}%
\end{pgfscope}%
\begin{pgfscope}%
\pgfpathrectangle{\pgfqpoint{0.100000in}{0.209042in}}{\pgfqpoint{2.906250in}{2.906250in}}%
\pgfusepath{clip}%
\pgfsetbuttcap%
\pgfsetroundjoin%
\definecolor{currentfill}{rgb}{0.632805,0.187893,0.495332}%
\pgfsetfillcolor{currentfill}%
\pgfsetfillopacity{0.800000}%
\pgfsetlinewidth{1.003750pt}%
\definecolor{currentstroke}{rgb}{0.632805,0.187893,0.495332}%
\pgfsetstrokecolor{currentstroke}%
\pgfsetstrokeopacity{0.800000}%
\pgfsetdash{}{0pt}%
\pgfpathmoveto{\pgfqpoint{1.466421in}{1.394561in}}%
\pgfpathcurveto{\pgfqpoint{1.472245in}{1.394561in}}{\pgfqpoint{1.477831in}{1.396875in}}{\pgfqpoint{1.481949in}{1.400993in}}%
\pgfpathcurveto{\pgfqpoint{1.486067in}{1.405111in}}{\pgfqpoint{1.488381in}{1.410697in}}{\pgfqpoint{1.488381in}{1.416521in}}%
\pgfpathcurveto{\pgfqpoint{1.488381in}{1.422345in}}{\pgfqpoint{1.486067in}{1.427931in}}{\pgfqpoint{1.481949in}{1.432049in}}%
\pgfpathcurveto{\pgfqpoint{1.477831in}{1.436167in}}{\pgfqpoint{1.472245in}{1.438481in}}{\pgfqpoint{1.466421in}{1.438481in}}%
\pgfpathcurveto{\pgfqpoint{1.460597in}{1.438481in}}{\pgfqpoint{1.455011in}{1.436167in}}{\pgfqpoint{1.450892in}{1.432049in}}%
\pgfpathcurveto{\pgfqpoint{1.446774in}{1.427931in}}{\pgfqpoint{1.444460in}{1.422345in}}{\pgfqpoint{1.444460in}{1.416521in}}%
\pgfpathcurveto{\pgfqpoint{1.444460in}{1.410697in}}{\pgfqpoint{1.446774in}{1.405111in}}{\pgfqpoint{1.450892in}{1.400993in}}%
\pgfpathcurveto{\pgfqpoint{1.455011in}{1.396875in}}{\pgfqpoint{1.460597in}{1.394561in}}{\pgfqpoint{1.466421in}{1.394561in}}%
\pgfpathlineto{\pgfqpoint{1.466421in}{1.394561in}}%
\pgfpathclose%
\pgfusepath{stroke,fill}%
\end{pgfscope}%
\begin{pgfscope}%
\pgfpathrectangle{\pgfqpoint{0.100000in}{0.209042in}}{\pgfqpoint{2.906250in}{2.906250in}}%
\pgfusepath{clip}%
\pgfsetbuttcap%
\pgfsetroundjoin%
\definecolor{currentfill}{rgb}{0.531507,0.154739,0.506895}%
\pgfsetfillcolor{currentfill}%
\pgfsetfillopacity{0.800000}%
\pgfsetlinewidth{1.003750pt}%
\definecolor{currentstroke}{rgb}{0.531507,0.154739,0.506895}%
\pgfsetstrokecolor{currentstroke}%
\pgfsetstrokeopacity{0.800000}%
\pgfsetdash{}{0pt}%
\pgfpathmoveto{\pgfqpoint{1.394527in}{1.566863in}}%
\pgfpathcurveto{\pgfqpoint{1.400350in}{1.566863in}}{\pgfqpoint{1.405937in}{1.569177in}}{\pgfqpoint{1.410055in}{1.573295in}}%
\pgfpathcurveto{\pgfqpoint{1.414173in}{1.577413in}}{\pgfqpoint{1.416487in}{1.583000in}}{\pgfqpoint{1.416487in}{1.588823in}}%
\pgfpathcurveto{\pgfqpoint{1.416487in}{1.594647in}}{\pgfqpoint{1.414173in}{1.600234in}}{\pgfqpoint{1.410055in}{1.604352in}}%
\pgfpathcurveto{\pgfqpoint{1.405937in}{1.608470in}}{\pgfqpoint{1.400350in}{1.610784in}}{\pgfqpoint{1.394527in}{1.610784in}}%
\pgfpathcurveto{\pgfqpoint{1.388703in}{1.610784in}}{\pgfqpoint{1.383116in}{1.608470in}}{\pgfqpoint{1.378998in}{1.604352in}}%
\pgfpathcurveto{\pgfqpoint{1.374880in}{1.600234in}}{\pgfqpoint{1.372566in}{1.594647in}}{\pgfqpoint{1.372566in}{1.588823in}}%
\pgfpathcurveto{\pgfqpoint{1.372566in}{1.583000in}}{\pgfqpoint{1.374880in}{1.577413in}}{\pgfqpoint{1.378998in}{1.573295in}}%
\pgfpathcurveto{\pgfqpoint{1.383116in}{1.569177in}}{\pgfqpoint{1.388703in}{1.566863in}}{\pgfqpoint{1.394527in}{1.566863in}}%
\pgfpathlineto{\pgfqpoint{1.394527in}{1.566863in}}%
\pgfpathclose%
\pgfusepath{stroke,fill}%
\end{pgfscope}%
\begin{pgfscope}%
\pgfpathrectangle{\pgfqpoint{0.100000in}{0.209042in}}{\pgfqpoint{2.906250in}{2.906250in}}%
\pgfusepath{clip}%
\pgfsetbuttcap%
\pgfsetroundjoin%
\definecolor{currentfill}{rgb}{0.822926,0.259016,0.433573}%
\pgfsetfillcolor{currentfill}%
\pgfsetfillopacity{0.800000}%
\pgfsetlinewidth{1.003750pt}%
\definecolor{currentstroke}{rgb}{0.822926,0.259016,0.433573}%
\pgfsetstrokecolor{currentstroke}%
\pgfsetstrokeopacity{0.800000}%
\pgfsetdash{}{0pt}%
\pgfpathmoveto{\pgfqpoint{1.685304in}{2.106695in}}%
\pgfpathcurveto{\pgfqpoint{1.691128in}{2.106695in}}{\pgfqpoint{1.696714in}{2.109008in}}{\pgfqpoint{1.700832in}{2.113127in}}%
\pgfpathcurveto{\pgfqpoint{1.704950in}{2.117245in}}{\pgfqpoint{1.707264in}{2.122831in}}{\pgfqpoint{1.707264in}{2.128655in}}%
\pgfpathcurveto{\pgfqpoint{1.707264in}{2.134479in}}{\pgfqpoint{1.704950in}{2.140065in}}{\pgfqpoint{1.700832in}{2.144183in}}%
\pgfpathcurveto{\pgfqpoint{1.696714in}{2.148301in}}{\pgfqpoint{1.691128in}{2.150615in}}{\pgfqpoint{1.685304in}{2.150615in}}%
\pgfpathcurveto{\pgfqpoint{1.679480in}{2.150615in}}{\pgfqpoint{1.673894in}{2.148301in}}{\pgfqpoint{1.669776in}{2.144183in}}%
\pgfpathcurveto{\pgfqpoint{1.665658in}{2.140065in}}{\pgfqpoint{1.663344in}{2.134479in}}{\pgfqpoint{1.663344in}{2.128655in}}%
\pgfpathcurveto{\pgfqpoint{1.663344in}{2.122831in}}{\pgfqpoint{1.665658in}{2.117245in}}{\pgfqpoint{1.669776in}{2.113127in}}%
\pgfpathcurveto{\pgfqpoint{1.673894in}{2.109008in}}{\pgfqpoint{1.679480in}{2.106695in}}{\pgfqpoint{1.685304in}{2.106695in}}%
\pgfpathlineto{\pgfqpoint{1.685304in}{2.106695in}}%
\pgfpathclose%
\pgfusepath{stroke,fill}%
\end{pgfscope}%
\begin{pgfscope}%
\pgfpathrectangle{\pgfqpoint{0.100000in}{0.209042in}}{\pgfqpoint{2.906250in}{2.906250in}}%
\pgfusepath{clip}%
\pgfsetbuttcap%
\pgfsetroundjoin%
\definecolor{currentfill}{rgb}{0.445163,0.122724,0.506901}%
\pgfsetfillcolor{currentfill}%
\pgfsetfillopacity{0.800000}%
\pgfsetlinewidth{1.003750pt}%
\definecolor{currentstroke}{rgb}{0.445163,0.122724,0.506901}%
\pgfsetstrokecolor{currentstroke}%
\pgfsetstrokeopacity{0.800000}%
\pgfsetdash{}{0pt}%
\pgfpathmoveto{\pgfqpoint{1.606974in}{1.810748in}}%
\pgfpathcurveto{\pgfqpoint{1.612798in}{1.810748in}}{\pgfqpoint{1.618384in}{1.813062in}}{\pgfqpoint{1.622502in}{1.817180in}}%
\pgfpathcurveto{\pgfqpoint{1.626620in}{1.821299in}}{\pgfqpoint{1.628934in}{1.826885in}}{\pgfqpoint{1.628934in}{1.832709in}}%
\pgfpathcurveto{\pgfqpoint{1.628934in}{1.838533in}}{\pgfqpoint{1.626620in}{1.844119in}}{\pgfqpoint{1.622502in}{1.848237in}}%
\pgfpathcurveto{\pgfqpoint{1.618384in}{1.852355in}}{\pgfqpoint{1.612798in}{1.854669in}}{\pgfqpoint{1.606974in}{1.854669in}}%
\pgfpathcurveto{\pgfqpoint{1.601150in}{1.854669in}}{\pgfqpoint{1.595563in}{1.852355in}}{\pgfqpoint{1.591445in}{1.848237in}}%
\pgfpathcurveto{\pgfqpoint{1.587327in}{1.844119in}}{\pgfqpoint{1.585013in}{1.838533in}}{\pgfqpoint{1.585013in}{1.832709in}}%
\pgfpathcurveto{\pgfqpoint{1.585013in}{1.826885in}}{\pgfqpoint{1.587327in}{1.821299in}}{\pgfqpoint{1.591445in}{1.817180in}}%
\pgfpathcurveto{\pgfqpoint{1.595563in}{1.813062in}}{\pgfqpoint{1.601150in}{1.810748in}}{\pgfqpoint{1.606974in}{1.810748in}}%
\pgfpathlineto{\pgfqpoint{1.606974in}{1.810748in}}%
\pgfpathclose%
\pgfusepath{stroke,fill}%
\end{pgfscope}%
\begin{pgfscope}%
\pgfpathrectangle{\pgfqpoint{0.100000in}{0.209042in}}{\pgfqpoint{2.906250in}{2.906250in}}%
\pgfusepath{clip}%
\pgfsetbuttcap%
\pgfsetroundjoin%
\definecolor{currentfill}{rgb}{0.575490,0.169530,0.503466}%
\pgfsetfillcolor{currentfill}%
\pgfsetfillopacity{0.800000}%
\pgfsetlinewidth{1.003750pt}%
\definecolor{currentstroke}{rgb}{0.575490,0.169530,0.503466}%
\pgfsetstrokecolor{currentstroke}%
\pgfsetstrokeopacity{0.800000}%
\pgfsetdash{}{0pt}%
\pgfpathmoveto{\pgfqpoint{1.948497in}{1.794932in}}%
\pgfpathcurveto{\pgfqpoint{1.954320in}{1.794932in}}{\pgfqpoint{1.959907in}{1.797246in}}{\pgfqpoint{1.964025in}{1.801364in}}%
\pgfpathcurveto{\pgfqpoint{1.968143in}{1.805482in}}{\pgfqpoint{1.970457in}{1.811068in}}{\pgfqpoint{1.970457in}{1.816892in}}%
\pgfpathcurveto{\pgfqpoint{1.970457in}{1.822716in}}{\pgfqpoint{1.968143in}{1.828302in}}{\pgfqpoint{1.964025in}{1.832421in}}%
\pgfpathcurveto{\pgfqpoint{1.959907in}{1.836539in}}{\pgfqpoint{1.954320in}{1.838853in}}{\pgfqpoint{1.948497in}{1.838853in}}%
\pgfpathcurveto{\pgfqpoint{1.942673in}{1.838853in}}{\pgfqpoint{1.937086in}{1.836539in}}{\pgfqpoint{1.932968in}{1.832421in}}%
\pgfpathcurveto{\pgfqpoint{1.928850in}{1.828302in}}{\pgfqpoint{1.926536in}{1.822716in}}{\pgfqpoint{1.926536in}{1.816892in}}%
\pgfpathcurveto{\pgfqpoint{1.926536in}{1.811068in}}{\pgfqpoint{1.928850in}{1.805482in}}{\pgfqpoint{1.932968in}{1.801364in}}%
\pgfpathcurveto{\pgfqpoint{1.937086in}{1.797246in}}{\pgfqpoint{1.942673in}{1.794932in}}{\pgfqpoint{1.948497in}{1.794932in}}%
\pgfpathlineto{\pgfqpoint{1.948497in}{1.794932in}}%
\pgfpathclose%
\pgfusepath{stroke,fill}%
\end{pgfscope}%
\begin{pgfscope}%
\pgfpathrectangle{\pgfqpoint{0.100000in}{0.209042in}}{\pgfqpoint{2.906250in}{2.906250in}}%
\pgfusepath{clip}%
\pgfsetbuttcap%
\pgfsetroundjoin%
\definecolor{currentfill}{rgb}{0.426877,0.115395,0.505714}%
\pgfsetfillcolor{currentfill}%
\pgfsetfillopacity{0.800000}%
\pgfsetlinewidth{1.003750pt}%
\definecolor{currentstroke}{rgb}{0.426877,0.115395,0.505714}%
\pgfsetstrokecolor{currentstroke}%
\pgfsetstrokeopacity{0.800000}%
\pgfsetdash{}{0pt}%
\pgfpathmoveto{\pgfqpoint{1.719052in}{1.544607in}}%
\pgfpathcurveto{\pgfqpoint{1.724876in}{1.544607in}}{\pgfqpoint{1.730462in}{1.546921in}}{\pgfqpoint{1.734580in}{1.551039in}}%
\pgfpathcurveto{\pgfqpoint{1.738698in}{1.555158in}}{\pgfqpoint{1.741012in}{1.560744in}}{\pgfqpoint{1.741012in}{1.566568in}}%
\pgfpathcurveto{\pgfqpoint{1.741012in}{1.572392in}}{\pgfqpoint{1.738698in}{1.577978in}}{\pgfqpoint{1.734580in}{1.582096in}}%
\pgfpathcurveto{\pgfqpoint{1.730462in}{1.586214in}}{\pgfqpoint{1.724876in}{1.588528in}}{\pgfqpoint{1.719052in}{1.588528in}}%
\pgfpathcurveto{\pgfqpoint{1.713228in}{1.588528in}}{\pgfqpoint{1.707642in}{1.586214in}}{\pgfqpoint{1.703524in}{1.582096in}}%
\pgfpathcurveto{\pgfqpoint{1.699406in}{1.577978in}}{\pgfqpoint{1.697092in}{1.572392in}}{\pgfqpoint{1.697092in}{1.566568in}}%
\pgfpathcurveto{\pgfqpoint{1.697092in}{1.560744in}}{\pgfqpoint{1.699406in}{1.555158in}}{\pgfqpoint{1.703524in}{1.551039in}}%
\pgfpathcurveto{\pgfqpoint{1.707642in}{1.546921in}}{\pgfqpoint{1.713228in}{1.544607in}}{\pgfqpoint{1.719052in}{1.544607in}}%
\pgfpathlineto{\pgfqpoint{1.719052in}{1.544607in}}%
\pgfpathclose%
\pgfusepath{stroke,fill}%
\end{pgfscope}%
\begin{pgfscope}%
\pgfpathrectangle{\pgfqpoint{0.100000in}{0.209042in}}{\pgfqpoint{2.906250in}{2.906250in}}%
\pgfusepath{clip}%
\pgfsetbuttcap%
\pgfsetroundjoin%
\definecolor{currentfill}{rgb}{0.613617,0.181811,0.498536}%
\pgfsetfillcolor{currentfill}%
\pgfsetfillopacity{0.800000}%
\pgfsetlinewidth{1.003750pt}%
\definecolor{currentstroke}{rgb}{0.613617,0.181811,0.498536}%
\pgfsetstrokecolor{currentstroke}%
\pgfsetstrokeopacity{0.800000}%
\pgfsetdash{}{0pt}%
\pgfpathmoveto{\pgfqpoint{2.004658in}{1.764294in}}%
\pgfpathcurveto{\pgfqpoint{2.010482in}{1.764294in}}{\pgfqpoint{2.016068in}{1.766608in}}{\pgfqpoint{2.020187in}{1.770726in}}%
\pgfpathcurveto{\pgfqpoint{2.024305in}{1.774844in}}{\pgfqpoint{2.026619in}{1.780431in}}{\pgfqpoint{2.026619in}{1.786254in}}%
\pgfpathcurveto{\pgfqpoint{2.026619in}{1.792078in}}{\pgfqpoint{2.024305in}{1.797665in}}{\pgfqpoint{2.020187in}{1.801783in}}%
\pgfpathcurveto{\pgfqpoint{2.016068in}{1.805901in}}{\pgfqpoint{2.010482in}{1.808215in}}{\pgfqpoint{2.004658in}{1.808215in}}%
\pgfpathcurveto{\pgfqpoint{1.998834in}{1.808215in}}{\pgfqpoint{1.993248in}{1.805901in}}{\pgfqpoint{1.989130in}{1.801783in}}%
\pgfpathcurveto{\pgfqpoint{1.985012in}{1.797665in}}{\pgfqpoint{1.982698in}{1.792078in}}{\pgfqpoint{1.982698in}{1.786254in}}%
\pgfpathcurveto{\pgfqpoint{1.982698in}{1.780431in}}{\pgfqpoint{1.985012in}{1.774844in}}{\pgfqpoint{1.989130in}{1.770726in}}%
\pgfpathcurveto{\pgfqpoint{1.993248in}{1.766608in}}{\pgfqpoint{1.998834in}{1.764294in}}{\pgfqpoint{2.004658in}{1.764294in}}%
\pgfpathlineto{\pgfqpoint{2.004658in}{1.764294in}}%
\pgfpathclose%
\pgfusepath{stroke,fill}%
\end{pgfscope}%
\begin{pgfscope}%
\pgfpathrectangle{\pgfqpoint{0.100000in}{0.209042in}}{\pgfqpoint{2.906250in}{2.906250in}}%
\pgfusepath{clip}%
\pgfsetbuttcap%
\pgfsetroundjoin%
\definecolor{currentfill}{rgb}{0.569172,0.167454,0.504105}%
\pgfsetfillcolor{currentfill}%
\pgfsetfillopacity{0.800000}%
\pgfsetlinewidth{1.003750pt}%
\definecolor{currentstroke}{rgb}{0.569172,0.167454,0.504105}%
\pgfsetstrokecolor{currentstroke}%
\pgfsetstrokeopacity{0.800000}%
\pgfsetdash{}{0pt}%
\pgfpathmoveto{\pgfqpoint{1.840069in}{1.432771in}}%
\pgfpathcurveto{\pgfqpoint{1.845893in}{1.432771in}}{\pgfqpoint{1.851479in}{1.435085in}}{\pgfqpoint{1.855598in}{1.439203in}}%
\pgfpathcurveto{\pgfqpoint{1.859716in}{1.443321in}}{\pgfqpoint{1.862030in}{1.448907in}}{\pgfqpoint{1.862030in}{1.454731in}}%
\pgfpathcurveto{\pgfqpoint{1.862030in}{1.460555in}}{\pgfqpoint{1.859716in}{1.466141in}}{\pgfqpoint{1.855598in}{1.470259in}}%
\pgfpathcurveto{\pgfqpoint{1.851479in}{1.474377in}}{\pgfqpoint{1.845893in}{1.476691in}}{\pgfqpoint{1.840069in}{1.476691in}}%
\pgfpathcurveto{\pgfqpoint{1.834245in}{1.476691in}}{\pgfqpoint{1.828659in}{1.474377in}}{\pgfqpoint{1.824541in}{1.470259in}}%
\pgfpathcurveto{\pgfqpoint{1.820423in}{1.466141in}}{\pgfqpoint{1.818109in}{1.460555in}}{\pgfqpoint{1.818109in}{1.454731in}}%
\pgfpathcurveto{\pgfqpoint{1.818109in}{1.448907in}}{\pgfqpoint{1.820423in}{1.443321in}}{\pgfqpoint{1.824541in}{1.439203in}}%
\pgfpathcurveto{\pgfqpoint{1.828659in}{1.435085in}}{\pgfqpoint{1.834245in}{1.432771in}}{\pgfqpoint{1.840069in}{1.432771in}}%
\pgfpathlineto{\pgfqpoint{1.840069in}{1.432771in}}%
\pgfpathclose%
\pgfusepath{stroke,fill}%
\end{pgfscope}%
\begin{pgfscope}%
\pgfpathrectangle{\pgfqpoint{0.100000in}{0.209042in}}{\pgfqpoint{2.906250in}{2.906250in}}%
\pgfusepath{clip}%
\pgfsetbuttcap%
\pgfsetroundjoin%
\definecolor{currentfill}{rgb}{0.426877,0.115395,0.505714}%
\pgfsetfillcolor{currentfill}%
\pgfsetfillopacity{0.800000}%
\pgfsetlinewidth{1.003750pt}%
\definecolor{currentstroke}{rgb}{0.426877,0.115395,0.505714}%
\pgfsetstrokecolor{currentstroke}%
\pgfsetstrokeopacity{0.800000}%
\pgfsetdash{}{0pt}%
\pgfpathmoveto{\pgfqpoint{1.588623in}{1.600820in}}%
\pgfpathcurveto{\pgfqpoint{1.594447in}{1.600820in}}{\pgfqpoint{1.600034in}{1.603134in}}{\pgfqpoint{1.604152in}{1.607252in}}%
\pgfpathcurveto{\pgfqpoint{1.608270in}{1.611370in}}{\pgfqpoint{1.610584in}{1.616956in}}{\pgfqpoint{1.610584in}{1.622780in}}%
\pgfpathcurveto{\pgfqpoint{1.610584in}{1.628604in}}{\pgfqpoint{1.608270in}{1.634191in}}{\pgfqpoint{1.604152in}{1.638309in}}%
\pgfpathcurveto{\pgfqpoint{1.600034in}{1.642427in}}{\pgfqpoint{1.594447in}{1.644741in}}{\pgfqpoint{1.588623in}{1.644741in}}%
\pgfpathcurveto{\pgfqpoint{1.582800in}{1.644741in}}{\pgfqpoint{1.577213in}{1.642427in}}{\pgfqpoint{1.573095in}{1.638309in}}%
\pgfpathcurveto{\pgfqpoint{1.568977in}{1.634191in}}{\pgfqpoint{1.566663in}{1.628604in}}{\pgfqpoint{1.566663in}{1.622780in}}%
\pgfpathcurveto{\pgfqpoint{1.566663in}{1.616956in}}{\pgfqpoint{1.568977in}{1.611370in}}{\pgfqpoint{1.573095in}{1.607252in}}%
\pgfpathcurveto{\pgfqpoint{1.577213in}{1.603134in}}{\pgfqpoint{1.582800in}{1.600820in}}{\pgfqpoint{1.588623in}{1.600820in}}%
\pgfpathlineto{\pgfqpoint{1.588623in}{1.600820in}}%
\pgfpathclose%
\pgfusepath{stroke,fill}%
\end{pgfscope}%
\begin{pgfscope}%
\pgfpathrectangle{\pgfqpoint{0.100000in}{0.209042in}}{\pgfqpoint{2.906250in}{2.906250in}}%
\pgfusepath{clip}%
\pgfsetbuttcap%
\pgfsetroundjoin%
\definecolor{currentfill}{rgb}{0.955849,0.404400,0.360619}%
\pgfsetfillcolor{currentfill}%
\pgfsetfillopacity{0.800000}%
\pgfsetlinewidth{1.003750pt}%
\definecolor{currentstroke}{rgb}{0.955849,0.404400,0.360619}%
\pgfsetstrokecolor{currentstroke}%
\pgfsetstrokeopacity{0.800000}%
\pgfsetdash{}{0pt}%
\pgfpathmoveto{\pgfqpoint{2.070823in}{1.199267in}}%
\pgfpathcurveto{\pgfqpoint{2.076647in}{1.199267in}}{\pgfqpoint{2.082233in}{1.201581in}}{\pgfqpoint{2.086351in}{1.205699in}}%
\pgfpathcurveto{\pgfqpoint{2.090469in}{1.209817in}}{\pgfqpoint{2.092783in}{1.215403in}}{\pgfqpoint{2.092783in}{1.221227in}}%
\pgfpathcurveto{\pgfqpoint{2.092783in}{1.227051in}}{\pgfqpoint{2.090469in}{1.232637in}}{\pgfqpoint{2.086351in}{1.236755in}}%
\pgfpathcurveto{\pgfqpoint{2.082233in}{1.240873in}}{\pgfqpoint{2.076647in}{1.243187in}}{\pgfqpoint{2.070823in}{1.243187in}}%
\pgfpathcurveto{\pgfqpoint{2.064999in}{1.243187in}}{\pgfqpoint{2.059413in}{1.240873in}}{\pgfqpoint{2.055295in}{1.236755in}}%
\pgfpathcurveto{\pgfqpoint{2.051177in}{1.232637in}}{\pgfqpoint{2.048863in}{1.227051in}}{\pgfqpoint{2.048863in}{1.221227in}}%
\pgfpathcurveto{\pgfqpoint{2.048863in}{1.215403in}}{\pgfqpoint{2.051177in}{1.209817in}}{\pgfqpoint{2.055295in}{1.205699in}}%
\pgfpathcurveto{\pgfqpoint{2.059413in}{1.201581in}}{\pgfqpoint{2.064999in}{1.199267in}}{\pgfqpoint{2.070823in}{1.199267in}}%
\pgfpathlineto{\pgfqpoint{2.070823in}{1.199267in}}%
\pgfpathclose%
\pgfusepath{stroke,fill}%
\end{pgfscope}%
\begin{pgfscope}%
\pgfpathrectangle{\pgfqpoint{0.100000in}{0.209042in}}{\pgfqpoint{2.906250in}{2.906250in}}%
\pgfusepath{clip}%
\pgfsetbuttcap%
\pgfsetroundjoin%
\definecolor{currentfill}{rgb}{0.463508,0.129893,0.507652}%
\pgfsetfillcolor{currentfill}%
\pgfsetfillopacity{0.800000}%
\pgfsetlinewidth{1.003750pt}%
\definecolor{currentstroke}{rgb}{0.463508,0.129893,0.507652}%
\pgfsetstrokecolor{currentstroke}%
\pgfsetstrokeopacity{0.800000}%
\pgfsetdash{}{0pt}%
\pgfpathmoveto{\pgfqpoint{1.581982in}{1.784647in}}%
\pgfpathcurveto{\pgfqpoint{1.587806in}{1.784647in}}{\pgfqpoint{1.593393in}{1.786961in}}{\pgfqpoint{1.597511in}{1.791079in}}%
\pgfpathcurveto{\pgfqpoint{1.601629in}{1.795198in}}{\pgfqpoint{1.603943in}{1.800784in}}{\pgfqpoint{1.603943in}{1.806608in}}%
\pgfpathcurveto{\pgfqpoint{1.603943in}{1.812432in}}{\pgfqpoint{1.601629in}{1.818018in}}{\pgfqpoint{1.597511in}{1.822136in}}%
\pgfpathcurveto{\pgfqpoint{1.593393in}{1.826254in}}{\pgfqpoint{1.587806in}{1.828568in}}{\pgfqpoint{1.581982in}{1.828568in}}%
\pgfpathcurveto{\pgfqpoint{1.576159in}{1.828568in}}{\pgfqpoint{1.570572in}{1.826254in}}{\pgfqpoint{1.566454in}{1.822136in}}%
\pgfpathcurveto{\pgfqpoint{1.562336in}{1.818018in}}{\pgfqpoint{1.560022in}{1.812432in}}{\pgfqpoint{1.560022in}{1.806608in}}%
\pgfpathcurveto{\pgfqpoint{1.560022in}{1.800784in}}{\pgfqpoint{1.562336in}{1.795198in}}{\pgfqpoint{1.566454in}{1.791079in}}%
\pgfpathcurveto{\pgfqpoint{1.570572in}{1.786961in}}{\pgfqpoint{1.576159in}{1.784647in}}{\pgfqpoint{1.581982in}{1.784647in}}%
\pgfpathlineto{\pgfqpoint{1.581982in}{1.784647in}}%
\pgfpathclose%
\pgfusepath{stroke,fill}%
\end{pgfscope}%
\begin{pgfscope}%
\pgfpathrectangle{\pgfqpoint{0.100000in}{0.209042in}}{\pgfqpoint{2.906250in}{2.906250in}}%
\pgfusepath{clip}%
\pgfsetbuttcap%
\pgfsetroundjoin%
\definecolor{currentfill}{rgb}{0.562866,0.165368,0.504692}%
\pgfsetfillcolor{currentfill}%
\pgfsetfillopacity{0.800000}%
\pgfsetlinewidth{1.003750pt}%
\definecolor{currentstroke}{rgb}{0.562866,0.165368,0.504692}%
\pgfsetstrokecolor{currentstroke}%
\pgfsetstrokeopacity{0.800000}%
\pgfsetdash{}{0pt}%
\pgfpathmoveto{\pgfqpoint{1.432823in}{1.811669in}}%
\pgfpathcurveto{\pgfqpoint{1.438647in}{1.811669in}}{\pgfqpoint{1.444234in}{1.813983in}}{\pgfqpoint{1.448352in}{1.818101in}}%
\pgfpathcurveto{\pgfqpoint{1.452470in}{1.822219in}}{\pgfqpoint{1.454784in}{1.827805in}}{\pgfqpoint{1.454784in}{1.833629in}}%
\pgfpathcurveto{\pgfqpoint{1.454784in}{1.839453in}}{\pgfqpoint{1.452470in}{1.845039in}}{\pgfqpoint{1.448352in}{1.849157in}}%
\pgfpathcurveto{\pgfqpoint{1.444234in}{1.853275in}}{\pgfqpoint{1.438647in}{1.855589in}}{\pgfqpoint{1.432823in}{1.855589in}}%
\pgfpathcurveto{\pgfqpoint{1.427000in}{1.855589in}}{\pgfqpoint{1.421413in}{1.853275in}}{\pgfqpoint{1.417295in}{1.849157in}}%
\pgfpathcurveto{\pgfqpoint{1.413177in}{1.845039in}}{\pgfqpoint{1.410863in}{1.839453in}}{\pgfqpoint{1.410863in}{1.833629in}}%
\pgfpathcurveto{\pgfqpoint{1.410863in}{1.827805in}}{\pgfqpoint{1.413177in}{1.822219in}}{\pgfqpoint{1.417295in}{1.818101in}}%
\pgfpathcurveto{\pgfqpoint{1.421413in}{1.813983in}}{\pgfqpoint{1.427000in}{1.811669in}}{\pgfqpoint{1.432823in}{1.811669in}}%
\pgfpathlineto{\pgfqpoint{1.432823in}{1.811669in}}%
\pgfpathclose%
\pgfusepath{stroke,fill}%
\end{pgfscope}%
\begin{pgfscope}%
\pgfpathrectangle{\pgfqpoint{0.100000in}{0.209042in}}{\pgfqpoint{2.906250in}{2.906250in}}%
\pgfusepath{clip}%
\pgfsetbuttcap%
\pgfsetroundjoin%
\definecolor{currentfill}{rgb}{0.613617,0.181811,0.498536}%
\pgfsetfillcolor{currentfill}%
\pgfsetfillopacity{0.800000}%
\pgfsetlinewidth{1.003750pt}%
\definecolor{currentstroke}{rgb}{0.613617,0.181811,0.498536}%
\pgfsetstrokecolor{currentstroke}%
\pgfsetstrokeopacity{0.800000}%
\pgfsetdash{}{0pt}%
\pgfpathmoveto{\pgfqpoint{1.377968in}{1.522927in}}%
\pgfpathcurveto{\pgfqpoint{1.383792in}{1.522927in}}{\pgfqpoint{1.389378in}{1.525241in}}{\pgfqpoint{1.393496in}{1.529359in}}%
\pgfpathcurveto{\pgfqpoint{1.397614in}{1.533477in}}{\pgfqpoint{1.399928in}{1.539063in}}{\pgfqpoint{1.399928in}{1.544887in}}%
\pgfpathcurveto{\pgfqpoint{1.399928in}{1.550711in}}{\pgfqpoint{1.397614in}{1.556297in}}{\pgfqpoint{1.393496in}{1.560416in}}%
\pgfpathcurveto{\pgfqpoint{1.389378in}{1.564534in}}{\pgfqpoint{1.383792in}{1.566848in}}{\pgfqpoint{1.377968in}{1.566848in}}%
\pgfpathcurveto{\pgfqpoint{1.372144in}{1.566848in}}{\pgfqpoint{1.366558in}{1.564534in}}{\pgfqpoint{1.362440in}{1.560416in}}%
\pgfpathcurveto{\pgfqpoint{1.358322in}{1.556297in}}{\pgfqpoint{1.356008in}{1.550711in}}{\pgfqpoint{1.356008in}{1.544887in}}%
\pgfpathcurveto{\pgfqpoint{1.356008in}{1.539063in}}{\pgfqpoint{1.358322in}{1.533477in}}{\pgfqpoint{1.362440in}{1.529359in}}%
\pgfpathcurveto{\pgfqpoint{1.366558in}{1.525241in}}{\pgfqpoint{1.372144in}{1.522927in}}{\pgfqpoint{1.377968in}{1.522927in}}%
\pgfpathlineto{\pgfqpoint{1.377968in}{1.522927in}}%
\pgfpathclose%
\pgfusepath{stroke,fill}%
\end{pgfscope}%
\begin{pgfscope}%
\pgfpathrectangle{\pgfqpoint{0.100000in}{0.209042in}}{\pgfqpoint{2.906250in}{2.906250in}}%
\pgfusepath{clip}%
\pgfsetbuttcap%
\pgfsetroundjoin%
\definecolor{currentfill}{rgb}{0.690661,0.206384,0.482558}%
\pgfsetfillcolor{currentfill}%
\pgfsetfillopacity{0.800000}%
\pgfsetlinewidth{1.003750pt}%
\definecolor{currentstroke}{rgb}{0.690661,0.206384,0.482558}%
\pgfsetstrokecolor{currentstroke}%
\pgfsetstrokeopacity{0.800000}%
\pgfsetdash{}{0pt}%
\pgfpathmoveto{\pgfqpoint{2.073152in}{1.730669in}}%
\pgfpathcurveto{\pgfqpoint{2.078976in}{1.730669in}}{\pgfqpoint{2.084562in}{1.732982in}}{\pgfqpoint{2.088680in}{1.737101in}}%
\pgfpathcurveto{\pgfqpoint{2.092799in}{1.741219in}}{\pgfqpoint{2.095112in}{1.746805in}}{\pgfqpoint{2.095112in}{1.752629in}}%
\pgfpathcurveto{\pgfqpoint{2.095112in}{1.758453in}}{\pgfqpoint{2.092799in}{1.764039in}}{\pgfqpoint{2.088680in}{1.768157in}}%
\pgfpathcurveto{\pgfqpoint{2.084562in}{1.772275in}}{\pgfqpoint{2.078976in}{1.774589in}}{\pgfqpoint{2.073152in}{1.774589in}}%
\pgfpathcurveto{\pgfqpoint{2.067328in}{1.774589in}}{\pgfqpoint{2.061742in}{1.772275in}}{\pgfqpoint{2.057624in}{1.768157in}}%
\pgfpathcurveto{\pgfqpoint{2.053506in}{1.764039in}}{\pgfqpoint{2.051192in}{1.758453in}}{\pgfqpoint{2.051192in}{1.752629in}}%
\pgfpathcurveto{\pgfqpoint{2.051192in}{1.746805in}}{\pgfqpoint{2.053506in}{1.741219in}}{\pgfqpoint{2.057624in}{1.737101in}}%
\pgfpathcurveto{\pgfqpoint{2.061742in}{1.732982in}}{\pgfqpoint{2.067328in}{1.730669in}}{\pgfqpoint{2.073152in}{1.730669in}}%
\pgfpathlineto{\pgfqpoint{2.073152in}{1.730669in}}%
\pgfpathclose%
\pgfusepath{stroke,fill}%
\end{pgfscope}%
\begin{pgfscope}%
\pgfpathrectangle{\pgfqpoint{0.100000in}{0.209042in}}{\pgfqpoint{2.906250in}{2.906250in}}%
\pgfusepath{clip}%
\pgfsetbuttcap%
\pgfsetroundjoin%
\definecolor{currentfill}{rgb}{0.767398,0.233705,0.457755}%
\pgfsetfillcolor{currentfill}%
\pgfsetfillopacity{0.800000}%
\pgfsetlinewidth{1.003750pt}%
\definecolor{currentstroke}{rgb}{0.767398,0.233705,0.457755}%
\pgfsetstrokecolor{currentstroke}%
\pgfsetstrokeopacity{0.800000}%
\pgfsetdash{}{0pt}%
\pgfpathmoveto{\pgfqpoint{1.349113in}{1.953669in}}%
\pgfpathcurveto{\pgfqpoint{1.354937in}{1.953669in}}{\pgfqpoint{1.360523in}{1.955983in}}{\pgfqpoint{1.364641in}{1.960101in}}%
\pgfpathcurveto{\pgfqpoint{1.368760in}{1.964219in}}{\pgfqpoint{1.371073in}{1.969805in}}{\pgfqpoint{1.371073in}{1.975629in}}%
\pgfpathcurveto{\pgfqpoint{1.371073in}{1.981453in}}{\pgfqpoint{1.368760in}{1.987039in}}{\pgfqpoint{1.364641in}{1.991157in}}%
\pgfpathcurveto{\pgfqpoint{1.360523in}{1.995276in}}{\pgfqpoint{1.354937in}{1.997589in}}{\pgfqpoint{1.349113in}{1.997589in}}%
\pgfpathcurveto{\pgfqpoint{1.343289in}{1.997589in}}{\pgfqpoint{1.337703in}{1.995276in}}{\pgfqpoint{1.333585in}{1.991157in}}%
\pgfpathcurveto{\pgfqpoint{1.329467in}{1.987039in}}{\pgfqpoint{1.327153in}{1.981453in}}{\pgfqpoint{1.327153in}{1.975629in}}%
\pgfpathcurveto{\pgfqpoint{1.327153in}{1.969805in}}{\pgfqpoint{1.329467in}{1.964219in}}{\pgfqpoint{1.333585in}{1.960101in}}%
\pgfpathcurveto{\pgfqpoint{1.337703in}{1.955983in}}{\pgfqpoint{1.343289in}{1.953669in}}{\pgfqpoint{1.349113in}{1.953669in}}%
\pgfpathlineto{\pgfqpoint{1.349113in}{1.953669in}}%
\pgfpathclose%
\pgfusepath{stroke,fill}%
\end{pgfscope}%
\begin{pgfscope}%
\pgfpathrectangle{\pgfqpoint{0.100000in}{0.209042in}}{\pgfqpoint{2.906250in}{2.906250in}}%
\pgfusepath{clip}%
\pgfsetbuttcap%
\pgfsetroundjoin%
\definecolor{currentfill}{rgb}{0.697098,0.208501,0.480835}%
\pgfsetfillcolor{currentfill}%
\pgfsetfillopacity{0.800000}%
\pgfsetlinewidth{1.003750pt}%
\definecolor{currentstroke}{rgb}{0.697098,0.208501,0.480835}%
\pgfsetstrokecolor{currentstroke}%
\pgfsetstrokeopacity{0.800000}%
\pgfsetdash{}{0pt}%
\pgfpathmoveto{\pgfqpoint{1.689229in}{1.992961in}}%
\pgfpathcurveto{\pgfqpoint{1.695053in}{1.992961in}}{\pgfqpoint{1.700639in}{1.995275in}}{\pgfqpoint{1.704757in}{1.999393in}}%
\pgfpathcurveto{\pgfqpoint{1.708875in}{2.003511in}}{\pgfqpoint{1.711189in}{2.009097in}}{\pgfqpoint{1.711189in}{2.014921in}}%
\pgfpathcurveto{\pgfqpoint{1.711189in}{2.020745in}}{\pgfqpoint{1.708875in}{2.026331in}}{\pgfqpoint{1.704757in}{2.030449in}}%
\pgfpathcurveto{\pgfqpoint{1.700639in}{2.034567in}}{\pgfqpoint{1.695053in}{2.036881in}}{\pgfqpoint{1.689229in}{2.036881in}}%
\pgfpathcurveto{\pgfqpoint{1.683405in}{2.036881in}}{\pgfqpoint{1.677819in}{2.034567in}}{\pgfqpoint{1.673701in}{2.030449in}}%
\pgfpathcurveto{\pgfqpoint{1.669583in}{2.026331in}}{\pgfqpoint{1.667269in}{2.020745in}}{\pgfqpoint{1.667269in}{2.014921in}}%
\pgfpathcurveto{\pgfqpoint{1.667269in}{2.009097in}}{\pgfqpoint{1.669583in}{2.003511in}}{\pgfqpoint{1.673701in}{1.999393in}}%
\pgfpathcurveto{\pgfqpoint{1.677819in}{1.995275in}}{\pgfqpoint{1.683405in}{1.992961in}}{\pgfqpoint{1.689229in}{1.992961in}}%
\pgfpathlineto{\pgfqpoint{1.689229in}{1.992961in}}%
\pgfpathclose%
\pgfusepath{stroke,fill}%
\end{pgfscope}%
\begin{pgfscope}%
\pgfpathrectangle{\pgfqpoint{0.100000in}{0.209042in}}{\pgfqpoint{2.906250in}{2.906250in}}%
\pgfusepath{clip}%
\pgfsetbuttcap%
\pgfsetroundjoin%
\definecolor{currentfill}{rgb}{0.494258,0.141462,0.507988}%
\pgfsetfillcolor{currentfill}%
\pgfsetfillopacity{0.800000}%
\pgfsetlinewidth{1.003750pt}%
\definecolor{currentstroke}{rgb}{0.494258,0.141462,0.507988}%
\pgfsetstrokecolor{currentstroke}%
\pgfsetstrokeopacity{0.800000}%
\pgfsetdash{}{0pt}%
\pgfpathmoveto{\pgfqpoint{1.544542in}{1.558435in}}%
\pgfpathcurveto{\pgfqpoint{1.550366in}{1.558435in}}{\pgfqpoint{1.555952in}{1.560749in}}{\pgfqpoint{1.560070in}{1.564867in}}%
\pgfpathcurveto{\pgfqpoint{1.564188in}{1.568985in}}{\pgfqpoint{1.566502in}{1.574571in}}{\pgfqpoint{1.566502in}{1.580395in}}%
\pgfpathcurveto{\pgfqpoint{1.566502in}{1.586219in}}{\pgfqpoint{1.564188in}{1.591805in}}{\pgfqpoint{1.560070in}{1.595923in}}%
\pgfpathcurveto{\pgfqpoint{1.555952in}{1.600042in}}{\pgfqpoint{1.550366in}{1.602355in}}{\pgfqpoint{1.544542in}{1.602355in}}%
\pgfpathcurveto{\pgfqpoint{1.538718in}{1.602355in}}{\pgfqpoint{1.533132in}{1.600042in}}{\pgfqpoint{1.529014in}{1.595923in}}%
\pgfpathcurveto{\pgfqpoint{1.524896in}{1.591805in}}{\pgfqpoint{1.522582in}{1.586219in}}{\pgfqpoint{1.522582in}{1.580395in}}%
\pgfpathcurveto{\pgfqpoint{1.522582in}{1.574571in}}{\pgfqpoint{1.524896in}{1.568985in}}{\pgfqpoint{1.529014in}{1.564867in}}%
\pgfpathcurveto{\pgfqpoint{1.533132in}{1.560749in}}{\pgfqpoint{1.538718in}{1.558435in}}{\pgfqpoint{1.544542in}{1.558435in}}%
\pgfpathlineto{\pgfqpoint{1.544542in}{1.558435in}}%
\pgfpathclose%
\pgfusepath{stroke,fill}%
\end{pgfscope}%
\begin{pgfscope}%
\pgfpathrectangle{\pgfqpoint{0.100000in}{0.209042in}}{\pgfqpoint{2.906250in}{2.906250in}}%
\pgfusepath{clip}%
\pgfsetbuttcap%
\pgfsetroundjoin%
\definecolor{currentfill}{rgb}{0.512831,0.148179,0.507648}%
\pgfsetfillcolor{currentfill}%
\pgfsetfillopacity{0.800000}%
\pgfsetlinewidth{1.003750pt}%
\definecolor{currentstroke}{rgb}{0.512831,0.148179,0.507648}%
\pgfsetstrokecolor{currentstroke}%
\pgfsetstrokeopacity{0.800000}%
\pgfsetdash{}{0pt}%
\pgfpathmoveto{\pgfqpoint{1.478420in}{1.607841in}}%
\pgfpathcurveto{\pgfqpoint{1.484244in}{1.607841in}}{\pgfqpoint{1.489830in}{1.610155in}}{\pgfqpoint{1.493948in}{1.614273in}}%
\pgfpathcurveto{\pgfqpoint{1.498066in}{1.618391in}}{\pgfqpoint{1.500380in}{1.623977in}}{\pgfqpoint{1.500380in}{1.629801in}}%
\pgfpathcurveto{\pgfqpoint{1.500380in}{1.635625in}}{\pgfqpoint{1.498066in}{1.641211in}}{\pgfqpoint{1.493948in}{1.645330in}}%
\pgfpathcurveto{\pgfqpoint{1.489830in}{1.649448in}}{\pgfqpoint{1.484244in}{1.651762in}}{\pgfqpoint{1.478420in}{1.651762in}}%
\pgfpathcurveto{\pgfqpoint{1.472596in}{1.651762in}}{\pgfqpoint{1.467010in}{1.649448in}}{\pgfqpoint{1.462892in}{1.645330in}}%
\pgfpathcurveto{\pgfqpoint{1.458774in}{1.641211in}}{\pgfqpoint{1.456460in}{1.635625in}}{\pgfqpoint{1.456460in}{1.629801in}}%
\pgfpathcurveto{\pgfqpoint{1.456460in}{1.623977in}}{\pgfqpoint{1.458774in}{1.618391in}}{\pgfqpoint{1.462892in}{1.614273in}}%
\pgfpathcurveto{\pgfqpoint{1.467010in}{1.610155in}}{\pgfqpoint{1.472596in}{1.607841in}}{\pgfqpoint{1.478420in}{1.607841in}}%
\pgfpathlineto{\pgfqpoint{1.478420in}{1.607841in}}%
\pgfpathclose%
\pgfusepath{stroke,fill}%
\end{pgfscope}%
\begin{pgfscope}%
\pgfpathrectangle{\pgfqpoint{0.100000in}{0.209042in}}{\pgfqpoint{2.906250in}{2.906250in}}%
\pgfusepath{clip}%
\pgfsetbuttcap%
\pgfsetroundjoin%
\definecolor{currentfill}{rgb}{0.697098,0.208501,0.480835}%
\pgfsetfillcolor{currentfill}%
\pgfsetfillopacity{0.800000}%
\pgfsetlinewidth{1.003750pt}%
\definecolor{currentstroke}{rgb}{0.697098,0.208501,0.480835}%
\pgfsetstrokecolor{currentstroke}%
\pgfsetstrokeopacity{0.800000}%
\pgfsetdash{}{0pt}%
\pgfpathmoveto{\pgfqpoint{1.824367in}{1.340099in}}%
\pgfpathcurveto{\pgfqpoint{1.830191in}{1.340099in}}{\pgfqpoint{1.835777in}{1.342412in}}{\pgfqpoint{1.839895in}{1.346531in}}%
\pgfpathcurveto{\pgfqpoint{1.844014in}{1.350649in}}{\pgfqpoint{1.846327in}{1.356235in}}{\pgfqpoint{1.846327in}{1.362059in}}%
\pgfpathcurveto{\pgfqpoint{1.846327in}{1.367883in}}{\pgfqpoint{1.844014in}{1.373469in}}{\pgfqpoint{1.839895in}{1.377587in}}%
\pgfpathcurveto{\pgfqpoint{1.835777in}{1.381705in}}{\pgfqpoint{1.830191in}{1.384019in}}{\pgfqpoint{1.824367in}{1.384019in}}%
\pgfpathcurveto{\pgfqpoint{1.818543in}{1.384019in}}{\pgfqpoint{1.812957in}{1.381705in}}{\pgfqpoint{1.808839in}{1.377587in}}%
\pgfpathcurveto{\pgfqpoint{1.804721in}{1.373469in}}{\pgfqpoint{1.802407in}{1.367883in}}{\pgfqpoint{1.802407in}{1.362059in}}%
\pgfpathcurveto{\pgfqpoint{1.802407in}{1.356235in}}{\pgfqpoint{1.804721in}{1.350649in}}{\pgfqpoint{1.808839in}{1.346531in}}%
\pgfpathcurveto{\pgfqpoint{1.812957in}{1.342412in}}{\pgfqpoint{1.818543in}{1.340099in}}{\pgfqpoint{1.824367in}{1.340099in}}%
\pgfpathlineto{\pgfqpoint{1.824367in}{1.340099in}}%
\pgfpathclose%
\pgfusepath{stroke,fill}%
\end{pgfscope}%
\begin{pgfscope}%
\pgfpathrectangle{\pgfqpoint{0.100000in}{0.209042in}}{\pgfqpoint{2.906250in}{2.906250in}}%
\pgfusepath{clip}%
\pgfsetbuttcap%
\pgfsetroundjoin%
\definecolor{currentfill}{rgb}{0.626401,0.185867,0.496456}%
\pgfsetfillcolor{currentfill}%
\pgfsetfillopacity{0.800000}%
\pgfsetlinewidth{1.003750pt}%
\definecolor{currentstroke}{rgb}{0.626401,0.185867,0.496456}%
\pgfsetstrokecolor{currentstroke}%
\pgfsetstrokeopacity{0.800000}%
\pgfsetdash{}{0pt}%
\pgfpathmoveto{\pgfqpoint{1.430643in}{1.859274in}}%
\pgfpathcurveto{\pgfqpoint{1.436467in}{1.859274in}}{\pgfqpoint{1.442053in}{1.861588in}}{\pgfqpoint{1.446171in}{1.865706in}}%
\pgfpathcurveto{\pgfqpoint{1.450289in}{1.869824in}}{\pgfqpoint{1.452603in}{1.875411in}}{\pgfqpoint{1.452603in}{1.881235in}}%
\pgfpathcurveto{\pgfqpoint{1.452603in}{1.887058in}}{\pgfqpoint{1.450289in}{1.892645in}}{\pgfqpoint{1.446171in}{1.896763in}}%
\pgfpathcurveto{\pgfqpoint{1.442053in}{1.900881in}}{\pgfqpoint{1.436467in}{1.903195in}}{\pgfqpoint{1.430643in}{1.903195in}}%
\pgfpathcurveto{\pgfqpoint{1.424819in}{1.903195in}}{\pgfqpoint{1.419233in}{1.900881in}}{\pgfqpoint{1.415114in}{1.896763in}}%
\pgfpathcurveto{\pgfqpoint{1.410996in}{1.892645in}}{\pgfqpoint{1.408682in}{1.887058in}}{\pgfqpoint{1.408682in}{1.881235in}}%
\pgfpathcurveto{\pgfqpoint{1.408682in}{1.875411in}}{\pgfqpoint{1.410996in}{1.869824in}}{\pgfqpoint{1.415114in}{1.865706in}}%
\pgfpathcurveto{\pgfqpoint{1.419233in}{1.861588in}}{\pgfqpoint{1.424819in}{1.859274in}}{\pgfqpoint{1.430643in}{1.859274in}}%
\pgfpathlineto{\pgfqpoint{1.430643in}{1.859274in}}%
\pgfpathclose%
\pgfusepath{stroke,fill}%
\end{pgfscope}%
\begin{pgfscope}%
\pgfpathrectangle{\pgfqpoint{0.100000in}{0.209042in}}{\pgfqpoint{2.906250in}{2.906250in}}%
\pgfusepath{clip}%
\pgfsetbuttcap%
\pgfsetroundjoin%
\definecolor{currentfill}{rgb}{0.600868,0.177743,0.500394}%
\pgfsetfillcolor{currentfill}%
\pgfsetfillopacity{0.800000}%
\pgfsetlinewidth{1.003750pt}%
\definecolor{currentstroke}{rgb}{0.600868,0.177743,0.500394}%
\pgfsetstrokecolor{currentstroke}%
\pgfsetstrokeopacity{0.800000}%
\pgfsetdash{}{0pt}%
\pgfpathmoveto{\pgfqpoint{1.422223in}{1.536520in}}%
\pgfpathcurveto{\pgfqpoint{1.428047in}{1.536520in}}{\pgfqpoint{1.433634in}{1.538834in}}{\pgfqpoint{1.437752in}{1.542952in}}%
\pgfpathcurveto{\pgfqpoint{1.441870in}{1.547070in}}{\pgfqpoint{1.444184in}{1.552656in}}{\pgfqpoint{1.444184in}{1.558480in}}%
\pgfpathcurveto{\pgfqpoint{1.444184in}{1.564304in}}{\pgfqpoint{1.441870in}{1.569890in}}{\pgfqpoint{1.437752in}{1.574008in}}%
\pgfpathcurveto{\pgfqpoint{1.433634in}{1.578126in}}{\pgfqpoint{1.428047in}{1.580440in}}{\pgfqpoint{1.422223in}{1.580440in}}%
\pgfpathcurveto{\pgfqpoint{1.416399in}{1.580440in}}{\pgfqpoint{1.410813in}{1.578126in}}{\pgfqpoint{1.406695in}{1.574008in}}%
\pgfpathcurveto{\pgfqpoint{1.402577in}{1.569890in}}{\pgfqpoint{1.400263in}{1.564304in}}{\pgfqpoint{1.400263in}{1.558480in}}%
\pgfpathcurveto{\pgfqpoint{1.400263in}{1.552656in}}{\pgfqpoint{1.402577in}{1.547070in}}{\pgfqpoint{1.406695in}{1.542952in}}%
\pgfpathcurveto{\pgfqpoint{1.410813in}{1.538834in}}{\pgfqpoint{1.416399in}{1.536520in}}{\pgfqpoint{1.422223in}{1.536520in}}%
\pgfpathlineto{\pgfqpoint{1.422223in}{1.536520in}}%
\pgfpathclose%
\pgfusepath{stroke,fill}%
\end{pgfscope}%
\begin{pgfscope}%
\pgfpathrectangle{\pgfqpoint{0.100000in}{0.209042in}}{\pgfqpoint{2.906250in}{2.906250in}}%
\pgfusepath{clip}%
\pgfsetbuttcap%
\pgfsetroundjoin%
\definecolor{currentfill}{rgb}{0.798608,0.247040,0.444848}%
\pgfsetfillcolor{currentfill}%
\pgfsetfillopacity{0.800000}%
\pgfsetlinewidth{1.003750pt}%
\definecolor{currentstroke}{rgb}{0.798608,0.247040,0.444848}%
\pgfsetstrokecolor{currentstroke}%
\pgfsetstrokeopacity{0.800000}%
\pgfsetdash{}{0pt}%
\pgfpathmoveto{\pgfqpoint{2.054315in}{1.881527in}}%
\pgfpathcurveto{\pgfqpoint{2.060139in}{1.881527in}}{\pgfqpoint{2.065725in}{1.883841in}}{\pgfqpoint{2.069843in}{1.887959in}}%
\pgfpathcurveto{\pgfqpoint{2.073962in}{1.892077in}}{\pgfqpoint{2.076275in}{1.897663in}}{\pgfqpoint{2.076275in}{1.903487in}}%
\pgfpathcurveto{\pgfqpoint{2.076275in}{1.909311in}}{\pgfqpoint{2.073962in}{1.914897in}}{\pgfqpoint{2.069843in}{1.919015in}}%
\pgfpathcurveto{\pgfqpoint{2.065725in}{1.923134in}}{\pgfqpoint{2.060139in}{1.925447in}}{\pgfqpoint{2.054315in}{1.925447in}}%
\pgfpathcurveto{\pgfqpoint{2.048491in}{1.925447in}}{\pgfqpoint{2.042905in}{1.923134in}}{\pgfqpoint{2.038787in}{1.919015in}}%
\pgfpathcurveto{\pgfqpoint{2.034669in}{1.914897in}}{\pgfqpoint{2.032355in}{1.909311in}}{\pgfqpoint{2.032355in}{1.903487in}}%
\pgfpathcurveto{\pgfqpoint{2.032355in}{1.897663in}}{\pgfqpoint{2.034669in}{1.892077in}}{\pgfqpoint{2.038787in}{1.887959in}}%
\pgfpathcurveto{\pgfqpoint{2.042905in}{1.883841in}}{\pgfqpoint{2.048491in}{1.881527in}}{\pgfqpoint{2.054315in}{1.881527in}}%
\pgfpathlineto{\pgfqpoint{2.054315in}{1.881527in}}%
\pgfpathclose%
\pgfusepath{stroke,fill}%
\end{pgfscope}%
\begin{pgfscope}%
\pgfpathrectangle{\pgfqpoint{0.100000in}{0.209042in}}{\pgfqpoint{2.906250in}{2.906250in}}%
\pgfusepath{clip}%
\pgfsetbuttcap%
\pgfsetroundjoin%
\definecolor{currentfill}{rgb}{0.792427,0.244242,0.447543}%
\pgfsetfillcolor{currentfill}%
\pgfsetfillopacity{0.800000}%
\pgfsetlinewidth{1.003750pt}%
\definecolor{currentstroke}{rgb}{0.792427,0.244242,0.447543}%
\pgfsetstrokecolor{currentstroke}%
\pgfsetstrokeopacity{0.800000}%
\pgfsetdash{}{0pt}%
\pgfpathmoveto{\pgfqpoint{1.710391in}{2.036972in}}%
\pgfpathcurveto{\pgfqpoint{1.716215in}{2.036972in}}{\pgfqpoint{1.721801in}{2.039286in}}{\pgfqpoint{1.725919in}{2.043404in}}%
\pgfpathcurveto{\pgfqpoint{1.730038in}{2.047522in}}{\pgfqpoint{1.732351in}{2.053108in}}{\pgfqpoint{1.732351in}{2.058932in}}%
\pgfpathcurveto{\pgfqpoint{1.732351in}{2.064756in}}{\pgfqpoint{1.730038in}{2.070342in}}{\pgfqpoint{1.725919in}{2.074461in}}%
\pgfpathcurveto{\pgfqpoint{1.721801in}{2.078579in}}{\pgfqpoint{1.716215in}{2.080893in}}{\pgfqpoint{1.710391in}{2.080893in}}%
\pgfpathcurveto{\pgfqpoint{1.704567in}{2.080893in}}{\pgfqpoint{1.698981in}{2.078579in}}{\pgfqpoint{1.694863in}{2.074461in}}%
\pgfpathcurveto{\pgfqpoint{1.690745in}{2.070342in}}{\pgfqpoint{1.688431in}{2.064756in}}{\pgfqpoint{1.688431in}{2.058932in}}%
\pgfpathcurveto{\pgfqpoint{1.688431in}{2.053108in}}{\pgfqpoint{1.690745in}{2.047522in}}{\pgfqpoint{1.694863in}{2.043404in}}%
\pgfpathcurveto{\pgfqpoint{1.698981in}{2.039286in}}{\pgfqpoint{1.704567in}{2.036972in}}{\pgfqpoint{1.710391in}{2.036972in}}%
\pgfpathlineto{\pgfqpoint{1.710391in}{2.036972in}}%
\pgfpathclose%
\pgfusepath{stroke,fill}%
\end{pgfscope}%
\begin{pgfscope}%
\pgfpathrectangle{\pgfqpoint{0.100000in}{0.209042in}}{\pgfqpoint{2.906250in}{2.906250in}}%
\pgfusepath{clip}%
\pgfsetbuttcap%
\pgfsetroundjoin%
\definecolor{currentfill}{rgb}{0.652056,0.193986,0.491611}%
\pgfsetfillcolor{currentfill}%
\pgfsetfillopacity{0.800000}%
\pgfsetlinewidth{1.003750pt}%
\definecolor{currentstroke}{rgb}{0.652056,0.193986,0.491611}%
\pgfsetstrokecolor{currentstroke}%
\pgfsetstrokeopacity{0.800000}%
\pgfsetdash{}{0pt}%
\pgfpathmoveto{\pgfqpoint{1.482821in}{1.441625in}}%
\pgfpathcurveto{\pgfqpoint{1.488645in}{1.441625in}}{\pgfqpoint{1.494231in}{1.443939in}}{\pgfqpoint{1.498349in}{1.448057in}}%
\pgfpathcurveto{\pgfqpoint{1.502467in}{1.452175in}}{\pgfqpoint{1.504781in}{1.457761in}}{\pgfqpoint{1.504781in}{1.463585in}}%
\pgfpathcurveto{\pgfqpoint{1.504781in}{1.469409in}}{\pgfqpoint{1.502467in}{1.474995in}}{\pgfqpoint{1.498349in}{1.479113in}}%
\pgfpathcurveto{\pgfqpoint{1.494231in}{1.483231in}}{\pgfqpoint{1.488645in}{1.485545in}}{\pgfqpoint{1.482821in}{1.485545in}}%
\pgfpathcurveto{\pgfqpoint{1.476997in}{1.485545in}}{\pgfqpoint{1.471411in}{1.483231in}}{\pgfqpoint{1.467293in}{1.479113in}}%
\pgfpathcurveto{\pgfqpoint{1.463174in}{1.474995in}}{\pgfqpoint{1.460861in}{1.469409in}}{\pgfqpoint{1.460861in}{1.463585in}}%
\pgfpathcurveto{\pgfqpoint{1.460861in}{1.457761in}}{\pgfqpoint{1.463174in}{1.452175in}}{\pgfqpoint{1.467293in}{1.448057in}}%
\pgfpathcurveto{\pgfqpoint{1.471411in}{1.443939in}}{\pgfqpoint{1.476997in}{1.441625in}}{\pgfqpoint{1.482821in}{1.441625in}}%
\pgfpathlineto{\pgfqpoint{1.482821in}{1.441625in}}%
\pgfpathclose%
\pgfusepath{stroke,fill}%
\end{pgfscope}%
\begin{pgfscope}%
\pgfpathrectangle{\pgfqpoint{0.100000in}{0.209042in}}{\pgfqpoint{2.906250in}{2.906250in}}%
\pgfusepath{clip}%
\pgfsetbuttcap%
\pgfsetroundjoin%
\definecolor{currentfill}{rgb}{0.953099,0.397563,0.361438}%
\pgfsetfillcolor{currentfill}%
\pgfsetfillopacity{0.800000}%
\pgfsetlinewidth{1.003750pt}%
\definecolor{currentstroke}{rgb}{0.953099,0.397563,0.361438}%
\pgfsetstrokecolor{currentstroke}%
\pgfsetstrokeopacity{0.800000}%
\pgfsetdash{}{0pt}%
\pgfpathmoveto{\pgfqpoint{1.629096in}{2.173302in}}%
\pgfpathcurveto{\pgfqpoint{1.634920in}{2.173302in}}{\pgfqpoint{1.640506in}{2.175616in}}{\pgfqpoint{1.644625in}{2.179734in}}%
\pgfpathcurveto{\pgfqpoint{1.648743in}{2.183852in}}{\pgfqpoint{1.651057in}{2.189438in}}{\pgfqpoint{1.651057in}{2.195262in}}%
\pgfpathcurveto{\pgfqpoint{1.651057in}{2.201086in}}{\pgfqpoint{1.648743in}{2.206672in}}{\pgfqpoint{1.644625in}{2.210790in}}%
\pgfpathcurveto{\pgfqpoint{1.640506in}{2.214908in}}{\pgfqpoint{1.634920in}{2.217222in}}{\pgfqpoint{1.629096in}{2.217222in}}%
\pgfpathcurveto{\pgfqpoint{1.623272in}{2.217222in}}{\pgfqpoint{1.617686in}{2.214908in}}{\pgfqpoint{1.613568in}{2.210790in}}%
\pgfpathcurveto{\pgfqpoint{1.609450in}{2.206672in}}{\pgfqpoint{1.607136in}{2.201086in}}{\pgfqpoint{1.607136in}{2.195262in}}%
\pgfpathcurveto{\pgfqpoint{1.607136in}{2.189438in}}{\pgfqpoint{1.609450in}{2.183852in}}{\pgfqpoint{1.613568in}{2.179734in}}%
\pgfpathcurveto{\pgfqpoint{1.617686in}{2.175616in}}{\pgfqpoint{1.623272in}{2.173302in}}{\pgfqpoint{1.629096in}{2.173302in}}%
\pgfpathlineto{\pgfqpoint{1.629096in}{2.173302in}}%
\pgfpathclose%
\pgfusepath{stroke,fill}%
\end{pgfscope}%
\begin{pgfscope}%
\pgfpathrectangle{\pgfqpoint{0.100000in}{0.209042in}}{\pgfqpoint{2.906250in}{2.906250in}}%
\pgfusepath{clip}%
\pgfsetbuttcap%
\pgfsetroundjoin%
\definecolor{currentfill}{rgb}{0.575490,0.169530,0.503466}%
\pgfsetfillcolor{currentfill}%
\pgfsetfillopacity{0.800000}%
\pgfsetlinewidth{1.003750pt}%
\definecolor{currentstroke}{rgb}{0.575490,0.169530,0.503466}%
\pgfsetstrokecolor{currentstroke}%
\pgfsetstrokeopacity{0.800000}%
\pgfsetdash{}{0pt}%
\pgfpathmoveto{\pgfqpoint{1.790855in}{1.470974in}}%
\pgfpathcurveto{\pgfqpoint{1.796679in}{1.470974in}}{\pgfqpoint{1.802265in}{1.473288in}}{\pgfqpoint{1.806383in}{1.477406in}}%
\pgfpathcurveto{\pgfqpoint{1.810502in}{1.481524in}}{\pgfqpoint{1.812815in}{1.487110in}}{\pgfqpoint{1.812815in}{1.492934in}}%
\pgfpathcurveto{\pgfqpoint{1.812815in}{1.498758in}}{\pgfqpoint{1.810502in}{1.504344in}}{\pgfqpoint{1.806383in}{1.508462in}}%
\pgfpathcurveto{\pgfqpoint{1.802265in}{1.512581in}}{\pgfqpoint{1.796679in}{1.514894in}}{\pgfqpoint{1.790855in}{1.514894in}}%
\pgfpathcurveto{\pgfqpoint{1.785031in}{1.514894in}}{\pgfqpoint{1.779445in}{1.512581in}}{\pgfqpoint{1.775327in}{1.508462in}}%
\pgfpathcurveto{\pgfqpoint{1.771209in}{1.504344in}}{\pgfqpoint{1.768895in}{1.498758in}}{\pgfqpoint{1.768895in}{1.492934in}}%
\pgfpathcurveto{\pgfqpoint{1.768895in}{1.487110in}}{\pgfqpoint{1.771209in}{1.481524in}}{\pgfqpoint{1.775327in}{1.477406in}}%
\pgfpathcurveto{\pgfqpoint{1.779445in}{1.473288in}}{\pgfqpoint{1.785031in}{1.470974in}}{\pgfqpoint{1.790855in}{1.470974in}}%
\pgfpathlineto{\pgfqpoint{1.790855in}{1.470974in}}%
\pgfpathclose%
\pgfusepath{stroke,fill}%
\end{pgfscope}%
\begin{pgfscope}%
\pgfpathrectangle{\pgfqpoint{0.100000in}{0.209042in}}{\pgfqpoint{2.906250in}{2.906250in}}%
\pgfusepath{clip}%
\pgfsetbuttcap%
\pgfsetroundjoin%
\definecolor{currentfill}{rgb}{0.908884,0.324755,0.385308}%
\pgfsetfillcolor{currentfill}%
\pgfsetfillopacity{0.800000}%
\pgfsetlinewidth{1.003750pt}%
\definecolor{currentstroke}{rgb}{0.908884,0.324755,0.385308}%
\pgfsetstrokecolor{currentstroke}%
\pgfsetstrokeopacity{0.800000}%
\pgfsetdash{}{0pt}%
\pgfpathmoveto{\pgfqpoint{2.250081in}{1.683192in}}%
\pgfpathcurveto{\pgfqpoint{2.255905in}{1.683192in}}{\pgfqpoint{2.261491in}{1.685506in}}{\pgfqpoint{2.265609in}{1.689624in}}%
\pgfpathcurveto{\pgfqpoint{2.269727in}{1.693742in}}{\pgfqpoint{2.272041in}{1.699329in}}{\pgfqpoint{2.272041in}{1.705152in}}%
\pgfpathcurveto{\pgfqpoint{2.272041in}{1.710976in}}{\pgfqpoint{2.269727in}{1.716563in}}{\pgfqpoint{2.265609in}{1.720681in}}%
\pgfpathcurveto{\pgfqpoint{2.261491in}{1.724799in}}{\pgfqpoint{2.255905in}{1.727113in}}{\pgfqpoint{2.250081in}{1.727113in}}%
\pgfpathcurveto{\pgfqpoint{2.244257in}{1.727113in}}{\pgfqpoint{2.238670in}{1.724799in}}{\pgfqpoint{2.234552in}{1.720681in}}%
\pgfpathcurveto{\pgfqpoint{2.230434in}{1.716563in}}{\pgfqpoint{2.228120in}{1.710976in}}{\pgfqpoint{2.228120in}{1.705152in}}%
\pgfpathcurveto{\pgfqpoint{2.228120in}{1.699329in}}{\pgfqpoint{2.230434in}{1.693742in}}{\pgfqpoint{2.234552in}{1.689624in}}%
\pgfpathcurveto{\pgfqpoint{2.238670in}{1.685506in}}{\pgfqpoint{2.244257in}{1.683192in}}{\pgfqpoint{2.250081in}{1.683192in}}%
\pgfpathlineto{\pgfqpoint{2.250081in}{1.683192in}}%
\pgfpathclose%
\pgfusepath{stroke,fill}%
\end{pgfscope}%
\begin{pgfscope}%
\pgfpathrectangle{\pgfqpoint{0.100000in}{0.209042in}}{\pgfqpoint{2.906250in}{2.906250in}}%
\pgfusepath{clip}%
\pgfsetbuttcap%
\pgfsetroundjoin%
\definecolor{currentfill}{rgb}{0.494258,0.141462,0.507988}%
\pgfsetfillcolor{currentfill}%
\pgfsetfillopacity{0.800000}%
\pgfsetlinewidth{1.003750pt}%
\definecolor{currentstroke}{rgb}{0.494258,0.141462,0.507988}%
\pgfsetstrokecolor{currentstroke}%
\pgfsetstrokeopacity{0.800000}%
\pgfsetdash{}{0pt}%
\pgfpathmoveto{\pgfqpoint{1.717129in}{1.698216in}}%
\pgfpathcurveto{\pgfqpoint{1.722953in}{1.698216in}}{\pgfqpoint{1.728539in}{1.700530in}}{\pgfqpoint{1.732657in}{1.704648in}}%
\pgfpathcurveto{\pgfqpoint{1.736776in}{1.708766in}}{\pgfqpoint{1.739089in}{1.714352in}}{\pgfqpoint{1.739089in}{1.720176in}}%
\pgfpathcurveto{\pgfqpoint{1.739089in}{1.726000in}}{\pgfqpoint{1.736776in}{1.731586in}}{\pgfqpoint{1.732657in}{1.735704in}}%
\pgfpathcurveto{\pgfqpoint{1.728539in}{1.739822in}}{\pgfqpoint{1.722953in}{1.742136in}}{\pgfqpoint{1.717129in}{1.742136in}}%
\pgfpathcurveto{\pgfqpoint{1.711305in}{1.742136in}}{\pgfqpoint{1.705719in}{1.739822in}}{\pgfqpoint{1.701601in}{1.735704in}}%
\pgfpathcurveto{\pgfqpoint{1.697483in}{1.731586in}}{\pgfqpoint{1.695169in}{1.726000in}}{\pgfqpoint{1.695169in}{1.720176in}}%
\pgfpathcurveto{\pgfqpoint{1.695169in}{1.714352in}}{\pgfqpoint{1.697483in}{1.708766in}}{\pgfqpoint{1.701601in}{1.704648in}}%
\pgfpathcurveto{\pgfqpoint{1.705719in}{1.700530in}}{\pgfqpoint{1.711305in}{1.698216in}}{\pgfqpoint{1.717129in}{1.698216in}}%
\pgfpathlineto{\pgfqpoint{1.717129in}{1.698216in}}%
\pgfpathclose%
\pgfusepath{stroke,fill}%
\end{pgfscope}%
\begin{pgfscope}%
\pgfpathrectangle{\pgfqpoint{0.100000in}{0.209042in}}{\pgfqpoint{2.906250in}{2.906250in}}%
\pgfusepath{clip}%
\pgfsetbuttcap%
\pgfsetroundjoin%
\definecolor{currentfill}{rgb}{0.607238,0.179779,0.499492}%
\pgfsetfillcolor{currentfill}%
\pgfsetfillopacity{0.800000}%
\pgfsetlinewidth{1.003750pt}%
\definecolor{currentstroke}{rgb}{0.607238,0.179779,0.499492}%
\pgfsetstrokecolor{currentstroke}%
\pgfsetstrokeopacity{0.800000}%
\pgfsetdash{}{0pt}%
\pgfpathmoveto{\pgfqpoint{1.613920in}{1.448080in}}%
\pgfpathcurveto{\pgfqpoint{1.619744in}{1.448080in}}{\pgfqpoint{1.625330in}{1.450394in}}{\pgfqpoint{1.629449in}{1.454512in}}%
\pgfpathcurveto{\pgfqpoint{1.633567in}{1.458630in}}{\pgfqpoint{1.635881in}{1.464216in}}{\pgfqpoint{1.635881in}{1.470040in}}%
\pgfpathcurveto{\pgfqpoint{1.635881in}{1.475864in}}{\pgfqpoint{1.633567in}{1.481450in}}{\pgfqpoint{1.629449in}{1.485568in}}%
\pgfpathcurveto{\pgfqpoint{1.625330in}{1.489687in}}{\pgfqpoint{1.619744in}{1.492000in}}{\pgfqpoint{1.613920in}{1.492000in}}%
\pgfpathcurveto{\pgfqpoint{1.608096in}{1.492000in}}{\pgfqpoint{1.602510in}{1.489687in}}{\pgfqpoint{1.598392in}{1.485568in}}%
\pgfpathcurveto{\pgfqpoint{1.594274in}{1.481450in}}{\pgfqpoint{1.591960in}{1.475864in}}{\pgfqpoint{1.591960in}{1.470040in}}%
\pgfpathcurveto{\pgfqpoint{1.591960in}{1.464216in}}{\pgfqpoint{1.594274in}{1.458630in}}{\pgfqpoint{1.598392in}{1.454512in}}%
\pgfpathcurveto{\pgfqpoint{1.602510in}{1.450394in}}{\pgfqpoint{1.608096in}{1.448080in}}{\pgfqpoint{1.613920in}{1.448080in}}%
\pgfpathlineto{\pgfqpoint{1.613920in}{1.448080in}}%
\pgfpathclose%
\pgfusepath{stroke,fill}%
\end{pgfscope}%
\begin{pgfscope}%
\pgfpathrectangle{\pgfqpoint{0.100000in}{0.209042in}}{\pgfqpoint{2.906250in}{2.906250in}}%
\pgfusepath{clip}%
\pgfsetbuttcap%
\pgfsetroundjoin%
\definecolor{currentfill}{rgb}{0.575490,0.169530,0.503466}%
\pgfsetfillcolor{currentfill}%
\pgfsetfillopacity{0.800000}%
\pgfsetlinewidth{1.003750pt}%
\definecolor{currentstroke}{rgb}{0.575490,0.169530,0.503466}%
\pgfsetstrokecolor{currentstroke}%
\pgfsetstrokeopacity{0.800000}%
\pgfsetdash{}{0pt}%
\pgfpathmoveto{\pgfqpoint{1.609307in}{1.492338in}}%
\pgfpathcurveto{\pgfqpoint{1.615131in}{1.492338in}}{\pgfqpoint{1.620717in}{1.494652in}}{\pgfqpoint{1.624836in}{1.498770in}}%
\pgfpathcurveto{\pgfqpoint{1.628954in}{1.502889in}}{\pgfqpoint{1.631268in}{1.508475in}}{\pgfqpoint{1.631268in}{1.514299in}}%
\pgfpathcurveto{\pgfqpoint{1.631268in}{1.520123in}}{\pgfqpoint{1.628954in}{1.525709in}}{\pgfqpoint{1.624836in}{1.529827in}}%
\pgfpathcurveto{\pgfqpoint{1.620717in}{1.533945in}}{\pgfqpoint{1.615131in}{1.536259in}}{\pgfqpoint{1.609307in}{1.536259in}}%
\pgfpathcurveto{\pgfqpoint{1.603483in}{1.536259in}}{\pgfqpoint{1.597897in}{1.533945in}}{\pgfqpoint{1.593779in}{1.529827in}}%
\pgfpathcurveto{\pgfqpoint{1.589661in}{1.525709in}}{\pgfqpoint{1.587347in}{1.520123in}}{\pgfqpoint{1.587347in}{1.514299in}}%
\pgfpathcurveto{\pgfqpoint{1.587347in}{1.508475in}}{\pgfqpoint{1.589661in}{1.502889in}}{\pgfqpoint{1.593779in}{1.498770in}}%
\pgfpathcurveto{\pgfqpoint{1.597897in}{1.494652in}}{\pgfqpoint{1.603483in}{1.492338in}}{\pgfqpoint{1.609307in}{1.492338in}}%
\pgfpathlineto{\pgfqpoint{1.609307in}{1.492338in}}%
\pgfpathclose%
\pgfusepath{stroke,fill}%
\end{pgfscope}%
\begin{pgfscope}%
\pgfpathrectangle{\pgfqpoint{0.100000in}{0.209042in}}{\pgfqpoint{2.906250in}{2.906250in}}%
\pgfusepath{clip}%
\pgfsetbuttcap%
\pgfsetroundjoin%
\definecolor{currentfill}{rgb}{0.937221,0.364929,0.368567}%
\pgfsetfillcolor{currentfill}%
\pgfsetfillopacity{0.800000}%
\pgfsetlinewidth{1.003750pt}%
\definecolor{currentstroke}{rgb}{0.937221,0.364929,0.368567}%
\pgfsetstrokecolor{currentstroke}%
\pgfsetstrokeopacity{0.800000}%
\pgfsetdash{}{0pt}%
\pgfpathmoveto{\pgfqpoint{1.101605in}{1.818143in}}%
\pgfpathcurveto{\pgfqpoint{1.107429in}{1.818143in}}{\pgfqpoint{1.113015in}{1.820457in}}{\pgfqpoint{1.117133in}{1.824575in}}%
\pgfpathcurveto{\pgfqpoint{1.121251in}{1.828693in}}{\pgfqpoint{1.123565in}{1.834279in}}{\pgfqpoint{1.123565in}{1.840103in}}%
\pgfpathcurveto{\pgfqpoint{1.123565in}{1.845927in}}{\pgfqpoint{1.121251in}{1.851513in}}{\pgfqpoint{1.117133in}{1.855632in}}%
\pgfpathcurveto{\pgfqpoint{1.113015in}{1.859750in}}{\pgfqpoint{1.107429in}{1.862064in}}{\pgfqpoint{1.101605in}{1.862064in}}%
\pgfpathcurveto{\pgfqpoint{1.095781in}{1.862064in}}{\pgfqpoint{1.090195in}{1.859750in}}{\pgfqpoint{1.086077in}{1.855632in}}%
\pgfpathcurveto{\pgfqpoint{1.081958in}{1.851513in}}{\pgfqpoint{1.079645in}{1.845927in}}{\pgfqpoint{1.079645in}{1.840103in}}%
\pgfpathcurveto{\pgfqpoint{1.079645in}{1.834279in}}{\pgfqpoint{1.081958in}{1.828693in}}{\pgfqpoint{1.086077in}{1.824575in}}%
\pgfpathcurveto{\pgfqpoint{1.090195in}{1.820457in}}{\pgfqpoint{1.095781in}{1.818143in}}{\pgfqpoint{1.101605in}{1.818143in}}%
\pgfpathlineto{\pgfqpoint{1.101605in}{1.818143in}}%
\pgfpathclose%
\pgfusepath{stroke,fill}%
\end{pgfscope}%
\begin{pgfscope}%
\pgfpathrectangle{\pgfqpoint{0.100000in}{0.209042in}}{\pgfqpoint{2.906250in}{2.906250in}}%
\pgfusepath{clip}%
\pgfsetbuttcap%
\pgfsetroundjoin%
\definecolor{currentfill}{rgb}{0.671349,0.200133,0.487358}%
\pgfsetfillcolor{currentfill}%
\pgfsetfillopacity{0.800000}%
\pgfsetlinewidth{1.003750pt}%
\definecolor{currentstroke}{rgb}{0.671349,0.200133,0.487358}%
\pgfsetstrokecolor{currentstroke}%
\pgfsetstrokeopacity{0.800000}%
\pgfsetdash{}{0pt}%
\pgfpathmoveto{\pgfqpoint{1.922393in}{1.493533in}}%
\pgfpathcurveto{\pgfqpoint{1.928217in}{1.493533in}}{\pgfqpoint{1.933804in}{1.495846in}}{\pgfqpoint{1.937922in}{1.499965in}}%
\pgfpathcurveto{\pgfqpoint{1.942040in}{1.504083in}}{\pgfqpoint{1.944354in}{1.509669in}}{\pgfqpoint{1.944354in}{1.515493in}}%
\pgfpathcurveto{\pgfqpoint{1.944354in}{1.521317in}}{\pgfqpoint{1.942040in}{1.526903in}}{\pgfqpoint{1.937922in}{1.531021in}}%
\pgfpathcurveto{\pgfqpoint{1.933804in}{1.535139in}}{\pgfqpoint{1.928217in}{1.537453in}}{\pgfqpoint{1.922393in}{1.537453in}}%
\pgfpathcurveto{\pgfqpoint{1.916570in}{1.537453in}}{\pgfqpoint{1.910983in}{1.535139in}}{\pgfqpoint{1.906865in}{1.531021in}}%
\pgfpathcurveto{\pgfqpoint{1.902747in}{1.526903in}}{\pgfqpoint{1.900433in}{1.521317in}}{\pgfqpoint{1.900433in}{1.515493in}}%
\pgfpathcurveto{\pgfqpoint{1.900433in}{1.509669in}}{\pgfqpoint{1.902747in}{1.504083in}}{\pgfqpoint{1.906865in}{1.499965in}}%
\pgfpathcurveto{\pgfqpoint{1.910983in}{1.495846in}}{\pgfqpoint{1.916570in}{1.493533in}}{\pgfqpoint{1.922393in}{1.493533in}}%
\pgfpathlineto{\pgfqpoint{1.922393in}{1.493533in}}%
\pgfpathclose%
\pgfusepath{stroke,fill}%
\end{pgfscope}%
\begin{pgfscope}%
\pgfpathrectangle{\pgfqpoint{0.100000in}{0.209042in}}{\pgfqpoint{2.906250in}{2.906250in}}%
\pgfusepath{clip}%
\pgfsetbuttcap%
\pgfsetroundjoin%
\definecolor{currentfill}{rgb}{0.658483,0.196027,0.490253}%
\pgfsetfillcolor{currentfill}%
\pgfsetfillopacity{0.800000}%
\pgfsetlinewidth{1.003750pt}%
\definecolor{currentstroke}{rgb}{0.658483,0.196027,0.490253}%
\pgfsetstrokecolor{currentstroke}%
\pgfsetstrokeopacity{0.800000}%
\pgfsetdash{}{0pt}%
\pgfpathmoveto{\pgfqpoint{1.762578in}{1.422570in}}%
\pgfpathcurveto{\pgfqpoint{1.768402in}{1.422570in}}{\pgfqpoint{1.773988in}{1.424884in}}{\pgfqpoint{1.778106in}{1.429002in}}%
\pgfpathcurveto{\pgfqpoint{1.782224in}{1.433120in}}{\pgfqpoint{1.784538in}{1.438706in}}{\pgfqpoint{1.784538in}{1.444530in}}%
\pgfpathcurveto{\pgfqpoint{1.784538in}{1.450354in}}{\pgfqpoint{1.782224in}{1.455941in}}{\pgfqpoint{1.778106in}{1.460059in}}%
\pgfpathcurveto{\pgfqpoint{1.773988in}{1.464177in}}{\pgfqpoint{1.768402in}{1.466491in}}{\pgfqpoint{1.762578in}{1.466491in}}%
\pgfpathcurveto{\pgfqpoint{1.756754in}{1.466491in}}{\pgfqpoint{1.751168in}{1.464177in}}{\pgfqpoint{1.747050in}{1.460059in}}%
\pgfpathcurveto{\pgfqpoint{1.742932in}{1.455941in}}{\pgfqpoint{1.740618in}{1.450354in}}{\pgfqpoint{1.740618in}{1.444530in}}%
\pgfpathcurveto{\pgfqpoint{1.740618in}{1.438706in}}{\pgfqpoint{1.742932in}{1.433120in}}{\pgfqpoint{1.747050in}{1.429002in}}%
\pgfpathcurveto{\pgfqpoint{1.751168in}{1.424884in}}{\pgfqpoint{1.756754in}{1.422570in}}{\pgfqpoint{1.762578in}{1.422570in}}%
\pgfpathlineto{\pgfqpoint{1.762578in}{1.422570in}}%
\pgfpathclose%
\pgfusepath{stroke,fill}%
\end{pgfscope}%
\begin{pgfscope}%
\pgfpathrectangle{\pgfqpoint{0.100000in}{0.209042in}}{\pgfqpoint{2.906250in}{2.906250in}}%
\pgfusepath{clip}%
\pgfsetbuttcap%
\pgfsetroundjoin%
\definecolor{currentfill}{rgb}{0.645633,0.191952,0.492910}%
\pgfsetfillcolor{currentfill}%
\pgfsetfillopacity{0.800000}%
\pgfsetlinewidth{1.003750pt}%
\definecolor{currentstroke}{rgb}{0.645633,0.191952,0.492910}%
\pgfsetstrokecolor{currentstroke}%
\pgfsetstrokeopacity{0.800000}%
\pgfsetdash{}{0pt}%
\pgfpathmoveto{\pgfqpoint{1.538247in}{1.820670in}}%
\pgfpathcurveto{\pgfqpoint{1.544071in}{1.820670in}}{\pgfqpoint{1.549658in}{1.822984in}}{\pgfqpoint{1.553776in}{1.827102in}}%
\pgfpathcurveto{\pgfqpoint{1.557894in}{1.831220in}}{\pgfqpoint{1.560208in}{1.836807in}}{\pgfqpoint{1.560208in}{1.842630in}}%
\pgfpathcurveto{\pgfqpoint{1.560208in}{1.848454in}}{\pgfqpoint{1.557894in}{1.854041in}}{\pgfqpoint{1.553776in}{1.858159in}}%
\pgfpathcurveto{\pgfqpoint{1.549658in}{1.862277in}}{\pgfqpoint{1.544071in}{1.864591in}}{\pgfqpoint{1.538247in}{1.864591in}}%
\pgfpathcurveto{\pgfqpoint{1.532424in}{1.864591in}}{\pgfqpoint{1.526837in}{1.862277in}}{\pgfqpoint{1.522719in}{1.858159in}}%
\pgfpathcurveto{\pgfqpoint{1.518601in}{1.854041in}}{\pgfqpoint{1.516287in}{1.848454in}}{\pgfqpoint{1.516287in}{1.842630in}}%
\pgfpathcurveto{\pgfqpoint{1.516287in}{1.836807in}}{\pgfqpoint{1.518601in}{1.831220in}}{\pgfqpoint{1.522719in}{1.827102in}}%
\pgfpathcurveto{\pgfqpoint{1.526837in}{1.822984in}}{\pgfqpoint{1.532424in}{1.820670in}}{\pgfqpoint{1.538247in}{1.820670in}}%
\pgfpathlineto{\pgfqpoint{1.538247in}{1.820670in}}%
\pgfpathclose%
\pgfusepath{stroke,fill}%
\end{pgfscope}%
\begin{pgfscope}%
\pgfpathrectangle{\pgfqpoint{0.100000in}{0.209042in}}{\pgfqpoint{2.906250in}{2.906250in}}%
\pgfusepath{clip}%
\pgfsetbuttcap%
\pgfsetroundjoin%
\definecolor{currentfill}{rgb}{0.985693,0.528148,0.379371}%
\pgfsetfillcolor{currentfill}%
\pgfsetfillopacity{0.800000}%
\pgfsetlinewidth{1.003750pt}%
\definecolor{currentstroke}{rgb}{0.985693,0.528148,0.379371}%
\pgfsetstrokecolor{currentstroke}%
\pgfsetstrokeopacity{0.800000}%
\pgfsetdash{}{0pt}%
\pgfpathmoveto{\pgfqpoint{1.902364in}{2.174112in}}%
\pgfpathcurveto{\pgfqpoint{1.908188in}{2.174112in}}{\pgfqpoint{1.913774in}{2.176426in}}{\pgfqpoint{1.917893in}{2.180544in}}%
\pgfpathcurveto{\pgfqpoint{1.922011in}{2.184663in}}{\pgfqpoint{1.924325in}{2.190249in}}{\pgfqpoint{1.924325in}{2.196073in}}%
\pgfpathcurveto{\pgfqpoint{1.924325in}{2.201897in}}{\pgfqpoint{1.922011in}{2.207483in}}{\pgfqpoint{1.917893in}{2.211601in}}%
\pgfpathcurveto{\pgfqpoint{1.913774in}{2.215719in}}{\pgfqpoint{1.908188in}{2.218033in}}{\pgfqpoint{1.902364in}{2.218033in}}%
\pgfpathcurveto{\pgfqpoint{1.896540in}{2.218033in}}{\pgfqpoint{1.890954in}{2.215719in}}{\pgfqpoint{1.886836in}{2.211601in}}%
\pgfpathcurveto{\pgfqpoint{1.882718in}{2.207483in}}{\pgfqpoint{1.880404in}{2.201897in}}{\pgfqpoint{1.880404in}{2.196073in}}%
\pgfpathcurveto{\pgfqpoint{1.880404in}{2.190249in}}{\pgfqpoint{1.882718in}{2.184663in}}{\pgfqpoint{1.886836in}{2.180544in}}%
\pgfpathcurveto{\pgfqpoint{1.890954in}{2.176426in}}{\pgfqpoint{1.896540in}{2.174112in}}{\pgfqpoint{1.902364in}{2.174112in}}%
\pgfpathlineto{\pgfqpoint{1.902364in}{2.174112in}}%
\pgfpathclose%
\pgfusepath{stroke,fill}%
\end{pgfscope}%
\begin{pgfscope}%
\pgfpathrectangle{\pgfqpoint{0.100000in}{0.209042in}}{\pgfqpoint{2.906250in}{2.906250in}}%
\pgfusepath{clip}%
\pgfsetbuttcap%
\pgfsetroundjoin%
\definecolor{currentfill}{rgb}{0.773695,0.236249,0.455289}%
\pgfsetfillcolor{currentfill}%
\pgfsetfillopacity{0.800000}%
\pgfsetlinewidth{1.003750pt}%
\definecolor{currentstroke}{rgb}{0.773695,0.236249,0.455289}%
\pgfsetstrokecolor{currentstroke}%
\pgfsetstrokeopacity{0.800000}%
\pgfsetdash{}{0pt}%
\pgfpathmoveto{\pgfqpoint{1.402031in}{1.444450in}}%
\pgfpathcurveto{\pgfqpoint{1.407855in}{1.444450in}}{\pgfqpoint{1.413441in}{1.446764in}}{\pgfqpoint{1.417560in}{1.450882in}}%
\pgfpathcurveto{\pgfqpoint{1.421678in}{1.455000in}}{\pgfqpoint{1.423992in}{1.460586in}}{\pgfqpoint{1.423992in}{1.466410in}}%
\pgfpathcurveto{\pgfqpoint{1.423992in}{1.472234in}}{\pgfqpoint{1.421678in}{1.477820in}}{\pgfqpoint{1.417560in}{1.481938in}}%
\pgfpathcurveto{\pgfqpoint{1.413441in}{1.486056in}}{\pgfqpoint{1.407855in}{1.488370in}}{\pgfqpoint{1.402031in}{1.488370in}}%
\pgfpathcurveto{\pgfqpoint{1.396207in}{1.488370in}}{\pgfqpoint{1.390621in}{1.486056in}}{\pgfqpoint{1.386503in}{1.481938in}}%
\pgfpathcurveto{\pgfqpoint{1.382385in}{1.477820in}}{\pgfqpoint{1.380071in}{1.472234in}}{\pgfqpoint{1.380071in}{1.466410in}}%
\pgfpathcurveto{\pgfqpoint{1.380071in}{1.460586in}}{\pgfqpoint{1.382385in}{1.455000in}}{\pgfqpoint{1.386503in}{1.450882in}}%
\pgfpathcurveto{\pgfqpoint{1.390621in}{1.446764in}}{\pgfqpoint{1.396207in}{1.444450in}}{\pgfqpoint{1.402031in}{1.444450in}}%
\pgfpathlineto{\pgfqpoint{1.402031in}{1.444450in}}%
\pgfpathclose%
\pgfusepath{stroke,fill}%
\end{pgfscope}%
\begin{pgfscope}%
\pgfpathrectangle{\pgfqpoint{0.100000in}{0.209042in}}{\pgfqpoint{2.906250in}{2.906250in}}%
\pgfusepath{clip}%
\pgfsetbuttcap%
\pgfsetroundjoin%
\definecolor{currentfill}{rgb}{0.684224,0.204286,0.484219}%
\pgfsetfillcolor{currentfill}%
\pgfsetfillopacity{0.800000}%
\pgfsetlinewidth{1.003750pt}%
\definecolor{currentstroke}{rgb}{0.684224,0.204286,0.484219}%
\pgfsetstrokecolor{currentstroke}%
\pgfsetstrokeopacity{0.800000}%
\pgfsetdash{}{0pt}%
\pgfpathmoveto{\pgfqpoint{1.443010in}{1.570048in}}%
\pgfpathcurveto{\pgfqpoint{1.448834in}{1.570048in}}{\pgfqpoint{1.454420in}{1.572361in}}{\pgfqpoint{1.458539in}{1.576480in}}%
\pgfpathcurveto{\pgfqpoint{1.462657in}{1.580598in}}{\pgfqpoint{1.464971in}{1.586184in}}{\pgfqpoint{1.464971in}{1.592008in}}%
\pgfpathcurveto{\pgfqpoint{1.464971in}{1.597832in}}{\pgfqpoint{1.462657in}{1.603418in}}{\pgfqpoint{1.458539in}{1.607536in}}%
\pgfpathcurveto{\pgfqpoint{1.454420in}{1.611654in}}{\pgfqpoint{1.448834in}{1.613968in}}{\pgfqpoint{1.443010in}{1.613968in}}%
\pgfpathcurveto{\pgfqpoint{1.437186in}{1.613968in}}{\pgfqpoint{1.431600in}{1.611654in}}{\pgfqpoint{1.427482in}{1.607536in}}%
\pgfpathcurveto{\pgfqpoint{1.423364in}{1.603418in}}{\pgfqpoint{1.421050in}{1.597832in}}{\pgfqpoint{1.421050in}{1.592008in}}%
\pgfpathcurveto{\pgfqpoint{1.421050in}{1.586184in}}{\pgfqpoint{1.423364in}{1.580598in}}{\pgfqpoint{1.427482in}{1.576480in}}%
\pgfpathcurveto{\pgfqpoint{1.431600in}{1.572361in}}{\pgfqpoint{1.437186in}{1.570048in}}{\pgfqpoint{1.443010in}{1.570048in}}%
\pgfpathlineto{\pgfqpoint{1.443010in}{1.570048in}}%
\pgfpathclose%
\pgfusepath{stroke,fill}%
\end{pgfscope}%
\begin{pgfscope}%
\pgfpathrectangle{\pgfqpoint{0.100000in}{0.209042in}}{\pgfqpoint{2.906250in}{2.906250in}}%
\pgfusepath{clip}%
\pgfsetbuttcap%
\pgfsetroundjoin%
\definecolor{currentfill}{rgb}{0.652056,0.193986,0.491611}%
\pgfsetfillcolor{currentfill}%
\pgfsetfillopacity{0.800000}%
\pgfsetlinewidth{1.003750pt}%
\definecolor{currentstroke}{rgb}{0.652056,0.193986,0.491611}%
\pgfsetstrokecolor{currentstroke}%
\pgfsetstrokeopacity{0.800000}%
\pgfsetdash{}{0pt}%
\pgfpathmoveto{\pgfqpoint{1.883079in}{1.580832in}}%
\pgfpathcurveto{\pgfqpoint{1.888903in}{1.580832in}}{\pgfqpoint{1.894489in}{1.583146in}}{\pgfqpoint{1.898608in}{1.587264in}}%
\pgfpathcurveto{\pgfqpoint{1.902726in}{1.591382in}}{\pgfqpoint{1.905040in}{1.596968in}}{\pgfqpoint{1.905040in}{1.602792in}}%
\pgfpathcurveto{\pgfqpoint{1.905040in}{1.608616in}}{\pgfqpoint{1.902726in}{1.614202in}}{\pgfqpoint{1.898608in}{1.618320in}}%
\pgfpathcurveto{\pgfqpoint{1.894489in}{1.622438in}}{\pgfqpoint{1.888903in}{1.624752in}}{\pgfqpoint{1.883079in}{1.624752in}}%
\pgfpathcurveto{\pgfqpoint{1.877255in}{1.624752in}}{\pgfqpoint{1.871669in}{1.622438in}}{\pgfqpoint{1.867551in}{1.618320in}}%
\pgfpathcurveto{\pgfqpoint{1.863433in}{1.614202in}}{\pgfqpoint{1.861119in}{1.608616in}}{\pgfqpoint{1.861119in}{1.602792in}}%
\pgfpathcurveto{\pgfqpoint{1.861119in}{1.596968in}}{\pgfqpoint{1.863433in}{1.591382in}}{\pgfqpoint{1.867551in}{1.587264in}}%
\pgfpathcurveto{\pgfqpoint{1.871669in}{1.583146in}}{\pgfqpoint{1.877255in}{1.580832in}}{\pgfqpoint{1.883079in}{1.580832in}}%
\pgfpathlineto{\pgfqpoint{1.883079in}{1.580832in}}%
\pgfpathclose%
\pgfusepath{stroke,fill}%
\end{pgfscope}%
\begin{pgfscope}%
\pgfpathrectangle{\pgfqpoint{0.100000in}{0.209042in}}{\pgfqpoint{2.906250in}{2.906250in}}%
\pgfusepath{clip}%
\pgfsetbuttcap%
\pgfsetroundjoin%
\definecolor{currentfill}{rgb}{0.976690,0.476226,0.364466}%
\pgfsetfillcolor{currentfill}%
\pgfsetfillopacity{0.800000}%
\pgfsetlinewidth{1.003750pt}%
\definecolor{currentstroke}{rgb}{0.976690,0.476226,0.364466}%
\pgfsetstrokecolor{currentstroke}%
\pgfsetstrokeopacity{0.800000}%
\pgfsetdash{}{0pt}%
\pgfpathmoveto{\pgfqpoint{2.257798in}{1.367442in}}%
\pgfpathcurveto{\pgfqpoint{2.263622in}{1.367442in}}{\pgfqpoint{2.269208in}{1.369756in}}{\pgfqpoint{2.273326in}{1.373874in}}%
\pgfpathcurveto{\pgfqpoint{2.277444in}{1.377993in}}{\pgfqpoint{2.279758in}{1.383579in}}{\pgfqpoint{2.279758in}{1.389403in}}%
\pgfpathcurveto{\pgfqpoint{2.279758in}{1.395227in}}{\pgfqpoint{2.277444in}{1.400813in}}{\pgfqpoint{2.273326in}{1.404931in}}%
\pgfpathcurveto{\pgfqpoint{2.269208in}{1.409049in}}{\pgfqpoint{2.263622in}{1.411363in}}{\pgfqpoint{2.257798in}{1.411363in}}%
\pgfpathcurveto{\pgfqpoint{2.251974in}{1.411363in}}{\pgfqpoint{2.246388in}{1.409049in}}{\pgfqpoint{2.242269in}{1.404931in}}%
\pgfpathcurveto{\pgfqpoint{2.238151in}{1.400813in}}{\pgfqpoint{2.235837in}{1.395227in}}{\pgfqpoint{2.235837in}{1.389403in}}%
\pgfpathcurveto{\pgfqpoint{2.235837in}{1.383579in}}{\pgfqpoint{2.238151in}{1.377993in}}{\pgfqpoint{2.242269in}{1.373874in}}%
\pgfpathcurveto{\pgfqpoint{2.246388in}{1.369756in}}{\pgfqpoint{2.251974in}{1.367442in}}{\pgfqpoint{2.257798in}{1.367442in}}%
\pgfpathlineto{\pgfqpoint{2.257798in}{1.367442in}}%
\pgfpathclose%
\pgfusepath{stroke,fill}%
\end{pgfscope}%
\begin{pgfscope}%
\pgfpathrectangle{\pgfqpoint{0.100000in}{0.209042in}}{\pgfqpoint{2.906250in}{2.906250in}}%
\pgfusepath{clip}%
\pgfsetbuttcap%
\pgfsetroundjoin%
\definecolor{currentfill}{rgb}{0.874176,0.291859,0.406205}%
\pgfsetfillcolor{currentfill}%
\pgfsetfillopacity{0.800000}%
\pgfsetlinewidth{1.003750pt}%
\definecolor{currentstroke}{rgb}{0.874176,0.291859,0.406205}%
\pgfsetstrokecolor{currentstroke}%
\pgfsetstrokeopacity{0.800000}%
\pgfsetdash{}{0pt}%
\pgfpathmoveto{\pgfqpoint{1.248269in}{1.522266in}}%
\pgfpathcurveto{\pgfqpoint{1.254093in}{1.522266in}}{\pgfqpoint{1.259679in}{1.524580in}}{\pgfqpoint{1.263797in}{1.528698in}}%
\pgfpathcurveto{\pgfqpoint{1.267915in}{1.532816in}}{\pgfqpoint{1.270229in}{1.538402in}}{\pgfqpoint{1.270229in}{1.544226in}}%
\pgfpathcurveto{\pgfqpoint{1.270229in}{1.550050in}}{\pgfqpoint{1.267915in}{1.555636in}}{\pgfqpoint{1.263797in}{1.559754in}}%
\pgfpathcurveto{\pgfqpoint{1.259679in}{1.563873in}}{\pgfqpoint{1.254093in}{1.566186in}}{\pgfqpoint{1.248269in}{1.566186in}}%
\pgfpathcurveto{\pgfqpoint{1.242445in}{1.566186in}}{\pgfqpoint{1.236859in}{1.563873in}}{\pgfqpoint{1.232741in}{1.559754in}}%
\pgfpathcurveto{\pgfqpoint{1.228622in}{1.555636in}}{\pgfqpoint{1.226309in}{1.550050in}}{\pgfqpoint{1.226309in}{1.544226in}}%
\pgfpathcurveto{\pgfqpoint{1.226309in}{1.538402in}}{\pgfqpoint{1.228622in}{1.532816in}}{\pgfqpoint{1.232741in}{1.528698in}}%
\pgfpathcurveto{\pgfqpoint{1.236859in}{1.524580in}}{\pgfqpoint{1.242445in}{1.522266in}}{\pgfqpoint{1.248269in}{1.522266in}}%
\pgfpathlineto{\pgfqpoint{1.248269in}{1.522266in}}%
\pgfpathclose%
\pgfusepath{stroke,fill}%
\end{pgfscope}%
\begin{pgfscope}%
\pgfpathrectangle{\pgfqpoint{0.100000in}{0.209042in}}{\pgfqpoint{2.906250in}{2.906250in}}%
\pgfusepath{clip}%
\pgfsetbuttcap%
\pgfsetroundjoin%
\definecolor{currentfill}{rgb}{0.958464,0.411324,0.360014}%
\pgfsetfillcolor{currentfill}%
\pgfsetfillopacity{0.800000}%
\pgfsetlinewidth{1.003750pt}%
\definecolor{currentstroke}{rgb}{0.958464,0.411324,0.360014}%
\pgfsetstrokecolor{currentstroke}%
\pgfsetstrokeopacity{0.800000}%
\pgfsetdash{}{0pt}%
\pgfpathmoveto{\pgfqpoint{1.138072in}{1.792947in}}%
\pgfpathcurveto{\pgfqpoint{1.143896in}{1.792947in}}{\pgfqpoint{1.149483in}{1.795261in}}{\pgfqpoint{1.153601in}{1.799379in}}%
\pgfpathcurveto{\pgfqpoint{1.157719in}{1.803497in}}{\pgfqpoint{1.160033in}{1.809083in}}{\pgfqpoint{1.160033in}{1.814907in}}%
\pgfpathcurveto{\pgfqpoint{1.160033in}{1.820731in}}{\pgfqpoint{1.157719in}{1.826317in}}{\pgfqpoint{1.153601in}{1.830435in}}%
\pgfpathcurveto{\pgfqpoint{1.149483in}{1.834553in}}{\pgfqpoint{1.143896in}{1.836867in}}{\pgfqpoint{1.138072in}{1.836867in}}%
\pgfpathcurveto{\pgfqpoint{1.132248in}{1.836867in}}{\pgfqpoint{1.126662in}{1.834553in}}{\pgfqpoint{1.122544in}{1.830435in}}%
\pgfpathcurveto{\pgfqpoint{1.118426in}{1.826317in}}{\pgfqpoint{1.116112in}{1.820731in}}{\pgfqpoint{1.116112in}{1.814907in}}%
\pgfpathcurveto{\pgfqpoint{1.116112in}{1.809083in}}{\pgfqpoint{1.118426in}{1.803497in}}{\pgfqpoint{1.122544in}{1.799379in}}%
\pgfpathcurveto{\pgfqpoint{1.126662in}{1.795261in}}{\pgfqpoint{1.132248in}{1.792947in}}{\pgfqpoint{1.138072in}{1.792947in}}%
\pgfpathlineto{\pgfqpoint{1.138072in}{1.792947in}}%
\pgfpathclose%
\pgfusepath{stroke,fill}%
\end{pgfscope}%
\begin{pgfscope}%
\pgfpathrectangle{\pgfqpoint{0.100000in}{0.209042in}}{\pgfqpoint{2.906250in}{2.906250in}}%
\pgfusepath{clip}%
\pgfsetbuttcap%
\pgfsetroundjoin%
\definecolor{currentfill}{rgb}{0.947180,0.384178,0.363701}%
\pgfsetfillcolor{currentfill}%
\pgfsetfillopacity{0.800000}%
\pgfsetlinewidth{1.003750pt}%
\definecolor{currentstroke}{rgb}{0.947180,0.384178,0.363701}%
\pgfsetstrokecolor{currentstroke}%
\pgfsetstrokeopacity{0.800000}%
\pgfsetdash{}{0pt}%
\pgfpathmoveto{\pgfqpoint{2.059879in}{1.304779in}}%
\pgfpathcurveto{\pgfqpoint{2.065703in}{1.304779in}}{\pgfqpoint{2.071289in}{1.307093in}}{\pgfqpoint{2.075407in}{1.311211in}}%
\pgfpathcurveto{\pgfqpoint{2.079525in}{1.315329in}}{\pgfqpoint{2.081839in}{1.320915in}}{\pgfqpoint{2.081839in}{1.326739in}}%
\pgfpathcurveto{\pgfqpoint{2.081839in}{1.332563in}}{\pgfqpoint{2.079525in}{1.338149in}}{\pgfqpoint{2.075407in}{1.342267in}}%
\pgfpathcurveto{\pgfqpoint{2.071289in}{1.346386in}}{\pgfqpoint{2.065703in}{1.348699in}}{\pgfqpoint{2.059879in}{1.348699in}}%
\pgfpathcurveto{\pgfqpoint{2.054055in}{1.348699in}}{\pgfqpoint{2.048469in}{1.346386in}}{\pgfqpoint{2.044351in}{1.342267in}}%
\pgfpathcurveto{\pgfqpoint{2.040232in}{1.338149in}}{\pgfqpoint{2.037919in}{1.332563in}}{\pgfqpoint{2.037919in}{1.326739in}}%
\pgfpathcurveto{\pgfqpoint{2.037919in}{1.320915in}}{\pgfqpoint{2.040232in}{1.315329in}}{\pgfqpoint{2.044351in}{1.311211in}}%
\pgfpathcurveto{\pgfqpoint{2.048469in}{1.307093in}}{\pgfqpoint{2.054055in}{1.304779in}}{\pgfqpoint{2.059879in}{1.304779in}}%
\pgfpathlineto{\pgfqpoint{2.059879in}{1.304779in}}%
\pgfpathclose%
\pgfusepath{stroke,fill}%
\end{pgfscope}%
\begin{pgfscope}%
\pgfpathrectangle{\pgfqpoint{0.100000in}{0.209042in}}{\pgfqpoint{2.906250in}{2.906250in}}%
\pgfusepath{clip}%
\pgfsetbuttcap%
\pgfsetroundjoin%
\definecolor{currentfill}{rgb}{0.697098,0.208501,0.480835}%
\pgfsetfillcolor{currentfill}%
\pgfsetfillopacity{0.800000}%
\pgfsetlinewidth{1.003750pt}%
\definecolor{currentstroke}{rgb}{0.697098,0.208501,0.480835}%
\pgfsetstrokecolor{currentstroke}%
\pgfsetstrokeopacity{0.800000}%
\pgfsetdash{}{0pt}%
\pgfpathmoveto{\pgfqpoint{1.627661in}{1.731604in}}%
\pgfpathcurveto{\pgfqpoint{1.633485in}{1.731604in}}{\pgfqpoint{1.639072in}{1.733918in}}{\pgfqpoint{1.643190in}{1.738036in}}%
\pgfpathcurveto{\pgfqpoint{1.647308in}{1.742154in}}{\pgfqpoint{1.649622in}{1.747741in}}{\pgfqpoint{1.649622in}{1.753565in}}%
\pgfpathcurveto{\pgfqpoint{1.649622in}{1.759388in}}{\pgfqpoint{1.647308in}{1.764975in}}{\pgfqpoint{1.643190in}{1.769093in}}%
\pgfpathcurveto{\pgfqpoint{1.639072in}{1.773211in}}{\pgfqpoint{1.633485in}{1.775525in}}{\pgfqpoint{1.627661in}{1.775525in}}%
\pgfpathcurveto{\pgfqpoint{1.621838in}{1.775525in}}{\pgfqpoint{1.616251in}{1.773211in}}{\pgfqpoint{1.612133in}{1.769093in}}%
\pgfpathcurveto{\pgfqpoint{1.608015in}{1.764975in}}{\pgfqpoint{1.605701in}{1.759388in}}{\pgfqpoint{1.605701in}{1.753565in}}%
\pgfpathcurveto{\pgfqpoint{1.605701in}{1.747741in}}{\pgfqpoint{1.608015in}{1.742154in}}{\pgfqpoint{1.612133in}{1.738036in}}%
\pgfpathcurveto{\pgfqpoint{1.616251in}{1.733918in}}{\pgfqpoint{1.621838in}{1.731604in}}{\pgfqpoint{1.627661in}{1.731604in}}%
\pgfpathlineto{\pgfqpoint{1.627661in}{1.731604in}}%
\pgfpathclose%
\pgfusepath{stroke,fill}%
\end{pgfscope}%
\begin{pgfscope}%
\pgfpathrectangle{\pgfqpoint{0.100000in}{0.209042in}}{\pgfqpoint{2.906250in}{2.906250in}}%
\pgfusepath{clip}%
\pgfsetbuttcap%
\pgfsetroundjoin%
\definecolor{currentfill}{rgb}{0.904281,0.319610,0.388137}%
\pgfsetfillcolor{currentfill}%
\pgfsetfillopacity{0.800000}%
\pgfsetlinewidth{1.003750pt}%
\definecolor{currentstroke}{rgb}{0.904281,0.319610,0.388137}%
\pgfsetstrokecolor{currentstroke}%
\pgfsetstrokeopacity{0.800000}%
\pgfsetdash{}{0pt}%
\pgfpathmoveto{\pgfqpoint{1.264340in}{1.557146in}}%
\pgfpathcurveto{\pgfqpoint{1.270164in}{1.557146in}}{\pgfqpoint{1.275751in}{1.559460in}}{\pgfqpoint{1.279869in}{1.563578in}}%
\pgfpathcurveto{\pgfqpoint{1.283987in}{1.567696in}}{\pgfqpoint{1.286301in}{1.573282in}}{\pgfqpoint{1.286301in}{1.579106in}}%
\pgfpathcurveto{\pgfqpoint{1.286301in}{1.584930in}}{\pgfqpoint{1.283987in}{1.590516in}}{\pgfqpoint{1.279869in}{1.594634in}}%
\pgfpathcurveto{\pgfqpoint{1.275751in}{1.598753in}}{\pgfqpoint{1.270164in}{1.601066in}}{\pgfqpoint{1.264340in}{1.601066in}}%
\pgfpathcurveto{\pgfqpoint{1.258516in}{1.601066in}}{\pgfqpoint{1.252930in}{1.598753in}}{\pgfqpoint{1.248812in}{1.594634in}}%
\pgfpathcurveto{\pgfqpoint{1.244694in}{1.590516in}}{\pgfqpoint{1.242380in}{1.584930in}}{\pgfqpoint{1.242380in}{1.579106in}}%
\pgfpathcurveto{\pgfqpoint{1.242380in}{1.573282in}}{\pgfqpoint{1.244694in}{1.567696in}}{\pgfqpoint{1.248812in}{1.563578in}}%
\pgfpathcurveto{\pgfqpoint{1.252930in}{1.559460in}}{\pgfqpoint{1.258516in}{1.557146in}}{\pgfqpoint{1.264340in}{1.557146in}}%
\pgfpathlineto{\pgfqpoint{1.264340in}{1.557146in}}%
\pgfpathclose%
\pgfusepath{stroke,fill}%
\end{pgfscope}%
\begin{pgfscope}%
\pgfpathrectangle{\pgfqpoint{0.100000in}{0.209042in}}{\pgfqpoint{2.906250in}{2.906250in}}%
\pgfusepath{clip}%
\pgfsetbuttcap%
\pgfsetroundjoin%
\definecolor{currentfill}{rgb}{0.846416,0.272473,0.421631}%
\pgfsetfillcolor{currentfill}%
\pgfsetfillopacity{0.800000}%
\pgfsetlinewidth{1.003750pt}%
\definecolor{currentstroke}{rgb}{0.846416,0.272473,0.421631}%
\pgfsetstrokecolor{currentstroke}%
\pgfsetstrokeopacity{0.800000}%
\pgfsetdash{}{0pt}%
\pgfpathmoveto{\pgfqpoint{1.784493in}{1.356555in}}%
\pgfpathcurveto{\pgfqpoint{1.790317in}{1.356555in}}{\pgfqpoint{1.795903in}{1.358868in}}{\pgfqpoint{1.800021in}{1.362987in}}%
\pgfpathcurveto{\pgfqpoint{1.804139in}{1.367105in}}{\pgfqpoint{1.806453in}{1.372691in}}{\pgfqpoint{1.806453in}{1.378515in}}%
\pgfpathcurveto{\pgfqpoint{1.806453in}{1.384339in}}{\pgfqpoint{1.804139in}{1.389925in}}{\pgfqpoint{1.800021in}{1.394043in}}%
\pgfpathcurveto{\pgfqpoint{1.795903in}{1.398161in}}{\pgfqpoint{1.790317in}{1.400475in}}{\pgfqpoint{1.784493in}{1.400475in}}%
\pgfpathcurveto{\pgfqpoint{1.778669in}{1.400475in}}{\pgfqpoint{1.773083in}{1.398161in}}{\pgfqpoint{1.768964in}{1.394043in}}%
\pgfpathcurveto{\pgfqpoint{1.764846in}{1.389925in}}{\pgfqpoint{1.762532in}{1.384339in}}{\pgfqpoint{1.762532in}{1.378515in}}%
\pgfpathcurveto{\pgfqpoint{1.762532in}{1.372691in}}{\pgfqpoint{1.764846in}{1.367105in}}{\pgfqpoint{1.768964in}{1.362987in}}%
\pgfpathcurveto{\pgfqpoint{1.773083in}{1.358868in}}{\pgfqpoint{1.778669in}{1.356555in}}{\pgfqpoint{1.784493in}{1.356555in}}%
\pgfpathlineto{\pgfqpoint{1.784493in}{1.356555in}}%
\pgfpathclose%
\pgfusepath{stroke,fill}%
\end{pgfscope}%
\begin{pgfscope}%
\pgfpathrectangle{\pgfqpoint{0.100000in}{0.209042in}}{\pgfqpoint{2.906250in}{2.906250in}}%
\pgfusepath{clip}%
\pgfsetbuttcap%
\pgfsetroundjoin%
\definecolor{currentfill}{rgb}{0.984622,0.520713,0.376698}%
\pgfsetfillcolor{currentfill}%
\pgfsetfillopacity{0.800000}%
\pgfsetlinewidth{1.003750pt}%
\definecolor{currentstroke}{rgb}{0.984622,0.520713,0.376698}%
\pgfsetstrokecolor{currentstroke}%
\pgfsetstrokeopacity{0.800000}%
\pgfsetdash{}{0pt}%
\pgfpathmoveto{\pgfqpoint{1.412903in}{1.201842in}}%
\pgfpathcurveto{\pgfqpoint{1.418727in}{1.201842in}}{\pgfqpoint{1.424314in}{1.204156in}}{\pgfqpoint{1.428432in}{1.208274in}}%
\pgfpathcurveto{\pgfqpoint{1.432550in}{1.212392in}}{\pgfqpoint{1.434864in}{1.217978in}}{\pgfqpoint{1.434864in}{1.223802in}}%
\pgfpathcurveto{\pgfqpoint{1.434864in}{1.229626in}}{\pgfqpoint{1.432550in}{1.235212in}}{\pgfqpoint{1.428432in}{1.239331in}}%
\pgfpathcurveto{\pgfqpoint{1.424314in}{1.243449in}}{\pgfqpoint{1.418727in}{1.245763in}}{\pgfqpoint{1.412903in}{1.245763in}}%
\pgfpathcurveto{\pgfqpoint{1.407079in}{1.245763in}}{\pgfqpoint{1.401493in}{1.243449in}}{\pgfqpoint{1.397375in}{1.239331in}}%
\pgfpathcurveto{\pgfqpoint{1.393257in}{1.235212in}}{\pgfqpoint{1.390943in}{1.229626in}}{\pgfqpoint{1.390943in}{1.223802in}}%
\pgfpathcurveto{\pgfqpoint{1.390943in}{1.217978in}}{\pgfqpoint{1.393257in}{1.212392in}}{\pgfqpoint{1.397375in}{1.208274in}}%
\pgfpathcurveto{\pgfqpoint{1.401493in}{1.204156in}}{\pgfqpoint{1.407079in}{1.201842in}}{\pgfqpoint{1.412903in}{1.201842in}}%
\pgfpathlineto{\pgfqpoint{1.412903in}{1.201842in}}%
\pgfpathclose%
\pgfusepath{stroke,fill}%
\end{pgfscope}%
\begin{pgfscope}%
\pgfpathrectangle{\pgfqpoint{0.100000in}{0.209042in}}{\pgfqpoint{2.906250in}{2.906250in}}%
\pgfusepath{clip}%
\pgfsetbuttcap%
\pgfsetroundjoin%
\definecolor{currentfill}{rgb}{0.904281,0.319610,0.388137}%
\pgfsetfillcolor{currentfill}%
\pgfsetfillopacity{0.800000}%
\pgfsetlinewidth{1.003750pt}%
\definecolor{currentstroke}{rgb}{0.904281,0.319610,0.388137}%
\pgfsetstrokecolor{currentstroke}%
\pgfsetstrokeopacity{0.800000}%
\pgfsetdash{}{0pt}%
\pgfpathmoveto{\pgfqpoint{1.355604in}{1.841943in}}%
\pgfpathcurveto{\pgfqpoint{1.361428in}{1.841943in}}{\pgfqpoint{1.367014in}{1.844257in}}{\pgfqpoint{1.371132in}{1.848375in}}%
\pgfpathcurveto{\pgfqpoint{1.375250in}{1.852493in}}{\pgfqpoint{1.377564in}{1.858080in}}{\pgfqpoint{1.377564in}{1.863904in}}%
\pgfpathcurveto{\pgfqpoint{1.377564in}{1.869727in}}{\pgfqpoint{1.375250in}{1.875314in}}{\pgfqpoint{1.371132in}{1.879432in}}%
\pgfpathcurveto{\pgfqpoint{1.367014in}{1.883550in}}{\pgfqpoint{1.361428in}{1.885864in}}{\pgfqpoint{1.355604in}{1.885864in}}%
\pgfpathcurveto{\pgfqpoint{1.349780in}{1.885864in}}{\pgfqpoint{1.344194in}{1.883550in}}{\pgfqpoint{1.340076in}{1.879432in}}%
\pgfpathcurveto{\pgfqpoint{1.335958in}{1.875314in}}{\pgfqpoint{1.333644in}{1.869727in}}{\pgfqpoint{1.333644in}{1.863904in}}%
\pgfpathcurveto{\pgfqpoint{1.333644in}{1.858080in}}{\pgfqpoint{1.335958in}{1.852493in}}{\pgfqpoint{1.340076in}{1.848375in}}%
\pgfpathcurveto{\pgfqpoint{1.344194in}{1.844257in}}{\pgfqpoint{1.349780in}{1.841943in}}{\pgfqpoint{1.355604in}{1.841943in}}%
\pgfpathlineto{\pgfqpoint{1.355604in}{1.841943in}}%
\pgfpathclose%
\pgfusepath{stroke,fill}%
\end{pgfscope}%
\begin{pgfscope}%
\pgfpathrectangle{\pgfqpoint{0.100000in}{0.209042in}}{\pgfqpoint{2.906250in}{2.906250in}}%
\pgfusepath{clip}%
\pgfsetbuttcap%
\pgfsetroundjoin%
\definecolor{currentfill}{rgb}{0.944006,0.377643,0.365136}%
\pgfsetfillcolor{currentfill}%
\pgfsetfillopacity{0.800000}%
\pgfsetlinewidth{1.003750pt}%
\definecolor{currentstroke}{rgb}{0.944006,0.377643,0.365136}%
\pgfsetstrokecolor{currentstroke}%
\pgfsetstrokeopacity{0.800000}%
\pgfsetdash{}{0pt}%
\pgfpathmoveto{\pgfqpoint{1.226659in}{1.580608in}}%
\pgfpathcurveto{\pgfqpoint{1.232483in}{1.580608in}}{\pgfqpoint{1.238069in}{1.582922in}}{\pgfqpoint{1.242187in}{1.587040in}}%
\pgfpathcurveto{\pgfqpoint{1.246305in}{1.591158in}}{\pgfqpoint{1.248619in}{1.596744in}}{\pgfqpoint{1.248619in}{1.602568in}}%
\pgfpathcurveto{\pgfqpoint{1.248619in}{1.608392in}}{\pgfqpoint{1.246305in}{1.613979in}}{\pgfqpoint{1.242187in}{1.618097in}}%
\pgfpathcurveto{\pgfqpoint{1.238069in}{1.622215in}}{\pgfqpoint{1.232483in}{1.624529in}}{\pgfqpoint{1.226659in}{1.624529in}}%
\pgfpathcurveto{\pgfqpoint{1.220835in}{1.624529in}}{\pgfqpoint{1.215248in}{1.622215in}}{\pgfqpoint{1.211130in}{1.618097in}}%
\pgfpathcurveto{\pgfqpoint{1.207012in}{1.613979in}}{\pgfqpoint{1.204698in}{1.608392in}}{\pgfqpoint{1.204698in}{1.602568in}}%
\pgfpathcurveto{\pgfqpoint{1.204698in}{1.596744in}}{\pgfqpoint{1.207012in}{1.591158in}}{\pgfqpoint{1.211130in}{1.587040in}}%
\pgfpathcurveto{\pgfqpoint{1.215248in}{1.582922in}}{\pgfqpoint{1.220835in}{1.580608in}}{\pgfqpoint{1.226659in}{1.580608in}}%
\pgfpathlineto{\pgfqpoint{1.226659in}{1.580608in}}%
\pgfpathclose%
\pgfusepath{stroke,fill}%
\end{pgfscope}%
\begin{pgfscope}%
\pgfpathrectangle{\pgfqpoint{0.100000in}{0.209042in}}{\pgfqpoint{2.906250in}{2.906250in}}%
\pgfusepath{clip}%
\pgfsetbuttcap%
\pgfsetroundjoin%
\definecolor{currentfill}{rgb}{0.953099,0.397563,0.361438}%
\pgfsetfillcolor{currentfill}%
\pgfsetfillopacity{0.800000}%
\pgfsetlinewidth{1.003750pt}%
\definecolor{currentstroke}{rgb}{0.953099,0.397563,0.361438}%
\pgfsetstrokecolor{currentstroke}%
\pgfsetstrokeopacity{0.800000}%
\pgfsetdash{}{0pt}%
\pgfpathmoveto{\pgfqpoint{1.415384in}{1.953243in}}%
\pgfpathcurveto{\pgfqpoint{1.421207in}{1.953243in}}{\pgfqpoint{1.426794in}{1.955557in}}{\pgfqpoint{1.430912in}{1.959675in}}%
\pgfpathcurveto{\pgfqpoint{1.435030in}{1.963793in}}{\pgfqpoint{1.437344in}{1.969379in}}{\pgfqpoint{1.437344in}{1.975203in}}%
\pgfpathcurveto{\pgfqpoint{1.437344in}{1.981027in}}{\pgfqpoint{1.435030in}{1.986613in}}{\pgfqpoint{1.430912in}{1.990732in}}%
\pgfpathcurveto{\pgfqpoint{1.426794in}{1.994850in}}{\pgfqpoint{1.421207in}{1.997164in}}{\pgfqpoint{1.415384in}{1.997164in}}%
\pgfpathcurveto{\pgfqpoint{1.409560in}{1.997164in}}{\pgfqpoint{1.403973in}{1.994850in}}{\pgfqpoint{1.399855in}{1.990732in}}%
\pgfpathcurveto{\pgfqpoint{1.395737in}{1.986613in}}{\pgfqpoint{1.393423in}{1.981027in}}{\pgfqpoint{1.393423in}{1.975203in}}%
\pgfpathcurveto{\pgfqpoint{1.393423in}{1.969379in}}{\pgfqpoint{1.395737in}{1.963793in}}{\pgfqpoint{1.399855in}{1.959675in}}%
\pgfpathcurveto{\pgfqpoint{1.403973in}{1.955557in}}{\pgfqpoint{1.409560in}{1.953243in}}{\pgfqpoint{1.415384in}{1.953243in}}%
\pgfpathlineto{\pgfqpoint{1.415384in}{1.953243in}}%
\pgfpathclose%
\pgfusepath{stroke,fill}%
\end{pgfscope}%
\begin{pgfscope}%
\pgfpathrectangle{\pgfqpoint{0.100000in}{0.209042in}}{\pgfqpoint{2.906250in}{2.906250in}}%
\pgfusepath{clip}%
\pgfsetbuttcap%
\pgfsetroundjoin%
\definecolor{currentfill}{rgb}{0.863320,0.283729,0.412403}%
\pgfsetfillcolor{currentfill}%
\pgfsetfillopacity{0.800000}%
\pgfsetlinewidth{1.003750pt}%
\definecolor{currentstroke}{rgb}{0.863320,0.283729,0.412403}%
\pgfsetstrokecolor{currentstroke}%
\pgfsetstrokeopacity{0.800000}%
\pgfsetdash{}{0pt}%
\pgfpathmoveto{\pgfqpoint{1.436723in}{1.523312in}}%
\pgfpathcurveto{\pgfqpoint{1.442547in}{1.523312in}}{\pgfqpoint{1.448133in}{1.525625in}}{\pgfqpoint{1.452251in}{1.529744in}}%
\pgfpathcurveto{\pgfqpoint{1.456369in}{1.533862in}}{\pgfqpoint{1.458683in}{1.539448in}}{\pgfqpoint{1.458683in}{1.545272in}}%
\pgfpathcurveto{\pgfqpoint{1.458683in}{1.551096in}}{\pgfqpoint{1.456369in}{1.556682in}}{\pgfqpoint{1.452251in}{1.560800in}}%
\pgfpathcurveto{\pgfqpoint{1.448133in}{1.564918in}}{\pgfqpoint{1.442547in}{1.567232in}}{\pgfqpoint{1.436723in}{1.567232in}}%
\pgfpathcurveto{\pgfqpoint{1.430899in}{1.567232in}}{\pgfqpoint{1.425313in}{1.564918in}}{\pgfqpoint{1.421194in}{1.560800in}}%
\pgfpathcurveto{\pgfqpoint{1.417076in}{1.556682in}}{\pgfqpoint{1.414762in}{1.551096in}}{\pgfqpoint{1.414762in}{1.545272in}}%
\pgfpathcurveto{\pgfqpoint{1.414762in}{1.539448in}}{\pgfqpoint{1.417076in}{1.533862in}}{\pgfqpoint{1.421194in}{1.529744in}}%
\pgfpathcurveto{\pgfqpoint{1.425313in}{1.525625in}}{\pgfqpoint{1.430899in}{1.523312in}}{\pgfqpoint{1.436723in}{1.523312in}}%
\pgfpathlineto{\pgfqpoint{1.436723in}{1.523312in}}%
\pgfpathclose%
\pgfusepath{stroke,fill}%
\end{pgfscope}%
\begin{pgfscope}%
\pgfpathrectangle{\pgfqpoint{0.100000in}{0.209042in}}{\pgfqpoint{2.906250in}{2.906250in}}%
\pgfusepath{clip}%
\pgfsetbuttcap%
\pgfsetroundjoin%
\definecolor{currentfill}{rgb}{0.899552,0.314616,0.391037}%
\pgfsetfillcolor{currentfill}%
\pgfsetfillopacity{0.800000}%
\pgfsetlinewidth{1.003750pt}%
\definecolor{currentstroke}{rgb}{0.899552,0.314616,0.391037}%
\pgfsetstrokecolor{currentstroke}%
\pgfsetstrokeopacity{0.800000}%
\pgfsetdash{}{0pt}%
\pgfpathmoveto{\pgfqpoint{1.757674in}{1.376561in}}%
\pgfpathcurveto{\pgfqpoint{1.763497in}{1.376561in}}{\pgfqpoint{1.769084in}{1.378874in}}{\pgfqpoint{1.773202in}{1.382993in}}%
\pgfpathcurveto{\pgfqpoint{1.777320in}{1.387111in}}{\pgfqpoint{1.779634in}{1.392697in}}{\pgfqpoint{1.779634in}{1.398521in}}%
\pgfpathcurveto{\pgfqpoint{1.779634in}{1.404345in}}{\pgfqpoint{1.777320in}{1.409931in}}{\pgfqpoint{1.773202in}{1.414049in}}%
\pgfpathcurveto{\pgfqpoint{1.769084in}{1.418167in}}{\pgfqpoint{1.763497in}{1.420481in}}{\pgfqpoint{1.757674in}{1.420481in}}%
\pgfpathcurveto{\pgfqpoint{1.751850in}{1.420481in}}{\pgfqpoint{1.746263in}{1.418167in}}{\pgfqpoint{1.742145in}{1.414049in}}%
\pgfpathcurveto{\pgfqpoint{1.738027in}{1.409931in}}{\pgfqpoint{1.735713in}{1.404345in}}{\pgfqpoint{1.735713in}{1.398521in}}%
\pgfpathcurveto{\pgfqpoint{1.735713in}{1.392697in}}{\pgfqpoint{1.738027in}{1.387111in}}{\pgfqpoint{1.742145in}{1.382993in}}%
\pgfpathcurveto{\pgfqpoint{1.746263in}{1.378874in}}{\pgfqpoint{1.751850in}{1.376561in}}{\pgfqpoint{1.757674in}{1.376561in}}%
\pgfpathlineto{\pgfqpoint{1.757674in}{1.376561in}}%
\pgfpathclose%
\pgfusepath{stroke,fill}%
\end{pgfscope}%
\begin{pgfscope}%
\pgfpathrectangle{\pgfqpoint{0.100000in}{0.209042in}}{\pgfqpoint{2.906250in}{2.906250in}}%
\pgfusepath{clip}%
\pgfsetbuttcap%
\pgfsetroundjoin%
\definecolor{currentfill}{rgb}{0.925937,0.346844,0.374959}%
\pgfsetfillcolor{currentfill}%
\pgfsetfillopacity{0.800000}%
\pgfsetlinewidth{1.003750pt}%
\definecolor{currentstroke}{rgb}{0.925937,0.346844,0.374959}%
\pgfsetstrokecolor{currentstroke}%
\pgfsetstrokeopacity{0.800000}%
\pgfsetdash{}{0pt}%
\pgfpathmoveto{\pgfqpoint{1.955619in}{1.642330in}}%
\pgfpathcurveto{\pgfqpoint{1.961443in}{1.642330in}}{\pgfqpoint{1.967029in}{1.644644in}}{\pgfqpoint{1.971148in}{1.648762in}}%
\pgfpathcurveto{\pgfqpoint{1.975266in}{1.652880in}}{\pgfqpoint{1.977580in}{1.658466in}}{\pgfqpoint{1.977580in}{1.664290in}}%
\pgfpathcurveto{\pgfqpoint{1.977580in}{1.670114in}}{\pgfqpoint{1.975266in}{1.675700in}}{\pgfqpoint{1.971148in}{1.679818in}}%
\pgfpathcurveto{\pgfqpoint{1.967029in}{1.683936in}}{\pgfqpoint{1.961443in}{1.686250in}}{\pgfqpoint{1.955619in}{1.686250in}}%
\pgfpathcurveto{\pgfqpoint{1.949795in}{1.686250in}}{\pgfqpoint{1.944209in}{1.683936in}}{\pgfqpoint{1.940091in}{1.679818in}}%
\pgfpathcurveto{\pgfqpoint{1.935973in}{1.675700in}}{\pgfqpoint{1.933659in}{1.670114in}}{\pgfqpoint{1.933659in}{1.664290in}}%
\pgfpathcurveto{\pgfqpoint{1.933659in}{1.658466in}}{\pgfqpoint{1.935973in}{1.652880in}}{\pgfqpoint{1.940091in}{1.648762in}}%
\pgfpathcurveto{\pgfqpoint{1.944209in}{1.644644in}}{\pgfqpoint{1.949795in}{1.642330in}}{\pgfqpoint{1.955619in}{1.642330in}}%
\pgfpathlineto{\pgfqpoint{1.955619in}{1.642330in}}%
\pgfpathclose%
\pgfusepath{stroke,fill}%
\end{pgfscope}%
\begin{pgfscope}%
\pgfpathrectangle{\pgfqpoint{0.100000in}{0.209042in}}{\pgfqpoint{2.906250in}{2.906250in}}%
\pgfusepath{clip}%
\pgfsetbuttcap%
\pgfsetroundjoin%
\definecolor{currentfill}{rgb}{0.995932,0.805527,0.566202}%
\pgfsetfillcolor{currentfill}%
\pgfsetfillopacity{0.800000}%
\pgfsetlinewidth{1.003750pt}%
\definecolor{currentstroke}{rgb}{0.995932,0.805527,0.566202}%
\pgfsetstrokecolor{currentstroke}%
\pgfsetstrokeopacity{0.800000}%
\pgfsetdash{}{0pt}%
\pgfpathmoveto{\pgfqpoint{2.080133in}{1.211584in}}%
\pgfpathcurveto{\pgfqpoint{2.085957in}{1.211584in}}{\pgfqpoint{2.091544in}{1.213897in}}{\pgfqpoint{2.095662in}{1.218016in}}%
\pgfpathcurveto{\pgfqpoint{2.099780in}{1.222134in}}{\pgfqpoint{2.102094in}{1.227720in}}{\pgfqpoint{2.102094in}{1.233544in}}%
\pgfpathcurveto{\pgfqpoint{2.102094in}{1.239368in}}{\pgfqpoint{2.099780in}{1.244954in}}{\pgfqpoint{2.095662in}{1.249072in}}%
\pgfpathcurveto{\pgfqpoint{2.091544in}{1.253190in}}{\pgfqpoint{2.085957in}{1.255504in}}{\pgfqpoint{2.080133in}{1.255504in}}%
\pgfpathcurveto{\pgfqpoint{2.074309in}{1.255504in}}{\pgfqpoint{2.068723in}{1.253190in}}{\pgfqpoint{2.064605in}{1.249072in}}%
\pgfpathcurveto{\pgfqpoint{2.060487in}{1.244954in}}{\pgfqpoint{2.058173in}{1.239368in}}{\pgfqpoint{2.058173in}{1.233544in}}%
\pgfpathcurveto{\pgfqpoint{2.058173in}{1.227720in}}{\pgfqpoint{2.060487in}{1.222134in}}{\pgfqpoint{2.064605in}{1.218016in}}%
\pgfpathcurveto{\pgfqpoint{2.068723in}{1.213897in}}{\pgfqpoint{2.074309in}{1.211584in}}{\pgfqpoint{2.080133in}{1.211584in}}%
\pgfpathlineto{\pgfqpoint{2.080133in}{1.211584in}}%
\pgfpathclose%
\pgfusepath{stroke,fill}%
\end{pgfscope}%
\begin{pgfscope}%
\pgfpathrectangle{\pgfqpoint{0.100000in}{0.209042in}}{\pgfqpoint{2.906250in}{2.906250in}}%
\pgfusepath{clip}%
\pgfsetbuttcap%
\pgfsetroundjoin%
\definecolor{currentfill}{rgb}{0.995131,0.827052,0.585701}%
\pgfsetfillcolor{currentfill}%
\pgfsetfillopacity{0.800000}%
\pgfsetlinewidth{1.003750pt}%
\definecolor{currentstroke}{rgb}{0.995131,0.827052,0.585701}%
\pgfsetstrokecolor{currentstroke}%
\pgfsetstrokeopacity{0.800000}%
\pgfsetdash{}{0pt}%
\pgfpathmoveto{\pgfqpoint{1.929718in}{1.899939in}}%
\pgfpathcurveto{\pgfqpoint{1.935542in}{1.899939in}}{\pgfqpoint{1.941128in}{1.902253in}}{\pgfqpoint{1.945247in}{1.906371in}}%
\pgfpathcurveto{\pgfqpoint{1.949365in}{1.910489in}}{\pgfqpoint{1.951679in}{1.916075in}}{\pgfqpoint{1.951679in}{1.921899in}}%
\pgfpathcurveto{\pgfqpoint{1.951679in}{1.927723in}}{\pgfqpoint{1.949365in}{1.933309in}}{\pgfqpoint{1.945247in}{1.937428in}}%
\pgfpathcurveto{\pgfqpoint{1.941128in}{1.941546in}}{\pgfqpoint{1.935542in}{1.943860in}}{\pgfqpoint{1.929718in}{1.943860in}}%
\pgfpathcurveto{\pgfqpoint{1.923894in}{1.943860in}}{\pgfqpoint{1.918308in}{1.941546in}}{\pgfqpoint{1.914190in}{1.937428in}}%
\pgfpathcurveto{\pgfqpoint{1.910072in}{1.933309in}}{\pgfqpoint{1.907758in}{1.927723in}}{\pgfqpoint{1.907758in}{1.921899in}}%
\pgfpathcurveto{\pgfqpoint{1.907758in}{1.916075in}}{\pgfqpoint{1.910072in}{1.910489in}}{\pgfqpoint{1.914190in}{1.906371in}}%
\pgfpathcurveto{\pgfqpoint{1.918308in}{1.902253in}}{\pgfqpoint{1.923894in}{1.899939in}}{\pgfqpoint{1.929718in}{1.899939in}}%
\pgfpathlineto{\pgfqpoint{1.929718in}{1.899939in}}%
\pgfpathclose%
\pgfusepath{stroke,fill}%
\end{pgfscope}%
\begin{pgfscope}%
\pgfpathrectangle{\pgfqpoint{0.100000in}{0.209042in}}{\pgfqpoint{2.906250in}{2.906250in}}%
\pgfusepath{clip}%
\pgfsetbuttcap%
\pgfsetroundjoin%
\definecolor{currentfill}{rgb}{0.989815,0.934329,0.690198}%
\pgfsetfillcolor{currentfill}%
\pgfsetfillopacity{0.800000}%
\pgfsetlinewidth{1.003750pt}%
\definecolor{currentstroke}{rgb}{0.989815,0.934329,0.690198}%
\pgfsetstrokecolor{currentstroke}%
\pgfsetstrokeopacity{0.800000}%
\pgfsetdash{}{0pt}%
\pgfpathmoveto{\pgfqpoint{1.542706in}{1.275574in}}%
\pgfpathcurveto{\pgfqpoint{1.548530in}{1.275574in}}{\pgfqpoint{1.554116in}{1.277887in}}{\pgfqpoint{1.558235in}{1.282006in}}%
\pgfpathcurveto{\pgfqpoint{1.562353in}{1.286124in}}{\pgfqpoint{1.564667in}{1.291710in}}{\pgfqpoint{1.564667in}{1.297534in}}%
\pgfpathcurveto{\pgfqpoint{1.564667in}{1.303358in}}{\pgfqpoint{1.562353in}{1.308944in}}{\pgfqpoint{1.558235in}{1.313062in}}%
\pgfpathcurveto{\pgfqpoint{1.554116in}{1.317180in}}{\pgfqpoint{1.548530in}{1.319494in}}{\pgfqpoint{1.542706in}{1.319494in}}%
\pgfpathcurveto{\pgfqpoint{1.536882in}{1.319494in}}{\pgfqpoint{1.531296in}{1.317180in}}{\pgfqpoint{1.527178in}{1.313062in}}%
\pgfpathcurveto{\pgfqpoint{1.523060in}{1.308944in}}{\pgfqpoint{1.520746in}{1.303358in}}{\pgfqpoint{1.520746in}{1.297534in}}%
\pgfpathcurveto{\pgfqpoint{1.520746in}{1.291710in}}{\pgfqpoint{1.523060in}{1.286124in}}{\pgfqpoint{1.527178in}{1.282006in}}%
\pgfpathcurveto{\pgfqpoint{1.531296in}{1.277887in}}{\pgfqpoint{1.536882in}{1.275574in}}{\pgfqpoint{1.542706in}{1.275574in}}%
\pgfpathlineto{\pgfqpoint{1.542706in}{1.275574in}}%
\pgfpathclose%
\pgfusepath{stroke,fill}%
\end{pgfscope}%
\begin{pgfscope}%
\definecolor{textcolor}{rgb}{0.000000,0.000000,0.000000}%
\pgfsetstrokecolor{textcolor}%
\pgfsetfillcolor{textcolor}%
\pgftext[x=1.553125in,y=3.198625in,,base]{\color{textcolor}{\rmfamily\fontsize{12.000000}{14.400000}\selectfont\catcode`\^=\active\def^{\ifmmode\sp\else\^{}\fi}\catcode`\%=\active\def%{\%}Nube de Puntos Multivariable}}%
\end{pgfscope}%
\begin{pgfscope}%
\pgfsetbuttcap%
\pgfsetmiterjoin%
\definecolor{currentfill}{rgb}{1.000000,1.000000,1.000000}%
\pgfsetfillcolor{currentfill}%
\pgfsetlinewidth{0.000000pt}%
\definecolor{currentstroke}{rgb}{0.000000,0.000000,0.000000}%
\pgfsetstrokecolor{currentstroke}%
\pgfsetstrokeopacity{0.000000}%
\pgfsetdash{}{0pt}%
\pgfpathmoveto{\pgfqpoint{3.393750in}{0.699667in}}%
\pgfpathlineto{\pgfqpoint{3.490000in}{0.699667in}}%
\pgfpathlineto{\pgfqpoint{3.490000in}{2.624667in}}%
\pgfpathlineto{\pgfqpoint{3.393750in}{2.624667in}}%
\pgfpathlineto{\pgfqpoint{3.393750in}{0.699667in}}%
\pgfpathclose%
\pgfusepath{fill}%
\end{pgfscope}%
\begin{pgfscope}%
\pgfpathrectangle{\pgfqpoint{3.393750in}{0.699667in}}{\pgfqpoint{0.096250in}{1.925000in}}%
\pgfusepath{clip}%
\pgfsetbuttcap%
\pgfsetroundjoin%
\definecolor{currentfill}{rgb}{0.001462,0.000466,0.013866}%
\pgfsetfillcolor{currentfill}%
\pgfsetfillopacity{0.800000}%
\pgfsetlinewidth{0.000000pt}%
\definecolor{currentstroke}{rgb}{0.000000,0.000000,0.000000}%
\pgfsetstrokecolor{currentstroke}%
\pgfsetdash{}{0pt}%
\pgfpathmoveto{\pgfqpoint{3.393750in}{0.699667in}}%
\pgfpathlineto{\pgfqpoint{3.490000in}{0.699667in}}%
\pgfpathlineto{\pgfqpoint{3.490000in}{0.707187in}}%
\pgfpathlineto{\pgfqpoint{3.393750in}{0.707187in}}%
\pgfpathlineto{\pgfqpoint{3.393750in}{0.699667in}}%
\pgfusepath{fill}%
\end{pgfscope}%
\begin{pgfscope}%
\pgfpathrectangle{\pgfqpoint{3.393750in}{0.699667in}}{\pgfqpoint{0.096250in}{1.925000in}}%
\pgfusepath{clip}%
\pgfsetbuttcap%
\pgfsetroundjoin%
\definecolor{currentfill}{rgb}{0.002258,0.001295,0.018331}%
\pgfsetfillcolor{currentfill}%
\pgfsetfillopacity{0.800000}%
\pgfsetlinewidth{0.000000pt}%
\definecolor{currentstroke}{rgb}{0.000000,0.000000,0.000000}%
\pgfsetstrokecolor{currentstroke}%
\pgfsetdash{}{0pt}%
\pgfpathmoveto{\pgfqpoint{3.393750in}{0.707187in}}%
\pgfpathlineto{\pgfqpoint{3.490000in}{0.707187in}}%
\pgfpathlineto{\pgfqpoint{3.490000in}{0.714706in}}%
\pgfpathlineto{\pgfqpoint{3.393750in}{0.714706in}}%
\pgfpathlineto{\pgfqpoint{3.393750in}{0.707187in}}%
\pgfusepath{fill}%
\end{pgfscope}%
\begin{pgfscope}%
\pgfpathrectangle{\pgfqpoint{3.393750in}{0.699667in}}{\pgfqpoint{0.096250in}{1.925000in}}%
\pgfusepath{clip}%
\pgfsetbuttcap%
\pgfsetroundjoin%
\definecolor{currentfill}{rgb}{0.003279,0.002305,0.023708}%
\pgfsetfillcolor{currentfill}%
\pgfsetfillopacity{0.800000}%
\pgfsetlinewidth{0.000000pt}%
\definecolor{currentstroke}{rgb}{0.000000,0.000000,0.000000}%
\pgfsetstrokecolor{currentstroke}%
\pgfsetdash{}{0pt}%
\pgfpathmoveto{\pgfqpoint{3.393750in}{0.714706in}}%
\pgfpathlineto{\pgfqpoint{3.490000in}{0.714706in}}%
\pgfpathlineto{\pgfqpoint{3.490000in}{0.722226in}}%
\pgfpathlineto{\pgfqpoint{3.393750in}{0.722226in}}%
\pgfpathlineto{\pgfqpoint{3.393750in}{0.714706in}}%
\pgfusepath{fill}%
\end{pgfscope}%
\begin{pgfscope}%
\pgfpathrectangle{\pgfqpoint{3.393750in}{0.699667in}}{\pgfqpoint{0.096250in}{1.925000in}}%
\pgfusepath{clip}%
\pgfsetbuttcap%
\pgfsetroundjoin%
\definecolor{currentfill}{rgb}{0.004512,0.003490,0.029965}%
\pgfsetfillcolor{currentfill}%
\pgfsetfillopacity{0.800000}%
\pgfsetlinewidth{0.000000pt}%
\definecolor{currentstroke}{rgb}{0.000000,0.000000,0.000000}%
\pgfsetstrokecolor{currentstroke}%
\pgfsetdash{}{0pt}%
\pgfpathmoveto{\pgfqpoint{3.393750in}{0.722226in}}%
\pgfpathlineto{\pgfqpoint{3.490000in}{0.722226in}}%
\pgfpathlineto{\pgfqpoint{3.490000in}{0.729745in}}%
\pgfpathlineto{\pgfqpoint{3.393750in}{0.729745in}}%
\pgfpathlineto{\pgfqpoint{3.393750in}{0.722226in}}%
\pgfusepath{fill}%
\end{pgfscope}%
\begin{pgfscope}%
\pgfpathrectangle{\pgfqpoint{3.393750in}{0.699667in}}{\pgfqpoint{0.096250in}{1.925000in}}%
\pgfusepath{clip}%
\pgfsetbuttcap%
\pgfsetroundjoin%
\definecolor{currentfill}{rgb}{0.005950,0.004843,0.037130}%
\pgfsetfillcolor{currentfill}%
\pgfsetfillopacity{0.800000}%
\pgfsetlinewidth{0.000000pt}%
\definecolor{currentstroke}{rgb}{0.000000,0.000000,0.000000}%
\pgfsetstrokecolor{currentstroke}%
\pgfsetdash{}{0pt}%
\pgfpathmoveto{\pgfqpoint{3.393750in}{0.729745in}}%
\pgfpathlineto{\pgfqpoint{3.490000in}{0.729745in}}%
\pgfpathlineto{\pgfqpoint{3.490000in}{0.737265in}}%
\pgfpathlineto{\pgfqpoint{3.393750in}{0.737265in}}%
\pgfpathlineto{\pgfqpoint{3.393750in}{0.729745in}}%
\pgfusepath{fill}%
\end{pgfscope}%
\begin{pgfscope}%
\pgfpathrectangle{\pgfqpoint{3.393750in}{0.699667in}}{\pgfqpoint{0.096250in}{1.925000in}}%
\pgfusepath{clip}%
\pgfsetbuttcap%
\pgfsetroundjoin%
\definecolor{currentfill}{rgb}{0.007588,0.006356,0.044973}%
\pgfsetfillcolor{currentfill}%
\pgfsetfillopacity{0.800000}%
\pgfsetlinewidth{0.000000pt}%
\definecolor{currentstroke}{rgb}{0.000000,0.000000,0.000000}%
\pgfsetstrokecolor{currentstroke}%
\pgfsetdash{}{0pt}%
\pgfpathmoveto{\pgfqpoint{3.393750in}{0.737265in}}%
\pgfpathlineto{\pgfqpoint{3.490000in}{0.737265in}}%
\pgfpathlineto{\pgfqpoint{3.490000in}{0.744784in}}%
\pgfpathlineto{\pgfqpoint{3.393750in}{0.744784in}}%
\pgfpathlineto{\pgfqpoint{3.393750in}{0.737265in}}%
\pgfusepath{fill}%
\end{pgfscope}%
\begin{pgfscope}%
\pgfpathrectangle{\pgfqpoint{3.393750in}{0.699667in}}{\pgfqpoint{0.096250in}{1.925000in}}%
\pgfusepath{clip}%
\pgfsetbuttcap%
\pgfsetroundjoin%
\definecolor{currentfill}{rgb}{0.009426,0.008022,0.052844}%
\pgfsetfillcolor{currentfill}%
\pgfsetfillopacity{0.800000}%
\pgfsetlinewidth{0.000000pt}%
\definecolor{currentstroke}{rgb}{0.000000,0.000000,0.000000}%
\pgfsetstrokecolor{currentstroke}%
\pgfsetdash{}{0pt}%
\pgfpathmoveto{\pgfqpoint{3.393750in}{0.744784in}}%
\pgfpathlineto{\pgfqpoint{3.490000in}{0.744784in}}%
\pgfpathlineto{\pgfqpoint{3.490000in}{0.752304in}}%
\pgfpathlineto{\pgfqpoint{3.393750in}{0.752304in}}%
\pgfpathlineto{\pgfqpoint{3.393750in}{0.744784in}}%
\pgfusepath{fill}%
\end{pgfscope}%
\begin{pgfscope}%
\pgfpathrectangle{\pgfqpoint{3.393750in}{0.699667in}}{\pgfqpoint{0.096250in}{1.925000in}}%
\pgfusepath{clip}%
\pgfsetbuttcap%
\pgfsetroundjoin%
\definecolor{currentfill}{rgb}{0.011465,0.009828,0.060750}%
\pgfsetfillcolor{currentfill}%
\pgfsetfillopacity{0.800000}%
\pgfsetlinewidth{0.000000pt}%
\definecolor{currentstroke}{rgb}{0.000000,0.000000,0.000000}%
\pgfsetstrokecolor{currentstroke}%
\pgfsetdash{}{0pt}%
\pgfpathmoveto{\pgfqpoint{3.393750in}{0.752304in}}%
\pgfpathlineto{\pgfqpoint{3.490000in}{0.752304in}}%
\pgfpathlineto{\pgfqpoint{3.490000in}{0.759823in}}%
\pgfpathlineto{\pgfqpoint{3.393750in}{0.759823in}}%
\pgfpathlineto{\pgfqpoint{3.393750in}{0.752304in}}%
\pgfusepath{fill}%
\end{pgfscope}%
\begin{pgfscope}%
\pgfpathrectangle{\pgfqpoint{3.393750in}{0.699667in}}{\pgfqpoint{0.096250in}{1.925000in}}%
\pgfusepath{clip}%
\pgfsetbuttcap%
\pgfsetroundjoin%
\definecolor{currentfill}{rgb}{0.013708,0.011771,0.068667}%
\pgfsetfillcolor{currentfill}%
\pgfsetfillopacity{0.800000}%
\pgfsetlinewidth{0.000000pt}%
\definecolor{currentstroke}{rgb}{0.000000,0.000000,0.000000}%
\pgfsetstrokecolor{currentstroke}%
\pgfsetdash{}{0pt}%
\pgfpathmoveto{\pgfqpoint{3.393750in}{0.759823in}}%
\pgfpathlineto{\pgfqpoint{3.490000in}{0.759823in}}%
\pgfpathlineto{\pgfqpoint{3.490000in}{0.767343in}}%
\pgfpathlineto{\pgfqpoint{3.393750in}{0.767343in}}%
\pgfpathlineto{\pgfqpoint{3.393750in}{0.759823in}}%
\pgfusepath{fill}%
\end{pgfscope}%
\begin{pgfscope}%
\pgfpathrectangle{\pgfqpoint{3.393750in}{0.699667in}}{\pgfqpoint{0.096250in}{1.925000in}}%
\pgfusepath{clip}%
\pgfsetbuttcap%
\pgfsetroundjoin%
\definecolor{currentfill}{rgb}{0.016156,0.013840,0.076603}%
\pgfsetfillcolor{currentfill}%
\pgfsetfillopacity{0.800000}%
\pgfsetlinewidth{0.000000pt}%
\definecolor{currentstroke}{rgb}{0.000000,0.000000,0.000000}%
\pgfsetstrokecolor{currentstroke}%
\pgfsetdash{}{0pt}%
\pgfpathmoveto{\pgfqpoint{3.393750in}{0.767343in}}%
\pgfpathlineto{\pgfqpoint{3.490000in}{0.767343in}}%
\pgfpathlineto{\pgfqpoint{3.490000in}{0.774862in}}%
\pgfpathlineto{\pgfqpoint{3.393750in}{0.774862in}}%
\pgfpathlineto{\pgfqpoint{3.393750in}{0.767343in}}%
\pgfusepath{fill}%
\end{pgfscope}%
\begin{pgfscope}%
\pgfpathrectangle{\pgfqpoint{3.393750in}{0.699667in}}{\pgfqpoint{0.096250in}{1.925000in}}%
\pgfusepath{clip}%
\pgfsetbuttcap%
\pgfsetroundjoin%
\definecolor{currentfill}{rgb}{0.018815,0.016026,0.084584}%
\pgfsetfillcolor{currentfill}%
\pgfsetfillopacity{0.800000}%
\pgfsetlinewidth{0.000000pt}%
\definecolor{currentstroke}{rgb}{0.000000,0.000000,0.000000}%
\pgfsetstrokecolor{currentstroke}%
\pgfsetdash{}{0pt}%
\pgfpathmoveto{\pgfqpoint{3.393750in}{0.774862in}}%
\pgfpathlineto{\pgfqpoint{3.490000in}{0.774862in}}%
\pgfpathlineto{\pgfqpoint{3.490000in}{0.782382in}}%
\pgfpathlineto{\pgfqpoint{3.393750in}{0.782382in}}%
\pgfpathlineto{\pgfqpoint{3.393750in}{0.774862in}}%
\pgfusepath{fill}%
\end{pgfscope}%
\begin{pgfscope}%
\pgfpathrectangle{\pgfqpoint{3.393750in}{0.699667in}}{\pgfqpoint{0.096250in}{1.925000in}}%
\pgfusepath{clip}%
\pgfsetbuttcap%
\pgfsetroundjoin%
\definecolor{currentfill}{rgb}{0.021692,0.018320,0.092610}%
\pgfsetfillcolor{currentfill}%
\pgfsetfillopacity{0.800000}%
\pgfsetlinewidth{0.000000pt}%
\definecolor{currentstroke}{rgb}{0.000000,0.000000,0.000000}%
\pgfsetstrokecolor{currentstroke}%
\pgfsetdash{}{0pt}%
\pgfpathmoveto{\pgfqpoint{3.393750in}{0.782382in}}%
\pgfpathlineto{\pgfqpoint{3.490000in}{0.782382in}}%
\pgfpathlineto{\pgfqpoint{3.490000in}{0.789902in}}%
\pgfpathlineto{\pgfqpoint{3.393750in}{0.789902in}}%
\pgfpathlineto{\pgfqpoint{3.393750in}{0.782382in}}%
\pgfusepath{fill}%
\end{pgfscope}%
\begin{pgfscope}%
\pgfpathrectangle{\pgfqpoint{3.393750in}{0.699667in}}{\pgfqpoint{0.096250in}{1.925000in}}%
\pgfusepath{clip}%
\pgfsetbuttcap%
\pgfsetroundjoin%
\definecolor{currentfill}{rgb}{0.024792,0.020715,0.100676}%
\pgfsetfillcolor{currentfill}%
\pgfsetfillopacity{0.800000}%
\pgfsetlinewidth{0.000000pt}%
\definecolor{currentstroke}{rgb}{0.000000,0.000000,0.000000}%
\pgfsetstrokecolor{currentstroke}%
\pgfsetdash{}{0pt}%
\pgfpathmoveto{\pgfqpoint{3.393750in}{0.789902in}}%
\pgfpathlineto{\pgfqpoint{3.490000in}{0.789902in}}%
\pgfpathlineto{\pgfqpoint{3.490000in}{0.797421in}}%
\pgfpathlineto{\pgfqpoint{3.393750in}{0.797421in}}%
\pgfpathlineto{\pgfqpoint{3.393750in}{0.789902in}}%
\pgfusepath{fill}%
\end{pgfscope}%
\begin{pgfscope}%
\pgfpathrectangle{\pgfqpoint{3.393750in}{0.699667in}}{\pgfqpoint{0.096250in}{1.925000in}}%
\pgfusepath{clip}%
\pgfsetbuttcap%
\pgfsetroundjoin%
\definecolor{currentfill}{rgb}{0.028123,0.023201,0.108787}%
\pgfsetfillcolor{currentfill}%
\pgfsetfillopacity{0.800000}%
\pgfsetlinewidth{0.000000pt}%
\definecolor{currentstroke}{rgb}{0.000000,0.000000,0.000000}%
\pgfsetstrokecolor{currentstroke}%
\pgfsetdash{}{0pt}%
\pgfpathmoveto{\pgfqpoint{3.393750in}{0.797421in}}%
\pgfpathlineto{\pgfqpoint{3.490000in}{0.797421in}}%
\pgfpathlineto{\pgfqpoint{3.490000in}{0.804941in}}%
\pgfpathlineto{\pgfqpoint{3.393750in}{0.804941in}}%
\pgfpathlineto{\pgfqpoint{3.393750in}{0.797421in}}%
\pgfusepath{fill}%
\end{pgfscope}%
\begin{pgfscope}%
\pgfpathrectangle{\pgfqpoint{3.393750in}{0.699667in}}{\pgfqpoint{0.096250in}{1.925000in}}%
\pgfusepath{clip}%
\pgfsetbuttcap%
\pgfsetroundjoin%
\definecolor{currentfill}{rgb}{0.031696,0.025765,0.116965}%
\pgfsetfillcolor{currentfill}%
\pgfsetfillopacity{0.800000}%
\pgfsetlinewidth{0.000000pt}%
\definecolor{currentstroke}{rgb}{0.000000,0.000000,0.000000}%
\pgfsetstrokecolor{currentstroke}%
\pgfsetdash{}{0pt}%
\pgfpathmoveto{\pgfqpoint{3.393750in}{0.804941in}}%
\pgfpathlineto{\pgfqpoint{3.490000in}{0.804941in}}%
\pgfpathlineto{\pgfqpoint{3.490000in}{0.812460in}}%
\pgfpathlineto{\pgfqpoint{3.393750in}{0.812460in}}%
\pgfpathlineto{\pgfqpoint{3.393750in}{0.804941in}}%
\pgfusepath{fill}%
\end{pgfscope}%
\begin{pgfscope}%
\pgfpathrectangle{\pgfqpoint{3.393750in}{0.699667in}}{\pgfqpoint{0.096250in}{1.925000in}}%
\pgfusepath{clip}%
\pgfsetbuttcap%
\pgfsetroundjoin%
\definecolor{currentfill}{rgb}{0.035520,0.028397,0.125209}%
\pgfsetfillcolor{currentfill}%
\pgfsetfillopacity{0.800000}%
\pgfsetlinewidth{0.000000pt}%
\definecolor{currentstroke}{rgb}{0.000000,0.000000,0.000000}%
\pgfsetstrokecolor{currentstroke}%
\pgfsetdash{}{0pt}%
\pgfpathmoveto{\pgfqpoint{3.393750in}{0.812460in}}%
\pgfpathlineto{\pgfqpoint{3.490000in}{0.812460in}}%
\pgfpathlineto{\pgfqpoint{3.490000in}{0.819980in}}%
\pgfpathlineto{\pgfqpoint{3.393750in}{0.819980in}}%
\pgfpathlineto{\pgfqpoint{3.393750in}{0.812460in}}%
\pgfusepath{fill}%
\end{pgfscope}%
\begin{pgfscope}%
\pgfpathrectangle{\pgfqpoint{3.393750in}{0.699667in}}{\pgfqpoint{0.096250in}{1.925000in}}%
\pgfusepath{clip}%
\pgfsetbuttcap%
\pgfsetroundjoin%
\definecolor{currentfill}{rgb}{0.039608,0.031090,0.133515}%
\pgfsetfillcolor{currentfill}%
\pgfsetfillopacity{0.800000}%
\pgfsetlinewidth{0.000000pt}%
\definecolor{currentstroke}{rgb}{0.000000,0.000000,0.000000}%
\pgfsetstrokecolor{currentstroke}%
\pgfsetdash{}{0pt}%
\pgfpathmoveto{\pgfqpoint{3.393750in}{0.819980in}}%
\pgfpathlineto{\pgfqpoint{3.490000in}{0.819980in}}%
\pgfpathlineto{\pgfqpoint{3.490000in}{0.827499in}}%
\pgfpathlineto{\pgfqpoint{3.393750in}{0.827499in}}%
\pgfpathlineto{\pgfqpoint{3.393750in}{0.819980in}}%
\pgfusepath{fill}%
\end{pgfscope}%
\begin{pgfscope}%
\pgfpathrectangle{\pgfqpoint{3.393750in}{0.699667in}}{\pgfqpoint{0.096250in}{1.925000in}}%
\pgfusepath{clip}%
\pgfsetbuttcap%
\pgfsetroundjoin%
\definecolor{currentfill}{rgb}{0.043830,0.033830,0.141886}%
\pgfsetfillcolor{currentfill}%
\pgfsetfillopacity{0.800000}%
\pgfsetlinewidth{0.000000pt}%
\definecolor{currentstroke}{rgb}{0.000000,0.000000,0.000000}%
\pgfsetstrokecolor{currentstroke}%
\pgfsetdash{}{0pt}%
\pgfpathmoveto{\pgfqpoint{3.393750in}{0.827499in}}%
\pgfpathlineto{\pgfqpoint{3.490000in}{0.827499in}}%
\pgfpathlineto{\pgfqpoint{3.490000in}{0.835019in}}%
\pgfpathlineto{\pgfqpoint{3.393750in}{0.835019in}}%
\pgfpathlineto{\pgfqpoint{3.393750in}{0.827499in}}%
\pgfusepath{fill}%
\end{pgfscope}%
\begin{pgfscope}%
\pgfpathrectangle{\pgfqpoint{3.393750in}{0.699667in}}{\pgfqpoint{0.096250in}{1.925000in}}%
\pgfusepath{clip}%
\pgfsetbuttcap%
\pgfsetroundjoin%
\definecolor{currentfill}{rgb}{0.048062,0.036607,0.150327}%
\pgfsetfillcolor{currentfill}%
\pgfsetfillopacity{0.800000}%
\pgfsetlinewidth{0.000000pt}%
\definecolor{currentstroke}{rgb}{0.000000,0.000000,0.000000}%
\pgfsetstrokecolor{currentstroke}%
\pgfsetdash{}{0pt}%
\pgfpathmoveto{\pgfqpoint{3.393750in}{0.835019in}}%
\pgfpathlineto{\pgfqpoint{3.490000in}{0.835019in}}%
\pgfpathlineto{\pgfqpoint{3.490000in}{0.842538in}}%
\pgfpathlineto{\pgfqpoint{3.393750in}{0.842538in}}%
\pgfpathlineto{\pgfqpoint{3.393750in}{0.835019in}}%
\pgfusepath{fill}%
\end{pgfscope}%
\begin{pgfscope}%
\pgfpathrectangle{\pgfqpoint{3.393750in}{0.699667in}}{\pgfqpoint{0.096250in}{1.925000in}}%
\pgfusepath{clip}%
\pgfsetbuttcap%
\pgfsetroundjoin%
\definecolor{currentfill}{rgb}{0.052320,0.039407,0.158841}%
\pgfsetfillcolor{currentfill}%
\pgfsetfillopacity{0.800000}%
\pgfsetlinewidth{0.000000pt}%
\definecolor{currentstroke}{rgb}{0.000000,0.000000,0.000000}%
\pgfsetstrokecolor{currentstroke}%
\pgfsetdash{}{0pt}%
\pgfpathmoveto{\pgfqpoint{3.393750in}{0.842538in}}%
\pgfpathlineto{\pgfqpoint{3.490000in}{0.842538in}}%
\pgfpathlineto{\pgfqpoint{3.490000in}{0.850058in}}%
\pgfpathlineto{\pgfqpoint{3.393750in}{0.850058in}}%
\pgfpathlineto{\pgfqpoint{3.393750in}{0.842538in}}%
\pgfusepath{fill}%
\end{pgfscope}%
\begin{pgfscope}%
\pgfpathrectangle{\pgfqpoint{3.393750in}{0.699667in}}{\pgfqpoint{0.096250in}{1.925000in}}%
\pgfusepath{clip}%
\pgfsetbuttcap%
\pgfsetroundjoin%
\definecolor{currentfill}{rgb}{0.056615,0.042160,0.167446}%
\pgfsetfillcolor{currentfill}%
\pgfsetfillopacity{0.800000}%
\pgfsetlinewidth{0.000000pt}%
\definecolor{currentstroke}{rgb}{0.000000,0.000000,0.000000}%
\pgfsetstrokecolor{currentstroke}%
\pgfsetdash{}{0pt}%
\pgfpathmoveto{\pgfqpoint{3.393750in}{0.850058in}}%
\pgfpathlineto{\pgfqpoint{3.490000in}{0.850058in}}%
\pgfpathlineto{\pgfqpoint{3.490000in}{0.857577in}}%
\pgfpathlineto{\pgfqpoint{3.393750in}{0.857577in}}%
\pgfpathlineto{\pgfqpoint{3.393750in}{0.850058in}}%
\pgfusepath{fill}%
\end{pgfscope}%
\begin{pgfscope}%
\pgfpathrectangle{\pgfqpoint{3.393750in}{0.699667in}}{\pgfqpoint{0.096250in}{1.925000in}}%
\pgfusepath{clip}%
\pgfsetbuttcap%
\pgfsetroundjoin%
\definecolor{currentfill}{rgb}{0.060949,0.044794,0.176129}%
\pgfsetfillcolor{currentfill}%
\pgfsetfillopacity{0.800000}%
\pgfsetlinewidth{0.000000pt}%
\definecolor{currentstroke}{rgb}{0.000000,0.000000,0.000000}%
\pgfsetstrokecolor{currentstroke}%
\pgfsetdash{}{0pt}%
\pgfpathmoveto{\pgfqpoint{3.393750in}{0.857577in}}%
\pgfpathlineto{\pgfqpoint{3.490000in}{0.857577in}}%
\pgfpathlineto{\pgfqpoint{3.490000in}{0.865097in}}%
\pgfpathlineto{\pgfqpoint{3.393750in}{0.865097in}}%
\pgfpathlineto{\pgfqpoint{3.393750in}{0.857577in}}%
\pgfusepath{fill}%
\end{pgfscope}%
\begin{pgfscope}%
\pgfpathrectangle{\pgfqpoint{3.393750in}{0.699667in}}{\pgfqpoint{0.096250in}{1.925000in}}%
\pgfusepath{clip}%
\pgfsetbuttcap%
\pgfsetroundjoin%
\definecolor{currentfill}{rgb}{0.065330,0.047318,0.184892}%
\pgfsetfillcolor{currentfill}%
\pgfsetfillopacity{0.800000}%
\pgfsetlinewidth{0.000000pt}%
\definecolor{currentstroke}{rgb}{0.000000,0.000000,0.000000}%
\pgfsetstrokecolor{currentstroke}%
\pgfsetdash{}{0pt}%
\pgfpathmoveto{\pgfqpoint{3.393750in}{0.865097in}}%
\pgfpathlineto{\pgfqpoint{3.490000in}{0.865097in}}%
\pgfpathlineto{\pgfqpoint{3.490000in}{0.872616in}}%
\pgfpathlineto{\pgfqpoint{3.393750in}{0.872616in}}%
\pgfpathlineto{\pgfqpoint{3.393750in}{0.865097in}}%
\pgfusepath{fill}%
\end{pgfscope}%
\begin{pgfscope}%
\pgfpathrectangle{\pgfqpoint{3.393750in}{0.699667in}}{\pgfqpoint{0.096250in}{1.925000in}}%
\pgfusepath{clip}%
\pgfsetbuttcap%
\pgfsetroundjoin%
\definecolor{currentfill}{rgb}{0.069764,0.049726,0.193735}%
\pgfsetfillcolor{currentfill}%
\pgfsetfillopacity{0.800000}%
\pgfsetlinewidth{0.000000pt}%
\definecolor{currentstroke}{rgb}{0.000000,0.000000,0.000000}%
\pgfsetstrokecolor{currentstroke}%
\pgfsetdash{}{0pt}%
\pgfpathmoveto{\pgfqpoint{3.393750in}{0.872616in}}%
\pgfpathlineto{\pgfqpoint{3.490000in}{0.872616in}}%
\pgfpathlineto{\pgfqpoint{3.490000in}{0.880136in}}%
\pgfpathlineto{\pgfqpoint{3.393750in}{0.880136in}}%
\pgfpathlineto{\pgfqpoint{3.393750in}{0.872616in}}%
\pgfusepath{fill}%
\end{pgfscope}%
\begin{pgfscope}%
\pgfpathrectangle{\pgfqpoint{3.393750in}{0.699667in}}{\pgfqpoint{0.096250in}{1.925000in}}%
\pgfusepath{clip}%
\pgfsetbuttcap%
\pgfsetroundjoin%
\definecolor{currentfill}{rgb}{0.074257,0.052017,0.202660}%
\pgfsetfillcolor{currentfill}%
\pgfsetfillopacity{0.800000}%
\pgfsetlinewidth{0.000000pt}%
\definecolor{currentstroke}{rgb}{0.000000,0.000000,0.000000}%
\pgfsetstrokecolor{currentstroke}%
\pgfsetdash{}{0pt}%
\pgfpathmoveto{\pgfqpoint{3.393750in}{0.880136in}}%
\pgfpathlineto{\pgfqpoint{3.490000in}{0.880136in}}%
\pgfpathlineto{\pgfqpoint{3.490000in}{0.887655in}}%
\pgfpathlineto{\pgfqpoint{3.393750in}{0.887655in}}%
\pgfpathlineto{\pgfqpoint{3.393750in}{0.880136in}}%
\pgfusepath{fill}%
\end{pgfscope}%
\begin{pgfscope}%
\pgfpathrectangle{\pgfqpoint{3.393750in}{0.699667in}}{\pgfqpoint{0.096250in}{1.925000in}}%
\pgfusepath{clip}%
\pgfsetbuttcap%
\pgfsetroundjoin%
\definecolor{currentfill}{rgb}{0.078815,0.054184,0.211667}%
\pgfsetfillcolor{currentfill}%
\pgfsetfillopacity{0.800000}%
\pgfsetlinewidth{0.000000pt}%
\definecolor{currentstroke}{rgb}{0.000000,0.000000,0.000000}%
\pgfsetstrokecolor{currentstroke}%
\pgfsetdash{}{0pt}%
\pgfpathmoveto{\pgfqpoint{3.393750in}{0.887655in}}%
\pgfpathlineto{\pgfqpoint{3.490000in}{0.887655in}}%
\pgfpathlineto{\pgfqpoint{3.490000in}{0.895175in}}%
\pgfpathlineto{\pgfqpoint{3.393750in}{0.895175in}}%
\pgfpathlineto{\pgfqpoint{3.393750in}{0.887655in}}%
\pgfusepath{fill}%
\end{pgfscope}%
\begin{pgfscope}%
\pgfpathrectangle{\pgfqpoint{3.393750in}{0.699667in}}{\pgfqpoint{0.096250in}{1.925000in}}%
\pgfusepath{clip}%
\pgfsetbuttcap%
\pgfsetroundjoin%
\definecolor{currentfill}{rgb}{0.083446,0.056225,0.220755}%
\pgfsetfillcolor{currentfill}%
\pgfsetfillopacity{0.800000}%
\pgfsetlinewidth{0.000000pt}%
\definecolor{currentstroke}{rgb}{0.000000,0.000000,0.000000}%
\pgfsetstrokecolor{currentstroke}%
\pgfsetdash{}{0pt}%
\pgfpathmoveto{\pgfqpoint{3.393750in}{0.895175in}}%
\pgfpathlineto{\pgfqpoint{3.490000in}{0.895175in}}%
\pgfpathlineto{\pgfqpoint{3.490000in}{0.902695in}}%
\pgfpathlineto{\pgfqpoint{3.393750in}{0.902695in}}%
\pgfpathlineto{\pgfqpoint{3.393750in}{0.895175in}}%
\pgfusepath{fill}%
\end{pgfscope}%
\begin{pgfscope}%
\pgfpathrectangle{\pgfqpoint{3.393750in}{0.699667in}}{\pgfqpoint{0.096250in}{1.925000in}}%
\pgfusepath{clip}%
\pgfsetbuttcap%
\pgfsetroundjoin%
\definecolor{currentfill}{rgb}{0.088155,0.058133,0.229922}%
\pgfsetfillcolor{currentfill}%
\pgfsetfillopacity{0.800000}%
\pgfsetlinewidth{0.000000pt}%
\definecolor{currentstroke}{rgb}{0.000000,0.000000,0.000000}%
\pgfsetstrokecolor{currentstroke}%
\pgfsetdash{}{0pt}%
\pgfpathmoveto{\pgfqpoint{3.393750in}{0.902695in}}%
\pgfpathlineto{\pgfqpoint{3.490000in}{0.902695in}}%
\pgfpathlineto{\pgfqpoint{3.490000in}{0.910214in}}%
\pgfpathlineto{\pgfqpoint{3.393750in}{0.910214in}}%
\pgfpathlineto{\pgfqpoint{3.393750in}{0.902695in}}%
\pgfusepath{fill}%
\end{pgfscope}%
\begin{pgfscope}%
\pgfpathrectangle{\pgfqpoint{3.393750in}{0.699667in}}{\pgfqpoint{0.096250in}{1.925000in}}%
\pgfusepath{clip}%
\pgfsetbuttcap%
\pgfsetroundjoin%
\definecolor{currentfill}{rgb}{0.092949,0.059904,0.239164}%
\pgfsetfillcolor{currentfill}%
\pgfsetfillopacity{0.800000}%
\pgfsetlinewidth{0.000000pt}%
\definecolor{currentstroke}{rgb}{0.000000,0.000000,0.000000}%
\pgfsetstrokecolor{currentstroke}%
\pgfsetdash{}{0pt}%
\pgfpathmoveto{\pgfqpoint{3.393750in}{0.910214in}}%
\pgfpathlineto{\pgfqpoint{3.490000in}{0.910214in}}%
\pgfpathlineto{\pgfqpoint{3.490000in}{0.917734in}}%
\pgfpathlineto{\pgfqpoint{3.393750in}{0.917734in}}%
\pgfpathlineto{\pgfqpoint{3.393750in}{0.910214in}}%
\pgfusepath{fill}%
\end{pgfscope}%
\begin{pgfscope}%
\pgfpathrectangle{\pgfqpoint{3.393750in}{0.699667in}}{\pgfqpoint{0.096250in}{1.925000in}}%
\pgfusepath{clip}%
\pgfsetbuttcap%
\pgfsetroundjoin%
\definecolor{currentfill}{rgb}{0.097833,0.061531,0.248477}%
\pgfsetfillcolor{currentfill}%
\pgfsetfillopacity{0.800000}%
\pgfsetlinewidth{0.000000pt}%
\definecolor{currentstroke}{rgb}{0.000000,0.000000,0.000000}%
\pgfsetstrokecolor{currentstroke}%
\pgfsetdash{}{0pt}%
\pgfpathmoveto{\pgfqpoint{3.393750in}{0.917734in}}%
\pgfpathlineto{\pgfqpoint{3.490000in}{0.917734in}}%
\pgfpathlineto{\pgfqpoint{3.490000in}{0.925253in}}%
\pgfpathlineto{\pgfqpoint{3.393750in}{0.925253in}}%
\pgfpathlineto{\pgfqpoint{3.393750in}{0.917734in}}%
\pgfusepath{fill}%
\end{pgfscope}%
\begin{pgfscope}%
\pgfpathrectangle{\pgfqpoint{3.393750in}{0.699667in}}{\pgfqpoint{0.096250in}{1.925000in}}%
\pgfusepath{clip}%
\pgfsetbuttcap%
\pgfsetroundjoin%
\definecolor{currentfill}{rgb}{0.102815,0.063010,0.257854}%
\pgfsetfillcolor{currentfill}%
\pgfsetfillopacity{0.800000}%
\pgfsetlinewidth{0.000000pt}%
\definecolor{currentstroke}{rgb}{0.000000,0.000000,0.000000}%
\pgfsetstrokecolor{currentstroke}%
\pgfsetdash{}{0pt}%
\pgfpathmoveto{\pgfqpoint{3.393750in}{0.925253in}}%
\pgfpathlineto{\pgfqpoint{3.490000in}{0.925253in}}%
\pgfpathlineto{\pgfqpoint{3.490000in}{0.932773in}}%
\pgfpathlineto{\pgfqpoint{3.393750in}{0.932773in}}%
\pgfpathlineto{\pgfqpoint{3.393750in}{0.925253in}}%
\pgfusepath{fill}%
\end{pgfscope}%
\begin{pgfscope}%
\pgfpathrectangle{\pgfqpoint{3.393750in}{0.699667in}}{\pgfqpoint{0.096250in}{1.925000in}}%
\pgfusepath{clip}%
\pgfsetbuttcap%
\pgfsetroundjoin%
\definecolor{currentfill}{rgb}{0.107899,0.064335,0.267289}%
\pgfsetfillcolor{currentfill}%
\pgfsetfillopacity{0.800000}%
\pgfsetlinewidth{0.000000pt}%
\definecolor{currentstroke}{rgb}{0.000000,0.000000,0.000000}%
\pgfsetstrokecolor{currentstroke}%
\pgfsetdash{}{0pt}%
\pgfpathmoveto{\pgfqpoint{3.393750in}{0.932773in}}%
\pgfpathlineto{\pgfqpoint{3.490000in}{0.932773in}}%
\pgfpathlineto{\pgfqpoint{3.490000in}{0.940292in}}%
\pgfpathlineto{\pgfqpoint{3.393750in}{0.940292in}}%
\pgfpathlineto{\pgfqpoint{3.393750in}{0.932773in}}%
\pgfusepath{fill}%
\end{pgfscope}%
\begin{pgfscope}%
\pgfpathrectangle{\pgfqpoint{3.393750in}{0.699667in}}{\pgfqpoint{0.096250in}{1.925000in}}%
\pgfusepath{clip}%
\pgfsetbuttcap%
\pgfsetroundjoin%
\definecolor{currentfill}{rgb}{0.113094,0.065492,0.276784}%
\pgfsetfillcolor{currentfill}%
\pgfsetfillopacity{0.800000}%
\pgfsetlinewidth{0.000000pt}%
\definecolor{currentstroke}{rgb}{0.000000,0.000000,0.000000}%
\pgfsetstrokecolor{currentstroke}%
\pgfsetdash{}{0pt}%
\pgfpathmoveto{\pgfqpoint{3.393750in}{0.940292in}}%
\pgfpathlineto{\pgfqpoint{3.490000in}{0.940292in}}%
\pgfpathlineto{\pgfqpoint{3.490000in}{0.947812in}}%
\pgfpathlineto{\pgfqpoint{3.393750in}{0.947812in}}%
\pgfpathlineto{\pgfqpoint{3.393750in}{0.940292in}}%
\pgfusepath{fill}%
\end{pgfscope}%
\begin{pgfscope}%
\pgfpathrectangle{\pgfqpoint{3.393750in}{0.699667in}}{\pgfqpoint{0.096250in}{1.925000in}}%
\pgfusepath{clip}%
\pgfsetbuttcap%
\pgfsetroundjoin%
\definecolor{currentfill}{rgb}{0.118405,0.066479,0.286321}%
\pgfsetfillcolor{currentfill}%
\pgfsetfillopacity{0.800000}%
\pgfsetlinewidth{0.000000pt}%
\definecolor{currentstroke}{rgb}{0.000000,0.000000,0.000000}%
\pgfsetstrokecolor{currentstroke}%
\pgfsetdash{}{0pt}%
\pgfpathmoveto{\pgfqpoint{3.393750in}{0.947812in}}%
\pgfpathlineto{\pgfqpoint{3.490000in}{0.947812in}}%
\pgfpathlineto{\pgfqpoint{3.490000in}{0.955331in}}%
\pgfpathlineto{\pgfqpoint{3.393750in}{0.955331in}}%
\pgfpathlineto{\pgfqpoint{3.393750in}{0.947812in}}%
\pgfusepath{fill}%
\end{pgfscope}%
\begin{pgfscope}%
\pgfpathrectangle{\pgfqpoint{3.393750in}{0.699667in}}{\pgfqpoint{0.096250in}{1.925000in}}%
\pgfusepath{clip}%
\pgfsetbuttcap%
\pgfsetroundjoin%
\definecolor{currentfill}{rgb}{0.123833,0.067295,0.295879}%
\pgfsetfillcolor{currentfill}%
\pgfsetfillopacity{0.800000}%
\pgfsetlinewidth{0.000000pt}%
\definecolor{currentstroke}{rgb}{0.000000,0.000000,0.000000}%
\pgfsetstrokecolor{currentstroke}%
\pgfsetdash{}{0pt}%
\pgfpathmoveto{\pgfqpoint{3.393750in}{0.955331in}}%
\pgfpathlineto{\pgfqpoint{3.490000in}{0.955331in}}%
\pgfpathlineto{\pgfqpoint{3.490000in}{0.962851in}}%
\pgfpathlineto{\pgfqpoint{3.393750in}{0.962851in}}%
\pgfpathlineto{\pgfqpoint{3.393750in}{0.955331in}}%
\pgfusepath{fill}%
\end{pgfscope}%
\begin{pgfscope}%
\pgfpathrectangle{\pgfqpoint{3.393750in}{0.699667in}}{\pgfqpoint{0.096250in}{1.925000in}}%
\pgfusepath{clip}%
\pgfsetbuttcap%
\pgfsetroundjoin%
\definecolor{currentfill}{rgb}{0.129380,0.067935,0.305443}%
\pgfsetfillcolor{currentfill}%
\pgfsetfillopacity{0.800000}%
\pgfsetlinewidth{0.000000pt}%
\definecolor{currentstroke}{rgb}{0.000000,0.000000,0.000000}%
\pgfsetstrokecolor{currentstroke}%
\pgfsetdash{}{0pt}%
\pgfpathmoveto{\pgfqpoint{3.393750in}{0.962851in}}%
\pgfpathlineto{\pgfqpoint{3.490000in}{0.962851in}}%
\pgfpathlineto{\pgfqpoint{3.490000in}{0.970370in}}%
\pgfpathlineto{\pgfqpoint{3.393750in}{0.970370in}}%
\pgfpathlineto{\pgfqpoint{3.393750in}{0.962851in}}%
\pgfusepath{fill}%
\end{pgfscope}%
\begin{pgfscope}%
\pgfpathrectangle{\pgfqpoint{3.393750in}{0.699667in}}{\pgfqpoint{0.096250in}{1.925000in}}%
\pgfusepath{clip}%
\pgfsetbuttcap%
\pgfsetroundjoin%
\definecolor{currentfill}{rgb}{0.135053,0.068391,0.315000}%
\pgfsetfillcolor{currentfill}%
\pgfsetfillopacity{0.800000}%
\pgfsetlinewidth{0.000000pt}%
\definecolor{currentstroke}{rgb}{0.000000,0.000000,0.000000}%
\pgfsetstrokecolor{currentstroke}%
\pgfsetdash{}{0pt}%
\pgfpathmoveto{\pgfqpoint{3.393750in}{0.970370in}}%
\pgfpathlineto{\pgfqpoint{3.490000in}{0.970370in}}%
\pgfpathlineto{\pgfqpoint{3.490000in}{0.977890in}}%
\pgfpathlineto{\pgfqpoint{3.393750in}{0.977890in}}%
\pgfpathlineto{\pgfqpoint{3.393750in}{0.970370in}}%
\pgfusepath{fill}%
\end{pgfscope}%
\begin{pgfscope}%
\pgfpathrectangle{\pgfqpoint{3.393750in}{0.699667in}}{\pgfqpoint{0.096250in}{1.925000in}}%
\pgfusepath{clip}%
\pgfsetbuttcap%
\pgfsetroundjoin%
\definecolor{currentfill}{rgb}{0.140858,0.068654,0.324538}%
\pgfsetfillcolor{currentfill}%
\pgfsetfillopacity{0.800000}%
\pgfsetlinewidth{0.000000pt}%
\definecolor{currentstroke}{rgb}{0.000000,0.000000,0.000000}%
\pgfsetstrokecolor{currentstroke}%
\pgfsetdash{}{0pt}%
\pgfpathmoveto{\pgfqpoint{3.393750in}{0.977890in}}%
\pgfpathlineto{\pgfqpoint{3.490000in}{0.977890in}}%
\pgfpathlineto{\pgfqpoint{3.490000in}{0.985409in}}%
\pgfpathlineto{\pgfqpoint{3.393750in}{0.985409in}}%
\pgfpathlineto{\pgfqpoint{3.393750in}{0.977890in}}%
\pgfusepath{fill}%
\end{pgfscope}%
\begin{pgfscope}%
\pgfpathrectangle{\pgfqpoint{3.393750in}{0.699667in}}{\pgfqpoint{0.096250in}{1.925000in}}%
\pgfusepath{clip}%
\pgfsetbuttcap%
\pgfsetroundjoin%
\definecolor{currentfill}{rgb}{0.146785,0.068738,0.334011}%
\pgfsetfillcolor{currentfill}%
\pgfsetfillopacity{0.800000}%
\pgfsetlinewidth{0.000000pt}%
\definecolor{currentstroke}{rgb}{0.000000,0.000000,0.000000}%
\pgfsetstrokecolor{currentstroke}%
\pgfsetdash{}{0pt}%
\pgfpathmoveto{\pgfqpoint{3.393750in}{0.985409in}}%
\pgfpathlineto{\pgfqpoint{3.490000in}{0.985409in}}%
\pgfpathlineto{\pgfqpoint{3.490000in}{0.992929in}}%
\pgfpathlineto{\pgfqpoint{3.393750in}{0.992929in}}%
\pgfpathlineto{\pgfqpoint{3.393750in}{0.985409in}}%
\pgfusepath{fill}%
\end{pgfscope}%
\begin{pgfscope}%
\pgfpathrectangle{\pgfqpoint{3.393750in}{0.699667in}}{\pgfqpoint{0.096250in}{1.925000in}}%
\pgfusepath{clip}%
\pgfsetbuttcap%
\pgfsetroundjoin%
\definecolor{currentfill}{rgb}{0.152839,0.068637,0.343404}%
\pgfsetfillcolor{currentfill}%
\pgfsetfillopacity{0.800000}%
\pgfsetlinewidth{0.000000pt}%
\definecolor{currentstroke}{rgb}{0.000000,0.000000,0.000000}%
\pgfsetstrokecolor{currentstroke}%
\pgfsetdash{}{0pt}%
\pgfpathmoveto{\pgfqpoint{3.393750in}{0.992929in}}%
\pgfpathlineto{\pgfqpoint{3.490000in}{0.992929in}}%
\pgfpathlineto{\pgfqpoint{3.490000in}{1.000448in}}%
\pgfpathlineto{\pgfqpoint{3.393750in}{1.000448in}}%
\pgfpathlineto{\pgfqpoint{3.393750in}{0.992929in}}%
\pgfusepath{fill}%
\end{pgfscope}%
\begin{pgfscope}%
\pgfpathrectangle{\pgfqpoint{3.393750in}{0.699667in}}{\pgfqpoint{0.096250in}{1.925000in}}%
\pgfusepath{clip}%
\pgfsetbuttcap%
\pgfsetroundjoin%
\definecolor{currentfill}{rgb}{0.159018,0.068354,0.352688}%
\pgfsetfillcolor{currentfill}%
\pgfsetfillopacity{0.800000}%
\pgfsetlinewidth{0.000000pt}%
\definecolor{currentstroke}{rgb}{0.000000,0.000000,0.000000}%
\pgfsetstrokecolor{currentstroke}%
\pgfsetdash{}{0pt}%
\pgfpathmoveto{\pgfqpoint{3.393750in}{1.000448in}}%
\pgfpathlineto{\pgfqpoint{3.490000in}{1.000448in}}%
\pgfpathlineto{\pgfqpoint{3.490000in}{1.007968in}}%
\pgfpathlineto{\pgfqpoint{3.393750in}{1.007968in}}%
\pgfpathlineto{\pgfqpoint{3.393750in}{1.000448in}}%
\pgfusepath{fill}%
\end{pgfscope}%
\begin{pgfscope}%
\pgfpathrectangle{\pgfqpoint{3.393750in}{0.699667in}}{\pgfqpoint{0.096250in}{1.925000in}}%
\pgfusepath{clip}%
\pgfsetbuttcap%
\pgfsetroundjoin%
\definecolor{currentfill}{rgb}{0.165308,0.067911,0.361816}%
\pgfsetfillcolor{currentfill}%
\pgfsetfillopacity{0.800000}%
\pgfsetlinewidth{0.000000pt}%
\definecolor{currentstroke}{rgb}{0.000000,0.000000,0.000000}%
\pgfsetstrokecolor{currentstroke}%
\pgfsetdash{}{0pt}%
\pgfpathmoveto{\pgfqpoint{3.393750in}{1.007968in}}%
\pgfpathlineto{\pgfqpoint{3.490000in}{1.007968in}}%
\pgfpathlineto{\pgfqpoint{3.490000in}{1.015487in}}%
\pgfpathlineto{\pgfqpoint{3.393750in}{1.015487in}}%
\pgfpathlineto{\pgfqpoint{3.393750in}{1.007968in}}%
\pgfusepath{fill}%
\end{pgfscope}%
\begin{pgfscope}%
\pgfpathrectangle{\pgfqpoint{3.393750in}{0.699667in}}{\pgfqpoint{0.096250in}{1.925000in}}%
\pgfusepath{clip}%
\pgfsetbuttcap%
\pgfsetroundjoin%
\definecolor{currentfill}{rgb}{0.171713,0.067305,0.370771}%
\pgfsetfillcolor{currentfill}%
\pgfsetfillopacity{0.800000}%
\pgfsetlinewidth{0.000000pt}%
\definecolor{currentstroke}{rgb}{0.000000,0.000000,0.000000}%
\pgfsetstrokecolor{currentstroke}%
\pgfsetdash{}{0pt}%
\pgfpathmoveto{\pgfqpoint{3.393750in}{1.015487in}}%
\pgfpathlineto{\pgfqpoint{3.490000in}{1.015487in}}%
\pgfpathlineto{\pgfqpoint{3.490000in}{1.023007in}}%
\pgfpathlineto{\pgfqpoint{3.393750in}{1.023007in}}%
\pgfpathlineto{\pgfqpoint{3.393750in}{1.015487in}}%
\pgfusepath{fill}%
\end{pgfscope}%
\begin{pgfscope}%
\pgfpathrectangle{\pgfqpoint{3.393750in}{0.699667in}}{\pgfqpoint{0.096250in}{1.925000in}}%
\pgfusepath{clip}%
\pgfsetbuttcap%
\pgfsetroundjoin%
\definecolor{currentfill}{rgb}{0.178212,0.066576,0.379497}%
\pgfsetfillcolor{currentfill}%
\pgfsetfillopacity{0.800000}%
\pgfsetlinewidth{0.000000pt}%
\definecolor{currentstroke}{rgb}{0.000000,0.000000,0.000000}%
\pgfsetstrokecolor{currentstroke}%
\pgfsetdash{}{0pt}%
\pgfpathmoveto{\pgfqpoint{3.393750in}{1.023007in}}%
\pgfpathlineto{\pgfqpoint{3.490000in}{1.023007in}}%
\pgfpathlineto{\pgfqpoint{3.490000in}{1.030527in}}%
\pgfpathlineto{\pgfqpoint{3.393750in}{1.030527in}}%
\pgfpathlineto{\pgfqpoint{3.393750in}{1.023007in}}%
\pgfusepath{fill}%
\end{pgfscope}%
\begin{pgfscope}%
\pgfpathrectangle{\pgfqpoint{3.393750in}{0.699667in}}{\pgfqpoint{0.096250in}{1.925000in}}%
\pgfusepath{clip}%
\pgfsetbuttcap%
\pgfsetroundjoin%
\definecolor{currentfill}{rgb}{0.184801,0.065732,0.387973}%
\pgfsetfillcolor{currentfill}%
\pgfsetfillopacity{0.800000}%
\pgfsetlinewidth{0.000000pt}%
\definecolor{currentstroke}{rgb}{0.000000,0.000000,0.000000}%
\pgfsetstrokecolor{currentstroke}%
\pgfsetdash{}{0pt}%
\pgfpathmoveto{\pgfqpoint{3.393750in}{1.030527in}}%
\pgfpathlineto{\pgfqpoint{3.490000in}{1.030527in}}%
\pgfpathlineto{\pgfqpoint{3.490000in}{1.038046in}}%
\pgfpathlineto{\pgfqpoint{3.393750in}{1.038046in}}%
\pgfpathlineto{\pgfqpoint{3.393750in}{1.030527in}}%
\pgfusepath{fill}%
\end{pgfscope}%
\begin{pgfscope}%
\pgfpathrectangle{\pgfqpoint{3.393750in}{0.699667in}}{\pgfqpoint{0.096250in}{1.925000in}}%
\pgfusepath{clip}%
\pgfsetbuttcap%
\pgfsetroundjoin%
\definecolor{currentfill}{rgb}{0.191460,0.064818,0.396152}%
\pgfsetfillcolor{currentfill}%
\pgfsetfillopacity{0.800000}%
\pgfsetlinewidth{0.000000pt}%
\definecolor{currentstroke}{rgb}{0.000000,0.000000,0.000000}%
\pgfsetstrokecolor{currentstroke}%
\pgfsetdash{}{0pt}%
\pgfpathmoveto{\pgfqpoint{3.393750in}{1.038046in}}%
\pgfpathlineto{\pgfqpoint{3.490000in}{1.038046in}}%
\pgfpathlineto{\pgfqpoint{3.490000in}{1.045566in}}%
\pgfpathlineto{\pgfqpoint{3.393750in}{1.045566in}}%
\pgfpathlineto{\pgfqpoint{3.393750in}{1.038046in}}%
\pgfusepath{fill}%
\end{pgfscope}%
\begin{pgfscope}%
\pgfpathrectangle{\pgfqpoint{3.393750in}{0.699667in}}{\pgfqpoint{0.096250in}{1.925000in}}%
\pgfusepath{clip}%
\pgfsetbuttcap%
\pgfsetroundjoin%
\definecolor{currentfill}{rgb}{0.198177,0.063862,0.404009}%
\pgfsetfillcolor{currentfill}%
\pgfsetfillopacity{0.800000}%
\pgfsetlinewidth{0.000000pt}%
\definecolor{currentstroke}{rgb}{0.000000,0.000000,0.000000}%
\pgfsetstrokecolor{currentstroke}%
\pgfsetdash{}{0pt}%
\pgfpathmoveto{\pgfqpoint{3.393750in}{1.045566in}}%
\pgfpathlineto{\pgfqpoint{3.490000in}{1.045566in}}%
\pgfpathlineto{\pgfqpoint{3.490000in}{1.053085in}}%
\pgfpathlineto{\pgfqpoint{3.393750in}{1.053085in}}%
\pgfpathlineto{\pgfqpoint{3.393750in}{1.045566in}}%
\pgfusepath{fill}%
\end{pgfscope}%
\begin{pgfscope}%
\pgfpathrectangle{\pgfqpoint{3.393750in}{0.699667in}}{\pgfqpoint{0.096250in}{1.925000in}}%
\pgfusepath{clip}%
\pgfsetbuttcap%
\pgfsetroundjoin%
\definecolor{currentfill}{rgb}{0.204935,0.062907,0.411514}%
\pgfsetfillcolor{currentfill}%
\pgfsetfillopacity{0.800000}%
\pgfsetlinewidth{0.000000pt}%
\definecolor{currentstroke}{rgb}{0.000000,0.000000,0.000000}%
\pgfsetstrokecolor{currentstroke}%
\pgfsetdash{}{0pt}%
\pgfpathmoveto{\pgfqpoint{3.393750in}{1.053085in}}%
\pgfpathlineto{\pgfqpoint{3.490000in}{1.053085in}}%
\pgfpathlineto{\pgfqpoint{3.490000in}{1.060605in}}%
\pgfpathlineto{\pgfqpoint{3.393750in}{1.060605in}}%
\pgfpathlineto{\pgfqpoint{3.393750in}{1.053085in}}%
\pgfusepath{fill}%
\end{pgfscope}%
\begin{pgfscope}%
\pgfpathrectangle{\pgfqpoint{3.393750in}{0.699667in}}{\pgfqpoint{0.096250in}{1.925000in}}%
\pgfusepath{clip}%
\pgfsetbuttcap%
\pgfsetroundjoin%
\definecolor{currentfill}{rgb}{0.211718,0.061992,0.418647}%
\pgfsetfillcolor{currentfill}%
\pgfsetfillopacity{0.800000}%
\pgfsetlinewidth{0.000000pt}%
\definecolor{currentstroke}{rgb}{0.000000,0.000000,0.000000}%
\pgfsetstrokecolor{currentstroke}%
\pgfsetdash{}{0pt}%
\pgfpathmoveto{\pgfqpoint{3.393750in}{1.060605in}}%
\pgfpathlineto{\pgfqpoint{3.490000in}{1.060605in}}%
\pgfpathlineto{\pgfqpoint{3.490000in}{1.068124in}}%
\pgfpathlineto{\pgfqpoint{3.393750in}{1.068124in}}%
\pgfpathlineto{\pgfqpoint{3.393750in}{1.060605in}}%
\pgfusepath{fill}%
\end{pgfscope}%
\begin{pgfscope}%
\pgfpathrectangle{\pgfqpoint{3.393750in}{0.699667in}}{\pgfqpoint{0.096250in}{1.925000in}}%
\pgfusepath{clip}%
\pgfsetbuttcap%
\pgfsetroundjoin%
\definecolor{currentfill}{rgb}{0.218512,0.061158,0.425392}%
\pgfsetfillcolor{currentfill}%
\pgfsetfillopacity{0.800000}%
\pgfsetlinewidth{0.000000pt}%
\definecolor{currentstroke}{rgb}{0.000000,0.000000,0.000000}%
\pgfsetstrokecolor{currentstroke}%
\pgfsetdash{}{0pt}%
\pgfpathmoveto{\pgfqpoint{3.393750in}{1.068124in}}%
\pgfpathlineto{\pgfqpoint{3.490000in}{1.068124in}}%
\pgfpathlineto{\pgfqpoint{3.490000in}{1.075644in}}%
\pgfpathlineto{\pgfqpoint{3.393750in}{1.075644in}}%
\pgfpathlineto{\pgfqpoint{3.393750in}{1.068124in}}%
\pgfusepath{fill}%
\end{pgfscope}%
\begin{pgfscope}%
\pgfpathrectangle{\pgfqpoint{3.393750in}{0.699667in}}{\pgfqpoint{0.096250in}{1.925000in}}%
\pgfusepath{clip}%
\pgfsetbuttcap%
\pgfsetroundjoin%
\definecolor{currentfill}{rgb}{0.225302,0.060445,0.431742}%
\pgfsetfillcolor{currentfill}%
\pgfsetfillopacity{0.800000}%
\pgfsetlinewidth{0.000000pt}%
\definecolor{currentstroke}{rgb}{0.000000,0.000000,0.000000}%
\pgfsetstrokecolor{currentstroke}%
\pgfsetdash{}{0pt}%
\pgfpathmoveto{\pgfqpoint{3.393750in}{1.075644in}}%
\pgfpathlineto{\pgfqpoint{3.490000in}{1.075644in}}%
\pgfpathlineto{\pgfqpoint{3.490000in}{1.083163in}}%
\pgfpathlineto{\pgfqpoint{3.393750in}{1.083163in}}%
\pgfpathlineto{\pgfqpoint{3.393750in}{1.075644in}}%
\pgfusepath{fill}%
\end{pgfscope}%
\begin{pgfscope}%
\pgfpathrectangle{\pgfqpoint{3.393750in}{0.699667in}}{\pgfqpoint{0.096250in}{1.925000in}}%
\pgfusepath{clip}%
\pgfsetbuttcap%
\pgfsetroundjoin%
\definecolor{currentfill}{rgb}{0.232077,0.059889,0.437695}%
\pgfsetfillcolor{currentfill}%
\pgfsetfillopacity{0.800000}%
\pgfsetlinewidth{0.000000pt}%
\definecolor{currentstroke}{rgb}{0.000000,0.000000,0.000000}%
\pgfsetstrokecolor{currentstroke}%
\pgfsetdash{}{0pt}%
\pgfpathmoveto{\pgfqpoint{3.393750in}{1.083163in}}%
\pgfpathlineto{\pgfqpoint{3.490000in}{1.083163in}}%
\pgfpathlineto{\pgfqpoint{3.490000in}{1.090683in}}%
\pgfpathlineto{\pgfqpoint{3.393750in}{1.090683in}}%
\pgfpathlineto{\pgfqpoint{3.393750in}{1.083163in}}%
\pgfusepath{fill}%
\end{pgfscope}%
\begin{pgfscope}%
\pgfpathrectangle{\pgfqpoint{3.393750in}{0.699667in}}{\pgfqpoint{0.096250in}{1.925000in}}%
\pgfusepath{clip}%
\pgfsetbuttcap%
\pgfsetroundjoin%
\definecolor{currentfill}{rgb}{0.238826,0.059517,0.443256}%
\pgfsetfillcolor{currentfill}%
\pgfsetfillopacity{0.800000}%
\pgfsetlinewidth{0.000000pt}%
\definecolor{currentstroke}{rgb}{0.000000,0.000000,0.000000}%
\pgfsetstrokecolor{currentstroke}%
\pgfsetdash{}{0pt}%
\pgfpathmoveto{\pgfqpoint{3.393750in}{1.090683in}}%
\pgfpathlineto{\pgfqpoint{3.490000in}{1.090683in}}%
\pgfpathlineto{\pgfqpoint{3.490000in}{1.098202in}}%
\pgfpathlineto{\pgfqpoint{3.393750in}{1.098202in}}%
\pgfpathlineto{\pgfqpoint{3.393750in}{1.090683in}}%
\pgfusepath{fill}%
\end{pgfscope}%
\begin{pgfscope}%
\pgfpathrectangle{\pgfqpoint{3.393750in}{0.699667in}}{\pgfqpoint{0.096250in}{1.925000in}}%
\pgfusepath{clip}%
\pgfsetbuttcap%
\pgfsetroundjoin%
\definecolor{currentfill}{rgb}{0.245543,0.059352,0.448436}%
\pgfsetfillcolor{currentfill}%
\pgfsetfillopacity{0.800000}%
\pgfsetlinewidth{0.000000pt}%
\definecolor{currentstroke}{rgb}{0.000000,0.000000,0.000000}%
\pgfsetstrokecolor{currentstroke}%
\pgfsetdash{}{0pt}%
\pgfpathmoveto{\pgfqpoint{3.393750in}{1.098202in}}%
\pgfpathlineto{\pgfqpoint{3.490000in}{1.098202in}}%
\pgfpathlineto{\pgfqpoint{3.490000in}{1.105722in}}%
\pgfpathlineto{\pgfqpoint{3.393750in}{1.105722in}}%
\pgfpathlineto{\pgfqpoint{3.393750in}{1.098202in}}%
\pgfusepath{fill}%
\end{pgfscope}%
\begin{pgfscope}%
\pgfpathrectangle{\pgfqpoint{3.393750in}{0.699667in}}{\pgfqpoint{0.096250in}{1.925000in}}%
\pgfusepath{clip}%
\pgfsetbuttcap%
\pgfsetroundjoin%
\definecolor{currentfill}{rgb}{0.252220,0.059415,0.453248}%
\pgfsetfillcolor{currentfill}%
\pgfsetfillopacity{0.800000}%
\pgfsetlinewidth{0.000000pt}%
\definecolor{currentstroke}{rgb}{0.000000,0.000000,0.000000}%
\pgfsetstrokecolor{currentstroke}%
\pgfsetdash{}{0pt}%
\pgfpathmoveto{\pgfqpoint{3.393750in}{1.105722in}}%
\pgfpathlineto{\pgfqpoint{3.490000in}{1.105722in}}%
\pgfpathlineto{\pgfqpoint{3.490000in}{1.113241in}}%
\pgfpathlineto{\pgfqpoint{3.393750in}{1.113241in}}%
\pgfpathlineto{\pgfqpoint{3.393750in}{1.105722in}}%
\pgfusepath{fill}%
\end{pgfscope}%
\begin{pgfscope}%
\pgfpathrectangle{\pgfqpoint{3.393750in}{0.699667in}}{\pgfqpoint{0.096250in}{1.925000in}}%
\pgfusepath{clip}%
\pgfsetbuttcap%
\pgfsetroundjoin%
\definecolor{currentfill}{rgb}{0.258857,0.059706,0.457710}%
\pgfsetfillcolor{currentfill}%
\pgfsetfillopacity{0.800000}%
\pgfsetlinewidth{0.000000pt}%
\definecolor{currentstroke}{rgb}{0.000000,0.000000,0.000000}%
\pgfsetstrokecolor{currentstroke}%
\pgfsetdash{}{0pt}%
\pgfpathmoveto{\pgfqpoint{3.393750in}{1.113241in}}%
\pgfpathlineto{\pgfqpoint{3.490000in}{1.113241in}}%
\pgfpathlineto{\pgfqpoint{3.490000in}{1.120761in}}%
\pgfpathlineto{\pgfqpoint{3.393750in}{1.120761in}}%
\pgfpathlineto{\pgfqpoint{3.393750in}{1.113241in}}%
\pgfusepath{fill}%
\end{pgfscope}%
\begin{pgfscope}%
\pgfpathrectangle{\pgfqpoint{3.393750in}{0.699667in}}{\pgfqpoint{0.096250in}{1.925000in}}%
\pgfusepath{clip}%
\pgfsetbuttcap%
\pgfsetroundjoin%
\definecolor{currentfill}{rgb}{0.265447,0.060237,0.461840}%
\pgfsetfillcolor{currentfill}%
\pgfsetfillopacity{0.800000}%
\pgfsetlinewidth{0.000000pt}%
\definecolor{currentstroke}{rgb}{0.000000,0.000000,0.000000}%
\pgfsetstrokecolor{currentstroke}%
\pgfsetdash{}{0pt}%
\pgfpathmoveto{\pgfqpoint{3.393750in}{1.120761in}}%
\pgfpathlineto{\pgfqpoint{3.490000in}{1.120761in}}%
\pgfpathlineto{\pgfqpoint{3.490000in}{1.128280in}}%
\pgfpathlineto{\pgfqpoint{3.393750in}{1.128280in}}%
\pgfpathlineto{\pgfqpoint{3.393750in}{1.120761in}}%
\pgfusepath{fill}%
\end{pgfscope}%
\begin{pgfscope}%
\pgfpathrectangle{\pgfqpoint{3.393750in}{0.699667in}}{\pgfqpoint{0.096250in}{1.925000in}}%
\pgfusepath{clip}%
\pgfsetbuttcap%
\pgfsetroundjoin%
\definecolor{currentfill}{rgb}{0.271994,0.060994,0.465660}%
\pgfsetfillcolor{currentfill}%
\pgfsetfillopacity{0.800000}%
\pgfsetlinewidth{0.000000pt}%
\definecolor{currentstroke}{rgb}{0.000000,0.000000,0.000000}%
\pgfsetstrokecolor{currentstroke}%
\pgfsetdash{}{0pt}%
\pgfpathmoveto{\pgfqpoint{3.393750in}{1.128280in}}%
\pgfpathlineto{\pgfqpoint{3.490000in}{1.128280in}}%
\pgfpathlineto{\pgfqpoint{3.490000in}{1.135800in}}%
\pgfpathlineto{\pgfqpoint{3.393750in}{1.135800in}}%
\pgfpathlineto{\pgfqpoint{3.393750in}{1.128280in}}%
\pgfusepath{fill}%
\end{pgfscope}%
\begin{pgfscope}%
\pgfpathrectangle{\pgfqpoint{3.393750in}{0.699667in}}{\pgfqpoint{0.096250in}{1.925000in}}%
\pgfusepath{clip}%
\pgfsetbuttcap%
\pgfsetroundjoin%
\definecolor{currentfill}{rgb}{0.278493,0.061978,0.469190}%
\pgfsetfillcolor{currentfill}%
\pgfsetfillopacity{0.800000}%
\pgfsetlinewidth{0.000000pt}%
\definecolor{currentstroke}{rgb}{0.000000,0.000000,0.000000}%
\pgfsetstrokecolor{currentstroke}%
\pgfsetdash{}{0pt}%
\pgfpathmoveto{\pgfqpoint{3.393750in}{1.135800in}}%
\pgfpathlineto{\pgfqpoint{3.490000in}{1.135800in}}%
\pgfpathlineto{\pgfqpoint{3.490000in}{1.143320in}}%
\pgfpathlineto{\pgfqpoint{3.393750in}{1.143320in}}%
\pgfpathlineto{\pgfqpoint{3.393750in}{1.135800in}}%
\pgfusepath{fill}%
\end{pgfscope}%
\begin{pgfscope}%
\pgfpathrectangle{\pgfqpoint{3.393750in}{0.699667in}}{\pgfqpoint{0.096250in}{1.925000in}}%
\pgfusepath{clip}%
\pgfsetbuttcap%
\pgfsetroundjoin%
\definecolor{currentfill}{rgb}{0.284951,0.063168,0.472451}%
\pgfsetfillcolor{currentfill}%
\pgfsetfillopacity{0.800000}%
\pgfsetlinewidth{0.000000pt}%
\definecolor{currentstroke}{rgb}{0.000000,0.000000,0.000000}%
\pgfsetstrokecolor{currentstroke}%
\pgfsetdash{}{0pt}%
\pgfpathmoveto{\pgfqpoint{3.393750in}{1.143320in}}%
\pgfpathlineto{\pgfqpoint{3.490000in}{1.143320in}}%
\pgfpathlineto{\pgfqpoint{3.490000in}{1.150839in}}%
\pgfpathlineto{\pgfqpoint{3.393750in}{1.150839in}}%
\pgfpathlineto{\pgfqpoint{3.393750in}{1.143320in}}%
\pgfusepath{fill}%
\end{pgfscope}%
\begin{pgfscope}%
\pgfpathrectangle{\pgfqpoint{3.393750in}{0.699667in}}{\pgfqpoint{0.096250in}{1.925000in}}%
\pgfusepath{clip}%
\pgfsetbuttcap%
\pgfsetroundjoin%
\definecolor{currentfill}{rgb}{0.291366,0.064553,0.475462}%
\pgfsetfillcolor{currentfill}%
\pgfsetfillopacity{0.800000}%
\pgfsetlinewidth{0.000000pt}%
\definecolor{currentstroke}{rgb}{0.000000,0.000000,0.000000}%
\pgfsetstrokecolor{currentstroke}%
\pgfsetdash{}{0pt}%
\pgfpathmoveto{\pgfqpoint{3.393750in}{1.150839in}}%
\pgfpathlineto{\pgfqpoint{3.490000in}{1.150839in}}%
\pgfpathlineto{\pgfqpoint{3.490000in}{1.158359in}}%
\pgfpathlineto{\pgfqpoint{3.393750in}{1.158359in}}%
\pgfpathlineto{\pgfqpoint{3.393750in}{1.150839in}}%
\pgfusepath{fill}%
\end{pgfscope}%
\begin{pgfscope}%
\pgfpathrectangle{\pgfqpoint{3.393750in}{0.699667in}}{\pgfqpoint{0.096250in}{1.925000in}}%
\pgfusepath{clip}%
\pgfsetbuttcap%
\pgfsetroundjoin%
\definecolor{currentfill}{rgb}{0.297740,0.066117,0.478243}%
\pgfsetfillcolor{currentfill}%
\pgfsetfillopacity{0.800000}%
\pgfsetlinewidth{0.000000pt}%
\definecolor{currentstroke}{rgb}{0.000000,0.000000,0.000000}%
\pgfsetstrokecolor{currentstroke}%
\pgfsetdash{}{0pt}%
\pgfpathmoveto{\pgfqpoint{3.393750in}{1.158359in}}%
\pgfpathlineto{\pgfqpoint{3.490000in}{1.158359in}}%
\pgfpathlineto{\pgfqpoint{3.490000in}{1.165878in}}%
\pgfpathlineto{\pgfqpoint{3.393750in}{1.165878in}}%
\pgfpathlineto{\pgfqpoint{3.393750in}{1.158359in}}%
\pgfusepath{fill}%
\end{pgfscope}%
\begin{pgfscope}%
\pgfpathrectangle{\pgfqpoint{3.393750in}{0.699667in}}{\pgfqpoint{0.096250in}{1.925000in}}%
\pgfusepath{clip}%
\pgfsetbuttcap%
\pgfsetroundjoin%
\definecolor{currentfill}{rgb}{0.304081,0.067835,0.480812}%
\pgfsetfillcolor{currentfill}%
\pgfsetfillopacity{0.800000}%
\pgfsetlinewidth{0.000000pt}%
\definecolor{currentstroke}{rgb}{0.000000,0.000000,0.000000}%
\pgfsetstrokecolor{currentstroke}%
\pgfsetdash{}{0pt}%
\pgfpathmoveto{\pgfqpoint{3.393750in}{1.165878in}}%
\pgfpathlineto{\pgfqpoint{3.490000in}{1.165878in}}%
\pgfpathlineto{\pgfqpoint{3.490000in}{1.173398in}}%
\pgfpathlineto{\pgfqpoint{3.393750in}{1.173398in}}%
\pgfpathlineto{\pgfqpoint{3.393750in}{1.165878in}}%
\pgfusepath{fill}%
\end{pgfscope}%
\begin{pgfscope}%
\pgfpathrectangle{\pgfqpoint{3.393750in}{0.699667in}}{\pgfqpoint{0.096250in}{1.925000in}}%
\pgfusepath{clip}%
\pgfsetbuttcap%
\pgfsetroundjoin%
\definecolor{currentfill}{rgb}{0.310382,0.069702,0.483186}%
\pgfsetfillcolor{currentfill}%
\pgfsetfillopacity{0.800000}%
\pgfsetlinewidth{0.000000pt}%
\definecolor{currentstroke}{rgb}{0.000000,0.000000,0.000000}%
\pgfsetstrokecolor{currentstroke}%
\pgfsetdash{}{0pt}%
\pgfpathmoveto{\pgfqpoint{3.393750in}{1.173398in}}%
\pgfpathlineto{\pgfqpoint{3.490000in}{1.173398in}}%
\pgfpathlineto{\pgfqpoint{3.490000in}{1.180917in}}%
\pgfpathlineto{\pgfqpoint{3.393750in}{1.180917in}}%
\pgfpathlineto{\pgfqpoint{3.393750in}{1.173398in}}%
\pgfusepath{fill}%
\end{pgfscope}%
\begin{pgfscope}%
\pgfpathrectangle{\pgfqpoint{3.393750in}{0.699667in}}{\pgfqpoint{0.096250in}{1.925000in}}%
\pgfusepath{clip}%
\pgfsetbuttcap%
\pgfsetroundjoin%
\definecolor{currentfill}{rgb}{0.316654,0.071690,0.485380}%
\pgfsetfillcolor{currentfill}%
\pgfsetfillopacity{0.800000}%
\pgfsetlinewidth{0.000000pt}%
\definecolor{currentstroke}{rgb}{0.000000,0.000000,0.000000}%
\pgfsetstrokecolor{currentstroke}%
\pgfsetdash{}{0pt}%
\pgfpathmoveto{\pgfqpoint{3.393750in}{1.180917in}}%
\pgfpathlineto{\pgfqpoint{3.490000in}{1.180917in}}%
\pgfpathlineto{\pgfqpoint{3.490000in}{1.188437in}}%
\pgfpathlineto{\pgfqpoint{3.393750in}{1.188437in}}%
\pgfpathlineto{\pgfqpoint{3.393750in}{1.180917in}}%
\pgfusepath{fill}%
\end{pgfscope}%
\begin{pgfscope}%
\pgfpathrectangle{\pgfqpoint{3.393750in}{0.699667in}}{\pgfqpoint{0.096250in}{1.925000in}}%
\pgfusepath{clip}%
\pgfsetbuttcap%
\pgfsetroundjoin%
\definecolor{currentfill}{rgb}{0.322899,0.073782,0.487408}%
\pgfsetfillcolor{currentfill}%
\pgfsetfillopacity{0.800000}%
\pgfsetlinewidth{0.000000pt}%
\definecolor{currentstroke}{rgb}{0.000000,0.000000,0.000000}%
\pgfsetstrokecolor{currentstroke}%
\pgfsetdash{}{0pt}%
\pgfpathmoveto{\pgfqpoint{3.393750in}{1.188437in}}%
\pgfpathlineto{\pgfqpoint{3.490000in}{1.188437in}}%
\pgfpathlineto{\pgfqpoint{3.490000in}{1.195956in}}%
\pgfpathlineto{\pgfqpoint{3.393750in}{1.195956in}}%
\pgfpathlineto{\pgfqpoint{3.393750in}{1.188437in}}%
\pgfusepath{fill}%
\end{pgfscope}%
\begin{pgfscope}%
\pgfpathrectangle{\pgfqpoint{3.393750in}{0.699667in}}{\pgfqpoint{0.096250in}{1.925000in}}%
\pgfusepath{clip}%
\pgfsetbuttcap%
\pgfsetroundjoin%
\definecolor{currentfill}{rgb}{0.329114,0.075972,0.489287}%
\pgfsetfillcolor{currentfill}%
\pgfsetfillopacity{0.800000}%
\pgfsetlinewidth{0.000000pt}%
\definecolor{currentstroke}{rgb}{0.000000,0.000000,0.000000}%
\pgfsetstrokecolor{currentstroke}%
\pgfsetdash{}{0pt}%
\pgfpathmoveto{\pgfqpoint{3.393750in}{1.195956in}}%
\pgfpathlineto{\pgfqpoint{3.490000in}{1.195956in}}%
\pgfpathlineto{\pgfqpoint{3.490000in}{1.203476in}}%
\pgfpathlineto{\pgfqpoint{3.393750in}{1.203476in}}%
\pgfpathlineto{\pgfqpoint{3.393750in}{1.195956in}}%
\pgfusepath{fill}%
\end{pgfscope}%
\begin{pgfscope}%
\pgfpathrectangle{\pgfqpoint{3.393750in}{0.699667in}}{\pgfqpoint{0.096250in}{1.925000in}}%
\pgfusepath{clip}%
\pgfsetbuttcap%
\pgfsetroundjoin%
\definecolor{currentfill}{rgb}{0.335308,0.078236,0.491024}%
\pgfsetfillcolor{currentfill}%
\pgfsetfillopacity{0.800000}%
\pgfsetlinewidth{0.000000pt}%
\definecolor{currentstroke}{rgb}{0.000000,0.000000,0.000000}%
\pgfsetstrokecolor{currentstroke}%
\pgfsetdash{}{0pt}%
\pgfpathmoveto{\pgfqpoint{3.393750in}{1.203476in}}%
\pgfpathlineto{\pgfqpoint{3.490000in}{1.203476in}}%
\pgfpathlineto{\pgfqpoint{3.490000in}{1.210995in}}%
\pgfpathlineto{\pgfqpoint{3.393750in}{1.210995in}}%
\pgfpathlineto{\pgfqpoint{3.393750in}{1.203476in}}%
\pgfusepath{fill}%
\end{pgfscope}%
\begin{pgfscope}%
\pgfpathrectangle{\pgfqpoint{3.393750in}{0.699667in}}{\pgfqpoint{0.096250in}{1.925000in}}%
\pgfusepath{clip}%
\pgfsetbuttcap%
\pgfsetroundjoin%
\definecolor{currentfill}{rgb}{0.341482,0.080564,0.492631}%
\pgfsetfillcolor{currentfill}%
\pgfsetfillopacity{0.800000}%
\pgfsetlinewidth{0.000000pt}%
\definecolor{currentstroke}{rgb}{0.000000,0.000000,0.000000}%
\pgfsetstrokecolor{currentstroke}%
\pgfsetdash{}{0pt}%
\pgfpathmoveto{\pgfqpoint{3.393750in}{1.210995in}}%
\pgfpathlineto{\pgfqpoint{3.490000in}{1.210995in}}%
\pgfpathlineto{\pgfqpoint{3.490000in}{1.218515in}}%
\pgfpathlineto{\pgfqpoint{3.393750in}{1.218515in}}%
\pgfpathlineto{\pgfqpoint{3.393750in}{1.210995in}}%
\pgfusepath{fill}%
\end{pgfscope}%
\begin{pgfscope}%
\pgfpathrectangle{\pgfqpoint{3.393750in}{0.699667in}}{\pgfqpoint{0.096250in}{1.925000in}}%
\pgfusepath{clip}%
\pgfsetbuttcap%
\pgfsetroundjoin%
\definecolor{currentfill}{rgb}{0.347636,0.082946,0.494121}%
\pgfsetfillcolor{currentfill}%
\pgfsetfillopacity{0.800000}%
\pgfsetlinewidth{0.000000pt}%
\definecolor{currentstroke}{rgb}{0.000000,0.000000,0.000000}%
\pgfsetstrokecolor{currentstroke}%
\pgfsetdash{}{0pt}%
\pgfpathmoveto{\pgfqpoint{3.393750in}{1.218515in}}%
\pgfpathlineto{\pgfqpoint{3.490000in}{1.218515in}}%
\pgfpathlineto{\pgfqpoint{3.490000in}{1.226034in}}%
\pgfpathlineto{\pgfqpoint{3.393750in}{1.226034in}}%
\pgfpathlineto{\pgfqpoint{3.393750in}{1.218515in}}%
\pgfusepath{fill}%
\end{pgfscope}%
\begin{pgfscope}%
\pgfpathrectangle{\pgfqpoint{3.393750in}{0.699667in}}{\pgfqpoint{0.096250in}{1.925000in}}%
\pgfusepath{clip}%
\pgfsetbuttcap%
\pgfsetroundjoin%
\definecolor{currentfill}{rgb}{0.353773,0.085373,0.495501}%
\pgfsetfillcolor{currentfill}%
\pgfsetfillopacity{0.800000}%
\pgfsetlinewidth{0.000000pt}%
\definecolor{currentstroke}{rgb}{0.000000,0.000000,0.000000}%
\pgfsetstrokecolor{currentstroke}%
\pgfsetdash{}{0pt}%
\pgfpathmoveto{\pgfqpoint{3.393750in}{1.226034in}}%
\pgfpathlineto{\pgfqpoint{3.490000in}{1.226034in}}%
\pgfpathlineto{\pgfqpoint{3.490000in}{1.233554in}}%
\pgfpathlineto{\pgfqpoint{3.393750in}{1.233554in}}%
\pgfpathlineto{\pgfqpoint{3.393750in}{1.226034in}}%
\pgfusepath{fill}%
\end{pgfscope}%
\begin{pgfscope}%
\pgfpathrectangle{\pgfqpoint{3.393750in}{0.699667in}}{\pgfqpoint{0.096250in}{1.925000in}}%
\pgfusepath{clip}%
\pgfsetbuttcap%
\pgfsetroundjoin%
\definecolor{currentfill}{rgb}{0.359898,0.087831,0.496778}%
\pgfsetfillcolor{currentfill}%
\pgfsetfillopacity{0.800000}%
\pgfsetlinewidth{0.000000pt}%
\definecolor{currentstroke}{rgb}{0.000000,0.000000,0.000000}%
\pgfsetstrokecolor{currentstroke}%
\pgfsetdash{}{0pt}%
\pgfpathmoveto{\pgfqpoint{3.393750in}{1.233554in}}%
\pgfpathlineto{\pgfqpoint{3.490000in}{1.233554in}}%
\pgfpathlineto{\pgfqpoint{3.490000in}{1.241073in}}%
\pgfpathlineto{\pgfqpoint{3.393750in}{1.241073in}}%
\pgfpathlineto{\pgfqpoint{3.393750in}{1.233554in}}%
\pgfusepath{fill}%
\end{pgfscope}%
\begin{pgfscope}%
\pgfpathrectangle{\pgfqpoint{3.393750in}{0.699667in}}{\pgfqpoint{0.096250in}{1.925000in}}%
\pgfusepath{clip}%
\pgfsetbuttcap%
\pgfsetroundjoin%
\definecolor{currentfill}{rgb}{0.366012,0.090314,0.497960}%
\pgfsetfillcolor{currentfill}%
\pgfsetfillopacity{0.800000}%
\pgfsetlinewidth{0.000000pt}%
\definecolor{currentstroke}{rgb}{0.000000,0.000000,0.000000}%
\pgfsetstrokecolor{currentstroke}%
\pgfsetdash{}{0pt}%
\pgfpathmoveto{\pgfqpoint{3.393750in}{1.241073in}}%
\pgfpathlineto{\pgfqpoint{3.490000in}{1.241073in}}%
\pgfpathlineto{\pgfqpoint{3.490000in}{1.248593in}}%
\pgfpathlineto{\pgfqpoint{3.393750in}{1.248593in}}%
\pgfpathlineto{\pgfqpoint{3.393750in}{1.241073in}}%
\pgfusepath{fill}%
\end{pgfscope}%
\begin{pgfscope}%
\pgfpathrectangle{\pgfqpoint{3.393750in}{0.699667in}}{\pgfqpoint{0.096250in}{1.925000in}}%
\pgfusepath{clip}%
\pgfsetbuttcap%
\pgfsetroundjoin%
\definecolor{currentfill}{rgb}{0.372116,0.092816,0.499053}%
\pgfsetfillcolor{currentfill}%
\pgfsetfillopacity{0.800000}%
\pgfsetlinewidth{0.000000pt}%
\definecolor{currentstroke}{rgb}{0.000000,0.000000,0.000000}%
\pgfsetstrokecolor{currentstroke}%
\pgfsetdash{}{0pt}%
\pgfpathmoveto{\pgfqpoint{3.393750in}{1.248593in}}%
\pgfpathlineto{\pgfqpoint{3.490000in}{1.248593in}}%
\pgfpathlineto{\pgfqpoint{3.490000in}{1.256112in}}%
\pgfpathlineto{\pgfqpoint{3.393750in}{1.256112in}}%
\pgfpathlineto{\pgfqpoint{3.393750in}{1.248593in}}%
\pgfusepath{fill}%
\end{pgfscope}%
\begin{pgfscope}%
\pgfpathrectangle{\pgfqpoint{3.393750in}{0.699667in}}{\pgfqpoint{0.096250in}{1.925000in}}%
\pgfusepath{clip}%
\pgfsetbuttcap%
\pgfsetroundjoin%
\definecolor{currentfill}{rgb}{0.378211,0.095332,0.500067}%
\pgfsetfillcolor{currentfill}%
\pgfsetfillopacity{0.800000}%
\pgfsetlinewidth{0.000000pt}%
\definecolor{currentstroke}{rgb}{0.000000,0.000000,0.000000}%
\pgfsetstrokecolor{currentstroke}%
\pgfsetdash{}{0pt}%
\pgfpathmoveto{\pgfqpoint{3.393750in}{1.256112in}}%
\pgfpathlineto{\pgfqpoint{3.490000in}{1.256112in}}%
\pgfpathlineto{\pgfqpoint{3.490000in}{1.263632in}}%
\pgfpathlineto{\pgfqpoint{3.393750in}{1.263632in}}%
\pgfpathlineto{\pgfqpoint{3.393750in}{1.256112in}}%
\pgfusepath{fill}%
\end{pgfscope}%
\begin{pgfscope}%
\pgfpathrectangle{\pgfqpoint{3.393750in}{0.699667in}}{\pgfqpoint{0.096250in}{1.925000in}}%
\pgfusepath{clip}%
\pgfsetbuttcap%
\pgfsetroundjoin%
\definecolor{currentfill}{rgb}{0.384299,0.097855,0.501002}%
\pgfsetfillcolor{currentfill}%
\pgfsetfillopacity{0.800000}%
\pgfsetlinewidth{0.000000pt}%
\definecolor{currentstroke}{rgb}{0.000000,0.000000,0.000000}%
\pgfsetstrokecolor{currentstroke}%
\pgfsetdash{}{0pt}%
\pgfpathmoveto{\pgfqpoint{3.393750in}{1.263632in}}%
\pgfpathlineto{\pgfqpoint{3.490000in}{1.263632in}}%
\pgfpathlineto{\pgfqpoint{3.490000in}{1.271152in}}%
\pgfpathlineto{\pgfqpoint{3.393750in}{1.271152in}}%
\pgfpathlineto{\pgfqpoint{3.393750in}{1.263632in}}%
\pgfusepath{fill}%
\end{pgfscope}%
\begin{pgfscope}%
\pgfpathrectangle{\pgfqpoint{3.393750in}{0.699667in}}{\pgfqpoint{0.096250in}{1.925000in}}%
\pgfusepath{clip}%
\pgfsetbuttcap%
\pgfsetroundjoin%
\definecolor{currentfill}{rgb}{0.390384,0.100379,0.501864}%
\pgfsetfillcolor{currentfill}%
\pgfsetfillopacity{0.800000}%
\pgfsetlinewidth{0.000000pt}%
\definecolor{currentstroke}{rgb}{0.000000,0.000000,0.000000}%
\pgfsetstrokecolor{currentstroke}%
\pgfsetdash{}{0pt}%
\pgfpathmoveto{\pgfqpoint{3.393750in}{1.271152in}}%
\pgfpathlineto{\pgfqpoint{3.490000in}{1.271152in}}%
\pgfpathlineto{\pgfqpoint{3.490000in}{1.278671in}}%
\pgfpathlineto{\pgfqpoint{3.393750in}{1.278671in}}%
\pgfpathlineto{\pgfqpoint{3.393750in}{1.271152in}}%
\pgfusepath{fill}%
\end{pgfscope}%
\begin{pgfscope}%
\pgfpathrectangle{\pgfqpoint{3.393750in}{0.699667in}}{\pgfqpoint{0.096250in}{1.925000in}}%
\pgfusepath{clip}%
\pgfsetbuttcap%
\pgfsetroundjoin%
\definecolor{currentfill}{rgb}{0.396467,0.102902,0.502658}%
\pgfsetfillcolor{currentfill}%
\pgfsetfillopacity{0.800000}%
\pgfsetlinewidth{0.000000pt}%
\definecolor{currentstroke}{rgb}{0.000000,0.000000,0.000000}%
\pgfsetstrokecolor{currentstroke}%
\pgfsetdash{}{0pt}%
\pgfpathmoveto{\pgfqpoint{3.393750in}{1.278671in}}%
\pgfpathlineto{\pgfqpoint{3.490000in}{1.278671in}}%
\pgfpathlineto{\pgfqpoint{3.490000in}{1.286191in}}%
\pgfpathlineto{\pgfqpoint{3.393750in}{1.286191in}}%
\pgfpathlineto{\pgfqpoint{3.393750in}{1.278671in}}%
\pgfusepath{fill}%
\end{pgfscope}%
\begin{pgfscope}%
\pgfpathrectangle{\pgfqpoint{3.393750in}{0.699667in}}{\pgfqpoint{0.096250in}{1.925000in}}%
\pgfusepath{clip}%
\pgfsetbuttcap%
\pgfsetroundjoin%
\definecolor{currentfill}{rgb}{0.402548,0.105420,0.503386}%
\pgfsetfillcolor{currentfill}%
\pgfsetfillopacity{0.800000}%
\pgfsetlinewidth{0.000000pt}%
\definecolor{currentstroke}{rgb}{0.000000,0.000000,0.000000}%
\pgfsetstrokecolor{currentstroke}%
\pgfsetdash{}{0pt}%
\pgfpathmoveto{\pgfqpoint{3.393750in}{1.286191in}}%
\pgfpathlineto{\pgfqpoint{3.490000in}{1.286191in}}%
\pgfpathlineto{\pgfqpoint{3.490000in}{1.293710in}}%
\pgfpathlineto{\pgfqpoint{3.393750in}{1.293710in}}%
\pgfpathlineto{\pgfqpoint{3.393750in}{1.286191in}}%
\pgfusepath{fill}%
\end{pgfscope}%
\begin{pgfscope}%
\pgfpathrectangle{\pgfqpoint{3.393750in}{0.699667in}}{\pgfqpoint{0.096250in}{1.925000in}}%
\pgfusepath{clip}%
\pgfsetbuttcap%
\pgfsetroundjoin%
\definecolor{currentfill}{rgb}{0.408629,0.107930,0.504052}%
\pgfsetfillcolor{currentfill}%
\pgfsetfillopacity{0.800000}%
\pgfsetlinewidth{0.000000pt}%
\definecolor{currentstroke}{rgb}{0.000000,0.000000,0.000000}%
\pgfsetstrokecolor{currentstroke}%
\pgfsetdash{}{0pt}%
\pgfpathmoveto{\pgfqpoint{3.393750in}{1.293710in}}%
\pgfpathlineto{\pgfqpoint{3.490000in}{1.293710in}}%
\pgfpathlineto{\pgfqpoint{3.490000in}{1.301230in}}%
\pgfpathlineto{\pgfqpoint{3.393750in}{1.301230in}}%
\pgfpathlineto{\pgfqpoint{3.393750in}{1.293710in}}%
\pgfusepath{fill}%
\end{pgfscope}%
\begin{pgfscope}%
\pgfpathrectangle{\pgfqpoint{3.393750in}{0.699667in}}{\pgfqpoint{0.096250in}{1.925000in}}%
\pgfusepath{clip}%
\pgfsetbuttcap%
\pgfsetroundjoin%
\definecolor{currentfill}{rgb}{0.414709,0.110431,0.504662}%
\pgfsetfillcolor{currentfill}%
\pgfsetfillopacity{0.800000}%
\pgfsetlinewidth{0.000000pt}%
\definecolor{currentstroke}{rgb}{0.000000,0.000000,0.000000}%
\pgfsetstrokecolor{currentstroke}%
\pgfsetdash{}{0pt}%
\pgfpathmoveto{\pgfqpoint{3.393750in}{1.301230in}}%
\pgfpathlineto{\pgfqpoint{3.490000in}{1.301230in}}%
\pgfpathlineto{\pgfqpoint{3.490000in}{1.308749in}}%
\pgfpathlineto{\pgfqpoint{3.393750in}{1.308749in}}%
\pgfpathlineto{\pgfqpoint{3.393750in}{1.301230in}}%
\pgfusepath{fill}%
\end{pgfscope}%
\begin{pgfscope}%
\pgfpathrectangle{\pgfqpoint{3.393750in}{0.699667in}}{\pgfqpoint{0.096250in}{1.925000in}}%
\pgfusepath{clip}%
\pgfsetbuttcap%
\pgfsetroundjoin%
\definecolor{currentfill}{rgb}{0.420791,0.112920,0.505215}%
\pgfsetfillcolor{currentfill}%
\pgfsetfillopacity{0.800000}%
\pgfsetlinewidth{0.000000pt}%
\definecolor{currentstroke}{rgb}{0.000000,0.000000,0.000000}%
\pgfsetstrokecolor{currentstroke}%
\pgfsetdash{}{0pt}%
\pgfpathmoveto{\pgfqpoint{3.393750in}{1.308749in}}%
\pgfpathlineto{\pgfqpoint{3.490000in}{1.308749in}}%
\pgfpathlineto{\pgfqpoint{3.490000in}{1.316269in}}%
\pgfpathlineto{\pgfqpoint{3.393750in}{1.316269in}}%
\pgfpathlineto{\pgfqpoint{3.393750in}{1.308749in}}%
\pgfusepath{fill}%
\end{pgfscope}%
\begin{pgfscope}%
\pgfpathrectangle{\pgfqpoint{3.393750in}{0.699667in}}{\pgfqpoint{0.096250in}{1.925000in}}%
\pgfusepath{clip}%
\pgfsetbuttcap%
\pgfsetroundjoin%
\definecolor{currentfill}{rgb}{0.426877,0.115395,0.505714}%
\pgfsetfillcolor{currentfill}%
\pgfsetfillopacity{0.800000}%
\pgfsetlinewidth{0.000000pt}%
\definecolor{currentstroke}{rgb}{0.000000,0.000000,0.000000}%
\pgfsetstrokecolor{currentstroke}%
\pgfsetdash{}{0pt}%
\pgfpathmoveto{\pgfqpoint{3.393750in}{1.316269in}}%
\pgfpathlineto{\pgfqpoint{3.490000in}{1.316269in}}%
\pgfpathlineto{\pgfqpoint{3.490000in}{1.323788in}}%
\pgfpathlineto{\pgfqpoint{3.393750in}{1.323788in}}%
\pgfpathlineto{\pgfqpoint{3.393750in}{1.316269in}}%
\pgfusepath{fill}%
\end{pgfscope}%
\begin{pgfscope}%
\pgfpathrectangle{\pgfqpoint{3.393750in}{0.699667in}}{\pgfqpoint{0.096250in}{1.925000in}}%
\pgfusepath{clip}%
\pgfsetbuttcap%
\pgfsetroundjoin%
\definecolor{currentfill}{rgb}{0.432967,0.117855,0.506160}%
\pgfsetfillcolor{currentfill}%
\pgfsetfillopacity{0.800000}%
\pgfsetlinewidth{0.000000pt}%
\definecolor{currentstroke}{rgb}{0.000000,0.000000,0.000000}%
\pgfsetstrokecolor{currentstroke}%
\pgfsetdash{}{0pt}%
\pgfpathmoveto{\pgfqpoint{3.393750in}{1.323788in}}%
\pgfpathlineto{\pgfqpoint{3.490000in}{1.323788in}}%
\pgfpathlineto{\pgfqpoint{3.490000in}{1.331308in}}%
\pgfpathlineto{\pgfqpoint{3.393750in}{1.331308in}}%
\pgfpathlineto{\pgfqpoint{3.393750in}{1.323788in}}%
\pgfusepath{fill}%
\end{pgfscope}%
\begin{pgfscope}%
\pgfpathrectangle{\pgfqpoint{3.393750in}{0.699667in}}{\pgfqpoint{0.096250in}{1.925000in}}%
\pgfusepath{clip}%
\pgfsetbuttcap%
\pgfsetroundjoin%
\definecolor{currentfill}{rgb}{0.439062,0.120298,0.506555}%
\pgfsetfillcolor{currentfill}%
\pgfsetfillopacity{0.800000}%
\pgfsetlinewidth{0.000000pt}%
\definecolor{currentstroke}{rgb}{0.000000,0.000000,0.000000}%
\pgfsetstrokecolor{currentstroke}%
\pgfsetdash{}{0pt}%
\pgfpathmoveto{\pgfqpoint{3.393750in}{1.331308in}}%
\pgfpathlineto{\pgfqpoint{3.490000in}{1.331308in}}%
\pgfpathlineto{\pgfqpoint{3.490000in}{1.338827in}}%
\pgfpathlineto{\pgfqpoint{3.393750in}{1.338827in}}%
\pgfpathlineto{\pgfqpoint{3.393750in}{1.331308in}}%
\pgfusepath{fill}%
\end{pgfscope}%
\begin{pgfscope}%
\pgfpathrectangle{\pgfqpoint{3.393750in}{0.699667in}}{\pgfqpoint{0.096250in}{1.925000in}}%
\pgfusepath{clip}%
\pgfsetbuttcap%
\pgfsetroundjoin%
\definecolor{currentfill}{rgb}{0.445163,0.122724,0.506901}%
\pgfsetfillcolor{currentfill}%
\pgfsetfillopacity{0.800000}%
\pgfsetlinewidth{0.000000pt}%
\definecolor{currentstroke}{rgb}{0.000000,0.000000,0.000000}%
\pgfsetstrokecolor{currentstroke}%
\pgfsetdash{}{0pt}%
\pgfpathmoveto{\pgfqpoint{3.393750in}{1.338827in}}%
\pgfpathlineto{\pgfqpoint{3.490000in}{1.338827in}}%
\pgfpathlineto{\pgfqpoint{3.490000in}{1.346347in}}%
\pgfpathlineto{\pgfqpoint{3.393750in}{1.346347in}}%
\pgfpathlineto{\pgfqpoint{3.393750in}{1.338827in}}%
\pgfusepath{fill}%
\end{pgfscope}%
\begin{pgfscope}%
\pgfpathrectangle{\pgfqpoint{3.393750in}{0.699667in}}{\pgfqpoint{0.096250in}{1.925000in}}%
\pgfusepath{clip}%
\pgfsetbuttcap%
\pgfsetroundjoin%
\definecolor{currentfill}{rgb}{0.451271,0.125132,0.507198}%
\pgfsetfillcolor{currentfill}%
\pgfsetfillopacity{0.800000}%
\pgfsetlinewidth{0.000000pt}%
\definecolor{currentstroke}{rgb}{0.000000,0.000000,0.000000}%
\pgfsetstrokecolor{currentstroke}%
\pgfsetdash{}{0pt}%
\pgfpathmoveto{\pgfqpoint{3.393750in}{1.346347in}}%
\pgfpathlineto{\pgfqpoint{3.490000in}{1.346347in}}%
\pgfpathlineto{\pgfqpoint{3.490000in}{1.353866in}}%
\pgfpathlineto{\pgfqpoint{3.393750in}{1.353866in}}%
\pgfpathlineto{\pgfqpoint{3.393750in}{1.346347in}}%
\pgfusepath{fill}%
\end{pgfscope}%
\begin{pgfscope}%
\pgfpathrectangle{\pgfqpoint{3.393750in}{0.699667in}}{\pgfqpoint{0.096250in}{1.925000in}}%
\pgfusepath{clip}%
\pgfsetbuttcap%
\pgfsetroundjoin%
\definecolor{currentfill}{rgb}{0.457386,0.127522,0.507448}%
\pgfsetfillcolor{currentfill}%
\pgfsetfillopacity{0.800000}%
\pgfsetlinewidth{0.000000pt}%
\definecolor{currentstroke}{rgb}{0.000000,0.000000,0.000000}%
\pgfsetstrokecolor{currentstroke}%
\pgfsetdash{}{0pt}%
\pgfpathmoveto{\pgfqpoint{3.393750in}{1.353866in}}%
\pgfpathlineto{\pgfqpoint{3.490000in}{1.353866in}}%
\pgfpathlineto{\pgfqpoint{3.490000in}{1.361386in}}%
\pgfpathlineto{\pgfqpoint{3.393750in}{1.361386in}}%
\pgfpathlineto{\pgfqpoint{3.393750in}{1.353866in}}%
\pgfusepath{fill}%
\end{pgfscope}%
\begin{pgfscope}%
\pgfpathrectangle{\pgfqpoint{3.393750in}{0.699667in}}{\pgfqpoint{0.096250in}{1.925000in}}%
\pgfusepath{clip}%
\pgfsetbuttcap%
\pgfsetroundjoin%
\definecolor{currentfill}{rgb}{0.463508,0.129893,0.507652}%
\pgfsetfillcolor{currentfill}%
\pgfsetfillopacity{0.800000}%
\pgfsetlinewidth{0.000000pt}%
\definecolor{currentstroke}{rgb}{0.000000,0.000000,0.000000}%
\pgfsetstrokecolor{currentstroke}%
\pgfsetdash{}{0pt}%
\pgfpathmoveto{\pgfqpoint{3.393750in}{1.361386in}}%
\pgfpathlineto{\pgfqpoint{3.490000in}{1.361386in}}%
\pgfpathlineto{\pgfqpoint{3.490000in}{1.368905in}}%
\pgfpathlineto{\pgfqpoint{3.393750in}{1.368905in}}%
\pgfpathlineto{\pgfqpoint{3.393750in}{1.361386in}}%
\pgfusepath{fill}%
\end{pgfscope}%
\begin{pgfscope}%
\pgfpathrectangle{\pgfqpoint{3.393750in}{0.699667in}}{\pgfqpoint{0.096250in}{1.925000in}}%
\pgfusepath{clip}%
\pgfsetbuttcap%
\pgfsetroundjoin%
\definecolor{currentfill}{rgb}{0.469640,0.132245,0.507809}%
\pgfsetfillcolor{currentfill}%
\pgfsetfillopacity{0.800000}%
\pgfsetlinewidth{0.000000pt}%
\definecolor{currentstroke}{rgb}{0.000000,0.000000,0.000000}%
\pgfsetstrokecolor{currentstroke}%
\pgfsetdash{}{0pt}%
\pgfpathmoveto{\pgfqpoint{3.393750in}{1.368905in}}%
\pgfpathlineto{\pgfqpoint{3.490000in}{1.368905in}}%
\pgfpathlineto{\pgfqpoint{3.490000in}{1.376425in}}%
\pgfpathlineto{\pgfqpoint{3.393750in}{1.376425in}}%
\pgfpathlineto{\pgfqpoint{3.393750in}{1.368905in}}%
\pgfusepath{fill}%
\end{pgfscope}%
\begin{pgfscope}%
\pgfpathrectangle{\pgfqpoint{3.393750in}{0.699667in}}{\pgfqpoint{0.096250in}{1.925000in}}%
\pgfusepath{clip}%
\pgfsetbuttcap%
\pgfsetroundjoin%
\definecolor{currentfill}{rgb}{0.475780,0.134577,0.507921}%
\pgfsetfillcolor{currentfill}%
\pgfsetfillopacity{0.800000}%
\pgfsetlinewidth{0.000000pt}%
\definecolor{currentstroke}{rgb}{0.000000,0.000000,0.000000}%
\pgfsetstrokecolor{currentstroke}%
\pgfsetdash{}{0pt}%
\pgfpathmoveto{\pgfqpoint{3.393750in}{1.376425in}}%
\pgfpathlineto{\pgfqpoint{3.490000in}{1.376425in}}%
\pgfpathlineto{\pgfqpoint{3.490000in}{1.383945in}}%
\pgfpathlineto{\pgfqpoint{3.393750in}{1.383945in}}%
\pgfpathlineto{\pgfqpoint{3.393750in}{1.376425in}}%
\pgfusepath{fill}%
\end{pgfscope}%
\begin{pgfscope}%
\pgfpathrectangle{\pgfqpoint{3.393750in}{0.699667in}}{\pgfqpoint{0.096250in}{1.925000in}}%
\pgfusepath{clip}%
\pgfsetbuttcap%
\pgfsetroundjoin%
\definecolor{currentfill}{rgb}{0.481929,0.136891,0.507989}%
\pgfsetfillcolor{currentfill}%
\pgfsetfillopacity{0.800000}%
\pgfsetlinewidth{0.000000pt}%
\definecolor{currentstroke}{rgb}{0.000000,0.000000,0.000000}%
\pgfsetstrokecolor{currentstroke}%
\pgfsetdash{}{0pt}%
\pgfpathmoveto{\pgfqpoint{3.393750in}{1.383945in}}%
\pgfpathlineto{\pgfqpoint{3.490000in}{1.383945in}}%
\pgfpathlineto{\pgfqpoint{3.490000in}{1.391464in}}%
\pgfpathlineto{\pgfqpoint{3.393750in}{1.391464in}}%
\pgfpathlineto{\pgfqpoint{3.393750in}{1.383945in}}%
\pgfusepath{fill}%
\end{pgfscope}%
\begin{pgfscope}%
\pgfpathrectangle{\pgfqpoint{3.393750in}{0.699667in}}{\pgfqpoint{0.096250in}{1.925000in}}%
\pgfusepath{clip}%
\pgfsetbuttcap%
\pgfsetroundjoin%
\definecolor{currentfill}{rgb}{0.488088,0.139186,0.508011}%
\pgfsetfillcolor{currentfill}%
\pgfsetfillopacity{0.800000}%
\pgfsetlinewidth{0.000000pt}%
\definecolor{currentstroke}{rgb}{0.000000,0.000000,0.000000}%
\pgfsetstrokecolor{currentstroke}%
\pgfsetdash{}{0pt}%
\pgfpathmoveto{\pgfqpoint{3.393750in}{1.391464in}}%
\pgfpathlineto{\pgfqpoint{3.490000in}{1.391464in}}%
\pgfpathlineto{\pgfqpoint{3.490000in}{1.398984in}}%
\pgfpathlineto{\pgfqpoint{3.393750in}{1.398984in}}%
\pgfpathlineto{\pgfqpoint{3.393750in}{1.391464in}}%
\pgfusepath{fill}%
\end{pgfscope}%
\begin{pgfscope}%
\pgfpathrectangle{\pgfqpoint{3.393750in}{0.699667in}}{\pgfqpoint{0.096250in}{1.925000in}}%
\pgfusepath{clip}%
\pgfsetbuttcap%
\pgfsetroundjoin%
\definecolor{currentfill}{rgb}{0.494258,0.141462,0.507988}%
\pgfsetfillcolor{currentfill}%
\pgfsetfillopacity{0.800000}%
\pgfsetlinewidth{0.000000pt}%
\definecolor{currentstroke}{rgb}{0.000000,0.000000,0.000000}%
\pgfsetstrokecolor{currentstroke}%
\pgfsetdash{}{0pt}%
\pgfpathmoveto{\pgfqpoint{3.393750in}{1.398984in}}%
\pgfpathlineto{\pgfqpoint{3.490000in}{1.398984in}}%
\pgfpathlineto{\pgfqpoint{3.490000in}{1.406503in}}%
\pgfpathlineto{\pgfqpoint{3.393750in}{1.406503in}}%
\pgfpathlineto{\pgfqpoint{3.393750in}{1.398984in}}%
\pgfusepath{fill}%
\end{pgfscope}%
\begin{pgfscope}%
\pgfpathrectangle{\pgfqpoint{3.393750in}{0.699667in}}{\pgfqpoint{0.096250in}{1.925000in}}%
\pgfusepath{clip}%
\pgfsetbuttcap%
\pgfsetroundjoin%
\definecolor{currentfill}{rgb}{0.500438,0.143719,0.507920}%
\pgfsetfillcolor{currentfill}%
\pgfsetfillopacity{0.800000}%
\pgfsetlinewidth{0.000000pt}%
\definecolor{currentstroke}{rgb}{0.000000,0.000000,0.000000}%
\pgfsetstrokecolor{currentstroke}%
\pgfsetdash{}{0pt}%
\pgfpathmoveto{\pgfqpoint{3.393750in}{1.406503in}}%
\pgfpathlineto{\pgfqpoint{3.490000in}{1.406503in}}%
\pgfpathlineto{\pgfqpoint{3.490000in}{1.414023in}}%
\pgfpathlineto{\pgfqpoint{3.393750in}{1.414023in}}%
\pgfpathlineto{\pgfqpoint{3.393750in}{1.406503in}}%
\pgfusepath{fill}%
\end{pgfscope}%
\begin{pgfscope}%
\pgfpathrectangle{\pgfqpoint{3.393750in}{0.699667in}}{\pgfqpoint{0.096250in}{1.925000in}}%
\pgfusepath{clip}%
\pgfsetbuttcap%
\pgfsetroundjoin%
\definecolor{currentfill}{rgb}{0.506629,0.145958,0.507806}%
\pgfsetfillcolor{currentfill}%
\pgfsetfillopacity{0.800000}%
\pgfsetlinewidth{0.000000pt}%
\definecolor{currentstroke}{rgb}{0.000000,0.000000,0.000000}%
\pgfsetstrokecolor{currentstroke}%
\pgfsetdash{}{0pt}%
\pgfpathmoveto{\pgfqpoint{3.393750in}{1.414023in}}%
\pgfpathlineto{\pgfqpoint{3.490000in}{1.414023in}}%
\pgfpathlineto{\pgfqpoint{3.490000in}{1.421542in}}%
\pgfpathlineto{\pgfqpoint{3.393750in}{1.421542in}}%
\pgfpathlineto{\pgfqpoint{3.393750in}{1.414023in}}%
\pgfusepath{fill}%
\end{pgfscope}%
\begin{pgfscope}%
\pgfpathrectangle{\pgfqpoint{3.393750in}{0.699667in}}{\pgfqpoint{0.096250in}{1.925000in}}%
\pgfusepath{clip}%
\pgfsetbuttcap%
\pgfsetroundjoin%
\definecolor{currentfill}{rgb}{0.512831,0.148179,0.507648}%
\pgfsetfillcolor{currentfill}%
\pgfsetfillopacity{0.800000}%
\pgfsetlinewidth{0.000000pt}%
\definecolor{currentstroke}{rgb}{0.000000,0.000000,0.000000}%
\pgfsetstrokecolor{currentstroke}%
\pgfsetdash{}{0pt}%
\pgfpathmoveto{\pgfqpoint{3.393750in}{1.421542in}}%
\pgfpathlineto{\pgfqpoint{3.490000in}{1.421542in}}%
\pgfpathlineto{\pgfqpoint{3.490000in}{1.429062in}}%
\pgfpathlineto{\pgfqpoint{3.393750in}{1.429062in}}%
\pgfpathlineto{\pgfqpoint{3.393750in}{1.421542in}}%
\pgfusepath{fill}%
\end{pgfscope}%
\begin{pgfscope}%
\pgfpathrectangle{\pgfqpoint{3.393750in}{0.699667in}}{\pgfqpoint{0.096250in}{1.925000in}}%
\pgfusepath{clip}%
\pgfsetbuttcap%
\pgfsetroundjoin%
\definecolor{currentfill}{rgb}{0.519045,0.150383,0.507443}%
\pgfsetfillcolor{currentfill}%
\pgfsetfillopacity{0.800000}%
\pgfsetlinewidth{0.000000pt}%
\definecolor{currentstroke}{rgb}{0.000000,0.000000,0.000000}%
\pgfsetstrokecolor{currentstroke}%
\pgfsetdash{}{0pt}%
\pgfpathmoveto{\pgfqpoint{3.393750in}{1.429062in}}%
\pgfpathlineto{\pgfqpoint{3.490000in}{1.429062in}}%
\pgfpathlineto{\pgfqpoint{3.490000in}{1.436581in}}%
\pgfpathlineto{\pgfqpoint{3.393750in}{1.436581in}}%
\pgfpathlineto{\pgfqpoint{3.393750in}{1.429062in}}%
\pgfusepath{fill}%
\end{pgfscope}%
\begin{pgfscope}%
\pgfpathrectangle{\pgfqpoint{3.393750in}{0.699667in}}{\pgfqpoint{0.096250in}{1.925000in}}%
\pgfusepath{clip}%
\pgfsetbuttcap%
\pgfsetroundjoin%
\definecolor{currentfill}{rgb}{0.525270,0.152569,0.507192}%
\pgfsetfillcolor{currentfill}%
\pgfsetfillopacity{0.800000}%
\pgfsetlinewidth{0.000000pt}%
\definecolor{currentstroke}{rgb}{0.000000,0.000000,0.000000}%
\pgfsetstrokecolor{currentstroke}%
\pgfsetdash{}{0pt}%
\pgfpathmoveto{\pgfqpoint{3.393750in}{1.436581in}}%
\pgfpathlineto{\pgfqpoint{3.490000in}{1.436581in}}%
\pgfpathlineto{\pgfqpoint{3.490000in}{1.444101in}}%
\pgfpathlineto{\pgfqpoint{3.393750in}{1.444101in}}%
\pgfpathlineto{\pgfqpoint{3.393750in}{1.436581in}}%
\pgfusepath{fill}%
\end{pgfscope}%
\begin{pgfscope}%
\pgfpathrectangle{\pgfqpoint{3.393750in}{0.699667in}}{\pgfqpoint{0.096250in}{1.925000in}}%
\pgfusepath{clip}%
\pgfsetbuttcap%
\pgfsetroundjoin%
\definecolor{currentfill}{rgb}{0.531507,0.154739,0.506895}%
\pgfsetfillcolor{currentfill}%
\pgfsetfillopacity{0.800000}%
\pgfsetlinewidth{0.000000pt}%
\definecolor{currentstroke}{rgb}{0.000000,0.000000,0.000000}%
\pgfsetstrokecolor{currentstroke}%
\pgfsetdash{}{0pt}%
\pgfpathmoveto{\pgfqpoint{3.393750in}{1.444101in}}%
\pgfpathlineto{\pgfqpoint{3.490000in}{1.444101in}}%
\pgfpathlineto{\pgfqpoint{3.490000in}{1.451620in}}%
\pgfpathlineto{\pgfqpoint{3.393750in}{1.451620in}}%
\pgfpathlineto{\pgfqpoint{3.393750in}{1.444101in}}%
\pgfusepath{fill}%
\end{pgfscope}%
\begin{pgfscope}%
\pgfpathrectangle{\pgfqpoint{3.393750in}{0.699667in}}{\pgfqpoint{0.096250in}{1.925000in}}%
\pgfusepath{clip}%
\pgfsetbuttcap%
\pgfsetroundjoin%
\definecolor{currentfill}{rgb}{0.537755,0.156894,0.506551}%
\pgfsetfillcolor{currentfill}%
\pgfsetfillopacity{0.800000}%
\pgfsetlinewidth{0.000000pt}%
\definecolor{currentstroke}{rgb}{0.000000,0.000000,0.000000}%
\pgfsetstrokecolor{currentstroke}%
\pgfsetdash{}{0pt}%
\pgfpathmoveto{\pgfqpoint{3.393750in}{1.451620in}}%
\pgfpathlineto{\pgfqpoint{3.490000in}{1.451620in}}%
\pgfpathlineto{\pgfqpoint{3.490000in}{1.459140in}}%
\pgfpathlineto{\pgfqpoint{3.393750in}{1.459140in}}%
\pgfpathlineto{\pgfqpoint{3.393750in}{1.451620in}}%
\pgfusepath{fill}%
\end{pgfscope}%
\begin{pgfscope}%
\pgfpathrectangle{\pgfqpoint{3.393750in}{0.699667in}}{\pgfqpoint{0.096250in}{1.925000in}}%
\pgfusepath{clip}%
\pgfsetbuttcap%
\pgfsetroundjoin%
\definecolor{currentfill}{rgb}{0.544015,0.159033,0.506159}%
\pgfsetfillcolor{currentfill}%
\pgfsetfillopacity{0.800000}%
\pgfsetlinewidth{0.000000pt}%
\definecolor{currentstroke}{rgb}{0.000000,0.000000,0.000000}%
\pgfsetstrokecolor{currentstroke}%
\pgfsetdash{}{0pt}%
\pgfpathmoveto{\pgfqpoint{3.393750in}{1.459140in}}%
\pgfpathlineto{\pgfqpoint{3.490000in}{1.459140in}}%
\pgfpathlineto{\pgfqpoint{3.490000in}{1.466659in}}%
\pgfpathlineto{\pgfqpoint{3.393750in}{1.466659in}}%
\pgfpathlineto{\pgfqpoint{3.393750in}{1.459140in}}%
\pgfusepath{fill}%
\end{pgfscope}%
\begin{pgfscope}%
\pgfpathrectangle{\pgfqpoint{3.393750in}{0.699667in}}{\pgfqpoint{0.096250in}{1.925000in}}%
\pgfusepath{clip}%
\pgfsetbuttcap%
\pgfsetroundjoin%
\definecolor{currentfill}{rgb}{0.550287,0.161158,0.505719}%
\pgfsetfillcolor{currentfill}%
\pgfsetfillopacity{0.800000}%
\pgfsetlinewidth{0.000000pt}%
\definecolor{currentstroke}{rgb}{0.000000,0.000000,0.000000}%
\pgfsetstrokecolor{currentstroke}%
\pgfsetdash{}{0pt}%
\pgfpathmoveto{\pgfqpoint{3.393750in}{1.466659in}}%
\pgfpathlineto{\pgfqpoint{3.490000in}{1.466659in}}%
\pgfpathlineto{\pgfqpoint{3.490000in}{1.474179in}}%
\pgfpathlineto{\pgfqpoint{3.393750in}{1.474179in}}%
\pgfpathlineto{\pgfqpoint{3.393750in}{1.466659in}}%
\pgfusepath{fill}%
\end{pgfscope}%
\begin{pgfscope}%
\pgfpathrectangle{\pgfqpoint{3.393750in}{0.699667in}}{\pgfqpoint{0.096250in}{1.925000in}}%
\pgfusepath{clip}%
\pgfsetbuttcap%
\pgfsetroundjoin%
\definecolor{currentfill}{rgb}{0.556571,0.163269,0.505230}%
\pgfsetfillcolor{currentfill}%
\pgfsetfillopacity{0.800000}%
\pgfsetlinewidth{0.000000pt}%
\definecolor{currentstroke}{rgb}{0.000000,0.000000,0.000000}%
\pgfsetstrokecolor{currentstroke}%
\pgfsetdash{}{0pt}%
\pgfpathmoveto{\pgfqpoint{3.393750in}{1.474179in}}%
\pgfpathlineto{\pgfqpoint{3.490000in}{1.474179in}}%
\pgfpathlineto{\pgfqpoint{3.490000in}{1.481698in}}%
\pgfpathlineto{\pgfqpoint{3.393750in}{1.481698in}}%
\pgfpathlineto{\pgfqpoint{3.393750in}{1.474179in}}%
\pgfusepath{fill}%
\end{pgfscope}%
\begin{pgfscope}%
\pgfpathrectangle{\pgfqpoint{3.393750in}{0.699667in}}{\pgfqpoint{0.096250in}{1.925000in}}%
\pgfusepath{clip}%
\pgfsetbuttcap%
\pgfsetroundjoin%
\definecolor{currentfill}{rgb}{0.562866,0.165368,0.504692}%
\pgfsetfillcolor{currentfill}%
\pgfsetfillopacity{0.800000}%
\pgfsetlinewidth{0.000000pt}%
\definecolor{currentstroke}{rgb}{0.000000,0.000000,0.000000}%
\pgfsetstrokecolor{currentstroke}%
\pgfsetdash{}{0pt}%
\pgfpathmoveto{\pgfqpoint{3.393750in}{1.481698in}}%
\pgfpathlineto{\pgfqpoint{3.490000in}{1.481698in}}%
\pgfpathlineto{\pgfqpoint{3.490000in}{1.489218in}}%
\pgfpathlineto{\pgfqpoint{3.393750in}{1.489218in}}%
\pgfpathlineto{\pgfqpoint{3.393750in}{1.481698in}}%
\pgfusepath{fill}%
\end{pgfscope}%
\begin{pgfscope}%
\pgfpathrectangle{\pgfqpoint{3.393750in}{0.699667in}}{\pgfqpoint{0.096250in}{1.925000in}}%
\pgfusepath{clip}%
\pgfsetbuttcap%
\pgfsetroundjoin%
\definecolor{currentfill}{rgb}{0.569172,0.167454,0.504105}%
\pgfsetfillcolor{currentfill}%
\pgfsetfillopacity{0.800000}%
\pgfsetlinewidth{0.000000pt}%
\definecolor{currentstroke}{rgb}{0.000000,0.000000,0.000000}%
\pgfsetstrokecolor{currentstroke}%
\pgfsetdash{}{0pt}%
\pgfpathmoveto{\pgfqpoint{3.393750in}{1.489218in}}%
\pgfpathlineto{\pgfqpoint{3.490000in}{1.489218in}}%
\pgfpathlineto{\pgfqpoint{3.490000in}{1.496737in}}%
\pgfpathlineto{\pgfqpoint{3.393750in}{1.496737in}}%
\pgfpathlineto{\pgfqpoint{3.393750in}{1.489218in}}%
\pgfusepath{fill}%
\end{pgfscope}%
\begin{pgfscope}%
\pgfpathrectangle{\pgfqpoint{3.393750in}{0.699667in}}{\pgfqpoint{0.096250in}{1.925000in}}%
\pgfusepath{clip}%
\pgfsetbuttcap%
\pgfsetroundjoin%
\definecolor{currentfill}{rgb}{0.575490,0.169530,0.503466}%
\pgfsetfillcolor{currentfill}%
\pgfsetfillopacity{0.800000}%
\pgfsetlinewidth{0.000000pt}%
\definecolor{currentstroke}{rgb}{0.000000,0.000000,0.000000}%
\pgfsetstrokecolor{currentstroke}%
\pgfsetdash{}{0pt}%
\pgfpathmoveto{\pgfqpoint{3.393750in}{1.496737in}}%
\pgfpathlineto{\pgfqpoint{3.490000in}{1.496737in}}%
\pgfpathlineto{\pgfqpoint{3.490000in}{1.504257in}}%
\pgfpathlineto{\pgfqpoint{3.393750in}{1.504257in}}%
\pgfpathlineto{\pgfqpoint{3.393750in}{1.496737in}}%
\pgfusepath{fill}%
\end{pgfscope}%
\begin{pgfscope}%
\pgfpathrectangle{\pgfqpoint{3.393750in}{0.699667in}}{\pgfqpoint{0.096250in}{1.925000in}}%
\pgfusepath{clip}%
\pgfsetbuttcap%
\pgfsetroundjoin%
\definecolor{currentfill}{rgb}{0.581819,0.171596,0.502777}%
\pgfsetfillcolor{currentfill}%
\pgfsetfillopacity{0.800000}%
\pgfsetlinewidth{0.000000pt}%
\definecolor{currentstroke}{rgb}{0.000000,0.000000,0.000000}%
\pgfsetstrokecolor{currentstroke}%
\pgfsetdash{}{0pt}%
\pgfpathmoveto{\pgfqpoint{3.393750in}{1.504257in}}%
\pgfpathlineto{\pgfqpoint{3.490000in}{1.504257in}}%
\pgfpathlineto{\pgfqpoint{3.490000in}{1.511777in}}%
\pgfpathlineto{\pgfqpoint{3.393750in}{1.511777in}}%
\pgfpathlineto{\pgfqpoint{3.393750in}{1.504257in}}%
\pgfusepath{fill}%
\end{pgfscope}%
\begin{pgfscope}%
\pgfpathrectangle{\pgfqpoint{3.393750in}{0.699667in}}{\pgfqpoint{0.096250in}{1.925000in}}%
\pgfusepath{clip}%
\pgfsetbuttcap%
\pgfsetroundjoin%
\definecolor{currentfill}{rgb}{0.588158,0.173652,0.502035}%
\pgfsetfillcolor{currentfill}%
\pgfsetfillopacity{0.800000}%
\pgfsetlinewidth{0.000000pt}%
\definecolor{currentstroke}{rgb}{0.000000,0.000000,0.000000}%
\pgfsetstrokecolor{currentstroke}%
\pgfsetdash{}{0pt}%
\pgfpathmoveto{\pgfqpoint{3.393750in}{1.511777in}}%
\pgfpathlineto{\pgfqpoint{3.490000in}{1.511777in}}%
\pgfpathlineto{\pgfqpoint{3.490000in}{1.519296in}}%
\pgfpathlineto{\pgfqpoint{3.393750in}{1.519296in}}%
\pgfpathlineto{\pgfqpoint{3.393750in}{1.511777in}}%
\pgfusepath{fill}%
\end{pgfscope}%
\begin{pgfscope}%
\pgfpathrectangle{\pgfqpoint{3.393750in}{0.699667in}}{\pgfqpoint{0.096250in}{1.925000in}}%
\pgfusepath{clip}%
\pgfsetbuttcap%
\pgfsetroundjoin%
\definecolor{currentfill}{rgb}{0.594508,0.175701,0.501241}%
\pgfsetfillcolor{currentfill}%
\pgfsetfillopacity{0.800000}%
\pgfsetlinewidth{0.000000pt}%
\definecolor{currentstroke}{rgb}{0.000000,0.000000,0.000000}%
\pgfsetstrokecolor{currentstroke}%
\pgfsetdash{}{0pt}%
\pgfpathmoveto{\pgfqpoint{3.393750in}{1.519296in}}%
\pgfpathlineto{\pgfqpoint{3.490000in}{1.519296in}}%
\pgfpathlineto{\pgfqpoint{3.490000in}{1.526816in}}%
\pgfpathlineto{\pgfqpoint{3.393750in}{1.526816in}}%
\pgfpathlineto{\pgfqpoint{3.393750in}{1.519296in}}%
\pgfusepath{fill}%
\end{pgfscope}%
\begin{pgfscope}%
\pgfpathrectangle{\pgfqpoint{3.393750in}{0.699667in}}{\pgfqpoint{0.096250in}{1.925000in}}%
\pgfusepath{clip}%
\pgfsetbuttcap%
\pgfsetroundjoin%
\definecolor{currentfill}{rgb}{0.600868,0.177743,0.500394}%
\pgfsetfillcolor{currentfill}%
\pgfsetfillopacity{0.800000}%
\pgfsetlinewidth{0.000000pt}%
\definecolor{currentstroke}{rgb}{0.000000,0.000000,0.000000}%
\pgfsetstrokecolor{currentstroke}%
\pgfsetdash{}{0pt}%
\pgfpathmoveto{\pgfqpoint{3.393750in}{1.526816in}}%
\pgfpathlineto{\pgfqpoint{3.490000in}{1.526816in}}%
\pgfpathlineto{\pgfqpoint{3.490000in}{1.534335in}}%
\pgfpathlineto{\pgfqpoint{3.393750in}{1.534335in}}%
\pgfpathlineto{\pgfqpoint{3.393750in}{1.526816in}}%
\pgfusepath{fill}%
\end{pgfscope}%
\begin{pgfscope}%
\pgfpathrectangle{\pgfqpoint{3.393750in}{0.699667in}}{\pgfqpoint{0.096250in}{1.925000in}}%
\pgfusepath{clip}%
\pgfsetbuttcap%
\pgfsetroundjoin%
\definecolor{currentfill}{rgb}{0.607238,0.179779,0.499492}%
\pgfsetfillcolor{currentfill}%
\pgfsetfillopacity{0.800000}%
\pgfsetlinewidth{0.000000pt}%
\definecolor{currentstroke}{rgb}{0.000000,0.000000,0.000000}%
\pgfsetstrokecolor{currentstroke}%
\pgfsetdash{}{0pt}%
\pgfpathmoveto{\pgfqpoint{3.393750in}{1.534335in}}%
\pgfpathlineto{\pgfqpoint{3.490000in}{1.534335in}}%
\pgfpathlineto{\pgfqpoint{3.490000in}{1.541855in}}%
\pgfpathlineto{\pgfqpoint{3.393750in}{1.541855in}}%
\pgfpathlineto{\pgfqpoint{3.393750in}{1.534335in}}%
\pgfusepath{fill}%
\end{pgfscope}%
\begin{pgfscope}%
\pgfpathrectangle{\pgfqpoint{3.393750in}{0.699667in}}{\pgfqpoint{0.096250in}{1.925000in}}%
\pgfusepath{clip}%
\pgfsetbuttcap%
\pgfsetroundjoin%
\definecolor{currentfill}{rgb}{0.613617,0.181811,0.498536}%
\pgfsetfillcolor{currentfill}%
\pgfsetfillopacity{0.800000}%
\pgfsetlinewidth{0.000000pt}%
\definecolor{currentstroke}{rgb}{0.000000,0.000000,0.000000}%
\pgfsetstrokecolor{currentstroke}%
\pgfsetdash{}{0pt}%
\pgfpathmoveto{\pgfqpoint{3.393750in}{1.541855in}}%
\pgfpathlineto{\pgfqpoint{3.490000in}{1.541855in}}%
\pgfpathlineto{\pgfqpoint{3.490000in}{1.549374in}}%
\pgfpathlineto{\pgfqpoint{3.393750in}{1.549374in}}%
\pgfpathlineto{\pgfqpoint{3.393750in}{1.541855in}}%
\pgfusepath{fill}%
\end{pgfscope}%
\begin{pgfscope}%
\pgfpathrectangle{\pgfqpoint{3.393750in}{0.699667in}}{\pgfqpoint{0.096250in}{1.925000in}}%
\pgfusepath{clip}%
\pgfsetbuttcap%
\pgfsetroundjoin%
\definecolor{currentfill}{rgb}{0.620005,0.183840,0.497524}%
\pgfsetfillcolor{currentfill}%
\pgfsetfillopacity{0.800000}%
\pgfsetlinewidth{0.000000pt}%
\definecolor{currentstroke}{rgb}{0.000000,0.000000,0.000000}%
\pgfsetstrokecolor{currentstroke}%
\pgfsetdash{}{0pt}%
\pgfpathmoveto{\pgfqpoint{3.393750in}{1.549374in}}%
\pgfpathlineto{\pgfqpoint{3.490000in}{1.549374in}}%
\pgfpathlineto{\pgfqpoint{3.490000in}{1.556894in}}%
\pgfpathlineto{\pgfqpoint{3.393750in}{1.556894in}}%
\pgfpathlineto{\pgfqpoint{3.393750in}{1.549374in}}%
\pgfusepath{fill}%
\end{pgfscope}%
\begin{pgfscope}%
\pgfpathrectangle{\pgfqpoint{3.393750in}{0.699667in}}{\pgfqpoint{0.096250in}{1.925000in}}%
\pgfusepath{clip}%
\pgfsetbuttcap%
\pgfsetroundjoin%
\definecolor{currentfill}{rgb}{0.626401,0.185867,0.496456}%
\pgfsetfillcolor{currentfill}%
\pgfsetfillopacity{0.800000}%
\pgfsetlinewidth{0.000000pt}%
\definecolor{currentstroke}{rgb}{0.000000,0.000000,0.000000}%
\pgfsetstrokecolor{currentstroke}%
\pgfsetdash{}{0pt}%
\pgfpathmoveto{\pgfqpoint{3.393750in}{1.556894in}}%
\pgfpathlineto{\pgfqpoint{3.490000in}{1.556894in}}%
\pgfpathlineto{\pgfqpoint{3.490000in}{1.564413in}}%
\pgfpathlineto{\pgfqpoint{3.393750in}{1.564413in}}%
\pgfpathlineto{\pgfqpoint{3.393750in}{1.556894in}}%
\pgfusepath{fill}%
\end{pgfscope}%
\begin{pgfscope}%
\pgfpathrectangle{\pgfqpoint{3.393750in}{0.699667in}}{\pgfqpoint{0.096250in}{1.925000in}}%
\pgfusepath{clip}%
\pgfsetbuttcap%
\pgfsetroundjoin%
\definecolor{currentfill}{rgb}{0.632805,0.187893,0.495332}%
\pgfsetfillcolor{currentfill}%
\pgfsetfillopacity{0.800000}%
\pgfsetlinewidth{0.000000pt}%
\definecolor{currentstroke}{rgb}{0.000000,0.000000,0.000000}%
\pgfsetstrokecolor{currentstroke}%
\pgfsetdash{}{0pt}%
\pgfpathmoveto{\pgfqpoint{3.393750in}{1.564413in}}%
\pgfpathlineto{\pgfqpoint{3.490000in}{1.564413in}}%
\pgfpathlineto{\pgfqpoint{3.490000in}{1.571933in}}%
\pgfpathlineto{\pgfqpoint{3.393750in}{1.571933in}}%
\pgfpathlineto{\pgfqpoint{3.393750in}{1.564413in}}%
\pgfusepath{fill}%
\end{pgfscope}%
\begin{pgfscope}%
\pgfpathrectangle{\pgfqpoint{3.393750in}{0.699667in}}{\pgfqpoint{0.096250in}{1.925000in}}%
\pgfusepath{clip}%
\pgfsetbuttcap%
\pgfsetroundjoin%
\definecolor{currentfill}{rgb}{0.639216,0.189921,0.494150}%
\pgfsetfillcolor{currentfill}%
\pgfsetfillopacity{0.800000}%
\pgfsetlinewidth{0.000000pt}%
\definecolor{currentstroke}{rgb}{0.000000,0.000000,0.000000}%
\pgfsetstrokecolor{currentstroke}%
\pgfsetdash{}{0pt}%
\pgfpathmoveto{\pgfqpoint{3.393750in}{1.571933in}}%
\pgfpathlineto{\pgfqpoint{3.490000in}{1.571933in}}%
\pgfpathlineto{\pgfqpoint{3.490000in}{1.579452in}}%
\pgfpathlineto{\pgfqpoint{3.393750in}{1.579452in}}%
\pgfpathlineto{\pgfqpoint{3.393750in}{1.571933in}}%
\pgfusepath{fill}%
\end{pgfscope}%
\begin{pgfscope}%
\pgfpathrectangle{\pgfqpoint{3.393750in}{0.699667in}}{\pgfqpoint{0.096250in}{1.925000in}}%
\pgfusepath{clip}%
\pgfsetbuttcap%
\pgfsetroundjoin%
\definecolor{currentfill}{rgb}{0.645633,0.191952,0.492910}%
\pgfsetfillcolor{currentfill}%
\pgfsetfillopacity{0.800000}%
\pgfsetlinewidth{0.000000pt}%
\definecolor{currentstroke}{rgb}{0.000000,0.000000,0.000000}%
\pgfsetstrokecolor{currentstroke}%
\pgfsetdash{}{0pt}%
\pgfpathmoveto{\pgfqpoint{3.393750in}{1.579452in}}%
\pgfpathlineto{\pgfqpoint{3.490000in}{1.579452in}}%
\pgfpathlineto{\pgfqpoint{3.490000in}{1.586972in}}%
\pgfpathlineto{\pgfqpoint{3.393750in}{1.586972in}}%
\pgfpathlineto{\pgfqpoint{3.393750in}{1.579452in}}%
\pgfusepath{fill}%
\end{pgfscope}%
\begin{pgfscope}%
\pgfpathrectangle{\pgfqpoint{3.393750in}{0.699667in}}{\pgfqpoint{0.096250in}{1.925000in}}%
\pgfusepath{clip}%
\pgfsetbuttcap%
\pgfsetroundjoin%
\definecolor{currentfill}{rgb}{0.652056,0.193986,0.491611}%
\pgfsetfillcolor{currentfill}%
\pgfsetfillopacity{0.800000}%
\pgfsetlinewidth{0.000000pt}%
\definecolor{currentstroke}{rgb}{0.000000,0.000000,0.000000}%
\pgfsetstrokecolor{currentstroke}%
\pgfsetdash{}{0pt}%
\pgfpathmoveto{\pgfqpoint{3.393750in}{1.586972in}}%
\pgfpathlineto{\pgfqpoint{3.490000in}{1.586972in}}%
\pgfpathlineto{\pgfqpoint{3.490000in}{1.594491in}}%
\pgfpathlineto{\pgfqpoint{3.393750in}{1.594491in}}%
\pgfpathlineto{\pgfqpoint{3.393750in}{1.586972in}}%
\pgfusepath{fill}%
\end{pgfscope}%
\begin{pgfscope}%
\pgfpathrectangle{\pgfqpoint{3.393750in}{0.699667in}}{\pgfqpoint{0.096250in}{1.925000in}}%
\pgfusepath{clip}%
\pgfsetbuttcap%
\pgfsetroundjoin%
\definecolor{currentfill}{rgb}{0.658483,0.196027,0.490253}%
\pgfsetfillcolor{currentfill}%
\pgfsetfillopacity{0.800000}%
\pgfsetlinewidth{0.000000pt}%
\definecolor{currentstroke}{rgb}{0.000000,0.000000,0.000000}%
\pgfsetstrokecolor{currentstroke}%
\pgfsetdash{}{0pt}%
\pgfpathmoveto{\pgfqpoint{3.393750in}{1.594491in}}%
\pgfpathlineto{\pgfqpoint{3.490000in}{1.594491in}}%
\pgfpathlineto{\pgfqpoint{3.490000in}{1.602011in}}%
\pgfpathlineto{\pgfqpoint{3.393750in}{1.602011in}}%
\pgfpathlineto{\pgfqpoint{3.393750in}{1.594491in}}%
\pgfusepath{fill}%
\end{pgfscope}%
\begin{pgfscope}%
\pgfpathrectangle{\pgfqpoint{3.393750in}{0.699667in}}{\pgfqpoint{0.096250in}{1.925000in}}%
\pgfusepath{clip}%
\pgfsetbuttcap%
\pgfsetroundjoin%
\definecolor{currentfill}{rgb}{0.664915,0.198075,0.488836}%
\pgfsetfillcolor{currentfill}%
\pgfsetfillopacity{0.800000}%
\pgfsetlinewidth{0.000000pt}%
\definecolor{currentstroke}{rgb}{0.000000,0.000000,0.000000}%
\pgfsetstrokecolor{currentstroke}%
\pgfsetdash{}{0pt}%
\pgfpathmoveto{\pgfqpoint{3.393750in}{1.602011in}}%
\pgfpathlineto{\pgfqpoint{3.490000in}{1.602011in}}%
\pgfpathlineto{\pgfqpoint{3.490000in}{1.609530in}}%
\pgfpathlineto{\pgfqpoint{3.393750in}{1.609530in}}%
\pgfpathlineto{\pgfqpoint{3.393750in}{1.602011in}}%
\pgfusepath{fill}%
\end{pgfscope}%
\begin{pgfscope}%
\pgfpathrectangle{\pgfqpoint{3.393750in}{0.699667in}}{\pgfqpoint{0.096250in}{1.925000in}}%
\pgfusepath{clip}%
\pgfsetbuttcap%
\pgfsetroundjoin%
\definecolor{currentfill}{rgb}{0.671349,0.200133,0.487358}%
\pgfsetfillcolor{currentfill}%
\pgfsetfillopacity{0.800000}%
\pgfsetlinewidth{0.000000pt}%
\definecolor{currentstroke}{rgb}{0.000000,0.000000,0.000000}%
\pgfsetstrokecolor{currentstroke}%
\pgfsetdash{}{0pt}%
\pgfpathmoveto{\pgfqpoint{3.393750in}{1.609530in}}%
\pgfpathlineto{\pgfqpoint{3.490000in}{1.609530in}}%
\pgfpathlineto{\pgfqpoint{3.490000in}{1.617050in}}%
\pgfpathlineto{\pgfqpoint{3.393750in}{1.617050in}}%
\pgfpathlineto{\pgfqpoint{3.393750in}{1.609530in}}%
\pgfusepath{fill}%
\end{pgfscope}%
\begin{pgfscope}%
\pgfpathrectangle{\pgfqpoint{3.393750in}{0.699667in}}{\pgfqpoint{0.096250in}{1.925000in}}%
\pgfusepath{clip}%
\pgfsetbuttcap%
\pgfsetroundjoin%
\definecolor{currentfill}{rgb}{0.677786,0.202203,0.485819}%
\pgfsetfillcolor{currentfill}%
\pgfsetfillopacity{0.800000}%
\pgfsetlinewidth{0.000000pt}%
\definecolor{currentstroke}{rgb}{0.000000,0.000000,0.000000}%
\pgfsetstrokecolor{currentstroke}%
\pgfsetdash{}{0pt}%
\pgfpathmoveto{\pgfqpoint{3.393750in}{1.617050in}}%
\pgfpathlineto{\pgfqpoint{3.490000in}{1.617050in}}%
\pgfpathlineto{\pgfqpoint{3.490000in}{1.624570in}}%
\pgfpathlineto{\pgfqpoint{3.393750in}{1.624570in}}%
\pgfpathlineto{\pgfqpoint{3.393750in}{1.617050in}}%
\pgfusepath{fill}%
\end{pgfscope}%
\begin{pgfscope}%
\pgfpathrectangle{\pgfqpoint{3.393750in}{0.699667in}}{\pgfqpoint{0.096250in}{1.925000in}}%
\pgfusepath{clip}%
\pgfsetbuttcap%
\pgfsetroundjoin%
\definecolor{currentfill}{rgb}{0.684224,0.204286,0.484219}%
\pgfsetfillcolor{currentfill}%
\pgfsetfillopacity{0.800000}%
\pgfsetlinewidth{0.000000pt}%
\definecolor{currentstroke}{rgb}{0.000000,0.000000,0.000000}%
\pgfsetstrokecolor{currentstroke}%
\pgfsetdash{}{0pt}%
\pgfpathmoveto{\pgfqpoint{3.393750in}{1.624570in}}%
\pgfpathlineto{\pgfqpoint{3.490000in}{1.624570in}}%
\pgfpathlineto{\pgfqpoint{3.490000in}{1.632089in}}%
\pgfpathlineto{\pgfqpoint{3.393750in}{1.632089in}}%
\pgfpathlineto{\pgfqpoint{3.393750in}{1.624570in}}%
\pgfusepath{fill}%
\end{pgfscope}%
\begin{pgfscope}%
\pgfpathrectangle{\pgfqpoint{3.393750in}{0.699667in}}{\pgfqpoint{0.096250in}{1.925000in}}%
\pgfusepath{clip}%
\pgfsetbuttcap%
\pgfsetroundjoin%
\definecolor{currentfill}{rgb}{0.690661,0.206384,0.482558}%
\pgfsetfillcolor{currentfill}%
\pgfsetfillopacity{0.800000}%
\pgfsetlinewidth{0.000000pt}%
\definecolor{currentstroke}{rgb}{0.000000,0.000000,0.000000}%
\pgfsetstrokecolor{currentstroke}%
\pgfsetdash{}{0pt}%
\pgfpathmoveto{\pgfqpoint{3.393750in}{1.632089in}}%
\pgfpathlineto{\pgfqpoint{3.490000in}{1.632089in}}%
\pgfpathlineto{\pgfqpoint{3.490000in}{1.639609in}}%
\pgfpathlineto{\pgfqpoint{3.393750in}{1.639609in}}%
\pgfpathlineto{\pgfqpoint{3.393750in}{1.632089in}}%
\pgfusepath{fill}%
\end{pgfscope}%
\begin{pgfscope}%
\pgfpathrectangle{\pgfqpoint{3.393750in}{0.699667in}}{\pgfqpoint{0.096250in}{1.925000in}}%
\pgfusepath{clip}%
\pgfsetbuttcap%
\pgfsetroundjoin%
\definecolor{currentfill}{rgb}{0.697098,0.208501,0.480835}%
\pgfsetfillcolor{currentfill}%
\pgfsetfillopacity{0.800000}%
\pgfsetlinewidth{0.000000pt}%
\definecolor{currentstroke}{rgb}{0.000000,0.000000,0.000000}%
\pgfsetstrokecolor{currentstroke}%
\pgfsetdash{}{0pt}%
\pgfpathmoveto{\pgfqpoint{3.393750in}{1.639609in}}%
\pgfpathlineto{\pgfqpoint{3.490000in}{1.639609in}}%
\pgfpathlineto{\pgfqpoint{3.490000in}{1.647128in}}%
\pgfpathlineto{\pgfqpoint{3.393750in}{1.647128in}}%
\pgfpathlineto{\pgfqpoint{3.393750in}{1.639609in}}%
\pgfusepath{fill}%
\end{pgfscope}%
\begin{pgfscope}%
\pgfpathrectangle{\pgfqpoint{3.393750in}{0.699667in}}{\pgfqpoint{0.096250in}{1.925000in}}%
\pgfusepath{clip}%
\pgfsetbuttcap%
\pgfsetroundjoin%
\definecolor{currentfill}{rgb}{0.703532,0.210638,0.479049}%
\pgfsetfillcolor{currentfill}%
\pgfsetfillopacity{0.800000}%
\pgfsetlinewidth{0.000000pt}%
\definecolor{currentstroke}{rgb}{0.000000,0.000000,0.000000}%
\pgfsetstrokecolor{currentstroke}%
\pgfsetdash{}{0pt}%
\pgfpathmoveto{\pgfqpoint{3.393750in}{1.647128in}}%
\pgfpathlineto{\pgfqpoint{3.490000in}{1.647128in}}%
\pgfpathlineto{\pgfqpoint{3.490000in}{1.654648in}}%
\pgfpathlineto{\pgfqpoint{3.393750in}{1.654648in}}%
\pgfpathlineto{\pgfqpoint{3.393750in}{1.647128in}}%
\pgfusepath{fill}%
\end{pgfscope}%
\begin{pgfscope}%
\pgfpathrectangle{\pgfqpoint{3.393750in}{0.699667in}}{\pgfqpoint{0.096250in}{1.925000in}}%
\pgfusepath{clip}%
\pgfsetbuttcap%
\pgfsetroundjoin%
\definecolor{currentfill}{rgb}{0.709962,0.212797,0.477201}%
\pgfsetfillcolor{currentfill}%
\pgfsetfillopacity{0.800000}%
\pgfsetlinewidth{0.000000pt}%
\definecolor{currentstroke}{rgb}{0.000000,0.000000,0.000000}%
\pgfsetstrokecolor{currentstroke}%
\pgfsetdash{}{0pt}%
\pgfpathmoveto{\pgfqpoint{3.393750in}{1.654648in}}%
\pgfpathlineto{\pgfqpoint{3.490000in}{1.654648in}}%
\pgfpathlineto{\pgfqpoint{3.490000in}{1.662167in}}%
\pgfpathlineto{\pgfqpoint{3.393750in}{1.662167in}}%
\pgfpathlineto{\pgfqpoint{3.393750in}{1.654648in}}%
\pgfusepath{fill}%
\end{pgfscope}%
\begin{pgfscope}%
\pgfpathrectangle{\pgfqpoint{3.393750in}{0.699667in}}{\pgfqpoint{0.096250in}{1.925000in}}%
\pgfusepath{clip}%
\pgfsetbuttcap%
\pgfsetroundjoin%
\definecolor{currentfill}{rgb}{0.716387,0.214982,0.475290}%
\pgfsetfillcolor{currentfill}%
\pgfsetfillopacity{0.800000}%
\pgfsetlinewidth{0.000000pt}%
\definecolor{currentstroke}{rgb}{0.000000,0.000000,0.000000}%
\pgfsetstrokecolor{currentstroke}%
\pgfsetdash{}{0pt}%
\pgfpathmoveto{\pgfqpoint{3.393750in}{1.662167in}}%
\pgfpathlineto{\pgfqpoint{3.490000in}{1.662167in}}%
\pgfpathlineto{\pgfqpoint{3.490000in}{1.669687in}}%
\pgfpathlineto{\pgfqpoint{3.393750in}{1.669687in}}%
\pgfpathlineto{\pgfqpoint{3.393750in}{1.662167in}}%
\pgfusepath{fill}%
\end{pgfscope}%
\begin{pgfscope}%
\pgfpathrectangle{\pgfqpoint{3.393750in}{0.699667in}}{\pgfqpoint{0.096250in}{1.925000in}}%
\pgfusepath{clip}%
\pgfsetbuttcap%
\pgfsetroundjoin%
\definecolor{currentfill}{rgb}{0.722805,0.217194,0.473316}%
\pgfsetfillcolor{currentfill}%
\pgfsetfillopacity{0.800000}%
\pgfsetlinewidth{0.000000pt}%
\definecolor{currentstroke}{rgb}{0.000000,0.000000,0.000000}%
\pgfsetstrokecolor{currentstroke}%
\pgfsetdash{}{0pt}%
\pgfpathmoveto{\pgfqpoint{3.393750in}{1.669687in}}%
\pgfpathlineto{\pgfqpoint{3.490000in}{1.669687in}}%
\pgfpathlineto{\pgfqpoint{3.490000in}{1.677206in}}%
\pgfpathlineto{\pgfqpoint{3.393750in}{1.677206in}}%
\pgfpathlineto{\pgfqpoint{3.393750in}{1.669687in}}%
\pgfusepath{fill}%
\end{pgfscope}%
\begin{pgfscope}%
\pgfpathrectangle{\pgfqpoint{3.393750in}{0.699667in}}{\pgfqpoint{0.096250in}{1.925000in}}%
\pgfusepath{clip}%
\pgfsetbuttcap%
\pgfsetroundjoin%
\definecolor{currentfill}{rgb}{0.729216,0.219437,0.471279}%
\pgfsetfillcolor{currentfill}%
\pgfsetfillopacity{0.800000}%
\pgfsetlinewidth{0.000000pt}%
\definecolor{currentstroke}{rgb}{0.000000,0.000000,0.000000}%
\pgfsetstrokecolor{currentstroke}%
\pgfsetdash{}{0pt}%
\pgfpathmoveto{\pgfqpoint{3.393750in}{1.677206in}}%
\pgfpathlineto{\pgfqpoint{3.490000in}{1.677206in}}%
\pgfpathlineto{\pgfqpoint{3.490000in}{1.684726in}}%
\pgfpathlineto{\pgfqpoint{3.393750in}{1.684726in}}%
\pgfpathlineto{\pgfqpoint{3.393750in}{1.677206in}}%
\pgfusepath{fill}%
\end{pgfscope}%
\begin{pgfscope}%
\pgfpathrectangle{\pgfqpoint{3.393750in}{0.699667in}}{\pgfqpoint{0.096250in}{1.925000in}}%
\pgfusepath{clip}%
\pgfsetbuttcap%
\pgfsetroundjoin%
\definecolor{currentfill}{rgb}{0.735616,0.221713,0.469180}%
\pgfsetfillcolor{currentfill}%
\pgfsetfillopacity{0.800000}%
\pgfsetlinewidth{0.000000pt}%
\definecolor{currentstroke}{rgb}{0.000000,0.000000,0.000000}%
\pgfsetstrokecolor{currentstroke}%
\pgfsetdash{}{0pt}%
\pgfpathmoveto{\pgfqpoint{3.393750in}{1.684726in}}%
\pgfpathlineto{\pgfqpoint{3.490000in}{1.684726in}}%
\pgfpathlineto{\pgfqpoint{3.490000in}{1.692245in}}%
\pgfpathlineto{\pgfqpoint{3.393750in}{1.692245in}}%
\pgfpathlineto{\pgfqpoint{3.393750in}{1.684726in}}%
\pgfusepath{fill}%
\end{pgfscope}%
\begin{pgfscope}%
\pgfpathrectangle{\pgfqpoint{3.393750in}{0.699667in}}{\pgfqpoint{0.096250in}{1.925000in}}%
\pgfusepath{clip}%
\pgfsetbuttcap%
\pgfsetroundjoin%
\definecolor{currentfill}{rgb}{0.742004,0.224025,0.467018}%
\pgfsetfillcolor{currentfill}%
\pgfsetfillopacity{0.800000}%
\pgfsetlinewidth{0.000000pt}%
\definecolor{currentstroke}{rgb}{0.000000,0.000000,0.000000}%
\pgfsetstrokecolor{currentstroke}%
\pgfsetdash{}{0pt}%
\pgfpathmoveto{\pgfqpoint{3.393750in}{1.692245in}}%
\pgfpathlineto{\pgfqpoint{3.490000in}{1.692245in}}%
\pgfpathlineto{\pgfqpoint{3.490000in}{1.699765in}}%
\pgfpathlineto{\pgfqpoint{3.393750in}{1.699765in}}%
\pgfpathlineto{\pgfqpoint{3.393750in}{1.692245in}}%
\pgfusepath{fill}%
\end{pgfscope}%
\begin{pgfscope}%
\pgfpathrectangle{\pgfqpoint{3.393750in}{0.699667in}}{\pgfqpoint{0.096250in}{1.925000in}}%
\pgfusepath{clip}%
\pgfsetbuttcap%
\pgfsetroundjoin%
\definecolor{currentfill}{rgb}{0.748378,0.226377,0.464794}%
\pgfsetfillcolor{currentfill}%
\pgfsetfillopacity{0.800000}%
\pgfsetlinewidth{0.000000pt}%
\definecolor{currentstroke}{rgb}{0.000000,0.000000,0.000000}%
\pgfsetstrokecolor{currentstroke}%
\pgfsetdash{}{0pt}%
\pgfpathmoveto{\pgfqpoint{3.393750in}{1.699765in}}%
\pgfpathlineto{\pgfqpoint{3.490000in}{1.699765in}}%
\pgfpathlineto{\pgfqpoint{3.490000in}{1.707284in}}%
\pgfpathlineto{\pgfqpoint{3.393750in}{1.707284in}}%
\pgfpathlineto{\pgfqpoint{3.393750in}{1.699765in}}%
\pgfusepath{fill}%
\end{pgfscope}%
\begin{pgfscope}%
\pgfpathrectangle{\pgfqpoint{3.393750in}{0.699667in}}{\pgfqpoint{0.096250in}{1.925000in}}%
\pgfusepath{clip}%
\pgfsetbuttcap%
\pgfsetroundjoin%
\definecolor{currentfill}{rgb}{0.754737,0.228772,0.462509}%
\pgfsetfillcolor{currentfill}%
\pgfsetfillopacity{0.800000}%
\pgfsetlinewidth{0.000000pt}%
\definecolor{currentstroke}{rgb}{0.000000,0.000000,0.000000}%
\pgfsetstrokecolor{currentstroke}%
\pgfsetdash{}{0pt}%
\pgfpathmoveto{\pgfqpoint{3.393750in}{1.707284in}}%
\pgfpathlineto{\pgfqpoint{3.490000in}{1.707284in}}%
\pgfpathlineto{\pgfqpoint{3.490000in}{1.714804in}}%
\pgfpathlineto{\pgfqpoint{3.393750in}{1.714804in}}%
\pgfpathlineto{\pgfqpoint{3.393750in}{1.707284in}}%
\pgfusepath{fill}%
\end{pgfscope}%
\begin{pgfscope}%
\pgfpathrectangle{\pgfqpoint{3.393750in}{0.699667in}}{\pgfqpoint{0.096250in}{1.925000in}}%
\pgfusepath{clip}%
\pgfsetbuttcap%
\pgfsetroundjoin%
\definecolor{currentfill}{rgb}{0.761077,0.231214,0.460162}%
\pgfsetfillcolor{currentfill}%
\pgfsetfillopacity{0.800000}%
\pgfsetlinewidth{0.000000pt}%
\definecolor{currentstroke}{rgb}{0.000000,0.000000,0.000000}%
\pgfsetstrokecolor{currentstroke}%
\pgfsetdash{}{0pt}%
\pgfpathmoveto{\pgfqpoint{3.393750in}{1.714804in}}%
\pgfpathlineto{\pgfqpoint{3.490000in}{1.714804in}}%
\pgfpathlineto{\pgfqpoint{3.490000in}{1.722323in}}%
\pgfpathlineto{\pgfqpoint{3.393750in}{1.722323in}}%
\pgfpathlineto{\pgfqpoint{3.393750in}{1.714804in}}%
\pgfusepath{fill}%
\end{pgfscope}%
\begin{pgfscope}%
\pgfpathrectangle{\pgfqpoint{3.393750in}{0.699667in}}{\pgfqpoint{0.096250in}{1.925000in}}%
\pgfusepath{clip}%
\pgfsetbuttcap%
\pgfsetroundjoin%
\definecolor{currentfill}{rgb}{0.767398,0.233705,0.457755}%
\pgfsetfillcolor{currentfill}%
\pgfsetfillopacity{0.800000}%
\pgfsetlinewidth{0.000000pt}%
\definecolor{currentstroke}{rgb}{0.000000,0.000000,0.000000}%
\pgfsetstrokecolor{currentstroke}%
\pgfsetdash{}{0pt}%
\pgfpathmoveto{\pgfqpoint{3.393750in}{1.722323in}}%
\pgfpathlineto{\pgfqpoint{3.490000in}{1.722323in}}%
\pgfpathlineto{\pgfqpoint{3.490000in}{1.729843in}}%
\pgfpathlineto{\pgfqpoint{3.393750in}{1.729843in}}%
\pgfpathlineto{\pgfqpoint{3.393750in}{1.722323in}}%
\pgfusepath{fill}%
\end{pgfscope}%
\begin{pgfscope}%
\pgfpathrectangle{\pgfqpoint{3.393750in}{0.699667in}}{\pgfqpoint{0.096250in}{1.925000in}}%
\pgfusepath{clip}%
\pgfsetbuttcap%
\pgfsetroundjoin%
\definecolor{currentfill}{rgb}{0.773695,0.236249,0.455289}%
\pgfsetfillcolor{currentfill}%
\pgfsetfillopacity{0.800000}%
\pgfsetlinewidth{0.000000pt}%
\definecolor{currentstroke}{rgb}{0.000000,0.000000,0.000000}%
\pgfsetstrokecolor{currentstroke}%
\pgfsetdash{}{0pt}%
\pgfpathmoveto{\pgfqpoint{3.393750in}{1.729843in}}%
\pgfpathlineto{\pgfqpoint{3.490000in}{1.729843in}}%
\pgfpathlineto{\pgfqpoint{3.490000in}{1.737362in}}%
\pgfpathlineto{\pgfqpoint{3.393750in}{1.737362in}}%
\pgfpathlineto{\pgfqpoint{3.393750in}{1.729843in}}%
\pgfusepath{fill}%
\end{pgfscope}%
\begin{pgfscope}%
\pgfpathrectangle{\pgfqpoint{3.393750in}{0.699667in}}{\pgfqpoint{0.096250in}{1.925000in}}%
\pgfusepath{clip}%
\pgfsetbuttcap%
\pgfsetroundjoin%
\definecolor{currentfill}{rgb}{0.779968,0.238851,0.452765}%
\pgfsetfillcolor{currentfill}%
\pgfsetfillopacity{0.800000}%
\pgfsetlinewidth{0.000000pt}%
\definecolor{currentstroke}{rgb}{0.000000,0.000000,0.000000}%
\pgfsetstrokecolor{currentstroke}%
\pgfsetdash{}{0pt}%
\pgfpathmoveto{\pgfqpoint{3.393750in}{1.737362in}}%
\pgfpathlineto{\pgfqpoint{3.490000in}{1.737362in}}%
\pgfpathlineto{\pgfqpoint{3.490000in}{1.744882in}}%
\pgfpathlineto{\pgfqpoint{3.393750in}{1.744882in}}%
\pgfpathlineto{\pgfqpoint{3.393750in}{1.737362in}}%
\pgfusepath{fill}%
\end{pgfscope}%
\begin{pgfscope}%
\pgfpathrectangle{\pgfqpoint{3.393750in}{0.699667in}}{\pgfqpoint{0.096250in}{1.925000in}}%
\pgfusepath{clip}%
\pgfsetbuttcap%
\pgfsetroundjoin%
\definecolor{currentfill}{rgb}{0.786212,0.241514,0.450184}%
\pgfsetfillcolor{currentfill}%
\pgfsetfillopacity{0.800000}%
\pgfsetlinewidth{0.000000pt}%
\definecolor{currentstroke}{rgb}{0.000000,0.000000,0.000000}%
\pgfsetstrokecolor{currentstroke}%
\pgfsetdash{}{0pt}%
\pgfpathmoveto{\pgfqpoint{3.393750in}{1.744882in}}%
\pgfpathlineto{\pgfqpoint{3.490000in}{1.744882in}}%
\pgfpathlineto{\pgfqpoint{3.490000in}{1.752402in}}%
\pgfpathlineto{\pgfqpoint{3.393750in}{1.752402in}}%
\pgfpathlineto{\pgfqpoint{3.393750in}{1.744882in}}%
\pgfusepath{fill}%
\end{pgfscope}%
\begin{pgfscope}%
\pgfpathrectangle{\pgfqpoint{3.393750in}{0.699667in}}{\pgfqpoint{0.096250in}{1.925000in}}%
\pgfusepath{clip}%
\pgfsetbuttcap%
\pgfsetroundjoin%
\definecolor{currentfill}{rgb}{0.792427,0.244242,0.447543}%
\pgfsetfillcolor{currentfill}%
\pgfsetfillopacity{0.800000}%
\pgfsetlinewidth{0.000000pt}%
\definecolor{currentstroke}{rgb}{0.000000,0.000000,0.000000}%
\pgfsetstrokecolor{currentstroke}%
\pgfsetdash{}{0pt}%
\pgfpathmoveto{\pgfqpoint{3.393750in}{1.752402in}}%
\pgfpathlineto{\pgfqpoint{3.490000in}{1.752402in}}%
\pgfpathlineto{\pgfqpoint{3.490000in}{1.759921in}}%
\pgfpathlineto{\pgfqpoint{3.393750in}{1.759921in}}%
\pgfpathlineto{\pgfqpoint{3.393750in}{1.752402in}}%
\pgfusepath{fill}%
\end{pgfscope}%
\begin{pgfscope}%
\pgfpathrectangle{\pgfqpoint{3.393750in}{0.699667in}}{\pgfqpoint{0.096250in}{1.925000in}}%
\pgfusepath{clip}%
\pgfsetbuttcap%
\pgfsetroundjoin%
\definecolor{currentfill}{rgb}{0.798608,0.247040,0.444848}%
\pgfsetfillcolor{currentfill}%
\pgfsetfillopacity{0.800000}%
\pgfsetlinewidth{0.000000pt}%
\definecolor{currentstroke}{rgb}{0.000000,0.000000,0.000000}%
\pgfsetstrokecolor{currentstroke}%
\pgfsetdash{}{0pt}%
\pgfpathmoveto{\pgfqpoint{3.393750in}{1.759921in}}%
\pgfpathlineto{\pgfqpoint{3.490000in}{1.759921in}}%
\pgfpathlineto{\pgfqpoint{3.490000in}{1.767441in}}%
\pgfpathlineto{\pgfqpoint{3.393750in}{1.767441in}}%
\pgfpathlineto{\pgfqpoint{3.393750in}{1.759921in}}%
\pgfusepath{fill}%
\end{pgfscope}%
\begin{pgfscope}%
\pgfpathrectangle{\pgfqpoint{3.393750in}{0.699667in}}{\pgfqpoint{0.096250in}{1.925000in}}%
\pgfusepath{clip}%
\pgfsetbuttcap%
\pgfsetroundjoin%
\definecolor{currentfill}{rgb}{0.804752,0.249911,0.442102}%
\pgfsetfillcolor{currentfill}%
\pgfsetfillopacity{0.800000}%
\pgfsetlinewidth{0.000000pt}%
\definecolor{currentstroke}{rgb}{0.000000,0.000000,0.000000}%
\pgfsetstrokecolor{currentstroke}%
\pgfsetdash{}{0pt}%
\pgfpathmoveto{\pgfqpoint{3.393750in}{1.767441in}}%
\pgfpathlineto{\pgfqpoint{3.490000in}{1.767441in}}%
\pgfpathlineto{\pgfqpoint{3.490000in}{1.774960in}}%
\pgfpathlineto{\pgfqpoint{3.393750in}{1.774960in}}%
\pgfpathlineto{\pgfqpoint{3.393750in}{1.767441in}}%
\pgfusepath{fill}%
\end{pgfscope}%
\begin{pgfscope}%
\pgfpathrectangle{\pgfqpoint{3.393750in}{0.699667in}}{\pgfqpoint{0.096250in}{1.925000in}}%
\pgfusepath{clip}%
\pgfsetbuttcap%
\pgfsetroundjoin%
\definecolor{currentfill}{rgb}{0.810855,0.252861,0.439305}%
\pgfsetfillcolor{currentfill}%
\pgfsetfillopacity{0.800000}%
\pgfsetlinewidth{0.000000pt}%
\definecolor{currentstroke}{rgb}{0.000000,0.000000,0.000000}%
\pgfsetstrokecolor{currentstroke}%
\pgfsetdash{}{0pt}%
\pgfpathmoveto{\pgfqpoint{3.393750in}{1.774960in}}%
\pgfpathlineto{\pgfqpoint{3.490000in}{1.774960in}}%
\pgfpathlineto{\pgfqpoint{3.490000in}{1.782480in}}%
\pgfpathlineto{\pgfqpoint{3.393750in}{1.782480in}}%
\pgfpathlineto{\pgfqpoint{3.393750in}{1.774960in}}%
\pgfusepath{fill}%
\end{pgfscope}%
\begin{pgfscope}%
\pgfpathrectangle{\pgfqpoint{3.393750in}{0.699667in}}{\pgfqpoint{0.096250in}{1.925000in}}%
\pgfusepath{clip}%
\pgfsetbuttcap%
\pgfsetroundjoin%
\definecolor{currentfill}{rgb}{0.816914,0.255895,0.436461}%
\pgfsetfillcolor{currentfill}%
\pgfsetfillopacity{0.800000}%
\pgfsetlinewidth{0.000000pt}%
\definecolor{currentstroke}{rgb}{0.000000,0.000000,0.000000}%
\pgfsetstrokecolor{currentstroke}%
\pgfsetdash{}{0pt}%
\pgfpathmoveto{\pgfqpoint{3.393750in}{1.782480in}}%
\pgfpathlineto{\pgfqpoint{3.490000in}{1.782480in}}%
\pgfpathlineto{\pgfqpoint{3.490000in}{1.789999in}}%
\pgfpathlineto{\pgfqpoint{3.393750in}{1.789999in}}%
\pgfpathlineto{\pgfqpoint{3.393750in}{1.782480in}}%
\pgfusepath{fill}%
\end{pgfscope}%
\begin{pgfscope}%
\pgfpathrectangle{\pgfqpoint{3.393750in}{0.699667in}}{\pgfqpoint{0.096250in}{1.925000in}}%
\pgfusepath{clip}%
\pgfsetbuttcap%
\pgfsetroundjoin%
\definecolor{currentfill}{rgb}{0.822926,0.259016,0.433573}%
\pgfsetfillcolor{currentfill}%
\pgfsetfillopacity{0.800000}%
\pgfsetlinewidth{0.000000pt}%
\definecolor{currentstroke}{rgb}{0.000000,0.000000,0.000000}%
\pgfsetstrokecolor{currentstroke}%
\pgfsetdash{}{0pt}%
\pgfpathmoveto{\pgfqpoint{3.393750in}{1.789999in}}%
\pgfpathlineto{\pgfqpoint{3.490000in}{1.789999in}}%
\pgfpathlineto{\pgfqpoint{3.490000in}{1.797519in}}%
\pgfpathlineto{\pgfqpoint{3.393750in}{1.797519in}}%
\pgfpathlineto{\pgfqpoint{3.393750in}{1.789999in}}%
\pgfusepath{fill}%
\end{pgfscope}%
\begin{pgfscope}%
\pgfpathrectangle{\pgfqpoint{3.393750in}{0.699667in}}{\pgfqpoint{0.096250in}{1.925000in}}%
\pgfusepath{clip}%
\pgfsetbuttcap%
\pgfsetroundjoin%
\definecolor{currentfill}{rgb}{0.828886,0.262229,0.430644}%
\pgfsetfillcolor{currentfill}%
\pgfsetfillopacity{0.800000}%
\pgfsetlinewidth{0.000000pt}%
\definecolor{currentstroke}{rgb}{0.000000,0.000000,0.000000}%
\pgfsetstrokecolor{currentstroke}%
\pgfsetdash{}{0pt}%
\pgfpathmoveto{\pgfqpoint{3.393750in}{1.797519in}}%
\pgfpathlineto{\pgfqpoint{3.490000in}{1.797519in}}%
\pgfpathlineto{\pgfqpoint{3.490000in}{1.805038in}}%
\pgfpathlineto{\pgfqpoint{3.393750in}{1.805038in}}%
\pgfpathlineto{\pgfqpoint{3.393750in}{1.797519in}}%
\pgfusepath{fill}%
\end{pgfscope}%
\begin{pgfscope}%
\pgfpathrectangle{\pgfqpoint{3.393750in}{0.699667in}}{\pgfqpoint{0.096250in}{1.925000in}}%
\pgfusepath{clip}%
\pgfsetbuttcap%
\pgfsetroundjoin%
\definecolor{currentfill}{rgb}{0.834791,0.265540,0.427671}%
\pgfsetfillcolor{currentfill}%
\pgfsetfillopacity{0.800000}%
\pgfsetlinewidth{0.000000pt}%
\definecolor{currentstroke}{rgb}{0.000000,0.000000,0.000000}%
\pgfsetstrokecolor{currentstroke}%
\pgfsetdash{}{0pt}%
\pgfpathmoveto{\pgfqpoint{3.393750in}{1.805038in}}%
\pgfpathlineto{\pgfqpoint{3.490000in}{1.805038in}}%
\pgfpathlineto{\pgfqpoint{3.490000in}{1.812558in}}%
\pgfpathlineto{\pgfqpoint{3.393750in}{1.812558in}}%
\pgfpathlineto{\pgfqpoint{3.393750in}{1.805038in}}%
\pgfusepath{fill}%
\end{pgfscope}%
\begin{pgfscope}%
\pgfpathrectangle{\pgfqpoint{3.393750in}{0.699667in}}{\pgfqpoint{0.096250in}{1.925000in}}%
\pgfusepath{clip}%
\pgfsetbuttcap%
\pgfsetroundjoin%
\definecolor{currentfill}{rgb}{0.840636,0.268953,0.424666}%
\pgfsetfillcolor{currentfill}%
\pgfsetfillopacity{0.800000}%
\pgfsetlinewidth{0.000000pt}%
\definecolor{currentstroke}{rgb}{0.000000,0.000000,0.000000}%
\pgfsetstrokecolor{currentstroke}%
\pgfsetdash{}{0pt}%
\pgfpathmoveto{\pgfqpoint{3.393750in}{1.812558in}}%
\pgfpathlineto{\pgfqpoint{3.490000in}{1.812558in}}%
\pgfpathlineto{\pgfqpoint{3.490000in}{1.820077in}}%
\pgfpathlineto{\pgfqpoint{3.393750in}{1.820077in}}%
\pgfpathlineto{\pgfqpoint{3.393750in}{1.812558in}}%
\pgfusepath{fill}%
\end{pgfscope}%
\begin{pgfscope}%
\pgfpathrectangle{\pgfqpoint{3.393750in}{0.699667in}}{\pgfqpoint{0.096250in}{1.925000in}}%
\pgfusepath{clip}%
\pgfsetbuttcap%
\pgfsetroundjoin%
\definecolor{currentfill}{rgb}{0.846416,0.272473,0.421631}%
\pgfsetfillcolor{currentfill}%
\pgfsetfillopacity{0.800000}%
\pgfsetlinewidth{0.000000pt}%
\definecolor{currentstroke}{rgb}{0.000000,0.000000,0.000000}%
\pgfsetstrokecolor{currentstroke}%
\pgfsetdash{}{0pt}%
\pgfpathmoveto{\pgfqpoint{3.393750in}{1.820077in}}%
\pgfpathlineto{\pgfqpoint{3.490000in}{1.820077in}}%
\pgfpathlineto{\pgfqpoint{3.490000in}{1.827597in}}%
\pgfpathlineto{\pgfqpoint{3.393750in}{1.827597in}}%
\pgfpathlineto{\pgfqpoint{3.393750in}{1.820077in}}%
\pgfusepath{fill}%
\end{pgfscope}%
\begin{pgfscope}%
\pgfpathrectangle{\pgfqpoint{3.393750in}{0.699667in}}{\pgfqpoint{0.096250in}{1.925000in}}%
\pgfusepath{clip}%
\pgfsetbuttcap%
\pgfsetroundjoin%
\definecolor{currentfill}{rgb}{0.852126,0.276106,0.418573}%
\pgfsetfillcolor{currentfill}%
\pgfsetfillopacity{0.800000}%
\pgfsetlinewidth{0.000000pt}%
\definecolor{currentstroke}{rgb}{0.000000,0.000000,0.000000}%
\pgfsetstrokecolor{currentstroke}%
\pgfsetdash{}{0pt}%
\pgfpathmoveto{\pgfqpoint{3.393750in}{1.827597in}}%
\pgfpathlineto{\pgfqpoint{3.490000in}{1.827597in}}%
\pgfpathlineto{\pgfqpoint{3.490000in}{1.835116in}}%
\pgfpathlineto{\pgfqpoint{3.393750in}{1.835116in}}%
\pgfpathlineto{\pgfqpoint{3.393750in}{1.827597in}}%
\pgfusepath{fill}%
\end{pgfscope}%
\begin{pgfscope}%
\pgfpathrectangle{\pgfqpoint{3.393750in}{0.699667in}}{\pgfqpoint{0.096250in}{1.925000in}}%
\pgfusepath{clip}%
\pgfsetbuttcap%
\pgfsetroundjoin%
\definecolor{currentfill}{rgb}{0.857763,0.279857,0.415496}%
\pgfsetfillcolor{currentfill}%
\pgfsetfillopacity{0.800000}%
\pgfsetlinewidth{0.000000pt}%
\definecolor{currentstroke}{rgb}{0.000000,0.000000,0.000000}%
\pgfsetstrokecolor{currentstroke}%
\pgfsetdash{}{0pt}%
\pgfpathmoveto{\pgfqpoint{3.393750in}{1.835116in}}%
\pgfpathlineto{\pgfqpoint{3.490000in}{1.835116in}}%
\pgfpathlineto{\pgfqpoint{3.490000in}{1.842636in}}%
\pgfpathlineto{\pgfqpoint{3.393750in}{1.842636in}}%
\pgfpathlineto{\pgfqpoint{3.393750in}{1.835116in}}%
\pgfusepath{fill}%
\end{pgfscope}%
\begin{pgfscope}%
\pgfpathrectangle{\pgfqpoint{3.393750in}{0.699667in}}{\pgfqpoint{0.096250in}{1.925000in}}%
\pgfusepath{clip}%
\pgfsetbuttcap%
\pgfsetroundjoin%
\definecolor{currentfill}{rgb}{0.863320,0.283729,0.412403}%
\pgfsetfillcolor{currentfill}%
\pgfsetfillopacity{0.800000}%
\pgfsetlinewidth{0.000000pt}%
\definecolor{currentstroke}{rgb}{0.000000,0.000000,0.000000}%
\pgfsetstrokecolor{currentstroke}%
\pgfsetdash{}{0pt}%
\pgfpathmoveto{\pgfqpoint{3.393750in}{1.842636in}}%
\pgfpathlineto{\pgfqpoint{3.490000in}{1.842636in}}%
\pgfpathlineto{\pgfqpoint{3.490000in}{1.850155in}}%
\pgfpathlineto{\pgfqpoint{3.393750in}{1.850155in}}%
\pgfpathlineto{\pgfqpoint{3.393750in}{1.842636in}}%
\pgfusepath{fill}%
\end{pgfscope}%
\begin{pgfscope}%
\pgfpathrectangle{\pgfqpoint{3.393750in}{0.699667in}}{\pgfqpoint{0.096250in}{1.925000in}}%
\pgfusepath{clip}%
\pgfsetbuttcap%
\pgfsetroundjoin%
\definecolor{currentfill}{rgb}{0.868793,0.287728,0.409303}%
\pgfsetfillcolor{currentfill}%
\pgfsetfillopacity{0.800000}%
\pgfsetlinewidth{0.000000pt}%
\definecolor{currentstroke}{rgb}{0.000000,0.000000,0.000000}%
\pgfsetstrokecolor{currentstroke}%
\pgfsetdash{}{0pt}%
\pgfpathmoveto{\pgfqpoint{3.393750in}{1.850155in}}%
\pgfpathlineto{\pgfqpoint{3.490000in}{1.850155in}}%
\pgfpathlineto{\pgfqpoint{3.490000in}{1.857675in}}%
\pgfpathlineto{\pgfqpoint{3.393750in}{1.857675in}}%
\pgfpathlineto{\pgfqpoint{3.393750in}{1.850155in}}%
\pgfusepath{fill}%
\end{pgfscope}%
\begin{pgfscope}%
\pgfpathrectangle{\pgfqpoint{3.393750in}{0.699667in}}{\pgfqpoint{0.096250in}{1.925000in}}%
\pgfusepath{clip}%
\pgfsetbuttcap%
\pgfsetroundjoin%
\definecolor{currentfill}{rgb}{0.874176,0.291859,0.406205}%
\pgfsetfillcolor{currentfill}%
\pgfsetfillopacity{0.800000}%
\pgfsetlinewidth{0.000000pt}%
\definecolor{currentstroke}{rgb}{0.000000,0.000000,0.000000}%
\pgfsetstrokecolor{currentstroke}%
\pgfsetdash{}{0pt}%
\pgfpathmoveto{\pgfqpoint{3.393750in}{1.857675in}}%
\pgfpathlineto{\pgfqpoint{3.490000in}{1.857675in}}%
\pgfpathlineto{\pgfqpoint{3.490000in}{1.865195in}}%
\pgfpathlineto{\pgfqpoint{3.393750in}{1.865195in}}%
\pgfpathlineto{\pgfqpoint{3.393750in}{1.857675in}}%
\pgfusepath{fill}%
\end{pgfscope}%
\begin{pgfscope}%
\pgfpathrectangle{\pgfqpoint{3.393750in}{0.699667in}}{\pgfqpoint{0.096250in}{1.925000in}}%
\pgfusepath{clip}%
\pgfsetbuttcap%
\pgfsetroundjoin%
\definecolor{currentfill}{rgb}{0.879464,0.296125,0.403118}%
\pgfsetfillcolor{currentfill}%
\pgfsetfillopacity{0.800000}%
\pgfsetlinewidth{0.000000pt}%
\definecolor{currentstroke}{rgb}{0.000000,0.000000,0.000000}%
\pgfsetstrokecolor{currentstroke}%
\pgfsetdash{}{0pt}%
\pgfpathmoveto{\pgfqpoint{3.393750in}{1.865195in}}%
\pgfpathlineto{\pgfqpoint{3.490000in}{1.865195in}}%
\pgfpathlineto{\pgfqpoint{3.490000in}{1.872714in}}%
\pgfpathlineto{\pgfqpoint{3.393750in}{1.872714in}}%
\pgfpathlineto{\pgfqpoint{3.393750in}{1.865195in}}%
\pgfusepath{fill}%
\end{pgfscope}%
\begin{pgfscope}%
\pgfpathrectangle{\pgfqpoint{3.393750in}{0.699667in}}{\pgfqpoint{0.096250in}{1.925000in}}%
\pgfusepath{clip}%
\pgfsetbuttcap%
\pgfsetroundjoin%
\definecolor{currentfill}{rgb}{0.884651,0.300530,0.400047}%
\pgfsetfillcolor{currentfill}%
\pgfsetfillopacity{0.800000}%
\pgfsetlinewidth{0.000000pt}%
\definecolor{currentstroke}{rgb}{0.000000,0.000000,0.000000}%
\pgfsetstrokecolor{currentstroke}%
\pgfsetdash{}{0pt}%
\pgfpathmoveto{\pgfqpoint{3.393750in}{1.872714in}}%
\pgfpathlineto{\pgfqpoint{3.490000in}{1.872714in}}%
\pgfpathlineto{\pgfqpoint{3.490000in}{1.880234in}}%
\pgfpathlineto{\pgfqpoint{3.393750in}{1.880234in}}%
\pgfpathlineto{\pgfqpoint{3.393750in}{1.872714in}}%
\pgfusepath{fill}%
\end{pgfscope}%
\begin{pgfscope}%
\pgfpathrectangle{\pgfqpoint{3.393750in}{0.699667in}}{\pgfqpoint{0.096250in}{1.925000in}}%
\pgfusepath{clip}%
\pgfsetbuttcap%
\pgfsetroundjoin%
\definecolor{currentfill}{rgb}{0.889731,0.305079,0.397002}%
\pgfsetfillcolor{currentfill}%
\pgfsetfillopacity{0.800000}%
\pgfsetlinewidth{0.000000pt}%
\definecolor{currentstroke}{rgb}{0.000000,0.000000,0.000000}%
\pgfsetstrokecolor{currentstroke}%
\pgfsetdash{}{0pt}%
\pgfpathmoveto{\pgfqpoint{3.393750in}{1.880234in}}%
\pgfpathlineto{\pgfqpoint{3.490000in}{1.880234in}}%
\pgfpathlineto{\pgfqpoint{3.490000in}{1.887753in}}%
\pgfpathlineto{\pgfqpoint{3.393750in}{1.887753in}}%
\pgfpathlineto{\pgfqpoint{3.393750in}{1.880234in}}%
\pgfusepath{fill}%
\end{pgfscope}%
\begin{pgfscope}%
\pgfpathrectangle{\pgfqpoint{3.393750in}{0.699667in}}{\pgfqpoint{0.096250in}{1.925000in}}%
\pgfusepath{clip}%
\pgfsetbuttcap%
\pgfsetroundjoin%
\definecolor{currentfill}{rgb}{0.894700,0.309773,0.393995}%
\pgfsetfillcolor{currentfill}%
\pgfsetfillopacity{0.800000}%
\pgfsetlinewidth{0.000000pt}%
\definecolor{currentstroke}{rgb}{0.000000,0.000000,0.000000}%
\pgfsetstrokecolor{currentstroke}%
\pgfsetdash{}{0pt}%
\pgfpathmoveto{\pgfqpoint{3.393750in}{1.887753in}}%
\pgfpathlineto{\pgfqpoint{3.490000in}{1.887753in}}%
\pgfpathlineto{\pgfqpoint{3.490000in}{1.895273in}}%
\pgfpathlineto{\pgfqpoint{3.393750in}{1.895273in}}%
\pgfpathlineto{\pgfqpoint{3.393750in}{1.887753in}}%
\pgfusepath{fill}%
\end{pgfscope}%
\begin{pgfscope}%
\pgfpathrectangle{\pgfqpoint{3.393750in}{0.699667in}}{\pgfqpoint{0.096250in}{1.925000in}}%
\pgfusepath{clip}%
\pgfsetbuttcap%
\pgfsetroundjoin%
\definecolor{currentfill}{rgb}{0.899552,0.314616,0.391037}%
\pgfsetfillcolor{currentfill}%
\pgfsetfillopacity{0.800000}%
\pgfsetlinewidth{0.000000pt}%
\definecolor{currentstroke}{rgb}{0.000000,0.000000,0.000000}%
\pgfsetstrokecolor{currentstroke}%
\pgfsetdash{}{0pt}%
\pgfpathmoveto{\pgfqpoint{3.393750in}{1.895273in}}%
\pgfpathlineto{\pgfqpoint{3.490000in}{1.895273in}}%
\pgfpathlineto{\pgfqpoint{3.490000in}{1.902792in}}%
\pgfpathlineto{\pgfqpoint{3.393750in}{1.902792in}}%
\pgfpathlineto{\pgfqpoint{3.393750in}{1.895273in}}%
\pgfusepath{fill}%
\end{pgfscope}%
\begin{pgfscope}%
\pgfpathrectangle{\pgfqpoint{3.393750in}{0.699667in}}{\pgfqpoint{0.096250in}{1.925000in}}%
\pgfusepath{clip}%
\pgfsetbuttcap%
\pgfsetroundjoin%
\definecolor{currentfill}{rgb}{0.904281,0.319610,0.388137}%
\pgfsetfillcolor{currentfill}%
\pgfsetfillopacity{0.800000}%
\pgfsetlinewidth{0.000000pt}%
\definecolor{currentstroke}{rgb}{0.000000,0.000000,0.000000}%
\pgfsetstrokecolor{currentstroke}%
\pgfsetdash{}{0pt}%
\pgfpathmoveto{\pgfqpoint{3.393750in}{1.902792in}}%
\pgfpathlineto{\pgfqpoint{3.490000in}{1.902792in}}%
\pgfpathlineto{\pgfqpoint{3.490000in}{1.910312in}}%
\pgfpathlineto{\pgfqpoint{3.393750in}{1.910312in}}%
\pgfpathlineto{\pgfqpoint{3.393750in}{1.902792in}}%
\pgfusepath{fill}%
\end{pgfscope}%
\begin{pgfscope}%
\pgfpathrectangle{\pgfqpoint{3.393750in}{0.699667in}}{\pgfqpoint{0.096250in}{1.925000in}}%
\pgfusepath{clip}%
\pgfsetbuttcap%
\pgfsetroundjoin%
\definecolor{currentfill}{rgb}{0.908884,0.324755,0.385308}%
\pgfsetfillcolor{currentfill}%
\pgfsetfillopacity{0.800000}%
\pgfsetlinewidth{0.000000pt}%
\definecolor{currentstroke}{rgb}{0.000000,0.000000,0.000000}%
\pgfsetstrokecolor{currentstroke}%
\pgfsetdash{}{0pt}%
\pgfpathmoveto{\pgfqpoint{3.393750in}{1.910312in}}%
\pgfpathlineto{\pgfqpoint{3.490000in}{1.910312in}}%
\pgfpathlineto{\pgfqpoint{3.490000in}{1.917831in}}%
\pgfpathlineto{\pgfqpoint{3.393750in}{1.917831in}}%
\pgfpathlineto{\pgfqpoint{3.393750in}{1.910312in}}%
\pgfusepath{fill}%
\end{pgfscope}%
\begin{pgfscope}%
\pgfpathrectangle{\pgfqpoint{3.393750in}{0.699667in}}{\pgfqpoint{0.096250in}{1.925000in}}%
\pgfusepath{clip}%
\pgfsetbuttcap%
\pgfsetroundjoin%
\definecolor{currentfill}{rgb}{0.913354,0.330052,0.382563}%
\pgfsetfillcolor{currentfill}%
\pgfsetfillopacity{0.800000}%
\pgfsetlinewidth{0.000000pt}%
\definecolor{currentstroke}{rgb}{0.000000,0.000000,0.000000}%
\pgfsetstrokecolor{currentstroke}%
\pgfsetdash{}{0pt}%
\pgfpathmoveto{\pgfqpoint{3.393750in}{1.917831in}}%
\pgfpathlineto{\pgfqpoint{3.490000in}{1.917831in}}%
\pgfpathlineto{\pgfqpoint{3.490000in}{1.925351in}}%
\pgfpathlineto{\pgfqpoint{3.393750in}{1.925351in}}%
\pgfpathlineto{\pgfqpoint{3.393750in}{1.917831in}}%
\pgfusepath{fill}%
\end{pgfscope}%
\begin{pgfscope}%
\pgfpathrectangle{\pgfqpoint{3.393750in}{0.699667in}}{\pgfqpoint{0.096250in}{1.925000in}}%
\pgfusepath{clip}%
\pgfsetbuttcap%
\pgfsetroundjoin%
\definecolor{currentfill}{rgb}{0.917689,0.335500,0.379915}%
\pgfsetfillcolor{currentfill}%
\pgfsetfillopacity{0.800000}%
\pgfsetlinewidth{0.000000pt}%
\definecolor{currentstroke}{rgb}{0.000000,0.000000,0.000000}%
\pgfsetstrokecolor{currentstroke}%
\pgfsetdash{}{0pt}%
\pgfpathmoveto{\pgfqpoint{3.393750in}{1.925351in}}%
\pgfpathlineto{\pgfqpoint{3.490000in}{1.925351in}}%
\pgfpathlineto{\pgfqpoint{3.490000in}{1.932870in}}%
\pgfpathlineto{\pgfqpoint{3.393750in}{1.932870in}}%
\pgfpathlineto{\pgfqpoint{3.393750in}{1.925351in}}%
\pgfusepath{fill}%
\end{pgfscope}%
\begin{pgfscope}%
\pgfpathrectangle{\pgfqpoint{3.393750in}{0.699667in}}{\pgfqpoint{0.096250in}{1.925000in}}%
\pgfusepath{clip}%
\pgfsetbuttcap%
\pgfsetroundjoin%
\definecolor{currentfill}{rgb}{0.921884,0.341098,0.377376}%
\pgfsetfillcolor{currentfill}%
\pgfsetfillopacity{0.800000}%
\pgfsetlinewidth{0.000000pt}%
\definecolor{currentstroke}{rgb}{0.000000,0.000000,0.000000}%
\pgfsetstrokecolor{currentstroke}%
\pgfsetdash{}{0pt}%
\pgfpathmoveto{\pgfqpoint{3.393750in}{1.932870in}}%
\pgfpathlineto{\pgfqpoint{3.490000in}{1.932870in}}%
\pgfpathlineto{\pgfqpoint{3.490000in}{1.940390in}}%
\pgfpathlineto{\pgfqpoint{3.393750in}{1.940390in}}%
\pgfpathlineto{\pgfqpoint{3.393750in}{1.932870in}}%
\pgfusepath{fill}%
\end{pgfscope}%
\begin{pgfscope}%
\pgfpathrectangle{\pgfqpoint{3.393750in}{0.699667in}}{\pgfqpoint{0.096250in}{1.925000in}}%
\pgfusepath{clip}%
\pgfsetbuttcap%
\pgfsetroundjoin%
\definecolor{currentfill}{rgb}{0.925937,0.346844,0.374959}%
\pgfsetfillcolor{currentfill}%
\pgfsetfillopacity{0.800000}%
\pgfsetlinewidth{0.000000pt}%
\definecolor{currentstroke}{rgb}{0.000000,0.000000,0.000000}%
\pgfsetstrokecolor{currentstroke}%
\pgfsetdash{}{0pt}%
\pgfpathmoveto{\pgfqpoint{3.393750in}{1.940390in}}%
\pgfpathlineto{\pgfqpoint{3.490000in}{1.940390in}}%
\pgfpathlineto{\pgfqpoint{3.490000in}{1.947909in}}%
\pgfpathlineto{\pgfqpoint{3.393750in}{1.947909in}}%
\pgfpathlineto{\pgfqpoint{3.393750in}{1.940390in}}%
\pgfusepath{fill}%
\end{pgfscope}%
\begin{pgfscope}%
\pgfpathrectangle{\pgfqpoint{3.393750in}{0.699667in}}{\pgfqpoint{0.096250in}{1.925000in}}%
\pgfusepath{clip}%
\pgfsetbuttcap%
\pgfsetroundjoin%
\definecolor{currentfill}{rgb}{0.929845,0.352734,0.372677}%
\pgfsetfillcolor{currentfill}%
\pgfsetfillopacity{0.800000}%
\pgfsetlinewidth{0.000000pt}%
\definecolor{currentstroke}{rgb}{0.000000,0.000000,0.000000}%
\pgfsetstrokecolor{currentstroke}%
\pgfsetdash{}{0pt}%
\pgfpathmoveto{\pgfqpoint{3.393750in}{1.947909in}}%
\pgfpathlineto{\pgfqpoint{3.490000in}{1.947909in}}%
\pgfpathlineto{\pgfqpoint{3.490000in}{1.955429in}}%
\pgfpathlineto{\pgfqpoint{3.393750in}{1.955429in}}%
\pgfpathlineto{\pgfqpoint{3.393750in}{1.947909in}}%
\pgfusepath{fill}%
\end{pgfscope}%
\begin{pgfscope}%
\pgfpathrectangle{\pgfqpoint{3.393750in}{0.699667in}}{\pgfqpoint{0.096250in}{1.925000in}}%
\pgfusepath{clip}%
\pgfsetbuttcap%
\pgfsetroundjoin%
\definecolor{currentfill}{rgb}{0.933606,0.358764,0.370541}%
\pgfsetfillcolor{currentfill}%
\pgfsetfillopacity{0.800000}%
\pgfsetlinewidth{0.000000pt}%
\definecolor{currentstroke}{rgb}{0.000000,0.000000,0.000000}%
\pgfsetstrokecolor{currentstroke}%
\pgfsetdash{}{0pt}%
\pgfpathmoveto{\pgfqpoint{3.393750in}{1.955429in}}%
\pgfpathlineto{\pgfqpoint{3.490000in}{1.955429in}}%
\pgfpathlineto{\pgfqpoint{3.490000in}{1.962948in}}%
\pgfpathlineto{\pgfqpoint{3.393750in}{1.962948in}}%
\pgfpathlineto{\pgfqpoint{3.393750in}{1.955429in}}%
\pgfusepath{fill}%
\end{pgfscope}%
\begin{pgfscope}%
\pgfpathrectangle{\pgfqpoint{3.393750in}{0.699667in}}{\pgfqpoint{0.096250in}{1.925000in}}%
\pgfusepath{clip}%
\pgfsetbuttcap%
\pgfsetroundjoin%
\definecolor{currentfill}{rgb}{0.937221,0.364929,0.368567}%
\pgfsetfillcolor{currentfill}%
\pgfsetfillopacity{0.800000}%
\pgfsetlinewidth{0.000000pt}%
\definecolor{currentstroke}{rgb}{0.000000,0.000000,0.000000}%
\pgfsetstrokecolor{currentstroke}%
\pgfsetdash{}{0pt}%
\pgfpathmoveto{\pgfqpoint{3.393750in}{1.962948in}}%
\pgfpathlineto{\pgfqpoint{3.490000in}{1.962948in}}%
\pgfpathlineto{\pgfqpoint{3.490000in}{1.970468in}}%
\pgfpathlineto{\pgfqpoint{3.393750in}{1.970468in}}%
\pgfpathlineto{\pgfqpoint{3.393750in}{1.962948in}}%
\pgfusepath{fill}%
\end{pgfscope}%
\begin{pgfscope}%
\pgfpathrectangle{\pgfqpoint{3.393750in}{0.699667in}}{\pgfqpoint{0.096250in}{1.925000in}}%
\pgfusepath{clip}%
\pgfsetbuttcap%
\pgfsetroundjoin%
\definecolor{currentfill}{rgb}{0.940687,0.371224,0.366762}%
\pgfsetfillcolor{currentfill}%
\pgfsetfillopacity{0.800000}%
\pgfsetlinewidth{0.000000pt}%
\definecolor{currentstroke}{rgb}{0.000000,0.000000,0.000000}%
\pgfsetstrokecolor{currentstroke}%
\pgfsetdash{}{0pt}%
\pgfpathmoveto{\pgfqpoint{3.393750in}{1.970468in}}%
\pgfpathlineto{\pgfqpoint{3.490000in}{1.970468in}}%
\pgfpathlineto{\pgfqpoint{3.490000in}{1.977987in}}%
\pgfpathlineto{\pgfqpoint{3.393750in}{1.977987in}}%
\pgfpathlineto{\pgfqpoint{3.393750in}{1.970468in}}%
\pgfusepath{fill}%
\end{pgfscope}%
\begin{pgfscope}%
\pgfpathrectangle{\pgfqpoint{3.393750in}{0.699667in}}{\pgfqpoint{0.096250in}{1.925000in}}%
\pgfusepath{clip}%
\pgfsetbuttcap%
\pgfsetroundjoin%
\definecolor{currentfill}{rgb}{0.944006,0.377643,0.365136}%
\pgfsetfillcolor{currentfill}%
\pgfsetfillopacity{0.800000}%
\pgfsetlinewidth{0.000000pt}%
\definecolor{currentstroke}{rgb}{0.000000,0.000000,0.000000}%
\pgfsetstrokecolor{currentstroke}%
\pgfsetdash{}{0pt}%
\pgfpathmoveto{\pgfqpoint{3.393750in}{1.977987in}}%
\pgfpathlineto{\pgfqpoint{3.490000in}{1.977987in}}%
\pgfpathlineto{\pgfqpoint{3.490000in}{1.985507in}}%
\pgfpathlineto{\pgfqpoint{3.393750in}{1.985507in}}%
\pgfpathlineto{\pgfqpoint{3.393750in}{1.977987in}}%
\pgfusepath{fill}%
\end{pgfscope}%
\begin{pgfscope}%
\pgfpathrectangle{\pgfqpoint{3.393750in}{0.699667in}}{\pgfqpoint{0.096250in}{1.925000in}}%
\pgfusepath{clip}%
\pgfsetbuttcap%
\pgfsetroundjoin%
\definecolor{currentfill}{rgb}{0.947180,0.384178,0.363701}%
\pgfsetfillcolor{currentfill}%
\pgfsetfillopacity{0.800000}%
\pgfsetlinewidth{0.000000pt}%
\definecolor{currentstroke}{rgb}{0.000000,0.000000,0.000000}%
\pgfsetstrokecolor{currentstroke}%
\pgfsetdash{}{0pt}%
\pgfpathmoveto{\pgfqpoint{3.393750in}{1.985507in}}%
\pgfpathlineto{\pgfqpoint{3.490000in}{1.985507in}}%
\pgfpathlineto{\pgfqpoint{3.490000in}{1.993027in}}%
\pgfpathlineto{\pgfqpoint{3.393750in}{1.993027in}}%
\pgfpathlineto{\pgfqpoint{3.393750in}{1.985507in}}%
\pgfusepath{fill}%
\end{pgfscope}%
\begin{pgfscope}%
\pgfpathrectangle{\pgfqpoint{3.393750in}{0.699667in}}{\pgfqpoint{0.096250in}{1.925000in}}%
\pgfusepath{clip}%
\pgfsetbuttcap%
\pgfsetroundjoin%
\definecolor{currentfill}{rgb}{0.950210,0.390820,0.362468}%
\pgfsetfillcolor{currentfill}%
\pgfsetfillopacity{0.800000}%
\pgfsetlinewidth{0.000000pt}%
\definecolor{currentstroke}{rgb}{0.000000,0.000000,0.000000}%
\pgfsetstrokecolor{currentstroke}%
\pgfsetdash{}{0pt}%
\pgfpathmoveto{\pgfqpoint{3.393750in}{1.993027in}}%
\pgfpathlineto{\pgfqpoint{3.490000in}{1.993027in}}%
\pgfpathlineto{\pgfqpoint{3.490000in}{2.000546in}}%
\pgfpathlineto{\pgfqpoint{3.393750in}{2.000546in}}%
\pgfpathlineto{\pgfqpoint{3.393750in}{1.993027in}}%
\pgfusepath{fill}%
\end{pgfscope}%
\begin{pgfscope}%
\pgfpathrectangle{\pgfqpoint{3.393750in}{0.699667in}}{\pgfqpoint{0.096250in}{1.925000in}}%
\pgfusepath{clip}%
\pgfsetbuttcap%
\pgfsetroundjoin%
\definecolor{currentfill}{rgb}{0.953099,0.397563,0.361438}%
\pgfsetfillcolor{currentfill}%
\pgfsetfillopacity{0.800000}%
\pgfsetlinewidth{0.000000pt}%
\definecolor{currentstroke}{rgb}{0.000000,0.000000,0.000000}%
\pgfsetstrokecolor{currentstroke}%
\pgfsetdash{}{0pt}%
\pgfpathmoveto{\pgfqpoint{3.393750in}{2.000546in}}%
\pgfpathlineto{\pgfqpoint{3.490000in}{2.000546in}}%
\pgfpathlineto{\pgfqpoint{3.490000in}{2.008066in}}%
\pgfpathlineto{\pgfqpoint{3.393750in}{2.008066in}}%
\pgfpathlineto{\pgfqpoint{3.393750in}{2.000546in}}%
\pgfusepath{fill}%
\end{pgfscope}%
\begin{pgfscope}%
\pgfpathrectangle{\pgfqpoint{3.393750in}{0.699667in}}{\pgfqpoint{0.096250in}{1.925000in}}%
\pgfusepath{clip}%
\pgfsetbuttcap%
\pgfsetroundjoin%
\definecolor{currentfill}{rgb}{0.955849,0.404400,0.360619}%
\pgfsetfillcolor{currentfill}%
\pgfsetfillopacity{0.800000}%
\pgfsetlinewidth{0.000000pt}%
\definecolor{currentstroke}{rgb}{0.000000,0.000000,0.000000}%
\pgfsetstrokecolor{currentstroke}%
\pgfsetdash{}{0pt}%
\pgfpathmoveto{\pgfqpoint{3.393750in}{2.008066in}}%
\pgfpathlineto{\pgfqpoint{3.490000in}{2.008066in}}%
\pgfpathlineto{\pgfqpoint{3.490000in}{2.015585in}}%
\pgfpathlineto{\pgfqpoint{3.393750in}{2.015585in}}%
\pgfpathlineto{\pgfqpoint{3.393750in}{2.008066in}}%
\pgfusepath{fill}%
\end{pgfscope}%
\begin{pgfscope}%
\pgfpathrectangle{\pgfqpoint{3.393750in}{0.699667in}}{\pgfqpoint{0.096250in}{1.925000in}}%
\pgfusepath{clip}%
\pgfsetbuttcap%
\pgfsetroundjoin%
\definecolor{currentfill}{rgb}{0.958464,0.411324,0.360014}%
\pgfsetfillcolor{currentfill}%
\pgfsetfillopacity{0.800000}%
\pgfsetlinewidth{0.000000pt}%
\definecolor{currentstroke}{rgb}{0.000000,0.000000,0.000000}%
\pgfsetstrokecolor{currentstroke}%
\pgfsetdash{}{0pt}%
\pgfpathmoveto{\pgfqpoint{3.393750in}{2.015585in}}%
\pgfpathlineto{\pgfqpoint{3.490000in}{2.015585in}}%
\pgfpathlineto{\pgfqpoint{3.490000in}{2.023105in}}%
\pgfpathlineto{\pgfqpoint{3.393750in}{2.023105in}}%
\pgfpathlineto{\pgfqpoint{3.393750in}{2.015585in}}%
\pgfusepath{fill}%
\end{pgfscope}%
\begin{pgfscope}%
\pgfpathrectangle{\pgfqpoint{3.393750in}{0.699667in}}{\pgfqpoint{0.096250in}{1.925000in}}%
\pgfusepath{clip}%
\pgfsetbuttcap%
\pgfsetroundjoin%
\definecolor{currentfill}{rgb}{0.960949,0.418323,0.359630}%
\pgfsetfillcolor{currentfill}%
\pgfsetfillopacity{0.800000}%
\pgfsetlinewidth{0.000000pt}%
\definecolor{currentstroke}{rgb}{0.000000,0.000000,0.000000}%
\pgfsetstrokecolor{currentstroke}%
\pgfsetdash{}{0pt}%
\pgfpathmoveto{\pgfqpoint{3.393750in}{2.023105in}}%
\pgfpathlineto{\pgfqpoint{3.490000in}{2.023105in}}%
\pgfpathlineto{\pgfqpoint{3.490000in}{2.030624in}}%
\pgfpathlineto{\pgfqpoint{3.393750in}{2.030624in}}%
\pgfpathlineto{\pgfqpoint{3.393750in}{2.023105in}}%
\pgfusepath{fill}%
\end{pgfscope}%
\begin{pgfscope}%
\pgfpathrectangle{\pgfqpoint{3.393750in}{0.699667in}}{\pgfqpoint{0.096250in}{1.925000in}}%
\pgfusepath{clip}%
\pgfsetbuttcap%
\pgfsetroundjoin%
\definecolor{currentfill}{rgb}{0.963310,0.425390,0.359469}%
\pgfsetfillcolor{currentfill}%
\pgfsetfillopacity{0.800000}%
\pgfsetlinewidth{0.000000pt}%
\definecolor{currentstroke}{rgb}{0.000000,0.000000,0.000000}%
\pgfsetstrokecolor{currentstroke}%
\pgfsetdash{}{0pt}%
\pgfpathmoveto{\pgfqpoint{3.393750in}{2.030624in}}%
\pgfpathlineto{\pgfqpoint{3.490000in}{2.030624in}}%
\pgfpathlineto{\pgfqpoint{3.490000in}{2.038144in}}%
\pgfpathlineto{\pgfqpoint{3.393750in}{2.038144in}}%
\pgfpathlineto{\pgfqpoint{3.393750in}{2.030624in}}%
\pgfusepath{fill}%
\end{pgfscope}%
\begin{pgfscope}%
\pgfpathrectangle{\pgfqpoint{3.393750in}{0.699667in}}{\pgfqpoint{0.096250in}{1.925000in}}%
\pgfusepath{clip}%
\pgfsetbuttcap%
\pgfsetroundjoin%
\definecolor{currentfill}{rgb}{0.965549,0.432519,0.359529}%
\pgfsetfillcolor{currentfill}%
\pgfsetfillopacity{0.800000}%
\pgfsetlinewidth{0.000000pt}%
\definecolor{currentstroke}{rgb}{0.000000,0.000000,0.000000}%
\pgfsetstrokecolor{currentstroke}%
\pgfsetdash{}{0pt}%
\pgfpathmoveto{\pgfqpoint{3.393750in}{2.038144in}}%
\pgfpathlineto{\pgfqpoint{3.490000in}{2.038144in}}%
\pgfpathlineto{\pgfqpoint{3.490000in}{2.045663in}}%
\pgfpathlineto{\pgfqpoint{3.393750in}{2.045663in}}%
\pgfpathlineto{\pgfqpoint{3.393750in}{2.038144in}}%
\pgfusepath{fill}%
\end{pgfscope}%
\begin{pgfscope}%
\pgfpathrectangle{\pgfqpoint{3.393750in}{0.699667in}}{\pgfqpoint{0.096250in}{1.925000in}}%
\pgfusepath{clip}%
\pgfsetbuttcap%
\pgfsetroundjoin%
\definecolor{currentfill}{rgb}{0.967671,0.439703,0.359810}%
\pgfsetfillcolor{currentfill}%
\pgfsetfillopacity{0.800000}%
\pgfsetlinewidth{0.000000pt}%
\definecolor{currentstroke}{rgb}{0.000000,0.000000,0.000000}%
\pgfsetstrokecolor{currentstroke}%
\pgfsetdash{}{0pt}%
\pgfpathmoveto{\pgfqpoint{3.393750in}{2.045663in}}%
\pgfpathlineto{\pgfqpoint{3.490000in}{2.045663in}}%
\pgfpathlineto{\pgfqpoint{3.490000in}{2.053183in}}%
\pgfpathlineto{\pgfqpoint{3.393750in}{2.053183in}}%
\pgfpathlineto{\pgfqpoint{3.393750in}{2.045663in}}%
\pgfusepath{fill}%
\end{pgfscope}%
\begin{pgfscope}%
\pgfpathrectangle{\pgfqpoint{3.393750in}{0.699667in}}{\pgfqpoint{0.096250in}{1.925000in}}%
\pgfusepath{clip}%
\pgfsetbuttcap%
\pgfsetroundjoin%
\definecolor{currentfill}{rgb}{0.969680,0.446936,0.360311}%
\pgfsetfillcolor{currentfill}%
\pgfsetfillopacity{0.800000}%
\pgfsetlinewidth{0.000000pt}%
\definecolor{currentstroke}{rgb}{0.000000,0.000000,0.000000}%
\pgfsetstrokecolor{currentstroke}%
\pgfsetdash{}{0pt}%
\pgfpathmoveto{\pgfqpoint{3.393750in}{2.053183in}}%
\pgfpathlineto{\pgfqpoint{3.490000in}{2.053183in}}%
\pgfpathlineto{\pgfqpoint{3.490000in}{2.060702in}}%
\pgfpathlineto{\pgfqpoint{3.393750in}{2.060702in}}%
\pgfpathlineto{\pgfqpoint{3.393750in}{2.053183in}}%
\pgfusepath{fill}%
\end{pgfscope}%
\begin{pgfscope}%
\pgfpathrectangle{\pgfqpoint{3.393750in}{0.699667in}}{\pgfqpoint{0.096250in}{1.925000in}}%
\pgfusepath{clip}%
\pgfsetbuttcap%
\pgfsetroundjoin%
\definecolor{currentfill}{rgb}{0.971582,0.454210,0.361030}%
\pgfsetfillcolor{currentfill}%
\pgfsetfillopacity{0.800000}%
\pgfsetlinewidth{0.000000pt}%
\definecolor{currentstroke}{rgb}{0.000000,0.000000,0.000000}%
\pgfsetstrokecolor{currentstroke}%
\pgfsetdash{}{0pt}%
\pgfpathmoveto{\pgfqpoint{3.393750in}{2.060702in}}%
\pgfpathlineto{\pgfqpoint{3.490000in}{2.060702in}}%
\pgfpathlineto{\pgfqpoint{3.490000in}{2.068222in}}%
\pgfpathlineto{\pgfqpoint{3.393750in}{2.068222in}}%
\pgfpathlineto{\pgfqpoint{3.393750in}{2.060702in}}%
\pgfusepath{fill}%
\end{pgfscope}%
\begin{pgfscope}%
\pgfpathrectangle{\pgfqpoint{3.393750in}{0.699667in}}{\pgfqpoint{0.096250in}{1.925000in}}%
\pgfusepath{clip}%
\pgfsetbuttcap%
\pgfsetroundjoin%
\definecolor{currentfill}{rgb}{0.973381,0.461520,0.361965}%
\pgfsetfillcolor{currentfill}%
\pgfsetfillopacity{0.800000}%
\pgfsetlinewidth{0.000000pt}%
\definecolor{currentstroke}{rgb}{0.000000,0.000000,0.000000}%
\pgfsetstrokecolor{currentstroke}%
\pgfsetdash{}{0pt}%
\pgfpathmoveto{\pgfqpoint{3.393750in}{2.068222in}}%
\pgfpathlineto{\pgfqpoint{3.490000in}{2.068222in}}%
\pgfpathlineto{\pgfqpoint{3.490000in}{2.075741in}}%
\pgfpathlineto{\pgfqpoint{3.393750in}{2.075741in}}%
\pgfpathlineto{\pgfqpoint{3.393750in}{2.068222in}}%
\pgfusepath{fill}%
\end{pgfscope}%
\begin{pgfscope}%
\pgfpathrectangle{\pgfqpoint{3.393750in}{0.699667in}}{\pgfqpoint{0.096250in}{1.925000in}}%
\pgfusepath{clip}%
\pgfsetbuttcap%
\pgfsetroundjoin%
\definecolor{currentfill}{rgb}{0.975082,0.468861,0.363111}%
\pgfsetfillcolor{currentfill}%
\pgfsetfillopacity{0.800000}%
\pgfsetlinewidth{0.000000pt}%
\definecolor{currentstroke}{rgb}{0.000000,0.000000,0.000000}%
\pgfsetstrokecolor{currentstroke}%
\pgfsetdash{}{0pt}%
\pgfpathmoveto{\pgfqpoint{3.393750in}{2.075741in}}%
\pgfpathlineto{\pgfqpoint{3.490000in}{2.075741in}}%
\pgfpathlineto{\pgfqpoint{3.490000in}{2.083261in}}%
\pgfpathlineto{\pgfqpoint{3.393750in}{2.083261in}}%
\pgfpathlineto{\pgfqpoint{3.393750in}{2.075741in}}%
\pgfusepath{fill}%
\end{pgfscope}%
\begin{pgfscope}%
\pgfpathrectangle{\pgfqpoint{3.393750in}{0.699667in}}{\pgfqpoint{0.096250in}{1.925000in}}%
\pgfusepath{clip}%
\pgfsetbuttcap%
\pgfsetroundjoin%
\definecolor{currentfill}{rgb}{0.976690,0.476226,0.364466}%
\pgfsetfillcolor{currentfill}%
\pgfsetfillopacity{0.800000}%
\pgfsetlinewidth{0.000000pt}%
\definecolor{currentstroke}{rgb}{0.000000,0.000000,0.000000}%
\pgfsetstrokecolor{currentstroke}%
\pgfsetdash{}{0pt}%
\pgfpathmoveto{\pgfqpoint{3.393750in}{2.083261in}}%
\pgfpathlineto{\pgfqpoint{3.490000in}{2.083261in}}%
\pgfpathlineto{\pgfqpoint{3.490000in}{2.090780in}}%
\pgfpathlineto{\pgfqpoint{3.393750in}{2.090780in}}%
\pgfpathlineto{\pgfqpoint{3.393750in}{2.083261in}}%
\pgfusepath{fill}%
\end{pgfscope}%
\begin{pgfscope}%
\pgfpathrectangle{\pgfqpoint{3.393750in}{0.699667in}}{\pgfqpoint{0.096250in}{1.925000in}}%
\pgfusepath{clip}%
\pgfsetbuttcap%
\pgfsetroundjoin%
\definecolor{currentfill}{rgb}{0.978210,0.483612,0.366025}%
\pgfsetfillcolor{currentfill}%
\pgfsetfillopacity{0.800000}%
\pgfsetlinewidth{0.000000pt}%
\definecolor{currentstroke}{rgb}{0.000000,0.000000,0.000000}%
\pgfsetstrokecolor{currentstroke}%
\pgfsetdash{}{0pt}%
\pgfpathmoveto{\pgfqpoint{3.393750in}{2.090780in}}%
\pgfpathlineto{\pgfqpoint{3.490000in}{2.090780in}}%
\pgfpathlineto{\pgfqpoint{3.490000in}{2.098300in}}%
\pgfpathlineto{\pgfqpoint{3.393750in}{2.098300in}}%
\pgfpathlineto{\pgfqpoint{3.393750in}{2.090780in}}%
\pgfusepath{fill}%
\end{pgfscope}%
\begin{pgfscope}%
\pgfpathrectangle{\pgfqpoint{3.393750in}{0.699667in}}{\pgfqpoint{0.096250in}{1.925000in}}%
\pgfusepath{clip}%
\pgfsetbuttcap%
\pgfsetroundjoin%
\definecolor{currentfill}{rgb}{0.979645,0.491014,0.367783}%
\pgfsetfillcolor{currentfill}%
\pgfsetfillopacity{0.800000}%
\pgfsetlinewidth{0.000000pt}%
\definecolor{currentstroke}{rgb}{0.000000,0.000000,0.000000}%
\pgfsetstrokecolor{currentstroke}%
\pgfsetdash{}{0pt}%
\pgfpathmoveto{\pgfqpoint{3.393750in}{2.098300in}}%
\pgfpathlineto{\pgfqpoint{3.490000in}{2.098300in}}%
\pgfpathlineto{\pgfqpoint{3.490000in}{2.105820in}}%
\pgfpathlineto{\pgfqpoint{3.393750in}{2.105820in}}%
\pgfpathlineto{\pgfqpoint{3.393750in}{2.098300in}}%
\pgfusepath{fill}%
\end{pgfscope}%
\begin{pgfscope}%
\pgfpathrectangle{\pgfqpoint{3.393750in}{0.699667in}}{\pgfqpoint{0.096250in}{1.925000in}}%
\pgfusepath{clip}%
\pgfsetbuttcap%
\pgfsetroundjoin%
\definecolor{currentfill}{rgb}{0.981000,0.498428,0.369734}%
\pgfsetfillcolor{currentfill}%
\pgfsetfillopacity{0.800000}%
\pgfsetlinewidth{0.000000pt}%
\definecolor{currentstroke}{rgb}{0.000000,0.000000,0.000000}%
\pgfsetstrokecolor{currentstroke}%
\pgfsetdash{}{0pt}%
\pgfpathmoveto{\pgfqpoint{3.393750in}{2.105820in}}%
\pgfpathlineto{\pgfqpoint{3.490000in}{2.105820in}}%
\pgfpathlineto{\pgfqpoint{3.490000in}{2.113339in}}%
\pgfpathlineto{\pgfqpoint{3.393750in}{2.113339in}}%
\pgfpathlineto{\pgfqpoint{3.393750in}{2.105820in}}%
\pgfusepath{fill}%
\end{pgfscope}%
\begin{pgfscope}%
\pgfpathrectangle{\pgfqpoint{3.393750in}{0.699667in}}{\pgfqpoint{0.096250in}{1.925000in}}%
\pgfusepath{clip}%
\pgfsetbuttcap%
\pgfsetroundjoin%
\definecolor{currentfill}{rgb}{0.982279,0.505851,0.371874}%
\pgfsetfillcolor{currentfill}%
\pgfsetfillopacity{0.800000}%
\pgfsetlinewidth{0.000000pt}%
\definecolor{currentstroke}{rgb}{0.000000,0.000000,0.000000}%
\pgfsetstrokecolor{currentstroke}%
\pgfsetdash{}{0pt}%
\pgfpathmoveto{\pgfqpoint{3.393750in}{2.113339in}}%
\pgfpathlineto{\pgfqpoint{3.490000in}{2.113339in}}%
\pgfpathlineto{\pgfqpoint{3.490000in}{2.120859in}}%
\pgfpathlineto{\pgfqpoint{3.393750in}{2.120859in}}%
\pgfpathlineto{\pgfqpoint{3.393750in}{2.113339in}}%
\pgfusepath{fill}%
\end{pgfscope}%
\begin{pgfscope}%
\pgfpathrectangle{\pgfqpoint{3.393750in}{0.699667in}}{\pgfqpoint{0.096250in}{1.925000in}}%
\pgfusepath{clip}%
\pgfsetbuttcap%
\pgfsetroundjoin%
\definecolor{currentfill}{rgb}{0.983485,0.513280,0.374198}%
\pgfsetfillcolor{currentfill}%
\pgfsetfillopacity{0.800000}%
\pgfsetlinewidth{0.000000pt}%
\definecolor{currentstroke}{rgb}{0.000000,0.000000,0.000000}%
\pgfsetstrokecolor{currentstroke}%
\pgfsetdash{}{0pt}%
\pgfpathmoveto{\pgfqpoint{3.393750in}{2.120859in}}%
\pgfpathlineto{\pgfqpoint{3.490000in}{2.120859in}}%
\pgfpathlineto{\pgfqpoint{3.490000in}{2.128378in}}%
\pgfpathlineto{\pgfqpoint{3.393750in}{2.128378in}}%
\pgfpathlineto{\pgfqpoint{3.393750in}{2.120859in}}%
\pgfusepath{fill}%
\end{pgfscope}%
\begin{pgfscope}%
\pgfpathrectangle{\pgfqpoint{3.393750in}{0.699667in}}{\pgfqpoint{0.096250in}{1.925000in}}%
\pgfusepath{clip}%
\pgfsetbuttcap%
\pgfsetroundjoin%
\definecolor{currentfill}{rgb}{0.984622,0.520713,0.376698}%
\pgfsetfillcolor{currentfill}%
\pgfsetfillopacity{0.800000}%
\pgfsetlinewidth{0.000000pt}%
\definecolor{currentstroke}{rgb}{0.000000,0.000000,0.000000}%
\pgfsetstrokecolor{currentstroke}%
\pgfsetdash{}{0pt}%
\pgfpathmoveto{\pgfqpoint{3.393750in}{2.128378in}}%
\pgfpathlineto{\pgfqpoint{3.490000in}{2.128378in}}%
\pgfpathlineto{\pgfqpoint{3.490000in}{2.135898in}}%
\pgfpathlineto{\pgfqpoint{3.393750in}{2.135898in}}%
\pgfpathlineto{\pgfqpoint{3.393750in}{2.128378in}}%
\pgfusepath{fill}%
\end{pgfscope}%
\begin{pgfscope}%
\pgfpathrectangle{\pgfqpoint{3.393750in}{0.699667in}}{\pgfqpoint{0.096250in}{1.925000in}}%
\pgfusepath{clip}%
\pgfsetbuttcap%
\pgfsetroundjoin%
\definecolor{currentfill}{rgb}{0.985693,0.528148,0.379371}%
\pgfsetfillcolor{currentfill}%
\pgfsetfillopacity{0.800000}%
\pgfsetlinewidth{0.000000pt}%
\definecolor{currentstroke}{rgb}{0.000000,0.000000,0.000000}%
\pgfsetstrokecolor{currentstroke}%
\pgfsetdash{}{0pt}%
\pgfpathmoveto{\pgfqpoint{3.393750in}{2.135898in}}%
\pgfpathlineto{\pgfqpoint{3.490000in}{2.135898in}}%
\pgfpathlineto{\pgfqpoint{3.490000in}{2.143417in}}%
\pgfpathlineto{\pgfqpoint{3.393750in}{2.143417in}}%
\pgfpathlineto{\pgfqpoint{3.393750in}{2.135898in}}%
\pgfusepath{fill}%
\end{pgfscope}%
\begin{pgfscope}%
\pgfpathrectangle{\pgfqpoint{3.393750in}{0.699667in}}{\pgfqpoint{0.096250in}{1.925000in}}%
\pgfusepath{clip}%
\pgfsetbuttcap%
\pgfsetroundjoin%
\definecolor{currentfill}{rgb}{0.986700,0.535582,0.382210}%
\pgfsetfillcolor{currentfill}%
\pgfsetfillopacity{0.800000}%
\pgfsetlinewidth{0.000000pt}%
\definecolor{currentstroke}{rgb}{0.000000,0.000000,0.000000}%
\pgfsetstrokecolor{currentstroke}%
\pgfsetdash{}{0pt}%
\pgfpathmoveto{\pgfqpoint{3.393750in}{2.143417in}}%
\pgfpathlineto{\pgfqpoint{3.490000in}{2.143417in}}%
\pgfpathlineto{\pgfqpoint{3.490000in}{2.150937in}}%
\pgfpathlineto{\pgfqpoint{3.393750in}{2.150937in}}%
\pgfpathlineto{\pgfqpoint{3.393750in}{2.143417in}}%
\pgfusepath{fill}%
\end{pgfscope}%
\begin{pgfscope}%
\pgfpathrectangle{\pgfqpoint{3.393750in}{0.699667in}}{\pgfqpoint{0.096250in}{1.925000in}}%
\pgfusepath{clip}%
\pgfsetbuttcap%
\pgfsetroundjoin%
\definecolor{currentfill}{rgb}{0.987646,0.543015,0.385210}%
\pgfsetfillcolor{currentfill}%
\pgfsetfillopacity{0.800000}%
\pgfsetlinewidth{0.000000pt}%
\definecolor{currentstroke}{rgb}{0.000000,0.000000,0.000000}%
\pgfsetstrokecolor{currentstroke}%
\pgfsetdash{}{0pt}%
\pgfpathmoveto{\pgfqpoint{3.393750in}{2.150937in}}%
\pgfpathlineto{\pgfqpoint{3.490000in}{2.150937in}}%
\pgfpathlineto{\pgfqpoint{3.490000in}{2.158456in}}%
\pgfpathlineto{\pgfqpoint{3.393750in}{2.158456in}}%
\pgfpathlineto{\pgfqpoint{3.393750in}{2.150937in}}%
\pgfusepath{fill}%
\end{pgfscope}%
\begin{pgfscope}%
\pgfpathrectangle{\pgfqpoint{3.393750in}{0.699667in}}{\pgfqpoint{0.096250in}{1.925000in}}%
\pgfusepath{clip}%
\pgfsetbuttcap%
\pgfsetroundjoin%
\definecolor{currentfill}{rgb}{0.988533,0.550446,0.388365}%
\pgfsetfillcolor{currentfill}%
\pgfsetfillopacity{0.800000}%
\pgfsetlinewidth{0.000000pt}%
\definecolor{currentstroke}{rgb}{0.000000,0.000000,0.000000}%
\pgfsetstrokecolor{currentstroke}%
\pgfsetdash{}{0pt}%
\pgfpathmoveto{\pgfqpoint{3.393750in}{2.158456in}}%
\pgfpathlineto{\pgfqpoint{3.490000in}{2.158456in}}%
\pgfpathlineto{\pgfqpoint{3.490000in}{2.165976in}}%
\pgfpathlineto{\pgfqpoint{3.393750in}{2.165976in}}%
\pgfpathlineto{\pgfqpoint{3.393750in}{2.158456in}}%
\pgfusepath{fill}%
\end{pgfscope}%
\begin{pgfscope}%
\pgfpathrectangle{\pgfqpoint{3.393750in}{0.699667in}}{\pgfqpoint{0.096250in}{1.925000in}}%
\pgfusepath{clip}%
\pgfsetbuttcap%
\pgfsetroundjoin%
\definecolor{currentfill}{rgb}{0.989363,0.557873,0.391671}%
\pgfsetfillcolor{currentfill}%
\pgfsetfillopacity{0.800000}%
\pgfsetlinewidth{0.000000pt}%
\definecolor{currentstroke}{rgb}{0.000000,0.000000,0.000000}%
\pgfsetstrokecolor{currentstroke}%
\pgfsetdash{}{0pt}%
\pgfpathmoveto{\pgfqpoint{3.393750in}{2.165976in}}%
\pgfpathlineto{\pgfqpoint{3.490000in}{2.165976in}}%
\pgfpathlineto{\pgfqpoint{3.490000in}{2.173495in}}%
\pgfpathlineto{\pgfqpoint{3.393750in}{2.173495in}}%
\pgfpathlineto{\pgfqpoint{3.393750in}{2.165976in}}%
\pgfusepath{fill}%
\end{pgfscope}%
\begin{pgfscope}%
\pgfpathrectangle{\pgfqpoint{3.393750in}{0.699667in}}{\pgfqpoint{0.096250in}{1.925000in}}%
\pgfusepath{clip}%
\pgfsetbuttcap%
\pgfsetroundjoin%
\definecolor{currentfill}{rgb}{0.990138,0.565296,0.395122}%
\pgfsetfillcolor{currentfill}%
\pgfsetfillopacity{0.800000}%
\pgfsetlinewidth{0.000000pt}%
\definecolor{currentstroke}{rgb}{0.000000,0.000000,0.000000}%
\pgfsetstrokecolor{currentstroke}%
\pgfsetdash{}{0pt}%
\pgfpathmoveto{\pgfqpoint{3.393750in}{2.173495in}}%
\pgfpathlineto{\pgfqpoint{3.490000in}{2.173495in}}%
\pgfpathlineto{\pgfqpoint{3.490000in}{2.181015in}}%
\pgfpathlineto{\pgfqpoint{3.393750in}{2.181015in}}%
\pgfpathlineto{\pgfqpoint{3.393750in}{2.173495in}}%
\pgfusepath{fill}%
\end{pgfscope}%
\begin{pgfscope}%
\pgfpathrectangle{\pgfqpoint{3.393750in}{0.699667in}}{\pgfqpoint{0.096250in}{1.925000in}}%
\pgfusepath{clip}%
\pgfsetbuttcap%
\pgfsetroundjoin%
\definecolor{currentfill}{rgb}{0.990871,0.572706,0.398714}%
\pgfsetfillcolor{currentfill}%
\pgfsetfillopacity{0.800000}%
\pgfsetlinewidth{0.000000pt}%
\definecolor{currentstroke}{rgb}{0.000000,0.000000,0.000000}%
\pgfsetstrokecolor{currentstroke}%
\pgfsetdash{}{0pt}%
\pgfpathmoveto{\pgfqpoint{3.393750in}{2.181015in}}%
\pgfpathlineto{\pgfqpoint{3.490000in}{2.181015in}}%
\pgfpathlineto{\pgfqpoint{3.490000in}{2.188534in}}%
\pgfpathlineto{\pgfqpoint{3.393750in}{2.188534in}}%
\pgfpathlineto{\pgfqpoint{3.393750in}{2.181015in}}%
\pgfusepath{fill}%
\end{pgfscope}%
\begin{pgfscope}%
\pgfpathrectangle{\pgfqpoint{3.393750in}{0.699667in}}{\pgfqpoint{0.096250in}{1.925000in}}%
\pgfusepath{clip}%
\pgfsetbuttcap%
\pgfsetroundjoin%
\definecolor{currentfill}{rgb}{0.991558,0.580107,0.402441}%
\pgfsetfillcolor{currentfill}%
\pgfsetfillopacity{0.800000}%
\pgfsetlinewidth{0.000000pt}%
\definecolor{currentstroke}{rgb}{0.000000,0.000000,0.000000}%
\pgfsetstrokecolor{currentstroke}%
\pgfsetdash{}{0pt}%
\pgfpathmoveto{\pgfqpoint{3.393750in}{2.188534in}}%
\pgfpathlineto{\pgfqpoint{3.490000in}{2.188534in}}%
\pgfpathlineto{\pgfqpoint{3.490000in}{2.196054in}}%
\pgfpathlineto{\pgfqpoint{3.393750in}{2.196054in}}%
\pgfpathlineto{\pgfqpoint{3.393750in}{2.188534in}}%
\pgfusepath{fill}%
\end{pgfscope}%
\begin{pgfscope}%
\pgfpathrectangle{\pgfqpoint{3.393750in}{0.699667in}}{\pgfqpoint{0.096250in}{1.925000in}}%
\pgfusepath{clip}%
\pgfsetbuttcap%
\pgfsetroundjoin%
\definecolor{currentfill}{rgb}{0.992196,0.587502,0.406299}%
\pgfsetfillcolor{currentfill}%
\pgfsetfillopacity{0.800000}%
\pgfsetlinewidth{0.000000pt}%
\definecolor{currentstroke}{rgb}{0.000000,0.000000,0.000000}%
\pgfsetstrokecolor{currentstroke}%
\pgfsetdash{}{0pt}%
\pgfpathmoveto{\pgfqpoint{3.393750in}{2.196054in}}%
\pgfpathlineto{\pgfqpoint{3.490000in}{2.196054in}}%
\pgfpathlineto{\pgfqpoint{3.490000in}{2.203573in}}%
\pgfpathlineto{\pgfqpoint{3.393750in}{2.203573in}}%
\pgfpathlineto{\pgfqpoint{3.393750in}{2.196054in}}%
\pgfusepath{fill}%
\end{pgfscope}%
\begin{pgfscope}%
\pgfpathrectangle{\pgfqpoint{3.393750in}{0.699667in}}{\pgfqpoint{0.096250in}{1.925000in}}%
\pgfusepath{clip}%
\pgfsetbuttcap%
\pgfsetroundjoin%
\definecolor{currentfill}{rgb}{0.992785,0.594891,0.410283}%
\pgfsetfillcolor{currentfill}%
\pgfsetfillopacity{0.800000}%
\pgfsetlinewidth{0.000000pt}%
\definecolor{currentstroke}{rgb}{0.000000,0.000000,0.000000}%
\pgfsetstrokecolor{currentstroke}%
\pgfsetdash{}{0pt}%
\pgfpathmoveto{\pgfqpoint{3.393750in}{2.203573in}}%
\pgfpathlineto{\pgfqpoint{3.490000in}{2.203573in}}%
\pgfpathlineto{\pgfqpoint{3.490000in}{2.211093in}}%
\pgfpathlineto{\pgfqpoint{3.393750in}{2.211093in}}%
\pgfpathlineto{\pgfqpoint{3.393750in}{2.203573in}}%
\pgfusepath{fill}%
\end{pgfscope}%
\begin{pgfscope}%
\pgfpathrectangle{\pgfqpoint{3.393750in}{0.699667in}}{\pgfqpoint{0.096250in}{1.925000in}}%
\pgfusepath{clip}%
\pgfsetbuttcap%
\pgfsetroundjoin%
\definecolor{currentfill}{rgb}{0.993326,0.602275,0.414390}%
\pgfsetfillcolor{currentfill}%
\pgfsetfillopacity{0.800000}%
\pgfsetlinewidth{0.000000pt}%
\definecolor{currentstroke}{rgb}{0.000000,0.000000,0.000000}%
\pgfsetstrokecolor{currentstroke}%
\pgfsetdash{}{0pt}%
\pgfpathmoveto{\pgfqpoint{3.393750in}{2.211093in}}%
\pgfpathlineto{\pgfqpoint{3.490000in}{2.211093in}}%
\pgfpathlineto{\pgfqpoint{3.490000in}{2.218612in}}%
\pgfpathlineto{\pgfqpoint{3.393750in}{2.218612in}}%
\pgfpathlineto{\pgfqpoint{3.393750in}{2.211093in}}%
\pgfusepath{fill}%
\end{pgfscope}%
\begin{pgfscope}%
\pgfpathrectangle{\pgfqpoint{3.393750in}{0.699667in}}{\pgfqpoint{0.096250in}{1.925000in}}%
\pgfusepath{clip}%
\pgfsetbuttcap%
\pgfsetroundjoin%
\definecolor{currentfill}{rgb}{0.993834,0.609644,0.418613}%
\pgfsetfillcolor{currentfill}%
\pgfsetfillopacity{0.800000}%
\pgfsetlinewidth{0.000000pt}%
\definecolor{currentstroke}{rgb}{0.000000,0.000000,0.000000}%
\pgfsetstrokecolor{currentstroke}%
\pgfsetdash{}{0pt}%
\pgfpathmoveto{\pgfqpoint{3.393750in}{2.218612in}}%
\pgfpathlineto{\pgfqpoint{3.490000in}{2.218612in}}%
\pgfpathlineto{\pgfqpoint{3.490000in}{2.226132in}}%
\pgfpathlineto{\pgfqpoint{3.393750in}{2.226132in}}%
\pgfpathlineto{\pgfqpoint{3.393750in}{2.218612in}}%
\pgfusepath{fill}%
\end{pgfscope}%
\begin{pgfscope}%
\pgfpathrectangle{\pgfqpoint{3.393750in}{0.699667in}}{\pgfqpoint{0.096250in}{1.925000in}}%
\pgfusepath{clip}%
\pgfsetbuttcap%
\pgfsetroundjoin%
\definecolor{currentfill}{rgb}{0.994309,0.616999,0.422950}%
\pgfsetfillcolor{currentfill}%
\pgfsetfillopacity{0.800000}%
\pgfsetlinewidth{0.000000pt}%
\definecolor{currentstroke}{rgb}{0.000000,0.000000,0.000000}%
\pgfsetstrokecolor{currentstroke}%
\pgfsetdash{}{0pt}%
\pgfpathmoveto{\pgfqpoint{3.393750in}{2.226132in}}%
\pgfpathlineto{\pgfqpoint{3.490000in}{2.226132in}}%
\pgfpathlineto{\pgfqpoint{3.490000in}{2.233652in}}%
\pgfpathlineto{\pgfqpoint{3.393750in}{2.233652in}}%
\pgfpathlineto{\pgfqpoint{3.393750in}{2.226132in}}%
\pgfusepath{fill}%
\end{pgfscope}%
\begin{pgfscope}%
\pgfpathrectangle{\pgfqpoint{3.393750in}{0.699667in}}{\pgfqpoint{0.096250in}{1.925000in}}%
\pgfusepath{clip}%
\pgfsetbuttcap%
\pgfsetroundjoin%
\definecolor{currentfill}{rgb}{0.994738,0.624350,0.427397}%
\pgfsetfillcolor{currentfill}%
\pgfsetfillopacity{0.800000}%
\pgfsetlinewidth{0.000000pt}%
\definecolor{currentstroke}{rgb}{0.000000,0.000000,0.000000}%
\pgfsetstrokecolor{currentstroke}%
\pgfsetdash{}{0pt}%
\pgfpathmoveto{\pgfqpoint{3.393750in}{2.233652in}}%
\pgfpathlineto{\pgfqpoint{3.490000in}{2.233652in}}%
\pgfpathlineto{\pgfqpoint{3.490000in}{2.241171in}}%
\pgfpathlineto{\pgfqpoint{3.393750in}{2.241171in}}%
\pgfpathlineto{\pgfqpoint{3.393750in}{2.233652in}}%
\pgfusepath{fill}%
\end{pgfscope}%
\begin{pgfscope}%
\pgfpathrectangle{\pgfqpoint{3.393750in}{0.699667in}}{\pgfqpoint{0.096250in}{1.925000in}}%
\pgfusepath{clip}%
\pgfsetbuttcap%
\pgfsetroundjoin%
\definecolor{currentfill}{rgb}{0.995122,0.631696,0.431951}%
\pgfsetfillcolor{currentfill}%
\pgfsetfillopacity{0.800000}%
\pgfsetlinewidth{0.000000pt}%
\definecolor{currentstroke}{rgb}{0.000000,0.000000,0.000000}%
\pgfsetstrokecolor{currentstroke}%
\pgfsetdash{}{0pt}%
\pgfpathmoveto{\pgfqpoint{3.393750in}{2.241171in}}%
\pgfpathlineto{\pgfqpoint{3.490000in}{2.241171in}}%
\pgfpathlineto{\pgfqpoint{3.490000in}{2.248691in}}%
\pgfpathlineto{\pgfqpoint{3.393750in}{2.248691in}}%
\pgfpathlineto{\pgfqpoint{3.393750in}{2.241171in}}%
\pgfusepath{fill}%
\end{pgfscope}%
\begin{pgfscope}%
\pgfpathrectangle{\pgfqpoint{3.393750in}{0.699667in}}{\pgfqpoint{0.096250in}{1.925000in}}%
\pgfusepath{clip}%
\pgfsetbuttcap%
\pgfsetroundjoin%
\definecolor{currentfill}{rgb}{0.995480,0.639027,0.436607}%
\pgfsetfillcolor{currentfill}%
\pgfsetfillopacity{0.800000}%
\pgfsetlinewidth{0.000000pt}%
\definecolor{currentstroke}{rgb}{0.000000,0.000000,0.000000}%
\pgfsetstrokecolor{currentstroke}%
\pgfsetdash{}{0pt}%
\pgfpathmoveto{\pgfqpoint{3.393750in}{2.248691in}}%
\pgfpathlineto{\pgfqpoint{3.490000in}{2.248691in}}%
\pgfpathlineto{\pgfqpoint{3.490000in}{2.256210in}}%
\pgfpathlineto{\pgfqpoint{3.393750in}{2.256210in}}%
\pgfpathlineto{\pgfqpoint{3.393750in}{2.248691in}}%
\pgfusepath{fill}%
\end{pgfscope}%
\begin{pgfscope}%
\pgfpathrectangle{\pgfqpoint{3.393750in}{0.699667in}}{\pgfqpoint{0.096250in}{1.925000in}}%
\pgfusepath{clip}%
\pgfsetbuttcap%
\pgfsetroundjoin%
\definecolor{currentfill}{rgb}{0.995810,0.646344,0.441361}%
\pgfsetfillcolor{currentfill}%
\pgfsetfillopacity{0.800000}%
\pgfsetlinewidth{0.000000pt}%
\definecolor{currentstroke}{rgb}{0.000000,0.000000,0.000000}%
\pgfsetstrokecolor{currentstroke}%
\pgfsetdash{}{0pt}%
\pgfpathmoveto{\pgfqpoint{3.393750in}{2.256210in}}%
\pgfpathlineto{\pgfqpoint{3.490000in}{2.256210in}}%
\pgfpathlineto{\pgfqpoint{3.490000in}{2.263730in}}%
\pgfpathlineto{\pgfqpoint{3.393750in}{2.263730in}}%
\pgfpathlineto{\pgfqpoint{3.393750in}{2.256210in}}%
\pgfusepath{fill}%
\end{pgfscope}%
\begin{pgfscope}%
\pgfpathrectangle{\pgfqpoint{3.393750in}{0.699667in}}{\pgfqpoint{0.096250in}{1.925000in}}%
\pgfusepath{clip}%
\pgfsetbuttcap%
\pgfsetroundjoin%
\definecolor{currentfill}{rgb}{0.996096,0.653659,0.446213}%
\pgfsetfillcolor{currentfill}%
\pgfsetfillopacity{0.800000}%
\pgfsetlinewidth{0.000000pt}%
\definecolor{currentstroke}{rgb}{0.000000,0.000000,0.000000}%
\pgfsetstrokecolor{currentstroke}%
\pgfsetdash{}{0pt}%
\pgfpathmoveto{\pgfqpoint{3.393750in}{2.263730in}}%
\pgfpathlineto{\pgfqpoint{3.490000in}{2.263730in}}%
\pgfpathlineto{\pgfqpoint{3.490000in}{2.271249in}}%
\pgfpathlineto{\pgfqpoint{3.393750in}{2.271249in}}%
\pgfpathlineto{\pgfqpoint{3.393750in}{2.263730in}}%
\pgfusepath{fill}%
\end{pgfscope}%
\begin{pgfscope}%
\pgfpathrectangle{\pgfqpoint{3.393750in}{0.699667in}}{\pgfqpoint{0.096250in}{1.925000in}}%
\pgfusepath{clip}%
\pgfsetbuttcap%
\pgfsetroundjoin%
\definecolor{currentfill}{rgb}{0.996341,0.660969,0.451160}%
\pgfsetfillcolor{currentfill}%
\pgfsetfillopacity{0.800000}%
\pgfsetlinewidth{0.000000pt}%
\definecolor{currentstroke}{rgb}{0.000000,0.000000,0.000000}%
\pgfsetstrokecolor{currentstroke}%
\pgfsetdash{}{0pt}%
\pgfpathmoveto{\pgfqpoint{3.393750in}{2.271249in}}%
\pgfpathlineto{\pgfqpoint{3.490000in}{2.271249in}}%
\pgfpathlineto{\pgfqpoint{3.490000in}{2.278769in}}%
\pgfpathlineto{\pgfqpoint{3.393750in}{2.278769in}}%
\pgfpathlineto{\pgfqpoint{3.393750in}{2.271249in}}%
\pgfusepath{fill}%
\end{pgfscope}%
\begin{pgfscope}%
\pgfpathrectangle{\pgfqpoint{3.393750in}{0.699667in}}{\pgfqpoint{0.096250in}{1.925000in}}%
\pgfusepath{clip}%
\pgfsetbuttcap%
\pgfsetroundjoin%
\definecolor{currentfill}{rgb}{0.996580,0.668256,0.456192}%
\pgfsetfillcolor{currentfill}%
\pgfsetfillopacity{0.800000}%
\pgfsetlinewidth{0.000000pt}%
\definecolor{currentstroke}{rgb}{0.000000,0.000000,0.000000}%
\pgfsetstrokecolor{currentstroke}%
\pgfsetdash{}{0pt}%
\pgfpathmoveto{\pgfqpoint{3.393750in}{2.278769in}}%
\pgfpathlineto{\pgfqpoint{3.490000in}{2.278769in}}%
\pgfpathlineto{\pgfqpoint{3.490000in}{2.286288in}}%
\pgfpathlineto{\pgfqpoint{3.393750in}{2.286288in}}%
\pgfpathlineto{\pgfqpoint{3.393750in}{2.278769in}}%
\pgfusepath{fill}%
\end{pgfscope}%
\begin{pgfscope}%
\pgfpathrectangle{\pgfqpoint{3.393750in}{0.699667in}}{\pgfqpoint{0.096250in}{1.925000in}}%
\pgfusepath{clip}%
\pgfsetbuttcap%
\pgfsetroundjoin%
\definecolor{currentfill}{rgb}{0.996775,0.675541,0.461314}%
\pgfsetfillcolor{currentfill}%
\pgfsetfillopacity{0.800000}%
\pgfsetlinewidth{0.000000pt}%
\definecolor{currentstroke}{rgb}{0.000000,0.000000,0.000000}%
\pgfsetstrokecolor{currentstroke}%
\pgfsetdash{}{0pt}%
\pgfpathmoveto{\pgfqpoint{3.393750in}{2.286288in}}%
\pgfpathlineto{\pgfqpoint{3.490000in}{2.286288in}}%
\pgfpathlineto{\pgfqpoint{3.490000in}{2.293808in}}%
\pgfpathlineto{\pgfqpoint{3.393750in}{2.293808in}}%
\pgfpathlineto{\pgfqpoint{3.393750in}{2.286288in}}%
\pgfusepath{fill}%
\end{pgfscope}%
\begin{pgfscope}%
\pgfpathrectangle{\pgfqpoint{3.393750in}{0.699667in}}{\pgfqpoint{0.096250in}{1.925000in}}%
\pgfusepath{clip}%
\pgfsetbuttcap%
\pgfsetroundjoin%
\definecolor{currentfill}{rgb}{0.996925,0.682828,0.466526}%
\pgfsetfillcolor{currentfill}%
\pgfsetfillopacity{0.800000}%
\pgfsetlinewidth{0.000000pt}%
\definecolor{currentstroke}{rgb}{0.000000,0.000000,0.000000}%
\pgfsetstrokecolor{currentstroke}%
\pgfsetdash{}{0pt}%
\pgfpathmoveto{\pgfqpoint{3.393750in}{2.293808in}}%
\pgfpathlineto{\pgfqpoint{3.490000in}{2.293808in}}%
\pgfpathlineto{\pgfqpoint{3.490000in}{2.301327in}}%
\pgfpathlineto{\pgfqpoint{3.393750in}{2.301327in}}%
\pgfpathlineto{\pgfqpoint{3.393750in}{2.293808in}}%
\pgfusepath{fill}%
\end{pgfscope}%
\begin{pgfscope}%
\pgfpathrectangle{\pgfqpoint{3.393750in}{0.699667in}}{\pgfqpoint{0.096250in}{1.925000in}}%
\pgfusepath{clip}%
\pgfsetbuttcap%
\pgfsetroundjoin%
\definecolor{currentfill}{rgb}{0.997077,0.690088,0.471811}%
\pgfsetfillcolor{currentfill}%
\pgfsetfillopacity{0.800000}%
\pgfsetlinewidth{0.000000pt}%
\definecolor{currentstroke}{rgb}{0.000000,0.000000,0.000000}%
\pgfsetstrokecolor{currentstroke}%
\pgfsetdash{}{0pt}%
\pgfpathmoveto{\pgfqpoint{3.393750in}{2.301327in}}%
\pgfpathlineto{\pgfqpoint{3.490000in}{2.301327in}}%
\pgfpathlineto{\pgfqpoint{3.490000in}{2.308847in}}%
\pgfpathlineto{\pgfqpoint{3.393750in}{2.308847in}}%
\pgfpathlineto{\pgfqpoint{3.393750in}{2.301327in}}%
\pgfusepath{fill}%
\end{pgfscope}%
\begin{pgfscope}%
\pgfpathrectangle{\pgfqpoint{3.393750in}{0.699667in}}{\pgfqpoint{0.096250in}{1.925000in}}%
\pgfusepath{clip}%
\pgfsetbuttcap%
\pgfsetroundjoin%
\definecolor{currentfill}{rgb}{0.997186,0.697349,0.477182}%
\pgfsetfillcolor{currentfill}%
\pgfsetfillopacity{0.800000}%
\pgfsetlinewidth{0.000000pt}%
\definecolor{currentstroke}{rgb}{0.000000,0.000000,0.000000}%
\pgfsetstrokecolor{currentstroke}%
\pgfsetdash{}{0pt}%
\pgfpathmoveto{\pgfqpoint{3.393750in}{2.308847in}}%
\pgfpathlineto{\pgfqpoint{3.490000in}{2.308847in}}%
\pgfpathlineto{\pgfqpoint{3.490000in}{2.316366in}}%
\pgfpathlineto{\pgfqpoint{3.393750in}{2.316366in}}%
\pgfpathlineto{\pgfqpoint{3.393750in}{2.308847in}}%
\pgfusepath{fill}%
\end{pgfscope}%
\begin{pgfscope}%
\pgfpathrectangle{\pgfqpoint{3.393750in}{0.699667in}}{\pgfqpoint{0.096250in}{1.925000in}}%
\pgfusepath{clip}%
\pgfsetbuttcap%
\pgfsetroundjoin%
\definecolor{currentfill}{rgb}{0.997254,0.704611,0.482635}%
\pgfsetfillcolor{currentfill}%
\pgfsetfillopacity{0.800000}%
\pgfsetlinewidth{0.000000pt}%
\definecolor{currentstroke}{rgb}{0.000000,0.000000,0.000000}%
\pgfsetstrokecolor{currentstroke}%
\pgfsetdash{}{0pt}%
\pgfpathmoveto{\pgfqpoint{3.393750in}{2.316366in}}%
\pgfpathlineto{\pgfqpoint{3.490000in}{2.316366in}}%
\pgfpathlineto{\pgfqpoint{3.490000in}{2.323886in}}%
\pgfpathlineto{\pgfqpoint{3.393750in}{2.323886in}}%
\pgfpathlineto{\pgfqpoint{3.393750in}{2.316366in}}%
\pgfusepath{fill}%
\end{pgfscope}%
\begin{pgfscope}%
\pgfpathrectangle{\pgfqpoint{3.393750in}{0.699667in}}{\pgfqpoint{0.096250in}{1.925000in}}%
\pgfusepath{clip}%
\pgfsetbuttcap%
\pgfsetroundjoin%
\definecolor{currentfill}{rgb}{0.997325,0.711848,0.488154}%
\pgfsetfillcolor{currentfill}%
\pgfsetfillopacity{0.800000}%
\pgfsetlinewidth{0.000000pt}%
\definecolor{currentstroke}{rgb}{0.000000,0.000000,0.000000}%
\pgfsetstrokecolor{currentstroke}%
\pgfsetdash{}{0pt}%
\pgfpathmoveto{\pgfqpoint{3.393750in}{2.323886in}}%
\pgfpathlineto{\pgfqpoint{3.490000in}{2.323886in}}%
\pgfpathlineto{\pgfqpoint{3.490000in}{2.331405in}}%
\pgfpathlineto{\pgfqpoint{3.393750in}{2.331405in}}%
\pgfpathlineto{\pgfqpoint{3.393750in}{2.323886in}}%
\pgfusepath{fill}%
\end{pgfscope}%
\begin{pgfscope}%
\pgfpathrectangle{\pgfqpoint{3.393750in}{0.699667in}}{\pgfqpoint{0.096250in}{1.925000in}}%
\pgfusepath{clip}%
\pgfsetbuttcap%
\pgfsetroundjoin%
\definecolor{currentfill}{rgb}{0.997351,0.719089,0.493755}%
\pgfsetfillcolor{currentfill}%
\pgfsetfillopacity{0.800000}%
\pgfsetlinewidth{0.000000pt}%
\definecolor{currentstroke}{rgb}{0.000000,0.000000,0.000000}%
\pgfsetstrokecolor{currentstroke}%
\pgfsetdash{}{0pt}%
\pgfpathmoveto{\pgfqpoint{3.393750in}{2.331405in}}%
\pgfpathlineto{\pgfqpoint{3.490000in}{2.331405in}}%
\pgfpathlineto{\pgfqpoint{3.490000in}{2.338925in}}%
\pgfpathlineto{\pgfqpoint{3.393750in}{2.338925in}}%
\pgfpathlineto{\pgfqpoint{3.393750in}{2.331405in}}%
\pgfusepath{fill}%
\end{pgfscope}%
\begin{pgfscope}%
\pgfpathrectangle{\pgfqpoint{3.393750in}{0.699667in}}{\pgfqpoint{0.096250in}{1.925000in}}%
\pgfusepath{clip}%
\pgfsetbuttcap%
\pgfsetroundjoin%
\definecolor{currentfill}{rgb}{0.997351,0.726324,0.499428}%
\pgfsetfillcolor{currentfill}%
\pgfsetfillopacity{0.800000}%
\pgfsetlinewidth{0.000000pt}%
\definecolor{currentstroke}{rgb}{0.000000,0.000000,0.000000}%
\pgfsetstrokecolor{currentstroke}%
\pgfsetdash{}{0pt}%
\pgfpathmoveto{\pgfqpoint{3.393750in}{2.338925in}}%
\pgfpathlineto{\pgfqpoint{3.490000in}{2.338925in}}%
\pgfpathlineto{\pgfqpoint{3.490000in}{2.346445in}}%
\pgfpathlineto{\pgfqpoint{3.393750in}{2.346445in}}%
\pgfpathlineto{\pgfqpoint{3.393750in}{2.338925in}}%
\pgfusepath{fill}%
\end{pgfscope}%
\begin{pgfscope}%
\pgfpathrectangle{\pgfqpoint{3.393750in}{0.699667in}}{\pgfqpoint{0.096250in}{1.925000in}}%
\pgfusepath{clip}%
\pgfsetbuttcap%
\pgfsetroundjoin%
\definecolor{currentfill}{rgb}{0.997341,0.733545,0.505167}%
\pgfsetfillcolor{currentfill}%
\pgfsetfillopacity{0.800000}%
\pgfsetlinewidth{0.000000pt}%
\definecolor{currentstroke}{rgb}{0.000000,0.000000,0.000000}%
\pgfsetstrokecolor{currentstroke}%
\pgfsetdash{}{0pt}%
\pgfpathmoveto{\pgfqpoint{3.393750in}{2.346445in}}%
\pgfpathlineto{\pgfqpoint{3.490000in}{2.346445in}}%
\pgfpathlineto{\pgfqpoint{3.490000in}{2.353964in}}%
\pgfpathlineto{\pgfqpoint{3.393750in}{2.353964in}}%
\pgfpathlineto{\pgfqpoint{3.393750in}{2.346445in}}%
\pgfusepath{fill}%
\end{pgfscope}%
\begin{pgfscope}%
\pgfpathrectangle{\pgfqpoint{3.393750in}{0.699667in}}{\pgfqpoint{0.096250in}{1.925000in}}%
\pgfusepath{clip}%
\pgfsetbuttcap%
\pgfsetroundjoin%
\definecolor{currentfill}{rgb}{0.997285,0.740772,0.510983}%
\pgfsetfillcolor{currentfill}%
\pgfsetfillopacity{0.800000}%
\pgfsetlinewidth{0.000000pt}%
\definecolor{currentstroke}{rgb}{0.000000,0.000000,0.000000}%
\pgfsetstrokecolor{currentstroke}%
\pgfsetdash{}{0pt}%
\pgfpathmoveto{\pgfqpoint{3.393750in}{2.353964in}}%
\pgfpathlineto{\pgfqpoint{3.490000in}{2.353964in}}%
\pgfpathlineto{\pgfqpoint{3.490000in}{2.361484in}}%
\pgfpathlineto{\pgfqpoint{3.393750in}{2.361484in}}%
\pgfpathlineto{\pgfqpoint{3.393750in}{2.353964in}}%
\pgfusepath{fill}%
\end{pgfscope}%
\begin{pgfscope}%
\pgfpathrectangle{\pgfqpoint{3.393750in}{0.699667in}}{\pgfqpoint{0.096250in}{1.925000in}}%
\pgfusepath{clip}%
\pgfsetbuttcap%
\pgfsetroundjoin%
\definecolor{currentfill}{rgb}{0.997228,0.747981,0.516859}%
\pgfsetfillcolor{currentfill}%
\pgfsetfillopacity{0.800000}%
\pgfsetlinewidth{0.000000pt}%
\definecolor{currentstroke}{rgb}{0.000000,0.000000,0.000000}%
\pgfsetstrokecolor{currentstroke}%
\pgfsetdash{}{0pt}%
\pgfpathmoveto{\pgfqpoint{3.393750in}{2.361484in}}%
\pgfpathlineto{\pgfqpoint{3.490000in}{2.361484in}}%
\pgfpathlineto{\pgfqpoint{3.490000in}{2.369003in}}%
\pgfpathlineto{\pgfqpoint{3.393750in}{2.369003in}}%
\pgfpathlineto{\pgfqpoint{3.393750in}{2.361484in}}%
\pgfusepath{fill}%
\end{pgfscope}%
\begin{pgfscope}%
\pgfpathrectangle{\pgfqpoint{3.393750in}{0.699667in}}{\pgfqpoint{0.096250in}{1.925000in}}%
\pgfusepath{clip}%
\pgfsetbuttcap%
\pgfsetroundjoin%
\definecolor{currentfill}{rgb}{0.997138,0.755190,0.522806}%
\pgfsetfillcolor{currentfill}%
\pgfsetfillopacity{0.800000}%
\pgfsetlinewidth{0.000000pt}%
\definecolor{currentstroke}{rgb}{0.000000,0.000000,0.000000}%
\pgfsetstrokecolor{currentstroke}%
\pgfsetdash{}{0pt}%
\pgfpathmoveto{\pgfqpoint{3.393750in}{2.369003in}}%
\pgfpathlineto{\pgfqpoint{3.490000in}{2.369003in}}%
\pgfpathlineto{\pgfqpoint{3.490000in}{2.376523in}}%
\pgfpathlineto{\pgfqpoint{3.393750in}{2.376523in}}%
\pgfpathlineto{\pgfqpoint{3.393750in}{2.369003in}}%
\pgfusepath{fill}%
\end{pgfscope}%
\begin{pgfscope}%
\pgfpathrectangle{\pgfqpoint{3.393750in}{0.699667in}}{\pgfqpoint{0.096250in}{1.925000in}}%
\pgfusepath{clip}%
\pgfsetbuttcap%
\pgfsetroundjoin%
\definecolor{currentfill}{rgb}{0.997019,0.762398,0.528821}%
\pgfsetfillcolor{currentfill}%
\pgfsetfillopacity{0.800000}%
\pgfsetlinewidth{0.000000pt}%
\definecolor{currentstroke}{rgb}{0.000000,0.000000,0.000000}%
\pgfsetstrokecolor{currentstroke}%
\pgfsetdash{}{0pt}%
\pgfpathmoveto{\pgfqpoint{3.393750in}{2.376523in}}%
\pgfpathlineto{\pgfqpoint{3.490000in}{2.376523in}}%
\pgfpathlineto{\pgfqpoint{3.490000in}{2.384042in}}%
\pgfpathlineto{\pgfqpoint{3.393750in}{2.384042in}}%
\pgfpathlineto{\pgfqpoint{3.393750in}{2.376523in}}%
\pgfusepath{fill}%
\end{pgfscope}%
\begin{pgfscope}%
\pgfpathrectangle{\pgfqpoint{3.393750in}{0.699667in}}{\pgfqpoint{0.096250in}{1.925000in}}%
\pgfusepath{clip}%
\pgfsetbuttcap%
\pgfsetroundjoin%
\definecolor{currentfill}{rgb}{0.996898,0.769591,0.534892}%
\pgfsetfillcolor{currentfill}%
\pgfsetfillopacity{0.800000}%
\pgfsetlinewidth{0.000000pt}%
\definecolor{currentstroke}{rgb}{0.000000,0.000000,0.000000}%
\pgfsetstrokecolor{currentstroke}%
\pgfsetdash{}{0pt}%
\pgfpathmoveto{\pgfqpoint{3.393750in}{2.384042in}}%
\pgfpathlineto{\pgfqpoint{3.490000in}{2.384042in}}%
\pgfpathlineto{\pgfqpoint{3.490000in}{2.391562in}}%
\pgfpathlineto{\pgfqpoint{3.393750in}{2.391562in}}%
\pgfpathlineto{\pgfqpoint{3.393750in}{2.384042in}}%
\pgfusepath{fill}%
\end{pgfscope}%
\begin{pgfscope}%
\pgfpathrectangle{\pgfqpoint{3.393750in}{0.699667in}}{\pgfqpoint{0.096250in}{1.925000in}}%
\pgfusepath{clip}%
\pgfsetbuttcap%
\pgfsetroundjoin%
\definecolor{currentfill}{rgb}{0.996727,0.776795,0.541039}%
\pgfsetfillcolor{currentfill}%
\pgfsetfillopacity{0.800000}%
\pgfsetlinewidth{0.000000pt}%
\definecolor{currentstroke}{rgb}{0.000000,0.000000,0.000000}%
\pgfsetstrokecolor{currentstroke}%
\pgfsetdash{}{0pt}%
\pgfpathmoveto{\pgfqpoint{3.393750in}{2.391562in}}%
\pgfpathlineto{\pgfqpoint{3.490000in}{2.391562in}}%
\pgfpathlineto{\pgfqpoint{3.490000in}{2.399081in}}%
\pgfpathlineto{\pgfqpoint{3.393750in}{2.399081in}}%
\pgfpathlineto{\pgfqpoint{3.393750in}{2.391562in}}%
\pgfusepath{fill}%
\end{pgfscope}%
\begin{pgfscope}%
\pgfpathrectangle{\pgfqpoint{3.393750in}{0.699667in}}{\pgfqpoint{0.096250in}{1.925000in}}%
\pgfusepath{clip}%
\pgfsetbuttcap%
\pgfsetroundjoin%
\definecolor{currentfill}{rgb}{0.996571,0.783977,0.547233}%
\pgfsetfillcolor{currentfill}%
\pgfsetfillopacity{0.800000}%
\pgfsetlinewidth{0.000000pt}%
\definecolor{currentstroke}{rgb}{0.000000,0.000000,0.000000}%
\pgfsetstrokecolor{currentstroke}%
\pgfsetdash{}{0pt}%
\pgfpathmoveto{\pgfqpoint{3.393750in}{2.399081in}}%
\pgfpathlineto{\pgfqpoint{3.490000in}{2.399081in}}%
\pgfpathlineto{\pgfqpoint{3.490000in}{2.406601in}}%
\pgfpathlineto{\pgfqpoint{3.393750in}{2.406601in}}%
\pgfpathlineto{\pgfqpoint{3.393750in}{2.399081in}}%
\pgfusepath{fill}%
\end{pgfscope}%
\begin{pgfscope}%
\pgfpathrectangle{\pgfqpoint{3.393750in}{0.699667in}}{\pgfqpoint{0.096250in}{1.925000in}}%
\pgfusepath{clip}%
\pgfsetbuttcap%
\pgfsetroundjoin%
\definecolor{currentfill}{rgb}{0.996369,0.791167,0.553499}%
\pgfsetfillcolor{currentfill}%
\pgfsetfillopacity{0.800000}%
\pgfsetlinewidth{0.000000pt}%
\definecolor{currentstroke}{rgb}{0.000000,0.000000,0.000000}%
\pgfsetstrokecolor{currentstroke}%
\pgfsetdash{}{0pt}%
\pgfpathmoveto{\pgfqpoint{3.393750in}{2.406601in}}%
\pgfpathlineto{\pgfqpoint{3.490000in}{2.406601in}}%
\pgfpathlineto{\pgfqpoint{3.490000in}{2.414120in}}%
\pgfpathlineto{\pgfqpoint{3.393750in}{2.414120in}}%
\pgfpathlineto{\pgfqpoint{3.393750in}{2.406601in}}%
\pgfusepath{fill}%
\end{pgfscope}%
\begin{pgfscope}%
\pgfpathrectangle{\pgfqpoint{3.393750in}{0.699667in}}{\pgfqpoint{0.096250in}{1.925000in}}%
\pgfusepath{clip}%
\pgfsetbuttcap%
\pgfsetroundjoin%
\definecolor{currentfill}{rgb}{0.996162,0.798348,0.559820}%
\pgfsetfillcolor{currentfill}%
\pgfsetfillopacity{0.800000}%
\pgfsetlinewidth{0.000000pt}%
\definecolor{currentstroke}{rgb}{0.000000,0.000000,0.000000}%
\pgfsetstrokecolor{currentstroke}%
\pgfsetdash{}{0pt}%
\pgfpathmoveto{\pgfqpoint{3.393750in}{2.414120in}}%
\pgfpathlineto{\pgfqpoint{3.490000in}{2.414120in}}%
\pgfpathlineto{\pgfqpoint{3.490000in}{2.421640in}}%
\pgfpathlineto{\pgfqpoint{3.393750in}{2.421640in}}%
\pgfpathlineto{\pgfqpoint{3.393750in}{2.414120in}}%
\pgfusepath{fill}%
\end{pgfscope}%
\begin{pgfscope}%
\pgfpathrectangle{\pgfqpoint{3.393750in}{0.699667in}}{\pgfqpoint{0.096250in}{1.925000in}}%
\pgfusepath{clip}%
\pgfsetbuttcap%
\pgfsetroundjoin%
\definecolor{currentfill}{rgb}{0.995932,0.805527,0.566202}%
\pgfsetfillcolor{currentfill}%
\pgfsetfillopacity{0.800000}%
\pgfsetlinewidth{0.000000pt}%
\definecolor{currentstroke}{rgb}{0.000000,0.000000,0.000000}%
\pgfsetstrokecolor{currentstroke}%
\pgfsetdash{}{0pt}%
\pgfpathmoveto{\pgfqpoint{3.393750in}{2.421640in}}%
\pgfpathlineto{\pgfqpoint{3.490000in}{2.421640in}}%
\pgfpathlineto{\pgfqpoint{3.490000in}{2.429159in}}%
\pgfpathlineto{\pgfqpoint{3.393750in}{2.429159in}}%
\pgfpathlineto{\pgfqpoint{3.393750in}{2.421640in}}%
\pgfusepath{fill}%
\end{pgfscope}%
\begin{pgfscope}%
\pgfpathrectangle{\pgfqpoint{3.393750in}{0.699667in}}{\pgfqpoint{0.096250in}{1.925000in}}%
\pgfusepath{clip}%
\pgfsetbuttcap%
\pgfsetroundjoin%
\definecolor{currentfill}{rgb}{0.995680,0.812706,0.572645}%
\pgfsetfillcolor{currentfill}%
\pgfsetfillopacity{0.800000}%
\pgfsetlinewidth{0.000000pt}%
\definecolor{currentstroke}{rgb}{0.000000,0.000000,0.000000}%
\pgfsetstrokecolor{currentstroke}%
\pgfsetdash{}{0pt}%
\pgfpathmoveto{\pgfqpoint{3.393750in}{2.429159in}}%
\pgfpathlineto{\pgfqpoint{3.490000in}{2.429159in}}%
\pgfpathlineto{\pgfqpoint{3.490000in}{2.436679in}}%
\pgfpathlineto{\pgfqpoint{3.393750in}{2.436679in}}%
\pgfpathlineto{\pgfqpoint{3.393750in}{2.429159in}}%
\pgfusepath{fill}%
\end{pgfscope}%
\begin{pgfscope}%
\pgfpathrectangle{\pgfqpoint{3.393750in}{0.699667in}}{\pgfqpoint{0.096250in}{1.925000in}}%
\pgfusepath{clip}%
\pgfsetbuttcap%
\pgfsetroundjoin%
\definecolor{currentfill}{rgb}{0.995424,0.819875,0.579140}%
\pgfsetfillcolor{currentfill}%
\pgfsetfillopacity{0.800000}%
\pgfsetlinewidth{0.000000pt}%
\definecolor{currentstroke}{rgb}{0.000000,0.000000,0.000000}%
\pgfsetstrokecolor{currentstroke}%
\pgfsetdash{}{0pt}%
\pgfpathmoveto{\pgfqpoint{3.393750in}{2.436679in}}%
\pgfpathlineto{\pgfqpoint{3.490000in}{2.436679in}}%
\pgfpathlineto{\pgfqpoint{3.490000in}{2.444198in}}%
\pgfpathlineto{\pgfqpoint{3.393750in}{2.444198in}}%
\pgfpathlineto{\pgfqpoint{3.393750in}{2.436679in}}%
\pgfusepath{fill}%
\end{pgfscope}%
\begin{pgfscope}%
\pgfpathrectangle{\pgfqpoint{3.393750in}{0.699667in}}{\pgfqpoint{0.096250in}{1.925000in}}%
\pgfusepath{clip}%
\pgfsetbuttcap%
\pgfsetroundjoin%
\definecolor{currentfill}{rgb}{0.995131,0.827052,0.585701}%
\pgfsetfillcolor{currentfill}%
\pgfsetfillopacity{0.800000}%
\pgfsetlinewidth{0.000000pt}%
\definecolor{currentstroke}{rgb}{0.000000,0.000000,0.000000}%
\pgfsetstrokecolor{currentstroke}%
\pgfsetdash{}{0pt}%
\pgfpathmoveto{\pgfqpoint{3.393750in}{2.444198in}}%
\pgfpathlineto{\pgfqpoint{3.490000in}{2.444198in}}%
\pgfpathlineto{\pgfqpoint{3.490000in}{2.451718in}}%
\pgfpathlineto{\pgfqpoint{3.393750in}{2.451718in}}%
\pgfpathlineto{\pgfqpoint{3.393750in}{2.444198in}}%
\pgfusepath{fill}%
\end{pgfscope}%
\begin{pgfscope}%
\pgfpathrectangle{\pgfqpoint{3.393750in}{0.699667in}}{\pgfqpoint{0.096250in}{1.925000in}}%
\pgfusepath{clip}%
\pgfsetbuttcap%
\pgfsetroundjoin%
\definecolor{currentfill}{rgb}{0.994851,0.834213,0.592307}%
\pgfsetfillcolor{currentfill}%
\pgfsetfillopacity{0.800000}%
\pgfsetlinewidth{0.000000pt}%
\definecolor{currentstroke}{rgb}{0.000000,0.000000,0.000000}%
\pgfsetstrokecolor{currentstroke}%
\pgfsetdash{}{0pt}%
\pgfpathmoveto{\pgfqpoint{3.393750in}{2.451718in}}%
\pgfpathlineto{\pgfqpoint{3.490000in}{2.451718in}}%
\pgfpathlineto{\pgfqpoint{3.490000in}{2.459237in}}%
\pgfpathlineto{\pgfqpoint{3.393750in}{2.459237in}}%
\pgfpathlineto{\pgfqpoint{3.393750in}{2.451718in}}%
\pgfusepath{fill}%
\end{pgfscope}%
\begin{pgfscope}%
\pgfpathrectangle{\pgfqpoint{3.393750in}{0.699667in}}{\pgfqpoint{0.096250in}{1.925000in}}%
\pgfusepath{clip}%
\pgfsetbuttcap%
\pgfsetroundjoin%
\definecolor{currentfill}{rgb}{0.994524,0.841387,0.598983}%
\pgfsetfillcolor{currentfill}%
\pgfsetfillopacity{0.800000}%
\pgfsetlinewidth{0.000000pt}%
\definecolor{currentstroke}{rgb}{0.000000,0.000000,0.000000}%
\pgfsetstrokecolor{currentstroke}%
\pgfsetdash{}{0pt}%
\pgfpathmoveto{\pgfqpoint{3.393750in}{2.459237in}}%
\pgfpathlineto{\pgfqpoint{3.490000in}{2.459237in}}%
\pgfpathlineto{\pgfqpoint{3.490000in}{2.466757in}}%
\pgfpathlineto{\pgfqpoint{3.393750in}{2.466757in}}%
\pgfpathlineto{\pgfqpoint{3.393750in}{2.459237in}}%
\pgfusepath{fill}%
\end{pgfscope}%
\begin{pgfscope}%
\pgfpathrectangle{\pgfqpoint{3.393750in}{0.699667in}}{\pgfqpoint{0.096250in}{1.925000in}}%
\pgfusepath{clip}%
\pgfsetbuttcap%
\pgfsetroundjoin%
\definecolor{currentfill}{rgb}{0.994222,0.848540,0.605696}%
\pgfsetfillcolor{currentfill}%
\pgfsetfillopacity{0.800000}%
\pgfsetlinewidth{0.000000pt}%
\definecolor{currentstroke}{rgb}{0.000000,0.000000,0.000000}%
\pgfsetstrokecolor{currentstroke}%
\pgfsetdash{}{0pt}%
\pgfpathmoveto{\pgfqpoint{3.393750in}{2.466757in}}%
\pgfpathlineto{\pgfqpoint{3.490000in}{2.466757in}}%
\pgfpathlineto{\pgfqpoint{3.490000in}{2.474277in}}%
\pgfpathlineto{\pgfqpoint{3.393750in}{2.474277in}}%
\pgfpathlineto{\pgfqpoint{3.393750in}{2.466757in}}%
\pgfusepath{fill}%
\end{pgfscope}%
\begin{pgfscope}%
\pgfpathrectangle{\pgfqpoint{3.393750in}{0.699667in}}{\pgfqpoint{0.096250in}{1.925000in}}%
\pgfusepath{clip}%
\pgfsetbuttcap%
\pgfsetroundjoin%
\definecolor{currentfill}{rgb}{0.993866,0.855711,0.612482}%
\pgfsetfillcolor{currentfill}%
\pgfsetfillopacity{0.800000}%
\pgfsetlinewidth{0.000000pt}%
\definecolor{currentstroke}{rgb}{0.000000,0.000000,0.000000}%
\pgfsetstrokecolor{currentstroke}%
\pgfsetdash{}{0pt}%
\pgfpathmoveto{\pgfqpoint{3.393750in}{2.474277in}}%
\pgfpathlineto{\pgfqpoint{3.490000in}{2.474277in}}%
\pgfpathlineto{\pgfqpoint{3.490000in}{2.481796in}}%
\pgfpathlineto{\pgfqpoint{3.393750in}{2.481796in}}%
\pgfpathlineto{\pgfqpoint{3.393750in}{2.474277in}}%
\pgfusepath{fill}%
\end{pgfscope}%
\begin{pgfscope}%
\pgfpathrectangle{\pgfqpoint{3.393750in}{0.699667in}}{\pgfqpoint{0.096250in}{1.925000in}}%
\pgfusepath{clip}%
\pgfsetbuttcap%
\pgfsetroundjoin%
\definecolor{currentfill}{rgb}{0.993545,0.862859,0.619299}%
\pgfsetfillcolor{currentfill}%
\pgfsetfillopacity{0.800000}%
\pgfsetlinewidth{0.000000pt}%
\definecolor{currentstroke}{rgb}{0.000000,0.000000,0.000000}%
\pgfsetstrokecolor{currentstroke}%
\pgfsetdash{}{0pt}%
\pgfpathmoveto{\pgfqpoint{3.393750in}{2.481796in}}%
\pgfpathlineto{\pgfqpoint{3.490000in}{2.481796in}}%
\pgfpathlineto{\pgfqpoint{3.490000in}{2.489316in}}%
\pgfpathlineto{\pgfqpoint{3.393750in}{2.489316in}}%
\pgfpathlineto{\pgfqpoint{3.393750in}{2.481796in}}%
\pgfusepath{fill}%
\end{pgfscope}%
\begin{pgfscope}%
\pgfpathrectangle{\pgfqpoint{3.393750in}{0.699667in}}{\pgfqpoint{0.096250in}{1.925000in}}%
\pgfusepath{clip}%
\pgfsetbuttcap%
\pgfsetroundjoin%
\definecolor{currentfill}{rgb}{0.993170,0.870024,0.626189}%
\pgfsetfillcolor{currentfill}%
\pgfsetfillopacity{0.800000}%
\pgfsetlinewidth{0.000000pt}%
\definecolor{currentstroke}{rgb}{0.000000,0.000000,0.000000}%
\pgfsetstrokecolor{currentstroke}%
\pgfsetdash{}{0pt}%
\pgfpathmoveto{\pgfqpoint{3.393750in}{2.489316in}}%
\pgfpathlineto{\pgfqpoint{3.490000in}{2.489316in}}%
\pgfpathlineto{\pgfqpoint{3.490000in}{2.496835in}}%
\pgfpathlineto{\pgfqpoint{3.393750in}{2.496835in}}%
\pgfpathlineto{\pgfqpoint{3.393750in}{2.489316in}}%
\pgfusepath{fill}%
\end{pgfscope}%
\begin{pgfscope}%
\pgfpathrectangle{\pgfqpoint{3.393750in}{0.699667in}}{\pgfqpoint{0.096250in}{1.925000in}}%
\pgfusepath{clip}%
\pgfsetbuttcap%
\pgfsetroundjoin%
\definecolor{currentfill}{rgb}{0.992831,0.877168,0.633109}%
\pgfsetfillcolor{currentfill}%
\pgfsetfillopacity{0.800000}%
\pgfsetlinewidth{0.000000pt}%
\definecolor{currentstroke}{rgb}{0.000000,0.000000,0.000000}%
\pgfsetstrokecolor{currentstroke}%
\pgfsetdash{}{0pt}%
\pgfpathmoveto{\pgfqpoint{3.393750in}{2.496835in}}%
\pgfpathlineto{\pgfqpoint{3.490000in}{2.496835in}}%
\pgfpathlineto{\pgfqpoint{3.490000in}{2.504355in}}%
\pgfpathlineto{\pgfqpoint{3.393750in}{2.504355in}}%
\pgfpathlineto{\pgfqpoint{3.393750in}{2.496835in}}%
\pgfusepath{fill}%
\end{pgfscope}%
\begin{pgfscope}%
\pgfpathrectangle{\pgfqpoint{3.393750in}{0.699667in}}{\pgfqpoint{0.096250in}{1.925000in}}%
\pgfusepath{clip}%
\pgfsetbuttcap%
\pgfsetroundjoin%
\definecolor{currentfill}{rgb}{0.992440,0.884330,0.640099}%
\pgfsetfillcolor{currentfill}%
\pgfsetfillopacity{0.800000}%
\pgfsetlinewidth{0.000000pt}%
\definecolor{currentstroke}{rgb}{0.000000,0.000000,0.000000}%
\pgfsetstrokecolor{currentstroke}%
\pgfsetdash{}{0pt}%
\pgfpathmoveto{\pgfqpoint{3.393750in}{2.504355in}}%
\pgfpathlineto{\pgfqpoint{3.490000in}{2.504355in}}%
\pgfpathlineto{\pgfqpoint{3.490000in}{2.511874in}}%
\pgfpathlineto{\pgfqpoint{3.393750in}{2.511874in}}%
\pgfpathlineto{\pgfqpoint{3.393750in}{2.504355in}}%
\pgfusepath{fill}%
\end{pgfscope}%
\begin{pgfscope}%
\pgfpathrectangle{\pgfqpoint{3.393750in}{0.699667in}}{\pgfqpoint{0.096250in}{1.925000in}}%
\pgfusepath{clip}%
\pgfsetbuttcap%
\pgfsetroundjoin%
\definecolor{currentfill}{rgb}{0.992089,0.891470,0.647116}%
\pgfsetfillcolor{currentfill}%
\pgfsetfillopacity{0.800000}%
\pgfsetlinewidth{0.000000pt}%
\definecolor{currentstroke}{rgb}{0.000000,0.000000,0.000000}%
\pgfsetstrokecolor{currentstroke}%
\pgfsetdash{}{0pt}%
\pgfpathmoveto{\pgfqpoint{3.393750in}{2.511874in}}%
\pgfpathlineto{\pgfqpoint{3.490000in}{2.511874in}}%
\pgfpathlineto{\pgfqpoint{3.490000in}{2.519394in}}%
\pgfpathlineto{\pgfqpoint{3.393750in}{2.519394in}}%
\pgfpathlineto{\pgfqpoint{3.393750in}{2.511874in}}%
\pgfusepath{fill}%
\end{pgfscope}%
\begin{pgfscope}%
\pgfpathrectangle{\pgfqpoint{3.393750in}{0.699667in}}{\pgfqpoint{0.096250in}{1.925000in}}%
\pgfusepath{clip}%
\pgfsetbuttcap%
\pgfsetroundjoin%
\definecolor{currentfill}{rgb}{0.991688,0.898627,0.654202}%
\pgfsetfillcolor{currentfill}%
\pgfsetfillopacity{0.800000}%
\pgfsetlinewidth{0.000000pt}%
\definecolor{currentstroke}{rgb}{0.000000,0.000000,0.000000}%
\pgfsetstrokecolor{currentstroke}%
\pgfsetdash{}{0pt}%
\pgfpathmoveto{\pgfqpoint{3.393750in}{2.519394in}}%
\pgfpathlineto{\pgfqpoint{3.490000in}{2.519394in}}%
\pgfpathlineto{\pgfqpoint{3.490000in}{2.526913in}}%
\pgfpathlineto{\pgfqpoint{3.393750in}{2.526913in}}%
\pgfpathlineto{\pgfqpoint{3.393750in}{2.519394in}}%
\pgfusepath{fill}%
\end{pgfscope}%
\begin{pgfscope}%
\pgfpathrectangle{\pgfqpoint{3.393750in}{0.699667in}}{\pgfqpoint{0.096250in}{1.925000in}}%
\pgfusepath{clip}%
\pgfsetbuttcap%
\pgfsetroundjoin%
\definecolor{currentfill}{rgb}{0.991332,0.905763,0.661309}%
\pgfsetfillcolor{currentfill}%
\pgfsetfillopacity{0.800000}%
\pgfsetlinewidth{0.000000pt}%
\definecolor{currentstroke}{rgb}{0.000000,0.000000,0.000000}%
\pgfsetstrokecolor{currentstroke}%
\pgfsetdash{}{0pt}%
\pgfpathmoveto{\pgfqpoint{3.393750in}{2.526913in}}%
\pgfpathlineto{\pgfqpoint{3.490000in}{2.526913in}}%
\pgfpathlineto{\pgfqpoint{3.490000in}{2.534433in}}%
\pgfpathlineto{\pgfqpoint{3.393750in}{2.534433in}}%
\pgfpathlineto{\pgfqpoint{3.393750in}{2.526913in}}%
\pgfusepath{fill}%
\end{pgfscope}%
\begin{pgfscope}%
\pgfpathrectangle{\pgfqpoint{3.393750in}{0.699667in}}{\pgfqpoint{0.096250in}{1.925000in}}%
\pgfusepath{clip}%
\pgfsetbuttcap%
\pgfsetroundjoin%
\definecolor{currentfill}{rgb}{0.990930,0.912915,0.668481}%
\pgfsetfillcolor{currentfill}%
\pgfsetfillopacity{0.800000}%
\pgfsetlinewidth{0.000000pt}%
\definecolor{currentstroke}{rgb}{0.000000,0.000000,0.000000}%
\pgfsetstrokecolor{currentstroke}%
\pgfsetdash{}{0pt}%
\pgfpathmoveto{\pgfqpoint{3.393750in}{2.534433in}}%
\pgfpathlineto{\pgfqpoint{3.490000in}{2.534433in}}%
\pgfpathlineto{\pgfqpoint{3.490000in}{2.541952in}}%
\pgfpathlineto{\pgfqpoint{3.393750in}{2.541952in}}%
\pgfpathlineto{\pgfqpoint{3.393750in}{2.534433in}}%
\pgfusepath{fill}%
\end{pgfscope}%
\begin{pgfscope}%
\pgfpathrectangle{\pgfqpoint{3.393750in}{0.699667in}}{\pgfqpoint{0.096250in}{1.925000in}}%
\pgfusepath{clip}%
\pgfsetbuttcap%
\pgfsetroundjoin%
\definecolor{currentfill}{rgb}{0.990570,0.920049,0.675675}%
\pgfsetfillcolor{currentfill}%
\pgfsetfillopacity{0.800000}%
\pgfsetlinewidth{0.000000pt}%
\definecolor{currentstroke}{rgb}{0.000000,0.000000,0.000000}%
\pgfsetstrokecolor{currentstroke}%
\pgfsetdash{}{0pt}%
\pgfpathmoveto{\pgfqpoint{3.393750in}{2.541952in}}%
\pgfpathlineto{\pgfqpoint{3.490000in}{2.541952in}}%
\pgfpathlineto{\pgfqpoint{3.490000in}{2.549472in}}%
\pgfpathlineto{\pgfqpoint{3.393750in}{2.549472in}}%
\pgfpathlineto{\pgfqpoint{3.393750in}{2.541952in}}%
\pgfusepath{fill}%
\end{pgfscope}%
\begin{pgfscope}%
\pgfpathrectangle{\pgfqpoint{3.393750in}{0.699667in}}{\pgfqpoint{0.096250in}{1.925000in}}%
\pgfusepath{clip}%
\pgfsetbuttcap%
\pgfsetroundjoin%
\definecolor{currentfill}{rgb}{0.990175,0.927196,0.682926}%
\pgfsetfillcolor{currentfill}%
\pgfsetfillopacity{0.800000}%
\pgfsetlinewidth{0.000000pt}%
\definecolor{currentstroke}{rgb}{0.000000,0.000000,0.000000}%
\pgfsetstrokecolor{currentstroke}%
\pgfsetdash{}{0pt}%
\pgfpathmoveto{\pgfqpoint{3.393750in}{2.549472in}}%
\pgfpathlineto{\pgfqpoint{3.490000in}{2.549472in}}%
\pgfpathlineto{\pgfqpoint{3.490000in}{2.556991in}}%
\pgfpathlineto{\pgfqpoint{3.393750in}{2.556991in}}%
\pgfpathlineto{\pgfqpoint{3.393750in}{2.549472in}}%
\pgfusepath{fill}%
\end{pgfscope}%
\begin{pgfscope}%
\pgfpathrectangle{\pgfqpoint{3.393750in}{0.699667in}}{\pgfqpoint{0.096250in}{1.925000in}}%
\pgfusepath{clip}%
\pgfsetbuttcap%
\pgfsetroundjoin%
\definecolor{currentfill}{rgb}{0.989815,0.934329,0.690198}%
\pgfsetfillcolor{currentfill}%
\pgfsetfillopacity{0.800000}%
\pgfsetlinewidth{0.000000pt}%
\definecolor{currentstroke}{rgb}{0.000000,0.000000,0.000000}%
\pgfsetstrokecolor{currentstroke}%
\pgfsetdash{}{0pt}%
\pgfpathmoveto{\pgfqpoint{3.393750in}{2.556991in}}%
\pgfpathlineto{\pgfqpoint{3.490000in}{2.556991in}}%
\pgfpathlineto{\pgfqpoint{3.490000in}{2.564511in}}%
\pgfpathlineto{\pgfqpoint{3.393750in}{2.564511in}}%
\pgfpathlineto{\pgfqpoint{3.393750in}{2.556991in}}%
\pgfusepath{fill}%
\end{pgfscope}%
\begin{pgfscope}%
\pgfpathrectangle{\pgfqpoint{3.393750in}{0.699667in}}{\pgfqpoint{0.096250in}{1.925000in}}%
\pgfusepath{clip}%
\pgfsetbuttcap%
\pgfsetroundjoin%
\definecolor{currentfill}{rgb}{0.989434,0.941470,0.697519}%
\pgfsetfillcolor{currentfill}%
\pgfsetfillopacity{0.800000}%
\pgfsetlinewidth{0.000000pt}%
\definecolor{currentstroke}{rgb}{0.000000,0.000000,0.000000}%
\pgfsetstrokecolor{currentstroke}%
\pgfsetdash{}{0pt}%
\pgfpathmoveto{\pgfqpoint{3.393750in}{2.564511in}}%
\pgfpathlineto{\pgfqpoint{3.490000in}{2.564511in}}%
\pgfpathlineto{\pgfqpoint{3.490000in}{2.572030in}}%
\pgfpathlineto{\pgfqpoint{3.393750in}{2.572030in}}%
\pgfpathlineto{\pgfqpoint{3.393750in}{2.564511in}}%
\pgfusepath{fill}%
\end{pgfscope}%
\begin{pgfscope}%
\pgfpathrectangle{\pgfqpoint{3.393750in}{0.699667in}}{\pgfqpoint{0.096250in}{1.925000in}}%
\pgfusepath{clip}%
\pgfsetbuttcap%
\pgfsetroundjoin%
\definecolor{currentfill}{rgb}{0.989077,0.948604,0.704863}%
\pgfsetfillcolor{currentfill}%
\pgfsetfillopacity{0.800000}%
\pgfsetlinewidth{0.000000pt}%
\definecolor{currentstroke}{rgb}{0.000000,0.000000,0.000000}%
\pgfsetstrokecolor{currentstroke}%
\pgfsetdash{}{0pt}%
\pgfpathmoveto{\pgfqpoint{3.393750in}{2.572030in}}%
\pgfpathlineto{\pgfqpoint{3.490000in}{2.572030in}}%
\pgfpathlineto{\pgfqpoint{3.490000in}{2.579550in}}%
\pgfpathlineto{\pgfqpoint{3.393750in}{2.579550in}}%
\pgfpathlineto{\pgfqpoint{3.393750in}{2.572030in}}%
\pgfusepath{fill}%
\end{pgfscope}%
\begin{pgfscope}%
\pgfpathrectangle{\pgfqpoint{3.393750in}{0.699667in}}{\pgfqpoint{0.096250in}{1.925000in}}%
\pgfusepath{clip}%
\pgfsetbuttcap%
\pgfsetroundjoin%
\definecolor{currentfill}{rgb}{0.988717,0.955742,0.712242}%
\pgfsetfillcolor{currentfill}%
\pgfsetfillopacity{0.800000}%
\pgfsetlinewidth{0.000000pt}%
\definecolor{currentstroke}{rgb}{0.000000,0.000000,0.000000}%
\pgfsetstrokecolor{currentstroke}%
\pgfsetdash{}{0pt}%
\pgfpathmoveto{\pgfqpoint{3.393750in}{2.579550in}}%
\pgfpathlineto{\pgfqpoint{3.490000in}{2.579550in}}%
\pgfpathlineto{\pgfqpoint{3.490000in}{2.587070in}}%
\pgfpathlineto{\pgfqpoint{3.393750in}{2.587070in}}%
\pgfpathlineto{\pgfqpoint{3.393750in}{2.579550in}}%
\pgfusepath{fill}%
\end{pgfscope}%
\begin{pgfscope}%
\pgfpathrectangle{\pgfqpoint{3.393750in}{0.699667in}}{\pgfqpoint{0.096250in}{1.925000in}}%
\pgfusepath{clip}%
\pgfsetbuttcap%
\pgfsetroundjoin%
\definecolor{currentfill}{rgb}{0.988367,0.962878,0.719649}%
\pgfsetfillcolor{currentfill}%
\pgfsetfillopacity{0.800000}%
\pgfsetlinewidth{0.000000pt}%
\definecolor{currentstroke}{rgb}{0.000000,0.000000,0.000000}%
\pgfsetstrokecolor{currentstroke}%
\pgfsetdash{}{0pt}%
\pgfpathmoveto{\pgfqpoint{3.393750in}{2.587070in}}%
\pgfpathlineto{\pgfqpoint{3.490000in}{2.587070in}}%
\pgfpathlineto{\pgfqpoint{3.490000in}{2.594589in}}%
\pgfpathlineto{\pgfqpoint{3.393750in}{2.594589in}}%
\pgfpathlineto{\pgfqpoint{3.393750in}{2.587070in}}%
\pgfusepath{fill}%
\end{pgfscope}%
\begin{pgfscope}%
\pgfpathrectangle{\pgfqpoint{3.393750in}{0.699667in}}{\pgfqpoint{0.096250in}{1.925000in}}%
\pgfusepath{clip}%
\pgfsetbuttcap%
\pgfsetroundjoin%
\definecolor{currentfill}{rgb}{0.988033,0.970012,0.727077}%
\pgfsetfillcolor{currentfill}%
\pgfsetfillopacity{0.800000}%
\pgfsetlinewidth{0.000000pt}%
\definecolor{currentstroke}{rgb}{0.000000,0.000000,0.000000}%
\pgfsetstrokecolor{currentstroke}%
\pgfsetdash{}{0pt}%
\pgfpathmoveto{\pgfqpoint{3.393750in}{2.594589in}}%
\pgfpathlineto{\pgfqpoint{3.490000in}{2.594589in}}%
\pgfpathlineto{\pgfqpoint{3.490000in}{2.602109in}}%
\pgfpathlineto{\pgfqpoint{3.393750in}{2.602109in}}%
\pgfpathlineto{\pgfqpoint{3.393750in}{2.594589in}}%
\pgfusepath{fill}%
\end{pgfscope}%
\begin{pgfscope}%
\pgfpathrectangle{\pgfqpoint{3.393750in}{0.699667in}}{\pgfqpoint{0.096250in}{1.925000in}}%
\pgfusepath{clip}%
\pgfsetbuttcap%
\pgfsetroundjoin%
\definecolor{currentfill}{rgb}{0.987691,0.977154,0.734536}%
\pgfsetfillcolor{currentfill}%
\pgfsetfillopacity{0.800000}%
\pgfsetlinewidth{0.000000pt}%
\definecolor{currentstroke}{rgb}{0.000000,0.000000,0.000000}%
\pgfsetstrokecolor{currentstroke}%
\pgfsetdash{}{0pt}%
\pgfpathmoveto{\pgfqpoint{3.393750in}{2.602109in}}%
\pgfpathlineto{\pgfqpoint{3.490000in}{2.602109in}}%
\pgfpathlineto{\pgfqpoint{3.490000in}{2.609628in}}%
\pgfpathlineto{\pgfqpoint{3.393750in}{2.609628in}}%
\pgfpathlineto{\pgfqpoint{3.393750in}{2.602109in}}%
\pgfusepath{fill}%
\end{pgfscope}%
\begin{pgfscope}%
\pgfpathrectangle{\pgfqpoint{3.393750in}{0.699667in}}{\pgfqpoint{0.096250in}{1.925000in}}%
\pgfusepath{clip}%
\pgfsetbuttcap%
\pgfsetroundjoin%
\definecolor{currentfill}{rgb}{0.987387,0.984288,0.742002}%
\pgfsetfillcolor{currentfill}%
\pgfsetfillopacity{0.800000}%
\pgfsetlinewidth{0.000000pt}%
\definecolor{currentstroke}{rgb}{0.000000,0.000000,0.000000}%
\pgfsetstrokecolor{currentstroke}%
\pgfsetdash{}{0pt}%
\pgfpathmoveto{\pgfqpoint{3.393750in}{2.609628in}}%
\pgfpathlineto{\pgfqpoint{3.490000in}{2.609628in}}%
\pgfpathlineto{\pgfqpoint{3.490000in}{2.617148in}}%
\pgfpathlineto{\pgfqpoint{3.393750in}{2.617148in}}%
\pgfpathlineto{\pgfqpoint{3.393750in}{2.609628in}}%
\pgfusepath{fill}%
\end{pgfscope}%
\begin{pgfscope}%
\pgfpathrectangle{\pgfqpoint{3.393750in}{0.699667in}}{\pgfqpoint{0.096250in}{1.925000in}}%
\pgfusepath{clip}%
\pgfsetbuttcap%
\pgfsetroundjoin%
\definecolor{currentfill}{rgb}{0.987053,0.991438,0.749504}%
\pgfsetfillcolor{currentfill}%
\pgfsetfillopacity{0.800000}%
\pgfsetlinewidth{0.000000pt}%
\definecolor{currentstroke}{rgb}{0.000000,0.000000,0.000000}%
\pgfsetstrokecolor{currentstroke}%
\pgfsetdash{}{0pt}%
\pgfpathmoveto{\pgfqpoint{3.393750in}{2.617148in}}%
\pgfpathlineto{\pgfqpoint{3.490000in}{2.617148in}}%
\pgfpathlineto{\pgfqpoint{3.490000in}{2.624667in}}%
\pgfpathlineto{\pgfqpoint{3.393750in}{2.624667in}}%
\pgfpathlineto{\pgfqpoint{3.393750in}{2.617148in}}%
\pgfusepath{fill}%
\end{pgfscope}%
\begin{pgfscope}%
\pgfsetbuttcap%
\pgfsetroundjoin%
\definecolor{currentfill}{rgb}{0.000000,0.000000,0.000000}%
\pgfsetfillcolor{currentfill}%
\pgfsetlinewidth{0.803000pt}%
\definecolor{currentstroke}{rgb}{0.000000,0.000000,0.000000}%
\pgfsetstrokecolor{currentstroke}%
\pgfsetdash{}{0pt}%
\pgfsys@defobject{currentmarker}{\pgfqpoint{0.000000in}{0.000000in}}{\pgfqpoint{0.048611in}{0.000000in}}{%
\pgfpathmoveto{\pgfqpoint{0.000000in}{0.000000in}}%
\pgfpathlineto{\pgfqpoint{0.048611in}{0.000000in}}%
\pgfusepath{stroke,fill}%
}%
\begin{pgfscope}%
\pgfsys@transformshift{3.490000in}{1.096884in}%
\pgfsys@useobject{currentmarker}{}%
\end{pgfscope}%
\end{pgfscope}%
\begin{pgfscope}%
\definecolor{textcolor}{rgb}{0.000000,0.000000,0.000000}%
\pgfsetstrokecolor{textcolor}%
\pgfsetfillcolor{textcolor}%
\pgftext[x=3.587222in, y=1.048689in, left, base]{\color{textcolor}{\rmfamily\fontsize{10.000000}{12.000000}\selectfont\catcode`\^=\active\def^{\ifmmode\sp\else\^{}\fi}\catcode`\%=\active\def%{\%}$\mathdefault{1}$}}%
\end{pgfscope}%
\begin{pgfscope}%
\pgfsetbuttcap%
\pgfsetroundjoin%
\definecolor{currentfill}{rgb}{0.000000,0.000000,0.000000}%
\pgfsetfillcolor{currentfill}%
\pgfsetlinewidth{0.803000pt}%
\definecolor{currentstroke}{rgb}{0.000000,0.000000,0.000000}%
\pgfsetstrokecolor{currentstroke}%
\pgfsetdash{}{0pt}%
\pgfsys@defobject{currentmarker}{\pgfqpoint{0.000000in}{0.000000in}}{\pgfqpoint{0.048611in}{0.000000in}}{%
\pgfpathmoveto{\pgfqpoint{0.000000in}{0.000000in}}%
\pgfpathlineto{\pgfqpoint{0.048611in}{0.000000in}}%
\pgfusepath{stroke,fill}%
}%
\begin{pgfscope}%
\pgfsys@transformshift{3.490000in}{1.591523in}%
\pgfsys@useobject{currentmarker}{}%
\end{pgfscope}%
\end{pgfscope}%
\begin{pgfscope}%
\definecolor{textcolor}{rgb}{0.000000,0.000000,0.000000}%
\pgfsetstrokecolor{textcolor}%
\pgfsetfillcolor{textcolor}%
\pgftext[x=3.587222in, y=1.543329in, left, base]{\color{textcolor}{\rmfamily\fontsize{10.000000}{12.000000}\selectfont\catcode`\^=\active\def^{\ifmmode\sp\else\^{}\fi}\catcode`\%=\active\def%{\%}$\mathdefault{2}$}}%
\end{pgfscope}%
\begin{pgfscope}%
\pgfsetbuttcap%
\pgfsetroundjoin%
\definecolor{currentfill}{rgb}{0.000000,0.000000,0.000000}%
\pgfsetfillcolor{currentfill}%
\pgfsetlinewidth{0.803000pt}%
\definecolor{currentstroke}{rgb}{0.000000,0.000000,0.000000}%
\pgfsetstrokecolor{currentstroke}%
\pgfsetdash{}{0pt}%
\pgfsys@defobject{currentmarker}{\pgfqpoint{0.000000in}{0.000000in}}{\pgfqpoint{0.048611in}{0.000000in}}{%
\pgfpathmoveto{\pgfqpoint{0.000000in}{0.000000in}}%
\pgfpathlineto{\pgfqpoint{0.048611in}{0.000000in}}%
\pgfusepath{stroke,fill}%
}%
\begin{pgfscope}%
\pgfsys@transformshift{3.490000in}{2.086163in}%
\pgfsys@useobject{currentmarker}{}%
\end{pgfscope}%
\end{pgfscope}%
\begin{pgfscope}%
\definecolor{textcolor}{rgb}{0.000000,0.000000,0.000000}%
\pgfsetstrokecolor{textcolor}%
\pgfsetfillcolor{textcolor}%
\pgftext[x=3.587222in, y=2.037969in, left, base]{\color{textcolor}{\rmfamily\fontsize{10.000000}{12.000000}\selectfont\catcode`\^=\active\def^{\ifmmode\sp\else\^{}\fi}\catcode`\%=\active\def%{\%}$\mathdefault{3}$}}%
\end{pgfscope}%
\begin{pgfscope}%
\pgfsetbuttcap%
\pgfsetroundjoin%
\definecolor{currentfill}{rgb}{0.000000,0.000000,0.000000}%
\pgfsetfillcolor{currentfill}%
\pgfsetlinewidth{0.803000pt}%
\definecolor{currentstroke}{rgb}{0.000000,0.000000,0.000000}%
\pgfsetstrokecolor{currentstroke}%
\pgfsetdash{}{0pt}%
\pgfsys@defobject{currentmarker}{\pgfqpoint{0.000000in}{0.000000in}}{\pgfqpoint{0.048611in}{0.000000in}}{%
\pgfpathmoveto{\pgfqpoint{0.000000in}{0.000000in}}%
\pgfpathlineto{\pgfqpoint{0.048611in}{0.000000in}}%
\pgfusepath{stroke,fill}%
}%
\begin{pgfscope}%
\pgfsys@transformshift{3.490000in}{2.580803in}%
\pgfsys@useobject{currentmarker}{}%
\end{pgfscope}%
\end{pgfscope}%
\begin{pgfscope}%
\definecolor{textcolor}{rgb}{0.000000,0.000000,0.000000}%
\pgfsetstrokecolor{textcolor}%
\pgfsetfillcolor{textcolor}%
\pgftext[x=3.587222in, y=2.532608in, left, base]{\color{textcolor}{\rmfamily\fontsize{10.000000}{12.000000}\selectfont\catcode`\^=\active\def^{\ifmmode\sp\else\^{}\fi}\catcode`\%=\active\def%{\%}$\mathdefault{4}$}}%
\end{pgfscope}%
\begin{pgfscope}%
\definecolor{textcolor}{rgb}{0.000000,0.000000,0.000000}%
\pgfsetstrokecolor{textcolor}%
\pgfsetfillcolor{textcolor}%
\pgftext[x=3.712222in,y=1.662167in,,top,rotate=90.000000]{\color{textcolor}{\rmfamily\fontsize{10.000000}{12.000000}\selectfont\catcode`\^=\active\def^{\ifmmode\sp\else\^{}\fi}\catcode`\%=\active\def%{\%}Intensidad / Distancia}}%
\end{pgfscope}%
\begin{pgfscope}%
\pgfsetrectcap%
\pgfsetmiterjoin%
\pgfsetlinewidth{0.803000pt}%
\definecolor{currentstroke}{rgb}{0.000000,0.000000,0.000000}%
\pgfsetstrokecolor{currentstroke}%
\pgfsetdash{}{0pt}%
\pgfpathmoveto{\pgfqpoint{3.393750in}{0.699667in}}%
\pgfpathlineto{\pgfqpoint{3.441875in}{0.699667in}}%
\pgfpathlineto{\pgfqpoint{3.490000in}{0.699667in}}%
\pgfpathlineto{\pgfqpoint{3.490000in}{2.624667in}}%
\pgfpathlineto{\pgfqpoint{3.441875in}{2.624667in}}%
\pgfpathlineto{\pgfqpoint{3.393750in}{2.624667in}}%
\pgfpathlineto{\pgfqpoint{3.393750in}{0.699667in}}%
\pgfpathclose%
\pgfusepath{stroke}%
\end{pgfscope}%
\end{pgfpicture}%
\makeatother%
\endgroup%
 
    \caption{Gráfico 3D scatter}
    \label{fig:grafico_3D_scatter}
\end{figure}

\subsubsection{Superficie en malla regular (\textit{ax.plot\_surface})}

La función \texttt{plot\_surface} funciona como una red de pesca perfectamente cuadrada, cuando la lanza sobre una superficie (como una duna de arena o una escultura), la red se adapta a las curvas de ese objeto, donde cada nudo de la red tiene una coordenada (\(x,y\)) en el suelo y una altura (\(z\)).

Lo que hace \texttt{plot\_surface} es rellenar los huecos entre los hilos de la red con color, creando una piel continua.


\emph{¿Cuándo se usa?}
\begin{itemize}
    \item \textit{Topografía y Mapas de Elevación}: Si tiene las mediciones de altura de un terreno tomadas por un satélite en una cuadrícula fija, \texttt{plot\_surface} le permite reconstruir la montaña o el valle digitalmente.
    \item \textit{Paisajes de Optimización en Machine Learning}: En Inteligencia Artificial, queremos que el error de un modelo sea el mínimo posible. Este error suele visualizarse como una superficie con valles y colinas.
    \item \textit{Modelado de Ondas y Propagación}: Visualiza cómo se expande el sonido en una habitación o las ondas electromagnéticas de una antena wifi.Lo que nos permite identificar puntos muertos de señal en un diseño arquitectónico.
\end{itemize}

\begin{pycode}{superficie de malla regular}
import matplotlib.pyplot as plt
import numpy as np

# Datos
x = np.linspace(-3, 3, 200)
y = np.linspace(-3, 3, 200)
X, Y = np.meshgrid(x, y)
Z = np.sin(np.sqrt(X**2 + Y**2))

# Entrono 3D
fig, ax = plt.subplots(subplot_kw={"projection": "3d"}, figsize=(5, 5), layout="constrained")

# Malla:
surf = ax.plot_surface(X, Y, Z, cmap="viridis", rstride=1, cstride=1, linewidth=0, antialiased=True)

ax.set_title("Superficie 3D: z = sin(sqrt(x^2 + y^2))")
ax.set_xlabel("X")
ax.set_ylabel("Y")
ax.set_zlabel("Z")

# Ángulo de visualización
cb = fig.colorbar(surf, ax=ax, shrink=0.7, pad=0.05)
cb.set_label("Altura (Z)")
ax.view_init(elev=25, azim=45)
ax.set_box_aspect((1, 1, 0.5))   
\end{pycode}

\emph{Análisis del código: Parámetros clave}
\begin{itemize}
\item \texttt{np.meshgrid(x, y)}: Convierte dos vectores simples en matrices de coordenadas que definen todos los puntos de cruce de la rejilla. Sin esto, no hay ``malla''.
\item \texttt{rstride} y \texttt{cstride}: (\textit{Row/Column stride}). Controlan la densidad de la malla. Un valor de \texttt{1} usa todos los datos; valores más altos saltan puntos para que el gráfico sea más liviano de procesar.
\item \texttt{antialiased=True}: Suaviza los bordes de los polígonos para que la superficie no se vea "pixelada" o con bordes de sierra.
\item \texttt{ax.view\_init(elev=25, azim=45)}: Define el ángulo de cámara. \texttt{elev} es la altura del observador y \texttt{azim} es la rotación lateral. Es vital para encontrar el ángulo que mejor muestra la curvatura.
\end{itemize}

El gráfico de malla regular se puede observar en la figura \ref{fig:grafico_3D_malla_regular}

\begin{figure}[ht]
    \centering
    \input{Graficos/grafico_3D_malla_regular.pgf} 
    \caption{Gráfico de Malla regular}
    \label{fig:grafico_3D_malla_regular}
\end{figure}



\subsubsection{Superficie con datos dispersos (\texttt{ax.plot\_trisurf})}

A diferencia del gráfico anterior donde tenía una rejilla perfecta, aquí tiene una nube de puntos aleatorios. Imagine que tira puñados de arena sobre una mesa; los puntos no estáran alineados.    

Para crear una superficie con puntos \textit{desordenados}, \textit{Matplotlib} usa la Triangulación de Delaunay. Basicamente el algoritmo conecta los puntos más cercanos entre sí formando triángulos, de manera que no haya ángulos demasiado agudos. Esto crea una piel o malla que cubre todos los puntos, permitiendo ver una superficie continua donde antes solo había puntos sueltos.

\emph{¿Cuándo usarlo?}

Cuando no se tiene control sobre las coordenadas X e Y (ej: ubicación geográfica de clientes) o cuando tiene una nube de puntos y quiere unirlos para ver la tendencia general de la superficie.

\begin{pycode}{Superficie con datos dispersos}
import matplotlib.pyplot as plt
import numpy as np
from matplotlib.tri import Triangulation # Necesitamos importar esta libreria de matplot

# Datos
rng = np.random.default_rng(0)
n = 200
x = 6 * rng.random(n) - 3
y = 6 * rng.random(n) - 3
z = np.sin(np.sqrt(x**2 + y**2))

# Creamos la triangulación, donde matplot conecta los puntos para formar triángulos
tri = Triangulation(x, y)

# Creamos la figura 
fig, ax = plt.subplots(subplot_kw={"projection": "3d"}, figsize=(5, 5), layout="constrained")

# Utilizamos Trisurf
surf = ax.plot_trisurf(tri, z, cmap="plasma", linewidth=0.2, antialiased=True)

ax.set_title("Superficie 3D triangulada (datos dispersos)")
ax.set_xlabel("X")
ax.set_ylabel("Y")
ax.set_zlabel("Z")

# Barra de color a la derecha
fig.colorbar(surf, ax=ax, shrink=0.7, pad=0.05).set_label("Altura (Z)")
ax.set_box_aspect((1, 1, 0.6))
\end{pycode}

\emph{Análisis del código}
\begin{itemize}
    \item \texttt{Triangulation(x, y)}: Es el motor del gráfico. Toma los puntos sueltos y calcula matemáticamente cómo conectarlos con triángulos sin que se crucen (normalmente usando la triangulación de Delaunay).
    \item \texttt{cmap='plasma'}: Otro mapa de colores perceptualmente uniforme. A diferencia de \textit{viridis}, \textit{plasma} usa tonos púrpuras y amarillos intensos, muy útiles para resaltar relieves.
    \item \texttt{linewidth=0.2}: Define el grosor de las líneas que unen los triángulos. En este tipo de gráficos, un poco de línea ayuda a entender la estructura de la superficie.
    \item \texttt{ax.set\_box\_aspect((1, 1, 0.6))}: Nota que aquí el valor de Z (0.6) es menor. Esto aplasta un poco el gráfico verticalmente, lo cual es útil si la variación en Z es muy abrupta y quiere que el gráfico sea de una dimensión normal.
\end{itemize}

\emph{Diferencias entre \texttt{plot\_surface} y \texttt{plot\_trisurf}}:

\begin{itemize}
\item \texttt{plot\_surface} requiere que los datos estén ordenados en una matriz (grilla). Es una estrictura rectangular uniforme. Es rápido.
\item \texttt{plot\_trisurf} acepta listas de datos desordenados y él mismo construye la conexión entre ellos. Es una estructura Triangular irregular. Es Lento.
\end{itemize}

Para ver el gráfico de \texttt{plot\_trisurf} ver la figura \ref{fig:grafico_3D_Triangulado}
\begin{figure}[ht]
    \centering
    %% Creator: Matplotlib, PGF backend
%%
%% To include the figure in your LaTeX document, write
%%   \input{<filename>.pgf}
%%
%% Make sure the required packages are loaded in your preamble
%%   \usepackage{pgf}
%%
%% Also ensure that all the required font packages are loaded; for instance,
%% the lmodern package is sometimes necessary when using math font.
%%   \usepackage{lmodern}
%%
%% Figures using additional raster images can only be included by \input if
%% they are in the same directory as the main LaTeX file. For loading figures
%% from other directories you can use the `import` package
%%   \usepackage{import}
%%
%% and then include the figures with
%%   \import{<path to file>}{<filename>.pgf}
%%
%% Matplotlib used the following preamble
%%   \def\mathdefault#1{#1}
%%   \everymath=\expandafter{\the\everymath\displaystyle}
%%   \IfFileExists{scrextend.sty}{
%%     \usepackage[fontsize=10.000000pt]{scrextend}
%%   }{
%%     \renewcommand{\normalsize}{\fontsize{10.000000}{12.000000}\selectfont}
%%     \normalsize
%%   }
%%   \usepackage{fontspec}
%%   \setmainfont{CMU Serif}
%%   \makeatletter\@ifpackageloaded{underscore}{}{\usepackage[strings]{underscore}}\makeatother
%%
\begingroup%
\makeatletter%
\begin{pgfpicture}%
\pgfpathrectangle{\pgfpointorigin}{\pgfqpoint{5.104281in}{4.365512in}}%
\pgfusepath{use as bounding box, clip}%
\begin{pgfscope}%
\pgfsetbuttcap%
\pgfsetmiterjoin%
\definecolor{currentfill}{rgb}{1.000000,1.000000,1.000000}%
\pgfsetfillcolor{currentfill}%
\pgfsetlinewidth{0.000000pt}%
\definecolor{currentstroke}{rgb}{1.000000,1.000000,1.000000}%
\pgfsetstrokecolor{currentstroke}%
\pgfsetdash{}{0pt}%
\pgfpathmoveto{\pgfqpoint{0.000000in}{0.000000in}}%
\pgfpathlineto{\pgfqpoint{5.104281in}{0.000000in}}%
\pgfpathlineto{\pgfqpoint{5.104281in}{4.365512in}}%
\pgfpathlineto{\pgfqpoint{0.000000in}{4.365512in}}%
\pgfpathlineto{\pgfqpoint{0.000000in}{0.000000in}}%
\pgfpathclose%
\pgfusepath{fill}%
\end{pgfscope}%
\begin{pgfscope}%
\pgfsetbuttcap%
\pgfsetmiterjoin%
\definecolor{currentfill}{rgb}{1.000000,1.000000,1.000000}%
\pgfsetfillcolor{currentfill}%
\pgfsetlinewidth{0.000000pt}%
\definecolor{currentstroke}{rgb}{0.000000,0.000000,0.000000}%
\pgfsetstrokecolor{currentstroke}%
\pgfsetstrokeopacity{0.000000}%
\pgfsetdash{}{0pt}%
\pgfpathmoveto{\pgfqpoint{0.100000in}{0.100000in}}%
\pgfpathlineto{\pgfqpoint{4.057178in}{0.100000in}}%
\pgfpathlineto{\pgfqpoint{4.057178in}{4.057178in}}%
\pgfpathlineto{\pgfqpoint{0.100000in}{4.057178in}}%
\pgfpathlineto{\pgfqpoint{0.100000in}{0.100000in}}%
\pgfpathclose%
\pgfusepath{fill}%
\end{pgfscope}%
\begin{pgfscope}%
\pgfsetbuttcap%
\pgfsetmiterjoin%
\definecolor{currentfill}{rgb}{0.950000,0.950000,0.950000}%
\pgfsetfillcolor{currentfill}%
\pgfsetfillopacity{0.500000}%
\pgfsetlinewidth{1.003750pt}%
\definecolor{currentstroke}{rgb}{0.950000,0.950000,0.950000}%
\pgfsetstrokecolor{currentstroke}%
\pgfsetstrokeopacity{0.500000}%
\pgfsetdash{}{0pt}%
\pgfpathmoveto{\pgfqpoint{0.317814in}{1.198888in}}%
\pgfpathlineto{\pgfqpoint{1.687356in}{2.330587in}}%
\pgfpathlineto{\pgfqpoint{1.671665in}{3.644771in}}%
\pgfpathlineto{\pgfqpoint{0.247611in}{2.596413in}}%
\pgfusepath{stroke,fill}%
\end{pgfscope}%
\begin{pgfscope}%
\pgfsetbuttcap%
\pgfsetmiterjoin%
\definecolor{currentfill}{rgb}{0.900000,0.900000,0.900000}%
\pgfsetfillcolor{currentfill}%
\pgfsetfillopacity{0.500000}%
\pgfsetlinewidth{1.003750pt}%
\definecolor{currentstroke}{rgb}{0.900000,0.900000,0.900000}%
\pgfsetstrokecolor{currentstroke}%
\pgfsetstrokeopacity{0.500000}%
\pgfsetdash{}{0pt}%
\pgfpathmoveto{\pgfqpoint{1.687356in}{2.330587in}}%
\pgfpathlineto{\pgfqpoint{3.877604in}{1.701946in}}%
\pgfpathlineto{\pgfqpoint{3.942496in}{3.063277in}}%
\pgfpathlineto{\pgfqpoint{1.671665in}{3.644771in}}%
\pgfusepath{stroke,fill}%
\end{pgfscope}%
\begin{pgfscope}%
\pgfsetbuttcap%
\pgfsetmiterjoin%
\definecolor{currentfill}{rgb}{0.925000,0.925000,0.925000}%
\pgfsetfillcolor{currentfill}%
\pgfsetfillopacity{0.500000}%
\pgfsetlinewidth{1.003750pt}%
\definecolor{currentstroke}{rgb}{0.925000,0.925000,0.925000}%
\pgfsetstrokecolor{currentstroke}%
\pgfsetstrokeopacity{0.500000}%
\pgfsetdash{}{0pt}%
\pgfpathmoveto{\pgfqpoint{0.317814in}{1.198888in}}%
\pgfpathlineto{\pgfqpoint{2.645817in}{0.444059in}}%
\pgfpathlineto{\pgfqpoint{3.877604in}{1.701946in}}%
\pgfpathlineto{\pgfqpoint{1.687356in}{2.330587in}}%
\pgfusepath{stroke,fill}%
\end{pgfscope}%
\begin{pgfscope}%
\pgfsetbuttcap%
\pgfsetroundjoin%
\pgfsetlinewidth{0.501875pt}%
\definecolor{currentstroke}{rgb}{0.501961,0.501961,0.501961}%
\pgfsetstrokecolor{currentstroke}%
\pgfsetdash{{1.850000pt}{0.800000pt}}{0.000000pt}%
\pgfpathmoveto{\pgfqpoint{0.453146in}{1.155008in}}%
\pgfpathlineto{\pgfqpoint{1.815260in}{2.293876in}}%
\pgfpathlineto{\pgfqpoint{1.804047in}{3.610872in}}%
\pgfusepath{stroke}%
\end{pgfscope}%
\begin{pgfscope}%
\pgfsetbuttcap%
\pgfsetroundjoin%
\pgfsetlinewidth{0.501875pt}%
\definecolor{currentstroke}{rgb}{0.501961,0.501961,0.501961}%
\pgfsetstrokecolor{currentstroke}%
\pgfsetdash{{1.850000pt}{0.800000pt}}{0.000000pt}%
\pgfpathmoveto{\pgfqpoint{0.780049in}{1.049014in}}%
\pgfpathlineto{\pgfqpoint{2.123925in}{2.205284in}}%
\pgfpathlineto{\pgfqpoint{2.123634in}{3.529034in}}%
\pgfusepath{stroke}%
\end{pgfscope}%
\begin{pgfscope}%
\pgfsetbuttcap%
\pgfsetroundjoin%
\pgfsetlinewidth{0.501875pt}%
\definecolor{currentstroke}{rgb}{0.501961,0.501961,0.501961}%
\pgfsetstrokecolor{currentstroke}%
\pgfsetdash{{1.850000pt}{0.800000pt}}{0.000000pt}%
\pgfpathmoveto{\pgfqpoint{1.112194in}{0.941319in}}%
\pgfpathlineto{\pgfqpoint{2.437110in}{2.115394in}}%
\pgfpathlineto{\pgfqpoint{2.448070in}{3.445956in}}%
\pgfusepath{stroke}%
\end{pgfscope}%
\begin{pgfscope}%
\pgfsetbuttcap%
\pgfsetroundjoin%
\pgfsetlinewidth{0.501875pt}%
\definecolor{currentstroke}{rgb}{0.501961,0.501961,0.501961}%
\pgfsetstrokecolor{currentstroke}%
\pgfsetdash{{1.850000pt}{0.800000pt}}{0.000000pt}%
\pgfpathmoveto{\pgfqpoint{1.449708in}{0.831884in}}%
\pgfpathlineto{\pgfqpoint{2.754915in}{2.024178in}}%
\pgfpathlineto{\pgfqpoint{2.777465in}{3.361607in}}%
\pgfusepath{stroke}%
\end{pgfscope}%
\begin{pgfscope}%
\pgfsetbuttcap%
\pgfsetroundjoin%
\pgfsetlinewidth{0.501875pt}%
\definecolor{currentstroke}{rgb}{0.501961,0.501961,0.501961}%
\pgfsetstrokecolor{currentstroke}%
\pgfsetdash{{1.850000pt}{0.800000pt}}{0.000000pt}%
\pgfpathmoveto{\pgfqpoint{1.792723in}{0.720666in}}%
\pgfpathlineto{\pgfqpoint{3.077444in}{1.931606in}}%
\pgfpathlineto{\pgfqpoint{3.111935in}{3.275959in}}%
\pgfusepath{stroke}%
\end{pgfscope}%
\begin{pgfscope}%
\pgfsetbuttcap%
\pgfsetroundjoin%
\pgfsetlinewidth{0.501875pt}%
\definecolor{currentstroke}{rgb}{0.501961,0.501961,0.501961}%
\pgfsetstrokecolor{currentstroke}%
\pgfsetdash{{1.850000pt}{0.800000pt}}{0.000000pt}%
\pgfpathmoveto{\pgfqpoint{2.141373in}{0.607620in}}%
\pgfpathlineto{\pgfqpoint{3.404802in}{1.837649in}}%
\pgfpathlineto{\pgfqpoint{3.451596in}{3.188982in}}%
\pgfusepath{stroke}%
\end{pgfscope}%
\begin{pgfscope}%
\pgfsetbuttcap%
\pgfsetroundjoin%
\pgfsetlinewidth{0.501875pt}%
\definecolor{currentstroke}{rgb}{0.501961,0.501961,0.501961}%
\pgfsetstrokecolor{currentstroke}%
\pgfsetdash{{1.850000pt}{0.800000pt}}{0.000000pt}%
\pgfpathmoveto{\pgfqpoint{2.495798in}{0.492701in}}%
\pgfpathlineto{\pgfqpoint{3.737099in}{1.742274in}}%
\pgfpathlineto{\pgfqpoint{3.796572in}{3.100644in}}%
\pgfusepath{stroke}%
\end{pgfscope}%
\begin{pgfscope}%
\pgfsetbuttcap%
\pgfsetroundjoin%
\pgfsetlinewidth{0.501875pt}%
\definecolor{currentstroke}{rgb}{0.501961,0.501961,0.501961}%
\pgfsetstrokecolor{currentstroke}%
\pgfsetdash{{1.850000pt}{0.800000pt}}{0.000000pt}%
\pgfpathmoveto{\pgfqpoint{0.345985in}{2.668834in}}%
\pgfpathlineto{\pgfqpoint{0.412133in}{1.276827in}}%
\pgfpathlineto{\pgfqpoint{2.731031in}{0.531079in}}%
\pgfusepath{stroke}%
\end{pgfscope}%
\begin{pgfscope}%
\pgfsetbuttcap%
\pgfsetroundjoin%
\pgfsetlinewidth{0.501875pt}%
\definecolor{currentstroke}{rgb}{0.501961,0.501961,0.501961}%
\pgfsetstrokecolor{currentstroke}%
\pgfsetdash{{1.850000pt}{0.800000pt}}{0.000000pt}%
\pgfpathmoveto{\pgfqpoint{0.568023in}{2.832294in}}%
\pgfpathlineto{\pgfqpoint{0.625176in}{1.452872in}}%
\pgfpathlineto{\pgfqpoint{2.923300in}{0.727422in}}%
\pgfusepath{stroke}%
\end{pgfscope}%
\begin{pgfscope}%
\pgfsetbuttcap%
\pgfsetroundjoin%
\pgfsetlinewidth{0.501875pt}%
\definecolor{currentstroke}{rgb}{0.501961,0.501961,0.501961}%
\pgfsetstrokecolor{currentstroke}%
\pgfsetdash{{1.850000pt}{0.800000pt}}{0.000000pt}%
\pgfpathmoveto{\pgfqpoint{0.784001in}{2.991292in}}%
\pgfpathlineto{\pgfqpoint{0.832614in}{1.624284in}}%
\pgfpathlineto{\pgfqpoint{3.110232in}{0.918314in}}%
\pgfusepath{stroke}%
\end{pgfscope}%
\begin{pgfscope}%
\pgfsetbuttcap%
\pgfsetroundjoin%
\pgfsetlinewidth{0.501875pt}%
\definecolor{currentstroke}{rgb}{0.501961,0.501961,0.501961}%
\pgfsetstrokecolor{currentstroke}%
\pgfsetdash{{1.850000pt}{0.800000pt}}{0.000000pt}%
\pgfpathmoveto{\pgfqpoint{0.994162in}{3.146008in}}%
\pgfpathlineto{\pgfqpoint{1.034664in}{1.791245in}}%
\pgfpathlineto{\pgfqpoint{3.292046in}{1.103981in}}%
\pgfusepath{stroke}%
\end{pgfscope}%
\begin{pgfscope}%
\pgfsetbuttcap%
\pgfsetroundjoin%
\pgfsetlinewidth{0.501875pt}%
\definecolor{currentstroke}{rgb}{0.501961,0.501961,0.501961}%
\pgfsetstrokecolor{currentstroke}%
\pgfsetdash{{1.850000pt}{0.800000pt}}{0.000000pt}%
\pgfpathmoveto{\pgfqpoint{1.198739in}{3.296613in}}%
\pgfpathlineto{\pgfqpoint{1.231533in}{1.953925in}}%
\pgfpathlineto{\pgfqpoint{3.468950in}{1.284633in}}%
\pgfusepath{stroke}%
\end{pgfscope}%
\begin{pgfscope}%
\pgfsetbuttcap%
\pgfsetroundjoin%
\pgfsetlinewidth{0.501875pt}%
\definecolor{currentstroke}{rgb}{0.501961,0.501961,0.501961}%
\pgfsetstrokecolor{currentstroke}%
\pgfsetdash{{1.850000pt}{0.800000pt}}{0.000000pt}%
\pgfpathmoveto{\pgfqpoint{1.397952in}{3.443269in}}%
\pgfpathlineto{\pgfqpoint{1.423419in}{2.112487in}}%
\pgfpathlineto{\pgfqpoint{3.641139in}{1.460471in}}%
\pgfusepath{stroke}%
\end{pgfscope}%
\begin{pgfscope}%
\pgfsetbuttcap%
\pgfsetroundjoin%
\pgfsetlinewidth{0.501875pt}%
\definecolor{currentstroke}{rgb}{0.501961,0.501961,0.501961}%
\pgfsetstrokecolor{currentstroke}%
\pgfsetdash{{1.850000pt}{0.800000pt}}{0.000000pt}%
\pgfpathmoveto{\pgfqpoint{1.592008in}{3.586128in}}%
\pgfpathlineto{\pgfqpoint{1.610509in}{2.267086in}}%
\pgfpathlineto{\pgfqpoint{3.808801in}{1.631685in}}%
\pgfusepath{stroke}%
\end{pgfscope}%
\begin{pgfscope}%
\pgfsetbuttcap%
\pgfsetroundjoin%
\pgfsetlinewidth{0.501875pt}%
\definecolor{currentstroke}{rgb}{0.501961,0.501961,0.501961}%
\pgfsetstrokecolor{currentstroke}%
\pgfsetdash{{1.850000pt}{0.800000pt}}{0.000000pt}%
\pgfpathmoveto{\pgfqpoint{3.878701in}{1.724959in}}%
\pgfpathlineto{\pgfqpoint{1.687090in}{2.352843in}}%
\pgfpathlineto{\pgfqpoint{0.316629in}{1.222479in}}%
\pgfusepath{stroke}%
\end{pgfscope}%
\begin{pgfscope}%
\pgfsetbuttcap%
\pgfsetroundjoin%
\pgfsetlinewidth{0.501875pt}%
\definecolor{currentstroke}{rgb}{0.501961,0.501961,0.501961}%
\pgfsetstrokecolor{currentstroke}%
\pgfsetdash{{1.850000pt}{0.800000pt}}{0.000000pt}%
\pgfpathmoveto{\pgfqpoint{3.885426in}{1.866054in}}%
\pgfpathlineto{\pgfqpoint{1.685461in}{2.489266in}}%
\pgfpathlineto{\pgfqpoint{0.309362in}{1.367142in}}%
\pgfusepath{stroke}%
\end{pgfscope}%
\begin{pgfscope}%
\pgfsetbuttcap%
\pgfsetroundjoin%
\pgfsetlinewidth{0.501875pt}%
\definecolor{currentstroke}{rgb}{0.501961,0.501961,0.501961}%
\pgfsetstrokecolor{currentstroke}%
\pgfsetdash{{1.850000pt}{0.800000pt}}{0.000000pt}%
\pgfpathmoveto{\pgfqpoint{3.892204in}{2.008239in}}%
\pgfpathlineto{\pgfqpoint{1.683821in}{2.626691in}}%
\pgfpathlineto{\pgfqpoint{0.302037in}{1.512968in}}%
\pgfusepath{stroke}%
\end{pgfscope}%
\begin{pgfscope}%
\pgfsetbuttcap%
\pgfsetroundjoin%
\pgfsetlinewidth{0.501875pt}%
\definecolor{currentstroke}{rgb}{0.501961,0.501961,0.501961}%
\pgfsetstrokecolor{currentstroke}%
\pgfsetdash{{1.850000pt}{0.800000pt}}{0.000000pt}%
\pgfpathmoveto{\pgfqpoint{3.899034in}{2.151528in}}%
\pgfpathlineto{\pgfqpoint{1.682168in}{2.765130in}}%
\pgfpathlineto{\pgfqpoint{0.294652in}{1.659971in}}%
\pgfusepath{stroke}%
\end{pgfscope}%
\begin{pgfscope}%
\pgfsetbuttcap%
\pgfsetroundjoin%
\pgfsetlinewidth{0.501875pt}%
\definecolor{currentstroke}{rgb}{0.501961,0.501961,0.501961}%
\pgfsetstrokecolor{currentstroke}%
\pgfsetdash{{1.850000pt}{0.800000pt}}{0.000000pt}%
\pgfpathmoveto{\pgfqpoint{3.905918in}{2.295934in}}%
\pgfpathlineto{\pgfqpoint{1.680503in}{2.904594in}}%
\pgfpathlineto{\pgfqpoint{0.287208in}{1.808165in}}%
\pgfusepath{stroke}%
\end{pgfscope}%
\begin{pgfscope}%
\pgfsetbuttcap%
\pgfsetroundjoin%
\pgfsetlinewidth{0.501875pt}%
\definecolor{currentstroke}{rgb}{0.501961,0.501961,0.501961}%
\pgfsetstrokecolor{currentstroke}%
\pgfsetdash{{1.850000pt}{0.800000pt}}{0.000000pt}%
\pgfpathmoveto{\pgfqpoint{3.912855in}{2.441468in}}%
\pgfpathlineto{\pgfqpoint{1.678825in}{3.045093in}}%
\pgfpathlineto{\pgfqpoint{0.279703in}{1.957565in}}%
\pgfusepath{stroke}%
\end{pgfscope}%
\begin{pgfscope}%
\pgfsetbuttcap%
\pgfsetroundjoin%
\pgfsetlinewidth{0.501875pt}%
\definecolor{currentstroke}{rgb}{0.501961,0.501961,0.501961}%
\pgfsetstrokecolor{currentstroke}%
\pgfsetdash{{1.850000pt}{0.800000pt}}{0.000000pt}%
\pgfpathmoveto{\pgfqpoint{3.919847in}{2.588146in}}%
\pgfpathlineto{\pgfqpoint{1.677135in}{3.186641in}}%
\pgfpathlineto{\pgfqpoint{0.272136in}{2.108185in}}%
\pgfusepath{stroke}%
\end{pgfscope}%
\begin{pgfscope}%
\pgfsetbuttcap%
\pgfsetroundjoin%
\pgfsetlinewidth{0.501875pt}%
\definecolor{currentstroke}{rgb}{0.501961,0.501961,0.501961}%
\pgfsetstrokecolor{currentstroke}%
\pgfsetdash{{1.850000pt}{0.800000pt}}{0.000000pt}%
\pgfpathmoveto{\pgfqpoint{3.926894in}{2.735980in}}%
\pgfpathlineto{\pgfqpoint{1.675433in}{3.329248in}}%
\pgfpathlineto{\pgfqpoint{0.264508in}{2.260041in}}%
\pgfusepath{stroke}%
\end{pgfscope}%
\begin{pgfscope}%
\pgfsetbuttcap%
\pgfsetroundjoin%
\pgfsetlinewidth{0.501875pt}%
\definecolor{currentstroke}{rgb}{0.501961,0.501961,0.501961}%
\pgfsetstrokecolor{currentstroke}%
\pgfsetdash{{1.850000pt}{0.800000pt}}{0.000000pt}%
\pgfpathmoveto{\pgfqpoint{3.933997in}{2.884984in}}%
\pgfpathlineto{\pgfqpoint{1.673717in}{3.472926in}}%
\pgfpathlineto{\pgfqpoint{0.256817in}{2.413147in}}%
\pgfusepath{stroke}%
\end{pgfscope}%
\begin{pgfscope}%
\pgfsetbuttcap%
\pgfsetroundjoin%
\pgfsetlinewidth{0.501875pt}%
\definecolor{currentstroke}{rgb}{0.501961,0.501961,0.501961}%
\pgfsetstrokecolor{currentstroke}%
\pgfsetdash{{1.850000pt}{0.800000pt}}{0.000000pt}%
\pgfpathmoveto{\pgfqpoint{3.941156in}{3.035172in}}%
\pgfpathlineto{\pgfqpoint{1.671989in}{3.617688in}}%
\pgfpathlineto{\pgfqpoint{0.249062in}{2.567520in}}%
\pgfusepath{stroke}%
\end{pgfscope}%
\begin{pgfscope}%
\pgfsetrectcap%
\pgfsetroundjoin%
\pgfsetlinewidth{0.803000pt}%
\definecolor{currentstroke}{rgb}{0.000000,0.000000,0.000000}%
\pgfsetstrokecolor{currentstroke}%
\pgfsetdash{}{0pt}%
\pgfpathmoveto{\pgfqpoint{0.317814in}{1.198888in}}%
\pgfpathlineto{\pgfqpoint{2.645817in}{0.444059in}}%
\pgfusepath{stroke}%
\end{pgfscope}%
\begin{pgfscope}%
\pgfpathrectangle{\pgfqpoint{0.100000in}{0.100000in}}{\pgfqpoint{3.957178in}{3.957178in}}%
\pgfusepath{clip}%
\pgfsetbuttcap%
\pgfsetroundjoin%
\pgfsetlinewidth{0.501875pt}%
\definecolor{currentstroke}{rgb}{0.501961,0.501961,0.501961}%
\pgfsetstrokecolor{currentstroke}%
\pgfsetstrokeopacity{0.500000}%
\pgfsetdash{{1.850000pt}{0.800000pt}}{0.000000pt}%
\pgfpathmoveto{\pgfqpoint{2.132065in}{2.132065in}}%
\pgfusepath{stroke}%
\end{pgfscope}%
\begin{pgfscope}%
\pgfsetrectcap%
\pgfsetroundjoin%
\pgfsetlinewidth{0.803000pt}%
\definecolor{currentstroke}{rgb}{0.000000,0.000000,0.000000}%
\pgfsetstrokecolor{currentstroke}%
\pgfsetdash{}{0pt}%
\pgfpathmoveto{\pgfqpoint{0.465052in}{1.164963in}}%
\pgfpathlineto{\pgfqpoint{0.429280in}{1.135054in}}%
\pgfusepath{stroke}%
\end{pgfscope}%
\begin{pgfscope}%
\definecolor{textcolor}{rgb}{0.000000,0.000000,0.000000}%
\pgfsetstrokecolor{textcolor}%
\pgfsetfillcolor{textcolor}%
\pgftext[x=0.342104in,y=0.953141in,,top]{\color{textcolor}{\rmfamily\fontsize{10.000000}{12.000000}\selectfont\catcode`\^=\active\def^{\ifmmode\sp\else\^{}\fi}\catcode`\%=\active\def%{\%}$\mathdefault{\ensuremath{-}3}$}}%
\end{pgfscope}%
\begin{pgfscope}%
\pgfpathrectangle{\pgfqpoint{0.100000in}{0.100000in}}{\pgfqpoint{3.957178in}{3.957178in}}%
\pgfusepath{clip}%
\pgfsetbuttcap%
\pgfsetroundjoin%
\pgfsetlinewidth{0.501875pt}%
\definecolor{currentstroke}{rgb}{0.501961,0.501961,0.501961}%
\pgfsetstrokecolor{currentstroke}%
\pgfsetstrokeopacity{0.500000}%
\pgfsetdash{{1.850000pt}{0.800000pt}}{0.000000pt}%
\pgfpathmoveto{\pgfqpoint{2.132065in}{2.132065in}}%
\pgfusepath{stroke}%
\end{pgfscope}%
\begin{pgfscope}%
\pgfsetrectcap%
\pgfsetroundjoin%
\pgfsetlinewidth{0.803000pt}%
\definecolor{currentstroke}{rgb}{0.000000,0.000000,0.000000}%
\pgfsetstrokecolor{currentstroke}%
\pgfsetdash{}{0pt}%
\pgfpathmoveto{\pgfqpoint{0.791804in}{1.059127in}}%
\pgfpathlineto{\pgfqpoint{0.756486in}{1.028740in}}%
\pgfusepath{stroke}%
\end{pgfscope}%
\begin{pgfscope}%
\definecolor{textcolor}{rgb}{0.000000,0.000000,0.000000}%
\pgfsetstrokecolor{textcolor}%
\pgfsetfillcolor{textcolor}%
\pgftext[x=0.669558in,y=0.845067in,,top]{\color{textcolor}{\rmfamily\fontsize{10.000000}{12.000000}\selectfont\catcode`\^=\active\def^{\ifmmode\sp\else\^{}\fi}\catcode`\%=\active\def%{\%}$\mathdefault{\ensuremath{-}2}$}}%
\end{pgfscope}%
\begin{pgfscope}%
\pgfpathrectangle{\pgfqpoint{0.100000in}{0.100000in}}{\pgfqpoint{3.957178in}{3.957178in}}%
\pgfusepath{clip}%
\pgfsetbuttcap%
\pgfsetroundjoin%
\pgfsetlinewidth{0.501875pt}%
\definecolor{currentstroke}{rgb}{0.501961,0.501961,0.501961}%
\pgfsetstrokecolor{currentstroke}%
\pgfsetstrokeopacity{0.500000}%
\pgfsetdash{{1.850000pt}{0.800000pt}}{0.000000pt}%
\pgfpathmoveto{\pgfqpoint{2.132065in}{2.132065in}}%
\pgfusepath{stroke}%
\end{pgfscope}%
\begin{pgfscope}%
\pgfsetrectcap%
\pgfsetroundjoin%
\pgfsetlinewidth{0.803000pt}%
\definecolor{currentstroke}{rgb}{0.000000,0.000000,0.000000}%
\pgfsetstrokecolor{currentstroke}%
\pgfsetdash{}{0pt}%
\pgfpathmoveto{\pgfqpoint{1.123791in}{0.951596in}}%
\pgfpathlineto{\pgfqpoint{1.088948in}{0.920719in}}%
\pgfusepath{stroke}%
\end{pgfscope}%
\begin{pgfscope}%
\definecolor{textcolor}{rgb}{0.000000,0.000000,0.000000}%
\pgfsetstrokecolor{textcolor}%
\pgfsetfillcolor{textcolor}%
\pgftext[x=1.002286in,y=0.735253in,,top]{\color{textcolor}{\rmfamily\fontsize{10.000000}{12.000000}\selectfont\catcode`\^=\active\def^{\ifmmode\sp\else\^{}\fi}\catcode`\%=\active\def%{\%}$\mathdefault{\ensuremath{-}1}$}}%
\end{pgfscope}%
\begin{pgfscope}%
\pgfpathrectangle{\pgfqpoint{0.100000in}{0.100000in}}{\pgfqpoint{3.957178in}{3.957178in}}%
\pgfusepath{clip}%
\pgfsetbuttcap%
\pgfsetroundjoin%
\pgfsetlinewidth{0.501875pt}%
\definecolor{currentstroke}{rgb}{0.501961,0.501961,0.501961}%
\pgfsetstrokecolor{currentstroke}%
\pgfsetstrokeopacity{0.500000}%
\pgfsetdash{{1.850000pt}{0.800000pt}}{0.000000pt}%
\pgfpathmoveto{\pgfqpoint{2.132065in}{2.132065in}}%
\pgfusepath{stroke}%
\end{pgfscope}%
\begin{pgfscope}%
\pgfsetrectcap%
\pgfsetroundjoin%
\pgfsetlinewidth{0.803000pt}%
\definecolor{currentstroke}{rgb}{0.000000,0.000000,0.000000}%
\pgfsetstrokecolor{currentstroke}%
\pgfsetdash{}{0pt}%
\pgfpathmoveto{\pgfqpoint{1.461141in}{0.842327in}}%
\pgfpathlineto{\pgfqpoint{1.426791in}{0.810949in}}%
\pgfusepath{stroke}%
\end{pgfscope}%
\begin{pgfscope}%
\definecolor{textcolor}{rgb}{0.000000,0.000000,0.000000}%
\pgfsetstrokecolor{textcolor}%
\pgfsetfillcolor{textcolor}%
\pgftext[x=1.340417in,y=0.623655in,,top]{\color{textcolor}{\rmfamily\fontsize{10.000000}{12.000000}\selectfont\catcode`\^=\active\def^{\ifmmode\sp\else\^{}\fi}\catcode`\%=\active\def%{\%}$\mathdefault{0}$}}%
\end{pgfscope}%
\begin{pgfscope}%
\pgfpathrectangle{\pgfqpoint{0.100000in}{0.100000in}}{\pgfqpoint{3.957178in}{3.957178in}}%
\pgfusepath{clip}%
\pgfsetbuttcap%
\pgfsetroundjoin%
\pgfsetlinewidth{0.501875pt}%
\definecolor{currentstroke}{rgb}{0.501961,0.501961,0.501961}%
\pgfsetstrokecolor{currentstroke}%
\pgfsetstrokeopacity{0.500000}%
\pgfsetdash{{1.850000pt}{0.800000pt}}{0.000000pt}%
\pgfpathmoveto{\pgfqpoint{2.132065in}{2.132065in}}%
\pgfusepath{stroke}%
\end{pgfscope}%
\begin{pgfscope}%
\pgfsetrectcap%
\pgfsetroundjoin%
\pgfsetlinewidth{0.803000pt}%
\definecolor{currentstroke}{rgb}{0.000000,0.000000,0.000000}%
\pgfsetstrokecolor{currentstroke}%
\pgfsetdash{}{0pt}%
\pgfpathmoveto{\pgfqpoint{1.803983in}{0.731280in}}%
\pgfpathlineto{\pgfqpoint{1.770149in}{0.699388in}}%
\pgfusepath{stroke}%
\end{pgfscope}%
\begin{pgfscope}%
\definecolor{textcolor}{rgb}{0.000000,0.000000,0.000000}%
\pgfsetstrokecolor{textcolor}%
\pgfsetfillcolor{textcolor}%
\pgftext[x=1.684082in,y=0.510231in,,top]{\color{textcolor}{\rmfamily\fontsize{10.000000}{12.000000}\selectfont\catcode`\^=\active\def^{\ifmmode\sp\else\^{}\fi}\catcode`\%=\active\def%{\%}$\mathdefault{1}$}}%
\end{pgfscope}%
\begin{pgfscope}%
\pgfpathrectangle{\pgfqpoint{0.100000in}{0.100000in}}{\pgfqpoint{3.957178in}{3.957178in}}%
\pgfusepath{clip}%
\pgfsetbuttcap%
\pgfsetroundjoin%
\pgfsetlinewidth{0.501875pt}%
\definecolor{currentstroke}{rgb}{0.501961,0.501961,0.501961}%
\pgfsetstrokecolor{currentstroke}%
\pgfsetstrokeopacity{0.500000}%
\pgfsetdash{{1.850000pt}{0.800000pt}}{0.000000pt}%
\pgfpathmoveto{\pgfqpoint{2.132065in}{2.132065in}}%
\pgfusepath{stroke}%
\end{pgfscope}%
\begin{pgfscope}%
\pgfsetrectcap%
\pgfsetroundjoin%
\pgfsetlinewidth{0.803000pt}%
\definecolor{currentstroke}{rgb}{0.000000,0.000000,0.000000}%
\pgfsetstrokecolor{currentstroke}%
\pgfsetdash{}{0pt}%
\pgfpathmoveto{\pgfqpoint{2.152455in}{0.618409in}}%
\pgfpathlineto{\pgfqpoint{2.119157in}{0.585991in}}%
\pgfusepath{stroke}%
\end{pgfscope}%
\begin{pgfscope}%
\definecolor{textcolor}{rgb}{0.000000,0.000000,0.000000}%
\pgfsetstrokecolor{textcolor}%
\pgfsetfillcolor{textcolor}%
\pgftext[x=2.033420in,y=0.394935in,,top]{\color{textcolor}{\rmfamily\fontsize{10.000000}{12.000000}\selectfont\catcode`\^=\active\def^{\ifmmode\sp\else\^{}\fi}\catcode`\%=\active\def%{\%}$\mathdefault{2}$}}%
\end{pgfscope}%
\begin{pgfscope}%
\pgfpathrectangle{\pgfqpoint{0.100000in}{0.100000in}}{\pgfqpoint{3.957178in}{3.957178in}}%
\pgfusepath{clip}%
\pgfsetbuttcap%
\pgfsetroundjoin%
\pgfsetlinewidth{0.501875pt}%
\definecolor{currentstroke}{rgb}{0.501961,0.501961,0.501961}%
\pgfsetstrokecolor{currentstroke}%
\pgfsetstrokeopacity{0.500000}%
\pgfsetdash{{1.850000pt}{0.800000pt}}{0.000000pt}%
\pgfpathmoveto{\pgfqpoint{2.132065in}{2.132065in}}%
\pgfusepath{stroke}%
\end{pgfscope}%
\begin{pgfscope}%
\pgfsetrectcap%
\pgfsetroundjoin%
\pgfsetlinewidth{0.803000pt}%
\definecolor{currentstroke}{rgb}{0.000000,0.000000,0.000000}%
\pgfsetstrokecolor{currentstroke}%
\pgfsetdash{}{0pt}%
\pgfpathmoveto{\pgfqpoint{2.506694in}{0.503670in}}%
\pgfpathlineto{\pgfqpoint{2.473955in}{0.470712in}}%
\pgfusepath{stroke}%
\end{pgfscope}%
\begin{pgfscope}%
\definecolor{textcolor}{rgb}{0.000000,0.000000,0.000000}%
\pgfsetstrokecolor{textcolor}%
\pgfsetfillcolor{textcolor}%
\pgftext[x=2.388571in,y=0.277720in,,top]{\color{textcolor}{\rmfamily\fontsize{10.000000}{12.000000}\selectfont\catcode`\^=\active\def^{\ifmmode\sp\else\^{}\fi}\catcode`\%=\active\def%{\%}$\mathdefault{3}$}}%
\end{pgfscope}%
\begin{pgfscope}%
\definecolor{textcolor}{rgb}{0.000000,0.000000,0.000000}%
\pgfsetstrokecolor{textcolor}%
\pgfsetfillcolor{textcolor}%
\pgftext[x=1.210812in,y=0.376843in,,]{\color{textcolor}{\rmfamily\fontsize{10.000000}{12.000000}\selectfont\catcode`\^=\active\def^{\ifmmode\sp\else\^{}\fi}\catcode`\%=\active\def%{\%}X}}%
\end{pgfscope}%
\begin{pgfscope}%
\pgfsetrectcap%
\pgfsetroundjoin%
\pgfsetlinewidth{0.803000pt}%
\definecolor{currentstroke}{rgb}{0.000000,0.000000,0.000000}%
\pgfsetstrokecolor{currentstroke}%
\pgfsetdash{}{0pt}%
\pgfpathmoveto{\pgfqpoint{3.877604in}{1.701946in}}%
\pgfpathlineto{\pgfqpoint{2.645817in}{0.444059in}}%
\pgfusepath{stroke}%
\end{pgfscope}%
\begin{pgfscope}%
\pgfpathrectangle{\pgfqpoint{0.100000in}{0.100000in}}{\pgfqpoint{3.957178in}{3.957178in}}%
\pgfusepath{clip}%
\pgfsetbuttcap%
\pgfsetroundjoin%
\pgfsetlinewidth{0.501875pt}%
\definecolor{currentstroke}{rgb}{0.501961,0.501961,0.501961}%
\pgfsetstrokecolor{currentstroke}%
\pgfsetstrokeopacity{0.500000}%
\pgfsetdash{{1.850000pt}{0.800000pt}}{0.000000pt}%
\pgfpathmoveto{\pgfqpoint{2.132065in}{2.132065in}}%
\pgfusepath{stroke}%
\end{pgfscope}%
\begin{pgfscope}%
\pgfsetrectcap%
\pgfsetroundjoin%
\pgfsetlinewidth{0.803000pt}%
\definecolor{currentstroke}{rgb}{0.000000,0.000000,0.000000}%
\pgfsetstrokecolor{currentstroke}%
\pgfsetdash{}{0pt}%
\pgfpathmoveto{\pgfqpoint{2.711441in}{0.537379in}}%
\pgfpathlineto{\pgfqpoint{2.770264in}{0.518462in}}%
\pgfusepath{stroke}%
\end{pgfscope}%
\begin{pgfscope}%
\definecolor{textcolor}{rgb}{0.000000,0.000000,0.000000}%
\pgfsetstrokecolor{textcolor}%
\pgfsetfillcolor{textcolor}%
\pgftext[x=2.919323in,y=0.362310in,,top]{\color{textcolor}{\rmfamily\fontsize{10.000000}{12.000000}\selectfont\catcode`\^=\active\def^{\ifmmode\sp\else\^{}\fi}\catcode`\%=\active\def%{\%}$\mathdefault{\ensuremath{-}3}$}}%
\end{pgfscope}%
\begin{pgfscope}%
\pgfpathrectangle{\pgfqpoint{0.100000in}{0.100000in}}{\pgfqpoint{3.957178in}{3.957178in}}%
\pgfusepath{clip}%
\pgfsetbuttcap%
\pgfsetroundjoin%
\pgfsetlinewidth{0.501875pt}%
\definecolor{currentstroke}{rgb}{0.501961,0.501961,0.501961}%
\pgfsetstrokecolor{currentstroke}%
\pgfsetstrokeopacity{0.500000}%
\pgfsetdash{{1.850000pt}{0.800000pt}}{0.000000pt}%
\pgfpathmoveto{\pgfqpoint{2.132065in}{2.132065in}}%
\pgfusepath{stroke}%
\end{pgfscope}%
\begin{pgfscope}%
\pgfsetrectcap%
\pgfsetroundjoin%
\pgfsetlinewidth{0.803000pt}%
\definecolor{currentstroke}{rgb}{0.000000,0.000000,0.000000}%
\pgfsetstrokecolor{currentstroke}%
\pgfsetdash{}{0pt}%
\pgfpathmoveto{\pgfqpoint{2.903900in}{0.733546in}}%
\pgfpathlineto{\pgfqpoint{2.962152in}{0.715157in}}%
\pgfusepath{stroke}%
\end{pgfscope}%
\begin{pgfscope}%
\definecolor{textcolor}{rgb}{0.000000,0.000000,0.000000}%
\pgfsetstrokecolor{textcolor}%
\pgfsetfillcolor{textcolor}%
\pgftext[x=3.109193in,y=0.561292in,,top]{\color{textcolor}{\rmfamily\fontsize{10.000000}{12.000000}\selectfont\catcode`\^=\active\def^{\ifmmode\sp\else\^{}\fi}\catcode`\%=\active\def%{\%}$\mathdefault{\ensuremath{-}2}$}}%
\end{pgfscope}%
\begin{pgfscope}%
\pgfpathrectangle{\pgfqpoint{0.100000in}{0.100000in}}{\pgfqpoint{3.957178in}{3.957178in}}%
\pgfusepath{clip}%
\pgfsetbuttcap%
\pgfsetroundjoin%
\pgfsetlinewidth{0.501875pt}%
\definecolor{currentstroke}{rgb}{0.501961,0.501961,0.501961}%
\pgfsetstrokecolor{currentstroke}%
\pgfsetstrokeopacity{0.500000}%
\pgfsetdash{{1.850000pt}{0.800000pt}}{0.000000pt}%
\pgfpathmoveto{\pgfqpoint{2.132065in}{2.132065in}}%
\pgfusepath{stroke}%
\end{pgfscope}%
\begin{pgfscope}%
\pgfsetrectcap%
\pgfsetroundjoin%
\pgfsetlinewidth{0.803000pt}%
\definecolor{currentstroke}{rgb}{0.000000,0.000000,0.000000}%
\pgfsetstrokecolor{currentstroke}%
\pgfsetdash{}{0pt}%
\pgfpathmoveto{\pgfqpoint{3.091019in}{0.924270in}}%
\pgfpathlineto{\pgfqpoint{3.148708in}{0.906388in}}%
\pgfusepath{stroke}%
\end{pgfscope}%
\begin{pgfscope}%
\definecolor{textcolor}{rgb}{0.000000,0.000000,0.000000}%
\pgfsetstrokecolor{textcolor}%
\pgfsetfillcolor{textcolor}%
\pgftext[x=3.293787in,y=0.754743in,,top]{\color{textcolor}{\rmfamily\fontsize{10.000000}{12.000000}\selectfont\catcode`\^=\active\def^{\ifmmode\sp\else\^{}\fi}\catcode`\%=\active\def%{\%}$\mathdefault{\ensuremath{-}1}$}}%
\end{pgfscope}%
\begin{pgfscope}%
\pgfpathrectangle{\pgfqpoint{0.100000in}{0.100000in}}{\pgfqpoint{3.957178in}{3.957178in}}%
\pgfusepath{clip}%
\pgfsetbuttcap%
\pgfsetroundjoin%
\pgfsetlinewidth{0.501875pt}%
\definecolor{currentstroke}{rgb}{0.501961,0.501961,0.501961}%
\pgfsetstrokecolor{currentstroke}%
\pgfsetstrokeopacity{0.500000}%
\pgfsetdash{{1.850000pt}{0.800000pt}}{0.000000pt}%
\pgfpathmoveto{\pgfqpoint{2.132065in}{2.132065in}}%
\pgfusepath{stroke}%
\end{pgfscope}%
\begin{pgfscope}%
\pgfsetrectcap%
\pgfsetroundjoin%
\pgfsetlinewidth{0.803000pt}%
\definecolor{currentstroke}{rgb}{0.000000,0.000000,0.000000}%
\pgfsetstrokecolor{currentstroke}%
\pgfsetdash{}{0pt}%
\pgfpathmoveto{\pgfqpoint{3.273017in}{1.109774in}}%
\pgfpathlineto{\pgfqpoint{3.330153in}{1.092379in}}%
\pgfusepath{stroke}%
\end{pgfscope}%
\begin{pgfscope}%
\definecolor{textcolor}{rgb}{0.000000,0.000000,0.000000}%
\pgfsetstrokecolor{textcolor}%
\pgfsetfillcolor{textcolor}%
\pgftext[x=3.473320in,y=0.942891in,,top]{\color{textcolor}{\rmfamily\fontsize{10.000000}{12.000000}\selectfont\catcode`\^=\active\def^{\ifmmode\sp\else\^{}\fi}\catcode`\%=\active\def%{\%}$\mathdefault{0}$}}%
\end{pgfscope}%
\begin{pgfscope}%
\pgfpathrectangle{\pgfqpoint{0.100000in}{0.100000in}}{\pgfqpoint{3.957178in}{3.957178in}}%
\pgfusepath{clip}%
\pgfsetbuttcap%
\pgfsetroundjoin%
\pgfsetlinewidth{0.501875pt}%
\definecolor{currentstroke}{rgb}{0.501961,0.501961,0.501961}%
\pgfsetstrokecolor{currentstroke}%
\pgfsetstrokeopacity{0.500000}%
\pgfsetdash{{1.850000pt}{0.800000pt}}{0.000000pt}%
\pgfpathmoveto{\pgfqpoint{2.132065in}{2.132065in}}%
\pgfusepath{stroke}%
\end{pgfscope}%
\begin{pgfscope}%
\pgfsetrectcap%
\pgfsetroundjoin%
\pgfsetlinewidth{0.803000pt}%
\definecolor{currentstroke}{rgb}{0.000000,0.000000,0.000000}%
\pgfsetstrokecolor{currentstroke}%
\pgfsetdash{}{0pt}%
\pgfpathmoveto{\pgfqpoint{3.450102in}{1.290271in}}%
\pgfpathlineto{\pgfqpoint{3.506693in}{1.273343in}}%
\pgfusepath{stroke}%
\end{pgfscope}%
\begin{pgfscope}%
\definecolor{textcolor}{rgb}{0.000000,0.000000,0.000000}%
\pgfsetstrokecolor{textcolor}%
\pgfsetfillcolor{textcolor}%
\pgftext[x=3.647998in,y=1.125952in,,top]{\color{textcolor}{\rmfamily\fontsize{10.000000}{12.000000}\selectfont\catcode`\^=\active\def^{\ifmmode\sp\else\^{}\fi}\catcode`\%=\active\def%{\%}$\mathdefault{1}$}}%
\end{pgfscope}%
\begin{pgfscope}%
\pgfpathrectangle{\pgfqpoint{0.100000in}{0.100000in}}{\pgfqpoint{3.957178in}{3.957178in}}%
\pgfusepath{clip}%
\pgfsetbuttcap%
\pgfsetroundjoin%
\pgfsetlinewidth{0.501875pt}%
\definecolor{currentstroke}{rgb}{0.501961,0.501961,0.501961}%
\pgfsetstrokecolor{currentstroke}%
\pgfsetstrokeopacity{0.500000}%
\pgfsetdash{{1.850000pt}{0.800000pt}}{0.000000pt}%
\pgfpathmoveto{\pgfqpoint{2.132065in}{2.132065in}}%
\pgfusepath{stroke}%
\end{pgfscope}%
\begin{pgfscope}%
\pgfsetrectcap%
\pgfsetroundjoin%
\pgfsetlinewidth{0.803000pt}%
\definecolor{currentstroke}{rgb}{0.000000,0.000000,0.000000}%
\pgfsetstrokecolor{currentstroke}%
\pgfsetdash{}{0pt}%
\pgfpathmoveto{\pgfqpoint{3.622470in}{1.465960in}}%
\pgfpathlineto{\pgfqpoint{3.678525in}{1.449480in}}%
\pgfusepath{stroke}%
\end{pgfscope}%
\begin{pgfscope}%
\definecolor{textcolor}{rgb}{0.000000,0.000000,0.000000}%
\pgfsetstrokecolor{textcolor}%
\pgfsetfillcolor{textcolor}%
\pgftext[x=3.818015in,y=1.304128in,,top]{\color{textcolor}{\rmfamily\fontsize{10.000000}{12.000000}\selectfont\catcode`\^=\active\def^{\ifmmode\sp\else\^{}\fi}\catcode`\%=\active\def%{\%}$\mathdefault{2}$}}%
\end{pgfscope}%
\begin{pgfscope}%
\pgfpathrectangle{\pgfqpoint{0.100000in}{0.100000in}}{\pgfqpoint{3.957178in}{3.957178in}}%
\pgfusepath{clip}%
\pgfsetbuttcap%
\pgfsetroundjoin%
\pgfsetlinewidth{0.501875pt}%
\definecolor{currentstroke}{rgb}{0.501961,0.501961,0.501961}%
\pgfsetstrokecolor{currentstroke}%
\pgfsetstrokeopacity{0.500000}%
\pgfsetdash{{1.850000pt}{0.800000pt}}{0.000000pt}%
\pgfpathmoveto{\pgfqpoint{2.132065in}{2.132065in}}%
\pgfusepath{stroke}%
\end{pgfscope}%
\begin{pgfscope}%
\pgfsetrectcap%
\pgfsetroundjoin%
\pgfsetlinewidth{0.803000pt}%
\definecolor{currentstroke}{rgb}{0.000000,0.000000,0.000000}%
\pgfsetstrokecolor{currentstroke}%
\pgfsetdash{}{0pt}%
\pgfpathmoveto{\pgfqpoint{3.790307in}{1.637031in}}%
\pgfpathlineto{\pgfqpoint{3.845834in}{1.620981in}}%
\pgfusepath{stroke}%
\end{pgfscope}%
\begin{pgfscope}%
\definecolor{textcolor}{rgb}{0.000000,0.000000,0.000000}%
\pgfsetstrokecolor{textcolor}%
\pgfsetfillcolor{textcolor}%
\pgftext[x=3.983556in,y=1.477612in,,top]{\color{textcolor}{\rmfamily\fontsize{10.000000}{12.000000}\selectfont\catcode`\^=\active\def^{\ifmmode\sp\else\^{}\fi}\catcode`\%=\active\def%{\%}$\mathdefault{3}$}}%
\end{pgfscope}%
\begin{pgfscope}%
\definecolor{textcolor}{rgb}{0.000000,0.000000,0.000000}%
\pgfsetstrokecolor{textcolor}%
\pgfsetfillcolor{textcolor}%
\pgftext[x=3.685310in,y=0.751734in,,]{\color{textcolor}{\rmfamily\fontsize{10.000000}{12.000000}\selectfont\catcode`\^=\active\def^{\ifmmode\sp\else\^{}\fi}\catcode`\%=\active\def%{\%}Y}}%
\end{pgfscope}%
\begin{pgfscope}%
\pgfsetrectcap%
\pgfsetroundjoin%
\pgfsetlinewidth{0.803000pt}%
\definecolor{currentstroke}{rgb}{0.000000,0.000000,0.000000}%
\pgfsetstrokecolor{currentstroke}%
\pgfsetdash{}{0pt}%
\pgfpathmoveto{\pgfqpoint{3.877604in}{1.701946in}}%
\pgfpathlineto{\pgfqpoint{3.942496in}{3.063277in}}%
\pgfusepath{stroke}%
\end{pgfscope}%
\begin{pgfscope}%
\pgfsetbuttcap%
\pgfsetroundjoin%
\pgfsetlinewidth{0.501875pt}%
\definecolor{currentstroke}{rgb}{0.501961,0.501961,0.501961}%
\pgfsetstrokecolor{currentstroke}%
\pgfsetstrokeopacity{0.500000}%
\pgfsetdash{{1.850000pt}{0.800000pt}}{0.000000pt}%
\pgfpathmoveto{\pgfqpoint{2.132065in}{2.132065in}}%
\pgfusepath{stroke}%
\end{pgfscope}%
\begin{pgfscope}%
\pgfsetrectcap%
\pgfsetroundjoin%
\pgfsetlinewidth{0.803000pt}%
\definecolor{currentstroke}{rgb}{0.000000,0.000000,0.000000}%
\pgfsetstrokecolor{currentstroke}%
\pgfsetdash{}{0pt}%
\pgfpathmoveto{\pgfqpoint{3.860268in}{1.730240in}}%
\pgfpathlineto{\pgfqpoint{3.915612in}{1.714384in}}%
\pgfusepath{stroke}%
\end{pgfscope}%
\begin{pgfscope}%
\definecolor{textcolor}{rgb}{0.000000,0.000000,0.000000}%
\pgfsetstrokecolor{textcolor}%
\pgfsetfillcolor{textcolor}%
\pgftext[x=4.143509in,y=1.762314in,,top]{\color{textcolor}{\rmfamily\fontsize{10.000000}{12.000000}\selectfont\catcode`\^=\active\def^{\ifmmode\sp\else\^{}\fi}\catcode`\%=\active\def%{\%}$\mathdefault{\ensuremath{-}0.8}$}}%
\end{pgfscope}%
\begin{pgfscope}%
\pgfsetbuttcap%
\pgfsetroundjoin%
\pgfsetlinewidth{0.501875pt}%
\definecolor{currentstroke}{rgb}{0.501961,0.501961,0.501961}%
\pgfsetstrokecolor{currentstroke}%
\pgfsetstrokeopacity{0.500000}%
\pgfsetdash{{1.850000pt}{0.800000pt}}{0.000000pt}%
\pgfpathmoveto{\pgfqpoint{2.132065in}{2.132065in}}%
\pgfusepath{stroke}%
\end{pgfscope}%
\begin{pgfscope}%
\pgfsetrectcap%
\pgfsetroundjoin%
\pgfsetlinewidth{0.803000pt}%
\definecolor{currentstroke}{rgb}{0.000000,0.000000,0.000000}%
\pgfsetstrokecolor{currentstroke}%
\pgfsetdash{}{0pt}%
\pgfpathmoveto{\pgfqpoint{3.866920in}{1.871296in}}%
\pgfpathlineto{\pgfqpoint{3.922486in}{1.855555in}}%
\pgfusepath{stroke}%
\end{pgfscope}%
\begin{pgfscope}%
\definecolor{textcolor}{rgb}{0.000000,0.000000,0.000000}%
\pgfsetstrokecolor{textcolor}%
\pgfsetfillcolor{textcolor}%
\pgftext[x=4.151232in,y=1.903137in,,top]{\color{textcolor}{\rmfamily\fontsize{10.000000}{12.000000}\selectfont\catcode`\^=\active\def^{\ifmmode\sp\else\^{}\fi}\catcode`\%=\active\def%{\%}$\mathdefault{\ensuremath{-}0.6}$}}%
\end{pgfscope}%
\begin{pgfscope}%
\pgfsetbuttcap%
\pgfsetroundjoin%
\pgfsetlinewidth{0.501875pt}%
\definecolor{currentstroke}{rgb}{0.501961,0.501961,0.501961}%
\pgfsetstrokecolor{currentstroke}%
\pgfsetstrokeopacity{0.500000}%
\pgfsetdash{{1.850000pt}{0.800000pt}}{0.000000pt}%
\pgfpathmoveto{\pgfqpoint{2.132065in}{2.132065in}}%
\pgfusepath{stroke}%
\end{pgfscope}%
\begin{pgfscope}%
\pgfsetrectcap%
\pgfsetroundjoin%
\pgfsetlinewidth{0.803000pt}%
\definecolor{currentstroke}{rgb}{0.000000,0.000000,0.000000}%
\pgfsetstrokecolor{currentstroke}%
\pgfsetdash{}{0pt}%
\pgfpathmoveto{\pgfqpoint{3.873623in}{2.013443in}}%
\pgfpathlineto{\pgfqpoint{3.929413in}{1.997819in}}%
\pgfusepath{stroke}%
\end{pgfscope}%
\begin{pgfscope}%
\definecolor{textcolor}{rgb}{0.000000,0.000000,0.000000}%
\pgfsetstrokecolor{textcolor}%
\pgfsetfillcolor{textcolor}%
\pgftext[x=4.159014in,y=2.045046in,,top]{\color{textcolor}{\rmfamily\fontsize{10.000000}{12.000000}\selectfont\catcode`\^=\active\def^{\ifmmode\sp\else\^{}\fi}\catcode`\%=\active\def%{\%}$\mathdefault{\ensuremath{-}0.4}$}}%
\end{pgfscope}%
\begin{pgfscope}%
\pgfsetbuttcap%
\pgfsetroundjoin%
\pgfsetlinewidth{0.501875pt}%
\definecolor{currentstroke}{rgb}{0.501961,0.501961,0.501961}%
\pgfsetstrokecolor{currentstroke}%
\pgfsetstrokeopacity{0.500000}%
\pgfsetdash{{1.850000pt}{0.800000pt}}{0.000000pt}%
\pgfpathmoveto{\pgfqpoint{2.132065in}{2.132065in}}%
\pgfusepath{stroke}%
\end{pgfscope}%
\begin{pgfscope}%
\pgfsetrectcap%
\pgfsetroundjoin%
\pgfsetlinewidth{0.803000pt}%
\definecolor{currentstroke}{rgb}{0.000000,0.000000,0.000000}%
\pgfsetstrokecolor{currentstroke}%
\pgfsetdash{}{0pt}%
\pgfpathmoveto{\pgfqpoint{3.880378in}{2.156692in}}%
\pgfpathlineto{\pgfqpoint{3.936393in}{2.141188in}}%
\pgfusepath{stroke}%
\end{pgfscope}%
\begin{pgfscope}%
\definecolor{textcolor}{rgb}{0.000000,0.000000,0.000000}%
\pgfsetstrokecolor{textcolor}%
\pgfsetfillcolor{textcolor}%
\pgftext[x=4.166856in,y=2.188053in,,top]{\color{textcolor}{\rmfamily\fontsize{10.000000}{12.000000}\selectfont\catcode`\^=\active\def^{\ifmmode\sp\else\^{}\fi}\catcode`\%=\active\def%{\%}$\mathdefault{\ensuremath{-}0.2}$}}%
\end{pgfscope}%
\begin{pgfscope}%
\pgfsetbuttcap%
\pgfsetroundjoin%
\pgfsetlinewidth{0.501875pt}%
\definecolor{currentstroke}{rgb}{0.501961,0.501961,0.501961}%
\pgfsetstrokecolor{currentstroke}%
\pgfsetstrokeopacity{0.500000}%
\pgfsetdash{{1.850000pt}{0.800000pt}}{0.000000pt}%
\pgfpathmoveto{\pgfqpoint{2.132065in}{2.132065in}}%
\pgfusepath{stroke}%
\end{pgfscope}%
\begin{pgfscope}%
\pgfsetrectcap%
\pgfsetroundjoin%
\pgfsetlinewidth{0.803000pt}%
\definecolor{currentstroke}{rgb}{0.000000,0.000000,0.000000}%
\pgfsetstrokecolor{currentstroke}%
\pgfsetdash{}{0pt}%
\pgfpathmoveto{\pgfqpoint{3.887186in}{2.301057in}}%
\pgfpathlineto{\pgfqpoint{3.943428in}{2.285674in}}%
\pgfusepath{stroke}%
\end{pgfscope}%
\begin{pgfscope}%
\definecolor{textcolor}{rgb}{0.000000,0.000000,0.000000}%
\pgfsetstrokecolor{textcolor}%
\pgfsetfillcolor{textcolor}%
\pgftext[x=4.174759in,y=2.332171in,,top]{\color{textcolor}{\rmfamily\fontsize{10.000000}{12.000000}\selectfont\catcode`\^=\active\def^{\ifmmode\sp\else\^{}\fi}\catcode`\%=\active\def%{\%}$\mathdefault{0.0}$}}%
\end{pgfscope}%
\begin{pgfscope}%
\pgfsetbuttcap%
\pgfsetroundjoin%
\pgfsetlinewidth{0.501875pt}%
\definecolor{currentstroke}{rgb}{0.501961,0.501961,0.501961}%
\pgfsetstrokecolor{currentstroke}%
\pgfsetstrokeopacity{0.500000}%
\pgfsetdash{{1.850000pt}{0.800000pt}}{0.000000pt}%
\pgfpathmoveto{\pgfqpoint{2.132065in}{2.132065in}}%
\pgfusepath{stroke}%
\end{pgfscope}%
\begin{pgfscope}%
\pgfsetrectcap%
\pgfsetroundjoin%
\pgfsetlinewidth{0.803000pt}%
\definecolor{currentstroke}{rgb}{0.000000,0.000000,0.000000}%
\pgfsetstrokecolor{currentstroke}%
\pgfsetdash{}{0pt}%
\pgfpathmoveto{\pgfqpoint{3.894048in}{2.446550in}}%
\pgfpathlineto{\pgfqpoint{3.950518in}{2.431292in}}%
\pgfusepath{stroke}%
\end{pgfscope}%
\begin{pgfscope}%
\definecolor{textcolor}{rgb}{0.000000,0.000000,0.000000}%
\pgfsetstrokecolor{textcolor}%
\pgfsetfillcolor{textcolor}%
\pgftext[x=4.182724in,y=2.477413in,,top]{\color{textcolor}{\rmfamily\fontsize{10.000000}{12.000000}\selectfont\catcode`\^=\active\def^{\ifmmode\sp\else\^{}\fi}\catcode`\%=\active\def%{\%}$\mathdefault{0.2}$}}%
\end{pgfscope}%
\begin{pgfscope}%
\pgfsetbuttcap%
\pgfsetroundjoin%
\pgfsetlinewidth{0.501875pt}%
\definecolor{currentstroke}{rgb}{0.501961,0.501961,0.501961}%
\pgfsetstrokecolor{currentstroke}%
\pgfsetstrokeopacity{0.500000}%
\pgfsetdash{{1.850000pt}{0.800000pt}}{0.000000pt}%
\pgfpathmoveto{\pgfqpoint{2.132065in}{2.132065in}}%
\pgfusepath{stroke}%
\end{pgfscope}%
\begin{pgfscope}%
\pgfsetrectcap%
\pgfsetroundjoin%
\pgfsetlinewidth{0.803000pt}%
\definecolor{currentstroke}{rgb}{0.000000,0.000000,0.000000}%
\pgfsetstrokecolor{currentstroke}%
\pgfsetdash{}{0pt}%
\pgfpathmoveto{\pgfqpoint{3.900963in}{2.593186in}}%
\pgfpathlineto{\pgfqpoint{3.957664in}{2.578054in}}%
\pgfusepath{stroke}%
\end{pgfscope}%
\begin{pgfscope}%
\definecolor{textcolor}{rgb}{0.000000,0.000000,0.000000}%
\pgfsetstrokecolor{textcolor}%
\pgfsetfillcolor{textcolor}%
\pgftext[x=4.190752in,y=2.623791in,,top]{\color{textcolor}{\rmfamily\fontsize{10.000000}{12.000000}\selectfont\catcode`\^=\active\def^{\ifmmode\sp\else\^{}\fi}\catcode`\%=\active\def%{\%}$\mathdefault{0.4}$}}%
\end{pgfscope}%
\begin{pgfscope}%
\pgfsetbuttcap%
\pgfsetroundjoin%
\pgfsetlinewidth{0.501875pt}%
\definecolor{currentstroke}{rgb}{0.501961,0.501961,0.501961}%
\pgfsetstrokecolor{currentstroke}%
\pgfsetstrokeopacity{0.500000}%
\pgfsetdash{{1.850000pt}{0.800000pt}}{0.000000pt}%
\pgfpathmoveto{\pgfqpoint{2.132065in}{2.132065in}}%
\pgfusepath{stroke}%
\end{pgfscope}%
\begin{pgfscope}%
\pgfsetrectcap%
\pgfsetroundjoin%
\pgfsetlinewidth{0.803000pt}%
\definecolor{currentstroke}{rgb}{0.000000,0.000000,0.000000}%
\pgfsetstrokecolor{currentstroke}%
\pgfsetdash{}{0pt}%
\pgfpathmoveto{\pgfqpoint{3.907932in}{2.740977in}}%
\pgfpathlineto{\pgfqpoint{3.964866in}{2.725974in}}%
\pgfusepath{stroke}%
\end{pgfscope}%
\begin{pgfscope}%
\definecolor{textcolor}{rgb}{0.000000,0.000000,0.000000}%
\pgfsetstrokecolor{textcolor}%
\pgfsetfillcolor{textcolor}%
\pgftext[x=4.198842in,y=2.771321in,,top]{\color{textcolor}{\rmfamily\fontsize{10.000000}{12.000000}\selectfont\catcode`\^=\active\def^{\ifmmode\sp\else\^{}\fi}\catcode`\%=\active\def%{\%}$\mathdefault{0.6}$}}%
\end{pgfscope}%
\begin{pgfscope}%
\pgfsetbuttcap%
\pgfsetroundjoin%
\pgfsetlinewidth{0.501875pt}%
\definecolor{currentstroke}{rgb}{0.501961,0.501961,0.501961}%
\pgfsetstrokecolor{currentstroke}%
\pgfsetstrokeopacity{0.500000}%
\pgfsetdash{{1.850000pt}{0.800000pt}}{0.000000pt}%
\pgfpathmoveto{\pgfqpoint{2.132065in}{2.132065in}}%
\pgfusepath{stroke}%
\end{pgfscope}%
\begin{pgfscope}%
\pgfsetrectcap%
\pgfsetroundjoin%
\pgfsetlinewidth{0.803000pt}%
\definecolor{currentstroke}{rgb}{0.000000,0.000000,0.000000}%
\pgfsetstrokecolor{currentstroke}%
\pgfsetdash{}{0pt}%
\pgfpathmoveto{\pgfqpoint{3.914957in}{2.889937in}}%
\pgfpathlineto{\pgfqpoint{3.972125in}{2.875066in}}%
\pgfusepath{stroke}%
\end{pgfscope}%
\begin{pgfscope}%
\definecolor{textcolor}{rgb}{0.000000,0.000000,0.000000}%
\pgfsetstrokecolor{textcolor}%
\pgfsetfillcolor{textcolor}%
\pgftext[x=4.206996in,y=2.920014in,,top]{\color{textcolor}{\rmfamily\fontsize{10.000000}{12.000000}\selectfont\catcode`\^=\active\def^{\ifmmode\sp\else\^{}\fi}\catcode`\%=\active\def%{\%}$\mathdefault{0.8}$}}%
\end{pgfscope}%
\begin{pgfscope}%
\pgfsetbuttcap%
\pgfsetroundjoin%
\pgfsetlinewidth{0.501875pt}%
\definecolor{currentstroke}{rgb}{0.501961,0.501961,0.501961}%
\pgfsetstrokecolor{currentstroke}%
\pgfsetstrokeopacity{0.500000}%
\pgfsetdash{{1.850000pt}{0.800000pt}}{0.000000pt}%
\pgfpathmoveto{\pgfqpoint{2.132065in}{2.132065in}}%
\pgfusepath{stroke}%
\end{pgfscope}%
\begin{pgfscope}%
\pgfsetrectcap%
\pgfsetroundjoin%
\pgfsetlinewidth{0.803000pt}%
\definecolor{currentstroke}{rgb}{0.000000,0.000000,0.000000}%
\pgfsetstrokecolor{currentstroke}%
\pgfsetdash{}{0pt}%
\pgfpathmoveto{\pgfqpoint{3.922037in}{3.040080in}}%
\pgfpathlineto{\pgfqpoint{3.979442in}{3.025344in}}%
\pgfusepath{stroke}%
\end{pgfscope}%
\begin{pgfscope}%
\definecolor{textcolor}{rgb}{0.000000,0.000000,0.000000}%
\pgfsetstrokecolor{textcolor}%
\pgfsetfillcolor{textcolor}%
\pgftext[x=4.215215in,y=3.069886in,,top]{\color{textcolor}{\rmfamily\fontsize{10.000000}{12.000000}\selectfont\catcode`\^=\active\def^{\ifmmode\sp\else\^{}\fi}\catcode`\%=\active\def%{\%}$\mathdefault{1.0}$}}%
\end{pgfscope}%
\begin{pgfscope}%
\definecolor{textcolor}{rgb}{0.000000,0.000000,0.000000}%
\pgfsetstrokecolor{textcolor}%
\pgfsetfillcolor{textcolor}%
\pgftext[x=4.492994in,y=2.448369in,,]{\color{textcolor}{\rmfamily\fontsize{10.000000}{12.000000}\selectfont\catcode`\^=\active\def^{\ifmmode\sp\else\^{}\fi}\catcode`\%=\active\def%{\%}Z}}%
\end{pgfscope}%
\begin{pgfscope}%
\pgfpathrectangle{\pgfqpoint{0.100000in}{0.100000in}}{\pgfqpoint{3.957178in}{3.957178in}}%
\pgfusepath{clip}%
\pgfsetbuttcap%
\pgfsetroundjoin%
\definecolor{currentfill}{rgb}{0.050383,0.029803,0.527975}%
\pgfsetfillcolor{currentfill}%
\pgfsetlinewidth{0.000000pt}%
\definecolor{currentstroke}{rgb}{0.000000,0.000000,0.000000}%
\pgfsetstrokecolor{currentstroke}%
\pgfsetdash{}{0pt}%
\pgfpathmoveto{\pgfqpoint{1.671443in}{2.260097in}}%
\pgfpathlineto{\pgfqpoint{1.695477in}{2.304468in}}%
\pgfpathlineto{\pgfqpoint{1.810215in}{2.231164in}}%
\pgfpathlineto{\pgfqpoint{1.671443in}{2.260097in}}%
\pgfpathclose%
\pgfusepath{fill}%
\end{pgfscope}%
\begin{pgfscope}%
\pgfpathrectangle{\pgfqpoint{0.100000in}{0.100000in}}{\pgfqpoint{3.957178in}{3.957178in}}%
\pgfusepath{clip}%
\pgfsetbuttcap%
\pgfsetroundjoin%
\definecolor{currentfill}{rgb}{0.193374,0.018354,0.590330}%
\pgfsetfillcolor{currentfill}%
\pgfsetlinewidth{0.000000pt}%
\definecolor{currentstroke}{rgb}{0.000000,0.000000,0.000000}%
\pgfsetstrokecolor{currentstroke}%
\pgfsetdash{}{0pt}%
\pgfpathmoveto{\pgfqpoint{1.810215in}{2.231164in}}%
\pgfpathlineto{\pgfqpoint{1.695477in}{2.304468in}}%
\pgfpathlineto{\pgfqpoint{1.804775in}{2.437869in}}%
\pgfpathlineto{\pgfqpoint{1.810215in}{2.231164in}}%
\pgfpathclose%
\pgfusepath{fill}%
\end{pgfscope}%
\begin{pgfscope}%
\pgfpathrectangle{\pgfqpoint{0.100000in}{0.100000in}}{\pgfqpoint{3.957178in}{3.957178in}}%
\pgfusepath{clip}%
\pgfsetbuttcap%
\pgfsetroundjoin%
\definecolor{currentfill}{rgb}{0.299855,0.009561,0.631624}%
\pgfsetfillcolor{currentfill}%
\pgfsetlinewidth{0.000000pt}%
\definecolor{currentstroke}{rgb}{0.000000,0.000000,0.000000}%
\pgfsetstrokecolor{currentstroke}%
\pgfsetdash{}{0pt}%
\pgfpathmoveto{\pgfqpoint{1.804775in}{2.437869in}}%
\pgfpathlineto{\pgfqpoint{1.938957in}{2.483638in}}%
\pgfpathlineto{\pgfqpoint{1.810215in}{2.231164in}}%
\pgfpathlineto{\pgfqpoint{1.804775in}{2.437869in}}%
\pgfpathclose%
\pgfusepath{fill}%
\end{pgfscope}%
\begin{pgfscope}%
\pgfpathrectangle{\pgfqpoint{0.100000in}{0.100000in}}{\pgfqpoint{3.957178in}{3.957178in}}%
\pgfusepath{clip}%
\pgfsetbuttcap%
\pgfsetroundjoin%
\definecolor{currentfill}{rgb}{0.267703,0.012716,0.620346}%
\pgfsetfillcolor{currentfill}%
\pgfsetlinewidth{0.000000pt}%
\definecolor{currentstroke}{rgb}{0.000000,0.000000,0.000000}%
\pgfsetstrokecolor{currentstroke}%
\pgfsetdash{}{0pt}%
\pgfpathmoveto{\pgfqpoint{1.695477in}{2.304468in}}%
\pgfpathlineto{\pgfqpoint{1.671443in}{2.260097in}}%
\pgfpathlineto{\pgfqpoint{1.705806in}{2.473649in}}%
\pgfpathlineto{\pgfqpoint{1.695477in}{2.304468in}}%
\pgfpathclose%
\pgfusepath{fill}%
\end{pgfscope}%
\begin{pgfscope}%
\pgfpathrectangle{\pgfqpoint{0.100000in}{0.100000in}}{\pgfqpoint{3.957178in}{3.957178in}}%
\pgfusepath{clip}%
\pgfsetbuttcap%
\pgfsetroundjoin%
\definecolor{currentfill}{rgb}{0.374897,0.002245,0.651876}%
\pgfsetfillcolor{currentfill}%
\pgfsetlinewidth{0.000000pt}%
\definecolor{currentstroke}{rgb}{0.000000,0.000000,0.000000}%
\pgfsetstrokecolor{currentstroke}%
\pgfsetdash{}{0pt}%
\pgfpathmoveto{\pgfqpoint{1.705806in}{2.473649in}}%
\pgfpathlineto{\pgfqpoint{1.804775in}{2.437869in}}%
\pgfpathlineto{\pgfqpoint{1.695477in}{2.304468in}}%
\pgfpathlineto{\pgfqpoint{1.705806in}{2.473649in}}%
\pgfpathclose%
\pgfusepath{fill}%
\end{pgfscope}%
\begin{pgfscope}%
\pgfpathrectangle{\pgfqpoint{0.100000in}{0.100000in}}{\pgfqpoint{3.957178in}{3.957178in}}%
\pgfusepath{clip}%
\pgfsetbuttcap%
\pgfsetroundjoin%
\definecolor{currentfill}{rgb}{0.435734,0.001127,0.659797}%
\pgfsetfillcolor{currentfill}%
\pgfsetlinewidth{0.000000pt}%
\definecolor{currentstroke}{rgb}{0.000000,0.000000,0.000000}%
\pgfsetstrokecolor{currentstroke}%
\pgfsetdash{}{0pt}%
\pgfpathmoveto{\pgfqpoint{1.810215in}{2.231164in}}%
\pgfpathlineto{\pgfqpoint{1.938957in}{2.483638in}}%
\pgfpathlineto{\pgfqpoint{2.401955in}{2.671792in}}%
\pgfpathlineto{\pgfqpoint{1.810215in}{2.231164in}}%
\pgfpathclose%
\pgfusepath{fill}%
\end{pgfscope}%
\begin{pgfscope}%
\pgfpathrectangle{\pgfqpoint{0.100000in}{0.100000in}}{\pgfqpoint{3.957178in}{3.957178in}}%
\pgfusepath{clip}%
\pgfsetbuttcap%
\pgfsetroundjoin%
\definecolor{currentfill}{rgb}{0.500678,0.014055,0.657088}%
\pgfsetfillcolor{currentfill}%
\pgfsetlinewidth{0.000000pt}%
\definecolor{currentstroke}{rgb}{0.000000,0.000000,0.000000}%
\pgfsetstrokecolor{currentstroke}%
\pgfsetdash{}{0pt}%
\pgfpathmoveto{\pgfqpoint{1.862002in}{2.548801in}}%
\pgfpathlineto{\pgfqpoint{1.938957in}{2.483638in}}%
\pgfpathlineto{\pgfqpoint{1.804775in}{2.437869in}}%
\pgfpathlineto{\pgfqpoint{1.862002in}{2.548801in}}%
\pgfpathclose%
\pgfusepath{fill}%
\end{pgfscope}%
\begin{pgfscope}%
\pgfpathrectangle{\pgfqpoint{0.100000in}{0.100000in}}{\pgfqpoint{3.957178in}{3.957178in}}%
\pgfusepath{clip}%
\pgfsetbuttcap%
\pgfsetroundjoin%
\definecolor{currentfill}{rgb}{0.435734,0.001127,0.659797}%
\pgfsetfillcolor{currentfill}%
\pgfsetlinewidth{0.000000pt}%
\definecolor{currentstroke}{rgb}{0.000000,0.000000,0.000000}%
\pgfsetstrokecolor{currentstroke}%
\pgfsetdash{}{0pt}%
\pgfpathmoveto{\pgfqpoint{1.705806in}{2.473649in}}%
\pgfpathlineto{\pgfqpoint{1.671443in}{2.260097in}}%
\pgfpathlineto{\pgfqpoint{1.592075in}{2.547375in}}%
\pgfpathlineto{\pgfqpoint{1.705806in}{2.473649in}}%
\pgfpathclose%
\pgfusepath{fill}%
\end{pgfscope}%
\begin{pgfscope}%
\pgfpathrectangle{\pgfqpoint{0.100000in}{0.100000in}}{\pgfqpoint{3.957178in}{3.957178in}}%
\pgfusepath{clip}%
\pgfsetbuttcap%
\pgfsetroundjoin%
\definecolor{currentfill}{rgb}{0.512206,0.018833,0.655209}%
\pgfsetfillcolor{currentfill}%
\pgfsetlinewidth{0.000000pt}%
\definecolor{currentstroke}{rgb}{0.000000,0.000000,0.000000}%
\pgfsetstrokecolor{currentstroke}%
\pgfsetdash{}{0pt}%
\pgfpathmoveto{\pgfqpoint{1.862002in}{2.548801in}}%
\pgfpathlineto{\pgfqpoint{1.804775in}{2.437869in}}%
\pgfpathlineto{\pgfqpoint{1.705806in}{2.473649in}}%
\pgfpathlineto{\pgfqpoint{1.862002in}{2.548801in}}%
\pgfpathclose%
\pgfusepath{fill}%
\end{pgfscope}%
\begin{pgfscope}%
\pgfpathrectangle{\pgfqpoint{0.100000in}{0.100000in}}{\pgfqpoint{3.957178in}{3.957178in}}%
\pgfusepath{clip}%
\pgfsetbuttcap%
\pgfsetroundjoin%
\definecolor{currentfill}{rgb}{0.459623,0.003574,0.660277}%
\pgfsetfillcolor{currentfill}%
\pgfsetlinewidth{0.000000pt}%
\definecolor{currentstroke}{rgb}{0.000000,0.000000,0.000000}%
\pgfsetstrokecolor{currentstroke}%
\pgfsetdash{}{0pt}%
\pgfpathmoveto{\pgfqpoint{1.459637in}{2.452204in}}%
\pgfpathlineto{\pgfqpoint{1.592075in}{2.547375in}}%
\pgfpathlineto{\pgfqpoint{1.671443in}{2.260097in}}%
\pgfpathlineto{\pgfqpoint{1.459637in}{2.452204in}}%
\pgfpathclose%
\pgfusepath{fill}%
\end{pgfscope}%
\begin{pgfscope}%
\pgfpathrectangle{\pgfqpoint{0.100000in}{0.100000in}}{\pgfqpoint{3.957178in}{3.957178in}}%
\pgfusepath{clip}%
\pgfsetbuttcap%
\pgfsetroundjoin%
\definecolor{currentfill}{rgb}{0.350150,0.004382,0.646298}%
\pgfsetfillcolor{currentfill}%
\pgfsetlinewidth{0.000000pt}%
\definecolor{currentstroke}{rgb}{0.000000,0.000000,0.000000}%
\pgfsetstrokecolor{currentstroke}%
\pgfsetdash{}{0pt}%
\pgfpathmoveto{\pgfqpoint{2.401955in}{2.671792in}}%
\pgfpathlineto{\pgfqpoint{3.420890in}{1.935610in}}%
\pgfpathlineto{\pgfqpoint{1.810215in}{2.231164in}}%
\pgfpathlineto{\pgfqpoint{2.401955in}{2.671792in}}%
\pgfpathclose%
\pgfusepath{fill}%
\end{pgfscope}%
\begin{pgfscope}%
\pgfpathrectangle{\pgfqpoint{0.100000in}{0.100000in}}{\pgfqpoint{3.957178in}{3.957178in}}%
\pgfusepath{clip}%
\pgfsetbuttcap%
\pgfsetroundjoin%
\definecolor{currentfill}{rgb}{0.692840,0.165141,0.564522}%
\pgfsetfillcolor{currentfill}%
\pgfsetlinewidth{0.000000pt}%
\definecolor{currentstroke}{rgb}{0.000000,0.000000,0.000000}%
\pgfsetstrokecolor{currentstroke}%
\pgfsetdash{}{0pt}%
\pgfpathmoveto{\pgfqpoint{1.938957in}{2.483638in}}%
\pgfpathlineto{\pgfqpoint{2.304577in}{2.722084in}}%
\pgfpathlineto{\pgfqpoint{2.401955in}{2.671792in}}%
\pgfpathlineto{\pgfqpoint{1.938957in}{2.483638in}}%
\pgfpathclose%
\pgfusepath{fill}%
\end{pgfscope}%
\begin{pgfscope}%
\pgfpathrectangle{\pgfqpoint{0.100000in}{0.100000in}}{\pgfqpoint{3.957178in}{3.957178in}}%
\pgfusepath{clip}%
\pgfsetbuttcap%
\pgfsetroundjoin%
\definecolor{currentfill}{rgb}{0.697324,0.169573,0.560919}%
\pgfsetfillcolor{currentfill}%
\pgfsetlinewidth{0.000000pt}%
\definecolor{currentstroke}{rgb}{0.000000,0.000000,0.000000}%
\pgfsetstrokecolor{currentstroke}%
\pgfsetdash{}{0pt}%
\pgfpathmoveto{\pgfqpoint{1.938957in}{2.483638in}}%
\pgfpathlineto{\pgfqpoint{1.862002in}{2.548801in}}%
\pgfpathlineto{\pgfqpoint{1.916062in}{2.844087in}}%
\pgfpathlineto{\pgfqpoint{1.938957in}{2.483638in}}%
\pgfpathclose%
\pgfusepath{fill}%
\end{pgfscope}%
\begin{pgfscope}%
\pgfpathrectangle{\pgfqpoint{0.100000in}{0.100000in}}{\pgfqpoint{3.957178in}{3.957178in}}%
\pgfusepath{clip}%
\pgfsetbuttcap%
\pgfsetroundjoin%
\definecolor{currentfill}{rgb}{0.692840,0.165141,0.564522}%
\pgfsetfillcolor{currentfill}%
\pgfsetlinewidth{0.000000pt}%
\definecolor{currentstroke}{rgb}{0.000000,0.000000,0.000000}%
\pgfsetstrokecolor{currentstroke}%
\pgfsetdash{}{0pt}%
\pgfpathmoveto{\pgfqpoint{1.862002in}{2.548801in}}%
\pgfpathlineto{\pgfqpoint{1.705806in}{2.473649in}}%
\pgfpathlineto{\pgfqpoint{1.847025in}{2.811821in}}%
\pgfpathlineto{\pgfqpoint{1.862002in}{2.548801in}}%
\pgfpathclose%
\pgfusepath{fill}%
\end{pgfscope}%
\begin{pgfscope}%
\pgfpathrectangle{\pgfqpoint{0.100000in}{0.100000in}}{\pgfqpoint{3.957178in}{3.957178in}}%
\pgfusepath{clip}%
\pgfsetbuttcap%
\pgfsetroundjoin%
\definecolor{currentfill}{rgb}{0.655580,0.129725,0.592317}%
\pgfsetfillcolor{currentfill}%
\pgfsetlinewidth{0.000000pt}%
\definecolor{currentstroke}{rgb}{0.000000,0.000000,0.000000}%
\pgfsetstrokecolor{currentstroke}%
\pgfsetdash{}{0pt}%
\pgfpathmoveto{\pgfqpoint{1.672923in}{2.652240in}}%
\pgfpathlineto{\pgfqpoint{1.705806in}{2.473649in}}%
\pgfpathlineto{\pgfqpoint{1.592075in}{2.547375in}}%
\pgfpathlineto{\pgfqpoint{1.672923in}{2.652240in}}%
\pgfpathclose%
\pgfusepath{fill}%
\end{pgfscope}%
\begin{pgfscope}%
\pgfpathrectangle{\pgfqpoint{0.100000in}{0.100000in}}{\pgfqpoint{3.957178in}{3.957178in}}%
\pgfusepath{clip}%
\pgfsetbuttcap%
\pgfsetroundjoin%
\definecolor{currentfill}{rgb}{0.756304,0.231555,0.509468}%
\pgfsetfillcolor{currentfill}%
\pgfsetlinewidth{0.000000pt}%
\definecolor{currentstroke}{rgb}{0.000000,0.000000,0.000000}%
\pgfsetstrokecolor{currentstroke}%
\pgfsetdash{}{0pt}%
\pgfpathmoveto{\pgfqpoint{2.206315in}{2.843027in}}%
\pgfpathlineto{\pgfqpoint{2.304577in}{2.722084in}}%
\pgfpathlineto{\pgfqpoint{1.938957in}{2.483638in}}%
\pgfpathlineto{\pgfqpoint{2.206315in}{2.843027in}}%
\pgfpathclose%
\pgfusepath{fill}%
\end{pgfscope}%
\begin{pgfscope}%
\pgfpathrectangle{\pgfqpoint{0.100000in}{0.100000in}}{\pgfqpoint{3.957178in}{3.957178in}}%
\pgfusepath{clip}%
\pgfsetbuttcap%
\pgfsetroundjoin%
\definecolor{currentfill}{rgb}{0.645872,0.120898,0.598867}%
\pgfsetfillcolor{currentfill}%
\pgfsetlinewidth{0.000000pt}%
\definecolor{currentstroke}{rgb}{0.000000,0.000000,0.000000}%
\pgfsetstrokecolor{currentstroke}%
\pgfsetdash{}{0pt}%
\pgfpathmoveto{\pgfqpoint{1.594754in}{2.578389in}}%
\pgfpathlineto{\pgfqpoint{1.592075in}{2.547375in}}%
\pgfpathlineto{\pgfqpoint{1.459637in}{2.452204in}}%
\pgfpathlineto{\pgfqpoint{1.594754in}{2.578389in}}%
\pgfpathclose%
\pgfusepath{fill}%
\end{pgfscope}%
\begin{pgfscope}%
\pgfpathrectangle{\pgfqpoint{0.100000in}{0.100000in}}{\pgfqpoint{3.957178in}{3.957178in}}%
\pgfusepath{clip}%
\pgfsetbuttcap%
\pgfsetroundjoin%
\definecolor{currentfill}{rgb}{0.748289,0.222711,0.516834}%
\pgfsetfillcolor{currentfill}%
\pgfsetlinewidth{0.000000pt}%
\definecolor{currentstroke}{rgb}{0.000000,0.000000,0.000000}%
\pgfsetstrokecolor{currentstroke}%
\pgfsetdash{}{0pt}%
\pgfpathmoveto{\pgfqpoint{1.672923in}{2.652240in}}%
\pgfpathlineto{\pgfqpoint{1.847025in}{2.811821in}}%
\pgfpathlineto{\pgfqpoint{1.705806in}{2.473649in}}%
\pgfpathlineto{\pgfqpoint{1.672923in}{2.652240in}}%
\pgfpathclose%
\pgfusepath{fill}%
\end{pgfscope}%
\begin{pgfscope}%
\pgfpathrectangle{\pgfqpoint{0.100000in}{0.100000in}}{\pgfqpoint{3.957178in}{3.957178in}}%
\pgfusepath{clip}%
\pgfsetbuttcap%
\pgfsetroundjoin%
\definecolor{currentfill}{rgb}{0.714883,0.187299,0.546338}%
\pgfsetfillcolor{currentfill}%
\pgfsetlinewidth{0.000000pt}%
\definecolor{currentstroke}{rgb}{0.000000,0.000000,0.000000}%
\pgfsetstrokecolor{currentstroke}%
\pgfsetdash{}{0pt}%
\pgfpathmoveto{\pgfqpoint{1.594754in}{2.578389in}}%
\pgfpathlineto{\pgfqpoint{1.672923in}{2.652240in}}%
\pgfpathlineto{\pgfqpoint{1.592075in}{2.547375in}}%
\pgfpathlineto{\pgfqpoint{1.594754in}{2.578389in}}%
\pgfpathclose%
\pgfusepath{fill}%
\end{pgfscope}%
\begin{pgfscope}%
\pgfpathrectangle{\pgfqpoint{0.100000in}{0.100000in}}{\pgfqpoint{3.957178in}{3.957178in}}%
\pgfusepath{clip}%
\pgfsetbuttcap%
\pgfsetroundjoin%
\definecolor{currentfill}{rgb}{0.669845,0.142992,0.582154}%
\pgfsetfillcolor{currentfill}%
\pgfsetlinewidth{0.000000pt}%
\definecolor{currentstroke}{rgb}{0.000000,0.000000,0.000000}%
\pgfsetstrokecolor{currentstroke}%
\pgfsetdash{}{0pt}%
\pgfpathmoveto{\pgfqpoint{1.459637in}{2.452204in}}%
\pgfpathlineto{\pgfqpoint{1.482696in}{2.562498in}}%
\pgfpathlineto{\pgfqpoint{1.594754in}{2.578389in}}%
\pgfpathlineto{\pgfqpoint{1.459637in}{2.452204in}}%
\pgfpathclose%
\pgfusepath{fill}%
\end{pgfscope}%
\begin{pgfscope}%
\pgfpathrectangle{\pgfqpoint{0.100000in}{0.100000in}}{\pgfqpoint{3.957178in}{3.957178in}}%
\pgfusepath{clip}%
\pgfsetbuttcap%
\pgfsetroundjoin%
\definecolor{currentfill}{rgb}{0.801855,0.284626,0.465971}%
\pgfsetfillcolor{currentfill}%
\pgfsetlinewidth{0.000000pt}%
\definecolor{currentstroke}{rgb}{0.000000,0.000000,0.000000}%
\pgfsetstrokecolor{currentstroke}%
\pgfsetdash{}{0pt}%
\pgfpathmoveto{\pgfqpoint{2.206315in}{2.843027in}}%
\pgfpathlineto{\pgfqpoint{1.938957in}{2.483638in}}%
\pgfpathlineto{\pgfqpoint{1.916062in}{2.844087in}}%
\pgfpathlineto{\pgfqpoint{2.206315in}{2.843027in}}%
\pgfpathclose%
\pgfusepath{fill}%
\end{pgfscope}%
\begin{pgfscope}%
\pgfpathrectangle{\pgfqpoint{0.100000in}{0.100000in}}{\pgfqpoint{3.957178in}{3.957178in}}%
\pgfusepath{clip}%
\pgfsetbuttcap%
\pgfsetroundjoin%
\definecolor{currentfill}{rgb}{0.731862,0.205013,0.531601}%
\pgfsetfillcolor{currentfill}%
\pgfsetlinewidth{0.000000pt}%
\definecolor{currentstroke}{rgb}{0.000000,0.000000,0.000000}%
\pgfsetstrokecolor{currentstroke}%
\pgfsetdash{}{0pt}%
\pgfpathmoveto{\pgfqpoint{1.482696in}{2.562498in}}%
\pgfpathlineto{\pgfqpoint{1.672923in}{2.652240in}}%
\pgfpathlineto{\pgfqpoint{1.594754in}{2.578389in}}%
\pgfpathlineto{\pgfqpoint{1.482696in}{2.562498in}}%
\pgfpathclose%
\pgfusepath{fill}%
\end{pgfscope}%
\begin{pgfscope}%
\pgfpathrectangle{\pgfqpoint{0.100000in}{0.100000in}}{\pgfqpoint{3.957178in}{3.957178in}}%
\pgfusepath{clip}%
\pgfsetbuttcap%
\pgfsetroundjoin%
\definecolor{currentfill}{rgb}{0.823132,0.311261,0.444806}%
\pgfsetfillcolor{currentfill}%
\pgfsetlinewidth{0.000000pt}%
\definecolor{currentstroke}{rgb}{0.000000,0.000000,0.000000}%
\pgfsetstrokecolor{currentstroke}%
\pgfsetdash{}{0pt}%
\pgfpathmoveto{\pgfqpoint{1.916062in}{2.844087in}}%
\pgfpathlineto{\pgfqpoint{1.862002in}{2.548801in}}%
\pgfpathlineto{\pgfqpoint{1.847025in}{2.811821in}}%
\pgfpathlineto{\pgfqpoint{1.916062in}{2.844087in}}%
\pgfpathclose%
\pgfusepath{fill}%
\end{pgfscope}%
\begin{pgfscope}%
\pgfpathrectangle{\pgfqpoint{0.100000in}{0.100000in}}{\pgfqpoint{3.957178in}{3.957178in}}%
\pgfusepath{clip}%
\pgfsetbuttcap%
\pgfsetroundjoin%
\definecolor{currentfill}{rgb}{0.714883,0.187299,0.546338}%
\pgfsetfillcolor{currentfill}%
\pgfsetlinewidth{0.000000pt}%
\definecolor{currentstroke}{rgb}{0.000000,0.000000,0.000000}%
\pgfsetstrokecolor{currentstroke}%
\pgfsetdash{}{0pt}%
\pgfpathmoveto{\pgfqpoint{1.459637in}{2.452204in}}%
\pgfpathlineto{\pgfqpoint{1.432136in}{2.627995in}}%
\pgfpathlineto{\pgfqpoint{1.482696in}{2.562498in}}%
\pgfpathlineto{\pgfqpoint{1.459637in}{2.452204in}}%
\pgfpathclose%
\pgfusepath{fill}%
\end{pgfscope}%
\begin{pgfscope}%
\pgfpathrectangle{\pgfqpoint{0.100000in}{0.100000in}}{\pgfqpoint{3.957178in}{3.957178in}}%
\pgfusepath{clip}%
\pgfsetbuttcap%
\pgfsetroundjoin%
\definecolor{currentfill}{rgb}{0.679160,0.151848,0.575189}%
\pgfsetfillcolor{currentfill}%
\pgfsetlinewidth{0.000000pt}%
\definecolor{currentstroke}{rgb}{0.000000,0.000000,0.000000}%
\pgfsetstrokecolor{currentstroke}%
\pgfsetdash{}{0pt}%
\pgfpathmoveto{\pgfqpoint{2.401955in}{2.671792in}}%
\pgfpathlineto{\pgfqpoint{2.526588in}{2.720411in}}%
\pgfpathlineto{\pgfqpoint{3.328876in}{2.067735in}}%
\pgfpathlineto{\pgfqpoint{2.401955in}{2.671792in}}%
\pgfpathclose%
\pgfusepath{fill}%
\end{pgfscope}%
\begin{pgfscope}%
\pgfpathrectangle{\pgfqpoint{0.100000in}{0.100000in}}{\pgfqpoint{3.957178in}{3.957178in}}%
\pgfusepath{clip}%
\pgfsetbuttcap%
\pgfsetroundjoin%
\definecolor{currentfill}{rgb}{0.471457,0.005678,0.659897}%
\pgfsetfillcolor{currentfill}%
\pgfsetlinewidth{0.000000pt}%
\definecolor{currentstroke}{rgb}{0.000000,0.000000,0.000000}%
\pgfsetstrokecolor{currentstroke}%
\pgfsetdash{}{0pt}%
\pgfpathmoveto{\pgfqpoint{2.401955in}{2.671792in}}%
\pgfpathlineto{\pgfqpoint{3.328876in}{2.067735in}}%
\pgfpathlineto{\pgfqpoint{3.420890in}{1.935610in}}%
\pgfpathlineto{\pgfqpoint{2.401955in}{2.671792in}}%
\pgfpathclose%
\pgfusepath{fill}%
\end{pgfscope}%
\begin{pgfscope}%
\pgfpathrectangle{\pgfqpoint{0.100000in}{0.100000in}}{\pgfqpoint{3.957178in}{3.957178in}}%
\pgfusepath{clip}%
\pgfsetbuttcap%
\pgfsetroundjoin%
\definecolor{currentfill}{rgb}{0.881443,0.392529,0.383229}%
\pgfsetfillcolor{currentfill}%
\pgfsetlinewidth{0.000000pt}%
\definecolor{currentstroke}{rgb}{0.000000,0.000000,0.000000}%
\pgfsetstrokecolor{currentstroke}%
\pgfsetdash{}{0pt}%
\pgfpathmoveto{\pgfqpoint{2.304577in}{2.722084in}}%
\pgfpathlineto{\pgfqpoint{2.397217in}{2.994070in}}%
\pgfpathlineto{\pgfqpoint{2.401955in}{2.671792in}}%
\pgfpathlineto{\pgfqpoint{2.304577in}{2.722084in}}%
\pgfpathclose%
\pgfusepath{fill}%
\end{pgfscope}%
\begin{pgfscope}%
\pgfpathrectangle{\pgfqpoint{0.100000in}{0.100000in}}{\pgfqpoint{3.957178in}{3.957178in}}%
\pgfusepath{clip}%
\pgfsetbuttcap%
\pgfsetroundjoin%
\definecolor{currentfill}{rgb}{0.887402,0.401762,0.376494}%
\pgfsetfillcolor{currentfill}%
\pgfsetlinewidth{0.000000pt}%
\definecolor{currentstroke}{rgb}{0.000000,0.000000,0.000000}%
\pgfsetstrokecolor{currentstroke}%
\pgfsetdash{}{0pt}%
\pgfpathmoveto{\pgfqpoint{2.397217in}{2.994070in}}%
\pgfpathlineto{\pgfqpoint{2.526588in}{2.720411in}}%
\pgfpathlineto{\pgfqpoint{2.401955in}{2.671792in}}%
\pgfpathlineto{\pgfqpoint{2.397217in}{2.994070in}}%
\pgfpathclose%
\pgfusepath{fill}%
\end{pgfscope}%
\begin{pgfscope}%
\pgfpathrectangle{\pgfqpoint{0.100000in}{0.100000in}}{\pgfqpoint{3.957178in}{3.957178in}}%
\pgfusepath{clip}%
\pgfsetbuttcap%
\pgfsetroundjoin%
\definecolor{currentfill}{rgb}{0.856547,0.356066,0.410322}%
\pgfsetfillcolor{currentfill}%
\pgfsetlinewidth{0.000000pt}%
\definecolor{currentstroke}{rgb}{0.000000,0.000000,0.000000}%
\pgfsetstrokecolor{currentstroke}%
\pgfsetdash{}{0pt}%
\pgfpathmoveto{\pgfqpoint{1.672923in}{2.652240in}}%
\pgfpathlineto{\pgfqpoint{1.482696in}{2.562498in}}%
\pgfpathlineto{\pgfqpoint{1.791056in}{2.960511in}}%
\pgfpathlineto{\pgfqpoint{1.672923in}{2.652240in}}%
\pgfpathclose%
\pgfusepath{fill}%
\end{pgfscope}%
\begin{pgfscope}%
\pgfpathrectangle{\pgfqpoint{0.100000in}{0.100000in}}{\pgfqpoint{3.957178in}{3.957178in}}%
\pgfusepath{clip}%
\pgfsetbuttcap%
\pgfsetroundjoin%
\definecolor{currentfill}{rgb}{0.853319,0.351553,0.413734}%
\pgfsetfillcolor{currentfill}%
\pgfsetlinewidth{0.000000pt}%
\definecolor{currentstroke}{rgb}{0.000000,0.000000,0.000000}%
\pgfsetstrokecolor{currentstroke}%
\pgfsetdash{}{0pt}%
\pgfpathmoveto{\pgfqpoint{2.586901in}{2.744610in}}%
\pgfpathlineto{\pgfqpoint{2.864660in}{2.659680in}}%
\pgfpathlineto{\pgfqpoint{2.526588in}{2.720411in}}%
\pgfpathlineto{\pgfqpoint{2.586901in}{2.744610in}}%
\pgfpathclose%
\pgfusepath{fill}%
\end{pgfscope}%
\begin{pgfscope}%
\pgfpathrectangle{\pgfqpoint{0.100000in}{0.100000in}}{\pgfqpoint{3.957178in}{3.957178in}}%
\pgfusepath{clip}%
\pgfsetbuttcap%
\pgfsetroundjoin%
\definecolor{currentfill}{rgb}{0.920714,0.458417,0.336166}%
\pgfsetfillcolor{currentfill}%
\pgfsetlinewidth{0.000000pt}%
\definecolor{currentstroke}{rgb}{0.000000,0.000000,0.000000}%
\pgfsetstrokecolor{currentstroke}%
\pgfsetdash{}{0pt}%
\pgfpathmoveto{\pgfqpoint{2.206315in}{2.843027in}}%
\pgfpathlineto{\pgfqpoint{2.397217in}{2.994070in}}%
\pgfpathlineto{\pgfqpoint{2.304577in}{2.722084in}}%
\pgfpathlineto{\pgfqpoint{2.206315in}{2.843027in}}%
\pgfpathclose%
\pgfusepath{fill}%
\end{pgfscope}%
\begin{pgfscope}%
\pgfpathrectangle{\pgfqpoint{0.100000in}{0.100000in}}{\pgfqpoint{3.957178in}{3.957178in}}%
\pgfusepath{clip}%
\pgfsetbuttcap%
\pgfsetroundjoin%
\definecolor{currentfill}{rgb}{0.904601,0.429797,0.356329}%
\pgfsetfillcolor{currentfill}%
\pgfsetlinewidth{0.000000pt}%
\definecolor{currentstroke}{rgb}{0.000000,0.000000,0.000000}%
\pgfsetstrokecolor{currentstroke}%
\pgfsetdash{}{0pt}%
\pgfpathmoveto{\pgfqpoint{1.791056in}{2.960511in}}%
\pgfpathlineto{\pgfqpoint{1.847025in}{2.811821in}}%
\pgfpathlineto{\pgfqpoint{1.672923in}{2.652240in}}%
\pgfpathlineto{\pgfqpoint{1.791056in}{2.960511in}}%
\pgfpathclose%
\pgfusepath{fill}%
\end{pgfscope}%
\begin{pgfscope}%
\pgfpathrectangle{\pgfqpoint{0.100000in}{0.100000in}}{\pgfqpoint{3.957178in}{3.957178in}}%
\pgfusepath{clip}%
\pgfsetbuttcap%
\pgfsetroundjoin%
\definecolor{currentfill}{rgb}{0.723444,0.196158,0.538981}%
\pgfsetfillcolor{currentfill}%
\pgfsetlinewidth{0.000000pt}%
\definecolor{currentstroke}{rgb}{0.000000,0.000000,0.000000}%
\pgfsetstrokecolor{currentstroke}%
\pgfsetdash{}{0pt}%
\pgfpathmoveto{\pgfqpoint{2.864660in}{2.659680in}}%
\pgfpathlineto{\pgfqpoint{3.328876in}{2.067735in}}%
\pgfpathlineto{\pgfqpoint{2.526588in}{2.720411in}}%
\pgfpathlineto{\pgfqpoint{2.864660in}{2.659680in}}%
\pgfpathclose%
\pgfusepath{fill}%
\end{pgfscope}%
\begin{pgfscope}%
\pgfpathrectangle{\pgfqpoint{0.100000in}{0.100000in}}{\pgfqpoint{3.957178in}{3.957178in}}%
\pgfusepath{clip}%
\pgfsetbuttcap%
\pgfsetroundjoin%
\definecolor{currentfill}{rgb}{0.293478,0.010213,0.629490}%
\pgfsetfillcolor{currentfill}%
\pgfsetlinewidth{0.000000pt}%
\definecolor{currentstroke}{rgb}{0.000000,0.000000,0.000000}%
\pgfsetstrokecolor{currentstroke}%
\pgfsetdash{}{0pt}%
\pgfpathmoveto{\pgfqpoint{3.328876in}{2.067735in}}%
\pgfpathlineto{\pgfqpoint{3.413142in}{1.956666in}}%
\pgfpathlineto{\pgfqpoint{3.420890in}{1.935610in}}%
\pgfpathlineto{\pgfqpoint{3.328876in}{2.067735in}}%
\pgfpathclose%
\pgfusepath{fill}%
\end{pgfscope}%
\begin{pgfscope}%
\pgfpathrectangle{\pgfqpoint{0.100000in}{0.100000in}}{\pgfqpoint{3.957178in}{3.957178in}}%
\pgfusepath{clip}%
\pgfsetbuttcap%
\pgfsetroundjoin%
\definecolor{currentfill}{rgb}{0.186213,0.018803,0.587228}%
\pgfsetfillcolor{currentfill}%
\pgfsetlinewidth{0.000000pt}%
\definecolor{currentstroke}{rgb}{0.000000,0.000000,0.000000}%
\pgfsetstrokecolor{currentstroke}%
\pgfsetdash{}{0pt}%
\pgfpathmoveto{\pgfqpoint{3.413142in}{1.956666in}}%
\pgfpathlineto{\pgfqpoint{3.547673in}{1.762813in}}%
\pgfpathlineto{\pgfqpoint{3.420890in}{1.935610in}}%
\pgfpathlineto{\pgfqpoint{3.413142in}{1.956666in}}%
\pgfpathclose%
\pgfusepath{fill}%
\end{pgfscope}%
\begin{pgfscope}%
\pgfpathrectangle{\pgfqpoint{0.100000in}{0.100000in}}{\pgfqpoint{3.957178in}{3.957178in}}%
\pgfusepath{clip}%
\pgfsetbuttcap%
\pgfsetroundjoin%
\definecolor{currentfill}{rgb}{0.756304,0.231555,0.509468}%
\pgfsetfillcolor{currentfill}%
\pgfsetlinewidth{0.000000pt}%
\definecolor{currentstroke}{rgb}{0.000000,0.000000,0.000000}%
\pgfsetstrokecolor{currentstroke}%
\pgfsetdash{}{0pt}%
\pgfpathmoveto{\pgfqpoint{1.459637in}{2.452204in}}%
\pgfpathlineto{\pgfqpoint{1.145690in}{2.475609in}}%
\pgfpathlineto{\pgfqpoint{1.432136in}{2.627995in}}%
\pgfpathlineto{\pgfqpoint{1.459637in}{2.452204in}}%
\pgfpathclose%
\pgfusepath{fill}%
\end{pgfscope}%
\begin{pgfscope}%
\pgfpathrectangle{\pgfqpoint{0.100000in}{0.100000in}}{\pgfqpoint{3.957178in}{3.957178in}}%
\pgfusepath{clip}%
\pgfsetbuttcap%
\pgfsetroundjoin%
\definecolor{currentfill}{rgb}{0.915471,0.448807,0.342890}%
\pgfsetfillcolor{currentfill}%
\pgfsetlinewidth{0.000000pt}%
\definecolor{currentstroke}{rgb}{0.000000,0.000000,0.000000}%
\pgfsetstrokecolor{currentstroke}%
\pgfsetdash{}{0pt}%
\pgfpathmoveto{\pgfqpoint{2.526588in}{2.720411in}}%
\pgfpathlineto{\pgfqpoint{2.397217in}{2.994070in}}%
\pgfpathlineto{\pgfqpoint{2.586901in}{2.744610in}}%
\pgfpathlineto{\pgfqpoint{2.526588in}{2.720411in}}%
\pgfpathclose%
\pgfusepath{fill}%
\end{pgfscope}%
\begin{pgfscope}%
\pgfpathrectangle{\pgfqpoint{0.100000in}{0.100000in}}{\pgfqpoint{3.957178in}{3.957178in}}%
\pgfusepath{clip}%
\pgfsetbuttcap%
\pgfsetroundjoin%
\definecolor{currentfill}{rgb}{0.942598,0.502639,0.305816}%
\pgfsetfillcolor{currentfill}%
\pgfsetlinewidth{0.000000pt}%
\definecolor{currentstroke}{rgb}{0.000000,0.000000,0.000000}%
\pgfsetstrokecolor{currentstroke}%
\pgfsetdash{}{0pt}%
\pgfpathmoveto{\pgfqpoint{1.847025in}{2.811821in}}%
\pgfpathlineto{\pgfqpoint{1.791056in}{2.960511in}}%
\pgfpathlineto{\pgfqpoint{1.916062in}{2.844087in}}%
\pgfpathlineto{\pgfqpoint{1.847025in}{2.811821in}}%
\pgfpathclose%
\pgfusepath{fill}%
\end{pgfscope}%
\begin{pgfscope}%
\pgfpathrectangle{\pgfqpoint{0.100000in}{0.100000in}}{\pgfqpoint{3.957178in}{3.957178in}}%
\pgfusepath{clip}%
\pgfsetbuttcap%
\pgfsetroundjoin%
\definecolor{currentfill}{rgb}{0.148607,0.021154,0.570562}%
\pgfsetfillcolor{currentfill}%
\pgfsetlinewidth{0.000000pt}%
\definecolor{currentstroke}{rgb}{0.000000,0.000000,0.000000}%
\pgfsetstrokecolor{currentstroke}%
\pgfsetdash{}{0pt}%
\pgfpathmoveto{\pgfqpoint{3.413142in}{1.956666in}}%
\pgfpathlineto{\pgfqpoint{3.550789in}{1.750820in}}%
\pgfpathlineto{\pgfqpoint{3.547673in}{1.762813in}}%
\pgfpathlineto{\pgfqpoint{3.413142in}{1.956666in}}%
\pgfpathclose%
\pgfusepath{fill}%
\end{pgfscope}%
\begin{pgfscope}%
\pgfpathrectangle{\pgfqpoint{0.100000in}{0.100000in}}{\pgfqpoint{3.957178in}{3.957178in}}%
\pgfusepath{clip}%
\pgfsetbuttcap%
\pgfsetroundjoin%
\definecolor{currentfill}{rgb}{0.884436,0.397139,0.379860}%
\pgfsetfillcolor{currentfill}%
\pgfsetlinewidth{0.000000pt}%
\definecolor{currentstroke}{rgb}{0.000000,0.000000,0.000000}%
\pgfsetstrokecolor{currentstroke}%
\pgfsetdash{}{0pt}%
\pgfpathmoveto{\pgfqpoint{2.805314in}{2.745587in}}%
\pgfpathlineto{\pgfqpoint{2.864660in}{2.659680in}}%
\pgfpathlineto{\pgfqpoint{2.586901in}{2.744610in}}%
\pgfpathlineto{\pgfqpoint{2.805314in}{2.745587in}}%
\pgfpathclose%
\pgfusepath{fill}%
\end{pgfscope}%
\begin{pgfscope}%
\pgfpathrectangle{\pgfqpoint{0.100000in}{0.100000in}}{\pgfqpoint{3.957178in}{3.957178in}}%
\pgfusepath{clip}%
\pgfsetbuttcap%
\pgfsetroundjoin%
\definecolor{currentfill}{rgb}{0.881443,0.392529,0.383229}%
\pgfsetfillcolor{currentfill}%
\pgfsetlinewidth{0.000000pt}%
\definecolor{currentstroke}{rgb}{0.000000,0.000000,0.000000}%
\pgfsetstrokecolor{currentstroke}%
\pgfsetdash{}{0pt}%
\pgfpathmoveto{\pgfqpoint{1.432136in}{2.627995in}}%
\pgfpathlineto{\pgfqpoint{1.596907in}{2.947877in}}%
\pgfpathlineto{\pgfqpoint{1.482696in}{2.562498in}}%
\pgfpathlineto{\pgfqpoint{1.432136in}{2.627995in}}%
\pgfpathclose%
\pgfusepath{fill}%
\end{pgfscope}%
\begin{pgfscope}%
\pgfpathrectangle{\pgfqpoint{0.100000in}{0.100000in}}{\pgfqpoint{3.957178in}{3.957178in}}%
\pgfusepath{clip}%
\pgfsetbuttcap%
\pgfsetroundjoin%
\definecolor{currentfill}{rgb}{0.966798,0.564396,0.265118}%
\pgfsetfillcolor{currentfill}%
\pgfsetlinewidth{0.000000pt}%
\definecolor{currentstroke}{rgb}{0.000000,0.000000,0.000000}%
\pgfsetstrokecolor{currentstroke}%
\pgfsetdash{}{0pt}%
\pgfpathmoveto{\pgfqpoint{1.916062in}{2.844087in}}%
\pgfpathlineto{\pgfqpoint{2.264880in}{3.086674in}}%
\pgfpathlineto{\pgfqpoint{2.206315in}{2.843027in}}%
\pgfpathlineto{\pgfqpoint{1.916062in}{2.844087in}}%
\pgfpathclose%
\pgfusepath{fill}%
\end{pgfscope}%
\begin{pgfscope}%
\pgfpathrectangle{\pgfqpoint{0.100000in}{0.100000in}}{\pgfqpoint{3.957178in}{3.957178in}}%
\pgfusepath{clip}%
\pgfsetbuttcap%
\pgfsetroundjoin%
\definecolor{currentfill}{rgb}{0.928329,0.472975,0.326067}%
\pgfsetfillcolor{currentfill}%
\pgfsetlinewidth{0.000000pt}%
\definecolor{currentstroke}{rgb}{0.000000,0.000000,0.000000}%
\pgfsetstrokecolor{currentstroke}%
\pgfsetdash{}{0pt}%
\pgfpathmoveto{\pgfqpoint{2.586901in}{2.744610in}}%
\pgfpathlineto{\pgfqpoint{2.661682in}{2.887953in}}%
\pgfpathlineto{\pgfqpoint{2.805314in}{2.745587in}}%
\pgfpathlineto{\pgfqpoint{2.586901in}{2.744610in}}%
\pgfpathclose%
\pgfusepath{fill}%
\end{pgfscope}%
\begin{pgfscope}%
\pgfpathrectangle{\pgfqpoint{0.100000in}{0.100000in}}{\pgfqpoint{3.957178in}{3.957178in}}%
\pgfusepath{clip}%
\pgfsetbuttcap%
\pgfsetroundjoin%
\definecolor{currentfill}{rgb}{0.221197,0.016497,0.602083}%
\pgfsetfillcolor{currentfill}%
\pgfsetlinewidth{0.000000pt}%
\definecolor{currentstroke}{rgb}{0.000000,0.000000,0.000000}%
\pgfsetstrokecolor{currentstroke}%
\pgfsetdash{}{0pt}%
\pgfpathmoveto{\pgfqpoint{3.413142in}{1.956666in}}%
\pgfpathlineto{\pgfqpoint{3.505448in}{1.816784in}}%
\pgfpathlineto{\pgfqpoint{3.550789in}{1.750820in}}%
\pgfpathlineto{\pgfqpoint{3.413142in}{1.956666in}}%
\pgfpathclose%
\pgfusepath{fill}%
\end{pgfscope}%
\begin{pgfscope}%
\pgfpathrectangle{\pgfqpoint{0.100000in}{0.100000in}}{\pgfqpoint{3.957178in}{3.957178in}}%
\pgfusepath{clip}%
\pgfsetbuttcap%
\pgfsetroundjoin%
\definecolor{currentfill}{rgb}{0.787133,0.266922,0.480307}%
\pgfsetfillcolor{currentfill}%
\pgfsetlinewidth{0.000000pt}%
\definecolor{currentstroke}{rgb}{0.000000,0.000000,0.000000}%
\pgfsetstrokecolor{currentstroke}%
\pgfsetdash{}{0pt}%
\pgfpathmoveto{\pgfqpoint{2.848634in}{2.776341in}}%
\pgfpathlineto{\pgfqpoint{3.328876in}{2.067735in}}%
\pgfpathlineto{\pgfqpoint{2.864660in}{2.659680in}}%
\pgfpathlineto{\pgfqpoint{2.848634in}{2.776341in}}%
\pgfpathclose%
\pgfusepath{fill}%
\end{pgfscope}%
\begin{pgfscope}%
\pgfpathrectangle{\pgfqpoint{0.100000in}{0.100000in}}{\pgfqpoint{3.957178in}{3.957178in}}%
\pgfusepath{clip}%
\pgfsetbuttcap%
\pgfsetroundjoin%
\definecolor{currentfill}{rgb}{0.471457,0.005678,0.659897}%
\pgfsetfillcolor{currentfill}%
\pgfsetlinewidth{0.000000pt}%
\definecolor{currentstroke}{rgb}{0.000000,0.000000,0.000000}%
\pgfsetstrokecolor{currentstroke}%
\pgfsetdash{}{0pt}%
\pgfpathmoveto{\pgfqpoint{3.413142in}{1.956666in}}%
\pgfpathlineto{\pgfqpoint{3.328876in}{2.067735in}}%
\pgfpathlineto{\pgfqpoint{3.282578in}{2.232451in}}%
\pgfpathlineto{\pgfqpoint{3.413142in}{1.956666in}}%
\pgfpathclose%
\pgfusepath{fill}%
\end{pgfscope}%
\begin{pgfscope}%
\pgfpathrectangle{\pgfqpoint{0.100000in}{0.100000in}}{\pgfqpoint{3.957178in}{3.957178in}}%
\pgfusepath{clip}%
\pgfsetbuttcap%
\pgfsetroundjoin%
\definecolor{currentfill}{rgb}{0.731862,0.205013,0.531601}%
\pgfsetfillcolor{currentfill}%
\pgfsetlinewidth{0.000000pt}%
\definecolor{currentstroke}{rgb}{0.000000,0.000000,0.000000}%
\pgfsetstrokecolor{currentstroke}%
\pgfsetdash{}{0pt}%
\pgfpathmoveto{\pgfqpoint{1.145690in}{2.475609in}}%
\pgfpathlineto{\pgfqpoint{1.459637in}{2.452204in}}%
\pgfpathlineto{\pgfqpoint{1.017358in}{2.293891in}}%
\pgfpathlineto{\pgfqpoint{1.145690in}{2.475609in}}%
\pgfpathclose%
\pgfusepath{fill}%
\end{pgfscope}%
\begin{pgfscope}%
\pgfpathrectangle{\pgfqpoint{0.100000in}{0.100000in}}{\pgfqpoint{3.957178in}{3.957178in}}%
\pgfusepath{clip}%
\pgfsetbuttcap%
\pgfsetroundjoin%
\definecolor{currentfill}{rgb}{0.944844,0.507658,0.302433}%
\pgfsetfillcolor{currentfill}%
\pgfsetlinewidth{0.000000pt}%
\definecolor{currentstroke}{rgb}{0.000000,0.000000,0.000000}%
\pgfsetstrokecolor{currentstroke}%
\pgfsetdash{}{0pt}%
\pgfpathmoveto{\pgfqpoint{1.482696in}{2.562498in}}%
\pgfpathlineto{\pgfqpoint{1.596907in}{2.947877in}}%
\pgfpathlineto{\pgfqpoint{1.791056in}{2.960511in}}%
\pgfpathlineto{\pgfqpoint{1.482696in}{2.562498in}}%
\pgfpathclose%
\pgfusepath{fill}%
\end{pgfscope}%
\begin{pgfscope}%
\pgfpathrectangle{\pgfqpoint{0.100000in}{0.100000in}}{\pgfqpoint{3.957178in}{3.957178in}}%
\pgfusepath{clip}%
\pgfsetbuttcap%
\pgfsetroundjoin%
\definecolor{currentfill}{rgb}{0.959424,0.543431,0.278701}%
\pgfsetfillcolor{currentfill}%
\pgfsetlinewidth{0.000000pt}%
\definecolor{currentstroke}{rgb}{0.000000,0.000000,0.000000}%
\pgfsetstrokecolor{currentstroke}%
\pgfsetdash{}{0pt}%
\pgfpathmoveto{\pgfqpoint{2.586901in}{2.744610in}}%
\pgfpathlineto{\pgfqpoint{2.545300in}{2.957675in}}%
\pgfpathlineto{\pgfqpoint{2.661682in}{2.887953in}}%
\pgfpathlineto{\pgfqpoint{2.586901in}{2.744610in}}%
\pgfpathclose%
\pgfusepath{fill}%
\end{pgfscope}%
\begin{pgfscope}%
\pgfpathrectangle{\pgfqpoint{0.100000in}{0.100000in}}{\pgfqpoint{3.957178in}{3.957178in}}%
\pgfusepath{clip}%
\pgfsetbuttcap%
\pgfsetroundjoin%
\definecolor{currentfill}{rgb}{0.973416,0.585761,0.251540}%
\pgfsetfillcolor{currentfill}%
\pgfsetlinewidth{0.000000pt}%
\definecolor{currentstroke}{rgb}{0.000000,0.000000,0.000000}%
\pgfsetstrokecolor{currentstroke}%
\pgfsetdash{}{0pt}%
\pgfpathmoveto{\pgfqpoint{2.586901in}{2.744610in}}%
\pgfpathlineto{\pgfqpoint{2.397217in}{2.994070in}}%
\pgfpathlineto{\pgfqpoint{2.465371in}{3.006355in}}%
\pgfpathlineto{\pgfqpoint{2.586901in}{2.744610in}}%
\pgfpathclose%
\pgfusepath{fill}%
\end{pgfscope}%
\begin{pgfscope}%
\pgfpathrectangle{\pgfqpoint{0.100000in}{0.100000in}}{\pgfqpoint{3.957178in}{3.957178in}}%
\pgfusepath{clip}%
\pgfsetbuttcap%
\pgfsetroundjoin%
\definecolor{currentfill}{rgb}{0.973416,0.585761,0.251540}%
\pgfsetfillcolor{currentfill}%
\pgfsetlinewidth{0.000000pt}%
\definecolor{currentstroke}{rgb}{0.000000,0.000000,0.000000}%
\pgfsetstrokecolor{currentstroke}%
\pgfsetdash{}{0pt}%
\pgfpathmoveto{\pgfqpoint{2.465371in}{3.006355in}}%
\pgfpathlineto{\pgfqpoint{2.545300in}{2.957675in}}%
\pgfpathlineto{\pgfqpoint{2.586901in}{2.744610in}}%
\pgfpathlineto{\pgfqpoint{2.465371in}{3.006355in}}%
\pgfpathclose%
\pgfusepath{fill}%
\end{pgfscope}%
\begin{pgfscope}%
\pgfpathrectangle{\pgfqpoint{0.100000in}{0.100000in}}{\pgfqpoint{3.957178in}{3.957178in}}%
\pgfusepath{clip}%
\pgfsetbuttcap%
\pgfsetroundjoin%
\definecolor{currentfill}{rgb}{0.988260,0.652325,0.211364}%
\pgfsetfillcolor{currentfill}%
\pgfsetlinewidth{0.000000pt}%
\definecolor{currentstroke}{rgb}{0.000000,0.000000,0.000000}%
\pgfsetstrokecolor{currentstroke}%
\pgfsetdash{}{0pt}%
\pgfpathmoveto{\pgfqpoint{2.264880in}{3.086674in}}%
\pgfpathlineto{\pgfqpoint{2.397217in}{2.994070in}}%
\pgfpathlineto{\pgfqpoint{2.206315in}{2.843027in}}%
\pgfpathlineto{\pgfqpoint{2.264880in}{3.086674in}}%
\pgfpathclose%
\pgfusepath{fill}%
\end{pgfscope}%
\begin{pgfscope}%
\pgfpathrectangle{\pgfqpoint{0.100000in}{0.100000in}}{\pgfqpoint{3.957178in}{3.957178in}}%
\pgfusepath{clip}%
\pgfsetbuttcap%
\pgfsetroundjoin%
\definecolor{currentfill}{rgb}{0.918109,0.453603,0.339529}%
\pgfsetfillcolor{currentfill}%
\pgfsetlinewidth{0.000000pt}%
\definecolor{currentstroke}{rgb}{0.000000,0.000000,0.000000}%
\pgfsetstrokecolor{currentstroke}%
\pgfsetdash{}{0pt}%
\pgfpathmoveto{\pgfqpoint{2.864660in}{2.659680in}}%
\pgfpathlineto{\pgfqpoint{2.805314in}{2.745587in}}%
\pgfpathlineto{\pgfqpoint{2.848634in}{2.776341in}}%
\pgfpathlineto{\pgfqpoint{2.864660in}{2.659680in}}%
\pgfpathclose%
\pgfusepath{fill}%
\end{pgfscope}%
\begin{pgfscope}%
\pgfpathrectangle{\pgfqpoint{0.100000in}{0.100000in}}{\pgfqpoint{3.957178in}{3.957178in}}%
\pgfusepath{clip}%
\pgfsetbuttcap%
\pgfsetroundjoin%
\definecolor{currentfill}{rgb}{0.987332,0.646633,0.214648}%
\pgfsetfillcolor{currentfill}%
\pgfsetlinewidth{0.000000pt}%
\definecolor{currentstroke}{rgb}{0.000000,0.000000,0.000000}%
\pgfsetstrokecolor{currentstroke}%
\pgfsetdash{}{0pt}%
\pgfpathmoveto{\pgfqpoint{1.916062in}{2.844087in}}%
\pgfpathlineto{\pgfqpoint{1.791056in}{2.960511in}}%
\pgfpathlineto{\pgfqpoint{1.836764in}{3.067780in}}%
\pgfpathlineto{\pgfqpoint{1.916062in}{2.844087in}}%
\pgfpathclose%
\pgfusepath{fill}%
\end{pgfscope}%
\begin{pgfscope}%
\pgfpathrectangle{\pgfqpoint{0.100000in}{0.100000in}}{\pgfqpoint{3.957178in}{3.957178in}}%
\pgfusepath{clip}%
\pgfsetbuttcap%
\pgfsetroundjoin%
\definecolor{currentfill}{rgb}{0.878423,0.387932,0.386600}%
\pgfsetfillcolor{currentfill}%
\pgfsetlinewidth{0.000000pt}%
\definecolor{currentstroke}{rgb}{0.000000,0.000000,0.000000}%
\pgfsetstrokecolor{currentstroke}%
\pgfsetdash{}{0pt}%
\pgfpathmoveto{\pgfqpoint{1.145690in}{2.475609in}}%
\pgfpathlineto{\pgfqpoint{1.402175in}{2.771187in}}%
\pgfpathlineto{\pgfqpoint{1.432136in}{2.627995in}}%
\pgfpathlineto{\pgfqpoint{1.145690in}{2.475609in}}%
\pgfpathclose%
\pgfusepath{fill}%
\end{pgfscope}%
\begin{pgfscope}%
\pgfpathrectangle{\pgfqpoint{0.100000in}{0.100000in}}{\pgfqpoint{3.957178in}{3.957178in}}%
\pgfusepath{clip}%
\pgfsetbuttcap%
\pgfsetroundjoin%
\definecolor{currentfill}{rgb}{0.435734,0.001127,0.659797}%
\pgfsetfillcolor{currentfill}%
\pgfsetlinewidth{0.000000pt}%
\definecolor{currentstroke}{rgb}{0.000000,0.000000,0.000000}%
\pgfsetstrokecolor{currentstroke}%
\pgfsetdash{}{0pt}%
\pgfpathmoveto{\pgfqpoint{3.505448in}{1.816784in}}%
\pgfpathlineto{\pgfqpoint{3.413142in}{1.956666in}}%
\pgfpathlineto{\pgfqpoint{3.282578in}{2.232451in}}%
\pgfpathlineto{\pgfqpoint{3.505448in}{1.816784in}}%
\pgfpathclose%
\pgfusepath{fill}%
\end{pgfscope}%
\begin{pgfscope}%
\pgfpathrectangle{\pgfqpoint{0.100000in}{0.100000in}}{\pgfqpoint{3.957178in}{3.957178in}}%
\pgfusepath{clip}%
\pgfsetbuttcap%
\pgfsetroundjoin%
\definecolor{currentfill}{rgb}{0.949217,0.517763,0.295662}%
\pgfsetfillcolor{currentfill}%
\pgfsetlinewidth{0.000000pt}%
\definecolor{currentstroke}{rgb}{0.000000,0.000000,0.000000}%
\pgfsetstrokecolor{currentstroke}%
\pgfsetdash{}{0pt}%
\pgfpathmoveto{\pgfqpoint{1.596907in}{2.947877in}}%
\pgfpathlineto{\pgfqpoint{1.432136in}{2.627995in}}%
\pgfpathlineto{\pgfqpoint{1.402175in}{2.771187in}}%
\pgfpathlineto{\pgfqpoint{1.596907in}{2.947877in}}%
\pgfpathclose%
\pgfusepath{fill}%
\end{pgfscope}%
\begin{pgfscope}%
\pgfpathrectangle{\pgfqpoint{0.100000in}{0.100000in}}{\pgfqpoint{3.957178in}{3.957178in}}%
\pgfusepath{clip}%
\pgfsetbuttcap%
\pgfsetroundjoin%
\definecolor{currentfill}{rgb}{0.994561,0.734791,0.168938}%
\pgfsetfillcolor{currentfill}%
\pgfsetlinewidth{0.000000pt}%
\definecolor{currentstroke}{rgb}{0.000000,0.000000,0.000000}%
\pgfsetstrokecolor{currentstroke}%
\pgfsetdash{}{0pt}%
\pgfpathmoveto{\pgfqpoint{2.379181in}{3.032853in}}%
\pgfpathlineto{\pgfqpoint{2.465371in}{3.006355in}}%
\pgfpathlineto{\pgfqpoint{2.397217in}{2.994070in}}%
\pgfpathlineto{\pgfqpoint{2.379181in}{3.032853in}}%
\pgfpathclose%
\pgfusepath{fill}%
\end{pgfscope}%
\begin{pgfscope}%
\pgfpathrectangle{\pgfqpoint{0.100000in}{0.100000in}}{\pgfqpoint{3.957178in}{3.957178in}}%
\pgfusepath{clip}%
\pgfsetbuttcap%
\pgfsetroundjoin%
\definecolor{currentfill}{rgb}{0.274191,0.012109,0.622722}%
\pgfsetfillcolor{currentfill}%
\pgfsetlinewidth{0.000000pt}%
\definecolor{currentstroke}{rgb}{0.000000,0.000000,0.000000}%
\pgfsetstrokecolor{currentstroke}%
\pgfsetdash{}{0pt}%
\pgfpathmoveto{\pgfqpoint{3.505448in}{1.816784in}}%
\pgfpathlineto{\pgfqpoint{3.464151in}{1.860528in}}%
\pgfpathlineto{\pgfqpoint{3.550789in}{1.750820in}}%
\pgfpathlineto{\pgfqpoint{3.505448in}{1.816784in}}%
\pgfpathclose%
\pgfusepath{fill}%
\end{pgfscope}%
\begin{pgfscope}%
\pgfpathrectangle{\pgfqpoint{0.100000in}{0.100000in}}{\pgfqpoint{3.957178in}{3.957178in}}%
\pgfusepath{clip}%
\pgfsetbuttcap%
\pgfsetroundjoin%
\definecolor{currentfill}{rgb}{0.994561,0.734791,0.168938}%
\pgfsetfillcolor{currentfill}%
\pgfsetlinewidth{0.000000pt}%
\definecolor{currentstroke}{rgb}{0.000000,0.000000,0.000000}%
\pgfsetstrokecolor{currentstroke}%
\pgfsetdash{}{0pt}%
\pgfpathmoveto{\pgfqpoint{1.916062in}{2.844087in}}%
\pgfpathlineto{\pgfqpoint{1.877921in}{3.091334in}}%
\pgfpathlineto{\pgfqpoint{2.264880in}{3.086674in}}%
\pgfpathlineto{\pgfqpoint{1.916062in}{2.844087in}}%
\pgfpathclose%
\pgfusepath{fill}%
\end{pgfscope}%
\begin{pgfscope}%
\pgfpathrectangle{\pgfqpoint{0.100000in}{0.100000in}}{\pgfqpoint{3.957178in}{3.957178in}}%
\pgfusepath{clip}%
\pgfsetbuttcap%
\pgfsetroundjoin%
\definecolor{currentfill}{rgb}{0.846788,0.342551,0.420579}%
\pgfsetfillcolor{currentfill}%
\pgfsetlinewidth{0.000000pt}%
\definecolor{currentstroke}{rgb}{0.000000,0.000000,0.000000}%
\pgfsetstrokecolor{currentstroke}%
\pgfsetdash{}{0pt}%
\pgfpathmoveto{\pgfqpoint{3.328876in}{2.067735in}}%
\pgfpathlineto{\pgfqpoint{2.848634in}{2.776341in}}%
\pgfpathlineto{\pgfqpoint{2.968450in}{2.738710in}}%
\pgfpathlineto{\pgfqpoint{3.328876in}{2.067735in}}%
\pgfpathclose%
\pgfusepath{fill}%
\end{pgfscope}%
\begin{pgfscope}%
\pgfpathrectangle{\pgfqpoint{0.100000in}{0.100000in}}{\pgfqpoint{3.957178in}{3.957178in}}%
\pgfusepath{clip}%
\pgfsetbuttcap%
\pgfsetroundjoin%
\definecolor{currentfill}{rgb}{0.974947,0.591165,0.248151}%
\pgfsetfillcolor{currentfill}%
\pgfsetlinewidth{0.000000pt}%
\definecolor{currentstroke}{rgb}{0.000000,0.000000,0.000000}%
\pgfsetstrokecolor{currentstroke}%
\pgfsetdash{}{0pt}%
\pgfpathmoveto{\pgfqpoint{2.805314in}{2.745587in}}%
\pgfpathlineto{\pgfqpoint{2.661682in}{2.887953in}}%
\pgfpathlineto{\pgfqpoint{2.786897in}{2.893114in}}%
\pgfpathlineto{\pgfqpoint{2.805314in}{2.745587in}}%
\pgfpathclose%
\pgfusepath{fill}%
\end{pgfscope}%
\begin{pgfscope}%
\pgfpathrectangle{\pgfqpoint{0.100000in}{0.100000in}}{\pgfqpoint{3.957178in}{3.957178in}}%
\pgfusepath{clip}%
\pgfsetbuttcap%
\pgfsetroundjoin%
\definecolor{currentfill}{rgb}{0.994561,0.734791,0.168938}%
\pgfsetfillcolor{currentfill}%
\pgfsetlinewidth{0.000000pt}%
\definecolor{currentstroke}{rgb}{0.000000,0.000000,0.000000}%
\pgfsetstrokecolor{currentstroke}%
\pgfsetdash{}{0pt}%
\pgfpathmoveto{\pgfqpoint{1.836764in}{3.067780in}}%
\pgfpathlineto{\pgfqpoint{1.877921in}{3.091334in}}%
\pgfpathlineto{\pgfqpoint{1.916062in}{2.844087in}}%
\pgfpathlineto{\pgfqpoint{1.836764in}{3.067780in}}%
\pgfpathclose%
\pgfusepath{fill}%
\end{pgfscope}%
\begin{pgfscope}%
\pgfpathrectangle{\pgfqpoint{0.100000in}{0.100000in}}{\pgfqpoint{3.957178in}{3.957178in}}%
\pgfusepath{clip}%
\pgfsetbuttcap%
\pgfsetroundjoin%
\definecolor{currentfill}{rgb}{0.993033,0.771720,0.154808}%
\pgfsetfillcolor{currentfill}%
\pgfsetlinewidth{0.000000pt}%
\definecolor{currentstroke}{rgb}{0.000000,0.000000,0.000000}%
\pgfsetstrokecolor{currentstroke}%
\pgfsetdash{}{0pt}%
\pgfpathmoveto{\pgfqpoint{2.379181in}{3.032853in}}%
\pgfpathlineto{\pgfqpoint{2.397217in}{2.994070in}}%
\pgfpathlineto{\pgfqpoint{2.264880in}{3.086674in}}%
\pgfpathlineto{\pgfqpoint{2.379181in}{3.032853in}}%
\pgfpathclose%
\pgfusepath{fill}%
\end{pgfscope}%
\begin{pgfscope}%
\pgfpathrectangle{\pgfqpoint{0.100000in}{0.100000in}}{\pgfqpoint{3.957178in}{3.957178in}}%
\pgfusepath{clip}%
\pgfsetbuttcap%
\pgfsetroundjoin%
\definecolor{currentfill}{rgb}{0.752312,0.227133,0.513149}%
\pgfsetfillcolor{currentfill}%
\pgfsetlinewidth{0.000000pt}%
\definecolor{currentstroke}{rgb}{0.000000,0.000000,0.000000}%
\pgfsetstrokecolor{currentstroke}%
\pgfsetdash{}{0pt}%
\pgfpathmoveto{\pgfqpoint{2.968450in}{2.738710in}}%
\pgfpathlineto{\pgfqpoint{3.282578in}{2.232451in}}%
\pgfpathlineto{\pgfqpoint{3.328876in}{2.067735in}}%
\pgfpathlineto{\pgfqpoint{2.968450in}{2.738710in}}%
\pgfpathclose%
\pgfusepath{fill}%
\end{pgfscope}%
\begin{pgfscope}%
\pgfpathrectangle{\pgfqpoint{0.100000in}{0.100000in}}{\pgfqpoint{3.957178in}{3.957178in}}%
\pgfusepath{clip}%
\pgfsetbuttcap%
\pgfsetroundjoin%
\definecolor{currentfill}{rgb}{0.968526,0.569700,0.261721}%
\pgfsetfillcolor{currentfill}%
\pgfsetlinewidth{0.000000pt}%
\definecolor{currentstroke}{rgb}{0.000000,0.000000,0.000000}%
\pgfsetstrokecolor{currentstroke}%
\pgfsetdash{}{0pt}%
\pgfpathmoveto{\pgfqpoint{2.848634in}{2.776341in}}%
\pgfpathlineto{\pgfqpoint{2.805314in}{2.745587in}}%
\pgfpathlineto{\pgfqpoint{2.786897in}{2.893114in}}%
\pgfpathlineto{\pgfqpoint{2.848634in}{2.776341in}}%
\pgfpathclose%
\pgfusepath{fill}%
\end{pgfscope}%
\begin{pgfscope}%
\pgfpathrectangle{\pgfqpoint{0.100000in}{0.100000in}}{\pgfqpoint{3.957178in}{3.957178in}}%
\pgfusepath{clip}%
\pgfsetbuttcap%
\pgfsetroundjoin%
\definecolor{currentfill}{rgb}{0.994474,0.722691,0.174381}%
\pgfsetfillcolor{currentfill}%
\pgfsetlinewidth{0.000000pt}%
\definecolor{currentstroke}{rgb}{0.000000,0.000000,0.000000}%
\pgfsetstrokecolor{currentstroke}%
\pgfsetdash{}{0pt}%
\pgfpathmoveto{\pgfqpoint{2.586979in}{3.035021in}}%
\pgfpathlineto{\pgfqpoint{2.661682in}{2.887953in}}%
\pgfpathlineto{\pgfqpoint{2.545300in}{2.957675in}}%
\pgfpathlineto{\pgfqpoint{2.586979in}{3.035021in}}%
\pgfpathclose%
\pgfusepath{fill}%
\end{pgfscope}%
\begin{pgfscope}%
\pgfpathrectangle{\pgfqpoint{0.100000in}{0.100000in}}{\pgfqpoint{3.957178in}{3.957178in}}%
\pgfusepath{clip}%
\pgfsetbuttcap%
\pgfsetroundjoin%
\definecolor{currentfill}{rgb}{0.529306,0.027747,0.651586}%
\pgfsetfillcolor{currentfill}%
\pgfsetlinewidth{0.000000pt}%
\definecolor{currentstroke}{rgb}{0.000000,0.000000,0.000000}%
\pgfsetstrokecolor{currentstroke}%
\pgfsetdash{}{0pt}%
\pgfpathmoveto{\pgfqpoint{3.386039in}{2.033520in}}%
\pgfpathlineto{\pgfqpoint{3.505448in}{1.816784in}}%
\pgfpathlineto{\pgfqpoint{3.282578in}{2.232451in}}%
\pgfpathlineto{\pgfqpoint{3.386039in}{2.033520in}}%
\pgfpathclose%
\pgfusepath{fill}%
\end{pgfscope}%
\begin{pgfscope}%
\pgfpathrectangle{\pgfqpoint{0.100000in}{0.100000in}}{\pgfqpoint{3.957178in}{3.957178in}}%
\pgfusepath{clip}%
\pgfsetbuttcap%
\pgfsetroundjoin%
\definecolor{currentfill}{rgb}{0.992505,0.777967,0.152855}%
\pgfsetfillcolor{currentfill}%
\pgfsetlinewidth{0.000000pt}%
\definecolor{currentstroke}{rgb}{0.000000,0.000000,0.000000}%
\pgfsetstrokecolor{currentstroke}%
\pgfsetdash{}{0pt}%
\pgfpathmoveto{\pgfqpoint{2.545300in}{2.957675in}}%
\pgfpathlineto{\pgfqpoint{2.465371in}{3.006355in}}%
\pgfpathlineto{\pgfqpoint{2.502911in}{3.063259in}}%
\pgfpathlineto{\pgfqpoint{2.545300in}{2.957675in}}%
\pgfpathclose%
\pgfusepath{fill}%
\end{pgfscope}%
\begin{pgfscope}%
\pgfpathrectangle{\pgfqpoint{0.100000in}{0.100000in}}{\pgfqpoint{3.957178in}{3.957178in}}%
\pgfusepath{clip}%
\pgfsetbuttcap%
\pgfsetroundjoin%
\definecolor{currentfill}{rgb}{0.982653,0.841812,0.142303}%
\pgfsetfillcolor{currentfill}%
\pgfsetlinewidth{0.000000pt}%
\definecolor{currentstroke}{rgb}{0.000000,0.000000,0.000000}%
\pgfsetstrokecolor{currentstroke}%
\pgfsetdash{}{0pt}%
\pgfpathmoveto{\pgfqpoint{2.379181in}{3.032853in}}%
\pgfpathlineto{\pgfqpoint{2.264880in}{3.086674in}}%
\pgfpathlineto{\pgfqpoint{2.277221in}{3.095792in}}%
\pgfpathlineto{\pgfqpoint{2.379181in}{3.032853in}}%
\pgfpathclose%
\pgfusepath{fill}%
\end{pgfscope}%
\begin{pgfscope}%
\pgfpathrectangle{\pgfqpoint{0.100000in}{0.100000in}}{\pgfqpoint{3.957178in}{3.957178in}}%
\pgfusepath{clip}%
\pgfsetbuttcap%
\pgfsetroundjoin%
\definecolor{currentfill}{rgb}{0.993851,0.759304,0.159092}%
\pgfsetfillcolor{currentfill}%
\pgfsetlinewidth{0.000000pt}%
\definecolor{currentstroke}{rgb}{0.000000,0.000000,0.000000}%
\pgfsetstrokecolor{currentstroke}%
\pgfsetdash{}{0pt}%
\pgfpathmoveto{\pgfqpoint{1.787647in}{3.068063in}}%
\pgfpathlineto{\pgfqpoint{1.791056in}{2.960511in}}%
\pgfpathlineto{\pgfqpoint{1.596907in}{2.947877in}}%
\pgfpathlineto{\pgfqpoint{1.787647in}{3.068063in}}%
\pgfpathclose%
\pgfusepath{fill}%
\end{pgfscope}%
\begin{pgfscope}%
\pgfpathrectangle{\pgfqpoint{0.100000in}{0.100000in}}{\pgfqpoint{3.957178in}{3.957178in}}%
\pgfusepath{clip}%
\pgfsetbuttcap%
\pgfsetroundjoin%
\definecolor{currentfill}{rgb}{0.988648,0.809579,0.145357}%
\pgfsetfillcolor{currentfill}%
\pgfsetlinewidth{0.000000pt}%
\definecolor{currentstroke}{rgb}{0.000000,0.000000,0.000000}%
\pgfsetstrokecolor{currentstroke}%
\pgfsetdash{}{0pt}%
\pgfpathmoveto{\pgfqpoint{2.379181in}{3.032853in}}%
\pgfpathlineto{\pgfqpoint{2.502911in}{3.063259in}}%
\pgfpathlineto{\pgfqpoint{2.465371in}{3.006355in}}%
\pgfpathlineto{\pgfqpoint{2.379181in}{3.032853in}}%
\pgfpathclose%
\pgfusepath{fill}%
\end{pgfscope}%
\begin{pgfscope}%
\pgfpathrectangle{\pgfqpoint{0.100000in}{0.100000in}}{\pgfqpoint{3.957178in}{3.957178in}}%
\pgfusepath{clip}%
\pgfsetbuttcap%
\pgfsetroundjoin%
\definecolor{currentfill}{rgb}{0.423689,0.000646,0.658956}%
\pgfsetfillcolor{currentfill}%
\pgfsetlinewidth{0.000000pt}%
\definecolor{currentstroke}{rgb}{0.000000,0.000000,0.000000}%
\pgfsetstrokecolor{currentstroke}%
\pgfsetdash{}{0pt}%
\pgfpathmoveto{\pgfqpoint{3.386039in}{2.033520in}}%
\pgfpathlineto{\pgfqpoint{3.464151in}{1.860528in}}%
\pgfpathlineto{\pgfqpoint{3.505448in}{1.816784in}}%
\pgfpathlineto{\pgfqpoint{3.386039in}{2.033520in}}%
\pgfpathclose%
\pgfusepath{fill}%
\end{pgfscope}%
\begin{pgfscope}%
\pgfpathrectangle{\pgfqpoint{0.100000in}{0.100000in}}{\pgfqpoint{3.957178in}{3.957178in}}%
\pgfusepath{clip}%
\pgfsetbuttcap%
\pgfsetroundjoin%
\definecolor{currentfill}{rgb}{0.947051,0.512699,0.299049}%
\pgfsetfillcolor{currentfill}%
\pgfsetlinewidth{0.000000pt}%
\definecolor{currentstroke}{rgb}{0.000000,0.000000,0.000000}%
\pgfsetstrokecolor{currentstroke}%
\pgfsetdash{}{0pt}%
\pgfpathmoveto{\pgfqpoint{1.447386in}{2.869624in}}%
\pgfpathlineto{\pgfqpoint{1.402175in}{2.771187in}}%
\pgfpathlineto{\pgfqpoint{1.145690in}{2.475609in}}%
\pgfpathlineto{\pgfqpoint{1.447386in}{2.869624in}}%
\pgfpathclose%
\pgfusepath{fill}%
\end{pgfscope}%
\begin{pgfscope}%
\pgfpathrectangle{\pgfqpoint{0.100000in}{0.100000in}}{\pgfqpoint{3.957178in}{3.957178in}}%
\pgfusepath{clip}%
\pgfsetbuttcap%
\pgfsetroundjoin%
\definecolor{currentfill}{rgb}{0.805467,0.289057,0.462415}%
\pgfsetfillcolor{currentfill}%
\pgfsetlinewidth{0.000000pt}%
\definecolor{currentstroke}{rgb}{0.000000,0.000000,0.000000}%
\pgfsetstrokecolor{currentstroke}%
\pgfsetdash{}{0pt}%
\pgfpathmoveto{\pgfqpoint{1.103235in}{2.431049in}}%
\pgfpathlineto{\pgfqpoint{1.145690in}{2.475609in}}%
\pgfpathlineto{\pgfqpoint{1.017358in}{2.293891in}}%
\pgfpathlineto{\pgfqpoint{1.103235in}{2.431049in}}%
\pgfpathclose%
\pgfusepath{fill}%
\end{pgfscope}%
\begin{pgfscope}%
\pgfpathrectangle{\pgfqpoint{0.100000in}{0.100000in}}{\pgfqpoint{3.957178in}{3.957178in}}%
\pgfusepath{clip}%
\pgfsetbuttcap%
\pgfsetroundjoin%
\definecolor{currentfill}{rgb}{0.988648,0.809579,0.145357}%
\pgfsetfillcolor{currentfill}%
\pgfsetlinewidth{0.000000pt}%
\definecolor{currentstroke}{rgb}{0.000000,0.000000,0.000000}%
\pgfsetstrokecolor{currentstroke}%
\pgfsetdash{}{0pt}%
\pgfpathmoveto{\pgfqpoint{1.836764in}{3.067780in}}%
\pgfpathlineto{\pgfqpoint{1.791056in}{2.960511in}}%
\pgfpathlineto{\pgfqpoint{1.787647in}{3.068063in}}%
\pgfpathlineto{\pgfqpoint{1.836764in}{3.067780in}}%
\pgfpathclose%
\pgfusepath{fill}%
\end{pgfscope}%
\begin{pgfscope}%
\pgfpathrectangle{\pgfqpoint{0.100000in}{0.100000in}}{\pgfqpoint{3.957178in}{3.957178in}}%
\pgfusepath{clip}%
\pgfsetbuttcap%
\pgfsetroundjoin%
\definecolor{currentfill}{rgb}{0.977995,0.861432,0.142808}%
\pgfsetfillcolor{currentfill}%
\pgfsetlinewidth{0.000000pt}%
\definecolor{currentstroke}{rgb}{0.000000,0.000000,0.000000}%
\pgfsetstrokecolor{currentstroke}%
\pgfsetdash{}{0pt}%
\pgfpathmoveto{\pgfqpoint{2.277221in}{3.095792in}}%
\pgfpathlineto{\pgfqpoint{2.286897in}{3.104208in}}%
\pgfpathlineto{\pgfqpoint{2.379181in}{3.032853in}}%
\pgfpathlineto{\pgfqpoint{2.277221in}{3.095792in}}%
\pgfpathclose%
\pgfusepath{fill}%
\end{pgfscope}%
\begin{pgfscope}%
\pgfpathrectangle{\pgfqpoint{0.100000in}{0.100000in}}{\pgfqpoint{3.957178in}{3.957178in}}%
\pgfusepath{clip}%
\pgfsetbuttcap%
\pgfsetroundjoin%
\definecolor{currentfill}{rgb}{0.989128,0.658043,0.208100}%
\pgfsetfillcolor{currentfill}%
\pgfsetlinewidth{0.000000pt}%
\definecolor{currentstroke}{rgb}{0.000000,0.000000,0.000000}%
\pgfsetstrokecolor{currentstroke}%
\pgfsetdash{}{0pt}%
\pgfpathmoveto{\pgfqpoint{1.402175in}{2.771187in}}%
\pgfpathlineto{\pgfqpoint{1.447386in}{2.869624in}}%
\pgfpathlineto{\pgfqpoint{1.596907in}{2.947877in}}%
\pgfpathlineto{\pgfqpoint{1.402175in}{2.771187in}}%
\pgfpathclose%
\pgfusepath{fill}%
\end{pgfscope}%
\begin{pgfscope}%
\pgfpathrectangle{\pgfqpoint{0.100000in}{0.100000in}}{\pgfqpoint{3.957178in}{3.957178in}}%
\pgfusepath{clip}%
\pgfsetbuttcap%
\pgfsetroundjoin%
\definecolor{currentfill}{rgb}{0.968443,0.894564,0.147014}%
\pgfsetfillcolor{currentfill}%
\pgfsetlinewidth{0.000000pt}%
\definecolor{currentstroke}{rgb}{0.000000,0.000000,0.000000}%
\pgfsetstrokecolor{currentstroke}%
\pgfsetdash{}{0pt}%
\pgfpathmoveto{\pgfqpoint{2.277221in}{3.095792in}}%
\pgfpathlineto{\pgfqpoint{2.264880in}{3.086674in}}%
\pgfpathlineto{\pgfqpoint{2.286897in}{3.104208in}}%
\pgfpathlineto{\pgfqpoint{2.277221in}{3.095792in}}%
\pgfpathclose%
\pgfusepath{fill}%
\end{pgfscope}%
\begin{pgfscope}%
\pgfpathrectangle{\pgfqpoint{0.100000in}{0.100000in}}{\pgfqpoint{3.957178in}{3.957178in}}%
\pgfusepath{clip}%
\pgfsetbuttcap%
\pgfsetroundjoin%
\definecolor{currentfill}{rgb}{0.976265,0.868016,0.143351}%
\pgfsetfillcolor{currentfill}%
\pgfsetlinewidth{0.000000pt}%
\definecolor{currentstroke}{rgb}{0.000000,0.000000,0.000000}%
\pgfsetstrokecolor{currentstroke}%
\pgfsetdash{}{0pt}%
\pgfpathmoveto{\pgfqpoint{2.286897in}{3.104208in}}%
\pgfpathlineto{\pgfqpoint{2.502911in}{3.063259in}}%
\pgfpathlineto{\pgfqpoint{2.379181in}{3.032853in}}%
\pgfpathlineto{\pgfqpoint{2.286897in}{3.104208in}}%
\pgfpathclose%
\pgfusepath{fill}%
\end{pgfscope}%
\begin{pgfscope}%
\pgfpathrectangle{\pgfqpoint{0.100000in}{0.100000in}}{\pgfqpoint{3.957178in}{3.957178in}}%
\pgfusepath{clip}%
\pgfsetbuttcap%
\pgfsetroundjoin%
\definecolor{currentfill}{rgb}{0.994561,0.734791,0.168938}%
\pgfsetfillcolor{currentfill}%
\pgfsetlinewidth{0.000000pt}%
\definecolor{currentstroke}{rgb}{0.000000,0.000000,0.000000}%
\pgfsetstrokecolor{currentstroke}%
\pgfsetdash{}{0pt}%
\pgfpathmoveto{\pgfqpoint{2.786897in}{2.893114in}}%
\pgfpathlineto{\pgfqpoint{2.661682in}{2.887953in}}%
\pgfpathlineto{\pgfqpoint{2.630218in}{3.018539in}}%
\pgfpathlineto{\pgfqpoint{2.786897in}{2.893114in}}%
\pgfpathclose%
\pgfusepath{fill}%
\end{pgfscope}%
\begin{pgfscope}%
\pgfpathrectangle{\pgfqpoint{0.100000in}{0.100000in}}{\pgfqpoint{3.957178in}{3.957178in}}%
\pgfusepath{clip}%
\pgfsetbuttcap%
\pgfsetroundjoin%
\definecolor{currentfill}{rgb}{0.809052,0.293491,0.458870}%
\pgfsetfillcolor{currentfill}%
\pgfsetlinewidth{0.000000pt}%
\definecolor{currentstroke}{rgb}{0.000000,0.000000,0.000000}%
\pgfsetstrokecolor{currentstroke}%
\pgfsetdash{}{0pt}%
\pgfpathmoveto{\pgfqpoint{1.017358in}{2.293891in}}%
\pgfpathlineto{\pgfqpoint{1.068273in}{2.394651in}}%
\pgfpathlineto{\pgfqpoint{1.103235in}{2.431049in}}%
\pgfpathlineto{\pgfqpoint{1.017358in}{2.293891in}}%
\pgfpathclose%
\pgfusepath{fill}%
\end{pgfscope}%
\begin{pgfscope}%
\pgfpathrectangle{\pgfqpoint{0.100000in}{0.100000in}}{\pgfqpoint{3.957178in}{3.957178in}}%
\pgfusepath{clip}%
\pgfsetbuttcap%
\pgfsetroundjoin%
\definecolor{currentfill}{rgb}{0.983041,0.624131,0.227937}%
\pgfsetfillcolor{currentfill}%
\pgfsetlinewidth{0.000000pt}%
\definecolor{currentstroke}{rgb}{0.000000,0.000000,0.000000}%
\pgfsetstrokecolor{currentstroke}%
\pgfsetdash{}{0pt}%
\pgfpathmoveto{\pgfqpoint{2.968450in}{2.738710in}}%
\pgfpathlineto{\pgfqpoint{2.848634in}{2.776341in}}%
\pgfpathlineto{\pgfqpoint{2.786897in}{2.893114in}}%
\pgfpathlineto{\pgfqpoint{2.968450in}{2.738710in}}%
\pgfpathclose%
\pgfusepath{fill}%
\end{pgfscope}%
\begin{pgfscope}%
\pgfpathrectangle{\pgfqpoint{0.100000in}{0.100000in}}{\pgfqpoint{3.957178in}{3.957178in}}%
\pgfusepath{clip}%
\pgfsetbuttcap%
\pgfsetroundjoin%
\definecolor{currentfill}{rgb}{0.985314,0.828846,0.142945}%
\pgfsetfillcolor{currentfill}%
\pgfsetlinewidth{0.000000pt}%
\definecolor{currentstroke}{rgb}{0.000000,0.000000,0.000000}%
\pgfsetstrokecolor{currentstroke}%
\pgfsetdash{}{0pt}%
\pgfpathmoveto{\pgfqpoint{2.502911in}{3.063259in}}%
\pgfpathlineto{\pgfqpoint{2.586979in}{3.035021in}}%
\pgfpathlineto{\pgfqpoint{2.545300in}{2.957675in}}%
\pgfpathlineto{\pgfqpoint{2.502911in}{3.063259in}}%
\pgfpathclose%
\pgfusepath{fill}%
\end{pgfscope}%
\begin{pgfscope}%
\pgfpathrectangle{\pgfqpoint{0.100000in}{0.100000in}}{\pgfqpoint{3.957178in}{3.957178in}}%
\pgfusepath{clip}%
\pgfsetbuttcap%
\pgfsetroundjoin%
\definecolor{currentfill}{rgb}{0.961681,0.914672,0.150520}%
\pgfsetfillcolor{currentfill}%
\pgfsetlinewidth{0.000000pt}%
\definecolor{currentstroke}{rgb}{0.000000,0.000000,0.000000}%
\pgfsetstrokecolor{currentstroke}%
\pgfsetdash{}{0pt}%
\pgfpathmoveto{\pgfqpoint{2.286897in}{3.104208in}}%
\pgfpathlineto{\pgfqpoint{2.264880in}{3.086674in}}%
\pgfpathlineto{\pgfqpoint{1.877921in}{3.091334in}}%
\pgfpathlineto{\pgfqpoint{2.286897in}{3.104208in}}%
\pgfpathclose%
\pgfusepath{fill}%
\end{pgfscope}%
\begin{pgfscope}%
\pgfpathrectangle{\pgfqpoint{0.100000in}{0.100000in}}{\pgfqpoint{3.957178in}{3.957178in}}%
\pgfusepath{clip}%
\pgfsetbuttcap%
\pgfsetroundjoin%
\definecolor{currentfill}{rgb}{0.990439,0.796859,0.147870}%
\pgfsetfillcolor{currentfill}%
\pgfsetlinewidth{0.000000pt}%
\definecolor{currentstroke}{rgb}{0.000000,0.000000,0.000000}%
\pgfsetstrokecolor{currentstroke}%
\pgfsetdash{}{0pt}%
\pgfpathmoveto{\pgfqpoint{2.661682in}{2.887953in}}%
\pgfpathlineto{\pgfqpoint{2.586979in}{3.035021in}}%
\pgfpathlineto{\pgfqpoint{2.630218in}{3.018539in}}%
\pgfpathlineto{\pgfqpoint{2.661682in}{2.887953in}}%
\pgfpathclose%
\pgfusepath{fill}%
\end{pgfscope}%
\begin{pgfscope}%
\pgfpathrectangle{\pgfqpoint{0.100000in}{0.100000in}}{\pgfqpoint{3.957178in}{3.957178in}}%
\pgfusepath{clip}%
\pgfsetbuttcap%
\pgfsetroundjoin%
\definecolor{currentfill}{rgb}{0.920714,0.458417,0.336166}%
\pgfsetfillcolor{currentfill}%
\pgfsetlinewidth{0.000000pt}%
\definecolor{currentstroke}{rgb}{0.000000,0.000000,0.000000}%
\pgfsetstrokecolor{currentstroke}%
\pgfsetdash{}{0pt}%
\pgfpathmoveto{\pgfqpoint{1.103235in}{2.431049in}}%
\pgfpathlineto{\pgfqpoint{1.409186in}{2.875647in}}%
\pgfpathlineto{\pgfqpoint{1.145690in}{2.475609in}}%
\pgfpathlineto{\pgfqpoint{1.103235in}{2.431049in}}%
\pgfpathclose%
\pgfusepath{fill}%
\end{pgfscope}%
\begin{pgfscope}%
\pgfpathrectangle{\pgfqpoint{0.100000in}{0.100000in}}{\pgfqpoint{3.957178in}{3.957178in}}%
\pgfusepath{clip}%
\pgfsetbuttcap%
\pgfsetroundjoin%
\definecolor{currentfill}{rgb}{0.979644,0.854866,0.142453}%
\pgfsetfillcolor{currentfill}%
\pgfsetlinewidth{0.000000pt}%
\definecolor{currentstroke}{rgb}{0.000000,0.000000,0.000000}%
\pgfsetstrokecolor{currentstroke}%
\pgfsetdash{}{0pt}%
\pgfpathmoveto{\pgfqpoint{1.596907in}{2.947877in}}%
\pgfpathlineto{\pgfqpoint{1.766836in}{3.066415in}}%
\pgfpathlineto{\pgfqpoint{1.787647in}{3.068063in}}%
\pgfpathlineto{\pgfqpoint{1.596907in}{2.947877in}}%
\pgfpathclose%
\pgfusepath{fill}%
\end{pgfscope}%
\begin{pgfscope}%
\pgfpathrectangle{\pgfqpoint{0.100000in}{0.100000in}}{\pgfqpoint{3.957178in}{3.957178in}}%
\pgfusepath{clip}%
\pgfsetbuttcap%
\pgfsetroundjoin%
\definecolor{currentfill}{rgb}{0.964021,0.907950,0.149370}%
\pgfsetfillcolor{currentfill}%
\pgfsetlinewidth{0.000000pt}%
\definecolor{currentstroke}{rgb}{0.000000,0.000000,0.000000}%
\pgfsetstrokecolor{currentstroke}%
\pgfsetdash{}{0pt}%
\pgfpathmoveto{\pgfqpoint{1.787647in}{3.068063in}}%
\pgfpathlineto{\pgfqpoint{1.877921in}{3.091334in}}%
\pgfpathlineto{\pgfqpoint{1.836764in}{3.067780in}}%
\pgfpathlineto{\pgfqpoint{1.787647in}{3.068063in}}%
\pgfpathclose%
\pgfusepath{fill}%
\end{pgfscope}%
\begin{pgfscope}%
\pgfpathrectangle{\pgfqpoint{0.100000in}{0.100000in}}{\pgfqpoint{3.957178in}{3.957178in}}%
\pgfusepath{clip}%
\pgfsetbuttcap%
\pgfsetroundjoin%
\definecolor{currentfill}{rgb}{0.701769,0.174005,0.557296}%
\pgfsetfillcolor{currentfill}%
\pgfsetlinewidth{0.000000pt}%
\definecolor{currentstroke}{rgb}{0.000000,0.000000,0.000000}%
\pgfsetstrokecolor{currentstroke}%
\pgfsetdash{}{0pt}%
\pgfpathmoveto{\pgfqpoint{3.282578in}{2.232451in}}%
\pgfpathlineto{\pgfqpoint{3.270626in}{2.259309in}}%
\pgfpathlineto{\pgfqpoint{3.386039in}{2.033520in}}%
\pgfpathlineto{\pgfqpoint{3.282578in}{2.232451in}}%
\pgfpathclose%
\pgfusepath{fill}%
\end{pgfscope}%
\begin{pgfscope}%
\pgfpathrectangle{\pgfqpoint{0.100000in}{0.100000in}}{\pgfqpoint{3.957178in}{3.957178in}}%
\pgfusepath{clip}%
\pgfsetbuttcap%
\pgfsetroundjoin%
\definecolor{currentfill}{rgb}{0.557243,0.047331,0.643443}%
\pgfsetfillcolor{currentfill}%
\pgfsetlinewidth{0.000000pt}%
\definecolor{currentstroke}{rgb}{0.000000,0.000000,0.000000}%
\pgfsetstrokecolor{currentstroke}%
\pgfsetdash{}{0pt}%
\pgfpathmoveto{\pgfqpoint{3.356437in}{2.081890in}}%
\pgfpathlineto{\pgfqpoint{3.464151in}{1.860528in}}%
\pgfpathlineto{\pgfqpoint{3.386039in}{2.033520in}}%
\pgfpathlineto{\pgfqpoint{3.356437in}{2.081890in}}%
\pgfpathclose%
\pgfusepath{fill}%
\end{pgfscope}%
\begin{pgfscope}%
\pgfpathrectangle{\pgfqpoint{0.100000in}{0.100000in}}{\pgfqpoint{3.957178in}{3.957178in}}%
\pgfusepath{clip}%
\pgfsetbuttcap%
\pgfsetroundjoin%
\definecolor{currentfill}{rgb}{0.974947,0.591165,0.248151}%
\pgfsetfillcolor{currentfill}%
\pgfsetlinewidth{0.000000pt}%
\definecolor{currentstroke}{rgb}{0.000000,0.000000,0.000000}%
\pgfsetstrokecolor{currentstroke}%
\pgfsetdash{}{0pt}%
\pgfpathmoveto{\pgfqpoint{1.145690in}{2.475609in}}%
\pgfpathlineto{\pgfqpoint{1.409186in}{2.875647in}}%
\pgfpathlineto{\pgfqpoint{1.447386in}{2.869624in}}%
\pgfpathlineto{\pgfqpoint{1.145690in}{2.475609in}}%
\pgfpathclose%
\pgfusepath{fill}%
\end{pgfscope}%
\begin{pgfscope}%
\pgfpathrectangle{\pgfqpoint{0.100000in}{0.100000in}}{\pgfqpoint{3.957178in}{3.957178in}}%
\pgfusepath{clip}%
\pgfsetbuttcap%
\pgfsetroundjoin%
\definecolor{currentfill}{rgb}{0.878423,0.387932,0.386600}%
\pgfsetfillcolor{currentfill}%
\pgfsetlinewidth{0.000000pt}%
\definecolor{currentstroke}{rgb}{0.000000,0.000000,0.000000}%
\pgfsetstrokecolor{currentstroke}%
\pgfsetdash{}{0pt}%
\pgfpathmoveto{\pgfqpoint{2.968450in}{2.738710in}}%
\pgfpathlineto{\pgfqpoint{3.185542in}{2.423397in}}%
\pgfpathlineto{\pgfqpoint{3.282578in}{2.232451in}}%
\pgfpathlineto{\pgfqpoint{2.968450in}{2.738710in}}%
\pgfpathclose%
\pgfusepath{fill}%
\end{pgfscope}%
\begin{pgfscope}%
\pgfpathrectangle{\pgfqpoint{0.100000in}{0.100000in}}{\pgfqpoint{3.957178in}{3.957178in}}%
\pgfusepath{clip}%
\pgfsetbuttcap%
\pgfsetroundjoin%
\definecolor{currentfill}{rgb}{0.949151,0.948435,0.152178}%
\pgfsetfillcolor{currentfill}%
\pgfsetlinewidth{0.000000pt}%
\definecolor{currentstroke}{rgb}{0.000000,0.000000,0.000000}%
\pgfsetstrokecolor{currentstroke}%
\pgfsetdash{}{0pt}%
\pgfpathmoveto{\pgfqpoint{1.877921in}{3.091334in}}%
\pgfpathlineto{\pgfqpoint{2.285509in}{3.097456in}}%
\pgfpathlineto{\pgfqpoint{2.286897in}{3.104208in}}%
\pgfpathlineto{\pgfqpoint{1.877921in}{3.091334in}}%
\pgfpathclose%
\pgfusepath{fill}%
\end{pgfscope}%
\begin{pgfscope}%
\pgfpathrectangle{\pgfqpoint{0.100000in}{0.100000in}}{\pgfqpoint{3.957178in}{3.957178in}}%
\pgfusepath{clip}%
\pgfsetbuttcap%
\pgfsetroundjoin%
\definecolor{currentfill}{rgb}{0.956808,0.928152,0.152409}%
\pgfsetfillcolor{currentfill}%
\pgfsetlinewidth{0.000000pt}%
\definecolor{currentstroke}{rgb}{0.000000,0.000000,0.000000}%
\pgfsetstrokecolor{currentstroke}%
\pgfsetdash{}{0pt}%
\pgfpathmoveto{\pgfqpoint{1.766836in}{3.066415in}}%
\pgfpathlineto{\pgfqpoint{1.877921in}{3.091334in}}%
\pgfpathlineto{\pgfqpoint{1.787647in}{3.068063in}}%
\pgfpathlineto{\pgfqpoint{1.766836in}{3.066415in}}%
\pgfpathclose%
\pgfusepath{fill}%
\end{pgfscope}%
\begin{pgfscope}%
\pgfpathrectangle{\pgfqpoint{0.100000in}{0.100000in}}{\pgfqpoint{3.957178in}{3.957178in}}%
\pgfusepath{clip}%
\pgfsetbuttcap%
\pgfsetroundjoin%
\definecolor{currentfill}{rgb}{0.951726,0.941671,0.152925}%
\pgfsetfillcolor{currentfill}%
\pgfsetlinewidth{0.000000pt}%
\definecolor{currentstroke}{rgb}{0.000000,0.000000,0.000000}%
\pgfsetstrokecolor{currentstroke}%
\pgfsetdash{}{0pt}%
\pgfpathmoveto{\pgfqpoint{2.502911in}{3.063259in}}%
\pgfpathlineto{\pgfqpoint{2.286897in}{3.104208in}}%
\pgfpathlineto{\pgfqpoint{2.285509in}{3.097456in}}%
\pgfpathlineto{\pgfqpoint{2.502911in}{3.063259in}}%
\pgfpathclose%
\pgfusepath{fill}%
\end{pgfscope}%
\begin{pgfscope}%
\pgfpathrectangle{\pgfqpoint{0.100000in}{0.100000in}}{\pgfqpoint{3.957178in}{3.957178in}}%
\pgfusepath{clip}%
\pgfsetbuttcap%
\pgfsetroundjoin%
\definecolor{currentfill}{rgb}{0.994103,0.710698,0.180097}%
\pgfsetfillcolor{currentfill}%
\pgfsetlinewidth{0.000000pt}%
\definecolor{currentstroke}{rgb}{0.000000,0.000000,0.000000}%
\pgfsetstrokecolor{currentstroke}%
\pgfsetdash{}{0pt}%
\pgfpathmoveto{\pgfqpoint{2.786897in}{2.893114in}}%
\pgfpathlineto{\pgfqpoint{2.770237in}{2.929808in}}%
\pgfpathlineto{\pgfqpoint{2.968450in}{2.738710in}}%
\pgfpathlineto{\pgfqpoint{2.786897in}{2.893114in}}%
\pgfpathclose%
\pgfusepath{fill}%
\end{pgfscope}%
\begin{pgfscope}%
\pgfpathrectangle{\pgfqpoint{0.100000in}{0.100000in}}{\pgfqpoint{3.957178in}{3.957178in}}%
\pgfusepath{clip}%
\pgfsetbuttcap%
\pgfsetroundjoin%
\definecolor{currentfill}{rgb}{0.798216,0.280197,0.469538}%
\pgfsetfillcolor{currentfill}%
\pgfsetlinewidth{0.000000pt}%
\definecolor{currentstroke}{rgb}{0.000000,0.000000,0.000000}%
\pgfsetstrokecolor{currentstroke}%
\pgfsetdash{}{0pt}%
\pgfpathmoveto{\pgfqpoint{3.270626in}{2.259309in}}%
\pgfpathlineto{\pgfqpoint{3.282578in}{2.232451in}}%
\pgfpathlineto{\pgfqpoint{3.185542in}{2.423397in}}%
\pgfpathlineto{\pgfqpoint{3.270626in}{2.259309in}}%
\pgfpathclose%
\pgfusepath{fill}%
\end{pgfscope}%
\begin{pgfscope}%
\pgfpathrectangle{\pgfqpoint{0.100000in}{0.100000in}}{\pgfqpoint{3.957178in}{3.957178in}}%
\pgfusepath{clip}%
\pgfsetbuttcap%
\pgfsetroundjoin%
\definecolor{currentfill}{rgb}{0.683758,0.156278,0.571660}%
\pgfsetfillcolor{currentfill}%
\pgfsetlinewidth{0.000000pt}%
\definecolor{currentstroke}{rgb}{0.000000,0.000000,0.000000}%
\pgfsetstrokecolor{currentstroke}%
\pgfsetdash{}{0pt}%
\pgfpathmoveto{\pgfqpoint{3.386039in}{2.033520in}}%
\pgfpathlineto{\pgfqpoint{3.270626in}{2.259309in}}%
\pgfpathlineto{\pgfqpoint{3.356437in}{2.081890in}}%
\pgfpathlineto{\pgfqpoint{3.386039in}{2.033520in}}%
\pgfpathclose%
\pgfusepath{fill}%
\end{pgfscope}%
\begin{pgfscope}%
\pgfpathrectangle{\pgfqpoint{0.100000in}{0.100000in}}{\pgfqpoint{3.957178in}{3.957178in}}%
\pgfusepath{clip}%
\pgfsetbuttcap%
\pgfsetroundjoin%
\definecolor{currentfill}{rgb}{0.991209,0.790537,0.149377}%
\pgfsetfillcolor{currentfill}%
\pgfsetlinewidth{0.000000pt}%
\definecolor{currentstroke}{rgb}{0.000000,0.000000,0.000000}%
\pgfsetstrokecolor{currentstroke}%
\pgfsetdash{}{0pt}%
\pgfpathmoveto{\pgfqpoint{1.596907in}{2.947877in}}%
\pgfpathlineto{\pgfqpoint{1.447386in}{2.869624in}}%
\pgfpathlineto{\pgfqpoint{1.576036in}{2.988808in}}%
\pgfpathlineto{\pgfqpoint{1.596907in}{2.947877in}}%
\pgfpathclose%
\pgfusepath{fill}%
\end{pgfscope}%
\begin{pgfscope}%
\pgfpathrectangle{\pgfqpoint{0.100000in}{0.100000in}}{\pgfqpoint{3.957178in}{3.957178in}}%
\pgfusepath{clip}%
\pgfsetbuttcap%
\pgfsetroundjoin%
\definecolor{currentfill}{rgb}{0.989587,0.803205,0.146529}%
\pgfsetfillcolor{currentfill}%
\pgfsetlinewidth{0.000000pt}%
\definecolor{currentstroke}{rgb}{0.000000,0.000000,0.000000}%
\pgfsetstrokecolor{currentstroke}%
\pgfsetdash{}{0pt}%
\pgfpathmoveto{\pgfqpoint{2.630218in}{3.018539in}}%
\pgfpathlineto{\pgfqpoint{2.770237in}{2.929808in}}%
\pgfpathlineto{\pgfqpoint{2.786897in}{2.893114in}}%
\pgfpathlineto{\pgfqpoint{2.630218in}{3.018539in}}%
\pgfpathclose%
\pgfusepath{fill}%
\end{pgfscope}%
\begin{pgfscope}%
\pgfpathrectangle{\pgfqpoint{0.100000in}{0.100000in}}{\pgfqpoint{3.957178in}{3.957178in}}%
\pgfusepath{clip}%
\pgfsetbuttcap%
\pgfsetroundjoin%
\definecolor{currentfill}{rgb}{0.974443,0.874622,0.144061}%
\pgfsetfillcolor{currentfill}%
\pgfsetlinewidth{0.000000pt}%
\definecolor{currentstroke}{rgb}{0.000000,0.000000,0.000000}%
\pgfsetstrokecolor{currentstroke}%
\pgfsetdash{}{0pt}%
\pgfpathmoveto{\pgfqpoint{1.691543in}{3.028826in}}%
\pgfpathlineto{\pgfqpoint{1.766836in}{3.066415in}}%
\pgfpathlineto{\pgfqpoint{1.596907in}{2.947877in}}%
\pgfpathlineto{\pgfqpoint{1.691543in}{3.028826in}}%
\pgfpathclose%
\pgfusepath{fill}%
\end{pgfscope}%
\begin{pgfscope}%
\pgfpathrectangle{\pgfqpoint{0.100000in}{0.100000in}}{\pgfqpoint{3.957178in}{3.957178in}}%
\pgfusepath{clip}%
\pgfsetbuttcap%
\pgfsetroundjoin%
\definecolor{currentfill}{rgb}{0.809052,0.293491,0.458870}%
\pgfsetfillcolor{currentfill}%
\pgfsetlinewidth{0.000000pt}%
\definecolor{currentstroke}{rgb}{0.000000,0.000000,0.000000}%
\pgfsetstrokecolor{currentstroke}%
\pgfsetdash{}{0pt}%
\pgfpathmoveto{\pgfqpoint{1.017358in}{2.293891in}}%
\pgfpathlineto{\pgfqpoint{1.003470in}{2.252260in}}%
\pgfpathlineto{\pgfqpoint{1.068273in}{2.394651in}}%
\pgfpathlineto{\pgfqpoint{1.017358in}{2.293891in}}%
\pgfpathclose%
\pgfusepath{fill}%
\end{pgfscope}%
\begin{pgfscope}%
\pgfpathrectangle{\pgfqpoint{0.100000in}{0.100000in}}{\pgfqpoint{3.957178in}{3.957178in}}%
\pgfusepath{clip}%
\pgfsetbuttcap%
\pgfsetroundjoin%
\definecolor{currentfill}{rgb}{0.655580,0.129725,0.592317}%
\pgfsetfillcolor{currentfill}%
\pgfsetlinewidth{0.000000pt}%
\definecolor{currentstroke}{rgb}{0.000000,0.000000,0.000000}%
\pgfsetstrokecolor{currentstroke}%
\pgfsetdash{}{0pt}%
\pgfpathmoveto{\pgfqpoint{1.017358in}{2.293891in}}%
\pgfpathlineto{\pgfqpoint{0.812085in}{1.825153in}}%
\pgfpathlineto{\pgfqpoint{0.921128in}{2.073538in}}%
\pgfpathlineto{\pgfqpoint{1.017358in}{2.293891in}}%
\pgfpathclose%
\pgfusepath{fill}%
\end{pgfscope}%
\begin{pgfscope}%
\pgfpathrectangle{\pgfqpoint{0.100000in}{0.100000in}}{\pgfqpoint{3.957178in}{3.957178in}}%
\pgfusepath{clip}%
\pgfsetbuttcap%
\pgfsetroundjoin%
\definecolor{currentfill}{rgb}{0.631017,0.107699,0.608287}%
\pgfsetfillcolor{currentfill}%
\pgfsetlinewidth{0.000000pt}%
\definecolor{currentstroke}{rgb}{0.000000,0.000000,0.000000}%
\pgfsetstrokecolor{currentstroke}%
\pgfsetdash{}{0pt}%
\pgfpathmoveto{\pgfqpoint{3.464151in}{1.860528in}}%
\pgfpathlineto{\pgfqpoint{3.356437in}{2.081890in}}%
\pgfpathlineto{\pgfqpoint{3.313847in}{2.143572in}}%
\pgfpathlineto{\pgfqpoint{3.464151in}{1.860528in}}%
\pgfpathclose%
\pgfusepath{fill}%
\end{pgfscope}%
\begin{pgfscope}%
\pgfpathrectangle{\pgfqpoint{0.100000in}{0.100000in}}{\pgfqpoint{3.957178in}{3.957178in}}%
\pgfusepath{clip}%
\pgfsetbuttcap%
\pgfsetroundjoin%
\definecolor{currentfill}{rgb}{0.925825,0.468103,0.329435}%
\pgfsetfillcolor{currentfill}%
\pgfsetlinewidth{0.000000pt}%
\definecolor{currentstroke}{rgb}{0.000000,0.000000,0.000000}%
\pgfsetstrokecolor{currentstroke}%
\pgfsetdash{}{0pt}%
\pgfpathmoveto{\pgfqpoint{1.335079in}{2.804114in}}%
\pgfpathlineto{\pgfqpoint{1.103235in}{2.431049in}}%
\pgfpathlineto{\pgfqpoint{1.068273in}{2.394651in}}%
\pgfpathlineto{\pgfqpoint{1.335079in}{2.804114in}}%
\pgfpathclose%
\pgfusepath{fill}%
\end{pgfscope}%
\begin{pgfscope}%
\pgfpathrectangle{\pgfqpoint{0.100000in}{0.100000in}}{\pgfqpoint{3.957178in}{3.957178in}}%
\pgfusepath{clip}%
\pgfsetbuttcap%
\pgfsetroundjoin%
\definecolor{currentfill}{rgb}{0.930798,0.477867,0.322697}%
\pgfsetfillcolor{currentfill}%
\pgfsetlinewidth{0.000000pt}%
\definecolor{currentstroke}{rgb}{0.000000,0.000000,0.000000}%
\pgfsetstrokecolor{currentstroke}%
\pgfsetdash{}{0pt}%
\pgfpathmoveto{\pgfqpoint{3.154229in}{2.481387in}}%
\pgfpathlineto{\pgfqpoint{3.185542in}{2.423397in}}%
\pgfpathlineto{\pgfqpoint{2.968450in}{2.738710in}}%
\pgfpathlineto{\pgfqpoint{3.154229in}{2.481387in}}%
\pgfpathclose%
\pgfusepath{fill}%
\end{pgfscope}%
\begin{pgfscope}%
\pgfpathrectangle{\pgfqpoint{0.100000in}{0.100000in}}{\pgfqpoint{3.957178in}{3.957178in}}%
\pgfusepath{clip}%
\pgfsetbuttcap%
\pgfsetroundjoin%
\definecolor{currentfill}{rgb}{0.944152,0.961916,0.146861}%
\pgfsetfillcolor{currentfill}%
\pgfsetlinewidth{0.000000pt}%
\definecolor{currentstroke}{rgb}{0.000000,0.000000,0.000000}%
\pgfsetstrokecolor{currentstroke}%
\pgfsetdash{}{0pt}%
\pgfpathmoveto{\pgfqpoint{2.422027in}{3.048101in}}%
\pgfpathlineto{\pgfqpoint{2.502911in}{3.063259in}}%
\pgfpathlineto{\pgfqpoint{2.285509in}{3.097456in}}%
\pgfpathlineto{\pgfqpoint{2.422027in}{3.048101in}}%
\pgfpathclose%
\pgfusepath{fill}%
\end{pgfscope}%
\begin{pgfscope}%
\pgfpathrectangle{\pgfqpoint{0.100000in}{0.100000in}}{\pgfqpoint{3.957178in}{3.957178in}}%
\pgfusepath{clip}%
\pgfsetbuttcap%
\pgfsetroundjoin%
\definecolor{currentfill}{rgb}{0.954287,0.934908,0.152921}%
\pgfsetfillcolor{currentfill}%
\pgfsetlinewidth{0.000000pt}%
\definecolor{currentstroke}{rgb}{0.000000,0.000000,0.000000}%
\pgfsetstrokecolor{currentstroke}%
\pgfsetdash{}{0pt}%
\pgfpathmoveto{\pgfqpoint{2.502911in}{3.063259in}}%
\pgfpathlineto{\pgfqpoint{2.484697in}{3.037406in}}%
\pgfpathlineto{\pgfqpoint{2.586979in}{3.035021in}}%
\pgfpathlineto{\pgfqpoint{2.502911in}{3.063259in}}%
\pgfpathclose%
\pgfusepath{fill}%
\end{pgfscope}%
\begin{pgfscope}%
\pgfpathrectangle{\pgfqpoint{0.100000in}{0.100000in}}{\pgfqpoint{3.957178in}{3.957178in}}%
\pgfusepath{clip}%
\pgfsetbuttcap%
\pgfsetroundjoin%
\definecolor{currentfill}{rgb}{0.959276,0.921407,0.151566}%
\pgfsetfillcolor{currentfill}%
\pgfsetlinewidth{0.000000pt}%
\definecolor{currentstroke}{rgb}{0.000000,0.000000,0.000000}%
\pgfsetstrokecolor{currentstroke}%
\pgfsetdash{}{0pt}%
\pgfpathmoveto{\pgfqpoint{2.054934in}{2.904720in}}%
\pgfpathlineto{\pgfqpoint{2.285509in}{3.097456in}}%
\pgfpathlineto{\pgfqpoint{1.877921in}{3.091334in}}%
\pgfpathlineto{\pgfqpoint{2.054934in}{2.904720in}}%
\pgfpathclose%
\pgfusepath{fill}%
\end{pgfscope}%
\begin{pgfscope}%
\pgfpathrectangle{\pgfqpoint{0.100000in}{0.100000in}}{\pgfqpoint{3.957178in}{3.957178in}}%
\pgfusepath{clip}%
\pgfsetbuttcap%
\pgfsetroundjoin%
\definecolor{currentfill}{rgb}{0.862927,0.365119,0.403519}%
\pgfsetfillcolor{currentfill}%
\pgfsetlinewidth{0.000000pt}%
\definecolor{currentstroke}{rgb}{0.000000,0.000000,0.000000}%
\pgfsetstrokecolor{currentstroke}%
\pgfsetdash{}{0pt}%
\pgfpathmoveto{\pgfqpoint{3.185542in}{2.423397in}}%
\pgfpathlineto{\pgfqpoint{3.154229in}{2.481387in}}%
\pgfpathlineto{\pgfqpoint{3.270626in}{2.259309in}}%
\pgfpathlineto{\pgfqpoint{3.185542in}{2.423397in}}%
\pgfpathclose%
\pgfusepath{fill}%
\end{pgfscope}%
\begin{pgfscope}%
\pgfpathrectangle{\pgfqpoint{0.100000in}{0.100000in}}{\pgfqpoint{3.957178in}{3.957178in}}%
\pgfusepath{clip}%
\pgfsetbuttcap%
\pgfsetroundjoin%
\definecolor{currentfill}{rgb}{0.974443,0.874622,0.144061}%
\pgfsetfillcolor{currentfill}%
\pgfsetlinewidth{0.000000pt}%
\definecolor{currentstroke}{rgb}{0.000000,0.000000,0.000000}%
\pgfsetstrokecolor{currentstroke}%
\pgfsetdash{}{0pt}%
\pgfpathmoveto{\pgfqpoint{2.671623in}{3.002540in}}%
\pgfpathlineto{\pgfqpoint{2.770237in}{2.929808in}}%
\pgfpathlineto{\pgfqpoint{2.630218in}{3.018539in}}%
\pgfpathlineto{\pgfqpoint{2.671623in}{3.002540in}}%
\pgfpathclose%
\pgfusepath{fill}%
\end{pgfscope}%
\begin{pgfscope}%
\pgfpathrectangle{\pgfqpoint{0.100000in}{0.100000in}}{\pgfqpoint{3.957178in}{3.957178in}}%
\pgfusepath{clip}%
\pgfsetbuttcap%
\pgfsetroundjoin%
\definecolor{currentfill}{rgb}{0.949151,0.948435,0.152178}%
\pgfsetfillcolor{currentfill}%
\pgfsetlinewidth{0.000000pt}%
\definecolor{currentstroke}{rgb}{0.000000,0.000000,0.000000}%
\pgfsetstrokecolor{currentstroke}%
\pgfsetdash{}{0pt}%
\pgfpathmoveto{\pgfqpoint{1.877921in}{3.091334in}}%
\pgfpathlineto{\pgfqpoint{1.766836in}{3.066415in}}%
\pgfpathlineto{\pgfqpoint{1.762322in}{3.005647in}}%
\pgfpathlineto{\pgfqpoint{1.877921in}{3.091334in}}%
\pgfpathclose%
\pgfusepath{fill}%
\end{pgfscope}%
\begin{pgfscope}%
\pgfpathrectangle{\pgfqpoint{0.100000in}{0.100000in}}{\pgfqpoint{3.957178in}{3.957178in}}%
\pgfusepath{clip}%
\pgfsetbuttcap%
\pgfsetroundjoin%
\definecolor{currentfill}{rgb}{0.972530,0.881250,0.144923}%
\pgfsetfillcolor{currentfill}%
\pgfsetlinewidth{0.000000pt}%
\definecolor{currentstroke}{rgb}{0.000000,0.000000,0.000000}%
\pgfsetstrokecolor{currentstroke}%
\pgfsetdash{}{0pt}%
\pgfpathmoveto{\pgfqpoint{1.576036in}{2.988808in}}%
\pgfpathlineto{\pgfqpoint{1.691543in}{3.028826in}}%
\pgfpathlineto{\pgfqpoint{1.596907in}{2.947877in}}%
\pgfpathlineto{\pgfqpoint{1.576036in}{2.988808in}}%
\pgfpathclose%
\pgfusepath{fill}%
\end{pgfscope}%
\begin{pgfscope}%
\pgfpathrectangle{\pgfqpoint{0.100000in}{0.100000in}}{\pgfqpoint{3.957178in}{3.957178in}}%
\pgfusepath{clip}%
\pgfsetbuttcap%
\pgfsetroundjoin%
\definecolor{currentfill}{rgb}{0.768090,0.244817,0.498465}%
\pgfsetfillcolor{currentfill}%
\pgfsetlinewidth{0.000000pt}%
\definecolor{currentstroke}{rgb}{0.000000,0.000000,0.000000}%
\pgfsetstrokecolor{currentstroke}%
\pgfsetdash{}{0pt}%
\pgfpathmoveto{\pgfqpoint{0.921128in}{2.073538in}}%
\pgfpathlineto{\pgfqpoint{1.003470in}{2.252260in}}%
\pgfpathlineto{\pgfqpoint{1.017358in}{2.293891in}}%
\pgfpathlineto{\pgfqpoint{0.921128in}{2.073538in}}%
\pgfpathclose%
\pgfusepath{fill}%
\end{pgfscope}%
\begin{pgfscope}%
\pgfpathrectangle{\pgfqpoint{0.100000in}{0.100000in}}{\pgfqpoint{3.957178in}{3.957178in}}%
\pgfusepath{clip}%
\pgfsetbuttcap%
\pgfsetroundjoin%
\definecolor{currentfill}{rgb}{0.956808,0.928152,0.152409}%
\pgfsetfillcolor{currentfill}%
\pgfsetlinewidth{0.000000pt}%
\definecolor{currentstroke}{rgb}{0.000000,0.000000,0.000000}%
\pgfsetstrokecolor{currentstroke}%
\pgfsetdash{}{0pt}%
\pgfpathmoveto{\pgfqpoint{2.586979in}{3.035021in}}%
\pgfpathlineto{\pgfqpoint{2.484697in}{3.037406in}}%
\pgfpathlineto{\pgfqpoint{2.630218in}{3.018539in}}%
\pgfpathlineto{\pgfqpoint{2.586979in}{3.035021in}}%
\pgfpathclose%
\pgfusepath{fill}%
\end{pgfscope}%
\begin{pgfscope}%
\pgfpathrectangle{\pgfqpoint{0.100000in}{0.100000in}}{\pgfqpoint{3.957178in}{3.957178in}}%
\pgfusepath{clip}%
\pgfsetbuttcap%
\pgfsetroundjoin%
\definecolor{currentfill}{rgb}{0.987621,0.815978,0.144363}%
\pgfsetfillcolor{currentfill}%
\pgfsetlinewidth{0.000000pt}%
\definecolor{currentstroke}{rgb}{0.000000,0.000000,0.000000}%
\pgfsetstrokecolor{currentstroke}%
\pgfsetdash{}{0pt}%
\pgfpathmoveto{\pgfqpoint{1.447386in}{2.869624in}}%
\pgfpathlineto{\pgfqpoint{1.409186in}{2.875647in}}%
\pgfpathlineto{\pgfqpoint{1.576036in}{2.988808in}}%
\pgfpathlineto{\pgfqpoint{1.447386in}{2.869624in}}%
\pgfpathclose%
\pgfusepath{fill}%
\end{pgfscope}%
\begin{pgfscope}%
\pgfpathrectangle{\pgfqpoint{0.100000in}{0.100000in}}{\pgfqpoint{3.957178in}{3.957178in}}%
\pgfusepath{clip}%
\pgfsetbuttcap%
\pgfsetroundjoin%
\definecolor{currentfill}{rgb}{0.744232,0.218288,0.520524}%
\pgfsetfillcolor{currentfill}%
\pgfsetlinewidth{0.000000pt}%
\definecolor{currentstroke}{rgb}{0.000000,0.000000,0.000000}%
\pgfsetstrokecolor{currentstroke}%
\pgfsetdash{}{0pt}%
\pgfpathmoveto{\pgfqpoint{3.313847in}{2.143572in}}%
\pgfpathlineto{\pgfqpoint{3.356437in}{2.081890in}}%
\pgfpathlineto{\pgfqpoint{3.270626in}{2.259309in}}%
\pgfpathlineto{\pgfqpoint{3.313847in}{2.143572in}}%
\pgfpathclose%
\pgfusepath{fill}%
\end{pgfscope}%
\begin{pgfscope}%
\pgfpathrectangle{\pgfqpoint{0.100000in}{0.100000in}}{\pgfqpoint{3.957178in}{3.957178in}}%
\pgfusepath{clip}%
\pgfsetbuttcap%
\pgfsetroundjoin%
\definecolor{currentfill}{rgb}{0.994561,0.734791,0.168938}%
\pgfsetfillcolor{currentfill}%
\pgfsetlinewidth{0.000000pt}%
\definecolor{currentstroke}{rgb}{0.000000,0.000000,0.000000}%
\pgfsetstrokecolor{currentstroke}%
\pgfsetdash{}{0pt}%
\pgfpathmoveto{\pgfqpoint{2.770237in}{2.929808in}}%
\pgfpathlineto{\pgfqpoint{2.917573in}{2.822478in}}%
\pgfpathlineto{\pgfqpoint{2.968450in}{2.738710in}}%
\pgfpathlineto{\pgfqpoint{2.770237in}{2.929808in}}%
\pgfpathclose%
\pgfusepath{fill}%
\end{pgfscope}%
\begin{pgfscope}%
\pgfpathrectangle{\pgfqpoint{0.100000in}{0.100000in}}{\pgfqpoint{3.957178in}{3.957178in}}%
\pgfusepath{clip}%
\pgfsetbuttcap%
\pgfsetroundjoin%
\definecolor{currentfill}{rgb}{0.983041,0.624131,0.227937}%
\pgfsetfillcolor{currentfill}%
\pgfsetlinewidth{0.000000pt}%
\definecolor{currentstroke}{rgb}{0.000000,0.000000,0.000000}%
\pgfsetstrokecolor{currentstroke}%
\pgfsetdash{}{0pt}%
\pgfpathmoveto{\pgfqpoint{1.409186in}{2.875647in}}%
\pgfpathlineto{\pgfqpoint{1.103235in}{2.431049in}}%
\pgfpathlineto{\pgfqpoint{1.335079in}{2.804114in}}%
\pgfpathlineto{\pgfqpoint{1.409186in}{2.875647in}}%
\pgfpathclose%
\pgfusepath{fill}%
\end{pgfscope}%
\begin{pgfscope}%
\pgfpathrectangle{\pgfqpoint{0.100000in}{0.100000in}}{\pgfqpoint{3.957178in}{3.957178in}}%
\pgfusepath{clip}%
\pgfsetbuttcap%
\pgfsetroundjoin%
\definecolor{currentfill}{rgb}{0.944152,0.961916,0.146861}%
\pgfsetfillcolor{currentfill}%
\pgfsetlinewidth{0.000000pt}%
\definecolor{currentstroke}{rgb}{0.000000,0.000000,0.000000}%
\pgfsetstrokecolor{currentstroke}%
\pgfsetdash{}{0pt}%
\pgfpathmoveto{\pgfqpoint{2.484697in}{3.037406in}}%
\pgfpathlineto{\pgfqpoint{2.502911in}{3.063259in}}%
\pgfpathlineto{\pgfqpoint{2.422027in}{3.048101in}}%
\pgfpathlineto{\pgfqpoint{2.484697in}{3.037406in}}%
\pgfpathclose%
\pgfusepath{fill}%
\end{pgfscope}%
\begin{pgfscope}%
\pgfpathrectangle{\pgfqpoint{0.100000in}{0.100000in}}{\pgfqpoint{3.957178in}{3.957178in}}%
\pgfusepath{clip}%
\pgfsetbuttcap%
\pgfsetroundjoin%
\definecolor{currentfill}{rgb}{0.980556,0.613039,0.234646}%
\pgfsetfillcolor{currentfill}%
\pgfsetlinewidth{0.000000pt}%
\definecolor{currentstroke}{rgb}{0.000000,0.000000,0.000000}%
\pgfsetstrokecolor{currentstroke}%
\pgfsetdash{}{0pt}%
\pgfpathmoveto{\pgfqpoint{2.968450in}{2.738710in}}%
\pgfpathlineto{\pgfqpoint{2.917573in}{2.822478in}}%
\pgfpathlineto{\pgfqpoint{3.154229in}{2.481387in}}%
\pgfpathlineto{\pgfqpoint{2.968450in}{2.738710in}}%
\pgfpathclose%
\pgfusepath{fill}%
\end{pgfscope}%
\begin{pgfscope}%
\pgfpathrectangle{\pgfqpoint{0.100000in}{0.100000in}}{\pgfqpoint{3.957178in}{3.957178in}}%
\pgfusepath{clip}%
\pgfsetbuttcap%
\pgfsetroundjoin%
\definecolor{currentfill}{rgb}{0.954287,0.934908,0.152921}%
\pgfsetfillcolor{currentfill}%
\pgfsetlinewidth{0.000000pt}%
\definecolor{currentstroke}{rgb}{0.000000,0.000000,0.000000}%
\pgfsetstrokecolor{currentstroke}%
\pgfsetdash{}{0pt}%
\pgfpathmoveto{\pgfqpoint{2.484697in}{3.037406in}}%
\pgfpathlineto{\pgfqpoint{2.671623in}{3.002540in}}%
\pgfpathlineto{\pgfqpoint{2.630218in}{3.018539in}}%
\pgfpathlineto{\pgfqpoint{2.484697in}{3.037406in}}%
\pgfpathclose%
\pgfusepath{fill}%
\end{pgfscope}%
\begin{pgfscope}%
\pgfpathrectangle{\pgfqpoint{0.100000in}{0.100000in}}{\pgfqpoint{3.957178in}{3.957178in}}%
\pgfusepath{clip}%
\pgfsetbuttcap%
\pgfsetroundjoin%
\definecolor{currentfill}{rgb}{0.679160,0.151848,0.575189}%
\pgfsetfillcolor{currentfill}%
\pgfsetlinewidth{0.000000pt}%
\definecolor{currentstroke}{rgb}{0.000000,0.000000,0.000000}%
\pgfsetstrokecolor{currentstroke}%
\pgfsetdash{}{0pt}%
\pgfpathmoveto{\pgfqpoint{3.313847in}{2.143572in}}%
\pgfpathlineto{\pgfqpoint{3.296283in}{2.159348in}}%
\pgfpathlineto{\pgfqpoint{3.464151in}{1.860528in}}%
\pgfpathlineto{\pgfqpoint{3.313847in}{2.143572in}}%
\pgfpathclose%
\pgfusepath{fill}%
\end{pgfscope}%
\begin{pgfscope}%
\pgfpathrectangle{\pgfqpoint{0.100000in}{0.100000in}}{\pgfqpoint{3.957178in}{3.957178in}}%
\pgfusepath{clip}%
\pgfsetbuttcap%
\pgfsetroundjoin%
\definecolor{currentfill}{rgb}{0.898984,0.420392,0.363047}%
\pgfsetfillcolor{currentfill}%
\pgfsetlinewidth{0.000000pt}%
\definecolor{currentstroke}{rgb}{0.000000,0.000000,0.000000}%
\pgfsetstrokecolor{currentstroke}%
\pgfsetdash{}{0pt}%
\pgfpathmoveto{\pgfqpoint{3.270626in}{2.259309in}}%
\pgfpathlineto{\pgfqpoint{3.154229in}{2.481387in}}%
\pgfpathlineto{\pgfqpoint{3.111123in}{2.556569in}}%
\pgfpathlineto{\pgfqpoint{3.270626in}{2.259309in}}%
\pgfpathclose%
\pgfusepath{fill}%
\end{pgfscope}%
\begin{pgfscope}%
\pgfpathrectangle{\pgfqpoint{0.100000in}{0.100000in}}{\pgfqpoint{3.957178in}{3.957178in}}%
\pgfusepath{clip}%
\pgfsetbuttcap%
\pgfsetroundjoin%
\definecolor{currentfill}{rgb}{0.964021,0.907950,0.149370}%
\pgfsetfillcolor{currentfill}%
\pgfsetlinewidth{0.000000pt}%
\definecolor{currentstroke}{rgb}{0.000000,0.000000,0.000000}%
\pgfsetstrokecolor{currentstroke}%
\pgfsetdash{}{0pt}%
\pgfpathmoveto{\pgfqpoint{1.877921in}{3.091334in}}%
\pgfpathlineto{\pgfqpoint{1.842130in}{2.944323in}}%
\pgfpathlineto{\pgfqpoint{2.054934in}{2.904720in}}%
\pgfpathlineto{\pgfqpoint{1.877921in}{3.091334in}}%
\pgfpathclose%
\pgfusepath{fill}%
\end{pgfscope}%
\begin{pgfscope}%
\pgfpathrectangle{\pgfqpoint{0.100000in}{0.100000in}}{\pgfqpoint{3.957178in}{3.957178in}}%
\pgfusepath{clip}%
\pgfsetbuttcap%
\pgfsetroundjoin%
\definecolor{currentfill}{rgb}{0.944152,0.961916,0.146861}%
\pgfsetfillcolor{currentfill}%
\pgfsetlinewidth{0.000000pt}%
\definecolor{currentstroke}{rgb}{0.000000,0.000000,0.000000}%
\pgfsetstrokecolor{currentstroke}%
\pgfsetdash{}{0pt}%
\pgfpathmoveto{\pgfqpoint{1.762322in}{3.005647in}}%
\pgfpathlineto{\pgfqpoint{1.766836in}{3.066415in}}%
\pgfpathlineto{\pgfqpoint{1.691543in}{3.028826in}}%
\pgfpathlineto{\pgfqpoint{1.762322in}{3.005647in}}%
\pgfpathclose%
\pgfusepath{fill}%
\end{pgfscope}%
\begin{pgfscope}%
\pgfpathrectangle{\pgfqpoint{0.100000in}{0.100000in}}{\pgfqpoint{3.957178in}{3.957178in}}%
\pgfusepath{clip}%
\pgfsetbuttcap%
\pgfsetroundjoin%
\definecolor{currentfill}{rgb}{0.946602,0.955190,0.150328}%
\pgfsetfillcolor{currentfill}%
\pgfsetlinewidth{0.000000pt}%
\definecolor{currentstroke}{rgb}{0.000000,0.000000,0.000000}%
\pgfsetstrokecolor{currentstroke}%
\pgfsetdash{}{0pt}%
\pgfpathmoveto{\pgfqpoint{2.422027in}{3.048101in}}%
\pgfpathlineto{\pgfqpoint{2.285509in}{3.097456in}}%
\pgfpathlineto{\pgfqpoint{2.403569in}{2.912894in}}%
\pgfpathlineto{\pgfqpoint{2.422027in}{3.048101in}}%
\pgfpathclose%
\pgfusepath{fill}%
\end{pgfscope}%
\begin{pgfscope}%
\pgfpathrectangle{\pgfqpoint{0.100000in}{0.100000in}}{\pgfqpoint{3.957178in}{3.957178in}}%
\pgfusepath{clip}%
\pgfsetbuttcap%
\pgfsetroundjoin%
\definecolor{currentfill}{rgb}{0.650746,0.125309,0.595617}%
\pgfsetfillcolor{currentfill}%
\pgfsetlinewidth{0.000000pt}%
\definecolor{currentstroke}{rgb}{0.000000,0.000000,0.000000}%
\pgfsetstrokecolor{currentstroke}%
\pgfsetdash{}{0pt}%
\pgfpathmoveto{\pgfqpoint{0.812085in}{1.825153in}}%
\pgfpathlineto{\pgfqpoint{0.942330in}{2.102782in}}%
\pgfpathlineto{\pgfqpoint{0.921128in}{2.073538in}}%
\pgfpathlineto{\pgfqpoint{0.812085in}{1.825153in}}%
\pgfpathclose%
\pgfusepath{fill}%
\end{pgfscope}%
\begin{pgfscope}%
\pgfpathrectangle{\pgfqpoint{0.100000in}{0.100000in}}{\pgfqpoint{3.957178in}{3.957178in}}%
\pgfusepath{clip}%
\pgfsetbuttcap%
\pgfsetroundjoin%
\definecolor{currentfill}{rgb}{0.949151,0.948435,0.152178}%
\pgfsetfillcolor{currentfill}%
\pgfsetlinewidth{0.000000pt}%
\definecolor{currentstroke}{rgb}{0.000000,0.000000,0.000000}%
\pgfsetstrokecolor{currentstroke}%
\pgfsetdash{}{0pt}%
\pgfpathmoveto{\pgfqpoint{1.877921in}{3.091334in}}%
\pgfpathlineto{\pgfqpoint{1.762322in}{3.005647in}}%
\pgfpathlineto{\pgfqpoint{1.842130in}{2.944323in}}%
\pgfpathlineto{\pgfqpoint{1.877921in}{3.091334in}}%
\pgfpathclose%
\pgfusepath{fill}%
\end{pgfscope}%
\begin{pgfscope}%
\pgfpathrectangle{\pgfqpoint{0.100000in}{0.100000in}}{\pgfqpoint{3.957178in}{3.957178in}}%
\pgfusepath{clip}%
\pgfsetbuttcap%
\pgfsetroundjoin%
\definecolor{currentfill}{rgb}{0.966271,0.901249,0.148180}%
\pgfsetfillcolor{currentfill}%
\pgfsetlinewidth{0.000000pt}%
\definecolor{currentstroke}{rgb}{0.000000,0.000000,0.000000}%
\pgfsetstrokecolor{currentstroke}%
\pgfsetdash{}{0pt}%
\pgfpathmoveto{\pgfqpoint{2.770237in}{2.929808in}}%
\pgfpathlineto{\pgfqpoint{2.671623in}{3.002540in}}%
\pgfpathlineto{\pgfqpoint{2.680052in}{2.966420in}}%
\pgfpathlineto{\pgfqpoint{2.770237in}{2.929808in}}%
\pgfpathclose%
\pgfusepath{fill}%
\end{pgfscope}%
\begin{pgfscope}%
\pgfpathrectangle{\pgfqpoint{0.100000in}{0.100000in}}{\pgfqpoint{3.957178in}{3.957178in}}%
\pgfusepath{clip}%
\pgfsetbuttcap%
\pgfsetroundjoin%
\definecolor{currentfill}{rgb}{0.964021,0.907950,0.149370}%
\pgfsetfillcolor{currentfill}%
\pgfsetlinewidth{0.000000pt}%
\definecolor{currentstroke}{rgb}{0.000000,0.000000,0.000000}%
\pgfsetstrokecolor{currentstroke}%
\pgfsetdash{}{0pt}%
\pgfpathmoveto{\pgfqpoint{2.403569in}{2.912894in}}%
\pgfpathlineto{\pgfqpoint{2.285509in}{3.097456in}}%
\pgfpathlineto{\pgfqpoint{2.054934in}{2.904720in}}%
\pgfpathlineto{\pgfqpoint{2.403569in}{2.912894in}}%
\pgfpathclose%
\pgfusepath{fill}%
\end{pgfscope}%
\begin{pgfscope}%
\pgfpathrectangle{\pgfqpoint{0.100000in}{0.100000in}}{\pgfqpoint{3.957178in}{3.957178in}}%
\pgfusepath{clip}%
\pgfsetbuttcap%
\pgfsetroundjoin%
\definecolor{currentfill}{rgb}{0.904601,0.429797,0.356329}%
\pgfsetfillcolor{currentfill}%
\pgfsetlinewidth{0.000000pt}%
\definecolor{currentstroke}{rgb}{0.000000,0.000000,0.000000}%
\pgfsetstrokecolor{currentstroke}%
\pgfsetdash{}{0pt}%
\pgfpathmoveto{\pgfqpoint{1.212770in}{2.639946in}}%
\pgfpathlineto{\pgfqpoint{1.068273in}{2.394651in}}%
\pgfpathlineto{\pgfqpoint{1.003470in}{2.252260in}}%
\pgfpathlineto{\pgfqpoint{1.212770in}{2.639946in}}%
\pgfpathclose%
\pgfusepath{fill}%
\end{pgfscope}%
\begin{pgfscope}%
\pgfpathrectangle{\pgfqpoint{0.100000in}{0.100000in}}{\pgfqpoint{3.957178in}{3.957178in}}%
\pgfusepath{clip}%
\pgfsetbuttcap%
\pgfsetroundjoin%
\definecolor{currentfill}{rgb}{0.692840,0.165141,0.564522}%
\pgfsetfillcolor{currentfill}%
\pgfsetlinewidth{0.000000pt}%
\definecolor{currentstroke}{rgb}{0.000000,0.000000,0.000000}%
\pgfsetstrokecolor{currentstroke}%
\pgfsetdash{}{0pt}%
\pgfpathmoveto{\pgfqpoint{3.293789in}{2.110296in}}%
\pgfpathlineto{\pgfqpoint{3.464151in}{1.860528in}}%
\pgfpathlineto{\pgfqpoint{3.296283in}{2.159348in}}%
\pgfpathlineto{\pgfqpoint{3.293789in}{2.110296in}}%
\pgfpathclose%
\pgfusepath{fill}%
\end{pgfscope}%
\begin{pgfscope}%
\pgfpathrectangle{\pgfqpoint{0.100000in}{0.100000in}}{\pgfqpoint{3.957178in}{3.957178in}}%
\pgfusepath{clip}%
\pgfsetbuttcap%
\pgfsetroundjoin%
\definecolor{currentfill}{rgb}{0.973416,0.585761,0.251540}%
\pgfsetfillcolor{currentfill}%
\pgfsetlinewidth{0.000000pt}%
\definecolor{currentstroke}{rgb}{0.000000,0.000000,0.000000}%
\pgfsetstrokecolor{currentstroke}%
\pgfsetdash{}{0pt}%
\pgfpathmoveto{\pgfqpoint{2.917573in}{2.822478in}}%
\pgfpathlineto{\pgfqpoint{3.111123in}{2.556569in}}%
\pgfpathlineto{\pgfqpoint{3.154229in}{2.481387in}}%
\pgfpathlineto{\pgfqpoint{2.917573in}{2.822478in}}%
\pgfpathclose%
\pgfusepath{fill}%
\end{pgfscope}%
\begin{pgfscope}%
\pgfpathrectangle{\pgfqpoint{0.100000in}{0.100000in}}{\pgfqpoint{3.957178in}{3.957178in}}%
\pgfusepath{clip}%
\pgfsetbuttcap%
\pgfsetroundjoin%
\definecolor{currentfill}{rgb}{0.764193,0.240396,0.502126}%
\pgfsetfillcolor{currentfill}%
\pgfsetlinewidth{0.000000pt}%
\definecolor{currentstroke}{rgb}{0.000000,0.000000,0.000000}%
\pgfsetstrokecolor{currentstroke}%
\pgfsetdash{}{0pt}%
\pgfpathmoveto{\pgfqpoint{0.921128in}{2.073538in}}%
\pgfpathlineto{\pgfqpoint{0.942330in}{2.102782in}}%
\pgfpathlineto{\pgfqpoint{1.003470in}{2.252260in}}%
\pgfpathlineto{\pgfqpoint{0.921128in}{2.073538in}}%
\pgfpathclose%
\pgfusepath{fill}%
\end{pgfscope}%
\begin{pgfscope}%
\pgfpathrectangle{\pgfqpoint{0.100000in}{0.100000in}}{\pgfqpoint{3.957178in}{3.957178in}}%
\pgfusepath{clip}%
\pgfsetbuttcap%
\pgfsetroundjoin%
\definecolor{currentfill}{rgb}{0.946602,0.955190,0.150328}%
\pgfsetfillcolor{currentfill}%
\pgfsetlinewidth{0.000000pt}%
\definecolor{currentstroke}{rgb}{0.000000,0.000000,0.000000}%
\pgfsetstrokecolor{currentstroke}%
\pgfsetdash{}{0pt}%
\pgfpathmoveto{\pgfqpoint{2.403569in}{2.912894in}}%
\pgfpathlineto{\pgfqpoint{2.484697in}{3.037406in}}%
\pgfpathlineto{\pgfqpoint{2.422027in}{3.048101in}}%
\pgfpathlineto{\pgfqpoint{2.403569in}{2.912894in}}%
\pgfpathclose%
\pgfusepath{fill}%
\end{pgfscope}%
\begin{pgfscope}%
\pgfpathrectangle{\pgfqpoint{0.100000in}{0.100000in}}{\pgfqpoint{3.957178in}{3.957178in}}%
\pgfusepath{clip}%
\pgfsetbuttcap%
\pgfsetroundjoin%
\definecolor{currentfill}{rgb}{0.979644,0.854866,0.142453}%
\pgfsetfillcolor{currentfill}%
\pgfsetlinewidth{0.000000pt}%
\definecolor{currentstroke}{rgb}{0.000000,0.000000,0.000000}%
\pgfsetstrokecolor{currentstroke}%
\pgfsetdash{}{0pt}%
\pgfpathmoveto{\pgfqpoint{2.917573in}{2.822478in}}%
\pgfpathlineto{\pgfqpoint{2.770237in}{2.929808in}}%
\pgfpathlineto{\pgfqpoint{2.680052in}{2.966420in}}%
\pgfpathlineto{\pgfqpoint{2.917573in}{2.822478in}}%
\pgfpathclose%
\pgfusepath{fill}%
\end{pgfscope}%
\begin{pgfscope}%
\pgfpathrectangle{\pgfqpoint{0.100000in}{0.100000in}}{\pgfqpoint{3.957178in}{3.957178in}}%
\pgfusepath{clip}%
\pgfsetbuttcap%
\pgfsetroundjoin%
\definecolor{currentfill}{rgb}{0.974947,0.591165,0.248151}%
\pgfsetfillcolor{currentfill}%
\pgfsetlinewidth{0.000000pt}%
\definecolor{currentstroke}{rgb}{0.000000,0.000000,0.000000}%
\pgfsetstrokecolor{currentstroke}%
\pgfsetdash{}{0pt}%
\pgfpathmoveto{\pgfqpoint{1.212770in}{2.639946in}}%
\pgfpathlineto{\pgfqpoint{1.335079in}{2.804114in}}%
\pgfpathlineto{\pgfqpoint{1.068273in}{2.394651in}}%
\pgfpathlineto{\pgfqpoint{1.212770in}{2.639946in}}%
\pgfpathclose%
\pgfusepath{fill}%
\end{pgfscope}%
\begin{pgfscope}%
\pgfpathrectangle{\pgfqpoint{0.100000in}{0.100000in}}{\pgfqpoint{3.957178in}{3.957178in}}%
\pgfusepath{clip}%
\pgfsetbuttcap%
\pgfsetroundjoin%
\definecolor{currentfill}{rgb}{0.941896,0.968590,0.140956}%
\pgfsetfillcolor{currentfill}%
\pgfsetlinewidth{0.000000pt}%
\definecolor{currentstroke}{rgb}{0.000000,0.000000,0.000000}%
\pgfsetstrokecolor{currentstroke}%
\pgfsetdash{}{0pt}%
\pgfpathmoveto{\pgfqpoint{1.741847in}{2.997501in}}%
\pgfpathlineto{\pgfqpoint{1.691543in}{3.028826in}}%
\pgfpathlineto{\pgfqpoint{1.576036in}{2.988808in}}%
\pgfpathlineto{\pgfqpoint{1.741847in}{2.997501in}}%
\pgfpathclose%
\pgfusepath{fill}%
\end{pgfscope}%
\begin{pgfscope}%
\pgfpathrectangle{\pgfqpoint{0.100000in}{0.100000in}}{\pgfqpoint{3.957178in}{3.957178in}}%
\pgfusepath{clip}%
\pgfsetbuttcap%
\pgfsetroundjoin%
\definecolor{currentfill}{rgb}{0.944152,0.961916,0.146861}%
\pgfsetfillcolor{currentfill}%
\pgfsetlinewidth{0.000000pt}%
\definecolor{currentstroke}{rgb}{0.000000,0.000000,0.000000}%
\pgfsetstrokecolor{currentstroke}%
\pgfsetdash{}{0pt}%
\pgfpathmoveto{\pgfqpoint{2.572519in}{2.947479in}}%
\pgfpathlineto{\pgfqpoint{2.671623in}{3.002540in}}%
\pgfpathlineto{\pgfqpoint{2.484697in}{3.037406in}}%
\pgfpathlineto{\pgfqpoint{2.572519in}{2.947479in}}%
\pgfpathclose%
\pgfusepath{fill}%
\end{pgfscope}%
\begin{pgfscope}%
\pgfpathrectangle{\pgfqpoint{0.100000in}{0.100000in}}{\pgfqpoint{3.957178in}{3.957178in}}%
\pgfusepath{clip}%
\pgfsetbuttcap%
\pgfsetroundjoin%
\definecolor{currentfill}{rgb}{0.940015,0.975158,0.131326}%
\pgfsetfillcolor{currentfill}%
\pgfsetlinewidth{0.000000pt}%
\definecolor{currentstroke}{rgb}{0.000000,0.000000,0.000000}%
\pgfsetstrokecolor{currentstroke}%
\pgfsetdash{}{0pt}%
\pgfpathmoveto{\pgfqpoint{1.691543in}{3.028826in}}%
\pgfpathlineto{\pgfqpoint{1.741847in}{2.997501in}}%
\pgfpathlineto{\pgfqpoint{1.762322in}{3.005647in}}%
\pgfpathlineto{\pgfqpoint{1.691543in}{3.028826in}}%
\pgfpathclose%
\pgfusepath{fill}%
\end{pgfscope}%
\begin{pgfscope}%
\pgfpathrectangle{\pgfqpoint{0.100000in}{0.100000in}}{\pgfqpoint{3.957178in}{3.957178in}}%
\pgfusepath{clip}%
\pgfsetbuttcap%
\pgfsetroundjoin%
\definecolor{currentfill}{rgb}{0.984031,0.835315,0.142528}%
\pgfsetfillcolor{currentfill}%
\pgfsetlinewidth{0.000000pt}%
\definecolor{currentstroke}{rgb}{0.000000,0.000000,0.000000}%
\pgfsetstrokecolor{currentstroke}%
\pgfsetdash{}{0pt}%
\pgfpathmoveto{\pgfqpoint{2.054934in}{2.904720in}}%
\pgfpathlineto{\pgfqpoint{2.074507in}{2.788846in}}%
\pgfpathlineto{\pgfqpoint{2.403569in}{2.912894in}}%
\pgfpathlineto{\pgfqpoint{2.054934in}{2.904720in}}%
\pgfpathclose%
\pgfusepath{fill}%
\end{pgfscope}%
\begin{pgfscope}%
\pgfpathrectangle{\pgfqpoint{0.100000in}{0.100000in}}{\pgfqpoint{3.957178in}{3.957178in}}%
\pgfusepath{clip}%
\pgfsetbuttcap%
\pgfsetroundjoin%
\definecolor{currentfill}{rgb}{0.979644,0.854866,0.142453}%
\pgfsetfillcolor{currentfill}%
\pgfsetlinewidth{0.000000pt}%
\definecolor{currentstroke}{rgb}{0.000000,0.000000,0.000000}%
\pgfsetstrokecolor{currentstroke}%
\pgfsetdash{}{0pt}%
\pgfpathmoveto{\pgfqpoint{1.576036in}{2.988808in}}%
\pgfpathlineto{\pgfqpoint{1.409186in}{2.875647in}}%
\pgfpathlineto{\pgfqpoint{1.335079in}{2.804114in}}%
\pgfpathlineto{\pgfqpoint{1.576036in}{2.988808in}}%
\pgfpathclose%
\pgfusepath{fill}%
\end{pgfscope}%
\begin{pgfscope}%
\pgfpathrectangle{\pgfqpoint{0.100000in}{0.100000in}}{\pgfqpoint{3.957178in}{3.957178in}}%
\pgfusepath{clip}%
\pgfsetbuttcap%
\pgfsetroundjoin%
\definecolor{currentfill}{rgb}{0.992505,0.777967,0.152855}%
\pgfsetfillcolor{currentfill}%
\pgfsetlinewidth{0.000000pt}%
\definecolor{currentstroke}{rgb}{0.000000,0.000000,0.000000}%
\pgfsetstrokecolor{currentstroke}%
\pgfsetdash{}{0pt}%
\pgfpathmoveto{\pgfqpoint{1.937002in}{2.750602in}}%
\pgfpathlineto{\pgfqpoint{2.074507in}{2.788846in}}%
\pgfpathlineto{\pgfqpoint{2.054934in}{2.904720in}}%
\pgfpathlineto{\pgfqpoint{1.937002in}{2.750602in}}%
\pgfpathclose%
\pgfusepath{fill}%
\end{pgfscope}%
\begin{pgfscope}%
\pgfpathrectangle{\pgfqpoint{0.100000in}{0.100000in}}{\pgfqpoint{3.957178in}{3.957178in}}%
\pgfusepath{clip}%
\pgfsetbuttcap%
\pgfsetroundjoin%
\definecolor{currentfill}{rgb}{0.994355,0.746995,0.163821}%
\pgfsetfillcolor{currentfill}%
\pgfsetlinewidth{0.000000pt}%
\definecolor{currentstroke}{rgb}{0.000000,0.000000,0.000000}%
\pgfsetstrokecolor{currentstroke}%
\pgfsetdash{}{0pt}%
\pgfpathmoveto{\pgfqpoint{2.403569in}{2.912894in}}%
\pgfpathlineto{\pgfqpoint{2.074507in}{2.788846in}}%
\pgfpathlineto{\pgfqpoint{2.090788in}{2.637228in}}%
\pgfpathlineto{\pgfqpoint{2.403569in}{2.912894in}}%
\pgfpathclose%
\pgfusepath{fill}%
\end{pgfscope}%
\begin{pgfscope}%
\pgfpathrectangle{\pgfqpoint{0.100000in}{0.100000in}}{\pgfqpoint{3.957178in}{3.957178in}}%
\pgfusepath{clip}%
\pgfsetbuttcap%
\pgfsetroundjoin%
\definecolor{currentfill}{rgb}{0.993456,0.698810,0.186041}%
\pgfsetfillcolor{currentfill}%
\pgfsetlinewidth{0.000000pt}%
\definecolor{currentstroke}{rgb}{0.000000,0.000000,0.000000}%
\pgfsetstrokecolor{currentstroke}%
\pgfsetdash{}{0pt}%
\pgfpathmoveto{\pgfqpoint{2.090788in}{2.637228in}}%
\pgfpathlineto{\pgfqpoint{2.074507in}{2.788846in}}%
\pgfpathlineto{\pgfqpoint{1.937002in}{2.750602in}}%
\pgfpathlineto{\pgfqpoint{2.090788in}{2.637228in}}%
\pgfpathclose%
\pgfusepath{fill}%
\end{pgfscope}%
\begin{pgfscope}%
\pgfpathrectangle{\pgfqpoint{0.100000in}{0.100000in}}{\pgfqpoint{3.957178in}{3.957178in}}%
\pgfusepath{clip}%
\pgfsetbuttcap%
\pgfsetroundjoin%
\definecolor{currentfill}{rgb}{0.869203,0.374212,0.396738}%
\pgfsetfillcolor{currentfill}%
\pgfsetlinewidth{0.000000pt}%
\definecolor{currentstroke}{rgb}{0.000000,0.000000,0.000000}%
\pgfsetstrokecolor{currentstroke}%
\pgfsetdash{}{0pt}%
\pgfpathmoveto{\pgfqpoint{3.096663in}{2.540795in}}%
\pgfpathlineto{\pgfqpoint{3.313847in}{2.143572in}}%
\pgfpathlineto{\pgfqpoint{3.270626in}{2.259309in}}%
\pgfpathlineto{\pgfqpoint{3.096663in}{2.540795in}}%
\pgfpathclose%
\pgfusepath{fill}%
\end{pgfscope}%
\begin{pgfscope}%
\pgfpathrectangle{\pgfqpoint{0.100000in}{0.100000in}}{\pgfqpoint{3.957178in}{3.957178in}}%
\pgfusepath{clip}%
\pgfsetbuttcap%
\pgfsetroundjoin%
\definecolor{currentfill}{rgb}{0.974947,0.591165,0.248151}%
\pgfsetfillcolor{currentfill}%
\pgfsetlinewidth{0.000000pt}%
\definecolor{currentstroke}{rgb}{0.000000,0.000000,0.000000}%
\pgfsetstrokecolor{currentstroke}%
\pgfsetdash{}{0pt}%
\pgfpathmoveto{\pgfqpoint{2.090788in}{2.637228in}}%
\pgfpathlineto{\pgfqpoint{1.937002in}{2.750602in}}%
\pgfpathlineto{\pgfqpoint{2.055579in}{2.406772in}}%
\pgfpathlineto{\pgfqpoint{2.090788in}{2.637228in}}%
\pgfpathclose%
\pgfusepath{fill}%
\end{pgfscope}%
\begin{pgfscope}%
\pgfpathrectangle{\pgfqpoint{0.100000in}{0.100000in}}{\pgfqpoint{3.957178in}{3.957178in}}%
\pgfusepath{clip}%
\pgfsetbuttcap%
\pgfsetroundjoin%
\definecolor{currentfill}{rgb}{0.976265,0.868016,0.143351}%
\pgfsetfillcolor{currentfill}%
\pgfsetlinewidth{0.000000pt}%
\definecolor{currentstroke}{rgb}{0.000000,0.000000,0.000000}%
\pgfsetstrokecolor{currentstroke}%
\pgfsetdash{}{0pt}%
\pgfpathmoveto{\pgfqpoint{2.054934in}{2.904720in}}%
\pgfpathlineto{\pgfqpoint{1.842130in}{2.944323in}}%
\pgfpathlineto{\pgfqpoint{1.894028in}{2.839705in}}%
\pgfpathlineto{\pgfqpoint{2.054934in}{2.904720in}}%
\pgfpathclose%
\pgfusepath{fill}%
\end{pgfscope}%
\begin{pgfscope}%
\pgfpathrectangle{\pgfqpoint{0.100000in}{0.100000in}}{\pgfqpoint{3.957178in}{3.957178in}}%
\pgfusepath{clip}%
\pgfsetbuttcap%
\pgfsetroundjoin%
\definecolor{currentfill}{rgb}{0.944152,0.961916,0.146861}%
\pgfsetfillcolor{currentfill}%
\pgfsetlinewidth{0.000000pt}%
\definecolor{currentstroke}{rgb}{0.000000,0.000000,0.000000}%
\pgfsetstrokecolor{currentstroke}%
\pgfsetdash{}{0pt}%
\pgfpathmoveto{\pgfqpoint{1.842130in}{2.944323in}}%
\pgfpathlineto{\pgfqpoint{1.762322in}{3.005647in}}%
\pgfpathlineto{\pgfqpoint{1.741847in}{2.997501in}}%
\pgfpathlineto{\pgfqpoint{1.842130in}{2.944323in}}%
\pgfpathclose%
\pgfusepath{fill}%
\end{pgfscope}%
\begin{pgfscope}%
\pgfpathrectangle{\pgfqpoint{0.100000in}{0.100000in}}{\pgfqpoint{3.957178in}{3.957178in}}%
\pgfusepath{clip}%
\pgfsetbuttcap%
\pgfsetroundjoin%
\definecolor{currentfill}{rgb}{0.988648,0.809579,0.145357}%
\pgfsetfillcolor{currentfill}%
\pgfsetlinewidth{0.000000pt}%
\definecolor{currentstroke}{rgb}{0.000000,0.000000,0.000000}%
\pgfsetstrokecolor{currentstroke}%
\pgfsetdash{}{0pt}%
\pgfpathmoveto{\pgfqpoint{2.054934in}{2.904720in}}%
\pgfpathlineto{\pgfqpoint{1.894028in}{2.839705in}}%
\pgfpathlineto{\pgfqpoint{1.937002in}{2.750602in}}%
\pgfpathlineto{\pgfqpoint{2.054934in}{2.904720in}}%
\pgfpathclose%
\pgfusepath{fill}%
\end{pgfscope}%
\begin{pgfscope}%
\pgfpathrectangle{\pgfqpoint{0.100000in}{0.100000in}}{\pgfqpoint{3.957178in}{3.957178in}}%
\pgfusepath{clip}%
\pgfsetbuttcap%
\pgfsetroundjoin%
\definecolor{currentfill}{rgb}{0.930798,0.477867,0.322697}%
\pgfsetfillcolor{currentfill}%
\pgfsetlinewidth{0.000000pt}%
\definecolor{currentstroke}{rgb}{0.000000,0.000000,0.000000}%
\pgfsetstrokecolor{currentstroke}%
\pgfsetdash{}{0pt}%
\pgfpathmoveto{\pgfqpoint{3.270626in}{2.259309in}}%
\pgfpathlineto{\pgfqpoint{3.111123in}{2.556569in}}%
\pgfpathlineto{\pgfqpoint{3.096663in}{2.540795in}}%
\pgfpathlineto{\pgfqpoint{3.270626in}{2.259309in}}%
\pgfpathclose%
\pgfusepath{fill}%
\end{pgfscope}%
\begin{pgfscope}%
\pgfpathrectangle{\pgfqpoint{0.100000in}{0.100000in}}{\pgfqpoint{3.957178in}{3.957178in}}%
\pgfusepath{clip}%
\pgfsetbuttcap%
\pgfsetroundjoin%
\definecolor{currentfill}{rgb}{0.798216,0.280197,0.469538}%
\pgfsetfillcolor{currentfill}%
\pgfsetlinewidth{0.000000pt}%
\definecolor{currentstroke}{rgb}{0.000000,0.000000,0.000000}%
\pgfsetstrokecolor{currentstroke}%
\pgfsetdash{}{0pt}%
\pgfpathmoveto{\pgfqpoint{3.296283in}{2.159348in}}%
\pgfpathlineto{\pgfqpoint{3.313847in}{2.143572in}}%
\pgfpathlineto{\pgfqpoint{3.266506in}{2.205083in}}%
\pgfpathlineto{\pgfqpoint{3.296283in}{2.159348in}}%
\pgfpathclose%
\pgfusepath{fill}%
\end{pgfscope}%
\begin{pgfscope}%
\pgfpathrectangle{\pgfqpoint{0.100000in}{0.100000in}}{\pgfqpoint{3.957178in}{3.957178in}}%
\pgfusepath{clip}%
\pgfsetbuttcap%
\pgfsetroundjoin%
\definecolor{currentfill}{rgb}{0.640959,0.116492,0.602065}%
\pgfsetfillcolor{currentfill}%
\pgfsetlinewidth{0.000000pt}%
\definecolor{currentstroke}{rgb}{0.000000,0.000000,0.000000}%
\pgfsetstrokecolor{currentstroke}%
\pgfsetdash{}{0pt}%
\pgfpathmoveto{\pgfqpoint{0.812085in}{1.825153in}}%
\pgfpathlineto{\pgfqpoint{0.931340in}{1.874842in}}%
\pgfpathlineto{\pgfqpoint{0.942330in}{2.102782in}}%
\pgfpathlineto{\pgfqpoint{0.812085in}{1.825153in}}%
\pgfpathclose%
\pgfusepath{fill}%
\end{pgfscope}%
\begin{pgfscope}%
\pgfpathrectangle{\pgfqpoint{0.100000in}{0.100000in}}{\pgfqpoint{3.957178in}{3.957178in}}%
\pgfusepath{clip}%
\pgfsetbuttcap%
\pgfsetroundjoin%
\definecolor{currentfill}{rgb}{0.949151,0.948435,0.152178}%
\pgfsetfillcolor{currentfill}%
\pgfsetlinewidth{0.000000pt}%
\definecolor{currentstroke}{rgb}{0.000000,0.000000,0.000000}%
\pgfsetstrokecolor{currentstroke}%
\pgfsetdash{}{0pt}%
\pgfpathmoveto{\pgfqpoint{2.572519in}{2.947479in}}%
\pgfpathlineto{\pgfqpoint{2.484697in}{3.037406in}}%
\pgfpathlineto{\pgfqpoint{2.403569in}{2.912894in}}%
\pgfpathlineto{\pgfqpoint{2.572519in}{2.947479in}}%
\pgfpathclose%
\pgfusepath{fill}%
\end{pgfscope}%
\begin{pgfscope}%
\pgfpathrectangle{\pgfqpoint{0.100000in}{0.100000in}}{\pgfqpoint{3.957178in}{3.957178in}}%
\pgfusepath{clip}%
\pgfsetbuttcap%
\pgfsetroundjoin%
\definecolor{currentfill}{rgb}{0.944152,0.961916,0.146861}%
\pgfsetfillcolor{currentfill}%
\pgfsetlinewidth{0.000000pt}%
\definecolor{currentstroke}{rgb}{0.000000,0.000000,0.000000}%
\pgfsetstrokecolor{currentstroke}%
\pgfsetdash{}{0pt}%
\pgfpathmoveto{\pgfqpoint{2.680052in}{2.966420in}}%
\pgfpathlineto{\pgfqpoint{2.671623in}{3.002540in}}%
\pgfpathlineto{\pgfqpoint{2.572519in}{2.947479in}}%
\pgfpathlineto{\pgfqpoint{2.680052in}{2.966420in}}%
\pgfpathclose%
\pgfusepath{fill}%
\end{pgfscope}%
\begin{pgfscope}%
\pgfpathrectangle{\pgfqpoint{0.100000in}{0.100000in}}{\pgfqpoint{3.957178in}{3.957178in}}%
\pgfusepath{clip}%
\pgfsetbuttcap%
\pgfsetroundjoin%
\definecolor{currentfill}{rgb}{0.923287,0.463251,0.332801}%
\pgfsetfillcolor{currentfill}%
\pgfsetlinewidth{0.000000pt}%
\definecolor{currentstroke}{rgb}{0.000000,0.000000,0.000000}%
\pgfsetstrokecolor{currentstroke}%
\pgfsetdash{}{0pt}%
\pgfpathmoveto{\pgfqpoint{2.055579in}{2.406772in}}%
\pgfpathlineto{\pgfqpoint{2.222079in}{2.223854in}}%
\pgfpathlineto{\pgfqpoint{2.090788in}{2.637228in}}%
\pgfpathlineto{\pgfqpoint{2.055579in}{2.406772in}}%
\pgfpathclose%
\pgfusepath{fill}%
\end{pgfscope}%
\begin{pgfscope}%
\pgfpathrectangle{\pgfqpoint{0.100000in}{0.100000in}}{\pgfqpoint{3.957178in}{3.957178in}}%
\pgfusepath{clip}%
\pgfsetbuttcap%
\pgfsetroundjoin%
\definecolor{currentfill}{rgb}{0.957469,0.538250,0.282096}%
\pgfsetfillcolor{currentfill}%
\pgfsetlinewidth{0.000000pt}%
\definecolor{currentstroke}{rgb}{0.000000,0.000000,0.000000}%
\pgfsetstrokecolor{currentstroke}%
\pgfsetdash{}{0pt}%
\pgfpathmoveto{\pgfqpoint{1.974463in}{2.390365in}}%
\pgfpathlineto{\pgfqpoint{2.055579in}{2.406772in}}%
\pgfpathlineto{\pgfqpoint{1.937002in}{2.750602in}}%
\pgfpathlineto{\pgfqpoint{1.974463in}{2.390365in}}%
\pgfpathclose%
\pgfusepath{fill}%
\end{pgfscope}%
\begin{pgfscope}%
\pgfpathrectangle{\pgfqpoint{0.100000in}{0.100000in}}{\pgfqpoint{3.957178in}{3.957178in}}%
\pgfusepath{clip}%
\pgfsetbuttcap%
\pgfsetroundjoin%
\definecolor{currentfill}{rgb}{0.994561,0.734791,0.168938}%
\pgfsetfillcolor{currentfill}%
\pgfsetlinewidth{0.000000pt}%
\definecolor{currentstroke}{rgb}{0.000000,0.000000,0.000000}%
\pgfsetstrokecolor{currentstroke}%
\pgfsetdash{}{0pt}%
\pgfpathmoveto{\pgfqpoint{2.917573in}{2.822478in}}%
\pgfpathlineto{\pgfqpoint{2.873884in}{2.856740in}}%
\pgfpathlineto{\pgfqpoint{3.111123in}{2.556569in}}%
\pgfpathlineto{\pgfqpoint{2.917573in}{2.822478in}}%
\pgfpathclose%
\pgfusepath{fill}%
\end{pgfscope}%
\begin{pgfscope}%
\pgfpathrectangle{\pgfqpoint{0.100000in}{0.100000in}}{\pgfqpoint{3.957178in}{3.957178in}}%
\pgfusepath{clip}%
\pgfsetbuttcap%
\pgfsetroundjoin%
\definecolor{currentfill}{rgb}{0.318856,0.007576,0.637640}%
\pgfsetfillcolor{currentfill}%
\pgfsetlinewidth{0.000000pt}%
\definecolor{currentstroke}{rgb}{0.000000,0.000000,0.000000}%
\pgfsetstrokecolor{currentstroke}%
\pgfsetdash{}{0pt}%
\pgfpathmoveto{\pgfqpoint{0.783462in}{1.527029in}}%
\pgfpathlineto{\pgfqpoint{0.812085in}{1.825153in}}%
\pgfpathlineto{\pgfqpoint{0.747252in}{1.316186in}}%
\pgfpathlineto{\pgfqpoint{0.783462in}{1.527029in}}%
\pgfpathclose%
\pgfusepath{fill}%
\end{pgfscope}%
\begin{pgfscope}%
\pgfpathrectangle{\pgfqpoint{0.100000in}{0.100000in}}{\pgfqpoint{3.957178in}{3.957178in}}%
\pgfusepath{clip}%
\pgfsetbuttcap%
\pgfsetroundjoin%
\definecolor{currentfill}{rgb}{0.977856,0.602051,0.241387}%
\pgfsetfillcolor{currentfill}%
\pgfsetlinewidth{0.000000pt}%
\definecolor{currentstroke}{rgb}{0.000000,0.000000,0.000000}%
\pgfsetstrokecolor{currentstroke}%
\pgfsetdash{}{0pt}%
\pgfpathmoveto{\pgfqpoint{2.090788in}{2.637228in}}%
\pgfpathlineto{\pgfqpoint{2.222079in}{2.223854in}}%
\pgfpathlineto{\pgfqpoint{2.403569in}{2.912894in}}%
\pgfpathlineto{\pgfqpoint{2.090788in}{2.637228in}}%
\pgfpathclose%
\pgfusepath{fill}%
\end{pgfscope}%
\begin{pgfscope}%
\pgfpathrectangle{\pgfqpoint{0.100000in}{0.100000in}}{\pgfqpoint{3.957178in}{3.957178in}}%
\pgfusepath{clip}%
\pgfsetbuttcap%
\pgfsetroundjoin%
\definecolor{currentfill}{rgb}{0.878423,0.387932,0.386600}%
\pgfsetfillcolor{currentfill}%
\pgfsetlinewidth{0.000000pt}%
\definecolor{currentstroke}{rgb}{0.000000,0.000000,0.000000}%
\pgfsetstrokecolor{currentstroke}%
\pgfsetdash{}{0pt}%
\pgfpathmoveto{\pgfqpoint{2.083883in}{2.155265in}}%
\pgfpathlineto{\pgfqpoint{2.055579in}{2.406772in}}%
\pgfpathlineto{\pgfqpoint{1.974463in}{2.390365in}}%
\pgfpathlineto{\pgfqpoint{2.083883in}{2.155265in}}%
\pgfpathclose%
\pgfusepath{fill}%
\end{pgfscope}%
\begin{pgfscope}%
\pgfpathrectangle{\pgfqpoint{0.100000in}{0.100000in}}{\pgfqpoint{3.957178in}{3.957178in}}%
\pgfusepath{clip}%
\pgfsetbuttcap%
\pgfsetroundjoin%
\definecolor{currentfill}{rgb}{0.853319,0.351553,0.413734}%
\pgfsetfillcolor{currentfill}%
\pgfsetlinewidth{0.000000pt}%
\definecolor{currentstroke}{rgb}{0.000000,0.000000,0.000000}%
\pgfsetstrokecolor{currentstroke}%
\pgfsetdash{}{0pt}%
\pgfpathmoveto{\pgfqpoint{2.055579in}{2.406772in}}%
\pgfpathlineto{\pgfqpoint{2.083883in}{2.155265in}}%
\pgfpathlineto{\pgfqpoint{2.222079in}{2.223854in}}%
\pgfpathlineto{\pgfqpoint{2.055579in}{2.406772in}}%
\pgfpathclose%
\pgfusepath{fill}%
\end{pgfscope}%
\begin{pgfscope}%
\pgfpathrectangle{\pgfqpoint{0.100000in}{0.100000in}}{\pgfqpoint{3.957178in}{3.957178in}}%
\pgfusepath{clip}%
\pgfsetbuttcap%
\pgfsetroundjoin%
\definecolor{currentfill}{rgb}{0.961681,0.914672,0.150520}%
\pgfsetfillcolor{currentfill}%
\pgfsetlinewidth{0.000000pt}%
\definecolor{currentstroke}{rgb}{0.000000,0.000000,0.000000}%
\pgfsetstrokecolor{currentstroke}%
\pgfsetdash{}{0pt}%
\pgfpathmoveto{\pgfqpoint{1.576036in}{2.988808in}}%
\pgfpathlineto{\pgfqpoint{1.783624in}{2.759073in}}%
\pgfpathlineto{\pgfqpoint{1.741847in}{2.997501in}}%
\pgfpathlineto{\pgfqpoint{1.576036in}{2.988808in}}%
\pgfpathclose%
\pgfusepath{fill}%
\end{pgfscope}%
\begin{pgfscope}%
\pgfpathrectangle{\pgfqpoint{0.100000in}{0.100000in}}{\pgfqpoint{3.957178in}{3.957178in}}%
\pgfusepath{clip}%
\pgfsetbuttcap%
\pgfsetroundjoin%
\definecolor{currentfill}{rgb}{0.500678,0.014055,0.657088}%
\pgfsetfillcolor{currentfill}%
\pgfsetlinewidth{0.000000pt}%
\definecolor{currentstroke}{rgb}{0.000000,0.000000,0.000000}%
\pgfsetstrokecolor{currentstroke}%
\pgfsetdash{}{0pt}%
\pgfpathmoveto{\pgfqpoint{0.783462in}{1.527029in}}%
\pgfpathlineto{\pgfqpoint{0.931340in}{1.874842in}}%
\pgfpathlineto{\pgfqpoint{0.812085in}{1.825153in}}%
\pgfpathlineto{\pgfqpoint{0.783462in}{1.527029in}}%
\pgfpathclose%
\pgfusepath{fill}%
\end{pgfscope}%
\begin{pgfscope}%
\pgfpathrectangle{\pgfqpoint{0.100000in}{0.100000in}}{\pgfqpoint{3.957178in}{3.957178in}}%
\pgfusepath{clip}%
\pgfsetbuttcap%
\pgfsetroundjoin%
\definecolor{currentfill}{rgb}{0.964021,0.907950,0.149370}%
\pgfsetfillcolor{currentfill}%
\pgfsetlinewidth{0.000000pt}%
\definecolor{currentstroke}{rgb}{0.000000,0.000000,0.000000}%
\pgfsetstrokecolor{currentstroke}%
\pgfsetdash{}{0pt}%
\pgfpathmoveto{\pgfqpoint{1.842130in}{2.944323in}}%
\pgfpathlineto{\pgfqpoint{1.741847in}{2.997501in}}%
\pgfpathlineto{\pgfqpoint{1.783624in}{2.759073in}}%
\pgfpathlineto{\pgfqpoint{1.842130in}{2.944323in}}%
\pgfpathclose%
\pgfusepath{fill}%
\end{pgfscope}%
\begin{pgfscope}%
\pgfpathrectangle{\pgfqpoint{0.100000in}{0.100000in}}{\pgfqpoint{3.957178in}{3.957178in}}%
\pgfusepath{clip}%
\pgfsetbuttcap%
\pgfsetroundjoin%
\definecolor{currentfill}{rgb}{0.970533,0.887896,0.145919}%
\pgfsetfillcolor{currentfill}%
\pgfsetlinewidth{0.000000pt}%
\definecolor{currentstroke}{rgb}{0.000000,0.000000,0.000000}%
\pgfsetstrokecolor{currentstroke}%
\pgfsetdash{}{0pt}%
\pgfpathmoveto{\pgfqpoint{2.680052in}{2.966420in}}%
\pgfpathlineto{\pgfqpoint{2.873884in}{2.856740in}}%
\pgfpathlineto{\pgfqpoint{2.917573in}{2.822478in}}%
\pgfpathlineto{\pgfqpoint{2.680052in}{2.966420in}}%
\pgfpathclose%
\pgfusepath{fill}%
\end{pgfscope}%
\begin{pgfscope}%
\pgfpathrectangle{\pgfqpoint{0.100000in}{0.100000in}}{\pgfqpoint{3.957178in}{3.957178in}}%
\pgfusepath{clip}%
\pgfsetbuttcap%
\pgfsetroundjoin%
\definecolor{currentfill}{rgb}{0.812612,0.297928,0.455338}%
\pgfsetfillcolor{currentfill}%
\pgfsetlinewidth{0.000000pt}%
\definecolor{currentstroke}{rgb}{0.000000,0.000000,0.000000}%
\pgfsetstrokecolor{currentstroke}%
\pgfsetdash{}{0pt}%
\pgfpathmoveto{\pgfqpoint{3.266506in}{2.205083in}}%
\pgfpathlineto{\pgfqpoint{3.284979in}{2.149227in}}%
\pgfpathlineto{\pgfqpoint{3.296283in}{2.159348in}}%
\pgfpathlineto{\pgfqpoint{3.266506in}{2.205083in}}%
\pgfpathclose%
\pgfusepath{fill}%
\end{pgfscope}%
\begin{pgfscope}%
\pgfpathrectangle{\pgfqpoint{0.100000in}{0.100000in}}{\pgfqpoint{3.957178in}{3.957178in}}%
\pgfusepath{clip}%
\pgfsetbuttcap%
\pgfsetroundjoin%
\definecolor{currentfill}{rgb}{0.979644,0.854866,0.142453}%
\pgfsetfillcolor{currentfill}%
\pgfsetlinewidth{0.000000pt}%
\definecolor{currentstroke}{rgb}{0.000000,0.000000,0.000000}%
\pgfsetstrokecolor{currentstroke}%
\pgfsetdash{}{0pt}%
\pgfpathmoveto{\pgfqpoint{1.783624in}{2.759073in}}%
\pgfpathlineto{\pgfqpoint{1.894028in}{2.839705in}}%
\pgfpathlineto{\pgfqpoint{1.842130in}{2.944323in}}%
\pgfpathlineto{\pgfqpoint{1.783624in}{2.759073in}}%
\pgfpathclose%
\pgfusepath{fill}%
\end{pgfscope}%
\begin{pgfscope}%
\pgfpathrectangle{\pgfqpoint{0.100000in}{0.100000in}}{\pgfqpoint{3.957178in}{3.957178in}}%
\pgfusepath{clip}%
\pgfsetbuttcap%
\pgfsetroundjoin%
\definecolor{currentfill}{rgb}{0.959276,0.921407,0.151566}%
\pgfsetfillcolor{currentfill}%
\pgfsetlinewidth{0.000000pt}%
\definecolor{currentstroke}{rgb}{0.000000,0.000000,0.000000}%
\pgfsetstrokecolor{currentstroke}%
\pgfsetdash{}{0pt}%
\pgfpathmoveto{\pgfqpoint{2.403569in}{2.912894in}}%
\pgfpathlineto{\pgfqpoint{2.503502in}{2.880289in}}%
\pgfpathlineto{\pgfqpoint{2.572519in}{2.947479in}}%
\pgfpathlineto{\pgfqpoint{2.403569in}{2.912894in}}%
\pgfpathclose%
\pgfusepath{fill}%
\end{pgfscope}%
\begin{pgfscope}%
\pgfpathrectangle{\pgfqpoint{0.100000in}{0.100000in}}{\pgfqpoint{3.957178in}{3.957178in}}%
\pgfusepath{clip}%
\pgfsetbuttcap%
\pgfsetroundjoin%
\definecolor{currentfill}{rgb}{0.878423,0.387932,0.386600}%
\pgfsetfillcolor{currentfill}%
\pgfsetlinewidth{0.000000pt}%
\definecolor{currentstroke}{rgb}{0.000000,0.000000,0.000000}%
\pgfsetstrokecolor{currentstroke}%
\pgfsetdash{}{0pt}%
\pgfpathmoveto{\pgfqpoint{3.313847in}{2.143572in}}%
\pgfpathlineto{\pgfqpoint{3.096663in}{2.540795in}}%
\pgfpathlineto{\pgfqpoint{3.266506in}{2.205083in}}%
\pgfpathlineto{\pgfqpoint{3.313847in}{2.143572in}}%
\pgfpathclose%
\pgfusepath{fill}%
\end{pgfscope}%
\begin{pgfscope}%
\pgfpathrectangle{\pgfqpoint{0.100000in}{0.100000in}}{\pgfqpoint{3.957178in}{3.957178in}}%
\pgfusepath{clip}%
\pgfsetbuttcap%
\pgfsetroundjoin%
\definecolor{currentfill}{rgb}{0.798216,0.280197,0.469538}%
\pgfsetfillcolor{currentfill}%
\pgfsetlinewidth{0.000000pt}%
\definecolor{currentstroke}{rgb}{0.000000,0.000000,0.000000}%
\pgfsetstrokecolor{currentstroke}%
\pgfsetdash{}{0pt}%
\pgfpathmoveto{\pgfqpoint{3.293789in}{2.110296in}}%
\pgfpathlineto{\pgfqpoint{3.296283in}{2.159348in}}%
\pgfpathlineto{\pgfqpoint{3.284979in}{2.149227in}}%
\pgfpathlineto{\pgfqpoint{3.293789in}{2.110296in}}%
\pgfpathclose%
\pgfusepath{fill}%
\end{pgfscope}%
\begin{pgfscope}%
\pgfpathrectangle{\pgfqpoint{0.100000in}{0.100000in}}{\pgfqpoint{3.957178in}{3.957178in}}%
\pgfusepath{clip}%
\pgfsetbuttcap%
\pgfsetroundjoin%
\definecolor{currentfill}{rgb}{0.992541,0.687030,0.192170}%
\pgfsetfillcolor{currentfill}%
\pgfsetlinewidth{0.000000pt}%
\definecolor{currentstroke}{rgb}{0.000000,0.000000,0.000000}%
\pgfsetstrokecolor{currentstroke}%
\pgfsetdash{}{0pt}%
\pgfpathmoveto{\pgfqpoint{2.222079in}{2.223854in}}%
\pgfpathlineto{\pgfqpoint{2.503502in}{2.880289in}}%
\pgfpathlineto{\pgfqpoint{2.403569in}{2.912894in}}%
\pgfpathlineto{\pgfqpoint{2.222079in}{2.223854in}}%
\pgfpathclose%
\pgfusepath{fill}%
\end{pgfscope}%
\begin{pgfscope}%
\pgfpathrectangle{\pgfqpoint{0.100000in}{0.100000in}}{\pgfqpoint{3.957178in}{3.957178in}}%
\pgfusepath{clip}%
\pgfsetbuttcap%
\pgfsetroundjoin%
\definecolor{currentfill}{rgb}{0.990439,0.796859,0.147870}%
\pgfsetfillcolor{currentfill}%
\pgfsetlinewidth{0.000000pt}%
\definecolor{currentstroke}{rgb}{0.000000,0.000000,0.000000}%
\pgfsetstrokecolor{currentstroke}%
\pgfsetdash{}{0pt}%
\pgfpathmoveto{\pgfqpoint{1.783624in}{2.759073in}}%
\pgfpathlineto{\pgfqpoint{1.937002in}{2.750602in}}%
\pgfpathlineto{\pgfqpoint{1.894028in}{2.839705in}}%
\pgfpathlineto{\pgfqpoint{1.783624in}{2.759073in}}%
\pgfpathclose%
\pgfusepath{fill}%
\end{pgfscope}%
\begin{pgfscope}%
\pgfpathrectangle{\pgfqpoint{0.100000in}{0.100000in}}{\pgfqpoint{3.957178in}{3.957178in}}%
\pgfusepath{clip}%
\pgfsetbuttcap%
\pgfsetroundjoin%
\definecolor{currentfill}{rgb}{0.972530,0.881250,0.144923}%
\pgfsetfillcolor{currentfill}%
\pgfsetlinewidth{0.000000pt}%
\definecolor{currentstroke}{rgb}{0.000000,0.000000,0.000000}%
\pgfsetstrokecolor{currentstroke}%
\pgfsetdash{}{0pt}%
\pgfpathmoveto{\pgfqpoint{1.335079in}{2.804114in}}%
\pgfpathlineto{\pgfqpoint{1.350836in}{2.803045in}}%
\pgfpathlineto{\pgfqpoint{1.576036in}{2.988808in}}%
\pgfpathlineto{\pgfqpoint{1.335079in}{2.804114in}}%
\pgfpathclose%
\pgfusepath{fill}%
\end{pgfscope}%
\begin{pgfscope}%
\pgfpathrectangle{\pgfqpoint{0.100000in}{0.100000in}}{\pgfqpoint{3.957178in}{3.957178in}}%
\pgfusepath{clip}%
\pgfsetbuttcap%
\pgfsetroundjoin%
\definecolor{currentfill}{rgb}{0.994103,0.710698,0.180097}%
\pgfsetfillcolor{currentfill}%
\pgfsetlinewidth{0.000000pt}%
\definecolor{currentstroke}{rgb}{0.000000,0.000000,0.000000}%
\pgfsetstrokecolor{currentstroke}%
\pgfsetdash{}{0pt}%
\pgfpathmoveto{\pgfqpoint{3.033800in}{2.651652in}}%
\pgfpathlineto{\pgfqpoint{3.111123in}{2.556569in}}%
\pgfpathlineto{\pgfqpoint{2.873884in}{2.856740in}}%
\pgfpathlineto{\pgfqpoint{3.033800in}{2.651652in}}%
\pgfpathclose%
\pgfusepath{fill}%
\end{pgfscope}%
\begin{pgfscope}%
\pgfpathrectangle{\pgfqpoint{0.100000in}{0.100000in}}{\pgfqpoint{3.957178in}{3.957178in}}%
\pgfusepath{clip}%
\pgfsetbuttcap%
\pgfsetroundjoin%
\definecolor{currentfill}{rgb}{0.907365,0.434524,0.352970}%
\pgfsetfillcolor{currentfill}%
\pgfsetlinewidth{0.000000pt}%
\definecolor{currentstroke}{rgb}{0.000000,0.000000,0.000000}%
\pgfsetstrokecolor{currentstroke}%
\pgfsetdash{}{0pt}%
\pgfpathmoveto{\pgfqpoint{1.003470in}{2.252260in}}%
\pgfpathlineto{\pgfqpoint{0.942330in}{2.102782in}}%
\pgfpathlineto{\pgfqpoint{1.306609in}{2.699955in}}%
\pgfpathlineto{\pgfqpoint{1.003470in}{2.252260in}}%
\pgfpathclose%
\pgfusepath{fill}%
\end{pgfscope}%
\begin{pgfscope}%
\pgfpathrectangle{\pgfqpoint{0.100000in}{0.100000in}}{\pgfqpoint{3.957178in}{3.957178in}}%
\pgfusepath{clip}%
\pgfsetbuttcap%
\pgfsetroundjoin%
\definecolor{currentfill}{rgb}{0.979233,0.607532,0.238013}%
\pgfsetfillcolor{currentfill}%
\pgfsetlinewidth{0.000000pt}%
\definecolor{currentstroke}{rgb}{0.000000,0.000000,0.000000}%
\pgfsetstrokecolor{currentstroke}%
\pgfsetdash{}{0pt}%
\pgfpathmoveto{\pgfqpoint{1.937002in}{2.750602in}}%
\pgfpathlineto{\pgfqpoint{1.861239in}{2.565895in}}%
\pgfpathlineto{\pgfqpoint{1.974463in}{2.390365in}}%
\pgfpathlineto{\pgfqpoint{1.937002in}{2.750602in}}%
\pgfpathclose%
\pgfusepath{fill}%
\end{pgfscope}%
\begin{pgfscope}%
\pgfpathrectangle{\pgfqpoint{0.100000in}{0.100000in}}{\pgfqpoint{3.957178in}{3.957178in}}%
\pgfusepath{clip}%
\pgfsetbuttcap%
\pgfsetroundjoin%
\definecolor{currentfill}{rgb}{0.840155,0.333580,0.427455}%
\pgfsetfillcolor{currentfill}%
\pgfsetlinewidth{0.000000pt}%
\definecolor{currentstroke}{rgb}{0.000000,0.000000,0.000000}%
\pgfsetstrokecolor{currentstroke}%
\pgfsetdash{}{0pt}%
\pgfpathmoveto{\pgfqpoint{2.222079in}{2.223854in}}%
\pgfpathlineto{\pgfqpoint{2.083883in}{2.155265in}}%
\pgfpathlineto{\pgfqpoint{2.185536in}{2.224054in}}%
\pgfpathlineto{\pgfqpoint{2.222079in}{2.223854in}}%
\pgfpathclose%
\pgfusepath{fill}%
\end{pgfscope}%
\begin{pgfscope}%
\pgfpathrectangle{\pgfqpoint{0.100000in}{0.100000in}}{\pgfqpoint{3.957178in}{3.957178in}}%
\pgfusepath{clip}%
\pgfsetbuttcap%
\pgfsetroundjoin%
\definecolor{currentfill}{rgb}{0.840155,0.333580,0.427455}%
\pgfsetfillcolor{currentfill}%
\pgfsetlinewidth{0.000000pt}%
\definecolor{currentstroke}{rgb}{0.000000,0.000000,0.000000}%
\pgfsetstrokecolor{currentstroke}%
\pgfsetdash{}{0pt}%
\pgfpathmoveto{\pgfqpoint{3.240670in}{2.243694in}}%
\pgfpathlineto{\pgfqpoint{3.284979in}{2.149227in}}%
\pgfpathlineto{\pgfqpoint{3.266506in}{2.205083in}}%
\pgfpathlineto{\pgfqpoint{3.240670in}{2.243694in}}%
\pgfpathclose%
\pgfusepath{fill}%
\end{pgfscope}%
\begin{pgfscope}%
\pgfpathrectangle{\pgfqpoint{0.100000in}{0.100000in}}{\pgfqpoint{3.957178in}{3.957178in}}%
\pgfusepath{clip}%
\pgfsetbuttcap%
\pgfsetroundjoin%
\definecolor{currentfill}{rgb}{0.972530,0.881250,0.144923}%
\pgfsetfillcolor{currentfill}%
\pgfsetlinewidth{0.000000pt}%
\definecolor{currentstroke}{rgb}{0.000000,0.000000,0.000000}%
\pgfsetstrokecolor{currentstroke}%
\pgfsetdash{}{0pt}%
\pgfpathmoveto{\pgfqpoint{1.350836in}{2.803045in}}%
\pgfpathlineto{\pgfqpoint{1.783624in}{2.759073in}}%
\pgfpathlineto{\pgfqpoint{1.576036in}{2.988808in}}%
\pgfpathlineto{\pgfqpoint{1.350836in}{2.803045in}}%
\pgfpathclose%
\pgfusepath{fill}%
\end{pgfscope}%
\begin{pgfscope}%
\pgfpathrectangle{\pgfqpoint{0.100000in}{0.100000in}}{\pgfqpoint{3.957178in}{3.957178in}}%
\pgfusepath{clip}%
\pgfsetbuttcap%
\pgfsetroundjoin%
\definecolor{currentfill}{rgb}{0.826588,0.315714,0.441316}%
\pgfsetfillcolor{currentfill}%
\pgfsetlinewidth{0.000000pt}%
\definecolor{currentstroke}{rgb}{0.000000,0.000000,0.000000}%
\pgfsetstrokecolor{currentstroke}%
\pgfsetdash{}{0pt}%
\pgfpathmoveto{\pgfqpoint{3.284979in}{2.149227in}}%
\pgfpathlineto{\pgfqpoint{3.240670in}{2.243694in}}%
\pgfpathlineto{\pgfqpoint{3.293789in}{2.110296in}}%
\pgfpathlineto{\pgfqpoint{3.284979in}{2.149227in}}%
\pgfpathclose%
\pgfusepath{fill}%
\end{pgfscope}%
\begin{pgfscope}%
\pgfpathrectangle{\pgfqpoint{0.100000in}{0.100000in}}{\pgfqpoint{3.957178in}{3.957178in}}%
\pgfusepath{clip}%
\pgfsetbuttcap%
\pgfsetroundjoin%
\definecolor{currentfill}{rgb}{0.940015,0.975158,0.131326}%
\pgfsetfillcolor{currentfill}%
\pgfsetlinewidth{0.000000pt}%
\definecolor{currentstroke}{rgb}{0.000000,0.000000,0.000000}%
\pgfsetstrokecolor{currentstroke}%
\pgfsetdash{}{0pt}%
\pgfpathmoveto{\pgfqpoint{2.572519in}{2.947479in}}%
\pgfpathlineto{\pgfqpoint{2.674215in}{2.908477in}}%
\pgfpathlineto{\pgfqpoint{2.680052in}{2.966420in}}%
\pgfpathlineto{\pgfqpoint{2.572519in}{2.947479in}}%
\pgfpathclose%
\pgfusepath{fill}%
\end{pgfscope}%
\begin{pgfscope}%
\pgfpathrectangle{\pgfqpoint{0.100000in}{0.100000in}}{\pgfqpoint{3.957178in}{3.957178in}}%
\pgfusepath{clip}%
\pgfsetbuttcap%
\pgfsetroundjoin%
\definecolor{currentfill}{rgb}{0.984199,0.629718,0.224595}%
\pgfsetfillcolor{currentfill}%
\pgfsetlinewidth{0.000000pt}%
\definecolor{currentstroke}{rgb}{0.000000,0.000000,0.000000}%
\pgfsetstrokecolor{currentstroke}%
\pgfsetdash{}{0pt}%
\pgfpathmoveto{\pgfqpoint{3.096663in}{2.540795in}}%
\pgfpathlineto{\pgfqpoint{3.111123in}{2.556569in}}%
\pgfpathlineto{\pgfqpoint{3.033800in}{2.651652in}}%
\pgfpathlineto{\pgfqpoint{3.096663in}{2.540795in}}%
\pgfpathclose%
\pgfusepath{fill}%
\end{pgfscope}%
\begin{pgfscope}%
\pgfpathrectangle{\pgfqpoint{0.100000in}{0.100000in}}{\pgfqpoint{3.957178in}{3.957178in}}%
\pgfusepath{clip}%
\pgfsetbuttcap%
\pgfsetroundjoin%
\definecolor{currentfill}{rgb}{0.991897,0.784239,0.151042}%
\pgfsetfillcolor{currentfill}%
\pgfsetlinewidth{0.000000pt}%
\definecolor{currentstroke}{rgb}{0.000000,0.000000,0.000000}%
\pgfsetstrokecolor{currentstroke}%
\pgfsetdash{}{0pt}%
\pgfpathmoveto{\pgfqpoint{1.350836in}{2.803045in}}%
\pgfpathlineto{\pgfqpoint{1.335079in}{2.804114in}}%
\pgfpathlineto{\pgfqpoint{1.212770in}{2.639946in}}%
\pgfpathlineto{\pgfqpoint{1.350836in}{2.803045in}}%
\pgfpathclose%
\pgfusepath{fill}%
\end{pgfscope}%
\begin{pgfscope}%
\pgfpathrectangle{\pgfqpoint{0.100000in}{0.100000in}}{\pgfqpoint{3.957178in}{3.957178in}}%
\pgfusepath{clip}%
\pgfsetbuttcap%
\pgfsetroundjoin%
\definecolor{currentfill}{rgb}{0.949151,0.948435,0.152178}%
\pgfsetfillcolor{currentfill}%
\pgfsetlinewidth{0.000000pt}%
\definecolor{currentstroke}{rgb}{0.000000,0.000000,0.000000}%
\pgfsetstrokecolor{currentstroke}%
\pgfsetdash{}{0pt}%
\pgfpathmoveto{\pgfqpoint{2.572519in}{2.947479in}}%
\pgfpathlineto{\pgfqpoint{2.503502in}{2.880289in}}%
\pgfpathlineto{\pgfqpoint{2.674215in}{2.908477in}}%
\pgfpathlineto{\pgfqpoint{2.572519in}{2.947479in}}%
\pgfpathclose%
\pgfusepath{fill}%
\end{pgfscope}%
\begin{pgfscope}%
\pgfpathrectangle{\pgfqpoint{0.100000in}{0.100000in}}{\pgfqpoint{3.957178in}{3.957178in}}%
\pgfusepath{clip}%
\pgfsetbuttcap%
\pgfsetroundjoin%
\definecolor{currentfill}{rgb}{0.994561,0.734791,0.168938}%
\pgfsetfillcolor{currentfill}%
\pgfsetlinewidth{0.000000pt}%
\definecolor{currentstroke}{rgb}{0.000000,0.000000,0.000000}%
\pgfsetstrokecolor{currentstroke}%
\pgfsetdash{}{0pt}%
\pgfpathmoveto{\pgfqpoint{1.783624in}{2.759073in}}%
\pgfpathlineto{\pgfqpoint{1.861239in}{2.565895in}}%
\pgfpathlineto{\pgfqpoint{1.937002in}{2.750602in}}%
\pgfpathlineto{\pgfqpoint{1.783624in}{2.759073in}}%
\pgfpathclose%
\pgfusepath{fill}%
\end{pgfscope}%
\begin{pgfscope}%
\pgfpathrectangle{\pgfqpoint{0.100000in}{0.100000in}}{\pgfqpoint{3.957178in}{3.957178in}}%
\pgfusepath{clip}%
\pgfsetbuttcap%
\pgfsetroundjoin%
\definecolor{currentfill}{rgb}{0.697324,0.169573,0.560919}%
\pgfsetfillcolor{currentfill}%
\pgfsetlinewidth{0.000000pt}%
\definecolor{currentstroke}{rgb}{0.000000,0.000000,0.000000}%
\pgfsetstrokecolor{currentstroke}%
\pgfsetdash{}{0pt}%
\pgfpathmoveto{\pgfqpoint{3.464151in}{1.860528in}}%
\pgfpathlineto{\pgfqpoint{3.293789in}{2.110296in}}%
\pgfpathlineto{\pgfqpoint{3.157732in}{1.874966in}}%
\pgfpathlineto{\pgfqpoint{3.464151in}{1.860528in}}%
\pgfpathclose%
\pgfusepath{fill}%
\end{pgfscope}%
\begin{pgfscope}%
\pgfpathrectangle{\pgfqpoint{0.100000in}{0.100000in}}{\pgfqpoint{3.957178in}{3.957178in}}%
\pgfusepath{clip}%
\pgfsetbuttcap%
\pgfsetroundjoin%
\definecolor{currentfill}{rgb}{0.949151,0.948435,0.152178}%
\pgfsetfillcolor{currentfill}%
\pgfsetlinewidth{0.000000pt}%
\definecolor{currentstroke}{rgb}{0.000000,0.000000,0.000000}%
\pgfsetstrokecolor{currentstroke}%
\pgfsetdash{}{0pt}%
\pgfpathmoveto{\pgfqpoint{2.873884in}{2.856740in}}%
\pgfpathlineto{\pgfqpoint{2.680052in}{2.966420in}}%
\pgfpathlineto{\pgfqpoint{2.764569in}{2.899060in}}%
\pgfpathlineto{\pgfqpoint{2.873884in}{2.856740in}}%
\pgfpathclose%
\pgfusepath{fill}%
\end{pgfscope}%
\begin{pgfscope}%
\pgfpathrectangle{\pgfqpoint{0.100000in}{0.100000in}}{\pgfqpoint{3.957178in}{3.957178in}}%
\pgfusepath{clip}%
\pgfsetbuttcap%
\pgfsetroundjoin%
\definecolor{currentfill}{rgb}{0.910098,0.439268,0.349610}%
\pgfsetfillcolor{currentfill}%
\pgfsetlinewidth{0.000000pt}%
\definecolor{currentstroke}{rgb}{0.000000,0.000000,0.000000}%
\pgfsetstrokecolor{currentstroke}%
\pgfsetdash{}{0pt}%
\pgfpathmoveto{\pgfqpoint{3.266506in}{2.205083in}}%
\pgfpathlineto{\pgfqpoint{3.096663in}{2.540795in}}%
\pgfpathlineto{\pgfqpoint{3.240670in}{2.243694in}}%
\pgfpathlineto{\pgfqpoint{3.266506in}{2.205083in}}%
\pgfpathclose%
\pgfusepath{fill}%
\end{pgfscope}%
\begin{pgfscope}%
\pgfpathrectangle{\pgfqpoint{0.100000in}{0.100000in}}{\pgfqpoint{3.957178in}{3.957178in}}%
\pgfusepath{clip}%
\pgfsetbuttcap%
\pgfsetroundjoin%
\definecolor{currentfill}{rgb}{0.974947,0.591165,0.248151}%
\pgfsetfillcolor{currentfill}%
\pgfsetlinewidth{0.000000pt}%
\definecolor{currentstroke}{rgb}{0.000000,0.000000,0.000000}%
\pgfsetstrokecolor{currentstroke}%
\pgfsetdash{}{0pt}%
\pgfpathmoveto{\pgfqpoint{1.306609in}{2.699955in}}%
\pgfpathlineto{\pgfqpoint{1.212770in}{2.639946in}}%
\pgfpathlineto{\pgfqpoint{1.003470in}{2.252260in}}%
\pgfpathlineto{\pgfqpoint{1.306609in}{2.699955in}}%
\pgfpathclose%
\pgfusepath{fill}%
\end{pgfscope}%
\begin{pgfscope}%
\pgfpathrectangle{\pgfqpoint{0.100000in}{0.100000in}}{\pgfqpoint{3.957178in}{3.957178in}}%
\pgfusepath{clip}%
\pgfsetbuttcap%
\pgfsetroundjoin%
\definecolor{currentfill}{rgb}{0.940015,0.975158,0.131326}%
\pgfsetfillcolor{currentfill}%
\pgfsetlinewidth{0.000000pt}%
\definecolor{currentstroke}{rgb}{0.000000,0.000000,0.000000}%
\pgfsetstrokecolor{currentstroke}%
\pgfsetdash{}{0pt}%
\pgfpathmoveto{\pgfqpoint{2.764569in}{2.899060in}}%
\pgfpathlineto{\pgfqpoint{2.680052in}{2.966420in}}%
\pgfpathlineto{\pgfqpoint{2.674215in}{2.908477in}}%
\pgfpathlineto{\pgfqpoint{2.764569in}{2.899060in}}%
\pgfpathclose%
\pgfusepath{fill}%
\end{pgfscope}%
\begin{pgfscope}%
\pgfpathrectangle{\pgfqpoint{0.100000in}{0.100000in}}{\pgfqpoint{3.957178in}{3.957178in}}%
\pgfusepath{clip}%
\pgfsetbuttcap%
\pgfsetroundjoin%
\definecolor{currentfill}{rgb}{0.928329,0.472975,0.326067}%
\pgfsetfillcolor{currentfill}%
\pgfsetlinewidth{0.000000pt}%
\definecolor{currentstroke}{rgb}{0.000000,0.000000,0.000000}%
\pgfsetstrokecolor{currentstroke}%
\pgfsetdash{}{0pt}%
\pgfpathmoveto{\pgfqpoint{1.835085in}{2.543290in}}%
\pgfpathlineto{\pgfqpoint{2.083883in}{2.155265in}}%
\pgfpathlineto{\pgfqpoint{1.974463in}{2.390365in}}%
\pgfpathlineto{\pgfqpoint{1.835085in}{2.543290in}}%
\pgfpathclose%
\pgfusepath{fill}%
\end{pgfscope}%
\begin{pgfscope}%
\pgfpathrectangle{\pgfqpoint{0.100000in}{0.100000in}}{\pgfqpoint{3.957178in}{3.957178in}}%
\pgfusepath{clip}%
\pgfsetbuttcap%
\pgfsetroundjoin%
\definecolor{currentfill}{rgb}{0.881443,0.392529,0.383229}%
\pgfsetfillcolor{currentfill}%
\pgfsetlinewidth{0.000000pt}%
\definecolor{currentstroke}{rgb}{0.000000,0.000000,0.000000}%
\pgfsetstrokecolor{currentstroke}%
\pgfsetdash{}{0pt}%
\pgfpathmoveto{\pgfqpoint{3.240670in}{2.243694in}}%
\pgfpathlineto{\pgfqpoint{3.145728in}{2.407588in}}%
\pgfpathlineto{\pgfqpoint{3.293789in}{2.110296in}}%
\pgfpathlineto{\pgfqpoint{3.240670in}{2.243694in}}%
\pgfpathclose%
\pgfusepath{fill}%
\end{pgfscope}%
\begin{pgfscope}%
\pgfpathrectangle{\pgfqpoint{0.100000in}{0.100000in}}{\pgfqpoint{3.957178in}{3.957178in}}%
\pgfusepath{clip}%
\pgfsetbuttcap%
\pgfsetroundjoin%
\definecolor{currentfill}{rgb}{0.986345,0.640969,0.217948}%
\pgfsetfillcolor{currentfill}%
\pgfsetlinewidth{0.000000pt}%
\definecolor{currentstroke}{rgb}{0.000000,0.000000,0.000000}%
\pgfsetstrokecolor{currentstroke}%
\pgfsetdash{}{0pt}%
\pgfpathmoveto{\pgfqpoint{2.459372in}{2.593099in}}%
\pgfpathlineto{\pgfqpoint{2.503502in}{2.880289in}}%
\pgfpathlineto{\pgfqpoint{2.222079in}{2.223854in}}%
\pgfpathlineto{\pgfqpoint{2.459372in}{2.593099in}}%
\pgfpathclose%
\pgfusepath{fill}%
\end{pgfscope}%
\begin{pgfscope}%
\pgfpathrectangle{\pgfqpoint{0.100000in}{0.100000in}}{\pgfqpoint{3.957178in}{3.957178in}}%
\pgfusepath{clip}%
\pgfsetbuttcap%
\pgfsetroundjoin%
\definecolor{currentfill}{rgb}{0.994141,0.753137,0.161404}%
\pgfsetfillcolor{currentfill}%
\pgfsetlinewidth{0.000000pt}%
\definecolor{currentstroke}{rgb}{0.000000,0.000000,0.000000}%
\pgfsetstrokecolor{currentstroke}%
\pgfsetdash{}{0pt}%
\pgfpathmoveto{\pgfqpoint{1.782145in}{2.750558in}}%
\pgfpathlineto{\pgfqpoint{1.861239in}{2.565895in}}%
\pgfpathlineto{\pgfqpoint{1.783624in}{2.759073in}}%
\pgfpathlineto{\pgfqpoint{1.782145in}{2.750558in}}%
\pgfpathclose%
\pgfusepath{fill}%
\end{pgfscope}%
\begin{pgfscope}%
\pgfpathrectangle{\pgfqpoint{0.100000in}{0.100000in}}{\pgfqpoint{3.957178in}{3.957178in}}%
\pgfusepath{clip}%
\pgfsetbuttcap%
\pgfsetroundjoin%
\definecolor{currentfill}{rgb}{0.985314,0.828846,0.142945}%
\pgfsetfillcolor{currentfill}%
\pgfsetlinewidth{0.000000pt}%
\definecolor{currentstroke}{rgb}{0.000000,0.000000,0.000000}%
\pgfsetstrokecolor{currentstroke}%
\pgfsetdash{}{0pt}%
\pgfpathmoveto{\pgfqpoint{1.783624in}{2.759073in}}%
\pgfpathlineto{\pgfqpoint{1.350836in}{2.803045in}}%
\pgfpathlineto{\pgfqpoint{1.782145in}{2.750558in}}%
\pgfpathlineto{\pgfqpoint{1.783624in}{2.759073in}}%
\pgfpathclose%
\pgfusepath{fill}%
\end{pgfscope}%
\begin{pgfscope}%
\pgfpathrectangle{\pgfqpoint{0.100000in}{0.100000in}}{\pgfqpoint{3.957178in}{3.957178in}}%
\pgfusepath{clip}%
\pgfsetbuttcap%
\pgfsetroundjoin%
\definecolor{currentfill}{rgb}{0.423689,0.000646,0.658956}%
\pgfsetfillcolor{currentfill}%
\pgfsetlinewidth{0.000000pt}%
\definecolor{currentstroke}{rgb}{0.000000,0.000000,0.000000}%
\pgfsetstrokecolor{currentstroke}%
\pgfsetdash{}{0pt}%
\pgfpathmoveto{\pgfqpoint{0.971959in}{1.915950in}}%
\pgfpathlineto{\pgfqpoint{0.783462in}{1.527029in}}%
\pgfpathlineto{\pgfqpoint{0.747252in}{1.316186in}}%
\pgfpathlineto{\pgfqpoint{0.971959in}{1.915950in}}%
\pgfpathclose%
\pgfusepath{fill}%
\end{pgfscope}%
\begin{pgfscope}%
\pgfpathrectangle{\pgfqpoint{0.100000in}{0.100000in}}{\pgfqpoint{3.957178in}{3.957178in}}%
\pgfusepath{clip}%
\pgfsetbuttcap%
\pgfsetroundjoin%
\definecolor{currentfill}{rgb}{0.981190,0.848329,0.142279}%
\pgfsetfillcolor{currentfill}%
\pgfsetlinewidth{0.000000pt}%
\definecolor{currentstroke}{rgb}{0.000000,0.000000,0.000000}%
\pgfsetstrokecolor{currentstroke}%
\pgfsetdash{}{0pt}%
\pgfpathmoveto{\pgfqpoint{3.033800in}{2.651652in}}%
\pgfpathlineto{\pgfqpoint{2.873884in}{2.856740in}}%
\pgfpathlineto{\pgfqpoint{2.849017in}{2.840181in}}%
\pgfpathlineto{\pgfqpoint{3.033800in}{2.651652in}}%
\pgfpathclose%
\pgfusepath{fill}%
\end{pgfscope}%
\begin{pgfscope}%
\pgfpathrectangle{\pgfqpoint{0.100000in}{0.100000in}}{\pgfqpoint{3.957178in}{3.957178in}}%
\pgfusepath{clip}%
\pgfsetbuttcap%
\pgfsetroundjoin%
\definecolor{currentfill}{rgb}{0.954287,0.934908,0.152921}%
\pgfsetfillcolor{currentfill}%
\pgfsetlinewidth{0.000000pt}%
\definecolor{currentstroke}{rgb}{0.000000,0.000000,0.000000}%
\pgfsetstrokecolor{currentstroke}%
\pgfsetdash{}{0pt}%
\pgfpathmoveto{\pgfqpoint{2.849017in}{2.840181in}}%
\pgfpathlineto{\pgfqpoint{2.873884in}{2.856740in}}%
\pgfpathlineto{\pgfqpoint{2.764569in}{2.899060in}}%
\pgfpathlineto{\pgfqpoint{2.849017in}{2.840181in}}%
\pgfpathclose%
\pgfusepath{fill}%
\end{pgfscope}%
\begin{pgfscope}%
\pgfpathrectangle{\pgfqpoint{0.100000in}{0.100000in}}{\pgfqpoint{3.957178in}{3.957178in}}%
\pgfusepath{clip}%
\pgfsetbuttcap%
\pgfsetroundjoin%
\definecolor{currentfill}{rgb}{0.928329,0.472975,0.326067}%
\pgfsetfillcolor{currentfill}%
\pgfsetlinewidth{0.000000pt}%
\definecolor{currentstroke}{rgb}{0.000000,0.000000,0.000000}%
\pgfsetstrokecolor{currentstroke}%
\pgfsetdash{}{0pt}%
\pgfpathmoveto{\pgfqpoint{2.222079in}{2.223854in}}%
\pgfpathlineto{\pgfqpoint{2.185536in}{2.224054in}}%
\pgfpathlineto{\pgfqpoint{2.459372in}{2.593099in}}%
\pgfpathlineto{\pgfqpoint{2.222079in}{2.223854in}}%
\pgfpathclose%
\pgfusepath{fill}%
\end{pgfscope}%
\begin{pgfscope}%
\pgfpathrectangle{\pgfqpoint{0.100000in}{0.100000in}}{\pgfqpoint{3.957178in}{3.957178in}}%
\pgfusepath{clip}%
\pgfsetbuttcap%
\pgfsetroundjoin%
\definecolor{currentfill}{rgb}{0.947051,0.512699,0.299049}%
\pgfsetfillcolor{currentfill}%
\pgfsetlinewidth{0.000000pt}%
\definecolor{currentstroke}{rgb}{0.000000,0.000000,0.000000}%
\pgfsetstrokecolor{currentstroke}%
\pgfsetdash{}{0pt}%
\pgfpathmoveto{\pgfqpoint{3.145728in}{2.407588in}}%
\pgfpathlineto{\pgfqpoint{3.240670in}{2.243694in}}%
\pgfpathlineto{\pgfqpoint{3.096663in}{2.540795in}}%
\pgfpathlineto{\pgfqpoint{3.145728in}{2.407588in}}%
\pgfpathclose%
\pgfusepath{fill}%
\end{pgfscope}%
\begin{pgfscope}%
\pgfpathrectangle{\pgfqpoint{0.100000in}{0.100000in}}{\pgfqpoint{3.957178in}{3.957178in}}%
\pgfusepath{clip}%
\pgfsetbuttcap%
\pgfsetroundjoin%
\definecolor{currentfill}{rgb}{0.595011,0.077190,0.627917}%
\pgfsetfillcolor{currentfill}%
\pgfsetlinewidth{0.000000pt}%
\definecolor{currentstroke}{rgb}{0.000000,0.000000,0.000000}%
\pgfsetstrokecolor{currentstroke}%
\pgfsetdash{}{0pt}%
\pgfpathmoveto{\pgfqpoint{0.931340in}{1.874842in}}%
\pgfpathlineto{\pgfqpoint{0.783462in}{1.527029in}}%
\pgfpathlineto{\pgfqpoint{0.971959in}{1.915950in}}%
\pgfpathlineto{\pgfqpoint{0.931340in}{1.874842in}}%
\pgfpathclose%
\pgfusepath{fill}%
\end{pgfscope}%
\begin{pgfscope}%
\pgfpathrectangle{\pgfqpoint{0.100000in}{0.100000in}}{\pgfqpoint{3.957178in}{3.957178in}}%
\pgfusepath{clip}%
\pgfsetbuttcap%
\pgfsetroundjoin%
\definecolor{currentfill}{rgb}{0.994141,0.753137,0.161404}%
\pgfsetfillcolor{currentfill}%
\pgfsetlinewidth{0.000000pt}%
\definecolor{currentstroke}{rgb}{0.000000,0.000000,0.000000}%
\pgfsetstrokecolor{currentstroke}%
\pgfsetdash{}{0pt}%
\pgfpathmoveto{\pgfqpoint{1.782145in}{2.750558in}}%
\pgfpathlineto{\pgfqpoint{1.759958in}{2.718143in}}%
\pgfpathlineto{\pgfqpoint{1.861239in}{2.565895in}}%
\pgfpathlineto{\pgfqpoint{1.782145in}{2.750558in}}%
\pgfpathclose%
\pgfusepath{fill}%
\end{pgfscope}%
\begin{pgfscope}%
\pgfpathrectangle{\pgfqpoint{0.100000in}{0.100000in}}{\pgfqpoint{3.957178in}{3.957178in}}%
\pgfusepath{clip}%
\pgfsetbuttcap%
\pgfsetroundjoin%
\definecolor{currentfill}{rgb}{0.985314,0.828846,0.142945}%
\pgfsetfillcolor{currentfill}%
\pgfsetlinewidth{0.000000pt}%
\definecolor{currentstroke}{rgb}{0.000000,0.000000,0.000000}%
\pgfsetstrokecolor{currentstroke}%
\pgfsetdash{}{0pt}%
\pgfpathmoveto{\pgfqpoint{1.350836in}{2.803045in}}%
\pgfpathlineto{\pgfqpoint{1.759958in}{2.718143in}}%
\pgfpathlineto{\pgfqpoint{1.782145in}{2.750558in}}%
\pgfpathlineto{\pgfqpoint{1.350836in}{2.803045in}}%
\pgfpathclose%
\pgfusepath{fill}%
\end{pgfscope}%
\begin{pgfscope}%
\pgfpathrectangle{\pgfqpoint{0.100000in}{0.100000in}}{\pgfqpoint{3.957178in}{3.957178in}}%
\pgfusepath{clip}%
\pgfsetbuttcap%
\pgfsetroundjoin%
\definecolor{currentfill}{rgb}{0.976265,0.868016,0.143351}%
\pgfsetfillcolor{currentfill}%
\pgfsetlinewidth{0.000000pt}%
\definecolor{currentstroke}{rgb}{0.000000,0.000000,0.000000}%
\pgfsetstrokecolor{currentstroke}%
\pgfsetdash{}{0pt}%
\pgfpathmoveto{\pgfqpoint{2.459372in}{2.593099in}}%
\pgfpathlineto{\pgfqpoint{2.674215in}{2.908477in}}%
\pgfpathlineto{\pgfqpoint{2.503502in}{2.880289in}}%
\pgfpathlineto{\pgfqpoint{2.459372in}{2.593099in}}%
\pgfpathclose%
\pgfusepath{fill}%
\end{pgfscope}%
\begin{pgfscope}%
\pgfpathrectangle{\pgfqpoint{0.100000in}{0.100000in}}{\pgfqpoint{3.957178in}{3.957178in}}%
\pgfusepath{clip}%
\pgfsetbuttcap%
\pgfsetroundjoin%
\definecolor{currentfill}{rgb}{0.991209,0.790537,0.149377}%
\pgfsetfillcolor{currentfill}%
\pgfsetlinewidth{0.000000pt}%
\definecolor{currentstroke}{rgb}{0.000000,0.000000,0.000000}%
\pgfsetstrokecolor{currentstroke}%
\pgfsetdash{}{0pt}%
\pgfpathmoveto{\pgfqpoint{1.350836in}{2.803045in}}%
\pgfpathlineto{\pgfqpoint{1.212770in}{2.639946in}}%
\pgfpathlineto{\pgfqpoint{1.306609in}{2.699955in}}%
\pgfpathlineto{\pgfqpoint{1.350836in}{2.803045in}}%
\pgfpathclose%
\pgfusepath{fill}%
\end{pgfscope}%
\begin{pgfscope}%
\pgfpathrectangle{\pgfqpoint{0.100000in}{0.100000in}}{\pgfqpoint{3.957178in}{3.957178in}}%
\pgfusepath{clip}%
\pgfsetbuttcap%
\pgfsetroundjoin%
\definecolor{currentfill}{rgb}{0.977856,0.602051,0.241387}%
\pgfsetfillcolor{currentfill}%
\pgfsetlinewidth{0.000000pt}%
\definecolor{currentstroke}{rgb}{0.000000,0.000000,0.000000}%
\pgfsetstrokecolor{currentstroke}%
\pgfsetdash{}{0pt}%
\pgfpathmoveto{\pgfqpoint{1.974463in}{2.390365in}}%
\pgfpathlineto{\pgfqpoint{1.861239in}{2.565895in}}%
\pgfpathlineto{\pgfqpoint{1.835085in}{2.543290in}}%
\pgfpathlineto{\pgfqpoint{1.974463in}{2.390365in}}%
\pgfpathclose%
\pgfusepath{fill}%
\end{pgfscope}%
\begin{pgfscope}%
\pgfpathrectangle{\pgfqpoint{0.100000in}{0.100000in}}{\pgfqpoint{3.957178in}{3.957178in}}%
\pgfusepath{clip}%
\pgfsetbuttcap%
\pgfsetroundjoin%
\definecolor{currentfill}{rgb}{0.940015,0.975158,0.131326}%
\pgfsetfillcolor{currentfill}%
\pgfsetlinewidth{0.000000pt}%
\definecolor{currentstroke}{rgb}{0.000000,0.000000,0.000000}%
\pgfsetstrokecolor{currentstroke}%
\pgfsetdash{}{0pt}%
\pgfpathmoveto{\pgfqpoint{2.674215in}{2.908477in}}%
\pgfpathlineto{\pgfqpoint{2.716885in}{2.870430in}}%
\pgfpathlineto{\pgfqpoint{2.764569in}{2.899060in}}%
\pgfpathlineto{\pgfqpoint{2.674215in}{2.908477in}}%
\pgfpathclose%
\pgfusepath{fill}%
\end{pgfscope}%
\begin{pgfscope}%
\pgfpathrectangle{\pgfqpoint{0.100000in}{0.100000in}}{\pgfqpoint{3.957178in}{3.957178in}}%
\pgfusepath{clip}%
\pgfsetbuttcap%
\pgfsetroundjoin%
\definecolor{currentfill}{rgb}{0.941896,0.968590,0.140956}%
\pgfsetfillcolor{currentfill}%
\pgfsetlinewidth{0.000000pt}%
\definecolor{currentstroke}{rgb}{0.000000,0.000000,0.000000}%
\pgfsetstrokecolor{currentstroke}%
\pgfsetdash{}{0pt}%
\pgfpathmoveto{\pgfqpoint{2.764569in}{2.899060in}}%
\pgfpathlineto{\pgfqpoint{2.716885in}{2.870430in}}%
\pgfpathlineto{\pgfqpoint{2.849017in}{2.840181in}}%
\pgfpathlineto{\pgfqpoint{2.764569in}{2.899060in}}%
\pgfpathclose%
\pgfusepath{fill}%
\end{pgfscope}%
\begin{pgfscope}%
\pgfpathrectangle{\pgfqpoint{0.100000in}{0.100000in}}{\pgfqpoint{3.957178in}{3.957178in}}%
\pgfusepath{clip}%
\pgfsetbuttcap%
\pgfsetroundjoin%
\definecolor{currentfill}{rgb}{0.843484,0.338062,0.424013}%
\pgfsetfillcolor{currentfill}%
\pgfsetlinewidth{0.000000pt}%
\definecolor{currentstroke}{rgb}{0.000000,0.000000,0.000000}%
\pgfsetstrokecolor{currentstroke}%
\pgfsetdash{}{0pt}%
\pgfpathmoveto{\pgfqpoint{1.246560in}{2.412228in}}%
\pgfpathlineto{\pgfqpoint{0.942330in}{2.102782in}}%
\pgfpathlineto{\pgfqpoint{0.931340in}{1.874842in}}%
\pgfpathlineto{\pgfqpoint{1.246560in}{2.412228in}}%
\pgfpathclose%
\pgfusepath{fill}%
\end{pgfscope}%
\begin{pgfscope}%
\pgfpathrectangle{\pgfqpoint{0.100000in}{0.100000in}}{\pgfqpoint{3.957178in}{3.957178in}}%
\pgfusepath{clip}%
\pgfsetbuttcap%
\pgfsetroundjoin%
\definecolor{currentfill}{rgb}{0.989935,0.663787,0.204859}%
\pgfsetfillcolor{currentfill}%
\pgfsetlinewidth{0.000000pt}%
\definecolor{currentstroke}{rgb}{0.000000,0.000000,0.000000}%
\pgfsetstrokecolor{currentstroke}%
\pgfsetdash{}{0pt}%
\pgfpathmoveto{\pgfqpoint{3.033800in}{2.651652in}}%
\pgfpathlineto{\pgfqpoint{3.083252in}{2.463869in}}%
\pgfpathlineto{\pgfqpoint{3.096663in}{2.540795in}}%
\pgfpathlineto{\pgfqpoint{3.033800in}{2.651652in}}%
\pgfpathclose%
\pgfusepath{fill}%
\end{pgfscope}%
\begin{pgfscope}%
\pgfpathrectangle{\pgfqpoint{0.100000in}{0.100000in}}{\pgfqpoint{3.957178in}{3.957178in}}%
\pgfusepath{clip}%
\pgfsetbuttcap%
\pgfsetroundjoin%
\definecolor{currentfill}{rgb}{0.968443,0.894564,0.147014}%
\pgfsetfillcolor{currentfill}%
\pgfsetlinewidth{0.000000pt}%
\definecolor{currentstroke}{rgb}{0.000000,0.000000,0.000000}%
\pgfsetstrokecolor{currentstroke}%
\pgfsetdash{}{0pt}%
\pgfpathmoveto{\pgfqpoint{2.459372in}{2.593099in}}%
\pgfpathlineto{\pgfqpoint{2.716885in}{2.870430in}}%
\pgfpathlineto{\pgfqpoint{2.674215in}{2.908477in}}%
\pgfpathlineto{\pgfqpoint{2.459372in}{2.593099in}}%
\pgfpathclose%
\pgfusepath{fill}%
\end{pgfscope}%
\begin{pgfscope}%
\pgfpathrectangle{\pgfqpoint{0.100000in}{0.100000in}}{\pgfqpoint{3.957178in}{3.957178in}}%
\pgfusepath{clip}%
\pgfsetbuttcap%
\pgfsetroundjoin%
\definecolor{currentfill}{rgb}{0.994141,0.753137,0.161404}%
\pgfsetfillcolor{currentfill}%
\pgfsetlinewidth{0.000000pt}%
\definecolor{currentstroke}{rgb}{0.000000,0.000000,0.000000}%
\pgfsetstrokecolor{currentstroke}%
\pgfsetdash{}{0pt}%
\pgfpathmoveto{\pgfqpoint{1.759958in}{2.718143in}}%
\pgfpathlineto{\pgfqpoint{1.747542in}{2.664918in}}%
\pgfpathlineto{\pgfqpoint{1.861239in}{2.565895in}}%
\pgfpathlineto{\pgfqpoint{1.759958in}{2.718143in}}%
\pgfpathclose%
\pgfusepath{fill}%
\end{pgfscope}%
\begin{pgfscope}%
\pgfpathrectangle{\pgfqpoint{0.100000in}{0.100000in}}{\pgfqpoint{3.957178in}{3.957178in}}%
\pgfusepath{clip}%
\pgfsetbuttcap%
\pgfsetroundjoin%
\definecolor{currentfill}{rgb}{0.912800,0.444029,0.346251}%
\pgfsetfillcolor{currentfill}%
\pgfsetlinewidth{0.000000pt}%
\definecolor{currentstroke}{rgb}{0.000000,0.000000,0.000000}%
\pgfsetstrokecolor{currentstroke}%
\pgfsetdash{}{0pt}%
\pgfpathmoveto{\pgfqpoint{2.185536in}{2.224054in}}%
\pgfpathlineto{\pgfqpoint{2.083883in}{2.155265in}}%
\pgfpathlineto{\pgfqpoint{2.196589in}{2.449515in}}%
\pgfpathlineto{\pgfqpoint{2.185536in}{2.224054in}}%
\pgfpathclose%
\pgfusepath{fill}%
\end{pgfscope}%
\begin{pgfscope}%
\pgfpathrectangle{\pgfqpoint{0.100000in}{0.100000in}}{\pgfqpoint{3.957178in}{3.957178in}}%
\pgfusepath{clip}%
\pgfsetbuttcap%
\pgfsetroundjoin%
\definecolor{currentfill}{rgb}{0.920714,0.458417,0.336166}%
\pgfsetfillcolor{currentfill}%
\pgfsetlinewidth{0.000000pt}%
\definecolor{currentstroke}{rgb}{0.000000,0.000000,0.000000}%
\pgfsetstrokecolor{currentstroke}%
\pgfsetdash{}{0pt}%
\pgfpathmoveto{\pgfqpoint{3.293789in}{2.110296in}}%
\pgfpathlineto{\pgfqpoint{3.145728in}{2.407588in}}%
\pgfpathlineto{\pgfqpoint{3.104331in}{2.335002in}}%
\pgfpathlineto{\pgfqpoint{3.293789in}{2.110296in}}%
\pgfpathclose%
\pgfusepath{fill}%
\end{pgfscope}%
\begin{pgfscope}%
\pgfpathrectangle{\pgfqpoint{0.100000in}{0.100000in}}{\pgfqpoint{3.957178in}{3.957178in}}%
\pgfusepath{clip}%
\pgfsetbuttcap%
\pgfsetroundjoin%
\definecolor{currentfill}{rgb}{0.977856,0.602051,0.241387}%
\pgfsetfillcolor{currentfill}%
\pgfsetlinewidth{0.000000pt}%
\definecolor{currentstroke}{rgb}{0.000000,0.000000,0.000000}%
\pgfsetstrokecolor{currentstroke}%
\pgfsetdash{}{0pt}%
\pgfpathmoveto{\pgfqpoint{3.096663in}{2.540795in}}%
\pgfpathlineto{\pgfqpoint{3.083252in}{2.463869in}}%
\pgfpathlineto{\pgfqpoint{3.145728in}{2.407588in}}%
\pgfpathlineto{\pgfqpoint{3.096663in}{2.540795in}}%
\pgfpathclose%
\pgfusepath{fill}%
\end{pgfscope}%
\begin{pgfscope}%
\pgfpathrectangle{\pgfqpoint{0.100000in}{0.100000in}}{\pgfqpoint{3.957178in}{3.957178in}}%
\pgfusepath{clip}%
\pgfsetbuttcap%
\pgfsetroundjoin%
\definecolor{currentfill}{rgb}{0.970533,0.887896,0.145919}%
\pgfsetfillcolor{currentfill}%
\pgfsetlinewidth{0.000000pt}%
\definecolor{currentstroke}{rgb}{0.000000,0.000000,0.000000}%
\pgfsetstrokecolor{currentstroke}%
\pgfsetdash{}{0pt}%
\pgfpathmoveto{\pgfqpoint{1.350836in}{2.803045in}}%
\pgfpathlineto{\pgfqpoint{1.556620in}{2.782331in}}%
\pgfpathlineto{\pgfqpoint{1.759958in}{2.718143in}}%
\pgfpathlineto{\pgfqpoint{1.350836in}{2.803045in}}%
\pgfpathclose%
\pgfusepath{fill}%
\end{pgfscope}%
\begin{pgfscope}%
\pgfpathrectangle{\pgfqpoint{0.100000in}{0.100000in}}{\pgfqpoint{3.957178in}{3.957178in}}%
\pgfusepath{clip}%
\pgfsetbuttcap%
\pgfsetroundjoin%
\definecolor{currentfill}{rgb}{0.963203,0.553865,0.271909}%
\pgfsetfillcolor{currentfill}%
\pgfsetlinewidth{0.000000pt}%
\definecolor{currentstroke}{rgb}{0.000000,0.000000,0.000000}%
\pgfsetstrokecolor{currentstroke}%
\pgfsetdash{}{0pt}%
\pgfpathmoveto{\pgfqpoint{1.246560in}{2.412228in}}%
\pgfpathlineto{\pgfqpoint{1.306609in}{2.699955in}}%
\pgfpathlineto{\pgfqpoint{0.942330in}{2.102782in}}%
\pgfpathlineto{\pgfqpoint{1.246560in}{2.412228in}}%
\pgfpathclose%
\pgfusepath{fill}%
\end{pgfscope}%
\begin{pgfscope}%
\pgfpathrectangle{\pgfqpoint{0.100000in}{0.100000in}}{\pgfqpoint{3.957178in}{3.957178in}}%
\pgfusepath{clip}%
\pgfsetbuttcap%
\pgfsetroundjoin%
\definecolor{currentfill}{rgb}{0.994495,0.740880,0.166335}%
\pgfsetfillcolor{currentfill}%
\pgfsetlinewidth{0.000000pt}%
\definecolor{currentstroke}{rgb}{0.000000,0.000000,0.000000}%
\pgfsetstrokecolor{currentstroke}%
\pgfsetdash{}{0pt}%
\pgfpathmoveto{\pgfqpoint{1.790823in}{2.597170in}}%
\pgfpathlineto{\pgfqpoint{1.861239in}{2.565895in}}%
\pgfpathlineto{\pgfqpoint{1.747542in}{2.664918in}}%
\pgfpathlineto{\pgfqpoint{1.790823in}{2.597170in}}%
\pgfpathclose%
\pgfusepath{fill}%
\end{pgfscope}%
\begin{pgfscope}%
\pgfpathrectangle{\pgfqpoint{0.100000in}{0.100000in}}{\pgfqpoint{3.957178in}{3.957178in}}%
\pgfusepath{clip}%
\pgfsetbuttcap%
\pgfsetroundjoin%
\definecolor{currentfill}{rgb}{0.972530,0.881250,0.144923}%
\pgfsetfillcolor{currentfill}%
\pgfsetlinewidth{0.000000pt}%
\definecolor{currentstroke}{rgb}{0.000000,0.000000,0.000000}%
\pgfsetstrokecolor{currentstroke}%
\pgfsetdash{}{0pt}%
\pgfpathmoveto{\pgfqpoint{1.350836in}{2.803045in}}%
\pgfpathlineto{\pgfqpoint{1.306609in}{2.699955in}}%
\pgfpathlineto{\pgfqpoint{1.399155in}{2.753227in}}%
\pgfpathlineto{\pgfqpoint{1.350836in}{2.803045in}}%
\pgfpathclose%
\pgfusepath{fill}%
\end{pgfscope}%
\begin{pgfscope}%
\pgfpathrectangle{\pgfqpoint{0.100000in}{0.100000in}}{\pgfqpoint{3.957178in}{3.957178in}}%
\pgfusepath{clip}%
\pgfsetbuttcap%
\pgfsetroundjoin%
\definecolor{currentfill}{rgb}{0.994103,0.710698,0.180097}%
\pgfsetfillcolor{currentfill}%
\pgfsetlinewidth{0.000000pt}%
\definecolor{currentstroke}{rgb}{0.000000,0.000000,0.000000}%
\pgfsetstrokecolor{currentstroke}%
\pgfsetdash{}{0pt}%
\pgfpathmoveto{\pgfqpoint{1.835085in}{2.543290in}}%
\pgfpathlineto{\pgfqpoint{1.861239in}{2.565895in}}%
\pgfpathlineto{\pgfqpoint{1.790823in}{2.597170in}}%
\pgfpathlineto{\pgfqpoint{1.835085in}{2.543290in}}%
\pgfpathclose%
\pgfusepath{fill}%
\end{pgfscope}%
\begin{pgfscope}%
\pgfpathrectangle{\pgfqpoint{0.100000in}{0.100000in}}{\pgfqpoint{3.957178in}{3.957178in}}%
\pgfusepath{clip}%
\pgfsetbuttcap%
\pgfsetroundjoin%
\definecolor{currentfill}{rgb}{0.972530,0.881250,0.144923}%
\pgfsetfillcolor{currentfill}%
\pgfsetlinewidth{0.000000pt}%
\definecolor{currentstroke}{rgb}{0.000000,0.000000,0.000000}%
\pgfsetstrokecolor{currentstroke}%
\pgfsetdash{}{0pt}%
\pgfpathmoveto{\pgfqpoint{2.849017in}{2.840181in}}%
\pgfpathlineto{\pgfqpoint{2.771031in}{2.746435in}}%
\pgfpathlineto{\pgfqpoint{3.033800in}{2.651652in}}%
\pgfpathlineto{\pgfqpoint{2.849017in}{2.840181in}}%
\pgfpathclose%
\pgfusepath{fill}%
\end{pgfscope}%
\begin{pgfscope}%
\pgfpathrectangle{\pgfqpoint{0.100000in}{0.100000in}}{\pgfqpoint{3.957178in}{3.957178in}}%
\pgfusepath{clip}%
\pgfsetbuttcap%
\pgfsetroundjoin%
\definecolor{currentfill}{rgb}{0.963203,0.553865,0.271909}%
\pgfsetfillcolor{currentfill}%
\pgfsetlinewidth{0.000000pt}%
\definecolor{currentstroke}{rgb}{0.000000,0.000000,0.000000}%
\pgfsetstrokecolor{currentstroke}%
\pgfsetdash{}{0pt}%
\pgfpathmoveto{\pgfqpoint{1.835085in}{2.543290in}}%
\pgfpathlineto{\pgfqpoint{2.196589in}{2.449515in}}%
\pgfpathlineto{\pgfqpoint{2.083883in}{2.155265in}}%
\pgfpathlineto{\pgfqpoint{1.835085in}{2.543290in}}%
\pgfpathclose%
\pgfusepath{fill}%
\end{pgfscope}%
\begin{pgfscope}%
\pgfpathrectangle{\pgfqpoint{0.100000in}{0.100000in}}{\pgfqpoint{3.957178in}{3.957178in}}%
\pgfusepath{clip}%
\pgfsetbuttcap%
\pgfsetroundjoin%
\definecolor{currentfill}{rgb}{0.893250,0.411048,0.369768}%
\pgfsetfillcolor{currentfill}%
\pgfsetlinewidth{0.000000pt}%
\definecolor{currentstroke}{rgb}{0.000000,0.000000,0.000000}%
\pgfsetstrokecolor{currentstroke}%
\pgfsetdash{}{0pt}%
\pgfpathmoveto{\pgfqpoint{3.104331in}{2.335002in}}%
\pgfpathlineto{\pgfqpoint{3.176004in}{2.131408in}}%
\pgfpathlineto{\pgfqpoint{3.293789in}{2.110296in}}%
\pgfpathlineto{\pgfqpoint{3.104331in}{2.335002in}}%
\pgfpathclose%
\pgfusepath{fill}%
\end{pgfscope}%
\begin{pgfscope}%
\pgfpathrectangle{\pgfqpoint{0.100000in}{0.100000in}}{\pgfqpoint{3.957178in}{3.957178in}}%
\pgfusepath{clip}%
\pgfsetbuttcap%
\pgfsetroundjoin%
\definecolor{currentfill}{rgb}{0.976265,0.868016,0.143351}%
\pgfsetfillcolor{currentfill}%
\pgfsetlinewidth{0.000000pt}%
\definecolor{currentstroke}{rgb}{0.000000,0.000000,0.000000}%
\pgfsetstrokecolor{currentstroke}%
\pgfsetdash{}{0pt}%
\pgfpathmoveto{\pgfqpoint{1.556620in}{2.782331in}}%
\pgfpathlineto{\pgfqpoint{1.747542in}{2.664918in}}%
\pgfpathlineto{\pgfqpoint{1.759958in}{2.718143in}}%
\pgfpathlineto{\pgfqpoint{1.556620in}{2.782331in}}%
\pgfpathclose%
\pgfusepath{fill}%
\end{pgfscope}%
\begin{pgfscope}%
\pgfpathrectangle{\pgfqpoint{0.100000in}{0.100000in}}{\pgfqpoint{3.957178in}{3.957178in}}%
\pgfusepath{clip}%
\pgfsetbuttcap%
\pgfsetroundjoin%
\definecolor{currentfill}{rgb}{0.941896,0.968590,0.140956}%
\pgfsetfillcolor{currentfill}%
\pgfsetlinewidth{0.000000pt}%
\definecolor{currentstroke}{rgb}{0.000000,0.000000,0.000000}%
\pgfsetstrokecolor{currentstroke}%
\pgfsetdash{}{0pt}%
\pgfpathmoveto{\pgfqpoint{2.716885in}{2.870430in}}%
\pgfpathlineto{\pgfqpoint{2.771031in}{2.746435in}}%
\pgfpathlineto{\pgfqpoint{2.849017in}{2.840181in}}%
\pgfpathlineto{\pgfqpoint{2.716885in}{2.870430in}}%
\pgfpathclose%
\pgfusepath{fill}%
\end{pgfscope}%
\begin{pgfscope}%
\pgfpathrectangle{\pgfqpoint{0.100000in}{0.100000in}}{\pgfqpoint{3.957178in}{3.957178in}}%
\pgfusepath{clip}%
\pgfsetbuttcap%
\pgfsetroundjoin%
\definecolor{currentfill}{rgb}{0.956808,0.928152,0.152409}%
\pgfsetfillcolor{currentfill}%
\pgfsetlinewidth{0.000000pt}%
\definecolor{currentstroke}{rgb}{0.000000,0.000000,0.000000}%
\pgfsetstrokecolor{currentstroke}%
\pgfsetdash{}{0pt}%
\pgfpathmoveto{\pgfqpoint{1.399155in}{2.753227in}}%
\pgfpathlineto{\pgfqpoint{1.556620in}{2.782331in}}%
\pgfpathlineto{\pgfqpoint{1.350836in}{2.803045in}}%
\pgfpathlineto{\pgfqpoint{1.399155in}{2.753227in}}%
\pgfpathclose%
\pgfusepath{fill}%
\end{pgfscope}%
\begin{pgfscope}%
\pgfpathrectangle{\pgfqpoint{0.100000in}{0.100000in}}{\pgfqpoint{3.957178in}{3.957178in}}%
\pgfusepath{clip}%
\pgfsetbuttcap%
\pgfsetroundjoin%
\definecolor{currentfill}{rgb}{0.840155,0.333580,0.427455}%
\pgfsetfillcolor{currentfill}%
\pgfsetlinewidth{0.000000pt}%
\definecolor{currentstroke}{rgb}{0.000000,0.000000,0.000000}%
\pgfsetstrokecolor{currentstroke}%
\pgfsetdash{}{0pt}%
\pgfpathmoveto{\pgfqpoint{0.931340in}{1.874842in}}%
\pgfpathlineto{\pgfqpoint{0.971959in}{1.915950in}}%
\pgfpathlineto{\pgfqpoint{1.246560in}{2.412228in}}%
\pgfpathlineto{\pgfqpoint{0.931340in}{1.874842in}}%
\pgfpathclose%
\pgfusepath{fill}%
\end{pgfscope}%
\begin{pgfscope}%
\pgfpathrectangle{\pgfqpoint{0.100000in}{0.100000in}}{\pgfqpoint{3.957178in}{3.957178in}}%
\pgfusepath{clip}%
\pgfsetbuttcap%
\pgfsetroundjoin%
\definecolor{currentfill}{rgb}{0.973416,0.585761,0.251540}%
\pgfsetfillcolor{currentfill}%
\pgfsetlinewidth{0.000000pt}%
\definecolor{currentstroke}{rgb}{0.000000,0.000000,0.000000}%
\pgfsetstrokecolor{currentstroke}%
\pgfsetdash{}{0pt}%
\pgfpathmoveto{\pgfqpoint{3.145728in}{2.407588in}}%
\pgfpathlineto{\pgfqpoint{3.083252in}{2.463869in}}%
\pgfpathlineto{\pgfqpoint{3.104331in}{2.335002in}}%
\pgfpathlineto{\pgfqpoint{3.145728in}{2.407588in}}%
\pgfpathclose%
\pgfusepath{fill}%
\end{pgfscope}%
\begin{pgfscope}%
\pgfpathrectangle{\pgfqpoint{0.100000in}{0.100000in}}{\pgfqpoint{3.957178in}{3.957178in}}%
\pgfusepath{clip}%
\pgfsetbuttcap%
\pgfsetroundjoin%
\definecolor{currentfill}{rgb}{0.968443,0.894564,0.147014}%
\pgfsetfillcolor{currentfill}%
\pgfsetlinewidth{0.000000pt}%
\definecolor{currentstroke}{rgb}{0.000000,0.000000,0.000000}%
\pgfsetstrokecolor{currentstroke}%
\pgfsetdash{}{0pt}%
\pgfpathmoveto{\pgfqpoint{2.771031in}{2.746435in}}%
\pgfpathlineto{\pgfqpoint{2.716885in}{2.870430in}}%
\pgfpathlineto{\pgfqpoint{2.459372in}{2.593099in}}%
\pgfpathlineto{\pgfqpoint{2.771031in}{2.746435in}}%
\pgfpathclose%
\pgfusepath{fill}%
\end{pgfscope}%
\begin{pgfscope}%
\pgfpathrectangle{\pgfqpoint{0.100000in}{0.100000in}}{\pgfqpoint{3.957178in}{3.957178in}}%
\pgfusepath{clip}%
\pgfsetbuttcap%
\pgfsetroundjoin%
\definecolor{currentfill}{rgb}{0.562738,0.051545,0.641509}%
\pgfsetfillcolor{currentfill}%
\pgfsetlinewidth{0.000000pt}%
\definecolor{currentstroke}{rgb}{0.000000,0.000000,0.000000}%
\pgfsetstrokecolor{currentstroke}%
\pgfsetdash{}{0pt}%
\pgfpathmoveto{\pgfqpoint{1.160246in}{1.693711in}}%
\pgfpathlineto{\pgfqpoint{0.971959in}{1.915950in}}%
\pgfpathlineto{\pgfqpoint{0.747252in}{1.316186in}}%
\pgfpathlineto{\pgfqpoint{1.160246in}{1.693711in}}%
\pgfpathclose%
\pgfusepath{fill}%
\end{pgfscope}%
\begin{pgfscope}%
\pgfpathrectangle{\pgfqpoint{0.100000in}{0.100000in}}{\pgfqpoint{3.957178in}{3.957178in}}%
\pgfusepath{clip}%
\pgfsetbuttcap%
\pgfsetroundjoin%
\definecolor{currentfill}{rgb}{0.830018,0.320172,0.437836}%
\pgfsetfillcolor{currentfill}%
\pgfsetlinewidth{0.000000pt}%
\definecolor{currentstroke}{rgb}{0.000000,0.000000,0.000000}%
\pgfsetstrokecolor{currentstroke}%
\pgfsetdash{}{0pt}%
\pgfpathmoveto{\pgfqpoint{3.157732in}{1.874966in}}%
\pgfpathlineto{\pgfqpoint{3.293789in}{2.110296in}}%
\pgfpathlineto{\pgfqpoint{3.176004in}{2.131408in}}%
\pgfpathlineto{\pgfqpoint{3.157732in}{1.874966in}}%
\pgfpathclose%
\pgfusepath{fill}%
\end{pgfscope}%
\begin{pgfscope}%
\pgfpathrectangle{\pgfqpoint{0.100000in}{0.100000in}}{\pgfqpoint{3.957178in}{3.957178in}}%
\pgfusepath{clip}%
\pgfsetbuttcap%
\pgfsetroundjoin%
\definecolor{currentfill}{rgb}{0.993851,0.759304,0.159092}%
\pgfsetfillcolor{currentfill}%
\pgfsetlinewidth{0.000000pt}%
\definecolor{currentstroke}{rgb}{0.000000,0.000000,0.000000}%
\pgfsetstrokecolor{currentstroke}%
\pgfsetdash{}{0pt}%
\pgfpathmoveto{\pgfqpoint{2.874143in}{2.655359in}}%
\pgfpathlineto{\pgfqpoint{3.083252in}{2.463869in}}%
\pgfpathlineto{\pgfqpoint{3.033800in}{2.651652in}}%
\pgfpathlineto{\pgfqpoint{2.874143in}{2.655359in}}%
\pgfpathclose%
\pgfusepath{fill}%
\end{pgfscope}%
\begin{pgfscope}%
\pgfpathrectangle{\pgfqpoint{0.100000in}{0.100000in}}{\pgfqpoint{3.957178in}{3.957178in}}%
\pgfusepath{clip}%
\pgfsetbuttcap%
\pgfsetroundjoin%
\definecolor{currentfill}{rgb}{0.981190,0.848329,0.142279}%
\pgfsetfillcolor{currentfill}%
\pgfsetlinewidth{0.000000pt}%
\definecolor{currentstroke}{rgb}{0.000000,0.000000,0.000000}%
\pgfsetstrokecolor{currentstroke}%
\pgfsetdash{}{0pt}%
\pgfpathmoveto{\pgfqpoint{1.747542in}{2.664918in}}%
\pgfpathlineto{\pgfqpoint{1.556620in}{2.782331in}}%
\pgfpathlineto{\pgfqpoint{1.790823in}{2.597170in}}%
\pgfpathlineto{\pgfqpoint{1.747542in}{2.664918in}}%
\pgfpathclose%
\pgfusepath{fill}%
\end{pgfscope}%
\begin{pgfscope}%
\pgfpathrectangle{\pgfqpoint{0.100000in}{0.100000in}}{\pgfqpoint{3.957178in}{3.957178in}}%
\pgfusepath{clip}%
\pgfsetbuttcap%
\pgfsetroundjoin%
\definecolor{currentfill}{rgb}{0.964021,0.907950,0.149370}%
\pgfsetfillcolor{currentfill}%
\pgfsetlinewidth{0.000000pt}%
\definecolor{currentstroke}{rgb}{0.000000,0.000000,0.000000}%
\pgfsetstrokecolor{currentstroke}%
\pgfsetdash{}{0pt}%
\pgfpathmoveto{\pgfqpoint{1.391550in}{2.747707in}}%
\pgfpathlineto{\pgfqpoint{1.399155in}{2.753227in}}%
\pgfpathlineto{\pgfqpoint{1.306609in}{2.699955in}}%
\pgfpathlineto{\pgfqpoint{1.391550in}{2.747707in}}%
\pgfpathclose%
\pgfusepath{fill}%
\end{pgfscope}%
\begin{pgfscope}%
\pgfpathrectangle{\pgfqpoint{0.100000in}{0.100000in}}{\pgfqpoint{3.957178in}{3.957178in}}%
\pgfusepath{clip}%
\pgfsetbuttcap%
\pgfsetroundjoin%
\definecolor{currentfill}{rgb}{0.987332,0.646633,0.214648}%
\pgfsetfillcolor{currentfill}%
\pgfsetlinewidth{0.000000pt}%
\definecolor{currentstroke}{rgb}{0.000000,0.000000,0.000000}%
\pgfsetstrokecolor{currentstroke}%
\pgfsetdash{}{0pt}%
\pgfpathmoveto{\pgfqpoint{2.392659in}{2.559380in}}%
\pgfpathlineto{\pgfqpoint{2.459372in}{2.593099in}}%
\pgfpathlineto{\pgfqpoint{2.185536in}{2.224054in}}%
\pgfpathlineto{\pgfqpoint{2.392659in}{2.559380in}}%
\pgfpathclose%
\pgfusepath{fill}%
\end{pgfscope}%
\begin{pgfscope}%
\pgfpathrectangle{\pgfqpoint{0.100000in}{0.100000in}}{\pgfqpoint{3.957178in}{3.957178in}}%
\pgfusepath{clip}%
\pgfsetbuttcap%
\pgfsetroundjoin%
\definecolor{currentfill}{rgb}{0.976265,0.868016,0.143351}%
\pgfsetfillcolor{currentfill}%
\pgfsetlinewidth{0.000000pt}%
\definecolor{currentstroke}{rgb}{0.000000,0.000000,0.000000}%
\pgfsetstrokecolor{currentstroke}%
\pgfsetdash{}{0pt}%
\pgfpathmoveto{\pgfqpoint{3.033800in}{2.651652in}}%
\pgfpathlineto{\pgfqpoint{2.771031in}{2.746435in}}%
\pgfpathlineto{\pgfqpoint{2.874143in}{2.655359in}}%
\pgfpathlineto{\pgfqpoint{3.033800in}{2.651652in}}%
\pgfpathclose%
\pgfusepath{fill}%
\end{pgfscope}%
\begin{pgfscope}%
\pgfpathrectangle{\pgfqpoint{0.100000in}{0.100000in}}{\pgfqpoint{3.957178in}{3.957178in}}%
\pgfusepath{clip}%
\pgfsetbuttcap%
\pgfsetroundjoin%
\definecolor{currentfill}{rgb}{0.981826,0.618572,0.231287}%
\pgfsetfillcolor{currentfill}%
\pgfsetlinewidth{0.000000pt}%
\definecolor{currentstroke}{rgb}{0.000000,0.000000,0.000000}%
\pgfsetstrokecolor{currentstroke}%
\pgfsetdash{}{0pt}%
\pgfpathmoveto{\pgfqpoint{3.104331in}{2.335002in}}%
\pgfpathlineto{\pgfqpoint{3.083252in}{2.463869in}}%
\pgfpathlineto{\pgfqpoint{3.059676in}{2.426644in}}%
\pgfpathlineto{\pgfqpoint{3.104331in}{2.335002in}}%
\pgfpathclose%
\pgfusepath{fill}%
\end{pgfscope}%
\begin{pgfscope}%
\pgfpathrectangle{\pgfqpoint{0.100000in}{0.100000in}}{\pgfqpoint{3.957178in}{3.957178in}}%
\pgfusepath{clip}%
\pgfsetbuttcap%
\pgfsetroundjoin%
\definecolor{currentfill}{rgb}{0.964021,0.907950,0.149370}%
\pgfsetfillcolor{currentfill}%
\pgfsetlinewidth{0.000000pt}%
\definecolor{currentstroke}{rgb}{0.000000,0.000000,0.000000}%
\pgfsetstrokecolor{currentstroke}%
\pgfsetdash{}{0pt}%
\pgfpathmoveto{\pgfqpoint{1.790823in}{2.597170in}}%
\pgfpathlineto{\pgfqpoint{1.556620in}{2.782331in}}%
\pgfpathlineto{\pgfqpoint{1.495458in}{2.759711in}}%
\pgfpathlineto{\pgfqpoint{1.790823in}{2.597170in}}%
\pgfpathclose%
\pgfusepath{fill}%
\end{pgfscope}%
\begin{pgfscope}%
\pgfpathrectangle{\pgfqpoint{0.100000in}{0.100000in}}{\pgfqpoint{3.957178in}{3.957178in}}%
\pgfusepath{clip}%
\pgfsetbuttcap%
\pgfsetroundjoin%
\definecolor{currentfill}{rgb}{0.987621,0.815978,0.144363}%
\pgfsetfillcolor{currentfill}%
\pgfsetlinewidth{0.000000pt}%
\definecolor{currentstroke}{rgb}{0.000000,0.000000,0.000000}%
\pgfsetstrokecolor{currentstroke}%
\pgfsetdash{}{0pt}%
\pgfpathmoveto{\pgfqpoint{1.306609in}{2.699955in}}%
\pgfpathlineto{\pgfqpoint{1.246560in}{2.412228in}}%
\pgfpathlineto{\pgfqpoint{1.391550in}{2.747707in}}%
\pgfpathlineto{\pgfqpoint{1.306609in}{2.699955in}}%
\pgfpathclose%
\pgfusepath{fill}%
\end{pgfscope}%
\begin{pgfscope}%
\pgfpathrectangle{\pgfqpoint{0.100000in}{0.100000in}}{\pgfqpoint{3.957178in}{3.957178in}}%
\pgfusepath{clip}%
\pgfsetbuttcap%
\pgfsetroundjoin%
\definecolor{currentfill}{rgb}{0.949151,0.948435,0.152178}%
\pgfsetfillcolor{currentfill}%
\pgfsetlinewidth{0.000000pt}%
\definecolor{currentstroke}{rgb}{0.000000,0.000000,0.000000}%
\pgfsetstrokecolor{currentstroke}%
\pgfsetdash{}{0pt}%
\pgfpathmoveto{\pgfqpoint{1.391550in}{2.747707in}}%
\pgfpathlineto{\pgfqpoint{1.495458in}{2.759711in}}%
\pgfpathlineto{\pgfqpoint{1.399155in}{2.753227in}}%
\pgfpathlineto{\pgfqpoint{1.391550in}{2.747707in}}%
\pgfpathclose%
\pgfusepath{fill}%
\end{pgfscope}%
\begin{pgfscope}%
\pgfpathrectangle{\pgfqpoint{0.100000in}{0.100000in}}{\pgfqpoint{3.957178in}{3.957178in}}%
\pgfusepath{clip}%
\pgfsetbuttcap%
\pgfsetroundjoin%
\definecolor{currentfill}{rgb}{0.941896,0.968590,0.140956}%
\pgfsetfillcolor{currentfill}%
\pgfsetlinewidth{0.000000pt}%
\definecolor{currentstroke}{rgb}{0.000000,0.000000,0.000000}%
\pgfsetstrokecolor{currentstroke}%
\pgfsetdash{}{0pt}%
\pgfpathmoveto{\pgfqpoint{1.556620in}{2.782331in}}%
\pgfpathlineto{\pgfqpoint{1.399155in}{2.753227in}}%
\pgfpathlineto{\pgfqpoint{1.495458in}{2.759711in}}%
\pgfpathlineto{\pgfqpoint{1.556620in}{2.782331in}}%
\pgfpathclose%
\pgfusepath{fill}%
\end{pgfscope}%
\begin{pgfscope}%
\pgfpathrectangle{\pgfqpoint{0.100000in}{0.100000in}}{\pgfqpoint{3.957178in}{3.957178in}}%
\pgfusepath{clip}%
\pgfsetbuttcap%
\pgfsetroundjoin%
\definecolor{currentfill}{rgb}{0.994553,0.728728,0.171622}%
\pgfsetfillcolor{currentfill}%
\pgfsetlinewidth{0.000000pt}%
\definecolor{currentstroke}{rgb}{0.000000,0.000000,0.000000}%
\pgfsetstrokecolor{currentstroke}%
\pgfsetdash{}{0pt}%
\pgfpathmoveto{\pgfqpoint{3.059676in}{2.426644in}}%
\pgfpathlineto{\pgfqpoint{3.083252in}{2.463869in}}%
\pgfpathlineto{\pgfqpoint{2.889972in}{2.631947in}}%
\pgfpathlineto{\pgfqpoint{3.059676in}{2.426644in}}%
\pgfpathclose%
\pgfusepath{fill}%
\end{pgfscope}%
\begin{pgfscope}%
\pgfpathrectangle{\pgfqpoint{0.100000in}{0.100000in}}{\pgfqpoint{3.957178in}{3.957178in}}%
\pgfusepath{clip}%
\pgfsetbuttcap%
\pgfsetroundjoin%
\definecolor{currentfill}{rgb}{0.986345,0.640969,0.217948}%
\pgfsetfillcolor{currentfill}%
\pgfsetlinewidth{0.000000pt}%
\definecolor{currentstroke}{rgb}{0.000000,0.000000,0.000000}%
\pgfsetstrokecolor{currentstroke}%
\pgfsetdash{}{0pt}%
\pgfpathmoveto{\pgfqpoint{2.392659in}{2.559380in}}%
\pgfpathlineto{\pgfqpoint{2.185536in}{2.224054in}}%
\pgfpathlineto{\pgfqpoint{2.196589in}{2.449515in}}%
\pgfpathlineto{\pgfqpoint{2.392659in}{2.559380in}}%
\pgfpathclose%
\pgfusepath{fill}%
\end{pgfscope}%
\begin{pgfscope}%
\pgfpathrectangle{\pgfqpoint{0.100000in}{0.100000in}}{\pgfqpoint{3.957178in}{3.957178in}}%
\pgfusepath{clip}%
\pgfsetbuttcap%
\pgfsetroundjoin%
\definecolor{currentfill}{rgb}{0.925825,0.468103,0.329435}%
\pgfsetfillcolor{currentfill}%
\pgfsetlinewidth{0.000000pt}%
\definecolor{currentstroke}{rgb}{0.000000,0.000000,0.000000}%
\pgfsetstrokecolor{currentstroke}%
\pgfsetdash{}{0pt}%
\pgfpathmoveto{\pgfqpoint{1.246560in}{2.412228in}}%
\pgfpathlineto{\pgfqpoint{0.971959in}{1.915950in}}%
\pgfpathlineto{\pgfqpoint{1.230510in}{2.226458in}}%
\pgfpathlineto{\pgfqpoint{1.246560in}{2.412228in}}%
\pgfpathclose%
\pgfusepath{fill}%
\end{pgfscope}%
\begin{pgfscope}%
\pgfpathrectangle{\pgfqpoint{0.100000in}{0.100000in}}{\pgfqpoint{3.957178in}{3.957178in}}%
\pgfusepath{clip}%
\pgfsetbuttcap%
\pgfsetroundjoin%
\definecolor{currentfill}{rgb}{0.968443,0.894564,0.147014}%
\pgfsetfillcolor{currentfill}%
\pgfsetlinewidth{0.000000pt}%
\definecolor{currentstroke}{rgb}{0.000000,0.000000,0.000000}%
\pgfsetstrokecolor{currentstroke}%
\pgfsetdash{}{0pt}%
\pgfpathmoveto{\pgfqpoint{2.459372in}{2.593099in}}%
\pgfpathlineto{\pgfqpoint{2.673729in}{2.664530in}}%
\pgfpathlineto{\pgfqpoint{2.771031in}{2.746435in}}%
\pgfpathlineto{\pgfqpoint{2.459372in}{2.593099in}}%
\pgfpathclose%
\pgfusepath{fill}%
\end{pgfscope}%
\begin{pgfscope}%
\pgfpathrectangle{\pgfqpoint{0.100000in}{0.100000in}}{\pgfqpoint{3.957178in}{3.957178in}}%
\pgfusepath{clip}%
\pgfsetbuttcap%
\pgfsetroundjoin%
\definecolor{currentfill}{rgb}{0.947051,0.512699,0.299049}%
\pgfsetfillcolor{currentfill}%
\pgfsetlinewidth{0.000000pt}%
\definecolor{currentstroke}{rgb}{0.000000,0.000000,0.000000}%
\pgfsetstrokecolor{currentstroke}%
\pgfsetdash{}{0pt}%
\pgfpathmoveto{\pgfqpoint{3.176004in}{2.131408in}}%
\pgfpathlineto{\pgfqpoint{3.104331in}{2.335002in}}%
\pgfpathlineto{\pgfqpoint{3.040008in}{2.239939in}}%
\pgfpathlineto{\pgfqpoint{3.176004in}{2.131408in}}%
\pgfpathclose%
\pgfusepath{fill}%
\end{pgfscope}%
\begin{pgfscope}%
\pgfpathrectangle{\pgfqpoint{0.100000in}{0.100000in}}{\pgfqpoint{3.957178in}{3.957178in}}%
\pgfusepath{clip}%
\pgfsetbuttcap%
\pgfsetroundjoin%
\definecolor{currentfill}{rgb}{0.987621,0.815978,0.144363}%
\pgfsetfillcolor{currentfill}%
\pgfsetlinewidth{0.000000pt}%
\definecolor{currentstroke}{rgb}{0.000000,0.000000,0.000000}%
\pgfsetstrokecolor{currentstroke}%
\pgfsetdash{}{0pt}%
\pgfpathmoveto{\pgfqpoint{2.874143in}{2.655359in}}%
\pgfpathlineto{\pgfqpoint{2.889972in}{2.631947in}}%
\pgfpathlineto{\pgfqpoint{3.083252in}{2.463869in}}%
\pgfpathlineto{\pgfqpoint{2.874143in}{2.655359in}}%
\pgfpathclose%
\pgfusepath{fill}%
\end{pgfscope}%
\begin{pgfscope}%
\pgfpathrectangle{\pgfqpoint{0.100000in}{0.100000in}}{\pgfqpoint{3.957178in}{3.957178in}}%
\pgfusepath{clip}%
\pgfsetbuttcap%
\pgfsetroundjoin%
\definecolor{currentfill}{rgb}{0.976428,0.596595,0.244767}%
\pgfsetfillcolor{currentfill}%
\pgfsetlinewidth{0.000000pt}%
\definecolor{currentstroke}{rgb}{0.000000,0.000000,0.000000}%
\pgfsetstrokecolor{currentstroke}%
\pgfsetdash{}{0pt}%
\pgfpathmoveto{\pgfqpoint{3.040008in}{2.239939in}}%
\pgfpathlineto{\pgfqpoint{3.104331in}{2.335002in}}%
\pgfpathlineto{\pgfqpoint{3.059676in}{2.426644in}}%
\pgfpathlineto{\pgfqpoint{3.040008in}{2.239939in}}%
\pgfpathclose%
\pgfusepath{fill}%
\end{pgfscope}%
\begin{pgfscope}%
\pgfpathrectangle{\pgfqpoint{0.100000in}{0.100000in}}{\pgfqpoint{3.957178in}{3.957178in}}%
\pgfusepath{clip}%
\pgfsetbuttcap%
\pgfsetroundjoin%
\definecolor{currentfill}{rgb}{0.970533,0.887896,0.145919}%
\pgfsetfillcolor{currentfill}%
\pgfsetlinewidth{0.000000pt}%
\definecolor{currentstroke}{rgb}{0.000000,0.000000,0.000000}%
\pgfsetstrokecolor{currentstroke}%
\pgfsetdash{}{0pt}%
\pgfpathmoveto{\pgfqpoint{1.495458in}{2.759711in}}%
\pgfpathlineto{\pgfqpoint{1.430860in}{2.594652in}}%
\pgfpathlineto{\pgfqpoint{1.790823in}{2.597170in}}%
\pgfpathlineto{\pgfqpoint{1.495458in}{2.759711in}}%
\pgfpathclose%
\pgfusepath{fill}%
\end{pgfscope}%
\begin{pgfscope}%
\pgfpathrectangle{\pgfqpoint{0.100000in}{0.100000in}}{\pgfqpoint{3.957178in}{3.957178in}}%
\pgfusepath{clip}%
\pgfsetbuttcap%
\pgfsetroundjoin%
\definecolor{currentfill}{rgb}{0.823132,0.311261,0.444806}%
\pgfsetfillcolor{currentfill}%
\pgfsetlinewidth{0.000000pt}%
\definecolor{currentstroke}{rgb}{0.000000,0.000000,0.000000}%
\pgfsetstrokecolor{currentstroke}%
\pgfsetdash{}{0pt}%
\pgfpathmoveto{\pgfqpoint{1.230510in}{2.226458in}}%
\pgfpathlineto{\pgfqpoint{0.971959in}{1.915950in}}%
\pgfpathlineto{\pgfqpoint{1.160246in}{1.693711in}}%
\pgfpathlineto{\pgfqpoint{1.230510in}{2.226458in}}%
\pgfpathclose%
\pgfusepath{fill}%
\end{pgfscope}%
\begin{pgfscope}%
\pgfpathrectangle{\pgfqpoint{0.100000in}{0.100000in}}{\pgfqpoint{3.957178in}{3.957178in}}%
\pgfusepath{clip}%
\pgfsetbuttcap%
\pgfsetroundjoin%
\definecolor{currentfill}{rgb}{0.951726,0.941671,0.152925}%
\pgfsetfillcolor{currentfill}%
\pgfsetlinewidth{0.000000pt}%
\definecolor{currentstroke}{rgb}{0.000000,0.000000,0.000000}%
\pgfsetstrokecolor{currentstroke}%
\pgfsetdash{}{0pt}%
\pgfpathmoveto{\pgfqpoint{1.430860in}{2.594652in}}%
\pgfpathlineto{\pgfqpoint{1.495458in}{2.759711in}}%
\pgfpathlineto{\pgfqpoint{1.391550in}{2.747707in}}%
\pgfpathlineto{\pgfqpoint{1.430860in}{2.594652in}}%
\pgfpathclose%
\pgfusepath{fill}%
\end{pgfscope}%
\begin{pgfscope}%
\pgfpathrectangle{\pgfqpoint{0.100000in}{0.100000in}}{\pgfqpoint{3.957178in}{3.957178in}}%
\pgfusepath{clip}%
\pgfsetbuttcap%
\pgfsetroundjoin%
\definecolor{currentfill}{rgb}{0.893250,0.411048,0.369768}%
\pgfsetfillcolor{currentfill}%
\pgfsetlinewidth{0.000000pt}%
\definecolor{currentstroke}{rgb}{0.000000,0.000000,0.000000}%
\pgfsetstrokecolor{currentstroke}%
\pgfsetdash{}{0pt}%
\pgfpathmoveto{\pgfqpoint{3.157732in}{1.874966in}}%
\pgfpathlineto{\pgfqpoint{3.176004in}{2.131408in}}%
\pgfpathlineto{\pgfqpoint{3.040008in}{2.239939in}}%
\pgfpathlineto{\pgfqpoint{3.157732in}{1.874966in}}%
\pgfpathclose%
\pgfusepath{fill}%
\end{pgfscope}%
\begin{pgfscope}%
\pgfpathrectangle{\pgfqpoint{0.100000in}{0.100000in}}{\pgfqpoint{3.957178in}{3.957178in}}%
\pgfusepath{clip}%
\pgfsetbuttcap%
\pgfsetroundjoin%
\definecolor{currentfill}{rgb}{0.984031,0.835315,0.142528}%
\pgfsetfillcolor{currentfill}%
\pgfsetlinewidth{0.000000pt}%
\definecolor{currentstroke}{rgb}{0.000000,0.000000,0.000000}%
\pgfsetstrokecolor{currentstroke}%
\pgfsetdash{}{0pt}%
\pgfpathmoveto{\pgfqpoint{1.391550in}{2.747707in}}%
\pgfpathlineto{\pgfqpoint{1.246560in}{2.412228in}}%
\pgfpathlineto{\pgfqpoint{1.430860in}{2.594652in}}%
\pgfpathlineto{\pgfqpoint{1.391550in}{2.747707in}}%
\pgfpathclose%
\pgfusepath{fill}%
\end{pgfscope}%
\begin{pgfscope}%
\pgfpathrectangle{\pgfqpoint{0.100000in}{0.100000in}}{\pgfqpoint{3.957178in}{3.957178in}}%
\pgfusepath{clip}%
\pgfsetbuttcap%
\pgfsetroundjoin%
\definecolor{currentfill}{rgb}{0.993456,0.698810,0.186041}%
\pgfsetfillcolor{currentfill}%
\pgfsetlinewidth{0.000000pt}%
\definecolor{currentstroke}{rgb}{0.000000,0.000000,0.000000}%
\pgfsetstrokecolor{currentstroke}%
\pgfsetdash{}{0pt}%
\pgfpathmoveto{\pgfqpoint{2.889972in}{2.631947in}}%
\pgfpathlineto{\pgfqpoint{3.040008in}{2.239939in}}%
\pgfpathlineto{\pgfqpoint{3.059676in}{2.426644in}}%
\pgfpathlineto{\pgfqpoint{2.889972in}{2.631947in}}%
\pgfpathclose%
\pgfusepath{fill}%
\end{pgfscope}%
\begin{pgfscope}%
\pgfpathrectangle{\pgfqpoint{0.100000in}{0.100000in}}{\pgfqpoint{3.957178in}{3.957178in}}%
\pgfusepath{clip}%
\pgfsetbuttcap%
\pgfsetroundjoin%
\definecolor{currentfill}{rgb}{0.977995,0.861432,0.142808}%
\pgfsetfillcolor{currentfill}%
\pgfsetlinewidth{0.000000pt}%
\definecolor{currentstroke}{rgb}{0.000000,0.000000,0.000000}%
\pgfsetstrokecolor{currentstroke}%
\pgfsetdash{}{0pt}%
\pgfpathmoveto{\pgfqpoint{2.673729in}{2.664530in}}%
\pgfpathlineto{\pgfqpoint{2.459372in}{2.593099in}}%
\pgfpathlineto{\pgfqpoint{2.392659in}{2.559380in}}%
\pgfpathlineto{\pgfqpoint{2.673729in}{2.664530in}}%
\pgfpathclose%
\pgfusepath{fill}%
\end{pgfscope}%
\begin{pgfscope}%
\pgfpathrectangle{\pgfqpoint{0.100000in}{0.100000in}}{\pgfqpoint{3.957178in}{3.957178in}}%
\pgfusepath{clip}%
\pgfsetbuttcap%
\pgfsetroundjoin%
\definecolor{currentfill}{rgb}{0.988648,0.809579,0.145357}%
\pgfsetfillcolor{currentfill}%
\pgfsetlinewidth{0.000000pt}%
\definecolor{currentstroke}{rgb}{0.000000,0.000000,0.000000}%
\pgfsetstrokecolor{currentstroke}%
\pgfsetdash{}{0pt}%
\pgfpathmoveto{\pgfqpoint{1.790823in}{2.597170in}}%
\pgfpathlineto{\pgfqpoint{1.737144in}{2.491232in}}%
\pgfpathlineto{\pgfqpoint{1.835085in}{2.543290in}}%
\pgfpathlineto{\pgfqpoint{1.790823in}{2.597170in}}%
\pgfpathclose%
\pgfusepath{fill}%
\end{pgfscope}%
\begin{pgfscope}%
\pgfpathrectangle{\pgfqpoint{0.100000in}{0.100000in}}{\pgfqpoint{3.957178in}{3.957178in}}%
\pgfusepath{clip}%
\pgfsetbuttcap%
\pgfsetroundjoin%
\definecolor{currentfill}{rgb}{0.949151,0.948435,0.152178}%
\pgfsetfillcolor{currentfill}%
\pgfsetlinewidth{0.000000pt}%
\definecolor{currentstroke}{rgb}{0.000000,0.000000,0.000000}%
\pgfsetstrokecolor{currentstroke}%
\pgfsetdash{}{0pt}%
\pgfpathmoveto{\pgfqpoint{2.792567in}{2.611354in}}%
\pgfpathlineto{\pgfqpoint{2.874143in}{2.655359in}}%
\pgfpathlineto{\pgfqpoint{2.771031in}{2.746435in}}%
\pgfpathlineto{\pgfqpoint{2.792567in}{2.611354in}}%
\pgfpathclose%
\pgfusepath{fill}%
\end{pgfscope}%
\begin{pgfscope}%
\pgfpathrectangle{\pgfqpoint{0.100000in}{0.100000in}}{\pgfqpoint{3.957178in}{3.957178in}}%
\pgfusepath{clip}%
\pgfsetbuttcap%
\pgfsetroundjoin%
\definecolor{currentfill}{rgb}{0.987332,0.646633,0.214648}%
\pgfsetfillcolor{currentfill}%
\pgfsetlinewidth{0.000000pt}%
\definecolor{currentstroke}{rgb}{0.000000,0.000000,0.000000}%
\pgfsetstrokecolor{currentstroke}%
\pgfsetdash{}{0pt}%
\pgfpathmoveto{\pgfqpoint{1.246560in}{2.412228in}}%
\pgfpathlineto{\pgfqpoint{1.230510in}{2.226458in}}%
\pgfpathlineto{\pgfqpoint{1.283566in}{2.420857in}}%
\pgfpathlineto{\pgfqpoint{1.246560in}{2.412228in}}%
\pgfpathclose%
\pgfusepath{fill}%
\end{pgfscope}%
\begin{pgfscope}%
\pgfpathrectangle{\pgfqpoint{0.100000in}{0.100000in}}{\pgfqpoint{3.957178in}{3.957178in}}%
\pgfusepath{clip}%
\pgfsetbuttcap%
\pgfsetroundjoin%
\definecolor{currentfill}{rgb}{0.961681,0.914672,0.150520}%
\pgfsetfillcolor{currentfill}%
\pgfsetlinewidth{0.000000pt}%
\definecolor{currentstroke}{rgb}{0.000000,0.000000,0.000000}%
\pgfsetstrokecolor{currentstroke}%
\pgfsetdash{}{0pt}%
\pgfpathmoveto{\pgfqpoint{2.889972in}{2.631947in}}%
\pgfpathlineto{\pgfqpoint{2.874143in}{2.655359in}}%
\pgfpathlineto{\pgfqpoint{2.792567in}{2.611354in}}%
\pgfpathlineto{\pgfqpoint{2.889972in}{2.631947in}}%
\pgfpathclose%
\pgfusepath{fill}%
\end{pgfscope}%
\begin{pgfscope}%
\pgfpathrectangle{\pgfqpoint{0.100000in}{0.100000in}}{\pgfqpoint{3.957178in}{3.957178in}}%
\pgfusepath{clip}%
\pgfsetbuttcap%
\pgfsetroundjoin%
\definecolor{currentfill}{rgb}{0.605485,0.085854,0.622686}%
\pgfsetfillcolor{currentfill}%
\pgfsetlinewidth{0.000000pt}%
\definecolor{currentstroke}{rgb}{0.000000,0.000000,0.000000}%
\pgfsetstrokecolor{currentstroke}%
\pgfsetdash{}{0pt}%
\pgfpathmoveto{\pgfqpoint{0.747252in}{1.316186in}}%
\pgfpathlineto{\pgfqpoint{1.651923in}{1.660516in}}%
\pgfpathlineto{\pgfqpoint{1.160246in}{1.693711in}}%
\pgfpathlineto{\pgfqpoint{0.747252in}{1.316186in}}%
\pgfpathclose%
\pgfusepath{fill}%
\end{pgfscope}%
\begin{pgfscope}%
\pgfpathrectangle{\pgfqpoint{0.100000in}{0.100000in}}{\pgfqpoint{3.957178in}{3.957178in}}%
\pgfusepath{clip}%
\pgfsetbuttcap%
\pgfsetroundjoin%
\definecolor{currentfill}{rgb}{0.993482,0.765499,0.156891}%
\pgfsetfillcolor{currentfill}%
\pgfsetlinewidth{0.000000pt}%
\definecolor{currentstroke}{rgb}{0.000000,0.000000,0.000000}%
\pgfsetstrokecolor{currentstroke}%
\pgfsetdash{}{0pt}%
\pgfpathmoveto{\pgfqpoint{1.283566in}{2.420857in}}%
\pgfpathlineto{\pgfqpoint{1.430860in}{2.594652in}}%
\pgfpathlineto{\pgfqpoint{1.246560in}{2.412228in}}%
\pgfpathlineto{\pgfqpoint{1.283566in}{2.420857in}}%
\pgfpathclose%
\pgfusepath{fill}%
\end{pgfscope}%
\begin{pgfscope}%
\pgfpathrectangle{\pgfqpoint{0.100000in}{0.100000in}}{\pgfqpoint{3.957178in}{3.957178in}}%
\pgfusepath{clip}%
\pgfsetbuttcap%
\pgfsetroundjoin%
\definecolor{currentfill}{rgb}{0.918109,0.453603,0.339529}%
\pgfsetfillcolor{currentfill}%
\pgfsetlinewidth{0.000000pt}%
\definecolor{currentstroke}{rgb}{0.000000,0.000000,0.000000}%
\pgfsetstrokecolor{currentstroke}%
\pgfsetdash{}{0pt}%
\pgfpathmoveto{\pgfqpoint{3.157732in}{1.874966in}}%
\pgfpathlineto{\pgfqpoint{3.040008in}{2.239939in}}%
\pgfpathlineto{\pgfqpoint{3.022196in}{2.172137in}}%
\pgfpathlineto{\pgfqpoint{3.157732in}{1.874966in}}%
\pgfpathclose%
\pgfusepath{fill}%
\end{pgfscope}%
\begin{pgfscope}%
\pgfpathrectangle{\pgfqpoint{0.100000in}{0.100000in}}{\pgfqpoint{3.957178in}{3.957178in}}%
\pgfusepath{clip}%
\pgfsetbuttcap%
\pgfsetroundjoin%
\definecolor{currentfill}{rgb}{0.941896,0.968590,0.140956}%
\pgfsetfillcolor{currentfill}%
\pgfsetlinewidth{0.000000pt}%
\definecolor{currentstroke}{rgb}{0.000000,0.000000,0.000000}%
\pgfsetstrokecolor{currentstroke}%
\pgfsetdash{}{0pt}%
\pgfpathmoveto{\pgfqpoint{2.792567in}{2.611354in}}%
\pgfpathlineto{\pgfqpoint{2.771031in}{2.746435in}}%
\pgfpathlineto{\pgfqpoint{2.673729in}{2.664530in}}%
\pgfpathlineto{\pgfqpoint{2.792567in}{2.611354in}}%
\pgfpathclose%
\pgfusepath{fill}%
\end{pgfscope}%
\begin{pgfscope}%
\pgfpathrectangle{\pgfqpoint{0.100000in}{0.100000in}}{\pgfqpoint{3.957178in}{3.957178in}}%
\pgfusepath{clip}%
\pgfsetbuttcap%
\pgfsetroundjoin%
\definecolor{currentfill}{rgb}{0.991209,0.790537,0.149377}%
\pgfsetfillcolor{currentfill}%
\pgfsetlinewidth{0.000000pt}%
\definecolor{currentstroke}{rgb}{0.000000,0.000000,0.000000}%
\pgfsetstrokecolor{currentstroke}%
\pgfsetdash{}{0pt}%
\pgfpathmoveto{\pgfqpoint{3.040008in}{2.239939in}}%
\pgfpathlineto{\pgfqpoint{2.889972in}{2.631947in}}%
\pgfpathlineto{\pgfqpoint{2.792567in}{2.611354in}}%
\pgfpathlineto{\pgfqpoint{3.040008in}{2.239939in}}%
\pgfpathclose%
\pgfusepath{fill}%
\end{pgfscope}%
\begin{pgfscope}%
\pgfpathrectangle{\pgfqpoint{0.100000in}{0.100000in}}{\pgfqpoint{3.957178in}{3.957178in}}%
\pgfusepath{clip}%
\pgfsetbuttcap%
\pgfsetroundjoin%
\definecolor{currentfill}{rgb}{0.987332,0.646633,0.214648}%
\pgfsetfillcolor{currentfill}%
\pgfsetlinewidth{0.000000pt}%
\definecolor{currentstroke}{rgb}{0.000000,0.000000,0.000000}%
\pgfsetstrokecolor{currentstroke}%
\pgfsetdash{}{0pt}%
\pgfpathmoveto{\pgfqpoint{1.283566in}{2.420857in}}%
\pgfpathlineto{\pgfqpoint{1.230510in}{2.226458in}}%
\pgfpathlineto{\pgfqpoint{1.291600in}{2.354765in}}%
\pgfpathlineto{\pgfqpoint{1.283566in}{2.420857in}}%
\pgfpathclose%
\pgfusepath{fill}%
\end{pgfscope}%
\begin{pgfscope}%
\pgfpathrectangle{\pgfqpoint{0.100000in}{0.100000in}}{\pgfqpoint{3.957178in}{3.957178in}}%
\pgfusepath{clip}%
\pgfsetbuttcap%
\pgfsetroundjoin%
\definecolor{currentfill}{rgb}{0.872303,0.378774,0.393355}%
\pgfsetfillcolor{currentfill}%
\pgfsetlinewidth{0.000000pt}%
\definecolor{currentstroke}{rgb}{0.000000,0.000000,0.000000}%
\pgfsetstrokecolor{currentstroke}%
\pgfsetdash{}{0pt}%
\pgfpathmoveto{\pgfqpoint{3.157732in}{1.874966in}}%
\pgfpathlineto{\pgfqpoint{3.022196in}{2.172137in}}%
\pgfpathlineto{\pgfqpoint{3.040306in}{1.870487in}}%
\pgfpathlineto{\pgfqpoint{3.157732in}{1.874966in}}%
\pgfpathclose%
\pgfusepath{fill}%
\end{pgfscope}%
\begin{pgfscope}%
\pgfpathrectangle{\pgfqpoint{0.100000in}{0.100000in}}{\pgfqpoint{3.957178in}{3.957178in}}%
\pgfusepath{clip}%
\pgfsetbuttcap%
\pgfsetroundjoin%
\definecolor{currentfill}{rgb}{0.993482,0.765499,0.156891}%
\pgfsetfillcolor{currentfill}%
\pgfsetlinewidth{0.000000pt}%
\definecolor{currentstroke}{rgb}{0.000000,0.000000,0.000000}%
\pgfsetstrokecolor{currentstroke}%
\pgfsetdash{}{0pt}%
\pgfpathmoveto{\pgfqpoint{1.291600in}{2.354765in}}%
\pgfpathlineto{\pgfqpoint{1.430860in}{2.594652in}}%
\pgfpathlineto{\pgfqpoint{1.283566in}{2.420857in}}%
\pgfpathlineto{\pgfqpoint{1.291600in}{2.354765in}}%
\pgfpathclose%
\pgfusepath{fill}%
\end{pgfscope}%
\begin{pgfscope}%
\pgfpathrectangle{\pgfqpoint{0.100000in}{0.100000in}}{\pgfqpoint{3.957178in}{3.957178in}}%
\pgfusepath{clip}%
\pgfsetbuttcap%
\pgfsetroundjoin%
\definecolor{currentfill}{rgb}{0.809052,0.293491,0.458870}%
\pgfsetfillcolor{currentfill}%
\pgfsetlinewidth{0.000000pt}%
\definecolor{currentstroke}{rgb}{0.000000,0.000000,0.000000}%
\pgfsetstrokecolor{currentstroke}%
\pgfsetdash{}{0pt}%
\pgfpathmoveto{\pgfqpoint{3.040306in}{1.870487in}}%
\pgfpathlineto{\pgfqpoint{3.072225in}{1.783143in}}%
\pgfpathlineto{\pgfqpoint{3.157732in}{1.874966in}}%
\pgfpathlineto{\pgfqpoint{3.040306in}{1.870487in}}%
\pgfpathclose%
\pgfusepath{fill}%
\end{pgfscope}%
\begin{pgfscope}%
\pgfpathrectangle{\pgfqpoint{0.100000in}{0.100000in}}{\pgfqpoint{3.957178in}{3.957178in}}%
\pgfusepath{clip}%
\pgfsetbuttcap%
\pgfsetroundjoin%
\definecolor{currentfill}{rgb}{0.990439,0.796859,0.147870}%
\pgfsetfillcolor{currentfill}%
\pgfsetlinewidth{0.000000pt}%
\definecolor{currentstroke}{rgb}{0.000000,0.000000,0.000000}%
\pgfsetstrokecolor{currentstroke}%
\pgfsetdash{}{0pt}%
\pgfpathmoveto{\pgfqpoint{2.196589in}{2.449515in}}%
\pgfpathlineto{\pgfqpoint{1.835085in}{2.543290in}}%
\pgfpathlineto{\pgfqpoint{1.969490in}{2.479538in}}%
\pgfpathlineto{\pgfqpoint{2.196589in}{2.449515in}}%
\pgfpathclose%
\pgfusepath{fill}%
\end{pgfscope}%
\begin{pgfscope}%
\pgfpathrectangle{\pgfqpoint{0.100000in}{0.100000in}}{\pgfqpoint{3.957178in}{3.957178in}}%
\pgfusepath{clip}%
\pgfsetbuttcap%
\pgfsetroundjoin%
\definecolor{currentfill}{rgb}{0.783383,0.262500,0.483918}%
\pgfsetfillcolor{currentfill}%
\pgfsetlinewidth{0.000000pt}%
\definecolor{currentstroke}{rgb}{0.000000,0.000000,0.000000}%
\pgfsetstrokecolor{currentstroke}%
\pgfsetdash{}{0pt}%
\pgfpathmoveto{\pgfqpoint{3.072225in}{1.783143in}}%
\pgfpathlineto{\pgfqpoint{3.038250in}{1.728766in}}%
\pgfpathlineto{\pgfqpoint{3.157732in}{1.874966in}}%
\pgfpathlineto{\pgfqpoint{3.072225in}{1.783143in}}%
\pgfpathclose%
\pgfusepath{fill}%
\end{pgfscope}%
\begin{pgfscope}%
\pgfpathrectangle{\pgfqpoint{0.100000in}{0.100000in}}{\pgfqpoint{3.957178in}{3.957178in}}%
\pgfusepath{clip}%
\pgfsetbuttcap%
\pgfsetroundjoin%
\definecolor{currentfill}{rgb}{0.974443,0.874622,0.144061}%
\pgfsetfillcolor{currentfill}%
\pgfsetlinewidth{0.000000pt}%
\definecolor{currentstroke}{rgb}{0.000000,0.000000,0.000000}%
\pgfsetstrokecolor{currentstroke}%
\pgfsetdash{}{0pt}%
\pgfpathmoveto{\pgfqpoint{1.430860in}{2.594652in}}%
\pgfpathlineto{\pgfqpoint{1.737144in}{2.491232in}}%
\pgfpathlineto{\pgfqpoint{1.790823in}{2.597170in}}%
\pgfpathlineto{\pgfqpoint{1.430860in}{2.594652in}}%
\pgfpathclose%
\pgfusepath{fill}%
\end{pgfscope}%
\begin{pgfscope}%
\pgfpathrectangle{\pgfqpoint{0.100000in}{0.100000in}}{\pgfqpoint{3.957178in}{3.957178in}}%
\pgfusepath{clip}%
\pgfsetbuttcap%
\pgfsetroundjoin%
\definecolor{currentfill}{rgb}{0.968526,0.569700,0.261721}%
\pgfsetfillcolor{currentfill}%
\pgfsetlinewidth{0.000000pt}%
\definecolor{currentstroke}{rgb}{0.000000,0.000000,0.000000}%
\pgfsetstrokecolor{currentstroke}%
\pgfsetdash{}{0pt}%
\pgfpathmoveto{\pgfqpoint{3.022196in}{2.172137in}}%
\pgfpathlineto{\pgfqpoint{3.040008in}{2.239939in}}%
\pgfpathlineto{\pgfqpoint{2.948801in}{2.248544in}}%
\pgfpathlineto{\pgfqpoint{3.022196in}{2.172137in}}%
\pgfpathclose%
\pgfusepath{fill}%
\end{pgfscope}%
\begin{pgfscope}%
\pgfpathrectangle{\pgfqpoint{0.100000in}{0.100000in}}{\pgfqpoint{3.957178in}{3.957178in}}%
\pgfusepath{clip}%
\pgfsetbuttcap%
\pgfsetroundjoin%
\definecolor{currentfill}{rgb}{0.993456,0.698810,0.186041}%
\pgfsetfillcolor{currentfill}%
\pgfsetlinewidth{0.000000pt}%
\definecolor{currentstroke}{rgb}{0.000000,0.000000,0.000000}%
\pgfsetstrokecolor{currentstroke}%
\pgfsetdash{}{0pt}%
\pgfpathmoveto{\pgfqpoint{2.948801in}{2.248544in}}%
\pgfpathlineto{\pgfqpoint{3.040008in}{2.239939in}}%
\pgfpathlineto{\pgfqpoint{2.792567in}{2.611354in}}%
\pgfpathlineto{\pgfqpoint{2.948801in}{2.248544in}}%
\pgfpathclose%
\pgfusepath{fill}%
\end{pgfscope}%
\begin{pgfscope}%
\pgfpathrectangle{\pgfqpoint{0.100000in}{0.100000in}}{\pgfqpoint{3.957178in}{3.957178in}}%
\pgfusepath{clip}%
\pgfsetbuttcap%
\pgfsetroundjoin%
\definecolor{currentfill}{rgb}{0.935630,0.487712,0.315952}%
\pgfsetfillcolor{currentfill}%
\pgfsetlinewidth{0.000000pt}%
\definecolor{currentstroke}{rgb}{0.000000,0.000000,0.000000}%
\pgfsetstrokecolor{currentstroke}%
\pgfsetdash{}{0pt}%
\pgfpathmoveto{\pgfqpoint{1.160246in}{1.693711in}}%
\pgfpathlineto{\pgfqpoint{1.432459in}{2.373573in}}%
\pgfpathlineto{\pgfqpoint{1.230510in}{2.226458in}}%
\pgfpathlineto{\pgfqpoint{1.160246in}{1.693711in}}%
\pgfpathclose%
\pgfusepath{fill}%
\end{pgfscope}%
\begin{pgfscope}%
\pgfpathrectangle{\pgfqpoint{0.100000in}{0.100000in}}{\pgfqpoint{3.957178in}{3.957178in}}%
\pgfusepath{clip}%
\pgfsetbuttcap%
\pgfsetroundjoin%
\definecolor{currentfill}{rgb}{0.989935,0.663787,0.204859}%
\pgfsetfillcolor{currentfill}%
\pgfsetlinewidth{0.000000pt}%
\definecolor{currentstroke}{rgb}{0.000000,0.000000,0.000000}%
\pgfsetstrokecolor{currentstroke}%
\pgfsetdash{}{0pt}%
\pgfpathmoveto{\pgfqpoint{1.230510in}{2.226458in}}%
\pgfpathlineto{\pgfqpoint{1.432459in}{2.373573in}}%
\pgfpathlineto{\pgfqpoint{1.291600in}{2.354765in}}%
\pgfpathlineto{\pgfqpoint{1.230510in}{2.226458in}}%
\pgfpathclose%
\pgfusepath{fill}%
\end{pgfscope}%
\begin{pgfscope}%
\pgfpathrectangle{\pgfqpoint{0.100000in}{0.100000in}}{\pgfqpoint{3.957178in}{3.957178in}}%
\pgfusepath{clip}%
\pgfsetbuttcap%
\pgfsetroundjoin%
\definecolor{currentfill}{rgb}{0.819651,0.306812,0.448306}%
\pgfsetfillcolor{currentfill}%
\pgfsetlinewidth{0.000000pt}%
\definecolor{currentstroke}{rgb}{0.000000,0.000000,0.000000}%
\pgfsetstrokecolor{currentstroke}%
\pgfsetdash{}{0pt}%
\pgfpathmoveto{\pgfqpoint{3.072225in}{1.783143in}}%
\pgfpathlineto{\pgfqpoint{3.040306in}{1.870487in}}%
\pgfpathlineto{\pgfqpoint{3.043803in}{1.852214in}}%
\pgfpathlineto{\pgfqpoint{3.072225in}{1.783143in}}%
\pgfpathclose%
\pgfusepath{fill}%
\end{pgfscope}%
\begin{pgfscope}%
\pgfpathrectangle{\pgfqpoint{0.100000in}{0.100000in}}{\pgfqpoint{3.957178in}{3.957178in}}%
\pgfusepath{clip}%
\pgfsetbuttcap%
\pgfsetroundjoin%
\definecolor{currentfill}{rgb}{0.935630,0.487712,0.315952}%
\pgfsetfillcolor{currentfill}%
\pgfsetlinewidth{0.000000pt}%
\definecolor{currentstroke}{rgb}{0.000000,0.000000,0.000000}%
\pgfsetstrokecolor{currentstroke}%
\pgfsetdash{}{0pt}%
\pgfpathmoveto{\pgfqpoint{3.040306in}{1.870487in}}%
\pgfpathlineto{\pgfqpoint{3.022196in}{2.172137in}}%
\pgfpathlineto{\pgfqpoint{2.948801in}{2.248544in}}%
\pgfpathlineto{\pgfqpoint{3.040306in}{1.870487in}}%
\pgfpathclose%
\pgfusepath{fill}%
\end{pgfscope}%
\begin{pgfscope}%
\pgfpathrectangle{\pgfqpoint{0.100000in}{0.100000in}}{\pgfqpoint{3.957178in}{3.957178in}}%
\pgfusepath{clip}%
\pgfsetbuttcap%
\pgfsetroundjoin%
\definecolor{currentfill}{rgb}{0.794549,0.275770,0.473117}%
\pgfsetfillcolor{currentfill}%
\pgfsetlinewidth{0.000000pt}%
\definecolor{currentstroke}{rgb}{0.000000,0.000000,0.000000}%
\pgfsetstrokecolor{currentstroke}%
\pgfsetdash{}{0pt}%
\pgfpathmoveto{\pgfqpoint{3.038250in}{1.728766in}}%
\pgfpathlineto{\pgfqpoint{3.072225in}{1.783143in}}%
\pgfpathlineto{\pgfqpoint{3.043803in}{1.852214in}}%
\pgfpathlineto{\pgfqpoint{3.038250in}{1.728766in}}%
\pgfpathclose%
\pgfusepath{fill}%
\end{pgfscope}%
\begin{pgfscope}%
\pgfpathrectangle{\pgfqpoint{0.100000in}{0.100000in}}{\pgfqpoint{3.957178in}{3.957178in}}%
\pgfusepath{clip}%
\pgfsetbuttcap%
\pgfsetroundjoin%
\definecolor{currentfill}{rgb}{0.944152,0.961916,0.146861}%
\pgfsetfillcolor{currentfill}%
\pgfsetlinewidth{0.000000pt}%
\definecolor{currentstroke}{rgb}{0.000000,0.000000,0.000000}%
\pgfsetstrokecolor{currentstroke}%
\pgfsetdash{}{0pt}%
\pgfpathmoveto{\pgfqpoint{2.673729in}{2.664530in}}%
\pgfpathlineto{\pgfqpoint{2.701042in}{2.612520in}}%
\pgfpathlineto{\pgfqpoint{2.792567in}{2.611354in}}%
\pgfpathlineto{\pgfqpoint{2.673729in}{2.664530in}}%
\pgfpathclose%
\pgfusepath{fill}%
\end{pgfscope}%
\begin{pgfscope}%
\pgfpathrectangle{\pgfqpoint{0.100000in}{0.100000in}}{\pgfqpoint{3.957178in}{3.957178in}}%
\pgfusepath{clip}%
\pgfsetbuttcap%
\pgfsetroundjoin%
\definecolor{currentfill}{rgb}{0.991897,0.784239,0.151042}%
\pgfsetfillcolor{currentfill}%
\pgfsetlinewidth{0.000000pt}%
\definecolor{currentstroke}{rgb}{0.000000,0.000000,0.000000}%
\pgfsetstrokecolor{currentstroke}%
\pgfsetdash{}{0pt}%
\pgfpathmoveto{\pgfqpoint{1.432459in}{2.373573in}}%
\pgfpathlineto{\pgfqpoint{1.430860in}{2.594652in}}%
\pgfpathlineto{\pgfqpoint{1.291600in}{2.354765in}}%
\pgfpathlineto{\pgfqpoint{1.432459in}{2.373573in}}%
\pgfpathclose%
\pgfusepath{fill}%
\end{pgfscope}%
\begin{pgfscope}%
\pgfpathrectangle{\pgfqpoint{0.100000in}{0.100000in}}{\pgfqpoint{3.957178in}{3.957178in}}%
\pgfusepath{clip}%
\pgfsetbuttcap%
\pgfsetroundjoin%
\definecolor{currentfill}{rgb}{0.979644,0.854866,0.142453}%
\pgfsetfillcolor{currentfill}%
\pgfsetlinewidth{0.000000pt}%
\definecolor{currentstroke}{rgb}{0.000000,0.000000,0.000000}%
\pgfsetstrokecolor{currentstroke}%
\pgfsetdash{}{0pt}%
\pgfpathmoveto{\pgfqpoint{2.196589in}{2.449515in}}%
\pgfpathlineto{\pgfqpoint{2.203139in}{2.533068in}}%
\pgfpathlineto{\pgfqpoint{2.392659in}{2.559380in}}%
\pgfpathlineto{\pgfqpoint{2.196589in}{2.449515in}}%
\pgfpathclose%
\pgfusepath{fill}%
\end{pgfscope}%
\begin{pgfscope}%
\pgfpathrectangle{\pgfqpoint{0.100000in}{0.100000in}}{\pgfqpoint{3.957178in}{3.957178in}}%
\pgfusepath{clip}%
\pgfsetbuttcap%
\pgfsetroundjoin%
\definecolor{currentfill}{rgb}{0.946602,0.955190,0.150328}%
\pgfsetfillcolor{currentfill}%
\pgfsetlinewidth{0.000000pt}%
\definecolor{currentstroke}{rgb}{0.000000,0.000000,0.000000}%
\pgfsetstrokecolor{currentstroke}%
\pgfsetdash{}{0pt}%
\pgfpathmoveto{\pgfqpoint{2.564497in}{2.593307in}}%
\pgfpathlineto{\pgfqpoint{2.673729in}{2.664530in}}%
\pgfpathlineto{\pgfqpoint{2.392659in}{2.559380in}}%
\pgfpathlineto{\pgfqpoint{2.564497in}{2.593307in}}%
\pgfpathclose%
\pgfusepath{fill}%
\end{pgfscope}%
\begin{pgfscope}%
\pgfpathrectangle{\pgfqpoint{0.100000in}{0.100000in}}{\pgfqpoint{3.957178in}{3.957178in}}%
\pgfusepath{clip}%
\pgfsetbuttcap%
\pgfsetroundjoin%
\definecolor{currentfill}{rgb}{0.940015,0.975158,0.131326}%
\pgfsetfillcolor{currentfill}%
\pgfsetlinewidth{0.000000pt}%
\definecolor{currentstroke}{rgb}{0.000000,0.000000,0.000000}%
\pgfsetstrokecolor{currentstroke}%
\pgfsetdash{}{0pt}%
\pgfpathmoveto{\pgfqpoint{2.564497in}{2.593307in}}%
\pgfpathlineto{\pgfqpoint{2.701042in}{2.612520in}}%
\pgfpathlineto{\pgfqpoint{2.673729in}{2.664530in}}%
\pgfpathlineto{\pgfqpoint{2.564497in}{2.593307in}}%
\pgfpathclose%
\pgfusepath{fill}%
\end{pgfscope}%
\begin{pgfscope}%
\pgfpathrectangle{\pgfqpoint{0.100000in}{0.100000in}}{\pgfqpoint{3.957178in}{3.957178in}}%
\pgfusepath{clip}%
\pgfsetbuttcap%
\pgfsetroundjoin%
\definecolor{currentfill}{rgb}{0.875376,0.383347,0.389976}%
\pgfsetfillcolor{currentfill}%
\pgfsetlinewidth{0.000000pt}%
\definecolor{currentstroke}{rgb}{0.000000,0.000000,0.000000}%
\pgfsetstrokecolor{currentstroke}%
\pgfsetdash{}{0pt}%
\pgfpathmoveto{\pgfqpoint{3.043803in}{1.852214in}}%
\pgfpathlineto{\pgfqpoint{3.040306in}{1.870487in}}%
\pgfpathlineto{\pgfqpoint{2.897243in}{2.047433in}}%
\pgfpathlineto{\pgfqpoint{3.043803in}{1.852214in}}%
\pgfpathclose%
\pgfusepath{fill}%
\end{pgfscope}%
\begin{pgfscope}%
\pgfpathrectangle{\pgfqpoint{0.100000in}{0.100000in}}{\pgfqpoint{3.957178in}{3.957178in}}%
\pgfusepath{clip}%
\pgfsetbuttcap%
\pgfsetroundjoin%
\definecolor{currentfill}{rgb}{0.853319,0.351553,0.413734}%
\pgfsetfillcolor{currentfill}%
\pgfsetlinewidth{0.000000pt}%
\definecolor{currentstroke}{rgb}{0.000000,0.000000,0.000000}%
\pgfsetstrokecolor{currentstroke}%
\pgfsetdash{}{0pt}%
\pgfpathmoveto{\pgfqpoint{3.043803in}{1.852214in}}%
\pgfpathlineto{\pgfqpoint{2.897243in}{2.047433in}}%
\pgfpathlineto{\pgfqpoint{3.038250in}{1.728766in}}%
\pgfpathlineto{\pgfqpoint{3.043803in}{1.852214in}}%
\pgfpathclose%
\pgfusepath{fill}%
\end{pgfscope}%
\begin{pgfscope}%
\pgfpathrectangle{\pgfqpoint{0.100000in}{0.100000in}}{\pgfqpoint{3.957178in}{3.957178in}}%
\pgfusepath{clip}%
\pgfsetbuttcap%
\pgfsetroundjoin%
\definecolor{currentfill}{rgb}{0.993033,0.771720,0.154808}%
\pgfsetfillcolor{currentfill}%
\pgfsetlinewidth{0.000000pt}%
\definecolor{currentstroke}{rgb}{0.000000,0.000000,0.000000}%
\pgfsetstrokecolor{currentstroke}%
\pgfsetdash{}{0pt}%
\pgfpathmoveto{\pgfqpoint{2.808781in}{2.368821in}}%
\pgfpathlineto{\pgfqpoint{2.948801in}{2.248544in}}%
\pgfpathlineto{\pgfqpoint{2.792567in}{2.611354in}}%
\pgfpathlineto{\pgfqpoint{2.808781in}{2.368821in}}%
\pgfpathclose%
\pgfusepath{fill}%
\end{pgfscope}%
\begin{pgfscope}%
\pgfpathrectangle{\pgfqpoint{0.100000in}{0.100000in}}{\pgfqpoint{3.957178in}{3.957178in}}%
\pgfusepath{clip}%
\pgfsetbuttcap%
\pgfsetroundjoin%
\definecolor{currentfill}{rgb}{0.933232,0.482780,0.319325}%
\pgfsetfillcolor{currentfill}%
\pgfsetlinewidth{0.000000pt}%
\definecolor{currentstroke}{rgb}{0.000000,0.000000,0.000000}%
\pgfsetstrokecolor{currentstroke}%
\pgfsetdash{}{0pt}%
\pgfpathmoveto{\pgfqpoint{2.948801in}{2.248544in}}%
\pgfpathlineto{\pgfqpoint{2.897243in}{2.047433in}}%
\pgfpathlineto{\pgfqpoint{3.040306in}{1.870487in}}%
\pgfpathlineto{\pgfqpoint{2.948801in}{2.248544in}}%
\pgfpathclose%
\pgfusepath{fill}%
\end{pgfscope}%
\begin{pgfscope}%
\pgfpathrectangle{\pgfqpoint{0.100000in}{0.100000in}}{\pgfqpoint{3.957178in}{3.957178in}}%
\pgfusepath{clip}%
\pgfsetbuttcap%
\pgfsetroundjoin%
\definecolor{currentfill}{rgb}{0.968443,0.894564,0.147014}%
\pgfsetfillcolor{currentfill}%
\pgfsetlinewidth{0.000000pt}%
\definecolor{currentstroke}{rgb}{0.000000,0.000000,0.000000}%
\pgfsetstrokecolor{currentstroke}%
\pgfsetdash{}{0pt}%
\pgfpathmoveto{\pgfqpoint{2.792567in}{2.611354in}}%
\pgfpathlineto{\pgfqpoint{2.701042in}{2.612520in}}%
\pgfpathlineto{\pgfqpoint{2.808781in}{2.368821in}}%
\pgfpathlineto{\pgfqpoint{2.792567in}{2.611354in}}%
\pgfpathclose%
\pgfusepath{fill}%
\end{pgfscope}%
\begin{pgfscope}%
\pgfpathrectangle{\pgfqpoint{0.100000in}{0.100000in}}{\pgfqpoint{3.957178in}{3.957178in}}%
\pgfusepath{clip}%
\pgfsetbuttcap%
\pgfsetroundjoin%
\definecolor{currentfill}{rgb}{0.974443,0.874622,0.144061}%
\pgfsetfillcolor{currentfill}%
\pgfsetlinewidth{0.000000pt}%
\definecolor{currentstroke}{rgb}{0.000000,0.000000,0.000000}%
\pgfsetstrokecolor{currentstroke}%
\pgfsetdash{}{0pt}%
\pgfpathmoveto{\pgfqpoint{1.969490in}{2.479538in}}%
\pgfpathlineto{\pgfqpoint{1.835085in}{2.543290in}}%
\pgfpathlineto{\pgfqpoint{1.737144in}{2.491232in}}%
\pgfpathlineto{\pgfqpoint{1.969490in}{2.479538in}}%
\pgfpathclose%
\pgfusepath{fill}%
\end{pgfscope}%
\begin{pgfscope}%
\pgfpathrectangle{\pgfqpoint{0.100000in}{0.100000in}}{\pgfqpoint{3.957178in}{3.957178in}}%
\pgfusepath{clip}%
\pgfsetbuttcap%
\pgfsetroundjoin%
\definecolor{currentfill}{rgb}{0.859750,0.360588,0.406917}%
\pgfsetfillcolor{currentfill}%
\pgfsetlinewidth{0.000000pt}%
\definecolor{currentstroke}{rgb}{0.000000,0.000000,0.000000}%
\pgfsetstrokecolor{currentstroke}%
\pgfsetdash{}{0pt}%
\pgfpathmoveto{\pgfqpoint{3.038250in}{1.728766in}}%
\pgfpathlineto{\pgfqpoint{2.897243in}{2.047433in}}%
\pgfpathlineto{\pgfqpoint{2.915616in}{1.815810in}}%
\pgfpathlineto{\pgfqpoint{3.038250in}{1.728766in}}%
\pgfpathclose%
\pgfusepath{fill}%
\end{pgfscope}%
\begin{pgfscope}%
\pgfpathrectangle{\pgfqpoint{0.100000in}{0.100000in}}{\pgfqpoint{3.957178in}{3.957178in}}%
\pgfusepath{clip}%
\pgfsetbuttcap%
\pgfsetroundjoin%
\definecolor{currentfill}{rgb}{0.963203,0.553865,0.271909}%
\pgfsetfillcolor{currentfill}%
\pgfsetlinewidth{0.000000pt}%
\definecolor{currentstroke}{rgb}{0.000000,0.000000,0.000000}%
\pgfsetstrokecolor{currentstroke}%
\pgfsetdash{}{0pt}%
\pgfpathmoveto{\pgfqpoint{1.432459in}{2.373573in}}%
\pgfpathlineto{\pgfqpoint{1.160246in}{1.693711in}}%
\pgfpathlineto{\pgfqpoint{1.571221in}{2.335690in}}%
\pgfpathlineto{\pgfqpoint{1.432459in}{2.373573in}}%
\pgfpathclose%
\pgfusepath{fill}%
\end{pgfscope}%
\begin{pgfscope}%
\pgfpathrectangle{\pgfqpoint{0.100000in}{0.100000in}}{\pgfqpoint{3.957178in}{3.957178in}}%
\pgfusepath{clip}%
\pgfsetbuttcap%
\pgfsetroundjoin%
\definecolor{currentfill}{rgb}{0.984199,0.629718,0.224595}%
\pgfsetfillcolor{currentfill}%
\pgfsetlinewidth{0.000000pt}%
\definecolor{currentstroke}{rgb}{0.000000,0.000000,0.000000}%
\pgfsetstrokecolor{currentstroke}%
\pgfsetdash{}{0pt}%
\pgfpathmoveto{\pgfqpoint{2.897243in}{2.047433in}}%
\pgfpathlineto{\pgfqpoint{2.948801in}{2.248544in}}%
\pgfpathlineto{\pgfqpoint{2.808781in}{2.368821in}}%
\pgfpathlineto{\pgfqpoint{2.897243in}{2.047433in}}%
\pgfpathclose%
\pgfusepath{fill}%
\end{pgfscope}%
\begin{pgfscope}%
\pgfpathrectangle{\pgfqpoint{0.100000in}{0.100000in}}{\pgfqpoint{3.957178in}{3.957178in}}%
\pgfusepath{clip}%
\pgfsetbuttcap%
\pgfsetroundjoin%
\definecolor{currentfill}{rgb}{0.974443,0.874622,0.144061}%
\pgfsetfillcolor{currentfill}%
\pgfsetlinewidth{0.000000pt}%
\definecolor{currentstroke}{rgb}{0.000000,0.000000,0.000000}%
\pgfsetstrokecolor{currentstroke}%
\pgfsetdash{}{0pt}%
\pgfpathmoveto{\pgfqpoint{1.430860in}{2.594652in}}%
\pgfpathlineto{\pgfqpoint{1.432459in}{2.373573in}}%
\pgfpathlineto{\pgfqpoint{1.737144in}{2.491232in}}%
\pgfpathlineto{\pgfqpoint{1.430860in}{2.594652in}}%
\pgfpathclose%
\pgfusepath{fill}%
\end{pgfscope}%
\begin{pgfscope}%
\pgfpathrectangle{\pgfqpoint{0.100000in}{0.100000in}}{\pgfqpoint{3.957178in}{3.957178in}}%
\pgfusepath{clip}%
\pgfsetbuttcap%
\pgfsetroundjoin%
\definecolor{currentfill}{rgb}{0.949151,0.948435,0.152178}%
\pgfsetfillcolor{currentfill}%
\pgfsetlinewidth{0.000000pt}%
\definecolor{currentstroke}{rgb}{0.000000,0.000000,0.000000}%
\pgfsetstrokecolor{currentstroke}%
\pgfsetdash{}{0pt}%
\pgfpathmoveto{\pgfqpoint{2.392659in}{2.559380in}}%
\pgfpathlineto{\pgfqpoint{2.440526in}{2.523204in}}%
\pgfpathlineto{\pgfqpoint{2.564497in}{2.593307in}}%
\pgfpathlineto{\pgfqpoint{2.392659in}{2.559380in}}%
\pgfpathclose%
\pgfusepath{fill}%
\end{pgfscope}%
\begin{pgfscope}%
\pgfpathrectangle{\pgfqpoint{0.100000in}{0.100000in}}{\pgfqpoint{3.957178in}{3.957178in}}%
\pgfusepath{clip}%
\pgfsetbuttcap%
\pgfsetroundjoin%
\definecolor{currentfill}{rgb}{0.764193,0.240396,0.502126}%
\pgfsetfillcolor{currentfill}%
\pgfsetlinewidth{0.000000pt}%
\definecolor{currentstroke}{rgb}{0.000000,0.000000,0.000000}%
\pgfsetstrokecolor{currentstroke}%
\pgfsetdash{}{0pt}%
\pgfpathmoveto{\pgfqpoint{3.038250in}{1.728766in}}%
\pgfpathlineto{\pgfqpoint{2.915616in}{1.815810in}}%
\pgfpathlineto{\pgfqpoint{2.838663in}{1.457713in}}%
\pgfpathlineto{\pgfqpoint{3.038250in}{1.728766in}}%
\pgfpathclose%
\pgfusepath{fill}%
\end{pgfscope}%
\begin{pgfscope}%
\pgfpathrectangle{\pgfqpoint{0.100000in}{0.100000in}}{\pgfqpoint{3.957178in}{3.957178in}}%
\pgfusepath{clip}%
\pgfsetbuttcap%
\pgfsetroundjoin%
\definecolor{currentfill}{rgb}{0.972530,0.881250,0.144923}%
\pgfsetfillcolor{currentfill}%
\pgfsetlinewidth{0.000000pt}%
\definecolor{currentstroke}{rgb}{0.000000,0.000000,0.000000}%
\pgfsetstrokecolor{currentstroke}%
\pgfsetdash{}{0pt}%
\pgfpathmoveto{\pgfqpoint{2.203139in}{2.533068in}}%
\pgfpathlineto{\pgfqpoint{2.196589in}{2.449515in}}%
\pgfpathlineto{\pgfqpoint{1.969490in}{2.479538in}}%
\pgfpathlineto{\pgfqpoint{2.203139in}{2.533068in}}%
\pgfpathclose%
\pgfusepath{fill}%
\end{pgfscope}%
\begin{pgfscope}%
\pgfpathrectangle{\pgfqpoint{0.100000in}{0.100000in}}{\pgfqpoint{3.957178in}{3.957178in}}%
\pgfusepath{clip}%
\pgfsetbuttcap%
\pgfsetroundjoin%
\definecolor{currentfill}{rgb}{0.944152,0.961916,0.146861}%
\pgfsetfillcolor{currentfill}%
\pgfsetlinewidth{0.000000pt}%
\definecolor{currentstroke}{rgb}{0.000000,0.000000,0.000000}%
\pgfsetstrokecolor{currentstroke}%
\pgfsetdash{}{0pt}%
\pgfpathmoveto{\pgfqpoint{2.593864in}{2.492515in}}%
\pgfpathlineto{\pgfqpoint{2.701042in}{2.612520in}}%
\pgfpathlineto{\pgfqpoint{2.564497in}{2.593307in}}%
\pgfpathlineto{\pgfqpoint{2.593864in}{2.492515in}}%
\pgfpathclose%
\pgfusepath{fill}%
\end{pgfscope}%
\begin{pgfscope}%
\pgfpathrectangle{\pgfqpoint{0.100000in}{0.100000in}}{\pgfqpoint{3.957178in}{3.957178in}}%
\pgfusepath{clip}%
\pgfsetbuttcap%
\pgfsetroundjoin%
\definecolor{currentfill}{rgb}{0.951726,0.941671,0.152925}%
\pgfsetfillcolor{currentfill}%
\pgfsetlinewidth{0.000000pt}%
\definecolor{currentstroke}{rgb}{0.000000,0.000000,0.000000}%
\pgfsetstrokecolor{currentstroke}%
\pgfsetdash{}{0pt}%
\pgfpathmoveto{\pgfqpoint{2.392659in}{2.559380in}}%
\pgfpathlineto{\pgfqpoint{2.203139in}{2.533068in}}%
\pgfpathlineto{\pgfqpoint{2.229537in}{2.531680in}}%
\pgfpathlineto{\pgfqpoint{2.392659in}{2.559380in}}%
\pgfpathclose%
\pgfusepath{fill}%
\end{pgfscope}%
\begin{pgfscope}%
\pgfpathrectangle{\pgfqpoint{0.100000in}{0.100000in}}{\pgfqpoint{3.957178in}{3.957178in}}%
\pgfusepath{clip}%
\pgfsetbuttcap%
\pgfsetroundjoin%
\definecolor{currentfill}{rgb}{0.968443,0.894564,0.147014}%
\pgfsetfillcolor{currentfill}%
\pgfsetlinewidth{0.000000pt}%
\definecolor{currentstroke}{rgb}{0.000000,0.000000,0.000000}%
\pgfsetstrokecolor{currentstroke}%
\pgfsetdash{}{0pt}%
\pgfpathmoveto{\pgfqpoint{2.701042in}{2.612520in}}%
\pgfpathlineto{\pgfqpoint{2.593864in}{2.492515in}}%
\pgfpathlineto{\pgfqpoint{2.808781in}{2.368821in}}%
\pgfpathlineto{\pgfqpoint{2.701042in}{2.612520in}}%
\pgfpathclose%
\pgfusepath{fill}%
\end{pgfscope}%
\begin{pgfscope}%
\pgfpathrectangle{\pgfqpoint{0.100000in}{0.100000in}}{\pgfqpoint{3.957178in}{3.957178in}}%
\pgfusepath{clip}%
\pgfsetbuttcap%
\pgfsetroundjoin%
\definecolor{currentfill}{rgb}{0.949151,0.948435,0.152178}%
\pgfsetfillcolor{currentfill}%
\pgfsetlinewidth{0.000000pt}%
\definecolor{currentstroke}{rgb}{0.000000,0.000000,0.000000}%
\pgfsetstrokecolor{currentstroke}%
\pgfsetdash{}{0pt}%
\pgfpathmoveto{\pgfqpoint{2.392659in}{2.559380in}}%
\pgfpathlineto{\pgfqpoint{2.229537in}{2.531680in}}%
\pgfpathlineto{\pgfqpoint{2.440526in}{2.523204in}}%
\pgfpathlineto{\pgfqpoint{2.392659in}{2.559380in}}%
\pgfpathclose%
\pgfusepath{fill}%
\end{pgfscope}%
\begin{pgfscope}%
\pgfpathrectangle{\pgfqpoint{0.100000in}{0.100000in}}{\pgfqpoint{3.957178in}{3.957178in}}%
\pgfusepath{clip}%
\pgfsetbuttcap%
\pgfsetroundjoin%
\definecolor{currentfill}{rgb}{0.620919,0.098934,0.614257}%
\pgfsetfillcolor{currentfill}%
\pgfsetlinewidth{0.000000pt}%
\definecolor{currentstroke}{rgb}{0.000000,0.000000,0.000000}%
\pgfsetstrokecolor{currentstroke}%
\pgfsetdash{}{0pt}%
\pgfpathmoveto{\pgfqpoint{2.838663in}{1.457713in}}%
\pgfpathlineto{\pgfqpoint{2.737026in}{1.125597in}}%
\pgfpathlineto{\pgfqpoint{3.038250in}{1.728766in}}%
\pgfpathlineto{\pgfqpoint{2.838663in}{1.457713in}}%
\pgfpathclose%
\pgfusepath{fill}%
\end{pgfscope}%
\begin{pgfscope}%
\pgfpathrectangle{\pgfqpoint{0.100000in}{0.100000in}}{\pgfqpoint{3.957178in}{3.957178in}}%
\pgfusepath{clip}%
\pgfsetbuttcap%
\pgfsetroundjoin%
\definecolor{currentfill}{rgb}{0.955470,0.533093,0.285490}%
\pgfsetfillcolor{currentfill}%
\pgfsetlinewidth{0.000000pt}%
\definecolor{currentstroke}{rgb}{0.000000,0.000000,0.000000}%
\pgfsetstrokecolor{currentstroke}%
\pgfsetdash{}{0pt}%
\pgfpathmoveto{\pgfqpoint{1.679165in}{2.180486in}}%
\pgfpathlineto{\pgfqpoint{1.571221in}{2.335690in}}%
\pgfpathlineto{\pgfqpoint{1.160246in}{1.693711in}}%
\pgfpathlineto{\pgfqpoint{1.679165in}{2.180486in}}%
\pgfpathclose%
\pgfusepath{fill}%
\end{pgfscope}%
\begin{pgfscope}%
\pgfpathrectangle{\pgfqpoint{0.100000in}{0.100000in}}{\pgfqpoint{3.957178in}{3.957178in}}%
\pgfusepath{clip}%
\pgfsetbuttcap%
\pgfsetroundjoin%
\definecolor{currentfill}{rgb}{0.942598,0.502639,0.305816}%
\pgfsetfillcolor{currentfill}%
\pgfsetlinewidth{0.000000pt}%
\definecolor{currentstroke}{rgb}{0.000000,0.000000,0.000000}%
\pgfsetstrokecolor{currentstroke}%
\pgfsetdash{}{0pt}%
\pgfpathmoveto{\pgfqpoint{2.897243in}{2.047433in}}%
\pgfpathlineto{\pgfqpoint{2.709474in}{2.194762in}}%
\pgfpathlineto{\pgfqpoint{2.915616in}{1.815810in}}%
\pgfpathlineto{\pgfqpoint{2.897243in}{2.047433in}}%
\pgfpathclose%
\pgfusepath{fill}%
\end{pgfscope}%
\begin{pgfscope}%
\pgfpathrectangle{\pgfqpoint{0.100000in}{0.100000in}}{\pgfqpoint{3.957178in}{3.957178in}}%
\pgfusepath{clip}%
\pgfsetbuttcap%
\pgfsetroundjoin%
\definecolor{currentfill}{rgb}{0.987332,0.646633,0.214648}%
\pgfsetfillcolor{currentfill}%
\pgfsetlinewidth{0.000000pt}%
\definecolor{currentstroke}{rgb}{0.000000,0.000000,0.000000}%
\pgfsetstrokecolor{currentstroke}%
\pgfsetdash{}{0pt}%
\pgfpathmoveto{\pgfqpoint{2.709474in}{2.194762in}}%
\pgfpathlineto{\pgfqpoint{2.897243in}{2.047433in}}%
\pgfpathlineto{\pgfqpoint{2.808781in}{2.368821in}}%
\pgfpathlineto{\pgfqpoint{2.709474in}{2.194762in}}%
\pgfpathclose%
\pgfusepath{fill}%
\end{pgfscope}%
\begin{pgfscope}%
\pgfpathrectangle{\pgfqpoint{0.100000in}{0.100000in}}{\pgfqpoint{3.957178in}{3.957178in}}%
\pgfusepath{clip}%
\pgfsetbuttcap%
\pgfsetroundjoin%
\definecolor{currentfill}{rgb}{0.941896,0.968590,0.140956}%
\pgfsetfillcolor{currentfill}%
\pgfsetlinewidth{0.000000pt}%
\definecolor{currentstroke}{rgb}{0.000000,0.000000,0.000000}%
\pgfsetstrokecolor{currentstroke}%
\pgfsetdash{}{0pt}%
\pgfpathmoveto{\pgfqpoint{2.564497in}{2.593307in}}%
\pgfpathlineto{\pgfqpoint{2.440526in}{2.523204in}}%
\pgfpathlineto{\pgfqpoint{2.593864in}{2.492515in}}%
\pgfpathlineto{\pgfqpoint{2.564497in}{2.593307in}}%
\pgfpathclose%
\pgfusepath{fill}%
\end{pgfscope}%
\begin{pgfscope}%
\pgfpathrectangle{\pgfqpoint{0.100000in}{0.100000in}}{\pgfqpoint{3.957178in}{3.957178in}}%
\pgfusepath{clip}%
\pgfsetbuttcap%
\pgfsetroundjoin%
\definecolor{currentfill}{rgb}{0.866078,0.369660,0.400126}%
\pgfsetfillcolor{currentfill}%
\pgfsetlinewidth{0.000000pt}%
\definecolor{currentstroke}{rgb}{0.000000,0.000000,0.000000}%
\pgfsetstrokecolor{currentstroke}%
\pgfsetdash{}{0pt}%
\pgfpathmoveto{\pgfqpoint{2.838663in}{1.457713in}}%
\pgfpathlineto{\pgfqpoint{2.915616in}{1.815810in}}%
\pgfpathlineto{\pgfqpoint{2.709474in}{2.194762in}}%
\pgfpathlineto{\pgfqpoint{2.838663in}{1.457713in}}%
\pgfpathclose%
\pgfusepath{fill}%
\end{pgfscope}%
\begin{pgfscope}%
\pgfpathrectangle{\pgfqpoint{0.100000in}{0.100000in}}{\pgfqpoint{3.957178in}{3.957178in}}%
\pgfusepath{clip}%
\pgfsetbuttcap%
\pgfsetroundjoin%
\definecolor{currentfill}{rgb}{0.984031,0.835315,0.142528}%
\pgfsetfillcolor{currentfill}%
\pgfsetlinewidth{0.000000pt}%
\definecolor{currentstroke}{rgb}{0.000000,0.000000,0.000000}%
\pgfsetstrokecolor{currentstroke}%
\pgfsetdash{}{0pt}%
\pgfpathmoveto{\pgfqpoint{1.737144in}{2.491232in}}%
\pgfpathlineto{\pgfqpoint{1.432459in}{2.373573in}}%
\pgfpathlineto{\pgfqpoint{1.571221in}{2.335690in}}%
\pgfpathlineto{\pgfqpoint{1.737144in}{2.491232in}}%
\pgfpathclose%
\pgfusepath{fill}%
\end{pgfscope}%
\begin{pgfscope}%
\pgfpathrectangle{\pgfqpoint{0.100000in}{0.100000in}}{\pgfqpoint{3.957178in}{3.957178in}}%
\pgfusepath{clip}%
\pgfsetbuttcap%
\pgfsetroundjoin%
\definecolor{currentfill}{rgb}{0.875376,0.383347,0.389976}%
\pgfsetfillcolor{currentfill}%
\pgfsetlinewidth{0.000000pt}%
\definecolor{currentstroke}{rgb}{0.000000,0.000000,0.000000}%
\pgfsetstrokecolor{currentstroke}%
\pgfsetdash{}{0pt}%
\pgfpathmoveto{\pgfqpoint{1.679165in}{2.180486in}}%
\pgfpathlineto{\pgfqpoint{1.160246in}{1.693711in}}%
\pgfpathlineto{\pgfqpoint{1.651923in}{1.660516in}}%
\pgfpathlineto{\pgfqpoint{1.679165in}{2.180486in}}%
\pgfpathclose%
\pgfusepath{fill}%
\end{pgfscope}%
\begin{pgfscope}%
\pgfpathrectangle{\pgfqpoint{0.100000in}{0.100000in}}{\pgfqpoint{3.957178in}{3.957178in}}%
\pgfusepath{clip}%
\pgfsetbuttcap%
\pgfsetroundjoin%
\definecolor{currentfill}{rgb}{0.990439,0.796859,0.147870}%
\pgfsetfillcolor{currentfill}%
\pgfsetlinewidth{0.000000pt}%
\definecolor{currentstroke}{rgb}{0.000000,0.000000,0.000000}%
\pgfsetstrokecolor{currentstroke}%
\pgfsetdash{}{0pt}%
\pgfpathmoveto{\pgfqpoint{2.808781in}{2.368821in}}%
\pgfpathlineto{\pgfqpoint{2.593864in}{2.492515in}}%
\pgfpathlineto{\pgfqpoint{2.709474in}{2.194762in}}%
\pgfpathlineto{\pgfqpoint{2.808781in}{2.368821in}}%
\pgfpathclose%
\pgfusepath{fill}%
\end{pgfscope}%
\begin{pgfscope}%
\pgfpathrectangle{\pgfqpoint{0.100000in}{0.100000in}}{\pgfqpoint{3.957178in}{3.957178in}}%
\pgfusepath{clip}%
\pgfsetbuttcap%
\pgfsetroundjoin%
\definecolor{currentfill}{rgb}{0.573632,0.060028,0.637349}%
\pgfsetfillcolor{currentfill}%
\pgfsetlinewidth{0.000000pt}%
\definecolor{currentstroke}{rgb}{0.000000,0.000000,0.000000}%
\pgfsetstrokecolor{currentstroke}%
\pgfsetdash{}{0pt}%
\pgfpathmoveto{\pgfqpoint{2.762632in}{1.426075in}}%
\pgfpathlineto{\pgfqpoint{2.737026in}{1.125597in}}%
\pgfpathlineto{\pgfqpoint{2.838663in}{1.457713in}}%
\pgfpathlineto{\pgfqpoint{2.762632in}{1.426075in}}%
\pgfpathclose%
\pgfusepath{fill}%
\end{pgfscope}%
\begin{pgfscope}%
\pgfpathrectangle{\pgfqpoint{0.100000in}{0.100000in}}{\pgfqpoint{3.957178in}{3.957178in}}%
\pgfusepath{clip}%
\pgfsetbuttcap%
\pgfsetroundjoin%
\definecolor{currentfill}{rgb}{0.809052,0.293491,0.458870}%
\pgfsetfillcolor{currentfill}%
\pgfsetlinewidth{0.000000pt}%
\definecolor{currentstroke}{rgb}{0.000000,0.000000,0.000000}%
\pgfsetstrokecolor{currentstroke}%
\pgfsetdash{}{0pt}%
\pgfpathmoveto{\pgfqpoint{2.838663in}{1.457713in}}%
\pgfpathlineto{\pgfqpoint{2.709474in}{2.194762in}}%
\pgfpathlineto{\pgfqpoint{2.762632in}{1.426075in}}%
\pgfpathlineto{\pgfqpoint{2.838663in}{1.457713in}}%
\pgfpathclose%
\pgfusepath{fill}%
\end{pgfscope}%
\begin{pgfscope}%
\pgfpathrectangle{\pgfqpoint{0.100000in}{0.100000in}}{\pgfqpoint{3.957178in}{3.957178in}}%
\pgfusepath{clip}%
\pgfsetbuttcap%
\pgfsetroundjoin%
\definecolor{currentfill}{rgb}{0.941896,0.968590,0.140956}%
\pgfsetfillcolor{currentfill}%
\pgfsetlinewidth{0.000000pt}%
\definecolor{currentstroke}{rgb}{0.000000,0.000000,0.000000}%
\pgfsetstrokecolor{currentstroke}%
\pgfsetdash{}{0pt}%
\pgfpathmoveto{\pgfqpoint{1.969490in}{2.479538in}}%
\pgfpathlineto{\pgfqpoint{2.229537in}{2.531680in}}%
\pgfpathlineto{\pgfqpoint{2.203139in}{2.533068in}}%
\pgfpathlineto{\pgfqpoint{1.969490in}{2.479538in}}%
\pgfpathclose%
\pgfusepath{fill}%
\end{pgfscope}%
\begin{pgfscope}%
\pgfpathrectangle{\pgfqpoint{0.100000in}{0.100000in}}{\pgfqpoint{3.957178in}{3.957178in}}%
\pgfusepath{clip}%
\pgfsetbuttcap%
\pgfsetroundjoin%
\definecolor{currentfill}{rgb}{0.977995,0.861432,0.142808}%
\pgfsetfillcolor{currentfill}%
\pgfsetlinewidth{0.000000pt}%
\definecolor{currentstroke}{rgb}{0.000000,0.000000,0.000000}%
\pgfsetstrokecolor{currentstroke}%
\pgfsetdash{}{0pt}%
\pgfpathmoveto{\pgfqpoint{2.593864in}{2.492515in}}%
\pgfpathlineto{\pgfqpoint{2.440526in}{2.523204in}}%
\pgfpathlineto{\pgfqpoint{2.709474in}{2.194762in}}%
\pgfpathlineto{\pgfqpoint{2.593864in}{2.492515in}}%
\pgfpathclose%
\pgfusepath{fill}%
\end{pgfscope}%
\begin{pgfscope}%
\pgfpathrectangle{\pgfqpoint{0.100000in}{0.100000in}}{\pgfqpoint{3.957178in}{3.957178in}}%
\pgfusepath{clip}%
\pgfsetbuttcap%
\pgfsetroundjoin%
\definecolor{currentfill}{rgb}{0.940015,0.975158,0.131326}%
\pgfsetfillcolor{currentfill}%
\pgfsetlinewidth{0.000000pt}%
\definecolor{currentstroke}{rgb}{0.000000,0.000000,0.000000}%
\pgfsetstrokecolor{currentstroke}%
\pgfsetdash{}{0pt}%
\pgfpathmoveto{\pgfqpoint{2.440526in}{2.523204in}}%
\pgfpathlineto{\pgfqpoint{2.229537in}{2.531680in}}%
\pgfpathlineto{\pgfqpoint{2.333197in}{2.501845in}}%
\pgfpathlineto{\pgfqpoint{2.440526in}{2.523204in}}%
\pgfpathclose%
\pgfusepath{fill}%
\end{pgfscope}%
\begin{pgfscope}%
\pgfpathrectangle{\pgfqpoint{0.100000in}{0.100000in}}{\pgfqpoint{3.957178in}{3.957178in}}%
\pgfusepath{clip}%
\pgfsetbuttcap%
\pgfsetroundjoin%
\definecolor{currentfill}{rgb}{0.318856,0.007576,0.637640}%
\pgfsetfillcolor{currentfill}%
\pgfsetlinewidth{0.000000pt}%
\definecolor{currentstroke}{rgb}{0.000000,0.000000,0.000000}%
\pgfsetstrokecolor{currentstroke}%
\pgfsetdash{}{0pt}%
\pgfpathmoveto{\pgfqpoint{2.737026in}{1.125597in}}%
\pgfpathlineto{\pgfqpoint{2.595041in}{1.089175in}}%
\pgfpathlineto{\pgfqpoint{2.602110in}{0.852544in}}%
\pgfpathlineto{\pgfqpoint{2.737026in}{1.125597in}}%
\pgfpathclose%
\pgfusepath{fill}%
\end{pgfscope}%
\begin{pgfscope}%
\pgfpathrectangle{\pgfqpoint{0.100000in}{0.100000in}}{\pgfqpoint{3.957178in}{3.957178in}}%
\pgfusepath{clip}%
\pgfsetbuttcap%
\pgfsetroundjoin%
\definecolor{currentfill}{rgb}{0.620919,0.098934,0.614257}%
\pgfsetfillcolor{currentfill}%
\pgfsetlinewidth{0.000000pt}%
\definecolor{currentstroke}{rgb}{0.000000,0.000000,0.000000}%
\pgfsetstrokecolor{currentstroke}%
\pgfsetdash{}{0pt}%
\pgfpathmoveto{\pgfqpoint{2.558352in}{1.557738in}}%
\pgfpathlineto{\pgfqpoint{2.737026in}{1.125597in}}%
\pgfpathlineto{\pgfqpoint{2.762632in}{1.426075in}}%
\pgfpathlineto{\pgfqpoint{2.558352in}{1.557738in}}%
\pgfpathclose%
\pgfusepath{fill}%
\end{pgfscope}%
\begin{pgfscope}%
\pgfpathrectangle{\pgfqpoint{0.100000in}{0.100000in}}{\pgfqpoint{3.957178in}{3.957178in}}%
\pgfusepath{clip}%
\pgfsetbuttcap%
\pgfsetroundjoin%
\definecolor{currentfill}{rgb}{0.941896,0.968590,0.140956}%
\pgfsetfillcolor{currentfill}%
\pgfsetlinewidth{0.000000pt}%
\definecolor{currentstroke}{rgb}{0.000000,0.000000,0.000000}%
\pgfsetstrokecolor{currentstroke}%
\pgfsetdash{}{0pt}%
\pgfpathmoveto{\pgfqpoint{2.328755in}{2.465158in}}%
\pgfpathlineto{\pgfqpoint{2.333197in}{2.501845in}}%
\pgfpathlineto{\pgfqpoint{2.229537in}{2.531680in}}%
\pgfpathlineto{\pgfqpoint{2.328755in}{2.465158in}}%
\pgfpathclose%
\pgfusepath{fill}%
\end{pgfscope}%
\begin{pgfscope}%
\pgfpathrectangle{\pgfqpoint{0.100000in}{0.100000in}}{\pgfqpoint{3.957178in}{3.957178in}}%
\pgfusepath{clip}%
\pgfsetbuttcap%
\pgfsetroundjoin%
\definecolor{currentfill}{rgb}{0.984031,0.835315,0.142528}%
\pgfsetfillcolor{currentfill}%
\pgfsetlinewidth{0.000000pt}%
\definecolor{currentstroke}{rgb}{0.000000,0.000000,0.000000}%
\pgfsetstrokecolor{currentstroke}%
\pgfsetdash{}{0pt}%
\pgfpathmoveto{\pgfqpoint{1.737144in}{2.491232in}}%
\pgfpathlineto{\pgfqpoint{1.571221in}{2.335690in}}%
\pgfpathlineto{\pgfqpoint{1.707645in}{2.261003in}}%
\pgfpathlineto{\pgfqpoint{1.737144in}{2.491232in}}%
\pgfpathclose%
\pgfusepath{fill}%
\end{pgfscope}%
\begin{pgfscope}%
\pgfpathrectangle{\pgfqpoint{0.100000in}{0.100000in}}{\pgfqpoint{3.957178in}{3.957178in}}%
\pgfusepath{clip}%
\pgfsetbuttcap%
\pgfsetroundjoin%
\definecolor{currentfill}{rgb}{0.976265,0.868016,0.143351}%
\pgfsetfillcolor{currentfill}%
\pgfsetlinewidth{0.000000pt}%
\definecolor{currentstroke}{rgb}{0.000000,0.000000,0.000000}%
\pgfsetstrokecolor{currentstroke}%
\pgfsetdash{}{0pt}%
\pgfpathmoveto{\pgfqpoint{2.709474in}{2.194762in}}%
\pgfpathlineto{\pgfqpoint{2.440526in}{2.523204in}}%
\pgfpathlineto{\pgfqpoint{2.328755in}{2.465158in}}%
\pgfpathlineto{\pgfqpoint{2.709474in}{2.194762in}}%
\pgfpathclose%
\pgfusepath{fill}%
\end{pgfscope}%
\begin{pgfscope}%
\pgfpathrectangle{\pgfqpoint{0.100000in}{0.100000in}}{\pgfqpoint{3.957178in}{3.957178in}}%
\pgfusepath{clip}%
\pgfsetbuttcap%
\pgfsetroundjoin%
\definecolor{currentfill}{rgb}{0.941896,0.968590,0.140956}%
\pgfsetfillcolor{currentfill}%
\pgfsetlinewidth{0.000000pt}%
\definecolor{currentstroke}{rgb}{0.000000,0.000000,0.000000}%
\pgfsetstrokecolor{currentstroke}%
\pgfsetdash{}{0pt}%
\pgfpathmoveto{\pgfqpoint{2.328755in}{2.465158in}}%
\pgfpathlineto{\pgfqpoint{2.229537in}{2.531680in}}%
\pgfpathlineto{\pgfqpoint{1.969490in}{2.479538in}}%
\pgfpathlineto{\pgfqpoint{2.328755in}{2.465158in}}%
\pgfpathclose%
\pgfusepath{fill}%
\end{pgfscope}%
\begin{pgfscope}%
\pgfpathrectangle{\pgfqpoint{0.100000in}{0.100000in}}{\pgfqpoint{3.957178in}{3.957178in}}%
\pgfusepath{clip}%
\pgfsetbuttcap%
\pgfsetroundjoin%
\definecolor{currentfill}{rgb}{0.940015,0.975158,0.131326}%
\pgfsetfillcolor{currentfill}%
\pgfsetlinewidth{0.000000pt}%
\definecolor{currentstroke}{rgb}{0.000000,0.000000,0.000000}%
\pgfsetstrokecolor{currentstroke}%
\pgfsetdash{}{0pt}%
\pgfpathmoveto{\pgfqpoint{2.440526in}{2.523204in}}%
\pgfpathlineto{\pgfqpoint{2.333197in}{2.501845in}}%
\pgfpathlineto{\pgfqpoint{2.328755in}{2.465158in}}%
\pgfpathlineto{\pgfqpoint{2.440526in}{2.523204in}}%
\pgfpathclose%
\pgfusepath{fill}%
\end{pgfscope}%
\begin{pgfscope}%
\pgfpathrectangle{\pgfqpoint{0.100000in}{0.100000in}}{\pgfqpoint{3.957178in}{3.957178in}}%
\pgfusepath{clip}%
\pgfsetbuttcap%
\pgfsetroundjoin%
\definecolor{currentfill}{rgb}{0.869203,0.374212,0.396738}%
\pgfsetfillcolor{currentfill}%
\pgfsetlinewidth{0.000000pt}%
\definecolor{currentstroke}{rgb}{0.000000,0.000000,0.000000}%
\pgfsetstrokecolor{currentstroke}%
\pgfsetdash{}{0pt}%
\pgfpathmoveto{\pgfqpoint{2.525070in}{1.721761in}}%
\pgfpathlineto{\pgfqpoint{2.762632in}{1.426075in}}%
\pgfpathlineto{\pgfqpoint{2.709474in}{2.194762in}}%
\pgfpathlineto{\pgfqpoint{2.525070in}{1.721761in}}%
\pgfpathclose%
\pgfusepath{fill}%
\end{pgfscope}%
\begin{pgfscope}%
\pgfpathrectangle{\pgfqpoint{0.100000in}{0.100000in}}{\pgfqpoint{3.957178in}{3.957178in}}%
\pgfusepath{clip}%
\pgfsetbuttcap%
\pgfsetroundjoin%
\definecolor{currentfill}{rgb}{0.540570,0.034950,0.648640}%
\pgfsetfillcolor{currentfill}%
\pgfsetlinewidth{0.000000pt}%
\definecolor{currentstroke}{rgb}{0.000000,0.000000,0.000000}%
\pgfsetstrokecolor{currentstroke}%
\pgfsetdash{}{0pt}%
\pgfpathmoveto{\pgfqpoint{2.595041in}{1.089175in}}%
\pgfpathlineto{\pgfqpoint{2.737026in}{1.125597in}}%
\pgfpathlineto{\pgfqpoint{2.558352in}{1.557738in}}%
\pgfpathlineto{\pgfqpoint{2.595041in}{1.089175in}}%
\pgfpathclose%
\pgfusepath{fill}%
\end{pgfscope}%
\begin{pgfscope}%
\pgfpathrectangle{\pgfqpoint{0.100000in}{0.100000in}}{\pgfqpoint{3.957178in}{3.957178in}}%
\pgfusepath{clip}%
\pgfsetbuttcap%
\pgfsetroundjoin%
\definecolor{currentfill}{rgb}{0.946602,0.955190,0.150328}%
\pgfsetfillcolor{currentfill}%
\pgfsetlinewidth{0.000000pt}%
\definecolor{currentstroke}{rgb}{0.000000,0.000000,0.000000}%
\pgfsetstrokecolor{currentstroke}%
\pgfsetdash{}{0pt}%
\pgfpathmoveto{\pgfqpoint{1.737144in}{2.491232in}}%
\pgfpathlineto{\pgfqpoint{1.882422in}{2.450114in}}%
\pgfpathlineto{\pgfqpoint{1.969490in}{2.479538in}}%
\pgfpathlineto{\pgfqpoint{1.737144in}{2.491232in}}%
\pgfpathclose%
\pgfusepath{fill}%
\end{pgfscope}%
\begin{pgfscope}%
\pgfpathrectangle{\pgfqpoint{0.100000in}{0.100000in}}{\pgfqpoint{3.957178in}{3.957178in}}%
\pgfusepath{clip}%
\pgfsetbuttcap%
\pgfsetroundjoin%
\definecolor{currentfill}{rgb}{0.312543,0.008239,0.635700}%
\pgfsetfillcolor{currentfill}%
\pgfsetlinewidth{0.000000pt}%
\definecolor{currentstroke}{rgb}{0.000000,0.000000,0.000000}%
\pgfsetstrokecolor{currentstroke}%
\pgfsetdash{}{0pt}%
\pgfpathmoveto{\pgfqpoint{2.595041in}{1.089175in}}%
\pgfpathlineto{\pgfqpoint{2.480551in}{1.051058in}}%
\pgfpathlineto{\pgfqpoint{2.602110in}{0.852544in}}%
\pgfpathlineto{\pgfqpoint{2.595041in}{1.089175in}}%
\pgfpathclose%
\pgfusepath{fill}%
\end{pgfscope}%
\begin{pgfscope}%
\pgfpathrectangle{\pgfqpoint{0.100000in}{0.100000in}}{\pgfqpoint{3.957178in}{3.957178in}}%
\pgfusepath{clip}%
\pgfsetbuttcap%
\pgfsetroundjoin%
\definecolor{currentfill}{rgb}{0.994141,0.753137,0.161404}%
\pgfsetfillcolor{currentfill}%
\pgfsetlinewidth{0.000000pt}%
\definecolor{currentstroke}{rgb}{0.000000,0.000000,0.000000}%
\pgfsetstrokecolor{currentstroke}%
\pgfsetdash{}{0pt}%
\pgfpathmoveto{\pgfqpoint{1.679165in}{2.180486in}}%
\pgfpathlineto{\pgfqpoint{1.707645in}{2.261003in}}%
\pgfpathlineto{\pgfqpoint{1.571221in}{2.335690in}}%
\pgfpathlineto{\pgfqpoint{1.679165in}{2.180486in}}%
\pgfpathclose%
\pgfusepath{fill}%
\end{pgfscope}%
\begin{pgfscope}%
\pgfpathrectangle{\pgfqpoint{0.100000in}{0.100000in}}{\pgfqpoint{3.957178in}{3.957178in}}%
\pgfusepath{clip}%
\pgfsetbuttcap%
\pgfsetroundjoin%
\definecolor{currentfill}{rgb}{0.970533,0.887896,0.145919}%
\pgfsetfillcolor{currentfill}%
\pgfsetlinewidth{0.000000pt}%
\definecolor{currentstroke}{rgb}{0.000000,0.000000,0.000000}%
\pgfsetstrokecolor{currentstroke}%
\pgfsetdash{}{0pt}%
\pgfpathmoveto{\pgfqpoint{1.707645in}{2.261003in}}%
\pgfpathlineto{\pgfqpoint{1.882422in}{2.450114in}}%
\pgfpathlineto{\pgfqpoint{1.737144in}{2.491232in}}%
\pgfpathlineto{\pgfqpoint{1.707645in}{2.261003in}}%
\pgfpathclose%
\pgfusepath{fill}%
\end{pgfscope}%
\begin{pgfscope}%
\pgfpathrectangle{\pgfqpoint{0.100000in}{0.100000in}}{\pgfqpoint{3.957178in}{3.957178in}}%
\pgfusepath{clip}%
\pgfsetbuttcap%
\pgfsetroundjoin%
\definecolor{currentfill}{rgb}{0.768090,0.244817,0.498465}%
\pgfsetfillcolor{currentfill}%
\pgfsetlinewidth{0.000000pt}%
\definecolor{currentstroke}{rgb}{0.000000,0.000000,0.000000}%
\pgfsetstrokecolor{currentstroke}%
\pgfsetdash{}{0pt}%
\pgfpathmoveto{\pgfqpoint{2.558352in}{1.557738in}}%
\pgfpathlineto{\pgfqpoint{2.762632in}{1.426075in}}%
\pgfpathlineto{\pgfqpoint{2.525070in}{1.721761in}}%
\pgfpathlineto{\pgfqpoint{2.558352in}{1.557738in}}%
\pgfpathclose%
\pgfusepath{fill}%
\end{pgfscope}%
\begin{pgfscope}%
\pgfpathrectangle{\pgfqpoint{0.100000in}{0.100000in}}{\pgfqpoint{3.957178in}{3.957178in}}%
\pgfusepath{clip}%
\pgfsetbuttcap%
\pgfsetroundjoin%
\definecolor{currentfill}{rgb}{0.529306,0.027747,0.651586}%
\pgfsetfillcolor{currentfill}%
\pgfsetlinewidth{0.000000pt}%
\definecolor{currentstroke}{rgb}{0.000000,0.000000,0.000000}%
\pgfsetstrokecolor{currentstroke}%
\pgfsetdash{}{0pt}%
\pgfpathmoveto{\pgfqpoint{2.480551in}{1.051058in}}%
\pgfpathlineto{\pgfqpoint{2.595041in}{1.089175in}}%
\pgfpathlineto{\pgfqpoint{2.558352in}{1.557738in}}%
\pgfpathlineto{\pgfqpoint{2.480551in}{1.051058in}}%
\pgfpathclose%
\pgfusepath{fill}%
\end{pgfscope}%
\begin{pgfscope}%
\pgfpathrectangle{\pgfqpoint{0.100000in}{0.100000in}}{\pgfqpoint{3.957178in}{3.957178in}}%
\pgfusepath{clip}%
\pgfsetbuttcap%
\pgfsetroundjoin%
\definecolor{currentfill}{rgb}{0.318856,0.007576,0.637640}%
\pgfsetfillcolor{currentfill}%
\pgfsetlinewidth{0.000000pt}%
\definecolor{currentstroke}{rgb}{0.000000,0.000000,0.000000}%
\pgfsetstrokecolor{currentstroke}%
\pgfsetdash{}{0pt}%
\pgfpathmoveto{\pgfqpoint{2.480551in}{1.051058in}}%
\pgfpathlineto{\pgfqpoint{2.336674in}{1.076454in}}%
\pgfpathlineto{\pgfqpoint{2.602110in}{0.852544in}}%
\pgfpathlineto{\pgfqpoint{2.480551in}{1.051058in}}%
\pgfpathclose%
\pgfusepath{fill}%
\end{pgfscope}%
\begin{pgfscope}%
\pgfpathrectangle{\pgfqpoint{0.100000in}{0.100000in}}{\pgfqpoint{3.957178in}{3.957178in}}%
\pgfusepath{clip}%
\pgfsetbuttcap%
\pgfsetroundjoin%
\definecolor{currentfill}{rgb}{0.990681,0.669558,0.201642}%
\pgfsetfillcolor{currentfill}%
\pgfsetlinewidth{0.000000pt}%
\definecolor{currentstroke}{rgb}{0.000000,0.000000,0.000000}%
\pgfsetstrokecolor{currentstroke}%
\pgfsetdash{}{0pt}%
\pgfpathmoveto{\pgfqpoint{2.328755in}{2.465158in}}%
\pgfpathlineto{\pgfqpoint{2.478698in}{1.804673in}}%
\pgfpathlineto{\pgfqpoint{2.709474in}{2.194762in}}%
\pgfpathlineto{\pgfqpoint{2.328755in}{2.465158in}}%
\pgfpathclose%
\pgfusepath{fill}%
\end{pgfscope}%
\begin{pgfscope}%
\pgfpathrectangle{\pgfqpoint{0.100000in}{0.100000in}}{\pgfqpoint{3.957178in}{3.957178in}}%
\pgfusepath{clip}%
\pgfsetbuttcap%
\pgfsetroundjoin%
\definecolor{currentfill}{rgb}{0.955470,0.533093,0.285490}%
\pgfsetfillcolor{currentfill}%
\pgfsetlinewidth{0.000000pt}%
\definecolor{currentstroke}{rgb}{0.000000,0.000000,0.000000}%
\pgfsetstrokecolor{currentstroke}%
\pgfsetdash{}{0pt}%
\pgfpathmoveto{\pgfqpoint{1.691922in}{2.103469in}}%
\pgfpathlineto{\pgfqpoint{1.679165in}{2.180486in}}%
\pgfpathlineto{\pgfqpoint{1.651923in}{1.660516in}}%
\pgfpathlineto{\pgfqpoint{1.691922in}{2.103469in}}%
\pgfpathclose%
\pgfusepath{fill}%
\end{pgfscope}%
\begin{pgfscope}%
\pgfpathrectangle{\pgfqpoint{0.100000in}{0.100000in}}{\pgfqpoint{3.957178in}{3.957178in}}%
\pgfusepath{clip}%
\pgfsetbuttcap%
\pgfsetroundjoin%
\definecolor{currentfill}{rgb}{0.951726,0.941671,0.152925}%
\pgfsetfillcolor{currentfill}%
\pgfsetlinewidth{0.000000pt}%
\definecolor{currentstroke}{rgb}{0.000000,0.000000,0.000000}%
\pgfsetstrokecolor{currentstroke}%
\pgfsetdash{}{0pt}%
\pgfpathmoveto{\pgfqpoint{1.882422in}{2.450114in}}%
\pgfpathlineto{\pgfqpoint{1.898991in}{2.394226in}}%
\pgfpathlineto{\pgfqpoint{1.969490in}{2.479538in}}%
\pgfpathlineto{\pgfqpoint{1.882422in}{2.450114in}}%
\pgfpathclose%
\pgfusepath{fill}%
\end{pgfscope}%
\begin{pgfscope}%
\pgfpathrectangle{\pgfqpoint{0.100000in}{0.100000in}}{\pgfqpoint{3.957178in}{3.957178in}}%
\pgfusepath{clip}%
\pgfsetbuttcap%
\pgfsetroundjoin%
\definecolor{currentfill}{rgb}{0.901807,0.425087,0.359688}%
\pgfsetfillcolor{currentfill}%
\pgfsetlinewidth{0.000000pt}%
\definecolor{currentstroke}{rgb}{0.000000,0.000000,0.000000}%
\pgfsetstrokecolor{currentstroke}%
\pgfsetdash{}{0pt}%
\pgfpathmoveto{\pgfqpoint{1.651923in}{1.660516in}}%
\pgfpathlineto{\pgfqpoint{1.682618in}{1.716885in}}%
\pgfpathlineto{\pgfqpoint{1.691922in}{2.103469in}}%
\pgfpathlineto{\pgfqpoint{1.651923in}{1.660516in}}%
\pgfpathclose%
\pgfusepath{fill}%
\end{pgfscope}%
\begin{pgfscope}%
\pgfpathrectangle{\pgfqpoint{0.100000in}{0.100000in}}{\pgfqpoint{3.957178in}{3.957178in}}%
\pgfusepath{clip}%
\pgfsetbuttcap%
\pgfsetroundjoin%
\definecolor{currentfill}{rgb}{0.933232,0.482780,0.319325}%
\pgfsetfillcolor{currentfill}%
\pgfsetlinewidth{0.000000pt}%
\definecolor{currentstroke}{rgb}{0.000000,0.000000,0.000000}%
\pgfsetstrokecolor{currentstroke}%
\pgfsetdash{}{0pt}%
\pgfpathmoveto{\pgfqpoint{2.709474in}{2.194762in}}%
\pgfpathlineto{\pgfqpoint{2.478698in}{1.804673in}}%
\pgfpathlineto{\pgfqpoint{2.525070in}{1.721761in}}%
\pgfpathlineto{\pgfqpoint{2.709474in}{2.194762in}}%
\pgfpathclose%
\pgfusepath{fill}%
\end{pgfscope}%
\begin{pgfscope}%
\pgfpathrectangle{\pgfqpoint{0.100000in}{0.100000in}}{\pgfqpoint{3.957178in}{3.957178in}}%
\pgfusepath{clip}%
\pgfsetbuttcap%
\pgfsetroundjoin%
\definecolor{currentfill}{rgb}{0.993456,0.698810,0.186041}%
\pgfsetfillcolor{currentfill}%
\pgfsetlinewidth{0.000000pt}%
\definecolor{currentstroke}{rgb}{0.000000,0.000000,0.000000}%
\pgfsetstrokecolor{currentstroke}%
\pgfsetdash{}{0pt}%
\pgfpathmoveto{\pgfqpoint{1.707645in}{2.261003in}}%
\pgfpathlineto{\pgfqpoint{1.679165in}{2.180486in}}%
\pgfpathlineto{\pgfqpoint{1.691922in}{2.103469in}}%
\pgfpathlineto{\pgfqpoint{1.707645in}{2.261003in}}%
\pgfpathclose%
\pgfusepath{fill}%
\end{pgfscope}%
\begin{pgfscope}%
\pgfpathrectangle{\pgfqpoint{0.100000in}{0.100000in}}{\pgfqpoint{3.957178in}{3.957178in}}%
\pgfusepath{clip}%
\pgfsetbuttcap%
\pgfsetroundjoin%
\definecolor{currentfill}{rgb}{0.568201,0.055778,0.639477}%
\pgfsetfillcolor{currentfill}%
\pgfsetlinewidth{0.000000pt}%
\definecolor{currentstroke}{rgb}{0.000000,0.000000,0.000000}%
\pgfsetstrokecolor{currentstroke}%
\pgfsetdash{}{0pt}%
\pgfpathmoveto{\pgfqpoint{2.558352in}{1.557738in}}%
\pgfpathlineto{\pgfqpoint{2.356873in}{1.187613in}}%
\pgfpathlineto{\pgfqpoint{2.480551in}{1.051058in}}%
\pgfpathlineto{\pgfqpoint{2.558352in}{1.557738in}}%
\pgfpathclose%
\pgfusepath{fill}%
\end{pgfscope}%
\begin{pgfscope}%
\pgfpathrectangle{\pgfqpoint{0.100000in}{0.100000in}}{\pgfqpoint{3.957178in}{3.957178in}}%
\pgfusepath{clip}%
\pgfsetbuttcap%
\pgfsetroundjoin%
\definecolor{currentfill}{rgb}{0.976265,0.868016,0.143351}%
\pgfsetfillcolor{currentfill}%
\pgfsetlinewidth{0.000000pt}%
\definecolor{currentstroke}{rgb}{0.000000,0.000000,0.000000}%
\pgfsetstrokecolor{currentstroke}%
\pgfsetdash{}{0pt}%
\pgfpathmoveto{\pgfqpoint{1.893950in}{2.383623in}}%
\pgfpathlineto{\pgfqpoint{1.882422in}{2.450114in}}%
\pgfpathlineto{\pgfqpoint{1.707645in}{2.261003in}}%
\pgfpathlineto{\pgfqpoint{1.893950in}{2.383623in}}%
\pgfpathclose%
\pgfusepath{fill}%
\end{pgfscope}%
\begin{pgfscope}%
\pgfpathrectangle{\pgfqpoint{0.100000in}{0.100000in}}{\pgfqpoint{3.957178in}{3.957178in}}%
\pgfusepath{clip}%
\pgfsetbuttcap%
\pgfsetroundjoin%
\definecolor{currentfill}{rgb}{0.435734,0.001127,0.659797}%
\pgfsetfillcolor{currentfill}%
\pgfsetlinewidth{0.000000pt}%
\definecolor{currentstroke}{rgb}{0.000000,0.000000,0.000000}%
\pgfsetstrokecolor{currentstroke}%
\pgfsetdash{}{0pt}%
\pgfpathmoveto{\pgfqpoint{2.336674in}{1.076454in}}%
\pgfpathlineto{\pgfqpoint{2.480551in}{1.051058in}}%
\pgfpathlineto{\pgfqpoint{2.356873in}{1.187613in}}%
\pgfpathlineto{\pgfqpoint{2.336674in}{1.076454in}}%
\pgfpathclose%
\pgfusepath{fill}%
\end{pgfscope}%
\begin{pgfscope}%
\pgfpathrectangle{\pgfqpoint{0.100000in}{0.100000in}}{\pgfqpoint{3.957178in}{3.957178in}}%
\pgfusepath{clip}%
\pgfsetbuttcap%
\pgfsetroundjoin%
\definecolor{currentfill}{rgb}{0.959276,0.921407,0.151566}%
\pgfsetfillcolor{currentfill}%
\pgfsetlinewidth{0.000000pt}%
\definecolor{currentstroke}{rgb}{0.000000,0.000000,0.000000}%
\pgfsetstrokecolor{currentstroke}%
\pgfsetdash{}{0pt}%
\pgfpathmoveto{\pgfqpoint{1.898991in}{2.394226in}}%
\pgfpathlineto{\pgfqpoint{1.882422in}{2.450114in}}%
\pgfpathlineto{\pgfqpoint{1.893950in}{2.383623in}}%
\pgfpathlineto{\pgfqpoint{1.898991in}{2.394226in}}%
\pgfpathclose%
\pgfusepath{fill}%
\end{pgfscope}%
\begin{pgfscope}%
\pgfpathrectangle{\pgfqpoint{0.100000in}{0.100000in}}{\pgfqpoint{3.957178in}{3.957178in}}%
\pgfusepath{clip}%
\pgfsetbuttcap%
\pgfsetroundjoin%
\definecolor{currentfill}{rgb}{0.625987,0.103312,0.611305}%
\pgfsetfillcolor{currentfill}%
\pgfsetlinewidth{0.000000pt}%
\definecolor{currentstroke}{rgb}{0.000000,0.000000,0.000000}%
\pgfsetstrokecolor{currentstroke}%
\pgfsetdash{}{0pt}%
\pgfpathmoveto{\pgfqpoint{1.651923in}{1.660516in}}%
\pgfpathlineto{\pgfqpoint{2.336674in}{1.076454in}}%
\pgfpathlineto{\pgfqpoint{2.212049in}{1.262956in}}%
\pgfpathlineto{\pgfqpoint{1.651923in}{1.660516in}}%
\pgfpathclose%
\pgfusepath{fill}%
\end{pgfscope}%
\begin{pgfscope}%
\pgfpathrectangle{\pgfqpoint{0.100000in}{0.100000in}}{\pgfqpoint{3.957178in}{3.957178in}}%
\pgfusepath{clip}%
\pgfsetbuttcap%
\pgfsetroundjoin%
\definecolor{currentfill}{rgb}{0.989935,0.663787,0.204859}%
\pgfsetfillcolor{currentfill}%
\pgfsetlinewidth{0.000000pt}%
\definecolor{currentstroke}{rgb}{0.000000,0.000000,0.000000}%
\pgfsetstrokecolor{currentstroke}%
\pgfsetdash{}{0pt}%
\pgfpathmoveto{\pgfqpoint{1.761100in}{2.035047in}}%
\pgfpathlineto{\pgfqpoint{1.707645in}{2.261003in}}%
\pgfpathlineto{\pgfqpoint{1.691922in}{2.103469in}}%
\pgfpathlineto{\pgfqpoint{1.761100in}{2.035047in}}%
\pgfpathclose%
\pgfusepath{fill}%
\end{pgfscope}%
\begin{pgfscope}%
\pgfpathrectangle{\pgfqpoint{0.100000in}{0.100000in}}{\pgfqpoint{3.957178in}{3.957178in}}%
\pgfusepath{clip}%
\pgfsetbuttcap%
\pgfsetroundjoin%
\definecolor{currentfill}{rgb}{0.949217,0.517763,0.295662}%
\pgfsetfillcolor{currentfill}%
\pgfsetlinewidth{0.000000pt}%
\definecolor{currentstroke}{rgb}{0.000000,0.000000,0.000000}%
\pgfsetstrokecolor{currentstroke}%
\pgfsetdash{}{0pt}%
\pgfpathmoveto{\pgfqpoint{1.691922in}{2.103469in}}%
\pgfpathlineto{\pgfqpoint{1.682618in}{1.716885in}}%
\pgfpathlineto{\pgfqpoint{1.761100in}{2.035047in}}%
\pgfpathlineto{\pgfqpoint{1.691922in}{2.103469in}}%
\pgfpathclose%
\pgfusepath{fill}%
\end{pgfscope}%
\begin{pgfscope}%
\pgfpathrectangle{\pgfqpoint{0.100000in}{0.100000in}}{\pgfqpoint{3.957178in}{3.957178in}}%
\pgfusepath{clip}%
\pgfsetbuttcap%
\pgfsetroundjoin%
\definecolor{currentfill}{rgb}{0.994355,0.746995,0.163821}%
\pgfsetfillcolor{currentfill}%
\pgfsetlinewidth{0.000000pt}%
\definecolor{currentstroke}{rgb}{0.000000,0.000000,0.000000}%
\pgfsetstrokecolor{currentstroke}%
\pgfsetdash{}{0pt}%
\pgfpathmoveto{\pgfqpoint{1.707645in}{2.261003in}}%
\pgfpathlineto{\pgfqpoint{1.761100in}{2.035047in}}%
\pgfpathlineto{\pgfqpoint{1.893950in}{2.383623in}}%
\pgfpathlineto{\pgfqpoint{1.707645in}{2.261003in}}%
\pgfpathclose%
\pgfusepath{fill}%
\end{pgfscope}%
\begin{pgfscope}%
\pgfpathrectangle{\pgfqpoint{0.100000in}{0.100000in}}{\pgfqpoint{3.957178in}{3.957178in}}%
\pgfusepath{clip}%
\pgfsetbuttcap%
\pgfsetroundjoin%
\definecolor{currentfill}{rgb}{0.968443,0.894564,0.147014}%
\pgfsetfillcolor{currentfill}%
\pgfsetlinewidth{0.000000pt}%
\definecolor{currentstroke}{rgb}{0.000000,0.000000,0.000000}%
\pgfsetstrokecolor{currentstroke}%
\pgfsetdash{}{0pt}%
\pgfpathmoveto{\pgfqpoint{1.969490in}{2.479538in}}%
\pgfpathlineto{\pgfqpoint{2.082218in}{2.187897in}}%
\pgfpathlineto{\pgfqpoint{2.328755in}{2.465158in}}%
\pgfpathlineto{\pgfqpoint{1.969490in}{2.479538in}}%
\pgfpathclose%
\pgfusepath{fill}%
\end{pgfscope}%
\begin{pgfscope}%
\pgfpathrectangle{\pgfqpoint{0.100000in}{0.100000in}}{\pgfqpoint{3.957178in}{3.957178in}}%
\pgfusepath{clip}%
\pgfsetbuttcap%
\pgfsetroundjoin%
\definecolor{currentfill}{rgb}{0.866078,0.369660,0.400126}%
\pgfsetfillcolor{currentfill}%
\pgfsetlinewidth{0.000000pt}%
\definecolor{currentstroke}{rgb}{0.000000,0.000000,0.000000}%
\pgfsetstrokecolor{currentstroke}%
\pgfsetdash{}{0pt}%
\pgfpathmoveto{\pgfqpoint{1.921272in}{1.823569in}}%
\pgfpathlineto{\pgfqpoint{1.682618in}{1.716885in}}%
\pgfpathlineto{\pgfqpoint{1.651923in}{1.660516in}}%
\pgfpathlineto{\pgfqpoint{1.921272in}{1.823569in}}%
\pgfpathclose%
\pgfusepath{fill}%
\end{pgfscope}%
\begin{pgfscope}%
\pgfpathrectangle{\pgfqpoint{0.100000in}{0.100000in}}{\pgfqpoint{3.957178in}{3.957178in}}%
\pgfusepath{clip}%
\pgfsetbuttcap%
\pgfsetroundjoin%
\definecolor{currentfill}{rgb}{0.964021,0.907950,0.149370}%
\pgfsetfillcolor{currentfill}%
\pgfsetlinewidth{0.000000pt}%
\definecolor{currentstroke}{rgb}{0.000000,0.000000,0.000000}%
\pgfsetstrokecolor{currentstroke}%
\pgfsetdash{}{0pt}%
\pgfpathmoveto{\pgfqpoint{1.969490in}{2.479538in}}%
\pgfpathlineto{\pgfqpoint{1.898991in}{2.394226in}}%
\pgfpathlineto{\pgfqpoint{1.974709in}{2.309071in}}%
\pgfpathlineto{\pgfqpoint{1.969490in}{2.479538in}}%
\pgfpathclose%
\pgfusepath{fill}%
\end{pgfscope}%
\begin{pgfscope}%
\pgfpathrectangle{\pgfqpoint{0.100000in}{0.100000in}}{\pgfqpoint{3.957178in}{3.957178in}}%
\pgfusepath{clip}%
\pgfsetbuttcap%
\pgfsetroundjoin%
\definecolor{currentfill}{rgb}{0.846788,0.342551,0.420579}%
\pgfsetfillcolor{currentfill}%
\pgfsetlinewidth{0.000000pt}%
\definecolor{currentstroke}{rgb}{0.000000,0.000000,0.000000}%
\pgfsetstrokecolor{currentstroke}%
\pgfsetdash{}{0pt}%
\pgfpathmoveto{\pgfqpoint{2.525070in}{1.721761in}}%
\pgfpathlineto{\pgfqpoint{2.478698in}{1.804673in}}%
\pgfpathlineto{\pgfqpoint{2.558352in}{1.557738in}}%
\pgfpathlineto{\pgfqpoint{2.525070in}{1.721761in}}%
\pgfpathclose%
\pgfusepath{fill}%
\end{pgfscope}%
\begin{pgfscope}%
\pgfpathrectangle{\pgfqpoint{0.100000in}{0.100000in}}{\pgfqpoint{3.957178in}{3.957178in}}%
\pgfusepath{clip}%
\pgfsetbuttcap%
\pgfsetroundjoin%
\definecolor{currentfill}{rgb}{0.506454,0.016333,0.656202}%
\pgfsetfillcolor{currentfill}%
\pgfsetlinewidth{0.000000pt}%
\definecolor{currentstroke}{rgb}{0.000000,0.000000,0.000000}%
\pgfsetstrokecolor{currentstroke}%
\pgfsetdash{}{0pt}%
\pgfpathmoveto{\pgfqpoint{2.356873in}{1.187613in}}%
\pgfpathlineto{\pgfqpoint{2.212049in}{1.262956in}}%
\pgfpathlineto{\pgfqpoint{2.336674in}{1.076454in}}%
\pgfpathlineto{\pgfqpoint{2.356873in}{1.187613in}}%
\pgfpathclose%
\pgfusepath{fill}%
\end{pgfscope}%
\begin{pgfscope}%
\pgfpathrectangle{\pgfqpoint{0.100000in}{0.100000in}}{\pgfqpoint{3.957178in}{3.957178in}}%
\pgfusepath{clip}%
\pgfsetbuttcap%
\pgfsetroundjoin%
\definecolor{currentfill}{rgb}{0.760264,0.235976,0.505794}%
\pgfsetfillcolor{currentfill}%
\pgfsetlinewidth{0.000000pt}%
\definecolor{currentstroke}{rgb}{0.000000,0.000000,0.000000}%
\pgfsetstrokecolor{currentstroke}%
\pgfsetdash{}{0pt}%
\pgfpathmoveto{\pgfqpoint{2.408786in}{1.844473in}}%
\pgfpathlineto{\pgfqpoint{2.356873in}{1.187613in}}%
\pgfpathlineto{\pgfqpoint{2.558352in}{1.557738in}}%
\pgfpathlineto{\pgfqpoint{2.408786in}{1.844473in}}%
\pgfpathclose%
\pgfusepath{fill}%
\end{pgfscope}%
\begin{pgfscope}%
\pgfpathrectangle{\pgfqpoint{0.100000in}{0.100000in}}{\pgfqpoint{3.957178in}{3.957178in}}%
\pgfusepath{clip}%
\pgfsetbuttcap%
\pgfsetroundjoin%
\definecolor{currentfill}{rgb}{0.976428,0.596595,0.244767}%
\pgfsetfillcolor{currentfill}%
\pgfsetlinewidth{0.000000pt}%
\definecolor{currentstroke}{rgb}{0.000000,0.000000,0.000000}%
\pgfsetstrokecolor{currentstroke}%
\pgfsetdash{}{0pt}%
\pgfpathmoveto{\pgfqpoint{2.328755in}{2.465158in}}%
\pgfpathlineto{\pgfqpoint{2.408786in}{1.844473in}}%
\pgfpathlineto{\pgfqpoint{2.478698in}{1.804673in}}%
\pgfpathlineto{\pgfqpoint{2.328755in}{2.465158in}}%
\pgfpathclose%
\pgfusepath{fill}%
\end{pgfscope}%
\begin{pgfscope}%
\pgfpathrectangle{\pgfqpoint{0.100000in}{0.100000in}}{\pgfqpoint{3.957178in}{3.957178in}}%
\pgfusepath{clip}%
\pgfsetbuttcap%
\pgfsetroundjoin%
\definecolor{currentfill}{rgb}{0.787133,0.266922,0.480307}%
\pgfsetfillcolor{currentfill}%
\pgfsetlinewidth{0.000000pt}%
\definecolor{currentstroke}{rgb}{0.000000,0.000000,0.000000}%
\pgfsetstrokecolor{currentstroke}%
\pgfsetdash{}{0pt}%
\pgfpathmoveto{\pgfqpoint{1.651923in}{1.660516in}}%
\pgfpathlineto{\pgfqpoint{2.212049in}{1.262956in}}%
\pgfpathlineto{\pgfqpoint{1.974809in}{1.790662in}}%
\pgfpathlineto{\pgfqpoint{1.651923in}{1.660516in}}%
\pgfpathclose%
\pgfusepath{fill}%
\end{pgfscope}%
\begin{pgfscope}%
\pgfpathrectangle{\pgfqpoint{0.100000in}{0.100000in}}{\pgfqpoint{3.957178in}{3.957178in}}%
\pgfusepath{clip}%
\pgfsetbuttcap%
\pgfsetroundjoin%
\definecolor{currentfill}{rgb}{0.970533,0.887896,0.145919}%
\pgfsetfillcolor{currentfill}%
\pgfsetlinewidth{0.000000pt}%
\definecolor{currentstroke}{rgb}{0.000000,0.000000,0.000000}%
\pgfsetstrokecolor{currentstroke}%
\pgfsetdash{}{0pt}%
\pgfpathmoveto{\pgfqpoint{1.893950in}{2.383623in}}%
\pgfpathlineto{\pgfqpoint{1.974709in}{2.309071in}}%
\pgfpathlineto{\pgfqpoint{1.898991in}{2.394226in}}%
\pgfpathlineto{\pgfqpoint{1.893950in}{2.383623in}}%
\pgfpathclose%
\pgfusepath{fill}%
\end{pgfscope}%
\begin{pgfscope}%
\pgfpathrectangle{\pgfqpoint{0.100000in}{0.100000in}}{\pgfqpoint{3.957178in}{3.957178in}}%
\pgfusepath{clip}%
\pgfsetbuttcap%
\pgfsetroundjoin%
\definecolor{currentfill}{rgb}{0.869203,0.374212,0.396738}%
\pgfsetfillcolor{currentfill}%
\pgfsetlinewidth{0.000000pt}%
\definecolor{currentstroke}{rgb}{0.000000,0.000000,0.000000}%
\pgfsetstrokecolor{currentstroke}%
\pgfsetdash{}{0pt}%
\pgfpathmoveto{\pgfqpoint{2.558352in}{1.557738in}}%
\pgfpathlineto{\pgfqpoint{2.478698in}{1.804673in}}%
\pgfpathlineto{\pgfqpoint{2.408786in}{1.844473in}}%
\pgfpathlineto{\pgfqpoint{2.558352in}{1.557738in}}%
\pgfpathclose%
\pgfusepath{fill}%
\end{pgfscope}%
\begin{pgfscope}%
\pgfpathrectangle{\pgfqpoint{0.100000in}{0.100000in}}{\pgfqpoint{3.957178in}{3.957178in}}%
\pgfusepath{clip}%
\pgfsetbuttcap%
\pgfsetroundjoin%
\definecolor{currentfill}{rgb}{0.923287,0.463251,0.332801}%
\pgfsetfillcolor{currentfill}%
\pgfsetlinewidth{0.000000pt}%
\definecolor{currentstroke}{rgb}{0.000000,0.000000,0.000000}%
\pgfsetstrokecolor{currentstroke}%
\pgfsetdash{}{0pt}%
\pgfpathmoveto{\pgfqpoint{1.761100in}{2.035047in}}%
\pgfpathlineto{\pgfqpoint{1.682618in}{1.716885in}}%
\pgfpathlineto{\pgfqpoint{1.921272in}{1.823569in}}%
\pgfpathlineto{\pgfqpoint{1.761100in}{2.035047in}}%
\pgfpathclose%
\pgfusepath{fill}%
\end{pgfscope}%
\begin{pgfscope}%
\pgfpathrectangle{\pgfqpoint{0.100000in}{0.100000in}}{\pgfqpoint{3.957178in}{3.957178in}}%
\pgfusepath{clip}%
\pgfsetbuttcap%
\pgfsetroundjoin%
\definecolor{currentfill}{rgb}{0.625987,0.103312,0.611305}%
\pgfsetfillcolor{currentfill}%
\pgfsetlinewidth{0.000000pt}%
\definecolor{currentstroke}{rgb}{0.000000,0.000000,0.000000}%
\pgfsetstrokecolor{currentstroke}%
\pgfsetdash{}{0pt}%
\pgfpathmoveto{\pgfqpoint{2.257180in}{1.507161in}}%
\pgfpathlineto{\pgfqpoint{2.212049in}{1.262956in}}%
\pgfpathlineto{\pgfqpoint{2.356873in}{1.187613in}}%
\pgfpathlineto{\pgfqpoint{2.257180in}{1.507161in}}%
\pgfpathclose%
\pgfusepath{fill}%
\end{pgfscope}%
\begin{pgfscope}%
\pgfpathrectangle{\pgfqpoint{0.100000in}{0.100000in}}{\pgfqpoint{3.957178in}{3.957178in}}%
\pgfusepath{clip}%
\pgfsetbuttcap%
\pgfsetroundjoin%
\definecolor{currentfill}{rgb}{0.992505,0.777967,0.152855}%
\pgfsetfillcolor{currentfill}%
\pgfsetlinewidth{0.000000pt}%
\definecolor{currentstroke}{rgb}{0.000000,0.000000,0.000000}%
\pgfsetstrokecolor{currentstroke}%
\pgfsetdash{}{0pt}%
\pgfpathmoveto{\pgfqpoint{1.893950in}{2.383623in}}%
\pgfpathlineto{\pgfqpoint{1.761100in}{2.035047in}}%
\pgfpathlineto{\pgfqpoint{1.974709in}{2.309071in}}%
\pgfpathlineto{\pgfqpoint{1.893950in}{2.383623in}}%
\pgfpathclose%
\pgfusepath{fill}%
\end{pgfscope}%
\begin{pgfscope}%
\pgfpathrectangle{\pgfqpoint{0.100000in}{0.100000in}}{\pgfqpoint{3.957178in}{3.957178in}}%
\pgfusepath{clip}%
\pgfsetbuttcap%
\pgfsetroundjoin%
\definecolor{currentfill}{rgb}{0.979644,0.854866,0.142453}%
\pgfsetfillcolor{currentfill}%
\pgfsetlinewidth{0.000000pt}%
\definecolor{currentstroke}{rgb}{0.000000,0.000000,0.000000}%
\pgfsetstrokecolor{currentstroke}%
\pgfsetdash{}{0pt}%
\pgfpathmoveto{\pgfqpoint{1.974709in}{2.309071in}}%
\pgfpathlineto{\pgfqpoint{2.082218in}{2.187897in}}%
\pgfpathlineto{\pgfqpoint{1.969490in}{2.479538in}}%
\pgfpathlineto{\pgfqpoint{1.974709in}{2.309071in}}%
\pgfpathclose%
\pgfusepath{fill}%
\end{pgfscope}%
\begin{pgfscope}%
\pgfpathrectangle{\pgfqpoint{0.100000in}{0.100000in}}{\pgfqpoint{3.957178in}{3.957178in}}%
\pgfusepath{clip}%
\pgfsetbuttcap%
\pgfsetroundjoin%
\definecolor{currentfill}{rgb}{0.884436,0.397139,0.379860}%
\pgfsetfillcolor{currentfill}%
\pgfsetlinewidth{0.000000pt}%
\definecolor{currentstroke}{rgb}{0.000000,0.000000,0.000000}%
\pgfsetstrokecolor{currentstroke}%
\pgfsetdash{}{0pt}%
\pgfpathmoveto{\pgfqpoint{1.974809in}{1.790662in}}%
\pgfpathlineto{\pgfqpoint{1.921272in}{1.823569in}}%
\pgfpathlineto{\pgfqpoint{1.651923in}{1.660516in}}%
\pgfpathlineto{\pgfqpoint{1.974809in}{1.790662in}}%
\pgfpathclose%
\pgfusepath{fill}%
\end{pgfscope}%
\begin{pgfscope}%
\pgfpathrectangle{\pgfqpoint{0.100000in}{0.100000in}}{\pgfqpoint{3.957178in}{3.957178in}}%
\pgfusepath{clip}%
\pgfsetbuttcap%
\pgfsetroundjoin%
\definecolor{currentfill}{rgb}{0.994103,0.710698,0.180097}%
\pgfsetfillcolor{currentfill}%
\pgfsetlinewidth{0.000000pt}%
\definecolor{currentstroke}{rgb}{0.000000,0.000000,0.000000}%
\pgfsetstrokecolor{currentstroke}%
\pgfsetdash{}{0pt}%
\pgfpathmoveto{\pgfqpoint{2.082218in}{2.187897in}}%
\pgfpathlineto{\pgfqpoint{2.408786in}{1.844473in}}%
\pgfpathlineto{\pgfqpoint{2.328755in}{2.465158in}}%
\pgfpathlineto{\pgfqpoint{2.082218in}{2.187897in}}%
\pgfpathclose%
\pgfusepath{fill}%
\end{pgfscope}%
\begin{pgfscope}%
\pgfpathrectangle{\pgfqpoint{0.100000in}{0.100000in}}{\pgfqpoint{3.957178in}{3.957178in}}%
\pgfusepath{clip}%
\pgfsetbuttcap%
\pgfsetroundjoin%
\definecolor{currentfill}{rgb}{0.756304,0.231555,0.509468}%
\pgfsetfillcolor{currentfill}%
\pgfsetlinewidth{0.000000pt}%
\definecolor{currentstroke}{rgb}{0.000000,0.000000,0.000000}%
\pgfsetstrokecolor{currentstroke}%
\pgfsetdash{}{0pt}%
\pgfpathmoveto{\pgfqpoint{2.408786in}{1.844473in}}%
\pgfpathlineto{\pgfqpoint{2.257180in}{1.507161in}}%
\pgfpathlineto{\pgfqpoint{2.356873in}{1.187613in}}%
\pgfpathlineto{\pgfqpoint{2.408786in}{1.844473in}}%
\pgfpathclose%
\pgfusepath{fill}%
\end{pgfscope}%
\begin{pgfscope}%
\pgfpathrectangle{\pgfqpoint{0.100000in}{0.100000in}}{\pgfqpoint{3.957178in}{3.957178in}}%
\pgfusepath{clip}%
\pgfsetbuttcap%
\pgfsetroundjoin%
\definecolor{currentfill}{rgb}{0.706178,0.178437,0.553657}%
\pgfsetfillcolor{currentfill}%
\pgfsetlinewidth{0.000000pt}%
\definecolor{currentstroke}{rgb}{0.000000,0.000000,0.000000}%
\pgfsetstrokecolor{currentstroke}%
\pgfsetdash{}{0pt}%
\pgfpathmoveto{\pgfqpoint{2.212049in}{1.262956in}}%
\pgfpathlineto{\pgfqpoint{2.257180in}{1.507161in}}%
\pgfpathlineto{\pgfqpoint{2.196255in}{1.504269in}}%
\pgfpathlineto{\pgfqpoint{2.212049in}{1.262956in}}%
\pgfpathclose%
\pgfusepath{fill}%
\end{pgfscope}%
\begin{pgfscope}%
\pgfpathrectangle{\pgfqpoint{0.100000in}{0.100000in}}{\pgfqpoint{3.957178in}{3.957178in}}%
\pgfusepath{clip}%
\pgfsetbuttcap%
\pgfsetroundjoin%
\definecolor{currentfill}{rgb}{0.981826,0.618572,0.231287}%
\pgfsetfillcolor{currentfill}%
\pgfsetlinewidth{0.000000pt}%
\definecolor{currentstroke}{rgb}{0.000000,0.000000,0.000000}%
\pgfsetstrokecolor{currentstroke}%
\pgfsetdash{}{0pt}%
\pgfpathmoveto{\pgfqpoint{1.921272in}{1.823569in}}%
\pgfpathlineto{\pgfqpoint{1.974709in}{2.309071in}}%
\pgfpathlineto{\pgfqpoint{1.761100in}{2.035047in}}%
\pgfpathlineto{\pgfqpoint{1.921272in}{1.823569in}}%
\pgfpathclose%
\pgfusepath{fill}%
\end{pgfscope}%
\begin{pgfscope}%
\pgfpathrectangle{\pgfqpoint{0.100000in}{0.100000in}}{\pgfqpoint{3.957178in}{3.957178in}}%
\pgfusepath{clip}%
\pgfsetbuttcap%
\pgfsetroundjoin%
\definecolor{currentfill}{rgb}{0.736019,0.209439,0.527908}%
\pgfsetfillcolor{currentfill}%
\pgfsetlinewidth{0.000000pt}%
\definecolor{currentstroke}{rgb}{0.000000,0.000000,0.000000}%
\pgfsetstrokecolor{currentstroke}%
\pgfsetdash{}{0pt}%
\pgfpathmoveto{\pgfqpoint{2.196255in}{1.504269in}}%
\pgfpathlineto{\pgfqpoint{2.097955in}{1.637802in}}%
\pgfpathlineto{\pgfqpoint{2.212049in}{1.262956in}}%
\pgfpathlineto{\pgfqpoint{2.196255in}{1.504269in}}%
\pgfpathclose%
\pgfusepath{fill}%
\end{pgfscope}%
\begin{pgfscope}%
\pgfpathrectangle{\pgfqpoint{0.100000in}{0.100000in}}{\pgfqpoint{3.957178in}{3.957178in}}%
\pgfusepath{clip}%
\pgfsetbuttcap%
\pgfsetroundjoin%
\definecolor{currentfill}{rgb}{0.833422,0.324635,0.434366}%
\pgfsetfillcolor{currentfill}%
\pgfsetlinewidth{0.000000pt}%
\definecolor{currentstroke}{rgb}{0.000000,0.000000,0.000000}%
\pgfsetstrokecolor{currentstroke}%
\pgfsetdash{}{0pt}%
\pgfpathmoveto{\pgfqpoint{2.206940in}{1.582485in}}%
\pgfpathlineto{\pgfqpoint{2.257180in}{1.507161in}}%
\pgfpathlineto{\pgfqpoint{2.408786in}{1.844473in}}%
\pgfpathlineto{\pgfqpoint{2.206940in}{1.582485in}}%
\pgfpathclose%
\pgfusepath{fill}%
\end{pgfscope}%
\begin{pgfscope}%
\pgfpathrectangle{\pgfqpoint{0.100000in}{0.100000in}}{\pgfqpoint{3.957178in}{3.957178in}}%
\pgfusepath{clip}%
\pgfsetbuttcap%
\pgfsetroundjoin%
\definecolor{currentfill}{rgb}{0.775796,0.253658,0.491171}%
\pgfsetfillcolor{currentfill}%
\pgfsetlinewidth{0.000000pt}%
\definecolor{currentstroke}{rgb}{0.000000,0.000000,0.000000}%
\pgfsetstrokecolor{currentstroke}%
\pgfsetdash{}{0pt}%
\pgfpathmoveto{\pgfqpoint{2.196255in}{1.504269in}}%
\pgfpathlineto{\pgfqpoint{2.257180in}{1.507161in}}%
\pgfpathlineto{\pgfqpoint{2.206940in}{1.582485in}}%
\pgfpathlineto{\pgfqpoint{2.196255in}{1.504269in}}%
\pgfpathclose%
\pgfusepath{fill}%
\end{pgfscope}%
\begin{pgfscope}%
\pgfpathrectangle{\pgfqpoint{0.100000in}{0.100000in}}{\pgfqpoint{3.957178in}{3.957178in}}%
\pgfusepath{clip}%
\pgfsetbuttcap%
\pgfsetroundjoin%
\definecolor{currentfill}{rgb}{0.794549,0.275770,0.473117}%
\pgfsetfillcolor{currentfill}%
\pgfsetlinewidth{0.000000pt}%
\definecolor{currentstroke}{rgb}{0.000000,0.000000,0.000000}%
\pgfsetstrokecolor{currentstroke}%
\pgfsetdash{}{0pt}%
\pgfpathmoveto{\pgfqpoint{1.974809in}{1.790662in}}%
\pgfpathlineto{\pgfqpoint{2.212049in}{1.262956in}}%
\pgfpathlineto{\pgfqpoint{2.097955in}{1.637802in}}%
\pgfpathlineto{\pgfqpoint{1.974809in}{1.790662in}}%
\pgfpathclose%
\pgfusepath{fill}%
\end{pgfscope}%
\begin{pgfscope}%
\pgfpathrectangle{\pgfqpoint{0.100000in}{0.100000in}}{\pgfqpoint{3.957178in}{3.957178in}}%
\pgfusepath{clip}%
\pgfsetbuttcap%
\pgfsetroundjoin%
\definecolor{currentfill}{rgb}{0.801855,0.284626,0.465971}%
\pgfsetfillcolor{currentfill}%
\pgfsetlinewidth{0.000000pt}%
\definecolor{currentstroke}{rgb}{0.000000,0.000000,0.000000}%
\pgfsetstrokecolor{currentstroke}%
\pgfsetdash{}{0pt}%
\pgfpathmoveto{\pgfqpoint{2.206940in}{1.582485in}}%
\pgfpathlineto{\pgfqpoint{2.097955in}{1.637802in}}%
\pgfpathlineto{\pgfqpoint{2.196255in}{1.504269in}}%
\pgfpathlineto{\pgfqpoint{2.206940in}{1.582485in}}%
\pgfpathclose%
\pgfusepath{fill}%
\end{pgfscope}%
\begin{pgfscope}%
\pgfpathrectangle{\pgfqpoint{0.100000in}{0.100000in}}{\pgfqpoint{3.957178in}{3.957178in}}%
\pgfusepath{clip}%
\pgfsetbuttcap%
\pgfsetroundjoin%
\definecolor{currentfill}{rgb}{0.970205,0.575028,0.258325}%
\pgfsetfillcolor{currentfill}%
\pgfsetlinewidth{0.000000pt}%
\definecolor{currentstroke}{rgb}{0.000000,0.000000,0.000000}%
\pgfsetstrokecolor{currentstroke}%
\pgfsetdash{}{0pt}%
\pgfpathmoveto{\pgfqpoint{1.980272in}{1.842883in}}%
\pgfpathlineto{\pgfqpoint{1.974709in}{2.309071in}}%
\pgfpathlineto{\pgfqpoint{1.921272in}{1.823569in}}%
\pgfpathlineto{\pgfqpoint{1.980272in}{1.842883in}}%
\pgfpathclose%
\pgfusepath{fill}%
\end{pgfscope}%
\begin{pgfscope}%
\pgfpathrectangle{\pgfqpoint{0.100000in}{0.100000in}}{\pgfqpoint{3.957178in}{3.957178in}}%
\pgfusepath{clip}%
\pgfsetbuttcap%
\pgfsetroundjoin%
\definecolor{currentfill}{rgb}{0.875376,0.383347,0.389976}%
\pgfsetfillcolor{currentfill}%
\pgfsetlinewidth{0.000000pt}%
\definecolor{currentstroke}{rgb}{0.000000,0.000000,0.000000}%
\pgfsetstrokecolor{currentstroke}%
\pgfsetdash{}{0pt}%
\pgfpathmoveto{\pgfqpoint{2.123266in}{1.746012in}}%
\pgfpathlineto{\pgfqpoint{2.206940in}{1.582485in}}%
\pgfpathlineto{\pgfqpoint{2.408786in}{1.844473in}}%
\pgfpathlineto{\pgfqpoint{2.123266in}{1.746012in}}%
\pgfpathclose%
\pgfusepath{fill}%
\end{pgfscope}%
\begin{pgfscope}%
\pgfpathrectangle{\pgfqpoint{0.100000in}{0.100000in}}{\pgfqpoint{3.957178in}{3.957178in}}%
\pgfusepath{clip}%
\pgfsetbuttcap%
\pgfsetroundjoin%
\definecolor{currentfill}{rgb}{0.840155,0.333580,0.427455}%
\pgfsetfillcolor{currentfill}%
\pgfsetlinewidth{0.000000pt}%
\definecolor{currentstroke}{rgb}{0.000000,0.000000,0.000000}%
\pgfsetstrokecolor{currentstroke}%
\pgfsetdash{}{0pt}%
\pgfpathmoveto{\pgfqpoint{2.097955in}{1.637802in}}%
\pgfpathlineto{\pgfqpoint{2.206940in}{1.582485in}}%
\pgfpathlineto{\pgfqpoint{2.117205in}{1.714526in}}%
\pgfpathlineto{\pgfqpoint{2.097955in}{1.637802in}}%
\pgfpathclose%
\pgfusepath{fill}%
\end{pgfscope}%
\begin{pgfscope}%
\pgfpathrectangle{\pgfqpoint{0.100000in}{0.100000in}}{\pgfqpoint{3.957178in}{3.957178in}}%
\pgfusepath{clip}%
\pgfsetbuttcap%
\pgfsetroundjoin%
\definecolor{currentfill}{rgb}{0.991985,0.681179,0.195295}%
\pgfsetfillcolor{currentfill}%
\pgfsetlinewidth{0.000000pt}%
\definecolor{currentstroke}{rgb}{0.000000,0.000000,0.000000}%
\pgfsetstrokecolor{currentstroke}%
\pgfsetdash{}{0pt}%
\pgfpathmoveto{\pgfqpoint{2.082218in}{2.187897in}}%
\pgfpathlineto{\pgfqpoint{1.974709in}{2.309071in}}%
\pgfpathlineto{\pgfqpoint{1.980272in}{1.842883in}}%
\pgfpathlineto{\pgfqpoint{2.082218in}{2.187897in}}%
\pgfpathclose%
\pgfusepath{fill}%
\end{pgfscope}%
\begin{pgfscope}%
\pgfpathrectangle{\pgfqpoint{0.100000in}{0.100000in}}{\pgfqpoint{3.957178in}{3.957178in}}%
\pgfusepath{clip}%
\pgfsetbuttcap%
\pgfsetroundjoin%
\definecolor{currentfill}{rgb}{0.915471,0.448807,0.342890}%
\pgfsetfillcolor{currentfill}%
\pgfsetlinewidth{0.000000pt}%
\definecolor{currentstroke}{rgb}{0.000000,0.000000,0.000000}%
\pgfsetstrokecolor{currentstroke}%
\pgfsetdash{}{0pt}%
\pgfpathmoveto{\pgfqpoint{1.980272in}{1.842883in}}%
\pgfpathlineto{\pgfqpoint{1.921272in}{1.823569in}}%
\pgfpathlineto{\pgfqpoint{1.974809in}{1.790662in}}%
\pgfpathlineto{\pgfqpoint{1.980272in}{1.842883in}}%
\pgfpathclose%
\pgfusepath{fill}%
\end{pgfscope}%
\begin{pgfscope}%
\pgfpathrectangle{\pgfqpoint{0.100000in}{0.100000in}}{\pgfqpoint{3.957178in}{3.957178in}}%
\pgfusepath{clip}%
\pgfsetbuttcap%
\pgfsetroundjoin%
\definecolor{currentfill}{rgb}{0.953428,0.527960,0.288883}%
\pgfsetfillcolor{currentfill}%
\pgfsetlinewidth{0.000000pt}%
\definecolor{currentstroke}{rgb}{0.000000,0.000000,0.000000}%
\pgfsetstrokecolor{currentstroke}%
\pgfsetdash{}{0pt}%
\pgfpathmoveto{\pgfqpoint{2.408786in}{1.844473in}}%
\pgfpathlineto{\pgfqpoint{2.082218in}{2.187897in}}%
\pgfpathlineto{\pgfqpoint{2.123266in}{1.746012in}}%
\pgfpathlineto{\pgfqpoint{2.408786in}{1.844473in}}%
\pgfpathclose%
\pgfusepath{fill}%
\end{pgfscope}%
\begin{pgfscope}%
\pgfpathrectangle{\pgfqpoint{0.100000in}{0.100000in}}{\pgfqpoint{3.957178in}{3.957178in}}%
\pgfusepath{clip}%
\pgfsetbuttcap%
\pgfsetroundjoin%
\definecolor{currentfill}{rgb}{0.872303,0.378774,0.393355}%
\pgfsetfillcolor{currentfill}%
\pgfsetlinewidth{0.000000pt}%
\definecolor{currentstroke}{rgb}{0.000000,0.000000,0.000000}%
\pgfsetstrokecolor{currentstroke}%
\pgfsetdash{}{0pt}%
\pgfpathmoveto{\pgfqpoint{2.117205in}{1.714526in}}%
\pgfpathlineto{\pgfqpoint{1.974809in}{1.790662in}}%
\pgfpathlineto{\pgfqpoint{2.097955in}{1.637802in}}%
\pgfpathlineto{\pgfqpoint{2.117205in}{1.714526in}}%
\pgfpathclose%
\pgfusepath{fill}%
\end{pgfscope}%
\begin{pgfscope}%
\pgfpathrectangle{\pgfqpoint{0.100000in}{0.100000in}}{\pgfqpoint{3.957178in}{3.957178in}}%
\pgfusepath{clip}%
\pgfsetbuttcap%
\pgfsetroundjoin%
\definecolor{currentfill}{rgb}{0.856547,0.356066,0.410322}%
\pgfsetfillcolor{currentfill}%
\pgfsetlinewidth{0.000000pt}%
\definecolor{currentstroke}{rgb}{0.000000,0.000000,0.000000}%
\pgfsetstrokecolor{currentstroke}%
\pgfsetdash{}{0pt}%
\pgfpathmoveto{\pgfqpoint{2.206940in}{1.582485in}}%
\pgfpathlineto{\pgfqpoint{2.123266in}{1.746012in}}%
\pgfpathlineto{\pgfqpoint{2.117205in}{1.714526in}}%
\pgfpathlineto{\pgfqpoint{2.206940in}{1.582485in}}%
\pgfpathclose%
\pgfusepath{fill}%
\end{pgfscope}%
\begin{pgfscope}%
\pgfpathrectangle{\pgfqpoint{0.100000in}{0.100000in}}{\pgfqpoint{3.957178in}{3.957178in}}%
\pgfusepath{clip}%
\pgfsetbuttcap%
\pgfsetroundjoin%
\definecolor{currentfill}{rgb}{0.901807,0.425087,0.359688}%
\pgfsetfillcolor{currentfill}%
\pgfsetlinewidth{0.000000pt}%
\definecolor{currentstroke}{rgb}{0.000000,0.000000,0.000000}%
\pgfsetstrokecolor{currentstroke}%
\pgfsetdash{}{0pt}%
\pgfpathmoveto{\pgfqpoint{1.980272in}{1.842883in}}%
\pgfpathlineto{\pgfqpoint{1.974809in}{1.790662in}}%
\pgfpathlineto{\pgfqpoint{2.117205in}{1.714526in}}%
\pgfpathlineto{\pgfqpoint{1.980272in}{1.842883in}}%
\pgfpathclose%
\pgfusepath{fill}%
\end{pgfscope}%
\begin{pgfscope}%
\pgfpathrectangle{\pgfqpoint{0.100000in}{0.100000in}}{\pgfqpoint{3.957178in}{3.957178in}}%
\pgfusepath{clip}%
\pgfsetbuttcap%
\pgfsetroundjoin%
\definecolor{currentfill}{rgb}{0.896131,0.415712,0.366407}%
\pgfsetfillcolor{currentfill}%
\pgfsetlinewidth{0.000000pt}%
\definecolor{currentstroke}{rgb}{0.000000,0.000000,0.000000}%
\pgfsetstrokecolor{currentstroke}%
\pgfsetdash{}{0pt}%
\pgfpathmoveto{\pgfqpoint{2.117205in}{1.714526in}}%
\pgfpathlineto{\pgfqpoint{2.123266in}{1.746012in}}%
\pgfpathlineto{\pgfqpoint{1.980272in}{1.842883in}}%
\pgfpathlineto{\pgfqpoint{2.117205in}{1.714526in}}%
\pgfpathclose%
\pgfusepath{fill}%
\end{pgfscope}%
\begin{pgfscope}%
\pgfpathrectangle{\pgfqpoint{0.100000in}{0.100000in}}{\pgfqpoint{3.957178in}{3.957178in}}%
\pgfusepath{clip}%
\pgfsetbuttcap%
\pgfsetroundjoin%
\definecolor{currentfill}{rgb}{0.955470,0.533093,0.285490}%
\pgfsetfillcolor{currentfill}%
\pgfsetlinewidth{0.000000pt}%
\definecolor{currentstroke}{rgb}{0.000000,0.000000,0.000000}%
\pgfsetstrokecolor{currentstroke}%
\pgfsetdash{}{0pt}%
\pgfpathmoveto{\pgfqpoint{2.123266in}{1.746012in}}%
\pgfpathlineto{\pgfqpoint{2.082218in}{2.187897in}}%
\pgfpathlineto{\pgfqpoint{1.980272in}{1.842883in}}%
\pgfpathlineto{\pgfqpoint{2.123266in}{1.746012in}}%
\pgfpathclose%
\pgfusepath{fill}%
\end{pgfscope}%
\begin{pgfscope}%
\definecolor{textcolor}{rgb}{0.000000,0.000000,0.000000}%
\pgfsetstrokecolor{textcolor}%
\pgfsetfillcolor{textcolor}%
\pgftext[x=2.078589in,y=4.140512in,,base]{\color{textcolor}{\rmfamily\fontsize{12.000000}{14.400000}\selectfont\catcode`\^=\active\def^{\ifmmode\sp\else\^{}\fi}\catcode`\%=\active\def%{\%}Superficie 3D triangulada (datos dispersos)}}%
\end{pgfscope}%
\begin{pgfscope}%
\pgfsetbuttcap%
\pgfsetmiterjoin%
\definecolor{currentfill}{rgb}{1.000000,1.000000,1.000000}%
\pgfsetfillcolor{currentfill}%
\pgfsetlinewidth{0.000000pt}%
\definecolor{currentstroke}{rgb}{0.000000,0.000000,0.000000}%
\pgfsetstrokecolor{currentstroke}%
\pgfsetstrokeopacity{0.000000}%
\pgfsetdash{}{0pt}%
\pgfpathmoveto{\pgfqpoint{4.255037in}{0.357758in}}%
\pgfpathlineto{\pgfqpoint{4.427120in}{0.357758in}}%
\pgfpathlineto{\pgfqpoint{4.427120in}{3.799420in}}%
\pgfpathlineto{\pgfqpoint{4.255037in}{3.799420in}}%
\pgfpathlineto{\pgfqpoint{4.255037in}{0.357758in}}%
\pgfpathclose%
\pgfusepath{fill}%
\end{pgfscope}%
\begin{pgfscope}%
\pgfpathrectangle{\pgfqpoint{4.255037in}{0.357758in}}{\pgfqpoint{0.172083in}{3.441662in}}%
\pgfusepath{clip}%
\pgfsetbuttcap%
\pgfsetroundjoin%
\definecolor{currentfill}{rgb}{0.050383,0.029803,0.527975}%
\pgfsetfillcolor{currentfill}%
\pgfsetlinewidth{0.000000pt}%
\definecolor{currentstroke}{rgb}{0.000000,0.000000,0.000000}%
\pgfsetstrokecolor{currentstroke}%
\pgfsetdash{}{0pt}%
\pgfpathmoveto{\pgfqpoint{4.255037in}{0.357758in}}%
\pgfpathlineto{\pgfqpoint{4.427120in}{0.357758in}}%
\pgfpathlineto{\pgfqpoint{4.427120in}{0.371202in}}%
\pgfpathlineto{\pgfqpoint{4.255037in}{0.371202in}}%
\pgfpathlineto{\pgfqpoint{4.255037in}{0.357758in}}%
\pgfusepath{fill}%
\end{pgfscope}%
\begin{pgfscope}%
\pgfpathrectangle{\pgfqpoint{4.255037in}{0.357758in}}{\pgfqpoint{0.172083in}{3.441662in}}%
\pgfusepath{clip}%
\pgfsetbuttcap%
\pgfsetroundjoin%
\definecolor{currentfill}{rgb}{0.063536,0.028426,0.533124}%
\pgfsetfillcolor{currentfill}%
\pgfsetlinewidth{0.000000pt}%
\definecolor{currentstroke}{rgb}{0.000000,0.000000,0.000000}%
\pgfsetstrokecolor{currentstroke}%
\pgfsetdash{}{0pt}%
\pgfpathmoveto{\pgfqpoint{4.255037in}{0.371202in}}%
\pgfpathlineto{\pgfqpoint{4.427120in}{0.371202in}}%
\pgfpathlineto{\pgfqpoint{4.427120in}{0.384646in}}%
\pgfpathlineto{\pgfqpoint{4.255037in}{0.384646in}}%
\pgfpathlineto{\pgfqpoint{4.255037in}{0.371202in}}%
\pgfusepath{fill}%
\end{pgfscope}%
\begin{pgfscope}%
\pgfpathrectangle{\pgfqpoint{4.255037in}{0.357758in}}{\pgfqpoint{0.172083in}{3.441662in}}%
\pgfusepath{clip}%
\pgfsetbuttcap%
\pgfsetroundjoin%
\definecolor{currentfill}{rgb}{0.075353,0.027206,0.538007}%
\pgfsetfillcolor{currentfill}%
\pgfsetlinewidth{0.000000pt}%
\definecolor{currentstroke}{rgb}{0.000000,0.000000,0.000000}%
\pgfsetstrokecolor{currentstroke}%
\pgfsetdash{}{0pt}%
\pgfpathmoveto{\pgfqpoint{4.255037in}{0.384646in}}%
\pgfpathlineto{\pgfqpoint{4.427120in}{0.384646in}}%
\pgfpathlineto{\pgfqpoint{4.427120in}{0.398090in}}%
\pgfpathlineto{\pgfqpoint{4.255037in}{0.398090in}}%
\pgfpathlineto{\pgfqpoint{4.255037in}{0.384646in}}%
\pgfusepath{fill}%
\end{pgfscope}%
\begin{pgfscope}%
\pgfpathrectangle{\pgfqpoint{4.255037in}{0.357758in}}{\pgfqpoint{0.172083in}{3.441662in}}%
\pgfusepath{clip}%
\pgfsetbuttcap%
\pgfsetroundjoin%
\definecolor{currentfill}{rgb}{0.086222,0.026125,0.542658}%
\pgfsetfillcolor{currentfill}%
\pgfsetlinewidth{0.000000pt}%
\definecolor{currentstroke}{rgb}{0.000000,0.000000,0.000000}%
\pgfsetstrokecolor{currentstroke}%
\pgfsetdash{}{0pt}%
\pgfpathmoveto{\pgfqpoint{4.255037in}{0.398090in}}%
\pgfpathlineto{\pgfqpoint{4.427120in}{0.398090in}}%
\pgfpathlineto{\pgfqpoint{4.427120in}{0.411534in}}%
\pgfpathlineto{\pgfqpoint{4.255037in}{0.411534in}}%
\pgfpathlineto{\pgfqpoint{4.255037in}{0.398090in}}%
\pgfusepath{fill}%
\end{pgfscope}%
\begin{pgfscope}%
\pgfpathrectangle{\pgfqpoint{4.255037in}{0.357758in}}{\pgfqpoint{0.172083in}{3.441662in}}%
\pgfusepath{clip}%
\pgfsetbuttcap%
\pgfsetroundjoin%
\definecolor{currentfill}{rgb}{0.096379,0.025165,0.547103}%
\pgfsetfillcolor{currentfill}%
\pgfsetlinewidth{0.000000pt}%
\definecolor{currentstroke}{rgb}{0.000000,0.000000,0.000000}%
\pgfsetstrokecolor{currentstroke}%
\pgfsetdash{}{0pt}%
\pgfpathmoveto{\pgfqpoint{4.255037in}{0.411534in}}%
\pgfpathlineto{\pgfqpoint{4.427120in}{0.411534in}}%
\pgfpathlineto{\pgfqpoint{4.427120in}{0.424978in}}%
\pgfpathlineto{\pgfqpoint{4.255037in}{0.424978in}}%
\pgfpathlineto{\pgfqpoint{4.255037in}{0.411534in}}%
\pgfusepath{fill}%
\end{pgfscope}%
\begin{pgfscope}%
\pgfpathrectangle{\pgfqpoint{4.255037in}{0.357758in}}{\pgfqpoint{0.172083in}{3.441662in}}%
\pgfusepath{clip}%
\pgfsetbuttcap%
\pgfsetroundjoin%
\definecolor{currentfill}{rgb}{0.105980,0.024309,0.551368}%
\pgfsetfillcolor{currentfill}%
\pgfsetlinewidth{0.000000pt}%
\definecolor{currentstroke}{rgb}{0.000000,0.000000,0.000000}%
\pgfsetstrokecolor{currentstroke}%
\pgfsetdash{}{0pt}%
\pgfpathmoveto{\pgfqpoint{4.255037in}{0.424978in}}%
\pgfpathlineto{\pgfqpoint{4.427120in}{0.424978in}}%
\pgfpathlineto{\pgfqpoint{4.427120in}{0.438422in}}%
\pgfpathlineto{\pgfqpoint{4.255037in}{0.438422in}}%
\pgfpathlineto{\pgfqpoint{4.255037in}{0.424978in}}%
\pgfusepath{fill}%
\end{pgfscope}%
\begin{pgfscope}%
\pgfpathrectangle{\pgfqpoint{4.255037in}{0.357758in}}{\pgfqpoint{0.172083in}{3.441662in}}%
\pgfusepath{clip}%
\pgfsetbuttcap%
\pgfsetroundjoin%
\definecolor{currentfill}{rgb}{0.115124,0.023556,0.555468}%
\pgfsetfillcolor{currentfill}%
\pgfsetlinewidth{0.000000pt}%
\definecolor{currentstroke}{rgb}{0.000000,0.000000,0.000000}%
\pgfsetstrokecolor{currentstroke}%
\pgfsetdash{}{0pt}%
\pgfpathmoveto{\pgfqpoint{4.255037in}{0.438422in}}%
\pgfpathlineto{\pgfqpoint{4.427120in}{0.438422in}}%
\pgfpathlineto{\pgfqpoint{4.427120in}{0.451866in}}%
\pgfpathlineto{\pgfqpoint{4.255037in}{0.451866in}}%
\pgfpathlineto{\pgfqpoint{4.255037in}{0.438422in}}%
\pgfusepath{fill}%
\end{pgfscope}%
\begin{pgfscope}%
\pgfpathrectangle{\pgfqpoint{4.255037in}{0.357758in}}{\pgfqpoint{0.172083in}{3.441662in}}%
\pgfusepath{clip}%
\pgfsetbuttcap%
\pgfsetroundjoin%
\definecolor{currentfill}{rgb}{0.123903,0.022878,0.559423}%
\pgfsetfillcolor{currentfill}%
\pgfsetlinewidth{0.000000pt}%
\definecolor{currentstroke}{rgb}{0.000000,0.000000,0.000000}%
\pgfsetstrokecolor{currentstroke}%
\pgfsetdash{}{0pt}%
\pgfpathmoveto{\pgfqpoint{4.255037in}{0.451866in}}%
\pgfpathlineto{\pgfqpoint{4.427120in}{0.451866in}}%
\pgfpathlineto{\pgfqpoint{4.427120in}{0.465310in}}%
\pgfpathlineto{\pgfqpoint{4.255037in}{0.465310in}}%
\pgfpathlineto{\pgfqpoint{4.255037in}{0.451866in}}%
\pgfusepath{fill}%
\end{pgfscope}%
\begin{pgfscope}%
\pgfpathrectangle{\pgfqpoint{4.255037in}{0.357758in}}{\pgfqpoint{0.172083in}{3.441662in}}%
\pgfusepath{clip}%
\pgfsetbuttcap%
\pgfsetroundjoin%
\definecolor{currentfill}{rgb}{0.132381,0.022258,0.563250}%
\pgfsetfillcolor{currentfill}%
\pgfsetlinewidth{0.000000pt}%
\definecolor{currentstroke}{rgb}{0.000000,0.000000,0.000000}%
\pgfsetstrokecolor{currentstroke}%
\pgfsetdash{}{0pt}%
\pgfpathmoveto{\pgfqpoint{4.255037in}{0.465310in}}%
\pgfpathlineto{\pgfqpoint{4.427120in}{0.465310in}}%
\pgfpathlineto{\pgfqpoint{4.427120in}{0.478754in}}%
\pgfpathlineto{\pgfqpoint{4.255037in}{0.478754in}}%
\pgfpathlineto{\pgfqpoint{4.255037in}{0.465310in}}%
\pgfusepath{fill}%
\end{pgfscope}%
\begin{pgfscope}%
\pgfpathrectangle{\pgfqpoint{4.255037in}{0.357758in}}{\pgfqpoint{0.172083in}{3.441662in}}%
\pgfusepath{clip}%
\pgfsetbuttcap%
\pgfsetroundjoin%
\definecolor{currentfill}{rgb}{0.140603,0.021687,0.566959}%
\pgfsetfillcolor{currentfill}%
\pgfsetlinewidth{0.000000pt}%
\definecolor{currentstroke}{rgb}{0.000000,0.000000,0.000000}%
\pgfsetstrokecolor{currentstroke}%
\pgfsetdash{}{0pt}%
\pgfpathmoveto{\pgfqpoint{4.255037in}{0.478754in}}%
\pgfpathlineto{\pgfqpoint{4.427120in}{0.478754in}}%
\pgfpathlineto{\pgfqpoint{4.427120in}{0.492198in}}%
\pgfpathlineto{\pgfqpoint{4.255037in}{0.492198in}}%
\pgfpathlineto{\pgfqpoint{4.255037in}{0.478754in}}%
\pgfusepath{fill}%
\end{pgfscope}%
\begin{pgfscope}%
\pgfpathrectangle{\pgfqpoint{4.255037in}{0.357758in}}{\pgfqpoint{0.172083in}{3.441662in}}%
\pgfusepath{clip}%
\pgfsetbuttcap%
\pgfsetroundjoin%
\definecolor{currentfill}{rgb}{0.148607,0.021154,0.570562}%
\pgfsetfillcolor{currentfill}%
\pgfsetlinewidth{0.000000pt}%
\definecolor{currentstroke}{rgb}{0.000000,0.000000,0.000000}%
\pgfsetstrokecolor{currentstroke}%
\pgfsetdash{}{0pt}%
\pgfpathmoveto{\pgfqpoint{4.255037in}{0.492198in}}%
\pgfpathlineto{\pgfqpoint{4.427120in}{0.492198in}}%
\pgfpathlineto{\pgfqpoint{4.427120in}{0.505642in}}%
\pgfpathlineto{\pgfqpoint{4.255037in}{0.505642in}}%
\pgfpathlineto{\pgfqpoint{4.255037in}{0.492198in}}%
\pgfusepath{fill}%
\end{pgfscope}%
\begin{pgfscope}%
\pgfpathrectangle{\pgfqpoint{4.255037in}{0.357758in}}{\pgfqpoint{0.172083in}{3.441662in}}%
\pgfusepath{clip}%
\pgfsetbuttcap%
\pgfsetroundjoin%
\definecolor{currentfill}{rgb}{0.156421,0.020651,0.574065}%
\pgfsetfillcolor{currentfill}%
\pgfsetlinewidth{0.000000pt}%
\definecolor{currentstroke}{rgb}{0.000000,0.000000,0.000000}%
\pgfsetstrokecolor{currentstroke}%
\pgfsetdash{}{0pt}%
\pgfpathmoveto{\pgfqpoint{4.255037in}{0.505642in}}%
\pgfpathlineto{\pgfqpoint{4.427120in}{0.505642in}}%
\pgfpathlineto{\pgfqpoint{4.427120in}{0.519086in}}%
\pgfpathlineto{\pgfqpoint{4.255037in}{0.519086in}}%
\pgfpathlineto{\pgfqpoint{4.255037in}{0.505642in}}%
\pgfusepath{fill}%
\end{pgfscope}%
\begin{pgfscope}%
\pgfpathrectangle{\pgfqpoint{4.255037in}{0.357758in}}{\pgfqpoint{0.172083in}{3.441662in}}%
\pgfusepath{clip}%
\pgfsetbuttcap%
\pgfsetroundjoin%
\definecolor{currentfill}{rgb}{0.164070,0.020171,0.577478}%
\pgfsetfillcolor{currentfill}%
\pgfsetlinewidth{0.000000pt}%
\definecolor{currentstroke}{rgb}{0.000000,0.000000,0.000000}%
\pgfsetstrokecolor{currentstroke}%
\pgfsetdash{}{0pt}%
\pgfpathmoveto{\pgfqpoint{4.255037in}{0.519086in}}%
\pgfpathlineto{\pgfqpoint{4.427120in}{0.519086in}}%
\pgfpathlineto{\pgfqpoint{4.427120in}{0.532530in}}%
\pgfpathlineto{\pgfqpoint{4.255037in}{0.532530in}}%
\pgfpathlineto{\pgfqpoint{4.255037in}{0.519086in}}%
\pgfusepath{fill}%
\end{pgfscope}%
\begin{pgfscope}%
\pgfpathrectangle{\pgfqpoint{4.255037in}{0.357758in}}{\pgfqpoint{0.172083in}{3.441662in}}%
\pgfusepath{clip}%
\pgfsetbuttcap%
\pgfsetroundjoin%
\definecolor{currentfill}{rgb}{0.171574,0.019706,0.580806}%
\pgfsetfillcolor{currentfill}%
\pgfsetlinewidth{0.000000pt}%
\definecolor{currentstroke}{rgb}{0.000000,0.000000,0.000000}%
\pgfsetstrokecolor{currentstroke}%
\pgfsetdash{}{0pt}%
\pgfpathmoveto{\pgfqpoint{4.255037in}{0.532530in}}%
\pgfpathlineto{\pgfqpoint{4.427120in}{0.532530in}}%
\pgfpathlineto{\pgfqpoint{4.427120in}{0.545974in}}%
\pgfpathlineto{\pgfqpoint{4.255037in}{0.545974in}}%
\pgfpathlineto{\pgfqpoint{4.255037in}{0.532530in}}%
\pgfusepath{fill}%
\end{pgfscope}%
\begin{pgfscope}%
\pgfpathrectangle{\pgfqpoint{4.255037in}{0.357758in}}{\pgfqpoint{0.172083in}{3.441662in}}%
\pgfusepath{clip}%
\pgfsetbuttcap%
\pgfsetroundjoin%
\definecolor{currentfill}{rgb}{0.178950,0.019252,0.584054}%
\pgfsetfillcolor{currentfill}%
\pgfsetlinewidth{0.000000pt}%
\definecolor{currentstroke}{rgb}{0.000000,0.000000,0.000000}%
\pgfsetstrokecolor{currentstroke}%
\pgfsetdash{}{0pt}%
\pgfpathmoveto{\pgfqpoint{4.255037in}{0.545974in}}%
\pgfpathlineto{\pgfqpoint{4.427120in}{0.545974in}}%
\pgfpathlineto{\pgfqpoint{4.427120in}{0.559418in}}%
\pgfpathlineto{\pgfqpoint{4.255037in}{0.559418in}}%
\pgfpathlineto{\pgfqpoint{4.255037in}{0.545974in}}%
\pgfusepath{fill}%
\end{pgfscope}%
\begin{pgfscope}%
\pgfpathrectangle{\pgfqpoint{4.255037in}{0.357758in}}{\pgfqpoint{0.172083in}{3.441662in}}%
\pgfusepath{clip}%
\pgfsetbuttcap%
\pgfsetroundjoin%
\definecolor{currentfill}{rgb}{0.186213,0.018803,0.587228}%
\pgfsetfillcolor{currentfill}%
\pgfsetlinewidth{0.000000pt}%
\definecolor{currentstroke}{rgb}{0.000000,0.000000,0.000000}%
\pgfsetstrokecolor{currentstroke}%
\pgfsetdash{}{0pt}%
\pgfpathmoveto{\pgfqpoint{4.255037in}{0.559418in}}%
\pgfpathlineto{\pgfqpoint{4.427120in}{0.559418in}}%
\pgfpathlineto{\pgfqpoint{4.427120in}{0.572862in}}%
\pgfpathlineto{\pgfqpoint{4.255037in}{0.572862in}}%
\pgfpathlineto{\pgfqpoint{4.255037in}{0.559418in}}%
\pgfusepath{fill}%
\end{pgfscope}%
\begin{pgfscope}%
\pgfpathrectangle{\pgfqpoint{4.255037in}{0.357758in}}{\pgfqpoint{0.172083in}{3.441662in}}%
\pgfusepath{clip}%
\pgfsetbuttcap%
\pgfsetroundjoin%
\definecolor{currentfill}{rgb}{0.193374,0.018354,0.590330}%
\pgfsetfillcolor{currentfill}%
\pgfsetlinewidth{0.000000pt}%
\definecolor{currentstroke}{rgb}{0.000000,0.000000,0.000000}%
\pgfsetstrokecolor{currentstroke}%
\pgfsetdash{}{0pt}%
\pgfpathmoveto{\pgfqpoint{4.255037in}{0.572862in}}%
\pgfpathlineto{\pgfqpoint{4.427120in}{0.572862in}}%
\pgfpathlineto{\pgfqpoint{4.427120in}{0.586306in}}%
\pgfpathlineto{\pgfqpoint{4.255037in}{0.586306in}}%
\pgfpathlineto{\pgfqpoint{4.255037in}{0.572862in}}%
\pgfusepath{fill}%
\end{pgfscope}%
\begin{pgfscope}%
\pgfpathrectangle{\pgfqpoint{4.255037in}{0.357758in}}{\pgfqpoint{0.172083in}{3.441662in}}%
\pgfusepath{clip}%
\pgfsetbuttcap%
\pgfsetroundjoin%
\definecolor{currentfill}{rgb}{0.200445,0.017902,0.593364}%
\pgfsetfillcolor{currentfill}%
\pgfsetlinewidth{0.000000pt}%
\definecolor{currentstroke}{rgb}{0.000000,0.000000,0.000000}%
\pgfsetstrokecolor{currentstroke}%
\pgfsetdash{}{0pt}%
\pgfpathmoveto{\pgfqpoint{4.255037in}{0.586306in}}%
\pgfpathlineto{\pgfqpoint{4.427120in}{0.586306in}}%
\pgfpathlineto{\pgfqpoint{4.427120in}{0.599750in}}%
\pgfpathlineto{\pgfqpoint{4.255037in}{0.599750in}}%
\pgfpathlineto{\pgfqpoint{4.255037in}{0.586306in}}%
\pgfusepath{fill}%
\end{pgfscope}%
\begin{pgfscope}%
\pgfpathrectangle{\pgfqpoint{4.255037in}{0.357758in}}{\pgfqpoint{0.172083in}{3.441662in}}%
\pgfusepath{clip}%
\pgfsetbuttcap%
\pgfsetroundjoin%
\definecolor{currentfill}{rgb}{0.207435,0.017442,0.596333}%
\pgfsetfillcolor{currentfill}%
\pgfsetlinewidth{0.000000pt}%
\definecolor{currentstroke}{rgb}{0.000000,0.000000,0.000000}%
\pgfsetstrokecolor{currentstroke}%
\pgfsetdash{}{0pt}%
\pgfpathmoveto{\pgfqpoint{4.255037in}{0.599750in}}%
\pgfpathlineto{\pgfqpoint{4.427120in}{0.599750in}}%
\pgfpathlineto{\pgfqpoint{4.427120in}{0.613194in}}%
\pgfpathlineto{\pgfqpoint{4.255037in}{0.613194in}}%
\pgfpathlineto{\pgfqpoint{4.255037in}{0.599750in}}%
\pgfusepath{fill}%
\end{pgfscope}%
\begin{pgfscope}%
\pgfpathrectangle{\pgfqpoint{4.255037in}{0.357758in}}{\pgfqpoint{0.172083in}{3.441662in}}%
\pgfusepath{clip}%
\pgfsetbuttcap%
\pgfsetroundjoin%
\definecolor{currentfill}{rgb}{0.214350,0.016973,0.599239}%
\pgfsetfillcolor{currentfill}%
\pgfsetlinewidth{0.000000pt}%
\definecolor{currentstroke}{rgb}{0.000000,0.000000,0.000000}%
\pgfsetstrokecolor{currentstroke}%
\pgfsetdash{}{0pt}%
\pgfpathmoveto{\pgfqpoint{4.255037in}{0.613194in}}%
\pgfpathlineto{\pgfqpoint{4.427120in}{0.613194in}}%
\pgfpathlineto{\pgfqpoint{4.427120in}{0.626638in}}%
\pgfpathlineto{\pgfqpoint{4.255037in}{0.626638in}}%
\pgfpathlineto{\pgfqpoint{4.255037in}{0.613194in}}%
\pgfusepath{fill}%
\end{pgfscope}%
\begin{pgfscope}%
\pgfpathrectangle{\pgfqpoint{4.255037in}{0.357758in}}{\pgfqpoint{0.172083in}{3.441662in}}%
\pgfusepath{clip}%
\pgfsetbuttcap%
\pgfsetroundjoin%
\definecolor{currentfill}{rgb}{0.221197,0.016497,0.602083}%
\pgfsetfillcolor{currentfill}%
\pgfsetlinewidth{0.000000pt}%
\definecolor{currentstroke}{rgb}{0.000000,0.000000,0.000000}%
\pgfsetstrokecolor{currentstroke}%
\pgfsetdash{}{0pt}%
\pgfpathmoveto{\pgfqpoint{4.255037in}{0.626638in}}%
\pgfpathlineto{\pgfqpoint{4.427120in}{0.626638in}}%
\pgfpathlineto{\pgfqpoint{4.427120in}{0.640082in}}%
\pgfpathlineto{\pgfqpoint{4.255037in}{0.640082in}}%
\pgfpathlineto{\pgfqpoint{4.255037in}{0.626638in}}%
\pgfusepath{fill}%
\end{pgfscope}%
\begin{pgfscope}%
\pgfpathrectangle{\pgfqpoint{4.255037in}{0.357758in}}{\pgfqpoint{0.172083in}{3.441662in}}%
\pgfusepath{clip}%
\pgfsetbuttcap%
\pgfsetroundjoin%
\definecolor{currentfill}{rgb}{0.227983,0.016007,0.604867}%
\pgfsetfillcolor{currentfill}%
\pgfsetlinewidth{0.000000pt}%
\definecolor{currentstroke}{rgb}{0.000000,0.000000,0.000000}%
\pgfsetstrokecolor{currentstroke}%
\pgfsetdash{}{0pt}%
\pgfpathmoveto{\pgfqpoint{4.255037in}{0.640082in}}%
\pgfpathlineto{\pgfqpoint{4.427120in}{0.640082in}}%
\pgfpathlineto{\pgfqpoint{4.427120in}{0.653526in}}%
\pgfpathlineto{\pgfqpoint{4.255037in}{0.653526in}}%
\pgfpathlineto{\pgfqpoint{4.255037in}{0.640082in}}%
\pgfusepath{fill}%
\end{pgfscope}%
\begin{pgfscope}%
\pgfpathrectangle{\pgfqpoint{4.255037in}{0.357758in}}{\pgfqpoint{0.172083in}{3.441662in}}%
\pgfusepath{clip}%
\pgfsetbuttcap%
\pgfsetroundjoin%
\definecolor{currentfill}{rgb}{0.234715,0.015502,0.607592}%
\pgfsetfillcolor{currentfill}%
\pgfsetlinewidth{0.000000pt}%
\definecolor{currentstroke}{rgb}{0.000000,0.000000,0.000000}%
\pgfsetstrokecolor{currentstroke}%
\pgfsetdash{}{0pt}%
\pgfpathmoveto{\pgfqpoint{4.255037in}{0.653526in}}%
\pgfpathlineto{\pgfqpoint{4.427120in}{0.653526in}}%
\pgfpathlineto{\pgfqpoint{4.427120in}{0.666970in}}%
\pgfpathlineto{\pgfqpoint{4.255037in}{0.666970in}}%
\pgfpathlineto{\pgfqpoint{4.255037in}{0.653526in}}%
\pgfusepath{fill}%
\end{pgfscope}%
\begin{pgfscope}%
\pgfpathrectangle{\pgfqpoint{4.255037in}{0.357758in}}{\pgfqpoint{0.172083in}{3.441662in}}%
\pgfusepath{clip}%
\pgfsetbuttcap%
\pgfsetroundjoin%
\definecolor{currentfill}{rgb}{0.241396,0.014979,0.610259}%
\pgfsetfillcolor{currentfill}%
\pgfsetlinewidth{0.000000pt}%
\definecolor{currentstroke}{rgb}{0.000000,0.000000,0.000000}%
\pgfsetstrokecolor{currentstroke}%
\pgfsetdash{}{0pt}%
\pgfpathmoveto{\pgfqpoint{4.255037in}{0.666970in}}%
\pgfpathlineto{\pgfqpoint{4.427120in}{0.666970in}}%
\pgfpathlineto{\pgfqpoint{4.427120in}{0.680414in}}%
\pgfpathlineto{\pgfqpoint{4.255037in}{0.680414in}}%
\pgfpathlineto{\pgfqpoint{4.255037in}{0.666970in}}%
\pgfusepath{fill}%
\end{pgfscope}%
\begin{pgfscope}%
\pgfpathrectangle{\pgfqpoint{4.255037in}{0.357758in}}{\pgfqpoint{0.172083in}{3.441662in}}%
\pgfusepath{clip}%
\pgfsetbuttcap%
\pgfsetroundjoin%
\definecolor{currentfill}{rgb}{0.248032,0.014439,0.612868}%
\pgfsetfillcolor{currentfill}%
\pgfsetlinewidth{0.000000pt}%
\definecolor{currentstroke}{rgb}{0.000000,0.000000,0.000000}%
\pgfsetstrokecolor{currentstroke}%
\pgfsetdash{}{0pt}%
\pgfpathmoveto{\pgfqpoint{4.255037in}{0.680414in}}%
\pgfpathlineto{\pgfqpoint{4.427120in}{0.680414in}}%
\pgfpathlineto{\pgfqpoint{4.427120in}{0.693858in}}%
\pgfpathlineto{\pgfqpoint{4.255037in}{0.693858in}}%
\pgfpathlineto{\pgfqpoint{4.255037in}{0.680414in}}%
\pgfusepath{fill}%
\end{pgfscope}%
\begin{pgfscope}%
\pgfpathrectangle{\pgfqpoint{4.255037in}{0.357758in}}{\pgfqpoint{0.172083in}{3.441662in}}%
\pgfusepath{clip}%
\pgfsetbuttcap%
\pgfsetroundjoin%
\definecolor{currentfill}{rgb}{0.254627,0.013882,0.615419}%
\pgfsetfillcolor{currentfill}%
\pgfsetlinewidth{0.000000pt}%
\definecolor{currentstroke}{rgb}{0.000000,0.000000,0.000000}%
\pgfsetstrokecolor{currentstroke}%
\pgfsetdash{}{0pt}%
\pgfpathmoveto{\pgfqpoint{4.255037in}{0.693858in}}%
\pgfpathlineto{\pgfqpoint{4.427120in}{0.693858in}}%
\pgfpathlineto{\pgfqpoint{4.427120in}{0.707302in}}%
\pgfpathlineto{\pgfqpoint{4.255037in}{0.707302in}}%
\pgfpathlineto{\pgfqpoint{4.255037in}{0.693858in}}%
\pgfusepath{fill}%
\end{pgfscope}%
\begin{pgfscope}%
\pgfpathrectangle{\pgfqpoint{4.255037in}{0.357758in}}{\pgfqpoint{0.172083in}{3.441662in}}%
\pgfusepath{clip}%
\pgfsetbuttcap%
\pgfsetroundjoin%
\definecolor{currentfill}{rgb}{0.261183,0.013308,0.617911}%
\pgfsetfillcolor{currentfill}%
\pgfsetlinewidth{0.000000pt}%
\definecolor{currentstroke}{rgb}{0.000000,0.000000,0.000000}%
\pgfsetstrokecolor{currentstroke}%
\pgfsetdash{}{0pt}%
\pgfpathmoveto{\pgfqpoint{4.255037in}{0.707302in}}%
\pgfpathlineto{\pgfqpoint{4.427120in}{0.707302in}}%
\pgfpathlineto{\pgfqpoint{4.427120in}{0.720746in}}%
\pgfpathlineto{\pgfqpoint{4.255037in}{0.720746in}}%
\pgfpathlineto{\pgfqpoint{4.255037in}{0.707302in}}%
\pgfusepath{fill}%
\end{pgfscope}%
\begin{pgfscope}%
\pgfpathrectangle{\pgfqpoint{4.255037in}{0.357758in}}{\pgfqpoint{0.172083in}{3.441662in}}%
\pgfusepath{clip}%
\pgfsetbuttcap%
\pgfsetroundjoin%
\definecolor{currentfill}{rgb}{0.267703,0.012716,0.620346}%
\pgfsetfillcolor{currentfill}%
\pgfsetlinewidth{0.000000pt}%
\definecolor{currentstroke}{rgb}{0.000000,0.000000,0.000000}%
\pgfsetstrokecolor{currentstroke}%
\pgfsetdash{}{0pt}%
\pgfpathmoveto{\pgfqpoint{4.255037in}{0.720746in}}%
\pgfpathlineto{\pgfqpoint{4.427120in}{0.720746in}}%
\pgfpathlineto{\pgfqpoint{4.427120in}{0.734190in}}%
\pgfpathlineto{\pgfqpoint{4.255037in}{0.734190in}}%
\pgfpathlineto{\pgfqpoint{4.255037in}{0.720746in}}%
\pgfusepath{fill}%
\end{pgfscope}%
\begin{pgfscope}%
\pgfpathrectangle{\pgfqpoint{4.255037in}{0.357758in}}{\pgfqpoint{0.172083in}{3.441662in}}%
\pgfusepath{clip}%
\pgfsetbuttcap%
\pgfsetroundjoin%
\definecolor{currentfill}{rgb}{0.274191,0.012109,0.622722}%
\pgfsetfillcolor{currentfill}%
\pgfsetlinewidth{0.000000pt}%
\definecolor{currentstroke}{rgb}{0.000000,0.000000,0.000000}%
\pgfsetstrokecolor{currentstroke}%
\pgfsetdash{}{0pt}%
\pgfpathmoveto{\pgfqpoint{4.255037in}{0.734190in}}%
\pgfpathlineto{\pgfqpoint{4.427120in}{0.734190in}}%
\pgfpathlineto{\pgfqpoint{4.427120in}{0.747634in}}%
\pgfpathlineto{\pgfqpoint{4.255037in}{0.747634in}}%
\pgfpathlineto{\pgfqpoint{4.255037in}{0.734190in}}%
\pgfusepath{fill}%
\end{pgfscope}%
\begin{pgfscope}%
\pgfpathrectangle{\pgfqpoint{4.255037in}{0.357758in}}{\pgfqpoint{0.172083in}{3.441662in}}%
\pgfusepath{clip}%
\pgfsetbuttcap%
\pgfsetroundjoin%
\definecolor{currentfill}{rgb}{0.280648,0.011488,0.625038}%
\pgfsetfillcolor{currentfill}%
\pgfsetlinewidth{0.000000pt}%
\definecolor{currentstroke}{rgb}{0.000000,0.000000,0.000000}%
\pgfsetstrokecolor{currentstroke}%
\pgfsetdash{}{0pt}%
\pgfpathmoveto{\pgfqpoint{4.255037in}{0.747634in}}%
\pgfpathlineto{\pgfqpoint{4.427120in}{0.747634in}}%
\pgfpathlineto{\pgfqpoint{4.427120in}{0.761078in}}%
\pgfpathlineto{\pgfqpoint{4.255037in}{0.761078in}}%
\pgfpathlineto{\pgfqpoint{4.255037in}{0.747634in}}%
\pgfusepath{fill}%
\end{pgfscope}%
\begin{pgfscope}%
\pgfpathrectangle{\pgfqpoint{4.255037in}{0.357758in}}{\pgfqpoint{0.172083in}{3.441662in}}%
\pgfusepath{clip}%
\pgfsetbuttcap%
\pgfsetroundjoin%
\definecolor{currentfill}{rgb}{0.287076,0.010855,0.627295}%
\pgfsetfillcolor{currentfill}%
\pgfsetlinewidth{0.000000pt}%
\definecolor{currentstroke}{rgb}{0.000000,0.000000,0.000000}%
\pgfsetstrokecolor{currentstroke}%
\pgfsetdash{}{0pt}%
\pgfpathmoveto{\pgfqpoint{4.255037in}{0.761078in}}%
\pgfpathlineto{\pgfqpoint{4.427120in}{0.761078in}}%
\pgfpathlineto{\pgfqpoint{4.427120in}{0.774522in}}%
\pgfpathlineto{\pgfqpoint{4.255037in}{0.774522in}}%
\pgfpathlineto{\pgfqpoint{4.255037in}{0.761078in}}%
\pgfusepath{fill}%
\end{pgfscope}%
\begin{pgfscope}%
\pgfpathrectangle{\pgfqpoint{4.255037in}{0.357758in}}{\pgfqpoint{0.172083in}{3.441662in}}%
\pgfusepath{clip}%
\pgfsetbuttcap%
\pgfsetroundjoin%
\definecolor{currentfill}{rgb}{0.293478,0.010213,0.629490}%
\pgfsetfillcolor{currentfill}%
\pgfsetlinewidth{0.000000pt}%
\definecolor{currentstroke}{rgb}{0.000000,0.000000,0.000000}%
\pgfsetstrokecolor{currentstroke}%
\pgfsetdash{}{0pt}%
\pgfpathmoveto{\pgfqpoint{4.255037in}{0.774522in}}%
\pgfpathlineto{\pgfqpoint{4.427120in}{0.774522in}}%
\pgfpathlineto{\pgfqpoint{4.427120in}{0.787966in}}%
\pgfpathlineto{\pgfqpoint{4.255037in}{0.787966in}}%
\pgfpathlineto{\pgfqpoint{4.255037in}{0.774522in}}%
\pgfusepath{fill}%
\end{pgfscope}%
\begin{pgfscope}%
\pgfpathrectangle{\pgfqpoint{4.255037in}{0.357758in}}{\pgfqpoint{0.172083in}{3.441662in}}%
\pgfusepath{clip}%
\pgfsetbuttcap%
\pgfsetroundjoin%
\definecolor{currentfill}{rgb}{0.299855,0.009561,0.631624}%
\pgfsetfillcolor{currentfill}%
\pgfsetlinewidth{0.000000pt}%
\definecolor{currentstroke}{rgb}{0.000000,0.000000,0.000000}%
\pgfsetstrokecolor{currentstroke}%
\pgfsetdash{}{0pt}%
\pgfpathmoveto{\pgfqpoint{4.255037in}{0.787966in}}%
\pgfpathlineto{\pgfqpoint{4.427120in}{0.787966in}}%
\pgfpathlineto{\pgfqpoint{4.427120in}{0.801410in}}%
\pgfpathlineto{\pgfqpoint{4.255037in}{0.801410in}}%
\pgfpathlineto{\pgfqpoint{4.255037in}{0.787966in}}%
\pgfusepath{fill}%
\end{pgfscope}%
\begin{pgfscope}%
\pgfpathrectangle{\pgfqpoint{4.255037in}{0.357758in}}{\pgfqpoint{0.172083in}{3.441662in}}%
\pgfusepath{clip}%
\pgfsetbuttcap%
\pgfsetroundjoin%
\definecolor{currentfill}{rgb}{0.306210,0.008902,0.633694}%
\pgfsetfillcolor{currentfill}%
\pgfsetlinewidth{0.000000pt}%
\definecolor{currentstroke}{rgb}{0.000000,0.000000,0.000000}%
\pgfsetstrokecolor{currentstroke}%
\pgfsetdash{}{0pt}%
\pgfpathmoveto{\pgfqpoint{4.255037in}{0.801410in}}%
\pgfpathlineto{\pgfqpoint{4.427120in}{0.801410in}}%
\pgfpathlineto{\pgfqpoint{4.427120in}{0.814854in}}%
\pgfpathlineto{\pgfqpoint{4.255037in}{0.814854in}}%
\pgfpathlineto{\pgfqpoint{4.255037in}{0.801410in}}%
\pgfusepath{fill}%
\end{pgfscope}%
\begin{pgfscope}%
\pgfpathrectangle{\pgfqpoint{4.255037in}{0.357758in}}{\pgfqpoint{0.172083in}{3.441662in}}%
\pgfusepath{clip}%
\pgfsetbuttcap%
\pgfsetroundjoin%
\definecolor{currentfill}{rgb}{0.312543,0.008239,0.635700}%
\pgfsetfillcolor{currentfill}%
\pgfsetlinewidth{0.000000pt}%
\definecolor{currentstroke}{rgb}{0.000000,0.000000,0.000000}%
\pgfsetstrokecolor{currentstroke}%
\pgfsetdash{}{0pt}%
\pgfpathmoveto{\pgfqpoint{4.255037in}{0.814854in}}%
\pgfpathlineto{\pgfqpoint{4.427120in}{0.814854in}}%
\pgfpathlineto{\pgfqpoint{4.427120in}{0.828298in}}%
\pgfpathlineto{\pgfqpoint{4.255037in}{0.828298in}}%
\pgfpathlineto{\pgfqpoint{4.255037in}{0.814854in}}%
\pgfusepath{fill}%
\end{pgfscope}%
\begin{pgfscope}%
\pgfpathrectangle{\pgfqpoint{4.255037in}{0.357758in}}{\pgfqpoint{0.172083in}{3.441662in}}%
\pgfusepath{clip}%
\pgfsetbuttcap%
\pgfsetroundjoin%
\definecolor{currentfill}{rgb}{0.318856,0.007576,0.637640}%
\pgfsetfillcolor{currentfill}%
\pgfsetlinewidth{0.000000pt}%
\definecolor{currentstroke}{rgb}{0.000000,0.000000,0.000000}%
\pgfsetstrokecolor{currentstroke}%
\pgfsetdash{}{0pt}%
\pgfpathmoveto{\pgfqpoint{4.255037in}{0.828298in}}%
\pgfpathlineto{\pgfqpoint{4.427120in}{0.828298in}}%
\pgfpathlineto{\pgfqpoint{4.427120in}{0.841742in}}%
\pgfpathlineto{\pgfqpoint{4.255037in}{0.841742in}}%
\pgfpathlineto{\pgfqpoint{4.255037in}{0.828298in}}%
\pgfusepath{fill}%
\end{pgfscope}%
\begin{pgfscope}%
\pgfpathrectangle{\pgfqpoint{4.255037in}{0.357758in}}{\pgfqpoint{0.172083in}{3.441662in}}%
\pgfusepath{clip}%
\pgfsetbuttcap%
\pgfsetroundjoin%
\definecolor{currentfill}{rgb}{0.325150,0.006915,0.639512}%
\pgfsetfillcolor{currentfill}%
\pgfsetlinewidth{0.000000pt}%
\definecolor{currentstroke}{rgb}{0.000000,0.000000,0.000000}%
\pgfsetstrokecolor{currentstroke}%
\pgfsetdash{}{0pt}%
\pgfpathmoveto{\pgfqpoint{4.255037in}{0.841742in}}%
\pgfpathlineto{\pgfqpoint{4.427120in}{0.841742in}}%
\pgfpathlineto{\pgfqpoint{4.427120in}{0.855186in}}%
\pgfpathlineto{\pgfqpoint{4.255037in}{0.855186in}}%
\pgfpathlineto{\pgfqpoint{4.255037in}{0.841742in}}%
\pgfusepath{fill}%
\end{pgfscope}%
\begin{pgfscope}%
\pgfpathrectangle{\pgfqpoint{4.255037in}{0.357758in}}{\pgfqpoint{0.172083in}{3.441662in}}%
\pgfusepath{clip}%
\pgfsetbuttcap%
\pgfsetroundjoin%
\definecolor{currentfill}{rgb}{0.331426,0.006261,0.641316}%
\pgfsetfillcolor{currentfill}%
\pgfsetlinewidth{0.000000pt}%
\definecolor{currentstroke}{rgb}{0.000000,0.000000,0.000000}%
\pgfsetstrokecolor{currentstroke}%
\pgfsetdash{}{0pt}%
\pgfpathmoveto{\pgfqpoint{4.255037in}{0.855186in}}%
\pgfpathlineto{\pgfqpoint{4.427120in}{0.855186in}}%
\pgfpathlineto{\pgfqpoint{4.427120in}{0.868630in}}%
\pgfpathlineto{\pgfqpoint{4.255037in}{0.868630in}}%
\pgfpathlineto{\pgfqpoint{4.255037in}{0.855186in}}%
\pgfusepath{fill}%
\end{pgfscope}%
\begin{pgfscope}%
\pgfpathrectangle{\pgfqpoint{4.255037in}{0.357758in}}{\pgfqpoint{0.172083in}{3.441662in}}%
\pgfusepath{clip}%
\pgfsetbuttcap%
\pgfsetroundjoin%
\definecolor{currentfill}{rgb}{0.337683,0.005618,0.643049}%
\pgfsetfillcolor{currentfill}%
\pgfsetlinewidth{0.000000pt}%
\definecolor{currentstroke}{rgb}{0.000000,0.000000,0.000000}%
\pgfsetstrokecolor{currentstroke}%
\pgfsetdash{}{0pt}%
\pgfpathmoveto{\pgfqpoint{4.255037in}{0.868630in}}%
\pgfpathlineto{\pgfqpoint{4.427120in}{0.868630in}}%
\pgfpathlineto{\pgfqpoint{4.427120in}{0.882074in}}%
\pgfpathlineto{\pgfqpoint{4.255037in}{0.882074in}}%
\pgfpathlineto{\pgfqpoint{4.255037in}{0.868630in}}%
\pgfusepath{fill}%
\end{pgfscope}%
\begin{pgfscope}%
\pgfpathrectangle{\pgfqpoint{4.255037in}{0.357758in}}{\pgfqpoint{0.172083in}{3.441662in}}%
\pgfusepath{clip}%
\pgfsetbuttcap%
\pgfsetroundjoin%
\definecolor{currentfill}{rgb}{0.343925,0.004991,0.644710}%
\pgfsetfillcolor{currentfill}%
\pgfsetlinewidth{0.000000pt}%
\definecolor{currentstroke}{rgb}{0.000000,0.000000,0.000000}%
\pgfsetstrokecolor{currentstroke}%
\pgfsetdash{}{0pt}%
\pgfpathmoveto{\pgfqpoint{4.255037in}{0.882074in}}%
\pgfpathlineto{\pgfqpoint{4.427120in}{0.882074in}}%
\pgfpathlineto{\pgfqpoint{4.427120in}{0.895518in}}%
\pgfpathlineto{\pgfqpoint{4.255037in}{0.895518in}}%
\pgfpathlineto{\pgfqpoint{4.255037in}{0.882074in}}%
\pgfusepath{fill}%
\end{pgfscope}%
\begin{pgfscope}%
\pgfpathrectangle{\pgfqpoint{4.255037in}{0.357758in}}{\pgfqpoint{0.172083in}{3.441662in}}%
\pgfusepath{clip}%
\pgfsetbuttcap%
\pgfsetroundjoin%
\definecolor{currentfill}{rgb}{0.350150,0.004382,0.646298}%
\pgfsetfillcolor{currentfill}%
\pgfsetlinewidth{0.000000pt}%
\definecolor{currentstroke}{rgb}{0.000000,0.000000,0.000000}%
\pgfsetstrokecolor{currentstroke}%
\pgfsetdash{}{0pt}%
\pgfpathmoveto{\pgfqpoint{4.255037in}{0.895518in}}%
\pgfpathlineto{\pgfqpoint{4.427120in}{0.895518in}}%
\pgfpathlineto{\pgfqpoint{4.427120in}{0.908962in}}%
\pgfpathlineto{\pgfqpoint{4.255037in}{0.908962in}}%
\pgfpathlineto{\pgfqpoint{4.255037in}{0.895518in}}%
\pgfusepath{fill}%
\end{pgfscope}%
\begin{pgfscope}%
\pgfpathrectangle{\pgfqpoint{4.255037in}{0.357758in}}{\pgfqpoint{0.172083in}{3.441662in}}%
\pgfusepath{clip}%
\pgfsetbuttcap%
\pgfsetroundjoin%
\definecolor{currentfill}{rgb}{0.356359,0.003798,0.647810}%
\pgfsetfillcolor{currentfill}%
\pgfsetlinewidth{0.000000pt}%
\definecolor{currentstroke}{rgb}{0.000000,0.000000,0.000000}%
\pgfsetstrokecolor{currentstroke}%
\pgfsetdash{}{0pt}%
\pgfpathmoveto{\pgfqpoint{4.255037in}{0.908962in}}%
\pgfpathlineto{\pgfqpoint{4.427120in}{0.908962in}}%
\pgfpathlineto{\pgfqpoint{4.427120in}{0.922406in}}%
\pgfpathlineto{\pgfqpoint{4.255037in}{0.922406in}}%
\pgfpathlineto{\pgfqpoint{4.255037in}{0.908962in}}%
\pgfusepath{fill}%
\end{pgfscope}%
\begin{pgfscope}%
\pgfpathrectangle{\pgfqpoint{4.255037in}{0.357758in}}{\pgfqpoint{0.172083in}{3.441662in}}%
\pgfusepath{clip}%
\pgfsetbuttcap%
\pgfsetroundjoin%
\definecolor{currentfill}{rgb}{0.362553,0.003243,0.649245}%
\pgfsetfillcolor{currentfill}%
\pgfsetlinewidth{0.000000pt}%
\definecolor{currentstroke}{rgb}{0.000000,0.000000,0.000000}%
\pgfsetstrokecolor{currentstroke}%
\pgfsetdash{}{0pt}%
\pgfpathmoveto{\pgfqpoint{4.255037in}{0.922406in}}%
\pgfpathlineto{\pgfqpoint{4.427120in}{0.922406in}}%
\pgfpathlineto{\pgfqpoint{4.427120in}{0.935850in}}%
\pgfpathlineto{\pgfqpoint{4.255037in}{0.935850in}}%
\pgfpathlineto{\pgfqpoint{4.255037in}{0.922406in}}%
\pgfusepath{fill}%
\end{pgfscope}%
\begin{pgfscope}%
\pgfpathrectangle{\pgfqpoint{4.255037in}{0.357758in}}{\pgfqpoint{0.172083in}{3.441662in}}%
\pgfusepath{clip}%
\pgfsetbuttcap%
\pgfsetroundjoin%
\definecolor{currentfill}{rgb}{0.368733,0.002724,0.650601}%
\pgfsetfillcolor{currentfill}%
\pgfsetlinewidth{0.000000pt}%
\definecolor{currentstroke}{rgb}{0.000000,0.000000,0.000000}%
\pgfsetstrokecolor{currentstroke}%
\pgfsetdash{}{0pt}%
\pgfpathmoveto{\pgfqpoint{4.255037in}{0.935850in}}%
\pgfpathlineto{\pgfqpoint{4.427120in}{0.935850in}}%
\pgfpathlineto{\pgfqpoint{4.427120in}{0.949294in}}%
\pgfpathlineto{\pgfqpoint{4.255037in}{0.949294in}}%
\pgfpathlineto{\pgfqpoint{4.255037in}{0.935850in}}%
\pgfusepath{fill}%
\end{pgfscope}%
\begin{pgfscope}%
\pgfpathrectangle{\pgfqpoint{4.255037in}{0.357758in}}{\pgfqpoint{0.172083in}{3.441662in}}%
\pgfusepath{clip}%
\pgfsetbuttcap%
\pgfsetroundjoin%
\definecolor{currentfill}{rgb}{0.374897,0.002245,0.651876}%
\pgfsetfillcolor{currentfill}%
\pgfsetlinewidth{0.000000pt}%
\definecolor{currentstroke}{rgb}{0.000000,0.000000,0.000000}%
\pgfsetstrokecolor{currentstroke}%
\pgfsetdash{}{0pt}%
\pgfpathmoveto{\pgfqpoint{4.255037in}{0.949294in}}%
\pgfpathlineto{\pgfqpoint{4.427120in}{0.949294in}}%
\pgfpathlineto{\pgfqpoint{4.427120in}{0.962738in}}%
\pgfpathlineto{\pgfqpoint{4.255037in}{0.962738in}}%
\pgfpathlineto{\pgfqpoint{4.255037in}{0.949294in}}%
\pgfusepath{fill}%
\end{pgfscope}%
\begin{pgfscope}%
\pgfpathrectangle{\pgfqpoint{4.255037in}{0.357758in}}{\pgfqpoint{0.172083in}{3.441662in}}%
\pgfusepath{clip}%
\pgfsetbuttcap%
\pgfsetroundjoin%
\definecolor{currentfill}{rgb}{0.381047,0.001814,0.653068}%
\pgfsetfillcolor{currentfill}%
\pgfsetlinewidth{0.000000pt}%
\definecolor{currentstroke}{rgb}{0.000000,0.000000,0.000000}%
\pgfsetstrokecolor{currentstroke}%
\pgfsetdash{}{0pt}%
\pgfpathmoveto{\pgfqpoint{4.255037in}{0.962738in}}%
\pgfpathlineto{\pgfqpoint{4.427120in}{0.962738in}}%
\pgfpathlineto{\pgfqpoint{4.427120in}{0.976182in}}%
\pgfpathlineto{\pgfqpoint{4.255037in}{0.976182in}}%
\pgfpathlineto{\pgfqpoint{4.255037in}{0.962738in}}%
\pgfusepath{fill}%
\end{pgfscope}%
\begin{pgfscope}%
\pgfpathrectangle{\pgfqpoint{4.255037in}{0.357758in}}{\pgfqpoint{0.172083in}{3.441662in}}%
\pgfusepath{clip}%
\pgfsetbuttcap%
\pgfsetroundjoin%
\definecolor{currentfill}{rgb}{0.387183,0.001434,0.654177}%
\pgfsetfillcolor{currentfill}%
\pgfsetlinewidth{0.000000pt}%
\definecolor{currentstroke}{rgb}{0.000000,0.000000,0.000000}%
\pgfsetstrokecolor{currentstroke}%
\pgfsetdash{}{0pt}%
\pgfpathmoveto{\pgfqpoint{4.255037in}{0.976182in}}%
\pgfpathlineto{\pgfqpoint{4.427120in}{0.976182in}}%
\pgfpathlineto{\pgfqpoint{4.427120in}{0.989626in}}%
\pgfpathlineto{\pgfqpoint{4.255037in}{0.989626in}}%
\pgfpathlineto{\pgfqpoint{4.255037in}{0.976182in}}%
\pgfusepath{fill}%
\end{pgfscope}%
\begin{pgfscope}%
\pgfpathrectangle{\pgfqpoint{4.255037in}{0.357758in}}{\pgfqpoint{0.172083in}{3.441662in}}%
\pgfusepath{clip}%
\pgfsetbuttcap%
\pgfsetroundjoin%
\definecolor{currentfill}{rgb}{0.393304,0.001114,0.655199}%
\pgfsetfillcolor{currentfill}%
\pgfsetlinewidth{0.000000pt}%
\definecolor{currentstroke}{rgb}{0.000000,0.000000,0.000000}%
\pgfsetstrokecolor{currentstroke}%
\pgfsetdash{}{0pt}%
\pgfpathmoveto{\pgfqpoint{4.255037in}{0.989626in}}%
\pgfpathlineto{\pgfqpoint{4.427120in}{0.989626in}}%
\pgfpathlineto{\pgfqpoint{4.427120in}{1.003070in}}%
\pgfpathlineto{\pgfqpoint{4.255037in}{1.003070in}}%
\pgfpathlineto{\pgfqpoint{4.255037in}{0.989626in}}%
\pgfusepath{fill}%
\end{pgfscope}%
\begin{pgfscope}%
\pgfpathrectangle{\pgfqpoint{4.255037in}{0.357758in}}{\pgfqpoint{0.172083in}{3.441662in}}%
\pgfusepath{clip}%
\pgfsetbuttcap%
\pgfsetroundjoin%
\definecolor{currentfill}{rgb}{0.399411,0.000859,0.656133}%
\pgfsetfillcolor{currentfill}%
\pgfsetlinewidth{0.000000pt}%
\definecolor{currentstroke}{rgb}{0.000000,0.000000,0.000000}%
\pgfsetstrokecolor{currentstroke}%
\pgfsetdash{}{0pt}%
\pgfpathmoveto{\pgfqpoint{4.255037in}{1.003070in}}%
\pgfpathlineto{\pgfqpoint{4.427120in}{1.003070in}}%
\pgfpathlineto{\pgfqpoint{4.427120in}{1.016514in}}%
\pgfpathlineto{\pgfqpoint{4.255037in}{1.016514in}}%
\pgfpathlineto{\pgfqpoint{4.255037in}{1.003070in}}%
\pgfusepath{fill}%
\end{pgfscope}%
\begin{pgfscope}%
\pgfpathrectangle{\pgfqpoint{4.255037in}{0.357758in}}{\pgfqpoint{0.172083in}{3.441662in}}%
\pgfusepath{clip}%
\pgfsetbuttcap%
\pgfsetroundjoin%
\definecolor{currentfill}{rgb}{0.405503,0.000678,0.656977}%
\pgfsetfillcolor{currentfill}%
\pgfsetlinewidth{0.000000pt}%
\definecolor{currentstroke}{rgb}{0.000000,0.000000,0.000000}%
\pgfsetstrokecolor{currentstroke}%
\pgfsetdash{}{0pt}%
\pgfpathmoveto{\pgfqpoint{4.255037in}{1.016514in}}%
\pgfpathlineto{\pgfqpoint{4.427120in}{1.016514in}}%
\pgfpathlineto{\pgfqpoint{4.427120in}{1.029958in}}%
\pgfpathlineto{\pgfqpoint{4.255037in}{1.029958in}}%
\pgfpathlineto{\pgfqpoint{4.255037in}{1.016514in}}%
\pgfusepath{fill}%
\end{pgfscope}%
\begin{pgfscope}%
\pgfpathrectangle{\pgfqpoint{4.255037in}{0.357758in}}{\pgfqpoint{0.172083in}{3.441662in}}%
\pgfusepath{clip}%
\pgfsetbuttcap%
\pgfsetroundjoin%
\definecolor{currentfill}{rgb}{0.411580,0.000577,0.657730}%
\pgfsetfillcolor{currentfill}%
\pgfsetlinewidth{0.000000pt}%
\definecolor{currentstroke}{rgb}{0.000000,0.000000,0.000000}%
\pgfsetstrokecolor{currentstroke}%
\pgfsetdash{}{0pt}%
\pgfpathmoveto{\pgfqpoint{4.255037in}{1.029958in}}%
\pgfpathlineto{\pgfqpoint{4.427120in}{1.029958in}}%
\pgfpathlineto{\pgfqpoint{4.427120in}{1.043402in}}%
\pgfpathlineto{\pgfqpoint{4.255037in}{1.043402in}}%
\pgfpathlineto{\pgfqpoint{4.255037in}{1.029958in}}%
\pgfusepath{fill}%
\end{pgfscope}%
\begin{pgfscope}%
\pgfpathrectangle{\pgfqpoint{4.255037in}{0.357758in}}{\pgfqpoint{0.172083in}{3.441662in}}%
\pgfusepath{clip}%
\pgfsetbuttcap%
\pgfsetroundjoin%
\definecolor{currentfill}{rgb}{0.417642,0.000564,0.658390}%
\pgfsetfillcolor{currentfill}%
\pgfsetlinewidth{0.000000pt}%
\definecolor{currentstroke}{rgb}{0.000000,0.000000,0.000000}%
\pgfsetstrokecolor{currentstroke}%
\pgfsetdash{}{0pt}%
\pgfpathmoveto{\pgfqpoint{4.255037in}{1.043402in}}%
\pgfpathlineto{\pgfqpoint{4.427120in}{1.043402in}}%
\pgfpathlineto{\pgfqpoint{4.427120in}{1.056846in}}%
\pgfpathlineto{\pgfqpoint{4.255037in}{1.056846in}}%
\pgfpathlineto{\pgfqpoint{4.255037in}{1.043402in}}%
\pgfusepath{fill}%
\end{pgfscope}%
\begin{pgfscope}%
\pgfpathrectangle{\pgfqpoint{4.255037in}{0.357758in}}{\pgfqpoint{0.172083in}{3.441662in}}%
\pgfusepath{clip}%
\pgfsetbuttcap%
\pgfsetroundjoin%
\definecolor{currentfill}{rgb}{0.423689,0.000646,0.658956}%
\pgfsetfillcolor{currentfill}%
\pgfsetlinewidth{0.000000pt}%
\definecolor{currentstroke}{rgb}{0.000000,0.000000,0.000000}%
\pgfsetstrokecolor{currentstroke}%
\pgfsetdash{}{0pt}%
\pgfpathmoveto{\pgfqpoint{4.255037in}{1.056846in}}%
\pgfpathlineto{\pgfqpoint{4.427120in}{1.056846in}}%
\pgfpathlineto{\pgfqpoint{4.427120in}{1.070290in}}%
\pgfpathlineto{\pgfqpoint{4.255037in}{1.070290in}}%
\pgfpathlineto{\pgfqpoint{4.255037in}{1.056846in}}%
\pgfusepath{fill}%
\end{pgfscope}%
\begin{pgfscope}%
\pgfpathrectangle{\pgfqpoint{4.255037in}{0.357758in}}{\pgfqpoint{0.172083in}{3.441662in}}%
\pgfusepath{clip}%
\pgfsetbuttcap%
\pgfsetroundjoin%
\definecolor{currentfill}{rgb}{0.429719,0.000831,0.659425}%
\pgfsetfillcolor{currentfill}%
\pgfsetlinewidth{0.000000pt}%
\definecolor{currentstroke}{rgb}{0.000000,0.000000,0.000000}%
\pgfsetstrokecolor{currentstroke}%
\pgfsetdash{}{0pt}%
\pgfpathmoveto{\pgfqpoint{4.255037in}{1.070290in}}%
\pgfpathlineto{\pgfqpoint{4.427120in}{1.070290in}}%
\pgfpathlineto{\pgfqpoint{4.427120in}{1.083734in}}%
\pgfpathlineto{\pgfqpoint{4.255037in}{1.083734in}}%
\pgfpathlineto{\pgfqpoint{4.255037in}{1.070290in}}%
\pgfusepath{fill}%
\end{pgfscope}%
\begin{pgfscope}%
\pgfpathrectangle{\pgfqpoint{4.255037in}{0.357758in}}{\pgfqpoint{0.172083in}{3.441662in}}%
\pgfusepath{clip}%
\pgfsetbuttcap%
\pgfsetroundjoin%
\definecolor{currentfill}{rgb}{0.435734,0.001127,0.659797}%
\pgfsetfillcolor{currentfill}%
\pgfsetlinewidth{0.000000pt}%
\definecolor{currentstroke}{rgb}{0.000000,0.000000,0.000000}%
\pgfsetstrokecolor{currentstroke}%
\pgfsetdash{}{0pt}%
\pgfpathmoveto{\pgfqpoint{4.255037in}{1.083734in}}%
\pgfpathlineto{\pgfqpoint{4.427120in}{1.083734in}}%
\pgfpathlineto{\pgfqpoint{4.427120in}{1.097178in}}%
\pgfpathlineto{\pgfqpoint{4.255037in}{1.097178in}}%
\pgfpathlineto{\pgfqpoint{4.255037in}{1.083734in}}%
\pgfusepath{fill}%
\end{pgfscope}%
\begin{pgfscope}%
\pgfpathrectangle{\pgfqpoint{4.255037in}{0.357758in}}{\pgfqpoint{0.172083in}{3.441662in}}%
\pgfusepath{clip}%
\pgfsetbuttcap%
\pgfsetroundjoin%
\definecolor{currentfill}{rgb}{0.441732,0.001540,0.660069}%
\pgfsetfillcolor{currentfill}%
\pgfsetlinewidth{0.000000pt}%
\definecolor{currentstroke}{rgb}{0.000000,0.000000,0.000000}%
\pgfsetstrokecolor{currentstroke}%
\pgfsetdash{}{0pt}%
\pgfpathmoveto{\pgfqpoint{4.255037in}{1.097178in}}%
\pgfpathlineto{\pgfqpoint{4.427120in}{1.097178in}}%
\pgfpathlineto{\pgfqpoint{4.427120in}{1.110622in}}%
\pgfpathlineto{\pgfqpoint{4.255037in}{1.110622in}}%
\pgfpathlineto{\pgfqpoint{4.255037in}{1.097178in}}%
\pgfusepath{fill}%
\end{pgfscope}%
\begin{pgfscope}%
\pgfpathrectangle{\pgfqpoint{4.255037in}{0.357758in}}{\pgfqpoint{0.172083in}{3.441662in}}%
\pgfusepath{clip}%
\pgfsetbuttcap%
\pgfsetroundjoin%
\definecolor{currentfill}{rgb}{0.447714,0.002080,0.660240}%
\pgfsetfillcolor{currentfill}%
\pgfsetlinewidth{0.000000pt}%
\definecolor{currentstroke}{rgb}{0.000000,0.000000,0.000000}%
\pgfsetstrokecolor{currentstroke}%
\pgfsetdash{}{0pt}%
\pgfpathmoveto{\pgfqpoint{4.255037in}{1.110622in}}%
\pgfpathlineto{\pgfqpoint{4.427120in}{1.110622in}}%
\pgfpathlineto{\pgfqpoint{4.427120in}{1.124066in}}%
\pgfpathlineto{\pgfqpoint{4.255037in}{1.124066in}}%
\pgfpathlineto{\pgfqpoint{4.255037in}{1.110622in}}%
\pgfusepath{fill}%
\end{pgfscope}%
\begin{pgfscope}%
\pgfpathrectangle{\pgfqpoint{4.255037in}{0.357758in}}{\pgfqpoint{0.172083in}{3.441662in}}%
\pgfusepath{clip}%
\pgfsetbuttcap%
\pgfsetroundjoin%
\definecolor{currentfill}{rgb}{0.453677,0.002755,0.660310}%
\pgfsetfillcolor{currentfill}%
\pgfsetlinewidth{0.000000pt}%
\definecolor{currentstroke}{rgb}{0.000000,0.000000,0.000000}%
\pgfsetstrokecolor{currentstroke}%
\pgfsetdash{}{0pt}%
\pgfpathmoveto{\pgfqpoint{4.255037in}{1.124066in}}%
\pgfpathlineto{\pgfqpoint{4.427120in}{1.124066in}}%
\pgfpathlineto{\pgfqpoint{4.427120in}{1.137510in}}%
\pgfpathlineto{\pgfqpoint{4.255037in}{1.137510in}}%
\pgfpathlineto{\pgfqpoint{4.255037in}{1.124066in}}%
\pgfusepath{fill}%
\end{pgfscope}%
\begin{pgfscope}%
\pgfpathrectangle{\pgfqpoint{4.255037in}{0.357758in}}{\pgfqpoint{0.172083in}{3.441662in}}%
\pgfusepath{clip}%
\pgfsetbuttcap%
\pgfsetroundjoin%
\definecolor{currentfill}{rgb}{0.459623,0.003574,0.660277}%
\pgfsetfillcolor{currentfill}%
\pgfsetlinewidth{0.000000pt}%
\definecolor{currentstroke}{rgb}{0.000000,0.000000,0.000000}%
\pgfsetstrokecolor{currentstroke}%
\pgfsetdash{}{0pt}%
\pgfpathmoveto{\pgfqpoint{4.255037in}{1.137510in}}%
\pgfpathlineto{\pgfqpoint{4.427120in}{1.137510in}}%
\pgfpathlineto{\pgfqpoint{4.427120in}{1.150954in}}%
\pgfpathlineto{\pgfqpoint{4.255037in}{1.150954in}}%
\pgfpathlineto{\pgfqpoint{4.255037in}{1.137510in}}%
\pgfusepath{fill}%
\end{pgfscope}%
\begin{pgfscope}%
\pgfpathrectangle{\pgfqpoint{4.255037in}{0.357758in}}{\pgfqpoint{0.172083in}{3.441662in}}%
\pgfusepath{clip}%
\pgfsetbuttcap%
\pgfsetroundjoin%
\definecolor{currentfill}{rgb}{0.465550,0.004545,0.660139}%
\pgfsetfillcolor{currentfill}%
\pgfsetlinewidth{0.000000pt}%
\definecolor{currentstroke}{rgb}{0.000000,0.000000,0.000000}%
\pgfsetstrokecolor{currentstroke}%
\pgfsetdash{}{0pt}%
\pgfpathmoveto{\pgfqpoint{4.255037in}{1.150954in}}%
\pgfpathlineto{\pgfqpoint{4.427120in}{1.150954in}}%
\pgfpathlineto{\pgfqpoint{4.427120in}{1.164398in}}%
\pgfpathlineto{\pgfqpoint{4.255037in}{1.164398in}}%
\pgfpathlineto{\pgfqpoint{4.255037in}{1.150954in}}%
\pgfusepath{fill}%
\end{pgfscope}%
\begin{pgfscope}%
\pgfpathrectangle{\pgfqpoint{4.255037in}{0.357758in}}{\pgfqpoint{0.172083in}{3.441662in}}%
\pgfusepath{clip}%
\pgfsetbuttcap%
\pgfsetroundjoin%
\definecolor{currentfill}{rgb}{0.471457,0.005678,0.659897}%
\pgfsetfillcolor{currentfill}%
\pgfsetlinewidth{0.000000pt}%
\definecolor{currentstroke}{rgb}{0.000000,0.000000,0.000000}%
\pgfsetstrokecolor{currentstroke}%
\pgfsetdash{}{0pt}%
\pgfpathmoveto{\pgfqpoint{4.255037in}{1.164398in}}%
\pgfpathlineto{\pgfqpoint{4.427120in}{1.164398in}}%
\pgfpathlineto{\pgfqpoint{4.427120in}{1.177842in}}%
\pgfpathlineto{\pgfqpoint{4.255037in}{1.177842in}}%
\pgfpathlineto{\pgfqpoint{4.255037in}{1.164398in}}%
\pgfusepath{fill}%
\end{pgfscope}%
\begin{pgfscope}%
\pgfpathrectangle{\pgfqpoint{4.255037in}{0.357758in}}{\pgfqpoint{0.172083in}{3.441662in}}%
\pgfusepath{clip}%
\pgfsetbuttcap%
\pgfsetroundjoin%
\definecolor{currentfill}{rgb}{0.477344,0.006980,0.659549}%
\pgfsetfillcolor{currentfill}%
\pgfsetlinewidth{0.000000pt}%
\definecolor{currentstroke}{rgb}{0.000000,0.000000,0.000000}%
\pgfsetstrokecolor{currentstroke}%
\pgfsetdash{}{0pt}%
\pgfpathmoveto{\pgfqpoint{4.255037in}{1.177842in}}%
\pgfpathlineto{\pgfqpoint{4.427120in}{1.177842in}}%
\pgfpathlineto{\pgfqpoint{4.427120in}{1.191286in}}%
\pgfpathlineto{\pgfqpoint{4.255037in}{1.191286in}}%
\pgfpathlineto{\pgfqpoint{4.255037in}{1.177842in}}%
\pgfusepath{fill}%
\end{pgfscope}%
\begin{pgfscope}%
\pgfpathrectangle{\pgfqpoint{4.255037in}{0.357758in}}{\pgfqpoint{0.172083in}{3.441662in}}%
\pgfusepath{clip}%
\pgfsetbuttcap%
\pgfsetroundjoin%
\definecolor{currentfill}{rgb}{0.483210,0.008460,0.659095}%
\pgfsetfillcolor{currentfill}%
\pgfsetlinewidth{0.000000pt}%
\definecolor{currentstroke}{rgb}{0.000000,0.000000,0.000000}%
\pgfsetstrokecolor{currentstroke}%
\pgfsetdash{}{0pt}%
\pgfpathmoveto{\pgfqpoint{4.255037in}{1.191286in}}%
\pgfpathlineto{\pgfqpoint{4.427120in}{1.191286in}}%
\pgfpathlineto{\pgfqpoint{4.427120in}{1.204730in}}%
\pgfpathlineto{\pgfqpoint{4.255037in}{1.204730in}}%
\pgfpathlineto{\pgfqpoint{4.255037in}{1.191286in}}%
\pgfusepath{fill}%
\end{pgfscope}%
\begin{pgfscope}%
\pgfpathrectangle{\pgfqpoint{4.255037in}{0.357758in}}{\pgfqpoint{0.172083in}{3.441662in}}%
\pgfusepath{clip}%
\pgfsetbuttcap%
\pgfsetroundjoin%
\definecolor{currentfill}{rgb}{0.489055,0.010127,0.658534}%
\pgfsetfillcolor{currentfill}%
\pgfsetlinewidth{0.000000pt}%
\definecolor{currentstroke}{rgb}{0.000000,0.000000,0.000000}%
\pgfsetstrokecolor{currentstroke}%
\pgfsetdash{}{0pt}%
\pgfpathmoveto{\pgfqpoint{4.255037in}{1.204730in}}%
\pgfpathlineto{\pgfqpoint{4.427120in}{1.204730in}}%
\pgfpathlineto{\pgfqpoint{4.427120in}{1.218174in}}%
\pgfpathlineto{\pgfqpoint{4.255037in}{1.218174in}}%
\pgfpathlineto{\pgfqpoint{4.255037in}{1.204730in}}%
\pgfusepath{fill}%
\end{pgfscope}%
\begin{pgfscope}%
\pgfpathrectangle{\pgfqpoint{4.255037in}{0.357758in}}{\pgfqpoint{0.172083in}{3.441662in}}%
\pgfusepath{clip}%
\pgfsetbuttcap%
\pgfsetroundjoin%
\definecolor{currentfill}{rgb}{0.494877,0.011990,0.657865}%
\pgfsetfillcolor{currentfill}%
\pgfsetlinewidth{0.000000pt}%
\definecolor{currentstroke}{rgb}{0.000000,0.000000,0.000000}%
\pgfsetstrokecolor{currentstroke}%
\pgfsetdash{}{0pt}%
\pgfpathmoveto{\pgfqpoint{4.255037in}{1.218174in}}%
\pgfpathlineto{\pgfqpoint{4.427120in}{1.218174in}}%
\pgfpathlineto{\pgfqpoint{4.427120in}{1.231618in}}%
\pgfpathlineto{\pgfqpoint{4.255037in}{1.231618in}}%
\pgfpathlineto{\pgfqpoint{4.255037in}{1.218174in}}%
\pgfusepath{fill}%
\end{pgfscope}%
\begin{pgfscope}%
\pgfpathrectangle{\pgfqpoint{4.255037in}{0.357758in}}{\pgfqpoint{0.172083in}{3.441662in}}%
\pgfusepath{clip}%
\pgfsetbuttcap%
\pgfsetroundjoin%
\definecolor{currentfill}{rgb}{0.500678,0.014055,0.657088}%
\pgfsetfillcolor{currentfill}%
\pgfsetlinewidth{0.000000pt}%
\definecolor{currentstroke}{rgb}{0.000000,0.000000,0.000000}%
\pgfsetstrokecolor{currentstroke}%
\pgfsetdash{}{0pt}%
\pgfpathmoveto{\pgfqpoint{4.255037in}{1.231618in}}%
\pgfpathlineto{\pgfqpoint{4.427120in}{1.231618in}}%
\pgfpathlineto{\pgfqpoint{4.427120in}{1.245062in}}%
\pgfpathlineto{\pgfqpoint{4.255037in}{1.245062in}}%
\pgfpathlineto{\pgfqpoint{4.255037in}{1.231618in}}%
\pgfusepath{fill}%
\end{pgfscope}%
\begin{pgfscope}%
\pgfpathrectangle{\pgfqpoint{4.255037in}{0.357758in}}{\pgfqpoint{0.172083in}{3.441662in}}%
\pgfusepath{clip}%
\pgfsetbuttcap%
\pgfsetroundjoin%
\definecolor{currentfill}{rgb}{0.506454,0.016333,0.656202}%
\pgfsetfillcolor{currentfill}%
\pgfsetlinewidth{0.000000pt}%
\definecolor{currentstroke}{rgb}{0.000000,0.000000,0.000000}%
\pgfsetstrokecolor{currentstroke}%
\pgfsetdash{}{0pt}%
\pgfpathmoveto{\pgfqpoint{4.255037in}{1.245062in}}%
\pgfpathlineto{\pgfqpoint{4.427120in}{1.245062in}}%
\pgfpathlineto{\pgfqpoint{4.427120in}{1.258506in}}%
\pgfpathlineto{\pgfqpoint{4.255037in}{1.258506in}}%
\pgfpathlineto{\pgfqpoint{4.255037in}{1.245062in}}%
\pgfusepath{fill}%
\end{pgfscope}%
\begin{pgfscope}%
\pgfpathrectangle{\pgfqpoint{4.255037in}{0.357758in}}{\pgfqpoint{0.172083in}{3.441662in}}%
\pgfusepath{clip}%
\pgfsetbuttcap%
\pgfsetroundjoin%
\definecolor{currentfill}{rgb}{0.512206,0.018833,0.655209}%
\pgfsetfillcolor{currentfill}%
\pgfsetlinewidth{0.000000pt}%
\definecolor{currentstroke}{rgb}{0.000000,0.000000,0.000000}%
\pgfsetstrokecolor{currentstroke}%
\pgfsetdash{}{0pt}%
\pgfpathmoveto{\pgfqpoint{4.255037in}{1.258506in}}%
\pgfpathlineto{\pgfqpoint{4.427120in}{1.258506in}}%
\pgfpathlineto{\pgfqpoint{4.427120in}{1.271950in}}%
\pgfpathlineto{\pgfqpoint{4.255037in}{1.271950in}}%
\pgfpathlineto{\pgfqpoint{4.255037in}{1.258506in}}%
\pgfusepath{fill}%
\end{pgfscope}%
\begin{pgfscope}%
\pgfpathrectangle{\pgfqpoint{4.255037in}{0.357758in}}{\pgfqpoint{0.172083in}{3.441662in}}%
\pgfusepath{clip}%
\pgfsetbuttcap%
\pgfsetroundjoin%
\definecolor{currentfill}{rgb}{0.517933,0.021563,0.654109}%
\pgfsetfillcolor{currentfill}%
\pgfsetlinewidth{0.000000pt}%
\definecolor{currentstroke}{rgb}{0.000000,0.000000,0.000000}%
\pgfsetstrokecolor{currentstroke}%
\pgfsetdash{}{0pt}%
\pgfpathmoveto{\pgfqpoint{4.255037in}{1.271950in}}%
\pgfpathlineto{\pgfqpoint{4.427120in}{1.271950in}}%
\pgfpathlineto{\pgfqpoint{4.427120in}{1.285394in}}%
\pgfpathlineto{\pgfqpoint{4.255037in}{1.285394in}}%
\pgfpathlineto{\pgfqpoint{4.255037in}{1.271950in}}%
\pgfusepath{fill}%
\end{pgfscope}%
\begin{pgfscope}%
\pgfpathrectangle{\pgfqpoint{4.255037in}{0.357758in}}{\pgfqpoint{0.172083in}{3.441662in}}%
\pgfusepath{clip}%
\pgfsetbuttcap%
\pgfsetroundjoin%
\definecolor{currentfill}{rgb}{0.523633,0.024532,0.652901}%
\pgfsetfillcolor{currentfill}%
\pgfsetlinewidth{0.000000pt}%
\definecolor{currentstroke}{rgb}{0.000000,0.000000,0.000000}%
\pgfsetstrokecolor{currentstroke}%
\pgfsetdash{}{0pt}%
\pgfpathmoveto{\pgfqpoint{4.255037in}{1.285394in}}%
\pgfpathlineto{\pgfqpoint{4.427120in}{1.285394in}}%
\pgfpathlineto{\pgfqpoint{4.427120in}{1.298838in}}%
\pgfpathlineto{\pgfqpoint{4.255037in}{1.298838in}}%
\pgfpathlineto{\pgfqpoint{4.255037in}{1.285394in}}%
\pgfusepath{fill}%
\end{pgfscope}%
\begin{pgfscope}%
\pgfpathrectangle{\pgfqpoint{4.255037in}{0.357758in}}{\pgfqpoint{0.172083in}{3.441662in}}%
\pgfusepath{clip}%
\pgfsetbuttcap%
\pgfsetroundjoin%
\definecolor{currentfill}{rgb}{0.529306,0.027747,0.651586}%
\pgfsetfillcolor{currentfill}%
\pgfsetlinewidth{0.000000pt}%
\definecolor{currentstroke}{rgb}{0.000000,0.000000,0.000000}%
\pgfsetstrokecolor{currentstroke}%
\pgfsetdash{}{0pt}%
\pgfpathmoveto{\pgfqpoint{4.255037in}{1.298838in}}%
\pgfpathlineto{\pgfqpoint{4.427120in}{1.298838in}}%
\pgfpathlineto{\pgfqpoint{4.427120in}{1.312282in}}%
\pgfpathlineto{\pgfqpoint{4.255037in}{1.312282in}}%
\pgfpathlineto{\pgfqpoint{4.255037in}{1.298838in}}%
\pgfusepath{fill}%
\end{pgfscope}%
\begin{pgfscope}%
\pgfpathrectangle{\pgfqpoint{4.255037in}{0.357758in}}{\pgfqpoint{0.172083in}{3.441662in}}%
\pgfusepath{clip}%
\pgfsetbuttcap%
\pgfsetroundjoin%
\definecolor{currentfill}{rgb}{0.534952,0.031217,0.650165}%
\pgfsetfillcolor{currentfill}%
\pgfsetlinewidth{0.000000pt}%
\definecolor{currentstroke}{rgb}{0.000000,0.000000,0.000000}%
\pgfsetstrokecolor{currentstroke}%
\pgfsetdash{}{0pt}%
\pgfpathmoveto{\pgfqpoint{4.255037in}{1.312282in}}%
\pgfpathlineto{\pgfqpoint{4.427120in}{1.312282in}}%
\pgfpathlineto{\pgfqpoint{4.427120in}{1.325726in}}%
\pgfpathlineto{\pgfqpoint{4.255037in}{1.325726in}}%
\pgfpathlineto{\pgfqpoint{4.255037in}{1.312282in}}%
\pgfusepath{fill}%
\end{pgfscope}%
\begin{pgfscope}%
\pgfpathrectangle{\pgfqpoint{4.255037in}{0.357758in}}{\pgfqpoint{0.172083in}{3.441662in}}%
\pgfusepath{clip}%
\pgfsetbuttcap%
\pgfsetroundjoin%
\definecolor{currentfill}{rgb}{0.540570,0.034950,0.648640}%
\pgfsetfillcolor{currentfill}%
\pgfsetlinewidth{0.000000pt}%
\definecolor{currentstroke}{rgb}{0.000000,0.000000,0.000000}%
\pgfsetstrokecolor{currentstroke}%
\pgfsetdash{}{0pt}%
\pgfpathmoveto{\pgfqpoint{4.255037in}{1.325726in}}%
\pgfpathlineto{\pgfqpoint{4.427120in}{1.325726in}}%
\pgfpathlineto{\pgfqpoint{4.427120in}{1.339170in}}%
\pgfpathlineto{\pgfqpoint{4.255037in}{1.339170in}}%
\pgfpathlineto{\pgfqpoint{4.255037in}{1.325726in}}%
\pgfusepath{fill}%
\end{pgfscope}%
\begin{pgfscope}%
\pgfpathrectangle{\pgfqpoint{4.255037in}{0.357758in}}{\pgfqpoint{0.172083in}{3.441662in}}%
\pgfusepath{clip}%
\pgfsetbuttcap%
\pgfsetroundjoin%
\definecolor{currentfill}{rgb}{0.546157,0.038954,0.647010}%
\pgfsetfillcolor{currentfill}%
\pgfsetlinewidth{0.000000pt}%
\definecolor{currentstroke}{rgb}{0.000000,0.000000,0.000000}%
\pgfsetstrokecolor{currentstroke}%
\pgfsetdash{}{0pt}%
\pgfpathmoveto{\pgfqpoint{4.255037in}{1.339170in}}%
\pgfpathlineto{\pgfqpoint{4.427120in}{1.339170in}}%
\pgfpathlineto{\pgfqpoint{4.427120in}{1.352614in}}%
\pgfpathlineto{\pgfqpoint{4.255037in}{1.352614in}}%
\pgfpathlineto{\pgfqpoint{4.255037in}{1.339170in}}%
\pgfusepath{fill}%
\end{pgfscope}%
\begin{pgfscope}%
\pgfpathrectangle{\pgfqpoint{4.255037in}{0.357758in}}{\pgfqpoint{0.172083in}{3.441662in}}%
\pgfusepath{clip}%
\pgfsetbuttcap%
\pgfsetroundjoin%
\definecolor{currentfill}{rgb}{0.551715,0.043136,0.645277}%
\pgfsetfillcolor{currentfill}%
\pgfsetlinewidth{0.000000pt}%
\definecolor{currentstroke}{rgb}{0.000000,0.000000,0.000000}%
\pgfsetstrokecolor{currentstroke}%
\pgfsetdash{}{0pt}%
\pgfpathmoveto{\pgfqpoint{4.255037in}{1.352614in}}%
\pgfpathlineto{\pgfqpoint{4.427120in}{1.352614in}}%
\pgfpathlineto{\pgfqpoint{4.427120in}{1.366058in}}%
\pgfpathlineto{\pgfqpoint{4.255037in}{1.366058in}}%
\pgfpathlineto{\pgfqpoint{4.255037in}{1.352614in}}%
\pgfusepath{fill}%
\end{pgfscope}%
\begin{pgfscope}%
\pgfpathrectangle{\pgfqpoint{4.255037in}{0.357758in}}{\pgfqpoint{0.172083in}{3.441662in}}%
\pgfusepath{clip}%
\pgfsetbuttcap%
\pgfsetroundjoin%
\definecolor{currentfill}{rgb}{0.557243,0.047331,0.643443}%
\pgfsetfillcolor{currentfill}%
\pgfsetlinewidth{0.000000pt}%
\definecolor{currentstroke}{rgb}{0.000000,0.000000,0.000000}%
\pgfsetstrokecolor{currentstroke}%
\pgfsetdash{}{0pt}%
\pgfpathmoveto{\pgfqpoint{4.255037in}{1.366058in}}%
\pgfpathlineto{\pgfqpoint{4.427120in}{1.366058in}}%
\pgfpathlineto{\pgfqpoint{4.427120in}{1.379502in}}%
\pgfpathlineto{\pgfqpoint{4.255037in}{1.379502in}}%
\pgfpathlineto{\pgfqpoint{4.255037in}{1.366058in}}%
\pgfusepath{fill}%
\end{pgfscope}%
\begin{pgfscope}%
\pgfpathrectangle{\pgfqpoint{4.255037in}{0.357758in}}{\pgfqpoint{0.172083in}{3.441662in}}%
\pgfusepath{clip}%
\pgfsetbuttcap%
\pgfsetroundjoin%
\definecolor{currentfill}{rgb}{0.562738,0.051545,0.641509}%
\pgfsetfillcolor{currentfill}%
\pgfsetlinewidth{0.000000pt}%
\definecolor{currentstroke}{rgb}{0.000000,0.000000,0.000000}%
\pgfsetstrokecolor{currentstroke}%
\pgfsetdash{}{0pt}%
\pgfpathmoveto{\pgfqpoint{4.255037in}{1.379502in}}%
\pgfpathlineto{\pgfqpoint{4.427120in}{1.379502in}}%
\pgfpathlineto{\pgfqpoint{4.427120in}{1.392946in}}%
\pgfpathlineto{\pgfqpoint{4.255037in}{1.392946in}}%
\pgfpathlineto{\pgfqpoint{4.255037in}{1.379502in}}%
\pgfusepath{fill}%
\end{pgfscope}%
\begin{pgfscope}%
\pgfpathrectangle{\pgfqpoint{4.255037in}{0.357758in}}{\pgfqpoint{0.172083in}{3.441662in}}%
\pgfusepath{clip}%
\pgfsetbuttcap%
\pgfsetroundjoin%
\definecolor{currentfill}{rgb}{0.568201,0.055778,0.639477}%
\pgfsetfillcolor{currentfill}%
\pgfsetlinewidth{0.000000pt}%
\definecolor{currentstroke}{rgb}{0.000000,0.000000,0.000000}%
\pgfsetstrokecolor{currentstroke}%
\pgfsetdash{}{0pt}%
\pgfpathmoveto{\pgfqpoint{4.255037in}{1.392946in}}%
\pgfpathlineto{\pgfqpoint{4.427120in}{1.392946in}}%
\pgfpathlineto{\pgfqpoint{4.427120in}{1.406390in}}%
\pgfpathlineto{\pgfqpoint{4.255037in}{1.406390in}}%
\pgfpathlineto{\pgfqpoint{4.255037in}{1.392946in}}%
\pgfusepath{fill}%
\end{pgfscope}%
\begin{pgfscope}%
\pgfpathrectangle{\pgfqpoint{4.255037in}{0.357758in}}{\pgfqpoint{0.172083in}{3.441662in}}%
\pgfusepath{clip}%
\pgfsetbuttcap%
\pgfsetroundjoin%
\definecolor{currentfill}{rgb}{0.573632,0.060028,0.637349}%
\pgfsetfillcolor{currentfill}%
\pgfsetlinewidth{0.000000pt}%
\definecolor{currentstroke}{rgb}{0.000000,0.000000,0.000000}%
\pgfsetstrokecolor{currentstroke}%
\pgfsetdash{}{0pt}%
\pgfpathmoveto{\pgfqpoint{4.255037in}{1.406390in}}%
\pgfpathlineto{\pgfqpoint{4.427120in}{1.406390in}}%
\pgfpathlineto{\pgfqpoint{4.427120in}{1.419834in}}%
\pgfpathlineto{\pgfqpoint{4.255037in}{1.419834in}}%
\pgfpathlineto{\pgfqpoint{4.255037in}{1.406390in}}%
\pgfusepath{fill}%
\end{pgfscope}%
\begin{pgfscope}%
\pgfpathrectangle{\pgfqpoint{4.255037in}{0.357758in}}{\pgfqpoint{0.172083in}{3.441662in}}%
\pgfusepath{clip}%
\pgfsetbuttcap%
\pgfsetroundjoin%
\definecolor{currentfill}{rgb}{0.579029,0.064296,0.635126}%
\pgfsetfillcolor{currentfill}%
\pgfsetlinewidth{0.000000pt}%
\definecolor{currentstroke}{rgb}{0.000000,0.000000,0.000000}%
\pgfsetstrokecolor{currentstroke}%
\pgfsetdash{}{0pt}%
\pgfpathmoveto{\pgfqpoint{4.255037in}{1.419834in}}%
\pgfpathlineto{\pgfqpoint{4.427120in}{1.419834in}}%
\pgfpathlineto{\pgfqpoint{4.427120in}{1.433278in}}%
\pgfpathlineto{\pgfqpoint{4.255037in}{1.433278in}}%
\pgfpathlineto{\pgfqpoint{4.255037in}{1.419834in}}%
\pgfusepath{fill}%
\end{pgfscope}%
\begin{pgfscope}%
\pgfpathrectangle{\pgfqpoint{4.255037in}{0.357758in}}{\pgfqpoint{0.172083in}{3.441662in}}%
\pgfusepath{clip}%
\pgfsetbuttcap%
\pgfsetroundjoin%
\definecolor{currentfill}{rgb}{0.584391,0.068579,0.632812}%
\pgfsetfillcolor{currentfill}%
\pgfsetlinewidth{0.000000pt}%
\definecolor{currentstroke}{rgb}{0.000000,0.000000,0.000000}%
\pgfsetstrokecolor{currentstroke}%
\pgfsetdash{}{0pt}%
\pgfpathmoveto{\pgfqpoint{4.255037in}{1.433278in}}%
\pgfpathlineto{\pgfqpoint{4.427120in}{1.433278in}}%
\pgfpathlineto{\pgfqpoint{4.427120in}{1.446722in}}%
\pgfpathlineto{\pgfqpoint{4.255037in}{1.446722in}}%
\pgfpathlineto{\pgfqpoint{4.255037in}{1.433278in}}%
\pgfusepath{fill}%
\end{pgfscope}%
\begin{pgfscope}%
\pgfpathrectangle{\pgfqpoint{4.255037in}{0.357758in}}{\pgfqpoint{0.172083in}{3.441662in}}%
\pgfusepath{clip}%
\pgfsetbuttcap%
\pgfsetroundjoin%
\definecolor{currentfill}{rgb}{0.589719,0.072878,0.630408}%
\pgfsetfillcolor{currentfill}%
\pgfsetlinewidth{0.000000pt}%
\definecolor{currentstroke}{rgb}{0.000000,0.000000,0.000000}%
\pgfsetstrokecolor{currentstroke}%
\pgfsetdash{}{0pt}%
\pgfpathmoveto{\pgfqpoint{4.255037in}{1.446722in}}%
\pgfpathlineto{\pgfqpoint{4.427120in}{1.446722in}}%
\pgfpathlineto{\pgfqpoint{4.427120in}{1.460166in}}%
\pgfpathlineto{\pgfqpoint{4.255037in}{1.460166in}}%
\pgfpathlineto{\pgfqpoint{4.255037in}{1.446722in}}%
\pgfusepath{fill}%
\end{pgfscope}%
\begin{pgfscope}%
\pgfpathrectangle{\pgfqpoint{4.255037in}{0.357758in}}{\pgfqpoint{0.172083in}{3.441662in}}%
\pgfusepath{clip}%
\pgfsetbuttcap%
\pgfsetroundjoin%
\definecolor{currentfill}{rgb}{0.595011,0.077190,0.627917}%
\pgfsetfillcolor{currentfill}%
\pgfsetlinewidth{0.000000pt}%
\definecolor{currentstroke}{rgb}{0.000000,0.000000,0.000000}%
\pgfsetstrokecolor{currentstroke}%
\pgfsetdash{}{0pt}%
\pgfpathmoveto{\pgfqpoint{4.255037in}{1.460166in}}%
\pgfpathlineto{\pgfqpoint{4.427120in}{1.460166in}}%
\pgfpathlineto{\pgfqpoint{4.427120in}{1.473610in}}%
\pgfpathlineto{\pgfqpoint{4.255037in}{1.473610in}}%
\pgfpathlineto{\pgfqpoint{4.255037in}{1.460166in}}%
\pgfusepath{fill}%
\end{pgfscope}%
\begin{pgfscope}%
\pgfpathrectangle{\pgfqpoint{4.255037in}{0.357758in}}{\pgfqpoint{0.172083in}{3.441662in}}%
\pgfusepath{clip}%
\pgfsetbuttcap%
\pgfsetroundjoin%
\definecolor{currentfill}{rgb}{0.600266,0.081516,0.625342}%
\pgfsetfillcolor{currentfill}%
\pgfsetlinewidth{0.000000pt}%
\definecolor{currentstroke}{rgb}{0.000000,0.000000,0.000000}%
\pgfsetstrokecolor{currentstroke}%
\pgfsetdash{}{0pt}%
\pgfpathmoveto{\pgfqpoint{4.255037in}{1.473610in}}%
\pgfpathlineto{\pgfqpoint{4.427120in}{1.473610in}}%
\pgfpathlineto{\pgfqpoint{4.427120in}{1.487054in}}%
\pgfpathlineto{\pgfqpoint{4.255037in}{1.487054in}}%
\pgfpathlineto{\pgfqpoint{4.255037in}{1.473610in}}%
\pgfusepath{fill}%
\end{pgfscope}%
\begin{pgfscope}%
\pgfpathrectangle{\pgfqpoint{4.255037in}{0.357758in}}{\pgfqpoint{0.172083in}{3.441662in}}%
\pgfusepath{clip}%
\pgfsetbuttcap%
\pgfsetroundjoin%
\definecolor{currentfill}{rgb}{0.605485,0.085854,0.622686}%
\pgfsetfillcolor{currentfill}%
\pgfsetlinewidth{0.000000pt}%
\definecolor{currentstroke}{rgb}{0.000000,0.000000,0.000000}%
\pgfsetstrokecolor{currentstroke}%
\pgfsetdash{}{0pt}%
\pgfpathmoveto{\pgfqpoint{4.255037in}{1.487054in}}%
\pgfpathlineto{\pgfqpoint{4.427120in}{1.487054in}}%
\pgfpathlineto{\pgfqpoint{4.427120in}{1.500498in}}%
\pgfpathlineto{\pgfqpoint{4.255037in}{1.500498in}}%
\pgfpathlineto{\pgfqpoint{4.255037in}{1.487054in}}%
\pgfusepath{fill}%
\end{pgfscope}%
\begin{pgfscope}%
\pgfpathrectangle{\pgfqpoint{4.255037in}{0.357758in}}{\pgfqpoint{0.172083in}{3.441662in}}%
\pgfusepath{clip}%
\pgfsetbuttcap%
\pgfsetroundjoin%
\definecolor{currentfill}{rgb}{0.610667,0.090204,0.619951}%
\pgfsetfillcolor{currentfill}%
\pgfsetlinewidth{0.000000pt}%
\definecolor{currentstroke}{rgb}{0.000000,0.000000,0.000000}%
\pgfsetstrokecolor{currentstroke}%
\pgfsetdash{}{0pt}%
\pgfpathmoveto{\pgfqpoint{4.255037in}{1.500498in}}%
\pgfpathlineto{\pgfqpoint{4.427120in}{1.500498in}}%
\pgfpathlineto{\pgfqpoint{4.427120in}{1.513942in}}%
\pgfpathlineto{\pgfqpoint{4.255037in}{1.513942in}}%
\pgfpathlineto{\pgfqpoint{4.255037in}{1.500498in}}%
\pgfusepath{fill}%
\end{pgfscope}%
\begin{pgfscope}%
\pgfpathrectangle{\pgfqpoint{4.255037in}{0.357758in}}{\pgfqpoint{0.172083in}{3.441662in}}%
\pgfusepath{clip}%
\pgfsetbuttcap%
\pgfsetroundjoin%
\definecolor{currentfill}{rgb}{0.615812,0.094564,0.617140}%
\pgfsetfillcolor{currentfill}%
\pgfsetlinewidth{0.000000pt}%
\definecolor{currentstroke}{rgb}{0.000000,0.000000,0.000000}%
\pgfsetstrokecolor{currentstroke}%
\pgfsetdash{}{0pt}%
\pgfpathmoveto{\pgfqpoint{4.255037in}{1.513942in}}%
\pgfpathlineto{\pgfqpoint{4.427120in}{1.513942in}}%
\pgfpathlineto{\pgfqpoint{4.427120in}{1.527385in}}%
\pgfpathlineto{\pgfqpoint{4.255037in}{1.527385in}}%
\pgfpathlineto{\pgfqpoint{4.255037in}{1.513942in}}%
\pgfusepath{fill}%
\end{pgfscope}%
\begin{pgfscope}%
\pgfpathrectangle{\pgfqpoint{4.255037in}{0.357758in}}{\pgfqpoint{0.172083in}{3.441662in}}%
\pgfusepath{clip}%
\pgfsetbuttcap%
\pgfsetroundjoin%
\definecolor{currentfill}{rgb}{0.620919,0.098934,0.614257}%
\pgfsetfillcolor{currentfill}%
\pgfsetlinewidth{0.000000pt}%
\definecolor{currentstroke}{rgb}{0.000000,0.000000,0.000000}%
\pgfsetstrokecolor{currentstroke}%
\pgfsetdash{}{0pt}%
\pgfpathmoveto{\pgfqpoint{4.255037in}{1.527385in}}%
\pgfpathlineto{\pgfqpoint{4.427120in}{1.527385in}}%
\pgfpathlineto{\pgfqpoint{4.427120in}{1.540829in}}%
\pgfpathlineto{\pgfqpoint{4.255037in}{1.540829in}}%
\pgfpathlineto{\pgfqpoint{4.255037in}{1.527385in}}%
\pgfusepath{fill}%
\end{pgfscope}%
\begin{pgfscope}%
\pgfpathrectangle{\pgfqpoint{4.255037in}{0.357758in}}{\pgfqpoint{0.172083in}{3.441662in}}%
\pgfusepath{clip}%
\pgfsetbuttcap%
\pgfsetroundjoin%
\definecolor{currentfill}{rgb}{0.625987,0.103312,0.611305}%
\pgfsetfillcolor{currentfill}%
\pgfsetlinewidth{0.000000pt}%
\definecolor{currentstroke}{rgb}{0.000000,0.000000,0.000000}%
\pgfsetstrokecolor{currentstroke}%
\pgfsetdash{}{0pt}%
\pgfpathmoveto{\pgfqpoint{4.255037in}{1.540829in}}%
\pgfpathlineto{\pgfqpoint{4.427120in}{1.540829in}}%
\pgfpathlineto{\pgfqpoint{4.427120in}{1.554273in}}%
\pgfpathlineto{\pgfqpoint{4.255037in}{1.554273in}}%
\pgfpathlineto{\pgfqpoint{4.255037in}{1.540829in}}%
\pgfusepath{fill}%
\end{pgfscope}%
\begin{pgfscope}%
\pgfpathrectangle{\pgfqpoint{4.255037in}{0.357758in}}{\pgfqpoint{0.172083in}{3.441662in}}%
\pgfusepath{clip}%
\pgfsetbuttcap%
\pgfsetroundjoin%
\definecolor{currentfill}{rgb}{0.631017,0.107699,0.608287}%
\pgfsetfillcolor{currentfill}%
\pgfsetlinewidth{0.000000pt}%
\definecolor{currentstroke}{rgb}{0.000000,0.000000,0.000000}%
\pgfsetstrokecolor{currentstroke}%
\pgfsetdash{}{0pt}%
\pgfpathmoveto{\pgfqpoint{4.255037in}{1.554273in}}%
\pgfpathlineto{\pgfqpoint{4.427120in}{1.554273in}}%
\pgfpathlineto{\pgfqpoint{4.427120in}{1.567717in}}%
\pgfpathlineto{\pgfqpoint{4.255037in}{1.567717in}}%
\pgfpathlineto{\pgfqpoint{4.255037in}{1.554273in}}%
\pgfusepath{fill}%
\end{pgfscope}%
\begin{pgfscope}%
\pgfpathrectangle{\pgfqpoint{4.255037in}{0.357758in}}{\pgfqpoint{0.172083in}{3.441662in}}%
\pgfusepath{clip}%
\pgfsetbuttcap%
\pgfsetroundjoin%
\definecolor{currentfill}{rgb}{0.636008,0.112092,0.605205}%
\pgfsetfillcolor{currentfill}%
\pgfsetlinewidth{0.000000pt}%
\definecolor{currentstroke}{rgb}{0.000000,0.000000,0.000000}%
\pgfsetstrokecolor{currentstroke}%
\pgfsetdash{}{0pt}%
\pgfpathmoveto{\pgfqpoint{4.255037in}{1.567717in}}%
\pgfpathlineto{\pgfqpoint{4.427120in}{1.567717in}}%
\pgfpathlineto{\pgfqpoint{4.427120in}{1.581161in}}%
\pgfpathlineto{\pgfqpoint{4.255037in}{1.581161in}}%
\pgfpathlineto{\pgfqpoint{4.255037in}{1.567717in}}%
\pgfusepath{fill}%
\end{pgfscope}%
\begin{pgfscope}%
\pgfpathrectangle{\pgfqpoint{4.255037in}{0.357758in}}{\pgfqpoint{0.172083in}{3.441662in}}%
\pgfusepath{clip}%
\pgfsetbuttcap%
\pgfsetroundjoin%
\definecolor{currentfill}{rgb}{0.640959,0.116492,0.602065}%
\pgfsetfillcolor{currentfill}%
\pgfsetlinewidth{0.000000pt}%
\definecolor{currentstroke}{rgb}{0.000000,0.000000,0.000000}%
\pgfsetstrokecolor{currentstroke}%
\pgfsetdash{}{0pt}%
\pgfpathmoveto{\pgfqpoint{4.255037in}{1.581161in}}%
\pgfpathlineto{\pgfqpoint{4.427120in}{1.581161in}}%
\pgfpathlineto{\pgfqpoint{4.427120in}{1.594605in}}%
\pgfpathlineto{\pgfqpoint{4.255037in}{1.594605in}}%
\pgfpathlineto{\pgfqpoint{4.255037in}{1.581161in}}%
\pgfusepath{fill}%
\end{pgfscope}%
\begin{pgfscope}%
\pgfpathrectangle{\pgfqpoint{4.255037in}{0.357758in}}{\pgfqpoint{0.172083in}{3.441662in}}%
\pgfusepath{clip}%
\pgfsetbuttcap%
\pgfsetroundjoin%
\definecolor{currentfill}{rgb}{0.645872,0.120898,0.598867}%
\pgfsetfillcolor{currentfill}%
\pgfsetlinewidth{0.000000pt}%
\definecolor{currentstroke}{rgb}{0.000000,0.000000,0.000000}%
\pgfsetstrokecolor{currentstroke}%
\pgfsetdash{}{0pt}%
\pgfpathmoveto{\pgfqpoint{4.255037in}{1.594605in}}%
\pgfpathlineto{\pgfqpoint{4.427120in}{1.594605in}}%
\pgfpathlineto{\pgfqpoint{4.427120in}{1.608049in}}%
\pgfpathlineto{\pgfqpoint{4.255037in}{1.608049in}}%
\pgfpathlineto{\pgfqpoint{4.255037in}{1.594605in}}%
\pgfusepath{fill}%
\end{pgfscope}%
\begin{pgfscope}%
\pgfpathrectangle{\pgfqpoint{4.255037in}{0.357758in}}{\pgfqpoint{0.172083in}{3.441662in}}%
\pgfusepath{clip}%
\pgfsetbuttcap%
\pgfsetroundjoin%
\definecolor{currentfill}{rgb}{0.650746,0.125309,0.595617}%
\pgfsetfillcolor{currentfill}%
\pgfsetlinewidth{0.000000pt}%
\definecolor{currentstroke}{rgb}{0.000000,0.000000,0.000000}%
\pgfsetstrokecolor{currentstroke}%
\pgfsetdash{}{0pt}%
\pgfpathmoveto{\pgfqpoint{4.255037in}{1.608049in}}%
\pgfpathlineto{\pgfqpoint{4.427120in}{1.608049in}}%
\pgfpathlineto{\pgfqpoint{4.427120in}{1.621493in}}%
\pgfpathlineto{\pgfqpoint{4.255037in}{1.621493in}}%
\pgfpathlineto{\pgfqpoint{4.255037in}{1.608049in}}%
\pgfusepath{fill}%
\end{pgfscope}%
\begin{pgfscope}%
\pgfpathrectangle{\pgfqpoint{4.255037in}{0.357758in}}{\pgfqpoint{0.172083in}{3.441662in}}%
\pgfusepath{clip}%
\pgfsetbuttcap%
\pgfsetroundjoin%
\definecolor{currentfill}{rgb}{0.655580,0.129725,0.592317}%
\pgfsetfillcolor{currentfill}%
\pgfsetlinewidth{0.000000pt}%
\definecolor{currentstroke}{rgb}{0.000000,0.000000,0.000000}%
\pgfsetstrokecolor{currentstroke}%
\pgfsetdash{}{0pt}%
\pgfpathmoveto{\pgfqpoint{4.255037in}{1.621493in}}%
\pgfpathlineto{\pgfqpoint{4.427120in}{1.621493in}}%
\pgfpathlineto{\pgfqpoint{4.427120in}{1.634937in}}%
\pgfpathlineto{\pgfqpoint{4.255037in}{1.634937in}}%
\pgfpathlineto{\pgfqpoint{4.255037in}{1.621493in}}%
\pgfusepath{fill}%
\end{pgfscope}%
\begin{pgfscope}%
\pgfpathrectangle{\pgfqpoint{4.255037in}{0.357758in}}{\pgfqpoint{0.172083in}{3.441662in}}%
\pgfusepath{clip}%
\pgfsetbuttcap%
\pgfsetroundjoin%
\definecolor{currentfill}{rgb}{0.660374,0.134144,0.588971}%
\pgfsetfillcolor{currentfill}%
\pgfsetlinewidth{0.000000pt}%
\definecolor{currentstroke}{rgb}{0.000000,0.000000,0.000000}%
\pgfsetstrokecolor{currentstroke}%
\pgfsetdash{}{0pt}%
\pgfpathmoveto{\pgfqpoint{4.255037in}{1.634937in}}%
\pgfpathlineto{\pgfqpoint{4.427120in}{1.634937in}}%
\pgfpathlineto{\pgfqpoint{4.427120in}{1.648381in}}%
\pgfpathlineto{\pgfqpoint{4.255037in}{1.648381in}}%
\pgfpathlineto{\pgfqpoint{4.255037in}{1.634937in}}%
\pgfusepath{fill}%
\end{pgfscope}%
\begin{pgfscope}%
\pgfpathrectangle{\pgfqpoint{4.255037in}{0.357758in}}{\pgfqpoint{0.172083in}{3.441662in}}%
\pgfusepath{clip}%
\pgfsetbuttcap%
\pgfsetroundjoin%
\definecolor{currentfill}{rgb}{0.665129,0.138566,0.585582}%
\pgfsetfillcolor{currentfill}%
\pgfsetlinewidth{0.000000pt}%
\definecolor{currentstroke}{rgb}{0.000000,0.000000,0.000000}%
\pgfsetstrokecolor{currentstroke}%
\pgfsetdash{}{0pt}%
\pgfpathmoveto{\pgfqpoint{4.255037in}{1.648381in}}%
\pgfpathlineto{\pgfqpoint{4.427120in}{1.648381in}}%
\pgfpathlineto{\pgfqpoint{4.427120in}{1.661825in}}%
\pgfpathlineto{\pgfqpoint{4.255037in}{1.661825in}}%
\pgfpathlineto{\pgfqpoint{4.255037in}{1.648381in}}%
\pgfusepath{fill}%
\end{pgfscope}%
\begin{pgfscope}%
\pgfpathrectangle{\pgfqpoint{4.255037in}{0.357758in}}{\pgfqpoint{0.172083in}{3.441662in}}%
\pgfusepath{clip}%
\pgfsetbuttcap%
\pgfsetroundjoin%
\definecolor{currentfill}{rgb}{0.669845,0.142992,0.582154}%
\pgfsetfillcolor{currentfill}%
\pgfsetlinewidth{0.000000pt}%
\definecolor{currentstroke}{rgb}{0.000000,0.000000,0.000000}%
\pgfsetstrokecolor{currentstroke}%
\pgfsetdash{}{0pt}%
\pgfpathmoveto{\pgfqpoint{4.255037in}{1.661825in}}%
\pgfpathlineto{\pgfqpoint{4.427120in}{1.661825in}}%
\pgfpathlineto{\pgfqpoint{4.427120in}{1.675269in}}%
\pgfpathlineto{\pgfqpoint{4.255037in}{1.675269in}}%
\pgfpathlineto{\pgfqpoint{4.255037in}{1.661825in}}%
\pgfusepath{fill}%
\end{pgfscope}%
\begin{pgfscope}%
\pgfpathrectangle{\pgfqpoint{4.255037in}{0.357758in}}{\pgfqpoint{0.172083in}{3.441662in}}%
\pgfusepath{clip}%
\pgfsetbuttcap%
\pgfsetroundjoin%
\definecolor{currentfill}{rgb}{0.674522,0.147419,0.578688}%
\pgfsetfillcolor{currentfill}%
\pgfsetlinewidth{0.000000pt}%
\definecolor{currentstroke}{rgb}{0.000000,0.000000,0.000000}%
\pgfsetstrokecolor{currentstroke}%
\pgfsetdash{}{0pt}%
\pgfpathmoveto{\pgfqpoint{4.255037in}{1.675269in}}%
\pgfpathlineto{\pgfqpoint{4.427120in}{1.675269in}}%
\pgfpathlineto{\pgfqpoint{4.427120in}{1.688713in}}%
\pgfpathlineto{\pgfqpoint{4.255037in}{1.688713in}}%
\pgfpathlineto{\pgfqpoint{4.255037in}{1.675269in}}%
\pgfusepath{fill}%
\end{pgfscope}%
\begin{pgfscope}%
\pgfpathrectangle{\pgfqpoint{4.255037in}{0.357758in}}{\pgfqpoint{0.172083in}{3.441662in}}%
\pgfusepath{clip}%
\pgfsetbuttcap%
\pgfsetroundjoin%
\definecolor{currentfill}{rgb}{0.679160,0.151848,0.575189}%
\pgfsetfillcolor{currentfill}%
\pgfsetlinewidth{0.000000pt}%
\definecolor{currentstroke}{rgb}{0.000000,0.000000,0.000000}%
\pgfsetstrokecolor{currentstroke}%
\pgfsetdash{}{0pt}%
\pgfpathmoveto{\pgfqpoint{4.255037in}{1.688713in}}%
\pgfpathlineto{\pgfqpoint{4.427120in}{1.688713in}}%
\pgfpathlineto{\pgfqpoint{4.427120in}{1.702157in}}%
\pgfpathlineto{\pgfqpoint{4.255037in}{1.702157in}}%
\pgfpathlineto{\pgfqpoint{4.255037in}{1.688713in}}%
\pgfusepath{fill}%
\end{pgfscope}%
\begin{pgfscope}%
\pgfpathrectangle{\pgfqpoint{4.255037in}{0.357758in}}{\pgfqpoint{0.172083in}{3.441662in}}%
\pgfusepath{clip}%
\pgfsetbuttcap%
\pgfsetroundjoin%
\definecolor{currentfill}{rgb}{0.683758,0.156278,0.571660}%
\pgfsetfillcolor{currentfill}%
\pgfsetlinewidth{0.000000pt}%
\definecolor{currentstroke}{rgb}{0.000000,0.000000,0.000000}%
\pgfsetstrokecolor{currentstroke}%
\pgfsetdash{}{0pt}%
\pgfpathmoveto{\pgfqpoint{4.255037in}{1.702157in}}%
\pgfpathlineto{\pgfqpoint{4.427120in}{1.702157in}}%
\pgfpathlineto{\pgfqpoint{4.427120in}{1.715601in}}%
\pgfpathlineto{\pgfqpoint{4.255037in}{1.715601in}}%
\pgfpathlineto{\pgfqpoint{4.255037in}{1.702157in}}%
\pgfusepath{fill}%
\end{pgfscope}%
\begin{pgfscope}%
\pgfpathrectangle{\pgfqpoint{4.255037in}{0.357758in}}{\pgfqpoint{0.172083in}{3.441662in}}%
\pgfusepath{clip}%
\pgfsetbuttcap%
\pgfsetroundjoin%
\definecolor{currentfill}{rgb}{0.688318,0.160709,0.568103}%
\pgfsetfillcolor{currentfill}%
\pgfsetlinewidth{0.000000pt}%
\definecolor{currentstroke}{rgb}{0.000000,0.000000,0.000000}%
\pgfsetstrokecolor{currentstroke}%
\pgfsetdash{}{0pt}%
\pgfpathmoveto{\pgfqpoint{4.255037in}{1.715601in}}%
\pgfpathlineto{\pgfqpoint{4.427120in}{1.715601in}}%
\pgfpathlineto{\pgfqpoint{4.427120in}{1.729045in}}%
\pgfpathlineto{\pgfqpoint{4.255037in}{1.729045in}}%
\pgfpathlineto{\pgfqpoint{4.255037in}{1.715601in}}%
\pgfusepath{fill}%
\end{pgfscope}%
\begin{pgfscope}%
\pgfpathrectangle{\pgfqpoint{4.255037in}{0.357758in}}{\pgfqpoint{0.172083in}{3.441662in}}%
\pgfusepath{clip}%
\pgfsetbuttcap%
\pgfsetroundjoin%
\definecolor{currentfill}{rgb}{0.692840,0.165141,0.564522}%
\pgfsetfillcolor{currentfill}%
\pgfsetlinewidth{0.000000pt}%
\definecolor{currentstroke}{rgb}{0.000000,0.000000,0.000000}%
\pgfsetstrokecolor{currentstroke}%
\pgfsetdash{}{0pt}%
\pgfpathmoveto{\pgfqpoint{4.255037in}{1.729045in}}%
\pgfpathlineto{\pgfqpoint{4.427120in}{1.729045in}}%
\pgfpathlineto{\pgfqpoint{4.427120in}{1.742489in}}%
\pgfpathlineto{\pgfqpoint{4.255037in}{1.742489in}}%
\pgfpathlineto{\pgfqpoint{4.255037in}{1.729045in}}%
\pgfusepath{fill}%
\end{pgfscope}%
\begin{pgfscope}%
\pgfpathrectangle{\pgfqpoint{4.255037in}{0.357758in}}{\pgfqpoint{0.172083in}{3.441662in}}%
\pgfusepath{clip}%
\pgfsetbuttcap%
\pgfsetroundjoin%
\definecolor{currentfill}{rgb}{0.697324,0.169573,0.560919}%
\pgfsetfillcolor{currentfill}%
\pgfsetlinewidth{0.000000pt}%
\definecolor{currentstroke}{rgb}{0.000000,0.000000,0.000000}%
\pgfsetstrokecolor{currentstroke}%
\pgfsetdash{}{0pt}%
\pgfpathmoveto{\pgfqpoint{4.255037in}{1.742489in}}%
\pgfpathlineto{\pgfqpoint{4.427120in}{1.742489in}}%
\pgfpathlineto{\pgfqpoint{4.427120in}{1.755933in}}%
\pgfpathlineto{\pgfqpoint{4.255037in}{1.755933in}}%
\pgfpathlineto{\pgfqpoint{4.255037in}{1.742489in}}%
\pgfusepath{fill}%
\end{pgfscope}%
\begin{pgfscope}%
\pgfpathrectangle{\pgfqpoint{4.255037in}{0.357758in}}{\pgfqpoint{0.172083in}{3.441662in}}%
\pgfusepath{clip}%
\pgfsetbuttcap%
\pgfsetroundjoin%
\definecolor{currentfill}{rgb}{0.701769,0.174005,0.557296}%
\pgfsetfillcolor{currentfill}%
\pgfsetlinewidth{0.000000pt}%
\definecolor{currentstroke}{rgb}{0.000000,0.000000,0.000000}%
\pgfsetstrokecolor{currentstroke}%
\pgfsetdash{}{0pt}%
\pgfpathmoveto{\pgfqpoint{4.255037in}{1.755933in}}%
\pgfpathlineto{\pgfqpoint{4.427120in}{1.755933in}}%
\pgfpathlineto{\pgfqpoint{4.427120in}{1.769377in}}%
\pgfpathlineto{\pgfqpoint{4.255037in}{1.769377in}}%
\pgfpathlineto{\pgfqpoint{4.255037in}{1.755933in}}%
\pgfusepath{fill}%
\end{pgfscope}%
\begin{pgfscope}%
\pgfpathrectangle{\pgfqpoint{4.255037in}{0.357758in}}{\pgfqpoint{0.172083in}{3.441662in}}%
\pgfusepath{clip}%
\pgfsetbuttcap%
\pgfsetroundjoin%
\definecolor{currentfill}{rgb}{0.706178,0.178437,0.553657}%
\pgfsetfillcolor{currentfill}%
\pgfsetlinewidth{0.000000pt}%
\definecolor{currentstroke}{rgb}{0.000000,0.000000,0.000000}%
\pgfsetstrokecolor{currentstroke}%
\pgfsetdash{}{0pt}%
\pgfpathmoveto{\pgfqpoint{4.255037in}{1.769377in}}%
\pgfpathlineto{\pgfqpoint{4.427120in}{1.769377in}}%
\pgfpathlineto{\pgfqpoint{4.427120in}{1.782821in}}%
\pgfpathlineto{\pgfqpoint{4.255037in}{1.782821in}}%
\pgfpathlineto{\pgfqpoint{4.255037in}{1.769377in}}%
\pgfusepath{fill}%
\end{pgfscope}%
\begin{pgfscope}%
\pgfpathrectangle{\pgfqpoint{4.255037in}{0.357758in}}{\pgfqpoint{0.172083in}{3.441662in}}%
\pgfusepath{clip}%
\pgfsetbuttcap%
\pgfsetroundjoin%
\definecolor{currentfill}{rgb}{0.710549,0.182868,0.550004}%
\pgfsetfillcolor{currentfill}%
\pgfsetlinewidth{0.000000pt}%
\definecolor{currentstroke}{rgb}{0.000000,0.000000,0.000000}%
\pgfsetstrokecolor{currentstroke}%
\pgfsetdash{}{0pt}%
\pgfpathmoveto{\pgfqpoint{4.255037in}{1.782821in}}%
\pgfpathlineto{\pgfqpoint{4.427120in}{1.782821in}}%
\pgfpathlineto{\pgfqpoint{4.427120in}{1.796265in}}%
\pgfpathlineto{\pgfqpoint{4.255037in}{1.796265in}}%
\pgfpathlineto{\pgfqpoint{4.255037in}{1.782821in}}%
\pgfusepath{fill}%
\end{pgfscope}%
\begin{pgfscope}%
\pgfpathrectangle{\pgfqpoint{4.255037in}{0.357758in}}{\pgfqpoint{0.172083in}{3.441662in}}%
\pgfusepath{clip}%
\pgfsetbuttcap%
\pgfsetroundjoin%
\definecolor{currentfill}{rgb}{0.714883,0.187299,0.546338}%
\pgfsetfillcolor{currentfill}%
\pgfsetlinewidth{0.000000pt}%
\definecolor{currentstroke}{rgb}{0.000000,0.000000,0.000000}%
\pgfsetstrokecolor{currentstroke}%
\pgfsetdash{}{0pt}%
\pgfpathmoveto{\pgfqpoint{4.255037in}{1.796265in}}%
\pgfpathlineto{\pgfqpoint{4.427120in}{1.796265in}}%
\pgfpathlineto{\pgfqpoint{4.427120in}{1.809709in}}%
\pgfpathlineto{\pgfqpoint{4.255037in}{1.809709in}}%
\pgfpathlineto{\pgfqpoint{4.255037in}{1.796265in}}%
\pgfusepath{fill}%
\end{pgfscope}%
\begin{pgfscope}%
\pgfpathrectangle{\pgfqpoint{4.255037in}{0.357758in}}{\pgfqpoint{0.172083in}{3.441662in}}%
\pgfusepath{clip}%
\pgfsetbuttcap%
\pgfsetroundjoin%
\definecolor{currentfill}{rgb}{0.719181,0.191729,0.542663}%
\pgfsetfillcolor{currentfill}%
\pgfsetlinewidth{0.000000pt}%
\definecolor{currentstroke}{rgb}{0.000000,0.000000,0.000000}%
\pgfsetstrokecolor{currentstroke}%
\pgfsetdash{}{0pt}%
\pgfpathmoveto{\pgfqpoint{4.255037in}{1.809709in}}%
\pgfpathlineto{\pgfqpoint{4.427120in}{1.809709in}}%
\pgfpathlineto{\pgfqpoint{4.427120in}{1.823153in}}%
\pgfpathlineto{\pgfqpoint{4.255037in}{1.823153in}}%
\pgfpathlineto{\pgfqpoint{4.255037in}{1.809709in}}%
\pgfusepath{fill}%
\end{pgfscope}%
\begin{pgfscope}%
\pgfpathrectangle{\pgfqpoint{4.255037in}{0.357758in}}{\pgfqpoint{0.172083in}{3.441662in}}%
\pgfusepath{clip}%
\pgfsetbuttcap%
\pgfsetroundjoin%
\definecolor{currentfill}{rgb}{0.723444,0.196158,0.538981}%
\pgfsetfillcolor{currentfill}%
\pgfsetlinewidth{0.000000pt}%
\definecolor{currentstroke}{rgb}{0.000000,0.000000,0.000000}%
\pgfsetstrokecolor{currentstroke}%
\pgfsetdash{}{0pt}%
\pgfpathmoveto{\pgfqpoint{4.255037in}{1.823153in}}%
\pgfpathlineto{\pgfqpoint{4.427120in}{1.823153in}}%
\pgfpathlineto{\pgfqpoint{4.427120in}{1.836597in}}%
\pgfpathlineto{\pgfqpoint{4.255037in}{1.836597in}}%
\pgfpathlineto{\pgfqpoint{4.255037in}{1.823153in}}%
\pgfusepath{fill}%
\end{pgfscope}%
\begin{pgfscope}%
\pgfpathrectangle{\pgfqpoint{4.255037in}{0.357758in}}{\pgfqpoint{0.172083in}{3.441662in}}%
\pgfusepath{clip}%
\pgfsetbuttcap%
\pgfsetroundjoin%
\definecolor{currentfill}{rgb}{0.727670,0.200586,0.535293}%
\pgfsetfillcolor{currentfill}%
\pgfsetlinewidth{0.000000pt}%
\definecolor{currentstroke}{rgb}{0.000000,0.000000,0.000000}%
\pgfsetstrokecolor{currentstroke}%
\pgfsetdash{}{0pt}%
\pgfpathmoveto{\pgfqpoint{4.255037in}{1.836597in}}%
\pgfpathlineto{\pgfqpoint{4.427120in}{1.836597in}}%
\pgfpathlineto{\pgfqpoint{4.427120in}{1.850041in}}%
\pgfpathlineto{\pgfqpoint{4.255037in}{1.850041in}}%
\pgfpathlineto{\pgfqpoint{4.255037in}{1.836597in}}%
\pgfusepath{fill}%
\end{pgfscope}%
\begin{pgfscope}%
\pgfpathrectangle{\pgfqpoint{4.255037in}{0.357758in}}{\pgfqpoint{0.172083in}{3.441662in}}%
\pgfusepath{clip}%
\pgfsetbuttcap%
\pgfsetroundjoin%
\definecolor{currentfill}{rgb}{0.731862,0.205013,0.531601}%
\pgfsetfillcolor{currentfill}%
\pgfsetlinewidth{0.000000pt}%
\definecolor{currentstroke}{rgb}{0.000000,0.000000,0.000000}%
\pgfsetstrokecolor{currentstroke}%
\pgfsetdash{}{0pt}%
\pgfpathmoveto{\pgfqpoint{4.255037in}{1.850041in}}%
\pgfpathlineto{\pgfqpoint{4.427120in}{1.850041in}}%
\pgfpathlineto{\pgfqpoint{4.427120in}{1.863485in}}%
\pgfpathlineto{\pgfqpoint{4.255037in}{1.863485in}}%
\pgfpathlineto{\pgfqpoint{4.255037in}{1.850041in}}%
\pgfusepath{fill}%
\end{pgfscope}%
\begin{pgfscope}%
\pgfpathrectangle{\pgfqpoint{4.255037in}{0.357758in}}{\pgfqpoint{0.172083in}{3.441662in}}%
\pgfusepath{clip}%
\pgfsetbuttcap%
\pgfsetroundjoin%
\definecolor{currentfill}{rgb}{0.736019,0.209439,0.527908}%
\pgfsetfillcolor{currentfill}%
\pgfsetlinewidth{0.000000pt}%
\definecolor{currentstroke}{rgb}{0.000000,0.000000,0.000000}%
\pgfsetstrokecolor{currentstroke}%
\pgfsetdash{}{0pt}%
\pgfpathmoveto{\pgfqpoint{4.255037in}{1.863485in}}%
\pgfpathlineto{\pgfqpoint{4.427120in}{1.863485in}}%
\pgfpathlineto{\pgfqpoint{4.427120in}{1.876929in}}%
\pgfpathlineto{\pgfqpoint{4.255037in}{1.876929in}}%
\pgfpathlineto{\pgfqpoint{4.255037in}{1.863485in}}%
\pgfusepath{fill}%
\end{pgfscope}%
\begin{pgfscope}%
\pgfpathrectangle{\pgfqpoint{4.255037in}{0.357758in}}{\pgfqpoint{0.172083in}{3.441662in}}%
\pgfusepath{clip}%
\pgfsetbuttcap%
\pgfsetroundjoin%
\definecolor{currentfill}{rgb}{0.740143,0.213864,0.524216}%
\pgfsetfillcolor{currentfill}%
\pgfsetlinewidth{0.000000pt}%
\definecolor{currentstroke}{rgb}{0.000000,0.000000,0.000000}%
\pgfsetstrokecolor{currentstroke}%
\pgfsetdash{}{0pt}%
\pgfpathmoveto{\pgfqpoint{4.255037in}{1.876929in}}%
\pgfpathlineto{\pgfqpoint{4.427120in}{1.876929in}}%
\pgfpathlineto{\pgfqpoint{4.427120in}{1.890373in}}%
\pgfpathlineto{\pgfqpoint{4.255037in}{1.890373in}}%
\pgfpathlineto{\pgfqpoint{4.255037in}{1.876929in}}%
\pgfusepath{fill}%
\end{pgfscope}%
\begin{pgfscope}%
\pgfpathrectangle{\pgfqpoint{4.255037in}{0.357758in}}{\pgfqpoint{0.172083in}{3.441662in}}%
\pgfusepath{clip}%
\pgfsetbuttcap%
\pgfsetroundjoin%
\definecolor{currentfill}{rgb}{0.744232,0.218288,0.520524}%
\pgfsetfillcolor{currentfill}%
\pgfsetlinewidth{0.000000pt}%
\definecolor{currentstroke}{rgb}{0.000000,0.000000,0.000000}%
\pgfsetstrokecolor{currentstroke}%
\pgfsetdash{}{0pt}%
\pgfpathmoveto{\pgfqpoint{4.255037in}{1.890373in}}%
\pgfpathlineto{\pgfqpoint{4.427120in}{1.890373in}}%
\pgfpathlineto{\pgfqpoint{4.427120in}{1.903817in}}%
\pgfpathlineto{\pgfqpoint{4.255037in}{1.903817in}}%
\pgfpathlineto{\pgfqpoint{4.255037in}{1.890373in}}%
\pgfusepath{fill}%
\end{pgfscope}%
\begin{pgfscope}%
\pgfpathrectangle{\pgfqpoint{4.255037in}{0.357758in}}{\pgfqpoint{0.172083in}{3.441662in}}%
\pgfusepath{clip}%
\pgfsetbuttcap%
\pgfsetroundjoin%
\definecolor{currentfill}{rgb}{0.748289,0.222711,0.516834}%
\pgfsetfillcolor{currentfill}%
\pgfsetlinewidth{0.000000pt}%
\definecolor{currentstroke}{rgb}{0.000000,0.000000,0.000000}%
\pgfsetstrokecolor{currentstroke}%
\pgfsetdash{}{0pt}%
\pgfpathmoveto{\pgfqpoint{4.255037in}{1.903817in}}%
\pgfpathlineto{\pgfqpoint{4.427120in}{1.903817in}}%
\pgfpathlineto{\pgfqpoint{4.427120in}{1.917261in}}%
\pgfpathlineto{\pgfqpoint{4.255037in}{1.917261in}}%
\pgfpathlineto{\pgfqpoint{4.255037in}{1.903817in}}%
\pgfusepath{fill}%
\end{pgfscope}%
\begin{pgfscope}%
\pgfpathrectangle{\pgfqpoint{4.255037in}{0.357758in}}{\pgfqpoint{0.172083in}{3.441662in}}%
\pgfusepath{clip}%
\pgfsetbuttcap%
\pgfsetroundjoin%
\definecolor{currentfill}{rgb}{0.752312,0.227133,0.513149}%
\pgfsetfillcolor{currentfill}%
\pgfsetlinewidth{0.000000pt}%
\definecolor{currentstroke}{rgb}{0.000000,0.000000,0.000000}%
\pgfsetstrokecolor{currentstroke}%
\pgfsetdash{}{0pt}%
\pgfpathmoveto{\pgfqpoint{4.255037in}{1.917261in}}%
\pgfpathlineto{\pgfqpoint{4.427120in}{1.917261in}}%
\pgfpathlineto{\pgfqpoint{4.427120in}{1.930705in}}%
\pgfpathlineto{\pgfqpoint{4.255037in}{1.930705in}}%
\pgfpathlineto{\pgfqpoint{4.255037in}{1.917261in}}%
\pgfusepath{fill}%
\end{pgfscope}%
\begin{pgfscope}%
\pgfpathrectangle{\pgfqpoint{4.255037in}{0.357758in}}{\pgfqpoint{0.172083in}{3.441662in}}%
\pgfusepath{clip}%
\pgfsetbuttcap%
\pgfsetroundjoin%
\definecolor{currentfill}{rgb}{0.756304,0.231555,0.509468}%
\pgfsetfillcolor{currentfill}%
\pgfsetlinewidth{0.000000pt}%
\definecolor{currentstroke}{rgb}{0.000000,0.000000,0.000000}%
\pgfsetstrokecolor{currentstroke}%
\pgfsetdash{}{0pt}%
\pgfpathmoveto{\pgfqpoint{4.255037in}{1.930705in}}%
\pgfpathlineto{\pgfqpoint{4.427120in}{1.930705in}}%
\pgfpathlineto{\pgfqpoint{4.427120in}{1.944149in}}%
\pgfpathlineto{\pgfqpoint{4.255037in}{1.944149in}}%
\pgfpathlineto{\pgfqpoint{4.255037in}{1.930705in}}%
\pgfusepath{fill}%
\end{pgfscope}%
\begin{pgfscope}%
\pgfpathrectangle{\pgfqpoint{4.255037in}{0.357758in}}{\pgfqpoint{0.172083in}{3.441662in}}%
\pgfusepath{clip}%
\pgfsetbuttcap%
\pgfsetroundjoin%
\definecolor{currentfill}{rgb}{0.760264,0.235976,0.505794}%
\pgfsetfillcolor{currentfill}%
\pgfsetlinewidth{0.000000pt}%
\definecolor{currentstroke}{rgb}{0.000000,0.000000,0.000000}%
\pgfsetstrokecolor{currentstroke}%
\pgfsetdash{}{0pt}%
\pgfpathmoveto{\pgfqpoint{4.255037in}{1.944149in}}%
\pgfpathlineto{\pgfqpoint{4.427120in}{1.944149in}}%
\pgfpathlineto{\pgfqpoint{4.427120in}{1.957593in}}%
\pgfpathlineto{\pgfqpoint{4.255037in}{1.957593in}}%
\pgfpathlineto{\pgfqpoint{4.255037in}{1.944149in}}%
\pgfusepath{fill}%
\end{pgfscope}%
\begin{pgfscope}%
\pgfpathrectangle{\pgfqpoint{4.255037in}{0.357758in}}{\pgfqpoint{0.172083in}{3.441662in}}%
\pgfusepath{clip}%
\pgfsetbuttcap%
\pgfsetroundjoin%
\definecolor{currentfill}{rgb}{0.764193,0.240396,0.502126}%
\pgfsetfillcolor{currentfill}%
\pgfsetlinewidth{0.000000pt}%
\definecolor{currentstroke}{rgb}{0.000000,0.000000,0.000000}%
\pgfsetstrokecolor{currentstroke}%
\pgfsetdash{}{0pt}%
\pgfpathmoveto{\pgfqpoint{4.255037in}{1.957593in}}%
\pgfpathlineto{\pgfqpoint{4.427120in}{1.957593in}}%
\pgfpathlineto{\pgfqpoint{4.427120in}{1.971037in}}%
\pgfpathlineto{\pgfqpoint{4.255037in}{1.971037in}}%
\pgfpathlineto{\pgfqpoint{4.255037in}{1.957593in}}%
\pgfusepath{fill}%
\end{pgfscope}%
\begin{pgfscope}%
\pgfpathrectangle{\pgfqpoint{4.255037in}{0.357758in}}{\pgfqpoint{0.172083in}{3.441662in}}%
\pgfusepath{clip}%
\pgfsetbuttcap%
\pgfsetroundjoin%
\definecolor{currentfill}{rgb}{0.768090,0.244817,0.498465}%
\pgfsetfillcolor{currentfill}%
\pgfsetlinewidth{0.000000pt}%
\definecolor{currentstroke}{rgb}{0.000000,0.000000,0.000000}%
\pgfsetstrokecolor{currentstroke}%
\pgfsetdash{}{0pt}%
\pgfpathmoveto{\pgfqpoint{4.255037in}{1.971037in}}%
\pgfpathlineto{\pgfqpoint{4.427120in}{1.971037in}}%
\pgfpathlineto{\pgfqpoint{4.427120in}{1.984481in}}%
\pgfpathlineto{\pgfqpoint{4.255037in}{1.984481in}}%
\pgfpathlineto{\pgfqpoint{4.255037in}{1.971037in}}%
\pgfusepath{fill}%
\end{pgfscope}%
\begin{pgfscope}%
\pgfpathrectangle{\pgfqpoint{4.255037in}{0.357758in}}{\pgfqpoint{0.172083in}{3.441662in}}%
\pgfusepath{clip}%
\pgfsetbuttcap%
\pgfsetroundjoin%
\definecolor{currentfill}{rgb}{0.771958,0.249237,0.494813}%
\pgfsetfillcolor{currentfill}%
\pgfsetlinewidth{0.000000pt}%
\definecolor{currentstroke}{rgb}{0.000000,0.000000,0.000000}%
\pgfsetstrokecolor{currentstroke}%
\pgfsetdash{}{0pt}%
\pgfpathmoveto{\pgfqpoint{4.255037in}{1.984481in}}%
\pgfpathlineto{\pgfqpoint{4.427120in}{1.984481in}}%
\pgfpathlineto{\pgfqpoint{4.427120in}{1.997925in}}%
\pgfpathlineto{\pgfqpoint{4.255037in}{1.997925in}}%
\pgfpathlineto{\pgfqpoint{4.255037in}{1.984481in}}%
\pgfusepath{fill}%
\end{pgfscope}%
\begin{pgfscope}%
\pgfpathrectangle{\pgfqpoint{4.255037in}{0.357758in}}{\pgfqpoint{0.172083in}{3.441662in}}%
\pgfusepath{clip}%
\pgfsetbuttcap%
\pgfsetroundjoin%
\definecolor{currentfill}{rgb}{0.775796,0.253658,0.491171}%
\pgfsetfillcolor{currentfill}%
\pgfsetlinewidth{0.000000pt}%
\definecolor{currentstroke}{rgb}{0.000000,0.000000,0.000000}%
\pgfsetstrokecolor{currentstroke}%
\pgfsetdash{}{0pt}%
\pgfpathmoveto{\pgfqpoint{4.255037in}{1.997925in}}%
\pgfpathlineto{\pgfqpoint{4.427120in}{1.997925in}}%
\pgfpathlineto{\pgfqpoint{4.427120in}{2.011369in}}%
\pgfpathlineto{\pgfqpoint{4.255037in}{2.011369in}}%
\pgfpathlineto{\pgfqpoint{4.255037in}{1.997925in}}%
\pgfusepath{fill}%
\end{pgfscope}%
\begin{pgfscope}%
\pgfpathrectangle{\pgfqpoint{4.255037in}{0.357758in}}{\pgfqpoint{0.172083in}{3.441662in}}%
\pgfusepath{clip}%
\pgfsetbuttcap%
\pgfsetroundjoin%
\definecolor{currentfill}{rgb}{0.779604,0.258078,0.487539}%
\pgfsetfillcolor{currentfill}%
\pgfsetlinewidth{0.000000pt}%
\definecolor{currentstroke}{rgb}{0.000000,0.000000,0.000000}%
\pgfsetstrokecolor{currentstroke}%
\pgfsetdash{}{0pt}%
\pgfpathmoveto{\pgfqpoint{4.255037in}{2.011369in}}%
\pgfpathlineto{\pgfqpoint{4.427120in}{2.011369in}}%
\pgfpathlineto{\pgfqpoint{4.427120in}{2.024813in}}%
\pgfpathlineto{\pgfqpoint{4.255037in}{2.024813in}}%
\pgfpathlineto{\pgfqpoint{4.255037in}{2.011369in}}%
\pgfusepath{fill}%
\end{pgfscope}%
\begin{pgfscope}%
\pgfpathrectangle{\pgfqpoint{4.255037in}{0.357758in}}{\pgfqpoint{0.172083in}{3.441662in}}%
\pgfusepath{clip}%
\pgfsetbuttcap%
\pgfsetroundjoin%
\definecolor{currentfill}{rgb}{0.783383,0.262500,0.483918}%
\pgfsetfillcolor{currentfill}%
\pgfsetlinewidth{0.000000pt}%
\definecolor{currentstroke}{rgb}{0.000000,0.000000,0.000000}%
\pgfsetstrokecolor{currentstroke}%
\pgfsetdash{}{0pt}%
\pgfpathmoveto{\pgfqpoint{4.255037in}{2.024813in}}%
\pgfpathlineto{\pgfqpoint{4.427120in}{2.024813in}}%
\pgfpathlineto{\pgfqpoint{4.427120in}{2.038257in}}%
\pgfpathlineto{\pgfqpoint{4.255037in}{2.038257in}}%
\pgfpathlineto{\pgfqpoint{4.255037in}{2.024813in}}%
\pgfusepath{fill}%
\end{pgfscope}%
\begin{pgfscope}%
\pgfpathrectangle{\pgfqpoint{4.255037in}{0.357758in}}{\pgfqpoint{0.172083in}{3.441662in}}%
\pgfusepath{clip}%
\pgfsetbuttcap%
\pgfsetroundjoin%
\definecolor{currentfill}{rgb}{0.787133,0.266922,0.480307}%
\pgfsetfillcolor{currentfill}%
\pgfsetlinewidth{0.000000pt}%
\definecolor{currentstroke}{rgb}{0.000000,0.000000,0.000000}%
\pgfsetstrokecolor{currentstroke}%
\pgfsetdash{}{0pt}%
\pgfpathmoveto{\pgfqpoint{4.255037in}{2.038257in}}%
\pgfpathlineto{\pgfqpoint{4.427120in}{2.038257in}}%
\pgfpathlineto{\pgfqpoint{4.427120in}{2.051701in}}%
\pgfpathlineto{\pgfqpoint{4.255037in}{2.051701in}}%
\pgfpathlineto{\pgfqpoint{4.255037in}{2.038257in}}%
\pgfusepath{fill}%
\end{pgfscope}%
\begin{pgfscope}%
\pgfpathrectangle{\pgfqpoint{4.255037in}{0.357758in}}{\pgfqpoint{0.172083in}{3.441662in}}%
\pgfusepath{clip}%
\pgfsetbuttcap%
\pgfsetroundjoin%
\definecolor{currentfill}{rgb}{0.790855,0.271345,0.476706}%
\pgfsetfillcolor{currentfill}%
\pgfsetlinewidth{0.000000pt}%
\definecolor{currentstroke}{rgb}{0.000000,0.000000,0.000000}%
\pgfsetstrokecolor{currentstroke}%
\pgfsetdash{}{0pt}%
\pgfpathmoveto{\pgfqpoint{4.255037in}{2.051701in}}%
\pgfpathlineto{\pgfqpoint{4.427120in}{2.051701in}}%
\pgfpathlineto{\pgfqpoint{4.427120in}{2.065145in}}%
\pgfpathlineto{\pgfqpoint{4.255037in}{2.065145in}}%
\pgfpathlineto{\pgfqpoint{4.255037in}{2.051701in}}%
\pgfusepath{fill}%
\end{pgfscope}%
\begin{pgfscope}%
\pgfpathrectangle{\pgfqpoint{4.255037in}{0.357758in}}{\pgfqpoint{0.172083in}{3.441662in}}%
\pgfusepath{clip}%
\pgfsetbuttcap%
\pgfsetroundjoin%
\definecolor{currentfill}{rgb}{0.794549,0.275770,0.473117}%
\pgfsetfillcolor{currentfill}%
\pgfsetlinewidth{0.000000pt}%
\definecolor{currentstroke}{rgb}{0.000000,0.000000,0.000000}%
\pgfsetstrokecolor{currentstroke}%
\pgfsetdash{}{0pt}%
\pgfpathmoveto{\pgfqpoint{4.255037in}{2.065145in}}%
\pgfpathlineto{\pgfqpoint{4.427120in}{2.065145in}}%
\pgfpathlineto{\pgfqpoint{4.427120in}{2.078589in}}%
\pgfpathlineto{\pgfqpoint{4.255037in}{2.078589in}}%
\pgfpathlineto{\pgfqpoint{4.255037in}{2.065145in}}%
\pgfusepath{fill}%
\end{pgfscope}%
\begin{pgfscope}%
\pgfpathrectangle{\pgfqpoint{4.255037in}{0.357758in}}{\pgfqpoint{0.172083in}{3.441662in}}%
\pgfusepath{clip}%
\pgfsetbuttcap%
\pgfsetroundjoin%
\definecolor{currentfill}{rgb}{0.798216,0.280197,0.469538}%
\pgfsetfillcolor{currentfill}%
\pgfsetlinewidth{0.000000pt}%
\definecolor{currentstroke}{rgb}{0.000000,0.000000,0.000000}%
\pgfsetstrokecolor{currentstroke}%
\pgfsetdash{}{0pt}%
\pgfpathmoveto{\pgfqpoint{4.255037in}{2.078589in}}%
\pgfpathlineto{\pgfqpoint{4.427120in}{2.078589in}}%
\pgfpathlineto{\pgfqpoint{4.427120in}{2.092033in}}%
\pgfpathlineto{\pgfqpoint{4.255037in}{2.092033in}}%
\pgfpathlineto{\pgfqpoint{4.255037in}{2.078589in}}%
\pgfusepath{fill}%
\end{pgfscope}%
\begin{pgfscope}%
\pgfpathrectangle{\pgfqpoint{4.255037in}{0.357758in}}{\pgfqpoint{0.172083in}{3.441662in}}%
\pgfusepath{clip}%
\pgfsetbuttcap%
\pgfsetroundjoin%
\definecolor{currentfill}{rgb}{0.801855,0.284626,0.465971}%
\pgfsetfillcolor{currentfill}%
\pgfsetlinewidth{0.000000pt}%
\definecolor{currentstroke}{rgb}{0.000000,0.000000,0.000000}%
\pgfsetstrokecolor{currentstroke}%
\pgfsetdash{}{0pt}%
\pgfpathmoveto{\pgfqpoint{4.255037in}{2.092033in}}%
\pgfpathlineto{\pgfqpoint{4.427120in}{2.092033in}}%
\pgfpathlineto{\pgfqpoint{4.427120in}{2.105477in}}%
\pgfpathlineto{\pgfqpoint{4.255037in}{2.105477in}}%
\pgfpathlineto{\pgfqpoint{4.255037in}{2.092033in}}%
\pgfusepath{fill}%
\end{pgfscope}%
\begin{pgfscope}%
\pgfpathrectangle{\pgfqpoint{4.255037in}{0.357758in}}{\pgfqpoint{0.172083in}{3.441662in}}%
\pgfusepath{clip}%
\pgfsetbuttcap%
\pgfsetroundjoin%
\definecolor{currentfill}{rgb}{0.805467,0.289057,0.462415}%
\pgfsetfillcolor{currentfill}%
\pgfsetlinewidth{0.000000pt}%
\definecolor{currentstroke}{rgb}{0.000000,0.000000,0.000000}%
\pgfsetstrokecolor{currentstroke}%
\pgfsetdash{}{0pt}%
\pgfpathmoveto{\pgfqpoint{4.255037in}{2.105477in}}%
\pgfpathlineto{\pgfqpoint{4.427120in}{2.105477in}}%
\pgfpathlineto{\pgfqpoint{4.427120in}{2.118921in}}%
\pgfpathlineto{\pgfqpoint{4.255037in}{2.118921in}}%
\pgfpathlineto{\pgfqpoint{4.255037in}{2.105477in}}%
\pgfusepath{fill}%
\end{pgfscope}%
\begin{pgfscope}%
\pgfpathrectangle{\pgfqpoint{4.255037in}{0.357758in}}{\pgfqpoint{0.172083in}{3.441662in}}%
\pgfusepath{clip}%
\pgfsetbuttcap%
\pgfsetroundjoin%
\definecolor{currentfill}{rgb}{0.809052,0.293491,0.458870}%
\pgfsetfillcolor{currentfill}%
\pgfsetlinewidth{0.000000pt}%
\definecolor{currentstroke}{rgb}{0.000000,0.000000,0.000000}%
\pgfsetstrokecolor{currentstroke}%
\pgfsetdash{}{0pt}%
\pgfpathmoveto{\pgfqpoint{4.255037in}{2.118921in}}%
\pgfpathlineto{\pgfqpoint{4.427120in}{2.118921in}}%
\pgfpathlineto{\pgfqpoint{4.427120in}{2.132365in}}%
\pgfpathlineto{\pgfqpoint{4.255037in}{2.132365in}}%
\pgfpathlineto{\pgfqpoint{4.255037in}{2.118921in}}%
\pgfusepath{fill}%
\end{pgfscope}%
\begin{pgfscope}%
\pgfpathrectangle{\pgfqpoint{4.255037in}{0.357758in}}{\pgfqpoint{0.172083in}{3.441662in}}%
\pgfusepath{clip}%
\pgfsetbuttcap%
\pgfsetroundjoin%
\definecolor{currentfill}{rgb}{0.812612,0.297928,0.455338}%
\pgfsetfillcolor{currentfill}%
\pgfsetlinewidth{0.000000pt}%
\definecolor{currentstroke}{rgb}{0.000000,0.000000,0.000000}%
\pgfsetstrokecolor{currentstroke}%
\pgfsetdash{}{0pt}%
\pgfpathmoveto{\pgfqpoint{4.255037in}{2.132365in}}%
\pgfpathlineto{\pgfqpoint{4.427120in}{2.132365in}}%
\pgfpathlineto{\pgfqpoint{4.427120in}{2.145809in}}%
\pgfpathlineto{\pgfqpoint{4.255037in}{2.145809in}}%
\pgfpathlineto{\pgfqpoint{4.255037in}{2.132365in}}%
\pgfusepath{fill}%
\end{pgfscope}%
\begin{pgfscope}%
\pgfpathrectangle{\pgfqpoint{4.255037in}{0.357758in}}{\pgfqpoint{0.172083in}{3.441662in}}%
\pgfusepath{clip}%
\pgfsetbuttcap%
\pgfsetroundjoin%
\definecolor{currentfill}{rgb}{0.816144,0.302368,0.451816}%
\pgfsetfillcolor{currentfill}%
\pgfsetlinewidth{0.000000pt}%
\definecolor{currentstroke}{rgb}{0.000000,0.000000,0.000000}%
\pgfsetstrokecolor{currentstroke}%
\pgfsetdash{}{0pt}%
\pgfpathmoveto{\pgfqpoint{4.255037in}{2.145809in}}%
\pgfpathlineto{\pgfqpoint{4.427120in}{2.145809in}}%
\pgfpathlineto{\pgfqpoint{4.427120in}{2.159253in}}%
\pgfpathlineto{\pgfqpoint{4.255037in}{2.159253in}}%
\pgfpathlineto{\pgfqpoint{4.255037in}{2.145809in}}%
\pgfusepath{fill}%
\end{pgfscope}%
\begin{pgfscope}%
\pgfpathrectangle{\pgfqpoint{4.255037in}{0.357758in}}{\pgfqpoint{0.172083in}{3.441662in}}%
\pgfusepath{clip}%
\pgfsetbuttcap%
\pgfsetroundjoin%
\definecolor{currentfill}{rgb}{0.819651,0.306812,0.448306}%
\pgfsetfillcolor{currentfill}%
\pgfsetlinewidth{0.000000pt}%
\definecolor{currentstroke}{rgb}{0.000000,0.000000,0.000000}%
\pgfsetstrokecolor{currentstroke}%
\pgfsetdash{}{0pt}%
\pgfpathmoveto{\pgfqpoint{4.255037in}{2.159253in}}%
\pgfpathlineto{\pgfqpoint{4.427120in}{2.159253in}}%
\pgfpathlineto{\pgfqpoint{4.427120in}{2.172697in}}%
\pgfpathlineto{\pgfqpoint{4.255037in}{2.172697in}}%
\pgfpathlineto{\pgfqpoint{4.255037in}{2.159253in}}%
\pgfusepath{fill}%
\end{pgfscope}%
\begin{pgfscope}%
\pgfpathrectangle{\pgfqpoint{4.255037in}{0.357758in}}{\pgfqpoint{0.172083in}{3.441662in}}%
\pgfusepath{clip}%
\pgfsetbuttcap%
\pgfsetroundjoin%
\definecolor{currentfill}{rgb}{0.823132,0.311261,0.444806}%
\pgfsetfillcolor{currentfill}%
\pgfsetlinewidth{0.000000pt}%
\definecolor{currentstroke}{rgb}{0.000000,0.000000,0.000000}%
\pgfsetstrokecolor{currentstroke}%
\pgfsetdash{}{0pt}%
\pgfpathmoveto{\pgfqpoint{4.255037in}{2.172697in}}%
\pgfpathlineto{\pgfqpoint{4.427120in}{2.172697in}}%
\pgfpathlineto{\pgfqpoint{4.427120in}{2.186141in}}%
\pgfpathlineto{\pgfqpoint{4.255037in}{2.186141in}}%
\pgfpathlineto{\pgfqpoint{4.255037in}{2.172697in}}%
\pgfusepath{fill}%
\end{pgfscope}%
\begin{pgfscope}%
\pgfpathrectangle{\pgfqpoint{4.255037in}{0.357758in}}{\pgfqpoint{0.172083in}{3.441662in}}%
\pgfusepath{clip}%
\pgfsetbuttcap%
\pgfsetroundjoin%
\definecolor{currentfill}{rgb}{0.826588,0.315714,0.441316}%
\pgfsetfillcolor{currentfill}%
\pgfsetlinewidth{0.000000pt}%
\definecolor{currentstroke}{rgb}{0.000000,0.000000,0.000000}%
\pgfsetstrokecolor{currentstroke}%
\pgfsetdash{}{0pt}%
\pgfpathmoveto{\pgfqpoint{4.255037in}{2.186141in}}%
\pgfpathlineto{\pgfqpoint{4.427120in}{2.186141in}}%
\pgfpathlineto{\pgfqpoint{4.427120in}{2.199585in}}%
\pgfpathlineto{\pgfqpoint{4.255037in}{2.199585in}}%
\pgfpathlineto{\pgfqpoint{4.255037in}{2.186141in}}%
\pgfusepath{fill}%
\end{pgfscope}%
\begin{pgfscope}%
\pgfpathrectangle{\pgfqpoint{4.255037in}{0.357758in}}{\pgfqpoint{0.172083in}{3.441662in}}%
\pgfusepath{clip}%
\pgfsetbuttcap%
\pgfsetroundjoin%
\definecolor{currentfill}{rgb}{0.830018,0.320172,0.437836}%
\pgfsetfillcolor{currentfill}%
\pgfsetlinewidth{0.000000pt}%
\definecolor{currentstroke}{rgb}{0.000000,0.000000,0.000000}%
\pgfsetstrokecolor{currentstroke}%
\pgfsetdash{}{0pt}%
\pgfpathmoveto{\pgfqpoint{4.255037in}{2.199585in}}%
\pgfpathlineto{\pgfqpoint{4.427120in}{2.199585in}}%
\pgfpathlineto{\pgfqpoint{4.427120in}{2.213029in}}%
\pgfpathlineto{\pgfqpoint{4.255037in}{2.213029in}}%
\pgfpathlineto{\pgfqpoint{4.255037in}{2.199585in}}%
\pgfusepath{fill}%
\end{pgfscope}%
\begin{pgfscope}%
\pgfpathrectangle{\pgfqpoint{4.255037in}{0.357758in}}{\pgfqpoint{0.172083in}{3.441662in}}%
\pgfusepath{clip}%
\pgfsetbuttcap%
\pgfsetroundjoin%
\definecolor{currentfill}{rgb}{0.833422,0.324635,0.434366}%
\pgfsetfillcolor{currentfill}%
\pgfsetlinewidth{0.000000pt}%
\definecolor{currentstroke}{rgb}{0.000000,0.000000,0.000000}%
\pgfsetstrokecolor{currentstroke}%
\pgfsetdash{}{0pt}%
\pgfpathmoveto{\pgfqpoint{4.255037in}{2.213029in}}%
\pgfpathlineto{\pgfqpoint{4.427120in}{2.213029in}}%
\pgfpathlineto{\pgfqpoint{4.427120in}{2.226473in}}%
\pgfpathlineto{\pgfqpoint{4.255037in}{2.226473in}}%
\pgfpathlineto{\pgfqpoint{4.255037in}{2.213029in}}%
\pgfusepath{fill}%
\end{pgfscope}%
\begin{pgfscope}%
\pgfpathrectangle{\pgfqpoint{4.255037in}{0.357758in}}{\pgfqpoint{0.172083in}{3.441662in}}%
\pgfusepath{clip}%
\pgfsetbuttcap%
\pgfsetroundjoin%
\definecolor{currentfill}{rgb}{0.836801,0.329105,0.430905}%
\pgfsetfillcolor{currentfill}%
\pgfsetlinewidth{0.000000pt}%
\definecolor{currentstroke}{rgb}{0.000000,0.000000,0.000000}%
\pgfsetstrokecolor{currentstroke}%
\pgfsetdash{}{0pt}%
\pgfpathmoveto{\pgfqpoint{4.255037in}{2.226473in}}%
\pgfpathlineto{\pgfqpoint{4.427120in}{2.226473in}}%
\pgfpathlineto{\pgfqpoint{4.427120in}{2.239917in}}%
\pgfpathlineto{\pgfqpoint{4.255037in}{2.239917in}}%
\pgfpathlineto{\pgfqpoint{4.255037in}{2.226473in}}%
\pgfusepath{fill}%
\end{pgfscope}%
\begin{pgfscope}%
\pgfpathrectangle{\pgfqpoint{4.255037in}{0.357758in}}{\pgfqpoint{0.172083in}{3.441662in}}%
\pgfusepath{clip}%
\pgfsetbuttcap%
\pgfsetroundjoin%
\definecolor{currentfill}{rgb}{0.840155,0.333580,0.427455}%
\pgfsetfillcolor{currentfill}%
\pgfsetlinewidth{0.000000pt}%
\definecolor{currentstroke}{rgb}{0.000000,0.000000,0.000000}%
\pgfsetstrokecolor{currentstroke}%
\pgfsetdash{}{0pt}%
\pgfpathmoveto{\pgfqpoint{4.255037in}{2.239917in}}%
\pgfpathlineto{\pgfqpoint{4.427120in}{2.239917in}}%
\pgfpathlineto{\pgfqpoint{4.427120in}{2.253361in}}%
\pgfpathlineto{\pgfqpoint{4.255037in}{2.253361in}}%
\pgfpathlineto{\pgfqpoint{4.255037in}{2.239917in}}%
\pgfusepath{fill}%
\end{pgfscope}%
\begin{pgfscope}%
\pgfpathrectangle{\pgfqpoint{4.255037in}{0.357758in}}{\pgfqpoint{0.172083in}{3.441662in}}%
\pgfusepath{clip}%
\pgfsetbuttcap%
\pgfsetroundjoin%
\definecolor{currentfill}{rgb}{0.843484,0.338062,0.424013}%
\pgfsetfillcolor{currentfill}%
\pgfsetlinewidth{0.000000pt}%
\definecolor{currentstroke}{rgb}{0.000000,0.000000,0.000000}%
\pgfsetstrokecolor{currentstroke}%
\pgfsetdash{}{0pt}%
\pgfpathmoveto{\pgfqpoint{4.255037in}{2.253361in}}%
\pgfpathlineto{\pgfqpoint{4.427120in}{2.253361in}}%
\pgfpathlineto{\pgfqpoint{4.427120in}{2.266805in}}%
\pgfpathlineto{\pgfqpoint{4.255037in}{2.266805in}}%
\pgfpathlineto{\pgfqpoint{4.255037in}{2.253361in}}%
\pgfusepath{fill}%
\end{pgfscope}%
\begin{pgfscope}%
\pgfpathrectangle{\pgfqpoint{4.255037in}{0.357758in}}{\pgfqpoint{0.172083in}{3.441662in}}%
\pgfusepath{clip}%
\pgfsetbuttcap%
\pgfsetroundjoin%
\definecolor{currentfill}{rgb}{0.846788,0.342551,0.420579}%
\pgfsetfillcolor{currentfill}%
\pgfsetlinewidth{0.000000pt}%
\definecolor{currentstroke}{rgb}{0.000000,0.000000,0.000000}%
\pgfsetstrokecolor{currentstroke}%
\pgfsetdash{}{0pt}%
\pgfpathmoveto{\pgfqpoint{4.255037in}{2.266805in}}%
\pgfpathlineto{\pgfqpoint{4.427120in}{2.266805in}}%
\pgfpathlineto{\pgfqpoint{4.427120in}{2.280249in}}%
\pgfpathlineto{\pgfqpoint{4.255037in}{2.280249in}}%
\pgfpathlineto{\pgfqpoint{4.255037in}{2.266805in}}%
\pgfusepath{fill}%
\end{pgfscope}%
\begin{pgfscope}%
\pgfpathrectangle{\pgfqpoint{4.255037in}{0.357758in}}{\pgfqpoint{0.172083in}{3.441662in}}%
\pgfusepath{clip}%
\pgfsetbuttcap%
\pgfsetroundjoin%
\definecolor{currentfill}{rgb}{0.850066,0.347048,0.417153}%
\pgfsetfillcolor{currentfill}%
\pgfsetlinewidth{0.000000pt}%
\definecolor{currentstroke}{rgb}{0.000000,0.000000,0.000000}%
\pgfsetstrokecolor{currentstroke}%
\pgfsetdash{}{0pt}%
\pgfpathmoveto{\pgfqpoint{4.255037in}{2.280249in}}%
\pgfpathlineto{\pgfqpoint{4.427120in}{2.280249in}}%
\pgfpathlineto{\pgfqpoint{4.427120in}{2.293693in}}%
\pgfpathlineto{\pgfqpoint{4.255037in}{2.293693in}}%
\pgfpathlineto{\pgfqpoint{4.255037in}{2.280249in}}%
\pgfusepath{fill}%
\end{pgfscope}%
\begin{pgfscope}%
\pgfpathrectangle{\pgfqpoint{4.255037in}{0.357758in}}{\pgfqpoint{0.172083in}{3.441662in}}%
\pgfusepath{clip}%
\pgfsetbuttcap%
\pgfsetroundjoin%
\definecolor{currentfill}{rgb}{0.853319,0.351553,0.413734}%
\pgfsetfillcolor{currentfill}%
\pgfsetlinewidth{0.000000pt}%
\definecolor{currentstroke}{rgb}{0.000000,0.000000,0.000000}%
\pgfsetstrokecolor{currentstroke}%
\pgfsetdash{}{0pt}%
\pgfpathmoveto{\pgfqpoint{4.255037in}{2.293693in}}%
\pgfpathlineto{\pgfqpoint{4.427120in}{2.293693in}}%
\pgfpathlineto{\pgfqpoint{4.427120in}{2.307137in}}%
\pgfpathlineto{\pgfqpoint{4.255037in}{2.307137in}}%
\pgfpathlineto{\pgfqpoint{4.255037in}{2.293693in}}%
\pgfusepath{fill}%
\end{pgfscope}%
\begin{pgfscope}%
\pgfpathrectangle{\pgfqpoint{4.255037in}{0.357758in}}{\pgfqpoint{0.172083in}{3.441662in}}%
\pgfusepath{clip}%
\pgfsetbuttcap%
\pgfsetroundjoin%
\definecolor{currentfill}{rgb}{0.856547,0.356066,0.410322}%
\pgfsetfillcolor{currentfill}%
\pgfsetlinewidth{0.000000pt}%
\definecolor{currentstroke}{rgb}{0.000000,0.000000,0.000000}%
\pgfsetstrokecolor{currentstroke}%
\pgfsetdash{}{0pt}%
\pgfpathmoveto{\pgfqpoint{4.255037in}{2.307137in}}%
\pgfpathlineto{\pgfqpoint{4.427120in}{2.307137in}}%
\pgfpathlineto{\pgfqpoint{4.427120in}{2.320581in}}%
\pgfpathlineto{\pgfqpoint{4.255037in}{2.320581in}}%
\pgfpathlineto{\pgfqpoint{4.255037in}{2.307137in}}%
\pgfusepath{fill}%
\end{pgfscope}%
\begin{pgfscope}%
\pgfpathrectangle{\pgfqpoint{4.255037in}{0.357758in}}{\pgfqpoint{0.172083in}{3.441662in}}%
\pgfusepath{clip}%
\pgfsetbuttcap%
\pgfsetroundjoin%
\definecolor{currentfill}{rgb}{0.859750,0.360588,0.406917}%
\pgfsetfillcolor{currentfill}%
\pgfsetlinewidth{0.000000pt}%
\definecolor{currentstroke}{rgb}{0.000000,0.000000,0.000000}%
\pgfsetstrokecolor{currentstroke}%
\pgfsetdash{}{0pt}%
\pgfpathmoveto{\pgfqpoint{4.255037in}{2.320581in}}%
\pgfpathlineto{\pgfqpoint{4.427120in}{2.320581in}}%
\pgfpathlineto{\pgfqpoint{4.427120in}{2.334025in}}%
\pgfpathlineto{\pgfqpoint{4.255037in}{2.334025in}}%
\pgfpathlineto{\pgfqpoint{4.255037in}{2.320581in}}%
\pgfusepath{fill}%
\end{pgfscope}%
\begin{pgfscope}%
\pgfpathrectangle{\pgfqpoint{4.255037in}{0.357758in}}{\pgfqpoint{0.172083in}{3.441662in}}%
\pgfusepath{clip}%
\pgfsetbuttcap%
\pgfsetroundjoin%
\definecolor{currentfill}{rgb}{0.862927,0.365119,0.403519}%
\pgfsetfillcolor{currentfill}%
\pgfsetlinewidth{0.000000pt}%
\definecolor{currentstroke}{rgb}{0.000000,0.000000,0.000000}%
\pgfsetstrokecolor{currentstroke}%
\pgfsetdash{}{0pt}%
\pgfpathmoveto{\pgfqpoint{4.255037in}{2.334025in}}%
\pgfpathlineto{\pgfqpoint{4.427120in}{2.334025in}}%
\pgfpathlineto{\pgfqpoint{4.427120in}{2.347469in}}%
\pgfpathlineto{\pgfqpoint{4.255037in}{2.347469in}}%
\pgfpathlineto{\pgfqpoint{4.255037in}{2.334025in}}%
\pgfusepath{fill}%
\end{pgfscope}%
\begin{pgfscope}%
\pgfpathrectangle{\pgfqpoint{4.255037in}{0.357758in}}{\pgfqpoint{0.172083in}{3.441662in}}%
\pgfusepath{clip}%
\pgfsetbuttcap%
\pgfsetroundjoin%
\definecolor{currentfill}{rgb}{0.866078,0.369660,0.400126}%
\pgfsetfillcolor{currentfill}%
\pgfsetlinewidth{0.000000pt}%
\definecolor{currentstroke}{rgb}{0.000000,0.000000,0.000000}%
\pgfsetstrokecolor{currentstroke}%
\pgfsetdash{}{0pt}%
\pgfpathmoveto{\pgfqpoint{4.255037in}{2.347469in}}%
\pgfpathlineto{\pgfqpoint{4.427120in}{2.347469in}}%
\pgfpathlineto{\pgfqpoint{4.427120in}{2.360913in}}%
\pgfpathlineto{\pgfqpoint{4.255037in}{2.360913in}}%
\pgfpathlineto{\pgfqpoint{4.255037in}{2.347469in}}%
\pgfusepath{fill}%
\end{pgfscope}%
\begin{pgfscope}%
\pgfpathrectangle{\pgfqpoint{4.255037in}{0.357758in}}{\pgfqpoint{0.172083in}{3.441662in}}%
\pgfusepath{clip}%
\pgfsetbuttcap%
\pgfsetroundjoin%
\definecolor{currentfill}{rgb}{0.869203,0.374212,0.396738}%
\pgfsetfillcolor{currentfill}%
\pgfsetlinewidth{0.000000pt}%
\definecolor{currentstroke}{rgb}{0.000000,0.000000,0.000000}%
\pgfsetstrokecolor{currentstroke}%
\pgfsetdash{}{0pt}%
\pgfpathmoveto{\pgfqpoint{4.255037in}{2.360913in}}%
\pgfpathlineto{\pgfqpoint{4.427120in}{2.360913in}}%
\pgfpathlineto{\pgfqpoint{4.427120in}{2.374357in}}%
\pgfpathlineto{\pgfqpoint{4.255037in}{2.374357in}}%
\pgfpathlineto{\pgfqpoint{4.255037in}{2.360913in}}%
\pgfusepath{fill}%
\end{pgfscope}%
\begin{pgfscope}%
\pgfpathrectangle{\pgfqpoint{4.255037in}{0.357758in}}{\pgfqpoint{0.172083in}{3.441662in}}%
\pgfusepath{clip}%
\pgfsetbuttcap%
\pgfsetroundjoin%
\definecolor{currentfill}{rgb}{0.872303,0.378774,0.393355}%
\pgfsetfillcolor{currentfill}%
\pgfsetlinewidth{0.000000pt}%
\definecolor{currentstroke}{rgb}{0.000000,0.000000,0.000000}%
\pgfsetstrokecolor{currentstroke}%
\pgfsetdash{}{0pt}%
\pgfpathmoveto{\pgfqpoint{4.255037in}{2.374357in}}%
\pgfpathlineto{\pgfqpoint{4.427120in}{2.374357in}}%
\pgfpathlineto{\pgfqpoint{4.427120in}{2.387801in}}%
\pgfpathlineto{\pgfqpoint{4.255037in}{2.387801in}}%
\pgfpathlineto{\pgfqpoint{4.255037in}{2.374357in}}%
\pgfusepath{fill}%
\end{pgfscope}%
\begin{pgfscope}%
\pgfpathrectangle{\pgfqpoint{4.255037in}{0.357758in}}{\pgfqpoint{0.172083in}{3.441662in}}%
\pgfusepath{clip}%
\pgfsetbuttcap%
\pgfsetroundjoin%
\definecolor{currentfill}{rgb}{0.875376,0.383347,0.389976}%
\pgfsetfillcolor{currentfill}%
\pgfsetlinewidth{0.000000pt}%
\definecolor{currentstroke}{rgb}{0.000000,0.000000,0.000000}%
\pgfsetstrokecolor{currentstroke}%
\pgfsetdash{}{0pt}%
\pgfpathmoveto{\pgfqpoint{4.255037in}{2.387801in}}%
\pgfpathlineto{\pgfqpoint{4.427120in}{2.387801in}}%
\pgfpathlineto{\pgfqpoint{4.427120in}{2.401245in}}%
\pgfpathlineto{\pgfqpoint{4.255037in}{2.401245in}}%
\pgfpathlineto{\pgfqpoint{4.255037in}{2.387801in}}%
\pgfusepath{fill}%
\end{pgfscope}%
\begin{pgfscope}%
\pgfpathrectangle{\pgfqpoint{4.255037in}{0.357758in}}{\pgfqpoint{0.172083in}{3.441662in}}%
\pgfusepath{clip}%
\pgfsetbuttcap%
\pgfsetroundjoin%
\definecolor{currentfill}{rgb}{0.878423,0.387932,0.386600}%
\pgfsetfillcolor{currentfill}%
\pgfsetlinewidth{0.000000pt}%
\definecolor{currentstroke}{rgb}{0.000000,0.000000,0.000000}%
\pgfsetstrokecolor{currentstroke}%
\pgfsetdash{}{0pt}%
\pgfpathmoveto{\pgfqpoint{4.255037in}{2.401245in}}%
\pgfpathlineto{\pgfqpoint{4.427120in}{2.401245in}}%
\pgfpathlineto{\pgfqpoint{4.427120in}{2.414689in}}%
\pgfpathlineto{\pgfqpoint{4.255037in}{2.414689in}}%
\pgfpathlineto{\pgfqpoint{4.255037in}{2.401245in}}%
\pgfusepath{fill}%
\end{pgfscope}%
\begin{pgfscope}%
\pgfpathrectangle{\pgfqpoint{4.255037in}{0.357758in}}{\pgfqpoint{0.172083in}{3.441662in}}%
\pgfusepath{clip}%
\pgfsetbuttcap%
\pgfsetroundjoin%
\definecolor{currentfill}{rgb}{0.881443,0.392529,0.383229}%
\pgfsetfillcolor{currentfill}%
\pgfsetlinewidth{0.000000pt}%
\definecolor{currentstroke}{rgb}{0.000000,0.000000,0.000000}%
\pgfsetstrokecolor{currentstroke}%
\pgfsetdash{}{0pt}%
\pgfpathmoveto{\pgfqpoint{4.255037in}{2.414689in}}%
\pgfpathlineto{\pgfqpoint{4.427120in}{2.414689in}}%
\pgfpathlineto{\pgfqpoint{4.427120in}{2.428133in}}%
\pgfpathlineto{\pgfqpoint{4.255037in}{2.428133in}}%
\pgfpathlineto{\pgfqpoint{4.255037in}{2.414689in}}%
\pgfusepath{fill}%
\end{pgfscope}%
\begin{pgfscope}%
\pgfpathrectangle{\pgfqpoint{4.255037in}{0.357758in}}{\pgfqpoint{0.172083in}{3.441662in}}%
\pgfusepath{clip}%
\pgfsetbuttcap%
\pgfsetroundjoin%
\definecolor{currentfill}{rgb}{0.884436,0.397139,0.379860}%
\pgfsetfillcolor{currentfill}%
\pgfsetlinewidth{0.000000pt}%
\definecolor{currentstroke}{rgb}{0.000000,0.000000,0.000000}%
\pgfsetstrokecolor{currentstroke}%
\pgfsetdash{}{0pt}%
\pgfpathmoveto{\pgfqpoint{4.255037in}{2.428133in}}%
\pgfpathlineto{\pgfqpoint{4.427120in}{2.428133in}}%
\pgfpathlineto{\pgfqpoint{4.427120in}{2.441577in}}%
\pgfpathlineto{\pgfqpoint{4.255037in}{2.441577in}}%
\pgfpathlineto{\pgfqpoint{4.255037in}{2.428133in}}%
\pgfusepath{fill}%
\end{pgfscope}%
\begin{pgfscope}%
\pgfpathrectangle{\pgfqpoint{4.255037in}{0.357758in}}{\pgfqpoint{0.172083in}{3.441662in}}%
\pgfusepath{clip}%
\pgfsetbuttcap%
\pgfsetroundjoin%
\definecolor{currentfill}{rgb}{0.887402,0.401762,0.376494}%
\pgfsetfillcolor{currentfill}%
\pgfsetlinewidth{0.000000pt}%
\definecolor{currentstroke}{rgb}{0.000000,0.000000,0.000000}%
\pgfsetstrokecolor{currentstroke}%
\pgfsetdash{}{0pt}%
\pgfpathmoveto{\pgfqpoint{4.255037in}{2.441577in}}%
\pgfpathlineto{\pgfqpoint{4.427120in}{2.441577in}}%
\pgfpathlineto{\pgfqpoint{4.427120in}{2.455021in}}%
\pgfpathlineto{\pgfqpoint{4.255037in}{2.455021in}}%
\pgfpathlineto{\pgfqpoint{4.255037in}{2.441577in}}%
\pgfusepath{fill}%
\end{pgfscope}%
\begin{pgfscope}%
\pgfpathrectangle{\pgfqpoint{4.255037in}{0.357758in}}{\pgfqpoint{0.172083in}{3.441662in}}%
\pgfusepath{clip}%
\pgfsetbuttcap%
\pgfsetroundjoin%
\definecolor{currentfill}{rgb}{0.890340,0.406398,0.373130}%
\pgfsetfillcolor{currentfill}%
\pgfsetlinewidth{0.000000pt}%
\definecolor{currentstroke}{rgb}{0.000000,0.000000,0.000000}%
\pgfsetstrokecolor{currentstroke}%
\pgfsetdash{}{0pt}%
\pgfpathmoveto{\pgfqpoint{4.255037in}{2.455021in}}%
\pgfpathlineto{\pgfqpoint{4.427120in}{2.455021in}}%
\pgfpathlineto{\pgfqpoint{4.427120in}{2.468465in}}%
\pgfpathlineto{\pgfqpoint{4.255037in}{2.468465in}}%
\pgfpathlineto{\pgfqpoint{4.255037in}{2.455021in}}%
\pgfusepath{fill}%
\end{pgfscope}%
\begin{pgfscope}%
\pgfpathrectangle{\pgfqpoint{4.255037in}{0.357758in}}{\pgfqpoint{0.172083in}{3.441662in}}%
\pgfusepath{clip}%
\pgfsetbuttcap%
\pgfsetroundjoin%
\definecolor{currentfill}{rgb}{0.893250,0.411048,0.369768}%
\pgfsetfillcolor{currentfill}%
\pgfsetlinewidth{0.000000pt}%
\definecolor{currentstroke}{rgb}{0.000000,0.000000,0.000000}%
\pgfsetstrokecolor{currentstroke}%
\pgfsetdash{}{0pt}%
\pgfpathmoveto{\pgfqpoint{4.255037in}{2.468465in}}%
\pgfpathlineto{\pgfqpoint{4.427120in}{2.468465in}}%
\pgfpathlineto{\pgfqpoint{4.427120in}{2.481909in}}%
\pgfpathlineto{\pgfqpoint{4.255037in}{2.481909in}}%
\pgfpathlineto{\pgfqpoint{4.255037in}{2.468465in}}%
\pgfusepath{fill}%
\end{pgfscope}%
\begin{pgfscope}%
\pgfpathrectangle{\pgfqpoint{4.255037in}{0.357758in}}{\pgfqpoint{0.172083in}{3.441662in}}%
\pgfusepath{clip}%
\pgfsetbuttcap%
\pgfsetroundjoin%
\definecolor{currentfill}{rgb}{0.896131,0.415712,0.366407}%
\pgfsetfillcolor{currentfill}%
\pgfsetlinewidth{0.000000pt}%
\definecolor{currentstroke}{rgb}{0.000000,0.000000,0.000000}%
\pgfsetstrokecolor{currentstroke}%
\pgfsetdash{}{0pt}%
\pgfpathmoveto{\pgfqpoint{4.255037in}{2.481909in}}%
\pgfpathlineto{\pgfqpoint{4.427120in}{2.481909in}}%
\pgfpathlineto{\pgfqpoint{4.427120in}{2.495353in}}%
\pgfpathlineto{\pgfqpoint{4.255037in}{2.495353in}}%
\pgfpathlineto{\pgfqpoint{4.255037in}{2.481909in}}%
\pgfusepath{fill}%
\end{pgfscope}%
\begin{pgfscope}%
\pgfpathrectangle{\pgfqpoint{4.255037in}{0.357758in}}{\pgfqpoint{0.172083in}{3.441662in}}%
\pgfusepath{clip}%
\pgfsetbuttcap%
\pgfsetroundjoin%
\definecolor{currentfill}{rgb}{0.898984,0.420392,0.363047}%
\pgfsetfillcolor{currentfill}%
\pgfsetlinewidth{0.000000pt}%
\definecolor{currentstroke}{rgb}{0.000000,0.000000,0.000000}%
\pgfsetstrokecolor{currentstroke}%
\pgfsetdash{}{0pt}%
\pgfpathmoveto{\pgfqpoint{4.255037in}{2.495353in}}%
\pgfpathlineto{\pgfqpoint{4.427120in}{2.495353in}}%
\pgfpathlineto{\pgfqpoint{4.427120in}{2.508797in}}%
\pgfpathlineto{\pgfqpoint{4.255037in}{2.508797in}}%
\pgfpathlineto{\pgfqpoint{4.255037in}{2.495353in}}%
\pgfusepath{fill}%
\end{pgfscope}%
\begin{pgfscope}%
\pgfpathrectangle{\pgfqpoint{4.255037in}{0.357758in}}{\pgfqpoint{0.172083in}{3.441662in}}%
\pgfusepath{clip}%
\pgfsetbuttcap%
\pgfsetroundjoin%
\definecolor{currentfill}{rgb}{0.901807,0.425087,0.359688}%
\pgfsetfillcolor{currentfill}%
\pgfsetlinewidth{0.000000pt}%
\definecolor{currentstroke}{rgb}{0.000000,0.000000,0.000000}%
\pgfsetstrokecolor{currentstroke}%
\pgfsetdash{}{0pt}%
\pgfpathmoveto{\pgfqpoint{4.255037in}{2.508797in}}%
\pgfpathlineto{\pgfqpoint{4.427120in}{2.508797in}}%
\pgfpathlineto{\pgfqpoint{4.427120in}{2.522241in}}%
\pgfpathlineto{\pgfqpoint{4.255037in}{2.522241in}}%
\pgfpathlineto{\pgfqpoint{4.255037in}{2.508797in}}%
\pgfusepath{fill}%
\end{pgfscope}%
\begin{pgfscope}%
\pgfpathrectangle{\pgfqpoint{4.255037in}{0.357758in}}{\pgfqpoint{0.172083in}{3.441662in}}%
\pgfusepath{clip}%
\pgfsetbuttcap%
\pgfsetroundjoin%
\definecolor{currentfill}{rgb}{0.904601,0.429797,0.356329}%
\pgfsetfillcolor{currentfill}%
\pgfsetlinewidth{0.000000pt}%
\definecolor{currentstroke}{rgb}{0.000000,0.000000,0.000000}%
\pgfsetstrokecolor{currentstroke}%
\pgfsetdash{}{0pt}%
\pgfpathmoveto{\pgfqpoint{4.255037in}{2.522241in}}%
\pgfpathlineto{\pgfqpoint{4.427120in}{2.522241in}}%
\pgfpathlineto{\pgfqpoint{4.427120in}{2.535685in}}%
\pgfpathlineto{\pgfqpoint{4.255037in}{2.535685in}}%
\pgfpathlineto{\pgfqpoint{4.255037in}{2.522241in}}%
\pgfusepath{fill}%
\end{pgfscope}%
\begin{pgfscope}%
\pgfpathrectangle{\pgfqpoint{4.255037in}{0.357758in}}{\pgfqpoint{0.172083in}{3.441662in}}%
\pgfusepath{clip}%
\pgfsetbuttcap%
\pgfsetroundjoin%
\definecolor{currentfill}{rgb}{0.907365,0.434524,0.352970}%
\pgfsetfillcolor{currentfill}%
\pgfsetlinewidth{0.000000pt}%
\definecolor{currentstroke}{rgb}{0.000000,0.000000,0.000000}%
\pgfsetstrokecolor{currentstroke}%
\pgfsetdash{}{0pt}%
\pgfpathmoveto{\pgfqpoint{4.255037in}{2.535685in}}%
\pgfpathlineto{\pgfqpoint{4.427120in}{2.535685in}}%
\pgfpathlineto{\pgfqpoint{4.427120in}{2.549129in}}%
\pgfpathlineto{\pgfqpoint{4.255037in}{2.549129in}}%
\pgfpathlineto{\pgfqpoint{4.255037in}{2.535685in}}%
\pgfusepath{fill}%
\end{pgfscope}%
\begin{pgfscope}%
\pgfpathrectangle{\pgfqpoint{4.255037in}{0.357758in}}{\pgfqpoint{0.172083in}{3.441662in}}%
\pgfusepath{clip}%
\pgfsetbuttcap%
\pgfsetroundjoin%
\definecolor{currentfill}{rgb}{0.910098,0.439268,0.349610}%
\pgfsetfillcolor{currentfill}%
\pgfsetlinewidth{0.000000pt}%
\definecolor{currentstroke}{rgb}{0.000000,0.000000,0.000000}%
\pgfsetstrokecolor{currentstroke}%
\pgfsetdash{}{0pt}%
\pgfpathmoveto{\pgfqpoint{4.255037in}{2.549129in}}%
\pgfpathlineto{\pgfqpoint{4.427120in}{2.549129in}}%
\pgfpathlineto{\pgfqpoint{4.427120in}{2.562573in}}%
\pgfpathlineto{\pgfqpoint{4.255037in}{2.562573in}}%
\pgfpathlineto{\pgfqpoint{4.255037in}{2.549129in}}%
\pgfusepath{fill}%
\end{pgfscope}%
\begin{pgfscope}%
\pgfpathrectangle{\pgfqpoint{4.255037in}{0.357758in}}{\pgfqpoint{0.172083in}{3.441662in}}%
\pgfusepath{clip}%
\pgfsetbuttcap%
\pgfsetroundjoin%
\definecolor{currentfill}{rgb}{0.912800,0.444029,0.346251}%
\pgfsetfillcolor{currentfill}%
\pgfsetlinewidth{0.000000pt}%
\definecolor{currentstroke}{rgb}{0.000000,0.000000,0.000000}%
\pgfsetstrokecolor{currentstroke}%
\pgfsetdash{}{0pt}%
\pgfpathmoveto{\pgfqpoint{4.255037in}{2.562573in}}%
\pgfpathlineto{\pgfqpoint{4.427120in}{2.562573in}}%
\pgfpathlineto{\pgfqpoint{4.427120in}{2.576017in}}%
\pgfpathlineto{\pgfqpoint{4.255037in}{2.576017in}}%
\pgfpathlineto{\pgfqpoint{4.255037in}{2.562573in}}%
\pgfusepath{fill}%
\end{pgfscope}%
\begin{pgfscope}%
\pgfpathrectangle{\pgfqpoint{4.255037in}{0.357758in}}{\pgfqpoint{0.172083in}{3.441662in}}%
\pgfusepath{clip}%
\pgfsetbuttcap%
\pgfsetroundjoin%
\definecolor{currentfill}{rgb}{0.915471,0.448807,0.342890}%
\pgfsetfillcolor{currentfill}%
\pgfsetlinewidth{0.000000pt}%
\definecolor{currentstroke}{rgb}{0.000000,0.000000,0.000000}%
\pgfsetstrokecolor{currentstroke}%
\pgfsetdash{}{0pt}%
\pgfpathmoveto{\pgfqpoint{4.255037in}{2.576017in}}%
\pgfpathlineto{\pgfqpoint{4.427120in}{2.576017in}}%
\pgfpathlineto{\pgfqpoint{4.427120in}{2.589461in}}%
\pgfpathlineto{\pgfqpoint{4.255037in}{2.589461in}}%
\pgfpathlineto{\pgfqpoint{4.255037in}{2.576017in}}%
\pgfusepath{fill}%
\end{pgfscope}%
\begin{pgfscope}%
\pgfpathrectangle{\pgfqpoint{4.255037in}{0.357758in}}{\pgfqpoint{0.172083in}{3.441662in}}%
\pgfusepath{clip}%
\pgfsetbuttcap%
\pgfsetroundjoin%
\definecolor{currentfill}{rgb}{0.918109,0.453603,0.339529}%
\pgfsetfillcolor{currentfill}%
\pgfsetlinewidth{0.000000pt}%
\definecolor{currentstroke}{rgb}{0.000000,0.000000,0.000000}%
\pgfsetstrokecolor{currentstroke}%
\pgfsetdash{}{0pt}%
\pgfpathmoveto{\pgfqpoint{4.255037in}{2.589461in}}%
\pgfpathlineto{\pgfqpoint{4.427120in}{2.589461in}}%
\pgfpathlineto{\pgfqpoint{4.427120in}{2.602905in}}%
\pgfpathlineto{\pgfqpoint{4.255037in}{2.602905in}}%
\pgfpathlineto{\pgfqpoint{4.255037in}{2.589461in}}%
\pgfusepath{fill}%
\end{pgfscope}%
\begin{pgfscope}%
\pgfpathrectangle{\pgfqpoint{4.255037in}{0.357758in}}{\pgfqpoint{0.172083in}{3.441662in}}%
\pgfusepath{clip}%
\pgfsetbuttcap%
\pgfsetroundjoin%
\definecolor{currentfill}{rgb}{0.920714,0.458417,0.336166}%
\pgfsetfillcolor{currentfill}%
\pgfsetlinewidth{0.000000pt}%
\definecolor{currentstroke}{rgb}{0.000000,0.000000,0.000000}%
\pgfsetstrokecolor{currentstroke}%
\pgfsetdash{}{0pt}%
\pgfpathmoveto{\pgfqpoint{4.255037in}{2.602905in}}%
\pgfpathlineto{\pgfqpoint{4.427120in}{2.602905in}}%
\pgfpathlineto{\pgfqpoint{4.427120in}{2.616349in}}%
\pgfpathlineto{\pgfqpoint{4.255037in}{2.616349in}}%
\pgfpathlineto{\pgfqpoint{4.255037in}{2.602905in}}%
\pgfusepath{fill}%
\end{pgfscope}%
\begin{pgfscope}%
\pgfpathrectangle{\pgfqpoint{4.255037in}{0.357758in}}{\pgfqpoint{0.172083in}{3.441662in}}%
\pgfusepath{clip}%
\pgfsetbuttcap%
\pgfsetroundjoin%
\definecolor{currentfill}{rgb}{0.923287,0.463251,0.332801}%
\pgfsetfillcolor{currentfill}%
\pgfsetlinewidth{0.000000pt}%
\definecolor{currentstroke}{rgb}{0.000000,0.000000,0.000000}%
\pgfsetstrokecolor{currentstroke}%
\pgfsetdash{}{0pt}%
\pgfpathmoveto{\pgfqpoint{4.255037in}{2.616349in}}%
\pgfpathlineto{\pgfqpoint{4.427120in}{2.616349in}}%
\pgfpathlineto{\pgfqpoint{4.427120in}{2.629793in}}%
\pgfpathlineto{\pgfqpoint{4.255037in}{2.629793in}}%
\pgfpathlineto{\pgfqpoint{4.255037in}{2.616349in}}%
\pgfusepath{fill}%
\end{pgfscope}%
\begin{pgfscope}%
\pgfpathrectangle{\pgfqpoint{4.255037in}{0.357758in}}{\pgfqpoint{0.172083in}{3.441662in}}%
\pgfusepath{clip}%
\pgfsetbuttcap%
\pgfsetroundjoin%
\definecolor{currentfill}{rgb}{0.925825,0.468103,0.329435}%
\pgfsetfillcolor{currentfill}%
\pgfsetlinewidth{0.000000pt}%
\definecolor{currentstroke}{rgb}{0.000000,0.000000,0.000000}%
\pgfsetstrokecolor{currentstroke}%
\pgfsetdash{}{0pt}%
\pgfpathmoveto{\pgfqpoint{4.255037in}{2.629793in}}%
\pgfpathlineto{\pgfqpoint{4.427120in}{2.629793in}}%
\pgfpathlineto{\pgfqpoint{4.427120in}{2.643237in}}%
\pgfpathlineto{\pgfqpoint{4.255037in}{2.643237in}}%
\pgfpathlineto{\pgfqpoint{4.255037in}{2.629793in}}%
\pgfusepath{fill}%
\end{pgfscope}%
\begin{pgfscope}%
\pgfpathrectangle{\pgfqpoint{4.255037in}{0.357758in}}{\pgfqpoint{0.172083in}{3.441662in}}%
\pgfusepath{clip}%
\pgfsetbuttcap%
\pgfsetroundjoin%
\definecolor{currentfill}{rgb}{0.928329,0.472975,0.326067}%
\pgfsetfillcolor{currentfill}%
\pgfsetlinewidth{0.000000pt}%
\definecolor{currentstroke}{rgb}{0.000000,0.000000,0.000000}%
\pgfsetstrokecolor{currentstroke}%
\pgfsetdash{}{0pt}%
\pgfpathmoveto{\pgfqpoint{4.255037in}{2.643237in}}%
\pgfpathlineto{\pgfqpoint{4.427120in}{2.643237in}}%
\pgfpathlineto{\pgfqpoint{4.427120in}{2.656681in}}%
\pgfpathlineto{\pgfqpoint{4.255037in}{2.656681in}}%
\pgfpathlineto{\pgfqpoint{4.255037in}{2.643237in}}%
\pgfusepath{fill}%
\end{pgfscope}%
\begin{pgfscope}%
\pgfpathrectangle{\pgfqpoint{4.255037in}{0.357758in}}{\pgfqpoint{0.172083in}{3.441662in}}%
\pgfusepath{clip}%
\pgfsetbuttcap%
\pgfsetroundjoin%
\definecolor{currentfill}{rgb}{0.930798,0.477867,0.322697}%
\pgfsetfillcolor{currentfill}%
\pgfsetlinewidth{0.000000pt}%
\definecolor{currentstroke}{rgb}{0.000000,0.000000,0.000000}%
\pgfsetstrokecolor{currentstroke}%
\pgfsetdash{}{0pt}%
\pgfpathmoveto{\pgfqpoint{4.255037in}{2.656681in}}%
\pgfpathlineto{\pgfqpoint{4.427120in}{2.656681in}}%
\pgfpathlineto{\pgfqpoint{4.427120in}{2.670125in}}%
\pgfpathlineto{\pgfqpoint{4.255037in}{2.670125in}}%
\pgfpathlineto{\pgfqpoint{4.255037in}{2.656681in}}%
\pgfusepath{fill}%
\end{pgfscope}%
\begin{pgfscope}%
\pgfpathrectangle{\pgfqpoint{4.255037in}{0.357758in}}{\pgfqpoint{0.172083in}{3.441662in}}%
\pgfusepath{clip}%
\pgfsetbuttcap%
\pgfsetroundjoin%
\definecolor{currentfill}{rgb}{0.933232,0.482780,0.319325}%
\pgfsetfillcolor{currentfill}%
\pgfsetlinewidth{0.000000pt}%
\definecolor{currentstroke}{rgb}{0.000000,0.000000,0.000000}%
\pgfsetstrokecolor{currentstroke}%
\pgfsetdash{}{0pt}%
\pgfpathmoveto{\pgfqpoint{4.255037in}{2.670125in}}%
\pgfpathlineto{\pgfqpoint{4.427120in}{2.670125in}}%
\pgfpathlineto{\pgfqpoint{4.427120in}{2.683569in}}%
\pgfpathlineto{\pgfqpoint{4.255037in}{2.683569in}}%
\pgfpathlineto{\pgfqpoint{4.255037in}{2.670125in}}%
\pgfusepath{fill}%
\end{pgfscope}%
\begin{pgfscope}%
\pgfpathrectangle{\pgfqpoint{4.255037in}{0.357758in}}{\pgfqpoint{0.172083in}{3.441662in}}%
\pgfusepath{clip}%
\pgfsetbuttcap%
\pgfsetroundjoin%
\definecolor{currentfill}{rgb}{0.935630,0.487712,0.315952}%
\pgfsetfillcolor{currentfill}%
\pgfsetlinewidth{0.000000pt}%
\definecolor{currentstroke}{rgb}{0.000000,0.000000,0.000000}%
\pgfsetstrokecolor{currentstroke}%
\pgfsetdash{}{0pt}%
\pgfpathmoveto{\pgfqpoint{4.255037in}{2.683569in}}%
\pgfpathlineto{\pgfqpoint{4.427120in}{2.683569in}}%
\pgfpathlineto{\pgfqpoint{4.427120in}{2.697013in}}%
\pgfpathlineto{\pgfqpoint{4.255037in}{2.697013in}}%
\pgfpathlineto{\pgfqpoint{4.255037in}{2.683569in}}%
\pgfusepath{fill}%
\end{pgfscope}%
\begin{pgfscope}%
\pgfpathrectangle{\pgfqpoint{4.255037in}{0.357758in}}{\pgfqpoint{0.172083in}{3.441662in}}%
\pgfusepath{clip}%
\pgfsetbuttcap%
\pgfsetroundjoin%
\definecolor{currentfill}{rgb}{0.937990,0.492667,0.312575}%
\pgfsetfillcolor{currentfill}%
\pgfsetlinewidth{0.000000pt}%
\definecolor{currentstroke}{rgb}{0.000000,0.000000,0.000000}%
\pgfsetstrokecolor{currentstroke}%
\pgfsetdash{}{0pt}%
\pgfpathmoveto{\pgfqpoint{4.255037in}{2.697013in}}%
\pgfpathlineto{\pgfqpoint{4.427120in}{2.697013in}}%
\pgfpathlineto{\pgfqpoint{4.427120in}{2.710457in}}%
\pgfpathlineto{\pgfqpoint{4.255037in}{2.710457in}}%
\pgfpathlineto{\pgfqpoint{4.255037in}{2.697013in}}%
\pgfusepath{fill}%
\end{pgfscope}%
\begin{pgfscope}%
\pgfpathrectangle{\pgfqpoint{4.255037in}{0.357758in}}{\pgfqpoint{0.172083in}{3.441662in}}%
\pgfusepath{clip}%
\pgfsetbuttcap%
\pgfsetroundjoin%
\definecolor{currentfill}{rgb}{0.940313,0.497642,0.309197}%
\pgfsetfillcolor{currentfill}%
\pgfsetlinewidth{0.000000pt}%
\definecolor{currentstroke}{rgb}{0.000000,0.000000,0.000000}%
\pgfsetstrokecolor{currentstroke}%
\pgfsetdash{}{0pt}%
\pgfpathmoveto{\pgfqpoint{4.255037in}{2.710457in}}%
\pgfpathlineto{\pgfqpoint{4.427120in}{2.710457in}}%
\pgfpathlineto{\pgfqpoint{4.427120in}{2.723901in}}%
\pgfpathlineto{\pgfqpoint{4.255037in}{2.723901in}}%
\pgfpathlineto{\pgfqpoint{4.255037in}{2.710457in}}%
\pgfusepath{fill}%
\end{pgfscope}%
\begin{pgfscope}%
\pgfpathrectangle{\pgfqpoint{4.255037in}{0.357758in}}{\pgfqpoint{0.172083in}{3.441662in}}%
\pgfusepath{clip}%
\pgfsetbuttcap%
\pgfsetroundjoin%
\definecolor{currentfill}{rgb}{0.942598,0.502639,0.305816}%
\pgfsetfillcolor{currentfill}%
\pgfsetlinewidth{0.000000pt}%
\definecolor{currentstroke}{rgb}{0.000000,0.000000,0.000000}%
\pgfsetstrokecolor{currentstroke}%
\pgfsetdash{}{0pt}%
\pgfpathmoveto{\pgfqpoint{4.255037in}{2.723901in}}%
\pgfpathlineto{\pgfqpoint{4.427120in}{2.723901in}}%
\pgfpathlineto{\pgfqpoint{4.427120in}{2.737345in}}%
\pgfpathlineto{\pgfqpoint{4.255037in}{2.737345in}}%
\pgfpathlineto{\pgfqpoint{4.255037in}{2.723901in}}%
\pgfusepath{fill}%
\end{pgfscope}%
\begin{pgfscope}%
\pgfpathrectangle{\pgfqpoint{4.255037in}{0.357758in}}{\pgfqpoint{0.172083in}{3.441662in}}%
\pgfusepath{clip}%
\pgfsetbuttcap%
\pgfsetroundjoin%
\definecolor{currentfill}{rgb}{0.944844,0.507658,0.302433}%
\pgfsetfillcolor{currentfill}%
\pgfsetlinewidth{0.000000pt}%
\definecolor{currentstroke}{rgb}{0.000000,0.000000,0.000000}%
\pgfsetstrokecolor{currentstroke}%
\pgfsetdash{}{0pt}%
\pgfpathmoveto{\pgfqpoint{4.255037in}{2.737345in}}%
\pgfpathlineto{\pgfqpoint{4.427120in}{2.737345in}}%
\pgfpathlineto{\pgfqpoint{4.427120in}{2.750789in}}%
\pgfpathlineto{\pgfqpoint{4.255037in}{2.750789in}}%
\pgfpathlineto{\pgfqpoint{4.255037in}{2.737345in}}%
\pgfusepath{fill}%
\end{pgfscope}%
\begin{pgfscope}%
\pgfpathrectangle{\pgfqpoint{4.255037in}{0.357758in}}{\pgfqpoint{0.172083in}{3.441662in}}%
\pgfusepath{clip}%
\pgfsetbuttcap%
\pgfsetroundjoin%
\definecolor{currentfill}{rgb}{0.947051,0.512699,0.299049}%
\pgfsetfillcolor{currentfill}%
\pgfsetlinewidth{0.000000pt}%
\definecolor{currentstroke}{rgb}{0.000000,0.000000,0.000000}%
\pgfsetstrokecolor{currentstroke}%
\pgfsetdash{}{0pt}%
\pgfpathmoveto{\pgfqpoint{4.255037in}{2.750789in}}%
\pgfpathlineto{\pgfqpoint{4.427120in}{2.750789in}}%
\pgfpathlineto{\pgfqpoint{4.427120in}{2.764233in}}%
\pgfpathlineto{\pgfqpoint{4.255037in}{2.764233in}}%
\pgfpathlineto{\pgfqpoint{4.255037in}{2.750789in}}%
\pgfusepath{fill}%
\end{pgfscope}%
\begin{pgfscope}%
\pgfpathrectangle{\pgfqpoint{4.255037in}{0.357758in}}{\pgfqpoint{0.172083in}{3.441662in}}%
\pgfusepath{clip}%
\pgfsetbuttcap%
\pgfsetroundjoin%
\definecolor{currentfill}{rgb}{0.949217,0.517763,0.295662}%
\pgfsetfillcolor{currentfill}%
\pgfsetlinewidth{0.000000pt}%
\definecolor{currentstroke}{rgb}{0.000000,0.000000,0.000000}%
\pgfsetstrokecolor{currentstroke}%
\pgfsetdash{}{0pt}%
\pgfpathmoveto{\pgfqpoint{4.255037in}{2.764233in}}%
\pgfpathlineto{\pgfqpoint{4.427120in}{2.764233in}}%
\pgfpathlineto{\pgfqpoint{4.427120in}{2.777677in}}%
\pgfpathlineto{\pgfqpoint{4.255037in}{2.777677in}}%
\pgfpathlineto{\pgfqpoint{4.255037in}{2.764233in}}%
\pgfusepath{fill}%
\end{pgfscope}%
\begin{pgfscope}%
\pgfpathrectangle{\pgfqpoint{4.255037in}{0.357758in}}{\pgfqpoint{0.172083in}{3.441662in}}%
\pgfusepath{clip}%
\pgfsetbuttcap%
\pgfsetroundjoin%
\definecolor{currentfill}{rgb}{0.951344,0.522850,0.292275}%
\pgfsetfillcolor{currentfill}%
\pgfsetlinewidth{0.000000pt}%
\definecolor{currentstroke}{rgb}{0.000000,0.000000,0.000000}%
\pgfsetstrokecolor{currentstroke}%
\pgfsetdash{}{0pt}%
\pgfpathmoveto{\pgfqpoint{4.255037in}{2.777677in}}%
\pgfpathlineto{\pgfqpoint{4.427120in}{2.777677in}}%
\pgfpathlineto{\pgfqpoint{4.427120in}{2.791121in}}%
\pgfpathlineto{\pgfqpoint{4.255037in}{2.791121in}}%
\pgfpathlineto{\pgfqpoint{4.255037in}{2.777677in}}%
\pgfusepath{fill}%
\end{pgfscope}%
\begin{pgfscope}%
\pgfpathrectangle{\pgfqpoint{4.255037in}{0.357758in}}{\pgfqpoint{0.172083in}{3.441662in}}%
\pgfusepath{clip}%
\pgfsetbuttcap%
\pgfsetroundjoin%
\definecolor{currentfill}{rgb}{0.953428,0.527960,0.288883}%
\pgfsetfillcolor{currentfill}%
\pgfsetlinewidth{0.000000pt}%
\definecolor{currentstroke}{rgb}{0.000000,0.000000,0.000000}%
\pgfsetstrokecolor{currentstroke}%
\pgfsetdash{}{0pt}%
\pgfpathmoveto{\pgfqpoint{4.255037in}{2.791121in}}%
\pgfpathlineto{\pgfqpoint{4.427120in}{2.791121in}}%
\pgfpathlineto{\pgfqpoint{4.427120in}{2.804565in}}%
\pgfpathlineto{\pgfqpoint{4.255037in}{2.804565in}}%
\pgfpathlineto{\pgfqpoint{4.255037in}{2.791121in}}%
\pgfusepath{fill}%
\end{pgfscope}%
\begin{pgfscope}%
\pgfpathrectangle{\pgfqpoint{4.255037in}{0.357758in}}{\pgfqpoint{0.172083in}{3.441662in}}%
\pgfusepath{clip}%
\pgfsetbuttcap%
\pgfsetroundjoin%
\definecolor{currentfill}{rgb}{0.955470,0.533093,0.285490}%
\pgfsetfillcolor{currentfill}%
\pgfsetlinewidth{0.000000pt}%
\definecolor{currentstroke}{rgb}{0.000000,0.000000,0.000000}%
\pgfsetstrokecolor{currentstroke}%
\pgfsetdash{}{0pt}%
\pgfpathmoveto{\pgfqpoint{4.255037in}{2.804565in}}%
\pgfpathlineto{\pgfqpoint{4.427120in}{2.804565in}}%
\pgfpathlineto{\pgfqpoint{4.427120in}{2.818009in}}%
\pgfpathlineto{\pgfqpoint{4.255037in}{2.818009in}}%
\pgfpathlineto{\pgfqpoint{4.255037in}{2.804565in}}%
\pgfusepath{fill}%
\end{pgfscope}%
\begin{pgfscope}%
\pgfpathrectangle{\pgfqpoint{4.255037in}{0.357758in}}{\pgfqpoint{0.172083in}{3.441662in}}%
\pgfusepath{clip}%
\pgfsetbuttcap%
\pgfsetroundjoin%
\definecolor{currentfill}{rgb}{0.957469,0.538250,0.282096}%
\pgfsetfillcolor{currentfill}%
\pgfsetlinewidth{0.000000pt}%
\definecolor{currentstroke}{rgb}{0.000000,0.000000,0.000000}%
\pgfsetstrokecolor{currentstroke}%
\pgfsetdash{}{0pt}%
\pgfpathmoveto{\pgfqpoint{4.255037in}{2.818009in}}%
\pgfpathlineto{\pgfqpoint{4.427120in}{2.818009in}}%
\pgfpathlineto{\pgfqpoint{4.427120in}{2.831453in}}%
\pgfpathlineto{\pgfqpoint{4.255037in}{2.831453in}}%
\pgfpathlineto{\pgfqpoint{4.255037in}{2.818009in}}%
\pgfusepath{fill}%
\end{pgfscope}%
\begin{pgfscope}%
\pgfpathrectangle{\pgfqpoint{4.255037in}{0.357758in}}{\pgfqpoint{0.172083in}{3.441662in}}%
\pgfusepath{clip}%
\pgfsetbuttcap%
\pgfsetroundjoin%
\definecolor{currentfill}{rgb}{0.959424,0.543431,0.278701}%
\pgfsetfillcolor{currentfill}%
\pgfsetlinewidth{0.000000pt}%
\definecolor{currentstroke}{rgb}{0.000000,0.000000,0.000000}%
\pgfsetstrokecolor{currentstroke}%
\pgfsetdash{}{0pt}%
\pgfpathmoveto{\pgfqpoint{4.255037in}{2.831453in}}%
\pgfpathlineto{\pgfqpoint{4.427120in}{2.831453in}}%
\pgfpathlineto{\pgfqpoint{4.427120in}{2.844897in}}%
\pgfpathlineto{\pgfqpoint{4.255037in}{2.844897in}}%
\pgfpathlineto{\pgfqpoint{4.255037in}{2.831453in}}%
\pgfusepath{fill}%
\end{pgfscope}%
\begin{pgfscope}%
\pgfpathrectangle{\pgfqpoint{4.255037in}{0.357758in}}{\pgfqpoint{0.172083in}{3.441662in}}%
\pgfusepath{clip}%
\pgfsetbuttcap%
\pgfsetroundjoin%
\definecolor{currentfill}{rgb}{0.961336,0.548636,0.275305}%
\pgfsetfillcolor{currentfill}%
\pgfsetlinewidth{0.000000pt}%
\definecolor{currentstroke}{rgb}{0.000000,0.000000,0.000000}%
\pgfsetstrokecolor{currentstroke}%
\pgfsetdash{}{0pt}%
\pgfpathmoveto{\pgfqpoint{4.255037in}{2.844897in}}%
\pgfpathlineto{\pgfqpoint{4.427120in}{2.844897in}}%
\pgfpathlineto{\pgfqpoint{4.427120in}{2.858341in}}%
\pgfpathlineto{\pgfqpoint{4.255037in}{2.858341in}}%
\pgfpathlineto{\pgfqpoint{4.255037in}{2.844897in}}%
\pgfusepath{fill}%
\end{pgfscope}%
\begin{pgfscope}%
\pgfpathrectangle{\pgfqpoint{4.255037in}{0.357758in}}{\pgfqpoint{0.172083in}{3.441662in}}%
\pgfusepath{clip}%
\pgfsetbuttcap%
\pgfsetroundjoin%
\definecolor{currentfill}{rgb}{0.963203,0.553865,0.271909}%
\pgfsetfillcolor{currentfill}%
\pgfsetlinewidth{0.000000pt}%
\definecolor{currentstroke}{rgb}{0.000000,0.000000,0.000000}%
\pgfsetstrokecolor{currentstroke}%
\pgfsetdash{}{0pt}%
\pgfpathmoveto{\pgfqpoint{4.255037in}{2.858341in}}%
\pgfpathlineto{\pgfqpoint{4.427120in}{2.858341in}}%
\pgfpathlineto{\pgfqpoint{4.427120in}{2.871785in}}%
\pgfpathlineto{\pgfqpoint{4.255037in}{2.871785in}}%
\pgfpathlineto{\pgfqpoint{4.255037in}{2.858341in}}%
\pgfusepath{fill}%
\end{pgfscope}%
\begin{pgfscope}%
\pgfpathrectangle{\pgfqpoint{4.255037in}{0.357758in}}{\pgfqpoint{0.172083in}{3.441662in}}%
\pgfusepath{clip}%
\pgfsetbuttcap%
\pgfsetroundjoin%
\definecolor{currentfill}{rgb}{0.965024,0.559118,0.268513}%
\pgfsetfillcolor{currentfill}%
\pgfsetlinewidth{0.000000pt}%
\definecolor{currentstroke}{rgb}{0.000000,0.000000,0.000000}%
\pgfsetstrokecolor{currentstroke}%
\pgfsetdash{}{0pt}%
\pgfpathmoveto{\pgfqpoint{4.255037in}{2.871785in}}%
\pgfpathlineto{\pgfqpoint{4.427120in}{2.871785in}}%
\pgfpathlineto{\pgfqpoint{4.427120in}{2.885229in}}%
\pgfpathlineto{\pgfqpoint{4.255037in}{2.885229in}}%
\pgfpathlineto{\pgfqpoint{4.255037in}{2.871785in}}%
\pgfusepath{fill}%
\end{pgfscope}%
\begin{pgfscope}%
\pgfpathrectangle{\pgfqpoint{4.255037in}{0.357758in}}{\pgfqpoint{0.172083in}{3.441662in}}%
\pgfusepath{clip}%
\pgfsetbuttcap%
\pgfsetroundjoin%
\definecolor{currentfill}{rgb}{0.966798,0.564396,0.265118}%
\pgfsetfillcolor{currentfill}%
\pgfsetlinewidth{0.000000pt}%
\definecolor{currentstroke}{rgb}{0.000000,0.000000,0.000000}%
\pgfsetstrokecolor{currentstroke}%
\pgfsetdash{}{0pt}%
\pgfpathmoveto{\pgfqpoint{4.255037in}{2.885229in}}%
\pgfpathlineto{\pgfqpoint{4.427120in}{2.885229in}}%
\pgfpathlineto{\pgfqpoint{4.427120in}{2.898673in}}%
\pgfpathlineto{\pgfqpoint{4.255037in}{2.898673in}}%
\pgfpathlineto{\pgfqpoint{4.255037in}{2.885229in}}%
\pgfusepath{fill}%
\end{pgfscope}%
\begin{pgfscope}%
\pgfpathrectangle{\pgfqpoint{4.255037in}{0.357758in}}{\pgfqpoint{0.172083in}{3.441662in}}%
\pgfusepath{clip}%
\pgfsetbuttcap%
\pgfsetroundjoin%
\definecolor{currentfill}{rgb}{0.968526,0.569700,0.261721}%
\pgfsetfillcolor{currentfill}%
\pgfsetlinewidth{0.000000pt}%
\definecolor{currentstroke}{rgb}{0.000000,0.000000,0.000000}%
\pgfsetstrokecolor{currentstroke}%
\pgfsetdash{}{0pt}%
\pgfpathmoveto{\pgfqpoint{4.255037in}{2.898673in}}%
\pgfpathlineto{\pgfqpoint{4.427120in}{2.898673in}}%
\pgfpathlineto{\pgfqpoint{4.427120in}{2.912117in}}%
\pgfpathlineto{\pgfqpoint{4.255037in}{2.912117in}}%
\pgfpathlineto{\pgfqpoint{4.255037in}{2.898673in}}%
\pgfusepath{fill}%
\end{pgfscope}%
\begin{pgfscope}%
\pgfpathrectangle{\pgfqpoint{4.255037in}{0.357758in}}{\pgfqpoint{0.172083in}{3.441662in}}%
\pgfusepath{clip}%
\pgfsetbuttcap%
\pgfsetroundjoin%
\definecolor{currentfill}{rgb}{0.970205,0.575028,0.258325}%
\pgfsetfillcolor{currentfill}%
\pgfsetlinewidth{0.000000pt}%
\definecolor{currentstroke}{rgb}{0.000000,0.000000,0.000000}%
\pgfsetstrokecolor{currentstroke}%
\pgfsetdash{}{0pt}%
\pgfpathmoveto{\pgfqpoint{4.255037in}{2.912117in}}%
\pgfpathlineto{\pgfqpoint{4.427120in}{2.912117in}}%
\pgfpathlineto{\pgfqpoint{4.427120in}{2.925561in}}%
\pgfpathlineto{\pgfqpoint{4.255037in}{2.925561in}}%
\pgfpathlineto{\pgfqpoint{4.255037in}{2.912117in}}%
\pgfusepath{fill}%
\end{pgfscope}%
\begin{pgfscope}%
\pgfpathrectangle{\pgfqpoint{4.255037in}{0.357758in}}{\pgfqpoint{0.172083in}{3.441662in}}%
\pgfusepath{clip}%
\pgfsetbuttcap%
\pgfsetroundjoin%
\definecolor{currentfill}{rgb}{0.971835,0.580382,0.254931}%
\pgfsetfillcolor{currentfill}%
\pgfsetlinewidth{0.000000pt}%
\definecolor{currentstroke}{rgb}{0.000000,0.000000,0.000000}%
\pgfsetstrokecolor{currentstroke}%
\pgfsetdash{}{0pt}%
\pgfpathmoveto{\pgfqpoint{4.255037in}{2.925561in}}%
\pgfpathlineto{\pgfqpoint{4.427120in}{2.925561in}}%
\pgfpathlineto{\pgfqpoint{4.427120in}{2.939005in}}%
\pgfpathlineto{\pgfqpoint{4.255037in}{2.939005in}}%
\pgfpathlineto{\pgfqpoint{4.255037in}{2.925561in}}%
\pgfusepath{fill}%
\end{pgfscope}%
\begin{pgfscope}%
\pgfpathrectangle{\pgfqpoint{4.255037in}{0.357758in}}{\pgfqpoint{0.172083in}{3.441662in}}%
\pgfusepath{clip}%
\pgfsetbuttcap%
\pgfsetroundjoin%
\definecolor{currentfill}{rgb}{0.973416,0.585761,0.251540}%
\pgfsetfillcolor{currentfill}%
\pgfsetlinewidth{0.000000pt}%
\definecolor{currentstroke}{rgb}{0.000000,0.000000,0.000000}%
\pgfsetstrokecolor{currentstroke}%
\pgfsetdash{}{0pt}%
\pgfpathmoveto{\pgfqpoint{4.255037in}{2.939005in}}%
\pgfpathlineto{\pgfqpoint{4.427120in}{2.939005in}}%
\pgfpathlineto{\pgfqpoint{4.427120in}{2.952449in}}%
\pgfpathlineto{\pgfqpoint{4.255037in}{2.952449in}}%
\pgfpathlineto{\pgfqpoint{4.255037in}{2.939005in}}%
\pgfusepath{fill}%
\end{pgfscope}%
\begin{pgfscope}%
\pgfpathrectangle{\pgfqpoint{4.255037in}{0.357758in}}{\pgfqpoint{0.172083in}{3.441662in}}%
\pgfusepath{clip}%
\pgfsetbuttcap%
\pgfsetroundjoin%
\definecolor{currentfill}{rgb}{0.974947,0.591165,0.248151}%
\pgfsetfillcolor{currentfill}%
\pgfsetlinewidth{0.000000pt}%
\definecolor{currentstroke}{rgb}{0.000000,0.000000,0.000000}%
\pgfsetstrokecolor{currentstroke}%
\pgfsetdash{}{0pt}%
\pgfpathmoveto{\pgfqpoint{4.255037in}{2.952449in}}%
\pgfpathlineto{\pgfqpoint{4.427120in}{2.952449in}}%
\pgfpathlineto{\pgfqpoint{4.427120in}{2.965893in}}%
\pgfpathlineto{\pgfqpoint{4.255037in}{2.965893in}}%
\pgfpathlineto{\pgfqpoint{4.255037in}{2.952449in}}%
\pgfusepath{fill}%
\end{pgfscope}%
\begin{pgfscope}%
\pgfpathrectangle{\pgfqpoint{4.255037in}{0.357758in}}{\pgfqpoint{0.172083in}{3.441662in}}%
\pgfusepath{clip}%
\pgfsetbuttcap%
\pgfsetroundjoin%
\definecolor{currentfill}{rgb}{0.976428,0.596595,0.244767}%
\pgfsetfillcolor{currentfill}%
\pgfsetlinewidth{0.000000pt}%
\definecolor{currentstroke}{rgb}{0.000000,0.000000,0.000000}%
\pgfsetstrokecolor{currentstroke}%
\pgfsetdash{}{0pt}%
\pgfpathmoveto{\pgfqpoint{4.255037in}{2.965893in}}%
\pgfpathlineto{\pgfqpoint{4.427120in}{2.965893in}}%
\pgfpathlineto{\pgfqpoint{4.427120in}{2.979337in}}%
\pgfpathlineto{\pgfqpoint{4.255037in}{2.979337in}}%
\pgfpathlineto{\pgfqpoint{4.255037in}{2.965893in}}%
\pgfusepath{fill}%
\end{pgfscope}%
\begin{pgfscope}%
\pgfpathrectangle{\pgfqpoint{4.255037in}{0.357758in}}{\pgfqpoint{0.172083in}{3.441662in}}%
\pgfusepath{clip}%
\pgfsetbuttcap%
\pgfsetroundjoin%
\definecolor{currentfill}{rgb}{0.977856,0.602051,0.241387}%
\pgfsetfillcolor{currentfill}%
\pgfsetlinewidth{0.000000pt}%
\definecolor{currentstroke}{rgb}{0.000000,0.000000,0.000000}%
\pgfsetstrokecolor{currentstroke}%
\pgfsetdash{}{0pt}%
\pgfpathmoveto{\pgfqpoint{4.255037in}{2.979337in}}%
\pgfpathlineto{\pgfqpoint{4.427120in}{2.979337in}}%
\pgfpathlineto{\pgfqpoint{4.427120in}{2.992781in}}%
\pgfpathlineto{\pgfqpoint{4.255037in}{2.992781in}}%
\pgfpathlineto{\pgfqpoint{4.255037in}{2.979337in}}%
\pgfusepath{fill}%
\end{pgfscope}%
\begin{pgfscope}%
\pgfpathrectangle{\pgfqpoint{4.255037in}{0.357758in}}{\pgfqpoint{0.172083in}{3.441662in}}%
\pgfusepath{clip}%
\pgfsetbuttcap%
\pgfsetroundjoin%
\definecolor{currentfill}{rgb}{0.979233,0.607532,0.238013}%
\pgfsetfillcolor{currentfill}%
\pgfsetlinewidth{0.000000pt}%
\definecolor{currentstroke}{rgb}{0.000000,0.000000,0.000000}%
\pgfsetstrokecolor{currentstroke}%
\pgfsetdash{}{0pt}%
\pgfpathmoveto{\pgfqpoint{4.255037in}{2.992781in}}%
\pgfpathlineto{\pgfqpoint{4.427120in}{2.992781in}}%
\pgfpathlineto{\pgfqpoint{4.427120in}{3.006225in}}%
\pgfpathlineto{\pgfqpoint{4.255037in}{3.006225in}}%
\pgfpathlineto{\pgfqpoint{4.255037in}{2.992781in}}%
\pgfusepath{fill}%
\end{pgfscope}%
\begin{pgfscope}%
\pgfpathrectangle{\pgfqpoint{4.255037in}{0.357758in}}{\pgfqpoint{0.172083in}{3.441662in}}%
\pgfusepath{clip}%
\pgfsetbuttcap%
\pgfsetroundjoin%
\definecolor{currentfill}{rgb}{0.980556,0.613039,0.234646}%
\pgfsetfillcolor{currentfill}%
\pgfsetlinewidth{0.000000pt}%
\definecolor{currentstroke}{rgb}{0.000000,0.000000,0.000000}%
\pgfsetstrokecolor{currentstroke}%
\pgfsetdash{}{0pt}%
\pgfpathmoveto{\pgfqpoint{4.255037in}{3.006225in}}%
\pgfpathlineto{\pgfqpoint{4.427120in}{3.006225in}}%
\pgfpathlineto{\pgfqpoint{4.427120in}{3.019669in}}%
\pgfpathlineto{\pgfqpoint{4.255037in}{3.019669in}}%
\pgfpathlineto{\pgfqpoint{4.255037in}{3.006225in}}%
\pgfusepath{fill}%
\end{pgfscope}%
\begin{pgfscope}%
\pgfpathrectangle{\pgfqpoint{4.255037in}{0.357758in}}{\pgfqpoint{0.172083in}{3.441662in}}%
\pgfusepath{clip}%
\pgfsetbuttcap%
\pgfsetroundjoin%
\definecolor{currentfill}{rgb}{0.981826,0.618572,0.231287}%
\pgfsetfillcolor{currentfill}%
\pgfsetlinewidth{0.000000pt}%
\definecolor{currentstroke}{rgb}{0.000000,0.000000,0.000000}%
\pgfsetstrokecolor{currentstroke}%
\pgfsetdash{}{0pt}%
\pgfpathmoveto{\pgfqpoint{4.255037in}{3.019669in}}%
\pgfpathlineto{\pgfqpoint{4.427120in}{3.019669in}}%
\pgfpathlineto{\pgfqpoint{4.427120in}{3.033113in}}%
\pgfpathlineto{\pgfqpoint{4.255037in}{3.033113in}}%
\pgfpathlineto{\pgfqpoint{4.255037in}{3.019669in}}%
\pgfusepath{fill}%
\end{pgfscope}%
\begin{pgfscope}%
\pgfpathrectangle{\pgfqpoint{4.255037in}{0.357758in}}{\pgfqpoint{0.172083in}{3.441662in}}%
\pgfusepath{clip}%
\pgfsetbuttcap%
\pgfsetroundjoin%
\definecolor{currentfill}{rgb}{0.983041,0.624131,0.227937}%
\pgfsetfillcolor{currentfill}%
\pgfsetlinewidth{0.000000pt}%
\definecolor{currentstroke}{rgb}{0.000000,0.000000,0.000000}%
\pgfsetstrokecolor{currentstroke}%
\pgfsetdash{}{0pt}%
\pgfpathmoveto{\pgfqpoint{4.255037in}{3.033113in}}%
\pgfpathlineto{\pgfqpoint{4.427120in}{3.033113in}}%
\pgfpathlineto{\pgfqpoint{4.427120in}{3.046557in}}%
\pgfpathlineto{\pgfqpoint{4.255037in}{3.046557in}}%
\pgfpathlineto{\pgfqpoint{4.255037in}{3.033113in}}%
\pgfusepath{fill}%
\end{pgfscope}%
\begin{pgfscope}%
\pgfpathrectangle{\pgfqpoint{4.255037in}{0.357758in}}{\pgfqpoint{0.172083in}{3.441662in}}%
\pgfusepath{clip}%
\pgfsetbuttcap%
\pgfsetroundjoin%
\definecolor{currentfill}{rgb}{0.984199,0.629718,0.224595}%
\pgfsetfillcolor{currentfill}%
\pgfsetlinewidth{0.000000pt}%
\definecolor{currentstroke}{rgb}{0.000000,0.000000,0.000000}%
\pgfsetstrokecolor{currentstroke}%
\pgfsetdash{}{0pt}%
\pgfpathmoveto{\pgfqpoint{4.255037in}{3.046557in}}%
\pgfpathlineto{\pgfqpoint{4.427120in}{3.046557in}}%
\pgfpathlineto{\pgfqpoint{4.427120in}{3.060001in}}%
\pgfpathlineto{\pgfqpoint{4.255037in}{3.060001in}}%
\pgfpathlineto{\pgfqpoint{4.255037in}{3.046557in}}%
\pgfusepath{fill}%
\end{pgfscope}%
\begin{pgfscope}%
\pgfpathrectangle{\pgfqpoint{4.255037in}{0.357758in}}{\pgfqpoint{0.172083in}{3.441662in}}%
\pgfusepath{clip}%
\pgfsetbuttcap%
\pgfsetroundjoin%
\definecolor{currentfill}{rgb}{0.985301,0.635330,0.221265}%
\pgfsetfillcolor{currentfill}%
\pgfsetlinewidth{0.000000pt}%
\definecolor{currentstroke}{rgb}{0.000000,0.000000,0.000000}%
\pgfsetstrokecolor{currentstroke}%
\pgfsetdash{}{0pt}%
\pgfpathmoveto{\pgfqpoint{4.255037in}{3.060001in}}%
\pgfpathlineto{\pgfqpoint{4.427120in}{3.060001in}}%
\pgfpathlineto{\pgfqpoint{4.427120in}{3.073445in}}%
\pgfpathlineto{\pgfqpoint{4.255037in}{3.073445in}}%
\pgfpathlineto{\pgfqpoint{4.255037in}{3.060001in}}%
\pgfusepath{fill}%
\end{pgfscope}%
\begin{pgfscope}%
\pgfpathrectangle{\pgfqpoint{4.255037in}{0.357758in}}{\pgfqpoint{0.172083in}{3.441662in}}%
\pgfusepath{clip}%
\pgfsetbuttcap%
\pgfsetroundjoin%
\definecolor{currentfill}{rgb}{0.986345,0.640969,0.217948}%
\pgfsetfillcolor{currentfill}%
\pgfsetlinewidth{0.000000pt}%
\definecolor{currentstroke}{rgb}{0.000000,0.000000,0.000000}%
\pgfsetstrokecolor{currentstroke}%
\pgfsetdash{}{0pt}%
\pgfpathmoveto{\pgfqpoint{4.255037in}{3.073445in}}%
\pgfpathlineto{\pgfqpoint{4.427120in}{3.073445in}}%
\pgfpathlineto{\pgfqpoint{4.427120in}{3.086889in}}%
\pgfpathlineto{\pgfqpoint{4.255037in}{3.086889in}}%
\pgfpathlineto{\pgfqpoint{4.255037in}{3.073445in}}%
\pgfusepath{fill}%
\end{pgfscope}%
\begin{pgfscope}%
\pgfpathrectangle{\pgfqpoint{4.255037in}{0.357758in}}{\pgfqpoint{0.172083in}{3.441662in}}%
\pgfusepath{clip}%
\pgfsetbuttcap%
\pgfsetroundjoin%
\definecolor{currentfill}{rgb}{0.987332,0.646633,0.214648}%
\pgfsetfillcolor{currentfill}%
\pgfsetlinewidth{0.000000pt}%
\definecolor{currentstroke}{rgb}{0.000000,0.000000,0.000000}%
\pgfsetstrokecolor{currentstroke}%
\pgfsetdash{}{0pt}%
\pgfpathmoveto{\pgfqpoint{4.255037in}{3.086889in}}%
\pgfpathlineto{\pgfqpoint{4.427120in}{3.086889in}}%
\pgfpathlineto{\pgfqpoint{4.427120in}{3.100333in}}%
\pgfpathlineto{\pgfqpoint{4.255037in}{3.100333in}}%
\pgfpathlineto{\pgfqpoint{4.255037in}{3.086889in}}%
\pgfusepath{fill}%
\end{pgfscope}%
\begin{pgfscope}%
\pgfpathrectangle{\pgfqpoint{4.255037in}{0.357758in}}{\pgfqpoint{0.172083in}{3.441662in}}%
\pgfusepath{clip}%
\pgfsetbuttcap%
\pgfsetroundjoin%
\definecolor{currentfill}{rgb}{0.988260,0.652325,0.211364}%
\pgfsetfillcolor{currentfill}%
\pgfsetlinewidth{0.000000pt}%
\definecolor{currentstroke}{rgb}{0.000000,0.000000,0.000000}%
\pgfsetstrokecolor{currentstroke}%
\pgfsetdash{}{0pt}%
\pgfpathmoveto{\pgfqpoint{4.255037in}{3.100333in}}%
\pgfpathlineto{\pgfqpoint{4.427120in}{3.100333in}}%
\pgfpathlineto{\pgfqpoint{4.427120in}{3.113777in}}%
\pgfpathlineto{\pgfqpoint{4.255037in}{3.113777in}}%
\pgfpathlineto{\pgfqpoint{4.255037in}{3.100333in}}%
\pgfusepath{fill}%
\end{pgfscope}%
\begin{pgfscope}%
\pgfpathrectangle{\pgfqpoint{4.255037in}{0.357758in}}{\pgfqpoint{0.172083in}{3.441662in}}%
\pgfusepath{clip}%
\pgfsetbuttcap%
\pgfsetroundjoin%
\definecolor{currentfill}{rgb}{0.989128,0.658043,0.208100}%
\pgfsetfillcolor{currentfill}%
\pgfsetlinewidth{0.000000pt}%
\definecolor{currentstroke}{rgb}{0.000000,0.000000,0.000000}%
\pgfsetstrokecolor{currentstroke}%
\pgfsetdash{}{0pt}%
\pgfpathmoveto{\pgfqpoint{4.255037in}{3.113777in}}%
\pgfpathlineto{\pgfqpoint{4.427120in}{3.113777in}}%
\pgfpathlineto{\pgfqpoint{4.427120in}{3.127221in}}%
\pgfpathlineto{\pgfqpoint{4.255037in}{3.127221in}}%
\pgfpathlineto{\pgfqpoint{4.255037in}{3.113777in}}%
\pgfusepath{fill}%
\end{pgfscope}%
\begin{pgfscope}%
\pgfpathrectangle{\pgfqpoint{4.255037in}{0.357758in}}{\pgfqpoint{0.172083in}{3.441662in}}%
\pgfusepath{clip}%
\pgfsetbuttcap%
\pgfsetroundjoin%
\definecolor{currentfill}{rgb}{0.989935,0.663787,0.204859}%
\pgfsetfillcolor{currentfill}%
\pgfsetlinewidth{0.000000pt}%
\definecolor{currentstroke}{rgb}{0.000000,0.000000,0.000000}%
\pgfsetstrokecolor{currentstroke}%
\pgfsetdash{}{0pt}%
\pgfpathmoveto{\pgfqpoint{4.255037in}{3.127221in}}%
\pgfpathlineto{\pgfqpoint{4.427120in}{3.127221in}}%
\pgfpathlineto{\pgfqpoint{4.427120in}{3.140665in}}%
\pgfpathlineto{\pgfqpoint{4.255037in}{3.140665in}}%
\pgfpathlineto{\pgfqpoint{4.255037in}{3.127221in}}%
\pgfusepath{fill}%
\end{pgfscope}%
\begin{pgfscope}%
\pgfpathrectangle{\pgfqpoint{4.255037in}{0.357758in}}{\pgfqpoint{0.172083in}{3.441662in}}%
\pgfusepath{clip}%
\pgfsetbuttcap%
\pgfsetroundjoin%
\definecolor{currentfill}{rgb}{0.990681,0.669558,0.201642}%
\pgfsetfillcolor{currentfill}%
\pgfsetlinewidth{0.000000pt}%
\definecolor{currentstroke}{rgb}{0.000000,0.000000,0.000000}%
\pgfsetstrokecolor{currentstroke}%
\pgfsetdash{}{0pt}%
\pgfpathmoveto{\pgfqpoint{4.255037in}{3.140665in}}%
\pgfpathlineto{\pgfqpoint{4.427120in}{3.140665in}}%
\pgfpathlineto{\pgfqpoint{4.427120in}{3.154109in}}%
\pgfpathlineto{\pgfqpoint{4.255037in}{3.154109in}}%
\pgfpathlineto{\pgfqpoint{4.255037in}{3.140665in}}%
\pgfusepath{fill}%
\end{pgfscope}%
\begin{pgfscope}%
\pgfpathrectangle{\pgfqpoint{4.255037in}{0.357758in}}{\pgfqpoint{0.172083in}{3.441662in}}%
\pgfusepath{clip}%
\pgfsetbuttcap%
\pgfsetroundjoin%
\definecolor{currentfill}{rgb}{0.991365,0.675355,0.198453}%
\pgfsetfillcolor{currentfill}%
\pgfsetlinewidth{0.000000pt}%
\definecolor{currentstroke}{rgb}{0.000000,0.000000,0.000000}%
\pgfsetstrokecolor{currentstroke}%
\pgfsetdash{}{0pt}%
\pgfpathmoveto{\pgfqpoint{4.255037in}{3.154109in}}%
\pgfpathlineto{\pgfqpoint{4.427120in}{3.154109in}}%
\pgfpathlineto{\pgfqpoint{4.427120in}{3.167553in}}%
\pgfpathlineto{\pgfqpoint{4.255037in}{3.167553in}}%
\pgfpathlineto{\pgfqpoint{4.255037in}{3.154109in}}%
\pgfusepath{fill}%
\end{pgfscope}%
\begin{pgfscope}%
\pgfpathrectangle{\pgfqpoint{4.255037in}{0.357758in}}{\pgfqpoint{0.172083in}{3.441662in}}%
\pgfusepath{clip}%
\pgfsetbuttcap%
\pgfsetroundjoin%
\definecolor{currentfill}{rgb}{0.991985,0.681179,0.195295}%
\pgfsetfillcolor{currentfill}%
\pgfsetlinewidth{0.000000pt}%
\definecolor{currentstroke}{rgb}{0.000000,0.000000,0.000000}%
\pgfsetstrokecolor{currentstroke}%
\pgfsetdash{}{0pt}%
\pgfpathmoveto{\pgfqpoint{4.255037in}{3.167553in}}%
\pgfpathlineto{\pgfqpoint{4.427120in}{3.167553in}}%
\pgfpathlineto{\pgfqpoint{4.427120in}{3.180997in}}%
\pgfpathlineto{\pgfqpoint{4.255037in}{3.180997in}}%
\pgfpathlineto{\pgfqpoint{4.255037in}{3.167553in}}%
\pgfusepath{fill}%
\end{pgfscope}%
\begin{pgfscope}%
\pgfpathrectangle{\pgfqpoint{4.255037in}{0.357758in}}{\pgfqpoint{0.172083in}{3.441662in}}%
\pgfusepath{clip}%
\pgfsetbuttcap%
\pgfsetroundjoin%
\definecolor{currentfill}{rgb}{0.992541,0.687030,0.192170}%
\pgfsetfillcolor{currentfill}%
\pgfsetlinewidth{0.000000pt}%
\definecolor{currentstroke}{rgb}{0.000000,0.000000,0.000000}%
\pgfsetstrokecolor{currentstroke}%
\pgfsetdash{}{0pt}%
\pgfpathmoveto{\pgfqpoint{4.255037in}{3.180997in}}%
\pgfpathlineto{\pgfqpoint{4.427120in}{3.180997in}}%
\pgfpathlineto{\pgfqpoint{4.427120in}{3.194441in}}%
\pgfpathlineto{\pgfqpoint{4.255037in}{3.194441in}}%
\pgfpathlineto{\pgfqpoint{4.255037in}{3.180997in}}%
\pgfusepath{fill}%
\end{pgfscope}%
\begin{pgfscope}%
\pgfpathrectangle{\pgfqpoint{4.255037in}{0.357758in}}{\pgfqpoint{0.172083in}{3.441662in}}%
\pgfusepath{clip}%
\pgfsetbuttcap%
\pgfsetroundjoin%
\definecolor{currentfill}{rgb}{0.993032,0.692907,0.189084}%
\pgfsetfillcolor{currentfill}%
\pgfsetlinewidth{0.000000pt}%
\definecolor{currentstroke}{rgb}{0.000000,0.000000,0.000000}%
\pgfsetstrokecolor{currentstroke}%
\pgfsetdash{}{0pt}%
\pgfpathmoveto{\pgfqpoint{4.255037in}{3.194441in}}%
\pgfpathlineto{\pgfqpoint{4.427120in}{3.194441in}}%
\pgfpathlineto{\pgfqpoint{4.427120in}{3.207885in}}%
\pgfpathlineto{\pgfqpoint{4.255037in}{3.207885in}}%
\pgfpathlineto{\pgfqpoint{4.255037in}{3.194441in}}%
\pgfusepath{fill}%
\end{pgfscope}%
\begin{pgfscope}%
\pgfpathrectangle{\pgfqpoint{4.255037in}{0.357758in}}{\pgfqpoint{0.172083in}{3.441662in}}%
\pgfusepath{clip}%
\pgfsetbuttcap%
\pgfsetroundjoin%
\definecolor{currentfill}{rgb}{0.993456,0.698810,0.186041}%
\pgfsetfillcolor{currentfill}%
\pgfsetlinewidth{0.000000pt}%
\definecolor{currentstroke}{rgb}{0.000000,0.000000,0.000000}%
\pgfsetstrokecolor{currentstroke}%
\pgfsetdash{}{0pt}%
\pgfpathmoveto{\pgfqpoint{4.255037in}{3.207885in}}%
\pgfpathlineto{\pgfqpoint{4.427120in}{3.207885in}}%
\pgfpathlineto{\pgfqpoint{4.427120in}{3.221329in}}%
\pgfpathlineto{\pgfqpoint{4.255037in}{3.221329in}}%
\pgfpathlineto{\pgfqpoint{4.255037in}{3.207885in}}%
\pgfusepath{fill}%
\end{pgfscope}%
\begin{pgfscope}%
\pgfpathrectangle{\pgfqpoint{4.255037in}{0.357758in}}{\pgfqpoint{0.172083in}{3.441662in}}%
\pgfusepath{clip}%
\pgfsetbuttcap%
\pgfsetroundjoin%
\definecolor{currentfill}{rgb}{0.993814,0.704741,0.183043}%
\pgfsetfillcolor{currentfill}%
\pgfsetlinewidth{0.000000pt}%
\definecolor{currentstroke}{rgb}{0.000000,0.000000,0.000000}%
\pgfsetstrokecolor{currentstroke}%
\pgfsetdash{}{0pt}%
\pgfpathmoveto{\pgfqpoint{4.255037in}{3.221329in}}%
\pgfpathlineto{\pgfqpoint{4.427120in}{3.221329in}}%
\pgfpathlineto{\pgfqpoint{4.427120in}{3.234773in}}%
\pgfpathlineto{\pgfqpoint{4.255037in}{3.234773in}}%
\pgfpathlineto{\pgfqpoint{4.255037in}{3.221329in}}%
\pgfusepath{fill}%
\end{pgfscope}%
\begin{pgfscope}%
\pgfpathrectangle{\pgfqpoint{4.255037in}{0.357758in}}{\pgfqpoint{0.172083in}{3.441662in}}%
\pgfusepath{clip}%
\pgfsetbuttcap%
\pgfsetroundjoin%
\definecolor{currentfill}{rgb}{0.994103,0.710698,0.180097}%
\pgfsetfillcolor{currentfill}%
\pgfsetlinewidth{0.000000pt}%
\definecolor{currentstroke}{rgb}{0.000000,0.000000,0.000000}%
\pgfsetstrokecolor{currentstroke}%
\pgfsetdash{}{0pt}%
\pgfpathmoveto{\pgfqpoint{4.255037in}{3.234773in}}%
\pgfpathlineto{\pgfqpoint{4.427120in}{3.234773in}}%
\pgfpathlineto{\pgfqpoint{4.427120in}{3.248216in}}%
\pgfpathlineto{\pgfqpoint{4.255037in}{3.248216in}}%
\pgfpathlineto{\pgfqpoint{4.255037in}{3.234773in}}%
\pgfusepath{fill}%
\end{pgfscope}%
\begin{pgfscope}%
\pgfpathrectangle{\pgfqpoint{4.255037in}{0.357758in}}{\pgfqpoint{0.172083in}{3.441662in}}%
\pgfusepath{clip}%
\pgfsetbuttcap%
\pgfsetroundjoin%
\definecolor{currentfill}{rgb}{0.994324,0.716681,0.177208}%
\pgfsetfillcolor{currentfill}%
\pgfsetlinewidth{0.000000pt}%
\definecolor{currentstroke}{rgb}{0.000000,0.000000,0.000000}%
\pgfsetstrokecolor{currentstroke}%
\pgfsetdash{}{0pt}%
\pgfpathmoveto{\pgfqpoint{4.255037in}{3.248216in}}%
\pgfpathlineto{\pgfqpoint{4.427120in}{3.248216in}}%
\pgfpathlineto{\pgfqpoint{4.427120in}{3.261660in}}%
\pgfpathlineto{\pgfqpoint{4.255037in}{3.261660in}}%
\pgfpathlineto{\pgfqpoint{4.255037in}{3.248216in}}%
\pgfusepath{fill}%
\end{pgfscope}%
\begin{pgfscope}%
\pgfpathrectangle{\pgfqpoint{4.255037in}{0.357758in}}{\pgfqpoint{0.172083in}{3.441662in}}%
\pgfusepath{clip}%
\pgfsetbuttcap%
\pgfsetroundjoin%
\definecolor{currentfill}{rgb}{0.994474,0.722691,0.174381}%
\pgfsetfillcolor{currentfill}%
\pgfsetlinewidth{0.000000pt}%
\definecolor{currentstroke}{rgb}{0.000000,0.000000,0.000000}%
\pgfsetstrokecolor{currentstroke}%
\pgfsetdash{}{0pt}%
\pgfpathmoveto{\pgfqpoint{4.255037in}{3.261660in}}%
\pgfpathlineto{\pgfqpoint{4.427120in}{3.261660in}}%
\pgfpathlineto{\pgfqpoint{4.427120in}{3.275104in}}%
\pgfpathlineto{\pgfqpoint{4.255037in}{3.275104in}}%
\pgfpathlineto{\pgfqpoint{4.255037in}{3.261660in}}%
\pgfusepath{fill}%
\end{pgfscope}%
\begin{pgfscope}%
\pgfpathrectangle{\pgfqpoint{4.255037in}{0.357758in}}{\pgfqpoint{0.172083in}{3.441662in}}%
\pgfusepath{clip}%
\pgfsetbuttcap%
\pgfsetroundjoin%
\definecolor{currentfill}{rgb}{0.994553,0.728728,0.171622}%
\pgfsetfillcolor{currentfill}%
\pgfsetlinewidth{0.000000pt}%
\definecolor{currentstroke}{rgb}{0.000000,0.000000,0.000000}%
\pgfsetstrokecolor{currentstroke}%
\pgfsetdash{}{0pt}%
\pgfpathmoveto{\pgfqpoint{4.255037in}{3.275104in}}%
\pgfpathlineto{\pgfqpoint{4.427120in}{3.275104in}}%
\pgfpathlineto{\pgfqpoint{4.427120in}{3.288548in}}%
\pgfpathlineto{\pgfqpoint{4.255037in}{3.288548in}}%
\pgfpathlineto{\pgfqpoint{4.255037in}{3.275104in}}%
\pgfusepath{fill}%
\end{pgfscope}%
\begin{pgfscope}%
\pgfpathrectangle{\pgfqpoint{4.255037in}{0.357758in}}{\pgfqpoint{0.172083in}{3.441662in}}%
\pgfusepath{clip}%
\pgfsetbuttcap%
\pgfsetroundjoin%
\definecolor{currentfill}{rgb}{0.994561,0.734791,0.168938}%
\pgfsetfillcolor{currentfill}%
\pgfsetlinewidth{0.000000pt}%
\definecolor{currentstroke}{rgb}{0.000000,0.000000,0.000000}%
\pgfsetstrokecolor{currentstroke}%
\pgfsetdash{}{0pt}%
\pgfpathmoveto{\pgfqpoint{4.255037in}{3.288548in}}%
\pgfpathlineto{\pgfqpoint{4.427120in}{3.288548in}}%
\pgfpathlineto{\pgfqpoint{4.427120in}{3.301992in}}%
\pgfpathlineto{\pgfqpoint{4.255037in}{3.301992in}}%
\pgfpathlineto{\pgfqpoint{4.255037in}{3.288548in}}%
\pgfusepath{fill}%
\end{pgfscope}%
\begin{pgfscope}%
\pgfpathrectangle{\pgfqpoint{4.255037in}{0.357758in}}{\pgfqpoint{0.172083in}{3.441662in}}%
\pgfusepath{clip}%
\pgfsetbuttcap%
\pgfsetroundjoin%
\definecolor{currentfill}{rgb}{0.994495,0.740880,0.166335}%
\pgfsetfillcolor{currentfill}%
\pgfsetlinewidth{0.000000pt}%
\definecolor{currentstroke}{rgb}{0.000000,0.000000,0.000000}%
\pgfsetstrokecolor{currentstroke}%
\pgfsetdash{}{0pt}%
\pgfpathmoveto{\pgfqpoint{4.255037in}{3.301992in}}%
\pgfpathlineto{\pgfqpoint{4.427120in}{3.301992in}}%
\pgfpathlineto{\pgfqpoint{4.427120in}{3.315436in}}%
\pgfpathlineto{\pgfqpoint{4.255037in}{3.315436in}}%
\pgfpathlineto{\pgfqpoint{4.255037in}{3.301992in}}%
\pgfusepath{fill}%
\end{pgfscope}%
\begin{pgfscope}%
\pgfpathrectangle{\pgfqpoint{4.255037in}{0.357758in}}{\pgfqpoint{0.172083in}{3.441662in}}%
\pgfusepath{clip}%
\pgfsetbuttcap%
\pgfsetroundjoin%
\definecolor{currentfill}{rgb}{0.994355,0.746995,0.163821}%
\pgfsetfillcolor{currentfill}%
\pgfsetlinewidth{0.000000pt}%
\definecolor{currentstroke}{rgb}{0.000000,0.000000,0.000000}%
\pgfsetstrokecolor{currentstroke}%
\pgfsetdash{}{0pt}%
\pgfpathmoveto{\pgfqpoint{4.255037in}{3.315436in}}%
\pgfpathlineto{\pgfqpoint{4.427120in}{3.315436in}}%
\pgfpathlineto{\pgfqpoint{4.427120in}{3.328880in}}%
\pgfpathlineto{\pgfqpoint{4.255037in}{3.328880in}}%
\pgfpathlineto{\pgfqpoint{4.255037in}{3.315436in}}%
\pgfusepath{fill}%
\end{pgfscope}%
\begin{pgfscope}%
\pgfpathrectangle{\pgfqpoint{4.255037in}{0.357758in}}{\pgfqpoint{0.172083in}{3.441662in}}%
\pgfusepath{clip}%
\pgfsetbuttcap%
\pgfsetroundjoin%
\definecolor{currentfill}{rgb}{0.994141,0.753137,0.161404}%
\pgfsetfillcolor{currentfill}%
\pgfsetlinewidth{0.000000pt}%
\definecolor{currentstroke}{rgb}{0.000000,0.000000,0.000000}%
\pgfsetstrokecolor{currentstroke}%
\pgfsetdash{}{0pt}%
\pgfpathmoveto{\pgfqpoint{4.255037in}{3.328880in}}%
\pgfpathlineto{\pgfqpoint{4.427120in}{3.328880in}}%
\pgfpathlineto{\pgfqpoint{4.427120in}{3.342324in}}%
\pgfpathlineto{\pgfqpoint{4.255037in}{3.342324in}}%
\pgfpathlineto{\pgfqpoint{4.255037in}{3.328880in}}%
\pgfusepath{fill}%
\end{pgfscope}%
\begin{pgfscope}%
\pgfpathrectangle{\pgfqpoint{4.255037in}{0.357758in}}{\pgfqpoint{0.172083in}{3.441662in}}%
\pgfusepath{clip}%
\pgfsetbuttcap%
\pgfsetroundjoin%
\definecolor{currentfill}{rgb}{0.993851,0.759304,0.159092}%
\pgfsetfillcolor{currentfill}%
\pgfsetlinewidth{0.000000pt}%
\definecolor{currentstroke}{rgb}{0.000000,0.000000,0.000000}%
\pgfsetstrokecolor{currentstroke}%
\pgfsetdash{}{0pt}%
\pgfpathmoveto{\pgfqpoint{4.255037in}{3.342324in}}%
\pgfpathlineto{\pgfqpoint{4.427120in}{3.342324in}}%
\pgfpathlineto{\pgfqpoint{4.427120in}{3.355768in}}%
\pgfpathlineto{\pgfqpoint{4.255037in}{3.355768in}}%
\pgfpathlineto{\pgfqpoint{4.255037in}{3.342324in}}%
\pgfusepath{fill}%
\end{pgfscope}%
\begin{pgfscope}%
\pgfpathrectangle{\pgfqpoint{4.255037in}{0.357758in}}{\pgfqpoint{0.172083in}{3.441662in}}%
\pgfusepath{clip}%
\pgfsetbuttcap%
\pgfsetroundjoin%
\definecolor{currentfill}{rgb}{0.993482,0.765499,0.156891}%
\pgfsetfillcolor{currentfill}%
\pgfsetlinewidth{0.000000pt}%
\definecolor{currentstroke}{rgb}{0.000000,0.000000,0.000000}%
\pgfsetstrokecolor{currentstroke}%
\pgfsetdash{}{0pt}%
\pgfpathmoveto{\pgfqpoint{4.255037in}{3.355768in}}%
\pgfpathlineto{\pgfqpoint{4.427120in}{3.355768in}}%
\pgfpathlineto{\pgfqpoint{4.427120in}{3.369212in}}%
\pgfpathlineto{\pgfqpoint{4.255037in}{3.369212in}}%
\pgfpathlineto{\pgfqpoint{4.255037in}{3.355768in}}%
\pgfusepath{fill}%
\end{pgfscope}%
\begin{pgfscope}%
\pgfpathrectangle{\pgfqpoint{4.255037in}{0.357758in}}{\pgfqpoint{0.172083in}{3.441662in}}%
\pgfusepath{clip}%
\pgfsetbuttcap%
\pgfsetroundjoin%
\definecolor{currentfill}{rgb}{0.993033,0.771720,0.154808}%
\pgfsetfillcolor{currentfill}%
\pgfsetlinewidth{0.000000pt}%
\definecolor{currentstroke}{rgb}{0.000000,0.000000,0.000000}%
\pgfsetstrokecolor{currentstroke}%
\pgfsetdash{}{0pt}%
\pgfpathmoveto{\pgfqpoint{4.255037in}{3.369212in}}%
\pgfpathlineto{\pgfqpoint{4.427120in}{3.369212in}}%
\pgfpathlineto{\pgfqpoint{4.427120in}{3.382656in}}%
\pgfpathlineto{\pgfqpoint{4.255037in}{3.382656in}}%
\pgfpathlineto{\pgfqpoint{4.255037in}{3.369212in}}%
\pgfusepath{fill}%
\end{pgfscope}%
\begin{pgfscope}%
\pgfpathrectangle{\pgfqpoint{4.255037in}{0.357758in}}{\pgfqpoint{0.172083in}{3.441662in}}%
\pgfusepath{clip}%
\pgfsetbuttcap%
\pgfsetroundjoin%
\definecolor{currentfill}{rgb}{0.992505,0.777967,0.152855}%
\pgfsetfillcolor{currentfill}%
\pgfsetlinewidth{0.000000pt}%
\definecolor{currentstroke}{rgb}{0.000000,0.000000,0.000000}%
\pgfsetstrokecolor{currentstroke}%
\pgfsetdash{}{0pt}%
\pgfpathmoveto{\pgfqpoint{4.255037in}{3.382656in}}%
\pgfpathlineto{\pgfqpoint{4.427120in}{3.382656in}}%
\pgfpathlineto{\pgfqpoint{4.427120in}{3.396100in}}%
\pgfpathlineto{\pgfqpoint{4.255037in}{3.396100in}}%
\pgfpathlineto{\pgfqpoint{4.255037in}{3.382656in}}%
\pgfusepath{fill}%
\end{pgfscope}%
\begin{pgfscope}%
\pgfpathrectangle{\pgfqpoint{4.255037in}{0.357758in}}{\pgfqpoint{0.172083in}{3.441662in}}%
\pgfusepath{clip}%
\pgfsetbuttcap%
\pgfsetroundjoin%
\definecolor{currentfill}{rgb}{0.991897,0.784239,0.151042}%
\pgfsetfillcolor{currentfill}%
\pgfsetlinewidth{0.000000pt}%
\definecolor{currentstroke}{rgb}{0.000000,0.000000,0.000000}%
\pgfsetstrokecolor{currentstroke}%
\pgfsetdash{}{0pt}%
\pgfpathmoveto{\pgfqpoint{4.255037in}{3.396100in}}%
\pgfpathlineto{\pgfqpoint{4.427120in}{3.396100in}}%
\pgfpathlineto{\pgfqpoint{4.427120in}{3.409544in}}%
\pgfpathlineto{\pgfqpoint{4.255037in}{3.409544in}}%
\pgfpathlineto{\pgfqpoint{4.255037in}{3.396100in}}%
\pgfusepath{fill}%
\end{pgfscope}%
\begin{pgfscope}%
\pgfpathrectangle{\pgfqpoint{4.255037in}{0.357758in}}{\pgfqpoint{0.172083in}{3.441662in}}%
\pgfusepath{clip}%
\pgfsetbuttcap%
\pgfsetroundjoin%
\definecolor{currentfill}{rgb}{0.991209,0.790537,0.149377}%
\pgfsetfillcolor{currentfill}%
\pgfsetlinewidth{0.000000pt}%
\definecolor{currentstroke}{rgb}{0.000000,0.000000,0.000000}%
\pgfsetstrokecolor{currentstroke}%
\pgfsetdash{}{0pt}%
\pgfpathmoveto{\pgfqpoint{4.255037in}{3.409544in}}%
\pgfpathlineto{\pgfqpoint{4.427120in}{3.409544in}}%
\pgfpathlineto{\pgfqpoint{4.427120in}{3.422988in}}%
\pgfpathlineto{\pgfqpoint{4.255037in}{3.422988in}}%
\pgfpathlineto{\pgfqpoint{4.255037in}{3.409544in}}%
\pgfusepath{fill}%
\end{pgfscope}%
\begin{pgfscope}%
\pgfpathrectangle{\pgfqpoint{4.255037in}{0.357758in}}{\pgfqpoint{0.172083in}{3.441662in}}%
\pgfusepath{clip}%
\pgfsetbuttcap%
\pgfsetroundjoin%
\definecolor{currentfill}{rgb}{0.990439,0.796859,0.147870}%
\pgfsetfillcolor{currentfill}%
\pgfsetlinewidth{0.000000pt}%
\definecolor{currentstroke}{rgb}{0.000000,0.000000,0.000000}%
\pgfsetstrokecolor{currentstroke}%
\pgfsetdash{}{0pt}%
\pgfpathmoveto{\pgfqpoint{4.255037in}{3.422988in}}%
\pgfpathlineto{\pgfqpoint{4.427120in}{3.422988in}}%
\pgfpathlineto{\pgfqpoint{4.427120in}{3.436432in}}%
\pgfpathlineto{\pgfqpoint{4.255037in}{3.436432in}}%
\pgfpathlineto{\pgfqpoint{4.255037in}{3.422988in}}%
\pgfusepath{fill}%
\end{pgfscope}%
\begin{pgfscope}%
\pgfpathrectangle{\pgfqpoint{4.255037in}{0.357758in}}{\pgfqpoint{0.172083in}{3.441662in}}%
\pgfusepath{clip}%
\pgfsetbuttcap%
\pgfsetroundjoin%
\definecolor{currentfill}{rgb}{0.989587,0.803205,0.146529}%
\pgfsetfillcolor{currentfill}%
\pgfsetlinewidth{0.000000pt}%
\definecolor{currentstroke}{rgb}{0.000000,0.000000,0.000000}%
\pgfsetstrokecolor{currentstroke}%
\pgfsetdash{}{0pt}%
\pgfpathmoveto{\pgfqpoint{4.255037in}{3.436432in}}%
\pgfpathlineto{\pgfqpoint{4.427120in}{3.436432in}}%
\pgfpathlineto{\pgfqpoint{4.427120in}{3.449876in}}%
\pgfpathlineto{\pgfqpoint{4.255037in}{3.449876in}}%
\pgfpathlineto{\pgfqpoint{4.255037in}{3.436432in}}%
\pgfusepath{fill}%
\end{pgfscope}%
\begin{pgfscope}%
\pgfpathrectangle{\pgfqpoint{4.255037in}{0.357758in}}{\pgfqpoint{0.172083in}{3.441662in}}%
\pgfusepath{clip}%
\pgfsetbuttcap%
\pgfsetroundjoin%
\definecolor{currentfill}{rgb}{0.988648,0.809579,0.145357}%
\pgfsetfillcolor{currentfill}%
\pgfsetlinewidth{0.000000pt}%
\definecolor{currentstroke}{rgb}{0.000000,0.000000,0.000000}%
\pgfsetstrokecolor{currentstroke}%
\pgfsetdash{}{0pt}%
\pgfpathmoveto{\pgfqpoint{4.255037in}{3.449876in}}%
\pgfpathlineto{\pgfqpoint{4.427120in}{3.449876in}}%
\pgfpathlineto{\pgfqpoint{4.427120in}{3.463320in}}%
\pgfpathlineto{\pgfqpoint{4.255037in}{3.463320in}}%
\pgfpathlineto{\pgfqpoint{4.255037in}{3.449876in}}%
\pgfusepath{fill}%
\end{pgfscope}%
\begin{pgfscope}%
\pgfpathrectangle{\pgfqpoint{4.255037in}{0.357758in}}{\pgfqpoint{0.172083in}{3.441662in}}%
\pgfusepath{clip}%
\pgfsetbuttcap%
\pgfsetroundjoin%
\definecolor{currentfill}{rgb}{0.987621,0.815978,0.144363}%
\pgfsetfillcolor{currentfill}%
\pgfsetlinewidth{0.000000pt}%
\definecolor{currentstroke}{rgb}{0.000000,0.000000,0.000000}%
\pgfsetstrokecolor{currentstroke}%
\pgfsetdash{}{0pt}%
\pgfpathmoveto{\pgfqpoint{4.255037in}{3.463320in}}%
\pgfpathlineto{\pgfqpoint{4.427120in}{3.463320in}}%
\pgfpathlineto{\pgfqpoint{4.427120in}{3.476764in}}%
\pgfpathlineto{\pgfqpoint{4.255037in}{3.476764in}}%
\pgfpathlineto{\pgfqpoint{4.255037in}{3.463320in}}%
\pgfusepath{fill}%
\end{pgfscope}%
\begin{pgfscope}%
\pgfpathrectangle{\pgfqpoint{4.255037in}{0.357758in}}{\pgfqpoint{0.172083in}{3.441662in}}%
\pgfusepath{clip}%
\pgfsetbuttcap%
\pgfsetroundjoin%
\definecolor{currentfill}{rgb}{0.986509,0.822401,0.143557}%
\pgfsetfillcolor{currentfill}%
\pgfsetlinewidth{0.000000pt}%
\definecolor{currentstroke}{rgb}{0.000000,0.000000,0.000000}%
\pgfsetstrokecolor{currentstroke}%
\pgfsetdash{}{0pt}%
\pgfpathmoveto{\pgfqpoint{4.255037in}{3.476764in}}%
\pgfpathlineto{\pgfqpoint{4.427120in}{3.476764in}}%
\pgfpathlineto{\pgfqpoint{4.427120in}{3.490208in}}%
\pgfpathlineto{\pgfqpoint{4.255037in}{3.490208in}}%
\pgfpathlineto{\pgfqpoint{4.255037in}{3.476764in}}%
\pgfusepath{fill}%
\end{pgfscope}%
\begin{pgfscope}%
\pgfpathrectangle{\pgfqpoint{4.255037in}{0.357758in}}{\pgfqpoint{0.172083in}{3.441662in}}%
\pgfusepath{clip}%
\pgfsetbuttcap%
\pgfsetroundjoin%
\definecolor{currentfill}{rgb}{0.985314,0.828846,0.142945}%
\pgfsetfillcolor{currentfill}%
\pgfsetlinewidth{0.000000pt}%
\definecolor{currentstroke}{rgb}{0.000000,0.000000,0.000000}%
\pgfsetstrokecolor{currentstroke}%
\pgfsetdash{}{0pt}%
\pgfpathmoveto{\pgfqpoint{4.255037in}{3.490208in}}%
\pgfpathlineto{\pgfqpoint{4.427120in}{3.490208in}}%
\pgfpathlineto{\pgfqpoint{4.427120in}{3.503652in}}%
\pgfpathlineto{\pgfqpoint{4.255037in}{3.503652in}}%
\pgfpathlineto{\pgfqpoint{4.255037in}{3.490208in}}%
\pgfusepath{fill}%
\end{pgfscope}%
\begin{pgfscope}%
\pgfpathrectangle{\pgfqpoint{4.255037in}{0.357758in}}{\pgfqpoint{0.172083in}{3.441662in}}%
\pgfusepath{clip}%
\pgfsetbuttcap%
\pgfsetroundjoin%
\definecolor{currentfill}{rgb}{0.984031,0.835315,0.142528}%
\pgfsetfillcolor{currentfill}%
\pgfsetlinewidth{0.000000pt}%
\definecolor{currentstroke}{rgb}{0.000000,0.000000,0.000000}%
\pgfsetstrokecolor{currentstroke}%
\pgfsetdash{}{0pt}%
\pgfpathmoveto{\pgfqpoint{4.255037in}{3.503652in}}%
\pgfpathlineto{\pgfqpoint{4.427120in}{3.503652in}}%
\pgfpathlineto{\pgfqpoint{4.427120in}{3.517096in}}%
\pgfpathlineto{\pgfqpoint{4.255037in}{3.517096in}}%
\pgfpathlineto{\pgfqpoint{4.255037in}{3.503652in}}%
\pgfusepath{fill}%
\end{pgfscope}%
\begin{pgfscope}%
\pgfpathrectangle{\pgfqpoint{4.255037in}{0.357758in}}{\pgfqpoint{0.172083in}{3.441662in}}%
\pgfusepath{clip}%
\pgfsetbuttcap%
\pgfsetroundjoin%
\definecolor{currentfill}{rgb}{0.982653,0.841812,0.142303}%
\pgfsetfillcolor{currentfill}%
\pgfsetlinewidth{0.000000pt}%
\definecolor{currentstroke}{rgb}{0.000000,0.000000,0.000000}%
\pgfsetstrokecolor{currentstroke}%
\pgfsetdash{}{0pt}%
\pgfpathmoveto{\pgfqpoint{4.255037in}{3.517096in}}%
\pgfpathlineto{\pgfqpoint{4.427120in}{3.517096in}}%
\pgfpathlineto{\pgfqpoint{4.427120in}{3.530540in}}%
\pgfpathlineto{\pgfqpoint{4.255037in}{3.530540in}}%
\pgfpathlineto{\pgfqpoint{4.255037in}{3.517096in}}%
\pgfusepath{fill}%
\end{pgfscope}%
\begin{pgfscope}%
\pgfpathrectangle{\pgfqpoint{4.255037in}{0.357758in}}{\pgfqpoint{0.172083in}{3.441662in}}%
\pgfusepath{clip}%
\pgfsetbuttcap%
\pgfsetroundjoin%
\definecolor{currentfill}{rgb}{0.981190,0.848329,0.142279}%
\pgfsetfillcolor{currentfill}%
\pgfsetlinewidth{0.000000pt}%
\definecolor{currentstroke}{rgb}{0.000000,0.000000,0.000000}%
\pgfsetstrokecolor{currentstroke}%
\pgfsetdash{}{0pt}%
\pgfpathmoveto{\pgfqpoint{4.255037in}{3.530540in}}%
\pgfpathlineto{\pgfqpoint{4.427120in}{3.530540in}}%
\pgfpathlineto{\pgfqpoint{4.427120in}{3.543984in}}%
\pgfpathlineto{\pgfqpoint{4.255037in}{3.543984in}}%
\pgfpathlineto{\pgfqpoint{4.255037in}{3.530540in}}%
\pgfusepath{fill}%
\end{pgfscope}%
\begin{pgfscope}%
\pgfpathrectangle{\pgfqpoint{4.255037in}{0.357758in}}{\pgfqpoint{0.172083in}{3.441662in}}%
\pgfusepath{clip}%
\pgfsetbuttcap%
\pgfsetroundjoin%
\definecolor{currentfill}{rgb}{0.979644,0.854866,0.142453}%
\pgfsetfillcolor{currentfill}%
\pgfsetlinewidth{0.000000pt}%
\definecolor{currentstroke}{rgb}{0.000000,0.000000,0.000000}%
\pgfsetstrokecolor{currentstroke}%
\pgfsetdash{}{0pt}%
\pgfpathmoveto{\pgfqpoint{4.255037in}{3.543984in}}%
\pgfpathlineto{\pgfqpoint{4.427120in}{3.543984in}}%
\pgfpathlineto{\pgfqpoint{4.427120in}{3.557428in}}%
\pgfpathlineto{\pgfqpoint{4.255037in}{3.557428in}}%
\pgfpathlineto{\pgfqpoint{4.255037in}{3.543984in}}%
\pgfusepath{fill}%
\end{pgfscope}%
\begin{pgfscope}%
\pgfpathrectangle{\pgfqpoint{4.255037in}{0.357758in}}{\pgfqpoint{0.172083in}{3.441662in}}%
\pgfusepath{clip}%
\pgfsetbuttcap%
\pgfsetroundjoin%
\definecolor{currentfill}{rgb}{0.977995,0.861432,0.142808}%
\pgfsetfillcolor{currentfill}%
\pgfsetlinewidth{0.000000pt}%
\definecolor{currentstroke}{rgb}{0.000000,0.000000,0.000000}%
\pgfsetstrokecolor{currentstroke}%
\pgfsetdash{}{0pt}%
\pgfpathmoveto{\pgfqpoint{4.255037in}{3.557428in}}%
\pgfpathlineto{\pgfqpoint{4.427120in}{3.557428in}}%
\pgfpathlineto{\pgfqpoint{4.427120in}{3.570872in}}%
\pgfpathlineto{\pgfqpoint{4.255037in}{3.570872in}}%
\pgfpathlineto{\pgfqpoint{4.255037in}{3.557428in}}%
\pgfusepath{fill}%
\end{pgfscope}%
\begin{pgfscope}%
\pgfpathrectangle{\pgfqpoint{4.255037in}{0.357758in}}{\pgfqpoint{0.172083in}{3.441662in}}%
\pgfusepath{clip}%
\pgfsetbuttcap%
\pgfsetroundjoin%
\definecolor{currentfill}{rgb}{0.976265,0.868016,0.143351}%
\pgfsetfillcolor{currentfill}%
\pgfsetlinewidth{0.000000pt}%
\definecolor{currentstroke}{rgb}{0.000000,0.000000,0.000000}%
\pgfsetstrokecolor{currentstroke}%
\pgfsetdash{}{0pt}%
\pgfpathmoveto{\pgfqpoint{4.255037in}{3.570872in}}%
\pgfpathlineto{\pgfqpoint{4.427120in}{3.570872in}}%
\pgfpathlineto{\pgfqpoint{4.427120in}{3.584316in}}%
\pgfpathlineto{\pgfqpoint{4.255037in}{3.584316in}}%
\pgfpathlineto{\pgfqpoint{4.255037in}{3.570872in}}%
\pgfusepath{fill}%
\end{pgfscope}%
\begin{pgfscope}%
\pgfpathrectangle{\pgfqpoint{4.255037in}{0.357758in}}{\pgfqpoint{0.172083in}{3.441662in}}%
\pgfusepath{clip}%
\pgfsetbuttcap%
\pgfsetroundjoin%
\definecolor{currentfill}{rgb}{0.974443,0.874622,0.144061}%
\pgfsetfillcolor{currentfill}%
\pgfsetlinewidth{0.000000pt}%
\definecolor{currentstroke}{rgb}{0.000000,0.000000,0.000000}%
\pgfsetstrokecolor{currentstroke}%
\pgfsetdash{}{0pt}%
\pgfpathmoveto{\pgfqpoint{4.255037in}{3.584316in}}%
\pgfpathlineto{\pgfqpoint{4.427120in}{3.584316in}}%
\pgfpathlineto{\pgfqpoint{4.427120in}{3.597760in}}%
\pgfpathlineto{\pgfqpoint{4.255037in}{3.597760in}}%
\pgfpathlineto{\pgfqpoint{4.255037in}{3.584316in}}%
\pgfusepath{fill}%
\end{pgfscope}%
\begin{pgfscope}%
\pgfpathrectangle{\pgfqpoint{4.255037in}{0.357758in}}{\pgfqpoint{0.172083in}{3.441662in}}%
\pgfusepath{clip}%
\pgfsetbuttcap%
\pgfsetroundjoin%
\definecolor{currentfill}{rgb}{0.972530,0.881250,0.144923}%
\pgfsetfillcolor{currentfill}%
\pgfsetlinewidth{0.000000pt}%
\definecolor{currentstroke}{rgb}{0.000000,0.000000,0.000000}%
\pgfsetstrokecolor{currentstroke}%
\pgfsetdash{}{0pt}%
\pgfpathmoveto{\pgfqpoint{4.255037in}{3.597760in}}%
\pgfpathlineto{\pgfqpoint{4.427120in}{3.597760in}}%
\pgfpathlineto{\pgfqpoint{4.427120in}{3.611204in}}%
\pgfpathlineto{\pgfqpoint{4.255037in}{3.611204in}}%
\pgfpathlineto{\pgfqpoint{4.255037in}{3.597760in}}%
\pgfusepath{fill}%
\end{pgfscope}%
\begin{pgfscope}%
\pgfpathrectangle{\pgfqpoint{4.255037in}{0.357758in}}{\pgfqpoint{0.172083in}{3.441662in}}%
\pgfusepath{clip}%
\pgfsetbuttcap%
\pgfsetroundjoin%
\definecolor{currentfill}{rgb}{0.970533,0.887896,0.145919}%
\pgfsetfillcolor{currentfill}%
\pgfsetlinewidth{0.000000pt}%
\definecolor{currentstroke}{rgb}{0.000000,0.000000,0.000000}%
\pgfsetstrokecolor{currentstroke}%
\pgfsetdash{}{0pt}%
\pgfpathmoveto{\pgfqpoint{4.255037in}{3.611204in}}%
\pgfpathlineto{\pgfqpoint{4.427120in}{3.611204in}}%
\pgfpathlineto{\pgfqpoint{4.427120in}{3.624648in}}%
\pgfpathlineto{\pgfqpoint{4.255037in}{3.624648in}}%
\pgfpathlineto{\pgfqpoint{4.255037in}{3.611204in}}%
\pgfusepath{fill}%
\end{pgfscope}%
\begin{pgfscope}%
\pgfpathrectangle{\pgfqpoint{4.255037in}{0.357758in}}{\pgfqpoint{0.172083in}{3.441662in}}%
\pgfusepath{clip}%
\pgfsetbuttcap%
\pgfsetroundjoin%
\definecolor{currentfill}{rgb}{0.968443,0.894564,0.147014}%
\pgfsetfillcolor{currentfill}%
\pgfsetlinewidth{0.000000pt}%
\definecolor{currentstroke}{rgb}{0.000000,0.000000,0.000000}%
\pgfsetstrokecolor{currentstroke}%
\pgfsetdash{}{0pt}%
\pgfpathmoveto{\pgfqpoint{4.255037in}{3.624648in}}%
\pgfpathlineto{\pgfqpoint{4.427120in}{3.624648in}}%
\pgfpathlineto{\pgfqpoint{4.427120in}{3.638092in}}%
\pgfpathlineto{\pgfqpoint{4.255037in}{3.638092in}}%
\pgfpathlineto{\pgfqpoint{4.255037in}{3.624648in}}%
\pgfusepath{fill}%
\end{pgfscope}%
\begin{pgfscope}%
\pgfpathrectangle{\pgfqpoint{4.255037in}{0.357758in}}{\pgfqpoint{0.172083in}{3.441662in}}%
\pgfusepath{clip}%
\pgfsetbuttcap%
\pgfsetroundjoin%
\definecolor{currentfill}{rgb}{0.966271,0.901249,0.148180}%
\pgfsetfillcolor{currentfill}%
\pgfsetlinewidth{0.000000pt}%
\definecolor{currentstroke}{rgb}{0.000000,0.000000,0.000000}%
\pgfsetstrokecolor{currentstroke}%
\pgfsetdash{}{0pt}%
\pgfpathmoveto{\pgfqpoint{4.255037in}{3.638092in}}%
\pgfpathlineto{\pgfqpoint{4.427120in}{3.638092in}}%
\pgfpathlineto{\pgfqpoint{4.427120in}{3.651536in}}%
\pgfpathlineto{\pgfqpoint{4.255037in}{3.651536in}}%
\pgfpathlineto{\pgfqpoint{4.255037in}{3.638092in}}%
\pgfusepath{fill}%
\end{pgfscope}%
\begin{pgfscope}%
\pgfpathrectangle{\pgfqpoint{4.255037in}{0.357758in}}{\pgfqpoint{0.172083in}{3.441662in}}%
\pgfusepath{clip}%
\pgfsetbuttcap%
\pgfsetroundjoin%
\definecolor{currentfill}{rgb}{0.964021,0.907950,0.149370}%
\pgfsetfillcolor{currentfill}%
\pgfsetlinewidth{0.000000pt}%
\definecolor{currentstroke}{rgb}{0.000000,0.000000,0.000000}%
\pgfsetstrokecolor{currentstroke}%
\pgfsetdash{}{0pt}%
\pgfpathmoveto{\pgfqpoint{4.255037in}{3.651536in}}%
\pgfpathlineto{\pgfqpoint{4.427120in}{3.651536in}}%
\pgfpathlineto{\pgfqpoint{4.427120in}{3.664980in}}%
\pgfpathlineto{\pgfqpoint{4.255037in}{3.664980in}}%
\pgfpathlineto{\pgfqpoint{4.255037in}{3.651536in}}%
\pgfusepath{fill}%
\end{pgfscope}%
\begin{pgfscope}%
\pgfpathrectangle{\pgfqpoint{4.255037in}{0.357758in}}{\pgfqpoint{0.172083in}{3.441662in}}%
\pgfusepath{clip}%
\pgfsetbuttcap%
\pgfsetroundjoin%
\definecolor{currentfill}{rgb}{0.961681,0.914672,0.150520}%
\pgfsetfillcolor{currentfill}%
\pgfsetlinewidth{0.000000pt}%
\definecolor{currentstroke}{rgb}{0.000000,0.000000,0.000000}%
\pgfsetstrokecolor{currentstroke}%
\pgfsetdash{}{0pt}%
\pgfpathmoveto{\pgfqpoint{4.255037in}{3.664980in}}%
\pgfpathlineto{\pgfqpoint{4.427120in}{3.664980in}}%
\pgfpathlineto{\pgfqpoint{4.427120in}{3.678424in}}%
\pgfpathlineto{\pgfqpoint{4.255037in}{3.678424in}}%
\pgfpathlineto{\pgfqpoint{4.255037in}{3.664980in}}%
\pgfusepath{fill}%
\end{pgfscope}%
\begin{pgfscope}%
\pgfpathrectangle{\pgfqpoint{4.255037in}{0.357758in}}{\pgfqpoint{0.172083in}{3.441662in}}%
\pgfusepath{clip}%
\pgfsetbuttcap%
\pgfsetroundjoin%
\definecolor{currentfill}{rgb}{0.959276,0.921407,0.151566}%
\pgfsetfillcolor{currentfill}%
\pgfsetlinewidth{0.000000pt}%
\definecolor{currentstroke}{rgb}{0.000000,0.000000,0.000000}%
\pgfsetstrokecolor{currentstroke}%
\pgfsetdash{}{0pt}%
\pgfpathmoveto{\pgfqpoint{4.255037in}{3.678424in}}%
\pgfpathlineto{\pgfqpoint{4.427120in}{3.678424in}}%
\pgfpathlineto{\pgfqpoint{4.427120in}{3.691868in}}%
\pgfpathlineto{\pgfqpoint{4.255037in}{3.691868in}}%
\pgfpathlineto{\pgfqpoint{4.255037in}{3.678424in}}%
\pgfusepath{fill}%
\end{pgfscope}%
\begin{pgfscope}%
\pgfpathrectangle{\pgfqpoint{4.255037in}{0.357758in}}{\pgfqpoint{0.172083in}{3.441662in}}%
\pgfusepath{clip}%
\pgfsetbuttcap%
\pgfsetroundjoin%
\definecolor{currentfill}{rgb}{0.956808,0.928152,0.152409}%
\pgfsetfillcolor{currentfill}%
\pgfsetlinewidth{0.000000pt}%
\definecolor{currentstroke}{rgb}{0.000000,0.000000,0.000000}%
\pgfsetstrokecolor{currentstroke}%
\pgfsetdash{}{0pt}%
\pgfpathmoveto{\pgfqpoint{4.255037in}{3.691868in}}%
\pgfpathlineto{\pgfqpoint{4.427120in}{3.691868in}}%
\pgfpathlineto{\pgfqpoint{4.427120in}{3.705312in}}%
\pgfpathlineto{\pgfqpoint{4.255037in}{3.705312in}}%
\pgfpathlineto{\pgfqpoint{4.255037in}{3.691868in}}%
\pgfusepath{fill}%
\end{pgfscope}%
\begin{pgfscope}%
\pgfpathrectangle{\pgfqpoint{4.255037in}{0.357758in}}{\pgfqpoint{0.172083in}{3.441662in}}%
\pgfusepath{clip}%
\pgfsetbuttcap%
\pgfsetroundjoin%
\definecolor{currentfill}{rgb}{0.954287,0.934908,0.152921}%
\pgfsetfillcolor{currentfill}%
\pgfsetlinewidth{0.000000pt}%
\definecolor{currentstroke}{rgb}{0.000000,0.000000,0.000000}%
\pgfsetstrokecolor{currentstroke}%
\pgfsetdash{}{0pt}%
\pgfpathmoveto{\pgfqpoint{4.255037in}{3.705312in}}%
\pgfpathlineto{\pgfqpoint{4.427120in}{3.705312in}}%
\pgfpathlineto{\pgfqpoint{4.427120in}{3.718756in}}%
\pgfpathlineto{\pgfqpoint{4.255037in}{3.718756in}}%
\pgfpathlineto{\pgfqpoint{4.255037in}{3.705312in}}%
\pgfusepath{fill}%
\end{pgfscope}%
\begin{pgfscope}%
\pgfpathrectangle{\pgfqpoint{4.255037in}{0.357758in}}{\pgfqpoint{0.172083in}{3.441662in}}%
\pgfusepath{clip}%
\pgfsetbuttcap%
\pgfsetroundjoin%
\definecolor{currentfill}{rgb}{0.951726,0.941671,0.152925}%
\pgfsetfillcolor{currentfill}%
\pgfsetlinewidth{0.000000pt}%
\definecolor{currentstroke}{rgb}{0.000000,0.000000,0.000000}%
\pgfsetstrokecolor{currentstroke}%
\pgfsetdash{}{0pt}%
\pgfpathmoveto{\pgfqpoint{4.255037in}{3.718756in}}%
\pgfpathlineto{\pgfqpoint{4.427120in}{3.718756in}}%
\pgfpathlineto{\pgfqpoint{4.427120in}{3.732200in}}%
\pgfpathlineto{\pgfqpoint{4.255037in}{3.732200in}}%
\pgfpathlineto{\pgfqpoint{4.255037in}{3.718756in}}%
\pgfusepath{fill}%
\end{pgfscope}%
\begin{pgfscope}%
\pgfpathrectangle{\pgfqpoint{4.255037in}{0.357758in}}{\pgfqpoint{0.172083in}{3.441662in}}%
\pgfusepath{clip}%
\pgfsetbuttcap%
\pgfsetroundjoin%
\definecolor{currentfill}{rgb}{0.949151,0.948435,0.152178}%
\pgfsetfillcolor{currentfill}%
\pgfsetlinewidth{0.000000pt}%
\definecolor{currentstroke}{rgb}{0.000000,0.000000,0.000000}%
\pgfsetstrokecolor{currentstroke}%
\pgfsetdash{}{0pt}%
\pgfpathmoveto{\pgfqpoint{4.255037in}{3.732200in}}%
\pgfpathlineto{\pgfqpoint{4.427120in}{3.732200in}}%
\pgfpathlineto{\pgfqpoint{4.427120in}{3.745644in}}%
\pgfpathlineto{\pgfqpoint{4.255037in}{3.745644in}}%
\pgfpathlineto{\pgfqpoint{4.255037in}{3.732200in}}%
\pgfusepath{fill}%
\end{pgfscope}%
\begin{pgfscope}%
\pgfpathrectangle{\pgfqpoint{4.255037in}{0.357758in}}{\pgfqpoint{0.172083in}{3.441662in}}%
\pgfusepath{clip}%
\pgfsetbuttcap%
\pgfsetroundjoin%
\definecolor{currentfill}{rgb}{0.946602,0.955190,0.150328}%
\pgfsetfillcolor{currentfill}%
\pgfsetlinewidth{0.000000pt}%
\definecolor{currentstroke}{rgb}{0.000000,0.000000,0.000000}%
\pgfsetstrokecolor{currentstroke}%
\pgfsetdash{}{0pt}%
\pgfpathmoveto{\pgfqpoint{4.255037in}{3.745644in}}%
\pgfpathlineto{\pgfqpoint{4.427120in}{3.745644in}}%
\pgfpathlineto{\pgfqpoint{4.427120in}{3.759088in}}%
\pgfpathlineto{\pgfqpoint{4.255037in}{3.759088in}}%
\pgfpathlineto{\pgfqpoint{4.255037in}{3.745644in}}%
\pgfusepath{fill}%
\end{pgfscope}%
\begin{pgfscope}%
\pgfpathrectangle{\pgfqpoint{4.255037in}{0.357758in}}{\pgfqpoint{0.172083in}{3.441662in}}%
\pgfusepath{clip}%
\pgfsetbuttcap%
\pgfsetroundjoin%
\definecolor{currentfill}{rgb}{0.944152,0.961916,0.146861}%
\pgfsetfillcolor{currentfill}%
\pgfsetlinewidth{0.000000pt}%
\definecolor{currentstroke}{rgb}{0.000000,0.000000,0.000000}%
\pgfsetstrokecolor{currentstroke}%
\pgfsetdash{}{0pt}%
\pgfpathmoveto{\pgfqpoint{4.255037in}{3.759088in}}%
\pgfpathlineto{\pgfqpoint{4.427120in}{3.759088in}}%
\pgfpathlineto{\pgfqpoint{4.427120in}{3.772532in}}%
\pgfpathlineto{\pgfqpoint{4.255037in}{3.772532in}}%
\pgfpathlineto{\pgfqpoint{4.255037in}{3.759088in}}%
\pgfusepath{fill}%
\end{pgfscope}%
\begin{pgfscope}%
\pgfpathrectangle{\pgfqpoint{4.255037in}{0.357758in}}{\pgfqpoint{0.172083in}{3.441662in}}%
\pgfusepath{clip}%
\pgfsetbuttcap%
\pgfsetroundjoin%
\definecolor{currentfill}{rgb}{0.941896,0.968590,0.140956}%
\pgfsetfillcolor{currentfill}%
\pgfsetlinewidth{0.000000pt}%
\definecolor{currentstroke}{rgb}{0.000000,0.000000,0.000000}%
\pgfsetstrokecolor{currentstroke}%
\pgfsetdash{}{0pt}%
\pgfpathmoveto{\pgfqpoint{4.255037in}{3.772532in}}%
\pgfpathlineto{\pgfqpoint{4.427120in}{3.772532in}}%
\pgfpathlineto{\pgfqpoint{4.427120in}{3.785976in}}%
\pgfpathlineto{\pgfqpoint{4.255037in}{3.785976in}}%
\pgfpathlineto{\pgfqpoint{4.255037in}{3.772532in}}%
\pgfusepath{fill}%
\end{pgfscope}%
\begin{pgfscope}%
\pgfpathrectangle{\pgfqpoint{4.255037in}{0.357758in}}{\pgfqpoint{0.172083in}{3.441662in}}%
\pgfusepath{clip}%
\pgfsetbuttcap%
\pgfsetroundjoin%
\definecolor{currentfill}{rgb}{0.940015,0.975158,0.131326}%
\pgfsetfillcolor{currentfill}%
\pgfsetlinewidth{0.000000pt}%
\definecolor{currentstroke}{rgb}{0.000000,0.000000,0.000000}%
\pgfsetstrokecolor{currentstroke}%
\pgfsetdash{}{0pt}%
\pgfpathmoveto{\pgfqpoint{4.255037in}{3.785976in}}%
\pgfpathlineto{\pgfqpoint{4.427120in}{3.785976in}}%
\pgfpathlineto{\pgfqpoint{4.427120in}{3.799420in}}%
\pgfpathlineto{\pgfqpoint{4.255037in}{3.799420in}}%
\pgfpathlineto{\pgfqpoint{4.255037in}{3.785976in}}%
\pgfusepath{fill}%
\end{pgfscope}%
\begin{pgfscope}%
\pgfsetbuttcap%
\pgfsetroundjoin%
\definecolor{currentfill}{rgb}{0.000000,0.000000,0.000000}%
\pgfsetfillcolor{currentfill}%
\pgfsetlinewidth{0.803000pt}%
\definecolor{currentstroke}{rgb}{0.000000,0.000000,0.000000}%
\pgfsetstrokecolor{currentstroke}%
\pgfsetdash{}{0pt}%
\pgfsys@defobject{currentmarker}{\pgfqpoint{0.000000in}{0.000000in}}{\pgfqpoint{0.048611in}{0.000000in}}{%
\pgfpathmoveto{\pgfqpoint{0.000000in}{0.000000in}}%
\pgfpathlineto{\pgfqpoint{0.048611in}{0.000000in}}%
\pgfusepath{stroke,fill}%
}%
\begin{pgfscope}%
\pgfsys@transformshift{4.427120in}{0.536447in}%
\pgfsys@useobject{currentmarker}{}%
\end{pgfscope}%
\end{pgfscope}%
\begin{pgfscope}%
\definecolor{textcolor}{rgb}{0.000000,0.000000,0.000000}%
\pgfsetstrokecolor{textcolor}%
\pgfsetfillcolor{textcolor}%
\pgftext[x=4.524343in, y=0.488253in, left, base]{\color{textcolor}{\rmfamily\fontsize{10.000000}{12.000000}\selectfont\catcode`\^=\active\def^{\ifmmode\sp\else\^{}\fi}\catcode`\%=\active\def%{\%}$\mathdefault{\ensuremath{-}0.6}$}}%
\end{pgfscope}%
\begin{pgfscope}%
\pgfsetbuttcap%
\pgfsetroundjoin%
\definecolor{currentfill}{rgb}{0.000000,0.000000,0.000000}%
\pgfsetfillcolor{currentfill}%
\pgfsetlinewidth{0.803000pt}%
\definecolor{currentstroke}{rgb}{0.000000,0.000000,0.000000}%
\pgfsetstrokecolor{currentstroke}%
\pgfsetdash{}{0pt}%
\pgfsys@defobject{currentmarker}{\pgfqpoint{0.000000in}{0.000000in}}{\pgfqpoint{0.048611in}{0.000000in}}{%
\pgfpathmoveto{\pgfqpoint{0.000000in}{0.000000in}}%
\pgfpathlineto{\pgfqpoint{0.048611in}{0.000000in}}%
\pgfusepath{stroke,fill}%
}%
\begin{pgfscope}%
\pgfsys@transformshift{4.427120in}{0.945312in}%
\pgfsys@useobject{currentmarker}{}%
\end{pgfscope}%
\end{pgfscope}%
\begin{pgfscope}%
\definecolor{textcolor}{rgb}{0.000000,0.000000,0.000000}%
\pgfsetstrokecolor{textcolor}%
\pgfsetfillcolor{textcolor}%
\pgftext[x=4.524343in, y=0.897118in, left, base]{\color{textcolor}{\rmfamily\fontsize{10.000000}{12.000000}\selectfont\catcode`\^=\active\def^{\ifmmode\sp\else\^{}\fi}\catcode`\%=\active\def%{\%}$\mathdefault{\ensuremath{-}0.4}$}}%
\end{pgfscope}%
\begin{pgfscope}%
\pgfsetbuttcap%
\pgfsetroundjoin%
\definecolor{currentfill}{rgb}{0.000000,0.000000,0.000000}%
\pgfsetfillcolor{currentfill}%
\pgfsetlinewidth{0.803000pt}%
\definecolor{currentstroke}{rgb}{0.000000,0.000000,0.000000}%
\pgfsetstrokecolor{currentstroke}%
\pgfsetdash{}{0pt}%
\pgfsys@defobject{currentmarker}{\pgfqpoint{0.000000in}{0.000000in}}{\pgfqpoint{0.048611in}{0.000000in}}{%
\pgfpathmoveto{\pgfqpoint{0.000000in}{0.000000in}}%
\pgfpathlineto{\pgfqpoint{0.048611in}{0.000000in}}%
\pgfusepath{stroke,fill}%
}%
\begin{pgfscope}%
\pgfsys@transformshift{4.427120in}{1.354177in}%
\pgfsys@useobject{currentmarker}{}%
\end{pgfscope}%
\end{pgfscope}%
\begin{pgfscope}%
\definecolor{textcolor}{rgb}{0.000000,0.000000,0.000000}%
\pgfsetstrokecolor{textcolor}%
\pgfsetfillcolor{textcolor}%
\pgftext[x=4.524343in, y=1.305983in, left, base]{\color{textcolor}{\rmfamily\fontsize{10.000000}{12.000000}\selectfont\catcode`\^=\active\def^{\ifmmode\sp\else\^{}\fi}\catcode`\%=\active\def%{\%}$\mathdefault{\ensuremath{-}0.2}$}}%
\end{pgfscope}%
\begin{pgfscope}%
\pgfsetbuttcap%
\pgfsetroundjoin%
\definecolor{currentfill}{rgb}{0.000000,0.000000,0.000000}%
\pgfsetfillcolor{currentfill}%
\pgfsetlinewidth{0.803000pt}%
\definecolor{currentstroke}{rgb}{0.000000,0.000000,0.000000}%
\pgfsetstrokecolor{currentstroke}%
\pgfsetdash{}{0pt}%
\pgfsys@defobject{currentmarker}{\pgfqpoint{0.000000in}{0.000000in}}{\pgfqpoint{0.048611in}{0.000000in}}{%
\pgfpathmoveto{\pgfqpoint{0.000000in}{0.000000in}}%
\pgfpathlineto{\pgfqpoint{0.048611in}{0.000000in}}%
\pgfusepath{stroke,fill}%
}%
\begin{pgfscope}%
\pgfsys@transformshift{4.427120in}{1.763042in}%
\pgfsys@useobject{currentmarker}{}%
\end{pgfscope}%
\end{pgfscope}%
\begin{pgfscope}%
\definecolor{textcolor}{rgb}{0.000000,0.000000,0.000000}%
\pgfsetstrokecolor{textcolor}%
\pgfsetfillcolor{textcolor}%
\pgftext[x=4.524343in, y=1.714848in, left, base]{\color{textcolor}{\rmfamily\fontsize{10.000000}{12.000000}\selectfont\catcode`\^=\active\def^{\ifmmode\sp\else\^{}\fi}\catcode`\%=\active\def%{\%}$\mathdefault{0.0}$}}%
\end{pgfscope}%
\begin{pgfscope}%
\pgfsetbuttcap%
\pgfsetroundjoin%
\definecolor{currentfill}{rgb}{0.000000,0.000000,0.000000}%
\pgfsetfillcolor{currentfill}%
\pgfsetlinewidth{0.803000pt}%
\definecolor{currentstroke}{rgb}{0.000000,0.000000,0.000000}%
\pgfsetstrokecolor{currentstroke}%
\pgfsetdash{}{0pt}%
\pgfsys@defobject{currentmarker}{\pgfqpoint{0.000000in}{0.000000in}}{\pgfqpoint{0.048611in}{0.000000in}}{%
\pgfpathmoveto{\pgfqpoint{0.000000in}{0.000000in}}%
\pgfpathlineto{\pgfqpoint{0.048611in}{0.000000in}}%
\pgfusepath{stroke,fill}%
}%
\begin{pgfscope}%
\pgfsys@transformshift{4.427120in}{2.171908in}%
\pgfsys@useobject{currentmarker}{}%
\end{pgfscope}%
\end{pgfscope}%
\begin{pgfscope}%
\definecolor{textcolor}{rgb}{0.000000,0.000000,0.000000}%
\pgfsetstrokecolor{textcolor}%
\pgfsetfillcolor{textcolor}%
\pgftext[x=4.524343in, y=2.123713in, left, base]{\color{textcolor}{\rmfamily\fontsize{10.000000}{12.000000}\selectfont\catcode`\^=\active\def^{\ifmmode\sp\else\^{}\fi}\catcode`\%=\active\def%{\%}$\mathdefault{0.2}$}}%
\end{pgfscope}%
\begin{pgfscope}%
\pgfsetbuttcap%
\pgfsetroundjoin%
\definecolor{currentfill}{rgb}{0.000000,0.000000,0.000000}%
\pgfsetfillcolor{currentfill}%
\pgfsetlinewidth{0.803000pt}%
\definecolor{currentstroke}{rgb}{0.000000,0.000000,0.000000}%
\pgfsetstrokecolor{currentstroke}%
\pgfsetdash{}{0pt}%
\pgfsys@defobject{currentmarker}{\pgfqpoint{0.000000in}{0.000000in}}{\pgfqpoint{0.048611in}{0.000000in}}{%
\pgfpathmoveto{\pgfqpoint{0.000000in}{0.000000in}}%
\pgfpathlineto{\pgfqpoint{0.048611in}{0.000000in}}%
\pgfusepath{stroke,fill}%
}%
\begin{pgfscope}%
\pgfsys@transformshift{4.427120in}{2.580773in}%
\pgfsys@useobject{currentmarker}{}%
\end{pgfscope}%
\end{pgfscope}%
\begin{pgfscope}%
\definecolor{textcolor}{rgb}{0.000000,0.000000,0.000000}%
\pgfsetstrokecolor{textcolor}%
\pgfsetfillcolor{textcolor}%
\pgftext[x=4.524343in, y=2.532578in, left, base]{\color{textcolor}{\rmfamily\fontsize{10.000000}{12.000000}\selectfont\catcode`\^=\active\def^{\ifmmode\sp\else\^{}\fi}\catcode`\%=\active\def%{\%}$\mathdefault{0.4}$}}%
\end{pgfscope}%
\begin{pgfscope}%
\pgfsetbuttcap%
\pgfsetroundjoin%
\definecolor{currentfill}{rgb}{0.000000,0.000000,0.000000}%
\pgfsetfillcolor{currentfill}%
\pgfsetlinewidth{0.803000pt}%
\definecolor{currentstroke}{rgb}{0.000000,0.000000,0.000000}%
\pgfsetstrokecolor{currentstroke}%
\pgfsetdash{}{0pt}%
\pgfsys@defobject{currentmarker}{\pgfqpoint{0.000000in}{0.000000in}}{\pgfqpoint{0.048611in}{0.000000in}}{%
\pgfpathmoveto{\pgfqpoint{0.000000in}{0.000000in}}%
\pgfpathlineto{\pgfqpoint{0.048611in}{0.000000in}}%
\pgfusepath{stroke,fill}%
}%
\begin{pgfscope}%
\pgfsys@transformshift{4.427120in}{2.989638in}%
\pgfsys@useobject{currentmarker}{}%
\end{pgfscope}%
\end{pgfscope}%
\begin{pgfscope}%
\definecolor{textcolor}{rgb}{0.000000,0.000000,0.000000}%
\pgfsetstrokecolor{textcolor}%
\pgfsetfillcolor{textcolor}%
\pgftext[x=4.524343in, y=2.941443in, left, base]{\color{textcolor}{\rmfamily\fontsize{10.000000}{12.000000}\selectfont\catcode`\^=\active\def^{\ifmmode\sp\else\^{}\fi}\catcode`\%=\active\def%{\%}$\mathdefault{0.6}$}}%
\end{pgfscope}%
\begin{pgfscope}%
\pgfsetbuttcap%
\pgfsetroundjoin%
\definecolor{currentfill}{rgb}{0.000000,0.000000,0.000000}%
\pgfsetfillcolor{currentfill}%
\pgfsetlinewidth{0.803000pt}%
\definecolor{currentstroke}{rgb}{0.000000,0.000000,0.000000}%
\pgfsetstrokecolor{currentstroke}%
\pgfsetdash{}{0pt}%
\pgfsys@defobject{currentmarker}{\pgfqpoint{0.000000in}{0.000000in}}{\pgfqpoint{0.048611in}{0.000000in}}{%
\pgfpathmoveto{\pgfqpoint{0.000000in}{0.000000in}}%
\pgfpathlineto{\pgfqpoint{0.048611in}{0.000000in}}%
\pgfusepath{stroke,fill}%
}%
\begin{pgfscope}%
\pgfsys@transformshift{4.427120in}{3.398503in}%
\pgfsys@useobject{currentmarker}{}%
\end{pgfscope}%
\end{pgfscope}%
\begin{pgfscope}%
\definecolor{textcolor}{rgb}{0.000000,0.000000,0.000000}%
\pgfsetstrokecolor{textcolor}%
\pgfsetfillcolor{textcolor}%
\pgftext[x=4.524343in, y=3.350308in, left, base]{\color{textcolor}{\rmfamily\fontsize{10.000000}{12.000000}\selectfont\catcode`\^=\active\def^{\ifmmode\sp\else\^{}\fi}\catcode`\%=\active\def%{\%}$\mathdefault{0.8}$}}%
\end{pgfscope}%
\begin{pgfscope}%
\definecolor{textcolor}{rgb}{0.000000,0.000000,0.000000}%
\pgfsetstrokecolor{textcolor}%
\pgfsetfillcolor{textcolor}%
\pgftext[x=4.865393in,y=2.078589in,,top,rotate=90.000000]{\color{textcolor}{\rmfamily\fontsize{10.000000}{12.000000}\selectfont\catcode`\^=\active\def^{\ifmmode\sp\else\^{}\fi}\catcode`\%=\active\def%{\%}Altura (Z)}}%
\end{pgfscope}%
\begin{pgfscope}%
\pgfsetrectcap%
\pgfsetmiterjoin%
\pgfsetlinewidth{0.803000pt}%
\definecolor{currentstroke}{rgb}{0.000000,0.000000,0.000000}%
\pgfsetstrokecolor{currentstroke}%
\pgfsetdash{}{0pt}%
\pgfpathmoveto{\pgfqpoint{4.255037in}{0.357758in}}%
\pgfpathlineto{\pgfqpoint{4.341079in}{0.357758in}}%
\pgfpathlineto{\pgfqpoint{4.427120in}{0.357758in}}%
\pgfpathlineto{\pgfqpoint{4.427120in}{3.799420in}}%
\pgfpathlineto{\pgfqpoint{4.341079in}{3.799420in}}%
\pgfpathlineto{\pgfqpoint{4.255037in}{3.799420in}}%
\pgfpathlineto{\pgfqpoint{4.255037in}{0.357758in}}%
\pgfpathclose%
\pgfusepath{stroke}%
\end{pgfscope}%
\end{pgfpicture}%
\makeatother%
\endgroup%
 
    \caption{Gráfico de superficie con datos dispersos}
    \label{fig:grafico_3D_Triangulado}
\end{figure}


\subsubsection{¿Cómo realizar gráficos Tridimensionales para \LaTeX ?}

Para integrar gráficos 3D en documentos profesionales sin perder calidad ni incrementar excesivamente el tamaño del archivo PDF, es necesario garantizar que el motor de renderizado no convierta la superficie en una imagen de mapa de bits (.png) interna. Para evitar esto utilizamos comandos como \texttt{.set\_rasterized(False)}. Vease en el código.

\begin{pycode}{Gráfico de malla regular para LaTeX}
import matplotlib.pyplot as plt
import numpy as np

# Para guardar en otra carpeta:
import os

#####################
#       Código      #
#####################

# ---- Para latex ---- #

plt.rcParams.update({
    "pgf.texsystem": "xelatex",    
    "font.family": "serif",
    "text.usetex": True,       
    "pgf.rcfonts": False,          
    "pgf.preamble": "\n".join([    
         r"\usepackage{fontspec}",
         r"\setmainfont{CMU Serif}",
    ])
})

# ---- Grilla ---- #

plt.rcParams.update({
    'axes.grid':True,
    'grid.linestyle':'--',
    'grid.linewidth':0.5,
    'grid.color':'gray',
    'grid.alpha':0.5
})
    
#####################
#       Datos       #
#####################

# Datos
x = np.linspace(-3, 3, 100) # Siempre reduzca la densidad, más de 100 es demasiado para una demostración.
y = np.linspace(-3, 3, 100)
X, Y = np.meshgrid(x, y)
Z = np.sin(np.sqrt(X**2 + Y**2))

# Entorno 3D
fig, ax = plt.subplots(subplot_kw={"projection": "3d"}, figsize=(5, 5), layout="constrained")

# 1: antialiased=False
# Al exportar de forma vectorial utilizar antialased=True puede crear lineas blancas entre poligonos. Desactivarlo hace que los vectores encajen de forma perfecta
surf = ax.plot_surface(X, Y, Z, cmap="viridis", rstride=1, cstride=1, 
                       linewidth=0, antialiased=False, rasterized=False)

# 2: Forzar que la superficie no sea rasterizada
surf.set_rasterized(False)

# Titulos con LaTeX
ax.set_title(r"Superficie 3D: $z = \sin(\sqrt{x^2 + y^2})$") 
ax.set_xlabel("X")
ax.set_ylabel("Y")
ax.set_zlabel("Z")

# Barra de color
cb = fig.colorbar(surf, ax=ax, shrink=0.7, pad=0.05)
cb.set_label("Altura (Z)")

# 3: Forzamos a que los elementos de la barra de color sean vectoriales. Buscamos que no se simplifique a una imagen .png
cb.solids.set_rasterized(False)

# Vista del gráfico
ax.view_init(elev=25, azim=45)
ax.set_box_aspect((1, 1, 0.5))

# 4. Guardar en la carpeta deseada
# Expandimos la tilde antes de guardar
ruta_tilde = "~/Documentos/.../carpeta_deseada/Nombre_del_documento.pfg"
ruta_real = os.path.expanduser(ruta_tilde)

plt.savefig(ruta_real, bbox_inches='tight')
\end{pycode}